\documentclass[10pt,a4paper]{article}
\usepackage{fontspec,amsmath,amssymb,graphicx}
\usepackage[margin=18mm]{geometry}
\setmainfont{Times New Roman}
\setlength{\parindent}{0pt}
\setlength{\parskip}{5pt}
\allowdisplaybreaks
\sloppy
\emergencystretch=2em
\begin{document}
\clearpage\noindent\texttt{\detokenize{VD_1_2_2_1_01}}\par
\section*{Lời giải bổ sung — Ví dụ 1 · Mục 1.2.2.1}

\subsection*{Kiến thức và trình tự giải}
\textbf{Mục tiêu.} Xác định cấp và dạng của phương trình trước khi chọn cách giải. Ở mục này, từ “bậc” trong đề được hiểu theo nghĩa cấp đạo hàm của phương trình.

\textbf{Bước 1 — Liệt kê đạo hàm và đếm cấp.} Với $\partial^{i+j}u/\partial x^i\partial y^j$, cấp là $i+j$. Vì vậy $u_x$ cấp 1, $u_{xy}$ cấp 2, $u_{txx}$ cấp 3. Hạng $u$, $e^u$, $\log u$ không chứa đạo hàm, được xem là cấp 0. $(u_y)^2$ là bình phương của một đạo hàm cấp 1, không phải đạo hàm cấp 2. Lấy giá trị lớn nhất trong các cấp vừa liệt kê.

\textbf{Bước 2 — Viết rõ phần chứa đạo hàm cấp cao nhất.} Gọi cấp đó là $m$. Tạm xem từng đạo hàm cấp $m$ như một ẩn đại số độc lập. Nếu chúng bị bình phương, nhân với nhau hoặc nằm trong hàm mũ/lượng giác, biểu thức đã cho hoàn toàn phi tuyến. Nếu chúng chỉ xuất hiện tuyến tính, chuyển sang bước 3. Khi khai triển $(f(u))_x$, phải dùng $(f(u))_x=f'(u)u_x$ trước khi phân loại.

\textbf{Bước 3 — Xét hệ số và các hạng cấp thấp.} Nếu mọi hạng chứa $u$ và đạo hàm đều tuyến tính, hệ số chỉ phụ thuộc biến độc lập, phương trình tuyến tính. Nếu chỉ phần cấp $m$ tuyến tính, hệ số của phần đó phụ thuộc biến độc lập, còn phần cấp thấp phi tuyến, phương trình nửa tuyến tính. Nếu hệ số phần cấp $m$ phụ thuộc $u$ hoặc đạo hàm cấp thấp, phương trình tựa tuyến tính. Khi ghi kết luận, chỉ rõ hạng nào quyết định phân loại.

\textbf{Bước 4 — Xét thuần nhất khi phương trình tuyến tính.} Gom toàn bộ hạng chứa ẩn về một vế, viết $Lu=f(x,y)$. Nếu $f\equiv0$, phương trình thuần nhất; nếu không, không thuần nhất. Ví dụ $u_{xy}=u$ phải viết $u_{xy}-u=0$, nên là thuần nhất.

\textbf{Bước 5 — Kiểm tra miền và yêu cầu đề.} $\log u$ đòi hỏi $u>0$ đối với nghiệm thực. Một phép biến đổi như lấy log có thể cho phương trình tương đương trên một miền, nhưng cần phân biệt dạng biểu thức in trong đề và dạng đã biến đổi. Nếu đề chỉ hỏi cấp, dừng sau bước 1 và giải thích đạo hàm quyết định cấp.
\subsection*{Lời giải áp dụng cho bài}
\textbf{Bước 1 — Xác định cấp.} Liệt kê $u_{xxx}$ cấp 3, các đạo hàm còn lại có cấp thấp hơn. Bình phương $(u_y)^2$ vẫn chỉ chứa đạo hàm cấp 1. Do đó cấp cao nhất là $m=3$.

\textbf{Bước 2 — Xét phần chính.} Đạo hàm $u_{xxx}$ xuất hiện ở lũy thừa 1 với hệ số 1, không phụ thuộc $u$ hay đạo hàm của $u$.

\textbf{Bước 3 — Xét phần cấp thấp.} Hạng $(u_y)^2$ phi tuyến nên toàn phương trình không tuyến tính. Vì phi tuyến chỉ ở cấp thấp và hệ số phần chính chỉ phụ thuộc biến độc lập, kết luận: \textbf{nửa tuyến tính, cấp 3}.

\clearpage\noindent\texttt{\detokenize{VD_1_2_2_1_02}}\par
\section*{Lời giải bổ sung — Ví dụ 2 · Mục 1.2.2.1}

\subsection*{Kiến thức và trình tự giải}
\textbf{Mục tiêu.} Xác định cấp và dạng của phương trình trước khi chọn cách giải. Ở mục này, từ “bậc” trong đề được hiểu theo nghĩa cấp đạo hàm của phương trình.

\textbf{Bước 1 — Liệt kê đạo hàm và đếm cấp.} Với $\partial^{i+j}u/\partial x^i\partial y^j$, cấp là $i+j$. Vì vậy $u_x$ cấp 1, $u_{xy}$ cấp 2, $u_{txx}$ cấp 3. Hạng $u$, $e^u$, $\log u$ không chứa đạo hàm, được xem là cấp 0. $(u_y)^2$ là bình phương của một đạo hàm cấp 1, không phải đạo hàm cấp 2. Lấy giá trị lớn nhất trong các cấp vừa liệt kê.

\textbf{Bước 2 — Viết rõ phần chứa đạo hàm cấp cao nhất.} Gọi cấp đó là $m$. Tạm xem từng đạo hàm cấp $m$ như một ẩn đại số độc lập. Nếu chúng bị bình phương, nhân với nhau hoặc nằm trong hàm mũ/lượng giác, biểu thức đã cho hoàn toàn phi tuyến. Nếu chúng chỉ xuất hiện tuyến tính, chuyển sang bước 3. Khi khai triển $(f(u))_x$, phải dùng $(f(u))_x=f'(u)u_x$ trước khi phân loại.

\textbf{Bước 3 — Xét hệ số và các hạng cấp thấp.} Nếu mọi hạng chứa $u$ và đạo hàm đều tuyến tính, hệ số chỉ phụ thuộc biến độc lập, phương trình tuyến tính. Nếu chỉ phần cấp $m$ tuyến tính, hệ số của phần đó phụ thuộc biến độc lập, còn phần cấp thấp phi tuyến, phương trình nửa tuyến tính. Nếu hệ số phần cấp $m$ phụ thuộc $u$ hoặc đạo hàm cấp thấp, phương trình tựa tuyến tính. Khi ghi kết luận, chỉ rõ hạng nào quyết định phân loại.

\textbf{Bước 4 — Xét thuần nhất khi phương trình tuyến tính.} Gom toàn bộ hạng chứa ẩn về một vế, viết $Lu=f(x,y)$. Nếu $f\equiv0$, phương trình thuần nhất; nếu không, không thuần nhất. Ví dụ $u_{xy}=u$ phải viết $u_{xy}-u=0$, nên là thuần nhất.

\textbf{Bước 5 — Kiểm tra miền và yêu cầu đề.} $\log u$ đòi hỏi $u>0$ đối với nghiệm thực. Một phép biến đổi như lấy log có thể cho phương trình tương đương trên một miền, nhưng cần phân biệt dạng biểu thức in trong đề và dạng đã biến đổi. Nếu đề chỉ hỏi cấp, dừng sau bước 1 và giải thích đạo hàm quyết định cấp.
\subsection*{Lời giải áp dụng cho bài}
\textbf{Bước 1 — Đếm cấp.} Đạo hàm cao nhất xuất hiện là $u_x$, nên $m=1$; $u^3$ là hạng không có đạo hàm.

\textbf{Bước 2 — Xét tính tuyến tính theo đạo hàm cấp 1.} Coi $p=u_x$ như một biến đại số. Biểu thức lượng giác và bình phương chứa $p$ không có dạng $A(x,y,u)p+B(x,y,u)$. Như vậy phi tuyến xuất hiện ngay ở đạo hàm cấp cao nhất.

\textbf{Kết luận.} Phương trình \textbf{hoàn toàn phi tuyến, cấp 1}. Hạng $u^3$ không làm tăng cấp và không phải lý do chính để phân biệt hoàn toàn phi tuyến với nửa tuyến tính.

\clearpage\noindent\texttt{\detokenize{VD_1_2_2_1_03}}\par
\section*{Lời giải bổ sung — Ví dụ 3 · Mục 1.2.2.1}

\subsection*{Kiến thức và trình tự giải}
\textbf{Mục tiêu.} Xác định cấp và dạng của phương trình trước khi chọn cách giải. Ở mục này, từ “bậc” trong đề được hiểu theo nghĩa cấp đạo hàm của phương trình.

\textbf{Bước 1 — Liệt kê đạo hàm và đếm cấp.} Với $\partial^{i+j}u/\partial x^i\partial y^j$, cấp là $i+j$. Vì vậy $u_x$ cấp 1, $u_{xy}$ cấp 2, $u_{txx}$ cấp 3. Hạng $u$, $e^u$, $\log u$ không chứa đạo hàm, được xem là cấp 0. $(u_y)^2$ là bình phương của một đạo hàm cấp 1, không phải đạo hàm cấp 2. Lấy giá trị lớn nhất trong các cấp vừa liệt kê.

\textbf{Bước 2 — Viết rõ phần chứa đạo hàm cấp cao nhất.} Gọi cấp đó là $m$. Tạm xem từng đạo hàm cấp $m$ như một ẩn đại số độc lập. Nếu chúng bị bình phương, nhân với nhau hoặc nằm trong hàm mũ/lượng giác, biểu thức đã cho hoàn toàn phi tuyến. Nếu chúng chỉ xuất hiện tuyến tính, chuyển sang bước 3. Khi khai triển $(f(u))_x$, phải dùng $(f(u))_x=f'(u)u_x$ trước khi phân loại.

\textbf{Bước 3 — Xét hệ số và các hạng cấp thấp.} Nếu mọi hạng chứa $u$ và đạo hàm đều tuyến tính, hệ số chỉ phụ thuộc biến độc lập, phương trình tuyến tính. Nếu chỉ phần cấp $m$ tuyến tính, hệ số của phần đó phụ thuộc biến độc lập, còn phần cấp thấp phi tuyến, phương trình nửa tuyến tính. Nếu hệ số phần cấp $m$ phụ thuộc $u$ hoặc đạo hàm cấp thấp, phương trình tựa tuyến tính. Khi ghi kết luận, chỉ rõ hạng nào quyết định phân loại.

\textbf{Bước 4 — Xét thuần nhất khi phương trình tuyến tính.} Gom toàn bộ hạng chứa ẩn về một vế, viết $Lu=f(x,y)$. Nếu $f\equiv0$, phương trình thuần nhất; nếu không, không thuần nhất. Ví dụ $u_{xy}=u$ phải viết $u_{xy}-u=0$, nên là thuần nhất.

\textbf{Bước 5 — Kiểm tra miền và yêu cầu đề.} $\log u$ đòi hỏi $u>0$ đối với nghiệm thực. Một phép biến đổi như lấy log có thể cho phương trình tương đương trên một miền, nhưng cần phân biệt dạng biểu thức in trong đề và dạng đã biến đổi. Nếu đề chỉ hỏi cấp, dừng sau bước 1 và giải thích đạo hàm quyết định cấp.
\subsection*{Lời giải áp dụng cho bài}
\textbf{Bước 1 — Xác định cấp.} Các đạo hàm $u_{xx},u_{xy},u_{yy}$ đều có cấp 2, nên $m=2$.

\textbf{Bước 2 — Đọc các hệ số.} Các hệ số tương ứng là $x,x-y,-y$, chỉ phụ thuộc $x,y$. Mỗi đạo hàm xuất hiện tuyến tính; không có tích giữa các đạo hàm hoặc tích với ẩn $u$.

\textbf{Bước 3 — Xét nguồn.} Phương trình đã có dạng $Lu=0$, nên thuần nhất. Kết luận: \textbf{tuyến tính thuần nhất, cấp 2}. Hệ số thay đổi theo vị trí không làm mất tính tuyến tính.

\clearpage\noindent\texttt{\detokenize{VD_1_2_2_1_04}}\par
\section*{Lời giải bổ sung — Ví dụ 4 · Mục 1.2.2.1}

\subsection*{Kiến thức và trình tự giải}
\textbf{Mục tiêu.} Xác định cấp và dạng của phương trình trước khi chọn cách giải. Ở mục này, từ “bậc” trong đề được hiểu theo nghĩa cấp đạo hàm của phương trình.

\textbf{Bước 1 — Liệt kê đạo hàm và đếm cấp.} Với $\partial^{i+j}u/\partial x^i\partial y^j$, cấp là $i+j$. Vì vậy $u_x$ cấp 1, $u_{xy}$ cấp 2, $u_{txx}$ cấp 3. Hạng $u$, $e^u$, $\log u$ không chứa đạo hàm, được xem là cấp 0. $(u_y)^2$ là bình phương của một đạo hàm cấp 1, không phải đạo hàm cấp 2. Lấy giá trị lớn nhất trong các cấp vừa liệt kê.

\textbf{Bước 2 — Viết rõ phần chứa đạo hàm cấp cao nhất.} Gọi cấp đó là $m$. Tạm xem từng đạo hàm cấp $m$ như một ẩn đại số độc lập. Nếu chúng bị bình phương, nhân với nhau hoặc nằm trong hàm mũ/lượng giác, biểu thức đã cho hoàn toàn phi tuyến. Nếu chúng chỉ xuất hiện tuyến tính, chuyển sang bước 3. Khi khai triển $(f(u))_x$, phải dùng $(f(u))_x=f'(u)u_x$ trước khi phân loại.

\textbf{Bước 3 — Xét hệ số và các hạng cấp thấp.} Nếu mọi hạng chứa $u$ và đạo hàm đều tuyến tính, hệ số chỉ phụ thuộc biến độc lập, phương trình tuyến tính. Nếu chỉ phần cấp $m$ tuyến tính, hệ số của phần đó phụ thuộc biến độc lập, còn phần cấp thấp phi tuyến, phương trình nửa tuyến tính. Nếu hệ số phần cấp $m$ phụ thuộc $u$ hoặc đạo hàm cấp thấp, phương trình tựa tuyến tính. Khi ghi kết luận, chỉ rõ hạng nào quyết định phân loại.

\textbf{Bước 4 — Xét thuần nhất khi phương trình tuyến tính.} Gom toàn bộ hạng chứa ẩn về một vế, viết $Lu=f(x,y)$. Nếu $f\equiv0$, phương trình thuần nhất; nếu không, không thuần nhất. Ví dụ $u_{xy}=u$ phải viết $u_{xy}-u=0$, nên là thuần nhất.

\textbf{Bước 5 — Kiểm tra miền và yêu cầu đề.} $\log u$ đòi hỏi $u>0$ đối với nghiệm thực. Một phép biến đổi như lấy log có thể cho phương trình tương đương trên một miền, nhưng cần phân biệt dạng biểu thức in trong đề và dạng đã biến đổi. Nếu đề chỉ hỏi cấp, dừng sau bước 1 và giải thích đạo hàm quyết định cấp.
\subsection*{Lời giải áp dụng cho bài}
\textbf{Bước 1 — Xác định cấp.} $u_{xx}$ là đạo hàm cao nhất, nên $m=2$.

\textbf{Bước 2 — Xét dạng của phần cấp 2.} Đạo hàm $u_{xx}$ xuất hiện tuyến tính, nhưng hệ số là $u^2$, phụ thuộc chính ẩn cần tìm. Vì thế không thể xem phương trình là tuyến tính hoặc nửa tuyến tính theo định nghĩa trong sách.

\textbf{Kết luận.} Phương trình \textbf{tựa tuyến tính, cấp 2}. Tại $u=0$, hệ số phần chính có thể triệt tiêu; cần lưu ý sự suy biến khi giải, nhưng phân loại biểu thức chung vẫn như trên.

\clearpage\noindent\texttt{\detokenize{VD_1_2_2_2_01}}\par
\section*{Lời giải bổ sung — Ví dụ 1 · Mục 1.2.2.2}

\subsection*{Kiến thức và trình tự giải}
\textbf{Mục tiêu.} Đưa phương trình về ODE bằng cách giữ cố định các biến không được lấy đạo hàm, hoặc đặt một đạo hàm/hợp thức trung gian.

\textbf{Bước 1 — Chọn biến và tham số.} Với $u_x=f(x,y)$, cố định $y$ và tích phân theo $x$. Khi đó “hằng số” tích phân có thể khác nhau ở mỗi $y$, nên phải viết $C(y)$. Với ba biến, tích phân theo $x$ cho một hàm $C(y,z)$.

\textbf{Bước 2 — Tích phân đúng số lần.} Chẳng hạn $u_{xx}=f$ cho $u_x=\int f\,dx+A(y)$, rồi $u=\int(\int f\,dx)dx+xA(y)+B(y)$. Với đạo hàm hỗn hợp, có thể đặt $v=u_x$ hoặc $v=yu_y+2u$ nếu PDE chứa đúng $v_x$. Mỗi lần tích phân phải giữ hàm tùy ý theo các biến còn lại.

\textbf{Bước 3 — Giải ODE tuyến tính nếu còn hạng chứa ẩn.} Với $u_x+P(x)u=Q(x,y)$, đặt $M(x)=e^{\int P(x)dx}$. Do $M'=PM$, quy tắc đạo hàm tích cho

\[
(Mu)_x=M(u_x+Pu)=MQ,\qquad
u=M^{-1}\left[\int MQ\,dx+C(y)\right].
\]

Ghi rõ $P$, $M$ và cách tính tích phân trong từng bài. Với ODE cấp 2 có hệ số hằng theo biến đang tích phân, thử $e^{\lambda x}$ và giải phương trình đặc trưng; nghiệm kép cần thêm $xe^{\lambda x}$.

\textbf{Bước 4 — Xét trường hợp suy biến trước khi chia.} Chia cho $x$, $y$ hoặc $u$ chỉ hợp lệ ở nơi chúng khác 0. Trên đường bị loại, dùng phương trình gốc để tìm ràng buộc. Nếu cần nghiệm trơn qua đường đó, xét giới hạn và đạo hàm của công thức vừa tìm.

\textbf{Bước 5 — Dùng dữ liệu và kiểm tra.} Thế điều kiện đầu/biên để xác định các hàm tùy ý, rồi lấy đạo hàm của nghiệm và thay lại phương trình. Nêu rõ miền và độ trơn cần thiết. Nếu chỉ tìm nghiệm dạng mũ, đó là một họ nghiệm thử; không tự gọi là toàn bộ nghiệm tổng quát.
\subsection*{Lời giải áp dụng cho bài}
\textbf{Bước 1 — Chọn biến tích phân.} Phương trình là $u_x=3x^2+y^2$. Giữ $y$ cố định và tích phân theo $x$.

\textbf{Bước 2 — Tính từng tích phân.}
\[u=\int3x^2\,dx+\int y^2\,dx+C(y)=x^3+xy^2+C(y).\]
Hàm $C(y)$ xuất hiện vì đạo hàm theo $x$ của mọi hàm chỉ phụ thuộc $y$ đều bằng 0.

\textbf{Bước 3 — Kiểm tra.} $u_x=3x^2+y^2$ đúng phương trình. Không có dữ liệu bổ sung để xác định $C$; chọn $C$ đủ trơn theo lớp nghiệm yêu cầu.

\clearpage\noindent\texttt{\detokenize{VD_1_2_2_2_02}}\par
\section*{Lời giải bổ sung — Ví dụ 2 · Mục 1.2.2.2}

\subsection*{Kiến thức và trình tự giải}
\textbf{Mục tiêu.} Đưa phương trình về ODE bằng cách giữ cố định các biến không được lấy đạo hàm, hoặc đặt một đạo hàm/hợp thức trung gian.

\textbf{Bước 1 — Chọn biến và tham số.} Với $u_x=f(x,y)$, cố định $y$ và tích phân theo $x$. Khi đó “hằng số” tích phân có thể khác nhau ở mỗi $y$, nên phải viết $C(y)$. Với ba biến, tích phân theo $x$ cho một hàm $C(y,z)$.

\textbf{Bước 2 — Tích phân đúng số lần.} Chẳng hạn $u_{xx}=f$ cho $u_x=\int f\,dx+A(y)$, rồi $u=\int(\int f\,dx)dx+xA(y)+B(y)$. Với đạo hàm hỗn hợp, có thể đặt $v=u_x$ hoặc $v=yu_y+2u$ nếu PDE chứa đúng $v_x$. Mỗi lần tích phân phải giữ hàm tùy ý theo các biến còn lại.

\textbf{Bước 3 — Giải ODE tuyến tính nếu còn hạng chứa ẩn.} Với $u_x+P(x)u=Q(x,y)$, đặt $M(x)=e^{\int P(x)dx}$. Do $M'=PM$, quy tắc đạo hàm tích cho

\[
(Mu)_x=M(u_x+Pu)=MQ,\qquad
u=M^{-1}\left[\int MQ\,dx+C(y)\right].
\]

Ghi rõ $P$, $M$ và cách tính tích phân trong từng bài. Với ODE cấp 2 có hệ số hằng theo biến đang tích phân, thử $e^{\lambda x}$ và giải phương trình đặc trưng; nghiệm kép cần thêm $xe^{\lambda x}$.

\textbf{Bước 4 — Xét trường hợp suy biến trước khi chia.} Chia cho $x$, $y$ hoặc $u$ chỉ hợp lệ ở nơi chúng khác 0. Trên đường bị loại, dùng phương trình gốc để tìm ràng buộc. Nếu cần nghiệm trơn qua đường đó, xét giới hạn và đạo hàm của công thức vừa tìm.

\textbf{Bước 5 — Dùng dữ liệu và kiểm tra.} Thế điều kiện đầu/biên để xác định các hàm tùy ý, rồi lấy đạo hàm của nghiệm và thay lại phương trình. Nêu rõ miền và độ trơn cần thiết. Nếu chỉ tìm nghiệm dạng mũ, đó là một họ nghiệm thử; không tự gọi là toàn bộ nghiệm tổng quát.
\subsection*{Lời giải áp dụng cho bài}
\textbf{Bước 1 — Giảm đạo hàm hỗn hợp.} Từ $u_{xyz}=0$, cố định $y,z$ và tích phân theo $x$, được $u_{yz}=a(y,z)$.

\textbf{Bước 2 — Tích phân theo $y$, rồi $z$.} Chọn một nguyên hàm $b(y,z)$ của $a$ theo $y$, được $u_z=b(y,z)+c(x,z)$. Tích phân theo $z$:
\[u(x,y,z)=B(y,z)+C(x,z)+A(x,y).\]
Mỗi hàm tùy ý bỏ một biến; điều này giữ đủ các hằng số tích phân của ba lần lấy nguyên hàm. Công thức có ý nghĩa trên miền tích hoặc cục bộ nơi các phép tích phân được thực hiện.

\textbf{Bước 3 — Kiểm tra.} $A$ không phụ thuộc $z$, $B$ không phụ thuộc $x$, $C$ không phụ thuộc $y$, nên đạo hàm $xyz$ của từng hạng bằng 0. Vì vậy tổng thỏa PDE.

\clearpage\noindent\texttt{\detokenize{VD_1_2_2_2_03}}\par
\section*{Lời giải bổ sung — Ví dụ 3 · Mục 1.2.2.2}

\subsection*{Kiến thức và trình tự giải}
\textbf{Mục tiêu.} Đưa phương trình về ODE bằng cách giữ cố định các biến không được lấy đạo hàm, hoặc đặt một đạo hàm/hợp thức trung gian.

\textbf{Bước 1 — Chọn biến và tham số.} Với $u_x=f(x,y)$, cố định $y$ và tích phân theo $x$. Khi đó “hằng số” tích phân có thể khác nhau ở mỗi $y$, nên phải viết $C(y)$. Với ba biến, tích phân theo $x$ cho một hàm $C(y,z)$.

\textbf{Bước 2 — Tích phân đúng số lần.} Chẳng hạn $u_{xx}=f$ cho $u_x=\int f\,dx+A(y)$, rồi $u=\int(\int f\,dx)dx+xA(y)+B(y)$. Với đạo hàm hỗn hợp, có thể đặt $v=u_x$ hoặc $v=yu_y+2u$ nếu PDE chứa đúng $v_x$. Mỗi lần tích phân phải giữ hàm tùy ý theo các biến còn lại.

\textbf{Bước 3 — Giải ODE tuyến tính nếu còn hạng chứa ẩn.} Với $u_x+P(x)u=Q(x,y)$, đặt $M(x)=e^{\int P(x)dx}$. Do $M'=PM$, quy tắc đạo hàm tích cho

\[
(Mu)_x=M(u_x+Pu)=MQ,\qquad
u=M^{-1}\left[\int MQ\,dx+C(y)\right].
\]

Ghi rõ $P$, $M$ và cách tính tích phân trong từng bài. Với ODE cấp 2 có hệ số hằng theo biến đang tích phân, thử $e^{\lambda x}$ và giải phương trình đặc trưng; nghiệm kép cần thêm $xe^{\lambda x}$.

\textbf{Bước 4 — Xét trường hợp suy biến trước khi chia.} Chia cho $x$, $y$ hoặc $u$ chỉ hợp lệ ở nơi chúng khác 0. Trên đường bị loại, dùng phương trình gốc để tìm ràng buộc. Nếu cần nghiệm trơn qua đường đó, xét giới hạn và đạo hàm của công thức vừa tìm.

\textbf{Bước 5 — Dùng dữ liệu và kiểm tra.} Thế điều kiện đầu/biên để xác định các hàm tùy ý, rồi lấy đạo hàm của nghiệm và thay lại phương trình. Nêu rõ miền và độ trơn cần thiết. Nếu chỉ tìm nghiệm dạng mũ, đó là một họ nghiệm thử; không tự gọi là toàn bộ nghiệm tổng quát.
\subsection*{Lời giải áp dụng cho bài}
\textbf{Bước 1 — Xét hệ số trước khi chia.} Với $yu_x=x^2y$, trên $y\ne0$ chia cho $y$ được $u_x=x^2$. Trên $y=0$ phương trình gốc là $0=0$.

\textbf{Bước 2 — Tích phân theo $x$.} Giữ $y$ cố định:
\[u(x,y)=\int x^2\,dx+C(y)=x^3/3+C(y).\]

\textbf{Bước 3 — Kéo dài và kiểm tra.} Nếu miền chứa trục $y=0$ và yêu cầu nghiệm $C^1$, giá trị trên trục phải là giới hạn của công thức ở hai phía; không gán độc lập một hàm khác trên trục. Với $C$ đủ trơn, $yu_x=yx^2$ đúng cả ở $y=0$.

\clearpage\noindent\texttt{\detokenize{VD_1_2_2_2_04}}\par
\section*{Lời giải bổ sung — Ví dụ 4 · Mục 1.2.2.2}

\subsection*{Kiến thức và trình tự giải}
\textbf{Mục tiêu.} Đưa phương trình về ODE bằng cách giữ cố định các biến không được lấy đạo hàm, hoặc đặt một đạo hàm/hợp thức trung gian.

\textbf{Bước 1 — Chọn biến và tham số.} Với $u_x=f(x,y)$, cố định $y$ và tích phân theo $x$. Khi đó “hằng số” tích phân có thể khác nhau ở mỗi $y$, nên phải viết $C(y)$. Với ba biến, tích phân theo $x$ cho một hàm $C(y,z)$.

\textbf{Bước 2 — Tích phân đúng số lần.} Chẳng hạn $u_{xx}=f$ cho $u_x=\int f\,dx+A(y)$, rồi $u=\int(\int f\,dx)dx+xA(y)+B(y)$. Với đạo hàm hỗn hợp, có thể đặt $v=u_x$ hoặc $v=yu_y+2u$ nếu PDE chứa đúng $v_x$. Mỗi lần tích phân phải giữ hàm tùy ý theo các biến còn lại.

\textbf{Bước 3 — Giải ODE tuyến tính nếu còn hạng chứa ẩn.} Với $u_x+P(x)u=Q(x,y)$, đặt $M(x)=e^{\int P(x)dx}$. Do $M'=PM$, quy tắc đạo hàm tích cho

\[
(Mu)_x=M(u_x+Pu)=MQ,\qquad
u=M^{-1}\left[\int MQ\,dx+C(y)\right].
\]

Ghi rõ $P$, $M$ và cách tính tích phân trong từng bài. Với ODE cấp 2 có hệ số hằng theo biến đang tích phân, thử $e^{\lambda x}$ và giải phương trình đặc trưng; nghiệm kép cần thêm $xe^{\lambda x}$.

\textbf{Bước 4 — Xét trường hợp suy biến trước khi chia.} Chia cho $x$, $y$ hoặc $u$ chỉ hợp lệ ở nơi chúng khác 0. Trên đường bị loại, dùng phương trình gốc để tìm ràng buộc. Nếu cần nghiệm trơn qua đường đó, xét giới hạn và đạo hàm của công thức vừa tìm.

\textbf{Bước 5 — Dùng dữ liệu và kiểm tra.} Thế điều kiện đầu/biên để xác định các hàm tùy ý, rồi lấy đạo hàm của nghiệm và thay lại phương trình. Nêu rõ miền và độ trơn cần thiết. Nếu chỉ tìm nghiệm dạng mũ, đó là một họ nghiệm thử; không tự gọi là toàn bộ nghiệm tổng quát.
\subsection*{Lời giải áp dụng cho bài}
\textbf{Bước 1 — Đưa về ODE.} Với $u_x=2u$, coi $y$ là tham số và viết $u_x-2u=0$.

\textbf{Bước 2 — Dùng thừa số tích phân.} $M=e^{-2x}$ và
\[(e^{-2x}u)_x=e^{-2x}(u_x-2u)=0.\]
Tích phân cho $e^{-2x}u=C(y)$, tức $u=e^{2x}C(y)$. Cách này giữ cả nghiệm $u=0$, không cần chia cho $u$.

\textbf{Bước 3 — Kiểm tra.} $u_x=2e^{2x}C(y)=2u$. Chưa có dữ liệu đầu nên $C$ tùy ý đủ trơn.

\clearpage\noindent\texttt{\detokenize{VD_1_2_2_2_05}}\par
\section*{Lời giải bổ sung — Ví dụ 5 · Mục 1.2.2.2}

\subsection*{Kiến thức và trình tự giải}
\textbf{Mục tiêu.} Đưa phương trình về ODE bằng cách giữ cố định các biến không được lấy đạo hàm, hoặc đặt một đạo hàm/hợp thức trung gian.

\textbf{Bước 1 — Chọn biến và tham số.} Với $u_x=f(x,y)$, cố định $y$ và tích phân theo $x$. Khi đó “hằng số” tích phân có thể khác nhau ở mỗi $y$, nên phải viết $C(y)$. Với ba biến, tích phân theo $x$ cho một hàm $C(y,z)$.

\textbf{Bước 2 — Tích phân đúng số lần.} Chẳng hạn $u_{xx}=f$ cho $u_x=\int f\,dx+A(y)$, rồi $u=\int(\int f\,dx)dx+xA(y)+B(y)$. Với đạo hàm hỗn hợp, có thể đặt $v=u_x$ hoặc $v=yu_y+2u$ nếu PDE chứa đúng $v_x$. Mỗi lần tích phân phải giữ hàm tùy ý theo các biến còn lại.

\textbf{Bước 3 — Giải ODE tuyến tính nếu còn hạng chứa ẩn.} Với $u_x+P(x)u=Q(x,y)$, đặt $M(x)=e^{\int P(x)dx}$. Do $M'=PM$, quy tắc đạo hàm tích cho

\[
(Mu)_x=M(u_x+Pu)=MQ,\qquad
u=M^{-1}\left[\int MQ\,dx+C(y)\right].
\]

Ghi rõ $P$, $M$ và cách tính tích phân trong từng bài. Với ODE cấp 2 có hệ số hằng theo biến đang tích phân, thử $e^{\lambda x}$ và giải phương trình đặc trưng; nghiệm kép cần thêm $xe^{\lambda x}$.

\textbf{Bước 4 — Xét trường hợp suy biến trước khi chia.} Chia cho $x$, $y$ hoặc $u$ chỉ hợp lệ ở nơi chúng khác 0. Trên đường bị loại, dùng phương trình gốc để tìm ràng buộc. Nếu cần nghiệm trơn qua đường đó, xét giới hạn và đạo hàm của công thức vừa tìm.

\textbf{Bước 5 — Dùng dữ liệu và kiểm tra.} Thế điều kiện đầu/biên để xác định các hàm tùy ý, rồi lấy đạo hàm của nghiệm và thay lại phương trình. Nêu rõ miền và độ trơn cần thiết. Nếu chỉ tìm nghiệm dạng mũ, đó là một họ nghiệm thử; không tự gọi là toàn bộ nghiệm tổng quát.
\subsection*{Lời giải áp dụng cho bài}
\textbf{Bước 1 — Nhận dạng ODE tuyến tính.} Phương trình $u_x+2xu=4xy$ có $P(x)=2x$, nên $M=e^{\int2x\,dx}=e^{x^2}$.

\textbf{Bước 2 — Nhân và tích phân.}
\[(e^{x^2}u)_x=4xy e^{x^2}.\]
Giữ $y$ cố định, dùng $d(e^{x^2})=2xe^{x^2}dx$:
\[\int4xy e^{x^2}dx=2y e^{x^2},\quad e^{x^2}u=2y e^{x^2}+C(y).\]
Suy ra $u=2y+e^{-x^2}C(y)$.

\textbf{Bước 3 — Kiểm tra.} $u_x=-2xe^{-x^2}C(y)$, nên $u_x+2xu=4xy$.

\clearpage\noindent\texttt{\detokenize{VD_1_2_2_2_06}}\par
\section*{Lời giải bổ sung — Ví dụ 6 · Mục 1.2.2.2}

\subsection*{Kiến thức và trình tự giải}
\textbf{Mục tiêu.} Đưa phương trình về ODE bằng cách giữ cố định các biến không được lấy đạo hàm, hoặc đặt một đạo hàm/hợp thức trung gian.

\textbf{Bước 1 — Chọn biến và tham số.} Với $u_x=f(x,y)$, cố định $y$ và tích phân theo $x$. Khi đó “hằng số” tích phân có thể khác nhau ở mỗi $y$, nên phải viết $C(y)$. Với ba biến, tích phân theo $x$ cho một hàm $C(y,z)$.

\textbf{Bước 2 — Tích phân đúng số lần.} Chẳng hạn $u_{xx}=f$ cho $u_x=\int f\,dx+A(y)$, rồi $u=\int(\int f\,dx)dx+xA(y)+B(y)$. Với đạo hàm hỗn hợp, có thể đặt $v=u_x$ hoặc $v=yu_y+2u$ nếu PDE chứa đúng $v_x$. Mỗi lần tích phân phải giữ hàm tùy ý theo các biến còn lại.

\textbf{Bước 3 — Giải ODE tuyến tính nếu còn hạng chứa ẩn.} Với $u_x+P(x)u=Q(x,y)$, đặt $M(x)=e^{\int P(x)dx}$. Do $M'=PM$, quy tắc đạo hàm tích cho

\[
(Mu)_x=M(u_x+Pu)=MQ,\qquad
u=M^{-1}\left[\int MQ\,dx+C(y)\right].
\]

Ghi rõ $P$, $M$ và cách tính tích phân trong từng bài. Với ODE cấp 2 có hệ số hằng theo biến đang tích phân, thử $e^{\lambda x}$ và giải phương trình đặc trưng; nghiệm kép cần thêm $xe^{\lambda x}$.

\textbf{Bước 4 — Xét trường hợp suy biến trước khi chia.} Chia cho $x$, $y$ hoặc $u$ chỉ hợp lệ ở nơi chúng khác 0. Trên đường bị loại, dùng phương trình gốc để tìm ràng buộc. Nếu cần nghiệm trơn qua đường đó, xét giới hạn và đạo hàm của công thức vừa tìm.

\textbf{Bước 5 — Dùng dữ liệu và kiểm tra.} Thế điều kiện đầu/biên để xác định các hàm tùy ý, rồi lấy đạo hàm của nghiệm và thay lại phương trình. Nêu rõ miền và độ trơn cần thiết. Nếu chỉ tìm nghiệm dạng mũ, đó là một họ nghiệm thử; không tự gọi là toàn bộ nghiệm tổng quát.
\subsection*{Lời giải áp dụng cho bài}
\textbf{Bước 1 — Gom đạo hàm hỗn hợp.} Với $yu_{xy}+2u_x=x$, đặt $v=yu_y+2u$. Vì $y$ không đổi khi lấy đạo hàm theo $x$, $v_x=yu_{xy}+2u_x=x$.

\textbf{Bước 2 — Tích phân theo $x$.} $v=x^2/2+A(y)$, tức $yu_y+2u=x^2/2+A(y)$.

\textbf{Bước 3 — Giải theo $y$ trên $y\ne0$.} Chia cho $y$ rồi dùng thừa số $y^2$:
\[(y^2u)_y=yx^2/2+yA(y).\]
Tích phân cho $y^2u=x^2y^2/4+\int yA(y)dy+C(x)$. Đặt $B(y)=y^{-2}\int yA(y)dy$, ta được
\[u=x^2/4+B(y)+C(x)/y^2.\]

\textbf{Bước 4 — Kiểm tra và trục suy biến.} $u_x=x/2+C'(x)/y^2$, $u_{xy}=-2C'(x)/y^3$, nên $yu_{xy}+2u_x=x$. Muốn kéo dài qua $y=0$, tổ hợp $B(y)+C(x)/y^2$ phải đủ trơn; không cho phép chọn tùy ý các hạng kỳ dị. Phương trình gốc trên trục yêu cầu $2u_x(x,0)=x$.

\clearpage\noindent\texttt{\detokenize{BT_1_2_3_01}}\par
\section*{Lời giải bổ sung — Bài tập 1 · Mục 1.2.3}

\subsection*{Kiến thức và trình tự giải}
\textbf{Mục tiêu.} Xác định cấp và dạng của phương trình trước khi chọn cách giải. Ở mục này, từ “bậc” trong đề được hiểu theo nghĩa cấp đạo hàm của phương trình.

\textbf{Bước 1 — Liệt kê đạo hàm và đếm cấp.} Với $\partial^{i+j}u/\partial x^i\partial y^j$, cấp là $i+j$. Vì vậy $u_x$ cấp 1, $u_{xy}$ cấp 2, $u_{txx}$ cấp 3. Hạng $u$, $e^u$, $\log u$ không chứa đạo hàm, được xem là cấp 0. $(u_y)^2$ là bình phương của một đạo hàm cấp 1, không phải đạo hàm cấp 2. Lấy giá trị lớn nhất trong các cấp vừa liệt kê.

\textbf{Bước 2 — Viết rõ phần chứa đạo hàm cấp cao nhất.} Gọi cấp đó là $m$. Tạm xem từng đạo hàm cấp $m$ như một ẩn đại số độc lập. Nếu chúng bị bình phương, nhân với nhau hoặc nằm trong hàm mũ/lượng giác, biểu thức đã cho hoàn toàn phi tuyến. Nếu chúng chỉ xuất hiện tuyến tính, chuyển sang bước 3. Khi khai triển $(f(u))_x$, phải dùng $(f(u))_x=f'(u)u_x$ trước khi phân loại.

\textbf{Bước 3 — Xét hệ số và các hạng cấp thấp.} Nếu mọi hạng chứa $u$ và đạo hàm đều tuyến tính, hệ số chỉ phụ thuộc biến độc lập, phương trình tuyến tính. Nếu chỉ phần cấp $m$ tuyến tính, hệ số của phần đó phụ thuộc biến độc lập, còn phần cấp thấp phi tuyến, phương trình nửa tuyến tính. Nếu hệ số phần cấp $m$ phụ thuộc $u$ hoặc đạo hàm cấp thấp, phương trình tựa tuyến tính. Khi ghi kết luận, chỉ rõ hạng nào quyết định phân loại.

\textbf{Bước 4 — Xét thuần nhất khi phương trình tuyến tính.} Gom toàn bộ hạng chứa ẩn về một vế, viết $Lu=f(x,y)$. Nếu $f\equiv0$, phương trình thuần nhất; nếu không, không thuần nhất. Ví dụ $u_{xy}=u$ phải viết $u_{xy}-u=0$, nên là thuần nhất.

\textbf{Bước 5 — Kiểm tra miền và yêu cầu đề.} $\log u$ đòi hỏi $u>0$ đối với nghiệm thực. Một phép biến đổi như lấy log có thể cho phương trình tương đương trên một miền, nhưng cần phân biệt dạng biểu thức in trong đề và dạng đã biến đổi. Nếu đề chỉ hỏi cấp, dừng sau bước 1 và giải thích đạo hàm quyết định cấp.
\subsection*{Lời giải áp dụng cho bài}
\textbf{Bước 1: xác định điều cần phân loại.} Theo phần lý thuyết trong sách, ``bậc'' của câu này được hiểu là \emph{cấp phương trình}: cấp cao nhất của đạo hàm riêng xuất hiện. Số mũ của đạo hàm và hàm như $\log u,e^u$ không làm tăng cấp đạo hàm.

\textbf{Bước 2: xét từng phương trình.}
(a) $u_{xx}+u_{yy}=0$: hai đạo hàm đều cấp 2, nên cấp 2.
(b) $u_{xxx}+u_{xy}+a(x)u_y+\log u=f(x,y)$: $u_{xxx}$ cấp 3, $u_{xy}$ cấp 2, $u_y$ cấp 1; $\log u$ cấp 0. Cấp cao nhất bằng 3. Nếu nghiệm thực, $\log u$ cần $u>0$ nhưng không thay đổi cấp.
(c) $u_x+cu_y=d$: các đạo hàm đều cấp 1, nên cấp 1.
(d) $uu_{xx}+u_y^2+e^u=0$: $u_{xx}$ cấp 2; $(u_y)^2$ vẫn chỉ là bình phương của đạo hàm cấp 1. Do đó cấp 2.

\textbf{Kết luận:} các cấp lần lượt là $\boxed{2,3,1,2}$.

\clearpage\noindent\texttt{\detokenize{BT_1_2_3_02}}\par
\section*{Lời giải bổ sung — Bài tập 2 · Mục 1.2.3}

\subsection*{Kiến thức và trình tự giải}
\textbf{Mục tiêu.} Xác định cấp và dạng của phương trình trước khi chọn cách giải. Ở mục này, từ “bậc” trong đề được hiểu theo nghĩa cấp đạo hàm của phương trình.

\textbf{Bước 1 — Liệt kê đạo hàm và đếm cấp.} Với $\partial^{i+j}u/\partial x^i\partial y^j$, cấp là $i+j$. Vì vậy $u_x$ cấp 1, $u_{xy}$ cấp 2, $u_{txx}$ cấp 3. Hạng $u$, $e^u$, $\log u$ không chứa đạo hàm, được xem là cấp 0. $(u_y)^2$ là bình phương của một đạo hàm cấp 1, không phải đạo hàm cấp 2. Lấy giá trị lớn nhất trong các cấp vừa liệt kê.

\textbf{Bước 2 — Viết rõ phần chứa đạo hàm cấp cao nhất.} Gọi cấp đó là $m$. Tạm xem từng đạo hàm cấp $m$ như một ẩn đại số độc lập. Nếu chúng bị bình phương, nhân với nhau hoặc nằm trong hàm mũ/lượng giác, biểu thức đã cho hoàn toàn phi tuyến. Nếu chúng chỉ xuất hiện tuyến tính, chuyển sang bước 3. Khi khai triển $(f(u))_x$, phải dùng $(f(u))_x=f'(u)u_x$ trước khi phân loại.

\textbf{Bước 3 — Xét hệ số và các hạng cấp thấp.} Nếu mọi hạng chứa $u$ và đạo hàm đều tuyến tính, hệ số chỉ phụ thuộc biến độc lập, phương trình tuyến tính. Nếu chỉ phần cấp $m$ tuyến tính, hệ số của phần đó phụ thuộc biến độc lập, còn phần cấp thấp phi tuyến, phương trình nửa tuyến tính. Nếu hệ số phần cấp $m$ phụ thuộc $u$ hoặc đạo hàm cấp thấp, phương trình tựa tuyến tính. Khi ghi kết luận, chỉ rõ hạng nào quyết định phân loại.

\textbf{Bước 4 — Xét thuần nhất khi phương trình tuyến tính.} Gom toàn bộ hạng chứa ẩn về một vế, viết $Lu=f(x,y)$. Nếu $f\equiv0$, phương trình thuần nhất; nếu không, không thuần nhất. Ví dụ $u_{xy}=u$ phải viết $u_{xy}-u=0$, nên là thuần nhất.

\textbf{Bước 5 — Kiểm tra miền và yêu cầu đề.} $\log u$ đòi hỏi $u>0$ đối với nghiệm thực. Một phép biến đổi như lấy log có thể cho phương trình tương đương trên một miền, nhưng cần phân biệt dạng biểu thức in trong đề và dạng đã biến đổi. Nếu đề chỉ hỏi cấp, dừng sau bước 1 và giải thích đạo hàm quyết định cấp.
\subsection*{Lời giải áp dụng cho bài}
\textbf{Bước 1: dùng định nghĩa theo cấp cao nhất $m$.} Tuyến tính: ẩn và mọi đạo hàm xuất hiện tuyến tính, hệ số chỉ phụ thuộc biến độc lập. Nửa tuyến tính: đạo hàm cấp $m$ tuyến tính với hệ số chỉ phụ thuộc biến độc lập, nhưng các hạng cấp thấp có thể phi tuyến. Tựa tuyến tính: đạo hàm cấp $m$ tuyến tính, hệ số có thể phụ thuộc $u$ và các đạo hàm cấp thấp hơn. Hoàn toàn phi tuyến: biểu thức không tuyến tính theo đạo hàm cấp $m$. Chỉ hỏi thuần nhất khi phương trình đã tuyến tính.

\textbf{Bước 2: áp dụng từng ý.}
(a) $e^{\Delta u}=x$. Vì $\Delta u=u_{xx}+u_{yy}$ nằm trong hàm mũ, biểu thức đã cho hoàn toàn phi tuyến cấp 2. Tuy nhiên, nghiệm thực chỉ có trên $x>0$, và ở đó lấy log được phương trình tương đương $\Delta u=\log x$, tuyến tính không thuần nhất. Cần phân biệt phân loại \emph{biểu thức được in} và dạng tương đương sau biến đổi.

(b) $u_t+u_{txx}+uu_x=0$. Đạo hàm cao nhất $u_{txx}$ cấp 3, tuyến tính với hệ số 1. Tích $uu_x$ phi tuyến ở cấp thấp. Kết luận: nửa tuyến tính cấp 3.

(c) $u_t+(f(u))_x=0$. Quy tắc dây chuyền cho $u_t+f'(u)u_x=0$. Hạng cấp 1 tuyến tính theo $u_t,u_x$ nhưng hệ số $f'(u)$ phụ thuộc $u$, nên tựa tuyến tính cấp 1. Ngoại lệ: nếu $f(u)=au+b$, được $u_t+au_x=0$, tuyến tính thuần nhất.

(d) $u_{xx}+u_{yy}-2u=x^2$. Các hạng của $u$ tuyến tính, hệ số hằng, vế phải khác 0: tuyến tính không thuần nhất cấp 2.

(e) $u_{xy}=u$. Chuyển vế thành $u_{xy}-u=0$: tuyến tính thuần nhất cấp 2. Việc vế phải ban đầu là $u$ không làm phương trình không thuần nhất; phải gom tất cả hạng chứa ẩn về cùng vế.

(f) $uu_x+xu_y=0$. Tuyến tính theo hai đạo hàm cấp 1, nhưng hệ số của $u_x$ là $u$. Kết luận: tựa tuyến tính cấp 1.

(g) $u_x(1+u_y)=u_{xx}$. Viết $u_{xx}-u_x-u_xu_y=0$. Đạo hàm cấp 2 $u_{xx}$ tuyến tính, hệ số 1; phi tuyến $u_xu_y$ chỉ ở cấp 1. Kết luận: nửa tuyến tính cấp 2.

(h) $(\sin u_x)u_x+u_y=e^x$. Đạo hàm cấp cao nhất là cấp 1 và $u_x\sin u_x$ không tuyến tính theo $u_x$. Kết luận: hoàn toàn phi tuyến cấp 1.

\textbf{Bước 3: kiểm tra cách phân biệt.} Các ý (b),(g) có phi tuyến nhưng chỉ ở cấp thấp, nên nửa tuyến tính. Các ý (c),(f) có hệ số đạo hàm cấp cao nhất phụ thuộc ẩn, nên tựa tuyến tính. Các ý (a),(h) phi tuyến ngay ở đạo hàm cấp cao nhất. Với các phương trình phi tuyến không dùng phân biệt ``thuần nhất/không thuần nhất'' theo định nghĩa toán tử tuyến tính.

\clearpage\noindent\texttt{\detokenize{BT_1_2_3_03}}\par
\section*{Lời giải bổ sung — Bài tập 3 · Mục 1.2.3}

\subsection*{Kiến thức và trình tự giải}
\textbf{Mục tiêu.} Tìm nghiệm dạng mũ cho PDE tuyến tính có hệ số hằng, theo dạng thử mà đề yêu cầu.

\textbf{Bước 1 — Đặt đủ các biến trong số mũ.} Hai biến dùng $u=e^{rx+sy}$; ba biến dùng $u=e^{rx+sy+tz}$. Các tham số $r,s,t$ cần được xác định, không phải biến độc lập.

\textbf{Bước 2 — Tính các đạo hàm.} Ta có $u_x=ru$, $u_y=su$, $u_{xx}=r^2u$, $u_{xy}=rsu$, $u_{xyz}=rst\,u$. Mỗi lần lấy đạo hàm tương ứng với nhân thêm tham số của biến ấy.

\textbf{Bước 3 — Thế và giải phương trình đại số.} Vì hàm mũ không bằng 0, chia PDE cho $u$ để thu được điều kiện theo tham số. Giải hoặc tham số hóa điều kiện này. Phân biệt tham số thực và phức khi phương trình đại số không có nghiệm thực khác 0.

\textbf{Bước 4 — Kiểm tra và hiểu phạm vi kết quả.} Thế một bộ tham số cụ thể để kiểm tra. Với phương trình tuyến tính thuần nhất, tổng hữu hạn các nghiệm cũng là nghiệm. Họ nghiệm mũ chưa nhất thiết mô tả toàn bộ nghiệm; nếu muốn nghiệm tổng quát phải dùng thêm đặc trưng hoặc phương pháp phù hợp.
\subsection*{Lời giải áp dụng cho bài}
\textbf{Bước 1: tính đạo hàm của nghiệm thử.} Với $u=e^{rx+sy}$,
\[u_x=ru,\quad u_y=su,\quad u_{xx}=r^2u,\quad u_{xy}=rsu,\quad u_{yy}=s^2u.\]
Vì $u\ne0$, sau khi thế vào PDE có thể chia cho $u$ để có phương trình đại số theo $r,s$. Đề dùng nghiệm dạng đặc biệt; một họ nghiệm mũ không tự động là nghiệm tổng quát của mọi PDE.

(a) $2u_x+3u_y-2u=0$ cho $2r+3s-2=0$. Với mỗi $s\in\mathbb R$, chọn $r=1-3s/2$, được một nghiệm $e^{(1-3s/2)x+sy}$. Để có nghiệm tổng quát, dùng bất biến $\xi=3x-2y$, đặt $u=e^xF(\xi)$. Tính đạo hàm cho $2u_x+3u_y=2u$, nên toàn bộ nghiệm là $u=e^xF(3x-2y)$.

(b) Vì PDE có ba biến, dùng $u=e^{rx+sy+tz}$. Khi đó $u_{xyz}=rst\,u$, nên điều kiện là $rst=1$. Chọn $r,s\ne0$ tùy ý và $t=1/(rs)$; ví dụ $u=e^{x+y+z}$. Tuyến tính cho phép cộng hữu hạn các mode này. Không thể bỏ biến $z$ vì nghiệm không phụ thuộc $z$ có $u_{xyz}=0$.

(c) $4u_{xx}-4u_{xy}+4u_{yy}=0$ cho $r^2-rs+s^2=0$. Viết $(r-s/2)^2+3s^2/4=0$: với $r,s$ thực chỉ có $r=s=0$, tức mode hằng. Với số phức, $r=(1\pm i\sqrt3)s/2$; phần thực và phần ảo của mode phức cho nghiệm thực. Nghiệm mũ thực không mô tả hết nghiệm điều hòa của phương trình elliptic này.

(d) $u_{xx}+u_{yy}=u$ cho $r^2+s^2=1$. Tham số hóa $r=\cos\theta,s=\sin\theta$, được $u_\theta=e^{x\cos\theta+y\sin\theta}$. Các ví dụ $e^x,e^y,e^{(x+y)/\sqrt2}$ đều đúng; cộng các nghiệm do tính tuyến tính vẫn cho nghiệm. Thế các đạo hàm ở bước 1 kiểm tra trực tiếp từng công thức.

\textbf{Chi tiết tìm họ nghiệm ở (a).} Đặc trưng là $X'=2,Y'=3,U'=2U$. Tích phân từ $(x_0,y_0,u_0)$ được $X=x_0+2s,Y=y_0+3s,U=u_0e^{2s}$. Vì $3X-2Y=3x_0-2y_0$ và $e^{-X}U=e^{-x_0}u_0$, đặt giá trị bất biến thứ hai là $F(3X-2Y)$. Đổi về $x,y$ được $u=e^xF(3x-2y)$. Kiểm tra bằng các đạo hàm $u_x=e^x(F+3F')$, $u_y=-2e^xF'$: $2u_x+3u_y-2u=e^x(2F+6F'-6F'-2F)=0$.

\clearpage\noindent\texttt{\detokenize{BT_1_2_3_04}}\par
\section*{Lời giải bổ sung — Bài tập 4 · Mục 1.2.3}

\subsection*{Kiến thức và trình tự giải}
\textbf{Mục tiêu của bài.} Đưa $au_{xx}+bu_{xy}+cu_{yy}=0$ với $\Delta=b^2-4ac>0$ về phương trình có thể tích phân. Ở đây hệ số hỗn hợp là $b$, không phải $2b$.

\textbf{Bước 1 — Tìm hướng đặc trưng.} Thử tọa độ $\xi=y-mx$. Hệ số của $v_{\xi\xi}$ sau đổi biến là $am^2-bm+c$. Giải $am^2-bm+c=0$; nếu $a\ne0$, $\Delta>0$ cho hai nghiệm thực phân biệt $m_1,m_2$. Nếu $a=0$, xét trực tiếp vì khi đó $b\ne0$.

\textbf{Bước 2 — Đổi biến và kiểm tra Jacobian.} Đặt $\xi=y-m_1x$, $\eta=y-m_2x$, $u=v(\xi,\eta)$. Jacobian là $m_2-m_1\ne0$. Viết $\partial_x=-m_1\partial_\xi-m_2\partial_\eta$, $\partial_y=\partial_\xi+\partial_\eta$, rồi áp dụng hai lần để tính $u_{xx},u_{xy},u_{yy}$.

\textbf{Bước 3 — Thế toàn bộ đạo hàm vào PDE.} Hai hệ số thuần bằng 0 theo phương trình của $m_i$. Hệ số hỗn hợp khác 0, nên chia được để có $v_{\xi\eta}=0$. Với đổi biến không tuyến tính trong bài khác, phải giữ cả các hạng đạo hàm cấp 1 phát sinh.

\textbf{Bước 4 — Tích phân hai lần.} $\partial_\eta(v_\xi)=0$ cho $v_\xi=A(\xi)$; tích phân theo $\xi$ cho $v=F(\xi)+G(\eta)$. Đổi về $x,y$ và thế lại PDE để kiểm tra. Các hàm $F,G$ cần đủ trơn, thường $C^2$.
\subsection*{Lời giải áp dụng cho bài}
\textbf{Bước 1: tìm hai họ đặc trưng.} Nếu $a\ne0$, phương trình $am^2-bm+c=0$ có hai nghiệm thực phân biệt $m_1,m_2$ vì $\Delta=b^2-4ac>0$. Ta có $m_1+m_2=b/a$, $m_1m_2=c/a$.

\textbf{Bước 2: đổi biến.} Đặt $\xi=y-m_1x,\eta=y-m_2x,u=v(\xi,\eta)$. Jacobian bằng $m_2-m_1\ne0$, nên phép đổi không suy biến. Quy tắc dây chuyền:
\begin{align*}
u_{xx}&=m_1^2v_{\xi\xi}+2m_1m_2v_{\xi\eta}+m_2^2v_{\eta\eta},\\
u_{xy}&=-m_1v_{\xi\xi}-(m_1+m_2)v_{\xi\eta}-m_2v_{\eta\eta},\\
u_{yy}&=v_{\xi\xi}+2v_{\xi\eta}+v_{\eta\eta}.
\end{align*}
Hai hệ số của $v_{\xi\xi},v_{\eta\eta}$ bằng 0. Hệ số hỗn hợp bằng $2am_1m_2-b(m_1+m_2)+2c=-\Delta/a\ne0$. Do đó PDE tương đương $v_{\xi\eta}=0$.

\textbf{Bước 3: tích phân.} Từ $\partial_\eta(v_\xi)=0$, $v_\xi=A(\xi)$. Tích phân theo $\xi$ cho $v=F(\xi)+G(\eta)$. Suy
\[\boxed{u=F(y-m_1x)+G(y-m_2x).}\]
Đây là dạng yêu cầu với $\alpha=-m_1,\beta=1,\gamma=-m_2,\delta=1$.

\textbf{Trường hợp $a=0$.} Khi đó $b\ne0$, PDE là $\partial_y(bu_x+cu_y)=0$. Đặt $\xi=x,\eta=y-(c/b)x$, tính được $bu_{xy}+cu_{yy}=b v_{\xi\eta}$. Vẫn có $u=F(x)+G(y-(c/b)x)$, hai dạng tuyến tính độc lập.

\textbf{Kiểm tra:} thế đạo hàm của $F(y-m_ix)$ cho hệ số $am_i^2-bm_i+c=0$; từng hạng giải PDE, và tổng cũng giải do tuyến tính.

\clearpage\noindent\texttt{\detokenize{BT_1_2_3_05}}\par
\section*{Lời giải bổ sung — Bài tập 5 · Mục 1.2.3}

\subsection*{Kiến thức và trình tự giải}
\textbf{Mục tiêu.} Đưa phương trình về ODE bằng cách giữ cố định các biến không được lấy đạo hàm, hoặc đặt một đạo hàm/hợp thức trung gian.

\textbf{Bước 1 — Chọn biến và tham số.} Với $u_x=f(x,y)$, cố định $y$ và tích phân theo $x$. Khi đó “hằng số” tích phân có thể khác nhau ở mỗi $y$, nên phải viết $C(y)$. Với ba biến, tích phân theo $x$ cho một hàm $C(y,z)$.

\textbf{Bước 2 — Tích phân đúng số lần.} Chẳng hạn $u_{xx}=f$ cho $u_x=\int f\,dx+A(y)$, rồi $u=\int(\int f\,dx)dx+xA(y)+B(y)$. Với đạo hàm hỗn hợp, có thể đặt $v=u_x$ hoặc $v=yu_y+2u$ nếu PDE chứa đúng $v_x$. Mỗi lần tích phân phải giữ hàm tùy ý theo các biến còn lại.

\textbf{Bước 3 — Giải ODE tuyến tính nếu còn hạng chứa ẩn.} Với $u_x+P(x)u=Q(x,y)$, đặt $M(x)=e^{\int P(x)dx}$. Do $M'=PM$, quy tắc đạo hàm tích cho

\[
(Mu)_x=M(u_x+Pu)=MQ,\qquad
u=M^{-1}\left[\int MQ\,dx+C(y)\right].
\]

Ghi rõ $P$, $M$ và cách tính tích phân trong từng bài. Với ODE cấp 2 có hệ số hằng theo biến đang tích phân, thử $e^{\lambda x}$ và giải phương trình đặc trưng; nghiệm kép cần thêm $xe^{\lambda x}$.

\textbf{Bước 4 — Xét trường hợp suy biến trước khi chia.} Chia cho $x$, $y$ hoặc $u$ chỉ hợp lệ ở nơi chúng khác 0. Trên đường bị loại, dùng phương trình gốc để tìm ràng buộc. Nếu cần nghiệm trơn qua đường đó, xét giới hạn và đạo hàm của công thức vừa tìm.

\textbf{Bước 5 — Dùng dữ liệu và kiểm tra.} Thế điều kiện đầu/biên để xác định các hàm tùy ý, rồi lấy đạo hàm của nghiệm và thay lại phương trình. Nêu rõ miền và độ trơn cần thiết. Nếu chỉ tìm nghiệm dạng mũ, đó là một họ nghiệm thử; không tự gọi là toàn bộ nghiệm tổng quát.

\textbf{Bổ sung cho ý tổng bình phương đạo hàm.} Với nghiệm thực, $u_x^2+u_y^2=0$ buộc từng đạo hàm bằng 0, nên $u$ hằng trên mỗi thành phần liên thông. Nếu vế phải là $k<0$, vô nghiệm thực; nếu $k>0$, đây là eikonal $|\nabla u|^2=k$. Nghiệm affine $u=ax+by+c$ cần $a^2+b^2=k$, nhưng chỉ là một họ. Muốn dựng nghiệm cục bộ từ dữ liệu trên đường $\gamma(s)$, chọn gradient ban đầu $P(s)$ thỏa $|P|^2=k$, $P\cdot\gamma'=h'$. Đặc trưng cho $X=\gamma(s)+2\tau P(s)$, $U=h(s)+2k\tau$; chỉ khử được tham số nơi Jacobian khác 0. Cách này được áp dụng riêng ở ý (d).
\subsection*{Lời giải áp dụng cho bài}
(a) Xem $y$ là tham số: $u_x=-u$. Nhân $e^x$ cho $(e^xu)_x=0$, nên $u=e^{-x}C(y)$. Kiểm tra $u_x+u=0$.

(b) Xem $x$ là tham số. ODE theo $y$ là $u_{yy}-x^2u=0$, có phương trình đặc trưng $\lambda^2-x^2=0$. Với $x\ne0$, $u=A(x)e^{xy}+B(x)e^{-xy}$. Với $x=0$, ODE thành $u_{yy}=0$, nên vết là affine theo $y$. Một dạng thống nhất, không mất nghiệm ở $x=0$, là
\[u=A(x)\cosh(xy)+B(x)\frac{\sinh(xy)}x,\]
trong đó $\sinh(xy)/x=y$ tại $x=0$ theo giới hạn. Hệ số $A,B$ đủ trơn cho lớp nghiệm mong muốn.

(c) Đặt $v=u_x$, được $v_x-v=0$, nên $v=B(y)e^x$. Tích phân theo $x$ cho $u=A(y)+B(y)e^x$. Đạo hàm lại cho $u_{xx}=u_x$.

(d) Với nghiệm thực, $u_x^2+u_y^2=0$ là tổng hai số không âm; buộc $u_x=u_y=0$. Do đó $u$ hằng trên từng thành phần liên thông. Nếu vế phải là $k\ne0$: $k<0$ thì vô nghiệm thực; $k>0$ là eikonal $|\nabla u|=\sqrt k$. Một họ nghiệm đơn giản $u=ax+by+c$, $a^2+b^2=k$, nhưng không phải toàn bộ nghiệm.

Để dựng nghiệm cục bộ tổng quát từ dữ liệu $u(\gamma(s))=h(s)$, chọn $P(s)$ thỏa $|P(s)|^2=k$ và $P(s)\cdot\gamma'(s)=h'(s)$. Hệ đặc trưng của $F=p^2+q^2-k$ có $P'=0$, $X'=2P$, $U'=2k$, nên
\[X=\gamma(s)+2tP(s),\qquad U=h(s)+2kt.\]
Khử $(s,t)$ ở nơi Jacobian của $X$ khác 0, được nghiệm $u(X)=U$. Đây là dạng áp dụng phương pháp đặc trưng, phân biệt với chỉ đưa ra ví dụ affine.

\textbf{Chi tiết xử lý điểm $x=0$ ở (b).} Chọn hai hệ số từ các vết tại $y=0$: $A(x)=u(x,0)$, $B(x)=u_y(x,0)$. Khi $x\ne0$, giải hệ $C+D=A$, $x(C-D)=B$ cho $C=(A+B/x)/2,D=(A-B/x)/2$, rồi gom hai hàm mũ thành $A\cosh(xy)+B\sinh(xy)/x$. Khai triển $\sinh(xy)/x=y+x^2y^3/3!+\cdots$ cho giới hạn $y$ và vết $u(0,y)=A(0)+B(0)y$; không loại nghiệm tuyến tính này khi hai nghiệm mũ trùng nhau.

\clearpage\noindent\texttt{\detokenize{VD_1_3_1_2_01}}\par
\section*{Lời giải bổ sung — Ví dụ 1 · Mục 1.3.1.2}

\subsection*{Kiến thức và trình tự giải}
\textbf{Mục tiêu.} Chứng minh tính chồng chất của các nghiệm phương trình tuyến tính thuần nhất.

\textbf{Bước 1 — Viết toán tử.} Gọi $L$ là biểu thức lấy đạo hàm và nhân hệ số của PDE. Chẳng hạn $L=a\partial_x+b\partial_t$. Tuyến tính nghĩa là $L(c_1u_1+c_2u_2)=c_1Lu_1+c_2Lu_2$ với $c_1,c_2$ hằng.

\textbf{Bước 2 — Tính đạo hàm của tổng.} Dùng $(c_1u_1+c_2u_2)_x=c_1(u_1)_x+c_2(u_2)_x$ và tương tự theo các biến khác, rồi nhóm thành $c_1Lu_1+c_2Lu_2$. Không cho các hệ số $c_i$ phụ thuộc vị trí, vì khi đó có thêm đạo hàm của hệ số.

\textbf{Bước 3 — Thế giả thiết nghiệm.} Nếu $Lu_i=0$, tổng có $L=0$ nên cũng là nghiệm. Nếu các nghiệm có cùng dữ liệu biên, phải kiểm tra riêng dữ liệu của tổng; tính chồng chất của PDE chưa tự giữ nguyên dữ liệu.

\textbf{Bước 4 — Mở rộng cho phương trình có nguồn.} Nếu $Lu_i=f$, hiệu $u_1-u_2$ thỏa thuần nhất, nhưng tổng có nguồn $2f$. Để giải $Lu=f$, dùng một nghiệm riêng $u_p$ cộng mọi nghiệm thuần nhất $u_h$.
\subsection*{Lời giải áp dụng cho bài}
\textbf{Bước 1 — Xác định toán tử.} Đặt $L=a\partial_x+b\partial_t$. Do lấy đạo hàm và nhân hệ số là các phép tuyến tính, với $c_1,c_2$ hằng,
\[L(c_1u_1+c_2u_2)=c_1Lu_1+c_2Lu_2.\]

\textbf{Bước 2 — Dùng giả thiết.} Nếu $Lu_1=Lu_2=0$, vế phải bằng 0, nên $c_1u_1+c_2u_2$ cũng là nghiệm. Điều này đúng cho mọi tổng hữu hạn các nghiệm.

\textbf{Bước 3 — Phạm vi chồng chất.} Với PDE có nguồn $Lu=f\ne0$, tổng hai nghiệm cùng nguồn có nguồn $2f$. Khi đó dùng $u=u_p+u_h$ với $Lu_p=f$, $Lu_h=0$; hiệu hai nghiệm cùng nguồn là nghiệm thuần nhất.

\clearpage\noindent\texttt{\detokenize{VD_1_3_1_2_02}}\par
\section*{Lời giải bổ sung — Ví dụ 2 · Mục 1.3.1.2}

\subsection*{Kiến thức và trình tự giải}
\textbf{Mục tiêu.} Với $A(x,y)u_x+B(x,y)u_y+C(x,y)u=R(x,y)$, tìm các đường mà trên đó PDE trở thành một ODE. Dùng chữ $R$ cho vế phải để không nhầm với hàm tùy ý $F$.

\textbf{Bước 1 — Đọc các hệ số và lập hệ đặc trưng.} Đặt $U(s)=u(X(s),Y(s))$. Quy tắc dây chuyền cho $U'=u_xX'+u_yY'$. Chọn $X'=A$, $Y'=B$, khi đó

\[
X'=A(X,Y),\qquad Y'=B(X,Y),\qquad U'+C(X,Y)U=R(X,Y).
\]

Nếu $A\ne0$, có thể dùng $x$ làm tham số: $dy/dx=B/A$, $du/dx+(C/A)u=R/A$. Nếu $A=0$, chọn biến khác hoặc giữ hệ theo $s$; không chia cho hệ số bằng 0.

\textbf{Bước 2 — Tìm bất biến để phân biệt các đặc trưng.} Tích phân hai phương trình đầu để tìm $\rho(X,Y)$ không đổi. Kiểm tra $A\rho_x+B\rho_y=0$. Hệ số hằng cho $\rho=Bx-Ay$. Ví dụ $X'=X,Y'=-Y$ cho $(XY)'=0$; $X'=-Y,Y'=X$ cho $(X^2+Y^2)'=0$. Khi đặc trưng là vòng kín, cần kiểm tra nghiệm sau một vòng trở về cùng giá trị.

\textbf{Bước 3 — Giải ODE của ẩn trên từng đường.} Đặt $M(s)=\exp(\int C(X(s),Y(s))ds)$. Nhân phương trình của $U$ với $M$:

\[
(MU)'=MR,\qquad
U(s)=M(s)^{-1}\left[K+\int M(s)R(X(s),Y(s))ds\right].
\]

$K$ hằng trên một đường, nhưng thay đổi giữa các đường, nên viết $K=F(\rho)$. Nếu tìm được nghiệm riêng $u_p$, có thể giải phương trình thuần nhất cho $u-u_p$ để giảm tính toán. Nếu đổi biến $(\rho,z)$, phải tính Jacobian và dùng quy tắc dây chuyền để thấy PDE chỉ còn đạo hàm theo $z$.

\textbf{Bước 4 — Thế dữ liệu và tìm hàm tùy ý.} Với đường dữ liệu $(x_0(r),y_0(r))$, $u=h(r)$, thế vào nghiệm tổng quát để thu được $F(\rho_0(r))=H(r)$, trong đó $\rho_0=\rho(x_0,y_0)$. Nếu $\rho_0$ khả nghịch, đặt $r=\rho_0^{-1}(z)$ rồi viết $F(z)=H(\rho_0^{-1}(z))$. Nếu hai giá trị $r_1,r_2$ có cùng $\rho_0$, phải có $H(r_1)=H(r_2)$. Nếu đường dữ liệu chính là đặc trưng, dữ liệu phải thỏa ODE trên đường đó và thường không xác định hết $F$.

\textbf{Bước 5 — Xác định miền tồn tại và duy nhất.} Điều kiện không đặc trưng là $x_0'B-y_0'A\ne0$ trên đường dữ liệu. Nó cho khả nghịch cục bộ dưới giả thiết trơn thích hợp; không tự cho duy nhất trên cả mặt phẳng. Chỉ những đặc trưng gặp dữ liệu mới được xác định. Tại gốc, trục hoặc đường hệ số suy biến, kiểm tra riêng và ghép các biểu thức đủ trơn.

\textbf{Bước 6 — Kiểm tra kết quả.} Tính $u_x,u_y$ bằng quy tắc dây chuyền, thay vào $Au_x+Bu_y+Cu$ và kiểm tra bằng $R$. Sau đó thế lại điều kiện đầu/biên. Nêu rõ nếu chỉ có nghiệm cục bộ, nghiệm từng miền hoặc không có nghiệm trong lớp trơn yêu cầu.

\textbf{Chi tiết bước tích phân ODE.} Sau khi tìm bất biến $\rho$, viết $U'+P(s,\rho)U=Q(s,\rho)$. Chọn $M=\exp(\int_{s_0}^{s}P(r,\rho)dr)$; tính $M'=PM$, nên $MU'+M'U=MQ$. Tích phân từ $s_0$: $M(s)U(s)-U(s_0)=\int_{s_0}^sM(r)Q(r)dr$. Hằng số $U(s_0)$ phụ thuộc $\rho$; thay dữ liệu để tìm nó rồi mới đổi về tọa độ gốc.
\subsection*{Lời giải áp dụng cho bài}
\textbf{Bước 1 — Đặc trưng.} Với $au_x+bu_t=0$ và $b\ne0$, dùng $t$ làm tham số: $X'=a/b$, $U'=0$. Do đó $X=\xi+at/b$, $u=F(x-at/b)$.

\textbf{Bước 2 — Áp các dữ liệu.} Với $u(x,0)=xe^{-x^2}$, $F(\xi)=\xi e^{-\xi^2}$, nên
\[u=(x-at/b)e^{-(x-at/b)^2}.\]
Với $u(0,t)=t$ và $a\ne0$, dùng $x$ làm tham số: $t-bx/a$ bất biến. Dữ liệu cho $u=t-bx/a$. Với $a=1,b=-2$ và dữ liệu $1/(1+x^2)$, có $u=1/[1+(x+t/2)^2]$.

\textbf{Bước 3 — Kiểm tra hệ số và ngoại lệ.} Với $u=F(x-at/b)$, $au_x+bu_t=aF'-aF'=0$. Với $u=t-bx/a$, $au_x+bu_t=-b+b=0$. Các hệ số $a,b$ trong hai kết quả in trong sách đã đổi vai trò; các công thức trên sửa theo PDE. Nếu hệ số đạo hàm pháp tuyến bằng 0, đường dữ liệu là đặc trưng và dữ liệu phải hằng dọc nó; các dữ liệu không hằng ở đây không phù hợp. Nếu $a=b=0$, PDE không ràng buộc nghiệm.

\clearpage\noindent\texttt{\detokenize{VD_1_3_1_2_03}}\par
\section*{Lời giải bổ sung — Ví dụ 3 · Mục 1.3.1.2}

\subsection*{Kiến thức và trình tự giải}
\textbf{Mục tiêu.} Với $A(x,y)u_x+B(x,y)u_y+C(x,y)u=R(x,y)$, tìm các đường mà trên đó PDE trở thành một ODE. Dùng chữ $R$ cho vế phải để không nhầm với hàm tùy ý $F$.

\textbf{Bước 1 — Đọc các hệ số và lập hệ đặc trưng.} Đặt $U(s)=u(X(s),Y(s))$. Quy tắc dây chuyền cho $U'=u_xX'+u_yY'$. Chọn $X'=A$, $Y'=B$, khi đó

\[
X'=A(X,Y),\qquad Y'=B(X,Y),\qquad U'+C(X,Y)U=R(X,Y).
\]

Nếu $A\ne0$, có thể dùng $x$ làm tham số: $dy/dx=B/A$, $du/dx+(C/A)u=R/A$. Nếu $A=0$, chọn biến khác hoặc giữ hệ theo $s$; không chia cho hệ số bằng 0.

\textbf{Bước 2 — Tìm bất biến để phân biệt các đặc trưng.} Tích phân hai phương trình đầu để tìm $\rho(X,Y)$ không đổi. Kiểm tra $A\rho_x+B\rho_y=0$. Hệ số hằng cho $\rho=Bx-Ay$. Ví dụ $X'=X,Y'=-Y$ cho $(XY)'=0$; $X'=-Y,Y'=X$ cho $(X^2+Y^2)'=0$. Khi đặc trưng là vòng kín, cần kiểm tra nghiệm sau một vòng trở về cùng giá trị.

\textbf{Bước 3 — Giải ODE của ẩn trên từng đường.} Đặt $M(s)=\exp(\int C(X(s),Y(s))ds)$. Nhân phương trình của $U$ với $M$:

\[
(MU)'=MR,\qquad
U(s)=M(s)^{-1}\left[K+\int M(s)R(X(s),Y(s))ds\right].
\]

$K$ hằng trên một đường, nhưng thay đổi giữa các đường, nên viết $K=F(\rho)$. Nếu tìm được nghiệm riêng $u_p$, có thể giải phương trình thuần nhất cho $u-u_p$ để giảm tính toán. Nếu đổi biến $(\rho,z)$, phải tính Jacobian và dùng quy tắc dây chuyền để thấy PDE chỉ còn đạo hàm theo $z$.

\textbf{Bước 4 — Thế dữ liệu và tìm hàm tùy ý.} Với đường dữ liệu $(x_0(r),y_0(r))$, $u=h(r)$, thế vào nghiệm tổng quát để thu được $F(\rho_0(r))=H(r)$, trong đó $\rho_0=\rho(x_0,y_0)$. Nếu $\rho_0$ khả nghịch, đặt $r=\rho_0^{-1}(z)$ rồi viết $F(z)=H(\rho_0^{-1}(z))$. Nếu hai giá trị $r_1,r_2$ có cùng $\rho_0$, phải có $H(r_1)=H(r_2)$. Nếu đường dữ liệu chính là đặc trưng, dữ liệu phải thỏa ODE trên đường đó và thường không xác định hết $F$.

\textbf{Bước 5 — Xác định miền tồn tại và duy nhất.} Điều kiện không đặc trưng là $x_0'B-y_0'A\ne0$ trên đường dữ liệu. Nó cho khả nghịch cục bộ dưới giả thiết trơn thích hợp; không tự cho duy nhất trên cả mặt phẳng. Chỉ những đặc trưng gặp dữ liệu mới được xác định. Tại gốc, trục hoặc đường hệ số suy biến, kiểm tra riêng và ghép các biểu thức đủ trơn.

\textbf{Bước 6 — Kiểm tra kết quả.} Tính $u_x,u_y$ bằng quy tắc dây chuyền, thay vào $Au_x+Bu_y+Cu$ và kiểm tra bằng $R$. Sau đó thế lại điều kiện đầu/biên. Nêu rõ nếu chỉ có nghiệm cục bộ, nghiệm từng miền hoặc không có nghiệm trong lớp trơn yêu cầu.

\textbf{Chi tiết bước tích phân ODE.} Sau khi tìm bất biến $\rho$, viết $U'+P(s,\rho)U=Q(s,\rho)$. Chọn $M=\exp(\int_{s_0}^{s}P(r,\rho)dr)$; tính $M'=PM$, nên $MU'+M'U=MQ$. Tích phân từ $s_0$: $M(s)U(s)-U(s_0)=\int_{s_0}^sM(r)Q(r)dr$. Hằng số $U(s_0)$ phụ thuộc $\rho$; thay dữ liệu để tìm nó rồi mới đổi về tọa độ gốc.
\subsection*{Lời giải áp dụng cho bài}
\textbf{Bước 1 — Lập đặc trưng.} $u_t-2u_x=2$ cho $X'=-2$, $U'=2$ với $t$ là tham số. Từ $X+2t=\xi$, có bất biến $x+2t$.

\textbf{Bước 2 — Tích phân ẩn.} $U=2t+F(\xi)$, nên $u=2t+F(x+2t)$. Cũng có thể chọn nghiệm riêng $2t$, rồi giải PDE thuần nhất cho phần còn lại.

\textbf{Bước 3 — Kiểm tra.} $u_t=2+2F'$, $u_x=F'$, nên $u_t-2u_x=2$. Không có dữ liệu để xác định $F$.

\clearpage\noindent\texttt{\detokenize{VD_1_3_1_2_04}}\par
\section*{Lời giải bổ sung — Ví dụ 4 · Mục 1.3.1.2}

\subsection*{Kiến thức và trình tự giải}
\textbf{Mục tiêu.} Với $A(x,y)u_x+B(x,y)u_y+C(x,y)u=R(x,y)$, tìm các đường mà trên đó PDE trở thành một ODE. Dùng chữ $R$ cho vế phải để không nhầm với hàm tùy ý $F$.

\textbf{Bước 1 — Đọc các hệ số và lập hệ đặc trưng.} Đặt $U(s)=u(X(s),Y(s))$. Quy tắc dây chuyền cho $U'=u_xX'+u_yY'$. Chọn $X'=A$, $Y'=B$, khi đó

\[
X'=A(X,Y),\qquad Y'=B(X,Y),\qquad U'+C(X,Y)U=R(X,Y).
\]

Nếu $A\ne0$, có thể dùng $x$ làm tham số: $dy/dx=B/A$, $du/dx+(C/A)u=R/A$. Nếu $A=0$, chọn biến khác hoặc giữ hệ theo $s$; không chia cho hệ số bằng 0.

\textbf{Bước 2 — Tìm bất biến để phân biệt các đặc trưng.} Tích phân hai phương trình đầu để tìm $\rho(X,Y)$ không đổi. Kiểm tra $A\rho_x+B\rho_y=0$. Hệ số hằng cho $\rho=Bx-Ay$. Ví dụ $X'=X,Y'=-Y$ cho $(XY)'=0$; $X'=-Y,Y'=X$ cho $(X^2+Y^2)'=0$. Khi đặc trưng là vòng kín, cần kiểm tra nghiệm sau một vòng trở về cùng giá trị.

\textbf{Bước 3 — Giải ODE của ẩn trên từng đường.} Đặt $M(s)=\exp(\int C(X(s),Y(s))ds)$. Nhân phương trình của $U$ với $M$:

\[
(MU)'=MR,\qquad
U(s)=M(s)^{-1}\left[K+\int M(s)R(X(s),Y(s))ds\right].
\]

$K$ hằng trên một đường, nhưng thay đổi giữa các đường, nên viết $K=F(\rho)$. Nếu tìm được nghiệm riêng $u_p$, có thể giải phương trình thuần nhất cho $u-u_p$ để giảm tính toán. Nếu đổi biến $(\rho,z)$, phải tính Jacobian và dùng quy tắc dây chuyền để thấy PDE chỉ còn đạo hàm theo $z$.

\textbf{Bước 4 — Thế dữ liệu và tìm hàm tùy ý.} Với đường dữ liệu $(x_0(r),y_0(r))$, $u=h(r)$, thế vào nghiệm tổng quát để thu được $F(\rho_0(r))=H(r)$, trong đó $\rho_0=\rho(x_0,y_0)$. Nếu $\rho_0$ khả nghịch, đặt $r=\rho_0^{-1}(z)$ rồi viết $F(z)=H(\rho_0^{-1}(z))$. Nếu hai giá trị $r_1,r_2$ có cùng $\rho_0$, phải có $H(r_1)=H(r_2)$. Nếu đường dữ liệu chính là đặc trưng, dữ liệu phải thỏa ODE trên đường đó và thường không xác định hết $F$.

\textbf{Bước 5 — Xác định miền tồn tại và duy nhất.} Điều kiện không đặc trưng là $x_0'B-y_0'A\ne0$ trên đường dữ liệu. Nó cho khả nghịch cục bộ dưới giả thiết trơn thích hợp; không tự cho duy nhất trên cả mặt phẳng. Chỉ những đặc trưng gặp dữ liệu mới được xác định. Tại gốc, trục hoặc đường hệ số suy biến, kiểm tra riêng và ghép các biểu thức đủ trơn.

\textbf{Bước 6 — Kiểm tra kết quả.} Tính $u_x,u_y$ bằng quy tắc dây chuyền, thay vào $Au_x+Bu_y+Cu$ và kiểm tra bằng $R$. Sau đó thế lại điều kiện đầu/biên. Nêu rõ nếu chỉ có nghiệm cục bộ, nghiệm từng miền hoặc không có nghiệm trong lớp trơn yêu cầu.

\textbf{Chi tiết bước tích phân ODE.} Sau khi tìm bất biến $\rho$, viết $U'+P(s,\rho)U=Q(s,\rho)$. Chọn $M=\exp(\int_{s_0}^{s}P(r,\rho)dr)$; tính $M'=PM$, nên $MU'+M'U=MQ$. Tích phân từ $s_0$: $M(s)U(s)-U(s_0)=\int_{s_0}^sM(r)Q(r)dr$. Hằng số $U(s_0)$ phụ thuộc $\rho$; thay dữ liệu để tìm nó rồi mới đổi về tọa độ gốc.
\subsection*{Lời giải áp dụng cho bài}
\textbf{Bước 1 — Đường đặc trưng.} Với $u_x+2u_y+u=0$, dùng $x$ làm tham số: $dy/dx=2$, nên $\rho=y-2x$ không đổi.

\textbf{Bước 2 — Tích phân ẩn.} Dọc đường, $dU/dx+U=0$, nên $(e^xU)'=0$. Vì vậy $u=e^{-x}F(y-2x)$.

\textbf{Bước 3 — Dùng $u(x,0)=x^2$.} $F(-2x)=x^2e^x$. Đặt $s=-2x$, được $F(s)=s^2e^{-s/2}/4$; suy ra
\[u(x,y)=\frac{(y-2x)^2}{4}e^{-y/2}.\]

\textbf{Bước 4 — Kiểm tra.} Trên $y=0$ công thức bằng $x^2$. Với dạng tổng quát, $u_x=e^{-x}(-F-2F')$, $u_y=e^{-x}F'$, nên $u_x+2u_y+u=0$. Mỗi đặc trưng gặp $y=0$ đúng một lần, nên nghiệm được xác định duy nhất.

\clearpage\noindent\texttt{\detokenize{VD_1_3_1_2_05}}\par
\section*{Lời giải bổ sung — Ví dụ 5 · Mục 1.3.1.2}

\subsection*{Kiến thức và trình tự giải}
\textbf{Mục tiêu.} Với $A(x,y)u_x+B(x,y)u_y+C(x,y)u=R(x,y)$, tìm các đường mà trên đó PDE trở thành một ODE. Dùng chữ $R$ cho vế phải để không nhầm với hàm tùy ý $F$.

\textbf{Bước 1 — Đọc các hệ số và lập hệ đặc trưng.} Đặt $U(s)=u(X(s),Y(s))$. Quy tắc dây chuyền cho $U'=u_xX'+u_yY'$. Chọn $X'=A$, $Y'=B$, khi đó

\[
X'=A(X,Y),\qquad Y'=B(X,Y),\qquad U'+C(X,Y)U=R(X,Y).
\]

Nếu $A\ne0$, có thể dùng $x$ làm tham số: $dy/dx=B/A$, $du/dx+(C/A)u=R/A$. Nếu $A=0$, chọn biến khác hoặc giữ hệ theo $s$; không chia cho hệ số bằng 0.

\textbf{Bước 2 — Tìm bất biến để phân biệt các đặc trưng.} Tích phân hai phương trình đầu để tìm $\rho(X,Y)$ không đổi. Kiểm tra $A\rho_x+B\rho_y=0$. Hệ số hằng cho $\rho=Bx-Ay$. Ví dụ $X'=X,Y'=-Y$ cho $(XY)'=0$; $X'=-Y,Y'=X$ cho $(X^2+Y^2)'=0$. Khi đặc trưng là vòng kín, cần kiểm tra nghiệm sau một vòng trở về cùng giá trị.

\textbf{Bước 3 — Giải ODE của ẩn trên từng đường.} Đặt $M(s)=\exp(\int C(X(s),Y(s))ds)$. Nhân phương trình của $U$ với $M$:

\[
(MU)'=MR,\qquad
U(s)=M(s)^{-1}\left[K+\int M(s)R(X(s),Y(s))ds\right].
\]

$K$ hằng trên một đường, nhưng thay đổi giữa các đường, nên viết $K=F(\rho)$. Nếu tìm được nghiệm riêng $u_p$, có thể giải phương trình thuần nhất cho $u-u_p$ để giảm tính toán. Nếu đổi biến $(\rho,z)$, phải tính Jacobian và dùng quy tắc dây chuyền để thấy PDE chỉ còn đạo hàm theo $z$.

\textbf{Bước 4 — Thế dữ liệu và tìm hàm tùy ý.} Với đường dữ liệu $(x_0(r),y_0(r))$, $u=h(r)$, thế vào nghiệm tổng quát để thu được $F(\rho_0(r))=H(r)$, trong đó $\rho_0=\rho(x_0,y_0)$. Nếu $\rho_0$ khả nghịch, đặt $r=\rho_0^{-1}(z)$ rồi viết $F(z)=H(\rho_0^{-1}(z))$. Nếu hai giá trị $r_1,r_2$ có cùng $\rho_0$, phải có $H(r_1)=H(r_2)$. Nếu đường dữ liệu chính là đặc trưng, dữ liệu phải thỏa ODE trên đường đó và thường không xác định hết $F$.

\textbf{Bước 5 — Xác định miền tồn tại và duy nhất.} Điều kiện không đặc trưng là $x_0'B-y_0'A\ne0$ trên đường dữ liệu. Nó cho khả nghịch cục bộ dưới giả thiết trơn thích hợp; không tự cho duy nhất trên cả mặt phẳng. Chỉ những đặc trưng gặp dữ liệu mới được xác định. Tại gốc, trục hoặc đường hệ số suy biến, kiểm tra riêng và ghép các biểu thức đủ trơn.

\textbf{Bước 6 — Kiểm tra kết quả.} Tính $u_x,u_y$ bằng quy tắc dây chuyền, thay vào $Au_x+Bu_y+Cu$ và kiểm tra bằng $R$. Sau đó thế lại điều kiện đầu/biên. Nêu rõ nếu chỉ có nghiệm cục bộ, nghiệm từng miền hoặc không có nghiệm trong lớp trơn yêu cầu.

\textbf{Chi tiết bước tích phân ODE.} Sau khi tìm bất biến $\rho$, viết $U'+P(s,\rho)U=Q(s,\rho)$. Chọn $M=\exp(\int_{s_0}^{s}P(r,\rho)dr)$; tính $M'=PM$, nên $MU'+M'U=MQ$. Tích phân từ $s_0$: $M(s)U(s)-U(s_0)=\int_{s_0}^sM(r)Q(r)dr$. Hằng số $U(s_0)$ phụ thuộc $\rho$; thay dữ liệu để tìm nó rồi mới đổi về tọa độ gốc.
\subsection*{Lời giải áp dụng cho bài}
\textbf{Bước 1 — Tìm bất biến.} Với $2u_x-3u_y=x$, $dy/dx=-3/2$, nên $\rho=3x+2y$.

\textbf{Bước 2 — Tích phân trên đường.} $dU/dx=x/2$, do đó $U=x^2/4+F(\rho)$ và
\[u=x^2/4+F(3x+2y).\]

\textbf{Bước 3 — Kiểm tra.} $u_x=x/2+3F'$, $u_y=2F'$, nên $2u_x-3u_y=x+6F'-6F'=x$.

\clearpage\noindent\texttt{\detokenize{VD_1_3_1_2_06}}\par
\section*{Lời giải bổ sung — Ví dụ 6 · Mục 1.3.1.2}

\subsection*{Kiến thức và trình tự giải}
\textbf{Mục tiêu.} Với $A(x,y)u_x+B(x,y)u_y+C(x,y)u=R(x,y)$, tìm các đường mà trên đó PDE trở thành một ODE. Dùng chữ $R$ cho vế phải để không nhầm với hàm tùy ý $F$.

\textbf{Bước 1 — Đọc các hệ số và lập hệ đặc trưng.} Đặt $U(s)=u(X(s),Y(s))$. Quy tắc dây chuyền cho $U'=u_xX'+u_yY'$. Chọn $X'=A$, $Y'=B$, khi đó

\[
X'=A(X,Y),\qquad Y'=B(X,Y),\qquad U'+C(X,Y)U=R(X,Y).
\]

Nếu $A\ne0$, có thể dùng $x$ làm tham số: $dy/dx=B/A$, $du/dx+(C/A)u=R/A$. Nếu $A=0$, chọn biến khác hoặc giữ hệ theo $s$; không chia cho hệ số bằng 0.

\textbf{Bước 2 — Tìm bất biến để phân biệt các đặc trưng.} Tích phân hai phương trình đầu để tìm $\rho(X,Y)$ không đổi. Kiểm tra $A\rho_x+B\rho_y=0$. Hệ số hằng cho $\rho=Bx-Ay$. Ví dụ $X'=X,Y'=-Y$ cho $(XY)'=0$; $X'=-Y,Y'=X$ cho $(X^2+Y^2)'=0$. Khi đặc trưng là vòng kín, cần kiểm tra nghiệm sau một vòng trở về cùng giá trị.

\textbf{Bước 3 — Giải ODE của ẩn trên từng đường.} Đặt $M(s)=\exp(\int C(X(s),Y(s))ds)$. Nhân phương trình của $U$ với $M$:

\[
(MU)'=MR,\qquad
U(s)=M(s)^{-1}\left[K+\int M(s)R(X(s),Y(s))ds\right].
\]

$K$ hằng trên một đường, nhưng thay đổi giữa các đường, nên viết $K=F(\rho)$. Nếu tìm được nghiệm riêng $u_p$, có thể giải phương trình thuần nhất cho $u-u_p$ để giảm tính toán. Nếu đổi biến $(\rho,z)$, phải tính Jacobian và dùng quy tắc dây chuyền để thấy PDE chỉ còn đạo hàm theo $z$.

\textbf{Bước 4 — Thế dữ liệu và tìm hàm tùy ý.} Với đường dữ liệu $(x_0(r),y_0(r))$, $u=h(r)$, thế vào nghiệm tổng quát để thu được $F(\rho_0(r))=H(r)$, trong đó $\rho_0=\rho(x_0,y_0)$. Nếu $\rho_0$ khả nghịch, đặt $r=\rho_0^{-1}(z)$ rồi viết $F(z)=H(\rho_0^{-1}(z))$. Nếu hai giá trị $r_1,r_2$ có cùng $\rho_0$, phải có $H(r_1)=H(r_2)$. Nếu đường dữ liệu chính là đặc trưng, dữ liệu phải thỏa ODE trên đường đó và thường không xác định hết $F$.

\textbf{Bước 5 — Xác định miền tồn tại và duy nhất.} Điều kiện không đặc trưng là $x_0'B-y_0'A\ne0$ trên đường dữ liệu. Nó cho khả nghịch cục bộ dưới giả thiết trơn thích hợp; không tự cho duy nhất trên cả mặt phẳng. Chỉ những đặc trưng gặp dữ liệu mới được xác định. Tại gốc, trục hoặc đường hệ số suy biến, kiểm tra riêng và ghép các biểu thức đủ trơn.

\textbf{Bước 6 — Kiểm tra kết quả.} Tính $u_x,u_y$ bằng quy tắc dây chuyền, thay vào $Au_x+Bu_y+Cu$ và kiểm tra bằng $R$. Sau đó thế lại điều kiện đầu/biên. Nêu rõ nếu chỉ có nghiệm cục bộ, nghiệm từng miền hoặc không có nghiệm trong lớp trơn yêu cầu.

\textbf{Chi tiết bước tích phân ODE.} Sau khi tìm bất biến $\rho$, viết $U'+P(s,\rho)U=Q(s,\rho)$. Chọn $M=\exp(\int_{s_0}^{s}P(r,\rho)dr)$; tính $M'=PM$, nên $MU'+M'U=MQ$. Tích phân từ $s_0$: $M(s)U(s)-U(s_0)=\int_{s_0}^sM(r)Q(r)dr$. Hằng số $U(s_0)$ phụ thuộc $\rho$; thay dữ liệu để tìm nó rồi mới đổi về tọa độ gốc.
\subsection*{Lời giải áp dụng cho bài}
\textbf{Bước 1 — Đường đặc trưng.} $u_x+2u_y-4u=e^{x+y}$ có $dy/dx=2$, nên $\rho=y-2x$. Trên một đường $y=2x+\rho$, ẩn thỏa $U'-4U=e^{3x+\rho}$.

\textbf{Bước 2 — Nhân thừa số tích phân.}
\[(e^{-4x}U)'=e^{\rho-x},\quad e^{-4x}U=-e^{\rho-x}+F(\rho).\]
Suy ra $u=-e^{x+y}+e^{4x}F(y-2x)$. Dạng này tương đương dạng mũ $e^{4(x+y)/3}$ trong sách sau khi đổi tên hàm tùy ý.

\textbf{Bước 3 — Xác định $F$ từ dữ liệu.} $u(x,0)=\sin(x^2)$ cho
\[F(-2x)=e^{-4x}[\sin(x^2)+e^x].\]
Đặt $s=-2x$, được $F(s)=e^{2s}\sin(s^2/4)+e^{3s/2}$, nên
\[u=e^{2y}\sin((y-2x)^2/4)+e^{x+3y/2}-e^{x+y}.\]

\textbf{Bước 4 — Kiểm tra.} Với $L=\partial_x+2\partial_y-4$, $L(e^{4x}F(\rho))=0$, còn $L(-e^{x+y})=e^{x+y}$. Trên $y=0$, hai hạng mũ triệt tiêu và còn $\sin(x^2)$. Hệ số theo $y$ khác 0 nên mỗi đặc trưng gặp dữ liệu đúng một lần.

\clearpage\noindent\texttt{\detokenize{VD_1_3_1_2_07}}\par
\section*{Lời giải bổ sung — Ví dụ 7 · Mục 1.3.1.2}

\subsection*{Kiến thức và trình tự giải}
\textbf{Mục tiêu.} Với $A(x,y)u_x+B(x,y)u_y+C(x,y)u=R(x,y)$, tìm các đường mà trên đó PDE trở thành một ODE. Dùng chữ $R$ cho vế phải để không nhầm với hàm tùy ý $F$.

\textbf{Bước 1 — Đọc các hệ số và lập hệ đặc trưng.} Đặt $U(s)=u(X(s),Y(s))$. Quy tắc dây chuyền cho $U'=u_xX'+u_yY'$. Chọn $X'=A$, $Y'=B$, khi đó

\[
X'=A(X,Y),\qquad Y'=B(X,Y),\qquad U'+C(X,Y)U=R(X,Y).
\]

Nếu $A\ne0$, có thể dùng $x$ làm tham số: $dy/dx=B/A$, $du/dx+(C/A)u=R/A$. Nếu $A=0$, chọn biến khác hoặc giữ hệ theo $s$; không chia cho hệ số bằng 0.

\textbf{Bước 2 — Tìm bất biến để phân biệt các đặc trưng.} Tích phân hai phương trình đầu để tìm $\rho(X,Y)$ không đổi. Kiểm tra $A\rho_x+B\rho_y=0$. Hệ số hằng cho $\rho=Bx-Ay$. Ví dụ $X'=X,Y'=-Y$ cho $(XY)'=0$; $X'=-Y,Y'=X$ cho $(X^2+Y^2)'=0$. Khi đặc trưng là vòng kín, cần kiểm tra nghiệm sau một vòng trở về cùng giá trị.

\textbf{Bước 3 — Giải ODE của ẩn trên từng đường.} Đặt $M(s)=\exp(\int C(X(s),Y(s))ds)$. Nhân phương trình của $U$ với $M$:

\[
(MU)'=MR,\qquad
U(s)=M(s)^{-1}\left[K+\int M(s)R(X(s),Y(s))ds\right].
\]

$K$ hằng trên một đường, nhưng thay đổi giữa các đường, nên viết $K=F(\rho)$. Nếu tìm được nghiệm riêng $u_p$, có thể giải phương trình thuần nhất cho $u-u_p$ để giảm tính toán. Nếu đổi biến $(\rho,z)$, phải tính Jacobian và dùng quy tắc dây chuyền để thấy PDE chỉ còn đạo hàm theo $z$.

\textbf{Bước 4 — Thế dữ liệu và tìm hàm tùy ý.} Với đường dữ liệu $(x_0(r),y_0(r))$, $u=h(r)$, thế vào nghiệm tổng quát để thu được $F(\rho_0(r))=H(r)$, trong đó $\rho_0=\rho(x_0,y_0)$. Nếu $\rho_0$ khả nghịch, đặt $r=\rho_0^{-1}(z)$ rồi viết $F(z)=H(\rho_0^{-1}(z))$. Nếu hai giá trị $r_1,r_2$ có cùng $\rho_0$, phải có $H(r_1)=H(r_2)$. Nếu đường dữ liệu chính là đặc trưng, dữ liệu phải thỏa ODE trên đường đó và thường không xác định hết $F$.

\textbf{Bước 5 — Xác định miền tồn tại và duy nhất.} Điều kiện không đặc trưng là $x_0'B-y_0'A\ne0$ trên đường dữ liệu. Nó cho khả nghịch cục bộ dưới giả thiết trơn thích hợp; không tự cho duy nhất trên cả mặt phẳng. Chỉ những đặc trưng gặp dữ liệu mới được xác định. Tại gốc, trục hoặc đường hệ số suy biến, kiểm tra riêng và ghép các biểu thức đủ trơn.

\textbf{Bước 6 — Kiểm tra kết quả.} Tính $u_x,u_y$ bằng quy tắc dây chuyền, thay vào $Au_x+Bu_y+Cu$ và kiểm tra bằng $R$. Sau đó thế lại điều kiện đầu/biên. Nêu rõ nếu chỉ có nghiệm cục bộ, nghiệm từng miền hoặc không có nghiệm trong lớp trơn yêu cầu.

\textbf{Chi tiết bước tích phân ODE.} Sau khi tìm bất biến $\rho$, viết $U'+P(s,\rho)U=Q(s,\rho)$. Chọn $M=\exp(\int_{s_0}^{s}P(r,\rho)dr)$; tính $M'=PM$, nên $MU'+M'U=MQ$. Tích phân từ $s_0$: $M(s)U(s)-U(s_0)=\int_{s_0}^sM(r)Q(r)dr$. Hằng số $U(s_0)$ phụ thuộc $\rho$; thay dữ liệu để tìm nó rồi mới đổi về tọa độ gốc.
\subsection*{Lời giải áp dụng cho bài}
\textbf{Bước 1 — Giải PDE.} $u_x+u_y-u=0$ cho $dy/dx=1$, $\rho=y-x$ và $U'=U$, nên $u=e^xF(y-x)$.

\textbf{Bước 2 — Thế dữ liệu trên đặc trưng.} Trên $y=x$, $\rho=0$, nên $u(x,x)=e^xF(0)$. Để bằng $\tan x$, phải có $F(0)=e^{-x}\tan x$ ở mọi điểm của khoảng dữ liệu.

\textbf{Bước 3 — Chứng minh mâu thuẫn.} Ở $x=0$, điều kiện cho $F(0)=0$, nhưng ở $x$ gần 0 khác 0, $e^{-x}\tan x\ne0$. Do đó không có nghiệm trên lân cận đường dữ liệu quanh gốc. Cũng có thể dùng điều kiện cần $g'=g$ trên đường; $g=\tan x$ không thỏa.

\clearpage\noindent\texttt{\detokenize{VD_1_3_1_2_08}}\par
\section*{Lời giải bổ sung — Ví dụ 8 · Mục 1.3.1.2}

\subsection*{Kiến thức và trình tự giải}
\textbf{Mục tiêu.} Với $A(x,y)u_x+B(x,y)u_y+C(x,y)u=R(x,y)$, tìm các đường mà trên đó PDE trở thành một ODE. Dùng chữ $R$ cho vế phải để không nhầm với hàm tùy ý $F$.

\textbf{Bước 1 — Đọc các hệ số và lập hệ đặc trưng.} Đặt $U(s)=u(X(s),Y(s))$. Quy tắc dây chuyền cho $U'=u_xX'+u_yY'$. Chọn $X'=A$, $Y'=B$, khi đó

\[
X'=A(X,Y),\qquad Y'=B(X,Y),\qquad U'+C(X,Y)U=R(X,Y).
\]

Nếu $A\ne0$, có thể dùng $x$ làm tham số: $dy/dx=B/A$, $du/dx+(C/A)u=R/A$. Nếu $A=0$, chọn biến khác hoặc giữ hệ theo $s$; không chia cho hệ số bằng 0.

\textbf{Bước 2 — Tìm bất biến để phân biệt các đặc trưng.} Tích phân hai phương trình đầu để tìm $\rho(X,Y)$ không đổi. Kiểm tra $A\rho_x+B\rho_y=0$. Hệ số hằng cho $\rho=Bx-Ay$. Ví dụ $X'=X,Y'=-Y$ cho $(XY)'=0$; $X'=-Y,Y'=X$ cho $(X^2+Y^2)'=0$. Khi đặc trưng là vòng kín, cần kiểm tra nghiệm sau một vòng trở về cùng giá trị.

\textbf{Bước 3 — Giải ODE của ẩn trên từng đường.} Đặt $M(s)=\exp(\int C(X(s),Y(s))ds)$. Nhân phương trình của $U$ với $M$:

\[
(MU)'=MR,\qquad
U(s)=M(s)^{-1}\left[K+\int M(s)R(X(s),Y(s))ds\right].
\]

$K$ hằng trên một đường, nhưng thay đổi giữa các đường, nên viết $K=F(\rho)$. Nếu tìm được nghiệm riêng $u_p$, có thể giải phương trình thuần nhất cho $u-u_p$ để giảm tính toán. Nếu đổi biến $(\rho,z)$, phải tính Jacobian và dùng quy tắc dây chuyền để thấy PDE chỉ còn đạo hàm theo $z$.

\textbf{Bước 4 — Thế dữ liệu và tìm hàm tùy ý.} Với đường dữ liệu $(x_0(r),y_0(r))$, $u=h(r)$, thế vào nghiệm tổng quát để thu được $F(\rho_0(r))=H(r)$, trong đó $\rho_0=\rho(x_0,y_0)$. Nếu $\rho_0$ khả nghịch, đặt $r=\rho_0^{-1}(z)$ rồi viết $F(z)=H(\rho_0^{-1}(z))$. Nếu hai giá trị $r_1,r_2$ có cùng $\rho_0$, phải có $H(r_1)=H(r_2)$. Nếu đường dữ liệu chính là đặc trưng, dữ liệu phải thỏa ODE trên đường đó và thường không xác định hết $F$.

\textbf{Bước 5 — Xác định miền tồn tại và duy nhất.} Điều kiện không đặc trưng là $x_0'B-y_0'A\ne0$ trên đường dữ liệu. Nó cho khả nghịch cục bộ dưới giả thiết trơn thích hợp; không tự cho duy nhất trên cả mặt phẳng. Chỉ những đặc trưng gặp dữ liệu mới được xác định. Tại gốc, trục hoặc đường hệ số suy biến, kiểm tra riêng và ghép các biểu thức đủ trơn.

\textbf{Bước 6 — Kiểm tra kết quả.} Tính $u_x,u_y$ bằng quy tắc dây chuyền, thay vào $Au_x+Bu_y+Cu$ và kiểm tra bằng $R$. Sau đó thế lại điều kiện đầu/biên. Nêu rõ nếu chỉ có nghiệm cục bộ, nghiệm từng miền hoặc không có nghiệm trong lớp trơn yêu cầu.

\textbf{Chi tiết bước tích phân ODE.} Sau khi tìm bất biến $\rho$, viết $U'+P(s,\rho)U=Q(s,\rho)$. Chọn $M=\exp(\int_{s_0}^{s}P(r,\rho)dr)$; tính $M'=PM$, nên $MU'+M'U=MQ$. Tích phân từ $s_0$: $M(s)U(s)-U(s_0)=\int_{s_0}^sM(r)Q(r)dr$. Hằng số $U(s_0)$ phụ thuộc $\rho$; thay dữ liệu để tìm nó rồi mới đổi về tọa độ gốc.
\subsection*{Lời giải áp dụng cho bài}
\textbf{Bước 1 — Dùng nghiệm tổng quát.} Như ví dụ trước, $u=e^xF(y-x)$. Trên $y=x$, chỉ số đặc trưng bằng 0, nên dữ liệu phải là $f(x)=e^xF(0)$.

\textbf{Bước 2 — Điều kiện cần và đủ.} Đặt $A=F(0)$, ta có $f(x)=Ae^x$. Ngược lại, nếu $f$ có dạng này, mọi hàm $F$ đủ trơn thỏa $F(0)=A$ đều cho nghiệm.

\textbf{Bước 3 — Giải thích không duy nhất.} Chọn $F_1(s)=A$ và $F_2(s)=A+s$ cho $u_1=Ae^x$, $u_2=e^x(A+y-x)$. Hai hàm cùng thỏa PDE và dữ liệu, nhưng khác nhau ngoài $y=x$. Dữ liệu trên một đặc trưng chỉ xác định giá trị của $F$ tại một điểm.

\textbf{Chi tiết dựng nghiệm tổng quát trước khi dùng vết.} Với $u_x+u_y=u$, chọn $x$ làm tham số: $Y'=1$, $U'=U$. Tích phân $Y=x+\rho$ và $(e^{-x}U)'=0$ cho $e^{-x}U=F(\rho)$. Vì $\rho=y-x$, thu $u=e^xF(y-x)$. Như vậy công thức dùng ở bước 1 xuất phát từ chính hệ đặc trưng của bài này.

\clearpage\noindent\texttt{\detokenize{VD_1_3_1_2_09}}\par
\section*{Lời giải bổ sung — Ví dụ 9 · Mục 1.3.1.2}

\subsection*{Kiến thức và trình tự giải}
\textbf{Mục tiêu.} Với $A(x,y)u_x+B(x,y)u_y+C(x,y)u=R(x,y)$, tìm các đường mà trên đó PDE trở thành một ODE. Dùng chữ $R$ cho vế phải để không nhầm với hàm tùy ý $F$.

\textbf{Bước 1 — Đọc các hệ số và lập hệ đặc trưng.} Đặt $U(s)=u(X(s),Y(s))$. Quy tắc dây chuyền cho $U'=u_xX'+u_yY'$. Chọn $X'=A$, $Y'=B$, khi đó

\[
X'=A(X,Y),\qquad Y'=B(X,Y),\qquad U'+C(X,Y)U=R(X,Y).
\]

Nếu $A\ne0$, có thể dùng $x$ làm tham số: $dy/dx=B/A$, $du/dx+(C/A)u=R/A$. Nếu $A=0$, chọn biến khác hoặc giữ hệ theo $s$; không chia cho hệ số bằng 0.

\textbf{Bước 2 — Tìm bất biến để phân biệt các đặc trưng.} Tích phân hai phương trình đầu để tìm $\rho(X,Y)$ không đổi. Kiểm tra $A\rho_x+B\rho_y=0$. Hệ số hằng cho $\rho=Bx-Ay$. Ví dụ $X'=X,Y'=-Y$ cho $(XY)'=0$; $X'=-Y,Y'=X$ cho $(X^2+Y^2)'=0$. Khi đặc trưng là vòng kín, cần kiểm tra nghiệm sau một vòng trở về cùng giá trị.

\textbf{Bước 3 — Giải ODE của ẩn trên từng đường.} Đặt $M(s)=\exp(\int C(X(s),Y(s))ds)$. Nhân phương trình của $U$ với $M$:

\[
(MU)'=MR,\qquad
U(s)=M(s)^{-1}\left[K+\int M(s)R(X(s),Y(s))ds\right].
\]

$K$ hằng trên một đường, nhưng thay đổi giữa các đường, nên viết $K=F(\rho)$. Nếu tìm được nghiệm riêng $u_p$, có thể giải phương trình thuần nhất cho $u-u_p$ để giảm tính toán. Nếu đổi biến $(\rho,z)$, phải tính Jacobian và dùng quy tắc dây chuyền để thấy PDE chỉ còn đạo hàm theo $z$.

\textbf{Bước 4 — Thế dữ liệu và tìm hàm tùy ý.} Với đường dữ liệu $(x_0(r),y_0(r))$, $u=h(r)$, thế vào nghiệm tổng quát để thu được $F(\rho_0(r))=H(r)$, trong đó $\rho_0=\rho(x_0,y_0)$. Nếu $\rho_0$ khả nghịch, đặt $r=\rho_0^{-1}(z)$ rồi viết $F(z)=H(\rho_0^{-1}(z))$. Nếu hai giá trị $r_1,r_2$ có cùng $\rho_0$, phải có $H(r_1)=H(r_2)$. Nếu đường dữ liệu chính là đặc trưng, dữ liệu phải thỏa ODE trên đường đó và thường không xác định hết $F$.

\textbf{Bước 5 — Xác định miền tồn tại và duy nhất.} Điều kiện không đặc trưng là $x_0'B-y_0'A\ne0$ trên đường dữ liệu. Nó cho khả nghịch cục bộ dưới giả thiết trơn thích hợp; không tự cho duy nhất trên cả mặt phẳng. Chỉ những đặc trưng gặp dữ liệu mới được xác định. Tại gốc, trục hoặc đường hệ số suy biến, kiểm tra riêng và ghép các biểu thức đủ trơn.

\textbf{Bước 6 — Kiểm tra kết quả.} Tính $u_x,u_y$ bằng quy tắc dây chuyền, thay vào $Au_x+Bu_y+Cu$ và kiểm tra bằng $R$. Sau đó thế lại điều kiện đầu/biên. Nêu rõ nếu chỉ có nghiệm cục bộ, nghiệm từng miền hoặc không có nghiệm trong lớp trơn yêu cầu.

\textbf{Chi tiết bước tích phân ODE.} Sau khi tìm bất biến $\rho$, viết $U'+P(s,\rho)U=Q(s,\rho)$. Chọn $M=\exp(\int_{s_0}^{s}P(r,\rho)dr)$; tính $M'=PM$, nên $MU'+M'U=MQ$. Tích phân từ $s_0$: $M(s)U(s)-U(s_0)=\int_{s_0}^sM(r)Q(r)dr$. Hằng số $U(s_0)$ phụ thuộc $\rho$; thay dữ liệu để tìm nó rồi mới đổi về tọa độ gốc.
\subsection*{Lời giải áp dụng cho bài}
\textbf{Bước 1 — Giải theo $x$.} Với $au_x+bu_y+cu=0$, $ab\ne0$, bất biến $\rho=bx-ay$ và $dU/dx+(c/a)U=0$. Suy ra $u=e^{-cx/a}F(\rho)$.

\textbf{Bước 2 — Giải theo $y$.} Tương tự $dU/dy+(c/b)U=0$, nên $u=e^{-cy/b}G(\rho)$. Cả hai cách đều đúng, với hàm tùy ý tương ứng.

\textbf{Bước 3 — Chuyển giữa hai biểu thức.} Vì $x/a-y/b=(bx-ay)/(ab)=\rho/(ab)$,
\[e^{-cx/a}F(\rho)=e^{-cy/b}e^{-c\rho/(ab)}F(\rho).\]
Do đó đặt $G(s)=e^{-cs/(ab)}F(s)$ để được cùng nghiệm. Không cần $F=G$; hai bạn đang dùng hai cách tham số hóa cùng họ nghiệm.

\clearpage\noindent\texttt{\detokenize{BT_1_3_1_3_01}}\par
\section*{Lời giải bổ sung — Bài tập 1 · Mục 1.3.1.3}

\subsection*{Kiến thức và trình tự giải}
\textbf{Mục tiêu.} Với $A(x,y)u_x+B(x,y)u_y+C(x,y)u=R(x,y)$, tìm các đường mà trên đó PDE trở thành một ODE. Dùng chữ $R$ cho vế phải để không nhầm với hàm tùy ý $F$.

\textbf{Bước 1 — Đọc các hệ số và lập hệ đặc trưng.} Đặt $U(s)=u(X(s),Y(s))$. Quy tắc dây chuyền cho $U'=u_xX'+u_yY'$. Chọn $X'=A$, $Y'=B$, khi đó

\[
X'=A(X,Y),\qquad Y'=B(X,Y),\qquad U'+C(X,Y)U=R(X,Y).
\]

Nếu $A\ne0$, có thể dùng $x$ làm tham số: $dy/dx=B/A$, $du/dx+(C/A)u=R/A$. Nếu $A=0$, chọn biến khác hoặc giữ hệ theo $s$; không chia cho hệ số bằng 0.

\textbf{Bước 2 — Tìm bất biến để phân biệt các đặc trưng.} Tích phân hai phương trình đầu để tìm $\rho(X,Y)$ không đổi. Kiểm tra $A\rho_x+B\rho_y=0$. Hệ số hằng cho $\rho=Bx-Ay$. Ví dụ $X'=X,Y'=-Y$ cho $(XY)'=0$; $X'=-Y,Y'=X$ cho $(X^2+Y^2)'=0$. Khi đặc trưng là vòng kín, cần kiểm tra nghiệm sau một vòng trở về cùng giá trị.

\textbf{Bước 3 — Giải ODE của ẩn trên từng đường.} Đặt $M(s)=\exp(\int C(X(s),Y(s))ds)$. Nhân phương trình của $U$ với $M$:

\[
(MU)'=MR,\qquad
U(s)=M(s)^{-1}\left[K+\int M(s)R(X(s),Y(s))ds\right].
\]

$K$ hằng trên một đường, nhưng thay đổi giữa các đường, nên viết $K=F(\rho)$. Nếu tìm được nghiệm riêng $u_p$, có thể giải phương trình thuần nhất cho $u-u_p$ để giảm tính toán. Nếu đổi biến $(\rho,z)$, phải tính Jacobian và dùng quy tắc dây chuyền để thấy PDE chỉ còn đạo hàm theo $z$.

\textbf{Bước 4 — Thế dữ liệu và tìm hàm tùy ý.} Với đường dữ liệu $(x_0(r),y_0(r))$, $u=h(r)$, thế vào nghiệm tổng quát để thu được $F(\rho_0(r))=H(r)$, trong đó $\rho_0=\rho(x_0,y_0)$. Nếu $\rho_0$ khả nghịch, đặt $r=\rho_0^{-1}(z)$ rồi viết $F(z)=H(\rho_0^{-1}(z))$. Nếu hai giá trị $r_1,r_2$ có cùng $\rho_0$, phải có $H(r_1)=H(r_2)$. Nếu đường dữ liệu chính là đặc trưng, dữ liệu phải thỏa ODE trên đường đó và thường không xác định hết $F$.

\textbf{Bước 5 — Xác định miền tồn tại và duy nhất.} Điều kiện không đặc trưng là $x_0'B-y_0'A\ne0$ trên đường dữ liệu. Nó cho khả nghịch cục bộ dưới giả thiết trơn thích hợp; không tự cho duy nhất trên cả mặt phẳng. Chỉ những đặc trưng gặp dữ liệu mới được xác định. Tại gốc, trục hoặc đường hệ số suy biến, kiểm tra riêng và ghép các biểu thức đủ trơn.

\textbf{Bước 6 — Kiểm tra kết quả.} Tính $u_x,u_y$ bằng quy tắc dây chuyền, thay vào $Au_x+Bu_y+Cu$ và kiểm tra bằng $R$. Sau đó thế lại điều kiện đầu/biên. Nêu rõ nếu chỉ có nghiệm cục bộ, nghiệm từng miền hoặc không có nghiệm trong lớp trơn yêu cầu.

\textbf{Chi tiết bước tích phân ODE.} Sau khi tìm bất biến $\rho$, viết $U'+P(s,\rho)U=Q(s,\rho)$. Chọn $M=\exp(\int_{s_0}^{s}P(r,\rho)dr)$; tính $M'=PM$, nên $MU'+M'U=MQ$. Tích phân từ $s_0$: $M(s)U(s)-U(s_0)=\int_{s_0}^sM(r)Q(r)dr$. Hằng số $U(s_0)$ phụ thuộc $\rho$; thay dữ liệu để tìm nó rồi mới đổi về tọa độ gốc.
\subsection*{Lời giải áp dụng cho bài}
\textbf{Ý (a) — Đạo hàm theo hướng.} Nếu $(a,b)\ne(0,0)$, đặt $\nu=(a,b)/\sqrt{a^2+b^2}$. Khi đó $D_\nu u=(au_x+bu_y)/\sqrt{a^2+b^2}=0$. Nghiệm không đổi khi di chuyển theo hướng $(a,b)$.

\textbf{Ý (b) — Đặc trưng.} Giải $X'=a,Y'=b$: $(X,Y)=(x_0,y_0)+s(a,b)$. Tính $(bX-aY)'=ba-ab=0$, nên các đường là $bx-ay=\text{hằng}$.

\textbf{Ý (c) — Nghiệm.} $U'=au_x+bu_y=0$ trên mỗi đường, nên $u=F(bx-ay)$. Kiểm tra $u_x=bF'$, $u_y=-aF'$, suy $au_x+bu_y=0$. Nếu $a=b=0$, PDE là $0=0$, mọi hàm $C^1$ là nghiệm. Trên toàn mặt phẳng, mỗi mức của bất biến là một đường liên thông; trên miền bị chia cắt cần xét từng phần đặc trưng liên thông.

\clearpage\noindent\texttt{\detokenize{BT_1_3_1_3_02}}\par
\section*{Lời giải bổ sung — Bài tập 2 · Mục 1.3.1.3}

\subsection*{Kiến thức và trình tự giải}
\textbf{Mục tiêu.} Với $A(x,y)u_x+B(x,y)u_y+C(x,y)u=R(x,y)$, tìm các đường mà trên đó PDE trở thành một ODE. Dùng chữ $R$ cho vế phải để không nhầm với hàm tùy ý $F$.

\textbf{Bước 1 — Đọc các hệ số và lập hệ đặc trưng.} Đặt $U(s)=u(X(s),Y(s))$. Quy tắc dây chuyền cho $U'=u_xX'+u_yY'$. Chọn $X'=A$, $Y'=B$, khi đó

\[
X'=A(X,Y),\qquad Y'=B(X,Y),\qquad U'+C(X,Y)U=R(X,Y).
\]

Nếu $A\ne0$, có thể dùng $x$ làm tham số: $dy/dx=B/A$, $du/dx+(C/A)u=R/A$. Nếu $A=0$, chọn biến khác hoặc giữ hệ theo $s$; không chia cho hệ số bằng 0.

\textbf{Bước 2 — Tìm bất biến để phân biệt các đặc trưng.} Tích phân hai phương trình đầu để tìm $\rho(X,Y)$ không đổi. Kiểm tra $A\rho_x+B\rho_y=0$. Hệ số hằng cho $\rho=Bx-Ay$. Ví dụ $X'=X,Y'=-Y$ cho $(XY)'=0$; $X'=-Y,Y'=X$ cho $(X^2+Y^2)'=0$. Khi đặc trưng là vòng kín, cần kiểm tra nghiệm sau một vòng trở về cùng giá trị.

\textbf{Bước 3 — Giải ODE của ẩn trên từng đường.} Đặt $M(s)=\exp(\int C(X(s),Y(s))ds)$. Nhân phương trình của $U$ với $M$:

\[
(MU)'=MR,\qquad
U(s)=M(s)^{-1}\left[K+\int M(s)R(X(s),Y(s))ds\right].
\]

$K$ hằng trên một đường, nhưng thay đổi giữa các đường, nên viết $K=F(\rho)$. Nếu tìm được nghiệm riêng $u_p$, có thể giải phương trình thuần nhất cho $u-u_p$ để giảm tính toán. Nếu đổi biến $(\rho,z)$, phải tính Jacobian và dùng quy tắc dây chuyền để thấy PDE chỉ còn đạo hàm theo $z$.

\textbf{Bước 4 — Thế dữ liệu và tìm hàm tùy ý.} Với đường dữ liệu $(x_0(r),y_0(r))$, $u=h(r)$, thế vào nghiệm tổng quát để thu được $F(\rho_0(r))=H(r)$, trong đó $\rho_0=\rho(x_0,y_0)$. Nếu $\rho_0$ khả nghịch, đặt $r=\rho_0^{-1}(z)$ rồi viết $F(z)=H(\rho_0^{-1}(z))$. Nếu hai giá trị $r_1,r_2$ có cùng $\rho_0$, phải có $H(r_1)=H(r_2)$. Nếu đường dữ liệu chính là đặc trưng, dữ liệu phải thỏa ODE trên đường đó và thường không xác định hết $F$.

\textbf{Bước 5 — Xác định miền tồn tại và duy nhất.} Điều kiện không đặc trưng là $x_0'B-y_0'A\ne0$ trên đường dữ liệu. Nó cho khả nghịch cục bộ dưới giả thiết trơn thích hợp; không tự cho duy nhất trên cả mặt phẳng. Chỉ những đặc trưng gặp dữ liệu mới được xác định. Tại gốc, trục hoặc đường hệ số suy biến, kiểm tra riêng và ghép các biểu thức đủ trơn.

\textbf{Bước 6 — Kiểm tra kết quả.} Tính $u_x,u_y$ bằng quy tắc dây chuyền, thay vào $Au_x+Bu_y+Cu$ và kiểm tra bằng $R$. Sau đó thế lại điều kiện đầu/biên. Nêu rõ nếu chỉ có nghiệm cục bộ, nghiệm từng miền hoặc không có nghiệm trong lớp trơn yêu cầu.

\textbf{Chi tiết bước tích phân ODE.} Sau khi tìm bất biến $\rho$, viết $U'+P(s,\rho)U=Q(s,\rho)$. Chọn $M=\exp(\int_{s_0}^{s}P(r,\rho)dr)$; tính $M'=PM$, nên $MU'+M'U=MQ$. Tích phân từ $s_0$: $M(s)U(s)-U(s_0)=\int_{s_0}^sM(r)Q(r)dr$. Hằng số $U(s_0)$ phụ thuộc $\rho$; thay dữ liệu để tìm nó rồi mới đổi về tọa độ gốc.
\subsection*{Lời giải áp dụng cho bài}
\textbf{Ý (a).} $u_x+u_y-u=0$. Bước 1: $dy/dx=1$ cho $\rho=y-x$. Bước 2: $U'=U$ cho $U=e^xF(\rho)$. Bước 3: $u_x=e^x(F-F')$, $u_y=e^xF'$, nên $u_x+u_y=u$.

\textbf{Ý (b).} $2u_t+3u_x=0$. Chia cho 2, dùng $t$ làm tham số: $X'=3/2$, $U'=0$. Bất biến $\rho=x-3t/2$, do đó $u=F(x-3t/2)$. Kiểm tra $2u_t+3u_x=2(-3F'/2)+3F'=0$.

\textbf{Ý (c).} $au_t+bu_x=u$ với $a,b\ne0$. $X'=b/a$ và $U'=U/a$ cho bất biến $\rho=ax-bt$, $U=e^{t/a}F(\rho)$. Do đó $u=e^{t/a}F(ax-bt)$. Đạo hàm cho $au_t=e^{t/a}(F-abF')$, $bu_x=ab e^{t/a}F'$, tổng bằng $u$. Các hàm $F$ tùy ý $C^1$.

\clearpage\noindent\texttt{\detokenize{BT_1_3_1_3_03}}\par
\section*{Lời giải bổ sung — Bài tập 3 · Mục 1.3.1.3}

\subsection*{Kiến thức và trình tự giải}
\textbf{Mục tiêu.} Với $A(x,y)u_x+B(x,y)u_y+C(x,y)u=R(x,y)$, tìm các đường mà trên đó PDE trở thành một ODE. Dùng chữ $R$ cho vế phải để không nhầm với hàm tùy ý $F$.

\textbf{Bước 1 — Đọc các hệ số và lập hệ đặc trưng.} Đặt $U(s)=u(X(s),Y(s))$. Quy tắc dây chuyền cho $U'=u_xX'+u_yY'$. Chọn $X'=A$, $Y'=B$, khi đó

\[
X'=A(X,Y),\qquad Y'=B(X,Y),\qquad U'+C(X,Y)U=R(X,Y).
\]

Nếu $A\ne0$, có thể dùng $x$ làm tham số: $dy/dx=B/A$, $du/dx+(C/A)u=R/A$. Nếu $A=0$, chọn biến khác hoặc giữ hệ theo $s$; không chia cho hệ số bằng 0.

\textbf{Bước 2 — Tìm bất biến để phân biệt các đặc trưng.} Tích phân hai phương trình đầu để tìm $\rho(X,Y)$ không đổi. Kiểm tra $A\rho_x+B\rho_y=0$. Hệ số hằng cho $\rho=Bx-Ay$. Ví dụ $X'=X,Y'=-Y$ cho $(XY)'=0$; $X'=-Y,Y'=X$ cho $(X^2+Y^2)'=0$. Khi đặc trưng là vòng kín, cần kiểm tra nghiệm sau một vòng trở về cùng giá trị.

\textbf{Bước 3 — Giải ODE của ẩn trên từng đường.} Đặt $M(s)=\exp(\int C(X(s),Y(s))ds)$. Nhân phương trình của $U$ với $M$:

\[
(MU)'=MR,\qquad
U(s)=M(s)^{-1}\left[K+\int M(s)R(X(s),Y(s))ds\right].
\]

$K$ hằng trên một đường, nhưng thay đổi giữa các đường, nên viết $K=F(\rho)$. Nếu tìm được nghiệm riêng $u_p$, có thể giải phương trình thuần nhất cho $u-u_p$ để giảm tính toán. Nếu đổi biến $(\rho,z)$, phải tính Jacobian và dùng quy tắc dây chuyền để thấy PDE chỉ còn đạo hàm theo $z$.

\textbf{Bước 4 — Thế dữ liệu và tìm hàm tùy ý.} Với đường dữ liệu $(x_0(r),y_0(r))$, $u=h(r)$, thế vào nghiệm tổng quát để thu được $F(\rho_0(r))=H(r)$, trong đó $\rho_0=\rho(x_0,y_0)$. Nếu $\rho_0$ khả nghịch, đặt $r=\rho_0^{-1}(z)$ rồi viết $F(z)=H(\rho_0^{-1}(z))$. Nếu hai giá trị $r_1,r_2$ có cùng $\rho_0$, phải có $H(r_1)=H(r_2)$. Nếu đường dữ liệu chính là đặc trưng, dữ liệu phải thỏa ODE trên đường đó và thường không xác định hết $F$.

\textbf{Bước 5 — Xác định miền tồn tại và duy nhất.} Điều kiện không đặc trưng là $x_0'B-y_0'A\ne0$ trên đường dữ liệu. Nó cho khả nghịch cục bộ dưới giả thiết trơn thích hợp; không tự cho duy nhất trên cả mặt phẳng. Chỉ những đặc trưng gặp dữ liệu mới được xác định. Tại gốc, trục hoặc đường hệ số suy biến, kiểm tra riêng và ghép các biểu thức đủ trơn.

\textbf{Bước 6 — Kiểm tra kết quả.} Tính $u_x,u_y$ bằng quy tắc dây chuyền, thay vào $Au_x+Bu_y+Cu$ và kiểm tra bằng $R$. Sau đó thế lại điều kiện đầu/biên. Nêu rõ nếu chỉ có nghiệm cục bộ, nghiệm từng miền hoặc không có nghiệm trong lớp trơn yêu cầu.

\textbf{Chi tiết bước tích phân ODE.} Sau khi tìm bất biến $\rho$, viết $U'+P(s,\rho)U=Q(s,\rho)$. Chọn $M=\exp(\int_{s_0}^{s}P(r,\rho)dr)$; tính $M'=PM$, nên $MU'+M'U=MQ$. Tích phân từ $s_0$: $M(s)U(s)-U(s_0)=\int_{s_0}^sM(r)Q(r)dr$. Hằng số $U(s_0)$ phụ thuộc $\rho$; thay dữ liệu để tìm nó rồi mới đổi về tọa độ gốc.
\subsection*{Lời giải áp dụng cho bài}
\textbf{Ý (a) — $au_x+bu_y=f$, $a,b,f$ hằng.} Nếu $a\ne0$, dùng $x$ làm tham số: $dy/dx=b/a$, nên $\rho=bx-ay$. Trên đường, $dU/dx=f/a$, tích phân cho
\[u=fx/a+F(bx-ay).\]
Kiểm tra $au_x+bu_y=a(f/a+bF')-abF'=f$. Nếu $a=0,b\ne0$, tích phân theo $y$ cho $u=fy/b+F(x)$. Nếu $a=b=0$, chỉ có nghiệm khi $f=0$, khi ấy mọi hàm đủ trơn đều thỏa.

\textbf{Ý (b) — $u_x+2u_y+2u=1$.} Bước 1: $dy/dx=2$ cho $\rho=y-2x$. Bước 2: $U'+2U=1$. Nhân $e^{2x}$:
\[(e^{2x}U)'=e^{2x},\quad e^{2x}U=e^{2x}/2+F(\rho).\]
Suy $u=1/2+e^{-2x}F(y-2x)$. Bước 3: $u_x=e^{-2x}(-2F-2F')$, $u_y=e^{-2x}F'$, nên $u_x+2u_y+2u=1$.

\textbf{Ý (c) — $3u_x-4u_y=x+e^x$.} $dy/dx=-4/3$ cho $\rho=4x+3y$. Dọc đường, $U'=(x+e^x)/3$; tích phân:
\[u=x^2/6+e^x/3+F(4x+3y).\]
Kiểm tra $3u_x=x+e^x+12F'$, $-4u_y=-12F'$, tổng bằng vế phải.

\textbf{Ý (d) — $u_x+3u_y=9y^2$.} Bất biến $\rho=y-3x$. Trên đường $y=3x+\rho$, có $U'=9(3x+\rho)^2$. Vì $d[(3x+\rho)^3]/dx=9(3x+\rho)^2$, được $u=y^3+F(y-3x)$. Kiểm tra $u_x=-3F'$, $u_y=3y^2+F'$, tổng bằng $9y^2$.

\textbf{Ý (e) — Dùng $h(x)=2x^2+3x$ trên $y=0$.} Thực hiện bước xác định hàm tùy ý riêng cho từng PDE:

Với (a), nếu $b\ne0$, đặc trưng từ $(\xi,0)$ có $x=\xi+ay/b$. Do đó $\xi=x-ay/b$, còn $dU/dy=f/b$, $U(0)=h(\xi)$. Suy
\[u=h(x-ay/b)+fy/b.\]
Trên $y=0$ bằng $h(x)$; $au_x+bu_y=ah'(\xi)+b[-ah'(\xi)/b+f/b]=f$.

Với (b), $h(x)=1/2+e^{-2x}F(-2x)$. Đặt $s=-2x$, được $F(s)=e^{-s}[h(-s/2)-1/2]$. Thay $s=y-2x$:
\[u=\tfrac12+e^{-y}[h(x-y/2)-\tfrac12].\]

Với (c), đặt $P(x)=x^2/6+e^x/3$. Dữ liệu cho $F(4x)=h(x)-P(x)$, nên $F(s)=h(s/4)-P(s/4)$. Suy
\[u=P(x)+h(x+3y/4)-P(x+3y/4).\]

Với (d), $F(-3x)=h(x)$, nên $F(s)=h(-s/3)$ và
\[u=y^3+h(x-y/3).\]
Ba công thức cuối thỏa PDE vì được lấy từ các nghiệm tổng quát đã kiểm tra; thế $y=0$ cho đúng $h(x)$.

\textbf{Trường hợp đặc trưng ở (a).} Nếu $b=0$, trên dữ liệu phải có $ah'(x)=f$, tức $a(4x+3)=f$ ở cả khoảng. Với các hằng số $a,f$, điều này chỉ có thể xảy ra khi $a=f=0$. Khi ấy PDE không ràng buộc và có vô số nghiệm cùng vết $h$. Trong các trường hợp còn lại không có nghiệm thỏa dữ liệu. Khi $b\ne0$, và ở (b),(c),(d), đường $y=0$ không đặc trưng, nên dữ liệu xác định duy nhất các nghiệm trên toàn mặt phẳng.

\textbf{Chi tiết thế đa thức dữ liệu vào bốn nghiệm ở (e).} Với $h(s)=2s^2+3s$, các công thức lần lượt là
\[\begin{aligned}
u_a&=2(x-ay/b)^2+3(x-ay/b)+fy/b &&(b\ne0),\\
u_b&=\tfrac12+e^{-y}[2(x-y/2)^2+3(x-y/2)-\tfrac12],\\
u_c&=\frac{x^2}{6}+\frac{e^x}{3}+\frac{11}{6}(x+3y/4)^2+3(x+3y/4)-\frac{e^{x+3y/4}}3,\\
u_d&=y^3+2(x-y/3)^2+3(x-y/3).
\end{aligned}\]
Ở (b), đặt $z=x-y/2,Q=h(z)-1/2$; $u_x=e^{-y}h'(z)$, $u_y=-e^{-y}Q-e^{-y}h'(z)/2$. Vậy $u_x+2u_y+2u=1$. Ở (c), đặt $z=x+3y/4$, $u_x=P'(x)+h'(z)-P'(z)$, $u_y=3[h'(z)-P'(z)]/4$; thế vào $3u_x-4u_y$ còn $3P'(x)=x+e^x$. Ở (d), $u_x=h'(x-y/3)$, $u_y=3y^2-h'(x-y/3)/3$, cho $u_x+3u_y=9y^2$.

\clearpage\noindent\texttt{\detokenize{BT_1_3_1_3_04}}\par
\section*{Lời giải bổ sung — Bài tập 4 · Mục 1.3.1.3}

\subsection*{Kiến thức và trình tự giải}
\textbf{Mục tiêu.} Với $A(x,y)u_x+B(x,y)u_y+C(x,y)u=R(x,y)$, tìm các đường mà trên đó PDE trở thành một ODE. Dùng chữ $R$ cho vế phải để không nhầm với hàm tùy ý $F$.

\textbf{Bước 1 — Đọc các hệ số và lập hệ đặc trưng.} Đặt $U(s)=u(X(s),Y(s))$. Quy tắc dây chuyền cho $U'=u_xX'+u_yY'$. Chọn $X'=A$, $Y'=B$, khi đó

\[
X'=A(X,Y),\qquad Y'=B(X,Y),\qquad U'+C(X,Y)U=R(X,Y).
\]

Nếu $A\ne0$, có thể dùng $x$ làm tham số: $dy/dx=B/A$, $du/dx+(C/A)u=R/A$. Nếu $A=0$, chọn biến khác hoặc giữ hệ theo $s$; không chia cho hệ số bằng 0.

\textbf{Bước 2 — Tìm bất biến để phân biệt các đặc trưng.} Tích phân hai phương trình đầu để tìm $\rho(X,Y)$ không đổi. Kiểm tra $A\rho_x+B\rho_y=0$. Hệ số hằng cho $\rho=Bx-Ay$. Ví dụ $X'=X,Y'=-Y$ cho $(XY)'=0$; $X'=-Y,Y'=X$ cho $(X^2+Y^2)'=0$. Khi đặc trưng là vòng kín, cần kiểm tra nghiệm sau một vòng trở về cùng giá trị.

\textbf{Bước 3 — Giải ODE của ẩn trên từng đường.} Đặt $M(s)=\exp(\int C(X(s),Y(s))ds)$. Nhân phương trình của $U$ với $M$:

\[
(MU)'=MR,\qquad
U(s)=M(s)^{-1}\left[K+\int M(s)R(X(s),Y(s))ds\right].
\]

$K$ hằng trên một đường, nhưng thay đổi giữa các đường, nên viết $K=F(\rho)$. Nếu tìm được nghiệm riêng $u_p$, có thể giải phương trình thuần nhất cho $u-u_p$ để giảm tính toán. Nếu đổi biến $(\rho,z)$, phải tính Jacobian và dùng quy tắc dây chuyền để thấy PDE chỉ còn đạo hàm theo $z$.

\textbf{Bước 4 — Thế dữ liệu và tìm hàm tùy ý.} Với đường dữ liệu $(x_0(r),y_0(r))$, $u=h(r)$, thế vào nghiệm tổng quát để thu được $F(\rho_0(r))=H(r)$, trong đó $\rho_0=\rho(x_0,y_0)$. Nếu $\rho_0$ khả nghịch, đặt $r=\rho_0^{-1}(z)$ rồi viết $F(z)=H(\rho_0^{-1}(z))$. Nếu hai giá trị $r_1,r_2$ có cùng $\rho_0$, phải có $H(r_1)=H(r_2)$. Nếu đường dữ liệu chính là đặc trưng, dữ liệu phải thỏa ODE trên đường đó và thường không xác định hết $F$.

\textbf{Bước 5 — Xác định miền tồn tại và duy nhất.} Điều kiện không đặc trưng là $x_0'B-y_0'A\ne0$ trên đường dữ liệu. Nó cho khả nghịch cục bộ dưới giả thiết trơn thích hợp; không tự cho duy nhất trên cả mặt phẳng. Chỉ những đặc trưng gặp dữ liệu mới được xác định. Tại gốc, trục hoặc đường hệ số suy biến, kiểm tra riêng và ghép các biểu thức đủ trơn.

\textbf{Bước 6 — Kiểm tra kết quả.} Tính $u_x,u_y$ bằng quy tắc dây chuyền, thay vào $Au_x+Bu_y+Cu$ và kiểm tra bằng $R$. Sau đó thế lại điều kiện đầu/biên. Nêu rõ nếu chỉ có nghiệm cục bộ, nghiệm từng miền hoặc không có nghiệm trong lớp trơn yêu cầu.

\textbf{Chi tiết bước tích phân ODE.} Sau khi tìm bất biến $\rho$, viết $U'+P(s,\rho)U=Q(s,\rho)$. Chọn $M=\exp(\int_{s_0}^{s}P(r,\rho)dr)$; tính $M'=PM$, nên $MU'+M'U=MQ$. Tích phân từ $s_0$: $M(s)U(s)-U(s_0)=\int_{s_0}^sM(r)Q(r)dr$. Hằng số $U(s_0)$ phụ thuộc $\rho$; thay dữ liệu để tìm nó rồi mới đổi về tọa độ gốc.
\subsection*{Lời giải áp dụng cho bài}
\textbf{Bước 1 — Giải PDE trước khi dùng dữ liệu.} $dy/dx=2$ cho $\rho=y-2x$. Trên đường, $U'-4U=e^{3x+\rho}$ và $(e^{-4x}U)'=e^{\rho-x}$. Tích phân cho
\[u=-e^{x+y}+e^{4x}F(y-2x).\]

\textbf{Bước 2 — Xác định $F$ từng ý.}
(a) $u(0,y)=y^2$ cho $-e^y+F(y)=y^2$, nên $F(s)=s^2+e^s$.

(b) $u(x,-x)=x$ cho $-1+e^{4x}F(-3x)=x$. Đặt $s=-3x$, suy $F(s)=(1-s/3)e^{4s/3}$.

(c) $u(x,4x)=x^2$ cho $-e^{5x}+e^{4x}F(2x)=x^2$. Đặt $s=2x$, suy $F(s)=s^2e^{-2s}/4+e^{s/2}$.

Trong mỗi ý, thay hàm $F$ vừa tìm vào công thức ở bước 1 để được nghiệm. Chú ý ở (b) cả hạng 1 và hạng $-s/3$ đều nhân $e^{4s/3}$.

(d) Dữ liệu trên parabol cho
\[F(x^2-2x)=H(x),\quad H(x)=e^{-4x}[f(x)+e^{x+x^2}].\]
Vì $x^2-2x=(x-1)^2-1$, hai điểm $x$ và $2-x$ nằm trên cùng đặc trưng. Cần $H(x)=H(2-x)$. Khi tương thích, đặt $F(s)=H(1+\sqrt{s+1})$ trên $s\ge-1$ và kiểm tra độ trơn tại $s=-1$; đối xứng đơn thuần chưa đủ nếu $H$ thiếu trơn. Miền dữ liệu phủ là $y-2x\ge-1$. Muốn nghiệm trên toàn mặt phẳng phải có kéo dài $F$ đủ trơn qua $-1$; trên $s<-1$ kéo dài không được dữ liệu quyết định.

\textbf{Bước 3 — Kiểm tra.} Với $L=\partial_x+2\partial_y-4$, phần $e^{4x}F(\rho)$ có $L=0$, phần $-e^{x+y}$ có $L=e^{x+y}$. Thế lại các đường dữ liệu vào $F$ kiểm tra đúng từng giá trị. Ở (a),(b),(c), chỉ số đặc trưng lần lượt là $y,-3x,2x$, phủ toàn trục nên cho duy nhất.

\textbf{Chi tiết thay $F$ và viết nghiệm cuối của (a),(b),(c).} Trong $u=-e^{x+y}+e^{4x}F(y-2x)$:
\[\begin{aligned}
(a)\quad u&=-e^{x+y}+e^{4x}(y-2x)^2+e^{2x+y},\\
(b)\quad u&=-e^{x+y}+\frac{3+2x-y}{3}e^{4(x+y)/3},\\
(c)\quad u&=-e^{x+y}+\frac{(y-2x)^2}{4}e^{8x-2y}+e^{3x+y/2}.
\end{aligned}\]
Ở (b), $4x+4(y-2x)/3=4(x+y)/3$; ở (c), $4x-2(y-2x)=8x-2y$ và $4x+(y-2x)/2=3x+y/2$. Các bước này giữ nguyên thừa số mũ khi thế $F(s)=s^2e^{-2s}/4+e^{s/2}$.
Thế $x=0$ vào (a): $-e^y+y^2+e^y=y^2$. Thế $y=-x$ vào (b): $-1+(3+3x)/3=x$. Thế $y=4x$ vào $u_c$: $-e^{5x}+x^2+e^{5x}=x^2$.

\clearpage\noindent\texttt{\detokenize{BT_1_3_1_3_05}}\par
\section*{Lời giải bổ sung — Bài tập 5 · Mục 1.3.1.3}

\subsection*{Kiến thức và trình tự giải}
\textbf{Mục tiêu.} Với $A(x,y)u_x+B(x,y)u_y+C(x,y)u=R(x,y)$, tìm các đường mà trên đó PDE trở thành một ODE. Dùng chữ $R$ cho vế phải để không nhầm với hàm tùy ý $F$.

\textbf{Bước 1 — Đọc các hệ số và lập hệ đặc trưng.} Đặt $U(s)=u(X(s),Y(s))$. Quy tắc dây chuyền cho $U'=u_xX'+u_yY'$. Chọn $X'=A$, $Y'=B$, khi đó

\[
X'=A(X,Y),\qquad Y'=B(X,Y),\qquad U'+C(X,Y)U=R(X,Y).
\]

Nếu $A\ne0$, có thể dùng $x$ làm tham số: $dy/dx=B/A$, $du/dx+(C/A)u=R/A$. Nếu $A=0$, chọn biến khác hoặc giữ hệ theo $s$; không chia cho hệ số bằng 0.

\textbf{Bước 2 — Tìm bất biến để phân biệt các đặc trưng.} Tích phân hai phương trình đầu để tìm $\rho(X,Y)$ không đổi. Kiểm tra $A\rho_x+B\rho_y=0$. Hệ số hằng cho $\rho=Bx-Ay$. Ví dụ $X'=X,Y'=-Y$ cho $(XY)'=0$; $X'=-Y,Y'=X$ cho $(X^2+Y^2)'=0$. Khi đặc trưng là vòng kín, cần kiểm tra nghiệm sau một vòng trở về cùng giá trị.

\textbf{Bước 3 — Giải ODE của ẩn trên từng đường.} Đặt $M(s)=\exp(\int C(X(s),Y(s))ds)$. Nhân phương trình của $U$ với $M$:

\[
(MU)'=MR,\qquad
U(s)=M(s)^{-1}\left[K+\int M(s)R(X(s),Y(s))ds\right].
\]

$K$ hằng trên một đường, nhưng thay đổi giữa các đường, nên viết $K=F(\rho)$. Nếu tìm được nghiệm riêng $u_p$, có thể giải phương trình thuần nhất cho $u-u_p$ để giảm tính toán. Nếu đổi biến $(\rho,z)$, phải tính Jacobian và dùng quy tắc dây chuyền để thấy PDE chỉ còn đạo hàm theo $z$.

\textbf{Bước 4 — Thế dữ liệu và tìm hàm tùy ý.} Với đường dữ liệu $(x_0(r),y_0(r))$, $u=h(r)$, thế vào nghiệm tổng quát để thu được $F(\rho_0(r))=H(r)$, trong đó $\rho_0=\rho(x_0,y_0)$. Nếu $\rho_0$ khả nghịch, đặt $r=\rho_0^{-1}(z)$ rồi viết $F(z)=H(\rho_0^{-1}(z))$. Nếu hai giá trị $r_1,r_2$ có cùng $\rho_0$, phải có $H(r_1)=H(r_2)$. Nếu đường dữ liệu chính là đặc trưng, dữ liệu phải thỏa ODE trên đường đó và thường không xác định hết $F$.

\textbf{Bước 5 — Xác định miền tồn tại và duy nhất.} Điều kiện không đặc trưng là $x_0'B-y_0'A\ne0$ trên đường dữ liệu. Nó cho khả nghịch cục bộ dưới giả thiết trơn thích hợp; không tự cho duy nhất trên cả mặt phẳng. Chỉ những đặc trưng gặp dữ liệu mới được xác định. Tại gốc, trục hoặc đường hệ số suy biến, kiểm tra riêng và ghép các biểu thức đủ trơn.

\textbf{Bước 6 — Kiểm tra kết quả.} Tính $u_x,u_y$ bằng quy tắc dây chuyền, thay vào $Au_x+Bu_y+Cu$ và kiểm tra bằng $R$. Sau đó thế lại điều kiện đầu/biên. Nêu rõ nếu chỉ có nghiệm cục bộ, nghiệm từng miền hoặc không có nghiệm trong lớp trơn yêu cầu.

\textbf{Chi tiết bước tích phân ODE.} Sau khi tìm bất biến $\rho$, viết $U'+P(s,\rho)U=Q(s,\rho)$. Chọn $M=\exp(\int_{s_0}^{s}P(r,\rho)dr)$; tính $M'=PM$, nên $MU'+M'U=MQ$. Tích phân từ $s_0$: $M(s)U(s)-U(s_0)=\int_{s_0}^sM(r)Q(r)dr$. Hằng số $U(s_0)$ phụ thuộc $\rho$; thay dữ liệu để tìm nó rồi mới đổi về tọa độ gốc.
\subsection*{Lời giải áp dụng cho bài}
\textbf{Bước 1 — Giải dọc đặc trưng.} $u_x+3u_y+u=1$ cho $\rho=y-3x$ và $U'+U=1$. Nhân $e^x$, tích phân được $u=1+e^{-x}F(y-3x)$.

\textbf{Bước 2 — Xét dữ liệu trên $y=3x$.} Trên đường này $\rho=0$, nên $g(x)=1+e^{-x}F(0)$. Do đó điều kiện cần và đủ là $g(x)=1+Ae^{-x}$ với $A$ hằng. Cũng có thể đạo hàm vết để có $g'+g=1$.

\textbf{Bước 3 — Tồn tại và không duy nhất.} Nếu dữ liệu phù hợp, chọn bất kỳ $F$ đủ trơn với $F(0)=A$. Hai lựa chọn $F=A$ và $F=A+s$ cho hai nghiệm khác nhau cùng vết.

\textbf{Ý (b) — Dữ liệu in trong đề.} Với $g=-1+2e^x$, ta có $g'+g=-1+4e^x$, không đồng nhất bằng 1. Vì thế vô nghiệm; không thể viết hai nghiệm cho dữ liệu này. Nếu sửa dữ liệu thành $1+2e^{-x}$, hai nghiệm là $1+2e^{-x}$ và $1+e^{-x}(2+y-3x)$. Đây là minh họa cho đề đã sửa, không phải đáp án của dữ liệu nguyên văn.

\clearpage\noindent\texttt{\detokenize{BT_1_3_1_3_06}}\par
\section*{Lời giải bổ sung — Bài tập 6 · Mục 1.3.1.3}

\subsection*{Kiến thức và trình tự giải}
\textbf{Mục tiêu.} Với $A(x,y)u_x+B(x,y)u_y+C(x,y)u=R(x,y)$, tìm các đường mà trên đó PDE trở thành một ODE. Dùng chữ $R$ cho vế phải để không nhầm với hàm tùy ý $F$.

\textbf{Bước 1 — Đọc các hệ số và lập hệ đặc trưng.} Đặt $U(s)=u(X(s),Y(s))$. Quy tắc dây chuyền cho $U'=u_xX'+u_yY'$. Chọn $X'=A$, $Y'=B$, khi đó

\[
X'=A(X,Y),\qquad Y'=B(X,Y),\qquad U'+C(X,Y)U=R(X,Y).
\]

Nếu $A\ne0$, có thể dùng $x$ làm tham số: $dy/dx=B/A$, $du/dx+(C/A)u=R/A$. Nếu $A=0$, chọn biến khác hoặc giữ hệ theo $s$; không chia cho hệ số bằng 0.

\textbf{Bước 2 — Tìm bất biến để phân biệt các đặc trưng.} Tích phân hai phương trình đầu để tìm $\rho(X,Y)$ không đổi. Kiểm tra $A\rho_x+B\rho_y=0$. Hệ số hằng cho $\rho=Bx-Ay$. Ví dụ $X'=X,Y'=-Y$ cho $(XY)'=0$; $X'=-Y,Y'=X$ cho $(X^2+Y^2)'=0$. Khi đặc trưng là vòng kín, cần kiểm tra nghiệm sau một vòng trở về cùng giá trị.

\textbf{Bước 3 — Giải ODE của ẩn trên từng đường.} Đặt $M(s)=\exp(\int C(X(s),Y(s))ds)$. Nhân phương trình của $U$ với $M$:

\[
(MU)'=MR,\qquad
U(s)=M(s)^{-1}\left[K+\int M(s)R(X(s),Y(s))ds\right].
\]

$K$ hằng trên một đường, nhưng thay đổi giữa các đường, nên viết $K=F(\rho)$. Nếu tìm được nghiệm riêng $u_p$, có thể giải phương trình thuần nhất cho $u-u_p$ để giảm tính toán. Nếu đổi biến $(\rho,z)$, phải tính Jacobian và dùng quy tắc dây chuyền để thấy PDE chỉ còn đạo hàm theo $z$.

\textbf{Bước 4 — Thế dữ liệu và tìm hàm tùy ý.} Với đường dữ liệu $(x_0(r),y_0(r))$, $u=h(r)$, thế vào nghiệm tổng quát để thu được $F(\rho_0(r))=H(r)$, trong đó $\rho_0=\rho(x_0,y_0)$. Nếu $\rho_0$ khả nghịch, đặt $r=\rho_0^{-1}(z)$ rồi viết $F(z)=H(\rho_0^{-1}(z))$. Nếu hai giá trị $r_1,r_2$ có cùng $\rho_0$, phải có $H(r_1)=H(r_2)$. Nếu đường dữ liệu chính là đặc trưng, dữ liệu phải thỏa ODE trên đường đó và thường không xác định hết $F$.

\textbf{Bước 5 — Xác định miền tồn tại và duy nhất.} Điều kiện không đặc trưng là $x_0'B-y_0'A\ne0$ trên đường dữ liệu. Nó cho khả nghịch cục bộ dưới giả thiết trơn thích hợp; không tự cho duy nhất trên cả mặt phẳng. Chỉ những đặc trưng gặp dữ liệu mới được xác định. Tại gốc, trục hoặc đường hệ số suy biến, kiểm tra riêng và ghép các biểu thức đủ trơn.

\textbf{Bước 6 — Kiểm tra kết quả.} Tính $u_x,u_y$ bằng quy tắc dây chuyền, thay vào $Au_x+Bu_y+Cu$ và kiểm tra bằng $R$. Sau đó thế lại điều kiện đầu/biên. Nêu rõ nếu chỉ có nghiệm cục bộ, nghiệm từng miền hoặc không có nghiệm trong lớp trơn yêu cầu.

\textbf{Chi tiết bước tích phân ODE.} Sau khi tìm bất biến $\rho$, viết $U'+P(s,\rho)U=Q(s,\rho)$. Chọn $M=\exp(\int_{s_0}^{s}P(r,\rho)dr)$; tính $M'=PM$, nên $MU'+M'U=MQ$. Tích phân từ $s_0$: $M(s)U(s)-U(s_0)=\int_{s_0}^sM(r)Q(r)dr$. Hằng số $U(s_0)$ phụ thuộc $\rho$; thay dữ liệu để tìm nó rồi mới đổi về tọa độ gốc.
\subsection*{Lời giải áp dụng cho bài}
\textbf{Bước 1 — Đặc trưng của $au_x+bu_y=0$.} Giải $X'=a,Y'=b$; tính $(bX-aY)'=0$, nên bất biến $\rho=bx-ay$.

\textbf{Bước 2 — Tích phân ẩn.} $U'=0$ cho $U$ hằng trên mỗi đặc trưng, nhưng hằng có thể khác giữa các đường. Suy $u=F(bx-ay)$ với $(a,b)\ne(0,0)$.

\textbf{Bước 3 — Kiểm tra và ngoại lệ.} $u_x=bF'$, $u_y=-aF'$ nên PDE đúng. Nếu $a=b=0$, phương trình không ràng buộc, mọi hàm đủ trơn là nghiệm; không dùng dạng $F(0)$ để thu hẹp sai họ nghiệm.

\clearpage\noindent\texttt{\detokenize{BT_1_3_1_3_07}}\par
\section*{Lời giải bổ sung — Bài tập 7 · Mục 1.3.1.3}

\subsection*{Kiến thức và trình tự giải}
\textbf{Mục tiêu.} Với $A(x,y)u_x+B(x,y)u_y+C(x,y)u=R(x,y)$, tìm các đường mà trên đó PDE trở thành một ODE. Dùng chữ $R$ cho vế phải để không nhầm với hàm tùy ý $F$.

\textbf{Bước 1 — Đọc các hệ số và lập hệ đặc trưng.} Đặt $U(s)=u(X(s),Y(s))$. Quy tắc dây chuyền cho $U'=u_xX'+u_yY'$. Chọn $X'=A$, $Y'=B$, khi đó

\[
X'=A(X,Y),\qquad Y'=B(X,Y),\qquad U'+C(X,Y)U=R(X,Y).
\]

Nếu $A\ne0$, có thể dùng $x$ làm tham số: $dy/dx=B/A$, $du/dx+(C/A)u=R/A$. Nếu $A=0$, chọn biến khác hoặc giữ hệ theo $s$; không chia cho hệ số bằng 0.

\textbf{Bước 2 — Tìm bất biến để phân biệt các đặc trưng.} Tích phân hai phương trình đầu để tìm $\rho(X,Y)$ không đổi. Kiểm tra $A\rho_x+B\rho_y=0$. Hệ số hằng cho $\rho=Bx-Ay$. Ví dụ $X'=X,Y'=-Y$ cho $(XY)'=0$; $X'=-Y,Y'=X$ cho $(X^2+Y^2)'=0$. Khi đặc trưng là vòng kín, cần kiểm tra nghiệm sau một vòng trở về cùng giá trị.

\textbf{Bước 3 — Giải ODE của ẩn trên từng đường.} Đặt $M(s)=\exp(\int C(X(s),Y(s))ds)$. Nhân phương trình của $U$ với $M$:

\[
(MU)'=MR,\qquad
U(s)=M(s)^{-1}\left[K+\int M(s)R(X(s),Y(s))ds\right].
\]

$K$ hằng trên một đường, nhưng thay đổi giữa các đường, nên viết $K=F(\rho)$. Nếu tìm được nghiệm riêng $u_p$, có thể giải phương trình thuần nhất cho $u-u_p$ để giảm tính toán. Nếu đổi biến $(\rho,z)$, phải tính Jacobian và dùng quy tắc dây chuyền để thấy PDE chỉ còn đạo hàm theo $z$.

\textbf{Bước 4 — Thế dữ liệu và tìm hàm tùy ý.} Với đường dữ liệu $(x_0(r),y_0(r))$, $u=h(r)$, thế vào nghiệm tổng quát để thu được $F(\rho_0(r))=H(r)$, trong đó $\rho_0=\rho(x_0,y_0)$. Nếu $\rho_0$ khả nghịch, đặt $r=\rho_0^{-1}(z)$ rồi viết $F(z)=H(\rho_0^{-1}(z))$. Nếu hai giá trị $r_1,r_2$ có cùng $\rho_0$, phải có $H(r_1)=H(r_2)$. Nếu đường dữ liệu chính là đặc trưng, dữ liệu phải thỏa ODE trên đường đó và thường không xác định hết $F$.

\textbf{Bước 5 — Xác định miền tồn tại và duy nhất.} Điều kiện không đặc trưng là $x_0'B-y_0'A\ne0$ trên đường dữ liệu. Nó cho khả nghịch cục bộ dưới giả thiết trơn thích hợp; không tự cho duy nhất trên cả mặt phẳng. Chỉ những đặc trưng gặp dữ liệu mới được xác định. Tại gốc, trục hoặc đường hệ số suy biến, kiểm tra riêng và ghép các biểu thức đủ trơn.

\textbf{Bước 6 — Kiểm tra kết quả.} Tính $u_x,u_y$ bằng quy tắc dây chuyền, thay vào $Au_x+Bu_y+Cu$ và kiểm tra bằng $R$. Sau đó thế lại điều kiện đầu/biên. Nêu rõ nếu chỉ có nghiệm cục bộ, nghiệm từng miền hoặc không có nghiệm trong lớp trơn yêu cầu.

\textbf{Chi tiết bước tích phân ODE.} Sau khi tìm bất biến $\rho$, viết $U'+P(s,\rho)U=Q(s,\rho)$. Chọn $M=\exp(\int_{s_0}^{s}P(r,\rho)dr)$; tính $M'=PM$, nên $MU'+M'U=MQ$. Tích phân từ $s_0$: $M(s)U(s)-U(s_0)=\int_{s_0}^sM(r)Q(r)dr$. Hằng số $U(s_0)$ phụ thuộc $\rho$; thay dữ liệu để tìm nó rồi mới đổi về tọa độ gốc.
\subsection*{Lời giải áp dụng cho bài}
\textbf{Bước 1 — Nghiệm tổng quát.} Với $(a,b)\ne(0,0)$, bài trước cho $u=F(bx-ay)$.

\textbf{Ý (a).} Trên $bx-ay=1$, nghiệm phải bằng hằng $F(1)$. Do đó dữ liệu $h(x)$ phải có cùng giá trị tại mọi điểm thực sự nằm trên đường. Nếu $a\ne0$, $x$ chạy toàn trục và cần $h\equiv H$. Khi đó mọi $F$ với $F(1)=H$ cho nghiệm, nên không duy nhất. Nếu $a=0,b\ne0$, đường là $x=1/b$, dữ liệu chỉ gán giá trị $h(1/b)$; không đòi $h$ hằng ở những $x$ không nằm trên đường.

\textbf{Ý (b).} Đường $x-y=1$ có $y=x-1$, nên $F((b-a)x+a)=h(x)$. Nếu $b\ne a$, đặt $s=(b-a)x+a$:
\[F(s)=h\left(\frac{s-a}{b-a}\right),\quad u=h\left(\frac{bx-ay-a}{b-a}\right).\]
Chỉ số $s$ phủ toàn trục, cho duy nhất. Nếu $b=a\ne0$, chỉ số bằng $a$, cần $h$ hằng; khi tương thích vẫn có vô số $F$ với cùng $F(a)$.

\textbf{Ngoại lệ $a=b=0$.} Ở (a), tập $bx-ay=1$ rỗng, nên không có điều kiện dữ liệu thực sự. Ở (b), chỉ cần một hàm đủ trơn có vết $h$ trên đường; PDE không cho duy nhất. Các kết luận trên giả sử dữ liệu có độ trơn phù hợp.

\clearpage\noindent\texttt{\detokenize{BT_1_3_1_3_08}}\par
\section*{Lời giải bổ sung — Bài tập 8 · Mục 1.3.1.3}

\subsection*{Kiến thức và trình tự giải}
\textbf{Mục tiêu.} Với $A(x,y)u_x+B(x,y)u_y+C(x,y)u=R(x,y)$, tìm các đường mà trên đó PDE trở thành một ODE. Dùng chữ $R$ cho vế phải để không nhầm với hàm tùy ý $F$.

\textbf{Bước 1 — Đọc các hệ số và lập hệ đặc trưng.} Đặt $U(s)=u(X(s),Y(s))$. Quy tắc dây chuyền cho $U'=u_xX'+u_yY'$. Chọn $X'=A$, $Y'=B$, khi đó

\[
X'=A(X,Y),\qquad Y'=B(X,Y),\qquad U'+C(X,Y)U=R(X,Y).
\]

Nếu $A\ne0$, có thể dùng $x$ làm tham số: $dy/dx=B/A$, $du/dx+(C/A)u=R/A$. Nếu $A=0$, chọn biến khác hoặc giữ hệ theo $s$; không chia cho hệ số bằng 0.

\textbf{Bước 2 — Tìm bất biến để phân biệt các đặc trưng.} Tích phân hai phương trình đầu để tìm $\rho(X,Y)$ không đổi. Kiểm tra $A\rho_x+B\rho_y=0$. Hệ số hằng cho $\rho=Bx-Ay$. Ví dụ $X'=X,Y'=-Y$ cho $(XY)'=0$; $X'=-Y,Y'=X$ cho $(X^2+Y^2)'=0$. Khi đặc trưng là vòng kín, cần kiểm tra nghiệm sau một vòng trở về cùng giá trị.

\textbf{Bước 3 — Giải ODE của ẩn trên từng đường.} Đặt $M(s)=\exp(\int C(X(s),Y(s))ds)$. Nhân phương trình của $U$ với $M$:

\[
(MU)'=MR,\qquad
U(s)=M(s)^{-1}\left[K+\int M(s)R(X(s),Y(s))ds\right].
\]

$K$ hằng trên một đường, nhưng thay đổi giữa các đường, nên viết $K=F(\rho)$. Nếu tìm được nghiệm riêng $u_p$, có thể giải phương trình thuần nhất cho $u-u_p$ để giảm tính toán. Nếu đổi biến $(\rho,z)$, phải tính Jacobian và dùng quy tắc dây chuyền để thấy PDE chỉ còn đạo hàm theo $z$.

\textbf{Bước 4 — Thế dữ liệu và tìm hàm tùy ý.} Với đường dữ liệu $(x_0(r),y_0(r))$, $u=h(r)$, thế vào nghiệm tổng quát để thu được $F(\rho_0(r))=H(r)$, trong đó $\rho_0=\rho(x_0,y_0)$. Nếu $\rho_0$ khả nghịch, đặt $r=\rho_0^{-1}(z)$ rồi viết $F(z)=H(\rho_0^{-1}(z))$. Nếu hai giá trị $r_1,r_2$ có cùng $\rho_0$, phải có $H(r_1)=H(r_2)$. Nếu đường dữ liệu chính là đặc trưng, dữ liệu phải thỏa ODE trên đường đó và thường không xác định hết $F$.

\textbf{Bước 5 — Xác định miền tồn tại và duy nhất.} Điều kiện không đặc trưng là $x_0'B-y_0'A\ne0$ trên đường dữ liệu. Nó cho khả nghịch cục bộ dưới giả thiết trơn thích hợp; không tự cho duy nhất trên cả mặt phẳng. Chỉ những đặc trưng gặp dữ liệu mới được xác định. Tại gốc, trục hoặc đường hệ số suy biến, kiểm tra riêng và ghép các biểu thức đủ trơn.

\textbf{Bước 6 — Kiểm tra kết quả.} Tính $u_x,u_y$ bằng quy tắc dây chuyền, thay vào $Au_x+Bu_y+Cu$ và kiểm tra bằng $R$. Sau đó thế lại điều kiện đầu/biên. Nêu rõ nếu chỉ có nghiệm cục bộ, nghiệm từng miền hoặc không có nghiệm trong lớp trơn yêu cầu.

\textbf{Chi tiết bước tích phân ODE.} Sau khi tìm bất biến $\rho$, viết $U'+P(s,\rho)U=Q(s,\rho)$. Chọn $M=\exp(\int_{s_0}^{s}P(r,\rho)dr)$; tính $M'=PM$, nên $MU'+M'U=MQ$. Tích phân từ $s_0$: $M(s)U(s)-U(s_0)=\int_{s_0}^sM(r)Q(r)dr$. Hằng số $U(s_0)$ phụ thuộc $\rho$; thay dữ liệu để tìm nó rồi mới đổi về tọa độ gốc.
\subsection*{Lời giải áp dụng cho bài}
\textbf{Bước 1 — Giải PDE.} $u_x=0$ nên $u=F(y)$: nghiệm hằng trên các đường ngang.

\textbf{Bước 2 — Thế dữ liệu.} $u(x,x^2)=x$ đòi $F(x^2)=x$. Với $r>0$, hai điểm $x=\sqrt r$ và $x=-\sqrt r$ cùng cho đối số $r$, nhưng yêu cầu $F(r)=\sqrt r$ và $F(r)=-\sqrt r$.

\textbf{Bước 3 — Kết luận và miền hẹp.} Không có nghiệm trên toàn đường dữ liệu. Nếu chỉ giữ nhánh $x>0$, có $F(y)=\sqrt y$ trên $y>0$. Nếu muốn gồm gốc và nghiệm $C^1$, $F'(y)=1/(2\sqrt y)$ không hữu hạn tại 0, nên cũng không có kéo dài $C^1$ theo công thức này.

\clearpage\noindent\texttt{\detokenize{BT_1_3_1_3_09}}\par
\section*{Lời giải bổ sung — Bài tập 9 · Mục 1.3.1.3}

\subsection*{Kiến thức và trình tự giải}
\textbf{Mục tiêu.} Với $A(x,y)u_x+B(x,y)u_y+C(x,y)u=R(x,y)$, tìm các đường mà trên đó PDE trở thành một ODE. Dùng chữ $R$ cho vế phải để không nhầm với hàm tùy ý $F$.

\textbf{Bước 1 — Đọc các hệ số và lập hệ đặc trưng.} Đặt $U(s)=u(X(s),Y(s))$. Quy tắc dây chuyền cho $U'=u_xX'+u_yY'$. Chọn $X'=A$, $Y'=B$, khi đó

\[
X'=A(X,Y),\qquad Y'=B(X,Y),\qquad U'+C(X,Y)U=R(X,Y).
\]

Nếu $A\ne0$, có thể dùng $x$ làm tham số: $dy/dx=B/A$, $du/dx+(C/A)u=R/A$. Nếu $A=0$, chọn biến khác hoặc giữ hệ theo $s$; không chia cho hệ số bằng 0.

\textbf{Bước 2 — Tìm bất biến để phân biệt các đặc trưng.} Tích phân hai phương trình đầu để tìm $\rho(X,Y)$ không đổi. Kiểm tra $A\rho_x+B\rho_y=0$. Hệ số hằng cho $\rho=Bx-Ay$. Ví dụ $X'=X,Y'=-Y$ cho $(XY)'=0$; $X'=-Y,Y'=X$ cho $(X^2+Y^2)'=0$. Khi đặc trưng là vòng kín, cần kiểm tra nghiệm sau một vòng trở về cùng giá trị.

\textbf{Bước 3 — Giải ODE của ẩn trên từng đường.} Đặt $M(s)=\exp(\int C(X(s),Y(s))ds)$. Nhân phương trình của $U$ với $M$:

\[
(MU)'=MR,\qquad
U(s)=M(s)^{-1}\left[K+\int M(s)R(X(s),Y(s))ds\right].
\]

$K$ hằng trên một đường, nhưng thay đổi giữa các đường, nên viết $K=F(\rho)$. Nếu tìm được nghiệm riêng $u_p$, có thể giải phương trình thuần nhất cho $u-u_p$ để giảm tính toán. Nếu đổi biến $(\rho,z)$, phải tính Jacobian và dùng quy tắc dây chuyền để thấy PDE chỉ còn đạo hàm theo $z$.

\textbf{Bước 4 — Thế dữ liệu và tìm hàm tùy ý.} Với đường dữ liệu $(x_0(r),y_0(r))$, $u=h(r)$, thế vào nghiệm tổng quát để thu được $F(\rho_0(r))=H(r)$, trong đó $\rho_0=\rho(x_0,y_0)$. Nếu $\rho_0$ khả nghịch, đặt $r=\rho_0^{-1}(z)$ rồi viết $F(z)=H(\rho_0^{-1}(z))$. Nếu hai giá trị $r_1,r_2$ có cùng $\rho_0$, phải có $H(r_1)=H(r_2)$. Nếu đường dữ liệu chính là đặc trưng, dữ liệu phải thỏa ODE trên đường đó và thường không xác định hết $F$.

\textbf{Bước 5 — Xác định miền tồn tại và duy nhất.} Điều kiện không đặc trưng là $x_0'B-y_0'A\ne0$ trên đường dữ liệu. Nó cho khả nghịch cục bộ dưới giả thiết trơn thích hợp; không tự cho duy nhất trên cả mặt phẳng. Chỉ những đặc trưng gặp dữ liệu mới được xác định. Tại gốc, trục hoặc đường hệ số suy biến, kiểm tra riêng và ghép các biểu thức đủ trơn.

\textbf{Bước 6 — Kiểm tra kết quả.} Tính $u_x,u_y$ bằng quy tắc dây chuyền, thay vào $Au_x+Bu_y+Cu$ và kiểm tra bằng $R$. Sau đó thế lại điều kiện đầu/biên. Nêu rõ nếu chỉ có nghiệm cục bộ, nghiệm từng miền hoặc không có nghiệm trong lớp trơn yêu cầu.

\textbf{Chi tiết bước tích phân ODE.} Sau khi tìm bất biến $\rho$, viết $U'+P(s,\rho)U=Q(s,\rho)$. Chọn $M=\exp(\int_{s_0}^{s}P(r,\rho)dr)$; tính $M'=PM$, nên $MU'+M'U=MQ$. Tích phân từ $s_0$: $M(s)U(s)-U(s_0)=\int_{s_0}^sM(r)Q(r)dr$. Hằng số $U(s_0)$ phụ thuộc $\rho$; thay dữ liệu để tìm nó rồi mới đổi về tọa độ gốc.
\subsection*{Lời giải áp dụng cho bài}
\textbf{Bước 1 — Đặc trưng.} Với $u_x-2u_y=0$, $dy/dx=-2$, nên $\rho=y+2x$ và $u=F(\rho)$.

\textbf{Bước 2 — Dùng dữ liệu.} Trên $y=e^x$, $F(e^x+2x)=(e^x+2x)^2$. Hàm $\phi(x)=e^x+2x$ có $\phi'=e^x+2>0$; khi $x\to-\infty$ thì $\phi\to-\infty$, khi $x\to+\infty$ thì $\phi\to+\infty$. Vì vậy mọi chỉ số đặc trưng được gặp đúng một lần, $F(s)=s^2$.

\textbf{Bước 3 — Kiểm tra.} $u=(y+2x)^2$ có $u_x=4(y+2x)$, $u_y=2(y+2x)$, nên $u_x-2u_y=0$; vết trên $y=e^x$ đúng dữ liệu. Nghiệm duy nhất trên cả mặt phẳng.

\clearpage\noindent\texttt{\detokenize{BT_1_3_1_3_10}}\par
\section*{Lời giải bổ sung — Bài tập 10 · Mục 1.3.1.3}

\subsection*{Kiến thức và trình tự giải}
\textbf{Mục tiêu.} Với $A(x,y)u_x+B(x,y)u_y+C(x,y)u=R(x,y)$, tìm các đường mà trên đó PDE trở thành một ODE. Dùng chữ $R$ cho vế phải để không nhầm với hàm tùy ý $F$.

\textbf{Bước 1 — Đọc các hệ số và lập hệ đặc trưng.} Đặt $U(s)=u(X(s),Y(s))$. Quy tắc dây chuyền cho $U'=u_xX'+u_yY'$. Chọn $X'=A$, $Y'=B$, khi đó

\[
X'=A(X,Y),\qquad Y'=B(X,Y),\qquad U'+C(X,Y)U=R(X,Y).
\]

Nếu $A\ne0$, có thể dùng $x$ làm tham số: $dy/dx=B/A$, $du/dx+(C/A)u=R/A$. Nếu $A=0$, chọn biến khác hoặc giữ hệ theo $s$; không chia cho hệ số bằng 0.

\textbf{Bước 2 — Tìm bất biến để phân biệt các đặc trưng.} Tích phân hai phương trình đầu để tìm $\rho(X,Y)$ không đổi. Kiểm tra $A\rho_x+B\rho_y=0$. Hệ số hằng cho $\rho=Bx-Ay$. Ví dụ $X'=X,Y'=-Y$ cho $(XY)'=0$; $X'=-Y,Y'=X$ cho $(X^2+Y^2)'=0$. Khi đặc trưng là vòng kín, cần kiểm tra nghiệm sau một vòng trở về cùng giá trị.

\textbf{Bước 3 — Giải ODE của ẩn trên từng đường.} Đặt $M(s)=\exp(\int C(X(s),Y(s))ds)$. Nhân phương trình của $U$ với $M$:

\[
(MU)'=MR,\qquad
U(s)=M(s)^{-1}\left[K+\int M(s)R(X(s),Y(s))ds\right].
\]

$K$ hằng trên một đường, nhưng thay đổi giữa các đường, nên viết $K=F(\rho)$. Nếu tìm được nghiệm riêng $u_p$, có thể giải phương trình thuần nhất cho $u-u_p$ để giảm tính toán. Nếu đổi biến $(\rho,z)$, phải tính Jacobian và dùng quy tắc dây chuyền để thấy PDE chỉ còn đạo hàm theo $z$.

\textbf{Bước 4 — Thế dữ liệu và tìm hàm tùy ý.} Với đường dữ liệu $(x_0(r),y_0(r))$, $u=h(r)$, thế vào nghiệm tổng quát để thu được $F(\rho_0(r))=H(r)$, trong đó $\rho_0=\rho(x_0,y_0)$. Nếu $\rho_0$ khả nghịch, đặt $r=\rho_0^{-1}(z)$ rồi viết $F(z)=H(\rho_0^{-1}(z))$. Nếu hai giá trị $r_1,r_2$ có cùng $\rho_0$, phải có $H(r_1)=H(r_2)$. Nếu đường dữ liệu chính là đặc trưng, dữ liệu phải thỏa ODE trên đường đó và thường không xác định hết $F$.

\textbf{Bước 5 — Xác định miền tồn tại và duy nhất.} Điều kiện không đặc trưng là $x_0'B-y_0'A\ne0$ trên đường dữ liệu. Nó cho khả nghịch cục bộ dưới giả thiết trơn thích hợp; không tự cho duy nhất trên cả mặt phẳng. Chỉ những đặc trưng gặp dữ liệu mới được xác định. Tại gốc, trục hoặc đường hệ số suy biến, kiểm tra riêng và ghép các biểu thức đủ trơn.

\textbf{Bước 6 — Kiểm tra kết quả.} Tính $u_x,u_y$ bằng quy tắc dây chuyền, thay vào $Au_x+Bu_y+Cu$ và kiểm tra bằng $R$. Sau đó thế lại điều kiện đầu/biên. Nêu rõ nếu chỉ có nghiệm cục bộ, nghiệm từng miền hoặc không có nghiệm trong lớp trơn yêu cầu.

\textbf{Chi tiết bước tích phân ODE.} Sau khi tìm bất biến $\rho$, viết $U'+P(s,\rho)U=Q(s,\rho)$. Chọn $M=\exp(\int_{s_0}^{s}P(r,\rho)dr)$; tính $M'=PM$, nên $MU'+M'U=MQ$. Tích phân từ $s_0$: $M(s)U(s)-U(s_0)=\int_{s_0}^sM(r)Q(r)dr$. Hằng số $U(s_0)$ phụ thuộc $\rho$; thay dữ liệu để tìm nó rồi mới đổi về tọa độ gốc.
\subsection*{Lời giải áp dụng cho bài}
\textbf{Bước 1 — Nghiệm riêng và thuần nhất.} Với $3u_x-2u_y+u=x$, thử $u_p=x-3$ và có $3+(x-3)=x$. Đặt $w=u-u_p$, thì $3w_x-2w_y+w=0$. Bất biến $\rho=2x+3y$, ODE theo $x$ là $W'+W/3=0$, nên
\[u=x-3+e^{-x/3}F(2x+3y).\]

\textbf{Ý (a).} Trên $y=x$, $u=x$ đòi $e^{-x/3}F(5x)=3$. Đặt $s=5x$, $F(s)=3e^{s/15}$, suy
\[u=x-3+3e^{(y-x)/5}.\]
Thế $y=x$ cho $x$; $3u_x-2u_y+u=x$ vì phần mũ có hệ số $-3/5-2/5+1=0$.

\textbf{Ý (b).} Trên $2x+3y=1$, dữ liệu $u=0$ đòi $F(1)=(3-x)e^{x/3}$. Vế trái hằng nhưng đạo hàm vế phải là $-xe^{x/3}/3$, không đồng nhất bằng 0. Do đó không có nghiệm trên toàn đường. Đây là đường đặc trưng, nên phải kiểm tra tính tương thích thay vì khử $F$ như ở (a).

\clearpage\noindent\texttt{\detokenize{BT_1_3_1_3_11}}\par
\section*{Lời giải bổ sung — Bài tập 11 · Mục 1.3.1.3}

\subsection*{Kiến thức và trình tự giải}
\textbf{Mục tiêu.} Với $A(x,y)u_x+B(x,y)u_y+C(x,y)u=R(x,y)$, tìm các đường mà trên đó PDE trở thành một ODE. Dùng chữ $R$ cho vế phải để không nhầm với hàm tùy ý $F$.

\textbf{Bước 1 — Đọc các hệ số và lập hệ đặc trưng.} Đặt $U(s)=u(X(s),Y(s))$. Quy tắc dây chuyền cho $U'=u_xX'+u_yY'$. Chọn $X'=A$, $Y'=B$, khi đó

\[
X'=A(X,Y),\qquad Y'=B(X,Y),\qquad U'+C(X,Y)U=R(X,Y).
\]

Nếu $A\ne0$, có thể dùng $x$ làm tham số: $dy/dx=B/A$, $du/dx+(C/A)u=R/A$. Nếu $A=0$, chọn biến khác hoặc giữ hệ theo $s$; không chia cho hệ số bằng 0.

\textbf{Bước 2 — Tìm bất biến để phân biệt các đặc trưng.} Tích phân hai phương trình đầu để tìm $\rho(X,Y)$ không đổi. Kiểm tra $A\rho_x+B\rho_y=0$. Hệ số hằng cho $\rho=Bx-Ay$. Ví dụ $X'=X,Y'=-Y$ cho $(XY)'=0$; $X'=-Y,Y'=X$ cho $(X^2+Y^2)'=0$. Khi đặc trưng là vòng kín, cần kiểm tra nghiệm sau một vòng trở về cùng giá trị.

\textbf{Bước 3 — Giải ODE của ẩn trên từng đường.} Đặt $M(s)=\exp(\int C(X(s),Y(s))ds)$. Nhân phương trình của $U$ với $M$:

\[
(MU)'=MR,\qquad
U(s)=M(s)^{-1}\left[K+\int M(s)R(X(s),Y(s))ds\right].
\]

$K$ hằng trên một đường, nhưng thay đổi giữa các đường, nên viết $K=F(\rho)$. Nếu tìm được nghiệm riêng $u_p$, có thể giải phương trình thuần nhất cho $u-u_p$ để giảm tính toán. Nếu đổi biến $(\rho,z)$, phải tính Jacobian và dùng quy tắc dây chuyền để thấy PDE chỉ còn đạo hàm theo $z$.

\textbf{Bước 4 — Thế dữ liệu và tìm hàm tùy ý.} Với đường dữ liệu $(x_0(r),y_0(r))$, $u=h(r)$, thế vào nghiệm tổng quát để thu được $F(\rho_0(r))=H(r)$, trong đó $\rho_0=\rho(x_0,y_0)$. Nếu $\rho_0$ khả nghịch, đặt $r=\rho_0^{-1}(z)$ rồi viết $F(z)=H(\rho_0^{-1}(z))$. Nếu hai giá trị $r_1,r_2$ có cùng $\rho_0$, phải có $H(r_1)=H(r_2)$. Nếu đường dữ liệu chính là đặc trưng, dữ liệu phải thỏa ODE trên đường đó và thường không xác định hết $F$.

\textbf{Bước 5 — Xác định miền tồn tại và duy nhất.} Điều kiện không đặc trưng là $x_0'B-y_0'A\ne0$ trên đường dữ liệu. Nó cho khả nghịch cục bộ dưới giả thiết trơn thích hợp; không tự cho duy nhất trên cả mặt phẳng. Chỉ những đặc trưng gặp dữ liệu mới được xác định. Tại gốc, trục hoặc đường hệ số suy biến, kiểm tra riêng và ghép các biểu thức đủ trơn.

\textbf{Bước 6 — Kiểm tra kết quả.} Tính $u_x,u_y$ bằng quy tắc dây chuyền, thay vào $Au_x+Bu_y+Cu$ và kiểm tra bằng $R$. Sau đó thế lại điều kiện đầu/biên. Nêu rõ nếu chỉ có nghiệm cục bộ, nghiệm từng miền hoặc không có nghiệm trong lớp trơn yêu cầu.

\textbf{Chi tiết bước tích phân ODE.} Sau khi tìm bất biến $\rho$, viết $U'+P(s,\rho)U=Q(s,\rho)$. Chọn $M=\exp(\int_{s_0}^{s}P(r,\rho)dr)$; tính $M'=PM$, nên $MU'+M'U=MQ$. Tích phân từ $s_0$: $M(s)U(s)-U(s_0)=\int_{s_0}^sM(r)Q(r)dr$. Hằng số $U(s_0)$ phụ thuộc $\rho$; thay dữ liệu để tìm nó rồi mới đổi về tọa độ gốc.
\subsection*{Lời giải áp dụng cho bài}
\textbf{Bước 1 — Nêu giả thiết mô hình.} Hiểu $y$ là tuổi, mỗi cá thể tăng tuổi với tốc độ 1. Mật độ ban đầu $P(y,0)=C(y)\ge0$, tổng dân số $N(t)=\int_0^\infty P(y,t)dy$, sinh ở tuổi 0 cho $P(0,t)=\alpha N(t)$. Nếu $y$ chỉ là tọa độ không gian thì đề còn thiếu luật vận chuyển; giả thiết tuổi là điều cần nói rõ.

\textbf{Bước 2 — Giữ cách ghi chết tuyệt đối $D=k$.} Cân bằng dọc tuổi cho $P_t+P_y=-k$. Đặc trưng từ dữ liệu đầu có $y=\xi+t$ và $P'= -k$, nên trên $y>t$,
\[P(y,t)=C(y-t)-kt.\]
Để mật độ không âm ở mọi tuổi với $t>0$ phải có $C(\xi)\ge kt$ ở mọi $\xi>0$. Khi đó $\int C\ge\int_0^\infty kt\,d\xi=\infty$, trái giả thiết hữu hạn. Vì vậy mô hình chết tuyệt đối không cho mật độ dân số không âm hữu hạn trên mọi tuổi.

\textbf{Bước 3 — Giải mô hình tỉ lệ chết nếu đó là ý định của đề.} Sửa thành $D=kP$, PDE là $P_t+P_y=-kP$. Trên một đặc trưng, $P$ giảm theo $e^{-kt}$. Từ dữ liệu đầu, $P=e^{-kt}C(y-t)$ khi $y\ge t$. Từ biên tuổi 0 ở thời điểm $t-y$, $P=e^{-ky}\alpha N(t-y)$ khi $y<t$.

\textbf{Bước 4 — Tính tổng dân số.} Thế hai nhánh vào tích phân:
\[N(t)=e^{-kt}N_0+\alpha\int_0^t e^{-k(t-s)}N(s)ds,\qquad N_0=\int_0^\infty C.\]
Nhân $e^{kt}$, đạo hàm được $(e^{kt}N)'=\alpha e^{kt}N$, nên $N=N_0e^{(\alpha-k)t}$. Suy
\[P(y,t)=\begin{cases}e^{-kt}C(y-t),&y\ge t,\\ \alpha N_0e^{(\alpha-k)t-\alpha y},&0\le y<t.\end{cases}\]

\textbf{Bước 5 — Kiểm tra và ghép.} Hai nhánh thỏa $P_t+P_y=-kP$; tại $t=0$ đúng $C$, tại $y=0$ đúng $\alpha N(t)$. Muốn liên tục qua $y=t$ cần $C(0)=\alpha N_0$; muốn $C^1$ qua đó cần thêm $C'(0)=-\alpha^2N_0$. Nếu không tương thích, công thức mô tả nghiệm từng miền/ nghiệm vận chuyển yếu. Đây là kết quả của mô hình đã sửa, không đổi ngầm dữ kiện gốc.

\clearpage\noindent\texttt{\detokenize{VD_1_3_2_2_01}}\par
\section*{Lời giải bổ sung — Ví dụ 1 · Mục 1.3.2.2}

\subsection*{Kiến thức và trình tự giải}
\textbf{Mục tiêu.} Với $A(x,y)u_x+B(x,y)u_y+C(x,y)u=R(x,y)$, tìm các đường mà trên đó PDE trở thành một ODE. Dùng chữ $R$ cho vế phải để không nhầm với hàm tùy ý $F$.

\textbf{Bước 1 — Đọc các hệ số và lập hệ đặc trưng.} Đặt $U(s)=u(X(s),Y(s))$. Quy tắc dây chuyền cho $U'=u_xX'+u_yY'$. Chọn $X'=A$, $Y'=B$, khi đó

\[
X'=A(X,Y),\qquad Y'=B(X,Y),\qquad U'+C(X,Y)U=R(X,Y).
\]

Nếu $A\ne0$, có thể dùng $x$ làm tham số: $dy/dx=B/A$, $du/dx+(C/A)u=R/A$. Nếu $A=0$, chọn biến khác hoặc giữ hệ theo $s$; không chia cho hệ số bằng 0.

\textbf{Bước 2 — Tìm bất biến để phân biệt các đặc trưng.} Tích phân hai phương trình đầu để tìm $\rho(X,Y)$ không đổi. Kiểm tra $A\rho_x+B\rho_y=0$. Hệ số hằng cho $\rho=Bx-Ay$. Ví dụ $X'=X,Y'=-Y$ cho $(XY)'=0$; $X'=-Y,Y'=X$ cho $(X^2+Y^2)'=0$. Khi đặc trưng là vòng kín, cần kiểm tra nghiệm sau một vòng trở về cùng giá trị.

\textbf{Bước 3 — Giải ODE của ẩn trên từng đường.} Đặt $M(s)=\exp(\int C(X(s),Y(s))ds)$. Nhân phương trình của $U$ với $M$:

\[
(MU)'=MR,\qquad
U(s)=M(s)^{-1}\left[K+\int M(s)R(X(s),Y(s))ds\right].
\]

$K$ hằng trên một đường, nhưng thay đổi giữa các đường, nên viết $K=F(\rho)$. Nếu tìm được nghiệm riêng $u_p$, có thể giải phương trình thuần nhất cho $u-u_p$ để giảm tính toán. Nếu đổi biến $(\rho,z)$, phải tính Jacobian và dùng quy tắc dây chuyền để thấy PDE chỉ còn đạo hàm theo $z$.

\textbf{Bước 4 — Thế dữ liệu và tìm hàm tùy ý.} Với đường dữ liệu $(x_0(r),y_0(r))$, $u=h(r)$, thế vào nghiệm tổng quát để thu được $F(\rho_0(r))=H(r)$, trong đó $\rho_0=\rho(x_0,y_0)$. Nếu $\rho_0$ khả nghịch, đặt $r=\rho_0^{-1}(z)$ rồi viết $F(z)=H(\rho_0^{-1}(z))$. Nếu hai giá trị $r_1,r_2$ có cùng $\rho_0$, phải có $H(r_1)=H(r_2)$. Nếu đường dữ liệu chính là đặc trưng, dữ liệu phải thỏa ODE trên đường đó và thường không xác định hết $F$.

\textbf{Bước 5 — Xác định miền tồn tại và duy nhất.} Điều kiện không đặc trưng là $x_0'B-y_0'A\ne0$ trên đường dữ liệu. Nó cho khả nghịch cục bộ dưới giả thiết trơn thích hợp; không tự cho duy nhất trên cả mặt phẳng. Chỉ những đặc trưng gặp dữ liệu mới được xác định. Tại gốc, trục hoặc đường hệ số suy biến, kiểm tra riêng và ghép các biểu thức đủ trơn.

\textbf{Bước 6 — Kiểm tra kết quả.} Tính $u_x,u_y$ bằng quy tắc dây chuyền, thay vào $Au_x+Bu_y+Cu$ và kiểm tra bằng $R$. Sau đó thế lại điều kiện đầu/biên. Nêu rõ nếu chỉ có nghiệm cục bộ, nghiệm từng miền hoặc không có nghiệm trong lớp trơn yêu cầu.

\textbf{Chi tiết bước tích phân ODE.} Sau khi tìm bất biến $\rho$, viết $U'+P(s,\rho)U=Q(s,\rho)$. Chọn $M=\exp(\int_{s_0}^{s}P(r,\rho)dr)$; tính $M'=PM$, nên $MU'+M'U=MQ$. Tích phân từ $s_0$: $M(s)U(s)-U(s_0)=\int_{s_0}^sM(r)Q(r)dr$. Hằng số $U(s_0)$ phụ thuộc $\rho$; thay dữ liệu để tìm nó rồi mới đổi về tọa độ gốc.
\subsection*{Lời giải áp dụng cho bài}
\textbf{Bước 1 — Tìm bất biến của trường hệ số.} Với $Au_x+Bu_y+Cu=R$, giải đặc trưng $X'=A,Y'=B$ để tìm $\rho$ thỏa $A\rho_x+B\rho_y=0$.

\textbf{Bước 2 — Chọn tọa độ còn lại.} Đặt $w=\rho(x,y)$, $z=y$, $u=v(w,z)$. Jacobian là $\rho_x$; chỉ dùng phép đổi ở nơi $\rho_x\ne0$.

\textbf{Bước 3 — Tính đạo hàm và thế.}
\[u_x=v_w\rho_x,\qquad u_y=v_w\rho_y+v_z.\]
Do đó $Au_x+Bu_y=(A\rho_x+B\rho_y)v_w+Bv_z=Bv_z$. PDE thành $Bv_z+Cv=R$, với mọi hệ số được viết theo tọa độ mới. Nếu $B\ne0$, chia và giải ODE tuyến tính theo $z$, xem $w$ là tham số.

\textbf{Bước 4 — Ngoại lệ và dữ liệu.} Nếu $\rho_x=0$, có thể chọn $z=x$ khi $\rho_y\ne0$, hoặc tìm tọa độ khác. Sau giải ODE, đổi về $x,y$, thế đường dữ liệu để tìm hàm tùy ý và kiểm tra PDE. Không dùng phép đổi suy biến để suy tính duy nhất.

\clearpage\noindent\texttt{\detokenize{VD_1_3_2_2_02}}\par
\section*{Lời giải bổ sung — Ví dụ 2 · Mục 1.3.2.2}

\subsection*{Kiến thức và trình tự giải}
\textbf{Mục tiêu.} Với $A(x,y)u_x+B(x,y)u_y+C(x,y)u=R(x,y)$, tìm các đường mà trên đó PDE trở thành một ODE. Dùng chữ $R$ cho vế phải để không nhầm với hàm tùy ý $F$.

\textbf{Bước 1 — Đọc các hệ số và lập hệ đặc trưng.} Đặt $U(s)=u(X(s),Y(s))$. Quy tắc dây chuyền cho $U'=u_xX'+u_yY'$. Chọn $X'=A$, $Y'=B$, khi đó

\[
X'=A(X,Y),\qquad Y'=B(X,Y),\qquad U'+C(X,Y)U=R(X,Y).
\]

Nếu $A\ne0$, có thể dùng $x$ làm tham số: $dy/dx=B/A$, $du/dx+(C/A)u=R/A$. Nếu $A=0$, chọn biến khác hoặc giữ hệ theo $s$; không chia cho hệ số bằng 0.

\textbf{Bước 2 — Tìm bất biến để phân biệt các đặc trưng.} Tích phân hai phương trình đầu để tìm $\rho(X,Y)$ không đổi. Kiểm tra $A\rho_x+B\rho_y=0$. Hệ số hằng cho $\rho=Bx-Ay$. Ví dụ $X'=X,Y'=-Y$ cho $(XY)'=0$; $X'=-Y,Y'=X$ cho $(X^2+Y^2)'=0$. Khi đặc trưng là vòng kín, cần kiểm tra nghiệm sau một vòng trở về cùng giá trị.

\textbf{Bước 3 — Giải ODE của ẩn trên từng đường.} Đặt $M(s)=\exp(\int C(X(s),Y(s))ds)$. Nhân phương trình của $U$ với $M$:

\[
(MU)'=MR,\qquad
U(s)=M(s)^{-1}\left[K+\int M(s)R(X(s),Y(s))ds\right].
\]

$K$ hằng trên một đường, nhưng thay đổi giữa các đường, nên viết $K=F(\rho)$. Nếu tìm được nghiệm riêng $u_p$, có thể giải phương trình thuần nhất cho $u-u_p$ để giảm tính toán. Nếu đổi biến $(\rho,z)$, phải tính Jacobian và dùng quy tắc dây chuyền để thấy PDE chỉ còn đạo hàm theo $z$.

\textbf{Bước 4 — Thế dữ liệu và tìm hàm tùy ý.} Với đường dữ liệu $(x_0(r),y_0(r))$, $u=h(r)$, thế vào nghiệm tổng quát để thu được $F(\rho_0(r))=H(r)$, trong đó $\rho_0=\rho(x_0,y_0)$. Nếu $\rho_0$ khả nghịch, đặt $r=\rho_0^{-1}(z)$ rồi viết $F(z)=H(\rho_0^{-1}(z))$. Nếu hai giá trị $r_1,r_2$ có cùng $\rho_0$, phải có $H(r_1)=H(r_2)$. Nếu đường dữ liệu chính là đặc trưng, dữ liệu phải thỏa ODE trên đường đó và thường không xác định hết $F$.

\textbf{Bước 5 — Xác định miền tồn tại và duy nhất.} Điều kiện không đặc trưng là $x_0'B-y_0'A\ne0$ trên đường dữ liệu. Nó cho khả nghịch cục bộ dưới giả thiết trơn thích hợp; không tự cho duy nhất trên cả mặt phẳng. Chỉ những đặc trưng gặp dữ liệu mới được xác định. Tại gốc, trục hoặc đường hệ số suy biến, kiểm tra riêng và ghép các biểu thức đủ trơn.

\textbf{Bước 6 — Kiểm tra kết quả.} Tính $u_x,u_y$ bằng quy tắc dây chuyền, thay vào $Au_x+Bu_y+Cu$ và kiểm tra bằng $R$. Sau đó thế lại điều kiện đầu/biên. Nêu rõ nếu chỉ có nghiệm cục bộ, nghiệm từng miền hoặc không có nghiệm trong lớp trơn yêu cầu.

\textbf{Chi tiết bước tích phân ODE.} Sau khi tìm bất biến $\rho$, viết $U'+P(s,\rho)U=Q(s,\rho)$. Chọn $M=\exp(\int_{s_0}^{s}P(r,\rho)dr)$; tính $M'=PM$, nên $MU'+M'U=MQ$. Tích phân từ $s_0$: $M(s)U(s)-U(s_0)=\int_{s_0}^sM(r)Q(r)dr$. Hằng số $U(s_0)$ phụ thuộc $\rho$; thay dữ liệu để tìm nó rồi mới đổi về tọa độ gốc.
\subsection*{Lời giải áp dụng cho bài}
\textbf{Ý (a) — $u_x+x^2u_y=0$.} Bước 1: $dy/dx=x^2$ cho $y=x^3/3+\rho$. Bước 2: $U'=0$, nên $u=F(y-x^3/3)$. Bước 3: $u_x=-x^2F'$, $u_y=F'$, thế vào PDE được 0.

\textbf{Ý (b) — $e^{x^2}u_x+xu_y=0$.} Vì $e^{x^2}>0$, chọn $x$ làm tham số. $dy/dx=xe^{-x^2}$; dùng $d(e^{-x^2})=-2xe^{-x^2}dx$, suy $y=-e^{-x^2}/2+\rho/2$. Do đó $u=F(2y+e^{-x^2})$. Kiểm tra $u_x=-2xe^{-x^2}F'$, $u_y=2F'$, nên $e^{x^2}u_x+xu_y=0$. Cả hai công thức hợp lệ trên toàn mặt phẳng với $F\in C^1$.

\clearpage\noindent\texttt{\detokenize{VD_1_3_2_2_03}}\par
\section*{Lời giải bổ sung — Ví dụ 3 · Mục 1.3.2.2}

\subsection*{Kiến thức và trình tự giải}
\textbf{Mục tiêu.} Với $A(x,y)u_x+B(x,y)u_y+C(x,y)u=R(x,y)$, tìm các đường mà trên đó PDE trở thành một ODE. Dùng chữ $R$ cho vế phải để không nhầm với hàm tùy ý $F$.

\textbf{Bước 1 — Đọc các hệ số và lập hệ đặc trưng.} Đặt $U(s)=u(X(s),Y(s))$. Quy tắc dây chuyền cho $U'=u_xX'+u_yY'$. Chọn $X'=A$, $Y'=B$, khi đó

\[
X'=A(X,Y),\qquad Y'=B(X,Y),\qquad U'+C(X,Y)U=R(X,Y).
\]

Nếu $A\ne0$, có thể dùng $x$ làm tham số: $dy/dx=B/A$, $du/dx+(C/A)u=R/A$. Nếu $A=0$, chọn biến khác hoặc giữ hệ theo $s$; không chia cho hệ số bằng 0.

\textbf{Bước 2 — Tìm bất biến để phân biệt các đặc trưng.} Tích phân hai phương trình đầu để tìm $\rho(X,Y)$ không đổi. Kiểm tra $A\rho_x+B\rho_y=0$. Hệ số hằng cho $\rho=Bx-Ay$. Ví dụ $X'=X,Y'=-Y$ cho $(XY)'=0$; $X'=-Y,Y'=X$ cho $(X^2+Y^2)'=0$. Khi đặc trưng là vòng kín, cần kiểm tra nghiệm sau một vòng trở về cùng giá trị.

\textbf{Bước 3 — Giải ODE của ẩn trên từng đường.} Đặt $M(s)=\exp(\int C(X(s),Y(s))ds)$. Nhân phương trình của $U$ với $M$:

\[
(MU)'=MR,\qquad
U(s)=M(s)^{-1}\left[K+\int M(s)R(X(s),Y(s))ds\right].
\]

$K$ hằng trên một đường, nhưng thay đổi giữa các đường, nên viết $K=F(\rho)$. Nếu tìm được nghiệm riêng $u_p$, có thể giải phương trình thuần nhất cho $u-u_p$ để giảm tính toán. Nếu đổi biến $(\rho,z)$, phải tính Jacobian và dùng quy tắc dây chuyền để thấy PDE chỉ còn đạo hàm theo $z$.

\textbf{Bước 4 — Thế dữ liệu và tìm hàm tùy ý.} Với đường dữ liệu $(x_0(r),y_0(r))$, $u=h(r)$, thế vào nghiệm tổng quát để thu được $F(\rho_0(r))=H(r)$, trong đó $\rho_0=\rho(x_0,y_0)$. Nếu $\rho_0$ khả nghịch, đặt $r=\rho_0^{-1}(z)$ rồi viết $F(z)=H(\rho_0^{-1}(z))$. Nếu hai giá trị $r_1,r_2$ có cùng $\rho_0$, phải có $H(r_1)=H(r_2)$. Nếu đường dữ liệu chính là đặc trưng, dữ liệu phải thỏa ODE trên đường đó và thường không xác định hết $F$.

\textbf{Bước 5 — Xác định miền tồn tại và duy nhất.} Điều kiện không đặc trưng là $x_0'B-y_0'A\ne0$ trên đường dữ liệu. Nó cho khả nghịch cục bộ dưới giả thiết trơn thích hợp; không tự cho duy nhất trên cả mặt phẳng. Chỉ những đặc trưng gặp dữ liệu mới được xác định. Tại gốc, trục hoặc đường hệ số suy biến, kiểm tra riêng và ghép các biểu thức đủ trơn.

\textbf{Bước 6 — Kiểm tra kết quả.} Tính $u_x,u_y$ bằng quy tắc dây chuyền, thay vào $Au_x+Bu_y+Cu$ và kiểm tra bằng $R$. Sau đó thế lại điều kiện đầu/biên. Nêu rõ nếu chỉ có nghiệm cục bộ, nghiệm từng miền hoặc không có nghiệm trong lớp trơn yêu cầu.

\textbf{Chi tiết bước tích phân ODE.} Sau khi tìm bất biến $\rho$, viết $U'+P(s,\rho)U=Q(s,\rho)$. Chọn $M=\exp(\int_{s_0}^{s}P(r,\rho)dr)$; tính $M'=PM$, nên $MU'+M'U=MQ$. Tích phân từ $s_0$: $M(s)U(s)-U(s_0)=\int_{s_0}^sM(r)Q(r)dr$. Hằng số $U(s_0)$ phụ thuộc $\rho$; thay dữ liệu để tìm nó rồi mới đổi về tọa độ gốc.
\subsection*{Lời giải áp dụng cho bài}
\textbf{Ý (a) — $u_x+x^2u_y=0$.} Giải $dy/dx=x^2$ cho $\rho=y-x^3/3$, rồi $U'=0$ cho $u=F(y-x^3/3)$. Kiểm tra $u_x=-x^2F'$, $u_y=F'$, hai hạng triệt tiêu.

\textbf{Ý (b) — $e^{x^2}u_x+xu_y=0$.} Vì hệ số $e^{x^2}>0$, giải $dy/dx=xe^{-x^2}$, được $\rho=y+e^{-x^2}/2$. Do $U'=0$, $u=F(\rho)$. Kiểm tra $u_x=-xe^{-x^2}F'$, $u_y=F'$, nên PDE bằng 0. Bất biến có thể nhân 2 và đổi tên $F$ để được dạng ở ví dụ 2.

\textbf{Ý (c) — $u_x+\sin x\,u_y=0$.} $dy/dx=\sin x$ cho $y=-\cos x+\rho$. Dọc đường $U'=0$, nên $u=F(y+\cos x)$. Đạo hàm $u_x=-\sin x F'$, $u_y=F'$ kiểm tra PDE.

\textbf{Ý (d) — $xu_x+yu_y=0$.} Giải hệ $X'=X,Y'=Y$: $X=x_0e^s,Y=y_0e^s$. Đặc trưng là các tia, không phải toàn đường thẳng qua gốc; chiều co giãn không đổi dấu của tọa độ. Trên mỗi nửa mặt phẳng $x\ne0$, $Y/X$ không đổi, nên $u=F(y/x)$; trên các miền $y\ne0$, dùng $u=G(x/y)$ để phủ các trục. Các hàm ở các miền cần ghép đủ trơn.

\textbf{Tại gốc.} Nếu yêu cầu nghiệm liên tục trên cả mặt phẳng, với mỗi $(x,y)$ ta có $u(x,y)=u(e^{-s}x,e^{-s}y)$. Cho $s\to\infty$, vế phải tiến tới $u(0,0)$, nên mọi điểm có cùng giá trị. Vì vậy nghiệm toàn mặt phẳng liên tục phải hằng.

\clearpage\noindent\texttt{\detokenize{VD_1_3_2_2_04}}\par
\section*{Lời giải bổ sung — Ví dụ 4 · Mục 1.3.2.2}

\subsection*{Kiến thức và trình tự giải}
\textbf{Mục tiêu.} Với $A(x,y)u_x+B(x,y)u_y+C(x,y)u=R(x,y)$, tìm các đường mà trên đó PDE trở thành một ODE. Dùng chữ $R$ cho vế phải để không nhầm với hàm tùy ý $F$.

\textbf{Bước 1 — Đọc các hệ số và lập hệ đặc trưng.} Đặt $U(s)=u(X(s),Y(s))$. Quy tắc dây chuyền cho $U'=u_xX'+u_yY'$. Chọn $X'=A$, $Y'=B$, khi đó

\[
X'=A(X,Y),\qquad Y'=B(X,Y),\qquad U'+C(X,Y)U=R(X,Y).
\]

Nếu $A\ne0$, có thể dùng $x$ làm tham số: $dy/dx=B/A$, $du/dx+(C/A)u=R/A$. Nếu $A=0$, chọn biến khác hoặc giữ hệ theo $s$; không chia cho hệ số bằng 0.

\textbf{Bước 2 — Tìm bất biến để phân biệt các đặc trưng.} Tích phân hai phương trình đầu để tìm $\rho(X,Y)$ không đổi. Kiểm tra $A\rho_x+B\rho_y=0$. Hệ số hằng cho $\rho=Bx-Ay$. Ví dụ $X'=X,Y'=-Y$ cho $(XY)'=0$; $X'=-Y,Y'=X$ cho $(X^2+Y^2)'=0$. Khi đặc trưng là vòng kín, cần kiểm tra nghiệm sau một vòng trở về cùng giá trị.

\textbf{Bước 3 — Giải ODE của ẩn trên từng đường.} Đặt $M(s)=\exp(\int C(X(s),Y(s))ds)$. Nhân phương trình của $U$ với $M$:

\[
(MU)'=MR,\qquad
U(s)=M(s)^{-1}\left[K+\int M(s)R(X(s),Y(s))ds\right].
\]

$K$ hằng trên một đường, nhưng thay đổi giữa các đường, nên viết $K=F(\rho)$. Nếu tìm được nghiệm riêng $u_p$, có thể giải phương trình thuần nhất cho $u-u_p$ để giảm tính toán. Nếu đổi biến $(\rho,z)$, phải tính Jacobian và dùng quy tắc dây chuyền để thấy PDE chỉ còn đạo hàm theo $z$.

\textbf{Bước 4 — Thế dữ liệu và tìm hàm tùy ý.} Với đường dữ liệu $(x_0(r),y_0(r))$, $u=h(r)$, thế vào nghiệm tổng quát để thu được $F(\rho_0(r))=H(r)$, trong đó $\rho_0=\rho(x_0,y_0)$. Nếu $\rho_0$ khả nghịch, đặt $r=\rho_0^{-1}(z)$ rồi viết $F(z)=H(\rho_0^{-1}(z))$. Nếu hai giá trị $r_1,r_2$ có cùng $\rho_0$, phải có $H(r_1)=H(r_2)$. Nếu đường dữ liệu chính là đặc trưng, dữ liệu phải thỏa ODE trên đường đó và thường không xác định hết $F$.

\textbf{Bước 5 — Xác định miền tồn tại và duy nhất.} Điều kiện không đặc trưng là $x_0'B-y_0'A\ne0$ trên đường dữ liệu. Nó cho khả nghịch cục bộ dưới giả thiết trơn thích hợp; không tự cho duy nhất trên cả mặt phẳng. Chỉ những đặc trưng gặp dữ liệu mới được xác định. Tại gốc, trục hoặc đường hệ số suy biến, kiểm tra riêng và ghép các biểu thức đủ trơn.

\textbf{Bước 6 — Kiểm tra kết quả.} Tính $u_x,u_y$ bằng quy tắc dây chuyền, thay vào $Au_x+Bu_y+Cu$ và kiểm tra bằng $R$. Sau đó thế lại điều kiện đầu/biên. Nêu rõ nếu chỉ có nghiệm cục bộ, nghiệm từng miền hoặc không có nghiệm trong lớp trơn yêu cầu.

\textbf{Chi tiết bước tích phân ODE.} Sau khi tìm bất biến $\rho$, viết $U'+P(s,\rho)U=Q(s,\rho)$. Chọn $M=\exp(\int_{s_0}^{s}P(r,\rho)dr)$; tính $M'=PM$, nên $MU'+M'U=MQ$. Tích phân từ $s_0$: $M(s)U(s)-U(s_0)=\int_{s_0}^sM(r)Q(r)dr$. Hằng số $U(s_0)$ phụ thuộc $\rho$; thay dữ liệu để tìm nó rồi mới đổi về tọa độ gốc.
\subsection*{Lời giải áp dụng cho bài}
\textbf{Bước 1 — Tìm đặc trưng.} $xu_x-2yu_y+u=e^x$, $x>0$. Hệ $X'=X,Y'=-2Y$ cho $(X^2Y)'=0$, nên $\rho=x^2y$. Cách tính này bao gồm $y=0$, không cần lấy log của $y$.

\textbf{Bước 2 — Giải ODE.} Dùng $x$ làm tham số, $U'+U/x=e^x/x$. Vì $x>0$, thừa số tích phân là $x$:
\[(xU)'=e^x,\quad xU=e^x+F(\rho).\]
Suy $u=[e^x+F(x^2y)]/x$.

\textbf{Bước 3 — Kiểm tra.} Đặt $V=x\partial_x-2y\partial_y$. $V\rho=0$, $V(1/x)=-1/x$; do đó $(V+1)(F(\rho)/x)=0$. Với $e^x/x$, $(V+1)(e^x/x)=e^x$. Công thức thỏa PDE trên miền $x>0$.

\clearpage\noindent\texttt{\detokenize{VD_1_3_2_2_05}}\par
\section*{Lời giải bổ sung — Ví dụ 5 · Mục 1.3.2.2}

\subsection*{Kiến thức và trình tự giải}
\textbf{Mục tiêu.} Với $A(x,y)u_x+B(x,y)u_y+C(x,y)u=R(x,y)$, tìm các đường mà trên đó PDE trở thành một ODE. Dùng chữ $R$ cho vế phải để không nhầm với hàm tùy ý $F$.

\textbf{Bước 1 — Đọc các hệ số và lập hệ đặc trưng.} Đặt $U(s)=u(X(s),Y(s))$. Quy tắc dây chuyền cho $U'=u_xX'+u_yY'$. Chọn $X'=A$, $Y'=B$, khi đó

\[
X'=A(X,Y),\qquad Y'=B(X,Y),\qquad U'+C(X,Y)U=R(X,Y).
\]

Nếu $A\ne0$, có thể dùng $x$ làm tham số: $dy/dx=B/A$, $du/dx+(C/A)u=R/A$. Nếu $A=0$, chọn biến khác hoặc giữ hệ theo $s$; không chia cho hệ số bằng 0.

\textbf{Bước 2 — Tìm bất biến để phân biệt các đặc trưng.} Tích phân hai phương trình đầu để tìm $\rho(X,Y)$ không đổi. Kiểm tra $A\rho_x+B\rho_y=0$. Hệ số hằng cho $\rho=Bx-Ay$. Ví dụ $X'=X,Y'=-Y$ cho $(XY)'=0$; $X'=-Y,Y'=X$ cho $(X^2+Y^2)'=0$. Khi đặc trưng là vòng kín, cần kiểm tra nghiệm sau một vòng trở về cùng giá trị.

\textbf{Bước 3 — Giải ODE của ẩn trên từng đường.} Đặt $M(s)=\exp(\int C(X(s),Y(s))ds)$. Nhân phương trình của $U$ với $M$:

\[
(MU)'=MR,\qquad
U(s)=M(s)^{-1}\left[K+\int M(s)R(X(s),Y(s))ds\right].
\]

$K$ hằng trên một đường, nhưng thay đổi giữa các đường, nên viết $K=F(\rho)$. Nếu tìm được nghiệm riêng $u_p$, có thể giải phương trình thuần nhất cho $u-u_p$ để giảm tính toán. Nếu đổi biến $(\rho,z)$, phải tính Jacobian và dùng quy tắc dây chuyền để thấy PDE chỉ còn đạo hàm theo $z$.

\textbf{Bước 4 — Thế dữ liệu và tìm hàm tùy ý.} Với đường dữ liệu $(x_0(r),y_0(r))$, $u=h(r)$, thế vào nghiệm tổng quát để thu được $F(\rho_0(r))=H(r)$, trong đó $\rho_0=\rho(x_0,y_0)$. Nếu $\rho_0$ khả nghịch, đặt $r=\rho_0^{-1}(z)$ rồi viết $F(z)=H(\rho_0^{-1}(z))$. Nếu hai giá trị $r_1,r_2$ có cùng $\rho_0$, phải có $H(r_1)=H(r_2)$. Nếu đường dữ liệu chính là đặc trưng, dữ liệu phải thỏa ODE trên đường đó và thường không xác định hết $F$.

\textbf{Bước 5 — Xác định miền tồn tại và duy nhất.} Điều kiện không đặc trưng là $x_0'B-y_0'A\ne0$ trên đường dữ liệu. Nó cho khả nghịch cục bộ dưới giả thiết trơn thích hợp; không tự cho duy nhất trên cả mặt phẳng. Chỉ những đặc trưng gặp dữ liệu mới được xác định. Tại gốc, trục hoặc đường hệ số suy biến, kiểm tra riêng và ghép các biểu thức đủ trơn.

\textbf{Bước 6 — Kiểm tra kết quả.} Tính $u_x,u_y$ bằng quy tắc dây chuyền, thay vào $Au_x+Bu_y+Cu$ và kiểm tra bằng $R$. Sau đó thế lại điều kiện đầu/biên. Nêu rõ nếu chỉ có nghiệm cục bộ, nghiệm từng miền hoặc không có nghiệm trong lớp trơn yêu cầu.

\textbf{Chi tiết bước tích phân ODE.} Sau khi tìm bất biến $\rho$, viết $U'+P(s,\rho)U=Q(s,\rho)$. Chọn $M=\exp(\int_{s_0}^{s}P(r,\rho)dr)$; tính $M'=PM$, nên $MU'+M'U=MQ$. Tích phân từ $s_0$: $M(s)U(s)-U(s_0)=\int_{s_0}^sM(r)Q(r)dr$. Hằng số $U(s_0)$ phụ thuộc $\rho$; thay dữ liệu để tìm nó rồi mới đổi về tọa độ gốc.
\subsection*{Lời giải áp dụng cho bài}
\textbf{Bước 1 — Dùng nghiệm tổng quát của ví dụ 4.} $u=[e^x+F(x^2y)]/x$ trên $x>0$.

\textbf{Bước 2 — Thế dữ liệu $u(1,y)=y$.} $e+F(y)=y$, nên $F(s)=s-e$. Thay lại:
\[u(x,y)=xy+(e^x-e)/x.\]

\textbf{Bước 3 — Kiểm tra và duy nhất.} Trên $x=1$, công thức bằng $y$. Phần $xy$ thỏa phương trình thuần nhất vì $x(xy)_x-2y(xy)_y+xy=0$; phần còn lại cho nguồn $e^x$. Mỗi đặc trưng $x^2y=\rho$ trong $x>0$ gặp $x=1$ đúng một lần, nên duy nhất trên miền này.

\textbf{Chi tiết giải lại PDE.} Hệ $X'=X,Y'=-2Y$ cho $X=x_0e^s,Y=y_0e^{-2s}$, nên $X^2Y=x_0^2y_0=\rho$. Vì $dX/ds=X>0$, phương trình của ẩn theo $x$ là $dU/dx+U/x=e^x/x$. Nhân $x$ rồi tích phân: $d(xU)/dx=e^x$, $xU=e^x+F(\rho)$. Do đó $u=(e^x+F(x^2y))/x$, trước khi thế $x=1$ như trong bước 2.

\clearpage\noindent\texttt{\detokenize{VD_1_3_2_2_06}}\par
\section*{Lời giải bổ sung — Ví dụ 6 · Mục 1.3.2.2}

\subsection*{Kiến thức và trình tự giải}
\textbf{Mục tiêu.} Với $A(x,y)u_x+B(x,y)u_y+C(x,y)u=R(x,y)$, tìm các đường mà trên đó PDE trở thành một ODE. Dùng chữ $R$ cho vế phải để không nhầm với hàm tùy ý $F$.

\textbf{Bước 1 — Đọc các hệ số và lập hệ đặc trưng.} Đặt $U(s)=u(X(s),Y(s))$. Quy tắc dây chuyền cho $U'=u_xX'+u_yY'$. Chọn $X'=A$, $Y'=B$, khi đó

\[
X'=A(X,Y),\qquad Y'=B(X,Y),\qquad U'+C(X,Y)U=R(X,Y).
\]

Nếu $A\ne0$, có thể dùng $x$ làm tham số: $dy/dx=B/A$, $du/dx+(C/A)u=R/A$. Nếu $A=0$, chọn biến khác hoặc giữ hệ theo $s$; không chia cho hệ số bằng 0.

\textbf{Bước 2 — Tìm bất biến để phân biệt các đặc trưng.} Tích phân hai phương trình đầu để tìm $\rho(X,Y)$ không đổi. Kiểm tra $A\rho_x+B\rho_y=0$. Hệ số hằng cho $\rho=Bx-Ay$. Ví dụ $X'=X,Y'=-Y$ cho $(XY)'=0$; $X'=-Y,Y'=X$ cho $(X^2+Y^2)'=0$. Khi đặc trưng là vòng kín, cần kiểm tra nghiệm sau một vòng trở về cùng giá trị.

\textbf{Bước 3 — Giải ODE của ẩn trên từng đường.} Đặt $M(s)=\exp(\int C(X(s),Y(s))ds)$. Nhân phương trình của $U$ với $M$:

\[
(MU)'=MR,\qquad
U(s)=M(s)^{-1}\left[K+\int M(s)R(X(s),Y(s))ds\right].
\]

$K$ hằng trên một đường, nhưng thay đổi giữa các đường, nên viết $K=F(\rho)$. Nếu tìm được nghiệm riêng $u_p$, có thể giải phương trình thuần nhất cho $u-u_p$ để giảm tính toán. Nếu đổi biến $(\rho,z)$, phải tính Jacobian và dùng quy tắc dây chuyền để thấy PDE chỉ còn đạo hàm theo $z$.

\textbf{Bước 4 — Thế dữ liệu và tìm hàm tùy ý.} Với đường dữ liệu $(x_0(r),y_0(r))$, $u=h(r)$, thế vào nghiệm tổng quát để thu được $F(\rho_0(r))=H(r)$, trong đó $\rho_0=\rho(x_0,y_0)$. Nếu $\rho_0$ khả nghịch, đặt $r=\rho_0^{-1}(z)$ rồi viết $F(z)=H(\rho_0^{-1}(z))$. Nếu hai giá trị $r_1,r_2$ có cùng $\rho_0$, phải có $H(r_1)=H(r_2)$. Nếu đường dữ liệu chính là đặc trưng, dữ liệu phải thỏa ODE trên đường đó và thường không xác định hết $F$.

\textbf{Bước 5 — Xác định miền tồn tại và duy nhất.} Điều kiện không đặc trưng là $x_0'B-y_0'A\ne0$ trên đường dữ liệu. Nó cho khả nghịch cục bộ dưới giả thiết trơn thích hợp; không tự cho duy nhất trên cả mặt phẳng. Chỉ những đặc trưng gặp dữ liệu mới được xác định. Tại gốc, trục hoặc đường hệ số suy biến, kiểm tra riêng và ghép các biểu thức đủ trơn.

\textbf{Bước 6 — Kiểm tra kết quả.} Tính $u_x,u_y$ bằng quy tắc dây chuyền, thay vào $Au_x+Bu_y+Cu$ và kiểm tra bằng $R$. Sau đó thế lại điều kiện đầu/biên. Nêu rõ nếu chỉ có nghiệm cục bộ, nghiệm từng miền hoặc không có nghiệm trong lớp trơn yêu cầu.

\textbf{Chi tiết bước tích phân ODE.} Sau khi tìm bất biến $\rho$, viết $U'+P(s,\rho)U=Q(s,\rho)$. Chọn $M=\exp(\int_{s_0}^{s}P(r,\rho)dr)$; tính $M'=PM$, nên $MU'+M'U=MQ$. Tích phân từ $s_0$: $M(s)U(s)-U(s_0)=\int_{s_0}^sM(r)Q(r)dr$. Hằng số $U(s_0)$ phụ thuộc $\rho$; thay dữ liệu để tìm nó rồi mới đổi về tọa độ gốc.
\subsection*{Lời giải áp dụng cho bài}
\textbf{Bước 1 — Nhận ra đường dữ liệu là đặc trưng.} Trong nghiệm $u=[e^x+F(x^2y)]/x$, trên $y=0$ chỉ số $x^2y=0$ không đổi.

\textbf{Bước 2 — Tìm điều kiện tương thích.} $u(x,0)=f(x)$ đòi $xf(x)-e^x=F(0)=C_0$, tức $f(x)=(e^x+C_0)/x$ trên $x>0$. Đây là điều kiện cần và đủ.

\textbf{Bước 3 — Chỉ ra vô số nghiệm.} Nếu dữ liệu tương thích, chọn $F_\lambda(s)=C_0+\lambda s$ với mọi $\lambda\in\mathbb R$. Khi đó
\[u_\lambda=(e^x+C_0)/x+\lambda xy.\]
Các nghiệm đều thỏa PDE và vết $f$, nhưng khác nhau ngoài $y=0$. Có thể dùng mọi $F\in C^1$ với $F(0)=C_0$; họ vừa viết đủ chứng minh không duy nhất.

\textbf{Chi tiết nguồn gốc của hàm tùy ý.} Dọc $X'=X,Y'=-2Y$, tính $(X^2Y)'=2X^2Y-2X^2Y=0$. Phương trình $dU/dx+U/x=e^x/x$ cho $(xU)'=e^x$, rồi $xU=e^x+F(X^2Y)$. Đường $y=0$ chỉ gặp chỉ số $\rho=0$, nên nó chỉ quyết định $F(0)$; các giá trị $F(\rho)$ với $\rho\ne0$ không được quyết định.

\clearpage\noindent\texttt{\detokenize{BT_1_3_2_3_01}}\par
\section*{Lời giải bổ sung — Bài tập 1 · Mục 1.3.2.3}

\subsection*{Kiến thức và trình tự giải}
\textbf{Mục tiêu.} Với $A(x,y)u_x+B(x,y)u_y+C(x,y)u=R(x,y)$, tìm các đường mà trên đó PDE trở thành một ODE. Dùng chữ $R$ cho vế phải để không nhầm với hàm tùy ý $F$.

\textbf{Bước 1 — Đọc các hệ số và lập hệ đặc trưng.} Đặt $U(s)=u(X(s),Y(s))$. Quy tắc dây chuyền cho $U'=u_xX'+u_yY'$. Chọn $X'=A$, $Y'=B$, khi đó

\[
X'=A(X,Y),\qquad Y'=B(X,Y),\qquad U'+C(X,Y)U=R(X,Y).
\]

Nếu $A\ne0$, có thể dùng $x$ làm tham số: $dy/dx=B/A$, $du/dx+(C/A)u=R/A$. Nếu $A=0$, chọn biến khác hoặc giữ hệ theo $s$; không chia cho hệ số bằng 0.

\textbf{Bước 2 — Tìm bất biến để phân biệt các đặc trưng.} Tích phân hai phương trình đầu để tìm $\rho(X,Y)$ không đổi. Kiểm tra $A\rho_x+B\rho_y=0$. Hệ số hằng cho $\rho=Bx-Ay$. Ví dụ $X'=X,Y'=-Y$ cho $(XY)'=0$; $X'=-Y,Y'=X$ cho $(X^2+Y^2)'=0$. Khi đặc trưng là vòng kín, cần kiểm tra nghiệm sau một vòng trở về cùng giá trị.

\textbf{Bước 3 — Giải ODE của ẩn trên từng đường.} Đặt $M(s)=\exp(\int C(X(s),Y(s))ds)$. Nhân phương trình của $U$ với $M$:

\[
(MU)'=MR,\qquad
U(s)=M(s)^{-1}\left[K+\int M(s)R(X(s),Y(s))ds\right].
\]

$K$ hằng trên một đường, nhưng thay đổi giữa các đường, nên viết $K=F(\rho)$. Nếu tìm được nghiệm riêng $u_p$, có thể giải phương trình thuần nhất cho $u-u_p$ để giảm tính toán. Nếu đổi biến $(\rho,z)$, phải tính Jacobian và dùng quy tắc dây chuyền để thấy PDE chỉ còn đạo hàm theo $z$.

\textbf{Bước 4 — Thế dữ liệu và tìm hàm tùy ý.} Với đường dữ liệu $(x_0(r),y_0(r))$, $u=h(r)$, thế vào nghiệm tổng quát để thu được $F(\rho_0(r))=H(r)$, trong đó $\rho_0=\rho(x_0,y_0)$. Nếu $\rho_0$ khả nghịch, đặt $r=\rho_0^{-1}(z)$ rồi viết $F(z)=H(\rho_0^{-1}(z))$. Nếu hai giá trị $r_1,r_2$ có cùng $\rho_0$, phải có $H(r_1)=H(r_2)$. Nếu đường dữ liệu chính là đặc trưng, dữ liệu phải thỏa ODE trên đường đó và thường không xác định hết $F$.

\textbf{Bước 5 — Xác định miền tồn tại và duy nhất.} Điều kiện không đặc trưng là $x_0'B-y_0'A\ne0$ trên đường dữ liệu. Nó cho khả nghịch cục bộ dưới giả thiết trơn thích hợp; không tự cho duy nhất trên cả mặt phẳng. Chỉ những đặc trưng gặp dữ liệu mới được xác định. Tại gốc, trục hoặc đường hệ số suy biến, kiểm tra riêng và ghép các biểu thức đủ trơn.

\textbf{Bước 6 — Kiểm tra kết quả.} Tính $u_x,u_y$ bằng quy tắc dây chuyền, thay vào $Au_x+Bu_y+Cu$ và kiểm tra bằng $R$. Sau đó thế lại điều kiện đầu/biên. Nêu rõ nếu chỉ có nghiệm cục bộ, nghiệm từng miền hoặc không có nghiệm trong lớp trơn yêu cầu.

\textbf{Chi tiết bước tích phân ODE.} Sau khi tìm bất biến $\rho$, viết $U'+P(s,\rho)U=Q(s,\rho)$. Chọn $M=\exp(\int_{s_0}^{s}P(r,\rho)dr)$; tính $M'=PM$, nên $MU'+M'U=MQ$. Tích phân từ $s_0$: $M(s)U(s)-U(s_0)=\int_{s_0}^sM(r)Q(r)dr$. Hằng số $U(s_0)$ phụ thuộc $\rho$; thay dữ liệu để tìm nó rồi mới đổi về tọa độ gốc.
\subsection*{Lời giải áp dụng cho bài}
\textbf{Ý (a).} Với $u_x+\sin x\,u_y=0$, tích phân $dy/dx=\sin x$ cho $\rho=y+\cos x$. $U'=0$ nên $u=F(\rho)$. Kiểm tra $u_x+\sin x\,u_y=-\sin xF'+\sin xF'=0$.

\textbf{Ý (b).} Với $xu_x+yu_y=0$, $X=x_0e^s,Y=y_0e^s$, nên $u$ hằng trên tia. Trên từng miền $x>0$ hoặc $x<0$, $u=F(y/x)$; dùng tọa độ $x/y$ trên miền $y\ne0$. Nếu cần liên tục trên toàn mặt phẳng, co giãn mỗi điểm về gốc cho $u(x,y)=u(0,0)$, nên chỉ có nghiệm hằng.

\textbf{Ý (c).} Với $xu_x+2yu_y=0$ trong $x,y>0$, $X'=X,Y'=2Y$. Tính $(Y/X^2)'=0$, nên $\rho=y/x^2$. $U'=0$ cho $u=F(y/x^2)$. Kiểm tra $u_x=-2yF'/x^3$, $u_y=F'/x^2$, hai hạng của PDE triệt tiêu.

\textbf{Ý (d).} Với $yu_x-4xu_y=2xy$, trường $V=y\partial_x-4x\partial_y$ có bất biến $r=y^2+4x^2$ vì $Vr=2y(-4x)+8xy=0$. Chọn nghiệm riêng $P=-y^2/4$, có $VP=2xy$. Đặt $w=u-P$, $Vw=0$. Các mức $r>0$ là ellipse liên thông và đặc trưng kín, nên
\[u=-y^2/4+F(y^2+4x^2).\]
Đạo hàm của phần $F$ cho $V=0$; phần riêng cho đúng nguồn. Để có nghiệm toàn mặt phẳng, tổ hợp xuyên tâm ellipse phải đủ trơn tại gốc; $F\in C^1([0,\infty))$ là một điều kiện đủ, nhưng không phải mô tả duy nhất của lớp trơn ở gốc.

\textbf{Chi tiết hệ đặc trưng ở (d).} Giải $X'=Y,Y'=-4X$: $X''=-4X$, nên $X=A\cos2s+B\sin2s$, $Y=-2A\sin2s+2B\cos2s$. Vì $4X^2+Y^2=4(A^2+B^2)$, hệ đi hết một ellipse sau chu kỳ $\pi$. Với $U'=2XY$, có $(Y^2/4)'=YY'/2=-2XY$, do đó $(U+Y^2/4)'=0$. Hằng số trên ellipse là $F(4X^2+Y^2)$, dẫn tới $u=-y^2/4+F(4x^2+y^2)$.

\clearpage\noindent\texttt{\detokenize{BT_1_3_2_3_02}}\par
\section*{Lời giải bổ sung — Bài tập 2 · Mục 1.3.2.3}

\subsection*{Kiến thức và trình tự giải}
\textbf{Mục tiêu.} Với $A(x,y)u_x+B(x,y)u_y+C(x,y)u=R(x,y)$, tìm các đường mà trên đó PDE trở thành một ODE. Dùng chữ $R$ cho vế phải để không nhầm với hàm tùy ý $F$.

\textbf{Bước 1 — Đọc các hệ số và lập hệ đặc trưng.} Đặt $U(s)=u(X(s),Y(s))$. Quy tắc dây chuyền cho $U'=u_xX'+u_yY'$. Chọn $X'=A$, $Y'=B$, khi đó

\[
X'=A(X,Y),\qquad Y'=B(X,Y),\qquad U'+C(X,Y)U=R(X,Y).
\]

Nếu $A\ne0$, có thể dùng $x$ làm tham số: $dy/dx=B/A$, $du/dx+(C/A)u=R/A$. Nếu $A=0$, chọn biến khác hoặc giữ hệ theo $s$; không chia cho hệ số bằng 0.

\textbf{Bước 2 — Tìm bất biến để phân biệt các đặc trưng.} Tích phân hai phương trình đầu để tìm $\rho(X,Y)$ không đổi. Kiểm tra $A\rho_x+B\rho_y=0$. Hệ số hằng cho $\rho=Bx-Ay$. Ví dụ $X'=X,Y'=-Y$ cho $(XY)'=0$; $X'=-Y,Y'=X$ cho $(X^2+Y^2)'=0$. Khi đặc trưng là vòng kín, cần kiểm tra nghiệm sau một vòng trở về cùng giá trị.

\textbf{Bước 3 — Giải ODE của ẩn trên từng đường.} Đặt $M(s)=\exp(\int C(X(s),Y(s))ds)$. Nhân phương trình của $U$ với $M$:

\[
(MU)'=MR,\qquad
U(s)=M(s)^{-1}\left[K+\int M(s)R(X(s),Y(s))ds\right].
\]

$K$ hằng trên một đường, nhưng thay đổi giữa các đường, nên viết $K=F(\rho)$. Nếu tìm được nghiệm riêng $u_p$, có thể giải phương trình thuần nhất cho $u-u_p$ để giảm tính toán. Nếu đổi biến $(\rho,z)$, phải tính Jacobian và dùng quy tắc dây chuyền để thấy PDE chỉ còn đạo hàm theo $z$.

\textbf{Bước 4 — Thế dữ liệu và tìm hàm tùy ý.} Với đường dữ liệu $(x_0(r),y_0(r))$, $u=h(r)$, thế vào nghiệm tổng quát để thu được $F(\rho_0(r))=H(r)$, trong đó $\rho_0=\rho(x_0,y_0)$. Nếu $\rho_0$ khả nghịch, đặt $r=\rho_0^{-1}(z)$ rồi viết $F(z)=H(\rho_0^{-1}(z))$. Nếu hai giá trị $r_1,r_2$ có cùng $\rho_0$, phải có $H(r_1)=H(r_2)$. Nếu đường dữ liệu chính là đặc trưng, dữ liệu phải thỏa ODE trên đường đó và thường không xác định hết $F$.

\textbf{Bước 5 — Xác định miền tồn tại và duy nhất.} Điều kiện không đặc trưng là $x_0'B-y_0'A\ne0$ trên đường dữ liệu. Nó cho khả nghịch cục bộ dưới giả thiết trơn thích hợp; không tự cho duy nhất trên cả mặt phẳng. Chỉ những đặc trưng gặp dữ liệu mới được xác định. Tại gốc, trục hoặc đường hệ số suy biến, kiểm tra riêng và ghép các biểu thức đủ trơn.

\textbf{Bước 6 — Kiểm tra kết quả.} Tính $u_x,u_y$ bằng quy tắc dây chuyền, thay vào $Au_x+Bu_y+Cu$ và kiểm tra bằng $R$. Sau đó thế lại điều kiện đầu/biên. Nêu rõ nếu chỉ có nghiệm cục bộ, nghiệm từng miền hoặc không có nghiệm trong lớp trơn yêu cầu.

\textbf{Chi tiết bước tích phân ODE.} Sau khi tìm bất biến $\rho$, viết $U'+P(s,\rho)U=Q(s,\rho)$. Chọn $M=\exp(\int_{s_0}^{s}P(r,\rho)dr)$; tính $M'=PM$, nên $MU'+M'U=MQ$. Tích phân từ $s_0$: $M(s)U(s)-U(s_0)=\int_{s_0}^sM(r)Q(r)dr$. Hằng số $U(s_0)$ phụ thuộc $\rho$; thay dữ liệu để tìm nó rồi mới đổi về tọa độ gốc.
\subsection*{Lời giải áp dụng cho bài}
\textbf{Bước 1 — Đặc trưng.} $-yu_x+xu_y=0$ có $X'=-Y,Y'=X$. Khi đó $X^2+Y^2$ không đổi, đặc trưng là các vòng tròn tâm gốc và $U'=0$.

\textbf{Bước 2 — So sánh dữ liệu trên cùng vòng.} Vòng bán kính $r>0$ gặp trục $x$ ở $(r,0)$ và $(-r,0)$. Nghiệm phải có hai giá trị bằng nhau, trong khi dữ liệu $u(x,0)=3x$ cho $3r$ và $-3r$.

\textbf{Bước 3 — Kết luận.} Không có nghiệm trên toàn mặt phẳng. Ngay trên trục cũng thấy điều kiện cần $f(r)=f(-r)$ bị vi phạm. Nếu chỉ xét miền cục bộ quanh một phần đường dữ liệu, phải nói rõ miền; kết luận toàn cục không tự áp cho mọi miền nhỏ.

\textbf{Chi tiết tích phân hệ quay.} Từ $X'=-Y,Y'=X$, lấy đạo hàm dòng đầu được $X''+X=0$. Với điểm đầu $(x_0,y_0)$, $X=x_0\cos s-y_0\sin s$, $Y=x_0\sin s+y_0\cos s$. Tính $X^2+Y^2=x_0^2+y_0^2$ và $U'=0$ cho $U=u_0$. Sau góc $\pi$, $(r,0)$ đi tới $(-r,0)$ mà $U$ không đổi; đây là lý do hai giá trị dữ liệu phải bằng nhau.

\clearpage\noindent\texttt{\detokenize{BT_1_3_2_3_03}}\par
\section*{Lời giải bổ sung — Bài tập 3 · Mục 1.3.2.3}

\subsection*{Kiến thức và trình tự giải}
\textbf{Mục tiêu.} Với $A(x,y)u_x+B(x,y)u_y+C(x,y)u=R(x,y)$, tìm các đường mà trên đó PDE trở thành một ODE. Dùng chữ $R$ cho vế phải để không nhầm với hàm tùy ý $F$.

\textbf{Bước 1 — Đọc các hệ số và lập hệ đặc trưng.} Đặt $U(s)=u(X(s),Y(s))$. Quy tắc dây chuyền cho $U'=u_xX'+u_yY'$. Chọn $X'=A$, $Y'=B$, khi đó

\[
X'=A(X,Y),\qquad Y'=B(X,Y),\qquad U'+C(X,Y)U=R(X,Y).
\]

Nếu $A\ne0$, có thể dùng $x$ làm tham số: $dy/dx=B/A$, $du/dx+(C/A)u=R/A$. Nếu $A=0$, chọn biến khác hoặc giữ hệ theo $s$; không chia cho hệ số bằng 0.

\textbf{Bước 2 — Tìm bất biến để phân biệt các đặc trưng.} Tích phân hai phương trình đầu để tìm $\rho(X,Y)$ không đổi. Kiểm tra $A\rho_x+B\rho_y=0$. Hệ số hằng cho $\rho=Bx-Ay$. Ví dụ $X'=X,Y'=-Y$ cho $(XY)'=0$; $X'=-Y,Y'=X$ cho $(X^2+Y^2)'=0$. Khi đặc trưng là vòng kín, cần kiểm tra nghiệm sau một vòng trở về cùng giá trị.

\textbf{Bước 3 — Giải ODE của ẩn trên từng đường.} Đặt $M(s)=\exp(\int C(X(s),Y(s))ds)$. Nhân phương trình của $U$ với $M$:

\[
(MU)'=MR,\qquad
U(s)=M(s)^{-1}\left[K+\int M(s)R(X(s),Y(s))ds\right].
\]

$K$ hằng trên một đường, nhưng thay đổi giữa các đường, nên viết $K=F(\rho)$. Nếu tìm được nghiệm riêng $u_p$, có thể giải phương trình thuần nhất cho $u-u_p$ để giảm tính toán. Nếu đổi biến $(\rho,z)$, phải tính Jacobian và dùng quy tắc dây chuyền để thấy PDE chỉ còn đạo hàm theo $z$.

\textbf{Bước 4 — Thế dữ liệu và tìm hàm tùy ý.} Với đường dữ liệu $(x_0(r),y_0(r))$, $u=h(r)$, thế vào nghiệm tổng quát để thu được $F(\rho_0(r))=H(r)$, trong đó $\rho_0=\rho(x_0,y_0)$. Nếu $\rho_0$ khả nghịch, đặt $r=\rho_0^{-1}(z)$ rồi viết $F(z)=H(\rho_0^{-1}(z))$. Nếu hai giá trị $r_1,r_2$ có cùng $\rho_0$, phải có $H(r_1)=H(r_2)$. Nếu đường dữ liệu chính là đặc trưng, dữ liệu phải thỏa ODE trên đường đó và thường không xác định hết $F$.

\textbf{Bước 5 — Xác định miền tồn tại và duy nhất.} Điều kiện không đặc trưng là $x_0'B-y_0'A\ne0$ trên đường dữ liệu. Nó cho khả nghịch cục bộ dưới giả thiết trơn thích hợp; không tự cho duy nhất trên cả mặt phẳng. Chỉ những đặc trưng gặp dữ liệu mới được xác định. Tại gốc, trục hoặc đường hệ số suy biến, kiểm tra riêng và ghép các biểu thức đủ trơn.

\textbf{Bước 6 — Kiểm tra kết quả.} Tính $u_x,u_y$ bằng quy tắc dây chuyền, thay vào $Au_x+Bu_y+Cu$ và kiểm tra bằng $R$. Sau đó thế lại điều kiện đầu/biên. Nêu rõ nếu chỉ có nghiệm cục bộ, nghiệm từng miền hoặc không có nghiệm trong lớp trơn yêu cầu.

\textbf{Chi tiết bước tích phân ODE.} Sau khi tìm bất biến $\rho$, viết $U'+P(s,\rho)U=Q(s,\rho)$. Chọn $M=\exp(\int_{s_0}^{s}P(r,\rho)dr)$; tính $M'=PM$, nên $MU'+M'U=MQ$. Tích phân từ $s_0$: $M(s)U(s)-U(s_0)=\int_{s_0}^sM(r)Q(r)dr$. Hằng số $U(s_0)$ phụ thuộc $\rho$; thay dữ liệu để tìm nó rồi mới đổi về tọa độ gốc.
\subsection*{Lời giải áp dụng cho bài}
\textbf{Bước 1 — Hình dạng nghiệm.} Như bài trước, nghiệm hằng trên mỗi vòng tròn nên $u=H(r)$, $r=\sqrt{x^2+y^2}$, ngoài gốc.

\textbf{Bước 2 — Điều kiện dữ liệu.} Trên trục $x$, $H(|x|)=f(x)$, nên cần $f$ chẵn. Nếu $f\in C^1(\mathbb R)$ chẵn thì $f'(0)=0$, và đặt $H(r)=f(r)$, $r\ge0$.

\textbf{Bước 3 — Kiểm tra độ trơn và PDE.} Với $r>0$, $u_x=f'(r)x/r$, $u_y=f'(r)y/r$, nên $-yu_x+xu_y=0$. Do $f'(r)\to0$, gradient tiến về 0 tại gốc; đạo hàm tại gốc cũng bằng 0, nên công thức cho nghiệm $C^1$ trên toàn mặt phẳng.

\textbf{Bước 4 — Duy nhất.} Mọi bán kính gặp dữ liệu nên $H$ được xác định hết. Ví dụ $f=x^2$ cho $u=x^2+y^2$; $f$ không chẵn cho vô nghiệm. Với miền toàn mặt phẳng của đề, không có trường hợp dữ liệu tương thích mà không duy nhất; điều đó chỉ có thể xảy ra trên miền hoặc dữ liệu bị thu hẹp.

\textbf{Chi tiết từ hệ đặc trưng đến nghiệm xuyên tâm.} Hệ $X'=-Y,Y'=X,U'=0$ có nghiệm $X=x_0\cos s-y_0\sin s$, $Y=x_0\sin s+y_0\cos s$, $U=u_0$. Mỗi quỹ đạo $r>0$ đi hết vòng $X^2+Y^2=r^2$, nên giá trị chỉ có thể phụ thuộc $r$. Ở gốc đặt $u(0,0)=f(0)$. Vì $f$ chẵn và khả vi, $f'(0)=-f'(0)$, suy $f'(0)=0$. Hơn nữa $|u(x,y)-u(0,0)|/r=|f(r)-f(0)|/r\to0$, nên $u$ khả vi ở gốc với gradient $0$; các công thức gradient ngoài gốc tiến về cùng giá trị ấy.

\clearpage\noindent\texttt{\detokenize{BT_1_3_2_3_04}}\par
\section*{Lời giải bổ sung — Bài tập 4 · Mục 1.3.2.3}

\subsection*{Kiến thức và trình tự giải}
\textbf{Mục tiêu.} Với $A(x,y)u_x+B(x,y)u_y+C(x,y)u=R(x,y)$, tìm các đường mà trên đó PDE trở thành một ODE. Dùng chữ $R$ cho vế phải để không nhầm với hàm tùy ý $F$.

\textbf{Bước 1 — Đọc các hệ số và lập hệ đặc trưng.} Đặt $U(s)=u(X(s),Y(s))$. Quy tắc dây chuyền cho $U'=u_xX'+u_yY'$. Chọn $X'=A$, $Y'=B$, khi đó

\[
X'=A(X,Y),\qquad Y'=B(X,Y),\qquad U'+C(X,Y)U=R(X,Y).
\]

Nếu $A\ne0$, có thể dùng $x$ làm tham số: $dy/dx=B/A$, $du/dx+(C/A)u=R/A$. Nếu $A=0$, chọn biến khác hoặc giữ hệ theo $s$; không chia cho hệ số bằng 0.

\textbf{Bước 2 — Tìm bất biến để phân biệt các đặc trưng.} Tích phân hai phương trình đầu để tìm $\rho(X,Y)$ không đổi. Kiểm tra $A\rho_x+B\rho_y=0$. Hệ số hằng cho $\rho=Bx-Ay$. Ví dụ $X'=X,Y'=-Y$ cho $(XY)'=0$; $X'=-Y,Y'=X$ cho $(X^2+Y^2)'=0$. Khi đặc trưng là vòng kín, cần kiểm tra nghiệm sau một vòng trở về cùng giá trị.

\textbf{Bước 3 — Giải ODE của ẩn trên từng đường.} Đặt $M(s)=\exp(\int C(X(s),Y(s))ds)$. Nhân phương trình của $U$ với $M$:

\[
(MU)'=MR,\qquad
U(s)=M(s)^{-1}\left[K+\int M(s)R(X(s),Y(s))ds\right].
\]

$K$ hằng trên một đường, nhưng thay đổi giữa các đường, nên viết $K=F(\rho)$. Nếu tìm được nghiệm riêng $u_p$, có thể giải phương trình thuần nhất cho $u-u_p$ để giảm tính toán. Nếu đổi biến $(\rho,z)$, phải tính Jacobian và dùng quy tắc dây chuyền để thấy PDE chỉ còn đạo hàm theo $z$.

\textbf{Bước 4 — Thế dữ liệu và tìm hàm tùy ý.} Với đường dữ liệu $(x_0(r),y_0(r))$, $u=h(r)$, thế vào nghiệm tổng quát để thu được $F(\rho_0(r))=H(r)$, trong đó $\rho_0=\rho(x_0,y_0)$. Nếu $\rho_0$ khả nghịch, đặt $r=\rho_0^{-1}(z)$ rồi viết $F(z)=H(\rho_0^{-1}(z))$. Nếu hai giá trị $r_1,r_2$ có cùng $\rho_0$, phải có $H(r_1)=H(r_2)$. Nếu đường dữ liệu chính là đặc trưng, dữ liệu phải thỏa ODE trên đường đó và thường không xác định hết $F$.

\textbf{Bước 5 — Xác định miền tồn tại và duy nhất.} Điều kiện không đặc trưng là $x_0'B-y_0'A\ne0$ trên đường dữ liệu. Nó cho khả nghịch cục bộ dưới giả thiết trơn thích hợp; không tự cho duy nhất trên cả mặt phẳng. Chỉ những đặc trưng gặp dữ liệu mới được xác định. Tại gốc, trục hoặc đường hệ số suy biến, kiểm tra riêng và ghép các biểu thức đủ trơn.

\textbf{Bước 6 — Kiểm tra kết quả.} Tính $u_x,u_y$ bằng quy tắc dây chuyền, thay vào $Au_x+Bu_y+Cu$ và kiểm tra bằng $R$. Sau đó thế lại điều kiện đầu/biên. Nêu rõ nếu chỉ có nghiệm cục bộ, nghiệm từng miền hoặc không có nghiệm trong lớp trơn yêu cầu.

\textbf{Chi tiết bước tích phân ODE.} Sau khi tìm bất biến $\rho$, viết $U'+P(s,\rho)U=Q(s,\rho)$. Chọn $M=\exp(\int_{s_0}^{s}P(r,\rho)dr)$; tính $M'=PM$, nên $MU'+M'U=MQ$. Tích phân từ $s_0$: $M(s)U(s)-U(s_0)=\int_{s_0}^sM(r)Q(r)dr$. Hằng số $U(s_0)$ phụ thuộc $\rho$; thay dữ liệu để tìm nó rồi mới đổi về tọa độ gốc.
\subsection*{Lời giải áp dụng cho bài}
\textbf{Bước 1 — Giữ đúng hệ số của đề.} PDE là $xu_x-yu_y+yu=0$, có hệ số của $u$ là $y$. Hệ đặc trưng $X'=X,Y'=-Y$ cho $r=XY$ không đổi.

\textbf{Bước 2 — Giải ẩn.} $U'=-YU$, còn $Y'=-Y$, nên $(e^{-Y}U)'=e^{-Y}(U'-Y'U)=0$. Do đó
\[u(x,y)=e^yF(xy).\]
Trên mỗi miền $x>0$ hoặc $x<0$, có thể chọn tọa độ $(r,\log|x|)$, nên công thức mô tả nghiệm trên các phần đặc trưng đó. Khi muốn qua trục, ghép đủ trơn.

\textbf{Bước 3 — Kiểm tra.} $u_x=e^y yF'$, $u_y=e^y(F+xF')$, nên $xu_x-yu_y+yu=e^y[xyF'-yF-xyF'+yF]=0$. Trên $x=0,y\ne0$, PDE cho $u_y=u$, phù hợp $u(0,y)=e^yF(0)$; trên $y=0$, PDE cho $xu_x=0$.

\clearpage\noindent\texttt{\detokenize{BT_1_3_2_3_05}}\par
\section*{Lời giải bổ sung — Bài tập 5 · Mục 1.3.2.3}

\subsection*{Kiến thức và trình tự giải}
\textbf{Mục tiêu.} Với $A(x,y)u_x+B(x,y)u_y+C(x,y)u=R(x,y)$, tìm các đường mà trên đó PDE trở thành một ODE. Dùng chữ $R$ cho vế phải để không nhầm với hàm tùy ý $F$.

\textbf{Bước 1 — Đọc các hệ số và lập hệ đặc trưng.} Đặt $U(s)=u(X(s),Y(s))$. Quy tắc dây chuyền cho $U'=u_xX'+u_yY'$. Chọn $X'=A$, $Y'=B$, khi đó

\[
X'=A(X,Y),\qquad Y'=B(X,Y),\qquad U'+C(X,Y)U=R(X,Y).
\]

Nếu $A\ne0$, có thể dùng $x$ làm tham số: $dy/dx=B/A$, $du/dx+(C/A)u=R/A$. Nếu $A=0$, chọn biến khác hoặc giữ hệ theo $s$; không chia cho hệ số bằng 0.

\textbf{Bước 2 — Tìm bất biến để phân biệt các đặc trưng.} Tích phân hai phương trình đầu để tìm $\rho(X,Y)$ không đổi. Kiểm tra $A\rho_x+B\rho_y=0$. Hệ số hằng cho $\rho=Bx-Ay$. Ví dụ $X'=X,Y'=-Y$ cho $(XY)'=0$; $X'=-Y,Y'=X$ cho $(X^2+Y^2)'=0$. Khi đặc trưng là vòng kín, cần kiểm tra nghiệm sau một vòng trở về cùng giá trị.

\textbf{Bước 3 — Giải ODE của ẩn trên từng đường.} Đặt $M(s)=\exp(\int C(X(s),Y(s))ds)$. Nhân phương trình của $U$ với $M$:

\[
(MU)'=MR,\qquad
U(s)=M(s)^{-1}\left[K+\int M(s)R(X(s),Y(s))ds\right].
\]

$K$ hằng trên một đường, nhưng thay đổi giữa các đường, nên viết $K=F(\rho)$. Nếu tìm được nghiệm riêng $u_p$, có thể giải phương trình thuần nhất cho $u-u_p$ để giảm tính toán. Nếu đổi biến $(\rho,z)$, phải tính Jacobian và dùng quy tắc dây chuyền để thấy PDE chỉ còn đạo hàm theo $z$.

\textbf{Bước 4 — Thế dữ liệu và tìm hàm tùy ý.} Với đường dữ liệu $(x_0(r),y_0(r))$, $u=h(r)$, thế vào nghiệm tổng quát để thu được $F(\rho_0(r))=H(r)$, trong đó $\rho_0=\rho(x_0,y_0)$. Nếu $\rho_0$ khả nghịch, đặt $r=\rho_0^{-1}(z)$ rồi viết $F(z)=H(\rho_0^{-1}(z))$. Nếu hai giá trị $r_1,r_2$ có cùng $\rho_0$, phải có $H(r_1)=H(r_2)$. Nếu đường dữ liệu chính là đặc trưng, dữ liệu phải thỏa ODE trên đường đó và thường không xác định hết $F$.

\textbf{Bước 5 — Xác định miền tồn tại và duy nhất.} Điều kiện không đặc trưng là $x_0'B-y_0'A\ne0$ trên đường dữ liệu. Nó cho khả nghịch cục bộ dưới giả thiết trơn thích hợp; không tự cho duy nhất trên cả mặt phẳng. Chỉ những đặc trưng gặp dữ liệu mới được xác định. Tại gốc, trục hoặc đường hệ số suy biến, kiểm tra riêng và ghép các biểu thức đủ trơn.

\textbf{Bước 6 — Kiểm tra kết quả.} Tính $u_x,u_y$ bằng quy tắc dây chuyền, thay vào $Au_x+Bu_y+Cu$ và kiểm tra bằng $R$. Sau đó thế lại điều kiện đầu/biên. Nêu rõ nếu chỉ có nghiệm cục bộ, nghiệm từng miền hoặc không có nghiệm trong lớp trơn yêu cầu.

\textbf{Chi tiết bước tích phân ODE.} Sau khi tìm bất biến $\rho$, viết $U'+P(s,\rho)U=Q(s,\rho)$. Chọn $M=\exp(\int_{s_0}^{s}P(r,\rho)dr)$; tính $M'=PM$, nên $MU'+M'U=MQ$. Tích phân từ $s_0$: $M(s)U(s)-U(s_0)=\int_{s_0}^sM(r)Q(r)dr$. Hằng số $U(s_0)$ phụ thuộc $\rho$; thay dữ liệu để tìm nó rồi mới đổi về tọa độ gốc.
\subsection*{Lời giải áp dụng cho bài}
\textbf{Bước 1 — Đặc trưng và ODE.} Với $xu_x-yu_y+yu=y^2$, bất biến vẫn là $r=xy$. Dọc đặc trưng $U'+YU=Y^2$, $Y'=-Y$.

\textbf{Bước 2 — Tích phân.} Nhân $e^{-Y}$:
\[(e^{-Y}U)'=Y^2e^{-Y}.\]
Vì $dY=-Yds$, tích phân vế phải thành $-\int Ye^{-Y}dY=(Y+1)e^{-Y}$. Suy
\[u=y+1+e^yF(xy).\]
Nghiệm riêng $y+1$ giúp thấy công thức vẫn hợp lệ ở $y=0$, dù bước đổi tham số theo $Y$ cần xét riêng nơi $Y=0$.

\textbf{Bước 3 — Kiểm tra.} $x(y+1)_x-y(y+1)_y+y(y+1)=-y+y^2+y=y^2$; phần $e^yF(xy)$ thỏa thuần nhất theo bài 4. Do đó tổng đúng PDE.

\textbf{Chi tiết tích phân từng phần.} Khi $Y\ne0$, $ds=-dY/Y$, nên $\int Y^2e^{-Y}ds=-\int Ye^{-Y}dY$. Chọn $a=Y,db=e^{-Y}dY$, được $\int Ye^{-Y}dY=-Ye^{-Y}-e^{-Y}=-(Y+1)e^{-Y}$. Vậy $e^{-Y}U=(Y+1)e^{-Y}+F(XY)$, nhân lại $e^Y$ cho công thức đã nêu. Với $u_h=e^yF(xy)$, $u_{h,x}=ye^yF'$, $u_{h,y}=e^yF+xe^yF'$ và $xu_{h,x}-yu_{h,y}+yu_h=0$ bằng triệt tiêu từng hạng.

\clearpage\noindent\texttt{\detokenize{BT_1_3_2_3_06}}\par
\section*{Lời giải bổ sung — Bài tập 6 · Mục 1.3.2.3}

\subsection*{Kiến thức và trình tự giải}
\textbf{Mục tiêu.} Với $A(x,y)u_x+B(x,y)u_y+C(x,y)u=R(x,y)$, tìm các đường mà trên đó PDE trở thành một ODE. Dùng chữ $R$ cho vế phải để không nhầm với hàm tùy ý $F$.

\textbf{Bước 1 — Đọc các hệ số và lập hệ đặc trưng.} Đặt $U(s)=u(X(s),Y(s))$. Quy tắc dây chuyền cho $U'=u_xX'+u_yY'$. Chọn $X'=A$, $Y'=B$, khi đó

\[
X'=A(X,Y),\qquad Y'=B(X,Y),\qquad U'+C(X,Y)U=R(X,Y).
\]

Nếu $A\ne0$, có thể dùng $x$ làm tham số: $dy/dx=B/A$, $du/dx+(C/A)u=R/A$. Nếu $A=0$, chọn biến khác hoặc giữ hệ theo $s$; không chia cho hệ số bằng 0.

\textbf{Bước 2 — Tìm bất biến để phân biệt các đặc trưng.} Tích phân hai phương trình đầu để tìm $\rho(X,Y)$ không đổi. Kiểm tra $A\rho_x+B\rho_y=0$. Hệ số hằng cho $\rho=Bx-Ay$. Ví dụ $X'=X,Y'=-Y$ cho $(XY)'=0$; $X'=-Y,Y'=X$ cho $(X^2+Y^2)'=0$. Khi đặc trưng là vòng kín, cần kiểm tra nghiệm sau một vòng trở về cùng giá trị.

\textbf{Bước 3 — Giải ODE của ẩn trên từng đường.} Đặt $M(s)=\exp(\int C(X(s),Y(s))ds)$. Nhân phương trình của $U$ với $M$:

\[
(MU)'=MR,\qquad
U(s)=M(s)^{-1}\left[K+\int M(s)R(X(s),Y(s))ds\right].
\]

$K$ hằng trên một đường, nhưng thay đổi giữa các đường, nên viết $K=F(\rho)$. Nếu tìm được nghiệm riêng $u_p$, có thể giải phương trình thuần nhất cho $u-u_p$ để giảm tính toán. Nếu đổi biến $(\rho,z)$, phải tính Jacobian và dùng quy tắc dây chuyền để thấy PDE chỉ còn đạo hàm theo $z$.

\textbf{Bước 4 — Thế dữ liệu và tìm hàm tùy ý.} Với đường dữ liệu $(x_0(r),y_0(r))$, $u=h(r)$, thế vào nghiệm tổng quát để thu được $F(\rho_0(r))=H(r)$, trong đó $\rho_0=\rho(x_0,y_0)$. Nếu $\rho_0$ khả nghịch, đặt $r=\rho_0^{-1}(z)$ rồi viết $F(z)=H(\rho_0^{-1}(z))$. Nếu hai giá trị $r_1,r_2$ có cùng $\rho_0$, phải có $H(r_1)=H(r_2)$. Nếu đường dữ liệu chính là đặc trưng, dữ liệu phải thỏa ODE trên đường đó và thường không xác định hết $F$.

\textbf{Bước 5 — Xác định miền tồn tại và duy nhất.} Điều kiện không đặc trưng là $x_0'B-y_0'A\ne0$ trên đường dữ liệu. Nó cho khả nghịch cục bộ dưới giả thiết trơn thích hợp; không tự cho duy nhất trên cả mặt phẳng. Chỉ những đặc trưng gặp dữ liệu mới được xác định. Tại gốc, trục hoặc đường hệ số suy biến, kiểm tra riêng và ghép các biểu thức đủ trơn.

\textbf{Bước 6 — Kiểm tra kết quả.} Tính $u_x,u_y$ bằng quy tắc dây chuyền, thay vào $Au_x+Bu_y+Cu$ và kiểm tra bằng $R$. Sau đó thế lại điều kiện đầu/biên. Nêu rõ nếu chỉ có nghiệm cục bộ, nghiệm từng miền hoặc không có nghiệm trong lớp trơn yêu cầu.

\textbf{Chi tiết bước tích phân ODE.} Sau khi tìm bất biến $\rho$, viết $U'+P(s,\rho)U=Q(s,\rho)$. Chọn $M=\exp(\int_{s_0}^{s}P(r,\rho)dr)$; tính $M'=PM$, nên $MU'+M'U=MQ$. Tích phân từ $s_0$: $M(s)U(s)-U(s_0)=\int_{s_0}^sM(r)Q(r)dr$. Hằng số $U(s_0)$ phụ thuộc $\rho$; thay dữ liệu để tìm nó rồi mới đổi về tọa độ gốc.
\subsection*{Lời giải áp dụng cho bài}
\textbf{Bước 1 — Dùng hai nghiệm tổng quát đã kiểm tra.} Bài 4: $u=e^yF(xy)$. Bài 5: $u=y+1+e^yF(xy)$. Xét dữ liệu cho từng phương trình.

\textbf{Ý (a) — $u(x,1/x)=x$, $x>0$.} Trên đường này $xy=1$. Với bài 4 phải có $F(1)=xe^{-1/x}$; đạo hàm vế phải bằng $e^{-1/x}(1+1/x)>0$, nên không thể hằng. Với bài 5 phải có $F(1)=e^{-1/x}(x-1-1/x)$; đạo hàm bằng $e^{-1/x}(1+1/x-1/x^3)$, không đồng nhất bằng 0. Cả hai đều vô nghiệm trên toàn đường dữ liệu.

\textbf{Ý (b) — $u(1,y)=y^2$.} Với bài 4, $F(y)=y^2e^{-y}$, nên
\[u=x^2y^2e^{y-xy}.\]
Với bài 5, $F(y)=e^{-y}(y^2-y-1)$, nên
\[u=y+1+e^{y-xy}(x^2y^2-xy-1).\]
Thế $x=1$ cho $y^2$ ở cả hai. Trong miền $x>0$, mọi đặc trưng gặp $x=1$ đúng một lần nên duy nhất ở miền đó; dữ liệu không tự xác định các đặc trưng trong $x<0$.

\textbf{Ý (c) — $u(x,0)=x^4$.} Với cả hai PDE, đặt $y=0$ trong phương trình gốc cho $xu_x(x,0)=0$. Dữ liệu có $u_x(x,0)=4x^3$, suy $4x^4=0$, mâu thuẫn trên mọi khoảng có $x\ne0$. Vì vậy vô nghiệm. Dùng phương trình gốc ở đường suy biến tránh chia cho $y$.

\textbf{Chi tiết hai công thức dùng ở bước 1.} Đặc trưng $X'=X,Y'=-Y$ cho $XY=\rho$. Với phương trình thuần nhất, $(e^{-Y}U)'=0$ nên $u=e^yF(\rho)$. Với phương trình có nguồn $y^2$, $(e^{-Y}U)'=Y^2e^{-Y}$; đổi $ds=-dY/Y$ và tích phân được $(Y+1)e^{-Y}+F(\rho)$, tức $u=y+1+e^yF(\rho)$. Ở (b), đường $x=1$ cho $\rho=y=s$, nên thay $s=xy$ vào $F(s)=s^2e^{-s}$ hoặc $F(s)=(s^2-s-1)e^{-s}$; đây là phép thế tạo ra hai nghiệm đã viết.

\clearpage\noindent\texttt{\detokenize{BT_1_3_2_3_07}}\par
\section*{Lời giải bổ sung — Bài tập 7 · Mục 1.3.2.3}

\subsection*{Kiến thức và trình tự giải}
\textbf{Mục tiêu.} Với $A(x,y)u_x+B(x,y)u_y+C(x,y)u=R(x,y)$, tìm các đường mà trên đó PDE trở thành một ODE. Dùng chữ $R$ cho vế phải để không nhầm với hàm tùy ý $F$.

\textbf{Bước 1 — Đọc các hệ số và lập hệ đặc trưng.} Đặt $U(s)=u(X(s),Y(s))$. Quy tắc dây chuyền cho $U'=u_xX'+u_yY'$. Chọn $X'=A$, $Y'=B$, khi đó

\[
X'=A(X,Y),\qquad Y'=B(X,Y),\qquad U'+C(X,Y)U=R(X,Y).
\]

Nếu $A\ne0$, có thể dùng $x$ làm tham số: $dy/dx=B/A$, $du/dx+(C/A)u=R/A$. Nếu $A=0$, chọn biến khác hoặc giữ hệ theo $s$; không chia cho hệ số bằng 0.

\textbf{Bước 2 — Tìm bất biến để phân biệt các đặc trưng.} Tích phân hai phương trình đầu để tìm $\rho(X,Y)$ không đổi. Kiểm tra $A\rho_x+B\rho_y=0$. Hệ số hằng cho $\rho=Bx-Ay$. Ví dụ $X'=X,Y'=-Y$ cho $(XY)'=0$; $X'=-Y,Y'=X$ cho $(X^2+Y^2)'=0$. Khi đặc trưng là vòng kín, cần kiểm tra nghiệm sau một vòng trở về cùng giá trị.

\textbf{Bước 3 — Giải ODE của ẩn trên từng đường.} Đặt $M(s)=\exp(\int C(X(s),Y(s))ds)$. Nhân phương trình của $U$ với $M$:

\[
(MU)'=MR,\qquad
U(s)=M(s)^{-1}\left[K+\int M(s)R(X(s),Y(s))ds\right].
\]

$K$ hằng trên một đường, nhưng thay đổi giữa các đường, nên viết $K=F(\rho)$. Nếu tìm được nghiệm riêng $u_p$, có thể giải phương trình thuần nhất cho $u-u_p$ để giảm tính toán. Nếu đổi biến $(\rho,z)$, phải tính Jacobian và dùng quy tắc dây chuyền để thấy PDE chỉ còn đạo hàm theo $z$.

\textbf{Bước 4 — Thế dữ liệu và tìm hàm tùy ý.} Với đường dữ liệu $(x_0(r),y_0(r))$, $u=h(r)$, thế vào nghiệm tổng quát để thu được $F(\rho_0(r))=H(r)$, trong đó $\rho_0=\rho(x_0,y_0)$. Nếu $\rho_0$ khả nghịch, đặt $r=\rho_0^{-1}(z)$ rồi viết $F(z)=H(\rho_0^{-1}(z))$. Nếu hai giá trị $r_1,r_2$ có cùng $\rho_0$, phải có $H(r_1)=H(r_2)$. Nếu đường dữ liệu chính là đặc trưng, dữ liệu phải thỏa ODE trên đường đó và thường không xác định hết $F$.

\textbf{Bước 5 — Xác định miền tồn tại và duy nhất.} Điều kiện không đặc trưng là $x_0'B-y_0'A\ne0$ trên đường dữ liệu. Nó cho khả nghịch cục bộ dưới giả thiết trơn thích hợp; không tự cho duy nhất trên cả mặt phẳng. Chỉ những đặc trưng gặp dữ liệu mới được xác định. Tại gốc, trục hoặc đường hệ số suy biến, kiểm tra riêng và ghép các biểu thức đủ trơn.

\textbf{Bước 6 — Kiểm tra kết quả.} Tính $u_x,u_y$ bằng quy tắc dây chuyền, thay vào $Au_x+Bu_y+Cu$ và kiểm tra bằng $R$. Sau đó thế lại điều kiện đầu/biên. Nêu rõ nếu chỉ có nghiệm cục bộ, nghiệm từng miền hoặc không có nghiệm trong lớp trơn yêu cầu.

\textbf{Chi tiết bước tích phân ODE.} Sau khi tìm bất biến $\rho$, viết $U'+P(s,\rho)U=Q(s,\rho)$. Chọn $M=\exp(\int_{s_0}^{s}P(r,\rho)dr)$; tính $M'=PM$, nên $MU'+M'U=MQ$. Tích phân từ $s_0$: $M(s)U(s)-U(s_0)=\int_{s_0}^sM(r)Q(r)dr$. Hằng số $U(s_0)$ phụ thuộc $\rho$; thay dữ liệu để tìm nó rồi mới đổi về tọa độ gốc.
\subsection*{Lời giải áp dụng cho bài}
\textbf{Bước 1 — Bất biến.} Với $-2yu_x+xu_y=0$, $X'=-2Y,Y'=X$. Tính $(X^2+2Y^2)'=2X(-2Y)+4YX=0$.

\textbf{Bước 2 — Giá trị trên ellipse.} $U'=0$, nên nghiệm hằng trên mỗi ellipse $x^2+2y^2=r>0$. Do vòng ellipse liên thông, $u=F(x^2+2y^2)$ ngoài gốc.

\textbf{Bước 3 — Kiểm tra và độ trơn.} $u_x=2xF'$, $u_y=4yF'$, nên $-2yu_x+xu_y=0$. Để mô tả đúng lớp $C^1$ ở gốc, cũng có thể viết $u=H(\sqrt{x^2+2y^2})$ với $H\in C^1$, $H'(0)=0$. Yêu cầu $F'$ hữu hạn tại 0 là điều kiện đủ, nhưng có thể loại một số nghiệm $C^1$ hợp lệ.

\textbf{Chi tiết giải hệ trên ellipse.} Từ $X'=-2Y,Y'=X$, suy $X''=-2X$. Với $\omega=\sqrt2$,
\[X=x_0\cos\omega s-\omega y_0\sin\omega s,\quad
Y=y_0\cos\omega s+\frac{x_0}{\omega}\sin\omega s.\]
Thế vào $X^2+2Y^2$ cho $x_0^2+2y_0^2$. Chu kỳ $2\pi/\sqrt2$ đi hết ellipse; $U'=0$ cho hằng số trên toàn ellipse, thay vì chỉ trên một cung.

\clearpage\noindent\texttt{\detokenize{BT_1_3_2_3_08}}\par
\section*{Lời giải bổ sung — Bài tập 8 · Mục 1.3.2.3}

\subsection*{Kiến thức và trình tự giải}
\textbf{Mục tiêu.} Với $A(x,y)u_x+B(x,y)u_y+C(x,y)u=R(x,y)$, tìm các đường mà trên đó PDE trở thành một ODE. Dùng chữ $R$ cho vế phải để không nhầm với hàm tùy ý $F$.

\textbf{Bước 1 — Đọc các hệ số và lập hệ đặc trưng.} Đặt $U(s)=u(X(s),Y(s))$. Quy tắc dây chuyền cho $U'=u_xX'+u_yY'$. Chọn $X'=A$, $Y'=B$, khi đó

\[
X'=A(X,Y),\qquad Y'=B(X,Y),\qquad U'+C(X,Y)U=R(X,Y).
\]

Nếu $A\ne0$, có thể dùng $x$ làm tham số: $dy/dx=B/A$, $du/dx+(C/A)u=R/A$. Nếu $A=0$, chọn biến khác hoặc giữ hệ theo $s$; không chia cho hệ số bằng 0.

\textbf{Bước 2 — Tìm bất biến để phân biệt các đặc trưng.} Tích phân hai phương trình đầu để tìm $\rho(X,Y)$ không đổi. Kiểm tra $A\rho_x+B\rho_y=0$. Hệ số hằng cho $\rho=Bx-Ay$. Ví dụ $X'=X,Y'=-Y$ cho $(XY)'=0$; $X'=-Y,Y'=X$ cho $(X^2+Y^2)'=0$. Khi đặc trưng là vòng kín, cần kiểm tra nghiệm sau một vòng trở về cùng giá trị.

\textbf{Bước 3 — Giải ODE của ẩn trên từng đường.} Đặt $M(s)=\exp(\int C(X(s),Y(s))ds)$. Nhân phương trình của $U$ với $M$:

\[
(MU)'=MR,\qquad
U(s)=M(s)^{-1}\left[K+\int M(s)R(X(s),Y(s))ds\right].
\]

$K$ hằng trên một đường, nhưng thay đổi giữa các đường, nên viết $K=F(\rho)$. Nếu tìm được nghiệm riêng $u_p$, có thể giải phương trình thuần nhất cho $u-u_p$ để giảm tính toán. Nếu đổi biến $(\rho,z)$, phải tính Jacobian và dùng quy tắc dây chuyền để thấy PDE chỉ còn đạo hàm theo $z$.

\textbf{Bước 4 — Thế dữ liệu và tìm hàm tùy ý.} Với đường dữ liệu $(x_0(r),y_0(r))$, $u=h(r)$, thế vào nghiệm tổng quát để thu được $F(\rho_0(r))=H(r)$, trong đó $\rho_0=\rho(x_0,y_0)$. Nếu $\rho_0$ khả nghịch, đặt $r=\rho_0^{-1}(z)$ rồi viết $F(z)=H(\rho_0^{-1}(z))$. Nếu hai giá trị $r_1,r_2$ có cùng $\rho_0$, phải có $H(r_1)=H(r_2)$. Nếu đường dữ liệu chính là đặc trưng, dữ liệu phải thỏa ODE trên đường đó và thường không xác định hết $F$.

\textbf{Bước 5 — Xác định miền tồn tại và duy nhất.} Điều kiện không đặc trưng là $x_0'B-y_0'A\ne0$ trên đường dữ liệu. Nó cho khả nghịch cục bộ dưới giả thiết trơn thích hợp; không tự cho duy nhất trên cả mặt phẳng. Chỉ những đặc trưng gặp dữ liệu mới được xác định. Tại gốc, trục hoặc đường hệ số suy biến, kiểm tra riêng và ghép các biểu thức đủ trơn.

\textbf{Bước 6 — Kiểm tra kết quả.} Tính $u_x,u_y$ bằng quy tắc dây chuyền, thay vào $Au_x+Bu_y+Cu$ và kiểm tra bằng $R$. Sau đó thế lại điều kiện đầu/biên. Nêu rõ nếu chỉ có nghiệm cục bộ, nghiệm từng miền hoặc không có nghiệm trong lớp trơn yêu cầu.

\textbf{Chi tiết bước tích phân ODE.} Sau khi tìm bất biến $\rho$, viết $U'+P(s,\rho)U=Q(s,\rho)$. Chọn $M=\exp(\int_{s_0}^{s}P(r,\rho)dr)$; tính $M'=PM$, nên $MU'+M'U=MQ$. Tích phân từ $s_0$: $M(s)U(s)-U(s_0)=\int_{s_0}^sM(r)Q(r)dr$. Hằng số $U(s_0)$ phụ thuộc $\rho$; thay dữ liệu để tìm nó rồi mới đổi về tọa độ gốc.
\subsection*{Lời giải áp dụng cho bài}
\textbf{Ý (a).} $yu_x-xu_y=u$ cho $X'=Y,Y'=-X,U'=U$. Đường đi là vòng tròn chu kỳ $2\pi$, nhưng $U(s+2\pi)=e^{2\pi}U(s)$. Một nghiệm đơn trị phải trở lại cùng giá trị, buộc $U=0$ trên mọi vòng. Dữ liệu $u(x,0)=x$ khác 0, nên vô nghiệm toàn cục. Cục bộ trong một góc không phủ vòng kín có thể có nghiệm, nhưng không kéo dài thành nghiệm đơn trị quanh vòng.

\textbf{Ý (b).} $u_t+xtu_x=0$. Dùng $t$: $X'=tX$, $U'=0$, $X(0)=\xi$, $U(0)=\sin\xi$. Tích phân $X=\xi e^{t^2/2}$, nên $\xi=xe^{-t^2/2}$ và
\[u=\sin(xe^{-t^2/2}).\]
Đặt $z=xe^{-t^2/2}$, $u_t=-tz\cos z$, $xtu_x=tz\cos z$, tổng bằng 0; tại $t=0$ đúng $\sin x$.

\textbf{Ý (c) — Giữ hệ số $-2xt$ trong đề.} $xu_t-2xtu_x=2tu$, $u(x,0)=x^2$. Trên $x\ne0$, chia cho $x$: $u_t-2tu_x=(2t/x)u$. Đặc trưng $X'= -2t$ cho $X=\xi-t^2$, nên $\xi=x+t^2$. Trên đường,
\[U'=\frac{2t}{\xi-t^2}U,\quad \frac{d}{dt}[(\xi-t^2)U]=0.\]
Ở $t=0$, tích này bằng $\xi^3$, do đó
\[u=\frac{(x+t^2)^3}{x}.\]
Đối với miền từ dữ liệu $x>0$, công thức hợp lệ khi $x>0$; từ dữ liệu $x<0$, cần $\xi=x+t^2<0$. Đặt $z=x+t^2$, $u_t=6tz^2/x$, $u_x=3z^2/x-z^3/x^2$, nên $xu_t-2xtu_x=2tz^3/x=2tu$. Trên $x=0,t>0$ phương trình gốc buộc $u=0$, nhưng công thức phía $x>0$ phân kỳ khi $x\downarrow0$. Vì thế không có nghiệm $C^1$ trên toàn mặt phẳng thời gian dương thỏa toàn bộ dữ liệu.

\textbf{Ý (d).} $yu_x+xu_y=0$ có $r=x^2-y^2$ bất biến và $U'=0$. Trên $x=0$, $r=-y^2$; dữ liệu cho $F(r)=e^r$ với $r\le0$. Vì thế trên miền các đặc trưng từ trục $y$ phủ, có $u=e^{x^2-y^2}$. Công thức cũng là một nghiệm toàn mặt phẳng, nhưng miền $r>0$ chưa được dữ liệu xác định. Chẳng hạn thêm $\psi(r)$ với $\psi(r)=0$ khi $r\le0$, $\psi(r)=e^{-1/r}$ khi $r>0$, vẫn cho nghiệm trơn cùng dữ liệu và khác trên $r>0$.

\textbf{Ý (e).} Với $V=y\partial_x+x\partial_y$, $Vr=0$ và $V(y^2/2)=xy$. Do đó $u=y^2/2+H(r)$. Trên $y=0$, $H(r)=e^{-r}$ với $r\ge0$. Trên $x=0$, $r=-y^2$, nên $H(r)=e^r+r/2$ với $r\le0$. Hai giá trị bằng nhau tại 0, nhưng $H'_+(0)=-1$, $H'_-(0)=3/2$. Trên đường $x=y>0$, $\nabla r\ne0$, nên gradient của $u$ không ghép liên tục. Kết luận: công thức ghép là nghiệm liên tục từng miền, nhưng không có nghiệm $C^1$ trên toàn góc phần tư thỏa hai vết đã cho.

\textbf{Chi tiết hệ hyperbola của (d),(e).} Với $X'=Y,Y'=X$, đặt $S=X+Y,T=X-Y$. Khi đó $S'=S,T'=-T$, nên $S=S_0e^s,T=T_0e^{-s}$ và $ST=X^2-Y^2=\rho$. Ở (d), $U'=0$ nên $U$ hằng trên từng thành phần hyperbola; trên trục $x=0$ có $\rho=-y^2$ và dữ liệu $e^{-y^2}=e^\rho$, quyết định các thành phần từ trục ấy. Ở (e), $U'=XY$ và $(Y^2/2)'=YY'=XY$, nên $(U-Y^2/2)'=0$. Bởi vậy $u=y^2/2+H(\rho)$ trước khi dùng hai vết. Đạo hàm hai biểu thức của $H$ tại $0$: $(e^{-r})'|_0=-1$, $(e^r+r/2)'|_0=1+1/2=3/2$, giải thích sự không ghép $C^1$.

\clearpage\noindent\texttt{\detokenize{VD_1_4_2_01}}\par
\section*{Lời giải bổ sung — Ví dụ 1 · Mục 1.4.2}

\subsection*{Kiến thức và trình tự giải}
\textbf{Mục tiêu.} Với $u_t+a(x,t,u)u_x=b(x,t,u)$, đặc trưng phụ thuộc cả giá trị nghiệm nên phải giải đồng thời đường đi và ẩn.

\textbf{Bước 1 — Gắn nhãn điểm xuất phát.} Gọi $\xi$ là tọa độ ở $t=0$, đặt $X(0;\xi)=\xi$, $U(0;\xi)=f(\xi)$. Dọc đường, $U(t;\xi)=u(X(t;\xi),t)$.

\textbf{Bước 2 — Lập và giải hệ.} Quy tắc dây chuyền cho

\[
X'=a(X,t,U),\qquad U'=b(X,t,U).
\]

Nếu ODE của $U$ giải được trước, thế kết quả vào $X'$ rồi tích phân từ 0 đến $t$, giữ đúng điều kiện $X(0)=\xi$. Chẳng hạn $U'=\lambda U$ cho $U=e^{\lambda t}f(\xi)$. Nếu chia cho $U$ hoặc $U+1$, xét riêng các nghiệm hằng bị loại.

\textbf{Bước 3 — Khử nhãn $\xi$.} Giải $x=X(t;\xi)$ để tìm $\xi=\xi(x,t)$ rồi thay vào $u=U(t;\xi)$. Khi không khử được bằng công thức đơn giản, ghi nghiệm tham số và phương trình xác định $\xi$.

\textbf{Bước 4 — Tìm miền còn nghiệm cổ điển.} Tính $X_\xi$. Khi $X_\xi\ne0$, có thể khử nhãn cục bộ. Vì $u_x=U_\xi/X_\xi$, mẫu số tiến về 0 có thể gây đạo hàm vô hạn. Với $u_t+a(u)u_x=0$, ta có $X=\xi+a(f(\xi))t$, $X_\xi=1+a'(f(\xi))f'(\xi)t$. Tìm thời điểm dương đầu tiên mẫu này bằng 0.

\textbf{Bước 5 — Nếu dữ liệu có bước nhảy.} Viết $F'(u)=a(u)$ để có dạng bảo toàn $u_t+(F(u))_x=0$. Đặc trưng hội tụ dẫn đến sốc, tách xa dẫn đến quạt giãn; dùng Rankine–Hugoniot và điều kiện entropy như phần kiến thức về sốc. Giá trị ở chỗ đặc trưng giao nhau không được chọn tùy ý.

\textbf{Bước 6 — Thử lại.} Kiểm tra phương trình trong vùng trơn, dữ liệu ở $t=0$, miền khả nghịch và mọi đường ghép. Với dữ liệu chỉ trên một đoạn, giữ điều kiện nhãn $\xi$ nằm trong đoạn đó.
\subsection*{Lời giải áp dụng cho bài}
\textbf{Ý (a) — $u_t+tu_x=u^2$.} Bước 1: $X'=t$ cho $\rho=x-t^2/2$. Bước 2: $U'=U^2$. Nếu $U\ne0$, $(1/U)'=-1$, nên $u=1/[F(\rho)-t]$, ở nơi mẫu khác 0. Phải thêm nghiệm $u=0$ bị mất khi chia. Một dạng từ vết $f(\rho)$ tại $t=0$ giữ cả điểm $f=0$ là $u=f(\rho)/[1-tf(\rho)]$. Dọc đặc trưng công thức có đạo hàm bằng $U^2$, nên kiểm tra được PDE.

\textbf{Ý (b) — $yu_x+xu_y=u^2+u$.} $V=y\partial_x+x\partial_y$ có bất biến $r=x^2-y^2$. Trên $x+y\ne0$, đặt $s=\log|x+y|$; $Vs=1$. ODE là $U'=U(U+1)$. Nếu $U\ne0,-1$, tích phân $dU/[U(U+1)]=ds$ cho $\log|U/(U+1)|=s+\text{hằng}$. Suy
\[u=\frac{A(r)e^s}{1-A(r)e^s}.\]
$u=0$ nằm trong họ khi $A=0$; thêm $u=-1$. Công thức trên mỗi biểu đồ có mẫu khác 0; trên $x+y=0$ dùng biểu đồ khác nếu trường không triệt tiêu, tại gốc PDE cho $u=0$ hoặc $-1$. Ghép đủ trơn trước khi kết luận toàn cục.

\textbf{Ý (c) — $u_x+xyu_y=1/u$.} $dy/dx=xy$ cho $r=ye^{-x^2/2}$. Dọc đường, $U'=1/U$, nên $(U^2)'=2$, $U^2=2x+A(r)$. Suy $u=\pm\sqrt{2x+A(ye^{-x^2/2})}$. Chọn miền trong căn dương và dấu liên tục. Kiểm tra bằng cách đạo hàm $u^2$: $(\partial_x+xy\partial_y)u^2=2$, suy PDE vì $u\ne0$. Không thêm $u=0$ vì phương trình gốc chứa $1/u$.

\textbf{Chi tiết phân thức ở (b).} $\frac1{U(U+1)}=\frac1U-\frac1{U+1}$, nên $\int dU/[U(U+1)]=\log|U|-\log|U+1|$. Sau khi tích phân theo $s$, viết $U/(U+1)=A(r)e^s$ với dấu được gộp vào $A$. Nhân chéo: $U=Ae^sU+Ae^s$, tức $(1-Ae^s)U=Ae^s$; chia khi mẫu khác $0$ được công thức. Hai giá trị cân bằng $0,-1$ phải được xét trước phép chia.

\clearpage\noindent\texttt{\detokenize{VD_1_4_2_02}}\par
\section*{Lời giải bổ sung — Ví dụ 2 · Mục 1.4.2}

\subsection*{Kiến thức và trình tự giải}
\textbf{Mục tiêu.} Với $u_t+a(x,t,u)u_x=b(x,t,u)$, đặc trưng phụ thuộc cả giá trị nghiệm nên phải giải đồng thời đường đi và ẩn.

\textbf{Bước 1 — Gắn nhãn điểm xuất phát.} Gọi $\xi$ là tọa độ ở $t=0$, đặt $X(0;\xi)=\xi$, $U(0;\xi)=f(\xi)$. Dọc đường, $U(t;\xi)=u(X(t;\xi),t)$.

\textbf{Bước 2 — Lập và giải hệ.} Quy tắc dây chuyền cho

\[
X'=a(X,t,U),\qquad U'=b(X,t,U).
\]

Nếu ODE của $U$ giải được trước, thế kết quả vào $X'$ rồi tích phân từ 0 đến $t$, giữ đúng điều kiện $X(0)=\xi$. Chẳng hạn $U'=\lambda U$ cho $U=e^{\lambda t}f(\xi)$. Nếu chia cho $U$ hoặc $U+1$, xét riêng các nghiệm hằng bị loại.

\textbf{Bước 3 — Khử nhãn $\xi$.} Giải $x=X(t;\xi)$ để tìm $\xi=\xi(x,t)$ rồi thay vào $u=U(t;\xi)$. Khi không khử được bằng công thức đơn giản, ghi nghiệm tham số và phương trình xác định $\xi$.

\textbf{Bước 4 — Tìm miền còn nghiệm cổ điển.} Tính $X_\xi$. Khi $X_\xi\ne0$, có thể khử nhãn cục bộ. Vì $u_x=U_\xi/X_\xi$, mẫu số tiến về 0 có thể gây đạo hàm vô hạn. Với $u_t+a(u)u_x=0$, ta có $X=\xi+a(f(\xi))t$, $X_\xi=1+a'(f(\xi))f'(\xi)t$. Tìm thời điểm dương đầu tiên mẫu này bằng 0.

\textbf{Bước 5 — Nếu dữ liệu có bước nhảy.} Viết $F'(u)=a(u)$ để có dạng bảo toàn $u_t+(F(u))_x=0$. Đặc trưng hội tụ dẫn đến sốc, tách xa dẫn đến quạt giãn; dùng Rankine–Hugoniot và điều kiện entropy như phần kiến thức về sốc. Giá trị ở chỗ đặc trưng giao nhau không được chọn tùy ý.

\textbf{Bước 6 — Thử lại.} Kiểm tra phương trình trong vùng trơn, dữ liệu ở $t=0$, miền khả nghịch và mọi đường ghép. Với dữ liệu chỉ trên một đoạn, giữ điều kiện nhãn $\xi$ nằm trong đoạn đó.
\subsection*{Lời giải áp dụng cho bài}
\textbf{Ý (a).} $u_t+uu_x=0$, $u(x,0)=3x$. $U'=0$ cho $U=3\xi$; $X'=U$ cho $X=\xi(1+3t)$. Khử $\xi=x/(1+3t)$:
\[u=3x/(1+3t).\]
$X_\xi=1+3t>0$ khi $t\ge0$, nên nghiệm cổ điển trên toàn miền thời gian dương. $u_t=-9x/(1+3t)^2$, $uu_x=9x/(1+3t)^2$.

\textbf{Ý (b).} $u_t-u^2u_x=3u$, $u(x,0)=f(x)$. Giải $U'=3U$ trước: $U=e^{3t}f(\xi)$. Sau đó
\[X'= -e^{6t}f(\xi)^2,\quad X=\xi-\frac{e^{6t}-1}{6}f(\xi)^2.\]
Nghiệm là $u=e^{3t}f(\xi)$, với $\xi$ xác định từ $x=X$. Tính $X_\xi=1-(e^{6t}-1)f(\xi)f'(\xi)/3$; chỉ khử được cục bộ ở nơi khác 0. Giữ nhãn và phương trình khử là cách giải phù hợp khi $f$ tùy ý.

\textbf{Ý (c).} $u_t-t^2uu_x=-u$. $U=e^{-t}f(\xi)$; $X'= -t^2e^{-t}f(\xi)$. Dùng tích phân từng phần hai lần:
\[\int_0^t s^2e^{-s}ds=2-(t^2+2t+2)e^{-t}=:J(t).\]
Suy $X=\xi-J(t)f(\xi)$, $u=e^{-t}f(\xi)$. Miền khử cần $1-J(t)f'(\xi)\ne0$; tại $t=0$, $J=0$, cho đúng dữ liệu.

\textbf{Ý (d).} $u_t+t^2uu_x=5$, $u(x,0)=x$. $U'=5$ cho $U=\xi+5t$. Tích phân $X'=t^2(\xi+5t)$:
\[X=\xi(1+t^3/3)+5t^4/4,\quad \xi=\frac{x-5t^4/4}{1+t^3/3}.\]
Do đó $u=(x-5t^4/4)/(1+t^3/3)+5t$. Với $t\ge0$, $X_\xi=1+t^3/3>0$. Trên mỗi nhánh trên, $dU/dt=u_t+X'u_x$ bằng đúng nguồn của PDE và dữ liệu tại 0 được giữ, nên công thức tham số thỏa bài.

\textbf{Chi tiết tích phân ở (c).} Đặt $J(t)=\int_0^t s^2e^{-s}ds$. Lần thứ nhất:
\[J=-t^2e^{-t}+2\int_0^t se^{-s}ds.\]
Lần thứ hai:
\[\int_0^t se^{-s}ds=-te^{-t}+\int_0^t e^{-s}ds=-te^{-t}+1-e^{-t}.\]
Thế lại được $J=2-(t^2+2t+2)e^{-t}$, rồi $X-\xi=-f(\xi)J(t)$. Ở (b), tích phân tương tự trực tiếp: $X-\xi=-f(\xi)^2\int_0^te^{6s}ds=-f(\xi)^2(e^{6t}-1)/6$.

\clearpage\noindent\texttt{\detokenize{VD_1_4_2_03}}\par
\section*{Lời giải bổ sung — Ví dụ 3 · Mục 1.4.2}

\subsection*{Kiến thức và trình tự giải}
\textbf{Mục tiêu.} Với dữ liệu gián đoạn hoặc khi đặc trưng giao nhau, tìm nghiệm yếu thỏa entropy của luật bảo toàn $u_t+(F(u))_x=0$.

\textbf{Bước 1 — Xác định thông lượng và đặc trưng.} Nếu PDE là $u_t+a(u)u_x=0$, tích phân $F'=a$. Trong vùng trơn, $u=u_0(\xi)$ và $x=\xi+F'(u_0(\xi))t$. Tính riêng tốc độ cho mỗi đoạn dữ liệu. Burgers dùng $F=u^2/2$; $u_t+u^2u_x=0$ dùng $F=u^3/3$.

\textbf{Bước 2 — Nhận diện sóng tại mỗi bước nhảy.} Gọi $u_L,u_R$ là giá trị trái, phải. So sánh $F'(u_L)$ và $F'(u_R)$. Với $F$ lồi trên khoảng trạng thái, $u_L>u_R$ tạo sốc, $u_L<u_R$ tạo quạt giãn; với $F$ lõm, thứ tự trạng thái đảo lại. Nếu độ lồi đổi dấu, phải xét thêm cấu trúc entropy, không dùng quy tắc lồi trên toàn khoảng.

\textbf{Bước 3 — Tính đường sốc từ bảo toàn.} Với $x=s(t)$,

\[
s'(t)=\frac{F(u_L)-F(u_R)}{u_L-u_R}.
\]

Nếu hai trạng thái hằng, tốc độ $\sigma$ hằng, nên $s=x_0+\sigma t$. Kiểm tra entropy bằng $F'(u_R)<\sigma<F'(u_L)$ đối với sốc nén có hai trạng thái phân biệt trong các bài này. Giá trị ngay trên đường sốc không ảnh hưởng nghiệm yếu.

\textbf{Bước 4 — Tính quạt giãn.} Từ $(x_0,0)$, đặt $z=(x-x_0)/t$, $u=V(z)$. Thế PDE cho $(F'(V)-z)V'=0$. Trong nhánh biến thiên, $F'(V)=z$, nên giải phương trình này tìm $V$. Mép quạt là $x=x_0+F'(u_L)t$ và $x=x_0+F'(u_R)t$, theo thứ tự tốc độ tăng; ghép liên tục với hai trạng thái ngoài. Chọn căn đúng dấu theo khoảng giá trị của dữ liệu.

\textbf{Bước 5 — Tìm mọi thời điểm tương tác.} Giải phương trình giao của hai sốc, sốc với mép quạt, hoặc quạt với quạt. Trước thời điểm gặp, ghép các sóng riêng. Sau khi gặp, tính lại hai trạng thái sát đường sốc. Nếu một phía nằm trong quạt, trạng thái đó phụ thuộc $s(t)$ và $t$, nên Rankine–Hugoniot trở thành ODE mới. Dùng $s(t_*)=x_*$ để tìm hằng số, rồi tiếp tục tìm lần tương tác sau.

\textbf{Bước 6 — Viết nghiệm theo từng khoảng thời gian.} Với mỗi khoảng, ghi rõ các miền $x$, thứ tự đường phân cách và công thức của $u$. Trên nửa trục, kiểm tra chiều truyền từ biên vào miền và thế điều kiện biên. Với dữ liệu trơn từng đoạn, tính $\partial x/\partial\xi$ để tìm thời điểm đầu tiên mất nghiệm cổ điển.

\textbf{Bước 7 — Kiểm tra.} Trong mỗi vùng trơn, thay vào PDE. Trên mỗi sốc, kiểm tra Rankine–Hugoniot và entropy; ở mép quạt, kiểm tra liên tục. Đối chiếu dữ liệu khi $t\downarrow0$, sự liên tục của vị trí sóng qua thời điểm gặp và khối lượng nếu dữ liệu có giá hữu hạn. Nhân PDE với một hàm của $u$ chỉ giữ tương đương cho nghiệm trơn; không tự áp luật đã nhân để tính cùng đường sốc.

\textbf{Chi tiết tích phân đường sốc nằm trong quạt.} Nếu $s'=s/(2t)$ và $s(t_*)=x_*>0$, chia cho $s$ rồi tích phân từ $t_*$ đến $t$: $\log(s/x_*)=\frac12\log(t/t_*)$, nên $s=x_*\sqrt{t/t_*}$. Nếu quạt có tâm $x_0$, thay $s$ bằng $s-x_0$ trước khi tách biến. Khi trạng thái sát sốc thay đổi tại lần gặp tiếp theo, lập ODE mới và lấy vị trí tại lần gặp làm dữ liệu; không giữ hằng số của giai đoạn trước.
\subsection*{Lời giải áp dụng cho bài}
\textbf{Bước 1 — Thông lượng và đặc trưng.} Burgers có $F=u^2/2$. Dữ liệu $u_L=1,u_R=2$ ở $x=0$ cho $X=\xi+t$ từ trái, $X=\xi+2t$ từ phải. Hai họ tách xa, để trống $t<x<2t$.

\textbf{Bước 2 — Lấp bằng quạt giãn.} Đặt $z=x/t$, $u=V(z)$. PDE cho $(V-z)V'=0$, nên nhánh biến thiên có $V=z$. Nó ghép liên tục với giá trị 1 tại $x=t$ và 2 tại $x=2t$.

\textbf{Bước 3 — Nghiệm và kiểm tra.}
\[u=\begin{cases}1,&x\le t,\\x/t,&t<x<2t,\\2,&x\ge2t,\end{cases}\quad t>0.\]
Trong quạt $u_t=-x/t^2$, $uu_x=x/t^2$, tổng bằng 0. Các mép liên tục và dữ liệu khi $t\downarrow0$ đúng ở $x\ne0$. Đây là nghiệm entropy; dữ liệu gián đoạn không cho nghiệm cổ điển qua $t=0$.

\clearpage\noindent\texttt{\detokenize{BT_1_4_3_01}}\par
\section*{Lời giải bổ sung — Bài tập 1 · Mục 1.4.3}

\subsection*{Kiến thức và trình tự giải}
\textbf{Mục tiêu.} Với dữ liệu gián đoạn hoặc khi đặc trưng giao nhau, tìm nghiệm yếu thỏa entropy của luật bảo toàn $u_t+(F(u))_x=0$.

\textbf{Bước 1 — Xác định thông lượng và đặc trưng.} Nếu PDE là $u_t+a(u)u_x=0$, tích phân $F'=a$. Trong vùng trơn, $u=u_0(\xi)$ và $x=\xi+F'(u_0(\xi))t$. Tính riêng tốc độ cho mỗi đoạn dữ liệu. Burgers dùng $F=u^2/2$; $u_t+u^2u_x=0$ dùng $F=u^3/3$.

\textbf{Bước 2 — Nhận diện sóng tại mỗi bước nhảy.} Gọi $u_L,u_R$ là giá trị trái, phải. So sánh $F'(u_L)$ và $F'(u_R)$. Với $F$ lồi trên khoảng trạng thái, $u_L>u_R$ tạo sốc, $u_L<u_R$ tạo quạt giãn; với $F$ lõm, thứ tự trạng thái đảo lại. Nếu độ lồi đổi dấu, phải xét thêm cấu trúc entropy, không dùng quy tắc lồi trên toàn khoảng.

\textbf{Bước 3 — Tính đường sốc từ bảo toàn.} Với $x=s(t)$,

\[
s'(t)=\frac{F(u_L)-F(u_R)}{u_L-u_R}.
\]

Nếu hai trạng thái hằng, tốc độ $\sigma$ hằng, nên $s=x_0+\sigma t$. Kiểm tra entropy bằng $F'(u_R)<\sigma<F'(u_L)$ đối với sốc nén có hai trạng thái phân biệt trong các bài này. Giá trị ngay trên đường sốc không ảnh hưởng nghiệm yếu.

\textbf{Bước 4 — Tính quạt giãn.} Từ $(x_0,0)$, đặt $z=(x-x_0)/t$, $u=V(z)$. Thế PDE cho $(F'(V)-z)V'=0$. Trong nhánh biến thiên, $F'(V)=z$, nên giải phương trình này tìm $V$. Mép quạt là $x=x_0+F'(u_L)t$ và $x=x_0+F'(u_R)t$, theo thứ tự tốc độ tăng; ghép liên tục với hai trạng thái ngoài. Chọn căn đúng dấu theo khoảng giá trị của dữ liệu.

\textbf{Bước 5 — Tìm mọi thời điểm tương tác.} Giải phương trình giao của hai sốc, sốc với mép quạt, hoặc quạt với quạt. Trước thời điểm gặp, ghép các sóng riêng. Sau khi gặp, tính lại hai trạng thái sát đường sốc. Nếu một phía nằm trong quạt, trạng thái đó phụ thuộc $s(t)$ và $t$, nên Rankine–Hugoniot trở thành ODE mới. Dùng $s(t_*)=x_*$ để tìm hằng số, rồi tiếp tục tìm lần tương tác sau.

\textbf{Bước 6 — Viết nghiệm theo từng khoảng thời gian.} Với mỗi khoảng, ghi rõ các miền $x$, thứ tự đường phân cách và công thức của $u$. Trên nửa trục, kiểm tra chiều truyền từ biên vào miền và thế điều kiện biên. Với dữ liệu trơn từng đoạn, tính $\partial x/\partial\xi$ để tìm thời điểm đầu tiên mất nghiệm cổ điển.

\textbf{Bước 7 — Kiểm tra.} Trong mỗi vùng trơn, thay vào PDE. Trên mỗi sốc, kiểm tra Rankine–Hugoniot và entropy; ở mép quạt, kiểm tra liên tục. Đối chiếu dữ liệu khi $t\downarrow0$, sự liên tục của vị trí sóng qua thời điểm gặp và khối lượng nếu dữ liệu có giá hữu hạn. Nhân PDE với một hàm của $u$ chỉ giữ tương đương cho nghiệm trơn; không tự áp luật đã nhân để tính cùng đường sốc.

\textbf{Chi tiết tích phân đường sốc nằm trong quạt.} Nếu $s'=s/(2t)$ và $s(t_*)=x_*>0$, chia cho $s$ rồi tích phân từ $t_*$ đến $t$: $\log(s/x_*)=\frac12\log(t/t_*)$, nên $s=x_*\sqrt{t/t_*}$. Nếu quạt có tâm $x_0$, thay $s$ bằng $s-x_0$ trước khi tách biến. Khi trạng thái sát sốc thay đổi tại lần gặp tiếp theo, lập ODE mới và lấy vị trí tại lần gặp làm dữ liệu; không giữ hằng số của giai đoạn trước.
\subsection*{Lời giải áp dụng cho bài}
\textbf{Bước 1 — Chọn phương pháp từng dữ liệu.} Trong vùng trơn dùng $U'=0$, $X'=F'(U)$; dữ liệu nhảy dùng sốc/quạt. Các thông lượng lần lượt là $u^2/2$ cho (a),(b), $u^3/3$ cho (c), $2u^2$ cho (d).

\textbf{Ý (a).} Dữ liệu nhảy từ 2 lên 12 tại $x=1$. Với Burgers, tốc độ trái 2, phải 12, nên hai họ tách xa và tạo quạt. Trong quạt, $F'(u)=u=(x-1)/t$. Do đó
\[u=\begin{cases}2,&x\le1+2t,\\(x-1)/t,&1+2t<x<1+12t,\\12,&x\ge1+12t.\end{cases}\]
Các mép cho giá trị 2,12; trong quạt $u_t=-(x-1)/t^2$, $uu_x=(x-1)/t^2$, nên PDE đúng.

\textbf{Ý (b).} Dữ liệu $u_0(\xi)=-\xi$ trơn. $U=-\xi$, $X=\xi(1-t)$. Khử $\xi=x/(1-t)$ cho $u=-x/(1-t)$ khi $0\le t<1$. $X_\xi=1-t$ triệt tiêu ở $t=1$, mọi đặc trưng hội tụ ở $x=0$. Vì vậy không kéo dài thành nghiệm cổ điển qua thời điểm đó.

\textbf{Ý (c).} Trạng thái 4 trái, 3 phải ở $x=0$; $F'=u^2$ có tốc độ 16 và 9, hội tụ. Rankine–Hugoniot:
\[\sigma=\frac{4^3/3-3^3/3}{4-3}=37/3.\]
Suy $s=37t/3$, $u=4$ khi $x<s$, $u=3$ khi $x>s$. Entropy: $9<37/3<16$.

\textbf{Ý (d).} Trạng thái 3 trái, 2 phải tại $x=1$, $F=2u^2$, tốc độ $F'=4u$ là 12 và 8. $\sigma=[2\cdot9-2\cdot4]/(3-2)=10$, nên $s=1+10t$, $u=3$ trái sốc, 2 phải sốc; $8<10<12$. Các trạng thái hằng thỏa PDE trong mỗi miền, đường sốc thỏa bảo toàn, và giới hạn $t\downarrow0$ khớp dữ liệu.

\clearpage\noindent\texttt{\detokenize{BT_1_4_3_02}}\par
\section*{Lời giải bổ sung — Bài tập 2 · Mục 1.4.3}

\subsection*{Kiến thức và trình tự giải}
\textbf{Mục tiêu.} Với dữ liệu gián đoạn hoặc khi đặc trưng giao nhau, tìm nghiệm yếu thỏa entropy của luật bảo toàn $u_t+(F(u))_x=0$.

\textbf{Bước 1 — Xác định thông lượng và đặc trưng.} Nếu PDE là $u_t+a(u)u_x=0$, tích phân $F'=a$. Trong vùng trơn, $u=u_0(\xi)$ và $x=\xi+F'(u_0(\xi))t$. Tính riêng tốc độ cho mỗi đoạn dữ liệu. Burgers dùng $F=u^2/2$; $u_t+u^2u_x=0$ dùng $F=u^3/3$.

\textbf{Bước 2 — Nhận diện sóng tại mỗi bước nhảy.} Gọi $u_L,u_R$ là giá trị trái, phải. So sánh $F'(u_L)$ và $F'(u_R)$. Với $F$ lồi trên khoảng trạng thái, $u_L>u_R$ tạo sốc, $u_L<u_R$ tạo quạt giãn; với $F$ lõm, thứ tự trạng thái đảo lại. Nếu độ lồi đổi dấu, phải xét thêm cấu trúc entropy, không dùng quy tắc lồi trên toàn khoảng.

\textbf{Bước 3 — Tính đường sốc từ bảo toàn.} Với $x=s(t)$,

\[
s'(t)=\frac{F(u_L)-F(u_R)}{u_L-u_R}.
\]

Nếu hai trạng thái hằng, tốc độ $\sigma$ hằng, nên $s=x_0+\sigma t$. Kiểm tra entropy bằng $F'(u_R)<\sigma<F'(u_L)$ đối với sốc nén có hai trạng thái phân biệt trong các bài này. Giá trị ngay trên đường sốc không ảnh hưởng nghiệm yếu.

\textbf{Bước 4 — Tính quạt giãn.} Từ $(x_0,0)$, đặt $z=(x-x_0)/t$, $u=V(z)$. Thế PDE cho $(F'(V)-z)V'=0$. Trong nhánh biến thiên, $F'(V)=z$, nên giải phương trình này tìm $V$. Mép quạt là $x=x_0+F'(u_L)t$ và $x=x_0+F'(u_R)t$, theo thứ tự tốc độ tăng; ghép liên tục với hai trạng thái ngoài. Chọn căn đúng dấu theo khoảng giá trị của dữ liệu.

\textbf{Bước 5 — Tìm mọi thời điểm tương tác.} Giải phương trình giao của hai sốc, sốc với mép quạt, hoặc quạt với quạt. Trước thời điểm gặp, ghép các sóng riêng. Sau khi gặp, tính lại hai trạng thái sát đường sốc. Nếu một phía nằm trong quạt, trạng thái đó phụ thuộc $s(t)$ và $t$, nên Rankine–Hugoniot trở thành ODE mới. Dùng $s(t_*)=x_*$ để tìm hằng số, rồi tiếp tục tìm lần tương tác sau.

\textbf{Bước 6 — Viết nghiệm theo từng khoảng thời gian.} Với mỗi khoảng, ghi rõ các miền $x$, thứ tự đường phân cách và công thức của $u$. Trên nửa trục, kiểm tra chiều truyền từ biên vào miền và thế điều kiện biên. Với dữ liệu trơn từng đoạn, tính $\partial x/\partial\xi$ để tìm thời điểm đầu tiên mất nghiệm cổ điển.

\textbf{Bước 7 — Kiểm tra.} Trong mỗi vùng trơn, thay vào PDE. Trên mỗi sốc, kiểm tra Rankine–Hugoniot và entropy; ở mép quạt, kiểm tra liên tục. Đối chiếu dữ liệu khi $t\downarrow0$, sự liên tục của vị trí sóng qua thời điểm gặp và khối lượng nếu dữ liệu có giá hữu hạn. Nhân PDE với một hàm của $u$ chỉ giữ tương đương cho nghiệm trơn; không tự áp luật đã nhân để tính cùng đường sốc.

\textbf{Chi tiết tích phân đường sốc nằm trong quạt.} Nếu $s'=s/(2t)$ và $s(t_*)=x_*>0$, chia cho $s$ rồi tích phân từ $t_*$ đến $t$: $\log(s/x_*)=\frac12\log(t/t_*)$, nên $s=x_*\sqrt{t/t_*}$. Nếu quạt có tâm $x_0$, thay $s$ bằng $s-x_0$ trước khi tách biến. Khi trạng thái sát sốc thay đổi tại lần gặp tiếp theo, lập ODE mới và lấy vị trí tại lần gặp làm dữ liệu; không giữ hằng số của giai đoạn trước.
\subsection*{Lời giải áp dụng cho bài}
\textbf{Bước 1 — Các sốc trước tương tác.} Burgers có dữ liệu lần lượt 2,1,0 với bước nhảy tại 0 và 1. Cả hai bước đều giảm nên tạo sốc. Tốc độ là trung bình hai trạng thái: $\sigma_1=(2+1)/2=3/2$, $\sigma_2=(1+0)/2=1/2$. Vì vậy $s_1=3t/2$, $s_2=1+t/2$.

\textbf{Bước 2 — Thời điểm gặp.} Giải $3t/2=1+t/2$ cho $t_*=1$, $x_*=3/2$. Khi $0<t<1$,
\[u=\begin{cases}2,&x<s_1,\\1,&s_1<x<s_2,\\0,&x>s_2.\end{cases}\]

\textbf{Bước 3 — Tính lại sốc sau gặp.} Trạng thái giữa biến mất; sốc nối 2 với 0 có tốc độ 1. Từ $s(1)=3/2$, $s=3/2+(t-1)=t+1/2$. Với $t\ge1$, $u=2$ khi $x<s$, $u=0$ khi $x>s$.

\textbf{Bước 4 — Kiểm tra.} Các tốc độ đều nằm giữa tốc độ đặc trưng hai phía: $1<3/2<2$, $0<1/2<1$, $0<1<2$. Vị trí sốc ghép liên tục ở $t=1$ và giới hạn ban đầu đúng các đoạn dữ liệu.

\clearpage\noindent\texttt{\detokenize{VD_1_5_2_01}}\par
\section*{Lời giải bổ sung — Ví dụ 1 · Mục 1.5.2}

\subsection*{Kiến thức và trình tự giải}
\textbf{Mục tiêu.} Giải PDE cấp 1 phi tuyến $F(x,y,u,p,q)=0$, với $p=u_x,q=u_y$. Chọn $s$ để gắn nhãn đường dữ liệu, $\tau$ là tham số chạy dọc đặc trưng; hai tham số này có vai trò khác nhau.

\textbf{Bước 1 — Viết đúng $F$ và tham số hóa dữ liệu.} Chuyển PDE về một vế, giữ mọi hệ số chứa $x,y,u$. Viết $X(s,0)=x_0(s)$, $Y(s,0)=y_0(s)$, $U(s,0)=h(s)$. Đạo hàm điều kiện theo $s$ cho

\[
p_0x_0'(s)+q_0y_0'(s)=h'(s),\qquad
F(x_0,y_0,h,p_0,q_0)=0.
\]

Giải hai phương trình này tìm tất cả nhánh thực $(p_0,q_0)$. Trên $y=0$, $p_0=h'$; trên $x=0$, $q_0=h'$; trên $y=x$, $p_0+q_0=h'$. Điều kiện dưới căn phải không âm. Các nhánh khác nhau có thể cho các nghiệm khác nhau cùng dữ liệu.

\textbf{Bước 2 — Kiểm tra không đặc trưng.} Tính $J_0=x_0'F_q-y_0'F_p$ ở dữ liệu của từng nhánh. Khi $J_0\ne0$ và các hàm đủ trơn, phép chuyển từ $(s,\tau)$ sang $(x,y)$ khả nghịch cục bộ. Khi $J_0=0$, cần xét riêng; không kết luận ngay vô nghiệm hoặc duy nhất.

\textbf{Bước 3 — Tính các đạo hàm riêng của $F$ rồi lập hệ Charpit.} Khi tính $F_x,F_y,F_u,F_p,F_q$, tạm coi năm đối số độc lập. Hệ là

\[
\begin{aligned}
X_\tau&=F_p,&Y_\tau&=F_q,&U_\tau&=pF_p+qF_q,\\
p_\tau&=-F_x-pF_u,&q_\tau&=-F_y-qF_u.
\end{aligned}
\]

Trong bài cụ thể phải thay các đạo hàm của $F$ vào hệ. Ví dụ $F=pq-1$ cho $X_\tau=q,Y_\tau=p,U_\tau=2,p_\tau=q_\tau=0$.

\textbf{Bước 4 — Tích phân hệ với điều kiện ban đầu.} Tìm phương trình giải được trước: $p'=0$, $q'=0$, hoặc tổ hợp như $(X-p)'=0$. Thế chúng vào phương trình còn lại. Nếu hệ tuyến tính $Z'=MZ$, dùng $Z=e^{\tau M}Z_0$; để tính cụ thể, tìm các tổ hợp tách hệ hoặc quan hệ $M^2, M^3$ giúp rút gọn chuỗi mũ. Ghi các phép tích phân và hằng số được dữ liệu xác định.

\textbf{Bước 5 — Khử tham số hoặc trình bày nghiệm tham số đã giải.} Từ $X(s,\tau),Y(s,\tau)$ tìm $s,\tau$ theo $x,y$ rồi thay vào $U$. Nếu giữ dạng tham số, phải có công thức $X,Y,U$ đã tích phân và miền $\det\partial(X,Y)/\partial(s,\tau)\ne0$; chỉ ghi hệ ODE chưa phải lời giải hoàn tất.

\textbf{Bước 6 — Dùng phép đổi ẩn khi thuận lợi.} Gần dữ liệu $u>0$, $p^2+q^2=u$ đổi $v=2\sqrt u$ cho $|\nabla v|^2=1$; $p^2+q^2=u^2$ đổi $v=\log u$ cũng cho eikonal. Phải tính $v_x,v_y$ để chứng minh phép đổi, đổi cả dữ liệu và giữ điều kiện dấu. Với dữ liệu hằng trên vòng tròn, đạo hàm tiếp tuyến bằng 0 nên gradient pháp tuyến có hai lựa chọn dấu.

\textbf{Bước 7 — Kiểm tra từng nhánh và miền.} Tính $p,q$ của công thức cuối, thế vào PDE và dữ liệu. Phân biệt một nghiệm đã tìm được, các nhánh không đặc trưng cục bộ và toàn bộ nghiệm toàn cục. Giữ điều kiện của nhãn dữ liệu nếu đề chỉ cho $0\le s\le1$; kiểm tra gốc, điểm căn bằng 0 và đường làm mẫu số bằng 0.

\textbf{Chi tiết tích phân và khử sau khi lập hệ.} Trong lúc tích phân theo $\tau$, nhãn $s$ và các gradient ban đầu $p_0(s),q_0(s)$ được giữ cố định. Với $p'=cp$, viết $p=p_0e^{c\tau}$; nếu $X'=dp$ thì $X=x_0+dp_0(e^{c\tau}-1)/c$ khi $c\ne0$, hoặc $X=x_0+dp_0\tau$ khi $c=0$. Với $Y'+cY=R(\tau)$, nhân $e^{c\tau}$ cho $(e^{c\tau}Y)'=e^{c\tau}R$, rồi dùng $Y(0)$ để xác định hằng số. Khử từ hai phương trình $x=X(s,\tau),y=Y(s,\tau)$; ghi rõ tổ hợp nào loại $\tau$ hoặc $s$. Nếu chỉ cho nghiệm tham số, vẫn phải có các tích phân đã tính và điều kiện Jacobian để khử cục bộ.
\subsection*{Lời giải áp dụng cho bài}
\textbf{Ý (a), bước 1 — Tương thích dữ liệu.} Đặt $F=p^2+q^2-1$, $X(s,0)=s,Y(s,0)=0,U(s,0)=f(s)$. Đạo hàm vết cho $p_0=f'(s)$, còn PDE cho $q_0^2=1-f'(s)^2$. Nếu $|f'(0)|>1$, không có nghiệm thực gần gốc.

\textbf{Bước 2 — Không đặc trưng và hệ Charpit.} Nếu $|f'(0)|<1$ và $f\in C^2$ gần gốc, thu hẹp khoảng để $|f'(s)|<1$. Có hai nhánh $q_0=\pm\sqrt{1-f'(s)^2}$ và $J_0=F_q=2q_0\ne0$. Hệ là
\[X'=2p,\quad Y'=2q,\quad U'=2,\quad p'=q'=0.\]
Do đó
\[X=s+2\tau f'(s),\quad Y=2\tau q_0(s),\quad U=f(s)+2\tau.\]
Trên từng nhánh, khử $s,\tau$ gần dữ liệu để được một nghiệm cục bộ. Nếu $f(s)=s/2$, khử trực tiếp cho $u=x/2\pm\sqrt3\,y/2$.

\textbf{Bước 3 — Trường hợp $|f'(0)|=1$.} $J_0=0$, nên không dùng định lý không đặc trưng để kết luận duy nhất. Dữ liệu $f(x)=x+x^2\operatorname{sgn}x$ có $f'(x)=1+2|x|>1$ khi $x\ne0$ gần gốc, nên vô nghiệm dù $f'(0)=1$. Dữ liệu $f(x)=x$ lại có vô số nghiệm $u_C=C-\sqrt{(x-C)^2+y^2}$ gần gốc, với $C>1$: trên $y=0,x<C$ có $u_C=x$, và $|\nabla u_C|=1$. Các ví dụ cho thấy phải xét cả dữ liệu lân cận, không chỉ giá trị đạo hàm tại gốc.

\textbf{Ý (b), bước 1 — Dữ liệu trên vòng tròn.} $x_0=\cos s,y_0=\sin s,h=0$. Đạo hàm tiếp tuyến cho $-p_0\sin s+q_0\cos s=0$; PDE cho $p_0^2+q_0^2=1$. Vì vậy gradient có hai lựa chọn $(p_0,q_0)=\varepsilon(\cos s,\sin s)$, $\varepsilon=\pm1$.

\textbf{Bước 2 — Tích phân và khử.} $X=(1+2\varepsilon\tau)\cos s$, $Y=(1+2\varepsilon\tau)\sin s$, $U=2\tau$. Gần vòng tròn giữ $1+2\varepsilon\tau>0$, nên $r=1+2\varepsilon\tau$ và
\[u=\varepsilon(r-1),\qquad r=\sqrt{x^2+y^2}.\]

\textbf{Bước 3 — Kiểm tra và sửa kết luận trong sách.} Hai hàm $r-1$ và $1-r$ đều có $|\nabla u|=1$ và vết 0 tại $r=1$. $J_0=-2\varepsilon\ne0$ cho duy nhất cục bộ \emph{trong từng nhánh gradient}, nhưng dữ liệu giá trị chưa chọn dấu nên có hai nhánh. Không kéo dài hai hàm xuyên tâm thành nghiệm $C^1$ qua gốc.

\textbf{Chi tiết khử tham số và chọn gradient.} Ở (a) với $f=s/2$, $q_0=\varepsilon\sqrt3/2$, $X=s+\tau$, $Y=\varepsilon\sqrt3\tau$, $U=s/2+2\tau$. Do đó $\tau=\varepsilon y/\sqrt3$, $s=x-\tau$, $U=x/2+3\tau/2=x/2+\varepsilon\sqrt3 y/2$. Ở (b), đặt tiếp tuyến $e_\theta=(-\sin s,\cos s)$ và pháp tuyến $e_r=(\cos s,\sin s)$. Viết $(p_0,q_0)=Ae_r+Be_\theta$; dữ liệu cho $B=0$, PDE cho $A^2+B^2=1$, nên $A=\varepsilon$. Từ $r=1+2\varepsilon\tau$ có $\tau=\varepsilon(r-1)/2$, $U=2\tau=\varepsilon(r-1)$.

\clearpage\noindent\texttt{\detokenize{VD_1_5_2_02}}\par
\section*{Lời giải bổ sung — Ví dụ 2 · Mục 1.5.2}

\subsection*{Kiến thức và trình tự giải}
\textbf{Mục tiêu.} Giải PDE cấp 1 phi tuyến $F(x,y,u,p,q)=0$, với $p=u_x,q=u_y$. Chọn $s$ để gắn nhãn đường dữ liệu, $\tau$ là tham số chạy dọc đặc trưng; hai tham số này có vai trò khác nhau.

\textbf{Bước 1 — Viết đúng $F$ và tham số hóa dữ liệu.} Chuyển PDE về một vế, giữ mọi hệ số chứa $x,y,u$. Viết $X(s,0)=x_0(s)$, $Y(s,0)=y_0(s)$, $U(s,0)=h(s)$. Đạo hàm điều kiện theo $s$ cho

\[
p_0x_0'(s)+q_0y_0'(s)=h'(s),\qquad
F(x_0,y_0,h,p_0,q_0)=0.
\]

Giải hai phương trình này tìm tất cả nhánh thực $(p_0,q_0)$. Trên $y=0$, $p_0=h'$; trên $x=0$, $q_0=h'$; trên $y=x$, $p_0+q_0=h'$. Điều kiện dưới căn phải không âm. Các nhánh khác nhau có thể cho các nghiệm khác nhau cùng dữ liệu.

\textbf{Bước 2 — Kiểm tra không đặc trưng.} Tính $J_0=x_0'F_q-y_0'F_p$ ở dữ liệu của từng nhánh. Khi $J_0\ne0$ và các hàm đủ trơn, phép chuyển từ $(s,\tau)$ sang $(x,y)$ khả nghịch cục bộ. Khi $J_0=0$, cần xét riêng; không kết luận ngay vô nghiệm hoặc duy nhất.

\textbf{Bước 3 — Tính các đạo hàm riêng của $F$ rồi lập hệ Charpit.} Khi tính $F_x,F_y,F_u,F_p,F_q$, tạm coi năm đối số độc lập. Hệ là

\[
\begin{aligned}
X_\tau&=F_p,&Y_\tau&=F_q,&U_\tau&=pF_p+qF_q,\\
p_\tau&=-F_x-pF_u,&q_\tau&=-F_y-qF_u.
\end{aligned}
\]

Trong bài cụ thể phải thay các đạo hàm của $F$ vào hệ. Ví dụ $F=pq-1$ cho $X_\tau=q,Y_\tau=p,U_\tau=2,p_\tau=q_\tau=0$.

\textbf{Bước 4 — Tích phân hệ với điều kiện ban đầu.} Tìm phương trình giải được trước: $p'=0$, $q'=0$, hoặc tổ hợp như $(X-p)'=0$. Thế chúng vào phương trình còn lại. Nếu hệ tuyến tính $Z'=MZ$, dùng $Z=e^{\tau M}Z_0$; để tính cụ thể, tìm các tổ hợp tách hệ hoặc quan hệ $M^2, M^3$ giúp rút gọn chuỗi mũ. Ghi các phép tích phân và hằng số được dữ liệu xác định.

\textbf{Bước 5 — Khử tham số hoặc trình bày nghiệm tham số đã giải.} Từ $X(s,\tau),Y(s,\tau)$ tìm $s,\tau$ theo $x,y$ rồi thay vào $U$. Nếu giữ dạng tham số, phải có công thức $X,Y,U$ đã tích phân và miền $\det\partial(X,Y)/\partial(s,\tau)\ne0$; chỉ ghi hệ ODE chưa phải lời giải hoàn tất.

\textbf{Bước 6 — Dùng phép đổi ẩn khi thuận lợi.} Gần dữ liệu $u>0$, $p^2+q^2=u$ đổi $v=2\sqrt u$ cho $|\nabla v|^2=1$; $p^2+q^2=u^2$ đổi $v=\log u$ cũng cho eikonal. Phải tính $v_x,v_y$ để chứng minh phép đổi, đổi cả dữ liệu và giữ điều kiện dấu. Với dữ liệu hằng trên vòng tròn, đạo hàm tiếp tuyến bằng 0 nên gradient pháp tuyến có hai lựa chọn dấu.

\textbf{Bước 7 — Kiểm tra từng nhánh và miền.} Tính $p,q$ của công thức cuối, thế vào PDE và dữ liệu. Phân biệt một nghiệm đã tìm được, các nhánh không đặc trưng cục bộ và toàn bộ nghiệm toàn cục. Giữ điều kiện của nhãn dữ liệu nếu đề chỉ cho $0\le s\le1$; kiểm tra gốc, điểm căn bằng 0 và đường làm mẫu số bằng 0.

\textbf{Chi tiết tích phân và khử sau khi lập hệ.} Trong lúc tích phân theo $\tau$, nhãn $s$ và các gradient ban đầu $p_0(s),q_0(s)$ được giữ cố định. Với $p'=cp$, viết $p=p_0e^{c\tau}$; nếu $X'=dp$ thì $X=x_0+dp_0(e^{c\tau}-1)/c$ khi $c\ne0$, hoặc $X=x_0+dp_0\tau$ khi $c=0$. Với $Y'+cY=R(\tau)$, nhân $e^{c\tau}$ cho $(e^{c\tau}Y)'=e^{c\tau}R$, rồi dùng $Y(0)$ để xác định hằng số. Khử từ hai phương trình $x=X(s,\tau),y=Y(s,\tau)$; ghi rõ tổ hợp nào loại $\tau$ hoặc $s$. Nếu chỉ cho nghiệm tham số, vẫn phải có các tích phân đã tính và điều kiện Jacobian để khử cục bộ.
\subsection*{Lời giải áp dụng cho bài}
\textbf{Bước 1 — Tương thích.} $F=p^2+q^2+2(p-x)(q-y)-2u$, dữ liệu $(s,0,0)$ cho $p_0=0$ và $q_0(q_0-2s)=0$. Có hai nhánh $q_0=0,2s$; $J_0=2(q_0-s)$ khác 0 khi $s>0$.

\textbf{Bước 2 — Hệ và các bất biến.} Tính đạo hàm của $F$:
\[X'=2(p+q-Y),\quad Y'=2(p+q-X),\quad p'=X',\quad q'=Y'.\]
Do đó $X-p=s$, $Y-q=-q_0$. Thế $p=X-s,q=Y+q_0$ vào hệ đường đi, được $X'=2X+2(q_0-s)$, $Y'=2Y+2(q_0-s)$. Tích phân với dữ liệu:
\[X=q_0e^{2\tau}+s-q_0,\qquad Y=(q_0-s)(e^{2\tau}-1).\]
Từ $F=0$,
\[U=\frac{X^2+Y^2+s^2-2Xs+2Yq_0-2sq_0+q_0^2}{2}.\]

\textbf{Bước 3 — Khử hai nhánh.} Nếu $q_0=0$, $X=s$, $Y=s(1-e^{2\tau})$, suy $p=0,q=Y$ và $u=y^2/2$. Nếu $q_0=2s$, $X=2se^{2\tau}-s$, $Y=s(e^{2\tau}-1)$, nên $s=X-2Y$, $p=2Y,q=2X-3Y$, suy $u=2xy-3y^2/2$.

\textbf{Bước 4 — Kiểm tra và miền.} Với $u=y^2/2$, $p=0,q=y$ cho vế trái $y^2=2u$. Với $u=2xy-3y^2/2$, $p=2y,q=2x-3y$ cho vế trái $4xy-3y^2=2u$. Cả hai có $u(x,0)=0$. Từ đoạn dữ liệu $0\le s\le1$, nhánh thứ nhất cần $s=x$ và $e^{2\tau}=1-y/x>0$; nhánh thứ hai cần $s=x-2y$ và $e^{2\tau}=(x-y)/(x-2y)>0$. Các đa thức là nghiệm kéo dài toàn cục, nhưng dữ liệu chỉ bảo đảm duy nhất cục bộ trên mỗi nhánh ở $s>0$; gốc suy biến. Lời giải sách nói duy nhất khi đã chọn gradient ban đầu; nếu chưa chọn nhánh thì có hai nghiệm.

\textbf{Chi tiết giải hai ODE tuyến tính.} Đặt $c=q_0-s$, hằng theo $\tau$. Từ $X'-2X=2c$, nhân $e^{-2\tau}$ và tích phân từ $0$:
\[e^{-2\tau}X-s=c(1-e^{-2\tau}),\quad X=(s+c)e^{2\tau}-c.\]
Tương tự $e^{-2\tau}Y=c(1-e^{-2\tau})$, nên $Y=c(e^{2\tau}-1)$. Vì $s+c=q_0$, đây đúng các công thức ở bước 2. Nhánh $q_0=2s$ cho $X-2Y=s$, $p=X-s=2Y$, $q=Y+2s=2X-3Y$. Tích phân $u_x=2y$ cho $u=2xy+C(y)$; phương trình $u_y=2x-3y$ cho $C'=-3y$, rồi $C=-3y^2/2$ nhờ vết ban đầu.

\clearpage\noindent\texttt{\detokenize{BT_1_5_3_01}}\par
\section*{Lời giải bổ sung — Bài tập 1 · Mục 1.5.3}

\subsection*{Kiến thức và trình tự giải}
\textbf{Mục tiêu.} Giải PDE cấp 1 phi tuyến $F(x,y,u,p,q)=0$, với $p=u_x,q=u_y$. Chọn $s$ để gắn nhãn đường dữ liệu, $\tau$ là tham số chạy dọc đặc trưng; hai tham số này có vai trò khác nhau.

\textbf{Bước 1 — Viết đúng $F$ và tham số hóa dữ liệu.} Chuyển PDE về một vế, giữ mọi hệ số chứa $x,y,u$. Viết $X(s,0)=x_0(s)$, $Y(s,0)=y_0(s)$, $U(s,0)=h(s)$. Đạo hàm điều kiện theo $s$ cho

\[
p_0x_0'(s)+q_0y_0'(s)=h'(s),\qquad
F(x_0,y_0,h,p_0,q_0)=0.
\]

Giải hai phương trình này tìm tất cả nhánh thực $(p_0,q_0)$. Trên $y=0$, $p_0=h'$; trên $x=0$, $q_0=h'$; trên $y=x$, $p_0+q_0=h'$. Điều kiện dưới căn phải không âm. Các nhánh khác nhau có thể cho các nghiệm khác nhau cùng dữ liệu.

\textbf{Bước 2 — Kiểm tra không đặc trưng.} Tính $J_0=x_0'F_q-y_0'F_p$ ở dữ liệu của từng nhánh. Khi $J_0\ne0$ và các hàm đủ trơn, phép chuyển từ $(s,\tau)$ sang $(x,y)$ khả nghịch cục bộ. Khi $J_0=0$, cần xét riêng; không kết luận ngay vô nghiệm hoặc duy nhất.

\textbf{Bước 3 — Tính các đạo hàm riêng của $F$ rồi lập hệ Charpit.} Khi tính $F_x,F_y,F_u,F_p,F_q$, tạm coi năm đối số độc lập. Hệ là

\[
\begin{aligned}
X_\tau&=F_p,&Y_\tau&=F_q,&U_\tau&=pF_p+qF_q,\\
p_\tau&=-F_x-pF_u,&q_\tau&=-F_y-qF_u.
\end{aligned}
\]

Trong bài cụ thể phải thay các đạo hàm của $F$ vào hệ. Ví dụ $F=pq-1$ cho $X_\tau=q,Y_\tau=p,U_\tau=2,p_\tau=q_\tau=0$.

\textbf{Bước 4 — Tích phân hệ với điều kiện ban đầu.} Tìm phương trình giải được trước: $p'=0$, $q'=0$, hoặc tổ hợp như $(X-p)'=0$. Thế chúng vào phương trình còn lại. Nếu hệ tuyến tính $Z'=MZ$, dùng $Z=e^{\tau M}Z_0$; để tính cụ thể, tìm các tổ hợp tách hệ hoặc quan hệ $M^2, M^3$ giúp rút gọn chuỗi mũ. Ghi các phép tích phân và hằng số được dữ liệu xác định.

\textbf{Bước 5 — Khử tham số hoặc trình bày nghiệm tham số đã giải.} Từ $X(s,\tau),Y(s,\tau)$ tìm $s,\tau$ theo $x,y$ rồi thay vào $U$. Nếu giữ dạng tham số, phải có công thức $X,Y,U$ đã tích phân và miền $\det\partial(X,Y)/\partial(s,\tau)\ne0$; chỉ ghi hệ ODE chưa phải lời giải hoàn tất.

\textbf{Bước 6 — Dùng phép đổi ẩn khi thuận lợi.} Gần dữ liệu $u>0$, $p^2+q^2=u$ đổi $v=2\sqrt u$ cho $|\nabla v|^2=1$; $p^2+q^2=u^2$ đổi $v=\log u$ cũng cho eikonal. Phải tính $v_x,v_y$ để chứng minh phép đổi, đổi cả dữ liệu và giữ điều kiện dấu. Với dữ liệu hằng trên vòng tròn, đạo hàm tiếp tuyến bằng 0 nên gradient pháp tuyến có hai lựa chọn dấu.

\textbf{Bước 7 — Kiểm tra từng nhánh và miền.} Tính $p,q$ của công thức cuối, thế vào PDE và dữ liệu. Phân biệt một nghiệm đã tìm được, các nhánh không đặc trưng cục bộ và toàn bộ nghiệm toàn cục. Giữ điều kiện của nhãn dữ liệu nếu đề chỉ cho $0\le s\le1$; kiểm tra gốc, điểm căn bằng 0 và đường làm mẫu số bằng 0.

\textbf{Chi tiết tích phân và khử sau khi lập hệ.} Trong lúc tích phân theo $\tau$, nhãn $s$ và các gradient ban đầu $p_0(s),q_0(s)$ được giữ cố định. Với $p'=cp$, viết $p=p_0e^{c\tau}$; nếu $X'=dp$ thì $X=x_0+dp_0(e^{c\tau}-1)/c$ khi $c\ne0$, hoặc $X=x_0+dp_0\tau$ khi $c=0$. Với $Y'+cY=R(\tau)$, nhân $e^{c\tau}$ cho $(e^{c\tau}Y)'=e^{c\tau}R$, rồi dùng $Y(0)$ để xác định hằng số. Khử từ hai phương trình $x=X(s,\tau),y=Y(s,\tau)$; ghi rõ tổ hợp nào loại $\tau$ hoặc $s$. Nếu chỉ cho nghiệm tham số, vẫn phải có các tích phân đã tính và điều kiện Jacobian để khử cục bộ.
\subsection*{Lời giải áp dụng cho bài}
\textbf{Bước 1 — Tìm gradient trên dữ liệu.} $F=p^2+pq-5u$, dữ liệu $u(0,s)=4s^2$ cho $q_0=8s$. PDE cho $p_0^2+8sp_0-20s^2=0=(p_0-2s)(p_0+10s)$, nên $p_0=2s$ hoặc $-10s$.

\textbf{Bước 2 — Giải Charpit.} $F_p=2p+q,F_q=p,F_u=-5$, nên $p'=5p,q'=5q,U'=10U$. Đặt $E=e^{5\tau}$:
\[p=p_0E,\quad q=8sE,\quad U=4s^2E^2,\quad X=(2p_0+8s)(E-1)/5,\quad Y=s+p_0(E-1)/5.\]

\textbf{Bước 3 — Khử từng nhánh.} Nếu $p_0=2s$, $X/2+2Y=2sE$, nên $u=(x/2+2y)^2$. Nếu $p_0=-10s$, $5X/2-2Y=-2sE$, nên $u=(5x/2-2y)^2$.

\textbf{Bước 4 — Kiểm tra.} Với $u=L^2$, $L=ax+by$, có $p^2+pq=4a(a+b)L^2$. Hai bộ $(a,b)=(1/2,2),(5/2,-2)$ đều có $4a(a+b)=5$, nên PDE đúng; ở $x=0$ đều bằng $4y^2$. Định thức trên dữ liệu $x=0$ là $-(2p_0+q_0)=\mp12s$, nên mỗi nhánh không đặc trưng ở $s\ne0$; không suy duy nhất ở gốc.

\textbf{Chi tiết tích phân và khử.} $p'=5p$ cho $p=p_0e^{5\tau}$, $q'=5q$ cho $q=8se^{5\tau}$; $U'=p(2p+q)=2(p^2+pq)=10U$ cho $U=4s^2e^{10\tau}$. Vì $\int_0^\tau e^{5r}dr=(E-1)/5$, tích phân $X',Y'$ được các công thức ở bước 2. Với $p_0=2s$: $X=12s(E-1)/5$, $Y=s+2s(E-1)/5$, nên $X/2+2Y=2s+2s(E-1)=2sE$. Với $p_0=-10s$: $X=-12s(E-1)/5$, $Y=s-2s(E-1)$, nên $5X/2-2Y=-2sE$. Bình phương cả hai cho $U$.

\clearpage\noindent\texttt{\detokenize{BT_1_5_3_02}}\par
\section*{Lời giải bổ sung — Bài tập 2 · Mục 1.5.3}

\subsection*{Kiến thức và trình tự giải}
\textbf{Mục tiêu.} Giải PDE cấp 1 phi tuyến $F(x,y,u,p,q)=0$, với $p=u_x,q=u_y$. Chọn $s$ để gắn nhãn đường dữ liệu, $\tau$ là tham số chạy dọc đặc trưng; hai tham số này có vai trò khác nhau.

\textbf{Bước 1 — Viết đúng $F$ và tham số hóa dữ liệu.} Chuyển PDE về một vế, giữ mọi hệ số chứa $x,y,u$. Viết $X(s,0)=x_0(s)$, $Y(s,0)=y_0(s)$, $U(s,0)=h(s)$. Đạo hàm điều kiện theo $s$ cho

\[
p_0x_0'(s)+q_0y_0'(s)=h'(s),\qquad
F(x_0,y_0,h,p_0,q_0)=0.
\]

Giải hai phương trình này tìm tất cả nhánh thực $(p_0,q_0)$. Trên $y=0$, $p_0=h'$; trên $x=0$, $q_0=h'$; trên $y=x$, $p_0+q_0=h'$. Điều kiện dưới căn phải không âm. Các nhánh khác nhau có thể cho các nghiệm khác nhau cùng dữ liệu.

\textbf{Bước 2 — Kiểm tra không đặc trưng.} Tính $J_0=x_0'F_q-y_0'F_p$ ở dữ liệu của từng nhánh. Khi $J_0\ne0$ và các hàm đủ trơn, phép chuyển từ $(s,\tau)$ sang $(x,y)$ khả nghịch cục bộ. Khi $J_0=0$, cần xét riêng; không kết luận ngay vô nghiệm hoặc duy nhất.

\textbf{Bước 3 — Tính các đạo hàm riêng của $F$ rồi lập hệ Charpit.} Khi tính $F_x,F_y,F_u,F_p,F_q$, tạm coi năm đối số độc lập. Hệ là

\[
\begin{aligned}
X_\tau&=F_p,&Y_\tau&=F_q,&U_\tau&=pF_p+qF_q,\\
p_\tau&=-F_x-pF_u,&q_\tau&=-F_y-qF_u.
\end{aligned}
\]

Trong bài cụ thể phải thay các đạo hàm của $F$ vào hệ. Ví dụ $F=pq-1$ cho $X_\tau=q,Y_\tau=p,U_\tau=2,p_\tau=q_\tau=0$.

\textbf{Bước 4 — Tích phân hệ với điều kiện ban đầu.} Tìm phương trình giải được trước: $p'=0$, $q'=0$, hoặc tổ hợp như $(X-p)'=0$. Thế chúng vào phương trình còn lại. Nếu hệ tuyến tính $Z'=MZ$, dùng $Z=e^{\tau M}Z_0$; để tính cụ thể, tìm các tổ hợp tách hệ hoặc quan hệ $M^2, M^3$ giúp rút gọn chuỗi mũ. Ghi các phép tích phân và hằng số được dữ liệu xác định.

\textbf{Bước 5 — Khử tham số hoặc trình bày nghiệm tham số đã giải.} Từ $X(s,\tau),Y(s,\tau)$ tìm $s,\tau$ theo $x,y$ rồi thay vào $U$. Nếu giữ dạng tham số, phải có công thức $X,Y,U$ đã tích phân và miền $\det\partial(X,Y)/\partial(s,\tau)\ne0$; chỉ ghi hệ ODE chưa phải lời giải hoàn tất.

\textbf{Bước 6 — Dùng phép đổi ẩn khi thuận lợi.} Gần dữ liệu $u>0$, $p^2+q^2=u$ đổi $v=2\sqrt u$ cho $|\nabla v|^2=1$; $p^2+q^2=u^2$ đổi $v=\log u$ cũng cho eikonal. Phải tính $v_x,v_y$ để chứng minh phép đổi, đổi cả dữ liệu và giữ điều kiện dấu. Với dữ liệu hằng trên vòng tròn, đạo hàm tiếp tuyến bằng 0 nên gradient pháp tuyến có hai lựa chọn dấu.

\textbf{Bước 7 — Kiểm tra từng nhánh và miền.} Tính $p,q$ của công thức cuối, thế vào PDE và dữ liệu. Phân biệt một nghiệm đã tìm được, các nhánh không đặc trưng cục bộ và toàn bộ nghiệm toàn cục. Giữ điều kiện của nhãn dữ liệu nếu đề chỉ cho $0\le s\le1$; kiểm tra gốc, điểm căn bằng 0 và đường làm mẫu số bằng 0.

\textbf{Chi tiết tích phân và khử sau khi lập hệ.} Trong lúc tích phân theo $\tau$, nhãn $s$ và các gradient ban đầu $p_0(s),q_0(s)$ được giữ cố định. Với $p'=cp$, viết $p=p_0e^{c\tau}$; nếu $X'=dp$ thì $X=x_0+dp_0(e^{c\tau}-1)/c$ khi $c\ne0$, hoặc $X=x_0+dp_0\tau$ khi $c=0$. Với $Y'+cY=R(\tau)$, nhân $e^{c\tau}$ cho $(e^{c\tau}Y)'=e^{c\tau}R$, rồi dùng $Y(0)$ để xác định hằng số. Khử từ hai phương trình $x=X(s,\tau),y=Y(s,\tau)$; ghi rõ tổ hợp nào loại $\tau$ hoặc $s$. Nếu chỉ cho nghiệm tham số, vẫn phải có các tích phân đã tính và điều kiện Jacobian để khử cục bộ.
\subsection*{Lời giải áp dụng cho bài}
\textbf{Bước 1 — Dữ liệu trên đường chéo.} $F=p^2-5q^2-u$, $u(s,s)=s^2$ cho $p_0+q_0=2s$. Thế $p_0=2s-q_0$:
\[(2s-q_0)^2-5q_0^2=s^2\ \Longleftrightarrow\ 4q_0^2+4sq_0-3s^2=0.\]
Hai nhánh là $(p_0,q_0)=(7s/2,-3s/2)$ và $(3s/2,s/2)$.

\textbf{Bước 2 — Hệ đặc trưng.} $X'=2p,Y'=-10q,p'=p,q'=q,U'=2U$. Đặt $E=e^\tau$, tích phân:
\[X=s+2p_0(E-1),\quad Y=s-10q_0(E-1),\quad U=s^2E^2.\]

\textbf{Bước 3 — Khử tham số.} Nhánh thứ nhất có $(7X-3Y)/4=sE$, nhánh thứ hai có $(3X+Y)/4=sE$. Vì vậy
\[u=((7x-3y)/4)^2,\qquad u=((3x+y)/4)^2.\]

\textbf{Bước 4 — Kiểm tra.} Với $u=(ax+by)^2$, PDE cần $4(a^2-5b^2)=1$, đúng cho $(7/4,-3/4)$ và $(3/4,1/4)$; $a+b=1$ cho vết $x^2$. $J_0=-10q_0-2p_0=\pm8s$, nên hai nhánh không đặc trưng ở $s\ne0$.

\textbf{Chi tiết nghiệm bậc hai và khử.} $4q_0^2+4sq_0-3s^2=(2q_0-s)(2q_0+3s)$, nên $q_0=s/2,-3s/2$; lấy $p_0=2s-q_0$. Nhánh $(7s/2,-3s/2)$ có $X=s+7s(E-1)$, $Y=s+15s(E-1)$, do đó $7X-3Y=4s+4s(E-1)=4sE$. Nhánh $(3s/2,s/2)$ có $X=s+3s(E-1)$, $Y=s-5s(E-1)$, nên $3X+Y=4sE$. Do $U=s^2E^2$, mỗi tổ hợp chia $4$ rồi bình phương cho nghiệm tương ứng.

\clearpage\noindent\texttt{\detokenize{BT_1_5_3_03}}\par
\section*{Lời giải bổ sung — Bài tập 3 · Mục 1.5.3}

\subsection*{Kiến thức và trình tự giải}
\textbf{Mục tiêu.} Giải PDE cấp 1 phi tuyến $F(x,y,u,p,q)=0$, với $p=u_x,q=u_y$. Chọn $s$ để gắn nhãn đường dữ liệu, $\tau$ là tham số chạy dọc đặc trưng; hai tham số này có vai trò khác nhau.

\textbf{Bước 1 — Viết đúng $F$ và tham số hóa dữ liệu.} Chuyển PDE về một vế, giữ mọi hệ số chứa $x,y,u$. Viết $X(s,0)=x_0(s)$, $Y(s,0)=y_0(s)$, $U(s,0)=h(s)$. Đạo hàm điều kiện theo $s$ cho

\[
p_0x_0'(s)+q_0y_0'(s)=h'(s),\qquad
F(x_0,y_0,h,p_0,q_0)=0.
\]

Giải hai phương trình này tìm tất cả nhánh thực $(p_0,q_0)$. Trên $y=0$, $p_0=h'$; trên $x=0$, $q_0=h'$; trên $y=x$, $p_0+q_0=h'$. Điều kiện dưới căn phải không âm. Các nhánh khác nhau có thể cho các nghiệm khác nhau cùng dữ liệu.

\textbf{Bước 2 — Kiểm tra không đặc trưng.} Tính $J_0=x_0'F_q-y_0'F_p$ ở dữ liệu của từng nhánh. Khi $J_0\ne0$ và các hàm đủ trơn, phép chuyển từ $(s,\tau)$ sang $(x,y)$ khả nghịch cục bộ. Khi $J_0=0$, cần xét riêng; không kết luận ngay vô nghiệm hoặc duy nhất.

\textbf{Bước 3 — Tính các đạo hàm riêng của $F$ rồi lập hệ Charpit.} Khi tính $F_x,F_y,F_u,F_p,F_q$, tạm coi năm đối số độc lập. Hệ là

\[
\begin{aligned}
X_\tau&=F_p,&Y_\tau&=F_q,&U_\tau&=pF_p+qF_q,\\
p_\tau&=-F_x-pF_u,&q_\tau&=-F_y-qF_u.
\end{aligned}
\]

Trong bài cụ thể phải thay các đạo hàm của $F$ vào hệ. Ví dụ $F=pq-1$ cho $X_\tau=q,Y_\tau=p,U_\tau=2,p_\tau=q_\tau=0$.

\textbf{Bước 4 — Tích phân hệ với điều kiện ban đầu.} Tìm phương trình giải được trước: $p'=0$, $q'=0$, hoặc tổ hợp như $(X-p)'=0$. Thế chúng vào phương trình còn lại. Nếu hệ tuyến tính $Z'=MZ$, dùng $Z=e^{\tau M}Z_0$; để tính cụ thể, tìm các tổ hợp tách hệ hoặc quan hệ $M^2, M^3$ giúp rút gọn chuỗi mũ. Ghi các phép tích phân và hằng số được dữ liệu xác định.

\textbf{Bước 5 — Khử tham số hoặc trình bày nghiệm tham số đã giải.} Từ $X(s,\tau),Y(s,\tau)$ tìm $s,\tau$ theo $x,y$ rồi thay vào $U$. Nếu giữ dạng tham số, phải có công thức $X,Y,U$ đã tích phân và miền $\det\partial(X,Y)/\partial(s,\tau)\ne0$; chỉ ghi hệ ODE chưa phải lời giải hoàn tất.

\textbf{Bước 6 — Dùng phép đổi ẩn khi thuận lợi.} Gần dữ liệu $u>0$, $p^2+q^2=u$ đổi $v=2\sqrt u$ cho $|\nabla v|^2=1$; $p^2+q^2=u^2$ đổi $v=\log u$ cũng cho eikonal. Phải tính $v_x,v_y$ để chứng minh phép đổi, đổi cả dữ liệu và giữ điều kiện dấu. Với dữ liệu hằng trên vòng tròn, đạo hàm tiếp tuyến bằng 0 nên gradient pháp tuyến có hai lựa chọn dấu.

\textbf{Bước 7 — Kiểm tra từng nhánh và miền.} Tính $p,q$ của công thức cuối, thế vào PDE và dữ liệu. Phân biệt một nghiệm đã tìm được, các nhánh không đặc trưng cục bộ và toàn bộ nghiệm toàn cục. Giữ điều kiện của nhãn dữ liệu nếu đề chỉ cho $0\le s\le1$; kiểm tra gốc, điểm căn bằng 0 và đường làm mẫu số bằng 0.

\textbf{Chi tiết tích phân và khử sau khi lập hệ.} Trong lúc tích phân theo $\tau$, nhãn $s$ và các gradient ban đầu $p_0(s),q_0(s)$ được giữ cố định. Với $p'=cp$, viết $p=p_0e^{c\tau}$; nếu $X'=dp$ thì $X=x_0+dp_0(e^{c\tau}-1)/c$ khi $c\ne0$, hoặc $X=x_0+dp_0\tau$ khi $c=0$. Với $Y'+cY=R(\tau)$, nhân $e^{c\tau}$ cho $(e^{c\tau}Y)'=e^{c\tau}R$, rồi dùng $Y(0)$ để xác định hằng số. Khử từ hai phương trình $x=X(s,\tau),y=Y(s,\tau)$; ghi rõ tổ hợp nào loại $\tau$ hoặc $s$. Nếu chỉ cho nghiệm tham số, vẫn phải có các tích phân đã tính và điều kiện Jacobian để khử cục bộ.
\subsection*{Lời giải áp dụng cho bài}
\textbf{Bước 1 — Giữ hệ số $y$ của đề.} $F=p^2+8qu-8yq^2$. Dữ liệu $(s,s,s^2)$ cho $p_0+q_0=2s$. Đặt $q_0=Q$, $p_0=2s-Q$, ta được
\[(1-8s)Q^2+(8s^2-4s)Q+4s^2=0.\]
Biệt thức là $\Delta=64s^3(s+1)$. Vì $\Delta<0$ trên $-1<s<0$, không có nghiệm thực trên lân cận toàn đường dữ liệu. Với $s>0$ hoặc $s<-1$, lấy các nhánh thực
\[Q=\frac{4s-8s^2\pm8\sqrt{s^3(s+1)}}{2(1-8s)}\quad(s\ne1/8).\]
Ở $s=1/8$, phương trình thành tuyến tính và cho $Q=1/6$; chọn nhánh có giới hạn hữu hạn qua điểm này.

\textbf{Bước 2 — Tính và giải hệ Charpit.} $F_x=0,F_y=-8q^2,F_u=8q,F_p=2p,F_q=8u-16yq$. Do đó
\[X'=2p,\quad Y'=8U-16Yq,\quad p'=-8pq,\quad q'=0.\]
Trên $F=0$, $U'=2p^2+8qU-16Yq^2=-8qU$. Vì $q=Q(s)$ hằng trên đường, đặt $E=e^{-8Q\tau}$:
\[p=p_0E,\quad q=Q,\quad U=s^2E,\quad X=s+\frac{p_0}{4Q}(1-E).\]
Giải $Y'+16QY=8s^2E$ bằng thừa số $e^{16Q\tau}$:
\[Y=sE^2+\frac{s^2}{Q}(E-E^2).\]

\textbf{Bước 3 — Nghiệm và miền.} Nghiệm là $u(X(s,\tau),Y(s,\tau))=s^2E$, với từng nhánh $Q$ trơn, tại nơi ánh xạ khả nghịch. Ở dữ liệu,
\[J_0=8s^2-16sQ-2p_0=8s^2-4s+2(1-8s)Q=\pm\sqrt\Delta.\]
Vì vậy có nghiệm cục bộ trên từng nhánh khi $s>0$ hoặc $s<-1$; $s=0,-1$ suy biến cần xét riêng. Trên các khoảng không đặc trưng này, $Q\ne0$ nên các công thức chia cho $Q$ hợp lệ.

\textbf{Bước 4 — Kiểm tra.} Tại $\tau=0$, $X=Y=s,U=s^2$. Thế các công thức vào $p^2+8qU-8Yq^2$ cho $E^2[p_0^2+8Qs^2-8sQ^2]=0$, theo phương trình tương thích. Hệ Charpit cùng điều kiện tiếp tuyến bảo đảm $p=u_x,q=u_y$ trên miền khả nghịch. Không bỏ hệ số $y$ để dùng công thức của một PDE khác.

\textbf{Chi tiết biệt thức và tích phân $Y$.} Với $a=1-8s,b=8s^2-4s,c=4s^2$,
\[\Delta=b^2-4ac=64s^4-64s^3+16s^2-16s^2+128s^3=64s^3(s+1).\]
Với $Q\ne0$, nhân $Y'+16QY=8s^2e^{-8Q\tau}$ bởi $e^{16Q\tau}$:
\[e^{16Q\tau}Y-s=8s^2\int_0^\tau e^{8Qr}dr=\frac{s^2}{Q}(e^{8Q\tau}-1).\]
Nhân $e^{-16Q\tau}$ được $Y=sE^2+(s^2/Q)(E-E^2)$. Đồng thời $X-s=2p_0\int_0^\tau e^{-8Qr}dr=p_0(1-E)/(4Q)$. Các hằng số đều lấy từ dữ liệu tại $\tau=0$.

\clearpage\noindent\texttt{\detokenize{BT_1_5_3_04}}\par
\section*{Lời giải bổ sung — Bài tập 4 · Mục 1.5.3}

\subsection*{Kiến thức và trình tự giải}
\textbf{Mục tiêu.} Giải PDE cấp 1 phi tuyến $F(x,y,u,p,q)=0$, với $p=u_x,q=u_y$. Chọn $s$ để gắn nhãn đường dữ liệu, $\tau$ là tham số chạy dọc đặc trưng; hai tham số này có vai trò khác nhau.

\textbf{Bước 1 — Viết đúng $F$ và tham số hóa dữ liệu.} Chuyển PDE về một vế, giữ mọi hệ số chứa $x,y,u$. Viết $X(s,0)=x_0(s)$, $Y(s,0)=y_0(s)$, $U(s,0)=h(s)$. Đạo hàm điều kiện theo $s$ cho

\[
p_0x_0'(s)+q_0y_0'(s)=h'(s),\qquad
F(x_0,y_0,h,p_0,q_0)=0.
\]

Giải hai phương trình này tìm tất cả nhánh thực $(p_0,q_0)$. Trên $y=0$, $p_0=h'$; trên $x=0$, $q_0=h'$; trên $y=x$, $p_0+q_0=h'$. Điều kiện dưới căn phải không âm. Các nhánh khác nhau có thể cho các nghiệm khác nhau cùng dữ liệu.

\textbf{Bước 2 — Kiểm tra không đặc trưng.} Tính $J_0=x_0'F_q-y_0'F_p$ ở dữ liệu của từng nhánh. Khi $J_0\ne0$ và các hàm đủ trơn, phép chuyển từ $(s,\tau)$ sang $(x,y)$ khả nghịch cục bộ. Khi $J_0=0$, cần xét riêng; không kết luận ngay vô nghiệm hoặc duy nhất.

\textbf{Bước 3 — Tính các đạo hàm riêng của $F$ rồi lập hệ Charpit.} Khi tính $F_x,F_y,F_u,F_p,F_q$, tạm coi năm đối số độc lập. Hệ là

\[
\begin{aligned}
X_\tau&=F_p,&Y_\tau&=F_q,&U_\tau&=pF_p+qF_q,\\
p_\tau&=-F_x-pF_u,&q_\tau&=-F_y-qF_u.
\end{aligned}
\]

Trong bài cụ thể phải thay các đạo hàm của $F$ vào hệ. Ví dụ $F=pq-1$ cho $X_\tau=q,Y_\tau=p,U_\tau=2,p_\tau=q_\tau=0$.

\textbf{Bước 4 — Tích phân hệ với điều kiện ban đầu.} Tìm phương trình giải được trước: $p'=0$, $q'=0$, hoặc tổ hợp như $(X-p)'=0$. Thế chúng vào phương trình còn lại. Nếu hệ tuyến tính $Z'=MZ$, dùng $Z=e^{\tau M}Z_0$; để tính cụ thể, tìm các tổ hợp tách hệ hoặc quan hệ $M^2, M^3$ giúp rút gọn chuỗi mũ. Ghi các phép tích phân và hằng số được dữ liệu xác định.

\textbf{Bước 5 — Khử tham số hoặc trình bày nghiệm tham số đã giải.} Từ $X(s,\tau),Y(s,\tau)$ tìm $s,\tau$ theo $x,y$ rồi thay vào $U$. Nếu giữ dạng tham số, phải có công thức $X,Y,U$ đã tích phân và miền $\det\partial(X,Y)/\partial(s,\tau)\ne0$; chỉ ghi hệ ODE chưa phải lời giải hoàn tất.

\textbf{Bước 6 — Dùng phép đổi ẩn khi thuận lợi.} Gần dữ liệu $u>0$, $p^2+q^2=u$ đổi $v=2\sqrt u$ cho $|\nabla v|^2=1$; $p^2+q^2=u^2$ đổi $v=\log u$ cũng cho eikonal. Phải tính $v_x,v_y$ để chứng minh phép đổi, đổi cả dữ liệu và giữ điều kiện dấu. Với dữ liệu hằng trên vòng tròn, đạo hàm tiếp tuyến bằng 0 nên gradient pháp tuyến có hai lựa chọn dấu.

\textbf{Bước 7 — Kiểm tra từng nhánh và miền.} Tính $p,q$ của công thức cuối, thế vào PDE và dữ liệu. Phân biệt một nghiệm đã tìm được, các nhánh không đặc trưng cục bộ và toàn bộ nghiệm toàn cục. Giữ điều kiện của nhãn dữ liệu nếu đề chỉ cho $0\le s\le1$; kiểm tra gốc, điểm căn bằng 0 và đường làm mẫu số bằng 0.

\textbf{Chi tiết tích phân và khử sau khi lập hệ.} Trong lúc tích phân theo $\tau$, nhãn $s$ và các gradient ban đầu $p_0(s),q_0(s)$ được giữ cố định. Với $p'=cp$, viết $p=p_0e^{c\tau}$; nếu $X'=dp$ thì $X=x_0+dp_0(e^{c\tau}-1)/c$ khi $c\ne0$, hoặc $X=x_0+dp_0\tau$ khi $c=0$. Với $Y'+cY=R(\tau)$, nhân $e^{c\tau}$ cho $(e^{c\tau}Y)'=e^{c\tau}R$, rồi dùng $Y(0)$ để xác định hằng số. Khử từ hai phương trình $x=X(s,\tau),y=Y(s,\tau)$; ghi rõ tổ hợp nào loại $\tau$ hoặc $s$. Nếu chỉ cho nghiệm tham số, vẫn phải có các tích phân đã tính và điều kiện Jacobian để khử cục bộ.
\subsection*{Lời giải áp dụng cho bài}
\textbf{Bước 1 — Tương thích.} $F=p^2+q^2+2xp+2yq-2u$. Dữ liệu $u(s,0)=(1-s^2)/2$ cho $p_0=-s$. Thế PDE: $s^2+q_0^2-2s^2=1-s^2$, suy $q_0^2=1$, $q_0=\pm1$.

\textbf{Bước 2 — Giải Charpit.} $p'=-2p+2p=0$, $q'=-2q+2q=0$. $X'=2p+2X=2(X-s)$, $X(0)=s$, nên $X=s$. $Y'=2q_0+2Y$ cho $Y=q_0(e^{2\tau}-1)$. $U'=2+2q_0Y=2e^{2\tau}$, nên $U=(1-s^2)/2+e^{2\tau}-1$.

\textbf{Bước 3 — Khử và kiểm tra.} $e^{2\tau}-1=q_0Y$, nên
\[u=(1-x^2)/2+y,\qquad u=(1-x^2)/2-y.\]
Với $p=-x,q=\pm1$, vế trái PDE bằng $1-x^2\pm2y=2u$ và vết đúng. $J_0=2q_0\ne0$ nên mỗi nhánh cho nghiệm cục bộ duy nhất khi gradient pháp tuyến đã chọn.

\clearpage\noindent\texttt{\detokenize{BT_1_5_3_05}}\par
\section*{Lời giải bổ sung — Bài tập 5 · Mục 1.5.3}

\subsection*{Kiến thức và trình tự giải}
\textbf{Mục tiêu.} Giải PDE cấp 1 phi tuyến $F(x,y,u,p,q)=0$, với $p=u_x,q=u_y$. Chọn $s$ để gắn nhãn đường dữ liệu, $\tau$ là tham số chạy dọc đặc trưng; hai tham số này có vai trò khác nhau.

\textbf{Bước 1 — Viết đúng $F$ và tham số hóa dữ liệu.} Chuyển PDE về một vế, giữ mọi hệ số chứa $x,y,u$. Viết $X(s,0)=x_0(s)$, $Y(s,0)=y_0(s)$, $U(s,0)=h(s)$. Đạo hàm điều kiện theo $s$ cho

\[
p_0x_0'(s)+q_0y_0'(s)=h'(s),\qquad
F(x_0,y_0,h,p_0,q_0)=0.
\]

Giải hai phương trình này tìm tất cả nhánh thực $(p_0,q_0)$. Trên $y=0$, $p_0=h'$; trên $x=0$, $q_0=h'$; trên $y=x$, $p_0+q_0=h'$. Điều kiện dưới căn phải không âm. Các nhánh khác nhau có thể cho các nghiệm khác nhau cùng dữ liệu.

\textbf{Bước 2 — Kiểm tra không đặc trưng.} Tính $J_0=x_0'F_q-y_0'F_p$ ở dữ liệu của từng nhánh. Khi $J_0\ne0$ và các hàm đủ trơn, phép chuyển từ $(s,\tau)$ sang $(x,y)$ khả nghịch cục bộ. Khi $J_0=0$, cần xét riêng; không kết luận ngay vô nghiệm hoặc duy nhất.

\textbf{Bước 3 — Tính các đạo hàm riêng của $F$ rồi lập hệ Charpit.} Khi tính $F_x,F_y,F_u,F_p,F_q$, tạm coi năm đối số độc lập. Hệ là

\[
\begin{aligned}
X_\tau&=F_p,&Y_\tau&=F_q,&U_\tau&=pF_p+qF_q,\\
p_\tau&=-F_x-pF_u,&q_\tau&=-F_y-qF_u.
\end{aligned}
\]

Trong bài cụ thể phải thay các đạo hàm của $F$ vào hệ. Ví dụ $F=pq-1$ cho $X_\tau=q,Y_\tau=p,U_\tau=2,p_\tau=q_\tau=0$.

\textbf{Bước 4 — Tích phân hệ với điều kiện ban đầu.} Tìm phương trình giải được trước: $p'=0$, $q'=0$, hoặc tổ hợp như $(X-p)'=0$. Thế chúng vào phương trình còn lại. Nếu hệ tuyến tính $Z'=MZ$, dùng $Z=e^{\tau M}Z_0$; để tính cụ thể, tìm các tổ hợp tách hệ hoặc quan hệ $M^2, M^3$ giúp rút gọn chuỗi mũ. Ghi các phép tích phân và hằng số được dữ liệu xác định.

\textbf{Bước 5 — Khử tham số hoặc trình bày nghiệm tham số đã giải.} Từ $X(s,\tau),Y(s,\tau)$ tìm $s,\tau$ theo $x,y$ rồi thay vào $U$. Nếu giữ dạng tham số, phải có công thức $X,Y,U$ đã tích phân và miền $\det\partial(X,Y)/\partial(s,\tau)\ne0$; chỉ ghi hệ ODE chưa phải lời giải hoàn tất.

\textbf{Bước 6 — Dùng phép đổi ẩn khi thuận lợi.} Gần dữ liệu $u>0$, $p^2+q^2=u$ đổi $v=2\sqrt u$ cho $|\nabla v|^2=1$; $p^2+q^2=u^2$ đổi $v=\log u$ cũng cho eikonal. Phải tính $v_x,v_y$ để chứng minh phép đổi, đổi cả dữ liệu và giữ điều kiện dấu. Với dữ liệu hằng trên vòng tròn, đạo hàm tiếp tuyến bằng 0 nên gradient pháp tuyến có hai lựa chọn dấu.

\textbf{Bước 7 — Kiểm tra từng nhánh và miền.} Tính $p,q$ của công thức cuối, thế vào PDE và dữ liệu. Phân biệt một nghiệm đã tìm được, các nhánh không đặc trưng cục bộ và toàn bộ nghiệm toàn cục. Giữ điều kiện của nhãn dữ liệu nếu đề chỉ cho $0\le s\le1$; kiểm tra gốc, điểm căn bằng 0 và đường làm mẫu số bằng 0.

\textbf{Chi tiết tích phân và khử sau khi lập hệ.} Trong lúc tích phân theo $\tau$, nhãn $s$ và các gradient ban đầu $p_0(s),q_0(s)$ được giữ cố định. Với $p'=cp$, viết $p=p_0e^{c\tau}$; nếu $X'=dp$ thì $X=x_0+dp_0(e^{c\tau}-1)/c$ khi $c\ne0$, hoặc $X=x_0+dp_0\tau$ khi $c=0$. Với $Y'+cY=R(\tau)$, nhân $e^{c\tau}$ cho $(e^{c\tau}Y)'=e^{c\tau}R$, rồi dùng $Y(0)$ để xác định hằng số. Khử từ hai phương trình $x=X(s,\tau),y=Y(s,\tau)$; ghi rõ tổ hợp nào loại $\tau$ hoặc $s$. Nếu chỉ cho nghiệm tham số, vẫn phải có các tích phân đã tính và điều kiện Jacobian để khử cục bộ.
\subsection*{Lời giải áp dụng cho bài}
\textbf{Bước 1 — Gradient ban đầu.} $F=p^2+q^2+2yp+2xq-2u$. Dữ liệu $u(s,0)=s^2/2$ cho $p_0=s$ và $q_0^2+2sq_0=0$, nên $q_0=0$ hoặc $-2s$.

\textbf{Bước 2 — Tách hệ Charpit.} Tính
\[X'=2p+2Y,\quad Y'=2q+2X,\quad p'=2p-2q,\quad q'=2q-2p.\]
Đặt $R=p+q$, $D=p-q$, $S=X+Y$, $T=X-Y$. Khi đó $R'=0,D'=4D$, $S'=2S+2R$, $T'=-2T+2D$. Với $R_0=s+q_0,D_0=s-q_0,S_0=T_0=s$, tích phân:
\[R=R_0,\quad D=D_0e^{4\tau},\quad S=(s+R_0)e^{2\tau}-R_0,\]
\[T=(s-D_0/3)e^{-2\tau}+(D_0/3)e^{4\tau}.\]
Lấy $X=(S+T)/2$, $Y=(S-T)/2$, $p=(R+D)/2$, $q=(R-D)/2$. Từ PDE,
\[U=(p^2+q^2+2Yp+2Xq)/2.\]
Đây là công thức tham số đã tích phân bằng các hàm mũ thông thường.

\textbf{Bước 3 — Hai nhánh.} Với $q_0=0$, $R_0=D_0=s$, dùng các công thức trên và khử $s,\tau$ ở nơi Jacobian khác 0; không cần ép thành một đa thức. Với $q_0=-2s$, $R_0=-s,D_0=3s$, nên $S=s,T=se^{4\tau}$, và
\[U=(3T^2-S^2)/4,\qquad u=(x^2+y^2)/2-2xy.\]

\textbf{Bước 4 — Kiểm tra và miền.} Nhánh đa thức có $p=x-2y,q=y-2x$; thế cho $p^2+q^2+2yp+2xq=x^2+y^2-4xy=2u$, vết đúng $x^2/2$. Với nhánh tham số, hệ được tích phân từ đúng Charpit và $U$ được khôi phục từ $F=0$; tại 0 đúng toàn bộ dữ liệu. $J_0=2(q_0+s)=\pm2s$, nên mỗi nhánh có thể khử cục bộ ở $s\ne0$; gốc không được định lý này bảo đảm.

\textbf{Chi tiết thừa số tích phân.} $R'=0$ cho $R=R_0$; $D'=4D$ cho $D=D_0e^{4\tau}$. Với $S'-2S=2R_0$, có $(e^{-2\tau}S)'=2R_0e^{-2\tau}$, nên $e^{-2\tau}S=s+R_0(1-e^{-2\tau})$. Với $T'+2T=2D_0e^{4\tau}$, có $(e^{2\tau}T)'=2D_0e^{6\tau}$, nên $e^{2\tau}T=s+D_0(e^{6\tau}-1)/3$. Đây là hai công thức $S,T$ trong bước 2. Ở nhánh $R=-S,D=3T$, $p=(-S+3T)/2,q=(-S-3T)/2$, và
\[p^2+q^2=(S^2+9T^2)/2,\quad 2Yp+2Xq=-S^2-3T^2.\]
Vậy $2U=(-S^2+3T^2)/2$, tức $U=(3T^2-S^2)/4$.

\clearpage\noindent\texttt{\detokenize{BT_1_5_3_06}}\par
\section*{Lời giải bổ sung — Bài tập 6 · Mục 1.5.3}

\subsection*{Kiến thức và trình tự giải}
\textbf{Mục tiêu.} Giải PDE cấp 1 phi tuyến $F(x,y,u,p,q)=0$, với $p=u_x,q=u_y$. Chọn $s$ để gắn nhãn đường dữ liệu, $\tau$ là tham số chạy dọc đặc trưng; hai tham số này có vai trò khác nhau.

\textbf{Bước 1 — Viết đúng $F$ và tham số hóa dữ liệu.} Chuyển PDE về một vế, giữ mọi hệ số chứa $x,y,u$. Viết $X(s,0)=x_0(s)$, $Y(s,0)=y_0(s)$, $U(s,0)=h(s)$. Đạo hàm điều kiện theo $s$ cho

\[
p_0x_0'(s)+q_0y_0'(s)=h'(s),\qquad
F(x_0,y_0,h,p_0,q_0)=0.
\]

Giải hai phương trình này tìm tất cả nhánh thực $(p_0,q_0)$. Trên $y=0$, $p_0=h'$; trên $x=0$, $q_0=h'$; trên $y=x$, $p_0+q_0=h'$. Điều kiện dưới căn phải không âm. Các nhánh khác nhau có thể cho các nghiệm khác nhau cùng dữ liệu.

\textbf{Bước 2 — Kiểm tra không đặc trưng.} Tính $J_0=x_0'F_q-y_0'F_p$ ở dữ liệu của từng nhánh. Khi $J_0\ne0$ và các hàm đủ trơn, phép chuyển từ $(s,\tau)$ sang $(x,y)$ khả nghịch cục bộ. Khi $J_0=0$, cần xét riêng; không kết luận ngay vô nghiệm hoặc duy nhất.

\textbf{Bước 3 — Tính các đạo hàm riêng của $F$ rồi lập hệ Charpit.} Khi tính $F_x,F_y,F_u,F_p,F_q$, tạm coi năm đối số độc lập. Hệ là

\[
\begin{aligned}
X_\tau&=F_p,&Y_\tau&=F_q,&U_\tau&=pF_p+qF_q,\\
p_\tau&=-F_x-pF_u,&q_\tau&=-F_y-qF_u.
\end{aligned}
\]

Trong bài cụ thể phải thay các đạo hàm của $F$ vào hệ. Ví dụ $F=pq-1$ cho $X_\tau=q,Y_\tau=p,U_\tau=2,p_\tau=q_\tau=0$.

\textbf{Bước 4 — Tích phân hệ với điều kiện ban đầu.} Tìm phương trình giải được trước: $p'=0$, $q'=0$, hoặc tổ hợp như $(X-p)'=0$. Thế chúng vào phương trình còn lại. Nếu hệ tuyến tính $Z'=MZ$, dùng $Z=e^{\tau M}Z_0$; để tính cụ thể, tìm các tổ hợp tách hệ hoặc quan hệ $M^2, M^3$ giúp rút gọn chuỗi mũ. Ghi các phép tích phân và hằng số được dữ liệu xác định.

\textbf{Bước 5 — Khử tham số hoặc trình bày nghiệm tham số đã giải.} Từ $X(s,\tau),Y(s,\tau)$ tìm $s,\tau$ theo $x,y$ rồi thay vào $U$. Nếu giữ dạng tham số, phải có công thức $X,Y,U$ đã tích phân và miền $\det\partial(X,Y)/\partial(s,\tau)\ne0$; chỉ ghi hệ ODE chưa phải lời giải hoàn tất.

\textbf{Bước 6 — Dùng phép đổi ẩn khi thuận lợi.} Gần dữ liệu $u>0$, $p^2+q^2=u$ đổi $v=2\sqrt u$ cho $|\nabla v|^2=1$; $p^2+q^2=u^2$ đổi $v=\log u$ cũng cho eikonal. Phải tính $v_x,v_y$ để chứng minh phép đổi, đổi cả dữ liệu và giữ điều kiện dấu. Với dữ liệu hằng trên vòng tròn, đạo hàm tiếp tuyến bằng 0 nên gradient pháp tuyến có hai lựa chọn dấu.

\textbf{Bước 7 — Kiểm tra từng nhánh và miền.} Tính $p,q$ của công thức cuối, thế vào PDE và dữ liệu. Phân biệt một nghiệm đã tìm được, các nhánh không đặc trưng cục bộ và toàn bộ nghiệm toàn cục. Giữ điều kiện của nhãn dữ liệu nếu đề chỉ cho $0\le s\le1$; kiểm tra gốc, điểm căn bằng 0 và đường làm mẫu số bằng 0.

\textbf{Chi tiết tích phân và khử sau khi lập hệ.} Trong lúc tích phân theo $\tau$, nhãn $s$ và các gradient ban đầu $p_0(s),q_0(s)$ được giữ cố định. Với $p'=cp$, viết $p=p_0e^{c\tau}$; nếu $X'=dp$ thì $X=x_0+dp_0(e^{c\tau}-1)/c$ khi $c\ne0$, hoặc $X=x_0+dp_0\tau$ khi $c=0$. Với $Y'+cY=R(\tau)$, nhân $e^{c\tau}$ cho $(e^{c\tau}Y)'=e^{c\tau}R$, rồi dùng $Y(0)$ để xác định hằng số. Khử từ hai phương trình $x=X(s,\tau),y=Y(s,\tau)$; ghi rõ tổ hợp nào loại $\tau$ hoặc $s$. Nếu chỉ cho nghiệm tham số, vẫn phải có các tích phân đã tính và điều kiện Jacobian để khử cục bộ.
\subsection*{Lời giải áp dụng cho bài}
\textbf{Bước 1 — Tương thích.} $F=4p^2+q^2+4(p-x)(q-4y)-8u$, dữ liệu $u(s,0)=0$ cho $p_0=0$ và $q_0(q_0-4s)=0$.

\textbf{Bước 2 — Hệ Charpit và tổ hợp bất biến.}
\[X'=8p+4q-16Y,\quad Y'=4p+2q-4X,\quad p'=X',\quad q'=4Y'.\]
Suy $X-p=s$, $q-4Y=q_0$. Thế $p=X-s$, $q=4Y+q_0$ vào hai phương trình đầu:
\[X'=8X-8s+4q_0,\quad Y'=8Y+2q_0-4s.\]
Tích phân với $X(0)=s,Y(0)=0$, đặt $E=e^{8\tau}$:
\[X=s+(q_0/2)(E-1),\quad Y=(q_0/4-s/2)(E-1).\]

\textbf{Bước 3 — Khử.} Với $q_0=0$, $X=s,p=0,q=4Y$, PDE cho $u=2y^2$. Với $q_0=4s$, $X=s+4Y$, nên $s=X-4Y$, $p=4Y,q=4X-12Y$, PDE cho $u=4xy-6y^2$.

\textbf{Bước 4 — Kiểm tra.} Nhánh đầu: $p=0,q=4y$ cho vế trái $16y^2=8u$. Nhánh sau: $p=4y,q=4x-12y$ cho vế trái $32xy-48y^2=8u$. Cả hai bằng 0 ở $y=0$. $J_0=2q_0-4s=\mp4s$ khác 0 khi $s\ne0$; hai nhánh trùng ở gốc không cho suy duy nhất tại đó.

\textbf{Chi tiết tích phân và rút gọn nhánh thứ hai.} Vì $p'=X'$, $p-X=-s$; vì $q'=4Y'$, $q-4Y=q_0$. Giải $X'-8X=-8s+4q_0$ bằng nhân $e^{-8\tau}$ cho $X=s+q_0(E-1)/2$. Tương tự $Y'-8Y=2q_0-4s$ cho $Y=(q_0/4-s/2)(E-1)$. Khi $q_0=4s$, $X-s=2s(E-1)=4Y$, nên $p=4Y,q=4X-12Y$. Thế vào PDE:
\[4p^2+q^2+4(p-X)(q-4Y)
=64Y^2+(4X-12Y)^2-16(X-4Y)^2
=32XY-48Y^2.\]
Chia $8$ cho $U=4XY-6Y^2$.

\clearpage\noindent\texttt{\detokenize{BT_1_5_3_07}}\par
\section*{Lời giải bổ sung — Bài tập 7 · Mục 1.5.3}

\subsection*{Kiến thức và trình tự giải}
\textbf{Mục tiêu.} Giải PDE cấp 1 phi tuyến $F(x,y,u,p,q)=0$, với $p=u_x,q=u_y$. Chọn $s$ để gắn nhãn đường dữ liệu, $\tau$ là tham số chạy dọc đặc trưng; hai tham số này có vai trò khác nhau.

\textbf{Bước 1 — Viết đúng $F$ và tham số hóa dữ liệu.} Chuyển PDE về một vế, giữ mọi hệ số chứa $x,y,u$. Viết $X(s,0)=x_0(s)$, $Y(s,0)=y_0(s)$, $U(s,0)=h(s)$. Đạo hàm điều kiện theo $s$ cho

\[
p_0x_0'(s)+q_0y_0'(s)=h'(s),\qquad
F(x_0,y_0,h,p_0,q_0)=0.
\]

Giải hai phương trình này tìm tất cả nhánh thực $(p_0,q_0)$. Trên $y=0$, $p_0=h'$; trên $x=0$, $q_0=h'$; trên $y=x$, $p_0+q_0=h'$. Điều kiện dưới căn phải không âm. Các nhánh khác nhau có thể cho các nghiệm khác nhau cùng dữ liệu.

\textbf{Bước 2 — Kiểm tra không đặc trưng.} Tính $J_0=x_0'F_q-y_0'F_p$ ở dữ liệu của từng nhánh. Khi $J_0\ne0$ và các hàm đủ trơn, phép chuyển từ $(s,\tau)$ sang $(x,y)$ khả nghịch cục bộ. Khi $J_0=0$, cần xét riêng; không kết luận ngay vô nghiệm hoặc duy nhất.

\textbf{Bước 3 — Tính các đạo hàm riêng của $F$ rồi lập hệ Charpit.} Khi tính $F_x,F_y,F_u,F_p,F_q$, tạm coi năm đối số độc lập. Hệ là

\[
\begin{aligned}
X_\tau&=F_p,&Y_\tau&=F_q,&U_\tau&=pF_p+qF_q,\\
p_\tau&=-F_x-pF_u,&q_\tau&=-F_y-qF_u.
\end{aligned}
\]

Trong bài cụ thể phải thay các đạo hàm của $F$ vào hệ. Ví dụ $F=pq-1$ cho $X_\tau=q,Y_\tau=p,U_\tau=2,p_\tau=q_\tau=0$.

\textbf{Bước 4 — Tích phân hệ với điều kiện ban đầu.} Tìm phương trình giải được trước: $p'=0$, $q'=0$, hoặc tổ hợp như $(X-p)'=0$. Thế chúng vào phương trình còn lại. Nếu hệ tuyến tính $Z'=MZ$, dùng $Z=e^{\tau M}Z_0$; để tính cụ thể, tìm các tổ hợp tách hệ hoặc quan hệ $M^2, M^3$ giúp rút gọn chuỗi mũ. Ghi các phép tích phân và hằng số được dữ liệu xác định.

\textbf{Bước 5 — Khử tham số hoặc trình bày nghiệm tham số đã giải.} Từ $X(s,\tau),Y(s,\tau)$ tìm $s,\tau$ theo $x,y$ rồi thay vào $U$. Nếu giữ dạng tham số, phải có công thức $X,Y,U$ đã tích phân và miền $\det\partial(X,Y)/\partial(s,\tau)\ne0$; chỉ ghi hệ ODE chưa phải lời giải hoàn tất.

\textbf{Bước 6 — Dùng phép đổi ẩn khi thuận lợi.} Gần dữ liệu $u>0$, $p^2+q^2=u$ đổi $v=2\sqrt u$ cho $|\nabla v|^2=1$; $p^2+q^2=u^2$ đổi $v=\log u$ cũng cho eikonal. Phải tính $v_x,v_y$ để chứng minh phép đổi, đổi cả dữ liệu và giữ điều kiện dấu. Với dữ liệu hằng trên vòng tròn, đạo hàm tiếp tuyến bằng 0 nên gradient pháp tuyến có hai lựa chọn dấu.

\textbf{Bước 7 — Kiểm tra từng nhánh và miền.} Tính $p,q$ của công thức cuối, thế vào PDE và dữ liệu. Phân biệt một nghiệm đã tìm được, các nhánh không đặc trưng cục bộ và toàn bộ nghiệm toàn cục. Giữ điều kiện của nhãn dữ liệu nếu đề chỉ cho $0\le s\le1$; kiểm tra gốc, điểm căn bằng 0 và đường làm mẫu số bằng 0.

\textbf{Chi tiết tích phân và khử sau khi lập hệ.} Trong lúc tích phân theo $\tau$, nhãn $s$ và các gradient ban đầu $p_0(s),q_0(s)$ được giữ cố định. Với $p'=cp$, viết $p=p_0e^{c\tau}$; nếu $X'=dp$ thì $X=x_0+dp_0(e^{c\tau}-1)/c$ khi $c\ne0$, hoặc $X=x_0+dp_0\tau$ khi $c=0$. Với $Y'+cY=R(\tau)$, nhân $e^{c\tau}$ cho $(e^{c\tau}Y)'=e^{c\tau}R$, rồi dùng $Y(0)$ để xác định hằng số. Khử từ hai phương trình $x=X(s,\tau),y=Y(s,\tau)$; ghi rõ tổ hợp nào loại $\tau$ hoặc $s$. Nếu chỉ cho nghiệm tham số, vẫn phải có các tích phân đã tính và điều kiện Jacobian để khử cục bộ.
\subsection*{Lời giải áp dụng cho bài}
\textbf{Bước 1 — Tương thích và miền log.} $F=pq-1$, $u(s,0)=\log s$ với $s>0$ cho $p_0=1/s$, $q_0=s$.

\textbf{Bước 2 — Giải Charpit.} $p'=q'=0$, $X'=q=s$, $Y'=p=1/s$, $U'=2pq=2$. Do đó $X=s+s\tau$, $Y=\tau/s$, $U=\log s+2\tau$.

\textbf{Bước 3 — Khử.} $\tau=sy$ và $x=s+s^2y$. Nghiệm của phương trình bậc hai liên tục về $s=x$ khi $y=0$ là
\[s=\frac{2x}{1+\sqrt{1+4xy}},\qquad u=\log s+2sy.\]
Chọn miền gần $x>0,y=0$ để $s>0$ và $1+4xy>0$.

\textbf{Bước 4 — Kiểm tra.} Từ $x=s+s^2y$, $s_x=1/(1+2sy)$ và $s_y=-s^2/(1+2sy)$. Đạo hàm $u$ cho $u_x=(1/s+2y)s_x=1/s$, $u_y=(1/s+2y)s_y+2s=s$, nên $pq=1$. Tại $y=0$, $s=x$, $u=\log x$. Mẫu $1+2sy$ cũng là điều kiện khử không suy biến.

\textbf{Chi tiết chọn nghiệm bậc hai.} Khi $y\ne0$, $ys^2+s-x=0$ cho $s=(-1\pm\sqrt{1+4xy})/(2y)$. Nhánh dấu cộng có giới hạn $x$ khi $y\to0$; nhân tử số và mẫu với $1+\sqrt{1+4xy}$ cho $s=2x/(1+\sqrt{1+4xy})$. Nhánh dấu trừ không về dữ liệu $s=x>0$. Sau đó $\tau=sy$ và $U=\log s+2\tau$ cho nghiệm.

\clearpage\noindent\texttt{\detokenize{BT_1_5_3_08}}\par
\section*{Lời giải bổ sung — Bài tập 8 · Mục 1.5.3}

\subsection*{Kiến thức và trình tự giải}
\textbf{Mục tiêu.} Giải PDE cấp 1 phi tuyến $F(x,y,u,p,q)=0$, với $p=u_x,q=u_y$. Chọn $s$ để gắn nhãn đường dữ liệu, $\tau$ là tham số chạy dọc đặc trưng; hai tham số này có vai trò khác nhau.

\textbf{Bước 1 — Viết đúng $F$ và tham số hóa dữ liệu.} Chuyển PDE về một vế, giữ mọi hệ số chứa $x,y,u$. Viết $X(s,0)=x_0(s)$, $Y(s,0)=y_0(s)$, $U(s,0)=h(s)$. Đạo hàm điều kiện theo $s$ cho

\[
p_0x_0'(s)+q_0y_0'(s)=h'(s),\qquad
F(x_0,y_0,h,p_0,q_0)=0.
\]

Giải hai phương trình này tìm tất cả nhánh thực $(p_0,q_0)$. Trên $y=0$, $p_0=h'$; trên $x=0$, $q_0=h'$; trên $y=x$, $p_0+q_0=h'$. Điều kiện dưới căn phải không âm. Các nhánh khác nhau có thể cho các nghiệm khác nhau cùng dữ liệu.

\textbf{Bước 2 — Kiểm tra không đặc trưng.} Tính $J_0=x_0'F_q-y_0'F_p$ ở dữ liệu của từng nhánh. Khi $J_0\ne0$ và các hàm đủ trơn, phép chuyển từ $(s,\tau)$ sang $(x,y)$ khả nghịch cục bộ. Khi $J_0=0$, cần xét riêng; không kết luận ngay vô nghiệm hoặc duy nhất.

\textbf{Bước 3 — Tính các đạo hàm riêng của $F$ rồi lập hệ Charpit.} Khi tính $F_x,F_y,F_u,F_p,F_q$, tạm coi năm đối số độc lập. Hệ là

\[
\begin{aligned}
X_\tau&=F_p,&Y_\tau&=F_q,&U_\tau&=pF_p+qF_q,\\
p_\tau&=-F_x-pF_u,&q_\tau&=-F_y-qF_u.
\end{aligned}
\]

Trong bài cụ thể phải thay các đạo hàm của $F$ vào hệ. Ví dụ $F=pq-1$ cho $X_\tau=q,Y_\tau=p,U_\tau=2,p_\tau=q_\tau=0$.

\textbf{Bước 4 — Tích phân hệ với điều kiện ban đầu.} Tìm phương trình giải được trước: $p'=0$, $q'=0$, hoặc tổ hợp như $(X-p)'=0$. Thế chúng vào phương trình còn lại. Nếu hệ tuyến tính $Z'=MZ$, dùng $Z=e^{\tau M}Z_0$; để tính cụ thể, tìm các tổ hợp tách hệ hoặc quan hệ $M^2, M^3$ giúp rút gọn chuỗi mũ. Ghi các phép tích phân và hằng số được dữ liệu xác định.

\textbf{Bước 5 — Khử tham số hoặc trình bày nghiệm tham số đã giải.} Từ $X(s,\tau),Y(s,\tau)$ tìm $s,\tau$ theo $x,y$ rồi thay vào $U$. Nếu giữ dạng tham số, phải có công thức $X,Y,U$ đã tích phân và miền $\det\partial(X,Y)/\partial(s,\tau)\ne0$; chỉ ghi hệ ODE chưa phải lời giải hoàn tất.

\textbf{Bước 6 — Dùng phép đổi ẩn khi thuận lợi.} Gần dữ liệu $u>0$, $p^2+q^2=u$ đổi $v=2\sqrt u$ cho $|\nabla v|^2=1$; $p^2+q^2=u^2$ đổi $v=\log u$ cũng cho eikonal. Phải tính $v_x,v_y$ để chứng minh phép đổi, đổi cả dữ liệu và giữ điều kiện dấu. Với dữ liệu hằng trên vòng tròn, đạo hàm tiếp tuyến bằng 0 nên gradient pháp tuyến có hai lựa chọn dấu.

\textbf{Bước 7 — Kiểm tra từng nhánh và miền.} Tính $p,q$ của công thức cuối, thế vào PDE và dữ liệu. Phân biệt một nghiệm đã tìm được, các nhánh không đặc trưng cục bộ và toàn bộ nghiệm toàn cục. Giữ điều kiện của nhãn dữ liệu nếu đề chỉ cho $0\le s\le1$; kiểm tra gốc, điểm căn bằng 0 và đường làm mẫu số bằng 0.

\textbf{Chi tiết tích phân và khử sau khi lập hệ.} Trong lúc tích phân theo $\tau$, nhãn $s$ và các gradient ban đầu $p_0(s),q_0(s)$ được giữ cố định. Với $p'=cp$, viết $p=p_0e^{c\tau}$; nếu $X'=dp$ thì $X=x_0+dp_0(e^{c\tau}-1)/c$ khi $c\ne0$, hoặc $X=x_0+dp_0\tau$ khi $c=0$. Với $Y'+cY=R(\tau)$, nhân $e^{c\tau}$ cho $(e^{c\tau}Y)'=e^{c\tau}R$, rồi dùng $Y(0)$ để xác định hằng số. Khử từ hai phương trình $x=X(s,\tau),y=Y(s,\tau)$; ghi rõ tổ hợp nào loại $\tau$ hoặc $s$. Nếu chỉ cho nghiệm tham số, vẫn phải có các tích phân đã tính và điều kiện Jacobian để khử cục bộ.
\subsection*{Lời giải áp dụng cho bài}
\textbf{Bước 1 — Dữ liệu trên $y=1$.} $F=pq-u$, $u(s,1)=s$ cho $p_0=1$, $q_0=s$.

\textbf{Bước 2 — Hệ đặc trưng.} $p'=p,q'=q$, $X'=q,Y'=p$, $U'=2U$. Tích phân:
\[p=e^\tau,\quad q=se^\tau,\quad X=se^\tau,\quad Y=e^\tau,\quad U=se^{2\tau}.\]

\textbf{Bước 3 — Khử và thử lại.} $U=XY$, nên $u=xy$. $u_x=y,u_y=x$ cho $pq=xy=u$, và $u(x,1)=x$. Vì dữ liệu chỉ có $0\le s\le1$, miền được đặc trưng phủ cần $y>0$ và $0\le x/y\le1$. Đa thức $xy$ là một kéo dài toàn cục; dữ liệu đoạn không tự xác định duy nhất ở mọi điểm ngoài miền phủ.

\clearpage\noindent\texttt{\detokenize{BT_1_5_3_09}}\par
\section*{Lời giải bổ sung — Bài tập 9 · Mục 1.5.3}

\subsection*{Kiến thức và trình tự giải}
\textbf{Mục tiêu.} Giải PDE cấp 1 phi tuyến $F(x,y,u,p,q)=0$, với $p=u_x,q=u_y$. Chọn $s$ để gắn nhãn đường dữ liệu, $\tau$ là tham số chạy dọc đặc trưng; hai tham số này có vai trò khác nhau.

\textbf{Bước 1 — Viết đúng $F$ và tham số hóa dữ liệu.} Chuyển PDE về một vế, giữ mọi hệ số chứa $x,y,u$. Viết $X(s,0)=x_0(s)$, $Y(s,0)=y_0(s)$, $U(s,0)=h(s)$. Đạo hàm điều kiện theo $s$ cho

\[
p_0x_0'(s)+q_0y_0'(s)=h'(s),\qquad
F(x_0,y_0,h,p_0,q_0)=0.
\]

Giải hai phương trình này tìm tất cả nhánh thực $(p_0,q_0)$. Trên $y=0$, $p_0=h'$; trên $x=0$, $q_0=h'$; trên $y=x$, $p_0+q_0=h'$. Điều kiện dưới căn phải không âm. Các nhánh khác nhau có thể cho các nghiệm khác nhau cùng dữ liệu.

\textbf{Bước 2 — Kiểm tra không đặc trưng.} Tính $J_0=x_0'F_q-y_0'F_p$ ở dữ liệu của từng nhánh. Khi $J_0\ne0$ và các hàm đủ trơn, phép chuyển từ $(s,\tau)$ sang $(x,y)$ khả nghịch cục bộ. Khi $J_0=0$, cần xét riêng; không kết luận ngay vô nghiệm hoặc duy nhất.

\textbf{Bước 3 — Tính các đạo hàm riêng của $F$ rồi lập hệ Charpit.} Khi tính $F_x,F_y,F_u,F_p,F_q$, tạm coi năm đối số độc lập. Hệ là

\[
\begin{aligned}
X_\tau&=F_p,&Y_\tau&=F_q,&U_\tau&=pF_p+qF_q,\\
p_\tau&=-F_x-pF_u,&q_\tau&=-F_y-qF_u.
\end{aligned}
\]

Trong bài cụ thể phải thay các đạo hàm của $F$ vào hệ. Ví dụ $F=pq-1$ cho $X_\tau=q,Y_\tau=p,U_\tau=2,p_\tau=q_\tau=0$.

\textbf{Bước 4 — Tích phân hệ với điều kiện ban đầu.} Tìm phương trình giải được trước: $p'=0$, $q'=0$, hoặc tổ hợp như $(X-p)'=0$. Thế chúng vào phương trình còn lại. Nếu hệ tuyến tính $Z'=MZ$, dùng $Z=e^{\tau M}Z_0$; để tính cụ thể, tìm các tổ hợp tách hệ hoặc quan hệ $M^2, M^3$ giúp rút gọn chuỗi mũ. Ghi các phép tích phân và hằng số được dữ liệu xác định.

\textbf{Bước 5 — Khử tham số hoặc trình bày nghiệm tham số đã giải.} Từ $X(s,\tau),Y(s,\tau)$ tìm $s,\tau$ theo $x,y$ rồi thay vào $U$. Nếu giữ dạng tham số, phải có công thức $X,Y,U$ đã tích phân và miền $\det\partial(X,Y)/\partial(s,\tau)\ne0$; chỉ ghi hệ ODE chưa phải lời giải hoàn tất.

\textbf{Bước 6 — Dùng phép đổi ẩn khi thuận lợi.} Gần dữ liệu $u>0$, $p^2+q^2=u$ đổi $v=2\sqrt u$ cho $|\nabla v|^2=1$; $p^2+q^2=u^2$ đổi $v=\log u$ cũng cho eikonal. Phải tính $v_x,v_y$ để chứng minh phép đổi, đổi cả dữ liệu và giữ điều kiện dấu. Với dữ liệu hằng trên vòng tròn, đạo hàm tiếp tuyến bằng 0 nên gradient pháp tuyến có hai lựa chọn dấu.

\textbf{Bước 7 — Kiểm tra từng nhánh và miền.} Tính $p,q$ của công thức cuối, thế vào PDE và dữ liệu. Phân biệt một nghiệm đã tìm được, các nhánh không đặc trưng cục bộ và toàn bộ nghiệm toàn cục. Giữ điều kiện của nhãn dữ liệu nếu đề chỉ cho $0\le s\le1$; kiểm tra gốc, điểm căn bằng 0 và đường làm mẫu số bằng 0.

\textbf{Chi tiết tích phân và khử sau khi lập hệ.} Trong lúc tích phân theo $\tau$, nhãn $s$ và các gradient ban đầu $p_0(s),q_0(s)$ được giữ cố định. Với $p'=cp$, viết $p=p_0e^{c\tau}$; nếu $X'=dp$ thì $X=x_0+dp_0(e^{c\tau}-1)/c$ khi $c\ne0$, hoặc $X=x_0+dp_0\tau$ khi $c=0$. Với $Y'+cY=R(\tau)$, nhân $e^{c\tau}$ cho $(e^{c\tau}Y)'=e^{c\tau}R$, rồi dùng $Y(0)$ để xác định hằng số. Khử từ hai phương trình $x=X(s,\tau),y=Y(s,\tau)$; ghi rõ tổ hợp nào loại $\tau$ hoặc $s$. Nếu chỉ cho nghiệm tham số, vẫn phải có các tích phân đã tính và điều kiện Jacobian để khử cục bộ.
\subsection*{Lời giải áp dụng cho bài}
\textbf{Bước 1 — Tương thích.} $F=pq-2$, $u(s,s)=3s$ cho $p_0+q_0=3$, $p_0q_0=2$. Hai nghiệm của $z^2-3z+2=0$ là 1,2, nên $(p_0,q_0)=(1,2)$ hoặc $(2,1)$.

\textbf{Bước 2 — Tích phân.} $p,q$ hằng; $X=s+q_0\tau$, $Y=s+p_0\tau$, $U=3s+4\tau$. Với nhánh $(1,2)$, $U=X+2Y$; nhánh $(2,1)$ cho $U=2X+Y$.

\textbf{Bước 3 — Nghiệm và miền.} $u=x+2y$ hoặc $u=2x+y$. Gradient có tích 2, vết bằng $3x$, nên cả hai đúng. $J_0=p_0-q_0=\mp1$ khác 0. Từ đoạn $0\le s\le1$, nhánh đầu cần $s=2y-x\in[0,1]$, nhánh sau cần $s=2x-y\in[0,1]$. Các công thức affine kéo dài được, nhưng kết luận duy nhất cục bộ áp cho từng nhánh đã chọn.

\clearpage\noindent\texttt{\detokenize{BT_1_5_3_10}}\par
\section*{Lời giải bổ sung — Bài tập 10 · Mục 1.5.3}

\subsection*{Kiến thức và trình tự giải}
\textbf{Mục tiêu.} Giải PDE cấp 1 phi tuyến $F(x,y,u,p,q)=0$, với $p=u_x,q=u_y$. Chọn $s$ để gắn nhãn đường dữ liệu, $\tau$ là tham số chạy dọc đặc trưng; hai tham số này có vai trò khác nhau.

\textbf{Bước 1 — Viết đúng $F$ và tham số hóa dữ liệu.} Chuyển PDE về một vế, giữ mọi hệ số chứa $x,y,u$. Viết $X(s,0)=x_0(s)$, $Y(s,0)=y_0(s)$, $U(s,0)=h(s)$. Đạo hàm điều kiện theo $s$ cho

\[
p_0x_0'(s)+q_0y_0'(s)=h'(s),\qquad
F(x_0,y_0,h,p_0,q_0)=0.
\]

Giải hai phương trình này tìm tất cả nhánh thực $(p_0,q_0)$. Trên $y=0$, $p_0=h'$; trên $x=0$, $q_0=h'$; trên $y=x$, $p_0+q_0=h'$. Điều kiện dưới căn phải không âm. Các nhánh khác nhau có thể cho các nghiệm khác nhau cùng dữ liệu.

\textbf{Bước 2 — Kiểm tra không đặc trưng.} Tính $J_0=x_0'F_q-y_0'F_p$ ở dữ liệu của từng nhánh. Khi $J_0\ne0$ và các hàm đủ trơn, phép chuyển từ $(s,\tau)$ sang $(x,y)$ khả nghịch cục bộ. Khi $J_0=0$, cần xét riêng; không kết luận ngay vô nghiệm hoặc duy nhất.

\textbf{Bước 3 — Tính các đạo hàm riêng của $F$ rồi lập hệ Charpit.} Khi tính $F_x,F_y,F_u,F_p,F_q$, tạm coi năm đối số độc lập. Hệ là

\[
\begin{aligned}
X_\tau&=F_p,&Y_\tau&=F_q,&U_\tau&=pF_p+qF_q,\\
p_\tau&=-F_x-pF_u,&q_\tau&=-F_y-qF_u.
\end{aligned}
\]

Trong bài cụ thể phải thay các đạo hàm của $F$ vào hệ. Ví dụ $F=pq-1$ cho $X_\tau=q,Y_\tau=p,U_\tau=2,p_\tau=q_\tau=0$.

\textbf{Bước 4 — Tích phân hệ với điều kiện ban đầu.} Tìm phương trình giải được trước: $p'=0$, $q'=0$, hoặc tổ hợp như $(X-p)'=0$. Thế chúng vào phương trình còn lại. Nếu hệ tuyến tính $Z'=MZ$, dùng $Z=e^{\tau M}Z_0$; để tính cụ thể, tìm các tổ hợp tách hệ hoặc quan hệ $M^2, M^3$ giúp rút gọn chuỗi mũ. Ghi các phép tích phân và hằng số được dữ liệu xác định.

\textbf{Bước 5 — Khử tham số hoặc trình bày nghiệm tham số đã giải.} Từ $X(s,\tau),Y(s,\tau)$ tìm $s,\tau$ theo $x,y$ rồi thay vào $U$. Nếu giữ dạng tham số, phải có công thức $X,Y,U$ đã tích phân và miền $\det\partial(X,Y)/\partial(s,\tau)\ne0$; chỉ ghi hệ ODE chưa phải lời giải hoàn tất.

\textbf{Bước 6 — Dùng phép đổi ẩn khi thuận lợi.} Gần dữ liệu $u>0$, $p^2+q^2=u$ đổi $v=2\sqrt u$ cho $|\nabla v|^2=1$; $p^2+q^2=u^2$ đổi $v=\log u$ cũng cho eikonal. Phải tính $v_x,v_y$ để chứng minh phép đổi, đổi cả dữ liệu và giữ điều kiện dấu. Với dữ liệu hằng trên vòng tròn, đạo hàm tiếp tuyến bằng 0 nên gradient pháp tuyến có hai lựa chọn dấu.

\textbf{Bước 7 — Kiểm tra từng nhánh và miền.} Tính $p,q$ của công thức cuối, thế vào PDE và dữ liệu. Phân biệt một nghiệm đã tìm được, các nhánh không đặc trưng cục bộ và toàn bộ nghiệm toàn cục. Giữ điều kiện của nhãn dữ liệu nếu đề chỉ cho $0\le s\le1$; kiểm tra gốc, điểm căn bằng 0 và đường làm mẫu số bằng 0.

\textbf{Chi tiết tích phân và khử sau khi lập hệ.} Trong lúc tích phân theo $\tau$, nhãn $s$ và các gradient ban đầu $p_0(s),q_0(s)$ được giữ cố định. Với $p'=cp$, viết $p=p_0e^{c\tau}$; nếu $X'=dp$ thì $X=x_0+dp_0(e^{c\tau}-1)/c$ khi $c\ne0$, hoặc $X=x_0+dp_0\tau$ khi $c=0$. Với $Y'+cY=R(\tau)$, nhân $e^{c\tau}$ cho $(e^{c\tau}Y)'=e^{c\tau}R$, rồi dùng $Y(0)$ để xác định hằng số. Khử từ hai phương trình $x=X(s,\tau),y=Y(s,\tau)$; ghi rõ tổ hợp nào loại $\tau$ hoặc $s$. Nếu chỉ cho nghiệm tham số, vẫn phải có các tích phân đã tính và điều kiện Jacobian để khử cục bộ.
\subsection*{Lời giải áp dụng cho bài}
\textbf{Bước 1 — Gradient đầu.} $F=p^3-q$, $u(s,0)=2s^{3/2}$, $0\le s\le1$, cho $p_0=3\sqrt s$, $q_0=p_0^3=27s^{3/2}$.

\textbf{Bước 2 — Charpit.} $p,q$ hằng, $X'=3p^2=27s$, $Y'=-1$, $U'=3p^3-q=54s^{3/2}$. Suy $X=s+27s\tau$, $Y=-\tau$, $U=2s^{3/2}(1+27\tau)$.

\textbf{Bước 3 — Khử.} $\tau=-y$, $s=x/(1-27y)$, nên
\[u=2x^{3/2}/\sqrt{1-27y}.\]
Chọn miền gần dữ liệu với $x>0$, $y<1/27$; miền từ đoạn dữ liệu còn cần $0\le x/(1-27y)\le1$.

\textbf{Bước 4 — Kiểm tra.} Đặt $D=1-27y$. $u_x=3\sqrt{x}/\sqrt D$, $u_y=27x^{3/2}/D^{3/2}$, nên $u_x^3=u_y$; tại $y=0$ đúng dữ liệu. Ở $x=0$ công thức kéo dài $C^1$ từ phía phải, nhưng $u_{xx}$ không hữu hạn, nên không gọi là nghiệm $C^2$ tại đầu mút này.

\clearpage\noindent\texttt{\detokenize{BT_1_5_3_11}}\par
\section*{Lời giải bổ sung — Bài tập 11 · Mục 1.5.3}

\subsection*{Kiến thức và trình tự giải}
\textbf{Mục tiêu.} Giải PDE cấp 1 phi tuyến $F(x,y,u,p,q)=0$, với $p=u_x,q=u_y$. Chọn $s$ để gắn nhãn đường dữ liệu, $\tau$ là tham số chạy dọc đặc trưng; hai tham số này có vai trò khác nhau.

\textbf{Bước 1 — Viết đúng $F$ và tham số hóa dữ liệu.} Chuyển PDE về một vế, giữ mọi hệ số chứa $x,y,u$. Viết $X(s,0)=x_0(s)$, $Y(s,0)=y_0(s)$, $U(s,0)=h(s)$. Đạo hàm điều kiện theo $s$ cho

\[
p_0x_0'(s)+q_0y_0'(s)=h'(s),\qquad
F(x_0,y_0,h,p_0,q_0)=0.
\]

Giải hai phương trình này tìm tất cả nhánh thực $(p_0,q_0)$. Trên $y=0$, $p_0=h'$; trên $x=0$, $q_0=h'$; trên $y=x$, $p_0+q_0=h'$. Điều kiện dưới căn phải không âm. Các nhánh khác nhau có thể cho các nghiệm khác nhau cùng dữ liệu.

\textbf{Bước 2 — Kiểm tra không đặc trưng.} Tính $J_0=x_0'F_q-y_0'F_p$ ở dữ liệu của từng nhánh. Khi $J_0\ne0$ và các hàm đủ trơn, phép chuyển từ $(s,\tau)$ sang $(x,y)$ khả nghịch cục bộ. Khi $J_0=0$, cần xét riêng; không kết luận ngay vô nghiệm hoặc duy nhất.

\textbf{Bước 3 — Tính các đạo hàm riêng của $F$ rồi lập hệ Charpit.} Khi tính $F_x,F_y,F_u,F_p,F_q$, tạm coi năm đối số độc lập. Hệ là

\[
\begin{aligned}
X_\tau&=F_p,&Y_\tau&=F_q,&U_\tau&=pF_p+qF_q,\\
p_\tau&=-F_x-pF_u,&q_\tau&=-F_y-qF_u.
\end{aligned}
\]

Trong bài cụ thể phải thay các đạo hàm của $F$ vào hệ. Ví dụ $F=pq-1$ cho $X_\tau=q,Y_\tau=p,U_\tau=2,p_\tau=q_\tau=0$.

\textbf{Bước 4 — Tích phân hệ với điều kiện ban đầu.} Tìm phương trình giải được trước: $p'=0$, $q'=0$, hoặc tổ hợp như $(X-p)'=0$. Thế chúng vào phương trình còn lại. Nếu hệ tuyến tính $Z'=MZ$, dùng $Z=e^{\tau M}Z_0$; để tính cụ thể, tìm các tổ hợp tách hệ hoặc quan hệ $M^2, M^3$ giúp rút gọn chuỗi mũ. Ghi các phép tích phân và hằng số được dữ liệu xác định.

\textbf{Bước 5 — Khử tham số hoặc trình bày nghiệm tham số đã giải.} Từ $X(s,\tau),Y(s,\tau)$ tìm $s,\tau$ theo $x,y$ rồi thay vào $U$. Nếu giữ dạng tham số, phải có công thức $X,Y,U$ đã tích phân và miền $\det\partial(X,Y)/\partial(s,\tau)\ne0$; chỉ ghi hệ ODE chưa phải lời giải hoàn tất.

\textbf{Bước 6 — Dùng phép đổi ẩn khi thuận lợi.} Gần dữ liệu $u>0$, $p^2+q^2=u$ đổi $v=2\sqrt u$ cho $|\nabla v|^2=1$; $p^2+q^2=u^2$ đổi $v=\log u$ cũng cho eikonal. Phải tính $v_x,v_y$ để chứng minh phép đổi, đổi cả dữ liệu và giữ điều kiện dấu. Với dữ liệu hằng trên vòng tròn, đạo hàm tiếp tuyến bằng 0 nên gradient pháp tuyến có hai lựa chọn dấu.

\textbf{Bước 7 — Kiểm tra từng nhánh và miền.} Tính $p,q$ của công thức cuối, thế vào PDE và dữ liệu. Phân biệt một nghiệm đã tìm được, các nhánh không đặc trưng cục bộ và toàn bộ nghiệm toàn cục. Giữ điều kiện của nhãn dữ liệu nếu đề chỉ cho $0\le s\le1$; kiểm tra gốc, điểm căn bằng 0 và đường làm mẫu số bằng 0.

\textbf{Chi tiết tích phân và khử sau khi lập hệ.} Trong lúc tích phân theo $\tau$, nhãn $s$ và các gradient ban đầu $p_0(s),q_0(s)$ được giữ cố định. Với $p'=cp$, viết $p=p_0e^{c\tau}$; nếu $X'=dp$ thì $X=x_0+dp_0(e^{c\tau}-1)/c$ khi $c\ne0$, hoặc $X=x_0+dp_0\tau$ khi $c=0$. Với $Y'+cY=R(\tau)$, nhân $e^{c\tau}$ cho $(e^{c\tau}Y)'=e^{c\tau}R$, rồi dùng $Y(0)$ để xác định hằng số. Khử từ hai phương trình $x=X(s,\tau),y=Y(s,\tau)$; ghi rõ tổ hợp nào loại $\tau$ hoặc $s$. Nếu chỉ cho nghiệm tham số, vẫn phải có các tích phân đã tính và điều kiện Jacobian để khử cục bộ.
\subsection*{Lời giải áp dụng cho bài}
\textbf{Bước 1 — Gradient đầu.} $F=p+q^2/2-1$, $u(0,s)=s^2$, $0\le s\le1$, cho $q_0=2s$, $p_0=1-2s^2$.

\textbf{Bước 2 — Tích phân.} $p,q$ hằng, $X'=1$, $Y'=q_0=2s$, $U'=p_0+q_0^2=1+2s^2$. Do đó $X=\tau$, $Y=s(1+2\tau)$, $U=s^2+\tau(1+2s^2)$.

\textbf{Bước 3 — Khử.} $\tau=x$, $s=y/(1+2x)$, suy $u=x+y^2/(1+2x)$. Miền cục bộ quanh dữ liệu lấy $1+2x>0$; còn cần $s\in[0,1]$. Công thức đại số cũng có nghĩa trên miền $x<-1/2$, nhưng miền này không nối với dữ liệu qua phép khử đang dùng.

\textbf{Bước 4 — Kiểm tra.} $p=1-2y^2/(1+2x)^2$, $q=2y/(1+2x)$, nên $p+q^2/2=1$; ở $x=0$ bằng $y^2$. Đường $x=-1/2$ làm Jacobian mất khả nghịch.

\clearpage\noindent\texttt{\detokenize{BT_1_5_3_12}}\par
\section*{Lời giải bổ sung — Bài tập 12 · Mục 1.5.3}

\subsection*{Kiến thức và trình tự giải}
\textbf{Mục tiêu.} Giải PDE cấp 1 phi tuyến $F(x,y,u,p,q)=0$, với $p=u_x,q=u_y$. Chọn $s$ để gắn nhãn đường dữ liệu, $\tau$ là tham số chạy dọc đặc trưng; hai tham số này có vai trò khác nhau.

\textbf{Bước 1 — Viết đúng $F$ và tham số hóa dữ liệu.} Chuyển PDE về một vế, giữ mọi hệ số chứa $x,y,u$. Viết $X(s,0)=x_0(s)$, $Y(s,0)=y_0(s)$, $U(s,0)=h(s)$. Đạo hàm điều kiện theo $s$ cho

\[
p_0x_0'(s)+q_0y_0'(s)=h'(s),\qquad
F(x_0,y_0,h,p_0,q_0)=0.
\]

Giải hai phương trình này tìm tất cả nhánh thực $(p_0,q_0)$. Trên $y=0$, $p_0=h'$; trên $x=0$, $q_0=h'$; trên $y=x$, $p_0+q_0=h'$. Điều kiện dưới căn phải không âm. Các nhánh khác nhau có thể cho các nghiệm khác nhau cùng dữ liệu.

\textbf{Bước 2 — Kiểm tra không đặc trưng.} Tính $J_0=x_0'F_q-y_0'F_p$ ở dữ liệu của từng nhánh. Khi $J_0\ne0$ và các hàm đủ trơn, phép chuyển từ $(s,\tau)$ sang $(x,y)$ khả nghịch cục bộ. Khi $J_0=0$, cần xét riêng; không kết luận ngay vô nghiệm hoặc duy nhất.

\textbf{Bước 3 — Tính các đạo hàm riêng của $F$ rồi lập hệ Charpit.} Khi tính $F_x,F_y,F_u,F_p,F_q$, tạm coi năm đối số độc lập. Hệ là

\[
\begin{aligned}
X_\tau&=F_p,&Y_\tau&=F_q,&U_\tau&=pF_p+qF_q,\\
p_\tau&=-F_x-pF_u,&q_\tau&=-F_y-qF_u.
\end{aligned}
\]

Trong bài cụ thể phải thay các đạo hàm của $F$ vào hệ. Ví dụ $F=pq-1$ cho $X_\tau=q,Y_\tau=p,U_\tau=2,p_\tau=q_\tau=0$.

\textbf{Bước 4 — Tích phân hệ với điều kiện ban đầu.} Tìm phương trình giải được trước: $p'=0$, $q'=0$, hoặc tổ hợp như $(X-p)'=0$. Thế chúng vào phương trình còn lại. Nếu hệ tuyến tính $Z'=MZ$, dùng $Z=e^{\tau M}Z_0$; để tính cụ thể, tìm các tổ hợp tách hệ hoặc quan hệ $M^2, M^3$ giúp rút gọn chuỗi mũ. Ghi các phép tích phân và hằng số được dữ liệu xác định.

\textbf{Bước 5 — Khử tham số hoặc trình bày nghiệm tham số đã giải.} Từ $X(s,\tau),Y(s,\tau)$ tìm $s,\tau$ theo $x,y$ rồi thay vào $U$. Nếu giữ dạng tham số, phải có công thức $X,Y,U$ đã tích phân và miền $\det\partial(X,Y)/\partial(s,\tau)\ne0$; chỉ ghi hệ ODE chưa phải lời giải hoàn tất.

\textbf{Bước 6 — Dùng phép đổi ẩn khi thuận lợi.} Gần dữ liệu $u>0$, $p^2+q^2=u$ đổi $v=2\sqrt u$ cho $|\nabla v|^2=1$; $p^2+q^2=u^2$ đổi $v=\log u$ cũng cho eikonal. Phải tính $v_x,v_y$ để chứng minh phép đổi, đổi cả dữ liệu và giữ điều kiện dấu. Với dữ liệu hằng trên vòng tròn, đạo hàm tiếp tuyến bằng 0 nên gradient pháp tuyến có hai lựa chọn dấu.

\textbf{Bước 7 — Kiểm tra từng nhánh và miền.} Tính $p,q$ của công thức cuối, thế vào PDE và dữ liệu. Phân biệt một nghiệm đã tìm được, các nhánh không đặc trưng cục bộ và toàn bộ nghiệm toàn cục. Giữ điều kiện của nhãn dữ liệu nếu đề chỉ cho $0\le s\le1$; kiểm tra gốc, điểm căn bằng 0 và đường làm mẫu số bằng 0.

\textbf{Chi tiết tích phân và khử sau khi lập hệ.} Trong lúc tích phân theo $\tau$, nhãn $s$ và các gradient ban đầu $p_0(s),q_0(s)$ được giữ cố định. Với $p'=cp$, viết $p=p_0e^{c\tau}$; nếu $X'=dp$ thì $X=x_0+dp_0(e^{c\tau}-1)/c$ khi $c\ne0$, hoặc $X=x_0+dp_0\tau$ khi $c=0$. Với $Y'+cY=R(\tau)$, nhân $e^{c\tau}$ cho $(e^{c\tau}Y)'=e^{c\tau}R$, rồi dùng $Y(0)$ để xác định hằng số. Khử từ hai phương trình $x=X(s,\tau),y=Y(s,\tau)$; ghi rõ tổ hợp nào loại $\tau$ hoặc $s$. Nếu chỉ cho nghiệm tham số, vẫn phải có các tích phân đã tính và điều kiện Jacobian để khử cục bộ.
\subsection*{Lời giải áp dụng cho bài}
\textbf{Bước 1 — Tương thích.} $F=2xp^2+yq-u$, $u(s,1)=s/2$ cho $p_0=1/2$. Thế PDE được $2s/4+q_0-s/2=0$, nên $q_0=0$.

\textbf{Bước 2 — Giải Charpit.} $p'=p-2p^2$, $q'=0$, nên với dữ liệu này $p=1/2,q=0$. $X'=4Xp=2X$, $Y'=Y$, $U'=4Xp^2+Yq=X$. Suy $X=se^{2\tau}$, $Y=e^\tau$, $U=se^{2\tau}/2$.

\textbf{Bước 3 — Khử và kiểm tra.} $U=X/2$, nên $u=x/2$. Gradient $p=1/2,q=0$ cho $2xp^2+yq=x/2=u$, và vết trên $y=1$ đúng. $J_0=F_q=1$ nên chỉ có một nhánh không đặc trưng cục bộ. Đây là trường hợp hệ có nghiệm đơn giản, không cần giữ mũ ma trận.

\clearpage\noindent\texttt{\detokenize{BT_1_5_3_13}}\par
\section*{Lời giải bổ sung — Bài tập 13 · Mục 1.5.3}

\subsection*{Kiến thức và trình tự giải}
\textbf{Mục tiêu.} Giải PDE cấp 1 phi tuyến $F(x,y,u,p,q)=0$, với $p=u_x,q=u_y$. Chọn $s$ để gắn nhãn đường dữ liệu, $\tau$ là tham số chạy dọc đặc trưng; hai tham số này có vai trò khác nhau.

\textbf{Bước 1 — Viết đúng $F$ và tham số hóa dữ liệu.} Chuyển PDE về một vế, giữ mọi hệ số chứa $x,y,u$. Viết $X(s,0)=x_0(s)$, $Y(s,0)=y_0(s)$, $U(s,0)=h(s)$. Đạo hàm điều kiện theo $s$ cho

\[
p_0x_0'(s)+q_0y_0'(s)=h'(s),\qquad
F(x_0,y_0,h,p_0,q_0)=0.
\]

Giải hai phương trình này tìm tất cả nhánh thực $(p_0,q_0)$. Trên $y=0$, $p_0=h'$; trên $x=0$, $q_0=h'$; trên $y=x$, $p_0+q_0=h'$. Điều kiện dưới căn phải không âm. Các nhánh khác nhau có thể cho các nghiệm khác nhau cùng dữ liệu.

\textbf{Bước 2 — Kiểm tra không đặc trưng.} Tính $J_0=x_0'F_q-y_0'F_p$ ở dữ liệu của từng nhánh. Khi $J_0\ne0$ và các hàm đủ trơn, phép chuyển từ $(s,\tau)$ sang $(x,y)$ khả nghịch cục bộ. Khi $J_0=0$, cần xét riêng; không kết luận ngay vô nghiệm hoặc duy nhất.

\textbf{Bước 3 — Tính các đạo hàm riêng của $F$ rồi lập hệ Charpit.} Khi tính $F_x,F_y,F_u,F_p,F_q$, tạm coi năm đối số độc lập. Hệ là

\[
\begin{aligned}
X_\tau&=F_p,&Y_\tau&=F_q,&U_\tau&=pF_p+qF_q,\\
p_\tau&=-F_x-pF_u,&q_\tau&=-F_y-qF_u.
\end{aligned}
\]

Trong bài cụ thể phải thay các đạo hàm của $F$ vào hệ. Ví dụ $F=pq-1$ cho $X_\tau=q,Y_\tau=p,U_\tau=2,p_\tau=q_\tau=0$.

\textbf{Bước 4 — Tích phân hệ với điều kiện ban đầu.} Tìm phương trình giải được trước: $p'=0$, $q'=0$, hoặc tổ hợp như $(X-p)'=0$. Thế chúng vào phương trình còn lại. Nếu hệ tuyến tính $Z'=MZ$, dùng $Z=e^{\tau M}Z_0$; để tính cụ thể, tìm các tổ hợp tách hệ hoặc quan hệ $M^2, M^3$ giúp rút gọn chuỗi mũ. Ghi các phép tích phân và hằng số được dữ liệu xác định.

\textbf{Bước 5 — Khử tham số hoặc trình bày nghiệm tham số đã giải.} Từ $X(s,\tau),Y(s,\tau)$ tìm $s,\tau$ theo $x,y$ rồi thay vào $U$. Nếu giữ dạng tham số, phải có công thức $X,Y,U$ đã tích phân và miền $\det\partial(X,Y)/\partial(s,\tau)\ne0$; chỉ ghi hệ ODE chưa phải lời giải hoàn tất.

\textbf{Bước 6 — Dùng phép đổi ẩn khi thuận lợi.} Gần dữ liệu $u>0$, $p^2+q^2=u$ đổi $v=2\sqrt u$ cho $|\nabla v|^2=1$; $p^2+q^2=u^2$ đổi $v=\log u$ cũng cho eikonal. Phải tính $v_x,v_y$ để chứng minh phép đổi, đổi cả dữ liệu và giữ điều kiện dấu. Với dữ liệu hằng trên vòng tròn, đạo hàm tiếp tuyến bằng 0 nên gradient pháp tuyến có hai lựa chọn dấu.

\textbf{Bước 7 — Kiểm tra từng nhánh và miền.} Tính $p,q$ của công thức cuối, thế vào PDE và dữ liệu. Phân biệt một nghiệm đã tìm được, các nhánh không đặc trưng cục bộ và toàn bộ nghiệm toàn cục. Giữ điều kiện của nhãn dữ liệu nếu đề chỉ cho $0\le s\le1$; kiểm tra gốc, điểm căn bằng 0 và đường làm mẫu số bằng 0.

\textbf{Chi tiết tích phân và khử sau khi lập hệ.} Trong lúc tích phân theo $\tau$, nhãn $s$ và các gradient ban đầu $p_0(s),q_0(s)$ được giữ cố định. Với $p'=cp$, viết $p=p_0e^{c\tau}$; nếu $X'=dp$ thì $X=x_0+dp_0(e^{c\tau}-1)/c$ khi $c\ne0$, hoặc $X=x_0+dp_0\tau$ khi $c=0$. Với $Y'+cY=R(\tau)$, nhân $e^{c\tau}$ cho $(e^{c\tau}Y)'=e^{c\tau}R$, rồi dùng $Y(0)$ để xác định hằng số. Khử từ hai phương trình $x=X(s,\tau),y=Y(s,\tau)$; ghi rõ tổ hợp nào loại $\tau$ hoặc $s$. Nếu chỉ cho nghiệm tham số, vẫn phải có các tích phân đã tính và điều kiện Jacobian để khử cục bộ.
\subsection*{Lời giải áp dụng cho bài}
\textbf{Bước 1 — Kiểm tra tồn tại thực.} $p^2+q^2=u$, dữ liệu $u(s,0)=s^2+1$ cho $p_0=2s$, nên $q_0^2=1-3s^2$. Vì vế phải âm khi $|s|>1/\sqrt3$, không có nghiệm thực trên lân cận toàn trục dữ liệu.

\textbf{Bước 2 — Đổi ẩn trên khoảng hợp lệ.} Gần dữ liệu $u>0$, đặt $v=2\sqrt u$. Vì $v_x=p/\sqrt u$, $v_y=q/\sqrt u$, có $v_x^2+v_y^2=1$ và $v(s,0)=2\sqrt{s^2+1}$. Đặt
\[a(s)=\frac{2s}{\sqrt{s^2+1}},\qquad b(s)=\frac{\pm\sqrt{1-3s^2}}{\sqrt{s^2+1}},\quad a^2+b^2=1.\]

\textbf{Bước 3 — Đặc trưng eikonal.} Dùng tham số độ dài: $X=s+\tau a$, $Y=\tau b$, $v=2\sqrt{s^2+1}+\tau$. Đổi về $u$:
\[u(X,Y)=(\sqrt{s^2+1}+\tau/2)^2.\]
Trên mỗi dấu, khử tham số nơi Jacobian khác 0; tại dữ liệu Jacobian bằng $b$, nên hợp lệ khi $|s|<1/\sqrt3$. Giữ $\sqrt{s^2+1}+\tau/2>0$ cho phép đổi căn.

\textbf{Bước 4 — Kiểm tra.} Công thức cho $v$ có gradient đơn vị trên miền khả nghịch, nên $|\nabla u|^2=(v^2/4)|\nabla v|^2=u$. Ở $\tau=0$ đúng dữ liệu. Hai đầu $|s|=1/\sqrt3$ có $b=0$, cần xét riêng, không áp kết luận không đặc trưng.

\textbf{Chi tiết xác nhận gradient của nghiệm tham số.} Đặt $h=2\sqrt{s^2+1}$, $a=h'$, $a^2+b^2=1$. Với $X=s+\tau a,Y=\tau b,v=h+\tau$, tính $dv=a\,ds+d\tau$. Mặt khác
\[a\,dX+b\,dY=a\,ds+(a^2+b^2)d\tau+\tau(aa'+bb')ds=a\,ds+d\tau,\]
vì đạo hàm $a^2+b^2=1$ cho $aa'+bb'=0$. Trên miền ánh xạ khả nghịch, so hệ số vi phân được $v_x=a,v_y=b$, nên $|\nabla v|^2=1$. Đổi $u=v^2/4$ cho $|\nabla u|^2=(v^2/4)|\nabla v|^2=u$.

\clearpage\noindent\texttt{\detokenize{BT_1_5_3_14}}\par
\section*{Lời giải bổ sung — Bài tập 14 · Mục 1.5.3}

\subsection*{Kiến thức và trình tự giải}
\textbf{Mục tiêu.} Giải PDE cấp 1 phi tuyến $F(x,y,u,p,q)=0$, với $p=u_x,q=u_y$. Chọn $s$ để gắn nhãn đường dữ liệu, $\tau$ là tham số chạy dọc đặc trưng; hai tham số này có vai trò khác nhau.

\textbf{Bước 1 — Viết đúng $F$ và tham số hóa dữ liệu.} Chuyển PDE về một vế, giữ mọi hệ số chứa $x,y,u$. Viết $X(s,0)=x_0(s)$, $Y(s,0)=y_0(s)$, $U(s,0)=h(s)$. Đạo hàm điều kiện theo $s$ cho

\[
p_0x_0'(s)+q_0y_0'(s)=h'(s),\qquad
F(x_0,y_0,h,p_0,q_0)=0.
\]

Giải hai phương trình này tìm tất cả nhánh thực $(p_0,q_0)$. Trên $y=0$, $p_0=h'$; trên $x=0$, $q_0=h'$; trên $y=x$, $p_0+q_0=h'$. Điều kiện dưới căn phải không âm. Các nhánh khác nhau có thể cho các nghiệm khác nhau cùng dữ liệu.

\textbf{Bước 2 — Kiểm tra không đặc trưng.} Tính $J_0=x_0'F_q-y_0'F_p$ ở dữ liệu của từng nhánh. Khi $J_0\ne0$ và các hàm đủ trơn, phép chuyển từ $(s,\tau)$ sang $(x,y)$ khả nghịch cục bộ. Khi $J_0=0$, cần xét riêng; không kết luận ngay vô nghiệm hoặc duy nhất.

\textbf{Bước 3 — Tính các đạo hàm riêng của $F$ rồi lập hệ Charpit.} Khi tính $F_x,F_y,F_u,F_p,F_q$, tạm coi năm đối số độc lập. Hệ là

\[
\begin{aligned}
X_\tau&=F_p,&Y_\tau&=F_q,&U_\tau&=pF_p+qF_q,\\
p_\tau&=-F_x-pF_u,&q_\tau&=-F_y-qF_u.
\end{aligned}
\]

Trong bài cụ thể phải thay các đạo hàm của $F$ vào hệ. Ví dụ $F=pq-1$ cho $X_\tau=q,Y_\tau=p,U_\tau=2,p_\tau=q_\tau=0$.

\textbf{Bước 4 — Tích phân hệ với điều kiện ban đầu.} Tìm phương trình giải được trước: $p'=0$, $q'=0$, hoặc tổ hợp như $(X-p)'=0$. Thế chúng vào phương trình còn lại. Nếu hệ tuyến tính $Z'=MZ$, dùng $Z=e^{\tau M}Z_0$; để tính cụ thể, tìm các tổ hợp tách hệ hoặc quan hệ $M^2, M^3$ giúp rút gọn chuỗi mũ. Ghi các phép tích phân và hằng số được dữ liệu xác định.

\textbf{Bước 5 — Khử tham số hoặc trình bày nghiệm tham số đã giải.} Từ $X(s,\tau),Y(s,\tau)$ tìm $s,\tau$ theo $x,y$ rồi thay vào $U$. Nếu giữ dạng tham số, phải có công thức $X,Y,U$ đã tích phân và miền $\det\partial(X,Y)/\partial(s,\tau)\ne0$; chỉ ghi hệ ODE chưa phải lời giải hoàn tất.

\textbf{Bước 6 — Dùng phép đổi ẩn khi thuận lợi.} Gần dữ liệu $u>0$, $p^2+q^2=u$ đổi $v=2\sqrt u$ cho $|\nabla v|^2=1$; $p^2+q^2=u^2$ đổi $v=\log u$ cũng cho eikonal. Phải tính $v_x,v_y$ để chứng minh phép đổi, đổi cả dữ liệu và giữ điều kiện dấu. Với dữ liệu hằng trên vòng tròn, đạo hàm tiếp tuyến bằng 0 nên gradient pháp tuyến có hai lựa chọn dấu.

\textbf{Bước 7 — Kiểm tra từng nhánh và miền.} Tính $p,q$ của công thức cuối, thế vào PDE và dữ liệu. Phân biệt một nghiệm đã tìm được, các nhánh không đặc trưng cục bộ và toàn bộ nghiệm toàn cục. Giữ điều kiện của nhãn dữ liệu nếu đề chỉ cho $0\le s\le1$; kiểm tra gốc, điểm căn bằng 0 và đường làm mẫu số bằng 0.

\textbf{Chi tiết tích phân và khử sau khi lập hệ.} Trong lúc tích phân theo $\tau$, nhãn $s$ và các gradient ban đầu $p_0(s),q_0(s)$ được giữ cố định. Với $p'=cp$, viết $p=p_0e^{c\tau}$; nếu $X'=dp$ thì $X=x_0+dp_0(e^{c\tau}-1)/c$ khi $c\ne0$, hoặc $X=x_0+dp_0\tau$ khi $c=0$. Với $Y'+cY=R(\tau)$, nhân $e^{c\tau}$ cho $(e^{c\tau}Y)'=e^{c\tau}R$, rồi dùng $Y(0)$ để xác định hằng số. Khử từ hai phương trình $x=X(s,\tau),y=Y(s,\tau)$; ghi rõ tổ hợp nào loại $\tau$ hoặc $s$. Nếu chỉ cho nghiệm tham số, vẫn phải có các tích phân đã tính và điều kiện Jacobian để khử cục bộ.
\subsection*{Lời giải áp dụng cho bài}
\textbf{Bước 1 — Đổi ẩn.} Dữ liệu $u(x,0)=1$ dương, nên trong lân cận đủ nhỏ đặt $v=\log u$. Khi đó $v_x=p/u$, $v_y=q/u$ và PDE $p^2+q^2=u^2$ thành $v_x^2+v_y^2=1$, dữ liệu $v(x,0)=0$.

\textbf{Bước 2 — Hai hướng pháp tuyến.} Đạo hàm theo $x$ của vết bằng 0, nên $v_x=0$ trên dữ liệu; PDE cho $v_y=\pm1$. Đặc trưng là đường thẳng pháp tuyến trục $x$, gradient giữ hằng, nên $v=\pm y$.

\textbf{Bước 3 — Đổi lại và thử.} $u=e^y$ hoặc $u=e^{-y}$. Với $p=0,q=\pm u$, có $p^2+q^2=u^2$; trên $y=0$ đều bằng 1. Có hai nhánh cục bộ khi chỉ cho giá trị trên dữ liệu; nếu chọn dấu đạo hàm pháp tuyến thì mỗi nhánh được xác định.

\clearpage\noindent\texttt{\detokenize{BT_1_5_3_15}}\par
\section*{Lời giải bổ sung — Bài tập 15 · Mục 1.5.3}

\subsection*{Kiến thức và trình tự giải}
\textbf{Mục tiêu.} Giải PDE cấp 1 phi tuyến $F(x,y,u,p,q)=0$, với $p=u_x,q=u_y$. Chọn $s$ để gắn nhãn đường dữ liệu, $\tau$ là tham số chạy dọc đặc trưng; hai tham số này có vai trò khác nhau.

\textbf{Bước 1 — Viết đúng $F$ và tham số hóa dữ liệu.} Chuyển PDE về một vế, giữ mọi hệ số chứa $x,y,u$. Viết $X(s,0)=x_0(s)$, $Y(s,0)=y_0(s)$, $U(s,0)=h(s)$. Đạo hàm điều kiện theo $s$ cho

\[
p_0x_0'(s)+q_0y_0'(s)=h'(s),\qquad
F(x_0,y_0,h,p_0,q_0)=0.
\]

Giải hai phương trình này tìm tất cả nhánh thực $(p_0,q_0)$. Trên $y=0$, $p_0=h'$; trên $x=0$, $q_0=h'$; trên $y=x$, $p_0+q_0=h'$. Điều kiện dưới căn phải không âm. Các nhánh khác nhau có thể cho các nghiệm khác nhau cùng dữ liệu.

\textbf{Bước 2 — Kiểm tra không đặc trưng.} Tính $J_0=x_0'F_q-y_0'F_p$ ở dữ liệu của từng nhánh. Khi $J_0\ne0$ và các hàm đủ trơn, phép chuyển từ $(s,\tau)$ sang $(x,y)$ khả nghịch cục bộ. Khi $J_0=0$, cần xét riêng; không kết luận ngay vô nghiệm hoặc duy nhất.

\textbf{Bước 3 — Tính các đạo hàm riêng của $F$ rồi lập hệ Charpit.} Khi tính $F_x,F_y,F_u,F_p,F_q$, tạm coi năm đối số độc lập. Hệ là

\[
\begin{aligned}
X_\tau&=F_p,&Y_\tau&=F_q,&U_\tau&=pF_p+qF_q,\\
p_\tau&=-F_x-pF_u,&q_\tau&=-F_y-qF_u.
\end{aligned}
\]

Trong bài cụ thể phải thay các đạo hàm của $F$ vào hệ. Ví dụ $F=pq-1$ cho $X_\tau=q,Y_\tau=p,U_\tau=2,p_\tau=q_\tau=0$.

\textbf{Bước 4 — Tích phân hệ với điều kiện ban đầu.} Tìm phương trình giải được trước: $p'=0$, $q'=0$, hoặc tổ hợp như $(X-p)'=0$. Thế chúng vào phương trình còn lại. Nếu hệ tuyến tính $Z'=MZ$, dùng $Z=e^{\tau M}Z_0$; để tính cụ thể, tìm các tổ hợp tách hệ hoặc quan hệ $M^2, M^3$ giúp rút gọn chuỗi mũ. Ghi các phép tích phân và hằng số được dữ liệu xác định.

\textbf{Bước 5 — Khử tham số hoặc trình bày nghiệm tham số đã giải.} Từ $X(s,\tau),Y(s,\tau)$ tìm $s,\tau$ theo $x,y$ rồi thay vào $U$. Nếu giữ dạng tham số, phải có công thức $X,Y,U$ đã tích phân và miền $\det\partial(X,Y)/\partial(s,\tau)\ne0$; chỉ ghi hệ ODE chưa phải lời giải hoàn tất.

\textbf{Bước 6 — Dùng phép đổi ẩn khi thuận lợi.} Gần dữ liệu $u>0$, $p^2+q^2=u$ đổi $v=2\sqrt u$ cho $|\nabla v|^2=1$; $p^2+q^2=u^2$ đổi $v=\log u$ cũng cho eikonal. Phải tính $v_x,v_y$ để chứng minh phép đổi, đổi cả dữ liệu và giữ điều kiện dấu. Với dữ liệu hằng trên vòng tròn, đạo hàm tiếp tuyến bằng 0 nên gradient pháp tuyến có hai lựa chọn dấu.

\textbf{Bước 7 — Kiểm tra từng nhánh và miền.} Tính $p,q$ của công thức cuối, thế vào PDE và dữ liệu. Phân biệt một nghiệm đã tìm được, các nhánh không đặc trưng cục bộ và toàn bộ nghiệm toàn cục. Giữ điều kiện của nhãn dữ liệu nếu đề chỉ cho $0\le s\le1$; kiểm tra gốc, điểm căn bằng 0 và đường làm mẫu số bằng 0.

\textbf{Chi tiết tích phân và khử sau khi lập hệ.} Trong lúc tích phân theo $\tau$, nhãn $s$ và các gradient ban đầu $p_0(s),q_0(s)$ được giữ cố định. Với $p'=cp$, viết $p=p_0e^{c\tau}$; nếu $X'=dp$ thì $X=x_0+dp_0(e^{c\tau}-1)/c$ khi $c\ne0$, hoặc $X=x_0+dp_0\tau$ khi $c=0$. Với $Y'+cY=R(\tau)$, nhân $e^{c\tau}$ cho $(e^{c\tau}Y)'=e^{c\tau}R$, rồi dùng $Y(0)$ để xác định hằng số. Khử từ hai phương trình $x=X(s,\tau),y=Y(s,\tau)$; ghi rõ tổ hợp nào loại $\tau$ hoặc $s$. Nếu chỉ cho nghiệm tham số, vẫn phải có các tích phân đã tính và điều kiện Jacobian để khử cục bộ.
\subsection*{Lời giải áp dụng cho bài}
\textbf{Bước 1 — Đổi ẩn và dữ liệu.} Đặt $v=\log u$ gần $u=1$. PDE thành $|\nabla v|^2=1$, dữ liệu $v=0$ trên $r=1$.

\textbf{Bước 2 — Chọn gradient ban đầu.} Đạo hàm tiếp tuyến trên vòng tròn bằng 0, nên gradient phải theo pháp tuyến. Độ dài 1 cho hai lựa chọn $\nabla v=\pm(\cos s,\sin s)$. Các đặc trưng là tia pháp tuyến và $v$ thay đổi với tốc độ $\pm1$ theo bán kính, nên $v=\pm(r-1)$.

\textbf{Bước 3 — Nghiệm và kiểm tra.} $u=e^{r-1}$ hoặc $u=e^{1-r}$ trong vành lân cận $r=1$. Vì $|\nabla r|=1$ khi $r>0$, $|\nabla u|^2=(u_r)^2=u^2$; ở $r=1$ bằng 1. Hai nhánh xuyên tâm không kéo dài $C^1$ qua gốc vì $u_r(0)$ khác 0. Dữ liệu không chọn dấu pháp tuyến nên chưa cho duy nhất giữa hai nhánh.

\clearpage\noindent\texttt{\detokenize{BT_1_5_3_16}}\par
\section*{Lời giải bổ sung — Bài tập 16 · Mục 1.5.3}

\subsection*{Kiến thức và trình tự giải}
\textbf{Mục tiêu.} Giải PDE cấp 1 phi tuyến $F(x,y,u,p,q)=0$, với $p=u_x,q=u_y$. Chọn $s$ để gắn nhãn đường dữ liệu, $\tau$ là tham số chạy dọc đặc trưng; hai tham số này có vai trò khác nhau.

\textbf{Bước 1 — Viết đúng $F$ và tham số hóa dữ liệu.} Chuyển PDE về một vế, giữ mọi hệ số chứa $x,y,u$. Viết $X(s,0)=x_0(s)$, $Y(s,0)=y_0(s)$, $U(s,0)=h(s)$. Đạo hàm điều kiện theo $s$ cho

\[
p_0x_0'(s)+q_0y_0'(s)=h'(s),\qquad
F(x_0,y_0,h,p_0,q_0)=0.
\]

Giải hai phương trình này tìm tất cả nhánh thực $(p_0,q_0)$. Trên $y=0$, $p_0=h'$; trên $x=0$, $q_0=h'$; trên $y=x$, $p_0+q_0=h'$. Điều kiện dưới căn phải không âm. Các nhánh khác nhau có thể cho các nghiệm khác nhau cùng dữ liệu.

\textbf{Bước 2 — Kiểm tra không đặc trưng.} Tính $J_0=x_0'F_q-y_0'F_p$ ở dữ liệu của từng nhánh. Khi $J_0\ne0$ và các hàm đủ trơn, phép chuyển từ $(s,\tau)$ sang $(x,y)$ khả nghịch cục bộ. Khi $J_0=0$, cần xét riêng; không kết luận ngay vô nghiệm hoặc duy nhất.

\textbf{Bước 3 — Tính các đạo hàm riêng của $F$ rồi lập hệ Charpit.} Khi tính $F_x,F_y,F_u,F_p,F_q$, tạm coi năm đối số độc lập. Hệ là

\[
\begin{aligned}
X_\tau&=F_p,&Y_\tau&=F_q,&U_\tau&=pF_p+qF_q,\\
p_\tau&=-F_x-pF_u,&q_\tau&=-F_y-qF_u.
\end{aligned}
\]

Trong bài cụ thể phải thay các đạo hàm của $F$ vào hệ. Ví dụ $F=pq-1$ cho $X_\tau=q,Y_\tau=p,U_\tau=2,p_\tau=q_\tau=0$.

\textbf{Bước 4 — Tích phân hệ với điều kiện ban đầu.} Tìm phương trình giải được trước: $p'=0$, $q'=0$, hoặc tổ hợp như $(X-p)'=0$. Thế chúng vào phương trình còn lại. Nếu hệ tuyến tính $Z'=MZ$, dùng $Z=e^{\tau M}Z_0$; để tính cụ thể, tìm các tổ hợp tách hệ hoặc quan hệ $M^2, M^3$ giúp rút gọn chuỗi mũ. Ghi các phép tích phân và hằng số được dữ liệu xác định.

\textbf{Bước 5 — Khử tham số hoặc trình bày nghiệm tham số đã giải.} Từ $X(s,\tau),Y(s,\tau)$ tìm $s,\tau$ theo $x,y$ rồi thay vào $U$. Nếu giữ dạng tham số, phải có công thức $X,Y,U$ đã tích phân và miền $\det\partial(X,Y)/\partial(s,\tau)\ne0$; chỉ ghi hệ ODE chưa phải lời giải hoàn tất.

\textbf{Bước 6 — Dùng phép đổi ẩn khi thuận lợi.} Gần dữ liệu $u>0$, $p^2+q^2=u$ đổi $v=2\sqrt u$ cho $|\nabla v|^2=1$; $p^2+q^2=u^2$ đổi $v=\log u$ cũng cho eikonal. Phải tính $v_x,v_y$ để chứng minh phép đổi, đổi cả dữ liệu và giữ điều kiện dấu. Với dữ liệu hằng trên vòng tròn, đạo hàm tiếp tuyến bằng 0 nên gradient pháp tuyến có hai lựa chọn dấu.

\textbf{Bước 7 — Kiểm tra từng nhánh và miền.} Tính $p,q$ của công thức cuối, thế vào PDE và dữ liệu. Phân biệt một nghiệm đã tìm được, các nhánh không đặc trưng cục bộ và toàn bộ nghiệm toàn cục. Giữ điều kiện của nhãn dữ liệu nếu đề chỉ cho $0\le s\le1$; kiểm tra gốc, điểm căn bằng 0 và đường làm mẫu số bằng 0.

\textbf{Chi tiết tích phân và khử sau khi lập hệ.} Trong lúc tích phân theo $\tau$, nhãn $s$ và các gradient ban đầu $p_0(s),q_0(s)$ được giữ cố định. Với $p'=cp$, viết $p=p_0e^{c\tau}$; nếu $X'=dp$ thì $X=x_0+dp_0(e^{c\tau}-1)/c$ khi $c\ne0$, hoặc $X=x_0+dp_0\tau$ khi $c=0$. Với $Y'+cY=R(\tau)$, nhân $e^{c\tau}$ cho $(e^{c\tau}Y)'=e^{c\tau}R$, rồi dùng $Y(0)$ để xác định hằng số. Khử từ hai phương trình $x=X(s,\tau),y=Y(s,\tau)$; ghi rõ tổ hợp nào loại $\tau$ hoặc $s$. Nếu chỉ cho nghiệm tham số, vẫn phải có các tích phân đã tính và điều kiện Jacobian để khử cục bộ.
\subsection*{Lời giải áp dụng cho bài}
\textbf{Bước 1 — Đổi ẩn.} Với $p^2+q^2=u$, gần dữ liệu $u=1>0$, đặt $v=2\sqrt u$; tính gradient như bài 13 cho $|\nabla v|^2=1$, $v|_{r=1}=2$.

\textbf{Bước 2 — Hai nhánh pháp tuyến.} Gradient tiếp tuyến bằng 0, gradient pháp tuyến là $\pm1$, nên $v=2\pm(r-1)$. Đổi lại:
\[u=(r+1)^2/4,\qquad u=(3-r)^2/4.\]

\textbf{Bước 3 — Kiểm tra và miền.} Với $r>0$, $|\nabla u|^2=(u_r)^2$; $u_r=(r+1)/2$ ở nhánh đầu, $(r-3)/2$ ở nhánh sau, nên đều có $(u_r)^2=u$. Tại $r=1$ bằng 1. Phép đổi căn cho nhánh sau dùng $v=3-r>0$, tức $r<3$ quanh dữ liệu; công thức đại số có thể giải PDE xa hơn sau khi kiểm tra trực tiếp. Cả hai nhánh không có kéo dài xuyên tâm $C^1$ qua gốc.

\textbf{Chi tiết đổi ẩn và tia pháp tuyến.} $v=2\sqrt u$ cho $v_x=u_x/\sqrt u,v_y=u_y/\sqrt u$, nên $|\nabla v|^2=(p^2+q^2)/u=1$. Lấy điểm đầu $(\cos s,\sin s)$, gradient $\varepsilon(\cos s,\sin s)$, tham số độ dài $\tau$. Hệ cho $(X,Y)=(1+\varepsilon\tau)(\cos s,\sin s)$, $v=2+\tau$. Gần vòng tròn, $r=1+\varepsilon\tau$, tức $\tau=\varepsilon(r-1)$, nên $v=2+\varepsilon(r-1)$. Bình phương rồi chia $4$ cho hai công thức của $u$.

\clearpage\noindent\texttt{\detokenize{BT_1_5_3_17}}\par
\section*{Lời giải bổ sung — Bài tập 17 · Mục 1.5.3}

\subsection*{Kiến thức và trình tự giải}
\textbf{Mục tiêu.} Giải PDE cấp 1 phi tuyến $F(x,y,u,p,q)=0$, với $p=u_x,q=u_y$. Chọn $s$ để gắn nhãn đường dữ liệu, $\tau$ là tham số chạy dọc đặc trưng; hai tham số này có vai trò khác nhau.

\textbf{Bước 1 — Viết đúng $F$ và tham số hóa dữ liệu.} Chuyển PDE về một vế, giữ mọi hệ số chứa $x,y,u$. Viết $X(s,0)=x_0(s)$, $Y(s,0)=y_0(s)$, $U(s,0)=h(s)$. Đạo hàm điều kiện theo $s$ cho

\[
p_0x_0'(s)+q_0y_0'(s)=h'(s),\qquad
F(x_0,y_0,h,p_0,q_0)=0.
\]

Giải hai phương trình này tìm tất cả nhánh thực $(p_0,q_0)$. Trên $y=0$, $p_0=h'$; trên $x=0$, $q_0=h'$; trên $y=x$, $p_0+q_0=h'$. Điều kiện dưới căn phải không âm. Các nhánh khác nhau có thể cho các nghiệm khác nhau cùng dữ liệu.

\textbf{Bước 2 — Kiểm tra không đặc trưng.} Tính $J_0=x_0'F_q-y_0'F_p$ ở dữ liệu của từng nhánh. Khi $J_0\ne0$ và các hàm đủ trơn, phép chuyển từ $(s,\tau)$ sang $(x,y)$ khả nghịch cục bộ. Khi $J_0=0$, cần xét riêng; không kết luận ngay vô nghiệm hoặc duy nhất.

\textbf{Bước 3 — Tính các đạo hàm riêng của $F$ rồi lập hệ Charpit.} Khi tính $F_x,F_y,F_u,F_p,F_q$, tạm coi năm đối số độc lập. Hệ là

\[
\begin{aligned}
X_\tau&=F_p,&Y_\tau&=F_q,&U_\tau&=pF_p+qF_q,\\
p_\tau&=-F_x-pF_u,&q_\tau&=-F_y-qF_u.
\end{aligned}
\]

Trong bài cụ thể phải thay các đạo hàm của $F$ vào hệ. Ví dụ $F=pq-1$ cho $X_\tau=q,Y_\tau=p,U_\tau=2,p_\tau=q_\tau=0$.

\textbf{Bước 4 — Tích phân hệ với điều kiện ban đầu.} Tìm phương trình giải được trước: $p'=0$, $q'=0$, hoặc tổ hợp như $(X-p)'=0$. Thế chúng vào phương trình còn lại. Nếu hệ tuyến tính $Z'=MZ$, dùng $Z=e^{\tau M}Z_0$; để tính cụ thể, tìm các tổ hợp tách hệ hoặc quan hệ $M^2, M^3$ giúp rút gọn chuỗi mũ. Ghi các phép tích phân và hằng số được dữ liệu xác định.

\textbf{Bước 5 — Khử tham số hoặc trình bày nghiệm tham số đã giải.} Từ $X(s,\tau),Y(s,\tau)$ tìm $s,\tau$ theo $x,y$ rồi thay vào $U$. Nếu giữ dạng tham số, phải có công thức $X,Y,U$ đã tích phân và miền $\det\partial(X,Y)/\partial(s,\tau)\ne0$; chỉ ghi hệ ODE chưa phải lời giải hoàn tất.

\textbf{Bước 6 — Dùng phép đổi ẩn khi thuận lợi.} Gần dữ liệu $u>0$, $p^2+q^2=u$ đổi $v=2\sqrt u$ cho $|\nabla v|^2=1$; $p^2+q^2=u^2$ đổi $v=\log u$ cũng cho eikonal. Phải tính $v_x,v_y$ để chứng minh phép đổi, đổi cả dữ liệu và giữ điều kiện dấu. Với dữ liệu hằng trên vòng tròn, đạo hàm tiếp tuyến bằng 0 nên gradient pháp tuyến có hai lựa chọn dấu.

\textbf{Bước 7 — Kiểm tra từng nhánh và miền.} Tính $p,q$ của công thức cuối, thế vào PDE và dữ liệu. Phân biệt một nghiệm đã tìm được, các nhánh không đặc trưng cục bộ và toàn bộ nghiệm toàn cục. Giữ điều kiện của nhãn dữ liệu nếu đề chỉ cho $0\le s\le1$; kiểm tra gốc, điểm căn bằng 0 và đường làm mẫu số bằng 0.

\textbf{Chi tiết tích phân và khử sau khi lập hệ.} Trong lúc tích phân theo $\tau$, nhãn $s$ và các gradient ban đầu $p_0(s),q_0(s)$ được giữ cố định. Với $p'=cp$, viết $p=p_0e^{c\tau}$; nếu $X'=dp$ thì $X=x_0+dp_0(e^{c\tau}-1)/c$ khi $c\ne0$, hoặc $X=x_0+dp_0\tau$ khi $c=0$. Với $Y'+cY=R(\tau)$, nhân $e^{c\tau}$ cho $(e^{c\tau}Y)'=e^{c\tau}R$, rồi dùng $Y(0)$ để xác định hằng số. Khử từ hai phương trình $x=X(s,\tau),y=Y(s,\tau)$; ghi rõ tổ hợp nào loại $\tau$ hoặc $s$. Nếu chỉ cho nghiệm tham số, vẫn phải có các tích phân đã tính và điều kiện Jacobian để khử cục bộ.
\subsection*{Lời giải áp dụng cho bài}
\textbf{Bước 1 — Điều kiện thực trên dữ liệu.} $F=q^2-p$, $u(s,0)=f(s)$ cho $p_0=f'(s)$ và $q_0^2=f'(s)$. Vì vậy cần $f'\ge0$. Nếu $f'<0$ tại một điểm thì vô nghiệm thực quanh điểm đó.

\textbf{Bước 2 — Hai nhánh không đặc trưng.} Khi $f'>0$ và $f$ đủ trơn, đặt $Q=\pm\sqrt{f'}$. $F_p=-1,F_q=2q$, $p'=q'=0$, $U'=-p+2q^2=p$. Tích phân:
\[X=s-\tau,\quad Y=2Q(s)\tau,\quad U=f(s)+f'(s)\tau.\]

\textbf{Bước 3 — Khử và miền.} Giải $x=s-\tau$, $y=2Q(s)\tau$ tìm $s,\tau$, rồi $u=U$. Jacobian là $2Q+2\tau Q'$, tại dữ liệu bằng $2Q\ne0$, nên có nghiệm cục bộ trên từng dấu. Không cần một công thức tường minh theo $x,y$ cho $f$ tùy ý.

\textbf{Bước 4 — Suy biến.} Nếu $f'=0$, dữ liệu đặc trưng và hai dấu trùng nhau. Dữ liệu hằng có nghiệm hằng, nhưng không suy duy nhất từ công thức ở bước 3. Trên miền khả nghịch, Charpit cho gradient $p=f'(s),q=Q(s)$, nên $q^2=p$ và tại $\tau=0$ đúng $f$.

\textbf{Chi tiết phương trình khử và gradient.} Thế $\tau=s-x$ vào $y=2Q(s)\tau$ được một phương trình một ẩn $y=2Q(s)(s-x)$. Đạo hàm theo $s$ là $2Q+2(s-x)Q'$, khác $0$ tại dữ liệu khi $f'>0$, nên giải cục bộ được $s=s(x,y)$. Đặt $P=f'=Q^2$; từ $X=s-\tau,Y=2Q\tau,U=f+P\tau$, tính
\[dU=(P+\tau P')ds+P\,d\tau=P\,dX+Q\,dY,\]
vì $P'=2QQ'$. So hệ số vi phân cho $u_x=P,u_y=Q$, xác nhận $u_y^2=u_x$.

\clearpage\noindent\texttt{\detokenize{BT_1_5_3_18}}\par
\section*{Lời giải bổ sung — Bài tập 18 · Mục 1.5.3}

\subsection*{Kiến thức và trình tự giải}
\textbf{Mục tiêu.} Giải PDE cấp 1 phi tuyến $F(x,y,u,p,q)=0$, với $p=u_x,q=u_y$. Chọn $s$ để gắn nhãn đường dữ liệu, $\tau$ là tham số chạy dọc đặc trưng; hai tham số này có vai trò khác nhau.

\textbf{Bước 1 — Viết đúng $F$ và tham số hóa dữ liệu.} Chuyển PDE về một vế, giữ mọi hệ số chứa $x,y,u$. Viết $X(s,0)=x_0(s)$, $Y(s,0)=y_0(s)$, $U(s,0)=h(s)$. Đạo hàm điều kiện theo $s$ cho

\[
p_0x_0'(s)+q_0y_0'(s)=h'(s),\qquad
F(x_0,y_0,h,p_0,q_0)=0.
\]

Giải hai phương trình này tìm tất cả nhánh thực $(p_0,q_0)$. Trên $y=0$, $p_0=h'$; trên $x=0$, $q_0=h'$; trên $y=x$, $p_0+q_0=h'$. Điều kiện dưới căn phải không âm. Các nhánh khác nhau có thể cho các nghiệm khác nhau cùng dữ liệu.

\textbf{Bước 2 — Kiểm tra không đặc trưng.} Tính $J_0=x_0'F_q-y_0'F_p$ ở dữ liệu của từng nhánh. Khi $J_0\ne0$ và các hàm đủ trơn, phép chuyển từ $(s,\tau)$ sang $(x,y)$ khả nghịch cục bộ. Khi $J_0=0$, cần xét riêng; không kết luận ngay vô nghiệm hoặc duy nhất.

\textbf{Bước 3 — Tính các đạo hàm riêng của $F$ rồi lập hệ Charpit.} Khi tính $F_x,F_y,F_u,F_p,F_q$, tạm coi năm đối số độc lập. Hệ là

\[
\begin{aligned}
X_\tau&=F_p,&Y_\tau&=F_q,&U_\tau&=pF_p+qF_q,\\
p_\tau&=-F_x-pF_u,&q_\tau&=-F_y-qF_u.
\end{aligned}
\]

Trong bài cụ thể phải thay các đạo hàm của $F$ vào hệ. Ví dụ $F=pq-1$ cho $X_\tau=q,Y_\tau=p,U_\tau=2,p_\tau=q_\tau=0$.

\textbf{Bước 4 — Tích phân hệ với điều kiện ban đầu.} Tìm phương trình giải được trước: $p'=0$, $q'=0$, hoặc tổ hợp như $(X-p)'=0$. Thế chúng vào phương trình còn lại. Nếu hệ tuyến tính $Z'=MZ$, dùng $Z=e^{\tau M}Z_0$; để tính cụ thể, tìm các tổ hợp tách hệ hoặc quan hệ $M^2, M^3$ giúp rút gọn chuỗi mũ. Ghi các phép tích phân và hằng số được dữ liệu xác định.

\textbf{Bước 5 — Khử tham số hoặc trình bày nghiệm tham số đã giải.} Từ $X(s,\tau),Y(s,\tau)$ tìm $s,\tau$ theo $x,y$ rồi thay vào $U$. Nếu giữ dạng tham số, phải có công thức $X,Y,U$ đã tích phân và miền $\det\partial(X,Y)/\partial(s,\tau)\ne0$; chỉ ghi hệ ODE chưa phải lời giải hoàn tất.

\textbf{Bước 6 — Dùng phép đổi ẩn khi thuận lợi.} Gần dữ liệu $u>0$, $p^2+q^2=u$ đổi $v=2\sqrt u$ cho $|\nabla v|^2=1$; $p^2+q^2=u^2$ đổi $v=\log u$ cũng cho eikonal. Phải tính $v_x,v_y$ để chứng minh phép đổi, đổi cả dữ liệu và giữ điều kiện dấu. Với dữ liệu hằng trên vòng tròn, đạo hàm tiếp tuyến bằng 0 nên gradient pháp tuyến có hai lựa chọn dấu.

\textbf{Bước 7 — Kiểm tra từng nhánh và miền.} Tính $p,q$ của công thức cuối, thế vào PDE và dữ liệu. Phân biệt một nghiệm đã tìm được, các nhánh không đặc trưng cục bộ và toàn bộ nghiệm toàn cục. Giữ điều kiện của nhãn dữ liệu nếu đề chỉ cho $0\le s\le1$; kiểm tra gốc, điểm căn bằng 0 và đường làm mẫu số bằng 0.

\textbf{Chi tiết tích phân và khử sau khi lập hệ.} Trong lúc tích phân theo $\tau$, nhãn $s$ và các gradient ban đầu $p_0(s),q_0(s)$ được giữ cố định. Với $p'=cp$, viết $p=p_0e^{c\tau}$; nếu $X'=dp$ thì $X=x_0+dp_0(e^{c\tau}-1)/c$ khi $c\ne0$, hoặc $X=x_0+dp_0\tau$ khi $c=0$. Với $Y'+cY=R(\tau)$, nhân $e^{c\tau}$ cho $(e^{c\tau}Y)'=e^{c\tau}R$, rồi dùng $Y(0)$ để xác định hằng số. Khử từ hai phương trình $x=X(s,\tau),y=Y(s,\tau)$; ghi rõ tổ hợp nào loại $\tau$ hoặc $s$. Nếu chỉ cho nghiệm tham số, vẫn phải có các tích phân đã tính và điều kiện Jacobian để khử cục bộ.
\subsection*{Lời giải áp dụng cho bài}
\textbf{Bước 1 — Tương thích.} $F=p^2-3q^2-u$, $u(s,0)=s^2$ cho $p_0=2s$. PDE cho $4s^2-3q_0^2=s^2$, nên $q_0=\pm s$.

\textbf{Bước 2 — Tích phân Charpit.} $p'=p,q'=q,X'=2p,Y'=-6q,U'=2U$. Đặt $E=e^\tau$:
\[X=s+4s(E-1),\quad Y=-6q_0(E-1),\quad U=s^2E^2.\]

\textbf{Bước 3 — Khử.} Với $q_0=s$, $X+Y/2=sE$; với $q_0=-s$, $X-Y/2=sE$. Vì thế $u=(x+y/2)^2$ hoặc $u=(x-y/2)^2$.

\textbf{Bước 4 — Kiểm tra.} Đặt $L=x\pm y/2$. $p=2L,q=\pm L$, nên $p^2-3q^2=L^2=u$, và $y=0$ cho $x^2$. $J_0=-6q_0=\mp6s$ khác 0 khi $s\ne0$; gốc là điểm suy biến.

\clearpage\noindent\texttt{\detokenize{VD_1_6_2_01}}\par
\section*{Lời giải bổ sung — Ví dụ 1 · Mục 1.6.2}

\subsection*{Kiến thức và trình tự giải}
\textbf{Mục tiêu.} Với dữ liệu gián đoạn hoặc khi đặc trưng giao nhau, tìm nghiệm yếu thỏa entropy của luật bảo toàn $u_t+(F(u))_x=0$.

\textbf{Bước 1 — Xác định thông lượng và đặc trưng.} Nếu PDE là $u_t+a(u)u_x=0$, tích phân $F'=a$. Trong vùng trơn, $u=u_0(\xi)$ và $x=\xi+F'(u_0(\xi))t$. Tính riêng tốc độ cho mỗi đoạn dữ liệu. Burgers dùng $F=u^2/2$; $u_t+u^2u_x=0$ dùng $F=u^3/3$.

\textbf{Bước 2 — Nhận diện sóng tại mỗi bước nhảy.} Gọi $u_L,u_R$ là giá trị trái, phải. So sánh $F'(u_L)$ và $F'(u_R)$. Với $F$ lồi trên khoảng trạng thái, $u_L>u_R$ tạo sốc, $u_L<u_R$ tạo quạt giãn; với $F$ lõm, thứ tự trạng thái đảo lại. Nếu độ lồi đổi dấu, phải xét thêm cấu trúc entropy, không dùng quy tắc lồi trên toàn khoảng.

\textbf{Bước 3 — Tính đường sốc từ bảo toàn.} Với $x=s(t)$,

\[
s'(t)=\frac{F(u_L)-F(u_R)}{u_L-u_R}.
\]

Nếu hai trạng thái hằng, tốc độ $\sigma$ hằng, nên $s=x_0+\sigma t$. Kiểm tra entropy bằng $F'(u_R)<\sigma<F'(u_L)$ đối với sốc nén có hai trạng thái phân biệt trong các bài này. Giá trị ngay trên đường sốc không ảnh hưởng nghiệm yếu.

\textbf{Bước 4 — Tính quạt giãn.} Từ $(x_0,0)$, đặt $z=(x-x_0)/t$, $u=V(z)$. Thế PDE cho $(F'(V)-z)V'=0$. Trong nhánh biến thiên, $F'(V)=z$, nên giải phương trình này tìm $V$. Mép quạt là $x=x_0+F'(u_L)t$ và $x=x_0+F'(u_R)t$, theo thứ tự tốc độ tăng; ghép liên tục với hai trạng thái ngoài. Chọn căn đúng dấu theo khoảng giá trị của dữ liệu.

\textbf{Bước 5 — Tìm mọi thời điểm tương tác.} Giải phương trình giao của hai sốc, sốc với mép quạt, hoặc quạt với quạt. Trước thời điểm gặp, ghép các sóng riêng. Sau khi gặp, tính lại hai trạng thái sát đường sốc. Nếu một phía nằm trong quạt, trạng thái đó phụ thuộc $s(t)$ và $t$, nên Rankine–Hugoniot trở thành ODE mới. Dùng $s(t_*)=x_*$ để tìm hằng số, rồi tiếp tục tìm lần tương tác sau.

\textbf{Bước 6 — Viết nghiệm theo từng khoảng thời gian.} Với mỗi khoảng, ghi rõ các miền $x$, thứ tự đường phân cách và công thức của $u$. Trên nửa trục, kiểm tra chiều truyền từ biên vào miền và thế điều kiện biên. Với dữ liệu trơn từng đoạn, tính $\partial x/\partial\xi$ để tìm thời điểm đầu tiên mất nghiệm cổ điển.

\textbf{Bước 7 — Kiểm tra.} Trong mỗi vùng trơn, thay vào PDE. Trên mỗi sốc, kiểm tra Rankine–Hugoniot và entropy; ở mép quạt, kiểm tra liên tục. Đối chiếu dữ liệu khi $t\downarrow0$, sự liên tục của vị trí sóng qua thời điểm gặp và khối lượng nếu dữ liệu có giá hữu hạn. Nhân PDE với một hàm của $u$ chỉ giữ tương đương cho nghiệm trơn; không tự áp luật đã nhân để tính cùng đường sốc.

\textbf{Chi tiết tích phân đường sốc nằm trong quạt.} Nếu $s'=s/(2t)$ và $s(t_*)=x_*>0$, chia cho $s$ rồi tích phân từ $t_*$ đến $t$: $\log(s/x_*)=\frac12\log(t/t_*)$, nên $s=x_*\sqrt{t/t_*}$. Nếu quạt có tâm $x_0$, thay $s$ bằng $s-x_0$ trước khi tách biến. Khi trạng thái sát sốc thay đổi tại lần gặp tiếp theo, lập ODE mới và lấy vị trí tại lần gặp làm dữ liệu; không giữ hằng số của giai đoạn trước.
\subsection*{Lời giải áp dụng cho bài}
\textbf{Bước 1 — Đặc trưng và vùng giao.} $F=u^2/2$, dữ liệu 1 khi $\xi<0$, 0 khi $\xi>0$. Vì $U'=0$, $X=\xi+t$ từ trái và $X=\xi$ từ phải. Vùng $x<0$ chỉ có đặc trưng trái, $x>t$ chỉ có đặc trưng phải; vùng $0<x<t$ hai họ giao nhau.

\textbf{Bước 2 — Tìm sốc.} Nối $u_L=1,u_R=0$, Rankine–Hugoniot cho $s'=(1/2-0)/(1-0)=1/2$. Với $s(0)=0$, $s=t/2$.

\textbf{Bước 3 — Nghiệm entropy.} $u=1$ khi $x<t/2$, $u=0$ khi $x>t/2$. Sốc nằm trong vùng giao $0<x<t$ và thỏa $F'(0)=0<s'<F'(1)=1$.

\textbf{Bước 4 — Đọc đồ thị.} Trong mặt phẳng $(x,t)$, vẽ họ $x=\xi+t$ ở trái, họ đứng $x=\xi$ ở phải và đường $x=t/2$ phân chia; đặc trưng đi vào sốc từ cả hai phía. Đồ thị theo $x$ ở mỗi thời điểm là bước nhảy 1 xuống 0 tại $t/2$. Giá trị trên chính sốc tùy chọn không đổi nghiệm yếu.

\clearpage\noindent\texttt{\detokenize{VD_1_6_2_02}}\par
\section*{Lời giải bổ sung — Ví dụ 2 · Mục 1.6.2}

\subsection*{Kiến thức và trình tự giải}
\textbf{Mục tiêu.} Với dữ liệu gián đoạn hoặc khi đặc trưng giao nhau, tìm nghiệm yếu thỏa entropy của luật bảo toàn $u_t+(F(u))_x=0$.

\textbf{Bước 1 — Xác định thông lượng và đặc trưng.} Nếu PDE là $u_t+a(u)u_x=0$, tích phân $F'=a$. Trong vùng trơn, $u=u_0(\xi)$ và $x=\xi+F'(u_0(\xi))t$. Tính riêng tốc độ cho mỗi đoạn dữ liệu. Burgers dùng $F=u^2/2$; $u_t+u^2u_x=0$ dùng $F=u^3/3$.

\textbf{Bước 2 — Nhận diện sóng tại mỗi bước nhảy.} Gọi $u_L,u_R$ là giá trị trái, phải. So sánh $F'(u_L)$ và $F'(u_R)$. Với $F$ lồi trên khoảng trạng thái, $u_L>u_R$ tạo sốc, $u_L<u_R$ tạo quạt giãn; với $F$ lõm, thứ tự trạng thái đảo lại. Nếu độ lồi đổi dấu, phải xét thêm cấu trúc entropy, không dùng quy tắc lồi trên toàn khoảng.

\textbf{Bước 3 — Tính đường sốc từ bảo toàn.} Với $x=s(t)$,

\[
s'(t)=\frac{F(u_L)-F(u_R)}{u_L-u_R}.
\]

Nếu hai trạng thái hằng, tốc độ $\sigma$ hằng, nên $s=x_0+\sigma t$. Kiểm tra entropy bằng $F'(u_R)<\sigma<F'(u_L)$ đối với sốc nén có hai trạng thái phân biệt trong các bài này. Giá trị ngay trên đường sốc không ảnh hưởng nghiệm yếu.

\textbf{Bước 4 — Tính quạt giãn.} Từ $(x_0,0)$, đặt $z=(x-x_0)/t$, $u=V(z)$. Thế PDE cho $(F'(V)-z)V'=0$. Trong nhánh biến thiên, $F'(V)=z$, nên giải phương trình này tìm $V$. Mép quạt là $x=x_0+F'(u_L)t$ và $x=x_0+F'(u_R)t$, theo thứ tự tốc độ tăng; ghép liên tục với hai trạng thái ngoài. Chọn căn đúng dấu theo khoảng giá trị của dữ liệu.

\textbf{Bước 5 — Tìm mọi thời điểm tương tác.} Giải phương trình giao của hai sốc, sốc với mép quạt, hoặc quạt với quạt. Trước thời điểm gặp, ghép các sóng riêng. Sau khi gặp, tính lại hai trạng thái sát đường sốc. Nếu một phía nằm trong quạt, trạng thái đó phụ thuộc $s(t)$ và $t$, nên Rankine–Hugoniot trở thành ODE mới. Dùng $s(t_*)=x_*$ để tìm hằng số, rồi tiếp tục tìm lần tương tác sau.

\textbf{Bước 6 — Viết nghiệm theo từng khoảng thời gian.} Với mỗi khoảng, ghi rõ các miền $x$, thứ tự đường phân cách và công thức của $u$. Trên nửa trục, kiểm tra chiều truyền từ biên vào miền và thế điều kiện biên. Với dữ liệu trơn từng đoạn, tính $\partial x/\partial\xi$ để tìm thời điểm đầu tiên mất nghiệm cổ điển.

\textbf{Bước 7 — Kiểm tra.} Trong mỗi vùng trơn, thay vào PDE. Trên mỗi sốc, kiểm tra Rankine–Hugoniot và entropy; ở mép quạt, kiểm tra liên tục. Đối chiếu dữ liệu khi $t\downarrow0$, sự liên tục của vị trí sóng qua thời điểm gặp và khối lượng nếu dữ liệu có giá hữu hạn. Nhân PDE với một hàm của $u$ chỉ giữ tương đương cho nghiệm trơn; không tự áp luật đã nhân để tính cùng đường sốc.

\textbf{Chi tiết tích phân đường sốc nằm trong quạt.} Nếu $s'=s/(2t)$ và $s(t_*)=x_*>0$, chia cho $s$ rồi tích phân từ $t_*$ đến $t$: $\log(s/x_*)=\frac12\log(t/t_*)$, nên $s=x_*\sqrt{t/t_*}$. Nếu quạt có tâm $x_0$, thay $s$ bằng $s-x_0$ trước khi tách biến. Khi trạng thái sát sốc thay đổi tại lần gặp tiếp theo, lập ODE mới và lấy vị trí tại lần gặp làm dữ liệu; không giữ hằng số của giai đoạn trước.
\subsection*{Lời giải áp dụng cho bài}
\textbf{Bước 1 — Đặc trưng tách xa.} Dữ liệu Burgers 0 khi $\xi<0$, 1 khi $\xi>0$ cho $X=\xi$ từ trái, $X=\xi+t$ từ phải. Vùng $0<x<t$ không có đặc trưng từ một điểm dữ liệu trơn.

\textbf{Bước 2 — Bổ sung quạt.} Tìm $u=V(z)$, $z=x/t$. Tính $u_t=-zV'/t$, $u_x=V'/t$, nên $(V-z)V'=0$. Nhánh nối liên tục 0 với 1 có $V=z$.

\textbf{Bước 3 — Công thức và kiểm tra.}
\[u=\begin{cases}0,&x\le0,\\x/t,&0<x<t,\\1,&x\ge t,\end{cases}\quad t>0.\]
PDE đúng trong quạt và hai vùng hằng; giá trị ghép tại mép là 0,1, nên không có bước nhảy cần tính sốc.

\textbf{Bước 4 — Đọc đồ thị.} Vẽ các tia $x=kt$, $0\le k\le1$ từ gốc để lấp vùng giữa hai họ ngoài. Trên mỗi tia $u=k$. Ở thời điểm cố định, đồ thị có đoạn thẳng từ $(0,0)$ đến $(t,1)$ giữa hai đoạn hằng. Giới hạn tại $t=0$ khớp dữ liệu ở $x\ne0$.

\clearpage\noindent\texttt{\detokenize{VD_1_6_2_03}}\par
\section*{Lời giải bổ sung — Ví dụ 3 · Mục 1.6.2}

\subsection*{Kiến thức và trình tự giải}
\textbf{Mục tiêu.} Với dữ liệu gián đoạn hoặc khi đặc trưng giao nhau, tìm nghiệm yếu thỏa entropy của luật bảo toàn $u_t+(F(u))_x=0$.

\textbf{Bước 1 — Xác định thông lượng và đặc trưng.} Nếu PDE là $u_t+a(u)u_x=0$, tích phân $F'=a$. Trong vùng trơn, $u=u_0(\xi)$ và $x=\xi+F'(u_0(\xi))t$. Tính riêng tốc độ cho mỗi đoạn dữ liệu. Burgers dùng $F=u^2/2$; $u_t+u^2u_x=0$ dùng $F=u^3/3$.

\textbf{Bước 2 — Nhận diện sóng tại mỗi bước nhảy.} Gọi $u_L,u_R$ là giá trị trái, phải. So sánh $F'(u_L)$ và $F'(u_R)$. Với $F$ lồi trên khoảng trạng thái, $u_L>u_R$ tạo sốc, $u_L<u_R$ tạo quạt giãn; với $F$ lõm, thứ tự trạng thái đảo lại. Nếu độ lồi đổi dấu, phải xét thêm cấu trúc entropy, không dùng quy tắc lồi trên toàn khoảng.

\textbf{Bước 3 — Tính đường sốc từ bảo toàn.} Với $x=s(t)$,

\[
s'(t)=\frac{F(u_L)-F(u_R)}{u_L-u_R}.
\]

Nếu hai trạng thái hằng, tốc độ $\sigma$ hằng, nên $s=x_0+\sigma t$. Kiểm tra entropy bằng $F'(u_R)<\sigma<F'(u_L)$ đối với sốc nén có hai trạng thái phân biệt trong các bài này. Giá trị ngay trên đường sốc không ảnh hưởng nghiệm yếu.

\textbf{Bước 4 — Tính quạt giãn.} Từ $(x_0,0)$, đặt $z=(x-x_0)/t$, $u=V(z)$. Thế PDE cho $(F'(V)-z)V'=0$. Trong nhánh biến thiên, $F'(V)=z$, nên giải phương trình này tìm $V$. Mép quạt là $x=x_0+F'(u_L)t$ và $x=x_0+F'(u_R)t$, theo thứ tự tốc độ tăng; ghép liên tục với hai trạng thái ngoài. Chọn căn đúng dấu theo khoảng giá trị của dữ liệu.

\textbf{Bước 5 — Tìm mọi thời điểm tương tác.} Giải phương trình giao của hai sốc, sốc với mép quạt, hoặc quạt với quạt. Trước thời điểm gặp, ghép các sóng riêng. Sau khi gặp, tính lại hai trạng thái sát đường sốc. Nếu một phía nằm trong quạt, trạng thái đó phụ thuộc $s(t)$ và $t$, nên Rankine–Hugoniot trở thành ODE mới. Dùng $s(t_*)=x_*$ để tìm hằng số, rồi tiếp tục tìm lần tương tác sau.

\textbf{Bước 6 — Viết nghiệm theo từng khoảng thời gian.} Với mỗi khoảng, ghi rõ các miền $x$, thứ tự đường phân cách và công thức của $u$. Trên nửa trục, kiểm tra chiều truyền từ biên vào miền và thế điều kiện biên. Với dữ liệu trơn từng đoạn, tính $\partial x/\partial\xi$ để tìm thời điểm đầu tiên mất nghiệm cổ điển.

\textbf{Bước 7 — Kiểm tra.} Trong mỗi vùng trơn, thay vào PDE. Trên mỗi sốc, kiểm tra Rankine–Hugoniot và entropy; ở mép quạt, kiểm tra liên tục. Đối chiếu dữ liệu khi $t\downarrow0$, sự liên tục của vị trí sóng qua thời điểm gặp và khối lượng nếu dữ liệu có giá hữu hạn. Nhân PDE với một hàm của $u$ chỉ giữ tương đương cho nghiệm trơn; không tự áp luật đã nhân để tính cùng đường sốc.

\textbf{Chi tiết tích phân đường sốc nằm trong quạt.} Nếu $s'=s/(2t)$ và $s(t_*)=x_*>0$, chia cho $s$ rồi tích phân từ $t_*$ đến $t$: $\log(s/x_*)=\frac12\log(t/t_*)$, nên $s=x_*\sqrt{t/t_*}$. Nếu quạt có tâm $x_0$, thay $s$ bằng $s-x_0$ trước khi tách biến. Khi trạng thái sát sốc thay đổi tại lần gặp tiếp theo, lập ODE mới và lấy vị trí tại lần gặp làm dữ liệu; không giữ hằng số của giai đoạn trước.
\subsection*{Lời giải áp dụng cho bài}
\textbf{Bước 1 — Sóng ban đầu.} Burgers có xung 1 trên $(0,1)$, 0 ở ngoài. Bước $0\to1$ tại 0 tạo quạt $u=x/t$ trên $0<x<t$; bước $1\to0$ tại 1 tạo sốc $s_1=1+t/2$.

\textbf{Bước 2 — Thời điểm sốc gặp quạt.} Giải $1+t/2=t$, được $t_*=2,x_*=2$. Trước đó,
\[u=\begin{cases}0,&x\le0,\\x/t,&0<x<t,\\1,&t\le x<1+t/2,\\0,&x>1+t/2.\end{cases}\quad0<t<2.\]

\textbf{Bước 3 — Sốc trong quạt.} Sau $t=2$, phía trái sốc là $u_L=s/t$, phải là 0. Rankine–Hugoniot cho $s'=s/(2t)$. Tích phân $d\log s=dt/(2t)$ cho $s=C\sqrt t$; dùng $s(2)=2$, suy $s=\sqrt{2t}$. Với $t\ge2$, $u=x/t$ trên $0<x<s$, bằng 0 ở ngoài.

\textbf{Bước 4 — Kiểm tra và hình học.} Sốc ghép liên tục ở 2, tốc độ sau gặp bằng nửa trạng thái trái nên thỏa entropy. Khối lượng sau gặp $\int_0^s(x/t)dx=s^2/(2t)=1$, đúng xung ban đầu. Đồ thị có quạt thẳng từ gốc và bước nhảy xuống 0 trên đường cong $x=\sqrt{2t}$; không kéo dài đường thẳng $1+t/2$ sau tương tác.

\clearpage\noindent\texttt{\detokenize{BT_1_6_3_01}}\par
\section*{Lời giải bổ sung — Bài tập 1 · Mục 1.6.3}

\subsection*{Kiến thức và trình tự giải}
\textbf{Mục tiêu.} Với dữ liệu gián đoạn hoặc khi đặc trưng giao nhau, tìm nghiệm yếu thỏa entropy của luật bảo toàn $u_t+(F(u))_x=0$.

\textbf{Bước 1 — Xác định thông lượng và đặc trưng.} Nếu PDE là $u_t+a(u)u_x=0$, tích phân $F'=a$. Trong vùng trơn, $u=u_0(\xi)$ và $x=\xi+F'(u_0(\xi))t$. Tính riêng tốc độ cho mỗi đoạn dữ liệu. Burgers dùng $F=u^2/2$; $u_t+u^2u_x=0$ dùng $F=u^3/3$.

\textbf{Bước 2 — Nhận diện sóng tại mỗi bước nhảy.} Gọi $u_L,u_R$ là giá trị trái, phải. So sánh $F'(u_L)$ và $F'(u_R)$. Với $F$ lồi trên khoảng trạng thái, $u_L>u_R$ tạo sốc, $u_L<u_R$ tạo quạt giãn; với $F$ lõm, thứ tự trạng thái đảo lại. Nếu độ lồi đổi dấu, phải xét thêm cấu trúc entropy, không dùng quy tắc lồi trên toàn khoảng.

\textbf{Bước 3 — Tính đường sốc từ bảo toàn.} Với $x=s(t)$,

\[
s'(t)=\frac{F(u_L)-F(u_R)}{u_L-u_R}.
\]

Nếu hai trạng thái hằng, tốc độ $\sigma$ hằng, nên $s=x_0+\sigma t$. Kiểm tra entropy bằng $F'(u_R)<\sigma<F'(u_L)$ đối với sốc nén có hai trạng thái phân biệt trong các bài này. Giá trị ngay trên đường sốc không ảnh hưởng nghiệm yếu.

\textbf{Bước 4 — Tính quạt giãn.} Từ $(x_0,0)$, đặt $z=(x-x_0)/t$, $u=V(z)$. Thế PDE cho $(F'(V)-z)V'=0$. Trong nhánh biến thiên, $F'(V)=z$, nên giải phương trình này tìm $V$. Mép quạt là $x=x_0+F'(u_L)t$ và $x=x_0+F'(u_R)t$, theo thứ tự tốc độ tăng; ghép liên tục với hai trạng thái ngoài. Chọn căn đúng dấu theo khoảng giá trị của dữ liệu.

\textbf{Bước 5 — Tìm mọi thời điểm tương tác.} Giải phương trình giao của hai sốc, sốc với mép quạt, hoặc quạt với quạt. Trước thời điểm gặp, ghép các sóng riêng. Sau khi gặp, tính lại hai trạng thái sát đường sốc. Nếu một phía nằm trong quạt, trạng thái đó phụ thuộc $s(t)$ và $t$, nên Rankine–Hugoniot trở thành ODE mới. Dùng $s(t_*)=x_*$ để tìm hằng số, rồi tiếp tục tìm lần tương tác sau.

\textbf{Bước 6 — Viết nghiệm theo từng khoảng thời gian.} Với mỗi khoảng, ghi rõ các miền $x$, thứ tự đường phân cách và công thức của $u$. Trên nửa trục, kiểm tra chiều truyền từ biên vào miền và thế điều kiện biên. Với dữ liệu trơn từng đoạn, tính $\partial x/\partial\xi$ để tìm thời điểm đầu tiên mất nghiệm cổ điển.

\textbf{Bước 7 — Kiểm tra.} Trong mỗi vùng trơn, thay vào PDE. Trên mỗi sốc, kiểm tra Rankine–Hugoniot và entropy; ở mép quạt, kiểm tra liên tục. Đối chiếu dữ liệu khi $t\downarrow0$, sự liên tục của vị trí sóng qua thời điểm gặp và khối lượng nếu dữ liệu có giá hữu hạn. Nhân PDE với một hàm của $u$ chỉ giữ tương đương cho nghiệm trơn; không tự áp luật đã nhân để tính cùng đường sốc.

\textbf{Chi tiết tích phân đường sốc nằm trong quạt.} Nếu $s'=s/(2t)$ và $s(t_*)=x_*>0$, chia cho $s$ rồi tích phân từ $t_*$ đến $t$: $\log(s/x_*)=\frac12\log(t/t_*)$, nên $s=x_*\sqrt{t/t_*}$. Nếu quạt có tâm $x_0$, thay $s$ bằng $s-x_0$ trước khi tách biến. Khi trạng thái sát sốc thay đổi tại lần gặp tiếp theo, lập ODE mới và lấy vị trí tại lần gặp làm dữ liệu; không giữ hằng số của giai đoạn trước.
\subsection*{Lời giải áp dụng cho bài}
\textbf{Bước 1 — Thông lượng và đặc trưng.} Burgers có $F=u^2/2$. Dữ liệu 2 trái, 1 phải tại $x=1$ cho $X=\xi+2t$ ở trái, $X=\xi+t$ ở phải. Hai họ giao trong $1+t<x<1+2t$.

\textbf{Bước 2 — Rankine–Hugoniot.} $s'=(2^2/2-1^2/2)/(2-1)=3/2$, $s(0)=1$, nên $s=1+3t/2$.

\textbf{Bước 3 — Nghiệm và kiểm tra.} $u=2$ khi $x<s$, $u=1$ khi $x>s$. Sốc nằm giữa hai mép vùng giao và $1<3/2<2$, nên entropy đúng. Các vùng hằng thỏa PDE; khi $t\downarrow0$ có đúng bước nhảy tại 1. Vẽ đặc trưng vào sốc và đồ thị bước nhảy tại $1+3t/2$ để đọc chuyển động.

\clearpage\noindent\texttt{\detokenize{BT_1_6_3_02}}\par
\section*{Lời giải bổ sung — Bài tập 2 · Mục 1.6.3}

\subsection*{Kiến thức và trình tự giải}
\textbf{Mục tiêu.} Với dữ liệu gián đoạn hoặc khi đặc trưng giao nhau, tìm nghiệm yếu thỏa entropy của luật bảo toàn $u_t+(F(u))_x=0$.

\textbf{Bước 1 — Xác định thông lượng và đặc trưng.} Nếu PDE là $u_t+a(u)u_x=0$, tích phân $F'=a$. Trong vùng trơn, $u=u_0(\xi)$ và $x=\xi+F'(u_0(\xi))t$. Tính riêng tốc độ cho mỗi đoạn dữ liệu. Burgers dùng $F=u^2/2$; $u_t+u^2u_x=0$ dùng $F=u^3/3$.

\textbf{Bước 2 — Nhận diện sóng tại mỗi bước nhảy.} Gọi $u_L,u_R$ là giá trị trái, phải. So sánh $F'(u_L)$ và $F'(u_R)$. Với $F$ lồi trên khoảng trạng thái, $u_L>u_R$ tạo sốc, $u_L<u_R$ tạo quạt giãn; với $F$ lõm, thứ tự trạng thái đảo lại. Nếu độ lồi đổi dấu, phải xét thêm cấu trúc entropy, không dùng quy tắc lồi trên toàn khoảng.

\textbf{Bước 3 — Tính đường sốc từ bảo toàn.} Với $x=s(t)$,

\[
s'(t)=\frac{F(u_L)-F(u_R)}{u_L-u_R}.
\]

Nếu hai trạng thái hằng, tốc độ $\sigma$ hằng, nên $s=x_0+\sigma t$. Kiểm tra entropy bằng $F'(u_R)<\sigma<F'(u_L)$ đối với sốc nén có hai trạng thái phân biệt trong các bài này. Giá trị ngay trên đường sốc không ảnh hưởng nghiệm yếu.

\textbf{Bước 4 — Tính quạt giãn.} Từ $(x_0,0)$, đặt $z=(x-x_0)/t$, $u=V(z)$. Thế PDE cho $(F'(V)-z)V'=0$. Trong nhánh biến thiên, $F'(V)=z$, nên giải phương trình này tìm $V$. Mép quạt là $x=x_0+F'(u_L)t$ và $x=x_0+F'(u_R)t$, theo thứ tự tốc độ tăng; ghép liên tục với hai trạng thái ngoài. Chọn căn đúng dấu theo khoảng giá trị của dữ liệu.

\textbf{Bước 5 — Tìm mọi thời điểm tương tác.} Giải phương trình giao của hai sốc, sốc với mép quạt, hoặc quạt với quạt. Trước thời điểm gặp, ghép các sóng riêng. Sau khi gặp, tính lại hai trạng thái sát đường sốc. Nếu một phía nằm trong quạt, trạng thái đó phụ thuộc $s(t)$ và $t$, nên Rankine–Hugoniot trở thành ODE mới. Dùng $s(t_*)=x_*$ để tìm hằng số, rồi tiếp tục tìm lần tương tác sau.

\textbf{Bước 6 — Viết nghiệm theo từng khoảng thời gian.} Với mỗi khoảng, ghi rõ các miền $x$, thứ tự đường phân cách và công thức của $u$. Trên nửa trục, kiểm tra chiều truyền từ biên vào miền và thế điều kiện biên. Với dữ liệu trơn từng đoạn, tính $\partial x/\partial\xi$ để tìm thời điểm đầu tiên mất nghiệm cổ điển.

\textbf{Bước 7 — Kiểm tra.} Trong mỗi vùng trơn, thay vào PDE. Trên mỗi sốc, kiểm tra Rankine–Hugoniot và entropy; ở mép quạt, kiểm tra liên tục. Đối chiếu dữ liệu khi $t\downarrow0$, sự liên tục của vị trí sóng qua thời điểm gặp và khối lượng nếu dữ liệu có giá hữu hạn. Nhân PDE với một hàm của $u$ chỉ giữ tương đương cho nghiệm trơn; không tự áp luật đã nhân để tính cùng đường sốc.

\textbf{Chi tiết tích phân đường sốc nằm trong quạt.} Nếu $s'=s/(2t)$ và $s(t_*)=x_*>0$, chia cho $s$ rồi tích phân từ $t_*$ đến $t$: $\log(s/x_*)=\frac12\log(t/t_*)$, nên $s=x_*\sqrt{t/t_*}$. Nếu quạt có tâm $x_0$, thay $s$ bằng $s-x_0$ trước khi tách biến. Khi trạng thái sát sốc thay đổi tại lần gặp tiếp theo, lập ODE mới và lấy vị trí tại lần gặp làm dữ liệu; không giữ hằng số của giai đoạn trước.
\subsection*{Lời giải áp dụng cho bài}
\textbf{Bước 1 — Thông lượng lõm.} $F=u-u^2$, $F'=1-2u$. Dữ liệu 0 trái, 1 phải tại 1 cho tốc độ 1 và $-1$. Hai họ hội tụ, tạo sốc dù bước nhảy tăng; nguyên nhân là $F''=-2<0$.

\textbf{Bước 2 — Tốc độ sốc.} $F(0)=F(1)=0$, nên $s'=[F(0)-F(1)]/(0-1)=0$. Từ $s(0)=1$, $s=1$.

\textbf{Bước 3 — Nghiệm và kiểm tra.} $u=0$ khi $x<1$, $u=1$ khi $x>1$. $-1<0<1$ kiểm tra entropy; đường sốc đứng yên nhưng đặc trưng vẫn đi vào từ cả hai phía. Nghiệm giữ nguyên dữ liệu ban đầu.

\clearpage\noindent\texttt{\detokenize{BT_1_6_3_03}}\par
\section*{Lời giải bổ sung — Bài tập 3 · Mục 1.6.3}

\subsection*{Kiến thức và trình tự giải}
\textbf{Mục tiêu.} Với dữ liệu gián đoạn hoặc khi đặc trưng giao nhau, tìm nghiệm yếu thỏa entropy của luật bảo toàn $u_t+(F(u))_x=0$.

\textbf{Bước 1 — Xác định thông lượng và đặc trưng.} Nếu PDE là $u_t+a(u)u_x=0$, tích phân $F'=a$. Trong vùng trơn, $u=u_0(\xi)$ và $x=\xi+F'(u_0(\xi))t$. Tính riêng tốc độ cho mỗi đoạn dữ liệu. Burgers dùng $F=u^2/2$; $u_t+u^2u_x=0$ dùng $F=u^3/3$.

\textbf{Bước 2 — Nhận diện sóng tại mỗi bước nhảy.} Gọi $u_L,u_R$ là giá trị trái, phải. So sánh $F'(u_L)$ và $F'(u_R)$. Với $F$ lồi trên khoảng trạng thái, $u_L>u_R$ tạo sốc, $u_L<u_R$ tạo quạt giãn; với $F$ lõm, thứ tự trạng thái đảo lại. Nếu độ lồi đổi dấu, phải xét thêm cấu trúc entropy, không dùng quy tắc lồi trên toàn khoảng.

\textbf{Bước 3 — Tính đường sốc từ bảo toàn.} Với $x=s(t)$,

\[
s'(t)=\frac{F(u_L)-F(u_R)}{u_L-u_R}.
\]

Nếu hai trạng thái hằng, tốc độ $\sigma$ hằng, nên $s=x_0+\sigma t$. Kiểm tra entropy bằng $F'(u_R)<\sigma<F'(u_L)$ đối với sốc nén có hai trạng thái phân biệt trong các bài này. Giá trị ngay trên đường sốc không ảnh hưởng nghiệm yếu.

\textbf{Bước 4 — Tính quạt giãn.} Từ $(x_0,0)$, đặt $z=(x-x_0)/t$, $u=V(z)$. Thế PDE cho $(F'(V)-z)V'=0$. Trong nhánh biến thiên, $F'(V)=z$, nên giải phương trình này tìm $V$. Mép quạt là $x=x_0+F'(u_L)t$ và $x=x_0+F'(u_R)t$, theo thứ tự tốc độ tăng; ghép liên tục với hai trạng thái ngoài. Chọn căn đúng dấu theo khoảng giá trị của dữ liệu.

\textbf{Bước 5 — Tìm mọi thời điểm tương tác.} Giải phương trình giao của hai sốc, sốc với mép quạt, hoặc quạt với quạt. Trước thời điểm gặp, ghép các sóng riêng. Sau khi gặp, tính lại hai trạng thái sát đường sốc. Nếu một phía nằm trong quạt, trạng thái đó phụ thuộc $s(t)$ và $t$, nên Rankine–Hugoniot trở thành ODE mới. Dùng $s(t_*)=x_*$ để tìm hằng số, rồi tiếp tục tìm lần tương tác sau.

\textbf{Bước 6 — Viết nghiệm theo từng khoảng thời gian.} Với mỗi khoảng, ghi rõ các miền $x$, thứ tự đường phân cách và công thức của $u$. Trên nửa trục, kiểm tra chiều truyền từ biên vào miền và thế điều kiện biên. Với dữ liệu trơn từng đoạn, tính $\partial x/\partial\xi$ để tìm thời điểm đầu tiên mất nghiệm cổ điển.

\textbf{Bước 7 — Kiểm tra.} Trong mỗi vùng trơn, thay vào PDE. Trên mỗi sốc, kiểm tra Rankine–Hugoniot và entropy; ở mép quạt, kiểm tra liên tục. Đối chiếu dữ liệu khi $t\downarrow0$, sự liên tục của vị trí sóng qua thời điểm gặp và khối lượng nếu dữ liệu có giá hữu hạn. Nhân PDE với một hàm của $u$ chỉ giữ tương đương cho nghiệm trơn; không tự áp luật đã nhân để tính cùng đường sốc.

\textbf{Chi tiết tích phân đường sốc nằm trong quạt.} Nếu $s'=s/(2t)$ và $s(t_*)=x_*>0$, chia cho $s$ rồi tích phân từ $t_*$ đến $t$: $\log(s/x_*)=\frac12\log(t/t_*)$, nên $s=x_*\sqrt{t/t_*}$. Nếu quạt có tâm $x_0$, thay $s$ bằng $s-x_0$ trước khi tách biến. Khi trạng thái sát sốc thay đổi tại lần gặp tiếp theo, lập ODE mới và lấy vị trí tại lần gặp làm dữ liệu; không giữ hằng số của giai đoạn trước.
\subsection*{Lời giải áp dụng cho bài}
\textbf{Bước 1 — Dữ liệu đầu và biên.} Trong $x<0$, trạng thái đầu là 1, tốc độ $F'(1)=-1$; trạng thái biên tại 0 là 2, tốc độ $F'(2)=-3$, truyền vào nửa trục âm. Hai trạng thái nối theo thứ tự 1 bên trái, 2 bên phải.

\textbf{Bước 2 — Tính sốc.} Với $F=u-u^2$, $F(1)=0,F(2)=-2$, nên $s'=[0-(-2)]/(1-2)=-2$. Sốc xuất phát ở góc $(0,0)$, $s=-2t$.

\textbf{Bước 3 — Công thức và thử lại.} $u=1$ khi $x<-2t$, $u=2$ khi $-2t<x<0$. Entropy: $-3<-2<-1$; biên $x\uparrow0$ cho 2, còn mỗi $x<0$ cố định khi $t\downarrow0$ cho 1. Đây là nghiệm yếu với dữ liệu góc không liên tục.

\clearpage\noindent\texttt{\detokenize{BT_1_6_3_04}}\par
\section*{Lời giải bổ sung — Bài tập 4 · Mục 1.6.3}

\subsection*{Kiến thức và trình tự giải}
\textbf{Mục tiêu.} Với dữ liệu gián đoạn hoặc khi đặc trưng giao nhau, tìm nghiệm yếu thỏa entropy của luật bảo toàn $u_t+(F(u))_x=0$.

\textbf{Bước 1 — Xác định thông lượng và đặc trưng.} Nếu PDE là $u_t+a(u)u_x=0$, tích phân $F'=a$. Trong vùng trơn, $u=u_0(\xi)$ và $x=\xi+F'(u_0(\xi))t$. Tính riêng tốc độ cho mỗi đoạn dữ liệu. Burgers dùng $F=u^2/2$; $u_t+u^2u_x=0$ dùng $F=u^3/3$.

\textbf{Bước 2 — Nhận diện sóng tại mỗi bước nhảy.} Gọi $u_L,u_R$ là giá trị trái, phải. So sánh $F'(u_L)$ và $F'(u_R)$. Với $F$ lồi trên khoảng trạng thái, $u_L>u_R$ tạo sốc, $u_L<u_R$ tạo quạt giãn; với $F$ lõm, thứ tự trạng thái đảo lại. Nếu độ lồi đổi dấu, phải xét thêm cấu trúc entropy, không dùng quy tắc lồi trên toàn khoảng.

\textbf{Bước 3 — Tính đường sốc từ bảo toàn.} Với $x=s(t)$,

\[
s'(t)=\frac{F(u_L)-F(u_R)}{u_L-u_R}.
\]

Nếu hai trạng thái hằng, tốc độ $\sigma$ hằng, nên $s=x_0+\sigma t$. Kiểm tra entropy bằng $F'(u_R)<\sigma<F'(u_L)$ đối với sốc nén có hai trạng thái phân biệt trong các bài này. Giá trị ngay trên đường sốc không ảnh hưởng nghiệm yếu.

\textbf{Bước 4 — Tính quạt giãn.} Từ $(x_0,0)$, đặt $z=(x-x_0)/t$, $u=V(z)$. Thế PDE cho $(F'(V)-z)V'=0$. Trong nhánh biến thiên, $F'(V)=z$, nên giải phương trình này tìm $V$. Mép quạt là $x=x_0+F'(u_L)t$ và $x=x_0+F'(u_R)t$, theo thứ tự tốc độ tăng; ghép liên tục với hai trạng thái ngoài. Chọn căn đúng dấu theo khoảng giá trị của dữ liệu.

\textbf{Bước 5 — Tìm mọi thời điểm tương tác.} Giải phương trình giao của hai sốc, sốc với mép quạt, hoặc quạt với quạt. Trước thời điểm gặp, ghép các sóng riêng. Sau khi gặp, tính lại hai trạng thái sát đường sốc. Nếu một phía nằm trong quạt, trạng thái đó phụ thuộc $s(t)$ và $t$, nên Rankine–Hugoniot trở thành ODE mới. Dùng $s(t_*)=x_*$ để tìm hằng số, rồi tiếp tục tìm lần tương tác sau.

\textbf{Bước 6 — Viết nghiệm theo từng khoảng thời gian.} Với mỗi khoảng, ghi rõ các miền $x$, thứ tự đường phân cách và công thức của $u$. Trên nửa trục, kiểm tra chiều truyền từ biên vào miền và thế điều kiện biên. Với dữ liệu trơn từng đoạn, tính $\partial x/\partial\xi$ để tìm thời điểm đầu tiên mất nghiệm cổ điển.

\textbf{Bước 7 — Kiểm tra.} Trong mỗi vùng trơn, thay vào PDE. Trên mỗi sốc, kiểm tra Rankine–Hugoniot và entropy; ở mép quạt, kiểm tra liên tục. Đối chiếu dữ liệu khi $t\downarrow0$, sự liên tục của vị trí sóng qua thời điểm gặp và khối lượng nếu dữ liệu có giá hữu hạn. Nhân PDE với một hàm của $u$ chỉ giữ tương đương cho nghiệm trơn; không tự áp luật đã nhân để tính cùng đường sốc.

\textbf{Chi tiết tích phân đường sốc nằm trong quạt.} Nếu $s'=s/(2t)$ và $s(t_*)=x_*>0$, chia cho $s$ rồi tích phân từ $t_*$ đến $t$: $\log(s/x_*)=\frac12\log(t/t_*)$, nên $s=x_*\sqrt{t/t_*}$. Nếu quạt có tâm $x_0$, thay $s$ bằng $s-x_0$ trước khi tách biến. Khi trạng thái sát sốc thay đổi tại lần gặp tiếp theo, lập ODE mới và lấy vị trí tại lần gặp làm dữ liệu; không giữ hằng số của giai đoạn trước.
\subsection*{Lời giải áp dụng cho bài}
\textbf{Bước 1 — Thông lượng và hướng truyền.} $F=u^3/3$, $F'=u^2$. Trong $x>0$, trạng thái biên 2 nằm trái, trạng thái đầu 1 nằm phải. Tốc độ là 4 và 1, đều hướng vào nửa trục dương.

\textbf{Bước 2 — Sốc từ góc.} $s'=(8/3-1/3)/(2-1)=7/3$, $s(0)=0$, nên $s=7t/3$.

\textbf{Bước 3 — Nghiệm và kiểm tra.} $u=2$ trên $0<x<7t/3$, $u=1$ trên $x>7t/3$. $1<7/3<4$ bảo đảm sốc nén; vết tại $x=0$ là 2 và giới hạn $t\downarrow0$ ở mọi $x>0$ là 1. Không dùng tốc độ $u$ của Burgers vì ở bài này tốc độ đặc trưng là $u^2$.

\clearpage\noindent\texttt{\detokenize{BT_1_6_3_05}}\par
\section*{Lời giải bổ sung — Bài tập 5 · Mục 1.6.3}

\subsection*{Kiến thức và trình tự giải}
\textbf{Mục tiêu.} Với dữ liệu gián đoạn hoặc khi đặc trưng giao nhau, tìm nghiệm yếu thỏa entropy của luật bảo toàn $u_t+(F(u))_x=0$.

\textbf{Bước 1 — Xác định thông lượng và đặc trưng.} Nếu PDE là $u_t+a(u)u_x=0$, tích phân $F'=a$. Trong vùng trơn, $u=u_0(\xi)$ và $x=\xi+F'(u_0(\xi))t$. Tính riêng tốc độ cho mỗi đoạn dữ liệu. Burgers dùng $F=u^2/2$; $u_t+u^2u_x=0$ dùng $F=u^3/3$.

\textbf{Bước 2 — Nhận diện sóng tại mỗi bước nhảy.} Gọi $u_L,u_R$ là giá trị trái, phải. So sánh $F'(u_L)$ và $F'(u_R)$. Với $F$ lồi trên khoảng trạng thái, $u_L>u_R$ tạo sốc, $u_L<u_R$ tạo quạt giãn; với $F$ lõm, thứ tự trạng thái đảo lại. Nếu độ lồi đổi dấu, phải xét thêm cấu trúc entropy, không dùng quy tắc lồi trên toàn khoảng.

\textbf{Bước 3 — Tính đường sốc từ bảo toàn.} Với $x=s(t)$,

\[
s'(t)=\frac{F(u_L)-F(u_R)}{u_L-u_R}.
\]

Nếu hai trạng thái hằng, tốc độ $\sigma$ hằng, nên $s=x_0+\sigma t$. Kiểm tra entropy bằng $F'(u_R)<\sigma<F'(u_L)$ đối với sốc nén có hai trạng thái phân biệt trong các bài này. Giá trị ngay trên đường sốc không ảnh hưởng nghiệm yếu.

\textbf{Bước 4 — Tính quạt giãn.} Từ $(x_0,0)$, đặt $z=(x-x_0)/t$, $u=V(z)$. Thế PDE cho $(F'(V)-z)V'=0$. Trong nhánh biến thiên, $F'(V)=z$, nên giải phương trình này tìm $V$. Mép quạt là $x=x_0+F'(u_L)t$ và $x=x_0+F'(u_R)t$, theo thứ tự tốc độ tăng; ghép liên tục với hai trạng thái ngoài. Chọn căn đúng dấu theo khoảng giá trị của dữ liệu.

\textbf{Bước 5 — Tìm mọi thời điểm tương tác.} Giải phương trình giao của hai sốc, sốc với mép quạt, hoặc quạt với quạt. Trước thời điểm gặp, ghép các sóng riêng. Sau khi gặp, tính lại hai trạng thái sát đường sốc. Nếu một phía nằm trong quạt, trạng thái đó phụ thuộc $s(t)$ và $t$, nên Rankine–Hugoniot trở thành ODE mới. Dùng $s(t_*)=x_*$ để tìm hằng số, rồi tiếp tục tìm lần tương tác sau.

\textbf{Bước 6 — Viết nghiệm theo từng khoảng thời gian.} Với mỗi khoảng, ghi rõ các miền $x$, thứ tự đường phân cách và công thức của $u$. Trên nửa trục, kiểm tra chiều truyền từ biên vào miền và thế điều kiện biên. Với dữ liệu trơn từng đoạn, tính $\partial x/\partial\xi$ để tìm thời điểm đầu tiên mất nghiệm cổ điển.

\textbf{Bước 7 — Kiểm tra.} Trong mỗi vùng trơn, thay vào PDE. Trên mỗi sốc, kiểm tra Rankine–Hugoniot và entropy; ở mép quạt, kiểm tra liên tục. Đối chiếu dữ liệu khi $t\downarrow0$, sự liên tục của vị trí sóng qua thời điểm gặp và khối lượng nếu dữ liệu có giá hữu hạn. Nhân PDE với một hàm của $u$ chỉ giữ tương đương cho nghiệm trơn; không tự áp luật đã nhân để tính cùng đường sốc.

\textbf{Chi tiết tích phân đường sốc nằm trong quạt.} Nếu $s'=s/(2t)$ và $s(t_*)=x_*>0$, chia cho $s$ rồi tích phân từ $t_*$ đến $t$: $\log(s/x_*)=\frac12\log(t/t_*)$, nên $s=x_*\sqrt{t/t_*}$. Nếu quạt có tâm $x_0$, thay $s$ bằng $s-x_0$ trước khi tách biến. Khi trạng thái sát sốc thay đổi tại lần gặp tiếp theo, lập ODE mới và lấy vị trí tại lần gặp làm dữ liệu; không giữ hằng số của giai đoạn trước.
\subsection*{Lời giải áp dụng cho bài}
\textbf{Bước 1 — Nhận dạng hai quạt.} Burgers có bước $-1\to1$ tại $-1$ và $1\to2$ tại 1. Cả hai tăng nên là quạt giãn. Quạt thứ nhất có mép $-1-t,-1+t$ và giá trị $(x+1)/t$. Quạt thứ hai có mép $1+t,1+2t$ và giá trị $(x-1)/t$.

\textbf{Bước 2 — Xét tương tác.} Mép phải quạt đầu và mép trái quạt sau cùng tốc độ 1, khoảng cách luôn $[1+t]-[-1+t]=2$. Vì vậy hai quạt không gặp nhau.

\textbf{Bước 3 — Ghép nghiệm.}
\[u=\begin{cases}-1,&x\le-1-t,\\(x+1)/t,&-1-t<x<-1+t,\\1,&-1+t\le x\le1+t,\\(x-1)/t,&1+t<x<1+2t,\\2,&x\ge1+2t.\end{cases}\]
Ở các mép, công thức biến thiên bằng trạng thái hằng kề bên. Mỗi quạt có $u_t+uu_x=0$ và dữ liệu khi $t\downarrow0$ đúng từng đoạn.

\clearpage\noindent\texttt{\detokenize{BT_1_6_3_06}}\par
\section*{Lời giải bổ sung — Bài tập 6 · Mục 1.6.3}

\subsection*{Kiến thức và trình tự giải}
\textbf{Mục tiêu.} Với dữ liệu gián đoạn hoặc khi đặc trưng giao nhau, tìm nghiệm yếu thỏa entropy của luật bảo toàn $u_t+(F(u))_x=0$.

\textbf{Bước 1 — Xác định thông lượng và đặc trưng.} Nếu PDE là $u_t+a(u)u_x=0$, tích phân $F'=a$. Trong vùng trơn, $u=u_0(\xi)$ và $x=\xi+F'(u_0(\xi))t$. Tính riêng tốc độ cho mỗi đoạn dữ liệu. Burgers dùng $F=u^2/2$; $u_t+u^2u_x=0$ dùng $F=u^3/3$.

\textbf{Bước 2 — Nhận diện sóng tại mỗi bước nhảy.} Gọi $u_L,u_R$ là giá trị trái, phải. So sánh $F'(u_L)$ và $F'(u_R)$. Với $F$ lồi trên khoảng trạng thái, $u_L>u_R$ tạo sốc, $u_L<u_R$ tạo quạt giãn; với $F$ lõm, thứ tự trạng thái đảo lại. Nếu độ lồi đổi dấu, phải xét thêm cấu trúc entropy, không dùng quy tắc lồi trên toàn khoảng.

\textbf{Bước 3 — Tính đường sốc từ bảo toàn.} Với $x=s(t)$,

\[
s'(t)=\frac{F(u_L)-F(u_R)}{u_L-u_R}.
\]

Nếu hai trạng thái hằng, tốc độ $\sigma$ hằng, nên $s=x_0+\sigma t$. Kiểm tra entropy bằng $F'(u_R)<\sigma<F'(u_L)$ đối với sốc nén có hai trạng thái phân biệt trong các bài này. Giá trị ngay trên đường sốc không ảnh hưởng nghiệm yếu.

\textbf{Bước 4 — Tính quạt giãn.} Từ $(x_0,0)$, đặt $z=(x-x_0)/t$, $u=V(z)$. Thế PDE cho $(F'(V)-z)V'=0$. Trong nhánh biến thiên, $F'(V)=z$, nên giải phương trình này tìm $V$. Mép quạt là $x=x_0+F'(u_L)t$ và $x=x_0+F'(u_R)t$, theo thứ tự tốc độ tăng; ghép liên tục với hai trạng thái ngoài. Chọn căn đúng dấu theo khoảng giá trị của dữ liệu.

\textbf{Bước 5 — Tìm mọi thời điểm tương tác.} Giải phương trình giao của hai sốc, sốc với mép quạt, hoặc quạt với quạt. Trước thời điểm gặp, ghép các sóng riêng. Sau khi gặp, tính lại hai trạng thái sát đường sốc. Nếu một phía nằm trong quạt, trạng thái đó phụ thuộc $s(t)$ và $t$, nên Rankine–Hugoniot trở thành ODE mới. Dùng $s(t_*)=x_*$ để tìm hằng số, rồi tiếp tục tìm lần tương tác sau.

\textbf{Bước 6 — Viết nghiệm theo từng khoảng thời gian.} Với mỗi khoảng, ghi rõ các miền $x$, thứ tự đường phân cách và công thức của $u$. Trên nửa trục, kiểm tra chiều truyền từ biên vào miền và thế điều kiện biên. Với dữ liệu trơn từng đoạn, tính $\partial x/\partial\xi$ để tìm thời điểm đầu tiên mất nghiệm cổ điển.

\textbf{Bước 7 — Kiểm tra.} Trong mỗi vùng trơn, thay vào PDE. Trên mỗi sốc, kiểm tra Rankine–Hugoniot và entropy; ở mép quạt, kiểm tra liên tục. Đối chiếu dữ liệu khi $t\downarrow0$, sự liên tục của vị trí sóng qua thời điểm gặp và khối lượng nếu dữ liệu có giá hữu hạn. Nhân PDE với một hàm của $u$ chỉ giữ tương đương cho nghiệm trơn; không tự áp luật đã nhân để tính cùng đường sốc.

\textbf{Chi tiết tích phân đường sốc nằm trong quạt.} Nếu $s'=s/(2t)$ và $s(t_*)=x_*>0$, chia cho $s$ rồi tích phân từ $t_*$ đến $t$: $\log(s/x_*)=\frac12\log(t/t_*)$, nên $s=x_*\sqrt{t/t_*}$. Nếu quạt có tâm $x_0$, thay $s$ bằng $s-x_0$ trước khi tách biến. Khi trạng thái sát sốc thay đổi tại lần gặp tiếp theo, lập ODE mới và lấy vị trí tại lần gặp làm dữ liệu; không giữ hằng số của giai đoạn trước.
\subsection*{Lời giải áp dụng cho bài}
\textbf{Bước 1 — Quạt của thông lượng lõm.} $F=u-u^2$, tốc độ $1-2u$. Trạng thái 2 bên trái cho tốc độ $-3$, trạng thái 1 bên phải cho $-1$, nên tách xa và tạo quạt ở nửa trục $x<0$.

\textbf{Bước 2 — Nhánh biến thiên.} Đặt $z=x/t$, giải $F'(u)=1-2u=z$, được $u=(1-z)/2$. Mép là $x=-3t,-t$.

\textbf{Bước 3 — Ghép và kiểm tra.}
\[u=\begin{cases}2,&x\le-3t,\\(1-x/t)/2,&-3t<x<-t,\\1,&-t\le x<0.\end{cases}\]
Trong quạt $u_t=x/(2t^2)$, $u_x=-1/(2t)$, $1-2u=x/t$, nên $u_t+(1-2u)u_x=0$. Vết biên là 1, còn trạng thái đầu ở $x<0$ là 2. Các tia trong quạt là $x=(1-2k)t$, $1\le k\le2$.

\clearpage\noindent\texttt{\detokenize{BT_1_6_3_07}}\par
\section*{Lời giải bổ sung — Bài tập 7 · Mục 1.6.3}

\subsection*{Kiến thức và trình tự giải}
\textbf{Mục tiêu.} Với dữ liệu gián đoạn hoặc khi đặc trưng giao nhau, tìm nghiệm yếu thỏa entropy của luật bảo toàn $u_t+(F(u))_x=0$.

\textbf{Bước 1 — Xác định thông lượng và đặc trưng.} Nếu PDE là $u_t+a(u)u_x=0$, tích phân $F'=a$. Trong vùng trơn, $u=u_0(\xi)$ và $x=\xi+F'(u_0(\xi))t$. Tính riêng tốc độ cho mỗi đoạn dữ liệu. Burgers dùng $F=u^2/2$; $u_t+u^2u_x=0$ dùng $F=u^3/3$.

\textbf{Bước 2 — Nhận diện sóng tại mỗi bước nhảy.} Gọi $u_L,u_R$ là giá trị trái, phải. So sánh $F'(u_L)$ và $F'(u_R)$. Với $F$ lồi trên khoảng trạng thái, $u_L>u_R$ tạo sốc, $u_L<u_R$ tạo quạt giãn; với $F$ lõm, thứ tự trạng thái đảo lại. Nếu độ lồi đổi dấu, phải xét thêm cấu trúc entropy, không dùng quy tắc lồi trên toàn khoảng.

\textbf{Bước 3 — Tính đường sốc từ bảo toàn.} Với $x=s(t)$,

\[
s'(t)=\frac{F(u_L)-F(u_R)}{u_L-u_R}.
\]

Nếu hai trạng thái hằng, tốc độ $\sigma$ hằng, nên $s=x_0+\sigma t$. Kiểm tra entropy bằng $F'(u_R)<\sigma<F'(u_L)$ đối với sốc nén có hai trạng thái phân biệt trong các bài này. Giá trị ngay trên đường sốc không ảnh hưởng nghiệm yếu.

\textbf{Bước 4 — Tính quạt giãn.} Từ $(x_0,0)$, đặt $z=(x-x_0)/t$, $u=V(z)$. Thế PDE cho $(F'(V)-z)V'=0$. Trong nhánh biến thiên, $F'(V)=z$, nên giải phương trình này tìm $V$. Mép quạt là $x=x_0+F'(u_L)t$ và $x=x_0+F'(u_R)t$, theo thứ tự tốc độ tăng; ghép liên tục với hai trạng thái ngoài. Chọn căn đúng dấu theo khoảng giá trị của dữ liệu.

\textbf{Bước 5 — Tìm mọi thời điểm tương tác.} Giải phương trình giao của hai sốc, sốc với mép quạt, hoặc quạt với quạt. Trước thời điểm gặp, ghép các sóng riêng. Sau khi gặp, tính lại hai trạng thái sát đường sốc. Nếu một phía nằm trong quạt, trạng thái đó phụ thuộc $s(t)$ và $t$, nên Rankine–Hugoniot trở thành ODE mới. Dùng $s(t_*)=x_*$ để tìm hằng số, rồi tiếp tục tìm lần tương tác sau.

\textbf{Bước 6 — Viết nghiệm theo từng khoảng thời gian.} Với mỗi khoảng, ghi rõ các miền $x$, thứ tự đường phân cách và công thức của $u$. Trên nửa trục, kiểm tra chiều truyền từ biên vào miền và thế điều kiện biên. Với dữ liệu trơn từng đoạn, tính $\partial x/\partial\xi$ để tìm thời điểm đầu tiên mất nghiệm cổ điển.

\textbf{Bước 7 — Kiểm tra.} Trong mỗi vùng trơn, thay vào PDE. Trên mỗi sốc, kiểm tra Rankine–Hugoniot và entropy; ở mép quạt, kiểm tra liên tục. Đối chiếu dữ liệu khi $t\downarrow0$, sự liên tục của vị trí sóng qua thời điểm gặp và khối lượng nếu dữ liệu có giá hữu hạn. Nhân PDE với một hàm của $u$ chỉ giữ tương đương cho nghiệm trơn; không tự áp luật đã nhân để tính cùng đường sốc.

\textbf{Chi tiết tích phân đường sốc nằm trong quạt.} Nếu $s'=s/(2t)$ và $s(t_*)=x_*>0$, chia cho $s$ rồi tích phân từ $t_*$ đến $t$: $\log(s/x_*)=\frac12\log(t/t_*)$, nên $s=x_*\sqrt{t/t_*}$. Nếu quạt có tâm $x_0$, thay $s$ bằng $s-x_0$ trước khi tách biến. Khi trạng thái sát sốc thay đổi tại lần gặp tiếp theo, lập ODE mới và lấy vị trí tại lần gặp làm dữ liệu; không giữ hằng số của giai đoạn trước.
\subsection*{Lời giải áp dụng cho bài}
\textbf{Bước 1 — Xét độ lõm trên khoảng dữ liệu.} $F=u-u^3$, $F'=1-3u^2$, $F''=-6u<0$ trên $[1,2]$. Trạng thái 2 trái có tốc độ $-11$, 1 phải có tốc độ $-2$, nên là quạt giãn.

\textbf{Bước 2 — Giải trong quạt.} $1-3u^2=x/t$ cho $u=\sqrt{(1-x/t)/3}$; chọn căn dương vì dữ liệu nằm trong $[1,2]$. Mép là $-11t,-2t$.

\textbf{Bước 3 — Nghiệm và kiểm tra.}
\[u=\begin{cases}2,&x\le-11t,\\\sqrt{(1-x/t)/3},&-11t<x<-2t,\\1,&-2t\le x<0.\end{cases}\]
Đẳng thức $F'(u)=x/t$ bảo đảm PDE tự đồng dạng; các mép cho 2 và 1. Biên tại 0 bằng 1 và khi $t\downarrow0$ mọi điểm $x<0$ trở về 2.

\clearpage\noindent\texttt{\detokenize{BT_1_6_3_08}}\par
\section*{Lời giải bổ sung — Bài tập 8 · Mục 1.6.3}

\subsection*{Kiến thức và trình tự giải}
\textbf{Mục tiêu.} Với dữ liệu gián đoạn hoặc khi đặc trưng giao nhau, tìm nghiệm yếu thỏa entropy của luật bảo toàn $u_t+(F(u))_x=0$.

\textbf{Bước 1 — Xác định thông lượng và đặc trưng.} Nếu PDE là $u_t+a(u)u_x=0$, tích phân $F'=a$. Trong vùng trơn, $u=u_0(\xi)$ và $x=\xi+F'(u_0(\xi))t$. Tính riêng tốc độ cho mỗi đoạn dữ liệu. Burgers dùng $F=u^2/2$; $u_t+u^2u_x=0$ dùng $F=u^3/3$.

\textbf{Bước 2 — Nhận diện sóng tại mỗi bước nhảy.} Gọi $u_L,u_R$ là giá trị trái, phải. So sánh $F'(u_L)$ và $F'(u_R)$. Với $F$ lồi trên khoảng trạng thái, $u_L>u_R$ tạo sốc, $u_L<u_R$ tạo quạt giãn; với $F$ lõm, thứ tự trạng thái đảo lại. Nếu độ lồi đổi dấu, phải xét thêm cấu trúc entropy, không dùng quy tắc lồi trên toàn khoảng.

\textbf{Bước 3 — Tính đường sốc từ bảo toàn.} Với $x=s(t)$,

\[
s'(t)=\frac{F(u_L)-F(u_R)}{u_L-u_R}.
\]

Nếu hai trạng thái hằng, tốc độ $\sigma$ hằng, nên $s=x_0+\sigma t$. Kiểm tra entropy bằng $F'(u_R)<\sigma<F'(u_L)$ đối với sốc nén có hai trạng thái phân biệt trong các bài này. Giá trị ngay trên đường sốc không ảnh hưởng nghiệm yếu.

\textbf{Bước 4 — Tính quạt giãn.} Từ $(x_0,0)$, đặt $z=(x-x_0)/t$, $u=V(z)$. Thế PDE cho $(F'(V)-z)V'=0$. Trong nhánh biến thiên, $F'(V)=z$, nên giải phương trình này tìm $V$. Mép quạt là $x=x_0+F'(u_L)t$ và $x=x_0+F'(u_R)t$, theo thứ tự tốc độ tăng; ghép liên tục với hai trạng thái ngoài. Chọn căn đúng dấu theo khoảng giá trị của dữ liệu.

\textbf{Bước 5 — Tìm mọi thời điểm tương tác.} Giải phương trình giao của hai sốc, sốc với mép quạt, hoặc quạt với quạt. Trước thời điểm gặp, ghép các sóng riêng. Sau khi gặp, tính lại hai trạng thái sát đường sốc. Nếu một phía nằm trong quạt, trạng thái đó phụ thuộc $s(t)$ và $t$, nên Rankine–Hugoniot trở thành ODE mới. Dùng $s(t_*)=x_*$ để tìm hằng số, rồi tiếp tục tìm lần tương tác sau.

\textbf{Bước 6 — Viết nghiệm theo từng khoảng thời gian.} Với mỗi khoảng, ghi rõ các miền $x$, thứ tự đường phân cách và công thức của $u$. Trên nửa trục, kiểm tra chiều truyền từ biên vào miền và thế điều kiện biên. Với dữ liệu trơn từng đoạn, tính $\partial x/\partial\xi$ để tìm thời điểm đầu tiên mất nghiệm cổ điển.

\textbf{Bước 7 — Kiểm tra.} Trong mỗi vùng trơn, thay vào PDE. Trên mỗi sốc, kiểm tra Rankine–Hugoniot và entropy; ở mép quạt, kiểm tra liên tục. Đối chiếu dữ liệu khi $t\downarrow0$, sự liên tục của vị trí sóng qua thời điểm gặp và khối lượng nếu dữ liệu có giá hữu hạn. Nhân PDE với một hàm của $u$ chỉ giữ tương đương cho nghiệm trơn; không tự áp luật đã nhân để tính cùng đường sốc.

\textbf{Chi tiết tích phân đường sốc nằm trong quạt.} Nếu $s'=s/(2t)$ và $s(t_*)=x_*>0$, chia cho $s$ rồi tích phân từ $t_*$ đến $t$: $\log(s/x_*)=\frac12\log(t/t_*)$, nên $s=x_*\sqrt{t/t_*}$. Nếu quạt có tâm $x_0$, thay $s$ bằng $s-x_0$ trước khi tách biến. Khi trạng thái sát sốc thay đổi tại lần gặp tiếp theo, lập ODE mới và lấy vị trí tại lần gặp làm dữ liệu; không giữ hằng số của giai đoạn trước.
\subsection*{Lời giải áp dụng cho bài}
\textbf{Bước 1 — Hai sốc của $F=u-u^2$.} Dữ liệu 0, $1/2$, 1 với bước tại 0 và 1. Công thức tốc độ rút gọn thành $\sigma=1-(u_L+u_R)$. Suy $\sigma_1=1/2$, $\sigma_2=-1/2$, $s_1=t/2$, $s_2=1-t/2$.

\textbf{Bước 2 — Thời điểm gặp.} $t/2=1-t/2$ cho $t_*=1,x_*=1/2$. Với $0<t<1$, nghiệm là 0 ở $x<s_1$, $1/2$ giữa hai sốc, 1 ở $x>s_2$.

\textbf{Bước 3 — Sốc sau gặp.} Nối trực tiếp 0 và 1, tốc độ mới $1-(0+1)=0$. Dùng $s(1)=1/2$, đường mới là $s=1/2$. Với $t\ge1$, $u=0$ trái đường, $u=1$ phải đường.

\textbf{Bước 4 — Kiểm tra.} Các tốc độ nằm giữa đặc trưng: $0<1/2<1$, $-1<-1/2<0$, $-1<0<1$. Đường sóng ghép đúng điểm gặp; không tiếp tục hai đường sốc cũ sau $t=1$.

\textbf{Chi tiết rút gọn tốc độ.} Với $F(u)=u-u^2$,
\[\frac{F(u_L)-F(u_R)}{u_L-u_R}
=\frac{(u_L-u_R)-(u_L^2-u_R^2)}{u_L-u_R}
=1-u_L-u_R.\]
Thế $(u_L,u_R)=(0,1/2)$ cho $1/2$, $(1/2,1)$ cho $-1/2$, và $(0,1)$ cho $0$ sau tương tác.

\clearpage\noindent\texttt{\detokenize{BT_1_6_3_09}}\par
\section*{Lời giải bổ sung — Bài tập 9 · Mục 1.6.3}

\subsection*{Kiến thức và trình tự giải}
\textbf{Mục tiêu.} Với dữ liệu gián đoạn hoặc khi đặc trưng giao nhau, tìm nghiệm yếu thỏa entropy của luật bảo toàn $u_t+(F(u))_x=0$.

\textbf{Bước 1 — Xác định thông lượng và đặc trưng.} Nếu PDE là $u_t+a(u)u_x=0$, tích phân $F'=a$. Trong vùng trơn, $u=u_0(\xi)$ và $x=\xi+F'(u_0(\xi))t$. Tính riêng tốc độ cho mỗi đoạn dữ liệu. Burgers dùng $F=u^2/2$; $u_t+u^2u_x=0$ dùng $F=u^3/3$.

\textbf{Bước 2 — Nhận diện sóng tại mỗi bước nhảy.} Gọi $u_L,u_R$ là giá trị trái, phải. So sánh $F'(u_L)$ và $F'(u_R)$. Với $F$ lồi trên khoảng trạng thái, $u_L>u_R$ tạo sốc, $u_L<u_R$ tạo quạt giãn; với $F$ lõm, thứ tự trạng thái đảo lại. Nếu độ lồi đổi dấu, phải xét thêm cấu trúc entropy, không dùng quy tắc lồi trên toàn khoảng.

\textbf{Bước 3 — Tính đường sốc từ bảo toàn.} Với $x=s(t)$,

\[
s'(t)=\frac{F(u_L)-F(u_R)}{u_L-u_R}.
\]

Nếu hai trạng thái hằng, tốc độ $\sigma$ hằng, nên $s=x_0+\sigma t$. Kiểm tra entropy bằng $F'(u_R)<\sigma<F'(u_L)$ đối với sốc nén có hai trạng thái phân biệt trong các bài này. Giá trị ngay trên đường sốc không ảnh hưởng nghiệm yếu.

\textbf{Bước 4 — Tính quạt giãn.} Từ $(x_0,0)$, đặt $z=(x-x_0)/t$, $u=V(z)$. Thế PDE cho $(F'(V)-z)V'=0$. Trong nhánh biến thiên, $F'(V)=z$, nên giải phương trình này tìm $V$. Mép quạt là $x=x_0+F'(u_L)t$ và $x=x_0+F'(u_R)t$, theo thứ tự tốc độ tăng; ghép liên tục với hai trạng thái ngoài. Chọn căn đúng dấu theo khoảng giá trị của dữ liệu.

\textbf{Bước 5 — Tìm mọi thời điểm tương tác.} Giải phương trình giao của hai sốc, sốc với mép quạt, hoặc quạt với quạt. Trước thời điểm gặp, ghép các sóng riêng. Sau khi gặp, tính lại hai trạng thái sát đường sốc. Nếu một phía nằm trong quạt, trạng thái đó phụ thuộc $s(t)$ và $t$, nên Rankine–Hugoniot trở thành ODE mới. Dùng $s(t_*)=x_*$ để tìm hằng số, rồi tiếp tục tìm lần tương tác sau.

\textbf{Bước 6 — Viết nghiệm theo từng khoảng thời gian.} Với mỗi khoảng, ghi rõ các miền $x$, thứ tự đường phân cách và công thức của $u$. Trên nửa trục, kiểm tra chiều truyền từ biên vào miền và thế điều kiện biên. Với dữ liệu trơn từng đoạn, tính $\partial x/\partial\xi$ để tìm thời điểm đầu tiên mất nghiệm cổ điển.

\textbf{Bước 7 — Kiểm tra.} Trong mỗi vùng trơn, thay vào PDE. Trên mỗi sốc, kiểm tra Rankine–Hugoniot và entropy; ở mép quạt, kiểm tra liên tục. Đối chiếu dữ liệu khi $t\downarrow0$, sự liên tục của vị trí sóng qua thời điểm gặp và khối lượng nếu dữ liệu có giá hữu hạn. Nhân PDE với một hàm của $u$ chỉ giữ tương đương cho nghiệm trơn; không tự áp luật đã nhân để tính cùng đường sốc.

\textbf{Chi tiết tích phân đường sốc nằm trong quạt.} Nếu $s'=s/(2t)$ và $s(t_*)=x_*>0$, chia cho $s$ rồi tích phân từ $t_*$ đến $t$: $\log(s/x_*)=\frac12\log(t/t_*)$, nên $s=x_*\sqrt{t/t_*}$. Nếu quạt có tâm $x_0$, thay $s$ bằng $s-x_0$ trước khi tách biến. Khi trạng thái sát sốc thay đổi tại lần gặp tiếp theo, lập ODE mới và lấy vị trí tại lần gặp làm dữ liệu; không giữ hằng số của giai đoạn trước.
\subsection*{Lời giải áp dụng cho bài}
\textbf{Bước 1 — Tốc độ hai sốc.} $F=u^3/3$, các trạng thái 2,1,0, bước tại 0,1. $\sigma_1=(8-1)/[3(2-1)]=7/3$, $\sigma_2=(1-0)/[3(1-0)]=1/3$. Vì vậy $s_1=7t/3$, $s_2=1+t/3$.

\textbf{Bước 2 — Tìm giao.} $7t/3=1+t/3$ cho $t_*=1/2$, $x_*=7/6$. Trước đó, $u=2$ trái $s_1$, 1 giữa hai sốc, 0 phải $s_2$.

\textbf{Bước 3 — Sau gặp.} Trạng thái 1 biến mất. Sốc nối 2 với 0 có $\sigma=8/[3\cdot2]=4/3$, nên
\[s=7/6+(4/3)(t-1/2)=4t/3+1/2.\]
Với $t\ge1/2$, $u=2$ trái sốc, 0 phải sốc.

\textbf{Bước 4 — Kiểm tra entropy.} Tốc độ đặc trưng $u^2$ cho $1<7/3<4$, $0<1/3<1$, $0<4/3<4$. Đường mới đi qua điểm gặp và các trạng thái hằng giải PDE.

\clearpage\noindent\texttt{\detokenize{BT_1_6_3_10}}\par
\section*{Lời giải bổ sung — Bài tập 10 · Mục 1.6.3}

\subsection*{Kiến thức và trình tự giải}
\textbf{Mục tiêu.} Với dữ liệu gián đoạn hoặc khi đặc trưng giao nhau, tìm nghiệm yếu thỏa entropy của luật bảo toàn $u_t+(F(u))_x=0$.

\textbf{Bước 1 — Xác định thông lượng và đặc trưng.} Nếu PDE là $u_t+a(u)u_x=0$, tích phân $F'=a$. Trong vùng trơn, $u=u_0(\xi)$ và $x=\xi+F'(u_0(\xi))t$. Tính riêng tốc độ cho mỗi đoạn dữ liệu. Burgers dùng $F=u^2/2$; $u_t+u^2u_x=0$ dùng $F=u^3/3$.

\textbf{Bước 2 — Nhận diện sóng tại mỗi bước nhảy.} Gọi $u_L,u_R$ là giá trị trái, phải. So sánh $F'(u_L)$ và $F'(u_R)$. Với $F$ lồi trên khoảng trạng thái, $u_L>u_R$ tạo sốc, $u_L<u_R$ tạo quạt giãn; với $F$ lõm, thứ tự trạng thái đảo lại. Nếu độ lồi đổi dấu, phải xét thêm cấu trúc entropy, không dùng quy tắc lồi trên toàn khoảng.

\textbf{Bước 3 — Tính đường sốc từ bảo toàn.} Với $x=s(t)$,

\[
s'(t)=\frac{F(u_L)-F(u_R)}{u_L-u_R}.
\]

Nếu hai trạng thái hằng, tốc độ $\sigma$ hằng, nên $s=x_0+\sigma t$. Kiểm tra entropy bằng $F'(u_R)<\sigma<F'(u_L)$ đối với sốc nén có hai trạng thái phân biệt trong các bài này. Giá trị ngay trên đường sốc không ảnh hưởng nghiệm yếu.

\textbf{Bước 4 — Tính quạt giãn.} Từ $(x_0,0)$, đặt $z=(x-x_0)/t$, $u=V(z)$. Thế PDE cho $(F'(V)-z)V'=0$. Trong nhánh biến thiên, $F'(V)=z$, nên giải phương trình này tìm $V$. Mép quạt là $x=x_0+F'(u_L)t$ và $x=x_0+F'(u_R)t$, theo thứ tự tốc độ tăng; ghép liên tục với hai trạng thái ngoài. Chọn căn đúng dấu theo khoảng giá trị của dữ liệu.

\textbf{Bước 5 — Tìm mọi thời điểm tương tác.} Giải phương trình giao của hai sốc, sốc với mép quạt, hoặc quạt với quạt. Trước thời điểm gặp, ghép các sóng riêng. Sau khi gặp, tính lại hai trạng thái sát đường sốc. Nếu một phía nằm trong quạt, trạng thái đó phụ thuộc $s(t)$ và $t$, nên Rankine–Hugoniot trở thành ODE mới. Dùng $s(t_*)=x_*$ để tìm hằng số, rồi tiếp tục tìm lần tương tác sau.

\textbf{Bước 6 — Viết nghiệm theo từng khoảng thời gian.} Với mỗi khoảng, ghi rõ các miền $x$, thứ tự đường phân cách và công thức của $u$. Trên nửa trục, kiểm tra chiều truyền từ biên vào miền và thế điều kiện biên. Với dữ liệu trơn từng đoạn, tính $\partial x/\partial\xi$ để tìm thời điểm đầu tiên mất nghiệm cổ điển.

\textbf{Bước 7 — Kiểm tra.} Trong mỗi vùng trơn, thay vào PDE. Trên mỗi sốc, kiểm tra Rankine–Hugoniot và entropy; ở mép quạt, kiểm tra liên tục. Đối chiếu dữ liệu khi $t\downarrow0$, sự liên tục của vị trí sóng qua thời điểm gặp và khối lượng nếu dữ liệu có giá hữu hạn. Nhân PDE với một hàm của $u$ chỉ giữ tương đương cho nghiệm trơn; không tự áp luật đã nhân để tính cùng đường sốc.

\textbf{Chi tiết tích phân đường sốc nằm trong quạt.} Nếu $s'=s/(2t)$ và $s(t_*)=x_*>0$, chia cho $s$ rồi tích phân từ $t_*$ đến $t$: $\log(s/x_*)=\frac12\log(t/t_*)$, nên $s=x_*\sqrt{t/t_*}$. Nếu quạt có tâm $x_0$, thay $s$ bằng $s-x_0$ trước khi tách biến. Khi trạng thái sát sốc thay đổi tại lần gặp tiếp theo, lập ODE mới và lấy vị trí tại lần gặp làm dữ liệu; không giữ hằng số của giai đoạn trước.
\subsection*{Lời giải áp dụng cho bài}
\textbf{Bước 1 — Sóng ban đầu.} $F=u^3/3$, xung 1 trên $(-1,1)$, 0 ngoài. Tại $-1$, bước $0\to1$ tạo quạt với $u^2=(x+1)/t$, nên $u=\sqrt{(x+1)/t}$ trên $-1<x<-1+t$. Tại 1, bước $1\to0$ tạo sốc $s_1=1+t/3$.

\textbf{Bước 2 — Thời điểm gặp.} $-1+t=1+t/3$ cho $t_*=3$, $x_*=2$. Khi $0<t<3$, $u=0$ ở $x\le-1$, bằng căn trong quạt, bằng 1 trên $-1+t\le x<s_1$, bằng 0 sau sốc.

\textbf{Bước 3 — Sốc khi phía trái là quạt.} $u_L=\sqrt{(s+1)/t}$, $u_R=0$, nên $s'=u_L^2/3=(s+1)/(3t)$. Tích phân:
\[\log(s+1)=\tfrac13\log t+\text{hằng},\quad s+1=Ct^{1/3}.\]
Dùng $s(3)=2$ cho $C=3^{2/3}$. Với $t\ge3$, $s=3^{2/3}t^{1/3}-1$, $u=\sqrt{(x+1)/t}$ trên $-1<x<s$, bằng 0 ngoài.

\textbf{Bước 4 — Kiểm tra.} Entropy $0<s'=u_L^2/3<u_L^2$. Khối lượng sau gặp là $2(s+1)^{3/2}/(3\sqrt t)=2$, bằng chiều dài xung ban đầu. Đường sốc liên tục ở $t=3$.

\textbf{Chi tiết ODE và khối lượng sau gặp.} Đặt $q=s+1>0$, có $q'=q/(3t)$, nên $q'/q=1/(3t)$. Tích phân từ $3$ đến $t$ cho $\log(q/3)=\frac13\log(t/3)$, tức $q=3^{2/3}t^{1/3}$. Khối lượng là
\[\int_{-1}^{s}\sqrt{\frac{x+1}{t}}dx
=\frac1{\sqrt t}\left[\frac23(x+1)^{3/2}\right]_{-1}^{s}
=\frac{2q^{3/2}}{3\sqrt t}=2.\]

\clearpage\noindent\texttt{\detokenize{BT_1_6_3_11}}\par
\section*{Lời giải bổ sung — Bài tập 11 · Mục 1.6.3}

\subsection*{Kiến thức và trình tự giải}
\textbf{Mục tiêu.} Với dữ liệu gián đoạn hoặc khi đặc trưng giao nhau, tìm nghiệm yếu thỏa entropy của luật bảo toàn $u_t+(F(u))_x=0$.

\textbf{Bước 1 — Xác định thông lượng và đặc trưng.} Nếu PDE là $u_t+a(u)u_x=0$, tích phân $F'=a$. Trong vùng trơn, $u=u_0(\xi)$ và $x=\xi+F'(u_0(\xi))t$. Tính riêng tốc độ cho mỗi đoạn dữ liệu. Burgers dùng $F=u^2/2$; $u_t+u^2u_x=0$ dùng $F=u^3/3$.

\textbf{Bước 2 — Nhận diện sóng tại mỗi bước nhảy.} Gọi $u_L,u_R$ là giá trị trái, phải. So sánh $F'(u_L)$ và $F'(u_R)$. Với $F$ lồi trên khoảng trạng thái, $u_L>u_R$ tạo sốc, $u_L<u_R$ tạo quạt giãn; với $F$ lõm, thứ tự trạng thái đảo lại. Nếu độ lồi đổi dấu, phải xét thêm cấu trúc entropy, không dùng quy tắc lồi trên toàn khoảng.

\textbf{Bước 3 — Tính đường sốc từ bảo toàn.} Với $x=s(t)$,

\[
s'(t)=\frac{F(u_L)-F(u_R)}{u_L-u_R}.
\]

Nếu hai trạng thái hằng, tốc độ $\sigma$ hằng, nên $s=x_0+\sigma t$. Kiểm tra entropy bằng $F'(u_R)<\sigma<F'(u_L)$ đối với sốc nén có hai trạng thái phân biệt trong các bài này. Giá trị ngay trên đường sốc không ảnh hưởng nghiệm yếu.

\textbf{Bước 4 — Tính quạt giãn.} Từ $(x_0,0)$, đặt $z=(x-x_0)/t$, $u=V(z)$. Thế PDE cho $(F'(V)-z)V'=0$. Trong nhánh biến thiên, $F'(V)=z$, nên giải phương trình này tìm $V$. Mép quạt là $x=x_0+F'(u_L)t$ và $x=x_0+F'(u_R)t$, theo thứ tự tốc độ tăng; ghép liên tục với hai trạng thái ngoài. Chọn căn đúng dấu theo khoảng giá trị của dữ liệu.

\textbf{Bước 5 — Tìm mọi thời điểm tương tác.} Giải phương trình giao của hai sốc, sốc với mép quạt, hoặc quạt với quạt. Trước thời điểm gặp, ghép các sóng riêng. Sau khi gặp, tính lại hai trạng thái sát đường sốc. Nếu một phía nằm trong quạt, trạng thái đó phụ thuộc $s(t)$ và $t$, nên Rankine–Hugoniot trở thành ODE mới. Dùng $s(t_*)=x_*$ để tìm hằng số, rồi tiếp tục tìm lần tương tác sau.

\textbf{Bước 6 — Viết nghiệm theo từng khoảng thời gian.} Với mỗi khoảng, ghi rõ các miền $x$, thứ tự đường phân cách và công thức của $u$. Trên nửa trục, kiểm tra chiều truyền từ biên vào miền và thế điều kiện biên. Với dữ liệu trơn từng đoạn, tính $\partial x/\partial\xi$ để tìm thời điểm đầu tiên mất nghiệm cổ điển.

\textbf{Bước 7 — Kiểm tra.} Trong mỗi vùng trơn, thay vào PDE. Trên mỗi sốc, kiểm tra Rankine–Hugoniot và entropy; ở mép quạt, kiểm tra liên tục. Đối chiếu dữ liệu khi $t\downarrow0$, sự liên tục của vị trí sóng qua thời điểm gặp và khối lượng nếu dữ liệu có giá hữu hạn. Nhân PDE với một hàm của $u$ chỉ giữ tương đương cho nghiệm trơn; không tự áp luật đã nhân để tính cùng đường sốc.

\textbf{Chi tiết tích phân đường sốc nằm trong quạt.} Nếu $s'=s/(2t)$ và $s(t_*)=x_*>0$, chia cho $s$ rồi tích phân từ $t_*$ đến $t$: $\log(s/x_*)=\frac12\log(t/t_*)$, nên $s=x_*\sqrt{t/t_*}$. Nếu quạt có tâm $x_0$, thay $s$ bằng $s-x_0$ trước khi tách biến. Khi trạng thái sát sốc thay đổi tại lần gặp tiếp theo, lập ODE mới và lấy vị trí tại lần gặp làm dữ liệu; không giữ hằng số của giai đoạn trước.
\subsection*{Lời giải áp dụng cho bài}
\textbf{Bước 1 — Sóng trước gặp.} $F=u-u^2/2$ lõm, $F'=1-u$. Bước $1\to2$ ở 0 tạo sốc tốc độ $1-(1+2)/2=-1/2$, $s_1=-t/2$. Bước $2\to1$ ở 1 tạo quạt có mép $1-t,1$ và nhánh $u=1+(1-x)/t$.

\textbf{Bước 2 — Thời điểm gặp và nghiệm trước gặp.} $-t/2=1-t$ cho $t_*=2,x_*=-1$. Với $0<t<2$,
\[u=\begin{cases}1,&x<s_1,\\2,&s_1<x\le1-t,\\1+(1-x)/t,&1-t<x<1,\\1,&x\ge1.\end{cases}\]

\textbf{Bước 3 — Sốc trong quạt.} Sau gặp, trạng thái trái là 1, phải là $1+(1-s)/t$. Tốc độ
\[s'=1-\frac{u_L+u_R}{2}=\frac{s-1}{2t}.\]
Tích phân $(1-s)=C\sqrt t$, dùng $s(2)=-1$ cho $C=\sqrt2$, nên $s=1-\sqrt{2t}$. Với $t\ge2$, $u=1$ khi $x<s$, bằng $1+(1-x)/t$ khi $s<x<1$, bằng 1 khi $x\ge1$.

\textbf{Bước 4 — Kiểm tra.} Ở sốc, tốc độ phải $-(1-s)/t$, trái 0, còn $s'=-(1-s)/(2t)$ nằm giữa. Mép tại 1 liên tục, vị trí sốc ghép ở 2. Khối lượng phần dư $u-1$ sau gặp bằng $(1-s)^2/(2t)=1$, đúng xung ban đầu.

\textbf{Chi tiết tích phân sau tương tác.} Đặt $q=1-s>0$. Từ $s'=(s-1)/(2t)$ suy $q'=q/(2t)$. Tích phân $q'/q=1/(2t)$ từ $2$ cho $\log(q/2)=\frac12\log(t/2)$; vậy $q=\sqrt{2t}$, $s=1-\sqrt{2t}$. Phần dư khối lượng là $\int_s^1(1-x)dx/t=(1-s)^2/(2t)=1$.

\clearpage\noindent\texttt{\detokenize{BT_1_6_3_12}}\par
\section*{Lời giải bổ sung — Bài tập 12 · Mục 1.6.3}

\subsection*{Kiến thức và trình tự giải}
\textbf{Mục tiêu.} Với dữ liệu gián đoạn hoặc khi đặc trưng giao nhau, tìm nghiệm yếu thỏa entropy của luật bảo toàn $u_t+(F(u))_x=0$.

\textbf{Bước 1 — Xác định thông lượng và đặc trưng.} Nếu PDE là $u_t+a(u)u_x=0$, tích phân $F'=a$. Trong vùng trơn, $u=u_0(\xi)$ và $x=\xi+F'(u_0(\xi))t$. Tính riêng tốc độ cho mỗi đoạn dữ liệu. Burgers dùng $F=u^2/2$; $u_t+u^2u_x=0$ dùng $F=u^3/3$.

\textbf{Bước 2 — Nhận diện sóng tại mỗi bước nhảy.} Gọi $u_L,u_R$ là giá trị trái, phải. So sánh $F'(u_L)$ và $F'(u_R)$. Với $F$ lồi trên khoảng trạng thái, $u_L>u_R$ tạo sốc, $u_L<u_R$ tạo quạt giãn; với $F$ lõm, thứ tự trạng thái đảo lại. Nếu độ lồi đổi dấu, phải xét thêm cấu trúc entropy, không dùng quy tắc lồi trên toàn khoảng.

\textbf{Bước 3 — Tính đường sốc từ bảo toàn.} Với $x=s(t)$,

\[
s'(t)=\frac{F(u_L)-F(u_R)}{u_L-u_R}.
\]

Nếu hai trạng thái hằng, tốc độ $\sigma$ hằng, nên $s=x_0+\sigma t$. Kiểm tra entropy bằng $F'(u_R)<\sigma<F'(u_L)$ đối với sốc nén có hai trạng thái phân biệt trong các bài này. Giá trị ngay trên đường sốc không ảnh hưởng nghiệm yếu.

\textbf{Bước 4 — Tính quạt giãn.} Từ $(x_0,0)$, đặt $z=(x-x_0)/t$, $u=V(z)$. Thế PDE cho $(F'(V)-z)V'=0$. Trong nhánh biến thiên, $F'(V)=z$, nên giải phương trình này tìm $V$. Mép quạt là $x=x_0+F'(u_L)t$ và $x=x_0+F'(u_R)t$, theo thứ tự tốc độ tăng; ghép liên tục với hai trạng thái ngoài. Chọn căn đúng dấu theo khoảng giá trị của dữ liệu.

\textbf{Bước 5 — Tìm mọi thời điểm tương tác.} Giải phương trình giao của hai sốc, sốc với mép quạt, hoặc quạt với quạt. Trước thời điểm gặp, ghép các sóng riêng. Sau khi gặp, tính lại hai trạng thái sát đường sốc. Nếu một phía nằm trong quạt, trạng thái đó phụ thuộc $s(t)$ và $t$, nên Rankine–Hugoniot trở thành ODE mới. Dùng $s(t_*)=x_*$ để tìm hằng số, rồi tiếp tục tìm lần tương tác sau.

\textbf{Bước 6 — Viết nghiệm theo từng khoảng thời gian.} Với mỗi khoảng, ghi rõ các miền $x$, thứ tự đường phân cách và công thức của $u$. Trên nửa trục, kiểm tra chiều truyền từ biên vào miền và thế điều kiện biên. Với dữ liệu trơn từng đoạn, tính $\partial x/\partial\xi$ để tìm thời điểm đầu tiên mất nghiệm cổ điển.

\textbf{Bước 7 — Kiểm tra.} Trong mỗi vùng trơn, thay vào PDE. Trên mỗi sốc, kiểm tra Rankine–Hugoniot và entropy; ở mép quạt, kiểm tra liên tục. Đối chiếu dữ liệu khi $t\downarrow0$, sự liên tục của vị trí sóng qua thời điểm gặp và khối lượng nếu dữ liệu có giá hữu hạn. Nhân PDE với một hàm của $u$ chỉ giữ tương đương cho nghiệm trơn; không tự áp luật đã nhân để tính cùng đường sốc.

\textbf{Chi tiết tích phân đường sốc nằm trong quạt.} Nếu $s'=s/(2t)$ và $s(t_*)=x_*>0$, chia cho $s$ rồi tích phân từ $t_*$ đến $t$: $\log(s/x_*)=\frac12\log(t/t_*)$, nên $s=x_*\sqrt{t/t_*}$. Nếu quạt có tâm $x_0$, thay $s$ bằng $s-x_0$ trước khi tách biến. Khi trạng thái sát sốc thay đổi tại lần gặp tiếp theo, lập ODE mới và lấy vị trí tại lần gặp làm dữ liệu; không giữ hằng số của giai đoạn trước.
\subsection*{Lời giải áp dụng cho bài}
\textbf{Bước 1 — Sóng trong nửa trục dương.} Burgers có giá trị biên 1, trạng thái đầu 2 trên $(0,1)$ rồi 0. Góc biên tạo quạt nối 1 đến 2: mép $t,2t$, $u=x/t$ trong quạt. Bước tại 1 tạo sốc 2 xuống 0 với $s_1=1+t$.

\textbf{Bước 2 — Tương tác thứ nhất.} $2t=1+t$ cho $t_1=1,x_1=2$. Khi $0<t<1$, $u=1$ trên $0<x\le t$, $u=x/t$ trên $t<x<2t$, $u=2$ trên $2t\le x<s_1$, 0 phía phải.

\textbf{Bước 3 — Sốc xuyên qua quạt.} Với $u_L=s/t,u_R=0$, $s'=s/(2t)$ cho $s=C\sqrt t$. Dùng $s(1)=2$, được $s=2\sqrt t$. Sốc gặp mép trái $x=t$ khi $2\sqrt t=t$, tức $t_2=4,x_2=4$. Với $1\le t<4$, $u=1$ trên $0<x\le t$, $u=x/t$ trên $t<x<s$, 0 phía phải.

\textbf{Bước 4 — Sau tương tác thứ hai.} Quạt đã bị sốc hấp thụ hết; trái sốc là 1, phải 0 nên tốc độ $1/2$. Từ $s(4)=4$, $s=4+(t-4)/2=t/2+2$. Với $t\ge4$, $u=1$ trên $0<x<s$, 0 phía phải.

\textbf{Bước 5 — Kiểm tra.} Hai vị trí ghép đúng tại $t=1,4$. Trong quạt PDE đúng; từng sốc có tốc độ nửa trạng thái trái nên thỏa entropy. Biên luôn bằng 1; khối lượng tăng theo thông lượng vào $F(1)=1/2$, phù hợp nhánh cuối có độ dài $t/2+2$.

\textbf{Chi tiết hai lần tích phân và kiểm tra khối lượng.} Trong quạt, $s'=s/(2t)$ cho $\log(s/2)=\frac12\log t$ bởi $s(1)=2$, nên $s=2\sqrt t$. Giải $2\sqrt t=t$ với $t>0$: chia $\sqrt t$ được $\sqrt t=2$, tức $t=4$ (loại $t=0$ ngoài giai đoạn). Sau đó $s'=1/2$ cho $s-4=(t-4)/2$. Khối lượng ở ba giai đoạn:
\[\begin{aligned}
M&=t+\int_t^{2t}\frac xtdx+2(1+t-2t)=2+t/2 &&(0<t<1),\\
M&=t+\int_t^{2\sqrt t}\frac xtdx=t+\frac{4t-t^2}{2t}=2+t/2 &&(1\le t<4),\\
M&=\int_0^{t/2+2}1\,dx=2+t/2 &&(t\ge4).
\end{aligned}\]
Do đó $M'=1/2$ ở mọi giai đoạn, đúng thông lượng vào $F(1)$.

\clearpage\noindent\texttt{\detokenize{NC_1_7_01}}\par
\section*{Lời giải bổ sung — Bài nâng cao 1 · Mục 1.7}

\subsection*{Kiến thức và trình tự giải}
\textbf{Mục tiêu.} Với dữ liệu gián đoạn hoặc khi đặc trưng giao nhau, tìm nghiệm yếu thỏa entropy của luật bảo toàn $u_t+(F(u))_x=0$.

\textbf{Bước 1 — Xác định thông lượng và đặc trưng.} Nếu PDE là $u_t+a(u)u_x=0$, tích phân $F'=a$. Trong vùng trơn, $u=u_0(\xi)$ và $x=\xi+F'(u_0(\xi))t$. Tính riêng tốc độ cho mỗi đoạn dữ liệu. Burgers dùng $F=u^2/2$; $u_t+u^2u_x=0$ dùng $F=u^3/3$.

\textbf{Bước 2 — Nhận diện sóng tại mỗi bước nhảy.} Gọi $u_L,u_R$ là giá trị trái, phải. So sánh $F'(u_L)$ và $F'(u_R)$. Với $F$ lồi trên khoảng trạng thái, $u_L>u_R$ tạo sốc, $u_L<u_R$ tạo quạt giãn; với $F$ lõm, thứ tự trạng thái đảo lại. Nếu độ lồi đổi dấu, phải xét thêm cấu trúc entropy, không dùng quy tắc lồi trên toàn khoảng.

\textbf{Bước 3 — Tính đường sốc từ bảo toàn.} Với $x=s(t)$,

\[
s'(t)=\frac{F(u_L)-F(u_R)}{u_L-u_R}.
\]

Nếu hai trạng thái hằng, tốc độ $\sigma$ hằng, nên $s=x_0+\sigma t$. Kiểm tra entropy bằng $F'(u_R)<\sigma<F'(u_L)$ đối với sốc nén có hai trạng thái phân biệt trong các bài này. Giá trị ngay trên đường sốc không ảnh hưởng nghiệm yếu.

\textbf{Bước 4 — Tính quạt giãn.} Từ $(x_0,0)$, đặt $z=(x-x_0)/t$, $u=V(z)$. Thế PDE cho $(F'(V)-z)V'=0$. Trong nhánh biến thiên, $F'(V)=z$, nên giải phương trình này tìm $V$. Mép quạt là $x=x_0+F'(u_L)t$ và $x=x_0+F'(u_R)t$, theo thứ tự tốc độ tăng; ghép liên tục với hai trạng thái ngoài. Chọn căn đúng dấu theo khoảng giá trị của dữ liệu.

\textbf{Bước 5 — Tìm mọi thời điểm tương tác.} Giải phương trình giao của hai sốc, sốc với mép quạt, hoặc quạt với quạt. Trước thời điểm gặp, ghép các sóng riêng. Sau khi gặp, tính lại hai trạng thái sát đường sốc. Nếu một phía nằm trong quạt, trạng thái đó phụ thuộc $s(t)$ và $t$, nên Rankine–Hugoniot trở thành ODE mới. Dùng $s(t_*)=x_*$ để tìm hằng số, rồi tiếp tục tìm lần tương tác sau.

\textbf{Bước 6 — Viết nghiệm theo từng khoảng thời gian.} Với mỗi khoảng, ghi rõ các miền $x$, thứ tự đường phân cách và công thức của $u$. Trên nửa trục, kiểm tra chiều truyền từ biên vào miền và thế điều kiện biên. Với dữ liệu trơn từng đoạn, tính $\partial x/\partial\xi$ để tìm thời điểm đầu tiên mất nghiệm cổ điển.

\textbf{Bước 7 — Kiểm tra.} Trong mỗi vùng trơn, thay vào PDE. Trên mỗi sốc, kiểm tra Rankine–Hugoniot và entropy; ở mép quạt, kiểm tra liên tục. Đối chiếu dữ liệu khi $t\downarrow0$, sự liên tục của vị trí sóng qua thời điểm gặp và khối lượng nếu dữ liệu có giá hữu hạn. Nhân PDE với một hàm của $u$ chỉ giữ tương đương cho nghiệm trơn; không tự áp luật đã nhân để tính cùng đường sốc.

\textbf{Chi tiết tích phân đường sốc nằm trong quạt.} Nếu $s'=s/(2t)$ và $s(t_*)=x_*>0$, chia cho $s$ rồi tích phân từ $t_*$ đến $t$: $\log(s/x_*)=\frac12\log(t/t_*)$, nên $s=x_*\sqrt{t/t_*}$. Nếu quạt có tâm $x_0$, thay $s$ bằng $s-x_0$ trước khi tách biến. Khi trạng thái sát sốc thay đổi tại lần gặp tiếp theo, lập ODE mới và lấy vị trí tại lần gặp làm dữ liệu; không giữ hằng số của giai đoạn trước.
\subsection*{Lời giải áp dụng cho bài}
\textbf{Bước 1 — Đặc trưng ở đoạn dốc.} Burgers có $u_0=1$ bên trái 0, $u_0=1-\xi$ trên $(0,1)$, 0 bên phải 1. Trên đoạn dốc, $X=\xi+(1-\xi)t=t+\xi(1-t)$, $U=1-\xi$.

\textbf{Bước 2 — Khử trước thời điểm vỡ.} Với $t<1$, $X_\xi=1-t>0$, $\xi=(x-t)/(1-t)$, nên $u=(1-x)/(1-t)$ trên $t<x<1$. Nghiệm bằng 1 khi $x\le t$, bằng 0 khi $x\ge1$. Tại $t=1$ cả đoạn hội tụ ở $x=1$.

\textbf{Bước 3 — Sau thời điểm vỡ.} Sốc nối 1 và 0 của luật Burgers gốc có tốc độ $1/2$, đi từ $(1,1)$, nên $s=1+(t-1)/2=(t+1)/2$. $u=1$ bên trái, 0 bên phải sốc; $0<1/2<1$ cho entropy.

\textbf{Bước 4 — Giải thích phép nhân $u^\alpha$.} Trong vùng trơn và nơi phép tính lũy thừa hợp lệ, nhân PDE rồi lấy nguyên hàm cho
\[\partial_t\frac{u^{\alpha+1}}{\alpha+1}+\partial_x\frac{u^{\alpha+2}}{\alpha+2}=0.\]
Với $\alpha\ge0$, sốc $1\to0$ của luật này có tốc độ
\[\sigma_\alpha=\frac{1/(\alpha+2)}{1/(\alpha+1)}=\frac{\alpha+1}{\alpha+2}.\]
Nó chỉ bằng $1/2$ khi $\alpha=0$. Vì vậy phép nhân giữ tương đương cho nghiệm trơn, không giữ cùng điều kiện sốc của nghiệm yếu. Không thể yêu cầu cùng đường sốc thỏa mọi luật đã nhân; phải dùng luật gốc và entropy.

\clearpage\noindent\texttt{\detokenize{NC_1_7_02}}\par
\section*{Lời giải bổ sung — Bài nâng cao 2 · Mục 1.7}

\subsection*{Kiến thức và trình tự giải}
\textbf{Mục tiêu.} Với dữ liệu gián đoạn hoặc khi đặc trưng giao nhau, tìm nghiệm yếu thỏa entropy của luật bảo toàn $u_t+(F(u))_x=0$.

\textbf{Bước 1 — Xác định thông lượng và đặc trưng.} Nếu PDE là $u_t+a(u)u_x=0$, tích phân $F'=a$. Trong vùng trơn, $u=u_0(\xi)$ và $x=\xi+F'(u_0(\xi))t$. Tính riêng tốc độ cho mỗi đoạn dữ liệu. Burgers dùng $F=u^2/2$; $u_t+u^2u_x=0$ dùng $F=u^3/3$.

\textbf{Bước 2 — Nhận diện sóng tại mỗi bước nhảy.} Gọi $u_L,u_R$ là giá trị trái, phải. So sánh $F'(u_L)$ và $F'(u_R)$. Với $F$ lồi trên khoảng trạng thái, $u_L>u_R$ tạo sốc, $u_L<u_R$ tạo quạt giãn; với $F$ lõm, thứ tự trạng thái đảo lại. Nếu độ lồi đổi dấu, phải xét thêm cấu trúc entropy, không dùng quy tắc lồi trên toàn khoảng.

\textbf{Bước 3 — Tính đường sốc từ bảo toàn.} Với $x=s(t)$,

\[
s'(t)=\frac{F(u_L)-F(u_R)}{u_L-u_R}.
\]

Nếu hai trạng thái hằng, tốc độ $\sigma$ hằng, nên $s=x_0+\sigma t$. Kiểm tra entropy bằng $F'(u_R)<\sigma<F'(u_L)$ đối với sốc nén có hai trạng thái phân biệt trong các bài này. Giá trị ngay trên đường sốc không ảnh hưởng nghiệm yếu.

\textbf{Bước 4 — Tính quạt giãn.} Từ $(x_0,0)$, đặt $z=(x-x_0)/t$, $u=V(z)$. Thế PDE cho $(F'(V)-z)V'=0$. Trong nhánh biến thiên, $F'(V)=z$, nên giải phương trình này tìm $V$. Mép quạt là $x=x_0+F'(u_L)t$ và $x=x_0+F'(u_R)t$, theo thứ tự tốc độ tăng; ghép liên tục với hai trạng thái ngoài. Chọn căn đúng dấu theo khoảng giá trị của dữ liệu.

\textbf{Bước 5 — Tìm mọi thời điểm tương tác.} Giải phương trình giao của hai sốc, sốc với mép quạt, hoặc quạt với quạt. Trước thời điểm gặp, ghép các sóng riêng. Sau khi gặp, tính lại hai trạng thái sát đường sốc. Nếu một phía nằm trong quạt, trạng thái đó phụ thuộc $s(t)$ và $t$, nên Rankine–Hugoniot trở thành ODE mới. Dùng $s(t_*)=x_*$ để tìm hằng số, rồi tiếp tục tìm lần tương tác sau.

\textbf{Bước 6 — Viết nghiệm theo từng khoảng thời gian.} Với mỗi khoảng, ghi rõ các miền $x$, thứ tự đường phân cách và công thức của $u$. Trên nửa trục, kiểm tra chiều truyền từ biên vào miền và thế điều kiện biên. Với dữ liệu trơn từng đoạn, tính $\partial x/\partial\xi$ để tìm thời điểm đầu tiên mất nghiệm cổ điển.

\textbf{Bước 7 — Kiểm tra.} Trong mỗi vùng trơn, thay vào PDE. Trên mỗi sốc, kiểm tra Rankine–Hugoniot và entropy; ở mép quạt, kiểm tra liên tục. Đối chiếu dữ liệu khi $t\downarrow0$, sự liên tục của vị trí sóng qua thời điểm gặp và khối lượng nếu dữ liệu có giá hữu hạn. Nhân PDE với một hàm của $u$ chỉ giữ tương đương cho nghiệm trơn; không tự áp luật đã nhân để tính cùng đường sốc.

\textbf{Chi tiết tích phân đường sốc nằm trong quạt.} Nếu $s'=s/(2t)$ và $s(t_*)=x_*>0$, chia cho $s$ rồi tích phân từ $t_*$ đến $t$: $\log(s/x_*)=\frac12\log(t/t_*)$, nên $s=x_*\sqrt{t/t_*}$. Nếu quạt có tâm $x_0$, thay $s$ bằng $s-x_0$ trước khi tách biến. Khi trạng thái sát sốc thay đổi tại lần gặp tiếp theo, lập ODE mới và lấy vị trí tại lần gặp làm dữ liệu; không giữ hằng số của giai đoạn trước.
\subsection*{Lời giải áp dụng cho bài}
\textbf{Bước 1 — Đặc trưng và thời điểm đầu mất khả nghịch.} Burgers có $u_0=1$ trên $\xi<-1$, $u_0=\xi^2$ trên $-1<\xi<0$, 0 trên $\xi>0$. Đoạn biến thiên cho $X=\xi+\xi^2t$, $X_\xi=1+2\xi t$. Giá trị nhỏ nhất tiến về 0 tại $t=1/2,\xi=-1$, nên sốc xuất hiện từ $(-1/2,1/2)$.

\textbf{Bước 2 — Nghiệm trước sốc.} Giải $t\xi^2+\xi-x=0$ và chọn nhánh về $\xi=x$ khi $t\downarrow0$:
\[\xi(x,t)=\frac{-1+\sqrt{1+4tx}}{2t}.\]
Với $0<t\le1/2$, $u=1$ khi $x<-1+t$, $u=\xi(x,t)^2$ khi $-1+t<x<0$, $u=0$ khi $x>0$.

\textbf{Bước 3 — Sốc trong đoạn biến thiên.} Gọi $\xi_s\in(-1,0)$ là nhãn sát phải sốc, $s=\xi_s+\xi_s^2t$. Trái là 1, phải $\xi_s^2$, nên $s'=(1+\xi_s^2)/2$. Đạo hàm biểu thức $s$:
\[(1+2t\xi_s)\xi_s'=(1-\xi_s^2)/2.\]
Đổi coi $t$ là hàm của $\xi_s$:
\[\frac{dt}{d\xi_s}-\frac{4\xi_s}{1-\xi_s^2}t=\frac2{1-\xi_s^2}.\]
Thừa số tích phân $(1-\xi_s^2)^2$ cho
\[\frac{d}{d\xi_s}[(1-\xi_s^2)^2t]=2(1-\xi_s^2).\]
Tích phân và dùng giới hạn ở $\xi_s=-1,t=1/2$ cho
\[t=\frac{2(2-\xi_s)}{3(1-\xi_s)^2},\qquad s=\xi_s+\xi_s^2t.\]
Vế phải tăng từ $1/2$ đến $4/3$ khi $\xi_s$ tăng từ $-1$ đến 0, nên xác định nhãn duy nhất. Với $1/2<t<4/3$, $u=1$ trái $s$, $u=\xi(x,t)^2$ giữa $s$ và 0, $u=0$ phải 0; chọn nhánh căn như bước 2.

\textbf{Bước 4 — Sau sốc tới $x=0$.} Tại $\xi_s=0$, $t=4/3,s=0$. Sau đó sốc nối 1 với 0, $s' =1/2$, nên $s=(t-4/3)/2$; $u=1$ trái, 0 phải. Mọi vị trí ghép liên tục; trong giai đoạn giữa, $\xi_s^2<(1+\xi_s^2)/2<1$ cho entropy.

\textbf{Chi tiết tích phân và khử nhân tử.} Gọi $z=\xi_s$. Tích phân $( (1-z^2)^2t)'=2(1-z^2)$ theo $z$ cho $(1-z^2)^2t=2z-2z^3/3+C$. Khi $z\to-1$, vế trái về $0$, nên $C=4/3$. Phân tích
\[2z-\tfrac23z^3+\tfrac43=\tfrac23(2-z)(z+1)^2,\quad
(1-z^2)^2=(1-z)^2(1+z)^2.\]
Trên $-1<z<0$, khử $(1+z)^2$ được $t=2(2-z)/[3(1-z)^2]$. Đạo hàm $dt/dz=2(3-z)/[3(1-z)^3]>0$, và giới hạn tại $z=-1,0$ là $1/2,4/3$, chứng minh nhãn xác định duy nhất.

\clearpage\noindent\texttt{\detokenize{NC_1_7_03}}\par
\section*{Lời giải bổ sung — Bài nâng cao 3 · Mục 1.7}

\subsection*{Kiến thức và trình tự giải}
\textbf{Mục tiêu.} Với dữ liệu gián đoạn hoặc khi đặc trưng giao nhau, tìm nghiệm yếu thỏa entropy của luật bảo toàn $u_t+(F(u))_x=0$.

\textbf{Bước 1 — Xác định thông lượng và đặc trưng.} Nếu PDE là $u_t+a(u)u_x=0$, tích phân $F'=a$. Trong vùng trơn, $u=u_0(\xi)$ và $x=\xi+F'(u_0(\xi))t$. Tính riêng tốc độ cho mỗi đoạn dữ liệu. Burgers dùng $F=u^2/2$; $u_t+u^2u_x=0$ dùng $F=u^3/3$.

\textbf{Bước 2 — Nhận diện sóng tại mỗi bước nhảy.} Gọi $u_L,u_R$ là giá trị trái, phải. So sánh $F'(u_L)$ và $F'(u_R)$. Với $F$ lồi trên khoảng trạng thái, $u_L>u_R$ tạo sốc, $u_L<u_R$ tạo quạt giãn; với $F$ lõm, thứ tự trạng thái đảo lại. Nếu độ lồi đổi dấu, phải xét thêm cấu trúc entropy, không dùng quy tắc lồi trên toàn khoảng.

\textbf{Bước 3 — Tính đường sốc từ bảo toàn.} Với $x=s(t)$,

\[
s'(t)=\frac{F(u_L)-F(u_R)}{u_L-u_R}.
\]

Nếu hai trạng thái hằng, tốc độ $\sigma$ hằng, nên $s=x_0+\sigma t$. Kiểm tra entropy bằng $F'(u_R)<\sigma<F'(u_L)$ đối với sốc nén có hai trạng thái phân biệt trong các bài này. Giá trị ngay trên đường sốc không ảnh hưởng nghiệm yếu.

\textbf{Bước 4 — Tính quạt giãn.} Từ $(x_0,0)$, đặt $z=(x-x_0)/t$, $u=V(z)$. Thế PDE cho $(F'(V)-z)V'=0$. Trong nhánh biến thiên, $F'(V)=z$, nên giải phương trình này tìm $V$. Mép quạt là $x=x_0+F'(u_L)t$ và $x=x_0+F'(u_R)t$, theo thứ tự tốc độ tăng; ghép liên tục với hai trạng thái ngoài. Chọn căn đúng dấu theo khoảng giá trị của dữ liệu.

\textbf{Bước 5 — Tìm mọi thời điểm tương tác.} Giải phương trình giao của hai sốc, sốc với mép quạt, hoặc quạt với quạt. Trước thời điểm gặp, ghép các sóng riêng. Sau khi gặp, tính lại hai trạng thái sát đường sốc. Nếu một phía nằm trong quạt, trạng thái đó phụ thuộc $s(t)$ và $t$, nên Rankine–Hugoniot trở thành ODE mới. Dùng $s(t_*)=x_*$ để tìm hằng số, rồi tiếp tục tìm lần tương tác sau.

\textbf{Bước 6 — Viết nghiệm theo từng khoảng thời gian.} Với mỗi khoảng, ghi rõ các miền $x$, thứ tự đường phân cách và công thức của $u$. Trên nửa trục, kiểm tra chiều truyền từ biên vào miền và thế điều kiện biên. Với dữ liệu trơn từng đoạn, tính $\partial x/\partial\xi$ để tìm thời điểm đầu tiên mất nghiệm cổ điển.

\textbf{Bước 7 — Kiểm tra.} Trong mỗi vùng trơn, thay vào PDE. Trên mỗi sốc, kiểm tra Rankine–Hugoniot và entropy; ở mép quạt, kiểm tra liên tục. Đối chiếu dữ liệu khi $t\downarrow0$, sự liên tục của vị trí sóng qua thời điểm gặp và khối lượng nếu dữ liệu có giá hữu hạn. Nhân PDE với một hàm của $u$ chỉ giữ tương đương cho nghiệm trơn; không tự áp luật đã nhân để tính cùng đường sốc.

\textbf{Chi tiết tích phân đường sốc nằm trong quạt.} Nếu $s'=s/(2t)$ và $s(t_*)=x_*>0$, chia cho $s$ rồi tích phân từ $t_*$ đến $t$: $\log(s/x_*)=\frac12\log(t/t_*)$, nên $s=x_*\sqrt{t/t_*}$. Nếu quạt có tâm $x_0$, thay $s$ bằng $s-x_0$ trước khi tách biến. Khi trạng thái sát sốc thay đổi tại lần gặp tiếp theo, lập ODE mới và lấy vị trí tại lần gặp làm dữ liệu; không giữ hằng số của giai đoạn trước.
\subsection*{Lời giải áp dụng cho bài}
\textbf{Bước 1 — Khử hai nhánh ban đầu.} Xung tam giác Burgers có $u_0=\xi+1$ trên $(-1,0)$, $1-\xi$ trên $(0,1)$, 0 ngoài. Nhánh tăng: $X=\xi+(\xi+1)t$ cho $u=(x+1)/(t+1)$ trên $-1<x<t$. Nhánh giảm: $X=\xi+(1-\xi)t$ cho $u=(1-x)/(1-t)$ trên $t<x<1$, khi $t<1$.

\textbf{Bước 2 — Nghiệm trước vỡ.}
\[u=\begin{cases}(x+1)/(t+1),&-1<x\le t,\\(1-x)/(1-t),&t<x<1,\\0,&x\le-1\ \text{hoặc}\ x\ge1,\end{cases}\quad0<t<1.\]
Hai nhánh bằng 1 tại $x=t$; nhánh giảm có $X_\xi=1-t$, hội tụ ở $(1,1)$.

\textbf{Bước 3 — Sốc sau vỡ.} Trái sốc là $u_L=(s+1)/(t+1)$, phải 0, nên $s'=(s+1)/[2(t+1)]$. Tích phân cho $s+1=C\sqrt{t+1}$; dùng $s(1)=1$, $C=\sqrt2$. Do đó $s=\sqrt{2(t+1)}-1$ và với $t\ge1$, $u=(x+1)/(t+1)$ trên $-1<x<s$, bằng 0 ngoài.

\textbf{Bước 4 — Kiểm tra.} Trong nhánh tăng, $u_t=-(x+1)/(t+1)^2$, $uu_x=(x+1)/(t+1)^2$. Sốc thỏa bảo toàn và entropy, ghép đúng $(1,1)$. Khối lượng sau vỡ $(s+1)^2/[2(t+1)]=1$, bằng diện tích tam giác ban đầu.

\clearpage\noindent\texttt{\detokenize{NC_1_7_04}}\par
\section*{Lời giải bổ sung — Bài nâng cao 4 · Mục 1.7}

\subsection*{Kiến thức và trình tự giải}
\textbf{Mục tiêu.} Với dữ liệu gián đoạn hoặc khi đặc trưng giao nhau, tìm nghiệm yếu thỏa entropy của luật bảo toàn $u_t+(F(u))_x=0$.

\textbf{Bước 1 — Xác định thông lượng và đặc trưng.} Nếu PDE là $u_t+a(u)u_x=0$, tích phân $F'=a$. Trong vùng trơn, $u=u_0(\xi)$ và $x=\xi+F'(u_0(\xi))t$. Tính riêng tốc độ cho mỗi đoạn dữ liệu. Burgers dùng $F=u^2/2$; $u_t+u^2u_x=0$ dùng $F=u^3/3$.

\textbf{Bước 2 — Nhận diện sóng tại mỗi bước nhảy.} Gọi $u_L,u_R$ là giá trị trái, phải. So sánh $F'(u_L)$ và $F'(u_R)$. Với $F$ lồi trên khoảng trạng thái, $u_L>u_R$ tạo sốc, $u_L<u_R$ tạo quạt giãn; với $F$ lõm, thứ tự trạng thái đảo lại. Nếu độ lồi đổi dấu, phải xét thêm cấu trúc entropy, không dùng quy tắc lồi trên toàn khoảng.

\textbf{Bước 3 — Tính đường sốc từ bảo toàn.} Với $x=s(t)$,

\[
s'(t)=\frac{F(u_L)-F(u_R)}{u_L-u_R}.
\]

Nếu hai trạng thái hằng, tốc độ $\sigma$ hằng, nên $s=x_0+\sigma t$. Kiểm tra entropy bằng $F'(u_R)<\sigma<F'(u_L)$ đối với sốc nén có hai trạng thái phân biệt trong các bài này. Giá trị ngay trên đường sốc không ảnh hưởng nghiệm yếu.

\textbf{Bước 4 — Tính quạt giãn.} Từ $(x_0,0)$, đặt $z=(x-x_0)/t$, $u=V(z)$. Thế PDE cho $(F'(V)-z)V'=0$. Trong nhánh biến thiên, $F'(V)=z$, nên giải phương trình này tìm $V$. Mép quạt là $x=x_0+F'(u_L)t$ và $x=x_0+F'(u_R)t$, theo thứ tự tốc độ tăng; ghép liên tục với hai trạng thái ngoài. Chọn căn đúng dấu theo khoảng giá trị của dữ liệu.

\textbf{Bước 5 — Tìm mọi thời điểm tương tác.} Giải phương trình giao của hai sốc, sốc với mép quạt, hoặc quạt với quạt. Trước thời điểm gặp, ghép các sóng riêng. Sau khi gặp, tính lại hai trạng thái sát đường sốc. Nếu một phía nằm trong quạt, trạng thái đó phụ thuộc $s(t)$ và $t$, nên Rankine–Hugoniot trở thành ODE mới. Dùng $s(t_*)=x_*$ để tìm hằng số, rồi tiếp tục tìm lần tương tác sau.

\textbf{Bước 6 — Viết nghiệm theo từng khoảng thời gian.} Với mỗi khoảng, ghi rõ các miền $x$, thứ tự đường phân cách và công thức của $u$. Trên nửa trục, kiểm tra chiều truyền từ biên vào miền và thế điều kiện biên. Với dữ liệu trơn từng đoạn, tính $\partial x/\partial\xi$ để tìm thời điểm đầu tiên mất nghiệm cổ điển.

\textbf{Bước 7 — Kiểm tra.} Trong mỗi vùng trơn, thay vào PDE. Trên mỗi sốc, kiểm tra Rankine–Hugoniot và entropy; ở mép quạt, kiểm tra liên tục. Đối chiếu dữ liệu khi $t\downarrow0$, sự liên tục của vị trí sóng qua thời điểm gặp và khối lượng nếu dữ liệu có giá hữu hạn. Nhân PDE với một hàm của $u$ chỉ giữ tương đương cho nghiệm trơn; không tự áp luật đã nhân để tính cùng đường sốc.

\textbf{Chi tiết tích phân đường sốc nằm trong quạt.} Nếu $s'=s/(2t)$ và $s(t_*)=x_*>0$, chia cho $s$ rồi tích phân từ $t_*$ đến $t$: $\log(s/x_*)=\frac12\log(t/t_*)$, nên $s=x_*\sqrt{t/t_*}$. Nếu quạt có tâm $x_0$, thay $s$ bằng $s-x_0$ trước khi tách biến. Khi trạng thái sát sốc thay đổi tại lần gặp tiếp theo, lập ODE mới và lấy vị trí tại lần gặp làm dữ liệu; không giữ hằng số của giai đoạn trước.
\subsection*{Lời giải áp dụng cho bài}
\textbf{Bước 1 — Xác định dữ liệu hợp lệ.} Dòng cuối nguyên văn ghi $u_0=0$ khi $|x|>0$, chồng lên các đoạn khác, nên không xác định một hàm. Phần dưới giải phiên bản sửa $u_0=0$ khi $|x|>2$: các trạng thái lần lượt là 0 trên $x<-2$, 1 trên $(-2,-1)$, 4 trên $(-1,1)$, 1 trên $(1,2)$, 0 trên $x>2$.

\textbf{Bước 2 — Sóng ban đầu.} Với Burgers, hai quạt là $R_0=(x+2)/t$ từ $-2$, mép $a=-2+t$; $R_1=(x+1)/t$ từ $-1$, mép $b=-1+t,c=-1+4t$. Hai sốc 4 xuống 1, 1 xuống 0 có $s_1=1+5t/2$, $s_2=2+t/2$.

\textbf{Bước 3 — Hai sốc gặp nhau.} $s_1=s_2$ cho $t_1=1/2,x_1=9/4$. Sau gặp, sốc 4 xuống 0 có tốc độ 2, nên $s=2t+5/4$. Nó gặp mép phải $c$ của quạt $R_1$ khi $2t+5/4=-1+4t$, tức $t_2=9/8,x_2=7/2$.

\textbf{Bước 4 — Sốc đi qua quạt $R_1$.} Trái sốc $u_L=(s+1)/t$, phải 0, nên $s'=(s+1)/(2t)$. Từ $s(9/8)=7/2$, được $s=3\sqrt{2t}-1$. Sốc tới mép trái $b=-1+t$ khi $3\sqrt{2t}=t$, tức $t_3=18,x_3=17$.

\textbf{Bước 5 — Từ trạng thái 1 sang quạt $R_0$.} Sau $18$, sốc nối 1 với 0 có tốc độ $1/2$: $s=t/2+8$. Nó gặp mép phải $a=-2+t$ khi $t/2+8=-2+t$, tức $t_4=20,x_4=18$. Từ đó $u_L=(s+2)/t$, nên $s'=(s+2)/(2t)$; dùng $s(20)=18$ cho $s=\sqrt{20t}-2$.

\textbf{Bước 6 — Viết đủ các miền.} Với $0<t<1/2$,
\[u=\begin{cases}0,&x\le-2,\\R_0,&-2<x<a,\\1,&a\le x\le b,\\R_1,&b<x<c,\\4,&c\le x<s_1,\\1,&s_1<x<s_2,\\0,&x>s_2.\end{cases}\]
Với $1/2\le t<9/8$, $s=2t+5/4$,
\[u=\begin{cases}0,&x\le-2,\\R_0,&-2<x<a,\\1,&a\le x\le b,\\R_1,&b<x<c,\\4,&c\le x<s,\\0,&x>s.\end{cases}\]
Với $9/8\le t<18$, $s=3\sqrt{2t}-1$,
\[u=\begin{cases}0,&x\le-2,\\R_0,&-2<x<a,\\1,&a\le x\le b,\\R_1,&b<x<s,\\0,&x>s.\end{cases}\]
Với $18\le t<20$, $s=t/2+8$, nghiệm bằng 0 khi $x\le-2$, bằng $R_0$ khi $-2<x<a$, bằng 1 khi $a\le x<s$, bằng 0 khi $x>s$. Với $t\ge20$, $s=\sqrt{20t}-2$, chỉ còn $u=R_0$ trên $-2<x<s$, bằng 0 ngoài.

\textbf{Bước 7 — Kiểm tra.} Các đường sốc ghép liên tục ở bốn thời điểm vừa tính. Tốc độ luôn bằng trung bình hai trạng thái, nên bảo toàn và entropy đúng. Khối lượng ban đầu $1+8+1=10$; giai đoạn cuối có $\int_{-2}^s(x+2)dx/t=(s+2)^2/(2t)=10$. Kết quả thuộc phiên bản dữ liệu đã sửa và phải ghi rõ điều đó.

\textbf{Chi tiết hằng số trên các nhánh cong.} Khi sốc đi qua quạt $R_1$, đặt $q=s+1$, $q'/q=1/(2t)$, nên $q=C\sqrt t$. Dùng $q(9/8)=9/2$ cho $C=(9/2)/\sqrt{9/8}=3\sqrt2$. Sau giai đoạn trạng thái $1$, tại $t=20$ có $s=18$; trong quạt $R_0$, đặt $q=s+2$, $q'/q=1/(2t)$ và $q(20)=20$, nên $q=\sqrt{20t}$. Vì vậy nhánh cuối $s=\sqrt{20t}-2$. Các hằng số khác nhau vì sốc đi qua hai quạt có tâm lần lượt $-1$ và $-2$.

\clearpage\noindent\texttt{\detokenize{VD_2_2_01}}\par
\section*{Lời giải bổ sung — Ví dụ 1 · Mục 2.2}

\subsection*{Kiến thức và trình tự giải}
\textbf{Mục tiêu.} Tính đạo hàm khi đổi tọa độ rồi thay vào PDE. Đổi tọa độ của điểm và đổi biểu thức của hàm là hai việc cần làm cùng nhau: đặt $q=q(x)$ và $v(q)=u(x(q))$.

\textbf{Bước 1 — Kiểm tra phép đổi biến.} Tính $J=\det(\partial q/\partial x)$. Chỉ dùng các công thức trên miền $J\ne0$ và chọn một nhánh liên tục cho các góc. Tọa độ cực/trụ loại trục $r=0$; tọa độ cầu còn loại $\sin\theta=0$. Đây là suy biến của tọa độ, chưa phải kết luận rằng hàm $u$ không tồn tại ở đó.

\textbf{Bước 2 — Tính đạo hàm cấp một bằng quy tắc dây chuyền.} Nếu $b_{ik}=\partial q_k/\partial x_i$, thì


\[
u_{x_i}=\sum_k b_{ik}v_{q_k},\qquad D_i=\sum_k b_{ik}\partial_{q_k}.
\]


Với $q=Ax$, $b_{ik}=a_{ki}$ là hằng số. Với tọa độ cong, lấy nghịch đảo ma trận $\partial x/\partial q$ để tìm $b_{ik}$.

\textbf{Bước 3 — Lấy đạo hàm cấp hai cả hàm lẫn hệ số.} Áp dụng đạo hàm tích cho từng hạng:


\[
u_{x_ix_j}=\sum_{k,l}b_{ik}b_{jl}v_{q_kq_l}
+\sum_{k,l}b_{ik}(\partial_{q_k}b_{jl})v_{q_l}.
\]


Hạng thứ hai chỉ mất đi khi phép đổi biến tuyến tính. Ví dụ $\partial_x=\cos\theta\,\partial_r-(\sin\theta/r)\partial_\theta$ thì lần đạo hàm tiếp theo phải tác động lên cả $\cos\theta$ và $\sin\theta/r$.

\textbf{Bước 4 — Thu gọn toán tử.} Gom hệ số của từng $v_{q_kq_l}$ và $v_{q_k}$ trước khi cộng. Với Laplace, lấy tổng các đạo hàm $u_{x_ix_i}$. Trong hệ trực giao có hệ số độ dài $h_1,h_2,h_3$, có thể đối chiếu bằng


\[
\Delta u=\frac1{h_1h_2h_3}\sum_{k=1}^3
\partial_{q_k}\left(\frac{h_1h_2h_3}{h_k^2}v_{q_k}\right).
\]


\textbf{Bước 5 — Kiểm tra bằng hàm đơn giản.} Thử $u=1$, $u=x_i$ và $u=x_1^2+x_2^2+x_3^2$. Công thức Laplace phải lần lượt cho $0$, $0$, $6$. Với Hessian biểu diễn trong cơ sở trực chuẩn, đổi lại cơ sở Descartes bằng $H_x=E H_q E^T$, không nhầm Hessian với ma trận chỉ gồm các đạo hàm $v_{q_kq_l}$.
\subsection*{Lời giải áp dụng cho bài}
\textbf{Bước 1. Đặt hàm mới và kiểm tra khả nghịch.} Với
\[\xi=ax+by,\quad\eta=cx+dy,\quad J=ad-bc\ne0,\]
ta có $x=(d\xi-b\eta)/J$, $y=(-c\xi+a\eta)/J$. Đặt $v(\xi,\eta)=u(x,y)$, tức $u(x,y)=v(ax+by,cx+dy)$.

\textbf{Bước 2. Đạo hàm cấp một.} Quy tắc dây chuyền cho
\[u_x=a v_\xi+c v_\eta,\qquad u_y=b v_\xi+d v_\eta.\]
Tương đương $\partial_x=a\partial_\xi+c\partial_\eta$, $\partial_y=b\partial_\xi+d\partial_\eta$ khi tác động lên hàm đã đổi biến.

\textbf{Bước 3. Đạo hàm cấp hai.} Các hệ số $a,b,c,d$ là hằng số, nên
\begin{align*}
u_{xx}&=a(av_{\xi\xi}+cv_{\xi\eta})+c(av_{\eta\xi}+cv_{\eta\eta})
=a^2v_{\xi\xi}+2acv_{\xi\eta}+c^2v_{\eta\eta},\\
u_{xy}&=a(bv_{\xi\xi}+dv_{\xi\eta})+c(bv_{\eta\xi}+dv_{\eta\eta})
=abv_{\xi\xi}+(ad+bc)v_{\xi\eta}+cdv_{\eta\eta},\\
u_{yy}&=b^2v_{\xi\xi}+2bdv_{\xi\eta}+d^2v_{\eta\eta}.
\end{align*}
Đã dùng $v_{\xi\eta}=v_{\eta\xi}$ vì $v\in C^2$. Nếu cần đổi một PDE, nhân ba dòng này với hệ số tương ứng rồi gom theo từng đạo hàm mới. Ma trận phần chính $M$ được biến thành $TMT^T$, với $T=\left(\begin{smallmatrix}a&b\\c&d\end{smallmatrix}\right)$.

\textbf{Kiểm tra.} Chọn $v=\xi\eta$, thì $u=(ax+by)(cx+dy)$. Lấy đạo hàm trực tiếp cho $u_{xx}=2ac$, $u_{xy}=ad+bc$, $u_{yy}=2bd$, đúng với ba công thức trên.

\clearpage\noindent\texttt{\detokenize{VD_2_2_02}}\par
\section*{Lời giải bổ sung — Ví dụ 2 · Mục 2.2}

\subsection*{Kiến thức và trình tự giải}
\textbf{Mục tiêu.} Phân loại và đưa phương trình tuyến tính cấp hai về dạng chính tắc, có giữ đầy đủ các hạng cấp thấp.

\textbf{Bước 1 — Chuẩn hóa hệ số.} Viết


\[
A u_{xx}+2B u_{xy}+C u_{yy}+D u_x+E u_y+K u=f.
\]


Hệ số $B$ bằng một nửa hệ số in trước $u_{xy}$. Tính $\delta=B^2-AC$: dương là hyperbolic, âm là elliptic, bằng $0$ là parabolic nếu phần cấp hai còn khác $0$. Phải xét theo từng miền khi hệ số thay đổi hoặc suy biến.

\textbf{Bước 2 — Chọn biến mới.} Với hyperbolic, giải $A(dy)^2-2Bdxdy+C(dx)^2=0$ để tìm hai bất biến $\xi,\eta$ của hai họ đặc trưng. Với parabolic, chỉ có một họ: lấy bất biến $s$ và một biến $z$ độc lập. Với elliptic, không có hai họ đặc trưng thực; hoàn thành bình phương của phần chính để loại hạng hỗn hợp rồi cân bằng hai hệ số dương. Mọi trường hợp đều kiểm tra Jacobian khác $0$.

\textbf{Bước 3 — Viết công thức thay đạo hàm.} Nếu $u(x,y)=v(\xi,\eta)$ thì


\[
\begin{aligned}
u_x&=\xi_xv_\xi+\eta_xv_\eta,\\
u_{xx}&=\xi_x^2v_{\xi\xi}+2\xi_x\eta_xv_{\xi\eta}
+\eta_x^2v_{\eta\eta}+\xi_{xx}v_\xi+\eta_{xx}v_\eta.
\end{aligned}
\]


Làm tương tự cho $u_{xy},u_{yy}$. Có thể dùng toán tử $\partial_x,\partial_y$ để rút ngắn, nhưng phải nêu các đạo hàm của hệ số nếu đổi biến phi tuyến.

\textbf{Bước 4 — Gom hệ số và chia hợp lệ.} Hệ số mới của $v_{\xi\xi}$ là $A\xi_x^2+2B\xi_x\xi_y+C\xi_y^2$; hệ số của $v_{\xi\eta}$ là $2[A\xi_x\eta_x+B(\xi_x\eta_y+\xi_y\eta_x)+C\xi_y\eta_y]$. Hệ số mới của $v_\xi$ là $A\xi_{xx}+2B\xi_{xy}+C\xi_{yy}+D\xi_x+E\xi_y$. Chỉ chia cho hệ số khác $0$ trên miền đang xét.

\textbf{Bước 5 — Với ba biến.} Lập ma trận đối xứng $M$ của phần cấp hai; hệ số hỗn hợp được chia đôi. Dùng hoàn thành bình phương để tìm số hướng dương, âm và bằng $0$. Nếu $\xi=T x$, ma trận mới là $TMT^T$. Cho rõ ma trận $T$ và kiểm tra nó khả nghịch; chỉ ghi dấu của dạng toàn phương chưa đủ để thực hiện phép đổi biến.

\textbf{Chi tiết gom hệ số sau đổi biến.} Với $Lu=Au_{xx}+2Bu_{xy}+Cu_{yy}+Du_x+Eu_y+Ku$, tính riêng


\[
\begin{aligned}
a&=A\xi_x^2+2B\xi_x\xi_y+C\xi_y^2,\\
2b&=2[A\xi_x\eta_x+B(\xi_x\eta_y+\xi_y\eta_x)+C\xi_y\eta_y],\\
c&=A\eta_x^2+2B\eta_x\eta_y+C\eta_y^2,\\
d&=A\xi_{xx}+2B\xi_{xy}+C\xi_{yy}+D\xi_x+E\xi_y,\\
e&=A\eta_{xx}+2B\eta_{xy}+C\eta_{yy}+D\eta_x+E\eta_y.
\end{aligned}
\]


Khi đó $Lu=av_{\xi\xi}+2bv_{\xi\eta}+cv_{\eta\eta}+dv_\xi+ev_\eta+Kv$. Trong lời giải từng ý, thế các số hoặc hàm hệ số vào từng dòng trước khi rút gọn. Với ba biến, tính $TM$ trước, rồi lấy tích vô hướng mỗi hàng của $TM$ với mỗi hàng của $T$ để có $TMT^T$.
\subsection*{Lời giải áp dụng cho bài}
\textbf{Bước 1. Phân loại.} Phương trình là $2u_{xx}+3u_{xy}+u_{yy}+u_x-u_y=0$. Lấy $A=2$, $B=3/2$, $C=1$, ta được $\delta=9/4-2=1/4>0$: hyperbolic trên $\mathbb R^2$.

\textbf{Bước 2. Hai họ đặc trưng.} Phương trình đặc trưng
\[2(dy)^2-3dxdy+(dx)^2=(dy-dx)(2dy-dx)=0\]
cho $y-x=\mathrm{const}$ và $2y-x=\mathrm{const}$. Chọn $\xi=x-y$, $\eta=x-2y$. Jacobian $J=-1\ne0$; phép đổi ngược là $x=2\xi-\eta$, $y=\xi-\eta$.

\textbf{Bước 3. Thay đạo hàm.} Vì $\partial_x=\partial_\xi+\partial_\eta$ và $\partial_y=-\partial_\xi-2\partial_\eta$, ta có
\begin{align*}
u_{xx}&=v_{\xi\xi}+2v_{\xi\eta}+v_{\eta\eta},\\
u_{xy}&=-v_{\xi\xi}-3v_{\xi\eta}-2v_{\eta\eta},\\
u_{yy}&=v_{\xi\xi}+4v_{\xi\eta}+4v_{\eta\eta},\\
u_x-u_y&=2v_\xi+3v_\eta.
\end{align*}
Các hệ số của $v_{\xi\xi}$, $v_{\eta\eta}$ lần lượt là $2-3+1=0$, $2-6+4=0$; hệ số hỗn hợp là $4-9+4=-1$.

\textbf{Kết quả.} Dạng chính tắc là
\[\boxed{-v_{\xi\eta}+2v_\xi+3v_\eta=0.}\]
Không bỏ các hạng cấp một khi phân loại phần chính. Có thể nhân cả phương trình với $-1$ để viết $v_{\xi\eta}-2v_\xi-3v_\eta=0$.

\clearpage\noindent\texttt{\detokenize{VD_2_2_03}}\par
\section*{Lời giải bổ sung — Ví dụ 3 · Mục 2.2}

\subsection*{Kiến thức và trình tự giải}
\textbf{Mục tiêu.} Phân loại và đưa phương trình tuyến tính cấp hai về dạng chính tắc, có giữ đầy đủ các hạng cấp thấp.

\textbf{Bước 1 — Chuẩn hóa hệ số.} Viết


\[
A u_{xx}+2B u_{xy}+C u_{yy}+D u_x+E u_y+K u=f.
\]


Hệ số $B$ bằng một nửa hệ số in trước $u_{xy}$. Tính $\delta=B^2-AC$: dương là hyperbolic, âm là elliptic, bằng $0$ là parabolic nếu phần cấp hai còn khác $0$. Phải xét theo từng miền khi hệ số thay đổi hoặc suy biến.

\textbf{Bước 2 — Chọn biến mới.} Với hyperbolic, giải $A(dy)^2-2Bdxdy+C(dx)^2=0$ để tìm hai bất biến $\xi,\eta$ của hai họ đặc trưng. Với parabolic, chỉ có một họ: lấy bất biến $s$ và một biến $z$ độc lập. Với elliptic, không có hai họ đặc trưng thực; hoàn thành bình phương của phần chính để loại hạng hỗn hợp rồi cân bằng hai hệ số dương. Mọi trường hợp đều kiểm tra Jacobian khác $0$.

\textbf{Bước 3 — Viết công thức thay đạo hàm.} Nếu $u(x,y)=v(\xi,\eta)$ thì


\[
\begin{aligned}
u_x&=\xi_xv_\xi+\eta_xv_\eta,\\
u_{xx}&=\xi_x^2v_{\xi\xi}+2\xi_x\eta_xv_{\xi\eta}
+\eta_x^2v_{\eta\eta}+\xi_{xx}v_\xi+\eta_{xx}v_\eta.
\end{aligned}
\]


Làm tương tự cho $u_{xy},u_{yy}$. Có thể dùng toán tử $\partial_x,\partial_y$ để rút ngắn, nhưng phải nêu các đạo hàm của hệ số nếu đổi biến phi tuyến.

\textbf{Bước 4 — Gom hệ số và chia hợp lệ.} Hệ số mới của $v_{\xi\xi}$ là $A\xi_x^2+2B\xi_x\xi_y+C\xi_y^2$; hệ số của $v_{\xi\eta}$ là $2[A\xi_x\eta_x+B(\xi_x\eta_y+\xi_y\eta_x)+C\xi_y\eta_y]$. Hệ số mới của $v_\xi$ là $A\xi_{xx}+2B\xi_{xy}+C\xi_{yy}+D\xi_x+E\xi_y$. Chỉ chia cho hệ số khác $0$ trên miền đang xét.

\textbf{Bước 5 — Với ba biến.} Lập ma trận đối xứng $M$ của phần cấp hai; hệ số hỗn hợp được chia đôi. Dùng hoàn thành bình phương để tìm số hướng dương, âm và bằng $0$. Nếu $\xi=T x$, ma trận mới là $TMT^T$. Cho rõ ma trận $T$ và kiểm tra nó khả nghịch; chỉ ghi dấu của dạng toàn phương chưa đủ để thực hiện phép đổi biến.

\textbf{Chi tiết gom hệ số sau đổi biến.} Với $Lu=Au_{xx}+2Bu_{xy}+Cu_{yy}+Du_x+Eu_y+Ku$, tính riêng


\[
\begin{aligned}
a&=A\xi_x^2+2B\xi_x\xi_y+C\xi_y^2,\\
2b&=2[A\xi_x\eta_x+B(\xi_x\eta_y+\xi_y\eta_x)+C\xi_y\eta_y],\\
c&=A\eta_x^2+2B\eta_x\eta_y+C\eta_y^2,\\
d&=A\xi_{xx}+2B\xi_{xy}+C\xi_{yy}+D\xi_x+E\xi_y,\\
e&=A\eta_{xx}+2B\eta_{xy}+C\eta_{yy}+D\eta_x+E\eta_y.
\end{aligned}
\]


Khi đó $Lu=av_{\xi\xi}+2bv_{\xi\eta}+cv_{\eta\eta}+dv_\xi+ev_\eta+Kv$. Trong lời giải từng ý, thế các số hoặc hàm hệ số vào từng dòng trước khi rút gọn. Với ba biến, tính $TM$ trước, rồi lấy tích vô hướng mỗi hàng của $TM$ với mỗi hàng của $T$ để có $TMT^T$.
\subsection*{Lời giải áp dụng cho bài}
\textbf{Bước 1. Phân loại.} Viết $2u_{xx}+2u_{xy}+u_{yy}+u_x+u_y-u=0$. Với $A=2$, $B=1$, $C=1$, $\delta=-1<0$, phương trình elliptic. Không có hai họ đặc trưng thực.

\textbf{Bước 2. Hoàn thành bình phương.} Phần chính có thể viết $(\partial_x+\partial_y)^2+\partial_x^2$. Chọn $\xi=x-y$, $\eta=y$, $J=1$. Khi đó $x=\xi+\eta$, $y=\eta$ và
\[\partial_x=\partial_\xi,\qquad\partial_y=-\partial_\xi+\partial_\eta.\]

\textbf{Bước 3. Thế cụ thể.}
\[u_{xx}=v_{\xi\xi},\quad u_{xy}=-v_{\xi\xi}+v_{\xi\eta},\quad
u_{yy}=v_{\xi\xi}-2v_{\xi\eta}+v_{\eta\eta}.\]
Do đó $2u_{xx}+2u_{xy}+u_{yy}=v_{\xi\xi}+v_{\eta\eta}$ và $u_x+u_y=v_\eta$.

\textbf{Kết quả.}
\[\boxed{v_{\xi\xi}+v_{\eta\eta}+v_\eta-v=0.}\]
Hai hệ số cấp hai đã bằng nhau và không còn đạo hàm hỗn hợp. Nếu muốn bỏ hạng $v_\eta$, đặt $v=e^{-\eta/2}w$; đạo hàm tích cho $w_{\xi\xi}+w_{\eta\eta}-\tfrac54w=0$. Đây là phép biến đổi phụ, không cần để hoàn thành yêu cầu đưa về chính tắc.

\textbf{Chi tiết gom hệ số và đổi ẩn phụ.} Ba hệ số cấp hai là $2-2+1=1$ của $v_{\xi\xi}$, $2-2=0$ của $v_{\xi\eta}$, và $1$ của $v_{\eta\eta}$. Với $v=e^{-\eta/2}w$,
\[v_{\xi\xi}=e^{-\eta/2}w_{\xi\xi},\quad v_\eta=e^{-\eta/2}(w_\eta-w/2),\quad
v_{\eta\eta}=e^{-\eta/2}(w_{\eta\eta}-w_\eta+w/4).\]
Thế rồi chia $e^{-\eta/2}$: hai hạng $w_\eta$ triệt tiêu, hệ số $w$ là $1/4-1/2-1=-5/4$.

\clearpage\noindent\texttt{\detokenize{VD_2_2_04}}\par
\section*{Lời giải bổ sung — Ví dụ 4 · Mục 2.2}

\subsection*{Kiến thức và trình tự giải}
\textbf{Mục tiêu.} Phân loại và đưa phương trình tuyến tính cấp hai về dạng chính tắc, có giữ đầy đủ các hạng cấp thấp.

\textbf{Bước 1 — Chuẩn hóa hệ số.} Viết


\[
A u_{xx}+2B u_{xy}+C u_{yy}+D u_x+E u_y+K u=f.
\]


Hệ số $B$ bằng một nửa hệ số in trước $u_{xy}$. Tính $\delta=B^2-AC$: dương là hyperbolic, âm là elliptic, bằng $0$ là parabolic nếu phần cấp hai còn khác $0$. Phải xét theo từng miền khi hệ số thay đổi hoặc suy biến.

\textbf{Bước 2 — Chọn biến mới.} Với hyperbolic, giải $A(dy)^2-2Bdxdy+C(dx)^2=0$ để tìm hai bất biến $\xi,\eta$ của hai họ đặc trưng. Với parabolic, chỉ có một họ: lấy bất biến $s$ và một biến $z$ độc lập. Với elliptic, không có hai họ đặc trưng thực; hoàn thành bình phương của phần chính để loại hạng hỗn hợp rồi cân bằng hai hệ số dương. Mọi trường hợp đều kiểm tra Jacobian khác $0$.

\textbf{Bước 3 — Viết công thức thay đạo hàm.} Nếu $u(x,y)=v(\xi,\eta)$ thì


\[
\begin{aligned}
u_x&=\xi_xv_\xi+\eta_xv_\eta,\\
u_{xx}&=\xi_x^2v_{\xi\xi}+2\xi_x\eta_xv_{\xi\eta}
+\eta_x^2v_{\eta\eta}+\xi_{xx}v_\xi+\eta_{xx}v_\eta.
\end{aligned}
\]


Làm tương tự cho $u_{xy},u_{yy}$. Có thể dùng toán tử $\partial_x,\partial_y$ để rút ngắn, nhưng phải nêu các đạo hàm của hệ số nếu đổi biến phi tuyến.

\textbf{Bước 4 — Gom hệ số và chia hợp lệ.} Hệ số mới của $v_{\xi\xi}$ là $A\xi_x^2+2B\xi_x\xi_y+C\xi_y^2$; hệ số của $v_{\xi\eta}$ là $2[A\xi_x\eta_x+B(\xi_x\eta_y+\xi_y\eta_x)+C\xi_y\eta_y]$. Hệ số mới của $v_\xi$ là $A\xi_{xx}+2B\xi_{xy}+C\xi_{yy}+D\xi_x+E\xi_y$. Chỉ chia cho hệ số khác $0$ trên miền đang xét.

\textbf{Bước 5 — Với ba biến.} Lập ma trận đối xứng $M$ của phần cấp hai; hệ số hỗn hợp được chia đôi. Dùng hoàn thành bình phương để tìm số hướng dương, âm và bằng $0$. Nếu $\xi=T x$, ma trận mới là $TMT^T$. Cho rõ ma trận $T$ và kiểm tra nó khả nghịch; chỉ ghi dấu của dạng toàn phương chưa đủ để thực hiện phép đổi biến.

\textbf{Chi tiết gom hệ số sau đổi biến.} Với $Lu=Au_{xx}+2Bu_{xy}+Cu_{yy}+Du_x+Eu_y+Ku$, tính riêng


\[
\begin{aligned}
a&=A\xi_x^2+2B\xi_x\xi_y+C\xi_y^2,\\
2b&=2[A\xi_x\eta_x+B(\xi_x\eta_y+\xi_y\eta_x)+C\xi_y\eta_y],\\
c&=A\eta_x^2+2B\eta_x\eta_y+C\eta_y^2,\\
d&=A\xi_{xx}+2B\xi_{xy}+C\xi_{yy}+D\xi_x+E\xi_y,\\
e&=A\eta_{xx}+2B\eta_{xy}+C\eta_{yy}+D\eta_x+E\eta_y.
\end{aligned}
\]


Khi đó $Lu=av_{\xi\xi}+2bv_{\xi\eta}+cv_{\eta\eta}+dv_\xi+ev_\eta+Kv$. Trong lời giải từng ý, thế các số hoặc hàm hệ số vào từng dòng trước khi rút gọn. Với ba biến, tính $TM$ trước, rồi lấy tích vô hướng mỗi hàng của $TM$ với mỗi hàng của $T$ để có $TMT^T$.
\subsection*{Lời giải áp dụng cho bài}
\textbf{Bước 1. Phân loại và đặc trưng.} Với $u_{xx}-2u_{xy}+u_{yy}+u_x-u_y+u=1$, ta có $A=C=1$, $B=-1$, $\delta=0$ và phần chính khác $0$. Phương trình parabolic; đặc trưng $(dy+dx)^2=0$ cho $x+y=\mathrm{const}$.

\textbf{Bước 2. Chọn tọa độ.} Đặt $s=x+y$, $z=x$, $J=-1$, $u=v(s,z)$. Ta có $\partial_x=\partial_s+\partial_z$, $\partial_y=\partial_s$, nên $D=\partial_x-\partial_y=\partial_z$. Vì vậy $u_x-u_y=v_z$ và $u_{xx}-2u_{xy}+u_{yy}=v_{zz}$.

\textbf{Bước 3. Dạng chính tắc.}
\[\boxed{v_{zz}+v_z+v=1.}\]
Chỉ còn đạo hàm theo $z$; $s$ giữ vai trò tham số. Đây cũng cho phép giải thêm: nghiệm riêng $v_p=1$, phương trình đặc trưng $\lambda^2+\lambda+1=0$ có nghiệm $-1/2\pm i\sqrt3/2$.

\textbf{Nghiệm tổng quát nếu cần.}
\[u=1+e^{-x/2}\left[F(x+y)\cos\frac{\sqrt3x}{2}+G(x+y)\sin\frac{\sqrt3x}{2}\right].\]
Với $F,G\in C^2$, áp dụng $D=\partial_z$ vào công thức cho $D^2u+Du+u=1$, kiểm tra lại phương trình gốc.

\textbf{Chi tiết ba đạo hàm cấp hai.} Từ $u_x=v_s+v_z,u_y=v_s$,
\[u_{xx}=v_{ss}+2v_{sz}+v_{zz},\quad u_{xy}=v_{ss}+v_{sz},\quad u_{yy}=v_{ss}.\]
Nhân với $1,-2,1$ và cộng, hệ số $v_{ss}$ là $1-2+1=0$, hệ số $v_{sz}$ là $2-2=0$, còn $v_{zz}$. Với $v=1+w$, phương trình thuần nhất $w_{zz}+w_z+w=0$ có $\lambda=(-1\pm\sqrt{-3})/2=-1/2\pm i\sqrt3/2$, cho hai nghiệm thực đã viết.

\clearpage\noindent\texttt{\detokenize{VD_2_2_05}}\par
\section*{Lời giải bổ sung — Ví dụ 5 · Mục 2.2}

\subsection*{Kiến thức và trình tự giải}
\textbf{Mục tiêu.} Tính đạo hàm khi đổi tọa độ rồi thay vào PDE. Đổi tọa độ của điểm và đổi biểu thức của hàm là hai việc cần làm cùng nhau: đặt $q=q(x)$ và $v(q)=u(x(q))$.

\textbf{Bước 1 — Kiểm tra phép đổi biến.} Tính $J=\det(\partial q/\partial x)$. Chỉ dùng các công thức trên miền $J\ne0$ và chọn một nhánh liên tục cho các góc. Tọa độ cực/trụ loại trục $r=0$; tọa độ cầu còn loại $\sin\theta=0$. Đây là suy biến của tọa độ, chưa phải kết luận rằng hàm $u$ không tồn tại ở đó.

\textbf{Bước 2 — Tính đạo hàm cấp một bằng quy tắc dây chuyền.} Nếu $b_{ik}=\partial q_k/\partial x_i$, thì


\[
u_{x_i}=\sum_k b_{ik}v_{q_k},\qquad D_i=\sum_k b_{ik}\partial_{q_k}.
\]


Với $q=Ax$, $b_{ik}=a_{ki}$ là hằng số. Với tọa độ cong, lấy nghịch đảo ma trận $\partial x/\partial q$ để tìm $b_{ik}$.

\textbf{Bước 3 — Lấy đạo hàm cấp hai cả hàm lẫn hệ số.} Áp dụng đạo hàm tích cho từng hạng:


\[
u_{x_ix_j}=\sum_{k,l}b_{ik}b_{jl}v_{q_kq_l}
+\sum_{k,l}b_{ik}(\partial_{q_k}b_{jl})v_{q_l}.
\]


Hạng thứ hai chỉ mất đi khi phép đổi biến tuyến tính. Ví dụ $\partial_x=\cos\theta\,\partial_r-(\sin\theta/r)\partial_\theta$ thì lần đạo hàm tiếp theo phải tác động lên cả $\cos\theta$ và $\sin\theta/r$.

\textbf{Bước 4 — Thu gọn toán tử.} Gom hệ số của từng $v_{q_kq_l}$ và $v_{q_k}$ trước khi cộng. Với Laplace, lấy tổng các đạo hàm $u_{x_ix_i}$. Trong hệ trực giao có hệ số độ dài $h_1,h_2,h_3$, có thể đối chiếu bằng


\[
\Delta u=\frac1{h_1h_2h_3}\sum_{k=1}^3
\partial_{q_k}\left(\frac{h_1h_2h_3}{h_k^2}v_{q_k}\right).
\]


\textbf{Bước 5 — Kiểm tra bằng hàm đơn giản.} Thử $u=1$, $u=x_i$ và $u=x_1^2+x_2^2+x_3^2$. Công thức Laplace phải lần lượt cho $0$, $0$, $6$. Với Hessian biểu diễn trong cơ sở trực chuẩn, đổi lại cơ sở Descartes bằng $H_x=E H_q E^T$, không nhầm Hessian với ma trận chỉ gồm các đạo hàm $v_{q_kq_l}$.
\subsection*{Lời giải áp dụng cho bài}
\textbf{Bước 1. Nghịch đảo vi phân tọa độ cực.} Viết $c=\cos\theta$, $s=\sin\theta$, $x=rc$, $y=rs$, $r>0$. Từ $dx=c\,dr-rs\,d\theta$, $dy=s\,dr+rc\,d\theta$ suy ra
\[dr=c\,dx+s\,dy,\qquad d\theta=-\frac{s}{r}dx+\frac{c}{r}dy.\]
Vì vậy $D_x=c\partial_r-s\partial_\theta/r$, $D_y=s\partial_r+c\partial_\theta/r$.

\textbf{Bước 2. Đạo hàm cấp một.}
\[u_x=cv_r-\frac{s}{r}v_\theta,\qquad u_y=sv_r+\frac{c}{r}v_\theta.\]

\textbf{Bước 3. Tính cả đạo hàm của hệ số.} Ta có $D_xc=s^2/r$, $D_x(-s/r)=2sc/r^2$, $D_yc=-sc/r$, $D_y(-s/r)=-(c^2-s^2)/r^2$. Ví dụ
\[u_{xx}=(D_xc)v_r+cD_xv_r+D_x(-s/r)v_\theta-(s/r)D_xv_\theta.\]
Khai triển cho cả ba đạo hàm:
\begin{align*}
u_{xx}&=c^2v_{rr}-\frac{2sc}{r}v_{r\theta}+\frac{s^2}{r^2}v_{\theta\theta}
+\frac{s^2}{r}v_r+\frac{2sc}{r^2}v_\theta,\\
u_{xy}&=scv_{rr}+\frac{c^2-s^2}{r}v_{r\theta}-\frac{sc}{r^2}v_{\theta\theta}
-\frac{sc}{r}v_r-\frac{c^2-s^2}{r^2}v_\theta,\\
u_{yy}&=s^2v_{rr}+\frac{2sc}{r}v_{r\theta}+\frac{c^2}{r^2}v_{\theta\theta}
+\frac{c^2}{r}v_r-\frac{2sc}{r^2}v_\theta.
\end{align*}

\textbf{Bước 4. Cộng để được Laplace.} Dùng $c^2+s^2=1$, các hạng $v_{r\theta}$ và $v_\theta$ triệt tiêu:
\[\boxed{\Delta u=v_{rr}+\frac1r v_r+\frac1{r^2}v_{\theta\theta}.}\]
Kiểm tra với $u=x$, tức $v=r\cos\theta$: kết quả bằng $0$. Với $u=x^2+y^2$, $v=r^2$, kết quả bằng $4$. Tại $r=0$, phải xét hàm gốc hoặc giới hạn trơn, không thế trực tiếp vào hệ số $1/r$.

\textbf{Chi tiết đạo hàm tích cho hai dòng còn lại.} Ngoài các hệ số đã tính, $D_ys=c^2/r$, $D_y(c/r)=-2sc/r^2$. Vì $D_yv_r=sv_{rr}+cv_{r\theta}/r$, $D_yv_\theta=sv_{r\theta}+cv_{\theta\theta}/r$,
\[\begin{aligned}
u_{xy}&=-scv_r/r+c(sv_{rr}+cv_{r\theta}/r)
-(c^2-s^2)v_\theta/r^2-(s/r)(sv_{r\theta}+cv_{\theta\theta}/r),\\
u_{yy}&=c^2v_r/r+s(sv_{rr}+cv_{r\theta}/r)
-2scv_\theta/r^2+(c/r)(sv_{r\theta}+cv_{\theta\theta}/r).
\end{aligned}\]
Gom các hạng cùng đạo hàm cho đúng hai dòng ở bước 3.

\clearpage\noindent\texttt{\detokenize{VD_2_2_06}}\par
\section*{Lời giải bổ sung — Ví dụ 6 · Mục 2.2}

\subsection*{Kiến thức và trình tự giải}
\textbf{Mục tiêu.} Giải phương trình hyperbolic sau khi đã tìm hai tọa độ đặc trưng độc lập.

\textbf{Bước 1 — Phân loại, tìm đặc trưng và thay đạo hàm.} Làm đủ các bước của phép đưa về chính tắc. Viết cả phép đổi ngược để biểu diễn vế phải $\widetilde f(\xi,\eta)=f(x(\xi,\eta),y(\xi,\eta))$.

\textbf{Bước 2 — Tích phân phương trình chính tắc.} Với $v_{\xi\eta}=g$, giữ biến kia cố định khi tích phân:


\[
v=F(\xi)+G(\eta)+\int_{\xi_0}^{\xi}\int_{\eta_0}^{\eta}g(a,b)\,db\,da.
\]


Nếu tìm được một nghiệm riêng $P(x,y)$ dễ hơn, viết $u=P+F(\xi)+G(\eta)$. Với $v_{\xi\eta}+\alpha v_\eta=0$, đặt $w=v_\eta$, giải $w_\xi+\alpha w=0$ rồi tích phân theo $\eta$; thu được $v=F(\xi)+e^{-\alpha\xi}G(\eta)$.

\textbf{Bước 3 — Đổi cả dữ liệu.} Thế tọa độ đường dữ liệu vào $\xi,\eta$, dùng quy tắc dây chuyền cho đạo hàm được cho. Trên đường không đặc trưng, hai dữ liệu thường xác định hai hàm $F,G$. Trên đường đặc trưng, có biến không thay đổi, nên dữ liệu chỉ xác định giá trị hoặc đạo hàm của một hàm tại một điểm; phải tìm điều kiện tương thích và nêu phần tùy ý còn lại.

\textbf{Bước 4 — Lập đẳng thức hình bình hành.} Với $v=F(\xi)+G(\eta)$, bốn đỉnh theo hai họ đặc trưng thỏa $v_{11}+v_{22}=v_{12}+v_{21}$. Nếu có số hạng mũ, trước hết nhân trọng số để được tổng hai hàm một biến. Với vế phải $g$, hiệu bốn đỉnh bằng tích phân kép của $g$ trên hình chữ nhật trong tọa độ mới.

\textbf{Bước 5 — Xác định miền được dữ liệu quyết định.} Hai đặc trưng đi qua điểm đang xét phải gặp đoạn dữ liệu phù hợp. Với phương trình sóng $v_{tt}=v_{zz}$ trên $z>0$, dữ liệu lúc $t=0$ xác định duy nhất miền $z>t$; ở $0<z<t$, cần thêm dữ liệu biên $z=0$. Với đặc trưng cong, dùng miền trong các tọa độ mới rồi đổi lại.

\textbf{Bước 6 — Kiểm tra.} Thế nghiệm vào cả PDE và từng điều kiện. Nếu đề yêu cầu hai nghiệm, chỉ rõ hai lựa chọn khác nhau cho hàm tùy ý và chứng minh phần chênh lệch thỏa phương trình thuần nhất, đồng thời triệt tiêu trên đường dữ liệu.
\subsection*{Lời giải áp dụng cho bài}
\textbf{(i) Phân loại và đổi biến.} $A=2$, $B=3/2$, $C=1$ cho $\delta=1/4>0$. Hai họ đặc trưng là $x-y=\mathrm{const}$, $x-2y=\mathrm{const}$. Đặt $\xi=x-y$, $\eta=x-2y$, $x=2\xi-\eta$, $y=\xi-\eta$. Như Ví dụ 2, phần chính bằng $-v_{\xi\eta}$; $au_x+bu_y=(a-b)v_\xi+(a-2b)v_\eta$. Dạng chính tắc là
\[-v_{\xi\eta}+(a-b)v_\xi+(a-2b)v_\eta+\widetilde f=0.\]

\textbf{(ii) Khi $a=b=0$: kiểm tra dấu của công thức trong đề.} Gọi $L=2\partial_x^2+3\partial_x\partial_y+\partial_y^2=(2\partial_x+\partial_y)(\partial_x+\partial_y)$. Với
\[F(x,y)=\int_0^y\int_{2t+x-2y}^{t+x-y}f(s,t)\,ds\,dt,\]
quy tắc Leibniz cho
\[(\partial_x+\partial_y)F=\int_0^y f(2t+x-2y,t)dt.\]
Ở đây hạng tại cận $t=y$ của tích phân ban đầu bằng $0$, vì hai cận trong đều bằng $x$. Tác động tiếp $2\partial_x+\partial_y$ lên dòng trên: đối số $2t+x-2y$ có đạo hàm theo hướng này bằng $0$, còn cận $y$ cho $f(x,y)$. Do đó $LF=f$.

\textbf{Lưu ý dấu in trong sách.} Phương trình được in là $Lu+f=0$, nên nghiệm riêng phải là $-F$, không phải $F$. Chẳng hạn $f=1$ cho $F=y^2/2$ và $LF+f=2$. Nếu viết PDE thành $Lu=f$ thì $F$ đúng. Giữ phương trình in ở đầu bài, nghiệm tổng quát là
\[\boxed{u=-F(x,y)+P(x-y)+Q(x-2y).}\]
Tương đương, từ $v_{\xi\eta}=\widetilde f$ có thể dùng tích phân kép trong tọa độ mới. Giả sử $f$ đủ trơn để các đạo hàm và tích phân vừa dùng có nghĩa.

\textbf{(iii) Khi $a=b=1$, $f=0$.} Phương trình thành $v_{\xi\eta}+v_\eta=0$. Đặt $w=v_\eta$, giải $w_\xi+w=0$ được $w=e^{-\xi}H(\eta)$. Tích phân theo $\eta$:
\[\boxed{u=P(x-y)+e^{-(x-y)}Q(x-2y).}\]
Nhân $v$ với $e^\xi$ để có $e^\xi v=e^\xi P(\xi)+Q(\eta)$. Vì là tổng hai hàm một biến, bốn đỉnh hình bình hành đặc trưng thỏa
\[e^{\xi_1}v_{11}+e^{\xi_2}v_{22}=e^{\xi_1}v_{12}+e^{\xi_2}v_{21},\]
với $v_{ij}=v(\xi_i,\eta_j)$. Thế lại $v_{\xi\eta}=-v_\eta$ kiểm tra ngay PDE.

\textbf{Chi tiết thế đạo hàm ở (i).} $u_{xx}=v_{\xi\xi}+2v_{\xi\eta}+v_{\eta\eta}$, $u_{xy}=-v_{\xi\xi}-3v_{\xi\eta}-2v_{\eta\eta}$, $u_{yy}=v_{\xi\xi}+4v_{\xi\eta}+4v_{\eta\eta}$. Phần chính có hệ số $2-3+1=0$, $4-9+4=-1$, $2-6+4=0$, lần lượt theo $v_{\xi\xi},v_{\xi\eta},v_{\eta\eta}$.

\textbf{Chi tiết Leibniz ở (ii).} Đặt $\ell=2t+x-2y$, $h=t+x-y$. Vì tích phân trong bằng $0$ tại $t=y$,
\[F_x=\int_0^y[f(h,t)-f(\ell,t)]dt,\quad
F_y=\int_0^y[-f(h,t)+2f(\ell,t)]dt.\]
Vậy $F_x+F_y=\int_0^yf(\ell,t)dt=:I$. Lấy $I_x=\int_0^yf_x(\ell,t)dt$ và $I_y=f(x,y)-2\int_0^yf_x(\ell,t)dt$, suy $2I_x+I_y=f(x,y)$. Đây là phép triệt tiêu tạo ra $LF=f$.

\clearpage\noindent\texttt{\detokenize{VD_2_2_07}}\par
\section*{Lời giải bổ sung — Ví dụ 7 · Mục 2.2}

\subsection*{Kiến thức và trình tự giải}
\textbf{Mục tiêu.} Tận dụng phần cấp hai có dạng bình phương một đạo hàm theo hướng để giảm PDE thành ODE có tham số.

\textbf{Bước 1 — Nhận ra hướng đạo hàm.} Nếu $D=a\partial_x+b\partial_y$ có hệ số hằng, $D^2u=a^2u_{xx}+2ab u_{xy}+b^2u_{yy}$. Với hệ số biến thiên phải dùng


\[
D^2u=a^2u_{xx}+2ab u_{xy}+b^2u_{yy}+(Da)u_x+(Db)u_y.
\]


So sánh cả hạng cấp một trước khi kết luận PDE bằng $D^2u+\lambda Du+\mu u=f$.

\textbf{Bước 2 — Tìm bất biến $s$ và biến tiến $z$.} Giải $Ds=0$, chọn $Dz=k\ne0$. Kiểm tra $J=\partial(s,z)/\partial(x,y)\ne0$. Khi $k$ hằng, $D=k\partial_z$; PDE trở thành ODE theo $z$, trong đó $s$ là tham số. Nếu còn $v_s$, phải giữ lại, chỉ các trường hợp nó triệt tiêu mới giải trực tiếp như ODE.

\textbf{Bước 3 — Giải ODE.} Với $v_{zz}=g(s,z)$, tích phân hai lần cho $v=zF(s)+G(s)+\int_{z_0}^z(z-t)g(s,t)dt$. Với $a v_{zz}+b v_z+c v=g$, giải đa thức $a\lambda^2+b\lambda+c=0$, giữ hai hàm tùy ý theo $s$. Nghiệm kép có thêm hệ số $z$; có thể tìm nghiệm riêng dạng mũ hoặc đa thức.

\textbf{Bước 4 — Dùng dữ liệu đúng biến.} Tính $u_x=s_xv_s+z_xv_z$ và $u_y=s_yv_s+z_yv_z$. Nếu dữ liệu ở $s=0$, chỉ biết $F(0),G(0)$ hoặc đạo hàm của chúng; không suy ra toàn bộ các hàm. So sánh hệ số của $e^z$, $z$, $z^2$ để tìm các hằng số tương thích.

\textbf{Bước 5 — Kiểm tra tính duy nhất và nghiệm cụ thể.} Khi đường dữ liệu gặp tất cả đặc trưng một cách độc lập, giải hệ cho $F,G$. Khi đường dữ liệu là một đặc trưng, mô tả các hàm còn tùy ý và đưa hai lựa chọn rõ ràng. Cuối cùng thế lại phương trình gốc, không chỉ ODE mới.
\subsection*{Lời giải áp dụng cho bài}
\textbf{(i) Phân loại và chính tắc.} $A=C=1$, $B=-1$, $\delta=0$: parabolic. Đặc trưng $x+y=\mathrm{const}$. Đặt $s=x+y$, $z=x$; $\partial_x=\partial_s+\partial_z$, $\partial_y=\partial_s$. Suy ra
\[u_{xx}-2u_{xy}+u_{yy}=v_{zz},\quad au_x+bu_y=(a+b)v_s+av_z.\]
Với $x=z$, $y=s-z$, phương trình mới là
\[\boxed{v_{zz}+av_z+(a+b)v_s+cv=f(z,s-z).}\]
Chỉ khi $a+b=0$ mới không còn đạo hàm theo tham số $s$.

\textbf{(ii.a) $a=b=c=0$.} Giữ $s$ cố định, tích phân $v_{zz}=f(z,s-z)$ hai lần. Với mốc $z=0$,
\[v=zP(s)+Q(s)+\int_0^z(z-t)f(t,s-t)dt.\]
Vì đạo hàm lần thứ nhất của tích phân là $\int_0^z f(t,s-t)dt$ và lần thứ hai là $f(z,s-z)$, công thức thỏa PDE. Đổi lại:
\[u=xP(x+y)+Q(x+y)+\int_0^x(x-t)f(t,x+y-t)dt.\]

\textbf{(ii.b) $a=b=0$, $c=1$, $f=0$.} ODE $v_{zz}+v=0$ có hai nghiệm độc lập $\cos z$, $\sin z$. Hằng số tích phân phụ thuộc $s$:
\[\boxed{u=P(x+y)\cos x+Q(x+y)\sin x.}\]

\textbf{(ii.c) $a=1$, $b=-1$, $c=0$, $f=1$.} ODE $v_{zz}+v_z=1$. Đặt $p=v_z$: $(e^zp)'=e^z$, nên $p=1+C(s)e^{-z}$. Tích phân lần nữa và đổi tên hàm tùy ý:
\[\boxed{u=x+P(x+y)+e^{-x}Q(x+y).}\]
Thật vậy, $v_z=1-e^{-z}Q$, $v_{zz}=e^{-z}Q$, tổng bằng $1$. Các hàm tùy ý lấy đủ trơn, chẳng hạn $C^2$.

\textbf{Chi tiết thế và hai lần tích phân.} Với $s=x+y,z=x$,
\[u_{xx}=v_{ss}+2v_{sz}+v_{zz},\quad u_{xy}=v_{ss}+v_{sz},\quad u_{yy}=v_{ss}.\]
Vậy $u_{xx}-2u_{xy}+u_{yy}=(1-2+1)v_{ss}+(2-2)v_{sz}+v_{zz}$. Ở (ii.a), lần đầu $v_z=P(s)+\int_0^zf(t,s-t)dt$; lần sau
\[v=Q(s)+zP(s)+\int_0^z\int_0^r f(t,s-t)dt\,dr.\]
Đổi thứ tự trên tam giác $0\le t\le r\le z$ được $\int_0^z(z-t)f(t,s-t)dt$ khi $z\ge0$. Với $z<0$, cùng công thức suy trực tiếp bằng hai đạo hàm và hai vết tại $z=0$, nên vẫn đúng với tích phân có hướng.

\clearpage\noindent\texttt{\detokenize{VD_2_2_08}}\par
\section*{Lời giải bổ sung — Ví dụ 8 · Mục 2.2}

\subsection*{Kiến thức và trình tự giải}
\textbf{Mục tiêu.} Giải phương trình hyperbolic sau khi đã tìm hai tọa độ đặc trưng độc lập.

\textbf{Bước 1 — Phân loại, tìm đặc trưng và thay đạo hàm.} Làm đủ các bước của phép đưa về chính tắc. Viết cả phép đổi ngược để biểu diễn vế phải $\widetilde f(\xi,\eta)=f(x(\xi,\eta),y(\xi,\eta))$.

\textbf{Bước 2 — Tích phân phương trình chính tắc.} Với $v_{\xi\eta}=g$, giữ biến kia cố định khi tích phân:


\[
v=F(\xi)+G(\eta)+\int_{\xi_0}^{\xi}\int_{\eta_0}^{\eta}g(a,b)\,db\,da.
\]


Nếu tìm được một nghiệm riêng $P(x,y)$ dễ hơn, viết $u=P+F(\xi)+G(\eta)$. Với $v_{\xi\eta}+\alpha v_\eta=0$, đặt $w=v_\eta$, giải $w_\xi+\alpha w=0$ rồi tích phân theo $\eta$; thu được $v=F(\xi)+e^{-\alpha\xi}G(\eta)$.

\textbf{Bước 3 — Đổi cả dữ liệu.} Thế tọa độ đường dữ liệu vào $\xi,\eta$, dùng quy tắc dây chuyền cho đạo hàm được cho. Trên đường không đặc trưng, hai dữ liệu thường xác định hai hàm $F,G$. Trên đường đặc trưng, có biến không thay đổi, nên dữ liệu chỉ xác định giá trị hoặc đạo hàm của một hàm tại một điểm; phải tìm điều kiện tương thích và nêu phần tùy ý còn lại.

\textbf{Bước 4 — Lập đẳng thức hình bình hành.} Với $v=F(\xi)+G(\eta)$, bốn đỉnh theo hai họ đặc trưng thỏa $v_{11}+v_{22}=v_{12}+v_{21}$. Nếu có số hạng mũ, trước hết nhân trọng số để được tổng hai hàm một biến. Với vế phải $g$, hiệu bốn đỉnh bằng tích phân kép của $g$ trên hình chữ nhật trong tọa độ mới.

\textbf{Bước 5 — Xác định miền được dữ liệu quyết định.} Hai đặc trưng đi qua điểm đang xét phải gặp đoạn dữ liệu phù hợp. Với phương trình sóng $v_{tt}=v_{zz}$ trên $z>0$, dữ liệu lúc $t=0$ xác định duy nhất miền $z>t$; ở $0<z<t$, cần thêm dữ liệu biên $z=0$. Với đặc trưng cong, dùng miền trong các tọa độ mới rồi đổi lại.

\textbf{Bước 6 — Kiểm tra.} Thế nghiệm vào cả PDE và từng điều kiện. Nếu đề yêu cầu hai nghiệm, chỉ rõ hai lựa chọn khác nhau cho hàm tùy ý và chứng minh phần chênh lệch thỏa phương trình thuần nhất, đồng thời triệt tiêu trên đường dữ liệu.
\subsection*{Lời giải áp dụng cho bài}
\textbf{Bước chung.} Đặt $\xi=x-y$, $\eta=x-2y$ như Ví dụ 6. Phương trình thành $v_{\xi\eta}+v_\eta=0$, nên
\[v=F(\xi)+e^{-\xi}G(\eta),\qquad
u_y=-F'(\xi)+e^{-\xi}G(\eta)-2e^{-\xi}G'(\eta).\]
Các công thức dữ liệu dưới đây đều được suy ra từ hai dòng này.

\textbf{(i) Cauchy trên $y=0$.} Lúc này $\xi=\eta=x$:
\[F(x)+e^{-x}G(x)=u_0(x),\qquad
-F'(x)+e^{-x}G(x)-2e^{-x}G'(x)=u_1(x).\]
Lấy đạo hàm phương trình thứ nhất theo $x$ rồi cộng với phương trình thứ hai cho $-e^{-x}G'=u_0'+u_1$. Do đó $G'=-e^x(u_0'+u_1)$ và $F=u_0-e^{-x}G$. Hằng số của $G$ triệt tiêu khi thế lại:
\[\boxed{u(x,y)=u_0(x-y)-e^{-(x-y)}\int_{x-y}^{x-2y}e^t[u_0'(t)+u_1(t)]dt.}\]
Hai đặc trưng qua $(x,y)$ gặp $y=0$ ở $x-y$ và $x-2y$, nên dữ liệu trên toàn trục xác định duy nhất toàn mặt phẳng, với giả thiết trơn thích hợp.

\textbf{(ii) Cauchy trên đặc trưng $y=x$.} Khi đó $\xi=0$, $\eta=-x$. Dữ liệu giá trị cho $F(0)+G(-x)=\varphi(x)$, suy ra $G'(-x)=-\varphi'(x)$. Thế vào dữ liệu đạo hàm:
\[\psi(x)=\varphi(x)+2\varphi'(x)-F(0)-F'(0).\]
Vậy điều kiện cần và đủ là
\[\boxed{\psi-\varphi-2\varphi'=C\quad\text{là hằng số}.}\]
Chuẩn hóa $F(0)=0$ bằng cách chuyển một số hạng giữa $F$ và $G$. Khi đó $G(\eta)=\varphi(-\eta)$, $F'(0)=-C$, và toàn bộ nghiệm là
\[u=F(x-y)+e^{-(x-y)}\varphi(2y-x),\quad F(0)=0,\quad F'(0)=-C.\]
Chọn $F(s)=-Cs$ hoặc $F(s)=-Cs+s^2$ cho hai nghiệm khác nhau. Đường dữ liệu là đặc trưng nên không có duy nhất.

\textbf{(iii) Darboux--Goursat.} Hai đường $y=x$ và $x=2y$ tương ứng $\xi=0$, $\eta=0$. Tại giao điểm, bắt buộc $u_0(0)=u_1(0)=k$. Chọn $F(0)=0$: dữ liệu thứ nhất cho $G(\eta)=u_0(-\eta)$; dữ liệu thứ hai cho $F(\xi)=u_1(\xi)-e^{-\xi}k$. Do đó
\[\boxed{u=u_1(x-y)+e^{-(x-y)}[u_0(2y-x)-k].}\]
Thế $(x,x)$ và $(2x,x)$ kiểm tra lần lượt được $u_0(x)$, $u_1(x)$. Khi điều kiện tại giao điểm đúng và dữ liệu đủ trơn, công thức là nghiệm duy nhất.

\textbf{Chi tiết giải PDE và khử hằng số.} Phần chính từ phép đổi $\xi=x-y,\eta=x-2y$ bằng $-v_{\xi\eta}$; $u_x+u_y=-v_\eta$. Do đó PDE là $v_{\xi\eta}+v_\eta=0$. Đặt $w=v_\eta$, nhân $e^\xi$ cho $(e^\xi w)_\xi=0$, nên $w=e^{-\xi}A(\eta)$; tích phân theo $\eta$ cho $v=F(\xi)+e^{-\xi}G(\eta)$. Ở (i),
\[G(\eta)-G(\xi)=-\int_\xi^\eta e^r[u_0'(r)+u_1(r)]dr,\quad
v=u_0(\xi)+e^{-\xi}[G(\eta)-G(\xi)].\]
Mọi hằng số cộng vào $G$ mất trong hiệu này. Ở (iii), tại $(x,x)$ có $\xi=0,\eta=-x$, nên $u=u_1(0)+u_0(x)-k=u_0(x)$; tại $(2x,x)$ có $\xi=x,\eta=0$, nên $u=u_1(x)+e^{-x}(u_0(0)-k)=u_1(x)$.

\clearpage\noindent\texttt{\detokenize{VD_2_2_09}}\par
\section*{Lời giải bổ sung — Ví dụ 9 · Mục 2.2}

\subsection*{Kiến thức và trình tự giải}
\textbf{Mục tiêu.} Giải phương trình hyperbolic sau khi đã tìm hai tọa độ đặc trưng độc lập.

\textbf{Bước 1 — Phân loại, tìm đặc trưng và thay đạo hàm.} Làm đủ các bước của phép đưa về chính tắc. Viết cả phép đổi ngược để biểu diễn vế phải $\widetilde f(\xi,\eta)=f(x(\xi,\eta),y(\xi,\eta))$.

\textbf{Bước 2 — Tích phân phương trình chính tắc.} Với $v_{\xi\eta}=g$, giữ biến kia cố định khi tích phân:


\[
v=F(\xi)+G(\eta)+\int_{\xi_0}^{\xi}\int_{\eta_0}^{\eta}g(a,b)\,db\,da.
\]


Nếu tìm được một nghiệm riêng $P(x,y)$ dễ hơn, viết $u=P+F(\xi)+G(\eta)$. Với $v_{\xi\eta}+\alpha v_\eta=0$, đặt $w=v_\eta$, giải $w_\xi+\alpha w=0$ rồi tích phân theo $\eta$; thu được $v=F(\xi)+e^{-\alpha\xi}G(\eta)$.

\textbf{Bước 3 — Đổi cả dữ liệu.} Thế tọa độ đường dữ liệu vào $\xi,\eta$, dùng quy tắc dây chuyền cho đạo hàm được cho. Trên đường không đặc trưng, hai dữ liệu thường xác định hai hàm $F,G$. Trên đường đặc trưng, có biến không thay đổi, nên dữ liệu chỉ xác định giá trị hoặc đạo hàm của một hàm tại một điểm; phải tìm điều kiện tương thích và nêu phần tùy ý còn lại.

\textbf{Bước 4 — Lập đẳng thức hình bình hành.} Với $v=F(\xi)+G(\eta)$, bốn đỉnh theo hai họ đặc trưng thỏa $v_{11}+v_{22}=v_{12}+v_{21}$. Nếu có số hạng mũ, trước hết nhân trọng số để được tổng hai hàm một biến. Với vế phải $g$, hiệu bốn đỉnh bằng tích phân kép của $g$ trên hình chữ nhật trong tọa độ mới.

\textbf{Bước 5 — Xác định miền được dữ liệu quyết định.} Hai đặc trưng đi qua điểm đang xét phải gặp đoạn dữ liệu phù hợp. Với phương trình sóng $v_{tt}=v_{zz}$ trên $z>0$, dữ liệu lúc $t=0$ xác định duy nhất miền $z>t$; ở $0<z<t$, cần thêm dữ liệu biên $z=0$. Với đặc trưng cong, dùng miền trong các tọa độ mới rồi đổi lại.

\textbf{Bước 6 — Kiểm tra.} Thế nghiệm vào cả PDE và từng điều kiện. Nếu đề yêu cầu hai nghiệm, chỉ rõ hai lựa chọn khác nhau cho hàm tùy ý và chứng minh phần chênh lệch thỏa phương trình thuần nhất, đồng thời triệt tiêu trên đường dữ liệu.

\textbf{Biến đổi riêng của bài.} Dùng tọa độ cực (hoặc cực sau khi co giãn trục). Khi thu được $r v_{r\theta}-v_\theta=0$, đặt $w=v/r$ thì $w_{r\theta}=0$. Vì vậy $v=rF(\theta)+G(r)$; phải chia nghiệm cho $r$ trước khi lập đẳng thức bốn đỉnh.
\subsection*{Lời giải áp dụng cho bài}
\textbf{(i) Phân loại và tọa độ cực.} Với $A=xy$, $2B=y^2-x^2$, $C=-xy$,
\[\delta=\frac{(y^2-x^2)^2}{4}+x^2y^2=\frac{(x^2+y^2)^2}{4}>0\quad(y>0).\]
Đặc trưng phân tích thành $(xdy-ydx)(ydy+xdx)=0$: các tia $x/y=\mathrm{const}$ và nửa đường tròn $x^2+y^2=\mathrm{const}$. Chọn $r=\sqrt{x^2+y^2}>0$, $\theta\in(0,\pi)$, $x=r\cos\theta$, $y=r\sin\theta$. Phép đổi ngược có Jacobian $r\ne0$.

\textbf{Tính toán phần chính.} Thế ba công thức đạo hàm cực ở Ví dụ 5; các hạng $v_{rr},v_{\theta\theta},v_r$ triệt tiêu, còn
\[xyu_{xx}+(y^2-x^2)u_{xy}-xyu_{yy}=-r v_{r\theta}+v_\theta.\]
Dạng chính tắc là $r v_{r\theta}-v_\theta=0$.

\textbf{(ii) Nghiệm tổng quát.} Đặt $w=v/r$. Vì $v_{r\theta}=w_\theta+r w_{r\theta}$, phương trình cho $r^2w_{r\theta}=0$. Tích phân hai lần: $w=F(\theta)+K(r)$, tức
\[\boxed{u=rF(\theta)+G(r).}\]

\textbf{(iii) Dữ liệu $y=1$.} Từ $u(x,1)=x^2$ có $u_x(x,1)=2x$; kết hợp $u_y(x,1)=4$, quy tắc dây chuyền cho
\[v_\theta=-y u_x+xu_y=2x\quad(y=1).\]
Trên đường này $r=1/\sin\theta$, $x=\cot\theta$, nên $rF'(\theta)=2\cot\theta$ và $F'=2\cos\theta$. Lấy $F=2\sin\theta$ (hằng số được hấp thụ vào $G$), dữ liệu giá trị cho $G(r)=r^2-3$, $r\ge1$. Nghiệm cụ thể
\[\boxed{u=x^2+y^2+2y-3}\]
thỏa PDE vì $u_{xx}=u_{yy}=2$, $u_{xy}=0$, và thỏa cả hai dữ liệu. Công thức đa thức xác định trên toàn $y>0$, nhưng dữ liệu chỉ quyết định duy nhất miền $x^2+y^2>1$, $y>0$ (và giá trị giới hạn tại $r=1$). Ở $r<1$ có thể đổi $G$ bằng một hàm $C^2$ triệt tiêu cùng hai đạo hàm tại $r=1$.

\textbf{(iv) Dữ liệu trên nửa đường tròn đơn vị.} Pháp tuyến ngoài là $\nu=(x,y)$, nên $\partial_\nu u=v_r$ tại $r=1$. Dữ liệu cho
\[F(\theta)+G(1)=\varphi(\cos\theta),\qquad F(\theta)+G'(1)=\psi(\cos\theta).\]
Điều kiện cần và đủ là $\psi(x)-\varphi(x)=C$ hằng trên $(-1,1)$. Toàn bộ nghiệm có thể viết
\[u=r\varphi(x/r)+H(r),\qquad H(1)=0,\quad H'(1)=C.\]
Ví dụ $H=C(r-1)$ hoặc $H=C(r-1)+(r-1)^2$; vì thế dữ liệu trên một đặc trưng không xác định duy nhất. Giả sử $\varphi$ và $H$ đủ trơn trên miền xét.

\textbf{(v) Hai dữ liệu giá trị và ghép nghiệm.} Trên $r=1$ có $F(\theta)+G(1)=1$ cho mọi $\theta$, nên $F$ hằng; hấp thụ nó vào $G$ được $u=H(r)$. Dữ liệu $u(x,0)=x$, $1\le x\le2$ cho $H(r)=r$ ở $1\le r\le2$. Vậy
\[\boxed{u=\sqrt{x^2+y^2}\quad\text{trên }1<x^2+y^2<4,\ y>0.}\]
Đây là miền duy nhất do hai dữ liệu quyết định. Với $0<r\le1$, ghép $u=h(r)$: các điều kiện $h(1)=1$, $h'(1)=1$, $h''(1)=0$ làm giá trị và các đạo hàm đến cấp hai khớp tại $r=1$. Mọi hàm chỉ phụ thuộc $r$ đều thỏa PDE, nên phép ghép cho nghiệm $C^2$ trên $0<r<2$, $y>0$, liên tục đến phần biên đã cho. Ngoài $r=2$, có thể tiếp tục bằng $u=r$ hoặc một phần ghép đủ trơn khác.

\textbf{Chi tiết gom năm hệ số trong tọa độ cực.} Viết $c=\cos\theta,s=\sin\theta$, $xy=r^2cs$, $y^2-x^2=r^2(s^2-c^2)$. Dùng ba đạo hàm cực, hệ số lần lượt là
\[\begin{aligned}
v_{rr}:&\quad r^2[cs\,c^2+(s^2-c^2)cs-cs\,s^2]=0,\\
v_{r\theta}:&\quad r[-2c^2s^2-(c^2-s^2)^2-2c^2s^2]=-r,\\
v_{\theta\theta}:&\quad cs\,s^2-(s^2-c^2)cs-cs\,c^2=0,\\
v_r:&\quad r[cs\,s^2-(s^2-c^2)cs-cs\,c^2]=0,\\
v_\theta:&\quad 2c^2s^2+(c^2-s^2)^2+2c^2s^2=1.
\end{aligned}\]
Đã dùng $4c^2s^2+(c^2-s^2)^2=(c^2+s^2)^2=1$. Vì vậy toán tử đúng bằng $-rv_{r\theta}+v_\theta$, kể cả tại các hướng có $x=0$ trong miền $y>0$.

\clearpage\noindent\texttt{\detokenize{BT_2_3_01}}\par
\section*{Lời giải bổ sung — Bài tập 1 · Mục 2.3}

\subsection*{Kiến thức và trình tự giải}
\textbf{Mục tiêu.} Tính đạo hàm khi đổi tọa độ rồi thay vào PDE. Đổi tọa độ của điểm và đổi biểu thức của hàm là hai việc cần làm cùng nhau: đặt $q=q(x)$ và $v(q)=u(x(q))$.

\textbf{Bước 1 — Kiểm tra phép đổi biến.} Tính $J=\det(\partial q/\partial x)$. Chỉ dùng các công thức trên miền $J\ne0$ và chọn một nhánh liên tục cho các góc. Tọa độ cực/trụ loại trục $r=0$; tọa độ cầu còn loại $\sin\theta=0$. Đây là suy biến của tọa độ, chưa phải kết luận rằng hàm $u$ không tồn tại ở đó.

\textbf{Bước 2 — Tính đạo hàm cấp một bằng quy tắc dây chuyền.} Nếu $b_{ik}=\partial q_k/\partial x_i$, thì


\[
u_{x_i}=\sum_k b_{ik}v_{q_k},\qquad D_i=\sum_k b_{ik}\partial_{q_k}.
\]


Với $q=Ax$, $b_{ik}=a_{ki}$ là hằng số. Với tọa độ cong, lấy nghịch đảo ma trận $\partial x/\partial q$ để tìm $b_{ik}$.

\textbf{Bước 3 — Lấy đạo hàm cấp hai cả hàm lẫn hệ số.} Áp dụng đạo hàm tích cho từng hạng:


\[
u_{x_ix_j}=\sum_{k,l}b_{ik}b_{jl}v_{q_kq_l}
+\sum_{k,l}b_{ik}(\partial_{q_k}b_{jl})v_{q_l}.
\]


Hạng thứ hai chỉ mất đi khi phép đổi biến tuyến tính. Ví dụ $\partial_x=\cos\theta\,\partial_r-(\sin\theta/r)\partial_\theta$ thì lần đạo hàm tiếp theo phải tác động lên cả $\cos\theta$ và $\sin\theta/r$.

\textbf{Bước 4 — Thu gọn toán tử.} Gom hệ số của từng $v_{q_kq_l}$ và $v_{q_k}$ trước khi cộng. Với Laplace, lấy tổng các đạo hàm $u_{x_ix_i}$. Trong hệ trực giao có hệ số độ dài $h_1,h_2,h_3$, có thể đối chiếu bằng


\[
\Delta u=\frac1{h_1h_2h_3}\sum_{k=1}^3
\partial_{q_k}\left(\frac{h_1h_2h_3}{h_k^2}v_{q_k}\right).
\]


\textbf{Bước 5 — Kiểm tra bằng hàm đơn giản.} Thử $u=1$, $u=x_i$ và $u=x_1^2+x_2^2+x_3^2$. Công thức Laplace phải lần lượt cho $0$, $0$, $6$. Với Hessian biểu diễn trong cơ sở trực chuẩn, đổi lại cơ sở Descartes bằng $H_x=E H_q E^T$, không nhầm Hessian với ma trận chỉ gồm các đạo hàm $v_{q_kq_l}$.
\subsection*{Lời giải áp dụng cho bài}
\textbf{(a) Phép đổi biến tuyến tính.} Đặt $\xi=Ax$, $\det A\ne0$, $v(\xi)=u(A^{-1}\xi)$. Vì $\partial\xi_k/\partial x_i=a_{ki}$ là hằng,
\[u=v,\qquad u_{x_i}=\sum_{k=1}^3a_{ki}v_{\xi_k},\qquad
u_{x_ix_j}=\sum_{k,l=1}^3a_{ki}a_{lj}v_{\xi_k\xi_l}.\]
Đây cho đủ mọi đa chỉ số $|\alpha|\le2$, gồm cả sáu đạo hàm cấp hai độc lập. Viết ma trận: $\nabla_xu=A^T\nabla_\xi v$, $H_xu=A^T H_\xi v A$.

\textbf{(b) Tọa độ trụ.} Viết $c=\cos\theta$, $s=\sin\theta$, $z=x_3$. Nghịch đảo vi phân $x_1=rc$, $x_2=rs$ cho
\[D_1=c\partial_r-s\partial_\theta/r,\quad D_2=s\partial_r+c\partial_\theta/r,\quad D_3=\partial_z.\]
Suy ra $u_1=cv_r-sv_\theta/r$, $u_2=sv_r+cv_\theta/r$, $u_3=v_z$. Đạo hàm tích thêm một lần:
\begin{align*}
u_{11}&=c^2v_{rr}-2scv_{r\theta}/r+s^2v_{\theta\theta}/r^2+s^2v_r/r+2scv_\theta/r^2,\\
u_{22}&=s^2v_{rr}+2scv_{r\theta}/r+c^2v_{\theta\theta}/r^2+c^2v_r/r-2scv_\theta/r^2,\\
u_{12}&=scv_{rr}+(c^2-s^2)v_{r\theta}/r-scv_{\theta\theta}/r^2-scv_r/r-(c^2-s^2)v_\theta/r^2,\\
u_{13}&=cv_{rz}-sv_{\theta z}/r,\quad
u_{23}=sv_{rz}+cv_{\theta z}/r,\quad u_{33}=v_{zz}.
\end{align*}

\textbf{(c), (d) Công thức tường minh chung cho cầu và xuyến.} Dùng $s=\sin\theta$, $c=\cos\theta$, $p=\cos\varphi$, $q=\sin\varphi$. Đặt $\rho=r s$ ở hệ cầu; đặt $\rho=R+rs$ ở hệ xuyến. Các đạo hàm cấp một là
\[u_1=spv_r+\frac{cp}{r}v_\theta-\frac q\rho v_\varphi,\quad
u_2=sqv_r+\frac{cq}{r}v_\theta+\frac p\rho v_\varphi,\quad
u_3=cv_r-\frac sr v_\theta.\]
Để trình bày đủ sáu đạo hàm cấp hai mà không lặp các biểu thức dài, đặt sáu tổ hợp đã tính hết hệ số:
\begin{align*}
H_{rr}&=v_{rr},& H_{r\theta}&=v_{r\theta}/r-v_\theta/r^2,\\
H_{r\varphi}&=v_{r\varphi}/\rho-sv_\varphi/\rho^2,&
H_{\theta\theta}&=v_{\theta\theta}/r^2+v_r/r,\\
H_{\theta\varphi}&=v_{\theta\varphi}/(r\rho)-cv_\varphi/\rho^2,&
H_{\varphi\varphi}&=v_{\varphi\varphi}/\rho^2+sv_r/\rho+cv_\theta/(r\rho).
\end{align*}
Đây là các thành phần Hessian trong cơ sở trực chuẩn
\[e_r=(sp,sq,c),\quad e_\theta=(cp,cq,-s),\quad e_\varphi=(-q,p,0).\]
Có thể suy ra, chẳng hạn, $H_{r\theta}=\partial_r(v_\theta/r)$ và $H_{r\varphi}=\partial_r(v_\varphi/\rho)$ vì cơ sở không đổi theo $r$. Với $H_{\varphi\varphi}$, đạo hàm của $e_\varphi$ theo cung tròn là $-(se_r+ce_\theta)/\rho$, tạo ra hai hạng cấp một. Đổi lại cơ sở Descartes cho các công thức cụ thể:
\begin{align*}
u_{11}&=s^2p^2H_{rr}+c^2p^2H_{\theta\theta}+q^2H_{\varphi\varphi}
+2scp^2H_{r\theta}-2spqH_{r\varphi}-2cpqH_{\theta\varphi},\\
u_{22}&=s^2q^2H_{rr}+c^2q^2H_{\theta\theta}+p^2H_{\varphi\varphi}
+2scq^2H_{r\theta}+2spqH_{r\varphi}+2cpqH_{\theta\varphi},\\
u_{33}&=c^2H_{rr}+s^2H_{\theta\theta}-2scH_{r\theta},\\
u_{12}&=s^2pqH_{rr}+c^2pqH_{\theta\theta}-pqH_{\varphi\varphi}
+2scpqH_{r\theta}+s(p^2-q^2)H_{r\varphi}+c(p^2-q^2)H_{\theta\varphi},\\
u_{13}&=scp(H_{rr}-H_{\theta\theta})+(c^2-s^2)pH_{r\theta}-cqH_{r\varphi}+sqH_{\theta\varphi},\\
u_{23}&=scq(H_{rr}-H_{\theta\theta})+(c^2-s^2)qH_{r\theta}+cpH_{r\varphi}-spH_{\theta\varphi}.
\end{align*}
Thế $\rho=rs$ hoặc $\rho=R+rs$ vào sáu $H$ cho đúng từng hệ tọa độ; không còn hệ số đạo hàm nào chưa tính. Ở hệ cầu, ví dụ $H_{\varphi\varphi}=v_{\varphi\varphi}/(r^2s^2)+v_r/r+\cot\theta\,v_\theta/r^2$.

\textbf{(e) Laplace và miền hợp lệ.} Lấy tổng $u_{11}+u_{22}+u_{33}$:
\begin{align*}
\Delta u\ \text{(tuyến tính)}&=\sum_{k,l}(AA^T)_{kl}v_{\xi_k\xi_l},\\
\Delta u\ \text{(trụ)}&=v_{rr}+v_r/r+v_{\theta\theta}/r^2+v_{zz},\\
\Delta u\ \text{(cầu)}&=v_{rr}+2v_r/r+(v_{\theta\theta}+\cot\theta\,v_\theta)/r^2+v_{\varphi\varphi}/(r^2\sin^2\theta),\\
\Delta u\ \text{(xuyến)}&=v_{rr}+\left(\frac1r+\frac{\sin\theta}{\rho}\right)v_r
+\frac{v_{\theta\theta}}{r^2}+\frac{\cos\theta}{r\rho}v_\theta+\frac{v_{\varphi\varphi}}{\rho^2}.
\end{align*}
Ở xuyến, cùng công thức viết dạng phân kỳ là $\Delta u=(r\rho)^{-1}[\partial_r(r\rho v_r)+\partial_\theta(\rho v_\theta/r)+\partial_\varphi(rv_\varphi/\rho)]$, do hệ số độ dài là $(1,r,\rho)$. Trụ cần $r>0$, cầu cần $r>0$, $\sin\theta\ne0$, xuyến cần $r>0$, $\rho>0$; do $r<a<R$, $\rho\ge R-r>0$. Chọn các nhánh góc trên từng miền tọa độ. Kiểm tra $u=x_i$ cho Hessian bằng $0$, và $u=|x|^2$ cho $\Delta u=6$ ở cả ba hệ.

\textbf{Chi tiết tạo Hessian và lấy vết.} Trong hệ trụ, $D_1c=s^2/r$, $D_1(-s/r)=2sc/r^2$, $D_2c=-sc/r$, $D_2(-s/r)=-(c^2-s^2)/r^2$. Thế các hệ số này vào $D_i(cv_r-sv_\theta/r)$ cho $u_{11},u_{12}$; làm tương tự với $sv_r+cv_\theta/r$ cho $u_{22}$. $D_3$ không tác động lên $r,\theta$, nên $u_{13},u_{23}$ chỉ lấy thêm đạo hàm $z$.

Với cầu và xuyến, lấy $e_r=(sp,sq,c)$, $e_\theta=(cp,cq,-s)$, $e_\varphi=(-q,p,0)$. Có $\partial_\theta e_r=e_\theta$, $\partial_\theta e_\theta=-e_r$ và $\partial_\varphi e_r=s e_\varphi$, $\partial_\varphi e_\theta=c e_\varphi$, $\partial_\varphi e_\varphi=-s e_r-c e_\theta$. Gradient là $v_re_r+(v_\theta/r)e_\theta+(v_\varphi/\rho)e_\varphi$. Đạo hàm gradient theo ba hướng $\partial_r,r^{-1}\partial_\theta,\rho^{-1}\partial_\varphi$ rồi chiếu lên cơ sở cho sáu $H$ đã ghi. Chẳng hạn
\[e_\theta\cdot(r^{-1}\partial_\theta\nabla u)=v_r/r+v_{\theta\theta}/r^2,\quad
e_\varphi\cdot(\rho^{-1}\partial_\varphi\nabla u)
=sv_r/\rho+cv_\theta/(r\rho)+v_{\varphi\varphi}/\rho^2.\]
Khi cộng $u_{11}+u_{22}+u_{33}$, hệ số của mỗi $H_{rr},H_{\theta\theta},H_{\varphi\varphi}$ là $1$, các hệ số ngoài đường chéo bằng $0$ bởi trực chuẩn. Vậy
\[\Delta u=v_{rr}+\left(\frac1r+\frac s\rho\right)v_r+
\frac{v_{\theta\theta}}{r^2}+\frac c{r\rho}v_\theta+\frac{v_{\varphi\varphi}}{\rho^2}.\]
Thế $\rho=rs$ cho $1/r+s/\rho=2/r$, $c/(r\rho)=\cot\theta/r^2$; thế $\rho=R+rs$ cho dạng xuyến.

\clearpage\noindent\texttt{\detokenize{BT_2_3_02}}\par
\section*{Lời giải bổ sung — Bài tập 2 · Mục 2.3}

\subsection*{Kiến thức và trình tự giải}
\textbf{Mục tiêu.} Phân loại và đưa phương trình tuyến tính cấp hai về dạng chính tắc, có giữ đầy đủ các hạng cấp thấp.

\textbf{Bước 1 — Chuẩn hóa hệ số.} Viết


\[
A u_{xx}+2B u_{xy}+C u_{yy}+D u_x+E u_y+K u=f.
\]


Hệ số $B$ bằng một nửa hệ số in trước $u_{xy}$. Tính $\delta=B^2-AC$: dương là hyperbolic, âm là elliptic, bằng $0$ là parabolic nếu phần cấp hai còn khác $0$. Phải xét theo từng miền khi hệ số thay đổi hoặc suy biến.

\textbf{Bước 2 — Chọn biến mới.} Với hyperbolic, giải $A(dy)^2-2Bdxdy+C(dx)^2=0$ để tìm hai bất biến $\xi,\eta$ của hai họ đặc trưng. Với parabolic, chỉ có một họ: lấy bất biến $s$ và một biến $z$ độc lập. Với elliptic, không có hai họ đặc trưng thực; hoàn thành bình phương của phần chính để loại hạng hỗn hợp rồi cân bằng hai hệ số dương. Mọi trường hợp đều kiểm tra Jacobian khác $0$.

\textbf{Bước 3 — Viết công thức thay đạo hàm.} Nếu $u(x,y)=v(\xi,\eta)$ thì


\[
\begin{aligned}
u_x&=\xi_xv_\xi+\eta_xv_\eta,\\
u_{xx}&=\xi_x^2v_{\xi\xi}+2\xi_x\eta_xv_{\xi\eta}
+\eta_x^2v_{\eta\eta}+\xi_{xx}v_\xi+\eta_{xx}v_\eta.
\end{aligned}
\]


Làm tương tự cho $u_{xy},u_{yy}$. Có thể dùng toán tử $\partial_x,\partial_y$ để rút ngắn, nhưng phải nêu các đạo hàm của hệ số nếu đổi biến phi tuyến.

\textbf{Bước 4 — Gom hệ số và chia hợp lệ.} Hệ số mới của $v_{\xi\xi}$ là $A\xi_x^2+2B\xi_x\xi_y+C\xi_y^2$; hệ số của $v_{\xi\eta}$ là $2[A\xi_x\eta_x+B(\xi_x\eta_y+\xi_y\eta_x)+C\xi_y\eta_y]$. Hệ số mới của $v_\xi$ là $A\xi_{xx}+2B\xi_{xy}+C\xi_{yy}+D\xi_x+E\xi_y$. Chỉ chia cho hệ số khác $0$ trên miền đang xét.

\textbf{Bước 5 — Với ba biến.} Lập ma trận đối xứng $M$ của phần cấp hai; hệ số hỗn hợp được chia đôi. Dùng hoàn thành bình phương để tìm số hướng dương, âm và bằng $0$. Nếu $\xi=T x$, ma trận mới là $TMT^T$. Cho rõ ma trận $T$ và kiểm tra nó khả nghịch; chỉ ghi dấu của dạng toàn phương chưa đủ để thực hiện phép đổi biến.

\textbf{Chi tiết gom hệ số sau đổi biến.} Với $Lu=Au_{xx}+2Bu_{xy}+Cu_{yy}+Du_x+Eu_y+Ku$, tính riêng


\[
\begin{aligned}
a&=A\xi_x^2+2B\xi_x\xi_y+C\xi_y^2,\\
2b&=2[A\xi_x\eta_x+B(\xi_x\eta_y+\xi_y\eta_x)+C\xi_y\eta_y],\\
c&=A\eta_x^2+2B\eta_x\eta_y+C\eta_y^2,\\
d&=A\xi_{xx}+2B\xi_{xy}+C\xi_{yy}+D\xi_x+E\xi_y,\\
e&=A\eta_{xx}+2B\eta_{xy}+C\eta_{yy}+D\eta_x+E\eta_y.
\end{aligned}
\]


Khi đó $Lu=av_{\xi\xi}+2bv_{\xi\eta}+cv_{\eta\eta}+dv_\xi+ev_\eta+Kv$. Trong lời giải từng ý, thế các số hoặc hàm hệ số vào từng dòng trước khi rút gọn. Với ba biến, tính $TM$ trước, rồi lấy tích vô hướng mỗi hàng của $TM$ với mỗi hàng của $T$ để có $TMT^T$.

\textbf{Lưu ý riêng của bài.} Các đường $\sin x=0$ hoặc $x=0$ có thể đổi loại hoặc làm phép đổi biến mất khả nghịch. Phân loại tại các đường ấy riêng; dạng chính tắc elliptic chỉ viết trên từng miền mở không chứa chúng.
\subsection*{Lời giải áp dụng cho bài}
\textbf{Quy ước.} Viết hệ số hỗn hợp là $2B$, tính $\delta=B^2-AC$; đặt $u=v(\xi,\eta)$. Với các đổi biến $\xi=h(x)$, $\eta=y+k(x)$, công thức chung là
\begin{align*}
u_x&=h'v_\xi+k'v_\eta,\quad u_y=v_\eta,\\
u_{xx}&=h'^2v_{\xi\xi}+2h'k'v_{\xi\eta}+k'^2v_{\eta\eta}+h''v_\xi+k''v_\eta,\\
u_{xy}&=h'v_{\xi\eta}+k'v_{\eta\eta},\quad u_{yy}=v_{\eta\eta}.
\end{align*}
Vì vậy chỉ cần nêu $h',k'$ và gom hệ số như dưới đây. Trong các dạng mới, $x,y$ được hiểu là hàm của $\xi,\eta$ trên miền đổi biến đã chọn.

\textbf{(i.a)} Phần chính $u_{xx}+4u_{xy}+5u_{yy}$ có $\delta=4-5=-1$. Chọn $\xi=x$, $\eta=y-2x$, $J=1$. $h'=1$, $k'=-2$ cho phần chính $v_{\xi\xi}+v_{\eta\eta}$. Đọc hạng cuối theo dạng $xu_x-u$, kết quả
\[v_{\xi\xi}+v_{\eta\eta}+\xi v_\xi-2\xi v_\eta-v=0.\]
\textit{Lưu ý bản in: cuối dòng (a) là chuỗi dấu $=u=0$. Nếu hiểu đúng từng dấu bằng thì còn yêu cầu $u=0$, chỉ có nghiệm không; dạng chính tắc trên dùng cách đọc phương trình $u_{xx}+4u_{xy}+5u_{yy}+xu_x=u$. Việc phân loại phần cấp hai giống nhau trong cả hai cách đọc.}

\textbf{(i.b)} $\delta=16x^2-(\sin^2x+16x^2)=-\sin^2x$. Elliptic trên mỗi miền $\sin x\ne0$, parabolic tại $x\in\pi\mathbb Z$. Chọn $\xi=-\cos x$, $\eta=y-2x^2$: $h'=\sin x$, $h''=\cos x$, $k'=-4x$, $k''=-4$, $J=\sin x$. Hệ số hỗn hợp triệt tiêu, hệ số hai đạo hàm thuần đều là $\sin^2x$:
\[\sin^2x(v_{\xi\xi}+v_{\eta\eta})+(\cos x+\sin^2x)v_\xi-(4+4x\sin x)v_\eta-v=0.\]
Có thể chia cho $\sin^2x$ chỉ trên các miền vừa nêu.

\textbf{(i.c)} $\delta=\sin^2x-(\sin^2x+4x^2)=-4x^2$. Chọn $\xi=x^2$, $\eta=y+\cos x$: $h'=2x$, $h''=2$, $k'=-\sin x$, $k''=-\cos x$, $J=2x$. Thu được
\[4x^2(v_{\xi\xi}+v_{\eta\eta})+(2+2x^2)v_\xi-(\cos x+x\sin x+1)v_\eta=0.\]
Elliptic trên $x>0$ và $x<0$; trên $x=0$ phần chính chỉ còn $u_{xx}$, nên parabolic. Không dùng $x^2$ như tọa độ khả nghịch qua $x=0$.

\textbf{(i.d)} Đặt $a(x)=9\sin^2x+16x^2$, $a'=18\sin x\cos x+32x$. $\delta=16x^2-a=-9\sin^2x$. Trên từng khoảng $\sin x\ne0$, chọn
\[\xi=\int_{x_*}^x\frac{3\sin t}{a(t)}dt,\qquad
\eta=y-\int_{x_*}^x\frac{4t}{a(t)}dt.\]
Đặt $h=\xi_x=3\sin x/a$, $k=\eta_x=-4x/a$. Hệ số hỗn hợp $2h(ak+4x)=0$; hai hệ số thuần đều $9\sin^2x/a$. Hệ số cấp một tính rõ:
\[a h'+h=3\cos x-\frac{3\sin x\,a'}a+\frac{3\sin x}a,\qquad
ak'+k=-4+\frac{4xa'}a-\frac{4x}a.\]
Dạng chính tắc:
\[\frac{9\sin^2x}a(v_{\xi\xi}+v_{\eta\eta})+(ah'+h)v_\xi+(ak'+k)v_\eta-v=0.\]
Ở ý này hệ số của $u_x$ trong đề là $1$, vì vậy không nhân thêm $x$ vào hai hệ số cấp một.
Jacobian $h\ne0$. Tại $x\in\pi\mathbb Z\setminus\{0\}$, phần chính là $(4x\partial_x+\partial_y)^2$ về các hạng cấp hai, hạng một không được tự lấy từ bình phương này; đó là parabolic. Tại $x=0$, phần chính $u_{yy}$ vẫn khác $0$, cũng parabolic.

\textbf{(i.e)} $\delta=0$ mọi nơi. Đặt $\xi=x$, $\eta=y-\cos x$, $k'=\sin x$, $k''=\cos x$, $J=1$. Khi thế vào phần chính $u_{xx}-2\sin x\,u_{xy}+\sin^2x\,u_{yy}$, còn $v_{\xi\xi}+\cos x\,v_\eta$. Dạng parabolic là
\[v_{\xi\xi}+y v_\xi+(\cos x+y\sin x-x)v_\eta=0.\]

\textbf{(i.f)} $\delta=\sin^2x+\cos^2x+3=4>0$. Đặc trưng $dy/dx=\sin x\pm2$ cho $y+\cos x\mp2x=\mathrm{const}$. Chọn $\xi=y+\cos x-2x$, $\eta=y+\cos x+2x$, $J=-4$. Tính $\xi_{xx}=\eta_{xx}=-\cos x$, $\xi_y=\eta_y=1$; phần chính thuần triệt tiêu, hỗn hợp bằng $-16v_{\xi\eta}$. Do đó
\[-16v_{\xi\eta}+(x-\cos x)(v_\xi+v_\eta)=xy.\]

\textbf{(ii) Ba biến: đưa về chính tắc bằng ma trận.} Gọi $M$ là ma trận đối xứng của phần chính. Nếu dạng toàn phương $\lambda^TM\lambda=\sum_j\varepsilon_j(b_j\cdot\lambda)^2$ và $B$ có các hàng $b_j$, thì $M=B^T\operatorname{diag}(\varepsilon_j)B$. Chọn tọa độ mới $\xi=B^{-T}(x,y,z)^T$ để ma trận mới đúng bằng $\operatorname{diag}(\varepsilon_j)$.

\textbf{(ii.a)} Hoàn thành bình phương:
\[Q=(\lambda_x+\lambda_y+\lambda_z)^2+(\lambda_y-\lambda_z)^2+\lambda_z^2.\]
Ba dấu dương: elliptic. Chọn $\xi=x$, $\eta=y-x$, $\zeta=z+y-2x$; định thức $1$. Dạng chính tắc $v_{\xi\xi}+v_{\eta\eta}+v_{\zeta\zeta}=0$.

\textbf{(ii.b)}
\[Q=(\lambda_x+2\lambda_y+6\lambda_z)^2+3\lambda_y^2-3(\lambda_y+3\lambda_z)^2.\]
Hai dấu dương, một dấu âm: hyperbolic. Chọn
\[\xi=x,\quad\eta=\frac{3y-z}{3\sqrt3},\quad\zeta=\frac{z-6x}{3\sqrt3}.\]
Định thức $1/9\ne0$. Dạng chính tắc $v_{\xi\xi}+v_{\eta\eta}-v_{\zeta\zeta}=0$.

\textbf{(ii.c)}
\[Q=3(\lambda_x-\lambda_y/3+\lambda_z/3)^2+\frac53(\lambda_y+7\lambda_z/5)^2-\frac35\lambda_z^2.\]
Hai dấu dương, một dấu âm: hyperbolic. Chọn
\[\xi=\frac{x}{\sqrt3},\quad\eta=\sqrt{\frac35}\left(y+\frac x3\right),\quad
\zeta=\sqrt{\frac53}\left(z-\frac75y-\frac45x\right).\]
Định thức $1/\sqrt3\ne0$. Dạng chính tắc cũng là $v_{\xi\xi}+v_{\eta\eta}-v_{\zeta\zeta}=0$. Với cả ba ý, có thể kiểm tra trực tiếp bằng phép nhân $TMT^T$; các hạng hỗn hợp đều bằng $0$.

\textbf{Chi tiết gom hệ số cho sáu phương trình hai biến.} Với $Lu=Au_{xx}+2Bu_{xy}+Cu_{yy}+Du_x+Eu_y+Ku$, phép đổi $\xi=h(x),\eta=y+k(x)$ cho
\[\begin{aligned}
L u={}&Ah'^2v_{\xi\xi}+2h'(Ak'+B)v_{\xi\eta}
+(Ak'^2+2Bk'+C)v_{\eta\eta}\\
&+(Ah''+Dh')v_\xi+(Ak''+Dk'+E)v_\eta+Kv.
\end{aligned}\]
Các hệ số dưới đây được tính trước khi chia bất kỳ biểu thức nào:

(a) $A=1,B=2,C=5,D=x$: hỗn hợp $2(1)(-2+2)=0$; thuần theo $\eta$ là $4-8+5=1$. Cấp một: $0+x=x$, $0-2x=-2x$.

(b) $A=1,B=4x,C=\sin^2x+16x^2,D=\sin x$: hỗn hợp $2\sin x(-4x+4x)=0$; thuần theo $\eta$ là $16x^2-32x^2+\sin^2x+16x^2=\sin^2x$. Cấp một: $\cos x+\sin x\sin x=\cos x+\sin^2x$, $-4+\sin x(-4x)=-4-4x\sin x$.

(c) $A=1,B=\sin x,C=\sin^2x+4x^2,D=x,E=-1$: hỗn hợp $4x(-\sin x+\sin x)=0$; thuần theo $\eta$ là $\sin^2x-2\sin^2x+\sin^2x+4x^2=4x^2$. Cấp một: $2+x(2x)=2+2x^2$, $-\cos x-x\sin x-1$.

(d) Giữ ký hiệu $h=\xi_x,k=\eta_x$ của ý này: $Ah^2=a(3\sin x/a)^2=9\sin^2x/a$; hỗn hợp $2h(ak+4x)=0$; thuần theo $\eta$ là
\[\frac{16x^2}a-\frac{32x^2}a+1=\frac{a-16x^2}a=\frac{9\sin^2x}a.\]
Đạo hàm thương $h'=3\cos x/a-3\sin x\,a'/a^2$, $k'=-4/a+4xa'/a^2$ rồi nhân $a$ và cộng $h,k$ cho hai hệ số cấp một đã ghi.

(e) $h'=1,k'=\sin x,A=1,B=-\sin x,C=\sin^2x,D=y,E=-x$: hỗn hợp $2(\sin x-\sin x)=0$, thuần theo $\eta$ là $\sin^2x-2\sin^2x+\sin^2x=0$. Cấp một $y$, $\cos x+y\sin x-x$.

(f) Đặt $p=-\sin x-2,q=-\sin x+2$. Ba đạo hàm là
\[u_{xx}=p^2v_{\xi\xi}+2pqv_{\xi\eta}+q^2v_{\eta\eta}-\cos x(v_\xi+v_\eta),\quad
u_{xy}=pv_{\xi\xi}+(p+q)v_{\xi\eta}+qv_{\eta\eta},\quad
u_{yy}=v_{\xi\xi}+2v_{\xi\eta}+v_{\eta\eta}.\]
Với $B=\sin x,C=\sin^2x-4$, các hệ số thuần là $(p+\sin x)^2-4=0$ và $(q+\sin x)^2-4=0$; hỗn hợp là
\[2[pq+\sin x(p+q)+\sin^2x-4]
=2[(\sin^2x-4)-2\sin^2x+\sin^2x-4]=-16.\]
Hạng $xu_y=x(v_\xi+v_\eta)$ cộng vào hệ số hai đạo hàm cấp một, cho $(x-\cos x)(v_\xi+v_\eta)$.

\textbf{Chi tiết nhân ma trận cho ba phương trình ba biến.} Gọi $T$ là ma trận các hàng gradient tọa độ mới, $N=TMT^T$.

(ii.a)
\[M=\begin{pmatrix}1&1&1\\1&2&0\\1&0&3\end{pmatrix},\quad
T=\begin{pmatrix}1&0&0\\-1&1&0\\-2&1&1\end{pmatrix},\quad
TM=\begin{pmatrix}1&1&1\\0&1&-1\\0&0&1\end{pmatrix}.\]
Nhân từng hàng của $TM$ với ba hàng của $T$: hàng đầu cho $1,-1+1,-2+1+1$; hàng hai cho $0,1,1-1$; hàng ba cho $0,0,1$. Vậy $N=\operatorname{diag}(1,1,1)$.

(ii.b) Đặt $b=1/(3\sqrt3)$, nên
\[M=\begin{pmatrix}1&2&6\\2&4&3\\6&3&9\end{pmatrix},\quad
T=\begin{pmatrix}1&0&0\\0&3b&-b\\-6b&0&b\end{pmatrix},\quad
TM=\begin{pmatrix}1&2&6\\0&9b&0\\0&-9b&-27b\end{pmatrix}.\]
Các tích trên đường chéo là $1$, $27b^2=1$, $-27b^2=-1$; ngoài đường chéo là $6b-6b=0$, $0$, $-27b^2+27b^2=0$. Vậy $N=\operatorname{diag}(1,1,-1)$.

(ii.c) Đặt $a=1/\sqrt3,b=\sqrt{3/5},c=\sqrt{5/3}$. Khi đó
\[M=\begin{pmatrix}3&-1&1\\-1&2&2\\1&2&3\end{pmatrix},\quad
T=\begin{pmatrix}a&0&0\\b/3&b&0\\-4c/5&-7c/5&c\end{pmatrix},\quad
TM=\begin{pmatrix}3a&-a&a\\0&5b/3&7b/3\\0&0&-3c/5\end{pmatrix}.\]
Đường chéo là $3a^2=1,5b^2/3=1,-3c^2/5=-1$. Hai tích hàng đầu với hàng sau là $ab-ab=0$ và $(-12+7+5)ac/5=0$; tích hàng hai với hàng ba là $-7bc/3+7bc/3=0$. Vậy $N=\operatorname{diag}(1,1,-1)$, xác nhận mọi hệ số hỗn hợp bằng $0$.

\clearpage\noindent\texttt{\detokenize{BT_2_3_03}}\par
\section*{Lời giải bổ sung — Bài tập 3 · Mục 2.3}

\subsection*{Kiến thức và trình tự giải}
\textbf{Mục tiêu.} Giải phương trình hyperbolic sau khi đã tìm hai tọa độ đặc trưng độc lập.

\textbf{Bước 1 — Phân loại, tìm đặc trưng và thay đạo hàm.} Làm đủ các bước của phép đưa về chính tắc. Viết cả phép đổi ngược để biểu diễn vế phải $\widetilde f(\xi,\eta)=f(x(\xi,\eta),y(\xi,\eta))$.

\textbf{Bước 2 — Tích phân phương trình chính tắc.} Với $v_{\xi\eta}=g$, giữ biến kia cố định khi tích phân:


\[
v=F(\xi)+G(\eta)+\int_{\xi_0}^{\xi}\int_{\eta_0}^{\eta}g(a,b)\,db\,da.
\]


Nếu tìm được một nghiệm riêng $P(x,y)$ dễ hơn, viết $u=P+F(\xi)+G(\eta)$. Với $v_{\xi\eta}+\alpha v_\eta=0$, đặt $w=v_\eta$, giải $w_\xi+\alpha w=0$ rồi tích phân theo $\eta$; thu được $v=F(\xi)+e^{-\alpha\xi}G(\eta)$.

\textbf{Bước 3 — Đổi cả dữ liệu.} Thế tọa độ đường dữ liệu vào $\xi,\eta$, dùng quy tắc dây chuyền cho đạo hàm được cho. Trên đường không đặc trưng, hai dữ liệu thường xác định hai hàm $F,G$. Trên đường đặc trưng, có biến không thay đổi, nên dữ liệu chỉ xác định giá trị hoặc đạo hàm của một hàm tại một điểm; phải tìm điều kiện tương thích và nêu phần tùy ý còn lại.

\textbf{Bước 4 — Lập đẳng thức hình bình hành.} Với $v=F(\xi)+G(\eta)$, bốn đỉnh theo hai họ đặc trưng thỏa $v_{11}+v_{22}=v_{12}+v_{21}$. Nếu có số hạng mũ, trước hết nhân trọng số để được tổng hai hàm một biến. Với vế phải $g$, hiệu bốn đỉnh bằng tích phân kép của $g$ trên hình chữ nhật trong tọa độ mới.

\textbf{Bước 5 — Xác định miền được dữ liệu quyết định.} Hai đặc trưng đi qua điểm đang xét phải gặp đoạn dữ liệu phù hợp. Với phương trình sóng $v_{tt}=v_{zz}$ trên $z>0$, dữ liệu lúc $t=0$ xác định duy nhất miền $z>t$; ở $0<z<t$, cần thêm dữ liệu biên $z=0$. Với đặc trưng cong, dùng miền trong các tọa độ mới rồi đổi lại.

\textbf{Bước 6 — Kiểm tra.} Thế nghiệm vào cả PDE và từng điều kiện. Nếu đề yêu cầu hai nghiệm, chỉ rõ hai lựa chọn khác nhau cho hàm tùy ý và chứng minh phần chênh lệch thỏa phương trình thuần nhất, đồng thời triệt tiêu trên đường dữ liệu.
\subsection*{Lời giải áp dụng cho bài}
\textbf{(a) Phân loại và đặc trưng.} $A=3$, $B=-1$, $C=-5$, $\delta=16>0$: hyperbolic. Phương trình đặc trưng $3(dy)^2+2dxdy-5(dx)^2=(dy-dx)(3dy+5dx)=0$ cho các đường $y-x=c_1$, $5x+3y=c_2$. Đặt $\xi=y-x$, $\eta=5x+3y$, $J=-8$; đổi ngược $x=(\eta-3\xi)/8$, $y=(\eta+5\xi)/8$. Có $\partial_x=-\partial_\xi+5\partial_\eta$, $\partial_y=\partial_\xi+3\partial_\eta$, nên
\[3u_{xx}-2u_{xy}-5u_{yy}=-64v_{\xi\eta}=\widetilde f.\]
Các đường đặc trưng cần vẽ là hai chùm thẳng có hệ số góc $1$ và $-5/3$.

\textbf{(b) Tích phân.} Giữ lần lượt $\xi,\eta$ cố định:
\[v=F(\xi)+G(\eta)-\frac1{64}\int_{\xi_0}^{\xi}\int_{\eta_0}^{\eta}\widetilde f(a,b)dbda.\]
Nếu $f=0$, $v=F+G$ và $v_{11}+v_{22}=v_{12}+v_{21}$ cho bốn đỉnh $(\xi_i,\eta_j)$, tức hình bình hành có cạnh theo hai họ đặc trưng trong mặt phẳng gốc.

\textbf{(c) Nghiệm riêng và dữ liệu.} Với $f=x+xy$, thử đa thức; lấy
\[P(x,y)=\frac{x^3(1+y)}{18}+\frac{x^4}{108}.\]
Ta có $P_{xx}=x(1+y)/3+x^2/9$, $P_{xy}=x^2/6$, $P_{yy}=0$, nên $3P_{xx}-2P_{xy}-5P_{yy}=x+xy$. Viết $u=P+F(\xi)+G(\eta)$. Trên $(r,r)$: $\xi=0$, $\eta=8r$, $F(0)+G(8r)=e^r-P(r,r)$. Trên $(3r,-5r)$: $\xi=-8r$, $\eta=0$, $F(-8r)+G(0)=r+1-P(3r,-5r)$. Hai dữ liệu tại gốc đều bằng $1$, nên tương thích. Lấy $F(0)=0$, $G(0)=1$ được
\[F(\xi)=-\xi/8-P(-3\xi/8,5\xi/8),\quad
G(\eta)=e^{\eta/8}-P(\eta/8,\eta/8).\]
Kết quả
\[\boxed{u=P(x,y)-\frac\xi8-P(-3\xi/8,5\xi/8)+e^{\eta/8}-P(\eta/8,\eta/8).}\]
Thế hai đường dữ liệu vào công thức lần lượt cho $e^r$ và $r+1$; phần $F+G$ có đạo hàm hỗn hợp bằng $0$, nên kiểm tra PDE qua $LP=f$. Hai đường là hai đặc trưng giao nhau, dữ liệu trên toàn bộ chúng xác định nghiệm duy nhất trên $\mathbb R^2$. Đồ thị kèm theo biểu diễn công thức này.

\textbf{Chi tiết thế ba đạo hàm cấp hai.} Với tọa độ đã chọn,
\[u_{xx}=v_{\xi\xi}-10v_{\xi\eta}+25v_{\eta\eta},\quad u_{xy}=-v_{\xi\xi}+2v_{\xi\eta}+15v_{\eta\eta},\quad u_{yy}=v_{\xi\xi}+6v_{\xi\eta}+9v_{\eta\eta}.\]
Nhân lần lượt với $3$,$-2$,$-5$, hệ số theo thứ tự $v_{\xi\xi},v_{\xi\eta},v_{\eta\eta}$ là
\[3+2-5=0,\qquad -30-4-30=-64,\qquad 75-30-45=0.\]
\begin{center}\includegraphics[width=.97\linewidth,height=11cm,keepaspectratio]{_xu_ly/graphs/BT_2_3_03.png}\end{center}
\textit{Đồ thị minh họa; với nghiệm số, xem phương pháp và tham số trong mã Python kèm theo.}

\clearpage\noindent\texttt{\detokenize{BT_2_3_04}}\par
\section*{Lời giải bổ sung — Bài tập 4 · Mục 2.3}

\subsection*{Kiến thức và trình tự giải}
\textbf{Mục tiêu.} Giải phương trình hyperbolic sau khi đã tìm hai tọa độ đặc trưng độc lập.

\textbf{Bước 1 — Phân loại, tìm đặc trưng và thay đạo hàm.} Làm đủ các bước của phép đưa về chính tắc. Viết cả phép đổi ngược để biểu diễn vế phải $\widetilde f(\xi,\eta)=f(x(\xi,\eta),y(\xi,\eta))$.

\textbf{Bước 2 — Tích phân phương trình chính tắc.} Với $v_{\xi\eta}=g$, giữ biến kia cố định khi tích phân:


\[
v=F(\xi)+G(\eta)+\int_{\xi_0}^{\xi}\int_{\eta_0}^{\eta}g(a,b)\,db\,da.
\]


Nếu tìm được một nghiệm riêng $P(x,y)$ dễ hơn, viết $u=P+F(\xi)+G(\eta)$. Với $v_{\xi\eta}+\alpha v_\eta=0$, đặt $w=v_\eta$, giải $w_\xi+\alpha w=0$ rồi tích phân theo $\eta$; thu được $v=F(\xi)+e^{-\alpha\xi}G(\eta)$.

\textbf{Bước 3 — Đổi cả dữ liệu.} Thế tọa độ đường dữ liệu vào $\xi,\eta$, dùng quy tắc dây chuyền cho đạo hàm được cho. Trên đường không đặc trưng, hai dữ liệu thường xác định hai hàm $F,G$. Trên đường đặc trưng, có biến không thay đổi, nên dữ liệu chỉ xác định giá trị hoặc đạo hàm của một hàm tại một điểm; phải tìm điều kiện tương thích và nêu phần tùy ý còn lại.

\textbf{Bước 4 — Lập đẳng thức hình bình hành.} Với $v=F(\xi)+G(\eta)$, bốn đỉnh theo hai họ đặc trưng thỏa $v_{11}+v_{22}=v_{12}+v_{21}$. Nếu có số hạng mũ, trước hết nhân trọng số để được tổng hai hàm một biến. Với vế phải $g$, hiệu bốn đỉnh bằng tích phân kép của $g$ trên hình chữ nhật trong tọa độ mới.

\textbf{Bước 5 — Xác định miền được dữ liệu quyết định.} Hai đặc trưng đi qua điểm đang xét phải gặp đoạn dữ liệu phù hợp. Với phương trình sóng $v_{tt}=v_{zz}$ trên $z>0$, dữ liệu lúc $t=0$ xác định duy nhất miền $z>t$; ở $0<z<t$, cần thêm dữ liệu biên $z=0$. Với đặc trưng cong, dùng miền trong các tọa độ mới rồi đổi lại.

\textbf{Bước 6 — Kiểm tra.} Thế nghiệm vào cả PDE và từng điều kiện. Nếu đề yêu cầu hai nghiệm, chỉ rõ hai lựa chọn khác nhau cho hàm tùy ý và chứng minh phần chênh lệch thỏa phương trình thuần nhất, đồng thời triệt tiêu trên đường dữ liệu.
\subsection*{Lời giải áp dụng cho bài}
\textbf{(a) Phân loại và chính tắc.} $A=15$, $B=-1$, $C=-1$, $\delta=16>0$. Đặc trưng $(5dy-dx)(3dy+dx)=0$ cho $x-5y=c_1$, $x+3y=c_2$. Chọn $\xi=x-5y$, $\eta=x+3y$, $J=8$, $x=(3\xi+5\eta)/8$, $y=(\eta-\xi)/8$. Với $\partial_x=\partial_\xi+\partial_\eta$, $\partial_y=-5\partial_\xi+3\partial_\eta$, phần chính là $64v_{\xi\eta}$. Dạng mới $64v_{\xi\eta}=\widetilde f$; hai chùm đặc trưng có hệ số góc $1/5$, $-1/3$.

\textbf{(b) Nghiệm tổng quát và bốn đỉnh.}
\[v=F(\xi)+G(\eta)+\frac1{64}\int_{\xi_0}^{\xi}\int_{\eta_0}^{\eta}\widetilde f(a,b)dbda.\]
Khi $f=0$, tổng ở hai đỉnh đối diện bằng tổng ở hai đỉnh còn lại, do mỗi $F(\xi_i)$ và $G(\eta_j)$ xuất hiện một lần ở mỗi vế.

\textbf{(c) $f=1$: dữ liệu trên hai đường không cùng là đặc trưng.} Chọn $P=x^2/30$ vì $15P_{xx}=1$. Dữ liệu $y=0$ cho $F(s)+G(s)=s-s^2/30$. Trên $(3t,t)$, $\xi=-2t$, $\eta=6t$; đặt $s=-2t$ rồi thay $G(-3s)=-3s-9s^2/30-F(-3s)$, được
\[F(s)-F(-3s)=R(s),\qquad R(s)=\sin(-s/2)+3s+\frac{9s^2}{40}.\]
Đặt $q=-1/3$. Vì $R(0)=0$, lặp ở $qs,q^2s,\ldots$ cho
\[F(s)=F(0)-\sum_{j=1}^{\infty}R(q^js).\]
Lấy $F(0)=0$ và $G(s)=s-s^2/30-F(s)$. Chuỗi hội tụ cùng hai đạo hàm trên mọi đoạn hữu hạn: mỗi lần đạo hàm thêm một hệ số $q^j$, còn $R(q^js)=O(|q|^j)$ ở gần $0$. Vì vậy
\[\boxed{u_c(x,y)=\frac{x^2}{30}+F(x-5y)+G(x+3y).}\]
Tính $F(s)-F(-3s)=R(s)$ bằng cách dịch chỉ số chuỗi kiểm tra dữ liệu thứ hai. Nếu hai nghiệm khác nhau, phần chênh $H$ thỏa $H(s)=H(-3s)$; lặp $s\mapsto qs$ và dùng liên tục ở $0$ cho $H$ hằng, rồi triệt tiêu với $G$. Do đó nghiệm liên tục, $C^2$ ở đây, là duy nhất.

\textbf{(e) $f=e^{x+3y}$: dữ liệu Cauchy trên $y=0$.} Bản gốc đánh số chuyển từ (c) sang (e), không có ý (d). Chọn $P=xe^{x+3y}/24$: trong tọa độ mới $P=(3\xi+5\eta)e^\eta/192$, nên $64P_{\xi\eta}=e^\eta$. Viết $u=P+F(\xi)+G(\eta)$ và đặt
\[H(s)=s-\frac{se^s}{24},\qquad K(s)=\sin s+1-\frac{se^s}{8}.\]
Hai dữ liệu cho $F+G=H$, $-5F'+3G'=K$. Lấy đạo hàm phương trình thứ nhất: $F'+G'=H'$. Giải hệ:
\[F'=\frac{3H'-K}{8},\quad G'=\frac{5H'+K}{8}.\]
Vậy chọn $F(s)=\int_0^s(3H'(r)-K(r))dr/8=s/4+(\cos s-1)/8-(e^s-1)/64$, $G=H-F$, được
\[\boxed{u_e=\frac{xe^{x+3y}}{24}+F(x-5y)+G(x+3y).}\]
Thế $y=0$ cho $u=x$, $u_y=\sin x+1$. Hai đặc trưng qua $(x,y)$ gặp trục ở $A=(x-5y,0)$, $B=(x+3y,0)$; tam giác phụ thuộc có ba đỉnh $A,B,(x,y)$. Công thức chỉ cần dữ liệu trên đoạn giữa $A$ và $B$. Các đồ thị kèm theo phân biệt nghiệm của (c) và (e).

\textbf{Chi tiết phương trình hàm ở (c).} Trên $(3t,t)$, $P=3t^2/10$, nên $F(-2t)+G(6t)=\sin t-3t^2/10$. Với $s=-2t$, $t=-s/2$, $6t=-3s$, ta có
\[F(s)+G(-3s)=\sin(-s/2)-3s^2/40.\]
Thế $G(-3s)=-3s-3s^2/10-F(-3s)$ rồi chuyển vế cho $R(s)$ đã ghi. Đặt $q=-1/3$, phương trình tại $qs$ là $F(qs)-F(s)=R(qs)$; lặp $N$ lần:
\[F(s)=F(q^Ns)-\sum_{j=1}^{N}R(q^js).\]
Cho $N\to\infty$ và dùng liên tục của $F$ tại $0$ được chuỗi.

\textbf{Chi tiết tích phân ở (e).} $H'=1-(1+s)e^s/24$, nên
\[3H'-K=2-\sin s-e^s/8,\qquad F'=1/4-\sin s/8-e^s/64.\]
Tích phân từ $0$ đến $s$ được $F=s/4+(\cos s-1)/8-(e^s-1)/64$, đúng biểu thức sử dụng.

\textbf{Chi tiết thế ba đạo hàm cấp hai.} Với tọa độ đã chọn,
\[u_{xx}=v_{\xi\xi}+2v_{\xi\eta}+v_{\eta\eta},\quad u_{xy}=-5v_{\xi\xi}-2v_{\xi\eta}+3v_{\eta\eta},\quad u_{yy}=25v_{\xi\xi}-30v_{\xi\eta}+9v_{\eta\eta}.\]
Nhân lần lượt với $15$,$-2$,$-1$, hệ số theo thứ tự $v_{\xi\xi},v_{\xi\eta},v_{\eta\eta}$ là
\[15+10-25=0,\qquad 30+4+30=64,\qquad 15-6-9=0.\]
\begin{center}\includegraphics[width=.97\linewidth,height=11cm,keepaspectratio]{_xu_ly/graphs/BT_2_3_04.png}\end{center}
\textit{Đồ thị minh họa; với nghiệm số, xem phương pháp và tham số trong mã Python kèm theo.}

\clearpage\noindent\texttt{\detokenize{BT_2_3_05}}\par
\section*{Lời giải bổ sung — Bài tập 5 · Mục 2.3}

\subsection*{Kiến thức và trình tự giải}
\textbf{Mục tiêu.} Giải phương trình hyperbolic sau khi đã tìm hai tọa độ đặc trưng độc lập.

\textbf{Bước 1 — Phân loại, tìm đặc trưng và thay đạo hàm.} Làm đủ các bước của phép đưa về chính tắc. Viết cả phép đổi ngược để biểu diễn vế phải $\widetilde f(\xi,\eta)=f(x(\xi,\eta),y(\xi,\eta))$.

\textbf{Bước 2 — Tích phân phương trình chính tắc.} Với $v_{\xi\eta}=g$, giữ biến kia cố định khi tích phân:


\[
v=F(\xi)+G(\eta)+\int_{\xi_0}^{\xi}\int_{\eta_0}^{\eta}g(a,b)\,db\,da.
\]


Nếu tìm được một nghiệm riêng $P(x,y)$ dễ hơn, viết $u=P+F(\xi)+G(\eta)$. Với $v_{\xi\eta}+\alpha v_\eta=0$, đặt $w=v_\eta$, giải $w_\xi+\alpha w=0$ rồi tích phân theo $\eta$; thu được $v=F(\xi)+e^{-\alpha\xi}G(\eta)$.

\textbf{Bước 3 — Đổi cả dữ liệu.} Thế tọa độ đường dữ liệu vào $\xi,\eta$, dùng quy tắc dây chuyền cho đạo hàm được cho. Trên đường không đặc trưng, hai dữ liệu thường xác định hai hàm $F,G$. Trên đường đặc trưng, có biến không thay đổi, nên dữ liệu chỉ xác định giá trị hoặc đạo hàm của một hàm tại một điểm; phải tìm điều kiện tương thích và nêu phần tùy ý còn lại.

\textbf{Bước 4 — Lập đẳng thức hình bình hành.} Với $v=F(\xi)+G(\eta)$, bốn đỉnh theo hai họ đặc trưng thỏa $v_{11}+v_{22}=v_{12}+v_{21}$. Nếu có số hạng mũ, trước hết nhân trọng số để được tổng hai hàm một biến. Với vế phải $g$, hiệu bốn đỉnh bằng tích phân kép của $g$ trên hình chữ nhật trong tọa độ mới.

\textbf{Bước 5 — Xác định miền được dữ liệu quyết định.} Hai đặc trưng đi qua điểm đang xét phải gặp đoạn dữ liệu phù hợp. Với phương trình sóng $v_{tt}=v_{zz}$ trên $z>0$, dữ liệu lúc $t=0$ xác định duy nhất miền $z>t$; ở $0<z<t$, cần thêm dữ liệu biên $z=0$. Với đặc trưng cong, dùng miền trong các tọa độ mới rồi đổi lại.

\textbf{Bước 6 — Kiểm tra.} Thế nghiệm vào cả PDE và từng điều kiện. Nếu đề yêu cầu hai nghiệm, chỉ rõ hai lựa chọn khác nhau cho hàm tùy ý và chứng minh phần chênh lệch thỏa phương trình thuần nhất, đồng thời triệt tiêu trên đường dữ liệu.
\subsection*{Lời giải áp dụng cho bài}
\textbf{(a) Phân loại, đặc trưng, đổi biến.} $A=3$, $B=1$, $C=-5$, $\delta=16>0$. Đặc trưng $(dy+dx)(3dy-5dx)=0$: $x+y=c_1$, $5x-3y=c_2$. Đặt $\xi=x+y$, $\eta=5x-3y$, $J=-8$; $x=(3\xi+\eta)/8$, $y=(5\xi-\eta)/8$. Có $\partial_x=\partial_\xi+5\partial_\eta$, $\partial_y=\partial_\xi-3\partial_\eta$, nên
\[3u_{xx}+2u_{xy}-5u_{yy}=64v_{\xi\eta},\quad 3u_x+5u_y=8v_\xi.\]
Dạng chính tắc $64v_{\xi\eta}+8v_\xi=0$.

\textbf{(b) Giải và hình bình hành.} Đặt $w=v_\xi$; ODE $w_\eta+w/8=0$ cho $w=e^{-\eta/8}A(\xi)$. Tích phân theo $\xi$:
\[\boxed{u=e^{-(5x-3y)/8}F(x+y)+G(5x-3y).}\]
Với $v_{ij}=v(\xi_i,\eta_j)$,
\[v_{22}-v_{12}=e^{-(\eta_2-\eta_1)/8}(v_{21}-v_{11}).\]
Thật vậy, lấy hiệu hai giá trị cùng $\eta_j$ làm $G$ triệt tiêu và chỉ còn trọng số mũ nhân $F(\xi_2)-F(\xi_1)$.

\textbf{(c) Điều kiện tương thích.} Trên $(3t,5t)$: $\xi=8t$, $\eta=0$. Dữ liệu giá trị cho $F(8t)+G(0)=8t$, nên $F(s)=s-g_0$, $g_0=G(0)$, $F'=1$. Từ
\[u_y=e^{-\eta/8}\left(F'+\frac38F\right)-3G'(\eta)\]
suy ra $u_y(3t,5t)=1+3t-3g_0/8-3G'(0)$. So với $Ct+1$ cho
\[\boxed{C=3,\qquad G'(0)=-g_0/8.}\]
Chỉ hai giá trị tại $0$ của $G$ bị ràng buộc, nên có vô số nghiệm. Chọn $G=0$ và $G(s)=s^2$, cùng $g_0=0$, cho
\[u_1=(x+y)e^{-(5x-3y)/8},\qquad u_2=u_1+(5x-3y)^2.\]
Phần chênh chỉ phụ thuộc $\eta$, nên thỏa PDE; nó và đạo hàm theo $y$ bằng $0$ khi $\eta=0$. Vì vậy cả hai thỏa đầy đủ dữ liệu. Đặc trưng cần vẽ có hệ số góc $-1$ và $5/3$.

\textbf{Chi tiết thế ba đạo hàm cấp hai.} Với tọa độ đã chọn,
\[u_{xx}=v_{\xi\xi}+10v_{\xi\eta}+25v_{\eta\eta},\quad u_{xy}=v_{\xi\xi}+2v_{\xi\eta}-15v_{\eta\eta},\quad u_{yy}=v_{\xi\xi}-6v_{\xi\eta}+9v_{\eta\eta}.\]
Nhân lần lượt với $3$,$2$,$-5$, hệ số theo thứ tự $v_{\xi\xi},v_{\xi\eta},v_{\eta\eta}$ là
\[3+2-5=0,\qquad 30+4+30=64,\qquad 75-30-45=0.\]
\begin{center}\includegraphics[width=.97\linewidth,height=11cm,keepaspectratio]{_xu_ly/graphs/BT_2_3_05.png}\end{center}
\textit{Đồ thị minh họa; với nghiệm số, xem phương pháp và tham số trong mã Python kèm theo.}

\clearpage\noindent\texttt{\detokenize{BT_2_3_06}}\par
\section*{Lời giải bổ sung — Bài tập 6 · Mục 2.3}

\subsection*{Kiến thức và trình tự giải}
\textbf{Mục tiêu.} Tận dụng phần cấp hai có dạng bình phương một đạo hàm theo hướng để giảm PDE thành ODE có tham số.

\textbf{Bước 1 — Nhận ra hướng đạo hàm.} Nếu $D=a\partial_x+b\partial_y$ có hệ số hằng, $D^2u=a^2u_{xx}+2ab u_{xy}+b^2u_{yy}$. Với hệ số biến thiên phải dùng


\[
D^2u=a^2u_{xx}+2ab u_{xy}+b^2u_{yy}+(Da)u_x+(Db)u_y.
\]


So sánh cả hạng cấp một trước khi kết luận PDE bằng $D^2u+\lambda Du+\mu u=f$.

\textbf{Bước 2 — Tìm bất biến $s$ và biến tiến $z$.} Giải $Ds=0$, chọn $Dz=k\ne0$. Kiểm tra $J=\partial(s,z)/\partial(x,y)\ne0$. Khi $k$ hằng, $D=k\partial_z$; PDE trở thành ODE theo $z$, trong đó $s$ là tham số. Nếu còn $v_s$, phải giữ lại, chỉ các trường hợp nó triệt tiêu mới giải trực tiếp như ODE.

\textbf{Bước 3 — Giải ODE.} Với $v_{zz}=g(s,z)$, tích phân hai lần cho $v=zF(s)+G(s)+\int_{z_0}^z(z-t)g(s,t)dt$. Với $a v_{zz}+b v_z+c v=g$, giải đa thức $a\lambda^2+b\lambda+c=0$, giữ hai hàm tùy ý theo $s$. Nghiệm kép có thêm hệ số $z$; có thể tìm nghiệm riêng dạng mũ hoặc đa thức.

\textbf{Bước 4 — Dùng dữ liệu đúng biến.} Tính $u_x=s_xv_s+z_xv_z$ và $u_y=s_yv_s+z_yv_z$. Nếu dữ liệu ở $s=0$, chỉ biết $F(0),G(0)$ hoặc đạo hàm của chúng; không suy ra toàn bộ các hàm. So sánh hệ số của $e^z$, $z$, $z^2$ để tìm các hằng số tương thích.

\textbf{Bước 5 — Kiểm tra tính duy nhất và nghiệm cụ thể.} Khi đường dữ liệu gặp tất cả đặc trưng một cách độc lập, giải hệ cho $F,G$. Khi đường dữ liệu là một đặc trưng, mô tả các hàm còn tùy ý và đưa hai lựa chọn rõ ràng. Cuối cùng thế lại phương trình gốc, không chỉ ODE mới.
\subsection*{Lời giải áp dụng cho bài}
\textbf{(a) Phân loại và đổi biến.} $A=1$, $B=-5$, $C=25$ cho $\delta=0$; phần chính là $D^2$, $D=\partial_x-5\partial_y$. Đặc trưng $5x+y=s$. Đặt $s=5x+y$, $z=x$, $J=-1$. Do $Ds=0$, $Dz=1$, ta có $D=\partial_z$, dạng chính tắc $v_{zz}=e^z$.

\textbf{(b) Tích phân hai lần.} Giữ $s$ cố định: $v_z=e^z+F(s)$, rồi $v=e^z+zF(s)+G(s)$. Vậy
\[u=e^x+xF(5x+y)+G(5x+y).\]

\textbf{(c) Xác định các hàm.} Dữ liệu $y=0$ và $u_y=xF'(s)+G'(s)$ cho
\[xF(5x)+G(5x)=-\cos5x,\qquad xF'(5x)+G'(5x)=\sin5x.\]
Lấy đạo hàm dòng đầu theo $x$ được $F(5x)+5xF'(5x)+5G'(5x)=5\sin5x$. Trừ $5$ lần dòng thứ hai, còn $F(5x)=0$ với mọi $x$. Suy ra $F=0$, $G(s)=-\cos s$:
\[\boxed{u=e^x-\cos(5x+y).}\]
Kiểm tra: $D(5x+y)=0$ nên $D^2[-\cos(5x+y)]=0$, $D^2e^x=e^x$. Thế $y=0$ cho đúng $u=e^x-\cos5x$, $u_y=\sin5x$. Toàn bộ $s\in\mathbb R$ đều xuất hiện trên dữ liệu, nên $F,G$ được xác định duy nhất. Đặc trưng là các đường thẳng hệ số góc $-5$.

\textbf{Chi tiết thế ba đạo hàm cấp hai.} Với tọa độ đã chọn,
\[u_{xx}=25v_{ss}+10v_{sz}+v_{zz},\quad u_{xy}=5v_{ss}+v_{sz},\quad u_{yy}=v_{ss}.\]
Nhân lần lượt với $1$,$-10$,$25$, hệ số theo thứ tự $v_{ss},v_{sz},v_{zz}$ là
\[25-50+25=0,\qquad 10-10+0=0,\qquad 1+0+0=1.\]
\begin{center}\includegraphics[width=.97\linewidth,height=11cm,keepaspectratio]{_xu_ly/graphs/BT_2_3_06.png}\end{center}
\textit{Đồ thị minh họa; với nghiệm số, xem phương pháp và tham số trong mã Python kèm theo.}

\clearpage\noindent\texttt{\detokenize{BT_2_3_07}}\par
\section*{Lời giải bổ sung — Bài tập 7 · Mục 2.3}

\subsection*{Kiến thức và trình tự giải}
\textbf{Mục tiêu.} Tận dụng phần cấp hai có dạng bình phương một đạo hàm theo hướng để giảm PDE thành ODE có tham số.

\textbf{Bước 1 — Nhận ra hướng đạo hàm.} Nếu $D=a\partial_x+b\partial_y$ có hệ số hằng, $D^2u=a^2u_{xx}+2ab u_{xy}+b^2u_{yy}$. Với hệ số biến thiên phải dùng


\[
D^2u=a^2u_{xx}+2ab u_{xy}+b^2u_{yy}+(Da)u_x+(Db)u_y.
\]


So sánh cả hạng cấp một trước khi kết luận PDE bằng $D^2u+\lambda Du+\mu u=f$.

\textbf{Bước 2 — Tìm bất biến $s$ và biến tiến $z$.} Giải $Ds=0$, chọn $Dz=k\ne0$. Kiểm tra $J=\partial(s,z)/\partial(x,y)\ne0$. Khi $k$ hằng, $D=k\partial_z$; PDE trở thành ODE theo $z$, trong đó $s$ là tham số. Nếu còn $v_s$, phải giữ lại, chỉ các trường hợp nó triệt tiêu mới giải trực tiếp như ODE.

\textbf{Bước 3 — Giải ODE.} Với $v_{zz}=g(s,z)$, tích phân hai lần cho $v=zF(s)+G(s)+\int_{z_0}^z(z-t)g(s,t)dt$. Với $a v_{zz}+b v_z+c v=g$, giải đa thức $a\lambda^2+b\lambda+c=0$, giữ hai hàm tùy ý theo $s$. Nghiệm kép có thêm hệ số $z$; có thể tìm nghiệm riêng dạng mũ hoặc đa thức.

\textbf{Bước 4 — Dùng dữ liệu đúng biến.} Tính $u_x=s_xv_s+z_xv_z$ và $u_y=s_yv_s+z_yv_z$. Nếu dữ liệu ở $s=0$, chỉ biết $F(0),G(0)$ hoặc đạo hàm của chúng; không suy ra toàn bộ các hàm. So sánh hệ số của $e^z$, $z$, $z^2$ để tìm các hằng số tương thích.

\textbf{Bước 5 — Kiểm tra tính duy nhất và nghiệm cụ thể.} Khi đường dữ liệu gặp tất cả đặc trưng một cách độc lập, giải hệ cho $F,G$. Khi đường dữ liệu là một đặc trưng, mô tả các hàm còn tùy ý và đưa hai lựa chọn rõ ràng. Cuối cùng thế lại phương trình gốc, không chỉ ODE mới.
\subsection*{Lời giải áp dụng cho bài}
\textbf{(a) Phân loại và tọa độ.} $A=1$, $B=-3$, $C=9$, $\delta=0$. Đặt $D=\partial_x-3\partial_y$, $s=3x+y$, $z=x$, $J=-1$. Vì $Ds=0$, $Dz=1$, PDE parabolic trở thành $v_{zz}=e^z$. Đặc trưng $y=-3x+s$.

\textbf{(b) Nghiệm tổng quát.} Tích phân theo $z$ hai lần:
\[u=e^x+xF(3x+y)+G(3x+y).\]

\textbf{(c) Dùng lần lượt hai đường dữ liệu.} Trên $x=0$, $s=y$, nên $1+G(y)=1+\sin y$: $G(s)=\sin s$. Trên $y=0$,
\[e^x+xF(3x)+\sin3x=e^x+\sin3x+x^3\]
cho $F(3x)=x^2$ khi $x\ne0$; tính liên tục xác định cả $F(0)=0$. Vì vậy $F(s)=s^2/9$ và
\[\boxed{u=e^x+\frac{x(3x+y)^2}{9}+\sin(3x+y).}\]
Giữ $s$ cố định thì $v_{zz}=e^z$, nên công thức thỏa PDE. Thế $y=0$ và $x=0$ kiểm tra đúng hai dữ liệu; chúng cùng bằng $1$ tại gốc. Hai đường dữ liệu đã xác định toàn bộ $F,G$, nên nghiệm là duy nhất trong lớp $C^2$ trên $\mathbb R^2$.

\textbf{Chi tiết tích phân theo biến còn lại.} Với $s=3x+y,z=x$, $u_x=3v_s+v_z,u_y=v_s$,
\[u_{xx}=9v_{ss}+6v_{sz}+v_{zz},\quad u_{xy}=3v_{ss}+v_{sz},\quad u_{yy}=v_{ss}.\]
Phần chính $u_{xx}-6u_{xy}+9u_{yy}$ có các hệ số $9-18+9=0$, $6-6=0$, còn $v_{zz}$. Tích phân $v_{zz}=e^z$ lần đầu được $v_z=e^z+F(s)$, lần sau $v=e^z+zF(s)+G(s)$.
\begin{center}\includegraphics[width=.97\linewidth,height=11cm,keepaspectratio]{_xu_ly/graphs/BT_2_3_07.png}\end{center}
\textit{Đồ thị minh họa; với nghiệm số, xem phương pháp và tham số trong mã Python kèm theo.}

\clearpage\noindent\texttt{\detokenize{BT_2_3_08}}\par
\section*{Lời giải bổ sung — Bài tập 8 · Mục 2.3}

\subsection*{Kiến thức và trình tự giải}
\textbf{Mục tiêu.} Tận dụng phần cấp hai có dạng bình phương một đạo hàm theo hướng để giảm PDE thành ODE có tham số.

\textbf{Bước 1 — Nhận ra hướng đạo hàm.} Nếu $D=a\partial_x+b\partial_y$ có hệ số hằng, $D^2u=a^2u_{xx}+2ab u_{xy}+b^2u_{yy}$. Với hệ số biến thiên phải dùng


\[
D^2u=a^2u_{xx}+2ab u_{xy}+b^2u_{yy}+(Da)u_x+(Db)u_y.
\]


So sánh cả hạng cấp một trước khi kết luận PDE bằng $D^2u+\lambda Du+\mu u=f$.

\textbf{Bước 2 — Tìm bất biến $s$ và biến tiến $z$.} Giải $Ds=0$, chọn $Dz=k\ne0$. Kiểm tra $J=\partial(s,z)/\partial(x,y)\ne0$. Khi $k$ hằng, $D=k\partial_z$; PDE trở thành ODE theo $z$, trong đó $s$ là tham số. Nếu còn $v_s$, phải giữ lại, chỉ các trường hợp nó triệt tiêu mới giải trực tiếp như ODE.

\textbf{Bước 3 — Giải ODE.} Với $v_{zz}=g(s,z)$, tích phân hai lần cho $v=zF(s)+G(s)+\int_{z_0}^z(z-t)g(s,t)dt$. Với $a v_{zz}+b v_z+c v=g$, giải đa thức $a\lambda^2+b\lambda+c=0$, giữ hai hàm tùy ý theo $s$. Nghiệm kép có thêm hệ số $z$; có thể tìm nghiệm riêng dạng mũ hoặc đa thức.

\textbf{Bước 4 — Dùng dữ liệu đúng biến.} Tính $u_x=s_xv_s+z_xv_z$ và $u_y=s_yv_s+z_yv_z$. Nếu dữ liệu ở $s=0$, chỉ biết $F(0),G(0)$ hoặc đạo hàm của chúng; không suy ra toàn bộ các hàm. So sánh hệ số của $e^z$, $z$, $z^2$ để tìm các hằng số tương thích.

\textbf{Bước 5 — Kiểm tra tính duy nhất và nghiệm cụ thể.} Khi đường dữ liệu gặp tất cả đặc trưng một cách độc lập, giải hệ cho $F,G$. Khi đường dữ liệu là một đặc trưng, mô tả các hàm còn tùy ý và đưa hai lựa chọn rõ ràng. Cuối cùng thế lại phương trình gốc, không chỉ ODE mới.
\subsection*{Lời giải áp dụng cho bài}
\textbf{(a) Phân loại và chính tắc.} $A=4$, $B=-2$, $C=1$, $\delta=0$. Phần chính $D^2$, $D=2\partial_x-\partial_y$. Bất biến $s=x+2y$; chọn $z=x$, $J=-2$. Khi đó $D=2\partial_z$, PDE trở thành $4v_{zz}=e^z$. Đặc trưng $y=-x/2+s/2$.

\textbf{(b) Tích phân.} Nghiệm riêng $e^z/4$; hai lần tích phân cho
\[u=e^x/4+xF(s)+G(s),\qquad s=x+2y.\]

\textbf{(c) Tương thích trên một đặc trưng.} Trên $(2t,-t)$, $s=0$:
\[e^{2t}/4+2tF(0)+G(0)=C_1e^{2t}.\]
Hàm mũ không thể bằng đa thức bậc nhất trên cả trục trừ khi các hệ số chênh đều bằng $0$; có thể lấy đạo hàm hai lần để thấy ngay $C_1=1/4$. Suy ra $F(0)=G(0)=0$. Đạo hàm theo $y$ là $u_y=2xF'(s)+2G'(s)$, nên
\[4tF'(0)+2G'(0)=C_2t^2+t+1.\]
So sánh ba hệ số cho
\[\boxed{C_1=1/4,\quad C_2=0,\quad F'(0)=1/4,\quad G'(0)=1/2.}\]
Chọn $F(s)=s/4$, $G(s)=s/2$ hoặc $G(s)=s/2+s^2$:
\[u_1=e^x/4+\frac{x(x+2y)}4+\frac{x+2y}2,\qquad u_2=u_1+(x+2y)^2.\]
Tại đường dữ liệu, $u_j=e^{2t}/4$ và $(u_j)_y=t+1$. Cả hai thỏa PDE vì $D(s)=0$, $D^2(xF(s)+G(s))=0$. Có vô số nghiệm: thêm các hàm $F,G$ với giá trị và đạo hàm tại $0$ giữ như trên.

\textbf{Chi tiết tích phân và so hệ số.} $4v_{zz}=e^z$ cho $v_z=e^z/4+F(s)$ rồi $v=e^z/4+zF(s)+G(s)$. Từ $(1/4-C_1)e^{2t}+2tF(0)+G(0)=0$, đạo hàm hai lần cho $4(1/4-C_1)e^{2t}=0$, suy $C_1=1/4$; còn $2tF(0)+G(0)=0$ cho $F(0)=G(0)=0$. Với $4tF'(0)+2G'(0)=C_2t^2+t+1$, so hệ số $t^2,t,1$ lần lượt cho $C_2=0,F'(0)=1/4,G'(0)=1/2$.

\textbf{Chi tiết thế ba đạo hàm cấp hai.} Với tọa độ đã chọn,
\[u_{xx}=v_{ss}+2v_{sz}+v_{zz},\quad u_{xy}=2v_{ss}+2v_{sz},\quad u_{yy}=4v_{ss}.\]
Nhân lần lượt với $4$,$-4$,$1$, hệ số theo thứ tự $v_{ss},v_{sz},v_{zz}$ là
\[4-8+4=0,\qquad 8-8+0=0,\qquad 4+0+0=4.\]
\begin{center}\includegraphics[width=.97\linewidth,height=11cm,keepaspectratio]{_xu_ly/graphs/BT_2_3_08.png}\end{center}
\textit{Đồ thị minh họa; với nghiệm số, xem phương pháp và tham số trong mã Python kèm theo.}

\clearpage\noindent\texttt{\detokenize{BT_2_3_09}}\par
\section*{Lời giải bổ sung — Bài tập 9 · Mục 2.3}

\subsection*{Kiến thức và trình tự giải}
\textbf{Mục tiêu.} Tận dụng phần cấp hai có dạng bình phương một đạo hàm theo hướng để giảm PDE thành ODE có tham số.

\textbf{Bước 1 — Nhận ra hướng đạo hàm.} Nếu $D=a\partial_x+b\partial_y$ có hệ số hằng, $D^2u=a^2u_{xx}+2ab u_{xy}+b^2u_{yy}$. Với hệ số biến thiên phải dùng


\[
D^2u=a^2u_{xx}+2ab u_{xy}+b^2u_{yy}+(Da)u_x+(Db)u_y.
\]


So sánh cả hạng cấp một trước khi kết luận PDE bằng $D^2u+\lambda Du+\mu u=f$.

\textbf{Bước 2 — Tìm bất biến $s$ và biến tiến $z$.} Giải $Ds=0$, chọn $Dz=k\ne0$. Kiểm tra $J=\partial(s,z)/\partial(x,y)\ne0$. Khi $k$ hằng, $D=k\partial_z$; PDE trở thành ODE theo $z$, trong đó $s$ là tham số. Nếu còn $v_s$, phải giữ lại, chỉ các trường hợp nó triệt tiêu mới giải trực tiếp như ODE.

\textbf{Bước 3 — Giải ODE.} Với $v_{zz}=g(s,z)$, tích phân hai lần cho $v=zF(s)+G(s)+\int_{z_0}^z(z-t)g(s,t)dt$. Với $a v_{zz}+b v_z+c v=g$, giải đa thức $a\lambda^2+b\lambda+c=0$, giữ hai hàm tùy ý theo $s$. Nghiệm kép có thêm hệ số $z$; có thể tìm nghiệm riêng dạng mũ hoặc đa thức.

\textbf{Bước 4 — Dùng dữ liệu đúng biến.} Tính $u_x=s_xv_s+z_xv_z$ và $u_y=s_yv_s+z_yv_z$. Nếu dữ liệu ở $s=0$, chỉ biết $F(0),G(0)$ hoặc đạo hàm của chúng; không suy ra toàn bộ các hàm. So sánh hệ số của $e^z$, $z$, $z^2$ để tìm các hằng số tương thích.

\textbf{Bước 5 — Kiểm tra tính duy nhất và nghiệm cụ thể.} Khi đường dữ liệu gặp tất cả đặc trưng một cách độc lập, giải hệ cho $F,G$. Khi đường dữ liệu là một đặc trưng, mô tả các hàm còn tùy ý và đưa hai lựa chọn rõ ràng. Cuối cùng thế lại phương trình gốc, không chỉ ODE mới.
\subsection*{Lời giải áp dụng cho bài}
\textbf{(a) Phân loại và chính tắc.} $A=25$, $B=-5$, $C=1$, $\delta=0$. Đặt $D=5\partial_x-\partial_y$. PDE là $D^2u-2Du=15e^x$. Chọn $s=x+5y$, $z=x$, $J=-5$; $Ds=0$, $Dz=5$, nên
\[25v_{zz}-10v_z=15e^z.\]
Đặc trưng là $y=-x/5+s/5$.

\textbf{(b) Giải ODE.} Đa thức đặc trưng $25\lambda^2-10\lambda=5\lambda(5\lambda-2)$ có nghiệm $0$, $2/5$. Nghiệm riêng $v_p=e^z$ vì $25-10=15$. Vì vậy
\[u=e^x+e^{2x/5}F(x+5y)+G(x+5y).\]

\textbf{(c) Tìm $C$ và hai nghiệm.} Trên $y=-x/5$, $s=0$, phương trình dữ liệu là
\[(1-C)e^x+(F(0)-1)e^{2x/5}+G(0)=0.\]
Ba hàm $e^x$, $e^{2x/5}$, $1$ độc lập tuyến tính, nên
\[\boxed{C=1,\qquad F(0)=1,\qquad G(0)=0.}\]
Lấy $F=1$, $G=0$ và $F=1$, $G(s)=s$ cho
\[u_1=e^x+e^{2x/5},\qquad u_2=e^x+e^{2x/5}+x+5y.\]
Kiểm tra: $D^2e^x-2De^x=15e^x$; $De^{2x/5}=2e^{2x/5}$ nên phần mũ thứ hai thỏa thuần nhất; $D(x+5y)=0$. Hai công thức bằng nhau trên đường dữ liệu nhưng khác ngoài nó. Dữ liệu trên một đặc trưng chỉ quy định hai giá trị $F(0),G(0)$, nên không duy nhất.

\textbf{Chi tiết tính độc lập của ba hàm.} Đặt $a=1-C,b=F(0)-1,c=G(0)$. Từ $ae^x+be^{2x/5}+c=0$, đạo hàm lần thứ nhất và thứ hai:
\[ae^x+\tfrac25be^{2x/5}=0,\quad ae^x+\tfrac4{25}be^{2x/5}=0.\]
Trừ hai dòng cho $(6/25)be^{2x/5}=0$, suy $b=0$; thế ngược cho $a=0$, rồi $c=0$. Do đó $C=1,F(0)=1,G(0)=0$ như kết luận.

\textbf{Chi tiết thế ba đạo hàm cấp hai.} Với tọa độ đã chọn,
\[u_{xx}=v_{ss}+2v_{sz}+v_{zz},\quad u_{xy}=5v_{ss}+5v_{sz},\quad u_{yy}=25v_{ss}.\]
Nhân lần lượt với $25$,$-10$,$1$, hệ số theo thứ tự $v_{ss},v_{sz},v_{zz}$ là
\[25-50+25=0,\qquad 50-50+0=0,\qquad 25+0+0=25.\]
\begin{center}\includegraphics[width=.97\linewidth,height=11cm,keepaspectratio]{_xu_ly/graphs/BT_2_3_09.png}\end{center}
\textit{Đồ thị minh họa; với nghiệm số, xem phương pháp và tham số trong mã Python kèm theo.}

\clearpage\noindent\texttt{\detokenize{BT_2_3_10}}\par
\section*{Lời giải bổ sung — Bài tập 10 · Mục 2.3}

\subsection*{Kiến thức và trình tự giải}
\textbf{Mục tiêu.} Tận dụng phần cấp hai có dạng bình phương một đạo hàm theo hướng để giảm PDE thành ODE có tham số.

\textbf{Bước 1 — Nhận ra hướng đạo hàm.} Nếu $D=a\partial_x+b\partial_y$ có hệ số hằng, $D^2u=a^2u_{xx}+2ab u_{xy}+b^2u_{yy}$. Với hệ số biến thiên phải dùng


\[
D^2u=a^2u_{xx}+2ab u_{xy}+b^2u_{yy}+(Da)u_x+(Db)u_y.
\]


So sánh cả hạng cấp một trước khi kết luận PDE bằng $D^2u+\lambda Du+\mu u=f$.

\textbf{Bước 2 — Tìm bất biến $s$ và biến tiến $z$.} Giải $Ds=0$, chọn $Dz=k\ne0$. Kiểm tra $J=\partial(s,z)/\partial(x,y)\ne0$. Khi $k$ hằng, $D=k\partial_z$; PDE trở thành ODE theo $z$, trong đó $s$ là tham số. Nếu còn $v_s$, phải giữ lại, chỉ các trường hợp nó triệt tiêu mới giải trực tiếp như ODE.

\textbf{Bước 3 — Giải ODE.} Với $v_{zz}=g(s,z)$, tích phân hai lần cho $v=zF(s)+G(s)+\int_{z_0}^z(z-t)g(s,t)dt$. Với $a v_{zz}+b v_z+c v=g$, giải đa thức $a\lambda^2+b\lambda+c=0$, giữ hai hàm tùy ý theo $s$. Nghiệm kép có thêm hệ số $z$; có thể tìm nghiệm riêng dạng mũ hoặc đa thức.

\textbf{Bước 4 — Dùng dữ liệu đúng biến.} Tính $u_x=s_xv_s+z_xv_z$ và $u_y=s_yv_s+z_yv_z$. Nếu dữ liệu ở $s=0$, chỉ biết $F(0),G(0)$ hoặc đạo hàm của chúng; không suy ra toàn bộ các hàm. So sánh hệ số của $e^z$, $z$, $z^2$ để tìm các hằng số tương thích.

\textbf{Bước 5 — Kiểm tra tính duy nhất và nghiệm cụ thể.} Khi đường dữ liệu gặp tất cả đặc trưng một cách độc lập, giải hệ cho $F,G$. Khi đường dữ liệu là một đặc trưng, mô tả các hàm còn tùy ý và đưa hai lựa chọn rõ ràng. Cuối cùng thế lại phương trình gốc, không chỉ ODE mới.
\subsection*{Lời giải áp dụng cho bài}
\textbf{(a) Phân loại và phép đổi biến.} $A=8$, $B=4$, $C=2$ cho $\delta=0$. Đặt $D=2\partial_x+\partial_y$; PDE là $2D^2u-Du=0$. Bất biến $s=x-2y$, chọn $z=x$, $J=2$. Do $D=2\partial_z$, dạng mới $8v_{zz}-2v_z=0$. Đặc trưng $y=x/2-s/2$.

\textbf{(b) Nghiệm tổng quát.} Phương trình đặc trưng $8\lambda^2-2\lambda=0$ cho $\lambda=0,1/4$, nên
\[u=e^{x/4}F(s)+G(s).\]

\textbf{(c) Dùng Cauchy.} Trên $y=0$, $s=x$:
\[e^{x/4}F(x)+G(x)=\cos x,\qquad
-2[e^{x/4}F'(x)+G'(x)]=\sin x.\]
Lấy đạo hàm phương trình thứ nhất: $e^{x/4}F/4+e^{x/4}F'+G'=-\sin x$. Dòng thứ hai cho $e^{x/4}F'+G'=-\sin x/2$, nên $F(x)=-2e^{-x/4}\sin x$, $G(x)=\cos x+2\sin x$. Thế lại:
\[\boxed{u=\cos(x-2y)+2(1-e^{y/2})\sin(x-2y).}\]
Để kiểm tra PDE, giữ $s$ cố định: phần biến thiên theo $z$ chỉ là một hằng cộng một bội của $e^{z/4}$, đúng ODE. Tại $y=0$, $u=\cos x$; tính $u_y$ cho $2\sin x-\sin x=\sin x$. Hai hàm $F,G$ được xác định trên toàn trục nên nghiệm duy nhất trên $\mathbb R^2$.

\textbf{Chi tiết thế ba đạo hàm cấp hai.} Với tọa độ đã chọn,
\[u_{xx}=v_{ss}+2v_{sz}+v_{zz},\quad u_{xy}=-2v_{ss}-2v_{sz},\quad u_{yy}=4v_{ss}.\]
Nhân lần lượt với $8$,$8$,$2$, hệ số theo thứ tự $v_{ss},v_{sz},v_{zz}$ là
\[8-16+8=0,\qquad 16-16+0=0,\qquad 8+0+0=8.\]
\begin{center}\includegraphics[width=.97\linewidth,height=11cm,keepaspectratio]{_xu_ly/graphs/BT_2_3_10.png}\end{center}
\textit{Đồ thị minh họa; với nghiệm số, xem phương pháp và tham số trong mã Python kèm theo.}

\clearpage\noindent\texttt{\detokenize{BT_2_3_11}}\par
\section*{Lời giải bổ sung — Bài tập 11 · Mục 2.3}

\subsection*{Kiến thức và trình tự giải}
\textbf{Mục tiêu.} Giải phương trình hyperbolic sau khi đã tìm hai tọa độ đặc trưng độc lập.

\textbf{Bước 1 — Phân loại, tìm đặc trưng và thay đạo hàm.} Làm đủ các bước của phép đưa về chính tắc. Viết cả phép đổi ngược để biểu diễn vế phải $\widetilde f(\xi,\eta)=f(x(\xi,\eta),y(\xi,\eta))$.

\textbf{Bước 2 — Tích phân phương trình chính tắc.} Với $v_{\xi\eta}=g$, giữ biến kia cố định khi tích phân:


\[
v=F(\xi)+G(\eta)+\int_{\xi_0}^{\xi}\int_{\eta_0}^{\eta}g(a,b)\,db\,da.
\]


Nếu tìm được một nghiệm riêng $P(x,y)$ dễ hơn, viết $u=P+F(\xi)+G(\eta)$. Với $v_{\xi\eta}+\alpha v_\eta=0$, đặt $w=v_\eta$, giải $w_\xi+\alpha w=0$ rồi tích phân theo $\eta$; thu được $v=F(\xi)+e^{-\alpha\xi}G(\eta)$.

\textbf{Bước 3 — Đổi cả dữ liệu.} Thế tọa độ đường dữ liệu vào $\xi,\eta$, dùng quy tắc dây chuyền cho đạo hàm được cho. Trên đường không đặc trưng, hai dữ liệu thường xác định hai hàm $F,G$. Trên đường đặc trưng, có biến không thay đổi, nên dữ liệu chỉ xác định giá trị hoặc đạo hàm của một hàm tại một điểm; phải tìm điều kiện tương thích và nêu phần tùy ý còn lại.

\textbf{Bước 4 — Lập đẳng thức hình bình hành.} Với $v=F(\xi)+G(\eta)$, bốn đỉnh theo hai họ đặc trưng thỏa $v_{11}+v_{22}=v_{12}+v_{21}$. Nếu có số hạng mũ, trước hết nhân trọng số để được tổng hai hàm một biến. Với vế phải $g$, hiệu bốn đỉnh bằng tích phân kép của $g$ trên hình chữ nhật trong tọa độ mới.

\textbf{Bước 5 — Xác định miền được dữ liệu quyết định.} Hai đặc trưng đi qua điểm đang xét phải gặp đoạn dữ liệu phù hợp. Với phương trình sóng $v_{tt}=v_{zz}$ trên $z>0$, dữ liệu lúc $t=0$ xác định duy nhất miền $z>t$; ở $0<z<t$, cần thêm dữ liệu biên $z=0$. Với đặc trưng cong, dùng miền trong các tọa độ mới rồi đổi lại.

\textbf{Bước 6 — Kiểm tra.} Thế nghiệm vào cả PDE và từng điều kiện. Nếu đề yêu cầu hai nghiệm, chỉ rõ hai lựa chọn khác nhau cho hàm tùy ý và chứng minh phần chênh lệch thỏa phương trình thuần nhất, đồng thời triệt tiêu trên đường dữ liệu.

\textbf{Biến đổi riêng của bài.} Dùng tọa độ cực (hoặc cực sau khi co giãn trục). Khi thu được $r v_{r\theta}-v_\theta=0$, đặt $w=v/r$ thì $w_{r\theta}=0$. Vì vậy $v=rF(\theta)+G(r)$; phải chia nghiệm cho $r$ trước khi lập đẳng thức bốn đỉnh.
\subsection*{Lời giải áp dụng cho bài}
\textbf{(a) Phân loại và đặc trưng cong.} $A=2xy$, $B=(2y^2-3x^2)/2$, $C=-3xy$ cho
\[\delta=(2y^2+3x^2)^2/4>0\quad(y>0).\]
Đặc trưng phân tích thành $(xdy-ydx)(2ydy+3xdx)=0$: các tia qua gốc và các nửa ellipse $3x^2+2y^2=\mathrm{const}$. Đặt
\[X=\sqrt3x,\quad Y=\sqrt2y,\quad r=\sqrt{3x^2+2y^2},\quad
X=r\cos\theta,\ Y=r\sin\theta,\quad0<\theta<\pi.\]
Trong biến $X,Y$, toán tử gốc bằng $\sqrt6[XY\partial_X^2+(Y^2-X^2)\partial_X\partial_Y-XY\partial_Y^2]$. Công thức cực ở Ví dụ 9 cho $-\sqrt6(r v_{r\theta}-v_\theta)$. Vì thế dạng chính tắc là $r v_{r\theta}-v_\theta=0$.

\textbf{(b) Giải và hình bình hành.} Đặt $w=v/r$ thì $w_{r\theta}=0$, suy ra
\[u=rF(\theta)+G(r).\]
Đẳng thức bốn đỉnh áp dụng cho $w$: $v(r_1,\theta_1)/r_1+v(r_2,\theta_2)/r_2=v(r_1,\theta_2)/r_1+v(r_2,\theta_1)/r_2$. Trong mặt phẳng gốc, bốn cạnh là hai tia và hai cung ellipse; đó là hình chữ nhật trong tọa độ $(r,\theta)$, không phải hình bình hành thẳng thông thường.

\textbf{(c) Dữ liệu và miền duy nhất.} Trên $r=\sqrt3$, dữ liệu $u=x$ cho
\[\sqrt3F(\theta)+G(\sqrt3)=\cos\theta.\]
Vì vậy $F(\theta)=\cos\theta/\sqrt3-k/\sqrt3$, $k=G(\sqrt3)$. Trên $y=0$, $1\le x\le2$, có $\theta=0$, $r=\sqrt3x$, dữ liệu $u=1$ cho $G(r)=1-rF(0)$ ở $\sqrt3\le r\le2\sqrt3$. Thế $F,G$ và rút gọn, mọi hằng $k$ triệt tiêu:
\[\boxed{u=x+1-\sqrt{x^2+\frac23y^2}.}\]
Công thức này xác định trên $y>0$ và có thể dùng như một nghiệm cụ thể ở toàn miền; nó được dữ liệu quyết định duy nhất trên
\[\boxed{y>0,\qquad3<3x^2+2y^2<12.}\]
Giá trị ở hai ellipse biên cũng được quyết định bằng liên tục. Ngoài khoảng bán kính này, $G$ chưa bị dữ liệu quy định. Kiểm tra: trên $r=\sqrt3$, căn bằng $1$, nên $u=x$; trên $y=0$, $x\in[1,2]$, căn bằng $x$, nên $u=1$. Nghiệm có dạng $rF(\theta)+G(r)$, do đó thỏa PDE.

\textbf{Chi tiết đổi tỉ lệ và khử hằng số.} Vì $\partial_x=\sqrt3\partial_X,\partial_y=\sqrt2\partial_Y$, $xy=XY/\sqrt6$, phần chính biến thành $\sqrt6[XYu_{XX}+(Y^2-X^2)u_{XY}-XYu_{YY}]$. Viết $X=rc,Y=rs$, thế đạo hàm cực: các hệ số $v_{rr},v_{\theta\theta},v_r$ bằng $0$; hệ số $v_{r\theta}$ là $-r[4c^2s^2+(c^2-s^2)^2]=-r$, hệ số $v_\theta$ là $1$. Với $w=v/r$, $v_\theta=rw_\theta$, $v_{r\theta}=w_\theta+rw_{r\theta}$, nên $rv_{r\theta}-v_\theta=r^2w_{r\theta}$.

Trong (c), $F=(\cos\theta-k)/\sqrt3$, $F(0)=(1-k)/\sqrt3$, $G=1-r(1-k)/\sqrt3$. Cộng:
\[rF+G=\frac r{\sqrt3}\cos\theta-\frac{rk}{\sqrt3}+1-\frac r{\sqrt3}+\frac{rk}{\sqrt3}
=x+1-\frac r{\sqrt3}.\]
Vì $r/\sqrt3=\sqrt{x^2+2y^2/3}$, thu nghiệm đã nêu.
\begin{center}\includegraphics[width=.97\linewidth,height=11cm,keepaspectratio]{_xu_ly/graphs/BT_2_3_11.png}\end{center}
\textit{Đồ thị minh họa; với nghiệm số, xem phương pháp và tham số trong mã Python kèm theo.}

\clearpage\noindent\texttt{\detokenize{BT_2_3_12}}\par
\section*{Lời giải bổ sung — Bài tập 12 · Mục 2.3}

\subsection*{Kiến thức và trình tự giải}
\textbf{Mục tiêu.} Giải phương trình hyperbolic sau khi đã tìm hai tọa độ đặc trưng độc lập.

\textbf{Bước 1 — Phân loại, tìm đặc trưng và thay đạo hàm.} Làm đủ các bước của phép đưa về chính tắc. Viết cả phép đổi ngược để biểu diễn vế phải $\widetilde f(\xi,\eta)=f(x(\xi,\eta),y(\xi,\eta))$.

\textbf{Bước 2 — Tích phân phương trình chính tắc.} Với $v_{\xi\eta}=g$, giữ biến kia cố định khi tích phân:


\[
v=F(\xi)+G(\eta)+\int_{\xi_0}^{\xi}\int_{\eta_0}^{\eta}g(a,b)\,db\,da.
\]


Nếu tìm được một nghiệm riêng $P(x,y)$ dễ hơn, viết $u=P+F(\xi)+G(\eta)$. Với $v_{\xi\eta}+\alpha v_\eta=0$, đặt $w=v_\eta$, giải $w_\xi+\alpha w=0$ rồi tích phân theo $\eta$; thu được $v=F(\xi)+e^{-\alpha\xi}G(\eta)$.

\textbf{Bước 3 — Đổi cả dữ liệu.} Thế tọa độ đường dữ liệu vào $\xi,\eta$, dùng quy tắc dây chuyền cho đạo hàm được cho. Trên đường không đặc trưng, hai dữ liệu thường xác định hai hàm $F,G$. Trên đường đặc trưng, có biến không thay đổi, nên dữ liệu chỉ xác định giá trị hoặc đạo hàm của một hàm tại một điểm; phải tìm điều kiện tương thích và nêu phần tùy ý còn lại.

\textbf{Bước 4 — Lập đẳng thức hình bình hành.} Với $v=F(\xi)+G(\eta)$, bốn đỉnh theo hai họ đặc trưng thỏa $v_{11}+v_{22}=v_{12}+v_{21}$. Nếu có số hạng mũ, trước hết nhân trọng số để được tổng hai hàm một biến. Với vế phải $g$, hiệu bốn đỉnh bằng tích phân kép của $g$ trên hình chữ nhật trong tọa độ mới.

\textbf{Bước 5 — Xác định miền được dữ liệu quyết định.} Hai đặc trưng đi qua điểm đang xét phải gặp đoạn dữ liệu phù hợp. Với phương trình sóng $v_{tt}=v_{zz}$ trên $z>0$, dữ liệu lúc $t=0$ xác định duy nhất miền $z>t$; ở $0<z<t$, cần thêm dữ liệu biên $z=0$. Với đặc trưng cong, dùng miền trong các tọa độ mới rồi đổi lại.

\textbf{Bước 6 — Kiểm tra.} Thế nghiệm vào cả PDE và từng điều kiện. Nếu đề yêu cầu hai nghiệm, chỉ rõ hai lựa chọn khác nhau cho hàm tùy ý và chứng minh phần chênh lệch thỏa phương trình thuần nhất, đồng thời triệt tiêu trên đường dữ liệu.
\subsection*{Lời giải áp dụng cho bài}
\textbf{(a) Phân loại và đổi biến.} Trong các biến $(x,t)$, $A=x$, $B=0$, $C=-x^5$, nên $\delta=x^6>0$ ở $x>0$. Chọn $z=x^3/3$, $u(x,t)=v(z,t)$, $z_x=x^2>0$. Tính
\[u_x=x^2v_z,\qquad u_{xx}=2xv_z+x^4v_{zz}.\]
Hai hạng cấp một triệt tiêu: $xu_{xx}-2u_x=x^5v_{zz}$. PDE trở thành $v_{tt}=v_{zz}$; dạng đặc trưng $(\xi,\eta)=(z+t,z-t)$ cho $v_{\xi\eta}=0$.

\textbf{(b) Tổng quát.} Tích phân đạo hàm hỗn hợp bằng $0$ cho $u=F(x^3/3+t)+G(x^3/3-t)$. Hai họ đặc trưng là $x^3/3\pm t=\mathrm{const}$.

\textbf{(c) Dữ liệu ban đầu.} Đặt $\Phi(z)=\varphi((3z)^{1/3})$, $\Psi(z)=\psi((3z)^{1/3})$, $z>0$. Tại $t=0$, $F+G=\Phi$, $F'-G'=\Psi$. Lấy đạo hàm phương trình đầu rồi giải hệ: $F'=(\Phi'+\Psi)/2$, $G'=(\Phi'-\Psi)/2$. Do đó công thức d'Alembert là
\[\boxed{u=\frac{\Phi(z+t)+\Phi(z-t)}2+\frac12\int_{z-t}^{z+t}\Psi(s)ds,\quad z=x^3/3>t.}\]
Miền duy nhất do dữ liệu ban đầu là $x>0$, $0<t<x^3/3$. Trong $0<z<t$, một đặc trưng gặp biên $z=0$ trước khi gặp trục $t=0$, nên cần thêm dữ liệu biên. Công thức kiểm tra $u(x,0)=\varphi(x)$, $u_t(x,0)=\psi(x)$ và $v_{tt}=v_{zz}$.

\textbf{(iv) Chọn điều khiển để $u(x,9)=0$ với $0\le x\le3$.} Đoạn đích tương ứng $0\le z\le9$. Giả sử $\Phi,\Psi$ có các giới hạn và độ trơn cần thiết tại $z=0$. Với $0\le z\le t$, dữ liệu biên $v(0,\tau)=h(\tau)$ cho
\[v(z,t)=h(t-z)+\frac{\Phi(t+z)-\Phi(t-z)}2+\frac12\int_{t-z}^{t+z}\Psi(s)ds.\]
Cách suy ra: $F(t+z)+G(z-t)$ và $h(t-z)=F(t-z)+G(z-t)$, nên hiệu bằng $F(t+z)-F(t-z)$; dùng $F'=(\Phi'+\Psi)/2$. Tại $t=9$, đặt $\tau=9-z\in[0,9]$, yêu cầu bằng $0$ cho
\[\boxed{h(\tau)=-\frac{\Phi(18-\tau)-\Phi(\tau)}2-\frac12\int_\tau^{18-\tau}\Psi(s)ds,\quad0\le\tau\le9.}\]
Đây là điều khiển duy nhất trên khoảng thời gian cần dùng; sau $9$ có thể kéo dài đủ trơn. Tại góc $(z,t)=(0,0)$, điều kiện giá trị $h(0)=\Phi(0)$ tương đương
\[\Phi(18)+\Phi(0)+\int_0^{18}\Psi(s)ds=0.\]
Muốn nghiệm $C^2$ qua đặc trưng $z=t$, còn cần $h'(0)=\Psi(0)$ và $h''(0)=\Phi''(0)$. Thay công thức $h$ cho hai điều kiện cụ thể:
\[\Phi'(18)+\Phi'(0)+\Psi(18)-\Psi(0)=0,\qquad
\Phi''(18)+\Phi''(0)+\Psi'(18)-\Psi'(0)=0.\]
Nếu dữ liệu không thỏa điều kiện góc tương ứng, không có nghiệm cổ điển với điều khiển này và đích đã cho trong lớp trơn ấy; không bỏ qua điểm $z=t$ khi kết luận. Với các điều kiện đúng, thế $\tau=9-z$ vào công thức vùng biên cho $v(z,9)=0$ trên toàn đoạn đích.

\textbf{Chi tiết tích phân hệ dữ liệu đầu.} Chọn mốc $z_0>0$ và $F(z_0)=c$. Từ $F'=(\Phi'+\Psi)/2$,
\[F(s)=c+\frac{\Phi(s)-\Phi(z_0)}2+\frac12\int_{z_0}^{s}\Psi(r)dr,\quad G(s)=\Phi(s)-F(s).\]
Thế $F(z+t)+G(z-t)=F(z+t)-F(z-t)+\Phi(z-t)$; hai hằng số mất, các tích phân gộp thành $\int_{z-t}^{z+t}\Psi$, cho đúng d'Alembert với $z-t>0$.

\clearpage\noindent\texttt{\detokenize{BT_2_3_13}}\par
\section*{Lời giải bổ sung — Bài tập 13 · Mục 2.3}

\subsection*{Kiến thức và trình tự giải}
\textbf{Mục tiêu.} Giải phương trình hyperbolic sau khi đã tìm hai tọa độ đặc trưng độc lập.

\textbf{Bước 1 — Phân loại, tìm đặc trưng và thay đạo hàm.} Làm đủ các bước của phép đưa về chính tắc. Viết cả phép đổi ngược để biểu diễn vế phải $\widetilde f(\xi,\eta)=f(x(\xi,\eta),y(\xi,\eta))$.

\textbf{Bước 2 — Tích phân phương trình chính tắc.} Với $v_{\xi\eta}=g$, giữ biến kia cố định khi tích phân:


\[
v=F(\xi)+G(\eta)+\int_{\xi_0}^{\xi}\int_{\eta_0}^{\eta}g(a,b)\,db\,da.
\]


Nếu tìm được một nghiệm riêng $P(x,y)$ dễ hơn, viết $u=P+F(\xi)+G(\eta)$. Với $v_{\xi\eta}+\alpha v_\eta=0$, đặt $w=v_\eta$, giải $w_\xi+\alpha w=0$ rồi tích phân theo $\eta$; thu được $v=F(\xi)+e^{-\alpha\xi}G(\eta)$.

\textbf{Bước 3 — Đổi cả dữ liệu.} Thế tọa độ đường dữ liệu vào $\xi,\eta$, dùng quy tắc dây chuyền cho đạo hàm được cho. Trên đường không đặc trưng, hai dữ liệu thường xác định hai hàm $F,G$. Trên đường đặc trưng, có biến không thay đổi, nên dữ liệu chỉ xác định giá trị hoặc đạo hàm của một hàm tại một điểm; phải tìm điều kiện tương thích và nêu phần tùy ý còn lại.

\textbf{Bước 4 — Lập đẳng thức hình bình hành.} Với $v=F(\xi)+G(\eta)$, bốn đỉnh theo hai họ đặc trưng thỏa $v_{11}+v_{22}=v_{12}+v_{21}$. Nếu có số hạng mũ, trước hết nhân trọng số để được tổng hai hàm một biến. Với vế phải $g$, hiệu bốn đỉnh bằng tích phân kép của $g$ trên hình chữ nhật trong tọa độ mới.

\textbf{Bước 5 — Xác định miền được dữ liệu quyết định.} Hai đặc trưng đi qua điểm đang xét phải gặp đoạn dữ liệu phù hợp. Với phương trình sóng $v_{tt}=v_{zz}$ trên $z>0$, dữ liệu lúc $t=0$ xác định duy nhất miền $z>t$; ở $0<z<t$, cần thêm dữ liệu biên $z=0$. Với đặc trưng cong, dùng miền trong các tọa độ mới rồi đổi lại.

\textbf{Bước 6 — Kiểm tra.} Thế nghiệm vào cả PDE và từng điều kiện. Nếu đề yêu cầu hai nghiệm, chỉ rõ hai lựa chọn khác nhau cho hàm tùy ý và chứng minh phần chênh lệch thỏa phương trình thuần nhất, đồng thời triệt tiêu trên đường dữ liệu.
\subsection*{Lời giải áp dụng cho bài}
\textbf{(a) Phân loại và đổi biến.} $A=(1+x^2)^2$, $B=0$, $C=-1$, $\delta=(1+x^2)^2>0$. Đặt $z=\arctan x\in(-\pi/2,\pi/2)$; $z_x=1/(1+x^2)$, $z_{xx}=-2x/(1+x^2)^2$. Thế vào
\[u_x=\frac{v_z}{1+x^2},\quad u_{xx}=\frac{v_{zz}-2xv_z}{(1+x^2)^2}\]
cho $v_{zz}-v_{yy}=0$: hai hạng $v_z$ triệt tiêu. Chọn $\xi=z+y$, $\eta=z-y$ cho $4v_{\xi\eta}=0$.

\textbf{(b) Nghiệm tổng quát.} $u=F(\arctan x+y)+G(\arctan x-y)$; hai họ đặc trưng $\arctan x\pm y=\mathrm{const}$.

\textbf{(c) Dữ liệu trên đặc trưng $y=z$.} Đặt $\Phi(z)=\varphi(\tan z)$, $\Psi(z)=\psi(\tan z)$. Tại $\eta=0$, $\xi=2z$,
\[F(2z)+G(0)=\Phi(z),\qquad F'(2z)-G'(0)=\Psi(z).\]
Lấy đạo hàm dòng đầu theo $z$ cho $2F'(2z)=\Phi'(z)$. Vậy điều kiện cần và đủ là
\[\boxed{\frac12\Phi'(z)-\Psi(z)=\lambda\quad\text{hằng số},\quad |z|<\pi/2.}\]
Trong biến gốc: $(1+x^2)\varphi'(x)/2-\psi(x)=\lambda$. Khi điều kiện đúng, $F(s)=\Phi(s/2)-G(0)$ ở $|s|<\pi$, còn $G$ tùy ý với $G'(0)=\lambda$. Lấy $G_1(s)=\lambda s$, $G_2(s)=\lambda s+s^2$:
\[u_1=\varphi\!\left(\tan\frac{\arctan x+y}{2}\right)+\lambda(\arctan x-y),\qquad
u_2=u_1+(\arctan x-y)^2.\]
Hai công thức xác định trên miền $|\arctan x+y|<\pi$. Khi $y=\arctan x$, $u_j=\varphi(x)$ và $(u_j)_y=(1+x^2)\varphi'(x)/2-\lambda=\psi(x)$. Trong tọa độ đặc trưng, mỗi nghiệm là tổng hai hàm một biến, nên thỏa PDE. Nếu yêu cầu nghiệm trên toàn $\mathbb R^2$, phải chọn kéo dài $C^2$ của $F$ ra ngoài $(-\pi,\pi)$; dữ liệu chỉ ở $x$ hữu hạn không bảo đảm kéo dài trơn tại hai đầu nếu $\varphi$ tăng không kiểm soát. Vì vậy công thức trên khẳng định tồn tại trên miền tự nhiên chứa toàn đường dữ liệu.

\clearpage\noindent\texttt{\detokenize{BT_2_3_14}}\par
\section*{Lời giải bổ sung — Bài tập 14 · Mục 2.3}

\subsection*{Kiến thức và trình tự giải}
\textbf{Mục tiêu.} Tận dụng phần cấp hai có dạng bình phương một đạo hàm theo hướng để giảm PDE thành ODE có tham số.

\textbf{Bước 1 — Nhận ra hướng đạo hàm.} Nếu $D=a\partial_x+b\partial_y$ có hệ số hằng, $D^2u=a^2u_{xx}+2ab u_{xy}+b^2u_{yy}$. Với hệ số biến thiên phải dùng


\[
D^2u=a^2u_{xx}+2ab u_{xy}+b^2u_{yy}+(Da)u_x+(Db)u_y.
\]


So sánh cả hạng cấp một trước khi kết luận PDE bằng $D^2u+\lambda Du+\mu u=f$.

\textbf{Bước 2 — Tìm bất biến $s$ và biến tiến $z$.} Giải $Ds=0$, chọn $Dz=k\ne0$. Kiểm tra $J=\partial(s,z)/\partial(x,y)\ne0$. Khi $k$ hằng, $D=k\partial_z$; PDE trở thành ODE theo $z$, trong đó $s$ là tham số. Nếu còn $v_s$, phải giữ lại, chỉ các trường hợp nó triệt tiêu mới giải trực tiếp như ODE.

\textbf{Bước 3 — Giải ODE.} Với $v_{zz}=g(s,z)$, tích phân hai lần cho $v=zF(s)+G(s)+\int_{z_0}^z(z-t)g(s,t)dt$. Với $a v_{zz}+b v_z+c v=g$, giải đa thức $a\lambda^2+b\lambda+c=0$, giữ hai hàm tùy ý theo $s$. Nghiệm kép có thêm hệ số $z$; có thể tìm nghiệm riêng dạng mũ hoặc đa thức.

\textbf{Bước 4 — Dùng dữ liệu đúng biến.} Tính $u_x=s_xv_s+z_xv_z$ và $u_y=s_yv_s+z_yv_z$. Nếu dữ liệu ở $s=0$, chỉ biết $F(0),G(0)$ hoặc đạo hàm của chúng; không suy ra toàn bộ các hàm. So sánh hệ số của $e^z$, $z$, $z^2$ để tìm các hằng số tương thích.

\textbf{Bước 5 — Kiểm tra tính duy nhất và nghiệm cụ thể.} Khi đường dữ liệu gặp tất cả đặc trưng một cách độc lập, giải hệ cho $F,G$. Khi đường dữ liệu là một đặc trưng, mô tả các hàm còn tùy ý và đưa hai lựa chọn rõ ràng. Cuối cùng thế lại phương trình gốc, không chỉ ODE mới.
\subsection*{Lời giải áp dụng cho bài}
\textbf{(a) Nhận dạng toán tử.} $A=e^{2x}$, $B=e^x$, $C=1$ cho $\delta=0$. Đặt $D=e^x\partial_x+\partial_y$. Do $De^x=e^{2x}$,
\[D^2u=e^{2x}u_{xx}+2e^xu_{xy}+u_{yy}+e^{2x}u_x,\]
đúng toàn bộ vế trái. Đặc trưng $dy/dx=e^{-x}$ có bất biến $s=y+e^{-x}$. Chọn $z=y$, $J=-e^{-x}\ne0$, $Ds=0$, $Dz=1$. Phép đổi ngược $x=-\log(s-z)$, miền $s-z>0$. Dạng chính tắc
\[v_{zz}=(s-z)^{-2}.\]

\textbf{(b) Tích phân.} Vì $[-\log(s-z)]_z=(s-z)^{-1}$ và đạo hàm tiếp bằng $(s-z)^{-2}$,
\[v=-\log(s-z)+zF(s)+G(s),\quad
u=x+yF(y+e^{-x})+G(y+e^{-x}).\]

\textbf{(c) Dữ liệu tại $x=0$.} Khi đó $s=y+1$. Dữ liệu giá trị cho $yF(y+1)+G(y+1)=2y+1$. Vì $s_x=-e^{-x}$,
\[u_x=1-e^{-x}[yF'(s)+G'(s)],\]
nên dữ liệu đạo hàm cho $yF'(y+1)+G'(y+1)=1$. Lấy đạo hàm dữ liệu giá trị theo $y$ rồi trừ dòng này, suy ra $F(y+1)=1$. Do đó $F=1$ và $G(s)=s$:
\[\boxed{u=x+2y+e^{-x}.}\]
Kiểm tra $u_x=1-e^{-x}$, $u_{xx}=e^{-x}$, $u_{xy}=u_{yy}=0$, nên vế trái bằng $e^{2x}$. Tại $x=0$, $u=2y+1$, $u_x=0$. Dữ liệu xác định $F,G$ trên toàn trục, nên nghiệm duy nhất trên $\mathbb R^2$.
\begin{center}\includegraphics[width=.97\linewidth,height=11cm,keepaspectratio]{_xu_ly/graphs/BT_2_3_14.png}\end{center}
\textit{Đồ thị minh họa; với nghiệm số, xem phương pháp và tham số trong mã Python kèm theo.}

\clearpage\noindent\texttt{\detokenize{BT_2_3_15}}\par
\section*{Lời giải bổ sung — Bài tập 15 · Mục 2.3}

\subsection*{Kiến thức và trình tự giải}
\textbf{Mục tiêu.} Tận dụng phần cấp hai có dạng bình phương một đạo hàm theo hướng để giảm PDE thành ODE có tham số.

\textbf{Bước 1 — Nhận ra hướng đạo hàm.} Nếu $D=a\partial_x+b\partial_y$ có hệ số hằng, $D^2u=a^2u_{xx}+2ab u_{xy}+b^2u_{yy}$. Với hệ số biến thiên phải dùng


\[
D^2u=a^2u_{xx}+2ab u_{xy}+b^2u_{yy}+(Da)u_x+(Db)u_y.
\]


So sánh cả hạng cấp một trước khi kết luận PDE bằng $D^2u+\lambda Du+\mu u=f$.

\textbf{Bước 2 — Tìm bất biến $s$ và biến tiến $z$.} Giải $Ds=0$, chọn $Dz=k\ne0$. Kiểm tra $J=\partial(s,z)/\partial(x,y)\ne0$. Khi $k$ hằng, $D=k\partial_z$; PDE trở thành ODE theo $z$, trong đó $s$ là tham số. Nếu còn $v_s$, phải giữ lại, chỉ các trường hợp nó triệt tiêu mới giải trực tiếp như ODE.

\textbf{Bước 3 — Giải ODE.} Với $v_{zz}=g(s,z)$, tích phân hai lần cho $v=zF(s)+G(s)+\int_{z_0}^z(z-t)g(s,t)dt$. Với $a v_{zz}+b v_z+c v=g$, giải đa thức $a\lambda^2+b\lambda+c=0$, giữ hai hàm tùy ý theo $s$. Nghiệm kép có thêm hệ số $z$; có thể tìm nghiệm riêng dạng mũ hoặc đa thức.

\textbf{Bước 4 — Dùng dữ liệu đúng biến.} Tính $u_x=s_xv_s+z_xv_z$ và $u_y=s_yv_s+z_yv_z$. Nếu dữ liệu ở $s=0$, chỉ biết $F(0),G(0)$ hoặc đạo hàm của chúng; không suy ra toàn bộ các hàm. So sánh hệ số của $e^z$, $z$, $z^2$ để tìm các hằng số tương thích.

\textbf{Bước 5 — Kiểm tra tính duy nhất và nghiệm cụ thể.} Khi đường dữ liệu gặp tất cả đặc trưng một cách độc lập, giải hệ cho $F,G$. Khi đường dữ liệu là một đặc trưng, mô tả các hàm còn tùy ý và đưa hai lựa chọn rõ ràng. Cuối cùng thế lại phương trình gốc, không chỉ ODE mới.
\subsection*{Lời giải áp dụng cho bài}
\textbf{(a) Phân loại và hướng đạo hàm.} Trên $0<x<\pi/2$, $y>0$, $A=\tan^2x$, $B=-y\tan x$, $C=y^2$, $\delta=0$. Đặt $D=\tan x\partial_x-y\partial_y$. Ta có
\[D^2u=\tan^2x\,u_{xx}-2y\tan x\,u_{xy}+y^2u_{yy}
+\tan x\sec^2x\,u_x+y u_y.\]
Trừ $Du$ cho hệ số $u_x$ là $\tan x(\sec^2x-1)=\tan^3x$, hệ số $u_y$ là $2y$, đúng đề. PDE là $(D^2-D)u=\tan^3x$.

\textbf{Tọa độ đặc trưng.} Từ $dy/dx=-y/\tan x$, tích phân $dy/y=-\cot x\,dx$ cho $s=y\sin x$ không đổi. Chọn $z=\log\sin x$: $Dz=1$, $Ds=0$, $J=-\cos x\ne0$. Miền mới $s>0$, $z<0$. Dạng chính tắc $v_{zz}-v_z=\tan^3x(z)$.

\textbf{(b) Giải ODE.} Nghiệm riêng $u_p=x$ vì $Dx=\tan x$, $D^2x-Dx=\tan x\sec^2x-\tan x=\tan^3x$. Phần thuần nhất có nghiệm $1,e^z$, nên
\[\boxed{u=x+F(y\sin x)+\sin x\,G(y\sin x).}\]

\textbf{(c) Tìm $a,b$ và hai nghiệm.} Giả sử nghiệm có vết liên tục đến $y=0$ với các hàm $F,G$ đủ trơn tại $s=0$. Trên biên, $x+F(0)+\sin xG(0)=ax+b\sin x$. Ba hàm $1,x,\sin x$ độc lập trên $(0,\pi/2)$, suy ra
\[\boxed{a=1,\quad F(0)=0,\quad G(0)=b,\quad b\in\mathbb R\text{ tùy ý}.}\]
Chọn $F=0$, $G=b$ và $F(s)=s$, $G=b$ cho
\[u_1=x+b\sin x,\qquad u_2=x+b\sin x+y\sin x.\]
Kiểm tra $D(\sin x)=\sin x$, nên $(D^2-D)(b\sin x)=0$; $D(y\sin x)=0$, nên phần chênh cũng thỏa thuần nhất. Cả hai có vết $x+b\sin x$ trên $y=0$. Có vô số nghiệm; dữ liệu biên là một đặc trưng, chỉ xác định giá trị hai hàm tùy ý tại $0$.

\textbf{Chi tiết chứng minh độc lập và dữ liệu biên.} Từ $(1-a)x+F(0)+(G(0)-b)\sin x=0$, đạo hàm hai lần theo $x$ cho $-(G(0)-b)\sin x=0$. Trên khoảng đang xét $\sin x>0$, nên $G(0)=b$; phương trình còn $(1-a)x+F(0)=0$, đạo hàm cho $a=1$ rồi $F(0)=0$. Không có điều kiện nào còn lại để cố định $b$.
\begin{center}\includegraphics[width=.97\linewidth,height=11cm,keepaspectratio]{_xu_ly/graphs/BT_2_3_15.png}\end{center}
\textit{Đồ thị minh họa; với nghiệm số, xem phương pháp và tham số trong mã Python kèm theo.}

\clearpage\noindent\texttt{\detokenize{BT_2_3_16}}\par
\section*{Lời giải bổ sung — Bài tập 16 · Mục 2.3}

\subsection*{Kiến thức và trình tự giải}
\textbf{Mục tiêu.} Đưa phương trình elliptic về dạng đơn giản, rồi tìm nghiệm giải tích từ dữ liệu Cauchy mà bài yêu cầu.

\textbf{Bước 1 — Loại hạng hỗn hợp và hạng cấp một do đổi biến.} Tính $\delta<0$, chọn $\xi,\eta$ thực có Jacobian khác $0$. Thay đầy đủ đạo hàm cấp hai; các đạo hàm của tọa độ có thể triệt tiêu với hạng cấp một trong đề. Kết quả thường là $v_{\xi\xi}+v_{\eta\eta}+c v=0$.

\textbf{Bước 2 — Chuyển dữ liệu sang đường tọa độ mới.} Nếu đường dữ liệu trở thành $\xi=0$, đặt $v(0,\eta)=f(\eta)$ và tính $v_\xi(0,\eta)=g(\eta)$ từ đạo hàm được cho. Nếu đường là $\eta=\eta_0$, đổi vai trò hai biến. Không lấy nguyên dữ liệu đạo hàm cũ làm dữ liệu mới khi hệ số dây chuyền khác $1$.

\textbf{Bước 3 — Khai triển và lập truy hồi.} Viết


\[
v(\xi,\eta)=\sum_{n=0}^{\infty}\frac{\xi^n}{n!}a_n(\eta),\qquad
a_0=f,\quad a_1=g,\quad a_{n+2}=-a_n''-c a_n.
\]


Tính từng $a_2,a_3,\ldots$. Nếu các hệ số sau một bậc đều bằng $0$, chuỗi kết thúc và cho một đa thức chính xác. Nếu không, cần xét hội tụ trong miền dữ liệu giải tích; chuỗi hình thức chưa tự là nghiệm.

\textbf{Bước 4 — Đổi ngược và kiểm tra.} Thế $\xi(x,y),\eta(x,y)$ vào nghiệm. Tính đạo hàm hoặc dùng toán tử đã được chứng minh tương đương để kiểm tra PDE, đồng thời kiểm tra từng dữ liệu gốc. Kết luận duy nhất trong lớp nghiệm giải tích trên miền liên thông mà phép đổi biến hợp lệ; dữ liệu Cauchy elliptic nói chung có thể không ổn định, nên không suy ra tồn tại cho dữ liệu tùy ý chỉ từ cách lập chuỗi.
\subsection*{Lời giải áp dụng cho bài}
\textbf{(a) Phân loại và đổi biến.} $A=1$, $B=-\sin x$, $C=2-\cos^2x=1+\sin^2x$, nên $\delta=-1<0$: elliptic trên toàn mặt phẳng. Đặt $\xi=x$, $\eta=y-\cos x$, $J=1$. Tính
\begin{align*}
u_x&=v_\xi+\sin x\,v_\eta,\quad u_y=v_\eta,\\
u_{xx}&=v_{\xi\xi}+2\sin x\,v_{\xi\eta}+\sin^2x\,v_{\eta\eta}+\cos x\,v_\eta,\\
u_{xy}&=v_{\xi\eta}+\sin x\,v_{\eta\eta},\quad u_{yy}=v_{\eta\eta}.
\end{align*}
Thế vào PDE, hạng hỗn hợp và hạng $v_\eta$ triệt tiêu; hệ số $v_{\eta\eta}$ bằng $\sin^2x-2\sin^2x+1+\sin^2x=1$. Dạng chính tắc $v_{\xi\xi}+v_{\eta\eta}=0$.

\textbf{(b) Nghiệm giải tích từ dữ liệu.} Tại $x=0$, $\xi=0$, $\eta=y-1$, $u_x=v_\xi$ vì $\sin0=0$. Vì vậy
\[v(0,\eta)=0,\qquad v_\xi(0,\eta)=\eta.\]
Viết $v=\sum_{n\ge0}a_n(\eta)\xi^n/n!$. PDE cho $a_{n+2}=-a_n''$, dữ liệu cho $a_0=0$, $a_1=\eta$. Do $a_0''=a_1''=0$, mọi $a_n$ với $n\ge2$ bằng $0$. Chuỗi kết thúc, nên $v=\xi\eta$, là nghiệm giải tích duy nhất.

\textbf{(c) Đổi ngược và kiểm tra trực tiếp.}
\[\boxed{u=x(y-\cos x).}\]
Tính $u_x=y-\cos x+x\sin x$, $u_{xx}=2\sin x+x\cos x$, $u_{xy}=1$, $u_{yy}=0$, $u_y=x$. Vế trái PDE bằng $(2\sin x+x\cos x)-2\sin x-x\cos x=0$. Trên $x=0$, $u=0$, $u_x=y-1$. Công thức giải tích trên $\mathbb R^2$ và duy nhất trong lớp giải tích trên miền liên thông chứa đường dữ liệu.
\begin{center}\includegraphics[width=.97\linewidth,height=11cm,keepaspectratio]{_xu_ly/graphs/BT_2_3_16.png}\end{center}
\textit{Đồ thị minh họa; với nghiệm số, xem phương pháp và tham số trong mã Python kèm theo.}

\clearpage\noindent\texttt{\detokenize{BT_2_3_17}}\par
\section*{Lời giải bổ sung — Bài tập 17 · Mục 2.3}

\subsection*{Kiến thức và trình tự giải}
\textbf{Mục tiêu.} Đưa phương trình elliptic về dạng đơn giản, rồi tìm nghiệm giải tích từ dữ liệu Cauchy mà bài yêu cầu.

\textbf{Bước 1 — Loại hạng hỗn hợp và hạng cấp một do đổi biến.} Tính $\delta<0$, chọn $\xi,\eta$ thực có Jacobian khác $0$. Thay đầy đủ đạo hàm cấp hai; các đạo hàm của tọa độ có thể triệt tiêu với hạng cấp một trong đề. Kết quả thường là $v_{\xi\xi}+v_{\eta\eta}+c v=0$.

\textbf{Bước 2 — Chuyển dữ liệu sang đường tọa độ mới.} Nếu đường dữ liệu trở thành $\xi=0$, đặt $v(0,\eta)=f(\eta)$ và tính $v_\xi(0,\eta)=g(\eta)$ từ đạo hàm được cho. Nếu đường là $\eta=\eta_0$, đổi vai trò hai biến. Không lấy nguyên dữ liệu đạo hàm cũ làm dữ liệu mới khi hệ số dây chuyền khác $1$.

\textbf{Bước 3 — Khai triển và lập truy hồi.} Viết


\[
v(\xi,\eta)=\sum_{n=0}^{\infty}\frac{\xi^n}{n!}a_n(\eta),\qquad
a_0=f,\quad a_1=g,\quad a_{n+2}=-a_n''-c a_n.
\]


Tính từng $a_2,a_3,\ldots$. Nếu các hệ số sau một bậc đều bằng $0$, chuỗi kết thúc và cho một đa thức chính xác. Nếu không, cần xét hội tụ trong miền dữ liệu giải tích; chuỗi hình thức chưa tự là nghiệm.

\textbf{Bước 4 — Đổi ngược và kiểm tra.} Thế $\xi(x,y),\eta(x,y)$ vào nghiệm. Tính đạo hàm hoặc dùng toán tử đã được chứng minh tương đương để kiểm tra PDE, đồng thời kiểm tra từng dữ liệu gốc. Kết luận duy nhất trong lớp nghiệm giải tích trên miền liên thông mà phép đổi biến hợp lệ; dữ liệu Cauchy elliptic nói chung có thể không ổn định, nên không suy ra tồn tại cho dữ liệu tùy ý chỉ từ cách lập chuỗi.
\subsection*{Lời giải áp dụng cho bài}
\textbf{(a) Phân loại và xây dựng biến mới.} $A=y^2$, $B=-y^2$, $C=y^2+1$, $\delta=-y^2<0$ khi $y>0$. Chọn $\xi=x+h(y)$, $\eta=k(y)$. Hệ số hỗn hợp bằng $2k'[-y^2+(y^2+1)h']$, nên lấy $h'=y^2/(1+y^2)$. Khi đó hệ số $v_{\xi\xi}$ bằng $y^2/(1+y^2)$; để cân bằng, lấy $k'=y/(1+y^2)$. Tích phân:
\[\xi=x+y-\arctan y,\qquad\eta=\tfrac12\log(1+y^2),\qquad J=\frac y{1+y^2}>0.\]
Có $h''=2y/(1+y^2)^2$, $k''=(1-y^2)/(1+y^2)^2$. Hệ số mới của $v_\xi$ bằng
\[(1+y^2)h''-y-\frac{1-y^2}{y}h'=0;\]
hệ số của $v_\eta$ bằng $(1+y^2)k''-(1-y^2)k'/y=0$. Vì vậy PDE là
\[\frac{y^2}{1+y^2}(v_{\xi\xi}+v_{\eta\eta})=0,\quad\text{tức }v_{\xi\xi}+v_{\eta\eta}=0.\]

\textbf{(b) Đổi dữ liệu rồi tính chuỗi.} Đặt $L=\log2$, $\eta_0=L/2$. Tại $y=1$, $\xi=x+1-\pi/4$; dữ liệu cho $v(\xi,\eta_0)=L\xi$. Đạo hàm theo $\xi$ trên đường này bằng $L$. Mặt khác $u_y=h'v_\xi+k'v_\eta=(v_\xi+v_\eta)/2$ tại $y=1$. Dữ liệu $u_y=\xi-L/2$ cho
\[v_\eta(\xi,\eta_0)=2\xi-2L.\]
Viết $v=\sum b_n(\xi)(\eta-\eta_0)^n/n!$. Có $b_0=L\xi$, $b_1=2\xi-2L$, $b_{n+2}=-b_n''$. Hai hàm đầu tuyến tính theo $\xi$, nên $b_n=0$ với $n\ge2$. Do đó
\[v=L\xi+(\eta-L/2)(2\xi-2L)=2\xi\eta-2L\eta+L^2.\]

\textbf{(c) Đổi ngược và kiểm tra.}
\[\boxed{u=(x+y-\arctan y-\log2)\log(1+y^2)+(\log2)^2.}\]
Trong biến mới $v_{\xi\xi}=v_{\eta\eta}=0$, nên PDE gốc bằng $0$ theo phép biến đổi đã tính. Tại $y=1$, giá trị là $L\xi$; đạo hàm là $(L+2\xi-2L)/2=\xi-L/2$, đúng đề. Có thể kiểm tra trực tiếp bằng $u_x=\log(1+y^2)$, $u_{xx}=0$, $u_{xy}=2y/(1+y^2)$ và
\[u_y=\frac{y^2}{1+y^2}\log(1+y^2)+\frac{2y}{1+y^2}(x+y-\arctan y-L).\]
Đạo hàm dòng này thêm một lần rồi thay vào PDE cho triệt tiêu từng nhóm. Nghiệm giải tích trên $y>0$; phép đổi biến khả nghịch toàn miền với $\eta>0$, nên kết luận duy nhất giải tích theo dữ liệu.

\textbf{Chi tiết đạo hàm cấp hai và triệt tiêu.} Đặt $h=y-\arctan y,k=\frac12\log(1+y^2)$, $u=v(x+h,k)$. Khi đó
\[u_{xx}=v_{\xi\xi},\quad u_{xy}=h'v_{\xi\xi}+k'v_{\xi\eta},\quad
u_{yy}=h'^2v_{\xi\xi}+2h'k'v_{\xi\eta}+k'^2v_{\eta\eta}+h''v_\xi+k''v_\eta.\]
Đặt $D=1+y^2$. Hệ số $v_{\xi\xi}$ là $y^2-2y^2(y^2/D)+D(y^4/D^2)=y^2/D$; hỗn hợp $2k'(-y^2+Dh')=0$; hệ số $v_{\eta\eta}$ là $D(y/D)^2=y^2/D$. Hai hệ số cấp một:
\[\frac{2y}D-y-\frac{(1-y^2)y}D
=\frac{2y-y-y^3-y+y^3}D=0,\quad
\frac{1-y^2}D-\frac{1-y^2}D=0.\]
Để kiểm tra trực tiếp, đặt $A=x+h-L$, có $u_y=h'\log D+2Ay/D$,
\[u_{yy}=h''\log D+\frac{4yh'}D+\frac{2A(1-y^2)}{D^2}.\]
Thế vào PDE: nhóm $\log D$ có hệ số $Dh''-y-(1-y^2)h'/y=0$; nhóm $A$ có hệ số $2(1-y^2)/D-2(1-y^2)/D=0$; nhóm còn lại $-4y^3/D+4yh'=0$.
\begin{center}\includegraphics[width=.97\linewidth,height=11cm,keepaspectratio]{_xu_ly/graphs/BT_2_3_17.png}\end{center}
\textit{Đồ thị minh họa; với nghiệm số, xem phương pháp và tham số trong mã Python kèm theo.}

\clearpage\noindent\texttt{\detokenize{BT_2_3_18}}\par
\section*{Lời giải bổ sung — Bài tập 18 · Mục 2.3}

\subsection*{Kiến thức và trình tự giải}
\textbf{Mục tiêu.} Đưa phương trình elliptic về dạng đơn giản, rồi tìm nghiệm giải tích từ dữ liệu Cauchy mà bài yêu cầu.

\textbf{Bước 1 — Loại hạng hỗn hợp và hạng cấp một do đổi biến.} Tính $\delta<0$, chọn $\xi,\eta$ thực có Jacobian khác $0$. Thay đầy đủ đạo hàm cấp hai; các đạo hàm của tọa độ có thể triệt tiêu với hạng cấp một trong đề. Kết quả thường là $v_{\xi\xi}+v_{\eta\eta}+c v=0$.

\textbf{Bước 2 — Chuyển dữ liệu sang đường tọa độ mới.} Nếu đường dữ liệu trở thành $\xi=0$, đặt $v(0,\eta)=f(\eta)$ và tính $v_\xi(0,\eta)=g(\eta)$ từ đạo hàm được cho. Nếu đường là $\eta=\eta_0$, đổi vai trò hai biến. Không lấy nguyên dữ liệu đạo hàm cũ làm dữ liệu mới khi hệ số dây chuyền khác $1$.

\textbf{Bước 3 — Khai triển và lập truy hồi.} Viết


\[
v(\xi,\eta)=\sum_{n=0}^{\infty}\frac{\xi^n}{n!}a_n(\eta),\qquad
a_0=f,\quad a_1=g,\quad a_{n+2}=-a_n''-c a_n.
\]


Tính từng $a_2,a_3,\ldots$. Nếu các hệ số sau một bậc đều bằng $0$, chuỗi kết thúc và cho một đa thức chính xác. Nếu không, cần xét hội tụ trong miền dữ liệu giải tích; chuỗi hình thức chưa tự là nghiệm.

\textbf{Bước 4 — Đổi ngược và kiểm tra.} Thế $\xi(x,y),\eta(x,y)$ vào nghiệm. Tính đạo hàm hoặc dùng toán tử đã được chứng minh tương đương để kiểm tra PDE, đồng thời kiểm tra từng dữ liệu gốc. Kết luận duy nhất trong lớp nghiệm giải tích trên miền liên thông mà phép đổi biến hợp lệ; dữ liệu Cauchy elliptic nói chung có thể không ổn định, nên không suy ra tồn tại cho dữ liệu tùy ý chỉ từ cách lập chuỗi.
\subsection*{Lời giải áp dụng cho bài}
\textbf{(a) Phân loại và nhận dạng toán tử.} $A=1$, $B=x^2$, $C=x^4+1$, $\delta=-1<0$. Đặt $D=\partial_x+x^2\partial_y$. Quy tắc đạo hàm tích cho $D^2u=u_{xx}+2x^2u_{xy}+x^4u_{yy}+2xu_y$, nên PDE là $D^2u+u_{yy}+u=0$. Chọn $\xi=x$, $\eta=y-x^3/3$: $D\xi=1$, $D\eta=0$, $J=1$. Do đó
\[\boxed{v_{\xi\xi}+v_{\eta\eta}+v=0.}\]

\textbf{(b) Chuỗi giải tích.} Tại $x=0$, $\eta=y$, $u_x=v_\xi$, nên $v(0,\eta)=\cos\eta$, $v_\xi(0,\eta)=0$. Khai triển $v=\sum a_n(\eta)\xi^n/n!$ cho
\[a_0=\cos\eta,\quad a_1=0,\quad a_{n+2}=-a_n''-a_n.\]
Vì $a_0''=-a_0$, có $a_2=0$; $a_1=0$ cho $a_3=0$, rồi toàn bộ các hệ số sau bằng $0$. Vậy $v=\cos\eta$, duy nhất trong lớp giải tích.

\textbf{(c) Đổi ngược và thế vào phương trình.}
\[\boxed{u=\cos(y-x^3/3).}\]
Đặt $s=y-x^3/3$. Ta có $u_x=x^2\sin s$, $u_y=-\sin s$, $u_{xx}=2x\sin s-x^4\cos s$, $u_{xy}=x^2\cos s$, $u_{yy}=-\cos s$. Tổng vế trái:
\[(2x\sin s-x^4\cos s)+2x^4\cos s-(x^4+1)\cos s-2x\sin s+\cos s=0.\]
Tại $x=0$, $u=\cos y$, $u_x=0$. Nghiệm xác định giải tích trên toàn $\mathbb R^2$.
\begin{center}\includegraphics[width=.97\linewidth,height=11cm,keepaspectratio]{_xu_ly/graphs/BT_2_3_18.png}\end{center}
\textit{Đồ thị minh họa; với nghiệm số, xem phương pháp và tham số trong mã Python kèm theo.}

\clearpage\noindent\texttt{\detokenize{VD_3_1_2_01}}\par
\section*{Lời giải bổ sung — Ví dụ 1 · Mục 3.1.2}

\subsection*{Kiến thức và trình tự giải}
\textbf{Bước 1 — Nhận dạng dữ liệu và tốc độ.} Viết $u_{tt}-a^2u_{xx}=F$, $u(x,0)=f(x)$, $u_t(x,0)=g(x)$; tốc độ là $a>0$, không phải $a^2$. Với dữ liệu trơn, cần $f\in C^2$, $g\in C^1$. Nếu dữ liệu có góc hoặc bước nhảy, công thức vẫn mô tả sóng từng miền nhưng phải xác định nơi mất tính trơn.

\textbf{Bước 2 — Tìm hai sóng chạy.} Đặt $\xi=x-at$, $\eta=x+at$. Quy tắc dây chuyền cho $u_{tt}-a^2u_{xx}=-4a^2u_{\xi\eta}$, nên khi $F=0$, $u=A(\xi)+B(\eta)$. Tại $t=0$, $A+B=f$, $a(B'-A')=g$. Chọn một nguyên hàm $H'=g$ rồi lấy $A=f/2-H/(2a)$, $B=f/2+H/(2a)$; hằng số của $H$ sẽ triệt tiêu.


\[
u_0(x,t)=\frac{f(x-at)+f(x+at)}2+\frac{H(x+at)-H(x-at)}{2a}.
\]


\textbf{Bước 3 — Tính tích phân theo dữ liệu thật.} Tìm nguyên hàm trên từng khoảng, nối liên tục nếu $g$ có bước nhảy; thay hai cận $x\pm at$. Dùng công thức cộng lượng giác trước khi kết luận chu kỳ. Với dữ liệu chỉ thị, tích phân bằng độ dài giao của $[x-at,x+at]$ với đoạn hỗ trợ, nhân hệ số của $g$.

\textbf{Bước 4 — Nếu có nguồn, thêm Duhamel.} Nghiệm riêng có dữ liệu đầu bằng không là


\[
u_F(x,t)=\frac1{2a}\int_0^t\int_{x-a(t-s)}^{x+a(t-s)}F(y,s)\,dy\,ds.
\]


Tính tích phân trong trước. Nếu $F$ là đa thức hoặc một mode lượng giác, có thể đoán $u_F$ rồi giải ODE theo $t$, nhưng vẫn phải buộc $u_F(x,0)=(u_F)_t(x,0)=0$. Nghiệm cuối là $u_0+u_F$.

\textbf{Bước 5 — Kiểm tra và đọc ý nghĩa.} Đạo hàm công thức: $u_x=(f'(\xi)+f'(\eta))/2+(g(\eta)-g(\xi))/(2a)$, $u_t=a(f'(\eta)-f'(\xi))/2+(g(\eta)+g(\xi))/2$. Tại $t=0$ thu lại $f,g$; đạo hàm tiếp cho PDE. Các điểm gãy $s=b$ phát sinh đường $x\pm at=b$. Chu kỳ phải đúng cho mọi $x,t$, không chỉ sau khi một mặt sóng đi qua một điểm cố định.
\subsection*{Lời giải áp dụng cho bài}
\textbf{Bước 1. Tách li độ và vận tốc.} Với $a=1$ ở (i), (iii), (iv), và $a=2$ ở (ii), dùng
\[u=\tfrac12[f(x-at)+f(x+at)]+\frac{H(x+at)-H(x-at)}{2a},\qquad H'=g.\]
\textbf{Bước 2. Tính từng trường hợp.}

(i) $g=0$, nên $H=0$. Công thức $\sin(x-t)+\sin(x+t)=2\sin x\cos t$ cho $u=\sin x\cos t$.

(ii) Chỉ đổi độ dịch thành $2t$: $u=\sin x\cos2t$. Hai chu kỳ cơ bản lần lượt là $2\pi$ và $\pi$.

(iii) $f=0$, $g(s)=s|s|$. Trên $s>0$, nguyên hàm là $s^3/3$; trên $s<0$, là $-s^3/3$. Nối tại $0$ được $H(s)=|s|^3/3$, do đó
\[u=\frac{|x+t|^3-|x-t|^3}{6}.\]
Hệ số $1/2$ ở nguyên hàm trong bản scan không đúng: $(|s|^3/3)'=s|s|$. Hàm $g$ thuộc $C^1$, nên nghiệm thuộc $C^2$, kể cả trên $x\pm t=0$; đạo hàm cấp ba có thể nhảy.

(iv) $g(s)=|s|$ có nguyên hàm $H=s|s|/2$, nên
\[u=\frac{(x+t)|x+t|-(x-t)|x-t|}{4}.\]
Vì $H\in C^1$ nhưng không $C^2$ tại $0$, nghiệm $C^1$ và không $C^2$ trên hai đặc trưng $x=\pm t$.

\textbf{Bước 3. Kiểm tra.} Tại $t=0$, hai hiệu nguyên hàm bằng không; đạo hàm theo $t$ cho $g(x)$. Mỗi hạng chỉ phụ thuộc $x\pm at$ nên thỏa PDE nơi đạo hàm tồn tại. Hai nghiệm (iii), (iv) tăng theo $t$ ở những điểm thích hợp, nên không tuần hoàn.

\clearpage\noindent\texttt{\detokenize{VD_3_1_2_02}}\par
\section*{Lời giải bổ sung — Ví dụ 2 · Mục 3.1.2}

\subsection*{Kiến thức và trình tự giải}
\textbf{Bước 1 — Nhận dạng dữ liệu và tốc độ.} Viết $u_{tt}-a^2u_{xx}=F$, $u(x,0)=f(x)$, $u_t(x,0)=g(x)$; tốc độ là $a>0$, không phải $a^2$. Với dữ liệu trơn, cần $f\in C^2$, $g\in C^1$. Nếu dữ liệu có góc hoặc bước nhảy, công thức vẫn mô tả sóng từng miền nhưng phải xác định nơi mất tính trơn.

\textbf{Bước 2 — Tìm hai sóng chạy.} Đặt $\xi=x-at$, $\eta=x+at$. Quy tắc dây chuyền cho $u_{tt}-a^2u_{xx}=-4a^2u_{\xi\eta}$, nên khi $F=0$, $u=A(\xi)+B(\eta)$. Tại $t=0$, $A+B=f$, $a(B'-A')=g$. Chọn một nguyên hàm $H'=g$ rồi lấy $A=f/2-H/(2a)$, $B=f/2+H/(2a)$; hằng số của $H$ sẽ triệt tiêu.


\[
u_0(x,t)=\frac{f(x-at)+f(x+at)}2+\frac{H(x+at)-H(x-at)}{2a}.
\]


\textbf{Bước 3 — Tính tích phân theo dữ liệu thật.} Tìm nguyên hàm trên từng khoảng, nối liên tục nếu $g$ có bước nhảy; thay hai cận $x\pm at$. Dùng công thức cộng lượng giác trước khi kết luận chu kỳ. Với dữ liệu chỉ thị, tích phân bằng độ dài giao của $[x-at,x+at]$ với đoạn hỗ trợ, nhân hệ số của $g$.

\textbf{Bước 4 — Nếu có nguồn, thêm Duhamel.} Nghiệm riêng có dữ liệu đầu bằng không là


\[
u_F(x,t)=\frac1{2a}\int_0^t\int_{x-a(t-s)}^{x+a(t-s)}F(y,s)\,dy\,ds.
\]


Tính tích phân trong trước. Nếu $F$ là đa thức hoặc một mode lượng giác, có thể đoán $u_F$ rồi giải ODE theo $t$, nhưng vẫn phải buộc $u_F(x,0)=(u_F)_t(x,0)=0$. Nghiệm cuối là $u_0+u_F$.

\textbf{Bước 5 — Kiểm tra và đọc ý nghĩa.} Đạo hàm công thức: $u_x=(f'(\xi)+f'(\eta))/2+(g(\eta)-g(\xi))/(2a)$, $u_t=a(f'(\eta)-f'(\xi))/2+(g(\eta)+g(\xi))/2$. Tại $t=0$ thu lại $f,g$; đạo hàm tiếp cho PDE. Các điểm gãy $s=b$ phát sinh đường $x\pm at=b$. Chu kỳ phải đúng cho mọi $x,t$, không chỉ sau khi một mặt sóng đi qua một điểm cố định.
\subsection*{Lời giải áp dụng cho bài}
\textbf{Bước 1. Giải phần có nguồn với dữ liệu không.} Nguồn $F(x,t)=xt$, $a=1$. Tích phân trong của Duhamel là
\[\int_{x-(t-s)}^{x+(t-s)}ys\,dy=2xs(t-s).\]
Suy $u_F=x\int_0^t s(t-s)ds=xt^3/6$.

\textbf{Bước 2. Thêm phần thuần nhất.} Ở (i), $f=g=0$, nên $u=xt^3/6$. Ở (ii), $f=\sin x$, $g=x|x|$, nguyên hàm đúng là $|x|^3/3$:
\[u=\sin x\cos t+\frac{|x+t|^3-|x-t|^3}{6}+\frac{xt^3}{6}.\]
\textbf{Bước 3. Kiểm tra.} $(u_F)_{tt}=xt$, $(u_F)_{xx}=0$, $u_F(x,0)=(u_F)_t(x,0)=0$. Phần còn lại thỏa phương trình thuần nhất và cho đúng $f,g$. Dữ liệu $g\in C^1$ bảo đảm nghiệm $C^2$; không cần loại hai đường $x=\pm t$ ở mức đạo hàm cấp hai.

\textbf{Chi tiết hai tích phân nguồn.} Đặt $h=t-s$. Tích phân trong bằng
\[s\left[\frac{y^2}{2}\right]_{x-h}^{x+h}
=\frac s2[(x+h)^2-(x-h)^2]=2xsh.\]
Nhân $1/2$ của Duhamel, tích phân ngoài là
\[x\int_0^t(ts-s^2)ds=x\left[\frac{ts^2}{2}-\frac{s^3}{3}\right]_0^t
=x(t^3/2-t^3/3)=xt^3/6.\]

\clearpage\noindent\texttt{\detokenize{VD_3_1_2_03}}\par
\section*{Lời giải bổ sung — Ví dụ 3 · Mục 3.1.2}

\subsection*{Kiến thức và trình tự giải}
\textbf{Bước 1 — Nhận dạng dữ liệu và tốc độ.} Viết $u_{tt}-a^2u_{xx}=F$, $u(x,0)=f(x)$, $u_t(x,0)=g(x)$; tốc độ là $a>0$, không phải $a^2$. Với dữ liệu trơn, cần $f\in C^2$, $g\in C^1$. Nếu dữ liệu có góc hoặc bước nhảy, công thức vẫn mô tả sóng từng miền nhưng phải xác định nơi mất tính trơn.

\textbf{Bước 2 — Tìm hai sóng chạy.} Đặt $\xi=x-at$, $\eta=x+at$. Quy tắc dây chuyền cho $u_{tt}-a^2u_{xx}=-4a^2u_{\xi\eta}$, nên khi $F=0$, $u=A(\xi)+B(\eta)$. Tại $t=0$, $A+B=f$, $a(B'-A')=g$. Chọn một nguyên hàm $H'=g$ rồi lấy $A=f/2-H/(2a)$, $B=f/2+H/(2a)$; hằng số của $H$ sẽ triệt tiêu.


\[
u_0(x,t)=\frac{f(x-at)+f(x+at)}2+\frac{H(x+at)-H(x-at)}{2a}.
\]


\textbf{Bước 3 — Tính tích phân theo dữ liệu thật.} Tìm nguyên hàm trên từng khoảng, nối liên tục nếu $g$ có bước nhảy; thay hai cận $x\pm at$. Dùng công thức cộng lượng giác trước khi kết luận chu kỳ. Với dữ liệu chỉ thị, tích phân bằng độ dài giao của $[x-at,x+at]$ với đoạn hỗ trợ, nhân hệ số của $g$.

\textbf{Bước 4 — Nếu có nguồn, thêm Duhamel.} Nghiệm riêng có dữ liệu đầu bằng không là


\[
u_F(x,t)=\frac1{2a}\int_0^t\int_{x-a(t-s)}^{x+a(t-s)}F(y,s)\,dy\,ds.
\]


Tính tích phân trong trước. Nếu $F$ là đa thức hoặc một mode lượng giác, có thể đoán $u_F$ rồi giải ODE theo $t$, nhưng vẫn phải buộc $u_F(x,0)=(u_F)_t(x,0)=0$. Nghiệm cuối là $u_0+u_F$.

\textbf{Bước 5 — Kiểm tra và đọc ý nghĩa.} Đạo hàm công thức: $u_x=(f'(\xi)+f'(\eta))/2+(g(\eta)-g(\xi))/(2a)$, $u_t=a(f'(\eta)-f'(\xi))/2+(g(\eta)+g(\xi))/2$. Tại $t=0$ thu lại $f,g$; đạo hàm tiếp cho PDE. Các điểm gãy $s=b$ phát sinh đường $x\pm at=b$. Chu kỳ phải đúng cho mọi $x,t$, không chỉ sau khi một mặt sóng đi qua một điểm cố định.
\subsection*{Lời giải áp dụng cho bài}
\textbf{Bước 1. Đặt tích phân trong.} Viết $V=(2a)^{-1}\int_0^t G(x,t,s)ds$, với
\[G=\int_{x-a(t-s)}^{x+a(t-s)}F(y,s)dy.\]
Ở đây $s$ là tham số khi lấy đạo hàm theo $t$.

\textbf{Bước 2. Đạo hàm các cận.} Với $F_\pm=F(x\pm a(t-s),s)$,
\[G_t=a(F_++F_-),\quad G_x=F_+-F_-,\quad G_{tt}=a^2(F_{x,+}-F_{x,-})=a^2G_{xx}.\]
Khi $s=t$, $G=0$ và $G_t=2aF(x,t)$.

\textbf{Bước 3. Đạo hàm tích phân ngoài.} Quy tắc Leibniz cho
\[V_t=\frac1{2a}\int_0^tG_tds,\qquad V_{tt}=F(x,t)+\frac1{2a}\int_0^tG_{tt}ds.\]
Do đó $V_{tt}-a^2V_{xx}=F$. Tại $t=0$ cả $V,V_t$ bằng không. Các đạo hàm trên hợp lệ nếu $F$ có các đạo hàm liên tục cần dùng (chẳng hạn $F\in C^1$ theo không gian và liên tục theo thời gian). Thêm nghiệm d'Alembert của dữ liệu $f,g$ được nghiệm đầy đủ.

\clearpage\noindent\texttt{\detokenize{VD_3_1_2_04}}\par
\section*{Lời giải bổ sung — Ví dụ 4 · Mục 3.1.2}

\subsection*{Kiến thức và trình tự giải}
\textbf{Bước 1 — Đọc đúng điều kiện biên.} Biên $u(0,t)=0$ dùng mở rộng lẻ; biên $u_x(0,t)=0$ dùng mở rộng chẵn. Với dữ liệu trên $x\ge0$, đặt $\widetilde f(-s)=-f(s)$ hoặc $f(s)$ tương ứng, làm tương tự cho $g$. Mở rộng lẻ là $C^2$ nếu $f(0)=f''(0)=0$, và mở rộng lẻ của $g$ là $C^1$ nếu $g(0)=0$. Mở rộng chẵn đòi $f'(0)=g'(0)=0$ cho mức trơn ấy.

\textbf{Bước 2 — Lập nguyên hàm sau khi phản xạ.} Đặt $H(s)=\int_0^s\widetilde g(r)dr$. Nếu $\widetilde g$ lẻ thì $H$ chẵn; nếu $\widetilde g$ chẵn thì $H$ lẻ. Chính nguyên hàm, chứ không phải vận tốc, xuất hiện trong


\[
u=\frac{\widetilde f(x-at)+\widetilde f(x+at)}2+\frac{H(x+at)-H(x-at)}{2a}.
\]


\textbf{Bước 3 — Tách miền trước và sau khi chạm biên.} Với $u(0,t)=h(t)$, ở $x\ge at$ dùng d'Alembert với dữ liệu gốc. Ở $0<x<at$:


\[
u=h(t-x/a)+\frac{f(x+at)-f(at-x)}2+\frac1{2a}\int_{at-x}^{at+x}g(s)ds.
\]


Với $u_x(0,t)=h(t)$, dùng nghiệm của các mở rộng chẵn rồi trừ $a\int_0^{t-x/a}h(s)ds$ trên miền $x<at$. Tính riêng đạo hàm theo $x$ của hạng này để kiểm tra biên.

\textbf{Bước 4 — Trên đoạn hữu hạn, phản xạ ở cả hai đầu.} Dirichlet ở hai đầu cho mở rộng lẻ tại $0,L$, chu kỳ không gian $2L$. Neumann ở hai đầu cho mở rộng chẵn tại $0,L$, cũng chu kỳ $2L$. Một đầu lẻ, một đầu chẵn cho $\widetilde f(s+2L)=-\widetilde f(s)$ và chu kỳ $4L$, tương tự $g$. Vì thế chu kỳ thời gian thường là $2L/a$ hoặc $4L/a$. Với Neumann hai đầu còn phải xét trung bình $g$: trung bình khác không sinh hạng tuyến tính theo $t$, nên nghiệm không tuần hoàn.

\textbf{Bước 5 — Kiểm tra góc và đường ghép.} Biên Dirichlet không thuần nhất cần $h(0)=f(0)$, $h'(0)=g(0)$, $h''(0)=a^2f''(0)$ để có nghiệm $C^2$ cho phương trình thuần nhất. So sánh hai vết trên $x=at$, gồm $u,u_x,u_t$ rồi đạo hàm cấp hai. Nếu không tương thích, nêu rõ mức trơn còn lại thay vì khẳng định nghiệm cổ điển trên toàn miền.
\subsection*{Lời giải áp dụng cho bài}
\textbf{Bước 1. Xác định hai miền.} Với biên $u(0,t)=h(t)$, ở $x\ge at$ dùng dữ liệu đầu; ở $x<at$ thêm tín hiệu biên tại thời điểm trễ $t-x/a$ và lấy hiệu hai giá trị $f$.

\textbf{Bước 2. Tính cụ thể.}

(i) $a=1$, $f=g=0$, $h=t^3$. Trước khi sóng biên đến, $u=0$; sau đó $u=(t-x)^3$. Gộp lại $u=(t-x)_+^3$. Hàm này $C^2$ trên đường $x=t$.

(ii) $a=2$, $f=\sin x$, $g=1$, $h=0$. Ngoài mặt sóng, phần vận tốc là $t$. Trong mặt sóng, đoạn tích phân $[2t-x,2t+x]$ dài $2x$, cho $x/2$. Phần li độ ở cả hai miền bằng $\sin x\cos2t$. Vì vậy
\[u=\sin x\cos2t+\min(t,x/2).\]
Nghiệm liên tục nhưng không $C^1$ trên $x=2t$: vận tốc đầu $g(0)=1$ trái với $h'(0)=0$.

(iii) $a=1$, $f=x^2$, $g=6x$, $h=t^2$. Ngoài mặt sóng, half-sum bằng $x^2+t^2$, tích phân vận tốc bằng $6xt$. Trong mặt sóng,
\[(t-x)^2+\tfrac12[(x+t)^2-(t-x)^2]+\tfrac12\int_{t-x}^{t+x}6sds=x^2+t^2+6xt.\]
\textbf{Bước 3. Kiểm tra.} Nghiệm (iii) có $u_{tt}=u_{xx}=2$, $u(x,0)=x^2$, $u_t(x,0)=6x$, $u(0,t)=t^2$. Ba điều kiện tương thích tại góc đều đúng.

\textbf{Chi tiết rút gọn câu (iii).} Ngoài mặt sóng, half-sum là $[(x-t)^2+(x+t)^2]/2=x^2+t^2$, còn tích phân là $3[(x+t)^2-(x-t)^2]/2=6xt$. Trong mặt sóng, half-difference là $[(x+t)^2-(t-x)^2]/2=2xt$, tích phân vận tốc là $3[(t+x)^2-(t-x)^2]/2=6xt$. Cộng thêm biên $(t-x)^2=x^2-2xt+t^2$ cho cùng $x^2+t^2+6xt$.

\clearpage\noindent\texttt{\detokenize{VD_3_1_2_05}}\par
\section*{Lời giải bổ sung — Ví dụ 5 · Mục 3.1.2}

\subsection*{Kiến thức và trình tự giải}
\textbf{Bước 1 — Đọc đúng điều kiện biên.} Biên $u(0,t)=0$ dùng mở rộng lẻ; biên $u_x(0,t)=0$ dùng mở rộng chẵn. Với dữ liệu trên $x\ge0$, đặt $\widetilde f(-s)=-f(s)$ hoặc $f(s)$ tương ứng, làm tương tự cho $g$. Mở rộng lẻ là $C^2$ nếu $f(0)=f''(0)=0$, và mở rộng lẻ của $g$ là $C^1$ nếu $g(0)=0$. Mở rộng chẵn đòi $f'(0)=g'(0)=0$ cho mức trơn ấy.

\textbf{Bước 2 — Lập nguyên hàm sau khi phản xạ.} Đặt $H(s)=\int_0^s\widetilde g(r)dr$. Nếu $\widetilde g$ lẻ thì $H$ chẵn; nếu $\widetilde g$ chẵn thì $H$ lẻ. Chính nguyên hàm, chứ không phải vận tốc, xuất hiện trong


\[
u=\frac{\widetilde f(x-at)+\widetilde f(x+at)}2+\frac{H(x+at)-H(x-at)}{2a}.
\]


\textbf{Bước 3 — Tách miền trước và sau khi chạm biên.} Với $u(0,t)=h(t)$, ở $x\ge at$ dùng d'Alembert với dữ liệu gốc. Ở $0<x<at$:


\[
u=h(t-x/a)+\frac{f(x+at)-f(at-x)}2+\frac1{2a}\int_{at-x}^{at+x}g(s)ds.
\]


Với $u_x(0,t)=h(t)$, dùng nghiệm của các mở rộng chẵn rồi trừ $a\int_0^{t-x/a}h(s)ds$ trên miền $x<at$. Tính riêng đạo hàm theo $x$ của hạng này để kiểm tra biên.

\textbf{Bước 4 — Trên đoạn hữu hạn, phản xạ ở cả hai đầu.} Dirichlet ở hai đầu cho mở rộng lẻ tại $0,L$, chu kỳ không gian $2L$. Neumann ở hai đầu cho mở rộng chẵn tại $0,L$, cũng chu kỳ $2L$. Một đầu lẻ, một đầu chẵn cho $\widetilde f(s+2L)=-\widetilde f(s)$ và chu kỳ $4L$, tương tự $g$. Vì thế chu kỳ thời gian thường là $2L/a$ hoặc $4L/a$. Với Neumann hai đầu còn phải xét trung bình $g$: trung bình khác không sinh hạng tuyến tính theo $t$, nên nghiệm không tuần hoàn.

\textbf{Bước 5 — Kiểm tra góc và đường ghép.} Biên Dirichlet không thuần nhất cần $h(0)=f(0)$, $h'(0)=g(0)$, $h''(0)=a^2f''(0)$ để có nghiệm $C^2$ cho phương trình thuần nhất. So sánh hai vết trên $x=at$, gồm $u,u_x,u_t$ rồi đạo hàm cấp hai. Nếu không tương thích, nêu rõ mức trơn còn lại thay vì khẳng định nghiệm cổ điển trên toàn miền.
\subsection*{Lời giải áp dụng cho bài}
\textbf{Bước 1. Điều kiện cần ở góc.} Nếu $u\in C^2$ đến $(0,0)$ thì $u(0,0)=f(0)=h(0)$, $u_t(0,0)=g(0)=h'(0)$. Từ PDE ở góc, $h''(0)=u_{tt}(0,0)=a^2f''(0)$.

\textbf{Bước 2. So sánh hai công thức trên $x=at$.} Đặt $s=x+at$, $r=at-x$. Công thức trong là $h(t-x/a)+(f(s)-f(r))/2+(2a)^{-1}\int_r^s g$. Ngoài là half-sum $f(x-at),f(s)$ cộng tích phân $g$. Tại $r=0$, chênh lệch vết trong trừ ngoài của $u$ là $h(0)-f(0)$.

\textbf{Bước 3. Đạo hàm cấp một.} Trong miền phản xạ,
\[u_x=-h'(t-x/a)/a+\tfrac12[f'(s)+f'(r)]+\frac{g(s)+g(r)}{2a}.\]
Ngoài miền phản xạ, $u_x=(f'(s)+f'(x-at))/2+(g(s)-g(x-at))/(2a)$. Hiệu vết là $(g(0)-h'(0))/a$; tương tự hiệu vết $u_t$ là $h'(0)-g(0)$.

\textbf{Bước 4. Đạo hàm cấp hai.} Hiệu vết $u_{xx}$ bằng $h''(0)/a^2-f''(0)$, vì các hạng $g'$ trùng nhau. Hiệu $u_{xt}$ bằng $af''(0)-h''(0)/a$; $u_{tt}$ bằng $h''(0)-a^2f''(0)$. Như vậy ba điều kiện đã nêu đủ để ghép thành $C^2$, nếu $f\in C^2,g\in C^1,h\in C^2$. Nghiệm thỏa dữ liệu và PDE theo công thức; tính duy nhất suy từ hai sóng chạy hoặc năng lượng của hiệu hai nghiệm.

\textbf{Chi tiết so các vết cấp hai.} Trong miền phản xạ,
\[\begin{aligned}
u_{xx}&=h''(\tau)/a^2+(f''(s)-f''(r))/2+(g'(s)-g'(r))/(2a),\\
u_{xt}&=-h''(\tau)/a+a(f''(s)+f''(r))/2+(g'(s)+g'(r))/2,\\
u_{tt}&=h''(\tau)+a^2(f''(s)-f''(r))/2+a(g'(s)-g'(r))/2.
\end{aligned}\]
Ngoài miền phản xạ, viết $\xi=x-at$, $\eta=x+at$:
\[\begin{aligned}
u_{xx}&=(f''(\xi)+f''(\eta))/2+(g'(\eta)-g'(\xi))/(2a),\\
u_{xt}&=a(f''(\eta)-f''(\xi))/2+(g'(\eta)+g'(\xi))/2,\\
u_{tt}&=a^2(f''(\xi)+f''(\eta))/2+a(g'(\eta)-g'(\xi))/2.
\end{aligned}\]
Trên $x=at$, $r=\xi=\tau=0$, $s=\eta=2at$. Trừ từng cặp: các hạng tại $2at$ và mọi hạng $g'$ mất, còn lần lượt $h''(0)/a^2-f''(0)$, $af''(0)-h''(0)/a$, $h''(0)-a^2f''(0)$.

\clearpage\noindent\texttt{\detokenize{BT_3_1_3_01}}\par
\section*{Lời giải bổ sung — Bài tập 1 · Mục 3.1.3}

\subsection*{Kiến thức và trình tự giải}
\textbf{Bước 1 — Chọn năng lượng theo toán tử.} Với sóng dùng $E=\frac12\int(u_t^2+a^2u_x^2)dx$. Với Klein–Gordon có thêm $u$, thêm $\frac12\int u^2dx$. Với dầm $u_{tt}=-a^2u_{xxxx}$ dùng $E=\frac12\int(u_t^2+a^2u_{xx}^2)dx$, không thay $u_{xx}$ bằng $u_x$.

\textbf{Bước 2 — Đạo hàm và tích phân từng phần.} Với $u_{tt}-a^2u_{xx}+\gamma u_t+\kappa u=F$, năng lượng có hạng $\kappa u^2$ cho


\[
E'=a^2[u_tu_x]_0^L-\gamma\int u_t^2dx+\int Fu_tdx.
\]


Với dầm, tích phân hai lần cho $E'=a^2[u_{tx}u_{xx}-u_tu_{xxx}]_0^L$. Viết riêng từng hạng biên rồi dùng điều kiện biên và đạo hàm theo thời gian của điều kiện đó để cho chúng bằng không.

\textbf{Bước 3 — Suy hệ quả.} Không cản, không nguồn thì $E$ bảo toàn. Có cản $\gamma>0$ thì $E'\le0$; nếu $E$ không có hệ số $1/2$, $E'=-2\gamma\int u_t^2\ge-2\gamma E$, suy $e^{-2\gamma t}E(0)\le E(t)\le E(0)$. Với hiệu hai nghiệm cùng dữ liệu, $E(0)=0$ nên có duy nhất.

\textbf{Bước 4 — Tách biến nếu phải tìm nghiệm cụ thể.} Sóng có cản $\gamma=1$ cho $T''+T'+a^2\lambda T=0$. Đặt $T=e^{-t/2}Y$, thu $Y''+(a^2\lambda-1/4)Y=0$, $Y(0)=A$, $Y'(0)=B+A/2$. Phân biệt dao động, nghiệm kép và hàm hyperbolic. Dầm cho $X''''=k^4X$, $T''+a^2k^4T=0$; dùng đủ bốn điều kiện biên để tìm phương trình của $kL$.

\textbf{Bước 5 — Kiểm tra hệ số và khẳng định của đề.} Năng lượng đầu chứa cả $g^2$ và các đạo hàm của $f$. Thử $f=0$, $g\ne0$ để phát hiện công thức bỏ vận tốc. Với dầm, nghiệm $k=0$ phải giải trực tiếp bằng đa thức bậc ba; nghiệm $0$ của phương trình lượng giác–hyperbolic có thể chỉ là nghiệm giả do cơ sở suy biến.

\textbf{Chi tiết phép triệt tiêu tích phân.} Đạo hàm động năng cho $\int u_tu_{tt}$; dùng PDE thay $u_{tt}$. Với sóng, $\int u_tu_{xx}=[u_tu_x]-\int u_{tx}u_x$, và tích phân cuối triệt tiêu đạo hàm thế năng $\int u_xu_{xt}$. Với dầm,


\[
\int u_tu_{xxxx}=[u_tu_{xxx}-u_{tx}u_{xx}]+\int u_{txx}u_{xx}.
\]


Tích phân cuối triệt tiêu đạo hàm $\frac12\int u_{xx}^2$ sau khi nhân hệ số và giữ đúng dấu PDE. Viết hạng biên trước, sau đó mới dùng điều kiện của từng đầu.
\subsection*{Lời giải áp dụng cho bài}
\textbf{Bước 1. Miền hỗ trợ.} Nếu $|x|>M+at$, hai cận $x\pm at$ và toàn đoạn giữa chúng cùng nằm ngoài $[-M,M]$. D'Alembert cho $u=0$; các đạo hàm cũng bằng không ngoài miền ấy khi dữ liệu đủ trơn.

\textbf{Bước 2. Đạo hàm năng lượng.} Đặt $I(t)=\int_{\mathbb R}(a^2u_x^2+u_t^2)dx$. Trên mỗi khoảng thời gian hữu hạn có thể tích phân trên một đoạn cố định đủ lớn, vì sóng có hỗ trợ hữu hạn. Dùng PDE:
\[I'=2a^2\int(u_xu_{xt}+u_tu_{xx})dx=2a^2[u_tu_x]_{-\infty}^{+\infty}=0.\]
\textbf{Bước 3. Giá trị đầu.} Tại $t=0$, $u_x=f'$, $u_t=g$, nên
\[I(t)=I(0)=\int_{\mathbb R}(a^2(f')^2+g^2)dx.\]
Điều kiện hỗ trợ bảo đảm hạng biên bằng không, không phải chỉ vì ký hiệu cận là vô hạn. Với dữ liệu năng lượng, kết quả tiếp tục đúng bằng xấp xỉ dữ liệu trơn.

\textbf{Chi tiết tích phân từng phần trong đạo hàm năng lượng.} Đặt $E=\frac12\int_{-R}^R(u_t^2+a^2u_x^2+\kappa u^2)dx$. Với $u_{tt}-a^2u_{xx}+\gamma u_t+\kappa u=F$,
\[\begin{aligned}
E'&=\int_{-R}^R(u_tu_{tt}+a^2u_xu_{xt}+\kappa uu_t)dx\\
&=\int_{-R}^Ru_t(a^2u_{xx}-\gamma u_t-\kappa u+F)dx
+a^2\int_{-R}^Ru_xu_{xt}dx+\kappa\int_{-R}^Ruu_tdx\\
&=a^2[u_tu_x]_{-R}^R-a^2\int_{-R}^Ru_{tx}u_xdx
+a^2\int_{-R}^Ru_xu_{xt}dx-\gamma\int_{-R}^Ru_t^2dx+\int_{-R}^RFu_tdx\\
&=a^2[u_tu_x]_{-R}^R-\gamma\int_{-R}^Ru_t^2dx+\int_{-R}^RFu_tdx.
\end{aligned}\]
Hai hạng $\kappa uu_t$ triệt tiêu ngay khi thế PDE; hai tích phân $u_xu_{xt}$ triệt tiêu sau tích phân từng phần. Chọn $R$ lớn hơn miền hỗ trợ trên khoảng thời gian xét, nên $u_t=u_x=0$ ở $x=\pm R$. Với $I=2E$, nhân đẳng thức đạo hàm với $2$ để có $I'=0$.

\clearpage\noindent\texttt{\detokenize{BT_3_1_3_02}}\par
\section*{Lời giải bổ sung — Bài tập 2 · Mục 3.1.3}

\subsection*{Kiến thức và trình tự giải}
\textbf{Bước 1 — Nhận dạng dữ liệu và tốc độ.} Viết $u_{tt}-a^2u_{xx}=F$, $u(x,0)=f(x)$, $u_t(x,0)=g(x)$; tốc độ là $a>0$, không phải $a^2$. Với dữ liệu trơn, cần $f\in C^2$, $g\in C^1$. Nếu dữ liệu có góc hoặc bước nhảy, công thức vẫn mô tả sóng từng miền nhưng phải xác định nơi mất tính trơn.

\textbf{Bước 2 — Tìm hai sóng chạy.} Đặt $\xi=x-at$, $\eta=x+at$. Quy tắc dây chuyền cho $u_{tt}-a^2u_{xx}=-4a^2u_{\xi\eta}$, nên khi $F=0$, $u=A(\xi)+B(\eta)$. Tại $t=0$, $A+B=f$, $a(B'-A')=g$. Chọn một nguyên hàm $H'=g$ rồi lấy $A=f/2-H/(2a)$, $B=f/2+H/(2a)$; hằng số của $H$ sẽ triệt tiêu.


\[
u_0(x,t)=\frac{f(x-at)+f(x+at)}2+\frac{H(x+at)-H(x-at)}{2a}.
\]


\textbf{Bước 3 — Tính tích phân theo dữ liệu thật.} Tìm nguyên hàm trên từng khoảng, nối liên tục nếu $g$ có bước nhảy; thay hai cận $x\pm at$. Dùng công thức cộng lượng giác trước khi kết luận chu kỳ. Với dữ liệu chỉ thị, tích phân bằng độ dài giao của $[x-at,x+at]$ với đoạn hỗ trợ, nhân hệ số của $g$.

\textbf{Bước 4 — Nếu có nguồn, thêm Duhamel.} Nghiệm riêng có dữ liệu đầu bằng không là


\[
u_F(x,t)=\frac1{2a}\int_0^t\int_{x-a(t-s)}^{x+a(t-s)}F(y,s)\,dy\,ds.
\]


Tính tích phân trong trước. Nếu $F$ là đa thức hoặc một mode lượng giác, có thể đoán $u_F$ rồi giải ODE theo $t$, nhưng vẫn phải buộc $u_F(x,0)=(u_F)_t(x,0)=0$. Nghiệm cuối là $u_0+u_F$.

\textbf{Bước 5 — Kiểm tra và đọc ý nghĩa.} Đạo hàm công thức: $u_x=(f'(\xi)+f'(\eta))/2+(g(\eta)-g(\xi))/(2a)$, $u_t=a(f'(\eta)-f'(\xi))/2+(g(\eta)+g(\xi))/2$. Tại $t=0$ thu lại $f,g$; đạo hàm tiếp cho PDE. Các điểm gãy $s=b$ phát sinh đường $x\pm at=b$. Chu kỳ phải đúng cho mọi $x,t$, không chỉ sau khi một mặt sóng đi qua một điểm cố định.
\subsection*{Lời giải áp dụng cho bài}
\textbf{Bước 1. Tính sóng bằng nguyên hàm.} Dùng $H'=g$ và half-sum của $f$. Công thức tổng quát trong phần kiến thức cho cách kiểm tra $u_t(x,0)=g(x)$; dưới đây ghi rõ nguyên hàm và kết quả từng câu.

(i) $a=3$, $H=-\cos2s/2$. Hiệu hai cos bằng $2\sin2x\sin6t$, nên
\[u=\cos x\cos3t+\tfrac16\sin2x\sin6t,\qquad T=2\pi/3.\]
(ii) $a=5$, $H=s^2/2$. Hiệu bình phương bằng $20xt$, cho $u=\sin3x\cos15t+xt$, không tuần hoàn.

(iii) $a=2$, $H=\sin s$. Half-sum của bình phương bằng $x^2+4t^2$, nên $u=x^2+4t^2+\cos x\sin2t/2$, không tuần hoàn.

(iv) $a=6$, $H=s/2+\sin2s/4$ vì $\cos^2s=(1+\cos2s)/2$:
\[u=\tfrac12(e^{-|x-6t|}+e^{-|x+6t|})+t/2+\cos2x\sin12t/24.\]
Hạng $t/2$ ngăn tuần hoàn; li độ có góc tại $0$ nên không $C^1$ trên $x=\pm6t$.

(v) $a=1$, $g=(3\sin s-\sin3s)/4$, $H=-3\cos s/4+\cos3s/12$. Suy
\[u=\sin x\cos t-\tfrac12-\tfrac12\cos2x\cos2t+\tfrac34\sin x\sin t-\tfrac1{12}\sin3x\sin3t.\]
Mode tần số $1$ khác không, nên chu kỳ cơ bản là $2\pi$.

(vi) $a=7$, $f=1$. Nguyên hàm của xung tam giác được nối liên tục:
\[H(s)=\begin{cases}0&s\le-1,\\(s+1)^2/2&-1<s\le0,\\1-(1-s)^2/2&0<s<1,\\1&s\ge1.\end{cases}\quad
u=1+\frac{H(x+7t)-H(x-7t)}{14}.\]
Tại mỗi $x$ cố định, sau khi cả hai cận ra ngoài hỗ trợ, $u=15/14$, khác giá trị đầu $1$; vì thế không tuần hoàn. Nghiệm $C^1$, các đạo hàm cấp hai gãy trên $x\pm7t=-1,0,1$.

\textbf{Bước 2. Dựng đồ thị.} Sóng phải $A=f/2-H/(2a)$ được dịch sang phải $at$, sóng trái $B=f/2+H/(2a)$ sang trái $at$. Các hình kèm theo gồm lát cắt thời gian; những vị trí gãy được suy từ các cận chạm điểm gãy của dữ liệu. Với các câu trơn, đạo hàm từng mode xác nhận PDE; với câu (iv), (vi), chỉ khẳng định cổ điển ngoài các đường gãy đã chỉ ra.

\textbf{Chi tiết phép thế nguyên hàm và lượng giác.} Ở (i),
\[\frac{H(x+3t)-H(x-3t)}6
=\frac{\cos(2x-6t)-\cos(2x+6t)}{12}
=\frac{\sin2x\sin6t}{6}.\]
Ở (ii), $[(x+5t)^2-(x-5t)^2]/20=xt$. Ở (iii), $[\sin(x+2t)-\sin(x-2t)]/4=\cos x\sin2t/2$. Ở (iv), phần tuyến tính của $H$ cho $(12t/2)/12=t/2$; phần $\sin2s/4$ cho $2\cos2x\sin12t/48=\cos2x\sin12t/24$.

Ở (v), $f(s)=\sin s-\cos^2s=\sin s-1/2-\cos2s/2$, nên half-sum là $\sin x\cos t-1/2-\cos2x\cos2t/2$. Từ $\sin^3s=(3\sin s-\sin3s)/4$, nguyên hàm đã ghi cho
\[\frac{H(x+t)-H(x-t)}2=\frac34\sin x\sin t-\frac1{12}\sin3x\sin3t.\]
Ở (vi), trên $-1<s<0$, $H=\int_{-1}^s(r+1)dr=(s+1)^2/2$; trên $0<s<1$, $H=H(0)+\int_0^s(1-r)dr=1/2+s-s^2/2=1-(1-s)^2/2$. Sau $1$ giữ $H(1)=1$, bảo đảm không có bước nhảy của nguyên hàm.
\begin{center}\includegraphics[width=.97\linewidth,height=11cm,keepaspectratio]{_xu_ly/graphs/BT_3_1_3_02.png}\end{center}
\textit{Đồ thị minh họa; với nghiệm số, xem phương pháp và tham số trong mã Python kèm theo.}
\begin{center}\includegraphics[width=.97\linewidth,height=11cm,keepaspectratio]{_xu_ly/graphs/BT_3_1_3_02_characteristics.png}\end{center}

\clearpage\noindent\texttt{\detokenize{BT_3_1_3_03}}\par
\section*{Lời giải bổ sung — Bài tập 3 · Mục 3.1.3}

\subsection*{Kiến thức và trình tự giải}
\textbf{Bước 1 — Nhận dạng dữ liệu và tốc độ.} Viết $u_{tt}-a^2u_{xx}=F$, $u(x,0)=f(x)$, $u_t(x,0)=g(x)$; tốc độ là $a>0$, không phải $a^2$. Với dữ liệu trơn, cần $f\in C^2$, $g\in C^1$. Nếu dữ liệu có góc hoặc bước nhảy, công thức vẫn mô tả sóng từng miền nhưng phải xác định nơi mất tính trơn.

\textbf{Bước 2 — Tìm hai sóng chạy.} Đặt $\xi=x-at$, $\eta=x+at$. Quy tắc dây chuyền cho $u_{tt}-a^2u_{xx}=-4a^2u_{\xi\eta}$, nên khi $F=0$, $u=A(\xi)+B(\eta)$. Tại $t=0$, $A+B=f$, $a(B'-A')=g$. Chọn một nguyên hàm $H'=g$ rồi lấy $A=f/2-H/(2a)$, $B=f/2+H/(2a)$; hằng số của $H$ sẽ triệt tiêu.


\[
u_0(x,t)=\frac{f(x-at)+f(x+at)}2+\frac{H(x+at)-H(x-at)}{2a}.
\]


\textbf{Bước 3 — Tính tích phân theo dữ liệu thật.} Tìm nguyên hàm trên từng khoảng, nối liên tục nếu $g$ có bước nhảy; thay hai cận $x\pm at$. Dùng công thức cộng lượng giác trước khi kết luận chu kỳ. Với dữ liệu chỉ thị, tích phân bằng độ dài giao của $[x-at,x+at]$ với đoạn hỗ trợ, nhân hệ số của $g$.

\textbf{Bước 4 — Nếu có nguồn, thêm Duhamel.} Nghiệm riêng có dữ liệu đầu bằng không là


\[
u_F(x,t)=\frac1{2a}\int_0^t\int_{x-a(t-s)}^{x+a(t-s)}F(y,s)\,dy\,ds.
\]


Tính tích phân trong trước. Nếu $F$ là đa thức hoặc một mode lượng giác, có thể đoán $u_F$ rồi giải ODE theo $t$, nhưng vẫn phải buộc $u_F(x,0)=(u_F)_t(x,0)=0$. Nghiệm cuối là $u_0+u_F$.

\textbf{Bước 5 — Kiểm tra và đọc ý nghĩa.} Đạo hàm công thức: $u_x=(f'(\xi)+f'(\eta))/2+(g(\eta)-g(\xi))/(2a)$, $u_t=a(f'(\eta)-f'(\xi))/2+(g(\eta)+g(\xi))/2$. Tại $t=0$ thu lại $f,g$; đạo hàm tiếp cho PDE. Các điểm gãy $s=b$ phát sinh đường $x\pm at=b$. Chu kỳ phải đúng cho mọi $x,t$, không chỉ sau khi một mặt sóng đi qua một điểm cố định.
\subsection*{Lời giải áp dụng cho bài}
\textbf{Bước 1. Tách sóng li độ và vận tốc.} $a=5$, $f=10\chi_{[-1,1]}$, $g=\chi_{[-1,1]}$. Tại $x=15$, cận phải $15+5t$ luôn ngoài hỗ trợ. Cận trái $15-5t$ nằm trong đoạn khi $14/5\le t\le16/5$. Đặt $\ell(t)=|[15-5t,15+5t]\cap[-1,1]|$.

\textbf{Bước 2. Viết áp suất theo từng giai đoạn.}
\[u(15,t)=\begin{cases}0&0\le t<14/5,\\5+(5t-14)/10&14/5\le t\le16/5,\\1/5&t>16/5.\end{cases}\]
Hạng $5$ là một nửa của li độ $10$; hạng thứ hai là $\ell/(2a)$. Ở hai đầu mút, giá trị phụ thuộc quy ước chỉ thị đóng; các vết một phía vẫn như trên.

\textbf{Bước 3. So sánh ngưỡng.} Trong khoảng giữa, hàm tăng, lớn nhất $5+2/10=26/5=5.2$, thấp hơn ngưỡng $8$. Sau xung chỉ còn $1/5$. Do đó áp suất không vượt ngưỡng. Không thể chỉ xét tích phân vận tốc mà bỏ sóng li độ gây bước nhảy $5$.
\begin{center}\includegraphics[width=.97\linewidth,height=11cm,keepaspectratio]{_xu_ly/graphs/BT_3_1_3_03.png}\end{center}
\textit{Đồ thị minh họa; với nghiệm số, xem phương pháp và tham số trong mã Python kèm theo.}

\clearpage\noindent\texttt{\detokenize{BT_3_1_3_04}}\par
\section*{Lời giải bổ sung — Bài tập 4 · Mục 3.1.3}

\subsection*{Kiến thức và trình tự giải}
\textbf{Bước 1 — Nhận dạng dữ liệu và tốc độ.} Viết $u_{tt}-a^2u_{xx}=F$, $u(x,0)=f(x)$, $u_t(x,0)=g(x)$; tốc độ là $a>0$, không phải $a^2$. Với dữ liệu trơn, cần $f\in C^2$, $g\in C^1$. Nếu dữ liệu có góc hoặc bước nhảy, công thức vẫn mô tả sóng từng miền nhưng phải xác định nơi mất tính trơn.

\textbf{Bước 2 — Tìm hai sóng chạy.} Đặt $\xi=x-at$, $\eta=x+at$. Quy tắc dây chuyền cho $u_{tt}-a^2u_{xx}=-4a^2u_{\xi\eta}$, nên khi $F=0$, $u=A(\xi)+B(\eta)$. Tại $t=0$, $A+B=f$, $a(B'-A')=g$. Chọn một nguyên hàm $H'=g$ rồi lấy $A=f/2-H/(2a)$, $B=f/2+H/(2a)$; hằng số của $H$ sẽ triệt tiêu.


\[
u_0(x,t)=\frac{f(x-at)+f(x+at)}2+\frac{H(x+at)-H(x-at)}{2a}.
\]


\textbf{Bước 3 — Tính tích phân theo dữ liệu thật.} Tìm nguyên hàm trên từng khoảng, nối liên tục nếu $g$ có bước nhảy; thay hai cận $x\pm at$. Dùng công thức cộng lượng giác trước khi kết luận chu kỳ. Với dữ liệu chỉ thị, tích phân bằng độ dài giao của $[x-at,x+at]$ với đoạn hỗ trợ, nhân hệ số của $g$.

\textbf{Bước 4 — Nếu có nguồn, thêm Duhamel.} Nghiệm riêng có dữ liệu đầu bằng không là


\[
u_F(x,t)=\frac1{2a}\int_0^t\int_{x-a(t-s)}^{x+a(t-s)}F(y,s)\,dy\,ds.
\]


Tính tích phân trong trước. Nếu $F$ là đa thức hoặc một mode lượng giác, có thể đoán $u_F$ rồi giải ODE theo $t$, nhưng vẫn phải buộc $u_F(x,0)=(u_F)_t(x,0)=0$. Nghiệm cuối là $u_0+u_F$.

\textbf{Bước 5 — Kiểm tra và đọc ý nghĩa.} Đạo hàm công thức: $u_x=(f'(\xi)+f'(\eta))/2+(g(\eta)-g(\xi))/(2a)$, $u_t=a(f'(\eta)-f'(\xi))/2+(g(\eta)+g(\xi))/2$. Tại $t=0$ thu lại $f,g$; đạo hàm tiếp cho PDE. Các điểm gãy $s=b$ phát sinh đường $x\pm at=b$. Chu kỳ phải đúng cho mọi $x,t$, không chỉ sau khi một mặt sóng đi qua một điểm cố định.
\subsection*{Lời giải áp dụng cho bài}
\textbf{Bước 1. Xung của nguồn tại thời điểm $T$.} Gọi $w(x,\tau;T)$ là nghiệm thuần nhất với $w(x,0;T)=0$, $w_\tau(x,0;T)=F(x,T)$. D'Alembert cho
\[w(x,\tau;T)=\frac1{2a}\int_{x-a\tau}^{x+a\tau}F(y,T)dy.\]
\textbf{Bước 2. Cộng các xung.} Với $V=\int_0^t w(x,t-T;T)dT$, Leibniz cho $V_{tt}=F(x,t)+\int_0^t w_{\tau\tau}dT$, còn $a^2V_{xx}=\int_0^ta^2w_{xx}dT$. Suy $V_{tt}-a^2V_{xx}=F$ và dữ liệu đầu bằng không.

\textbf{Bước 3. Khôi phục dữ liệu đầu.} Nếu đề có $f,g\ne0$ thì $\int w$ chỉ là phần nguồn; phải thêm $u_0$ của d'Alembert. Miền tích phân là tam giác ngược $0\le T\le t$, $|y-x|\le a(t-T)$; đổi thứ tự tích phân trên miền này cho dạng tích phân Green tương đương.

\textbf{Bước 4. Dạng Green trực tiếp.} Trên tam giác đó, đặt $p=u_T$, $q=-a^2u_y$, nên $p_T+q_y=F$. Tích phân và lấy thông lượng hai cạnh đặc trưng, trên đó $du/dT=u_T\mp au_y$, cho
\[\int_0^t\int_{x-a(t-T)}^{x+a(t-T)}F\,dy\,dT
=2au(x,t)-a[f(x-at)+f(x+at)]-\int_{x-at}^{x+at}g(y)dy.\]
Chia $2a$ thu đúng công thức. Công thức này đồng thời chỉ ra dấu và hệ số của từng dữ liệu.

\textbf{Chi tiết thông lượng trên từng cạnh tam giác Green.} Đặt $\ell(T)=x-a(t-T)$, $r(T)=x+a(t-T)$ và $M(T)=\int_{\ell(T)}^{r(T)}u_T(y,T)dy$. Quy tắc Leibniz cho
\[M'=\int_\ell^r u_{TT}dy-a u_T(r,T)-a u_T(\ell,T).\]
Tích phân $u_{TT}-a^2u_{yy}=F$ theo $y$, rồi thay dòng trên:
\[\int_\ell^r Fdy=M'+a(u_T+a u_y)(\ell,T)+a(u_T-a u_y)(r,T).\]
Cạnh trái có $\ell'=a$, nên $(u_T+a u_y)(\ell,T)=d[u(\ell(T),T)]/dT$; cạnh phải có $r'=-a$, nên $(u_T-a u_y)(r,T)=d[u(r(T),T)]/dT$. Tích phân mỗi cạnh từ $0$ đến $t$ cho
\[a[u(x,t)-f(x-at)],\qquad a[u(x,t)-f(x+at)].\]
Mặt đáy cho $M(t)-M(0)=0-\int_{x-at}^{x+at}g(y)dy$, vì ở đỉnh đoạn tích phân co về độ dài $0$. Cộng ba đóng góp được đẳng thức Green ở bước 4, rồi chia $2a$ thu half-sum, tích phân vận tốc và tích phân nguồn.

\clearpage\noindent\texttt{\detokenize{BT_3_1_3_05}}\par
\section*{Lời giải bổ sung — Bài tập 5 · Mục 3.1.3}

\subsection*{Kiến thức và trình tự giải}
\textbf{Bước 1 — Nhận dạng dữ liệu và tốc độ.} Viết $u_{tt}-a^2u_{xx}=F$, $u(x,0)=f(x)$, $u_t(x,0)=g(x)$; tốc độ là $a>0$, không phải $a^2$. Với dữ liệu trơn, cần $f\in C^2$, $g\in C^1$. Nếu dữ liệu có góc hoặc bước nhảy, công thức vẫn mô tả sóng từng miền nhưng phải xác định nơi mất tính trơn.

\textbf{Bước 2 — Tìm hai sóng chạy.} Đặt $\xi=x-at$, $\eta=x+at$. Quy tắc dây chuyền cho $u_{tt}-a^2u_{xx}=-4a^2u_{\xi\eta}$, nên khi $F=0$, $u=A(\xi)+B(\eta)$. Tại $t=0$, $A+B=f$, $a(B'-A')=g$. Chọn một nguyên hàm $H'=g$ rồi lấy $A=f/2-H/(2a)$, $B=f/2+H/(2a)$; hằng số của $H$ sẽ triệt tiêu.


\[
u_0(x,t)=\frac{f(x-at)+f(x+at)}2+\frac{H(x+at)-H(x-at)}{2a}.
\]


\textbf{Bước 3 — Tính tích phân theo dữ liệu thật.} Tìm nguyên hàm trên từng khoảng, nối liên tục nếu $g$ có bước nhảy; thay hai cận $x\pm at$. Dùng công thức cộng lượng giác trước khi kết luận chu kỳ. Với dữ liệu chỉ thị, tích phân bằng độ dài giao của $[x-at,x+at]$ với đoạn hỗ trợ, nhân hệ số của $g$.

\textbf{Bước 4 — Nếu có nguồn, thêm Duhamel.} Nghiệm riêng có dữ liệu đầu bằng không là


\[
u_F(x,t)=\frac1{2a}\int_0^t\int_{x-a(t-s)}^{x+a(t-s)}F(y,s)\,dy\,ds.
\]


Tính tích phân trong trước. Nếu $F$ là đa thức hoặc một mode lượng giác, có thể đoán $u_F$ rồi giải ODE theo $t$, nhưng vẫn phải buộc $u_F(x,0)=(u_F)_t(x,0)=0$. Nghiệm cuối là $u_0+u_F$.

\textbf{Bước 5 — Kiểm tra và đọc ý nghĩa.} Đạo hàm công thức: $u_x=(f'(\xi)+f'(\eta))/2+(g(\eta)-g(\xi))/(2a)$, $u_t=a(f'(\eta)-f'(\xi))/2+(g(\eta)+g(\xi))/2$. Tại $t=0$ thu lại $f,g$; đạo hàm tiếp cho PDE. Các điểm gãy $s=b$ phát sinh đường $x\pm at=b$. Chu kỳ phải đúng cho mọi $x,t$, không chỉ sau khi một mặt sóng đi qua một điểm cố định.
\subsection*{Lời giải áp dụng cho bài}
\textbf{Bước 1. Giải phần thuần nhất.} Thay $f,g$ vào d'Alembert. Nguyên hàm cần dùng lần lượt là $-\cos2s/2$, $\arctan s$, $-\cos s$, $e^s$, $\sin s$ và $H$ của xung tam giác. Sau đó tìm phần nguồn có hai dữ liệu đầu bằng không.

\textbf{Bước 2. Tính từng nguồn.}

(i) $a=2$, nguồn $x\cos t$ có nghiệm riêng $x(1-\cos t)$ vì đạo hàm không gian cấp hai bằng không. Suy
\[u=e^{-x}\cosh2t+\tfrac14\sin2x\sin4t+x(1-\cos t).\]
(ii) $a=3$. Nguồn $e^x$ dùng $e^xq$, $q''-9q=1$, $q(0)=q'(0)=0$, cho $q=(\cosh3t-1)/9$. Nguồn $\sin t$ độc lập $x$ cho $t-\sin t$:
\[u=\frac{\arctan(x+3t)-\arctan(x-3t)}6+\frac{e^x(\cosh3t-1)}9+t-\sin t.\]
(iii) $a=4$, nguồn $x+t$ cho $xt^2/2+t^3/6$:
\[u=x^2+16t^2+\tfrac14\sin x\sin4t+xt^2/2+t^3/6.\]
(iv) $a=7$, phần vận tốc $e^x$ là $e^x\sinh7t/7$, phần nguồn $xt$ là $xt^3/6$.

(v) $a=1$, $F=\sin(x+t)=\sin x\cos t+\cos x\sin t$. Hai ODE là $q''+q=\cos t$ và $r''+r=\sin t$ với dữ liệu không, ở cộng hưởng. Suy $q=t\sin t/2$, $r=(\sin t-t\cos t)/2$:
\[u=x+\cos x\sin t+\tfrac12\{\sin x\,t\sin t+\cos x(\sin t-t\cos t)\}.\]
(vi) $a=2$, $f$ là hình thang trên $[0,3]$ với nhánh $s,1,3-s$ tại các đoạn $[0,1],[1,2],[2,3]$. $H$ là nguyên hàm tam giác ở Bài 2(vi). Nguồn $-e^x$ cho
\[u=\tfrac12[f(x-2t)+f(x+2t)]+\frac{H(x+2t)-H(x-2t)}4-\frac{e^x(\cosh2t-1)}4.\]
\textbf{Bước 3. Kiểm tra.} Mỗi phần nguồn và đạo hàm thời gian của nó bằng không ở $t=0$. Các ODE vừa lập chứng minh PDE, không chỉ đoán dạng đáp án. Câu (vi) có góc của $f,g$, nên nghiệm được hiểu từng miền/ theo nghĩa yếu trên các đường $x\pm2t=0,1,2,3$ và $x\pm2t=-1,0,1$.

\textbf{Bước 4. Vẽ đặc trưng và các lát cắt.} Với mỗi tốc độ $a$ của từng câu, các đường xuất phát từ $(b,0)$ là $x=b+at$ và $x=b-at$. Nguồn thay đổi giá trị trên miền ảnh hưởng nhưng không đổi hai họ đặc trưng này. Hình kèm lấy $t=0,0.15,0.35,0.65$ và thay trực tiếp vào sáu công thức đã tính; hình thứ hai cho mặt phẳng $x,t$ với các đường đặc trưng mẫu.

\textbf{Chi tiết giải các ODE nguồn.} (i) Đặt $u_F=xq$, có $q''=\cos t$, $q'(0)=q(0)=0$. Tích phân $q'=\sin t$, $q=1-\cos t$. (ii) Với $u_F=e^xq$, $q''-9q=1$: thử $q_p=-1/9$, phần thuần nhất $A\cosh3t+B\sinh3t$. Dữ liệu cho $A=1/9,B=0$. Nguồn $\sin t$ cho $r''=\sin t$, $r'=1-\cos t$, $r=t-\sin t$.

(iii) $q''=x+t$, tích phân theo thời gian với $x$ cố định được $q'=xt+t^2/2$, $q=xt^2/2+t^3/6$. (iv) $q''=xt$ cho $q'=xt^2/2$, $q=xt^3/6$.

(v) Tính trực tiếp $q=t\sin t/2$: $q'=(\sin t+t\cos t)/2$, $q''=\cos t-t\sin t/2$, nên $q''+q=\cos t$. Với $r=(\sin t-t\cos t)/2$, $r'=t\sin t/2$, $r''=(\sin t+t\cos t)/2$, nên $r''+r=\sin t$. Cả $q,r$ và đạo hàm đầu bằng $0$. (vi) $u_F=e^xq$ cho $q''-4q=-1$, $q_p=1/4$, $q=A\cosh2t+B\sinh2t+1/4$; dữ liệu cho $A=-1/4,B=0$.
\begin{center}\includegraphics[width=.97\linewidth,height=11cm,keepaspectratio]{_xu_ly/graphs/BT_3_1_3_05.png}\end{center}
\textit{Đồ thị minh họa; với nghiệm số, xem phương pháp và tham số trong mã Python kèm theo.}
\begin{center}\includegraphics[width=.97\linewidth,height=11cm,keepaspectratio]{_xu_ly/graphs/BT_3_1_3_05_characteristics.png}\end{center}

\clearpage\noindent\texttt{\detokenize{BT_3_1_3_06}}\par
\section*{Lời giải bổ sung — Bài tập 6 · Mục 3.1.3}

\subsection*{Kiến thức và trình tự giải}
\textbf{Bước 1 — Đọc đúng điều kiện biên.} Biên $u(0,t)=0$ dùng mở rộng lẻ; biên $u_x(0,t)=0$ dùng mở rộng chẵn. Với dữ liệu trên $x\ge0$, đặt $\widetilde f(-s)=-f(s)$ hoặc $f(s)$ tương ứng, làm tương tự cho $g$. Mở rộng lẻ là $C^2$ nếu $f(0)=f''(0)=0$, và mở rộng lẻ của $g$ là $C^1$ nếu $g(0)=0$. Mở rộng chẵn đòi $f'(0)=g'(0)=0$ cho mức trơn ấy.

\textbf{Bước 2 — Lập nguyên hàm sau khi phản xạ.} Đặt $H(s)=\int_0^s\widetilde g(r)dr$. Nếu $\widetilde g$ lẻ thì $H$ chẵn; nếu $\widetilde g$ chẵn thì $H$ lẻ. Chính nguyên hàm, chứ không phải vận tốc, xuất hiện trong


\[
u=\frac{\widetilde f(x-at)+\widetilde f(x+at)}2+\frac{H(x+at)-H(x-at)}{2a}.
\]


\textbf{Bước 3 — Tách miền trước và sau khi chạm biên.} Với $u(0,t)=h(t)$, ở $x\ge at$ dùng d'Alembert với dữ liệu gốc. Ở $0<x<at$:


\[
u=h(t-x/a)+\frac{f(x+at)-f(at-x)}2+\frac1{2a}\int_{at-x}^{at+x}g(s)ds.
\]


Với $u_x(0,t)=h(t)$, dùng nghiệm của các mở rộng chẵn rồi trừ $a\int_0^{t-x/a}h(s)ds$ trên miền $x<at$. Tính riêng đạo hàm theo $x$ của hạng này để kiểm tra biên.

\textbf{Bước 4 — Trên đoạn hữu hạn, phản xạ ở cả hai đầu.} Dirichlet ở hai đầu cho mở rộng lẻ tại $0,L$, chu kỳ không gian $2L$. Neumann ở hai đầu cho mở rộng chẵn tại $0,L$, cũng chu kỳ $2L$. Một đầu lẻ, một đầu chẵn cho $\widetilde f(s+2L)=-\widetilde f(s)$ và chu kỳ $4L$, tương tự $g$. Vì thế chu kỳ thời gian thường là $2L/a$ hoặc $4L/a$. Với Neumann hai đầu còn phải xét trung bình $g$: trung bình khác không sinh hạng tuyến tính theo $t$, nên nghiệm không tuần hoàn.

\textbf{Bước 5 — Kiểm tra góc và đường ghép.} Biên Dirichlet không thuần nhất cần $h(0)=f(0)$, $h'(0)=g(0)$, $h''(0)=a^2f''(0)$ để có nghiệm $C^2$ cho phương trình thuần nhất. So sánh hai vết trên $x=at$, gồm $u,u_x,u_t$ rồi đạo hàm cấp hai. Nếu không tương thích, nêu rõ mức trơn còn lại thay vì khẳng định nghiệm cổ điển trên toàn miền.
\subsection*{Lời giải áp dụng cho bài}
\textbf{Bước 1. Đạo hàm công thức nửa trục.} Trong miền $x<at$, đặt $s=x+at$, $r=at-x$, $\tau=t-x/a$. Từ $u=h(\tau)+(f(s)-f(r))/2+(2a)^{-1}\int_r^s g$ thu
\[u_x=-h'(\tau)/a+\tfrac12(f'(s)+f'(r))+\frac{g(s)+g(r)}{2a}.\]
Đây là công thức (3.1.13); dấu cộng của $g(r)$ do cận dưới có đạo hàm $r_x=-1$.

\textbf{Bước 2. Đạo hàm thời gian hai lần.}
\[u_{tt}=h''(\tau)+\tfrac{a^2}2(f''(s)-f''(r))+\tfrac a2(g'(s)-g'(r)).\]
Ngoài mặt sóng, thay dấu trừ giữa hai $f''$ thành dấu cộng, và giữ cùng hiệu $g'$. Tại $x=at$, hiệu hai vết là $h''(0)-a^2f''(0)$, cho công thức ghép (3.1.15).

\textbf{Bước 3. Kết luận.} Cần $h(0)=f(0)$ để liên tục, $h'(0)=g(0)$ để ghép cấp một, và $h''(0)=a^2f''(0)$ để ghép cấp hai. Những đạo hàm tính ở đây cũng cho $u_{tt}=a^2u_{xx}$ trong miền phản xạ và đúng biên $u(0,t)=h(t)$.

\textbf{Chi tiết lần đạo hàm thứ nhất theo thời gian.} Với $s_t=r_t=a,\tau_t=1$,
\[u_t=h'(\tau)+\frac a2[f'(s)-f'(r)]+\frac12[g(s)-g(r)].\]
Lấy tiếp đạo hàm: $h'$ cho $h''$, hai hạng $f'$ cho $a^2(f''(s)-f''(r))/2$, hai hạng $g$ cho $a(g'(s)-g'(r))/2$. Tương tự, lấy đạo hàm $x$ của công thức $u_x$ cho
\[u_{xx}=h''(\tau)/a^2+(f''(s)-f''(r))/2+(g'(s)-g'(r))/(2a),\]
chứng minh trực tiếp $u_{tt}=a^2u_{xx}$.

\clearpage\noindent\texttt{\detokenize{BT_3_1_3_07}}\par
\section*{Lời giải bổ sung — Bài tập 7 · Mục 3.1.3}

\subsection*{Kiến thức và trình tự giải}
\textbf{Bước 1 — Đọc đúng điều kiện biên.} Biên $u(0,t)=0$ dùng mở rộng lẻ; biên $u_x(0,t)=0$ dùng mở rộng chẵn. Với dữ liệu trên $x\ge0$, đặt $\widetilde f(-s)=-f(s)$ hoặc $f(s)$ tương ứng, làm tương tự cho $g$. Mở rộng lẻ là $C^2$ nếu $f(0)=f''(0)=0$, và mở rộng lẻ của $g$ là $C^1$ nếu $g(0)=0$. Mở rộng chẵn đòi $f'(0)=g'(0)=0$ cho mức trơn ấy.

\textbf{Bước 2 — Lập nguyên hàm sau khi phản xạ.} Đặt $H(s)=\int_0^s\widetilde g(r)dr$. Nếu $\widetilde g$ lẻ thì $H$ chẵn; nếu $\widetilde g$ chẵn thì $H$ lẻ. Chính nguyên hàm, chứ không phải vận tốc, xuất hiện trong


\[
u=\frac{\widetilde f(x-at)+\widetilde f(x+at)}2+\frac{H(x+at)-H(x-at)}{2a}.
\]


\textbf{Bước 3 — Tách miền trước và sau khi chạm biên.} Với $u(0,t)=h(t)$, ở $x\ge at$ dùng d'Alembert với dữ liệu gốc. Ở $0<x<at$:


\[
u=h(t-x/a)+\frac{f(x+at)-f(at-x)}2+\frac1{2a}\int_{at-x}^{at+x}g(s)ds.
\]


Với $u_x(0,t)=h(t)$, dùng nghiệm của các mở rộng chẵn rồi trừ $a\int_0^{t-x/a}h(s)ds$ trên miền $x<at$. Tính riêng đạo hàm theo $x$ của hạng này để kiểm tra biên.

\textbf{Bước 4 — Trên đoạn hữu hạn, phản xạ ở cả hai đầu.} Dirichlet ở hai đầu cho mở rộng lẻ tại $0,L$, chu kỳ không gian $2L$. Neumann ở hai đầu cho mở rộng chẵn tại $0,L$, cũng chu kỳ $2L$. Một đầu lẻ, một đầu chẵn cho $\widetilde f(s+2L)=-\widetilde f(s)$ và chu kỳ $4L$, tương tự $g$. Vì thế chu kỳ thời gian thường là $2L/a$ hoặc $4L/a$. Với Neumann hai đầu còn phải xét trung bình $g$: trung bình khác không sinh hạng tuyến tính theo $t$, nên nghiệm không tuần hoàn.

\textbf{Bước 5 — Kiểm tra góc và đường ghép.} Biên Dirichlet không thuần nhất cần $h(0)=f(0)$, $h'(0)=g(0)$, $h''(0)=a^2f''(0)$ để có nghiệm $C^2$ cho phương trình thuần nhất. So sánh hai vết trên $x=at$, gồm $u,u_x,u_t$ rồi đạo hàm cấp hai. Nếu không tương thích, nêu rõ mức trơn còn lại thay vì khẳng định nghiệm cổ điển trên toàn miền.
\subsection*{Lời giải áp dụng cho bài}
\textbf{Bước 1. Mở rộng lẻ.} Đặt $\widetilde f(-s)=-f(s)$, $\widetilde g(-s)=-g(s)$. So vết tại $0$: $f(0)=0$ nối giá trị, $f'$ tự nối, $f''(0)=0$ nối cấp hai; $g(0)=0$ nối giá trị, $g'$ tự nối. Do đó $\widetilde f\in C^2$, $\widetilde g\in C^1$ dưới các điều kiện tương thích.

\textbf{Bước 2. Giải trên toàn trục.} Đặt $H(s)=\int_0^s\widetilde g(r)dr$; $H$ chẵn. D'Alembert cho $u=(\widetilde f(x-at)+\widetilde f(x+at))/2+(H(x+at)-H(x-at))/(2a)$. Tại $x=0$, hai giá trị $\widetilde f$ đối nhau và hai giá trị $H$ bằng nhau, nên $u(0,t)=0$.

\textbf{Bước 3. Viết lại bằng dữ liệu dương.} Khi $x\ge at$ giữ nguyên công thức. Khi $0<x<at$, $\widetilde f(x-at)=-f(at-x)$, còn $H(at+x)-H(at-x)=\int_{at-x}^{at+x}g$. Thu công thức phản xạ trong phần kiến thức với $h=0$. Ranh giới đúng là $x=at$; $x=t$ chỉ đúng nếu $a=1$. Dữ liệu đầu và tính $C^2$ suy trực tiếp từ mở rộng.

\clearpage\noindent\texttt{\detokenize{BT_3_1_3_08}}\par
\section*{Lời giải bổ sung — Bài tập 8 · Mục 3.1.3}

\subsection*{Kiến thức và trình tự giải}
\textbf{Bước 1 — Đọc đúng điều kiện biên.} Biên $u(0,t)=0$ dùng mở rộng lẻ; biên $u_x(0,t)=0$ dùng mở rộng chẵn. Với dữ liệu trên $x\ge0$, đặt $\widetilde f(-s)=-f(s)$ hoặc $f(s)$ tương ứng, làm tương tự cho $g$. Mở rộng lẻ là $C^2$ nếu $f(0)=f''(0)=0$, và mở rộng lẻ của $g$ là $C^1$ nếu $g(0)=0$. Mở rộng chẵn đòi $f'(0)=g'(0)=0$ cho mức trơn ấy.

\textbf{Bước 2 — Lập nguyên hàm sau khi phản xạ.} Đặt $H(s)=\int_0^s\widetilde g(r)dr$. Nếu $\widetilde g$ lẻ thì $H$ chẵn; nếu $\widetilde g$ chẵn thì $H$ lẻ. Chính nguyên hàm, chứ không phải vận tốc, xuất hiện trong


\[
u=\frac{\widetilde f(x-at)+\widetilde f(x+at)}2+\frac{H(x+at)-H(x-at)}{2a}.
\]


\textbf{Bước 3 — Tách miền trước và sau khi chạm biên.} Với $u(0,t)=h(t)$, ở $x\ge at$ dùng d'Alembert với dữ liệu gốc. Ở $0<x<at$:


\[
u=h(t-x/a)+\frac{f(x+at)-f(at-x)}2+\frac1{2a}\int_{at-x}^{at+x}g(s)ds.
\]


Với $u_x(0,t)=h(t)$, dùng nghiệm của các mở rộng chẵn rồi trừ $a\int_0^{t-x/a}h(s)ds$ trên miền $x<at$. Tính riêng đạo hàm theo $x$ của hạng này để kiểm tra biên.

\textbf{Bước 4 — Trên đoạn hữu hạn, phản xạ ở cả hai đầu.} Dirichlet ở hai đầu cho mở rộng lẻ tại $0,L$, chu kỳ không gian $2L$. Neumann ở hai đầu cho mở rộng chẵn tại $0,L$, cũng chu kỳ $2L$. Một đầu lẻ, một đầu chẵn cho $\widetilde f(s+2L)=-\widetilde f(s)$ và chu kỳ $4L$, tương tự $g$. Vì thế chu kỳ thời gian thường là $2L/a$ hoặc $4L/a$. Với Neumann hai đầu còn phải xét trung bình $g$: trung bình khác không sinh hạng tuyến tính theo $t$, nên nghiệm không tuần hoàn.

\textbf{Bước 5 — Kiểm tra góc và đường ghép.} Biên Dirichlet không thuần nhất cần $h(0)=f(0)$, $h'(0)=g(0)$, $h''(0)=a^2f''(0)$ để có nghiệm $C^2$ cho phương trình thuần nhất. So sánh hai vết trên $x=at$, gồm $u,u_x,u_t$ rồi đạo hàm cấp hai. Nếu không tương thích, nêu rõ mức trơn còn lại thay vì khẳng định nghiệm cổ điển trên toàn miền.
\subsection*{Lời giải áp dụng cho bài}
\textbf{Bước 1. Dùng công thức Dirichlet với thời điểm trễ.} Ở $x\ge at$ sóng chưa chạm biên; ở $x<at$ dùng $h(t-x/a)$ và phản xạ lẻ. Dưới đây giữ hai miền khi biên làm thay đổi vận tốc.

(i) $a=2$, $f=\sin x$, $g=\cos2x$, $h=0$. Mở rộng $g$ lẻ, $H(s)=\sin(2|s|)/2$. Do đó
\[u=\sin x\cos2t+\begin{cases}\cos2x\sin4t/4&x\ge2t,\\\sin2x\cos4t/4&0<x<2t.\end{cases}\]
$g(0)=1\ne h'(0)$ nên không $C^1$ trên mặt sóng. Sau khi mặt sóng đi qua mỗi $x$, công thức tuần hoàn chu kỳ $\pi$; toàn nghiệm không có chu kỳ chung vì hai nhánh khác nhau.

(ii) $a=3$, $f=g=0$, $h=t/(1+t)$. Đặt $\tau=t-x/3$: $u=0$ khi $\tau\le0$, $u=\tau/(1+\tau)$ khi $\tau>0$. Trên tia $x=ct$, giới hạn là $1$ nếu $0\le c<3$, và $0$ nếu $c\ge3$. Nghiệm không tuần hoàn: tại mỗi $x$ cố định, sau khi mặt sóng đến, $u$ tăng nghiêm ngặt về $1$.

(iii) $a=1$, mở rộng lẻ $f$ là $s|s|$, còn $g=\sin s$ đã lẻ:
\[u=\tfrac12[(x+t)|x+t|+(x-t)|x-t|]+\sin x\sin t.\]
Nghiệm $C^1$, không $C^2$ trên $x=t$ vì $f''(0)=2\ne0$, và không tuần hoàn.

(iv) $a=5$, $f=0$, $g=\sin x-\sin3x/3$ có $H=-\cos s+\cos3s/9$. Suy $u=\sin x\sin5t/5-\sin3x\sin15t/45$, trơn, chu kỳ $2\pi/5$.

(v) $a=2$, $f=\cos x-\sin x$, $g=0$, $h=\cos2t-\sin2t$. Ngoài mặt sóng, half-sum của $f$ là $(\cos x-\sin x)\cos2t$. Trong mặt sóng:
\[h(t-x/2)+\tfrac12[f(x+2t)-f(2t-x)]=\cos x(\cos2t-\sin2t).\]
Vậy
\[u=\begin{cases}(\cos x-\sin x)\cos2t&x\ge2t,\\\cos x(\cos2t-\sin2t)&0<x<2t.\end{cases}\]
Công thức cũ $\cos(x-2t)-\sin(x+2t)$ không thỏa $u_t(x,0)=0$. Nghiệm đúng liên tục, không $C^1$ vì $h'(0)=-2\ne0$. Chỉ nhánh sau mặt sóng có chu kỳ $\pi$ ở từng điểm.

(vi) $a=1$, $g=\chi_{[1,2]}-\chi_{[3,4]}$. Mở rộng lẻ có nguyên hàm chẵn $H(s)=[|s|-1]_0^1-[|s|-3]_0^1$, nên $u=(H(x+t)-H(x-t))/2$. Tại $x$ cố định, hai giá trị $H$ đều bằng không sau $t>x+4$, nên $u=0$; đây là xung quá độ, không tuần hoàn. Nghiệm liên tục nhưng có gãy cấp một khi các cận gặp $\pm1,\pm2,\pm3,\pm4$.

\textbf{Bước 2. Kiểm tra chung.} Thế $x=0$ vào nhánh trong cho đúng $h$; thế $t=0$ vào nhánh ngoài cho $f,g$. So hai vết trên $x=at$ theo ba điều kiện tương thích, không dùng tính tuần hoàn của nhánh trong để kết luận tuần hoàn toàn miền.

\textbf{Chi tiết hai nhánh vận tốc ở (i).} Ngoài mặt sóng,
\[\frac14\int_{x-2t}^{x+2t}\cos2s\,ds
=\frac{\sin(2x+4t)-\sin(2x-4t)}8=\frac14\cos2x\sin4t.\]
Trong mặt sóng, cận dưới là $2t-x$, nên hiệu sine bằng $\sin(4t+2x)-\sin(4t-2x)=2\cos4t\sin2x$; chia $8$ cho nhánh đã ghi.

\textbf{Chi tiết rút gọn (v).} Trong miền phản xạ, biên là $\cos(2t-x)-\sin(2t-x)$. Half-difference của $f$ bằng
\[-\sin2t\sin x-\cos2t\sin x.\]
Khai triển biên thành $\cos2t\cos x+\sin2t\sin x-\sin2t\cos x+\cos2t\sin x$ rồi cộng, hai hạng chứa $\sin x$ triệt tiêu, còn $\cos x(\cos2t-\sin2t)$.
\begin{center}\includegraphics[width=.97\linewidth,height=11cm,keepaspectratio]{_xu_ly/graphs/BT_3_1_3_08.png}\end{center}
\textit{Đồ thị minh họa; với nghiệm số, xem phương pháp và tham số trong mã Python kèm theo.}
\begin{center}\includegraphics[width=.97\linewidth,height=11cm,keepaspectratio]{_xu_ly/graphs/BT_3_1_3_08_characteristics.png}\end{center}

\clearpage\noindent\texttt{\detokenize{BT_3_1_3_09}}\par
\section*{Lời giải bổ sung — Bài tập 9 · Mục 3.1.3}

\subsection*{Kiến thức và trình tự giải}
\textbf{Bước 1 — Đọc đúng điều kiện biên.} Biên $u(0,t)=0$ dùng mở rộng lẻ; biên $u_x(0,t)=0$ dùng mở rộng chẵn. Với dữ liệu trên $x\ge0$, đặt $\widetilde f(-s)=-f(s)$ hoặc $f(s)$ tương ứng, làm tương tự cho $g$. Mở rộng lẻ là $C^2$ nếu $f(0)=f''(0)=0$, và mở rộng lẻ của $g$ là $C^1$ nếu $g(0)=0$. Mở rộng chẵn đòi $f'(0)=g'(0)=0$ cho mức trơn ấy.

\textbf{Bước 2 — Lập nguyên hàm sau khi phản xạ.} Đặt $H(s)=\int_0^s\widetilde g(r)dr$. Nếu $\widetilde g$ lẻ thì $H$ chẵn; nếu $\widetilde g$ chẵn thì $H$ lẻ. Chính nguyên hàm, chứ không phải vận tốc, xuất hiện trong


\[
u=\frac{\widetilde f(x-at)+\widetilde f(x+at)}2+\frac{H(x+at)-H(x-at)}{2a}.
\]


\textbf{Bước 3 — Tách miền trước và sau khi chạm biên.} Với $u(0,t)=h(t)$, ở $x\ge at$ dùng d'Alembert với dữ liệu gốc. Ở $0<x<at$:


\[
u=h(t-x/a)+\frac{f(x+at)-f(at-x)}2+\frac1{2a}\int_{at-x}^{at+x}g(s)ds.
\]


Với $u_x(0,t)=h(t)$, dùng nghiệm của các mở rộng chẵn rồi trừ $a\int_0^{t-x/a}h(s)ds$ trên miền $x<at$. Tính riêng đạo hàm theo $x$ của hạng này để kiểm tra biên.

\textbf{Bước 4 — Trên đoạn hữu hạn, phản xạ ở cả hai đầu.} Dirichlet ở hai đầu cho mở rộng lẻ tại $0,L$, chu kỳ không gian $2L$. Neumann ở hai đầu cho mở rộng chẵn tại $0,L$, cũng chu kỳ $2L$. Một đầu lẻ, một đầu chẵn cho $\widetilde f(s+2L)=-\widetilde f(s)$ và chu kỳ $4L$, tương tự $g$. Vì thế chu kỳ thời gian thường là $2L/a$ hoặc $4L/a$. Với Neumann hai đầu còn phải xét trung bình $g$: trung bình khác không sinh hạng tuyến tính theo $t$, nên nghiệm không tuần hoàn.

\textbf{Bước 5 — Kiểm tra góc và đường ghép.} Biên Dirichlet không thuần nhất cần $h(0)=f(0)$, $h'(0)=g(0)$, $h''(0)=a^2f''(0)$ để có nghiệm $C^2$ cho phương trình thuần nhất. So sánh hai vết trên $x=at$, gồm $u,u_x,u_t$ rồi đạo hàm cấp hai. Nếu không tương thích, nêu rõ mức trơn còn lại thay vì khẳng định nghiệm cổ điển trên toàn miền.
\subsection*{Lời giải áp dụng cho bài}
\textbf{Bước 1. Tìm các hàm sóng ở hai miền.} Biên nghiêng là $t=\alpha x$, không phải $x=0$. Viết $u=p+A(x+at)+B(x-at)$, trong đó $p$ là nghiệm riêng. Dữ liệu ở $t=0$ xác định $A,B$ trên đối số dương; dữ liệu biên xác định $B$ ở đối số âm.

(i) $a=1$, $\alpha=2$, $F=0$, $f=x$, $g=0$, vết biên $3x$. Trên $s>0$, $A(s)=B(s)=s/2$ (bỏ hai hằng số triệt tiêu). Tại $t=2x$, $A(3x)+B(-x)=3x$, nên $B(-x)=3x/2$, tức $B(s)=-3s/2$ khi $s<0$. Suy
\[u=\begin{cases}x&0\le t\le x,\\2t-x&x<t<2x.\end{cases}\]
Hai vết giá trị bằng nhau, nhưng $u_t$ đổi từ $0$ sang $2$; không có nghiệm $C^1$ trên toàn tam giác thỏa mọi dữ liệu này.

(ii) $a=2$, $\alpha=3$, $F=1$, $f=g=0$, biên bằng không. Chọn $p=-x^2/8$, vì $p_{tt}-4p_{xx}=1$. Dữ liệu đầu cho $A=B=s^2/16$ trên $s>0$. Tại $t=3x$, $A(7x)+B(-5x)-x^2/8=0$, suy $B(s)=-47s^2/400$ khi $s<0$:
\[u=\begin{cases}t^2/2&0\le t\le x/2,\\-x^2/8+(x+2t)^2/16-47(x-2t)^2/400&x/2<t<3x.\end{cases}\]
\textbf{Bước 2. Kiểm tra.} Tại đường $x=2t$, $B$ và $B'$ nối bằng không nên nghiệm $C^1$, nhưng $B''$ đổi từ $1/8$ sang $-47/200$, nên không $C^2$. Mỗi nhánh thỏa PDE; tại $t=3x$ đa thức triệt tiêu đúng. Hai kết quả phải hiểu từng miền khi đề yêu cầu mức trơn cao hơn mức dữ liệu cho phép.

\clearpage\noindent\texttt{\detokenize{BT_3_1_3_10}}\par
\section*{Lời giải bổ sung — Bài tập 10 · Mục 3.1.3}

\subsection*{Kiến thức và trình tự giải}
\textbf{Bước 1 — Đọc đúng điều kiện biên.} Biên $u(0,t)=0$ dùng mở rộng lẻ; biên $u_x(0,t)=0$ dùng mở rộng chẵn. Với dữ liệu trên $x\ge0$, đặt $\widetilde f(-s)=-f(s)$ hoặc $f(s)$ tương ứng, làm tương tự cho $g$. Mở rộng lẻ là $C^2$ nếu $f(0)=f''(0)=0$, và mở rộng lẻ của $g$ là $C^1$ nếu $g(0)=0$. Mở rộng chẵn đòi $f'(0)=g'(0)=0$ cho mức trơn ấy.

\textbf{Bước 2 — Lập nguyên hàm sau khi phản xạ.} Đặt $H(s)=\int_0^s\widetilde g(r)dr$. Nếu $\widetilde g$ lẻ thì $H$ chẵn; nếu $\widetilde g$ chẵn thì $H$ lẻ. Chính nguyên hàm, chứ không phải vận tốc, xuất hiện trong


\[
u=\frac{\widetilde f(x-at)+\widetilde f(x+at)}2+\frac{H(x+at)-H(x-at)}{2a}.
\]


\textbf{Bước 3 — Tách miền trước và sau khi chạm biên.} Với $u(0,t)=h(t)$, ở $x\ge at$ dùng d'Alembert với dữ liệu gốc. Ở $0<x<at$:


\[
u=h(t-x/a)+\frac{f(x+at)-f(at-x)}2+\frac1{2a}\int_{at-x}^{at+x}g(s)ds.
\]


Với $u_x(0,t)=h(t)$, dùng nghiệm của các mở rộng chẵn rồi trừ $a\int_0^{t-x/a}h(s)ds$ trên miền $x<at$. Tính riêng đạo hàm theo $x$ của hạng này để kiểm tra biên.

\textbf{Bước 4 — Trên đoạn hữu hạn, phản xạ ở cả hai đầu.} Dirichlet ở hai đầu cho mở rộng lẻ tại $0,L$, chu kỳ không gian $2L$. Neumann ở hai đầu cho mở rộng chẵn tại $0,L$, cũng chu kỳ $2L$. Một đầu lẻ, một đầu chẵn cho $\widetilde f(s+2L)=-\widetilde f(s)$ và chu kỳ $4L$, tương tự $g$. Vì thế chu kỳ thời gian thường là $2L/a$ hoặc $4L/a$. Với Neumann hai đầu còn phải xét trung bình $g$: trung bình khác không sinh hạng tuyến tính theo $t$, nên nghiệm không tuần hoàn.

\textbf{Bước 5 — Kiểm tra góc và đường ghép.} Biên Dirichlet không thuần nhất cần $h(0)=f(0)$, $h'(0)=g(0)$, $h''(0)=a^2f''(0)$ để có nghiệm $C^2$ cho phương trình thuần nhất. So sánh hai vết trên $x=at$, gồm $u,u_x,u_t$ rồi đạo hàm cấp hai. Nếu không tương thích, nêu rõ mức trơn còn lại thay vì khẳng định nghiệm cổ điển trên toàn miền.
\subsection*{Lời giải áp dụng cho bài}
\textbf{Bước 1. Điều kiện Neumann.} $u_x(0,t)=0$ kéo theo $f'(0)=0$, $g'(0)=0$ nếu nghiệm $C^2$. Mở rộng chẵn của $f,g$ có vết $f'$ và $g'$ đổi dấu, nên hai điều kiện này bảo đảm các vết cần thiết nối nhau.

\textbf{Bước 2. Áp dụng d'Alembert.} Với mở rộng chẵn $\widetilde f,\widetilde g$, nguyên hàm $H=\int_0^s\widetilde g$ là lẻ. Nghiệm là half-sum cộng hiệu $H/(2a)$. Nó chẵn theo $x$, do đó đạo hàm không gian ở $0$ bằng không. Trong miền $x<at$ có thể viết trực tiếp
\[u=\frac{f(at-x)+f(at+x)}2+\frac1a\int_0^{at-x}g(s)ds+\frac1{2a}\int_{at-x}^{at+x}g(s)ds.\]
Hạng tích phân đầu xuất hiện vì đoạn $[x-at,0]$ cũng góp vận tốc sau phản xạ chẵn.

\textbf{Bước 3. Ví dụ.} Với $a=2$, $f=x^2$, $g=\cos x$, cả hai đã chẵn và tương thích. Half-sum cho $x^2+4t^2$, nguyên hàm $\sin s$ cho $\cos x\sin2t/2$. Vậy $u=x^2+4t^2+\cos x\sin2t/2$. Kiểm tra $u_x(0,t)=0$, $u_t(x,0)=\cos x$ và $u_{tt}=4u_{xx}$ trực tiếp.

\clearpage\noindent\texttt{\detokenize{BT_3_1_3_11}}\par
\section*{Lời giải bổ sung — Bài tập 11 · Mục 3.1.3}

\subsection*{Kiến thức và trình tự giải}
\textbf{Bước 1 — Đọc đúng điều kiện biên.} Biên $u(0,t)=0$ dùng mở rộng lẻ; biên $u_x(0,t)=0$ dùng mở rộng chẵn. Với dữ liệu trên $x\ge0$, đặt $\widetilde f(-s)=-f(s)$ hoặc $f(s)$ tương ứng, làm tương tự cho $g$. Mở rộng lẻ là $C^2$ nếu $f(0)=f''(0)=0$, và mở rộng lẻ của $g$ là $C^1$ nếu $g(0)=0$. Mở rộng chẵn đòi $f'(0)=g'(0)=0$ cho mức trơn ấy.

\textbf{Bước 2 — Lập nguyên hàm sau khi phản xạ.} Đặt $H(s)=\int_0^s\widetilde g(r)dr$. Nếu $\widetilde g$ lẻ thì $H$ chẵn; nếu $\widetilde g$ chẵn thì $H$ lẻ. Chính nguyên hàm, chứ không phải vận tốc, xuất hiện trong


\[
u=\frac{\widetilde f(x-at)+\widetilde f(x+at)}2+\frac{H(x+at)-H(x-at)}{2a}.
\]


\textbf{Bước 3 — Tách miền trước và sau khi chạm biên.} Với $u(0,t)=h(t)$, ở $x\ge at$ dùng d'Alembert với dữ liệu gốc. Ở $0<x<at$:


\[
u=h(t-x/a)+\frac{f(x+at)-f(at-x)}2+\frac1{2a}\int_{at-x}^{at+x}g(s)ds.
\]


Với $u_x(0,t)=h(t)$, dùng nghiệm của các mở rộng chẵn rồi trừ $a\int_0^{t-x/a}h(s)ds$ trên miền $x<at$. Tính riêng đạo hàm theo $x$ của hạng này để kiểm tra biên.

\textbf{Bước 4 — Trên đoạn hữu hạn, phản xạ ở cả hai đầu.} Dirichlet ở hai đầu cho mở rộng lẻ tại $0,L$, chu kỳ không gian $2L$. Neumann ở hai đầu cho mở rộng chẵn tại $0,L$, cũng chu kỳ $2L$. Một đầu lẻ, một đầu chẵn cho $\widetilde f(s+2L)=-\widetilde f(s)$ và chu kỳ $4L$, tương tự $g$. Vì thế chu kỳ thời gian thường là $2L/a$ hoặc $4L/a$. Với Neumann hai đầu còn phải xét trung bình $g$: trung bình khác không sinh hạng tuyến tính theo $t$, nên nghiệm không tuần hoàn.

\textbf{Bước 5 — Kiểm tra góc và đường ghép.} Biên Dirichlet không thuần nhất cần $h(0)=f(0)$, $h'(0)=g(0)$, $h''(0)=a^2f''(0)$ để có nghiệm $C^2$ cho phương trình thuần nhất. So sánh hai vết trên $x=at$, gồm $u,u_x,u_t$ rồi đạo hàm cấp hai. Nếu không tương thích, nêu rõ mức trơn còn lại thay vì khẳng định nghiệm cổ điển trên toàn miền.
\subsection*{Lời giải áp dụng cho bài}
\textbf{Bước 1. Chuyển đạo hàm thành bài Dirichlet.} Đặt $v=u_x$. Khi dữ liệu đủ trơn, $v_{tt}=a^2v_{xx}$, $v(0,t)=h(t)$, $v(x,0)=f'(x)$, $v_t(x,0)=g'(x)$. Tương thích cấp hai cho bài này là $h(0)=f'(0)$, $h'(0)=g'(0)$, $h''(0)=a^2f'''(0)$.

\textbf{Bước 2. Tích phân lại và giữ hằng số theo thời gian.} $u(x,t)=C(t)+\int_0^xv(s,t)ds$. Thế vào PDE cho $C''(t)=a^2v_x(0,t)$, $C(0)=f(0)$, $C'(0)=g(0)$. Không được bỏ $C$ chỉ vì đã biết $u_x$.

\textbf{Bước 3. Công thức tiện dụng tương đương.} Gọi $u_e$ là nghiệm d'Alembert của các mở rộng chẵn $f,g$. Đặt
\[u=u_e-a\int_0^{(t-x/a)_+}h(s)ds.\]
Ở $x<at$, đạo hàm không gian của hạng cuối là $h(t-x/a)$, cho đúng $u_x(0,t)=h(t)$; hạng đó là sóng chạy nên thỏa PDE. Nếu $h(0)\ne0$, góc của $u_e$ do $f'(0)$ được hạng biên bù khi tương thích đúng.

\textbf{Bước 4. Viết hai nhánh.} Ngoài mặt sóng dùng d'Alembert với dữ liệu dương. Trong mặt sóng dùng công thức Bài 10 rồi trừ $a\int_0^{t-x/a}h$. Hai dữ liệu đầu giữ nguyên vì hạng bổ sung chưa đến $x>0$ tại $t=0$. Nếu chỉ cần $u\in C^2$, hai tương thích đầu đủ với mức trơn phù hợp; điều kiện thứ ba ở trên phục vụ $v\in C^2$, tức mức trơn cao hơn của $u$.

\clearpage\noindent\texttt{\detokenize{BT_3_1_3_12}}\par
\section*{Lời giải bổ sung — Bài tập 12 · Mục 3.1.3}

\subsection*{Kiến thức và trình tự giải}
\textbf{Bước 1 — Đọc đúng điều kiện biên.} Biên $u(0,t)=0$ dùng mở rộng lẻ; biên $u_x(0,t)=0$ dùng mở rộng chẵn. Với dữ liệu trên $x\ge0$, đặt $\widetilde f(-s)=-f(s)$ hoặc $f(s)$ tương ứng, làm tương tự cho $g$. Mở rộng lẻ là $C^2$ nếu $f(0)=f''(0)=0$, và mở rộng lẻ của $g$ là $C^1$ nếu $g(0)=0$. Mở rộng chẵn đòi $f'(0)=g'(0)=0$ cho mức trơn ấy.

\textbf{Bước 2 — Lập nguyên hàm sau khi phản xạ.} Đặt $H(s)=\int_0^s\widetilde g(r)dr$. Nếu $\widetilde g$ lẻ thì $H$ chẵn; nếu $\widetilde g$ chẵn thì $H$ lẻ. Chính nguyên hàm, chứ không phải vận tốc, xuất hiện trong


\[
u=\frac{\widetilde f(x-at)+\widetilde f(x+at)}2+\frac{H(x+at)-H(x-at)}{2a}.
\]


\textbf{Bước 3 — Tách miền trước và sau khi chạm biên.} Với $u(0,t)=h(t)$, ở $x\ge at$ dùng d'Alembert với dữ liệu gốc. Ở $0<x<at$:


\[
u=h(t-x/a)+\frac{f(x+at)-f(at-x)}2+\frac1{2a}\int_{at-x}^{at+x}g(s)ds.
\]


Với $u_x(0,t)=h(t)$, dùng nghiệm của các mở rộng chẵn rồi trừ $a\int_0^{t-x/a}h(s)ds$ trên miền $x<at$. Tính riêng đạo hàm theo $x$ của hạng này để kiểm tra biên.

\textbf{Bước 4 — Trên đoạn hữu hạn, phản xạ ở cả hai đầu.} Dirichlet ở hai đầu cho mở rộng lẻ tại $0,L$, chu kỳ không gian $2L$. Neumann ở hai đầu cho mở rộng chẵn tại $0,L$, cũng chu kỳ $2L$. Một đầu lẻ, một đầu chẵn cho $\widetilde f(s+2L)=-\widetilde f(s)$ và chu kỳ $4L$, tương tự $g$. Vì thế chu kỳ thời gian thường là $2L/a$ hoặc $4L/a$. Với Neumann hai đầu còn phải xét trung bình $g$: trung bình khác không sinh hạng tuyến tính theo $t$, nên nghiệm không tuần hoàn.

\textbf{Bước 5 — Kiểm tra góc và đường ghép.} Biên Dirichlet không thuần nhất cần $h(0)=f(0)$, $h'(0)=g(0)$, $h''(0)=a^2f''(0)$ để có nghiệm $C^2$ cho phương trình thuần nhất. So sánh hai vết trên $x=at$, gồm $u,u_x,u_t$ rồi đạo hàm cấp hai. Nếu không tương thích, nêu rõ mức trơn còn lại thay vì khẳng định nghiệm cổ điển trên toàn miền.
\subsection*{Lời giải áp dụng cho bài}
\textbf{Bước 1. Tách hai sóng trên dữ liệu gốc.} Đặt $H(s)=\int_0^sg$, $A=f/2-H/(2a)$, $B=f/2+H/(2a)$ trên $0\le s\le L$. Hai biên là $u(0,t)=\alpha(t)$ và $u(L,t)=\beta(t)$. Chỉ xét $t\le L/a$, nên mỗi chân đặc trưng phản xạ tối đa một lần.

\textbf{Bước 2. Viết đủ bốn miền.} Đặt $\xi=x-at$, $\eta=x+at$:
\[u=\begin{cases}
A(\xi)+B(\eta)&\xi\ge0,\ \eta\le L,\\
\alpha(t-x/a)-B(at-x)+B(x+at)&\xi<0,\ \eta\le L,\\
A(x-at)+\beta(t-(L-x)/a)-A(2L-x-at)&\xi\ge0,\ \eta>L,\\
\alpha(t-x/a)+\beta(t-(L-x)/a)-B(at-x)-A(2L-x-at)&\xi<0,\ \eta>L.
\end{cases}\]
Hai công thức phản xạ suy từ $A(-s)=\alpha(s/a)-B(s)$ và $B(L+s)=\beta(s/a)-A(L-s)$; mọi đối số của $A,B$ còn lại đều trong $[0,L]$.

\textbf{Bước 3. Trường hợp (i).} $a=2$, $L=\pi$, $f=\sin x$, $g=1$, $\alpha=\beta=0$. $A=\sin s/2-s/4$, $B=\sin s/2+s/4$. Trong bốn miền theo đúng thứ tự trên:
\[u=\sin x\cos2t+\begin{cases}t&\text{miền I},\\x/2&\text{miền II},\\(\pi-x)/2&\text{miền III},\\\pi/2-t&\text{miền IV}.\end{cases}\]
Phần vận tốc liên tục, nhưng có góc trên hai mặt sóng vì $g$ không tương thích với biên không.

\textbf{Bước 4. Trường hợp (ii).} $a=L=1$, $f=g=0$, $\alpha(t)=\sin\pi t$, $\beta(t)=t$. Vì $A=B=0$, bốn nhánh là $0$, $\sin\pi(t-x)$, $t-1+x$, và $\sin\pi(t-x)+t-1+x$. Thế $x=0,1$ trong khoảng $0\le t\le1$ cho đúng hai biên. Vận tốc góc không tương thích ($\alpha'(0)=\pi$, $\beta'(0)=1$), nên không có nghiệm $C^1$ trên toàn miền; các nhánh và đường ghép đã mô tả chính xác sự mất trơn.

\textbf{Chi tiết thế vào miền cả hai mặt sóng đã phản xạ ở (i).} Đặt $p=2t-x$, $q=2\pi-x-2t$. Khi $\xi<0,\eta>\pi$, $p,q\in[0,\pi]$ và $u=-B(p)-A(q)$:
\[u=-\frac{\sin(2t-x)}2-\frac{\sin(2\pi-x-2t)}2+\frac{q-p}{4}
=\frac{\sin(x+2t)+\sin(x-2t)}2+\frac{2\pi-4t}{4}.\]
Tổng sine bằng $2\sin x\cos2t$, hạng tuyến tính bằng $\pi/2-t$. Ở miền chỉ phản xạ trái, $-B(2t-x)+B(x+2t)$ cho hạng tuyến tính $[(x+2t)-(2t-x)]/4=x/2$. Ở miền chỉ phản xạ phải, $A(x-2t)-A(2\pi-x-2t)$ cho $[-(x-2t)+(2\pi-x-2t)]/4=(\pi-x)/2$.

\clearpage\noindent\texttt{\detokenize{VD_3_2_2_01}}\par
\section*{Lời giải bổ sung — Ví dụ 1 · Mục 3.2.2}

\subsection*{Kiến thức và trình tự giải}
\textbf{Bước 1 — Liệt kê các khoảng của dữ liệu.} Ghi đủ điểm đầu, điểm cuối hỗ trợ, điểm gãy và bước nhảy của $f,g$. Nếu có biên, phản xạ chẵn hoặc lẻ trước khi tích phân; trên đoạn hữu hạn tiếp tục phản xạ đến một chu kỳ hoàn chỉnh.

\textbf{Bước 2 — Tính nguyên hàm liên tục.} Lấy $H'=\widetilde g$ trên từng khoảng và chọn hằng số để hai vết của $H$ bằng nhau. Viết $[s]_p^q=\min(q,\max(p,s))$ hoặc $(s)_+=\max(s,0)$ khi cần gộp nhiều nhánh. Một nguyên hàm liên tục tránh tạo ra bước nhảy giả trong nghiệm.

\textbf{Bước 3 — Xác định sóng phải và sóng trái.} Đặt $A(s)=\widetilde f(s)/2-H(s)/(2a)$, $B(s)=\widetilde f(s)/2+H(s)/(2a)$; khi đó $u=A(x-at)+B(x+at)$. Với từng $t$ được yêu cầu, dịch đồ thị $A$ sang phải $at$ và $B$ sang trái $at$, rồi cộng tung độ. Các đường $x-at=b$ và $x+at=b$ chia mặt phẳng thành các miền có công thức đơn giản.

\textbf{Bước 4 — Xét liên tục, đạo hàm và giao của các mặt sóng.} Nếu $f$ có bước nhảy, $u$ thường có bước nhảy; nếu $f$ liên tục và $g$ bị chặn thì $u$ liên tục. Tại thời điểm cố định,


\[
u_x=\frac{\widetilde f'(x-at)+\widetilde f'(x+at)}2+\frac{\widetilde g(x+at)-\widetilde g(x-at)}{2a}.
\]


Tính giới hạn trái/phải của biểu thức này tại từng vị trí ứng viên. Có thể có triệt tiêu khi hai đường gặp nhau, nên không chỉ liệt kê đường rồi khẳng định mọi đạo hàm đều mất. Với tính $C^1$ theo cả $x,t$, còn kiểm tra $u_t$. Với $C^2$, xét tiếp các góc của $g$ và đạo hàm cấp hai của $f$.

\textbf{Bước 5 — Tính thời điểm đến, giới hạn và ngưỡng.} Sóng ở $(x,t)$ chỉ nhận dữ liệu trong $[x-at,x+at]$. Giải bất đẳng thức giao hỗ trợ để tìm thời điểm đến và đi qua. Nếu $H$ hằng ngoài đoạn hỗ trợ, dùng các giá trị hằng này để tính giới hạn khi $t\to\infty$. Nếu hỏi áp suất cực đại, viết hàm từng khoảng thời gian và so sánh cực trị, bước nhảy cùng hai giá trị một phía.
\subsection*{Lời giải áp dụng cho bài}
\textbf{Bước 1. Chia dữ liệu và hai sóng.} $a=1$, $f(s)=s(s+1)$ trên $[-1,0]$, bằng $0$ ngoài đoạn, $g=0$. Vì vậy $A=B=f/2$, không cần nguyên hàm vận tốc.

\textbf{Bước 2. Dịch từng đoạn.} Sóng phải $A(x-t)$ khác không khi $t-1<x<t$, sóng trái $B(x+t)$ khác không khi $-t-1<x<-t$. Nghiệm chính xác là
\[u(x,t)=\tfrac12(x-t)(x-t+1)\chi_{[-1,0]}(x-t)+\tfrac12(x+t)(x+t+1)\chi_{[-1,0]}(x+t).\]
Khi hai đoạn chồng nhau, cộng hai đa thức; khi chỉ một đoạn hiện diện lấy một hạng; ngoài cả hai đoạn lấy $0$. Đây cũng là cách dựng các hình trong sách.

\textbf{Bước 3. Xét tính trơn.} $f$ bằng $0$ tại hai đầu nên $u$ liên tục. $f'$ trong đoạn bằng $2s+1$, có bước nhảy tại $-1,0$, nên các đường gãy cấp một là $x\pm t=-1,0$. Ở giao điểm $(-1/2,1/2)$, lát cắt $u(x,1/2)$ có đạo hàm trái theo $x$ bằng $1/2$, đạo hàm phải bằng $-1/2$; vì vậy điểm ấy vẫn không $C^1$. Tại giao của các mặt sóng phải so trực tiếp các đạo hàm một phía, không mặc nhiên giả sử chúng triệt tiêu.

\textbf{Bước 4. Kiểm tra.} Thế $t=0$ cho $u=f$, còn đạo hàm theo $t$ của hai sóng đối nhau nên vận tốc đầu bằng $0$ nơi $f'$ tồn tại. Mỗi đa thức của $x\pm t$ thỏa PDE; toàn nghiệm đúng theo nghĩa yếu trên các đường gãy.

\textbf{Chi tiết đạo hàm tại giao điểm hai mặt sóng.} Với $t=1/2$, bên trái $x=-1/2$, chỉ $x+t=x+1/2\in[-1,0]$ góp:
\[u(x,1/2)=\tfrac12(x+1/2)(x+3/2),\quad u_x=x+1\to1/2.\]
Bên phải, chỉ $x-t=x-1/2\in[-1,0]$ góp:
\[u(x,1/2)=\tfrac12(x-1/2)(x+1/2),\quad u_x=x\to-1/2.\]
Hai vết giá trị đều $0$, nhưng hai đạo hàm không bằng nhau, nên không thể gọi điểm giao là điểm trơn.

\clearpage\noindent\texttt{\detokenize{VD_3_2_2_02}}\par
\section*{Lời giải bổ sung — Ví dụ 2 · Mục 3.2.2}

\subsection*{Kiến thức và trình tự giải}
\textbf{Bước 1 — Liệt kê các khoảng của dữ liệu.} Ghi đủ điểm đầu, điểm cuối hỗ trợ, điểm gãy và bước nhảy của $f,g$. Nếu có biên, phản xạ chẵn hoặc lẻ trước khi tích phân; trên đoạn hữu hạn tiếp tục phản xạ đến một chu kỳ hoàn chỉnh.

\textbf{Bước 2 — Tính nguyên hàm liên tục.} Lấy $H'=\widetilde g$ trên từng khoảng và chọn hằng số để hai vết của $H$ bằng nhau. Viết $[s]_p^q=\min(q,\max(p,s))$ hoặc $(s)_+=\max(s,0)$ khi cần gộp nhiều nhánh. Một nguyên hàm liên tục tránh tạo ra bước nhảy giả trong nghiệm.

\textbf{Bước 3 — Xác định sóng phải và sóng trái.} Đặt $A(s)=\widetilde f(s)/2-H(s)/(2a)$, $B(s)=\widetilde f(s)/2+H(s)/(2a)$; khi đó $u=A(x-at)+B(x+at)$. Với từng $t$ được yêu cầu, dịch đồ thị $A$ sang phải $at$ và $B$ sang trái $at$, rồi cộng tung độ. Các đường $x-at=b$ và $x+at=b$ chia mặt phẳng thành các miền có công thức đơn giản.

\textbf{Bước 4 — Xét liên tục, đạo hàm và giao của các mặt sóng.} Nếu $f$ có bước nhảy, $u$ thường có bước nhảy; nếu $f$ liên tục và $g$ bị chặn thì $u$ liên tục. Tại thời điểm cố định,


\[
u_x=\frac{\widetilde f'(x-at)+\widetilde f'(x+at)}2+\frac{\widetilde g(x+at)-\widetilde g(x-at)}{2a}.
\]


Tính giới hạn trái/phải của biểu thức này tại từng vị trí ứng viên. Có thể có triệt tiêu khi hai đường gặp nhau, nên không chỉ liệt kê đường rồi khẳng định mọi đạo hàm đều mất. Với tính $C^1$ theo cả $x,t$, còn kiểm tra $u_t$. Với $C^2$, xét tiếp các góc của $g$ và đạo hàm cấp hai của $f$.

\textbf{Bước 5 — Tính thời điểm đến, giới hạn và ngưỡng.} Sóng ở $(x,t)$ chỉ nhận dữ liệu trong $[x-at,x+at]$. Giải bất đẳng thức giao hỗ trợ để tìm thời điểm đến và đi qua. Nếu $H$ hằng ngoài đoạn hỗ trợ, dùng các giá trị hằng này để tính giới hạn khi $t\to\infty$. Nếu hỏi áp suất cực đại, viết hàm từng khoảng thời gian và so sánh cực trị, bước nhảy cùng hai giá trị một phía.
\subsection*{Lời giải áp dụng cho bài}
\textbf{Bước 1. Tích phân vận tốc.} $f=0$, $g=\chi_{[-1,0]}$, $a=1$. Chọn
\[H(s)=\begin{cases}0&s\le-1,\\s+1&-1<s<0,\\1&s\ge0.\end{cases}\]
Hai hằng số đuôi đã nối liên tục với nhánh giữa.

\textbf{Bước 2. Tách và ghép sóng.} $A=-H/2$, $B=H/2$, nên $u=(H(x+t)-H(x-t))/2$. Tương đương
\[u=\tfrac12\big[\min(x+t,0)-\max(x-t,-1)\big]_+.\]
Biểu thức trong ngoặc là độ dài giao của khoảng ảnh hưởng với hỗ trợ vận tốc. Trên mỗi miền mà các min/max chọn cùng nhánh, $u$ là affine theo $x,t$.

\textbf{Bước 3. Miền gãy và giá trị lâu dài.} $H$ liên tục nên $u$ liên tục. $H'$ nhảy tại $-1,0$, phát sinh các đường $x\pm t=-1,0$ mất trơn cấp một. Với mỗi $x$ cố định, khi $t>\max(|x|,|x+1|)$, toàn đoạn $[-1,0]$ nằm trong khoảng ảnh hưởng, nên $u=1/2$. Kiểm tra tại $t=0$ có hiệu $H=0$ và $u_t=g$ tại các điểm liên tục của dữ liệu.

\clearpage\noindent\texttt{\detokenize{VD_3_2_2_03}}\par
\section*{Lời giải bổ sung — Ví dụ 3 · Mục 3.2.2}

\subsection*{Kiến thức và trình tự giải}
\textbf{Bước 1 — Liệt kê các khoảng của dữ liệu.} Ghi đủ điểm đầu, điểm cuối hỗ trợ, điểm gãy và bước nhảy của $f,g$. Nếu có biên, phản xạ chẵn hoặc lẻ trước khi tích phân; trên đoạn hữu hạn tiếp tục phản xạ đến một chu kỳ hoàn chỉnh.

\textbf{Bước 2 — Tính nguyên hàm liên tục.} Lấy $H'=\widetilde g$ trên từng khoảng và chọn hằng số để hai vết của $H$ bằng nhau. Viết $[s]_p^q=\min(q,\max(p,s))$ hoặc $(s)_+=\max(s,0)$ khi cần gộp nhiều nhánh. Một nguyên hàm liên tục tránh tạo ra bước nhảy giả trong nghiệm.

\textbf{Bước 3 — Xác định sóng phải và sóng trái.} Đặt $A(s)=\widetilde f(s)/2-H(s)/(2a)$, $B(s)=\widetilde f(s)/2+H(s)/(2a)$; khi đó $u=A(x-at)+B(x+at)$. Với từng $t$ được yêu cầu, dịch đồ thị $A$ sang phải $at$ và $B$ sang trái $at$, rồi cộng tung độ. Các đường $x-at=b$ và $x+at=b$ chia mặt phẳng thành các miền có công thức đơn giản.

\textbf{Bước 4 — Xét liên tục, đạo hàm và giao của các mặt sóng.} Nếu $f$ có bước nhảy, $u$ thường có bước nhảy; nếu $f$ liên tục và $g$ bị chặn thì $u$ liên tục. Tại thời điểm cố định,


\[
u_x=\frac{\widetilde f'(x-at)+\widetilde f'(x+at)}2+\frac{\widetilde g(x+at)-\widetilde g(x-at)}{2a}.
\]


Tính giới hạn trái/phải của biểu thức này tại từng vị trí ứng viên. Có thể có triệt tiêu khi hai đường gặp nhau, nên không chỉ liệt kê đường rồi khẳng định mọi đạo hàm đều mất. Với tính $C^1$ theo cả $x,t$, còn kiểm tra $u_t$. Với $C^2$, xét tiếp các góc của $g$ và đạo hàm cấp hai của $f$.

\textbf{Bước 5 — Tính thời điểm đến, giới hạn và ngưỡng.} Sóng ở $(x,t)$ chỉ nhận dữ liệu trong $[x-at,x+at]$. Giải bất đẳng thức giao hỗ trợ để tìm thời điểm đến và đi qua. Nếu $H$ hằng ngoài đoạn hỗ trợ, dùng các giá trị hằng này để tính giới hạn khi $t\to\infty$. Nếu hỏi áp suất cực đại, viết hàm từng khoảng thời gian và so sánh cực trị, bước nhảy cùng hai giá trị một phía.

\textbf{Bước 1 — Đọc đúng điều kiện biên.} Biên $u(0,t)=0$ dùng mở rộng lẻ; biên $u_x(0,t)=0$ dùng mở rộng chẵn. Với dữ liệu trên $x\ge0$, đặt $\widetilde f(-s)=-f(s)$ hoặc $f(s)$ tương ứng, làm tương tự cho $g$. Mở rộng lẻ là $C^2$ nếu $f(0)=f''(0)=0$, và mở rộng lẻ của $g$ là $C^1$ nếu $g(0)=0$. Mở rộng chẵn đòi $f'(0)=g'(0)=0$ cho mức trơn ấy.

\textbf{Bước 2 — Lập nguyên hàm sau khi phản xạ.} Đặt $H(s)=\int_0^s\widetilde g(r)dr$. Nếu $\widetilde g$ lẻ thì $H$ chẵn; nếu $\widetilde g$ chẵn thì $H$ lẻ. Chính nguyên hàm, chứ không phải vận tốc, xuất hiện trong


\[
u=\frac{\widetilde f(x-at)+\widetilde f(x+at)}2+\frac{H(x+at)-H(x-at)}{2a}.
\]


\textbf{Bước 3 — Tách miền trước và sau khi chạm biên.} Với $u(0,t)=h(t)$, ở $x\ge at$ dùng d'Alembert với dữ liệu gốc. Ở $0<x<at$:


\[
u=h(t-x/a)+\frac{f(x+at)-f(at-x)}2+\frac1{2a}\int_{at-x}^{at+x}g(s)ds.
\]


Với $u_x(0,t)=h(t)$, dùng nghiệm của các mở rộng chẵn rồi trừ $a\int_0^{t-x/a}h(s)ds$ trên miền $x<at$. Tính riêng đạo hàm theo $x$ của hạng này để kiểm tra biên.

\textbf{Bước 4 — Trên đoạn hữu hạn, phản xạ ở cả hai đầu.} Dirichlet ở hai đầu cho mở rộng lẻ tại $0,L$, chu kỳ không gian $2L$. Neumann ở hai đầu cho mở rộng chẵn tại $0,L$, cũng chu kỳ $2L$. Một đầu lẻ, một đầu chẵn cho $\widetilde f(s+2L)=-\widetilde f(s)$ và chu kỳ $4L$, tương tự $g$. Vì thế chu kỳ thời gian thường là $2L/a$ hoặc $4L/a$. Với Neumann hai đầu còn phải xét trung bình $g$: trung bình khác không sinh hạng tuyến tính theo $t$, nên nghiệm không tuần hoàn.

\textbf{Bước 5 — Kiểm tra góc và đường ghép.} Biên Dirichlet không thuần nhất cần $h(0)=f(0)$, $h'(0)=g(0)$, $h''(0)=a^2f''(0)$ để có nghiệm $C^2$ cho phương trình thuần nhất. So sánh hai vết trên $x=at$, gồm $u,u_x,u_t$ rồi đạo hàm cấp hai. Nếu không tương thích, nêu rõ mức trơn còn lại thay vì khẳng định nghiệm cổ điển trên toàn miền.
\subsection*{Lời giải áp dụng cho bài}
\textbf{Bước 1. Phản xạ Dirichlet.} $a=2$, $f=0$, $g=\chi_{[1,2]}$ trên nửa trục, biên $u(0,t)=0$. Mở rộng vận tốc lẻ: bằng $1$ trên $(1,2)$, $-1$ trên $(-2,-1)$, bằng $0$ còn lại.

\textbf{Bước 2. Nguyên hàm chẵn.} Với $H(0)=0$, trên đối số dương $H=0$ trước $1$, $s-1$ trên $(1,2)$, $1$ sau $2$; phản xạ được $H(s)=[|s|-1]_0^1$. Vì thế
\[A=-H/4,\quad B=H/4,\qquad u=\frac{[|x+2t|-1]_0^1-[|x-2t|-1]_0^1}{4}.\]
\textbf{Bước 3. Kiểm tra biên và mặt sóng.} Tại $x=0$, hai cận có cùng trị tuyệt đối nên $u=0$. $u_t(x,0)=H'(x)=g(x)$ tại điểm liên tục. Nghiệm liên tục nhưng không $C^1$ trên các nhánh $x\pm2t=\pm1,\pm2$ nằm trong $x,t>0$. Chúng chia đồ thị thành các đoạn thẳng. Tại mỗi $x$ cố định, sau $t>(x+2)/2$, hai đuôi bằng nhau và $u=0$.

\clearpage\noindent\texttt{\detokenize{VD_3_2_2_04}}\par
\section*{Lời giải bổ sung — Ví dụ 4 · Mục 3.2.2}

\subsection*{Kiến thức và trình tự giải}
\textbf{Bước 1 — Liệt kê các khoảng của dữ liệu.} Ghi đủ điểm đầu, điểm cuối hỗ trợ, điểm gãy và bước nhảy của $f,g$. Nếu có biên, phản xạ chẵn hoặc lẻ trước khi tích phân; trên đoạn hữu hạn tiếp tục phản xạ đến một chu kỳ hoàn chỉnh.

\textbf{Bước 2 — Tính nguyên hàm liên tục.} Lấy $H'=\widetilde g$ trên từng khoảng và chọn hằng số để hai vết của $H$ bằng nhau. Viết $[s]_p^q=\min(q,\max(p,s))$ hoặc $(s)_+=\max(s,0)$ khi cần gộp nhiều nhánh. Một nguyên hàm liên tục tránh tạo ra bước nhảy giả trong nghiệm.

\textbf{Bước 3 — Xác định sóng phải và sóng trái.} Đặt $A(s)=\widetilde f(s)/2-H(s)/(2a)$, $B(s)=\widetilde f(s)/2+H(s)/(2a)$; khi đó $u=A(x-at)+B(x+at)$. Với từng $t$ được yêu cầu, dịch đồ thị $A$ sang phải $at$ và $B$ sang trái $at$, rồi cộng tung độ. Các đường $x-at=b$ và $x+at=b$ chia mặt phẳng thành các miền có công thức đơn giản.

\textbf{Bước 4 — Xét liên tục, đạo hàm và giao của các mặt sóng.} Nếu $f$ có bước nhảy, $u$ thường có bước nhảy; nếu $f$ liên tục và $g$ bị chặn thì $u$ liên tục. Tại thời điểm cố định,


\[
u_x=\frac{\widetilde f'(x-at)+\widetilde f'(x+at)}2+\frac{\widetilde g(x+at)-\widetilde g(x-at)}{2a}.
\]


Tính giới hạn trái/phải của biểu thức này tại từng vị trí ứng viên. Có thể có triệt tiêu khi hai đường gặp nhau, nên không chỉ liệt kê đường rồi khẳng định mọi đạo hàm đều mất. Với tính $C^1$ theo cả $x,t$, còn kiểm tra $u_t$. Với $C^2$, xét tiếp các góc của $g$ và đạo hàm cấp hai của $f$.

\textbf{Bước 5 — Tính thời điểm đến, giới hạn và ngưỡng.} Sóng ở $(x,t)$ chỉ nhận dữ liệu trong $[x-at,x+at]$. Giải bất đẳng thức giao hỗ trợ để tìm thời điểm đến và đi qua. Nếu $H$ hằng ngoài đoạn hỗ trợ, dùng các giá trị hằng này để tính giới hạn khi $t\to\infty$. Nếu hỏi áp suất cực đại, viết hàm từng khoảng thời gian và so sánh cực trị, bước nhảy cùng hai giá trị một phía.

\textbf{Bước 1 — Đọc đúng điều kiện biên.} Biên $u(0,t)=0$ dùng mở rộng lẻ; biên $u_x(0,t)=0$ dùng mở rộng chẵn. Với dữ liệu trên $x\ge0$, đặt $\widetilde f(-s)=-f(s)$ hoặc $f(s)$ tương ứng, làm tương tự cho $g$. Mở rộng lẻ là $C^2$ nếu $f(0)=f''(0)=0$, và mở rộng lẻ của $g$ là $C^1$ nếu $g(0)=0$. Mở rộng chẵn đòi $f'(0)=g'(0)=0$ cho mức trơn ấy.

\textbf{Bước 2 — Lập nguyên hàm sau khi phản xạ.} Đặt $H(s)=\int_0^s\widetilde g(r)dr$. Nếu $\widetilde g$ lẻ thì $H$ chẵn; nếu $\widetilde g$ chẵn thì $H$ lẻ. Chính nguyên hàm, chứ không phải vận tốc, xuất hiện trong


\[
u=\frac{\widetilde f(x-at)+\widetilde f(x+at)}2+\frac{H(x+at)-H(x-at)}{2a}.
\]


\textbf{Bước 3 — Tách miền trước và sau khi chạm biên.} Với $u(0,t)=h(t)$, ở $x\ge at$ dùng d'Alembert với dữ liệu gốc. Ở $0<x<at$:


\[
u=h(t-x/a)+\frac{f(x+at)-f(at-x)}2+\frac1{2a}\int_{at-x}^{at+x}g(s)ds.
\]


Với $u_x(0,t)=h(t)$, dùng nghiệm của các mở rộng chẵn rồi trừ $a\int_0^{t-x/a}h(s)ds$ trên miền $x<at$. Tính riêng đạo hàm theo $x$ của hạng này để kiểm tra biên.

\textbf{Bước 4 — Trên đoạn hữu hạn, phản xạ ở cả hai đầu.} Dirichlet ở hai đầu cho mở rộng lẻ tại $0,L$, chu kỳ không gian $2L$. Neumann ở hai đầu cho mở rộng chẵn tại $0,L$, cũng chu kỳ $2L$. Một đầu lẻ, một đầu chẵn cho $\widetilde f(s+2L)=-\widetilde f(s)$ và chu kỳ $4L$, tương tự $g$. Vì thế chu kỳ thời gian thường là $2L/a$ hoặc $4L/a$. Với Neumann hai đầu còn phải xét trung bình $g$: trung bình khác không sinh hạng tuyến tính theo $t$, nên nghiệm không tuần hoàn.

\textbf{Bước 5 — Kiểm tra góc và đường ghép.} Biên Dirichlet không thuần nhất cần $h(0)=f(0)$, $h'(0)=g(0)$, $h''(0)=a^2f''(0)$ để có nghiệm $C^2$ cho phương trình thuần nhất. So sánh hai vết trên $x=at$, gồm $u,u_x,u_t$ rồi đạo hàm cấp hai. Nếu không tương thích, nêu rõ mức trơn còn lại thay vì khẳng định nghiệm cổ điển trên toàn miền.
\subsection*{Lời giải áp dụng cho bài}
\textbf{Bước 1. Phản xạ Neumann.} Dữ liệu như Ví dụ 3, nhưng $u_x(0,t)=0$. Mở rộng vận tốc chẵn thành $\chi_{[1,2]}(|s|)$, nên nguyên hàm phải lẻ, không dùng nguyên hàm chẵn của ví dụ trước.

\textbf{Bước 2. Tính hai sóng.}
\[H(s)=\operatorname{sgn}(s)[|s|-1]_0^1,\quad A=-H/4,\quad B=H/4,\quad u=\frac{H(x+2t)-H(x-2t)}4.\]
Trước $x=2t$ vẫn có thể dùng các cận dương; sau đó giá trị $H(x-2t)$ âm, tạo xung phản xạ cùng dấu.

\textbf{Bước 3. Kiểm tra.} $u_x(0,t)=(H'(2t)-H'(-2t))/4=0$. $u$ liên tục, gãy cấp một trên $x\pm2t=\pm1,\pm2$. Độ dịch đúng là $2t$ vì PDE có hệ số $4$; những chỗ ghi $x\pm t$ trong hình/công thức cần hiểu theo tốc độ thật. Sau $t>(x+2)/2$, hai đuôi $H$ là $1,-1$, nên $u=1/2$, khác phản xạ Dirichlet trở về $0$.

\clearpage\noindent\texttt{\detokenize{VD_3_2_2_05}}\par
\section*{Lời giải bổ sung — Ví dụ 5 · Mục 3.2.2}

\subsection*{Kiến thức và trình tự giải}
\textbf{Bước 1 — Liệt kê các khoảng của dữ liệu.} Ghi đủ điểm đầu, điểm cuối hỗ trợ, điểm gãy và bước nhảy của $f,g$. Nếu có biên, phản xạ chẵn hoặc lẻ trước khi tích phân; trên đoạn hữu hạn tiếp tục phản xạ đến một chu kỳ hoàn chỉnh.

\textbf{Bước 2 — Tính nguyên hàm liên tục.} Lấy $H'=\widetilde g$ trên từng khoảng và chọn hằng số để hai vết của $H$ bằng nhau. Viết $[s]_p^q=\min(q,\max(p,s))$ hoặc $(s)_+=\max(s,0)$ khi cần gộp nhiều nhánh. Một nguyên hàm liên tục tránh tạo ra bước nhảy giả trong nghiệm.

\textbf{Bước 3 — Xác định sóng phải và sóng trái.} Đặt $A(s)=\widetilde f(s)/2-H(s)/(2a)$, $B(s)=\widetilde f(s)/2+H(s)/(2a)$; khi đó $u=A(x-at)+B(x+at)$. Với từng $t$ được yêu cầu, dịch đồ thị $A$ sang phải $at$ và $B$ sang trái $at$, rồi cộng tung độ. Các đường $x-at=b$ và $x+at=b$ chia mặt phẳng thành các miền có công thức đơn giản.

\textbf{Bước 4 — Xét liên tục, đạo hàm và giao của các mặt sóng.} Nếu $f$ có bước nhảy, $u$ thường có bước nhảy; nếu $f$ liên tục và $g$ bị chặn thì $u$ liên tục. Tại thời điểm cố định,


\[
u_x=\frac{\widetilde f'(x-at)+\widetilde f'(x+at)}2+\frac{\widetilde g(x+at)-\widetilde g(x-at)}{2a}.
\]


Tính giới hạn trái/phải của biểu thức này tại từng vị trí ứng viên. Có thể có triệt tiêu khi hai đường gặp nhau, nên không chỉ liệt kê đường rồi khẳng định mọi đạo hàm đều mất. Với tính $C^1$ theo cả $x,t$, còn kiểm tra $u_t$. Với $C^2$, xét tiếp các góc của $g$ và đạo hàm cấp hai của $f$.

\textbf{Bước 5 — Tính thời điểm đến, giới hạn và ngưỡng.} Sóng ở $(x,t)$ chỉ nhận dữ liệu trong $[x-at,x+at]$. Giải bất đẳng thức giao hỗ trợ để tìm thời điểm đến và đi qua. Nếu $H$ hằng ngoài đoạn hỗ trợ, dùng các giá trị hằng này để tính giới hạn khi $t\to\infty$. Nếu hỏi áp suất cực đại, viết hàm từng khoảng thời gian và so sánh cực trị, bước nhảy cùng hai giá trị một phía.

\textbf{Bước 1 — Đọc đúng điều kiện biên.} Biên $u(0,t)=0$ dùng mở rộng lẻ; biên $u_x(0,t)=0$ dùng mở rộng chẵn. Với dữ liệu trên $x\ge0$, đặt $\widetilde f(-s)=-f(s)$ hoặc $f(s)$ tương ứng, làm tương tự cho $g$. Mở rộng lẻ là $C^2$ nếu $f(0)=f''(0)=0$, và mở rộng lẻ của $g$ là $C^1$ nếu $g(0)=0$. Mở rộng chẵn đòi $f'(0)=g'(0)=0$ cho mức trơn ấy.

\textbf{Bước 2 — Lập nguyên hàm sau khi phản xạ.} Đặt $H(s)=\int_0^s\widetilde g(r)dr$. Nếu $\widetilde g$ lẻ thì $H$ chẵn; nếu $\widetilde g$ chẵn thì $H$ lẻ. Chính nguyên hàm, chứ không phải vận tốc, xuất hiện trong


\[
u=\frac{\widetilde f(x-at)+\widetilde f(x+at)}2+\frac{H(x+at)-H(x-at)}{2a}.
\]


\textbf{Bước 3 — Tách miền trước và sau khi chạm biên.} Với $u(0,t)=h(t)$, ở $x\ge at$ dùng d'Alembert với dữ liệu gốc. Ở $0<x<at$:


\[
u=h(t-x/a)+\frac{f(x+at)-f(at-x)}2+\frac1{2a}\int_{at-x}^{at+x}g(s)ds.
\]


Với $u_x(0,t)=h(t)$, dùng nghiệm của các mở rộng chẵn rồi trừ $a\int_0^{t-x/a}h(s)ds$ trên miền $x<at$. Tính riêng đạo hàm theo $x$ của hạng này để kiểm tra biên.

\textbf{Bước 4 — Trên đoạn hữu hạn, phản xạ ở cả hai đầu.} Dirichlet ở hai đầu cho mở rộng lẻ tại $0,L$, chu kỳ không gian $2L$. Neumann ở hai đầu cho mở rộng chẵn tại $0,L$, cũng chu kỳ $2L$. Một đầu lẻ, một đầu chẵn cho $\widetilde f(s+2L)=-\widetilde f(s)$ và chu kỳ $4L$, tương tự $g$. Vì thế chu kỳ thời gian thường là $2L/a$ hoặc $4L/a$. Với Neumann hai đầu còn phải xét trung bình $g$: trung bình khác không sinh hạng tuyến tính theo $t$, nên nghiệm không tuần hoàn.

\textbf{Bước 5 — Kiểm tra góc và đường ghép.} Biên Dirichlet không thuần nhất cần $h(0)=f(0)$, $h'(0)=g(0)$, $h''(0)=a^2f''(0)$ để có nghiệm $C^2$ cho phương trình thuần nhất. So sánh hai vết trên $x=at$, gồm $u,u_x,u_t$ rồi đạo hàm cấp hai. Nếu không tương thích, nêu rõ mức trơn còn lại thay vì khẳng định nghiệm cổ điển trên toàn miền.
\subsection*{Lời giải áp dụng cho bài}
\textbf{Bước 1. Phản xạ hai biên cố định.} $a=L=1$, $f=\chi_{[1/4,1/3]}$, $g=4\chi_{[2/3,3/4]}$. Mở rộng lẻ tại $0,1$ cho chu kỳ $2$. Đặt $r=\operatorname{mod}(s,2)$: $f_e(s)=f(r)$ nếu $r\le1$, $-f(2-r)$ nếu $r>1$; vận tốc mở rộng tương tự.

\textbf{Bước 2. Nguyên hàm trong một chu kỳ.} Trên $[0,1]$, $H(s)=4[s-2/3]_0^{1/12}$; nó tăng từ $0$ lên $1/3$. Trên $[1,2]$, $H(s)=H(2-s)$, sau đó mở rộng chu kỳ $2$. Sự giảm trên nửa sau do vận tốc phản xạ âm, nên $H(2)=0$.

\textbf{Bước 3. Ghép nghiệm.}
\[u=\tfrac12[f_e(x-t)+f_e(x+t)]+\tfrac12[H(x+t)-H(x-t)].\]
Hai sóng chạy là $A=f_e/2-H/2$, $B=f_e/2+H/2$. Tính từng hình bằng giảm các đối số về $[0,2]$, rồi chọn các nhánh trên; không cần viết một danh sách dài công thức thời gian.

\textbf{Bước 4. Chu kỳ và tính trơn.} Cả $f_e,H$ có chu kỳ $2$, nên $u(t+2)=u(t)$. Li độ bước nhảy gây gián đoạn trên các đường $x\pm t=\pm1/4+2k,\pm1/3+2k$; vận tốc gây gãy cấp một trên $x\pm t=\pm2/3+2k,\pm3/4+2k$, $k\in\mathbb Z$. Ngoài các đường ấy, nghiệm affine cộng hằng. Biên $u(0,t)=u(1,t)=0$ theo tính lẻ tại hai đầu; dữ liệu ban đầu đúng trong miền theo nghĩa yếu.

\clearpage\noindent\texttt{\detokenize{VD_3_2_2_06}}\par
\section*{Lời giải bổ sung — Ví dụ 6 · Mục 3.2.2}

\subsection*{Kiến thức và trình tự giải}
\textbf{Bước 1 — Liệt kê các khoảng của dữ liệu.} Ghi đủ điểm đầu, điểm cuối hỗ trợ, điểm gãy và bước nhảy của $f,g$. Nếu có biên, phản xạ chẵn hoặc lẻ trước khi tích phân; trên đoạn hữu hạn tiếp tục phản xạ đến một chu kỳ hoàn chỉnh.

\textbf{Bước 2 — Tính nguyên hàm liên tục.} Lấy $H'=\widetilde g$ trên từng khoảng và chọn hằng số để hai vết của $H$ bằng nhau. Viết $[s]_p^q=\min(q,\max(p,s))$ hoặc $(s)_+=\max(s,0)$ khi cần gộp nhiều nhánh. Một nguyên hàm liên tục tránh tạo ra bước nhảy giả trong nghiệm.

\textbf{Bước 3 — Xác định sóng phải và sóng trái.} Đặt $A(s)=\widetilde f(s)/2-H(s)/(2a)$, $B(s)=\widetilde f(s)/2+H(s)/(2a)$; khi đó $u=A(x-at)+B(x+at)$. Với từng $t$ được yêu cầu, dịch đồ thị $A$ sang phải $at$ và $B$ sang trái $at$, rồi cộng tung độ. Các đường $x-at=b$ và $x+at=b$ chia mặt phẳng thành các miền có công thức đơn giản.

\textbf{Bước 4 — Xét liên tục, đạo hàm và giao của các mặt sóng.} Nếu $f$ có bước nhảy, $u$ thường có bước nhảy; nếu $f$ liên tục và $g$ bị chặn thì $u$ liên tục. Tại thời điểm cố định,


\[
u_x=\frac{\widetilde f'(x-at)+\widetilde f'(x+at)}2+\frac{\widetilde g(x+at)-\widetilde g(x-at)}{2a}.
\]


Tính giới hạn trái/phải của biểu thức này tại từng vị trí ứng viên. Có thể có triệt tiêu khi hai đường gặp nhau, nên không chỉ liệt kê đường rồi khẳng định mọi đạo hàm đều mất. Với tính $C^1$ theo cả $x,t$, còn kiểm tra $u_t$. Với $C^2$, xét tiếp các góc của $g$ và đạo hàm cấp hai của $f$.

\textbf{Bước 5 — Tính thời điểm đến, giới hạn và ngưỡng.} Sóng ở $(x,t)$ chỉ nhận dữ liệu trong $[x-at,x+at]$. Giải bất đẳng thức giao hỗ trợ để tìm thời điểm đến và đi qua. Nếu $H$ hằng ngoài đoạn hỗ trợ, dùng các giá trị hằng này để tính giới hạn khi $t\to\infty$. Nếu hỏi áp suất cực đại, viết hàm từng khoảng thời gian và so sánh cực trị, bước nhảy cùng hai giá trị một phía.

\textbf{Bước 1 — Đọc đúng điều kiện biên.} Biên $u(0,t)=0$ dùng mở rộng lẻ; biên $u_x(0,t)=0$ dùng mở rộng chẵn. Với dữ liệu trên $x\ge0$, đặt $\widetilde f(-s)=-f(s)$ hoặc $f(s)$ tương ứng, làm tương tự cho $g$. Mở rộng lẻ là $C^2$ nếu $f(0)=f''(0)=0$, và mở rộng lẻ của $g$ là $C^1$ nếu $g(0)=0$. Mở rộng chẵn đòi $f'(0)=g'(0)=0$ cho mức trơn ấy.

\textbf{Bước 2 — Lập nguyên hàm sau khi phản xạ.} Đặt $H(s)=\int_0^s\widetilde g(r)dr$. Nếu $\widetilde g$ lẻ thì $H$ chẵn; nếu $\widetilde g$ chẵn thì $H$ lẻ. Chính nguyên hàm, chứ không phải vận tốc, xuất hiện trong


\[
u=\frac{\widetilde f(x-at)+\widetilde f(x+at)}2+\frac{H(x+at)-H(x-at)}{2a}.
\]


\textbf{Bước 3 — Tách miền trước và sau khi chạm biên.} Với $u(0,t)=h(t)$, ở $x\ge at$ dùng d'Alembert với dữ liệu gốc. Ở $0<x<at$:


\[
u=h(t-x/a)+\frac{f(x+at)-f(at-x)}2+\frac1{2a}\int_{at-x}^{at+x}g(s)ds.
\]


Với $u_x(0,t)=h(t)$, dùng nghiệm của các mở rộng chẵn rồi trừ $a\int_0^{t-x/a}h(s)ds$ trên miền $x<at$. Tính riêng đạo hàm theo $x$ của hạng này để kiểm tra biên.

\textbf{Bước 4 — Trên đoạn hữu hạn, phản xạ ở cả hai đầu.} Dirichlet ở hai đầu cho mở rộng lẻ tại $0,L$, chu kỳ không gian $2L$. Neumann ở hai đầu cho mở rộng chẵn tại $0,L$, cũng chu kỳ $2L$. Một đầu lẻ, một đầu chẵn cho $\widetilde f(s+2L)=-\widetilde f(s)$ và chu kỳ $4L$, tương tự $g$. Vì thế chu kỳ thời gian thường là $2L/a$ hoặc $4L/a$. Với Neumann hai đầu còn phải xét trung bình $g$: trung bình khác không sinh hạng tuyến tính theo $t$, nên nghiệm không tuần hoàn.

\textbf{Bước 5 — Kiểm tra góc và đường ghép.} Biên Dirichlet không thuần nhất cần $h(0)=f(0)$, $h'(0)=g(0)$, $h''(0)=a^2f''(0)$ để có nghiệm $C^2$ cho phương trình thuần nhất. So sánh hai vết trên $x=at$, gồm $u,u_x,u_t$ rồi đạo hàm cấp hai. Nếu không tương thích, nêu rõ mức trơn còn lại thay vì khẳng định nghiệm cổ điển trên toàn miền.
\subsection*{Lời giải áp dụng cho bài}
\textbf{Bước 1. Phản xạ biên hỗn hợp.} Dữ liệu giống Ví dụ 5, nhưng $u(0,t)=0$, $u_x(1,t)=0$. Trên $[0,2]$, mở rộng $f_e(s)=f(s)$ nếu $s\le1$, $f(2-s)$ nếu $s\ge1$; sau đó $f_e(s+2)=-f_e(s)$ và chu kỳ $4$. Làm tương tự với $g_e$.

\textbf{Bước 2. Tính nguyên hàm bốn nhánh.} Đặt $h(s)=4[s-2/3]_0^{1/12}$ trên $[0,1]$. Trên $[0,4]$:
\[H(s)=\begin{cases}h(s)&0\le s\le1,\\2/3-h(2-s)&1\le s\le2,\\2/3-H(s-2)&2\le s\le4.\end{cases}\]
Hạng $2/3$ là tích phân vận tốc trên $[0,2]$; toàn chu kỳ $4$ có trung bình bằng $0$, nên mở rộng $H$ chu kỳ $4$.

\textbf{Bước 3. Nghiệm và phản chu kỳ.} Vẫn dùng $u=(f_e(x-t)+f_e(x+t))/2+(H(x+t)-H(x-t))/2$. Ta có $H(s+2)=2/3-H(s)$, nên
\[u(x,t+2)=-u(x,t),\qquad u(x,t+4)=u(x,t).\]
Vì dữ liệu khác không, không có chu kỳ $2$; câu yêu cầu chu kỳ $2$ trong sách cần sửa thành phản chu kỳ $2$ và chu kỳ $4$. Mode đầu khác không (tích phân $f\sin(\pi x/2)>0$), xác nhận chu kỳ cơ bản $4$.

\textbf{Bước 4. Kiểm tra biên.} $f_e,g_e$ lẻ tại $0$, chẵn tại $1$, nên nghiệm có đúng vết Dirichlet và Neumann nơi vết đạo hàm tồn tại. Các bước nhảy dữ liệu lan dọc $x\pm t=b$, trong đó $b$ là các điểm biên hỗ trợ cùng các ảnh phản xạ đã xây dựng. Mức trơn là nghiệm từng miền/ nghiệm yếu, không $C^2$ trên toàn hình chữ nhật.

\textbf{Chi tiết nối nguyên hàm và phản chu kỳ.} Xung dương có diện tích $4(3/4-2/3)=1/3$. Phản xạ chẵn tại $1$ tạo thêm diện tích $1/3$ trên $(1,2)$, nên $H(2)=2/3$. Vì $g_e(s+2)=-g_e(s)$,
\[H(s+2)=H(2)+\int_2^{s+2}g_e(r)dr=2/3-\int_0^sg_e(r)dr=2/3-H(s).\]
Khi dịch $t$ thêm $2$, các đối số $x+t$ tăng $2$, $x-t$ giảm $2$; cả hai giá trị $H$ thành $2/3-H$ của giá trị cũ. Vì thế hiệu $H$ đổi dấu. Các giá trị $f_e$ cũng đổi dấu, cho $u(t+2)=-u(t)$.

\clearpage\noindent\texttt{\detokenize{BT_3_2_3_01}}\par
\section*{Lời giải bổ sung — Bài tập 1 · Mục 3.2.3}

\subsection*{Kiến thức và trình tự giải}
\textbf{Bước 1 — Liệt kê các khoảng của dữ liệu.} Ghi đủ điểm đầu, điểm cuối hỗ trợ, điểm gãy và bước nhảy của $f,g$. Nếu có biên, phản xạ chẵn hoặc lẻ trước khi tích phân; trên đoạn hữu hạn tiếp tục phản xạ đến một chu kỳ hoàn chỉnh.

\textbf{Bước 2 — Tính nguyên hàm liên tục.} Lấy $H'=\widetilde g$ trên từng khoảng và chọn hằng số để hai vết của $H$ bằng nhau. Viết $[s]_p^q=\min(q,\max(p,s))$ hoặc $(s)_+=\max(s,0)$ khi cần gộp nhiều nhánh. Một nguyên hàm liên tục tránh tạo ra bước nhảy giả trong nghiệm.

\textbf{Bước 3 — Xác định sóng phải và sóng trái.} Đặt $A(s)=\widetilde f(s)/2-H(s)/(2a)$, $B(s)=\widetilde f(s)/2+H(s)/(2a)$; khi đó $u=A(x-at)+B(x+at)$. Với từng $t$ được yêu cầu, dịch đồ thị $A$ sang phải $at$ và $B$ sang trái $at$, rồi cộng tung độ. Các đường $x-at=b$ và $x+at=b$ chia mặt phẳng thành các miền có công thức đơn giản.

\textbf{Bước 4 — Xét liên tục, đạo hàm và giao của các mặt sóng.} Nếu $f$ có bước nhảy, $u$ thường có bước nhảy; nếu $f$ liên tục và $g$ bị chặn thì $u$ liên tục. Tại thời điểm cố định,


\[
u_x=\frac{\widetilde f'(x-at)+\widetilde f'(x+at)}2+\frac{\widetilde g(x+at)-\widetilde g(x-at)}{2a}.
\]


Tính giới hạn trái/phải của biểu thức này tại từng vị trí ứng viên. Có thể có triệt tiêu khi hai đường gặp nhau, nên không chỉ liệt kê đường rồi khẳng định mọi đạo hàm đều mất. Với tính $C^1$ theo cả $x,t$, còn kiểm tra $u_t$. Với $C^2$, xét tiếp các góc của $g$ và đạo hàm cấp hai của $f$.

\textbf{Bước 5 — Tính thời điểm đến, giới hạn và ngưỡng.} Sóng ở $(x,t)$ chỉ nhận dữ liệu trong $[x-at,x+at]$. Giải bất đẳng thức giao hỗ trợ để tìm thời điểm đến và đi qua. Nếu $H$ hằng ngoài đoạn hỗ trợ, dùng các giá trị hằng này để tính giới hạn khi $t\to\infty$. Nếu hỏi áp suất cực đại, viết hàm từng khoảng thời gian và so sánh cực trị, bước nhảy cùng hai giá trị một phía.
\subsection*{Lời giải áp dụng cho bài}
\textbf{Bước 1. Đọc dữ liệu và chọn phản xạ.} Toàn trục, $a=2$, $f=0$, $g=s$ trên $[-1,0]$, bằng không bên ngoài.

\textbf{Bước 2. Tính từng nhánh và nối nguyên hàm.} Tích phân từ $-1$ đến $s$ cho $(s^2-1)/2$; trước $-1$ lấy $0$, sau $0$ giữ $-1/2$. Do đó $H$ liên tục ở cả hai đầu.
\[f_e(s)=0,\qquad H(s)=\tfrac12(\operatorname{clip}(s,-1,0)^2-1).\]
Ở đây $\operatorname{clip}(s,p,q)=\min(q,\max(p,s))$; $\operatorname{mod}(s,L)\in[0,L)$ kể cả với $s<0$.

\textbf{Bước 3. Xác định hai sóng chạy và ghép nghiệm.} $A=f_e/2-H/(2a)$, $B=f_e/2+H/(2a)$. Thay hai cận:
\[u(x,t)=\frac{f_e(x-2t)+f_e(x+2t)}2+\frac{H(x+2t)-H(x-2t)}{4}.\]

\textbf{Bước 4. Trả lời các yêu cầu và kiểm tra.} Tại $t=1$, cận $x+2=-1$ cho $x=-3$, cận $x-2=-1$ cho $x=1$. Bước nhảy của $g$ ở $-1$ khác không nên hai điểm này không có $u_x$. Ở $0$, $g$ nối về $0$ nên $u_x$ vẫn tồn tại; chỉ $u_{xx}$ gãy.

Tại $t=0$, hiệu nguyên hàm bằng $0$ và đạo hàm thời gian cho $H'(x)=g(x)$ trong miền dữ liệu; half-sum cho đúng li độ. Mỗi hạng là sóng của $x\pm at$, nên thỏa PDE ngoài các mặt gãy và theo nghĩa phân phối trên toàn miền. Đồ thị hai sóng ban đầu và các lát cắt thời gian được tính trực tiếp từ các nhánh trên.
\begin{center}\includegraphics[width=.97\linewidth,height=11cm,keepaspectratio]{_xu_ly/graphs/BT_3_2_3_01.png}\end{center}
\textit{Đồ thị minh họa; với nghiệm số, xem phương pháp và tham số trong mã Python kèm theo.}

\clearpage\noindent\texttt{\detokenize{BT_3_2_3_02}}\par
\section*{Lời giải bổ sung — Bài tập 2 · Mục 3.2.3}

\subsection*{Kiến thức và trình tự giải}
\textbf{Bước 1 — Liệt kê các khoảng của dữ liệu.} Ghi đủ điểm đầu, điểm cuối hỗ trợ, điểm gãy và bước nhảy của $f,g$. Nếu có biên, phản xạ chẵn hoặc lẻ trước khi tích phân; trên đoạn hữu hạn tiếp tục phản xạ đến một chu kỳ hoàn chỉnh.

\textbf{Bước 2 — Tính nguyên hàm liên tục.} Lấy $H'=\widetilde g$ trên từng khoảng và chọn hằng số để hai vết của $H$ bằng nhau. Viết $[s]_p^q=\min(q,\max(p,s))$ hoặc $(s)_+=\max(s,0)$ khi cần gộp nhiều nhánh. Một nguyên hàm liên tục tránh tạo ra bước nhảy giả trong nghiệm.

\textbf{Bước 3 — Xác định sóng phải và sóng trái.} Đặt $A(s)=\widetilde f(s)/2-H(s)/(2a)$, $B(s)=\widetilde f(s)/2+H(s)/(2a)$; khi đó $u=A(x-at)+B(x+at)$. Với từng $t$ được yêu cầu, dịch đồ thị $A$ sang phải $at$ và $B$ sang trái $at$, rồi cộng tung độ. Các đường $x-at=b$ và $x+at=b$ chia mặt phẳng thành các miền có công thức đơn giản.

\textbf{Bước 4 — Xét liên tục, đạo hàm và giao của các mặt sóng.} Nếu $f$ có bước nhảy, $u$ thường có bước nhảy; nếu $f$ liên tục và $g$ bị chặn thì $u$ liên tục. Tại thời điểm cố định,


\[
u_x=\frac{\widetilde f'(x-at)+\widetilde f'(x+at)}2+\frac{\widetilde g(x+at)-\widetilde g(x-at)}{2a}.
\]


Tính giới hạn trái/phải của biểu thức này tại từng vị trí ứng viên. Có thể có triệt tiêu khi hai đường gặp nhau, nên không chỉ liệt kê đường rồi khẳng định mọi đạo hàm đều mất. Với tính $C^1$ theo cả $x,t$, còn kiểm tra $u_t$. Với $C^2$, xét tiếp các góc của $g$ và đạo hàm cấp hai của $f$.

\textbf{Bước 5 — Tính thời điểm đến, giới hạn và ngưỡng.} Sóng ở $(x,t)$ chỉ nhận dữ liệu trong $[x-at,x+at]$. Giải bất đẳng thức giao hỗ trợ để tìm thời điểm đến và đi qua. Nếu $H$ hằng ngoài đoạn hỗ trợ, dùng các giá trị hằng này để tính giới hạn khi $t\to\infty$. Nếu hỏi áp suất cực đại, viết hàm từng khoảng thời gian và so sánh cực trị, bước nhảy cùng hai giá trị một phía.
\subsection*{Lời giải áp dụng cho bài}
\textbf{Bước 1. Đọc dữ liệu và chọn phản xạ.} Toàn trục, $a=2$, $f=(1-|s|)_+$, $g=\chi_{[1,2]}$.

\textbf{Bước 2. Tính từng nhánh và nối nguyên hàm.} $H=0$ trước $1$, $s-1$ trên $[1,2]$, $1$ sau $2$. Các nhánh của $f$ có hệ số góc $0,1,-1,0$ tại $-1,0,1$.
\[f_e(s)=(1-|s|)_+,\qquad H(s)=\operatorname{clip}(s-1,0,1).\]
Ở đây $\operatorname{clip}(s,p,q)=\min(q,\max(p,s))$; $\operatorname{mod}(s,L)\in[0,L)$ kể cả với $s<0$.

\textbf{Bước 3. Xác định hai sóng chạy và ghép nghiệm.} $A=f_e/2-H/(2a)$, $B=f_e/2+H/(2a)$. Thay hai cận:
\[u(x,t)=\frac{f_e(x-2t)+f_e(x+2t)}2+\frac{H(x+2t)-H(x-2t)}{4}.\]

\textbf{Bước 4. Trả lời các yêu cầu và kiểm tra.} $f,H$ liên tục nên $u$ liên tục ở mọi $(x,t)$, kể cả $t=1$. Ở $t=1$, vị trí ứng viên của $u_x$ là $-3,-2,-1,0,1,2,3,4$; cộng các bước nhảy theo công thức đạo hàm để xét từng lát cắt. Trong hai biến, các đường xuất phát từ $-1,0,1$ của $f$ và $1,2$ của $g$ là các mặt sóng cấp một.

Tại $t=0$, hiệu nguyên hàm bằng $0$ và đạo hàm thời gian cho $H'(x)=g(x)$ trong miền dữ liệu; half-sum cho đúng li độ. Mỗi hạng là sóng của $x\pm at$, nên thỏa PDE ngoài các mặt gãy và theo nghĩa phân phối trên toàn miền. Đồ thị hai sóng ban đầu và các lát cắt thời gian được tính trực tiếp từ các nhánh trên.
\begin{center}\includegraphics[width=.97\linewidth,height=11cm,keepaspectratio]{_xu_ly/graphs/BT_3_2_3_02.png}\end{center}
\textit{Đồ thị minh họa; với nghiệm số, xem phương pháp và tham số trong mã Python kèm theo.}

\clearpage\noindent\texttt{\detokenize{BT_3_2_3_03}}\par
\section*{Lời giải bổ sung — Bài tập 3 · Mục 3.2.3}

\subsection*{Kiến thức và trình tự giải}
\textbf{Bước 1 — Liệt kê các khoảng của dữ liệu.} Ghi đủ điểm đầu, điểm cuối hỗ trợ, điểm gãy và bước nhảy của $f,g$. Nếu có biên, phản xạ chẵn hoặc lẻ trước khi tích phân; trên đoạn hữu hạn tiếp tục phản xạ đến một chu kỳ hoàn chỉnh.

\textbf{Bước 2 — Tính nguyên hàm liên tục.} Lấy $H'=\widetilde g$ trên từng khoảng và chọn hằng số để hai vết của $H$ bằng nhau. Viết $[s]_p^q=\min(q,\max(p,s))$ hoặc $(s)_+=\max(s,0)$ khi cần gộp nhiều nhánh. Một nguyên hàm liên tục tránh tạo ra bước nhảy giả trong nghiệm.

\textbf{Bước 3 — Xác định sóng phải và sóng trái.} Đặt $A(s)=\widetilde f(s)/2-H(s)/(2a)$, $B(s)=\widetilde f(s)/2+H(s)/(2a)$; khi đó $u=A(x-at)+B(x+at)$. Với từng $t$ được yêu cầu, dịch đồ thị $A$ sang phải $at$ và $B$ sang trái $at$, rồi cộng tung độ. Các đường $x-at=b$ và $x+at=b$ chia mặt phẳng thành các miền có công thức đơn giản.

\textbf{Bước 4 — Xét liên tục, đạo hàm và giao của các mặt sóng.} Nếu $f$ có bước nhảy, $u$ thường có bước nhảy; nếu $f$ liên tục và $g$ bị chặn thì $u$ liên tục. Tại thời điểm cố định,


\[
u_x=\frac{\widetilde f'(x-at)+\widetilde f'(x+at)}2+\frac{\widetilde g(x+at)-\widetilde g(x-at)}{2a}.
\]


Tính giới hạn trái/phải của biểu thức này tại từng vị trí ứng viên. Có thể có triệt tiêu khi hai đường gặp nhau, nên không chỉ liệt kê đường rồi khẳng định mọi đạo hàm đều mất. Với tính $C^1$ theo cả $x,t$, còn kiểm tra $u_t$. Với $C^2$, xét tiếp các góc của $g$ và đạo hàm cấp hai của $f$.

\textbf{Bước 5 — Tính thời điểm đến, giới hạn và ngưỡng.} Sóng ở $(x,t)$ chỉ nhận dữ liệu trong $[x-at,x+at]$. Giải bất đẳng thức giao hỗ trợ để tìm thời điểm đến và đi qua. Nếu $H$ hằng ngoài đoạn hỗ trợ, dùng các giá trị hằng này để tính giới hạn khi $t\to\infty$. Nếu hỏi áp suất cực đại, viết hàm từng khoảng thời gian và so sánh cực trị, bước nhảy cùng hai giá trị một phía.
\subsection*{Lời giải áp dụng cho bài}
\textbf{Bước 1. Đọc dữ liệu và chọn phản xạ.} Toàn trục, $a=3$, $f=\chi_{[-1,1]}$, $g=s\chi_{[1,2]}$.

\textbf{Bước 2. Tính từng nhánh và nối nguyên hàm.} Trên $[1,2]$, $\int_1^s r,dr=(s^2-1)/2$; sau $2$ giữ $3/2$. Phần li độ là hai xung biên độ $1/2$, chạy ngược chiều.
\[f_e(s)=\chi_{[-1,1]}(s),\qquad H(s)=\tfrac12(\operatorname{clip}(s,1,2)^2-1).\]
Ở đây $\operatorname{clip}(s,p,q)=\min(q,\max(p,s))$; $\operatorname{mod}(s,L)\in[0,L)$ kể cả với $s<0$.

\textbf{Bước 3. Xác định hai sóng chạy và ghép nghiệm.} $A=f_e/2-H/(2a)$, $B=f_e/2+H/(2a)$. Thay hai cận:
\[u(x,t)=\frac{f_e(x-3t)+f_e(x+3t)}2+\frac{H(x+3t)-H(x-3t)}{6}.\]

\textbf{Bước 4. Trả lời các yêu cầu và kiểm tra.} Tại $t=1$, hai đoạn li độ là $[-4,-2]$ và $[2,4]$, có bước nhảy tại $-4,-2,2,4$. Phần $H$ liên tục nên không thể khử các bước nhảy ấy. Trong hai biến, ở giao điểm hai đường, cần so vết theo nhiều hướng; triệt tiêu trên một lát cắt không xóa mất kỳ dị hai chiều.

Tại $t=0$, hiệu nguyên hàm bằng $0$ và đạo hàm thời gian cho $H'(x)=g(x)$ trong miền dữ liệu; half-sum cho đúng li độ. Mỗi hạng là sóng của $x\pm at$, nên thỏa PDE ngoài các mặt gãy và theo nghĩa phân phối trên toàn miền. Đồ thị hai sóng ban đầu và các lát cắt thời gian được tính trực tiếp từ các nhánh trên.
\begin{center}\includegraphics[width=.97\linewidth,height=11cm,keepaspectratio]{_xu_ly/graphs/BT_3_2_3_03.png}\end{center}
\textit{Đồ thị minh họa; với nghiệm số, xem phương pháp và tham số trong mã Python kèm theo.}

\clearpage\noindent\texttt{\detokenize{BT_3_2_3_04}}\par
\section*{Lời giải bổ sung — Bài tập 4 · Mục 3.2.3}

\subsection*{Kiến thức và trình tự giải}
\textbf{Bước 1 — Liệt kê các khoảng của dữ liệu.} Ghi đủ điểm đầu, điểm cuối hỗ trợ, điểm gãy và bước nhảy của $f,g$. Nếu có biên, phản xạ chẵn hoặc lẻ trước khi tích phân; trên đoạn hữu hạn tiếp tục phản xạ đến một chu kỳ hoàn chỉnh.

\textbf{Bước 2 — Tính nguyên hàm liên tục.} Lấy $H'=\widetilde g$ trên từng khoảng và chọn hằng số để hai vết của $H$ bằng nhau. Viết $[s]_p^q=\min(q,\max(p,s))$ hoặc $(s)_+=\max(s,0)$ khi cần gộp nhiều nhánh. Một nguyên hàm liên tục tránh tạo ra bước nhảy giả trong nghiệm.

\textbf{Bước 3 — Xác định sóng phải và sóng trái.} Đặt $A(s)=\widetilde f(s)/2-H(s)/(2a)$, $B(s)=\widetilde f(s)/2+H(s)/(2a)$; khi đó $u=A(x-at)+B(x+at)$. Với từng $t$ được yêu cầu, dịch đồ thị $A$ sang phải $at$ và $B$ sang trái $at$, rồi cộng tung độ. Các đường $x-at=b$ và $x+at=b$ chia mặt phẳng thành các miền có công thức đơn giản.

\textbf{Bước 4 — Xét liên tục, đạo hàm và giao của các mặt sóng.} Nếu $f$ có bước nhảy, $u$ thường có bước nhảy; nếu $f$ liên tục và $g$ bị chặn thì $u$ liên tục. Tại thời điểm cố định,


\[
u_x=\frac{\widetilde f'(x-at)+\widetilde f'(x+at)}2+\frac{\widetilde g(x+at)-\widetilde g(x-at)}{2a}.
\]


Tính giới hạn trái/phải của biểu thức này tại từng vị trí ứng viên. Có thể có triệt tiêu khi hai đường gặp nhau, nên không chỉ liệt kê đường rồi khẳng định mọi đạo hàm đều mất. Với tính $C^1$ theo cả $x,t$, còn kiểm tra $u_t$. Với $C^2$, xét tiếp các góc của $g$ và đạo hàm cấp hai của $f$.

\textbf{Bước 5 — Tính thời điểm đến, giới hạn và ngưỡng.} Sóng ở $(x,t)$ chỉ nhận dữ liệu trong $[x-at,x+at]$. Giải bất đẳng thức giao hỗ trợ để tìm thời điểm đến và đi qua. Nếu $H$ hằng ngoài đoạn hỗ trợ, dùng các giá trị hằng này để tính giới hạn khi $t\to\infty$. Nếu hỏi áp suất cực đại, viết hàm từng khoảng thời gian và so sánh cực trị, bước nhảy cùng hai giá trị một phía.
\subsection*{Lời giải áp dụng cho bài}
\textbf{Bước 1. Đọc dữ liệu và chọn phản xạ.} Toàn trục, $a=2$, $f=\chi_{[-2,-1]}$, $g=s\chi_{[0,1]}$.

\textbf{Bước 2. Tính từng nhánh và nối nguyên hàm.} $H=0$ trước $0$, $s^2/2$ trên $[0,1]$, $1/2$ sau $1$. $g$ liên tục tại $0$, có bước nhảy từ $1$ về $0$ tại $1$.
\[f_e(s)=\chi_{[-2,-1]}(s),\qquad H(s)=\tfrac12\operatorname{clip}(s,0,1)^2.\]
Ở đây $\operatorname{clip}(s,p,q)=\min(q,\max(p,s))$; $\operatorname{mod}(s,L)\in[0,L)$ kể cả với $s<0$.

\textbf{Bước 3. Xác định hai sóng chạy và ghép nghiệm.} $A=f_e/2-H/(2a)$, $B=f_e/2+H/(2a)$. Thay hai cận:
\[u(x,t)=\frac{f_e(x-2t)+f_e(x+2t)}2+\frac{H(x+2t)-H(x-2t)}{4}.\]

\textbf{Bước 4. Trả lời các yêu cầu và kiểm tra.} $u$ gián đoạn trên các nhánh $x\pm2t=-2,-1$. Vận tốc chỉ làm gãy $u_x$ trên $x\pm2t=1$; tại $x\pm2t=0$ chỉ gãy đạo hàm cấp hai. Ở $t=1$, các bước nhảy li độ nằm tại $-4,-3,0,1$.

Tại $t=0$, hiệu nguyên hàm bằng $0$ và đạo hàm thời gian cho $H'(x)=g(x)$ trong miền dữ liệu; half-sum cho đúng li độ. Mỗi hạng là sóng của $x\pm at$, nên thỏa PDE ngoài các mặt gãy và theo nghĩa phân phối trên toàn miền. Đồ thị hai sóng ban đầu và các lát cắt thời gian được tính trực tiếp từ các nhánh trên.
\begin{center}\includegraphics[width=.97\linewidth,height=11cm,keepaspectratio]{_xu_ly/graphs/BT_3_2_3_04.png}\end{center}
\textit{Đồ thị minh họa; với nghiệm số, xem phương pháp và tham số trong mã Python kèm theo.}

\clearpage\noindent\texttt{\detokenize{BT_3_2_3_05}}\par
\section*{Lời giải bổ sung — Bài tập 5 · Mục 3.2.3}

\subsection*{Kiến thức và trình tự giải}
\textbf{Bước 1 — Liệt kê các khoảng của dữ liệu.} Ghi đủ điểm đầu, điểm cuối hỗ trợ, điểm gãy và bước nhảy của $f,g$. Nếu có biên, phản xạ chẵn hoặc lẻ trước khi tích phân; trên đoạn hữu hạn tiếp tục phản xạ đến một chu kỳ hoàn chỉnh.

\textbf{Bước 2 — Tính nguyên hàm liên tục.} Lấy $H'=\widetilde g$ trên từng khoảng và chọn hằng số để hai vết của $H$ bằng nhau. Viết $[s]_p^q=\min(q,\max(p,s))$ hoặc $(s)_+=\max(s,0)$ khi cần gộp nhiều nhánh. Một nguyên hàm liên tục tránh tạo ra bước nhảy giả trong nghiệm.

\textbf{Bước 3 — Xác định sóng phải và sóng trái.} Đặt $A(s)=\widetilde f(s)/2-H(s)/(2a)$, $B(s)=\widetilde f(s)/2+H(s)/(2a)$; khi đó $u=A(x-at)+B(x+at)$. Với từng $t$ được yêu cầu, dịch đồ thị $A$ sang phải $at$ và $B$ sang trái $at$, rồi cộng tung độ. Các đường $x-at=b$ và $x+at=b$ chia mặt phẳng thành các miền có công thức đơn giản.

\textbf{Bước 4 — Xét liên tục, đạo hàm và giao của các mặt sóng.} Nếu $f$ có bước nhảy, $u$ thường có bước nhảy; nếu $f$ liên tục và $g$ bị chặn thì $u$ liên tục. Tại thời điểm cố định,


\[
u_x=\frac{\widetilde f'(x-at)+\widetilde f'(x+at)}2+\frac{\widetilde g(x+at)-\widetilde g(x-at)}{2a}.
\]


Tính giới hạn trái/phải của biểu thức này tại từng vị trí ứng viên. Có thể có triệt tiêu khi hai đường gặp nhau, nên không chỉ liệt kê đường rồi khẳng định mọi đạo hàm đều mất. Với tính $C^1$ theo cả $x,t$, còn kiểm tra $u_t$. Với $C^2$, xét tiếp các góc của $g$ và đạo hàm cấp hai của $f$.

\textbf{Bước 5 — Tính thời điểm đến, giới hạn và ngưỡng.} Sóng ở $(x,t)$ chỉ nhận dữ liệu trong $[x-at,x+at]$. Giải bất đẳng thức giao hỗ trợ để tìm thời điểm đến và đi qua. Nếu $H$ hằng ngoài đoạn hỗ trợ, dùng các giá trị hằng này để tính giới hạn khi $t\to\infty$. Nếu hỏi áp suất cực đại, viết hàm từng khoảng thời gian và so sánh cực trị, bước nhảy cùng hai giá trị một phía.
\subsection*{Lời giải áp dụng cho bài}
\textbf{Bước 1. Đọc dữ liệu và chọn phản xạ.} Toàn trục, $a=4$, $f=10\chi_{[-1,1]}$, $g=\chi_{[-1,1]}$.

\textbf{Bước 2. Tính từng nhánh và nối nguyên hàm.} $H(s)=0$ khi $s\le-1$, $s+1$ khi $-1<s<1$, $2$ khi $s\ge1$. Tại $x=10$, sóng phải đến lúc $10-4t=1$ và rời lúc $10-4t=-1$.
\[f_e(s)=10\chi_{[-1,1]}(s),\qquad H(s)=\operatorname{clip}(s+1,0,2).\]
Ở đây $\operatorname{clip}(s,p,q)=\min(q,\max(p,s))$; $\operatorname{mod}(s,L)\in[0,L)$ kể cả với $s<0$.

\textbf{Bước 3. Xác định hai sóng chạy và ghép nghiệm.} $A=f_e/2-H/(2a)$, $B=f_e/2+H/(2a)$. Thay hai cận:
\[u(x,t)=\frac{f_e(x-4t)+f_e(x+4t)}2+\frac{H(x+4t)-H(x-4t)}{8}.\]

\textbf{Bước 4. Trả lời các yêu cầu và kiểm tra.} Đó là $t=9/4$ và $11/4$. Trong khoảng giữa, $u(10,t)=5+(4t-9)/8$; trước đó $0$, sau đó $1/4$. Cực đại $21/4<6$, kể cả vết một phía tại bước nhảy. Không đạt ngưỡng phá hủy.

Tại $t=0$, hiệu nguyên hàm bằng $0$ và đạo hàm thời gian cho $H'(x)=g(x)$ trong miền dữ liệu; half-sum cho đúng li độ. Mỗi hạng là sóng của $x\pm at$, nên thỏa PDE ngoài các mặt gãy và theo nghĩa phân phối trên toàn miền. Đồ thị hai sóng ban đầu và các lát cắt thời gian được tính trực tiếp từ các nhánh trên.
\begin{center}\includegraphics[width=.97\linewidth,height=11cm,keepaspectratio]{_xu_ly/graphs/BT_3_2_3_05.png}\end{center}
\textit{Đồ thị minh họa; với nghiệm số, xem phương pháp và tham số trong mã Python kèm theo.}

\clearpage\noindent\texttt{\detokenize{BT_3_2_3_06}}\par
\section*{Lời giải bổ sung — Bài tập 6 · Mục 3.2.3}

\subsection*{Kiến thức và trình tự giải}
\textbf{Bước 1 — Liệt kê các khoảng của dữ liệu.} Ghi đủ điểm đầu, điểm cuối hỗ trợ, điểm gãy và bước nhảy của $f,g$. Nếu có biên, phản xạ chẵn hoặc lẻ trước khi tích phân; trên đoạn hữu hạn tiếp tục phản xạ đến một chu kỳ hoàn chỉnh.

\textbf{Bước 2 — Tính nguyên hàm liên tục.} Lấy $H'=\widetilde g$ trên từng khoảng và chọn hằng số để hai vết của $H$ bằng nhau. Viết $[s]_p^q=\min(q,\max(p,s))$ hoặc $(s)_+=\max(s,0)$ khi cần gộp nhiều nhánh. Một nguyên hàm liên tục tránh tạo ra bước nhảy giả trong nghiệm.

\textbf{Bước 3 — Xác định sóng phải và sóng trái.} Đặt $A(s)=\widetilde f(s)/2-H(s)/(2a)$, $B(s)=\widetilde f(s)/2+H(s)/(2a)$; khi đó $u=A(x-at)+B(x+at)$. Với từng $t$ được yêu cầu, dịch đồ thị $A$ sang phải $at$ và $B$ sang trái $at$, rồi cộng tung độ. Các đường $x-at=b$ và $x+at=b$ chia mặt phẳng thành các miền có công thức đơn giản.

\textbf{Bước 4 — Xét liên tục, đạo hàm và giao của các mặt sóng.} Nếu $f$ có bước nhảy, $u$ thường có bước nhảy; nếu $f$ liên tục và $g$ bị chặn thì $u$ liên tục. Tại thời điểm cố định,


\[
u_x=\frac{\widetilde f'(x-at)+\widetilde f'(x+at)}2+\frac{\widetilde g(x+at)-\widetilde g(x-at)}{2a}.
\]


Tính giới hạn trái/phải của biểu thức này tại từng vị trí ứng viên. Có thể có triệt tiêu khi hai đường gặp nhau, nên không chỉ liệt kê đường rồi khẳng định mọi đạo hàm đều mất. Với tính $C^1$ theo cả $x,t$, còn kiểm tra $u_t$. Với $C^2$, xét tiếp các góc của $g$ và đạo hàm cấp hai của $f$.

\textbf{Bước 5 — Tính thời điểm đến, giới hạn và ngưỡng.} Sóng ở $(x,t)$ chỉ nhận dữ liệu trong $[x-at,x+at]$. Giải bất đẳng thức giao hỗ trợ để tìm thời điểm đến và đi qua. Nếu $H$ hằng ngoài đoạn hỗ trợ, dùng các giá trị hằng này để tính giới hạn khi $t\to\infty$. Nếu hỏi áp suất cực đại, viết hàm từng khoảng thời gian và so sánh cực trị, bước nhảy cùng hai giá trị một phía.

\textbf{Bước 1 — Đọc đúng điều kiện biên.} Biên $u(0,t)=0$ dùng mở rộng lẻ; biên $u_x(0,t)=0$ dùng mở rộng chẵn. Với dữ liệu trên $x\ge0$, đặt $\widetilde f(-s)=-f(s)$ hoặc $f(s)$ tương ứng, làm tương tự cho $g$. Mở rộng lẻ là $C^2$ nếu $f(0)=f''(0)=0$, và mở rộng lẻ của $g$ là $C^1$ nếu $g(0)=0$. Mở rộng chẵn đòi $f'(0)=g'(0)=0$ cho mức trơn ấy.

\textbf{Bước 2 — Lập nguyên hàm sau khi phản xạ.} Đặt $H(s)=\int_0^s\widetilde g(r)dr$. Nếu $\widetilde g$ lẻ thì $H$ chẵn; nếu $\widetilde g$ chẵn thì $H$ lẻ. Chính nguyên hàm, chứ không phải vận tốc, xuất hiện trong


\[
u=\frac{\widetilde f(x-at)+\widetilde f(x+at)}2+\frac{H(x+at)-H(x-at)}{2a}.
\]


\textbf{Bước 3 — Tách miền trước và sau khi chạm biên.} Với $u(0,t)=h(t)$, ở $x\ge at$ dùng d'Alembert với dữ liệu gốc. Ở $0<x<at$:


\[
u=h(t-x/a)+\frac{f(x+at)-f(at-x)}2+\frac1{2a}\int_{at-x}^{at+x}g(s)ds.
\]


Với $u_x(0,t)=h(t)$, dùng nghiệm của các mở rộng chẵn rồi trừ $a\int_0^{t-x/a}h(s)ds$ trên miền $x<at$. Tính riêng đạo hàm theo $x$ của hạng này để kiểm tra biên.

\textbf{Bước 4 — Trên đoạn hữu hạn, phản xạ ở cả hai đầu.} Dirichlet ở hai đầu cho mở rộng lẻ tại $0,L$, chu kỳ không gian $2L$. Neumann ở hai đầu cho mở rộng chẵn tại $0,L$, cũng chu kỳ $2L$. Một đầu lẻ, một đầu chẵn cho $\widetilde f(s+2L)=-\widetilde f(s)$ và chu kỳ $4L$, tương tự $g$. Vì thế chu kỳ thời gian thường là $2L/a$ hoặc $4L/a$. Với Neumann hai đầu còn phải xét trung bình $g$: trung bình khác không sinh hạng tuyến tính theo $t$, nên nghiệm không tuần hoàn.

\textbf{Bước 5 — Kiểm tra góc và đường ghép.} Biên Dirichlet không thuần nhất cần $h(0)=f(0)$, $h'(0)=g(0)$, $h''(0)=a^2f''(0)$ để có nghiệm $C^2$ cho phương trình thuần nhất. So sánh hai vết trên $x=at$, gồm $u,u_x,u_t$ rồi đạo hàm cấp hai. Nếu không tương thích, nêu rõ mức trơn còn lại thay vì khẳng định nghiệm cổ điển trên toàn miền.
\subsection*{Lời giải áp dụng cho bài}
\textbf{Bước 1. Đọc dữ liệu và chọn phản xạ.} Nửa trục Dirichlet, $a=3$, $f=s(s-1)$ trên $[0,1]$, $g=6$.

\textbf{Bước 2. Tính từng nhánh và nối nguyên hàm.} Mở rộng lẻ của $f$ là $\operatorname{sgn}(s)|s|(|s|-1)\chi_{[0,1]}(|s|)$. Mở rộng lẻ của $g$ là $6\operatorname{sgn}(s)$, nên $H=6|s|$.
\[f_e(s)=\operatorname{sgn}(s)|s|(|s|-1)\chi_{[0,1]}(|s|),\qquad H(s)=6|s|.\]
Ở đây $\operatorname{clip}(s,p,q)=\min(q,\max(p,s))$; $\operatorname{mod}(s,L)\in[0,L)$ kể cả với $s<0$.

\textbf{Bước 3. Xác định hai sóng chạy và ghép nghiệm.} $A=f_e/2-H/(2a)$, $B=f_e/2+H/(2a)$. Thay hai cận:
\[u(x,t)=\frac{f_e(x-3t)+f_e(x+3t)}2+\frac{H(x+3t)-H(x-3t)}{6}.\]

\textbf{Bước 4. Trả lời các yêu cầu và kiểm tra.} Ở $t=1/3$, $x-1=0$ cho bước nhảy của $g$, tức $x=1$. Góc của $f$ tại $s=1$ cho $x=2$; các góc khác nằm ngoài $x>0$. $f$ nối đạo hàm ở $0$, nên không khử bước nhảy vận tốc tại $x=1$.

Tại $t=0$, hiệu nguyên hàm bằng $0$ và đạo hàm thời gian cho $H'(x)=g(x)$ trong miền dữ liệu; half-sum cho đúng li độ. Mỗi hạng là sóng của $x\pm at$, nên thỏa PDE ngoài các mặt gãy và theo nghĩa phân phối trên toàn miền. Đồ thị hai sóng ban đầu và các lát cắt thời gian được tính trực tiếp từ các nhánh trên.
\begin{center}\includegraphics[width=.97\linewidth,height=11cm,keepaspectratio]{_xu_ly/graphs/BT_3_2_3_06.png}\end{center}
\textit{Đồ thị minh họa; với nghiệm số, xem phương pháp và tham số trong mã Python kèm theo.}

\clearpage\noindent\texttt{\detokenize{BT_3_2_3_07}}\par
\section*{Lời giải bổ sung — Bài tập 7 · Mục 3.2.3}

\subsection*{Kiến thức và trình tự giải}
\textbf{Bước 1 — Liệt kê các khoảng của dữ liệu.} Ghi đủ điểm đầu, điểm cuối hỗ trợ, điểm gãy và bước nhảy của $f,g$. Nếu có biên, phản xạ chẵn hoặc lẻ trước khi tích phân; trên đoạn hữu hạn tiếp tục phản xạ đến một chu kỳ hoàn chỉnh.

\textbf{Bước 2 — Tính nguyên hàm liên tục.} Lấy $H'=\widetilde g$ trên từng khoảng và chọn hằng số để hai vết của $H$ bằng nhau. Viết $[s]_p^q=\min(q,\max(p,s))$ hoặc $(s)_+=\max(s,0)$ khi cần gộp nhiều nhánh. Một nguyên hàm liên tục tránh tạo ra bước nhảy giả trong nghiệm.

\textbf{Bước 3 — Xác định sóng phải và sóng trái.} Đặt $A(s)=\widetilde f(s)/2-H(s)/(2a)$, $B(s)=\widetilde f(s)/2+H(s)/(2a)$; khi đó $u=A(x-at)+B(x+at)$. Với từng $t$ được yêu cầu, dịch đồ thị $A$ sang phải $at$ và $B$ sang trái $at$, rồi cộng tung độ. Các đường $x-at=b$ và $x+at=b$ chia mặt phẳng thành các miền có công thức đơn giản.

\textbf{Bước 4 — Xét liên tục, đạo hàm và giao của các mặt sóng.} Nếu $f$ có bước nhảy, $u$ thường có bước nhảy; nếu $f$ liên tục và $g$ bị chặn thì $u$ liên tục. Tại thời điểm cố định,


\[
u_x=\frac{\widetilde f'(x-at)+\widetilde f'(x+at)}2+\frac{\widetilde g(x+at)-\widetilde g(x-at)}{2a}.
\]


Tính giới hạn trái/phải của biểu thức này tại từng vị trí ứng viên. Có thể có triệt tiêu khi hai đường gặp nhau, nên không chỉ liệt kê đường rồi khẳng định mọi đạo hàm đều mất. Với tính $C^1$ theo cả $x,t$, còn kiểm tra $u_t$. Với $C^2$, xét tiếp các góc của $g$ và đạo hàm cấp hai của $f$.

\textbf{Bước 5 — Tính thời điểm đến, giới hạn và ngưỡng.} Sóng ở $(x,t)$ chỉ nhận dữ liệu trong $[x-at,x+at]$. Giải bất đẳng thức giao hỗ trợ để tìm thời điểm đến và đi qua. Nếu $H$ hằng ngoài đoạn hỗ trợ, dùng các giá trị hằng này để tính giới hạn khi $t\to\infty$. Nếu hỏi áp suất cực đại, viết hàm từng khoảng thời gian và so sánh cực trị, bước nhảy cùng hai giá trị một phía.

\textbf{Bước 1 — Đọc đúng điều kiện biên.} Biên $u(0,t)=0$ dùng mở rộng lẻ; biên $u_x(0,t)=0$ dùng mở rộng chẵn. Với dữ liệu trên $x\ge0$, đặt $\widetilde f(-s)=-f(s)$ hoặc $f(s)$ tương ứng, làm tương tự cho $g$. Mở rộng lẻ là $C^2$ nếu $f(0)=f''(0)=0$, và mở rộng lẻ của $g$ là $C^1$ nếu $g(0)=0$. Mở rộng chẵn đòi $f'(0)=g'(0)=0$ cho mức trơn ấy.

\textbf{Bước 2 — Lập nguyên hàm sau khi phản xạ.} Đặt $H(s)=\int_0^s\widetilde g(r)dr$. Nếu $\widetilde g$ lẻ thì $H$ chẵn; nếu $\widetilde g$ chẵn thì $H$ lẻ. Chính nguyên hàm, chứ không phải vận tốc, xuất hiện trong


\[
u=\frac{\widetilde f(x-at)+\widetilde f(x+at)}2+\frac{H(x+at)-H(x-at)}{2a}.
\]


\textbf{Bước 3 — Tách miền trước và sau khi chạm biên.} Với $u(0,t)=h(t)$, ở $x\ge at$ dùng d'Alembert với dữ liệu gốc. Ở $0<x<at$:


\[
u=h(t-x/a)+\frac{f(x+at)-f(at-x)}2+\frac1{2a}\int_{at-x}^{at+x}g(s)ds.
\]


Với $u_x(0,t)=h(t)$, dùng nghiệm của các mở rộng chẵn rồi trừ $a\int_0^{t-x/a}h(s)ds$ trên miền $x<at$. Tính riêng đạo hàm theo $x$ của hạng này để kiểm tra biên.

\textbf{Bước 4 — Trên đoạn hữu hạn, phản xạ ở cả hai đầu.} Dirichlet ở hai đầu cho mở rộng lẻ tại $0,L$, chu kỳ không gian $2L$. Neumann ở hai đầu cho mở rộng chẵn tại $0,L$, cũng chu kỳ $2L$. Một đầu lẻ, một đầu chẵn cho $\widetilde f(s+2L)=-\widetilde f(s)$ và chu kỳ $4L$, tương tự $g$. Vì thế chu kỳ thời gian thường là $2L/a$ hoặc $4L/a$. Với Neumann hai đầu còn phải xét trung bình $g$: trung bình khác không sinh hạng tuyến tính theo $t$, nên nghiệm không tuần hoàn.

\textbf{Bước 5 — Kiểm tra góc và đường ghép.} Biên Dirichlet không thuần nhất cần $h(0)=f(0)$, $h'(0)=g(0)$, $h''(0)=a^2f''(0)$ để có nghiệm $C^2$ cho phương trình thuần nhất. So sánh hai vết trên $x=at$, gồm $u,u_x,u_t$ rồi đạo hàm cấp hai. Nếu không tương thích, nêu rõ mức trơn còn lại thay vì khẳng định nghiệm cổ điển trên toàn miền.
\subsection*{Lời giải áp dụng cho bài}
\textbf{Bước 1. Đọc dữ liệu và chọn phản xạ.} Nửa trục Neumann, $a=1$, $f=0$, $g=s-1$ trên $[0,1]$, bằng không ngoài đoạn.

\textbf{Bước 2. Tính từng nhánh và nối nguyên hàm.} Mở rộng chẵn $\widetilde g=(|s|-1)\chi_{[0,1]}(|s|)$. Tích phân từ $0$ cho $s^2/2-s$ ở $s\in[0,1]$, sau đó giữ $-1/2$; phản xạ nguyên hàm lẻ.
\[f_e(s)=0,\qquad H(s)=\operatorname{sgn}(s)(m^2/2-m),\quad m=\min(|s|,1).\]
Ở đây $\operatorname{clip}(s,p,q)=\min(q,\max(p,s))$; $\operatorname{mod}(s,L)\in[0,L)$ kể cả với $s<0$.

\textbf{Bước 3. Xác định hai sóng chạy và ghép nghiệm.} $A=f_e/2-H/(2a)$, $B=f_e/2+H/(2a)$. Thay hai cận:
\[u(x,t)=\frac{f_e(x-1t)+f_e(x+1t)}2+\frac{H(x+1t)-H(x-1t)}{2}.\]

\textbf{Bước 4. Trả lời các yêu cầu và kiểm tra.} Vì $\widetilde g$ liên tục cả ở $0,\pm1$, $u_x$ tồn tại ở mọi $x$ tại $t=1$. Khi $t\ge2$ tại $x=1$, hai cận đã ở hai phía ngoài hỗ trợ, $H(1+t)=-1/2$, $H(1-t)=1/2$, nên $u=-1/2$.

Tại $t=0$, hiệu nguyên hàm bằng $0$ và đạo hàm thời gian cho $H'(x)=g(x)$ trong miền dữ liệu; half-sum cho đúng li độ. Mỗi hạng là sóng của $x\pm at$, nên thỏa PDE ngoài các mặt gãy và theo nghĩa phân phối trên toàn miền. Đồ thị hai sóng ban đầu và các lát cắt thời gian được tính trực tiếp từ các nhánh trên.
\begin{center}\includegraphics[width=.97\linewidth,height=11cm,keepaspectratio]{_xu_ly/graphs/BT_3_2_3_07.png}\end{center}
\textit{Đồ thị minh họa; với nghiệm số, xem phương pháp và tham số trong mã Python kèm theo.}

\clearpage\noindent\texttt{\detokenize{BT_3_2_3_08}}\par
\section*{Lời giải bổ sung — Bài tập 8 · Mục 3.2.3}

\subsection*{Kiến thức và trình tự giải}
\textbf{Bước 1 — Liệt kê các khoảng của dữ liệu.} Ghi đủ điểm đầu, điểm cuối hỗ trợ, điểm gãy và bước nhảy của $f,g$. Nếu có biên, phản xạ chẵn hoặc lẻ trước khi tích phân; trên đoạn hữu hạn tiếp tục phản xạ đến một chu kỳ hoàn chỉnh.

\textbf{Bước 2 — Tính nguyên hàm liên tục.} Lấy $H'=\widetilde g$ trên từng khoảng và chọn hằng số để hai vết của $H$ bằng nhau. Viết $[s]_p^q=\min(q,\max(p,s))$ hoặc $(s)_+=\max(s,0)$ khi cần gộp nhiều nhánh. Một nguyên hàm liên tục tránh tạo ra bước nhảy giả trong nghiệm.

\textbf{Bước 3 — Xác định sóng phải và sóng trái.} Đặt $A(s)=\widetilde f(s)/2-H(s)/(2a)$, $B(s)=\widetilde f(s)/2+H(s)/(2a)$; khi đó $u=A(x-at)+B(x+at)$. Với từng $t$ được yêu cầu, dịch đồ thị $A$ sang phải $at$ và $B$ sang trái $at$, rồi cộng tung độ. Các đường $x-at=b$ và $x+at=b$ chia mặt phẳng thành các miền có công thức đơn giản.

\textbf{Bước 4 — Xét liên tục, đạo hàm và giao của các mặt sóng.} Nếu $f$ có bước nhảy, $u$ thường có bước nhảy; nếu $f$ liên tục và $g$ bị chặn thì $u$ liên tục. Tại thời điểm cố định,


\[
u_x=\frac{\widetilde f'(x-at)+\widetilde f'(x+at)}2+\frac{\widetilde g(x+at)-\widetilde g(x-at)}{2a}.
\]


Tính giới hạn trái/phải của biểu thức này tại từng vị trí ứng viên. Có thể có triệt tiêu khi hai đường gặp nhau, nên không chỉ liệt kê đường rồi khẳng định mọi đạo hàm đều mất. Với tính $C^1$ theo cả $x,t$, còn kiểm tra $u_t$. Với $C^2$, xét tiếp các góc của $g$ và đạo hàm cấp hai của $f$.

\textbf{Bước 5 — Tính thời điểm đến, giới hạn và ngưỡng.} Sóng ở $(x,t)$ chỉ nhận dữ liệu trong $[x-at,x+at]$. Giải bất đẳng thức giao hỗ trợ để tìm thời điểm đến và đi qua. Nếu $H$ hằng ngoài đoạn hỗ trợ, dùng các giá trị hằng này để tính giới hạn khi $t\to\infty$. Nếu hỏi áp suất cực đại, viết hàm từng khoảng thời gian và so sánh cực trị, bước nhảy cùng hai giá trị một phía.

\textbf{Bước 1 — Đọc đúng điều kiện biên.} Biên $u(0,t)=0$ dùng mở rộng lẻ; biên $u_x(0,t)=0$ dùng mở rộng chẵn. Với dữ liệu trên $x\ge0$, đặt $\widetilde f(-s)=-f(s)$ hoặc $f(s)$ tương ứng, làm tương tự cho $g$. Mở rộng lẻ là $C^2$ nếu $f(0)=f''(0)=0$, và mở rộng lẻ của $g$ là $C^1$ nếu $g(0)=0$. Mở rộng chẵn đòi $f'(0)=g'(0)=0$ cho mức trơn ấy.

\textbf{Bước 2 — Lập nguyên hàm sau khi phản xạ.} Đặt $H(s)=\int_0^s\widetilde g(r)dr$. Nếu $\widetilde g$ lẻ thì $H$ chẵn; nếu $\widetilde g$ chẵn thì $H$ lẻ. Chính nguyên hàm, chứ không phải vận tốc, xuất hiện trong


\[
u=\frac{\widetilde f(x-at)+\widetilde f(x+at)}2+\frac{H(x+at)-H(x-at)}{2a}.
\]


\textbf{Bước 3 — Tách miền trước và sau khi chạm biên.} Với $u(0,t)=h(t)$, ở $x\ge at$ dùng d'Alembert với dữ liệu gốc. Ở $0<x<at$:


\[
u=h(t-x/a)+\frac{f(x+at)-f(at-x)}2+\frac1{2a}\int_{at-x}^{at+x}g(s)ds.
\]


Với $u_x(0,t)=h(t)$, dùng nghiệm của các mở rộng chẵn rồi trừ $a\int_0^{t-x/a}h(s)ds$ trên miền $x<at$. Tính riêng đạo hàm theo $x$ của hạng này để kiểm tra biên.

\textbf{Bước 4 — Trên đoạn hữu hạn, phản xạ ở cả hai đầu.} Dirichlet ở hai đầu cho mở rộng lẻ tại $0,L$, chu kỳ không gian $2L$. Neumann ở hai đầu cho mở rộng chẵn tại $0,L$, cũng chu kỳ $2L$. Một đầu lẻ, một đầu chẵn cho $\widetilde f(s+2L)=-\widetilde f(s)$ và chu kỳ $4L$, tương tự $g$. Vì thế chu kỳ thời gian thường là $2L/a$ hoặc $4L/a$. Với Neumann hai đầu còn phải xét trung bình $g$: trung bình khác không sinh hạng tuyến tính theo $t$, nên nghiệm không tuần hoàn.

\textbf{Bước 5 — Kiểm tra góc và đường ghép.} Biên Dirichlet không thuần nhất cần $h(0)=f(0)$, $h'(0)=g(0)$, $h''(0)=a^2f''(0)$ để có nghiệm $C^2$ cho phương trình thuần nhất. So sánh hai vết trên $x=at$, gồm $u,u_x,u_t$ rồi đạo hàm cấp hai. Nếu không tương thích, nêu rõ mức trơn còn lại thay vì khẳng định nghiệm cổ điển trên toàn miền.
\subsection*{Lời giải áp dụng cho bài}
\textbf{Bước 1. Đọc dữ liệu và chọn phản xạ.} Nửa trục Neumann, $a=3$, $f=s(2-s)$ trên $[0,2]$, $g=6$.

\textbf{Bước 2. Tính từng nhánh và nối nguyên hàm.} Mở rộng chẵn $f=|s|(2-|s|)\chi_{[0,2]}(|s|)$ và $g=6$, nên $H=6s$. Hiệu $H$ góp đúng $6t$.
\[f_e(s)=|s|(2-|s|)\chi_{[0,2]}(|s|),\qquad H(s)=6s.\]
Ở đây $\operatorname{clip}(s,p,q)=\min(q,\max(p,s))$; $\operatorname{mod}(s,L)\in[0,L)$ kể cả với $s<0$.

\textbf{Bước 3. Xác định hai sóng chạy và ghép nghiệm.} $A=f_e/2-H/(2a)$, $B=f_e/2+H/(2a)$. Thay hai cận:
\[u(x,t)=\frac{f_e(x-3t)+f_e(x+3t)}2+\frac{H(x+3t)-H(x-3t)}{6}.\]

\textbf{Bước 4. Trả lời các yêu cầu và kiểm tra.} Ở $t=1/3$, $x-1=0$ là góc của mở rộng chẵn $f$, cho $x=1$; $x-1=2$ cho $x=3$. Các góc ở cận phải không cho điểm dương khác. Sau $t>1$ tại $x=1$, phần li độ hết, $u=6t\to\infty$.

Tại $t=0$, hiệu nguyên hàm bằng $0$ và đạo hàm thời gian cho $H'(x)=g(x)$ trong miền dữ liệu; half-sum cho đúng li độ. Mỗi hạng là sóng của $x\pm at$, nên thỏa PDE ngoài các mặt gãy và theo nghĩa phân phối trên toàn miền. Đồ thị hai sóng ban đầu và các lát cắt thời gian được tính trực tiếp từ các nhánh trên.
\begin{center}\includegraphics[width=.97\linewidth,height=11cm,keepaspectratio]{_xu_ly/graphs/BT_3_2_3_08.png}\end{center}
\textit{Đồ thị minh họa; với nghiệm số, xem phương pháp và tham số trong mã Python kèm theo.}

\clearpage\noindent\texttt{\detokenize{BT_3_2_3_09}}\par
\section*{Lời giải bổ sung — Bài tập 9 · Mục 3.2.3}

\subsection*{Kiến thức và trình tự giải}
\textbf{Bước 1 — Liệt kê các khoảng của dữ liệu.} Ghi đủ điểm đầu, điểm cuối hỗ trợ, điểm gãy và bước nhảy của $f,g$. Nếu có biên, phản xạ chẵn hoặc lẻ trước khi tích phân; trên đoạn hữu hạn tiếp tục phản xạ đến một chu kỳ hoàn chỉnh.

\textbf{Bước 2 — Tính nguyên hàm liên tục.} Lấy $H'=\widetilde g$ trên từng khoảng và chọn hằng số để hai vết của $H$ bằng nhau. Viết $[s]_p^q=\min(q,\max(p,s))$ hoặc $(s)_+=\max(s,0)$ khi cần gộp nhiều nhánh. Một nguyên hàm liên tục tránh tạo ra bước nhảy giả trong nghiệm.

\textbf{Bước 3 — Xác định sóng phải và sóng trái.} Đặt $A(s)=\widetilde f(s)/2-H(s)/(2a)$, $B(s)=\widetilde f(s)/2+H(s)/(2a)$; khi đó $u=A(x-at)+B(x+at)$. Với từng $t$ được yêu cầu, dịch đồ thị $A$ sang phải $at$ và $B$ sang trái $at$, rồi cộng tung độ. Các đường $x-at=b$ và $x+at=b$ chia mặt phẳng thành các miền có công thức đơn giản.

\textbf{Bước 4 — Xét liên tục, đạo hàm và giao của các mặt sóng.} Nếu $f$ có bước nhảy, $u$ thường có bước nhảy; nếu $f$ liên tục và $g$ bị chặn thì $u$ liên tục. Tại thời điểm cố định,


\[
u_x=\frac{\widetilde f'(x-at)+\widetilde f'(x+at)}2+\frac{\widetilde g(x+at)-\widetilde g(x-at)}{2a}.
\]


Tính giới hạn trái/phải của biểu thức này tại từng vị trí ứng viên. Có thể có triệt tiêu khi hai đường gặp nhau, nên không chỉ liệt kê đường rồi khẳng định mọi đạo hàm đều mất. Với tính $C^1$ theo cả $x,t$, còn kiểm tra $u_t$. Với $C^2$, xét tiếp các góc của $g$ và đạo hàm cấp hai của $f$.

\textbf{Bước 5 — Tính thời điểm đến, giới hạn và ngưỡng.} Sóng ở $(x,t)$ chỉ nhận dữ liệu trong $[x-at,x+at]$. Giải bất đẳng thức giao hỗ trợ để tìm thời điểm đến và đi qua. Nếu $H$ hằng ngoài đoạn hỗ trợ, dùng các giá trị hằng này để tính giới hạn khi $t\to\infty$. Nếu hỏi áp suất cực đại, viết hàm từng khoảng thời gian và so sánh cực trị, bước nhảy cùng hai giá trị một phía.

\textbf{Bước 1 — Đọc đúng điều kiện biên.} Biên $u(0,t)=0$ dùng mở rộng lẻ; biên $u_x(0,t)=0$ dùng mở rộng chẵn. Với dữ liệu trên $x\ge0$, đặt $\widetilde f(-s)=-f(s)$ hoặc $f(s)$ tương ứng, làm tương tự cho $g$. Mở rộng lẻ là $C^2$ nếu $f(0)=f''(0)=0$, và mở rộng lẻ của $g$ là $C^1$ nếu $g(0)=0$. Mở rộng chẵn đòi $f'(0)=g'(0)=0$ cho mức trơn ấy.

\textbf{Bước 2 — Lập nguyên hàm sau khi phản xạ.} Đặt $H(s)=\int_0^s\widetilde g(r)dr$. Nếu $\widetilde g$ lẻ thì $H$ chẵn; nếu $\widetilde g$ chẵn thì $H$ lẻ. Chính nguyên hàm, chứ không phải vận tốc, xuất hiện trong


\[
u=\frac{\widetilde f(x-at)+\widetilde f(x+at)}2+\frac{H(x+at)-H(x-at)}{2a}.
\]


\textbf{Bước 3 — Tách miền trước và sau khi chạm biên.} Với $u(0,t)=h(t)$, ở $x\ge at$ dùng d'Alembert với dữ liệu gốc. Ở $0<x<at$:


\[
u=h(t-x/a)+\frac{f(x+at)-f(at-x)}2+\frac1{2a}\int_{at-x}^{at+x}g(s)ds.
\]


Với $u_x(0,t)=h(t)$, dùng nghiệm của các mở rộng chẵn rồi trừ $a\int_0^{t-x/a}h(s)ds$ trên miền $x<at$. Tính riêng đạo hàm theo $x$ của hạng này để kiểm tra biên.

\textbf{Bước 4 — Trên đoạn hữu hạn, phản xạ ở cả hai đầu.} Dirichlet ở hai đầu cho mở rộng lẻ tại $0,L$, chu kỳ không gian $2L$. Neumann ở hai đầu cho mở rộng chẵn tại $0,L$, cũng chu kỳ $2L$. Một đầu lẻ, một đầu chẵn cho $\widetilde f(s+2L)=-\widetilde f(s)$ và chu kỳ $4L$, tương tự $g$. Vì thế chu kỳ thời gian thường là $2L/a$ hoặc $4L/a$. Với Neumann hai đầu còn phải xét trung bình $g$: trung bình khác không sinh hạng tuyến tính theo $t$, nên nghiệm không tuần hoàn.

\textbf{Bước 5 — Kiểm tra góc và đường ghép.} Biên Dirichlet không thuần nhất cần $h(0)=f(0)$, $h'(0)=g(0)$, $h''(0)=a^2f''(0)$ để có nghiệm $C^2$ cho phương trình thuần nhất. So sánh hai vết trên $x=at$, gồm $u,u_x,u_t$ rồi đạo hàm cấp hai. Nếu không tương thích, nêu rõ mức trơn còn lại thay vì khẳng định nghiệm cổ điển trên toàn miền.
\subsection*{Lời giải áp dụng cho bài}
\textbf{Bước 1. Đọc dữ liệu và chọn phản xạ.} Nửa trục Dirichlet, $a=1$, $f=0$, $g=2s-2$ trên $[0,2]$.

\textbf{Bước 2. Tính từng nhánh và nối nguyên hàm.} Mở rộng lẻ là $\operatorname{sgn}(s)(2|s|-2)\chi_{[0,2]}(|s|)$. Tích phân dương cho $s^2-2s$; tại $2$ bằng $0$, sau đó giữ $0$. Nguyên hàm chẵn.
\[f_e(s)=0,\qquad H(s)=m^2-2m,\quad m=\min(|s|,2).\]
Ở đây $\operatorname{clip}(s,p,q)=\min(q,\max(p,s))$; $\operatorname{mod}(s,L)\in[0,L)$ kể cả với $s<0$.

\textbf{Bước 3. Xác định hai sóng chạy và ghép nghiệm.} $A=f_e/2-H/(2a)$, $B=f_e/2+H/(2a)$. Thay hai cận:
\[u(x,t)=\frac{f_e(x-1t)+f_e(x+1t)}2+\frac{H(x+1t)-H(x-1t)}{2}.\]

\textbf{Bước 4. Trả lời các yêu cầu và kiểm tra.} Tại $t=1$, các cận gặp bước nhảy ở $0,2$ tại $x=1,3$; các bước nhảy khác không cho vị trí dương mới. Sau $t\ge3$ ở $x=1$, cả hai giá trị nguyên hàm đều $0$, nên giới hạn bằng $0$.

Tại $t=0$, hiệu nguyên hàm bằng $0$ và đạo hàm thời gian cho $H'(x)=g(x)$ trong miền dữ liệu; half-sum cho đúng li độ. Mỗi hạng là sóng của $x\pm at$, nên thỏa PDE ngoài các mặt gãy và theo nghĩa phân phối trên toàn miền. Đồ thị hai sóng ban đầu và các lát cắt thời gian được tính trực tiếp từ các nhánh trên.
\begin{center}\includegraphics[width=.97\linewidth,height=11cm,keepaspectratio]{_xu_ly/graphs/BT_3_2_3_09.png}\end{center}
\textit{Đồ thị minh họa; với nghiệm số, xem phương pháp và tham số trong mã Python kèm theo.}

\clearpage\noindent\texttt{\detokenize{BT_3_2_3_10}}\par
\section*{Lời giải bổ sung — Bài tập 10 · Mục 3.2.3}

\subsection*{Kiến thức và trình tự giải}
\textbf{Bước 1 — Liệt kê các khoảng của dữ liệu.} Ghi đủ điểm đầu, điểm cuối hỗ trợ, điểm gãy và bước nhảy của $f,g$. Nếu có biên, phản xạ chẵn hoặc lẻ trước khi tích phân; trên đoạn hữu hạn tiếp tục phản xạ đến một chu kỳ hoàn chỉnh.

\textbf{Bước 2 — Tính nguyên hàm liên tục.} Lấy $H'=\widetilde g$ trên từng khoảng và chọn hằng số để hai vết của $H$ bằng nhau. Viết $[s]_p^q=\min(q,\max(p,s))$ hoặc $(s)_+=\max(s,0)$ khi cần gộp nhiều nhánh. Một nguyên hàm liên tục tránh tạo ra bước nhảy giả trong nghiệm.

\textbf{Bước 3 — Xác định sóng phải và sóng trái.} Đặt $A(s)=\widetilde f(s)/2-H(s)/(2a)$, $B(s)=\widetilde f(s)/2+H(s)/(2a)$; khi đó $u=A(x-at)+B(x+at)$. Với từng $t$ được yêu cầu, dịch đồ thị $A$ sang phải $at$ và $B$ sang trái $at$, rồi cộng tung độ. Các đường $x-at=b$ và $x+at=b$ chia mặt phẳng thành các miền có công thức đơn giản.

\textbf{Bước 4 — Xét liên tục, đạo hàm và giao của các mặt sóng.} Nếu $f$ có bước nhảy, $u$ thường có bước nhảy; nếu $f$ liên tục và $g$ bị chặn thì $u$ liên tục. Tại thời điểm cố định,


\[
u_x=\frac{\widetilde f'(x-at)+\widetilde f'(x+at)}2+\frac{\widetilde g(x+at)-\widetilde g(x-at)}{2a}.
\]


Tính giới hạn trái/phải của biểu thức này tại từng vị trí ứng viên. Có thể có triệt tiêu khi hai đường gặp nhau, nên không chỉ liệt kê đường rồi khẳng định mọi đạo hàm đều mất. Với tính $C^1$ theo cả $x,t$, còn kiểm tra $u_t$. Với $C^2$, xét tiếp các góc của $g$ và đạo hàm cấp hai của $f$.

\textbf{Bước 5 — Tính thời điểm đến, giới hạn và ngưỡng.} Sóng ở $(x,t)$ chỉ nhận dữ liệu trong $[x-at,x+at]$. Giải bất đẳng thức giao hỗ trợ để tìm thời điểm đến và đi qua. Nếu $H$ hằng ngoài đoạn hỗ trợ, dùng các giá trị hằng này để tính giới hạn khi $t\to\infty$. Nếu hỏi áp suất cực đại, viết hàm từng khoảng thời gian và so sánh cực trị, bước nhảy cùng hai giá trị một phía.

\textbf{Bước 1 — Đọc đúng điều kiện biên.} Biên $u(0,t)=0$ dùng mở rộng lẻ; biên $u_x(0,t)=0$ dùng mở rộng chẵn. Với dữ liệu trên $x\ge0$, đặt $\widetilde f(-s)=-f(s)$ hoặc $f(s)$ tương ứng, làm tương tự cho $g$. Mở rộng lẻ là $C^2$ nếu $f(0)=f''(0)=0$, và mở rộng lẻ của $g$ là $C^1$ nếu $g(0)=0$. Mở rộng chẵn đòi $f'(0)=g'(0)=0$ cho mức trơn ấy.

\textbf{Bước 2 — Lập nguyên hàm sau khi phản xạ.} Đặt $H(s)=\int_0^s\widetilde g(r)dr$. Nếu $\widetilde g$ lẻ thì $H$ chẵn; nếu $\widetilde g$ chẵn thì $H$ lẻ. Chính nguyên hàm, chứ không phải vận tốc, xuất hiện trong


\[
u=\frac{\widetilde f(x-at)+\widetilde f(x+at)}2+\frac{H(x+at)-H(x-at)}{2a}.
\]


\textbf{Bước 3 — Tách miền trước và sau khi chạm biên.} Với $u(0,t)=h(t)$, ở $x\ge at$ dùng d'Alembert với dữ liệu gốc. Ở $0<x<at$:


\[
u=h(t-x/a)+\frac{f(x+at)-f(at-x)}2+\frac1{2a}\int_{at-x}^{at+x}g(s)ds.
\]


Với $u_x(0,t)=h(t)$, dùng nghiệm của các mở rộng chẵn rồi trừ $a\int_0^{t-x/a}h(s)ds$ trên miền $x<at$. Tính riêng đạo hàm theo $x$ của hạng này để kiểm tra biên.

\textbf{Bước 4 — Trên đoạn hữu hạn, phản xạ ở cả hai đầu.} Dirichlet ở hai đầu cho mở rộng lẻ tại $0,L$, chu kỳ không gian $2L$. Neumann ở hai đầu cho mở rộng chẵn tại $0,L$, cũng chu kỳ $2L$. Một đầu lẻ, một đầu chẵn cho $\widetilde f(s+2L)=-\widetilde f(s)$ và chu kỳ $4L$, tương tự $g$. Vì thế chu kỳ thời gian thường là $2L/a$ hoặc $4L/a$. Với Neumann hai đầu còn phải xét trung bình $g$: trung bình khác không sinh hạng tuyến tính theo $t$, nên nghiệm không tuần hoàn.

\textbf{Bước 5 — Kiểm tra góc và đường ghép.} Biên Dirichlet không thuần nhất cần $h(0)=f(0)$, $h'(0)=g(0)$, $h''(0)=a^2f''(0)$ để có nghiệm $C^2$ cho phương trình thuần nhất. So sánh hai vết trên $x=at$, gồm $u,u_x,u_t$ rồi đạo hàm cấp hai. Nếu không tương thích, nêu rõ mức trơn còn lại thay vì khẳng định nghiệm cổ điển trên toàn miền.
\subsection*{Lời giải áp dụng cho bài}
\textbf{Bước 1. Đọc dữ liệu và chọn phản xạ.} Nửa trục Neumann, $a=1$, $f=0$, $g=2-s$ trên $[0,2]$.

\textbf{Bước 2. Tính từng nhánh và nối nguyên hàm.} Mở rộng chẵn là $(2-|s|)\chi_{[0,2]}(|s|)$. Tích phân dương cho $2s-s^2/2$, tại $2$ bằng $2$; nguyên hàm lẻ có hai đuôi $\pm2$.
\[f_e(s)=0,\qquad H(s)=\operatorname{sgn}(s)(2m-m^2/2),\quad m=\min(|s|,2).\]
Ở đây $\operatorname{clip}(s,p,q)=\min(q,\max(p,s))$; $\operatorname{mod}(s,L)\in[0,L)$ kể cả với $s<0$.

\textbf{Bước 3. Xác định hai sóng chạy và ghép nghiệm.} $A=f_e/2-H/(2a)$, $B=f_e/2+H/(2a)$. Thay hai cận:
\[u(x,t)=\frac{f_e(x-1t)+f_e(x+1t)}2+\frac{H(x+1t)-H(x-1t)}{2}.\]

\textbf{Bước 4. Trả lời các yêu cầu và kiểm tra.} $\widetilde g$ liên tục nên tại $t=1$, $u_x$ tồn tại mọi $x>0$. Ở $x=1/2$, khi $t\ge5/2$, hiệu hai đuôi là $4$, chia $2$ cho $u=2$. Góc của $\widetilde g$ tại $0$ ảnh hưởng cấp hai, không ảnh hưởng tồn tại $u_x$.

Tại $t=0$, hiệu nguyên hàm bằng $0$ và đạo hàm thời gian cho $H'(x)=g(x)$ trong miền dữ liệu; half-sum cho đúng li độ. Mỗi hạng là sóng của $x\pm at$, nên thỏa PDE ngoài các mặt gãy và theo nghĩa phân phối trên toàn miền. Đồ thị hai sóng ban đầu và các lát cắt thời gian được tính trực tiếp từ các nhánh trên.
\begin{center}\includegraphics[width=.97\linewidth,height=11cm,keepaspectratio]{_xu_ly/graphs/BT_3_2_3_10.png}\end{center}
\textit{Đồ thị minh họa; với nghiệm số, xem phương pháp và tham số trong mã Python kèm theo.}

\clearpage\noindent\texttt{\detokenize{BT_3_2_3_11}}\par
\section*{Lời giải bổ sung — Bài tập 11 · Mục 3.2.3}

\subsection*{Kiến thức và trình tự giải}
\textbf{Bước 1 — Liệt kê các khoảng của dữ liệu.} Ghi đủ điểm đầu, điểm cuối hỗ trợ, điểm gãy và bước nhảy của $f,g$. Nếu có biên, phản xạ chẵn hoặc lẻ trước khi tích phân; trên đoạn hữu hạn tiếp tục phản xạ đến một chu kỳ hoàn chỉnh.

\textbf{Bước 2 — Tính nguyên hàm liên tục.} Lấy $H'=\widetilde g$ trên từng khoảng và chọn hằng số để hai vết của $H$ bằng nhau. Viết $[s]_p^q=\min(q,\max(p,s))$ hoặc $(s)_+=\max(s,0)$ khi cần gộp nhiều nhánh. Một nguyên hàm liên tục tránh tạo ra bước nhảy giả trong nghiệm.

\textbf{Bước 3 — Xác định sóng phải và sóng trái.} Đặt $A(s)=\widetilde f(s)/2-H(s)/(2a)$, $B(s)=\widetilde f(s)/2+H(s)/(2a)$; khi đó $u=A(x-at)+B(x+at)$. Với từng $t$ được yêu cầu, dịch đồ thị $A$ sang phải $at$ và $B$ sang trái $at$, rồi cộng tung độ. Các đường $x-at=b$ và $x+at=b$ chia mặt phẳng thành các miền có công thức đơn giản.

\textbf{Bước 4 — Xét liên tục, đạo hàm và giao của các mặt sóng.} Nếu $f$ có bước nhảy, $u$ thường có bước nhảy; nếu $f$ liên tục và $g$ bị chặn thì $u$ liên tục. Tại thời điểm cố định,


\[
u_x=\frac{\widetilde f'(x-at)+\widetilde f'(x+at)}2+\frac{\widetilde g(x+at)-\widetilde g(x-at)}{2a}.
\]


Tính giới hạn trái/phải của biểu thức này tại từng vị trí ứng viên. Có thể có triệt tiêu khi hai đường gặp nhau, nên không chỉ liệt kê đường rồi khẳng định mọi đạo hàm đều mất. Với tính $C^1$ theo cả $x,t$, còn kiểm tra $u_t$. Với $C^2$, xét tiếp các góc của $g$ và đạo hàm cấp hai của $f$.

\textbf{Bước 5 — Tính thời điểm đến, giới hạn và ngưỡng.} Sóng ở $(x,t)$ chỉ nhận dữ liệu trong $[x-at,x+at]$. Giải bất đẳng thức giao hỗ trợ để tìm thời điểm đến và đi qua. Nếu $H$ hằng ngoài đoạn hỗ trợ, dùng các giá trị hằng này để tính giới hạn khi $t\to\infty$. Nếu hỏi áp suất cực đại, viết hàm từng khoảng thời gian và so sánh cực trị, bước nhảy cùng hai giá trị một phía.

\textbf{Bước 1 — Đọc đúng điều kiện biên.} Biên $u(0,t)=0$ dùng mở rộng lẻ; biên $u_x(0,t)=0$ dùng mở rộng chẵn. Với dữ liệu trên $x\ge0$, đặt $\widetilde f(-s)=-f(s)$ hoặc $f(s)$ tương ứng, làm tương tự cho $g$. Mở rộng lẻ là $C^2$ nếu $f(0)=f''(0)=0$, và mở rộng lẻ của $g$ là $C^1$ nếu $g(0)=0$. Mở rộng chẵn đòi $f'(0)=g'(0)=0$ cho mức trơn ấy.

\textbf{Bước 2 — Lập nguyên hàm sau khi phản xạ.} Đặt $H(s)=\int_0^s\widetilde g(r)dr$. Nếu $\widetilde g$ lẻ thì $H$ chẵn; nếu $\widetilde g$ chẵn thì $H$ lẻ. Chính nguyên hàm, chứ không phải vận tốc, xuất hiện trong


\[
u=\frac{\widetilde f(x-at)+\widetilde f(x+at)}2+\frac{H(x+at)-H(x-at)}{2a}.
\]


\textbf{Bước 3 — Tách miền trước và sau khi chạm biên.} Với $u(0,t)=h(t)$, ở $x\ge at$ dùng d'Alembert với dữ liệu gốc. Ở $0<x<at$:


\[
u=h(t-x/a)+\frac{f(x+at)-f(at-x)}2+\frac1{2a}\int_{at-x}^{at+x}g(s)ds.
\]


Với $u_x(0,t)=h(t)$, dùng nghiệm của các mở rộng chẵn rồi trừ $a\int_0^{t-x/a}h(s)ds$ trên miền $x<at$. Tính riêng đạo hàm theo $x$ của hạng này để kiểm tra biên.

\textbf{Bước 4 — Trên đoạn hữu hạn, phản xạ ở cả hai đầu.} Dirichlet ở hai đầu cho mở rộng lẻ tại $0,L$, chu kỳ không gian $2L$. Neumann ở hai đầu cho mở rộng chẵn tại $0,L$, cũng chu kỳ $2L$. Một đầu lẻ, một đầu chẵn cho $\widetilde f(s+2L)=-\widetilde f(s)$ và chu kỳ $4L$, tương tự $g$. Vì thế chu kỳ thời gian thường là $2L/a$ hoặc $4L/a$. Với Neumann hai đầu còn phải xét trung bình $g$: trung bình khác không sinh hạng tuyến tính theo $t$, nên nghiệm không tuần hoàn.

\textbf{Bước 5 — Kiểm tra góc và đường ghép.} Biên Dirichlet không thuần nhất cần $h(0)=f(0)$, $h'(0)=g(0)$, $h''(0)=a^2f''(0)$ để có nghiệm $C^2$ cho phương trình thuần nhất. So sánh hai vết trên $x=at$, gồm $u,u_x,u_t$ rồi đạo hàm cấp hai. Nếu không tương thích, nêu rõ mức trơn còn lại thay vì khẳng định nghiệm cổ điển trên toàn miền.
\subsection*{Lời giải áp dụng cho bài}
\textbf{Bước 1. Đọc dữ liệu và chọn phản xạ.} Nửa trục Neumann $u_x(0,t)=0$, $a=1$, $f=0$, $g=\chi_{[1,2]}$.

\textbf{Bước 2. Tính từng nhánh và nối nguyên hàm.} Biên trong bản scan là Neumann: phải mở rộng chẵn $\widetilde g=\chi_{[1,2]}(|s|)$. Tích phân dương cho $[s-1]_0^1$; nguyên hàm lẻ là $\operatorname{sgn}(s)[|s|-1]_0^1$. Công thức lẻ của vận tốc trong lời giải cũ là nhầm biên.
\[f_e(s)=0,\qquad H(s)=\operatorname{sgn}(s)\operatorname{clip}(|s|-1,0,1).\]
Ở đây $\operatorname{clip}(s,p,q)=\min(q,\max(p,s))$; $\operatorname{mod}(s,L)\in[0,L)$ kể cả với $s<0$.

\textbf{Bước 3. Xác định hai sóng chạy và ghép nghiệm.} $A=f_e/2-H/(2a)$, $B=f_e/2+H/(2a)$. Thay hai cận:
\[u(x,t)=\frac{f_e(x-1t)+f_e(x+1t)}2+\frac{H(x+1t)-H(x-1t)}{2}.\]

\textbf{Bước 4. Trả lời các yêu cầu và kiểm tra.} Vì $H$ lẻ, $H'(-s)=H'(s)$, nên $u_x(0,t)=(\widetilde g(t)-\widetilde g(-t))/2=0$. Nghiệm liên tục; đạo hàm cấp một gãy trên $x\pm t=\pm1,\pm2$ ở phần nằm trong nửa trục.

Tại $t=0$, hiệu nguyên hàm bằng $0$ và đạo hàm thời gian cho $H'(x)=g(x)$ trong miền dữ liệu; half-sum cho đúng li độ. Mỗi hạng là sóng của $x\pm at$, nên thỏa PDE ngoài các mặt gãy và theo nghĩa phân phối trên toàn miền. Đồ thị hai sóng ban đầu và các lát cắt thời gian được tính trực tiếp từ các nhánh trên.
\begin{center}\includegraphics[width=.97\linewidth,height=11cm,keepaspectratio]{_xu_ly/graphs/BT_3_2_3_11.png}\end{center}
\textit{Đồ thị minh họa; với nghiệm số, xem phương pháp và tham số trong mã Python kèm theo.}

\clearpage\noindent\texttt{\detokenize{BT_3_2_3_12}}\par
\section*{Lời giải bổ sung — Bài tập 12 · Mục 3.2.3}

\subsection*{Kiến thức và trình tự giải}
\textbf{Bước 1 — Liệt kê các khoảng của dữ liệu.} Ghi đủ điểm đầu, điểm cuối hỗ trợ, điểm gãy và bước nhảy của $f,g$. Nếu có biên, phản xạ chẵn hoặc lẻ trước khi tích phân; trên đoạn hữu hạn tiếp tục phản xạ đến một chu kỳ hoàn chỉnh.

\textbf{Bước 2 — Tính nguyên hàm liên tục.} Lấy $H'=\widetilde g$ trên từng khoảng và chọn hằng số để hai vết của $H$ bằng nhau. Viết $[s]_p^q=\min(q,\max(p,s))$ hoặc $(s)_+=\max(s,0)$ khi cần gộp nhiều nhánh. Một nguyên hàm liên tục tránh tạo ra bước nhảy giả trong nghiệm.

\textbf{Bước 3 — Xác định sóng phải và sóng trái.} Đặt $A(s)=\widetilde f(s)/2-H(s)/(2a)$, $B(s)=\widetilde f(s)/2+H(s)/(2a)$; khi đó $u=A(x-at)+B(x+at)$. Với từng $t$ được yêu cầu, dịch đồ thị $A$ sang phải $at$ và $B$ sang trái $at$, rồi cộng tung độ. Các đường $x-at=b$ và $x+at=b$ chia mặt phẳng thành các miền có công thức đơn giản.

\textbf{Bước 4 — Xét liên tục, đạo hàm và giao của các mặt sóng.} Nếu $f$ có bước nhảy, $u$ thường có bước nhảy; nếu $f$ liên tục và $g$ bị chặn thì $u$ liên tục. Tại thời điểm cố định,


\[
u_x=\frac{\widetilde f'(x-at)+\widetilde f'(x+at)}2+\frac{\widetilde g(x+at)-\widetilde g(x-at)}{2a}.
\]


Tính giới hạn trái/phải của biểu thức này tại từng vị trí ứng viên. Có thể có triệt tiêu khi hai đường gặp nhau, nên không chỉ liệt kê đường rồi khẳng định mọi đạo hàm đều mất. Với tính $C^1$ theo cả $x,t$, còn kiểm tra $u_t$. Với $C^2$, xét tiếp các góc của $g$ và đạo hàm cấp hai của $f$.

\textbf{Bước 5 — Tính thời điểm đến, giới hạn và ngưỡng.} Sóng ở $(x,t)$ chỉ nhận dữ liệu trong $[x-at,x+at]$. Giải bất đẳng thức giao hỗ trợ để tìm thời điểm đến và đi qua. Nếu $H$ hằng ngoài đoạn hỗ trợ, dùng các giá trị hằng này để tính giới hạn khi $t\to\infty$. Nếu hỏi áp suất cực đại, viết hàm từng khoảng thời gian và so sánh cực trị, bước nhảy cùng hai giá trị một phía.

\textbf{Bước 1 — Đọc đúng điều kiện biên.} Biên $u(0,t)=0$ dùng mở rộng lẻ; biên $u_x(0,t)=0$ dùng mở rộng chẵn. Với dữ liệu trên $x\ge0$, đặt $\widetilde f(-s)=-f(s)$ hoặc $f(s)$ tương ứng, làm tương tự cho $g$. Mở rộng lẻ là $C^2$ nếu $f(0)=f''(0)=0$, và mở rộng lẻ của $g$ là $C^1$ nếu $g(0)=0$. Mở rộng chẵn đòi $f'(0)=g'(0)=0$ cho mức trơn ấy.

\textbf{Bước 2 — Lập nguyên hàm sau khi phản xạ.} Đặt $H(s)=\int_0^s\widetilde g(r)dr$. Nếu $\widetilde g$ lẻ thì $H$ chẵn; nếu $\widetilde g$ chẵn thì $H$ lẻ. Chính nguyên hàm, chứ không phải vận tốc, xuất hiện trong


\[
u=\frac{\widetilde f(x-at)+\widetilde f(x+at)}2+\frac{H(x+at)-H(x-at)}{2a}.
\]


\textbf{Bước 3 — Tách miền trước và sau khi chạm biên.} Với $u(0,t)=h(t)$, ở $x\ge at$ dùng d'Alembert với dữ liệu gốc. Ở $0<x<at$:


\[
u=h(t-x/a)+\frac{f(x+at)-f(at-x)}2+\frac1{2a}\int_{at-x}^{at+x}g(s)ds.
\]


Với $u_x(0,t)=h(t)$, dùng nghiệm của các mở rộng chẵn rồi trừ $a\int_0^{t-x/a}h(s)ds$ trên miền $x<at$. Tính riêng đạo hàm theo $x$ của hạng này để kiểm tra biên.

\textbf{Bước 4 — Trên đoạn hữu hạn, phản xạ ở cả hai đầu.} Dirichlet ở hai đầu cho mở rộng lẻ tại $0,L$, chu kỳ không gian $2L$. Neumann ở hai đầu cho mở rộng chẵn tại $0,L$, cũng chu kỳ $2L$. Một đầu lẻ, một đầu chẵn cho $\widetilde f(s+2L)=-\widetilde f(s)$ và chu kỳ $4L$, tương tự $g$. Vì thế chu kỳ thời gian thường là $2L/a$ hoặc $4L/a$. Với Neumann hai đầu còn phải xét trung bình $g$: trung bình khác không sinh hạng tuyến tính theo $t$, nên nghiệm không tuần hoàn.

\textbf{Bước 5 — Kiểm tra góc và đường ghép.} Biên Dirichlet không thuần nhất cần $h(0)=f(0)$, $h'(0)=g(0)$, $h''(0)=a^2f''(0)$ để có nghiệm $C^2$ cho phương trình thuần nhất. So sánh hai vết trên $x=at$, gồm $u,u_x,u_t$ rồi đạo hàm cấp hai. Nếu không tương thích, nêu rõ mức trơn còn lại thay vì khẳng định nghiệm cổ điển trên toàn miền.
\subsection*{Lời giải áp dụng cho bài}
\textbf{Bước 1. Đọc dữ liệu và chọn phản xạ.} Nửa trục Neumann, $a=2$, $f=0$, $g=\chi_{[1,2]}$.

\textbf{Bước 2. Tính từng nhánh và nối nguyên hàm.} Mở rộng chẵn và nguyên hàm giống Bài 11. Chỉ tốc độ đổi từ $1$ thành $2$, nên độ dịch là $2t$ và hệ số tích phân là $1/4$.
\[f_e(s)=0,\qquad H(s)=\operatorname{sgn}(s)\operatorname{clip}(|s|-1,0,1).\]
Ở đây $\operatorname{clip}(s,p,q)=\min(q,\max(p,s))$; $\operatorname{mod}(s,L)\in[0,L)$ kể cả với $s<0$.

\textbf{Bước 3. Xác định hai sóng chạy và ghép nghiệm.} $A=f_e/2-H/(2a)$, $B=f_e/2+H/(2a)$. Thay hai cận:
\[u(x,t)=\frac{f_e(x-2t)+f_e(x+2t)}2+\frac{H(x+2t)-H(x-2t)}{4}.\]

\textbf{Bước 4. Trả lời các yêu cầu và kiểm tra.} $u_x(0,t)=0$ bởi tính chẵn của vận tốc mở rộng. $H$ liên tục nên $u$ liên tục; $H'$ nhảy ở $\pm1,\pm2$ nên các đường mất trơn cấp một là $x\pm2t=\pm1,\pm2$. Khi $t>(x+2)/2$, $u=1/2$ tại mỗi điểm cố định.

Tại $t=0$, hiệu nguyên hàm bằng $0$ và đạo hàm thời gian cho $H'(x)=g(x)$ trong miền dữ liệu; half-sum cho đúng li độ. Mỗi hạng là sóng của $x\pm at$, nên thỏa PDE ngoài các mặt gãy và theo nghĩa phân phối trên toàn miền. Đồ thị hai sóng ban đầu và các lát cắt thời gian được tính trực tiếp từ các nhánh trên.
\begin{center}\includegraphics[width=.97\linewidth,height=11cm,keepaspectratio]{_xu_ly/graphs/BT_3_2_3_12.png}\end{center}
\textit{Đồ thị minh họa; với nghiệm số, xem phương pháp và tham số trong mã Python kèm theo.}

\clearpage\noindent\texttt{\detokenize{BT_3_2_3_13}}\par
\section*{Lời giải bổ sung — Bài tập 13 · Mục 3.2.3}

\subsection*{Kiến thức và trình tự giải}
\textbf{Bước 1 — Liệt kê các khoảng của dữ liệu.} Ghi đủ điểm đầu, điểm cuối hỗ trợ, điểm gãy và bước nhảy của $f,g$. Nếu có biên, phản xạ chẵn hoặc lẻ trước khi tích phân; trên đoạn hữu hạn tiếp tục phản xạ đến một chu kỳ hoàn chỉnh.

\textbf{Bước 2 — Tính nguyên hàm liên tục.} Lấy $H'=\widetilde g$ trên từng khoảng và chọn hằng số để hai vết của $H$ bằng nhau. Viết $[s]_p^q=\min(q,\max(p,s))$ hoặc $(s)_+=\max(s,0)$ khi cần gộp nhiều nhánh. Một nguyên hàm liên tục tránh tạo ra bước nhảy giả trong nghiệm.

\textbf{Bước 3 — Xác định sóng phải và sóng trái.} Đặt $A(s)=\widetilde f(s)/2-H(s)/(2a)$, $B(s)=\widetilde f(s)/2+H(s)/(2a)$; khi đó $u=A(x-at)+B(x+at)$. Với từng $t$ được yêu cầu, dịch đồ thị $A$ sang phải $at$ và $B$ sang trái $at$, rồi cộng tung độ. Các đường $x-at=b$ và $x+at=b$ chia mặt phẳng thành các miền có công thức đơn giản.

\textbf{Bước 4 — Xét liên tục, đạo hàm và giao của các mặt sóng.} Nếu $f$ có bước nhảy, $u$ thường có bước nhảy; nếu $f$ liên tục và $g$ bị chặn thì $u$ liên tục. Tại thời điểm cố định,


\[
u_x=\frac{\widetilde f'(x-at)+\widetilde f'(x+at)}2+\frac{\widetilde g(x+at)-\widetilde g(x-at)}{2a}.
\]


Tính giới hạn trái/phải của biểu thức này tại từng vị trí ứng viên. Có thể có triệt tiêu khi hai đường gặp nhau, nên không chỉ liệt kê đường rồi khẳng định mọi đạo hàm đều mất. Với tính $C^1$ theo cả $x,t$, còn kiểm tra $u_t$. Với $C^2$, xét tiếp các góc của $g$ và đạo hàm cấp hai của $f$.

\textbf{Bước 5 — Tính thời điểm đến, giới hạn và ngưỡng.} Sóng ở $(x,t)$ chỉ nhận dữ liệu trong $[x-at,x+at]$. Giải bất đẳng thức giao hỗ trợ để tìm thời điểm đến và đi qua. Nếu $H$ hằng ngoài đoạn hỗ trợ, dùng các giá trị hằng này để tính giới hạn khi $t\to\infty$. Nếu hỏi áp suất cực đại, viết hàm từng khoảng thời gian và so sánh cực trị, bước nhảy cùng hai giá trị một phía.

\textbf{Bước 1 — Đọc đúng điều kiện biên.} Biên $u(0,t)=0$ dùng mở rộng lẻ; biên $u_x(0,t)=0$ dùng mở rộng chẵn. Với dữ liệu trên $x\ge0$, đặt $\widetilde f(-s)=-f(s)$ hoặc $f(s)$ tương ứng, làm tương tự cho $g$. Mở rộng lẻ là $C^2$ nếu $f(0)=f''(0)=0$, và mở rộng lẻ của $g$ là $C^1$ nếu $g(0)=0$. Mở rộng chẵn đòi $f'(0)=g'(0)=0$ cho mức trơn ấy.

\textbf{Bước 2 — Lập nguyên hàm sau khi phản xạ.} Đặt $H(s)=\int_0^s\widetilde g(r)dr$. Nếu $\widetilde g$ lẻ thì $H$ chẵn; nếu $\widetilde g$ chẵn thì $H$ lẻ. Chính nguyên hàm, chứ không phải vận tốc, xuất hiện trong


\[
u=\frac{\widetilde f(x-at)+\widetilde f(x+at)}2+\frac{H(x+at)-H(x-at)}{2a}.
\]


\textbf{Bước 3 — Tách miền trước và sau khi chạm biên.} Với $u(0,t)=h(t)$, ở $x\ge at$ dùng d'Alembert với dữ liệu gốc. Ở $0<x<at$:


\[
u=h(t-x/a)+\frac{f(x+at)-f(at-x)}2+\frac1{2a}\int_{at-x}^{at+x}g(s)ds.
\]


Với $u_x(0,t)=h(t)$, dùng nghiệm của các mở rộng chẵn rồi trừ $a\int_0^{t-x/a}h(s)ds$ trên miền $x<at$. Tính riêng đạo hàm theo $x$ của hạng này để kiểm tra biên.

\textbf{Bước 4 — Trên đoạn hữu hạn, phản xạ ở cả hai đầu.} Dirichlet ở hai đầu cho mở rộng lẻ tại $0,L$, chu kỳ không gian $2L$. Neumann ở hai đầu cho mở rộng chẵn tại $0,L$, cũng chu kỳ $2L$. Một đầu lẻ, một đầu chẵn cho $\widetilde f(s+2L)=-\widetilde f(s)$ và chu kỳ $4L$, tương tự $g$. Vì thế chu kỳ thời gian thường là $2L/a$ hoặc $4L/a$. Với Neumann hai đầu còn phải xét trung bình $g$: trung bình khác không sinh hạng tuyến tính theo $t$, nên nghiệm không tuần hoàn.

\textbf{Bước 5 — Kiểm tra góc và đường ghép.} Biên Dirichlet không thuần nhất cần $h(0)=f(0)$, $h'(0)=g(0)$, $h''(0)=a^2f''(0)$ để có nghiệm $C^2$ cho phương trình thuần nhất. So sánh hai vết trên $x=at$, gồm $u,u_x,u_t$ rồi đạo hàm cấp hai. Nếu không tương thích, nêu rõ mức trơn còn lại thay vì khẳng định nghiệm cổ điển trên toàn miền.
\subsection*{Lời giải áp dụng cho bài}
\textbf{Bước 1. Đọc dữ liệu và chọn phản xạ.} Nửa trục Dirichlet, $a=2$, $f=0$, $g=s\chi_{[1,2]}$.

\textbf{Bước 2. Tính từng nhánh và nối nguyên hàm.} Mở rộng lẻ; trên đối số dương, $H=0$ trước $1$, $(s^2-1)/2$ trên $[1,2]$, $3/2$ sau $2$. Phản xạ $H$ chẵn, nên dạng gộp dùng $|s|$.
\[f_e(s)=0,\qquad H(s)=\tfrac12(\operatorname{clip}(|s|,1,2)^2-1).\]
Ở đây $\operatorname{clip}(s,p,q)=\min(q,\max(p,s))$; $\operatorname{mod}(s,L)\in[0,L)$ kể cả với $s<0$.

\textbf{Bước 3. Xác định hai sóng chạy và ghép nghiệm.} $A=f_e/2-H/(2a)$, $B=f_e/2+H/(2a)$. Thay hai cận:
\[u(x,t)=\frac{f_e(x-2t)+f_e(x+2t)}2+\frac{H(x+2t)-H(x-2t)}{4}.\]

\textbf{Bước 4. Trả lời các yêu cầu và kiểm tra.} Tại $x=0$, hai giá trị $H$ bằng nhau nên $u=0$. Nghiệm liên tục nhưng các bước nhảy của $g$ ở $\pm1,\pm2$ phát sinh các mặt gãy cấp một. Sau $t>(x+2)/2$, cả hai đuôi của $H$ bằng $3/2$, nên nghiệm trở về $0$.

Tại $t=0$, hiệu nguyên hàm bằng $0$ và đạo hàm thời gian cho $H'(x)=g(x)$ trong miền dữ liệu; half-sum cho đúng li độ. Mỗi hạng là sóng của $x\pm at$, nên thỏa PDE ngoài các mặt gãy và theo nghĩa phân phối trên toàn miền. Đồ thị hai sóng ban đầu và các lát cắt thời gian được tính trực tiếp từ các nhánh trên.
\begin{center}\includegraphics[width=.97\linewidth,height=11cm,keepaspectratio]{_xu_ly/graphs/BT_3_2_3_13.png}\end{center}
\textit{Đồ thị minh họa; với nghiệm số, xem phương pháp và tham số trong mã Python kèm theo.}

\clearpage\noindent\texttt{\detokenize{BT_3_2_3_14}}\par
\section*{Lời giải bổ sung — Bài tập 14 · Mục 3.2.3}

\subsection*{Kiến thức và trình tự giải}
\textbf{Bước 1 — Liệt kê các khoảng của dữ liệu.} Ghi đủ điểm đầu, điểm cuối hỗ trợ, điểm gãy và bước nhảy của $f,g$. Nếu có biên, phản xạ chẵn hoặc lẻ trước khi tích phân; trên đoạn hữu hạn tiếp tục phản xạ đến một chu kỳ hoàn chỉnh.

\textbf{Bước 2 — Tính nguyên hàm liên tục.} Lấy $H'=\widetilde g$ trên từng khoảng và chọn hằng số để hai vết của $H$ bằng nhau. Viết $[s]_p^q=\min(q,\max(p,s))$ hoặc $(s)_+=\max(s,0)$ khi cần gộp nhiều nhánh. Một nguyên hàm liên tục tránh tạo ra bước nhảy giả trong nghiệm.

\textbf{Bước 3 — Xác định sóng phải và sóng trái.} Đặt $A(s)=\widetilde f(s)/2-H(s)/(2a)$, $B(s)=\widetilde f(s)/2+H(s)/(2a)$; khi đó $u=A(x-at)+B(x+at)$. Với từng $t$ được yêu cầu, dịch đồ thị $A$ sang phải $at$ và $B$ sang trái $at$, rồi cộng tung độ. Các đường $x-at=b$ và $x+at=b$ chia mặt phẳng thành các miền có công thức đơn giản.

\textbf{Bước 4 — Xét liên tục, đạo hàm và giao của các mặt sóng.} Nếu $f$ có bước nhảy, $u$ thường có bước nhảy; nếu $f$ liên tục và $g$ bị chặn thì $u$ liên tục. Tại thời điểm cố định,


\[
u_x=\frac{\widetilde f'(x-at)+\widetilde f'(x+at)}2+\frac{\widetilde g(x+at)-\widetilde g(x-at)}{2a}.
\]


Tính giới hạn trái/phải của biểu thức này tại từng vị trí ứng viên. Có thể có triệt tiêu khi hai đường gặp nhau, nên không chỉ liệt kê đường rồi khẳng định mọi đạo hàm đều mất. Với tính $C^1$ theo cả $x,t$, còn kiểm tra $u_t$. Với $C^2$, xét tiếp các góc của $g$ và đạo hàm cấp hai của $f$.

\textbf{Bước 5 — Tính thời điểm đến, giới hạn và ngưỡng.} Sóng ở $(x,t)$ chỉ nhận dữ liệu trong $[x-at,x+at]$. Giải bất đẳng thức giao hỗ trợ để tìm thời điểm đến và đi qua. Nếu $H$ hằng ngoài đoạn hỗ trợ, dùng các giá trị hằng này để tính giới hạn khi $t\to\infty$. Nếu hỏi áp suất cực đại, viết hàm từng khoảng thời gian và so sánh cực trị, bước nhảy cùng hai giá trị một phía.

\textbf{Bước 1 — Đọc đúng điều kiện biên.} Biên $u(0,t)=0$ dùng mở rộng lẻ; biên $u_x(0,t)=0$ dùng mở rộng chẵn. Với dữ liệu trên $x\ge0$, đặt $\widetilde f(-s)=-f(s)$ hoặc $f(s)$ tương ứng, làm tương tự cho $g$. Mở rộng lẻ là $C^2$ nếu $f(0)=f''(0)=0$, và mở rộng lẻ của $g$ là $C^1$ nếu $g(0)=0$. Mở rộng chẵn đòi $f'(0)=g'(0)=0$ cho mức trơn ấy.

\textbf{Bước 2 — Lập nguyên hàm sau khi phản xạ.} Đặt $H(s)=\int_0^s\widetilde g(r)dr$. Nếu $\widetilde g$ lẻ thì $H$ chẵn; nếu $\widetilde g$ chẵn thì $H$ lẻ. Chính nguyên hàm, chứ không phải vận tốc, xuất hiện trong


\[
u=\frac{\widetilde f(x-at)+\widetilde f(x+at)}2+\frac{H(x+at)-H(x-at)}{2a}.
\]


\textbf{Bước 3 — Tách miền trước và sau khi chạm biên.} Với $u(0,t)=h(t)$, ở $x\ge at$ dùng d'Alembert với dữ liệu gốc. Ở $0<x<at$:


\[
u=h(t-x/a)+\frac{f(x+at)-f(at-x)}2+\frac1{2a}\int_{at-x}^{at+x}g(s)ds.
\]


Với $u_x(0,t)=h(t)$, dùng nghiệm của các mở rộng chẵn rồi trừ $a\int_0^{t-x/a}h(s)ds$ trên miền $x<at$. Tính riêng đạo hàm theo $x$ của hạng này để kiểm tra biên.

\textbf{Bước 4 — Trên đoạn hữu hạn, phản xạ ở cả hai đầu.} Dirichlet ở hai đầu cho mở rộng lẻ tại $0,L$, chu kỳ không gian $2L$. Neumann ở hai đầu cho mở rộng chẵn tại $0,L$, cũng chu kỳ $2L$. Một đầu lẻ, một đầu chẵn cho $\widetilde f(s+2L)=-\widetilde f(s)$ và chu kỳ $4L$, tương tự $g$. Vì thế chu kỳ thời gian thường là $2L/a$ hoặc $4L/a$. Với Neumann hai đầu còn phải xét trung bình $g$: trung bình khác không sinh hạng tuyến tính theo $t$, nên nghiệm không tuần hoàn.

\textbf{Bước 5 — Kiểm tra góc và đường ghép.} Biên Dirichlet không thuần nhất cần $h(0)=f(0)$, $h'(0)=g(0)$, $h''(0)=a^2f''(0)$ để có nghiệm $C^2$ cho phương trình thuần nhất. So sánh hai vết trên $x=at$, gồm $u,u_x,u_t$ rồi đạo hàm cấp hai. Nếu không tương thích, nêu rõ mức trơn còn lại thay vì khẳng định nghiệm cổ điển trên toàn miền.
\subsection*{Lời giải áp dụng cho bài}
\textbf{Bước 1. Đọc dữ liệu và chọn phản xạ.} Nửa trục Neumann, $a=3$, $f=0$, $g=s\chi_{[2,3]}$.

\textbf{Bước 2. Tính từng nhánh và nối nguyên hàm.} Mở rộng chẵn. Trên $s>0$, nguyên hàm là $0$ trước $2$, $(s^2-4)/2$ trên $[2,3]$, $5/2$ sau $3$; phản xạ nguyên hàm lẻ.
\[f_e(s)=0,\qquad H(s)=\tfrac12\operatorname{sgn}(s)(\operatorname{clip}(|s|,2,3)^2-4).\]
Ở đây $\operatorname{clip}(s,p,q)=\min(q,\max(p,s))$; $\operatorname{mod}(s,L)\in[0,L)$ kể cả với $s<0$.

\textbf{Bước 3. Xác định hai sóng chạy và ghép nghiệm.} $A=f_e/2-H/(2a)$, $B=f_e/2+H/(2a)$. Thay hai cận:
\[u(x,t)=\frac{f_e(x-3t)+f_e(x+3t)}2+\frac{H(x+3t)-H(x-3t)}{6}.\]

\textbf{Bước 4. Trả lời các yêu cầu và kiểm tra.} Nghiệm liên tục, $u_x(0,t)=0$, các đường gãy cấp một là $x\pm3t=\pm2,\pm3$. Sau $t>(x+3)/3$, hiệu hai đuôi là $5$, nên $u=5/6$. Điều này khác trường hợp Dirichlet có hai đuôi nguyên hàm bằng nhau.

Tại $t=0$, hiệu nguyên hàm bằng $0$ và đạo hàm thời gian cho $H'(x)=g(x)$ trong miền dữ liệu; half-sum cho đúng li độ. Mỗi hạng là sóng của $x\pm at$, nên thỏa PDE ngoài các mặt gãy và theo nghĩa phân phối trên toàn miền. Đồ thị hai sóng ban đầu và các lát cắt thời gian được tính trực tiếp từ các nhánh trên.
\begin{center}\includegraphics[width=.97\linewidth,height=11cm,keepaspectratio]{_xu_ly/graphs/BT_3_2_3_14.png}\end{center}
\textit{Đồ thị minh họa; với nghiệm số, xem phương pháp và tham số trong mã Python kèm theo.}

\clearpage\noindent\texttt{\detokenize{BT_3_2_3_15}}\par
\section*{Lời giải bổ sung — Bài tập 15 · Mục 3.2.3}

\subsection*{Kiến thức và trình tự giải}
\textbf{Bước 1 — Liệt kê các khoảng của dữ liệu.} Ghi đủ điểm đầu, điểm cuối hỗ trợ, điểm gãy và bước nhảy của $f,g$. Nếu có biên, phản xạ chẵn hoặc lẻ trước khi tích phân; trên đoạn hữu hạn tiếp tục phản xạ đến một chu kỳ hoàn chỉnh.

\textbf{Bước 2 — Tính nguyên hàm liên tục.} Lấy $H'=\widetilde g$ trên từng khoảng và chọn hằng số để hai vết của $H$ bằng nhau. Viết $[s]_p^q=\min(q,\max(p,s))$ hoặc $(s)_+=\max(s,0)$ khi cần gộp nhiều nhánh. Một nguyên hàm liên tục tránh tạo ra bước nhảy giả trong nghiệm.

\textbf{Bước 3 — Xác định sóng phải và sóng trái.} Đặt $A(s)=\widetilde f(s)/2-H(s)/(2a)$, $B(s)=\widetilde f(s)/2+H(s)/(2a)$; khi đó $u=A(x-at)+B(x+at)$. Với từng $t$ được yêu cầu, dịch đồ thị $A$ sang phải $at$ và $B$ sang trái $at$, rồi cộng tung độ. Các đường $x-at=b$ và $x+at=b$ chia mặt phẳng thành các miền có công thức đơn giản.

\textbf{Bước 4 — Xét liên tục, đạo hàm và giao của các mặt sóng.} Nếu $f$ có bước nhảy, $u$ thường có bước nhảy; nếu $f$ liên tục và $g$ bị chặn thì $u$ liên tục. Tại thời điểm cố định,


\[
u_x=\frac{\widetilde f'(x-at)+\widetilde f'(x+at)}2+\frac{\widetilde g(x+at)-\widetilde g(x-at)}{2a}.
\]


Tính giới hạn trái/phải của biểu thức này tại từng vị trí ứng viên. Có thể có triệt tiêu khi hai đường gặp nhau, nên không chỉ liệt kê đường rồi khẳng định mọi đạo hàm đều mất. Với tính $C^1$ theo cả $x,t$, còn kiểm tra $u_t$. Với $C^2$, xét tiếp các góc của $g$ và đạo hàm cấp hai của $f$.

\textbf{Bước 5 — Tính thời điểm đến, giới hạn và ngưỡng.} Sóng ở $(x,t)$ chỉ nhận dữ liệu trong $[x-at,x+at]$. Giải bất đẳng thức giao hỗ trợ để tìm thời điểm đến và đi qua. Nếu $H$ hằng ngoài đoạn hỗ trợ, dùng các giá trị hằng này để tính giới hạn khi $t\to\infty$. Nếu hỏi áp suất cực đại, viết hàm từng khoảng thời gian và so sánh cực trị, bước nhảy cùng hai giá trị một phía.

\textbf{Bước 1 — Đọc đúng điều kiện biên.} Biên $u(0,t)=0$ dùng mở rộng lẻ; biên $u_x(0,t)=0$ dùng mở rộng chẵn. Với dữ liệu trên $x\ge0$, đặt $\widetilde f(-s)=-f(s)$ hoặc $f(s)$ tương ứng, làm tương tự cho $g$. Mở rộng lẻ là $C^2$ nếu $f(0)=f''(0)=0$, và mở rộng lẻ của $g$ là $C^1$ nếu $g(0)=0$. Mở rộng chẵn đòi $f'(0)=g'(0)=0$ cho mức trơn ấy.

\textbf{Bước 2 — Lập nguyên hàm sau khi phản xạ.} Đặt $H(s)=\int_0^s\widetilde g(r)dr$. Nếu $\widetilde g$ lẻ thì $H$ chẵn; nếu $\widetilde g$ chẵn thì $H$ lẻ. Chính nguyên hàm, chứ không phải vận tốc, xuất hiện trong


\[
u=\frac{\widetilde f(x-at)+\widetilde f(x+at)}2+\frac{H(x+at)-H(x-at)}{2a}.
\]


\textbf{Bước 3 — Tách miền trước và sau khi chạm biên.} Với $u(0,t)=h(t)$, ở $x\ge at$ dùng d'Alembert với dữ liệu gốc. Ở $0<x<at$:


\[
u=h(t-x/a)+\frac{f(x+at)-f(at-x)}2+\frac1{2a}\int_{at-x}^{at+x}g(s)ds.
\]


Với $u_x(0,t)=h(t)$, dùng nghiệm của các mở rộng chẵn rồi trừ $a\int_0^{t-x/a}h(s)ds$ trên miền $x<at$. Tính riêng đạo hàm theo $x$ của hạng này để kiểm tra biên.

\textbf{Bước 4 — Trên đoạn hữu hạn, phản xạ ở cả hai đầu.} Dirichlet ở hai đầu cho mở rộng lẻ tại $0,L$, chu kỳ không gian $2L$. Neumann ở hai đầu cho mở rộng chẵn tại $0,L$, cũng chu kỳ $2L$. Một đầu lẻ, một đầu chẵn cho $\widetilde f(s+2L)=-\widetilde f(s)$ và chu kỳ $4L$, tương tự $g$. Vì thế chu kỳ thời gian thường là $2L/a$ hoặc $4L/a$. Với Neumann hai đầu còn phải xét trung bình $g$: trung bình khác không sinh hạng tuyến tính theo $t$, nên nghiệm không tuần hoàn.

\textbf{Bước 5 — Kiểm tra góc và đường ghép.} Biên Dirichlet không thuần nhất cần $h(0)=f(0)$, $h'(0)=g(0)$, $h''(0)=a^2f''(0)$ để có nghiệm $C^2$ cho phương trình thuần nhất. So sánh hai vết trên $x=at$, gồm $u,u_x,u_t$ rồi đạo hàm cấp hai. Nếu không tương thích, nêu rõ mức trơn còn lại thay vì khẳng định nghiệm cổ điển trên toàn miền.
\subsection*{Lời giải áp dụng cho bài}
\textbf{Bước 1. Đọc dữ liệu và chọn phản xạ.} Đoạn $(0,1)$, Dirichlet hai đầu, $a=1$, $f=0$, $g=1$.

\textbf{Bước 2. Tính từng nhánh và nối nguyên hàm.} Mở rộng lẻ tại $0,1$ cho vận tốc bằng $1$ trên $(0,1)$, $-1$ trên $(1,2)$ rồi chu kỳ $2$. Tích phân cho $H=s$ ở $[0,1]$, $2-s$ ở $[1,2]$; trung bình vận tốc một chu kỳ bằng $0$.
\[f_e(s)=0,\qquad H(s)=1-|\operatorname{mod}(s,2)-1|.\]
Ở đây $\operatorname{clip}(s,p,q)=\min(q,\max(p,s))$; $\operatorname{mod}(s,L)\in[0,L)$ kể cả với $s<0$.

\textbf{Bước 3. Xác định hai sóng chạy và ghép nghiệm.} $A=f_e/2-H/(2a)$, $B=f_e/2+H/(2a)$. Thay hai cận:
\[u(x,t)=\frac{f_e(x-1t)+f_e(x+1t)}2+\frac{H(x+1t)-H(x-1t)}{2}.\]

\textbf{Bước 4. Trả lời các yêu cầu và kiểm tra.} $H(s+2)=H(s)$ nên nghiệm có chu kỳ thời gian $2$. Chiếu $g$ lên sine: $B_n=2\int_0^1\sin(n\pi x)dx=2(1-(-1)^n)/(n\pi)$. Chia $\omega_n=n\pi$ cho chuỗi ở dưới. $g$ khác không ở góc biên nên nghiệm không $C^1$ qua mặt sóng từ các góc.

Tách biến cho $u=\sum_{n\text{ lẻ}}4\sin(n\pi t)\sin(n\pi x)/(n^2\pi^2)$. Dữ liệu $g=1$ tại hai đầu không tương thích nên đây là nghiệm yếu; chu kỳ thời gian 2.

Tại $t=0$, hiệu nguyên hàm bằng $0$ và đạo hàm thời gian cho $H'(x)=g(x)$ trong miền dữ liệu; half-sum cho đúng li độ. Mỗi hạng là sóng của $x\pm at$, nên thỏa PDE ngoài các mặt gãy và theo nghĩa phân phối trên toàn miền. Đồ thị hai sóng ban đầu và các lát cắt thời gian được tính trực tiếp từ các nhánh trên.

\textbf{Chi tiết tích phân của dữ liệu hằng.} Với $k_n=n\pi/L$,
\[\int_0^L\sin k_nx\,dx=\left[-\frac{\cos k_nx}{k_n}\right]_0^L=\frac{1-(-1)^n}{k_n}.\]
Nhân $2/L$ cho hệ số của hằng $1$ là $2(1-(-1)^n)/(n\pi)$; hằng $-1$ đổi dấu. Khi chiếu vận tốc, còn chia $\omega_n$ trong hệ số $\sin\omega_nt$; khi chiếu li độ, hệ số $\cos\omega_nt$ giữ nguyên.
\begin{center}\includegraphics[width=.97\linewidth,height=11cm,keepaspectratio]{_xu_ly/graphs/BT_3_2_3_15.png}\end{center}
\textit{Đồ thị minh họa; với nghiệm số, xem phương pháp và tham số trong mã Python kèm theo.}

\clearpage\noindent\texttt{\detokenize{BT_3_2_3_16}}\par
\section*{Lời giải bổ sung — Bài tập 16 · Mục 3.2.3}

\subsection*{Kiến thức và trình tự giải}
\textbf{Bước 1 — Liệt kê các khoảng của dữ liệu.} Ghi đủ điểm đầu, điểm cuối hỗ trợ, điểm gãy và bước nhảy của $f,g$. Nếu có biên, phản xạ chẵn hoặc lẻ trước khi tích phân; trên đoạn hữu hạn tiếp tục phản xạ đến một chu kỳ hoàn chỉnh.

\textbf{Bước 2 — Tính nguyên hàm liên tục.} Lấy $H'=\widetilde g$ trên từng khoảng và chọn hằng số để hai vết của $H$ bằng nhau. Viết $[s]_p^q=\min(q,\max(p,s))$ hoặc $(s)_+=\max(s,0)$ khi cần gộp nhiều nhánh. Một nguyên hàm liên tục tránh tạo ra bước nhảy giả trong nghiệm.

\textbf{Bước 3 — Xác định sóng phải và sóng trái.} Đặt $A(s)=\widetilde f(s)/2-H(s)/(2a)$, $B(s)=\widetilde f(s)/2+H(s)/(2a)$; khi đó $u=A(x-at)+B(x+at)$. Với từng $t$ được yêu cầu, dịch đồ thị $A$ sang phải $at$ và $B$ sang trái $at$, rồi cộng tung độ. Các đường $x-at=b$ và $x+at=b$ chia mặt phẳng thành các miền có công thức đơn giản.

\textbf{Bước 4 — Xét liên tục, đạo hàm và giao của các mặt sóng.} Nếu $f$ có bước nhảy, $u$ thường có bước nhảy; nếu $f$ liên tục và $g$ bị chặn thì $u$ liên tục. Tại thời điểm cố định,


\[
u_x=\frac{\widetilde f'(x-at)+\widetilde f'(x+at)}2+\frac{\widetilde g(x+at)-\widetilde g(x-at)}{2a}.
\]


Tính giới hạn trái/phải của biểu thức này tại từng vị trí ứng viên. Có thể có triệt tiêu khi hai đường gặp nhau, nên không chỉ liệt kê đường rồi khẳng định mọi đạo hàm đều mất. Với tính $C^1$ theo cả $x,t$, còn kiểm tra $u_t$. Với $C^2$, xét tiếp các góc của $g$ và đạo hàm cấp hai của $f$.

\textbf{Bước 5 — Tính thời điểm đến, giới hạn và ngưỡng.} Sóng ở $(x,t)$ chỉ nhận dữ liệu trong $[x-at,x+at]$. Giải bất đẳng thức giao hỗ trợ để tìm thời điểm đến và đi qua. Nếu $H$ hằng ngoài đoạn hỗ trợ, dùng các giá trị hằng này để tính giới hạn khi $t\to\infty$. Nếu hỏi áp suất cực đại, viết hàm từng khoảng thời gian và so sánh cực trị, bước nhảy cùng hai giá trị một phía.

\textbf{Bước 1 — Đọc đúng điều kiện biên.} Biên $u(0,t)=0$ dùng mở rộng lẻ; biên $u_x(0,t)=0$ dùng mở rộng chẵn. Với dữ liệu trên $x\ge0$, đặt $\widetilde f(-s)=-f(s)$ hoặc $f(s)$ tương ứng, làm tương tự cho $g$. Mở rộng lẻ là $C^2$ nếu $f(0)=f''(0)=0$, và mở rộng lẻ của $g$ là $C^1$ nếu $g(0)=0$. Mở rộng chẵn đòi $f'(0)=g'(0)=0$ cho mức trơn ấy.

\textbf{Bước 2 — Lập nguyên hàm sau khi phản xạ.} Đặt $H(s)=\int_0^s\widetilde g(r)dr$. Nếu $\widetilde g$ lẻ thì $H$ chẵn; nếu $\widetilde g$ chẵn thì $H$ lẻ. Chính nguyên hàm, chứ không phải vận tốc, xuất hiện trong


\[
u=\frac{\widetilde f(x-at)+\widetilde f(x+at)}2+\frac{H(x+at)-H(x-at)}{2a}.
\]


\textbf{Bước 3 — Tách miền trước và sau khi chạm biên.} Với $u(0,t)=h(t)$, ở $x\ge at$ dùng d'Alembert với dữ liệu gốc. Ở $0<x<at$:


\[
u=h(t-x/a)+\frac{f(x+at)-f(at-x)}2+\frac1{2a}\int_{at-x}^{at+x}g(s)ds.
\]


Với $u_x(0,t)=h(t)$, dùng nghiệm của các mở rộng chẵn rồi trừ $a\int_0^{t-x/a}h(s)ds$ trên miền $x<at$. Tính riêng đạo hàm theo $x$ của hạng này để kiểm tra biên.

\textbf{Bước 4 — Trên đoạn hữu hạn, phản xạ ở cả hai đầu.} Dirichlet ở hai đầu cho mở rộng lẻ tại $0,L$, chu kỳ không gian $2L$. Neumann ở hai đầu cho mở rộng chẵn tại $0,L$, cũng chu kỳ $2L$. Một đầu lẻ, một đầu chẵn cho $\widetilde f(s+2L)=-\widetilde f(s)$ và chu kỳ $4L$, tương tự $g$. Vì thế chu kỳ thời gian thường là $2L/a$ hoặc $4L/a$. Với Neumann hai đầu còn phải xét trung bình $g$: trung bình khác không sinh hạng tuyến tính theo $t$, nên nghiệm không tuần hoàn.

\textbf{Bước 5 — Kiểm tra góc và đường ghép.} Biên Dirichlet không thuần nhất cần $h(0)=f(0)$, $h'(0)=g(0)$, $h''(0)=a^2f''(0)$ để có nghiệm $C^2$ cho phương trình thuần nhất. So sánh hai vết trên $x=at$, gồm $u,u_x,u_t$ rồi đạo hàm cấp hai. Nếu không tương thích, nêu rõ mức trơn còn lại thay vì khẳng định nghiệm cổ điển trên toàn miền.
\subsection*{Lời giải áp dụng cho bài}
\textbf{Bước 1. Đọc dữ liệu và chọn phản xạ.} Đoạn $(0,\pi)$, Neumann hai đầu, $a=2$, $f=x$, $g=0$.

\textbf{Bước 2. Tính từng nhánh và nối nguyên hàm.} Mở rộng chẵn tại $0,\pi$: trên $[0,2\pi]$, $f_e=s$ trước $\pi$ và $2\pi-s$ sau $\pi$. Đó là hàm tam giác chu kỳ $2\pi$; $H=0$.
\[f_e(s)=\pi-|\operatorname{mod}(s,2\pi)-\pi|,\qquad H(s)=0.\]
Ở đây $\operatorname{clip}(s,p,q)=\min(q,\max(p,s))$; $\operatorname{mod}(s,L)\in[0,L)$ kể cả với $s<0$.

\textbf{Bước 3. Xác định hai sóng chạy và ghép nghiệm.} $A=f_e/2-H/(2a)$, $B=f_e/2+H/(2a)$. Thay hai cận:
\[u(x,t)=\frac{f_e(x-2t)+f_e(x+2t)}2+\frac{H(x+2t)-H(x-2t)}{4}.\]

\textbf{Bước 4. Trả lời các yêu cầu và kiểm tra.} $A_0=\pi/2$. Tích phân từng phần cho $A_n=(2/\pi)((-1)^n-1)/n^2$, chỉ còn $n$ lẻ. Ghép với $\omega_n=2n$ cho chuỗi bên dưới và chu kỳ $\pi$. $f'(0)=f'(\pi)=1$ không thỏa biên Neumann, nên có góc phản xạ.

Tách biến cho $u=\pi/2-(4/\pi)\sum_{n\text{ lẻ}}\cos(2nt)\cos(nx)/n^2$. Chu kỳ $\pi$. Dữ liệu $f=x$ không có đạo hàm biên 0, nên không có nghiệm cổ điển tại các góc.

Tại $t=0$, hiệu nguyên hàm bằng $0$ và đạo hàm thời gian cho $H'(x)=g(x)$ trong miền dữ liệu; half-sum cho đúng li độ. Mỗi hạng là sóng của $x\pm at$, nên thỏa PDE ngoài các mặt gãy và theo nghĩa phân phối trên toàn miền. Đồ thị hai sóng ban đầu và các lát cắt thời gian được tính trực tiếp từ các nhánh trên.

\textbf{Chi tiết hệ số cos của dữ liệu $x$.} Trên $(0,\pi)$, trung bình là $A_0=\pi^{-1}\int_0^\pi x\,dx=\pi/2$. Với $n\ge1$,
\[\int_0^\pi x\cos nx\,dx=\left[\frac{x\sin nx}{n}\right]_0^\pi-\frac1n\int_0^\pi\sin nx\,dx
=\frac{(-1)^n-1}{n^2}.\]
Nhân $2/\pi$ cho $2((-1)^n-1)/(\pi n^2)$. Nếu dữ liệu này là li độ thì giữ hệ số với $\cos2nt$; nếu là vận tốc thì chia $2n$ và nhân $\sin2nt$. Vì $n$ lẻ có $(-1)^n-1=-2$, hệ số lần lượt là $-4/(\pi n^2)$ và $-2/(\pi n^3)$.
\begin{center}\includegraphics[width=.97\linewidth,height=11cm,keepaspectratio]{_xu_ly/graphs/BT_3_2_3_16.png}\end{center}
\textit{Đồ thị minh họa; với nghiệm số, xem phương pháp và tham số trong mã Python kèm theo.}

\clearpage\noindent\texttt{\detokenize{BT_3_2_3_17}}\par
\section*{Lời giải bổ sung — Bài tập 17 · Mục 3.2.3}

\subsection*{Kiến thức và trình tự giải}
\textbf{Bước 1 — Liệt kê các khoảng của dữ liệu.} Ghi đủ điểm đầu, điểm cuối hỗ trợ, điểm gãy và bước nhảy của $f,g$. Nếu có biên, phản xạ chẵn hoặc lẻ trước khi tích phân; trên đoạn hữu hạn tiếp tục phản xạ đến một chu kỳ hoàn chỉnh.

\textbf{Bước 2 — Tính nguyên hàm liên tục.} Lấy $H'=\widetilde g$ trên từng khoảng và chọn hằng số để hai vết của $H$ bằng nhau. Viết $[s]_p^q=\min(q,\max(p,s))$ hoặc $(s)_+=\max(s,0)$ khi cần gộp nhiều nhánh. Một nguyên hàm liên tục tránh tạo ra bước nhảy giả trong nghiệm.

\textbf{Bước 3 — Xác định sóng phải và sóng trái.} Đặt $A(s)=\widetilde f(s)/2-H(s)/(2a)$, $B(s)=\widetilde f(s)/2+H(s)/(2a)$; khi đó $u=A(x-at)+B(x+at)$. Với từng $t$ được yêu cầu, dịch đồ thị $A$ sang phải $at$ và $B$ sang trái $at$, rồi cộng tung độ. Các đường $x-at=b$ và $x+at=b$ chia mặt phẳng thành các miền có công thức đơn giản.

\textbf{Bước 4 — Xét liên tục, đạo hàm và giao của các mặt sóng.} Nếu $f$ có bước nhảy, $u$ thường có bước nhảy; nếu $f$ liên tục và $g$ bị chặn thì $u$ liên tục. Tại thời điểm cố định,


\[
u_x=\frac{\widetilde f'(x-at)+\widetilde f'(x+at)}2+\frac{\widetilde g(x+at)-\widetilde g(x-at)}{2a}.
\]


Tính giới hạn trái/phải của biểu thức này tại từng vị trí ứng viên. Có thể có triệt tiêu khi hai đường gặp nhau, nên không chỉ liệt kê đường rồi khẳng định mọi đạo hàm đều mất. Với tính $C^1$ theo cả $x,t$, còn kiểm tra $u_t$. Với $C^2$, xét tiếp các góc của $g$ và đạo hàm cấp hai của $f$.

\textbf{Bước 5 — Tính thời điểm đến, giới hạn và ngưỡng.} Sóng ở $(x,t)$ chỉ nhận dữ liệu trong $[x-at,x+at]$. Giải bất đẳng thức giao hỗ trợ để tìm thời điểm đến và đi qua. Nếu $H$ hằng ngoài đoạn hỗ trợ, dùng các giá trị hằng này để tính giới hạn khi $t\to\infty$. Nếu hỏi áp suất cực đại, viết hàm từng khoảng thời gian và so sánh cực trị, bước nhảy cùng hai giá trị một phía.

\textbf{Bước 1 — Đọc đúng điều kiện biên.} Biên $u(0,t)=0$ dùng mở rộng lẻ; biên $u_x(0,t)=0$ dùng mở rộng chẵn. Với dữ liệu trên $x\ge0$, đặt $\widetilde f(-s)=-f(s)$ hoặc $f(s)$ tương ứng, làm tương tự cho $g$. Mở rộng lẻ là $C^2$ nếu $f(0)=f''(0)=0$, và mở rộng lẻ của $g$ là $C^1$ nếu $g(0)=0$. Mở rộng chẵn đòi $f'(0)=g'(0)=0$ cho mức trơn ấy.

\textbf{Bước 2 — Lập nguyên hàm sau khi phản xạ.} Đặt $H(s)=\int_0^s\widetilde g(r)dr$. Nếu $\widetilde g$ lẻ thì $H$ chẵn; nếu $\widetilde g$ chẵn thì $H$ lẻ. Chính nguyên hàm, chứ không phải vận tốc, xuất hiện trong


\[
u=\frac{\widetilde f(x-at)+\widetilde f(x+at)}2+\frac{H(x+at)-H(x-at)}{2a}.
\]


\textbf{Bước 3 — Tách miền trước và sau khi chạm biên.} Với $u(0,t)=h(t)$, ở $x\ge at$ dùng d'Alembert với dữ liệu gốc. Ở $0<x<at$:


\[
u=h(t-x/a)+\frac{f(x+at)-f(at-x)}2+\frac1{2a}\int_{at-x}^{at+x}g(s)ds.
\]


Với $u_x(0,t)=h(t)$, dùng nghiệm của các mở rộng chẵn rồi trừ $a\int_0^{t-x/a}h(s)ds$ trên miền $x<at$. Tính riêng đạo hàm theo $x$ của hạng này để kiểm tra biên.

\textbf{Bước 4 — Trên đoạn hữu hạn, phản xạ ở cả hai đầu.} Dirichlet ở hai đầu cho mở rộng lẻ tại $0,L$, chu kỳ không gian $2L$. Neumann ở hai đầu cho mở rộng chẵn tại $0,L$, cũng chu kỳ $2L$. Một đầu lẻ, một đầu chẵn cho $\widetilde f(s+2L)=-\widetilde f(s)$ và chu kỳ $4L$, tương tự $g$. Vì thế chu kỳ thời gian thường là $2L/a$ hoặc $4L/a$. Với Neumann hai đầu còn phải xét trung bình $g$: trung bình khác không sinh hạng tuyến tính theo $t$, nên nghiệm không tuần hoàn.

\textbf{Bước 5 — Kiểm tra góc và đường ghép.} Biên Dirichlet không thuần nhất cần $h(0)=f(0)$, $h'(0)=g(0)$, $h''(0)=a^2f''(0)$ để có nghiệm $C^2$ cho phương trình thuần nhất. So sánh hai vết trên $x=at$, gồm $u,u_x,u_t$ rồi đạo hàm cấp hai. Nếu không tương thích, nêu rõ mức trơn còn lại thay vì khẳng định nghiệm cổ điển trên toàn miền.
\subsection*{Lời giải áp dụng cho bài}
\textbf{Bước 1. Đọc dữ liệu và chọn phản xạ.} Đoạn $(0,1)$, Neumann tại $0$, Dirichlet tại $1$, $a=2$, $f=1-x$, $g=0$.

\textbf{Bước 2. Tính từng nhánh và nối nguyên hàm.} Phản xạ chẵn tại $0$, lẻ tại $1$ cho trên $[0,4]$ hai nhánh $1-s$ ở $[0,2]$, $s-3$ ở $[2,4]$. Vì thế có công thức chính xác $f_e(s)=|\operatorname{mod}(s,4)-2|-1$.
\[f_e(s)=|\operatorname{mod}(s,4)-2|-1,\qquad H(s)=0.\]
Ở đây $\operatorname{clip}(s,p,q)=\min(q,\max(p,s))$; $\operatorname{mod}(s,L)\in[0,L)$ kể cả với $s<0$.

\textbf{Bước 3. Xác định hai sóng chạy và ghép nghiệm.} $A=f_e/2-H/(2a)$, $B=f_e/2+H/(2a)$. Thay hai cận:
\[u(x,t)=\frac{f_e(x-2t)+f_e(x+2t)}2+\frac{H(x+2t)-H(x-2t)}{4}.\]

\textbf{Bước 4. Trả lời các yêu cầu và kiểm tra.} $f_e(s+2)=-f_e(s)$ nên $u(t+1)=-u(t)$, chu kỳ $2$. Với $k_n=(n+1/2)\pi$, tích phân $2\int_0^1(1-x)\cos(k_nx)dx=2/k_n^2$ cho chuỗi. Góc tại phản xạ đầu $0$ xuất phát từ $f'(0)=-1$ không tương thích.

Mô hình $u_{tt}=4u_{xx}$, $u_x(0,t)=0,u(1,t)=0$, $f=1-x,g=0$. Với $k_n=(n+1/2)\pi$, nghiệm $u=\sum_{n\ge0}2\cos(2k_nt)\cos(k_nx)/k_n^2$. Chu kỳ 2; dữ liệu không tương thích đạo hàm tại đầu tự do.

Tại $t=0$, hiệu nguyên hàm bằng $0$ và đạo hàm thời gian cho $H'(x)=g(x)$ trong miền dữ liệu; half-sum cho đúng li độ. Mỗi hạng là sóng của $x\pm at$, nên thỏa PDE ngoài các mặt gãy và theo nghĩa phân phối trên toàn miền. Đồ thị hai sóng ban đầu và các lát cắt thời gian được tính trực tiếp từ các nhánh trên.

\textbf{Chi tiết chiếu li độ lên hệ cos nửa nguyên.} $k_n=(n+1/2)\pi$, $\cos k_n=0$, chuẩn $\int_0^1\cos^2k_nx\,dx=1/2$. Tích phân từng phần:
\[\int_0^1(1-x)\cos k_nx\,dx=\left[\frac{(1-x)\sin k_nx}{k_n}\right]_0^1+\frac1{k_n}\int_0^1\sin k_nx\,dx
=\frac{1-\cos k_n}{k_n^2}=\frac1{k_n^2}.\]
Chia chuẩn $1/2$ cho hệ số $2/k_n^2$; vận tốc đầu bằng $0$ nên không có phần sine theo thời gian.
\begin{center}\includegraphics[width=.97\linewidth,height=11cm,keepaspectratio]{_xu_ly/graphs/BT_3_2_3_17.png}\end{center}
\textit{Đồ thị minh họa; với nghiệm số, xem phương pháp và tham số trong mã Python kèm theo.}

\clearpage\noindent\texttt{\detokenize{BT_3_2_3_18}}\par
\section*{Lời giải bổ sung — Bài tập 18 · Mục 3.2.3}

\subsection*{Kiến thức và trình tự giải}
\textbf{Bước 1 — Liệt kê các khoảng của dữ liệu.} Ghi đủ điểm đầu, điểm cuối hỗ trợ, điểm gãy và bước nhảy của $f,g$. Nếu có biên, phản xạ chẵn hoặc lẻ trước khi tích phân; trên đoạn hữu hạn tiếp tục phản xạ đến một chu kỳ hoàn chỉnh.

\textbf{Bước 2 — Tính nguyên hàm liên tục.} Lấy $H'=\widetilde g$ trên từng khoảng và chọn hằng số để hai vết của $H$ bằng nhau. Viết $[s]_p^q=\min(q,\max(p,s))$ hoặc $(s)_+=\max(s,0)$ khi cần gộp nhiều nhánh. Một nguyên hàm liên tục tránh tạo ra bước nhảy giả trong nghiệm.

\textbf{Bước 3 — Xác định sóng phải và sóng trái.} Đặt $A(s)=\widetilde f(s)/2-H(s)/(2a)$, $B(s)=\widetilde f(s)/2+H(s)/(2a)$; khi đó $u=A(x-at)+B(x+at)$. Với từng $t$ được yêu cầu, dịch đồ thị $A$ sang phải $at$ và $B$ sang trái $at$, rồi cộng tung độ. Các đường $x-at=b$ và $x+at=b$ chia mặt phẳng thành các miền có công thức đơn giản.

\textbf{Bước 4 — Xét liên tục, đạo hàm và giao của các mặt sóng.} Nếu $f$ có bước nhảy, $u$ thường có bước nhảy; nếu $f$ liên tục và $g$ bị chặn thì $u$ liên tục. Tại thời điểm cố định,


\[
u_x=\frac{\widetilde f'(x-at)+\widetilde f'(x+at)}2+\frac{\widetilde g(x+at)-\widetilde g(x-at)}{2a}.
\]


Tính giới hạn trái/phải của biểu thức này tại từng vị trí ứng viên. Có thể có triệt tiêu khi hai đường gặp nhau, nên không chỉ liệt kê đường rồi khẳng định mọi đạo hàm đều mất. Với tính $C^1$ theo cả $x,t$, còn kiểm tra $u_t$. Với $C^2$, xét tiếp các góc của $g$ và đạo hàm cấp hai của $f$.

\textbf{Bước 5 — Tính thời điểm đến, giới hạn và ngưỡng.} Sóng ở $(x,t)$ chỉ nhận dữ liệu trong $[x-at,x+at]$. Giải bất đẳng thức giao hỗ trợ để tìm thời điểm đến và đi qua. Nếu $H$ hằng ngoài đoạn hỗ trợ, dùng các giá trị hằng này để tính giới hạn khi $t\to\infty$. Nếu hỏi áp suất cực đại, viết hàm từng khoảng thời gian và so sánh cực trị, bước nhảy cùng hai giá trị một phía.

\textbf{Bước 1 — Đọc đúng điều kiện biên.} Biên $u(0,t)=0$ dùng mở rộng lẻ; biên $u_x(0,t)=0$ dùng mở rộng chẵn. Với dữ liệu trên $x\ge0$, đặt $\widetilde f(-s)=-f(s)$ hoặc $f(s)$ tương ứng, làm tương tự cho $g$. Mở rộng lẻ là $C^2$ nếu $f(0)=f''(0)=0$, và mở rộng lẻ của $g$ là $C^1$ nếu $g(0)=0$. Mở rộng chẵn đòi $f'(0)=g'(0)=0$ cho mức trơn ấy.

\textbf{Bước 2 — Lập nguyên hàm sau khi phản xạ.} Đặt $H(s)=\int_0^s\widetilde g(r)dr$. Nếu $\widetilde g$ lẻ thì $H$ chẵn; nếu $\widetilde g$ chẵn thì $H$ lẻ. Chính nguyên hàm, chứ không phải vận tốc, xuất hiện trong


\[
u=\frac{\widetilde f(x-at)+\widetilde f(x+at)}2+\frac{H(x+at)-H(x-at)}{2a}.
\]


\textbf{Bước 3 — Tách miền trước và sau khi chạm biên.} Với $u(0,t)=h(t)$, ở $x\ge at$ dùng d'Alembert với dữ liệu gốc. Ở $0<x<at$:


\[
u=h(t-x/a)+\frac{f(x+at)-f(at-x)}2+\frac1{2a}\int_{at-x}^{at+x}g(s)ds.
\]


Với $u_x(0,t)=h(t)$, dùng nghiệm của các mở rộng chẵn rồi trừ $a\int_0^{t-x/a}h(s)ds$ trên miền $x<at$. Tính riêng đạo hàm theo $x$ của hạng này để kiểm tra biên.

\textbf{Bước 4 — Trên đoạn hữu hạn, phản xạ ở cả hai đầu.} Dirichlet ở hai đầu cho mở rộng lẻ tại $0,L$, chu kỳ không gian $2L$. Neumann ở hai đầu cho mở rộng chẵn tại $0,L$, cũng chu kỳ $2L$. Một đầu lẻ, một đầu chẵn cho $\widetilde f(s+2L)=-\widetilde f(s)$ và chu kỳ $4L$, tương tự $g$. Vì thế chu kỳ thời gian thường là $2L/a$ hoặc $4L/a$. Với Neumann hai đầu còn phải xét trung bình $g$: trung bình khác không sinh hạng tuyến tính theo $t$, nên nghiệm không tuần hoàn.

\textbf{Bước 5 — Kiểm tra góc và đường ghép.} Biên Dirichlet không thuần nhất cần $h(0)=f(0)$, $h'(0)=g(0)$, $h''(0)=a^2f''(0)$ để có nghiệm $C^2$ cho phương trình thuần nhất. So sánh hai vết trên $x=at$, gồm $u,u_x,u_t$ rồi đạo hàm cấp hai. Nếu không tương thích, nêu rõ mức trơn còn lại thay vì khẳng định nghiệm cổ điển trên toàn miền.
\subsection*{Lời giải áp dụng cho bài}
\textbf{Bước 1. Đọc dữ liệu và chọn phản xạ.} Đoạn $(0,1)$, $a=2$, hai biên $u=1$, dữ liệu đầu $f=0$, $g=1$.

\textbf{Bước 2. Tính từng nhánh và nối nguyên hàm.} Lấy hàm nâng $b=1$, $w=u-1$ có biên không, $w_0=-1$, $w_{t,0}=1$. Mở rộng lẻ chu kỳ $2$ của $-1$ và nguyên hàm tam giác của vận tốc được dùng trong công thức; cuối cùng phải cộng lại $1$.
\[f_e(s)=\widetilde f=-1\text{ trên }(0,1),\text{ thác triển lẻ chu kỳ }2,\qquad H(s)=1-|\operatorname{mod}(s,2)-1|.\]
Chọn giá trị mở rộng li độ bằng $0$ tại các số nguyên (trung bình hai vết), để công thức có đúng giá trị biên tại cả thời điểm hai mặt sóng gặp nhau.
Ở đây $\operatorname{clip}(s,p,q)=\min(q,\max(p,s))$; $\operatorname{mod}(s,L)\in[0,L)$ kể cả với $s<0$.

\textbf{Bước 3. Xác định hai sóng chạy và ghép nghiệm.} $A=f_e/2-H/(2a)$, $B=f_e/2+H/(2a)$. Thay hai cận:
\[u(x,t)=\frac{f_e(x-2t)+f_e(x+2t)}2+\frac{H(x+2t)-H(x-2t)}{4}+1.\]

\textbf{Bước 4. Trả lời các yêu cầu và kiểm tra.} Hệ số li độ $A_n=-2(1-(-1)^n)/(n\pi)$, vận tốc $B_n=2(1-(-1)^n)/(n\pi)$, tần số $2n\pi$. Dữ liệu li độ và biên mâu thuẫn ngay ở góc nên không có nghiệm cổ điển trên bao đóng; chuỗi và sóng chạy cho nghiệm trong miền với các bước nhảy.

Chọn $v=1$, $w=u-1$ có biên 0, $w(x,0)=-1,w_t(x,0)=1$. Nghiệm $u=1+\sum_{n\text{ lẻ}}[-4\cos(2n\pi t)/(n\pi)+2\sin(2n\pi t)/(n^2\pi^2)]\sin(n\pi x)$. Góc dữ liệu $u(x,0)=0$ mâu thuẫn biên 1: chỉ có nghiệm yếu với mặt sóng.

Tại $t=0$, hiệu nguyên hàm bằng $0$ và đạo hàm thời gian cho $H'(x)=g(x)$ trong miền dữ liệu; half-sum cho đúng li độ sau khi cộng hàm nâng biên. Mỗi hạng là sóng của $x\pm at$, nên thỏa PDE ngoài các mặt gãy và theo nghĩa phân phối trên toàn miền. Đồ thị hai sóng ban đầu và các lát cắt thời gian được tính trực tiếp từ các nhánh trên.

\textbf{Chi tiết tích phân của dữ liệu hằng.} Với $k_n=n\pi/L$,
\[\int_0^L\sin k_nx\,dx=\left[-\frac{\cos k_nx}{k_n}\right]_0^L=\frac{1-(-1)^n}{k_n}.\]
Nhân $2/L$ cho hệ số của hằng $1$ là $2(1-(-1)^n)/(n\pi)$; hằng $-1$ đổi dấu. Khi chiếu vận tốc, còn chia $\omega_n$ trong hệ số $\sin\omega_nt$; khi chiếu li độ, hệ số $\cos\omega_nt$ giữ nguyên.
\begin{center}\includegraphics[width=.97\linewidth,height=11cm,keepaspectratio]{_xu_ly/graphs/BT_3_2_3_18.png}\end{center}
\textit{Đồ thị minh họa; với nghiệm số, xem phương pháp và tham số trong mã Python kèm theo.}

\clearpage\noindent\texttt{\detokenize{BT_3_2_3_19}}\par
\section*{Lời giải bổ sung — Bài tập 19 · Mục 3.2.3}

\subsection*{Kiến thức và trình tự giải}
\textbf{Bước 1 — Liệt kê các khoảng của dữ liệu.} Ghi đủ điểm đầu, điểm cuối hỗ trợ, điểm gãy và bước nhảy của $f,g$. Nếu có biên, phản xạ chẵn hoặc lẻ trước khi tích phân; trên đoạn hữu hạn tiếp tục phản xạ đến một chu kỳ hoàn chỉnh.

\textbf{Bước 2 — Tính nguyên hàm liên tục.} Lấy $H'=\widetilde g$ trên từng khoảng và chọn hằng số để hai vết của $H$ bằng nhau. Viết $[s]_p^q=\min(q,\max(p,s))$ hoặc $(s)_+=\max(s,0)$ khi cần gộp nhiều nhánh. Một nguyên hàm liên tục tránh tạo ra bước nhảy giả trong nghiệm.

\textbf{Bước 3 — Xác định sóng phải và sóng trái.} Đặt $A(s)=\widetilde f(s)/2-H(s)/(2a)$, $B(s)=\widetilde f(s)/2+H(s)/(2a)$; khi đó $u=A(x-at)+B(x+at)$. Với từng $t$ được yêu cầu, dịch đồ thị $A$ sang phải $at$ và $B$ sang trái $at$, rồi cộng tung độ. Các đường $x-at=b$ và $x+at=b$ chia mặt phẳng thành các miền có công thức đơn giản.

\textbf{Bước 4 — Xét liên tục, đạo hàm và giao của các mặt sóng.} Nếu $f$ có bước nhảy, $u$ thường có bước nhảy; nếu $f$ liên tục và $g$ bị chặn thì $u$ liên tục. Tại thời điểm cố định,


\[
u_x=\frac{\widetilde f'(x-at)+\widetilde f'(x+at)}2+\frac{\widetilde g(x+at)-\widetilde g(x-at)}{2a}.
\]


Tính giới hạn trái/phải của biểu thức này tại từng vị trí ứng viên. Có thể có triệt tiêu khi hai đường gặp nhau, nên không chỉ liệt kê đường rồi khẳng định mọi đạo hàm đều mất. Với tính $C^1$ theo cả $x,t$, còn kiểm tra $u_t$. Với $C^2$, xét tiếp các góc của $g$ và đạo hàm cấp hai của $f$.

\textbf{Bước 5 — Tính thời điểm đến, giới hạn và ngưỡng.} Sóng ở $(x,t)$ chỉ nhận dữ liệu trong $[x-at,x+at]$. Giải bất đẳng thức giao hỗ trợ để tìm thời điểm đến và đi qua. Nếu $H$ hằng ngoài đoạn hỗ trợ, dùng các giá trị hằng này để tính giới hạn khi $t\to\infty$. Nếu hỏi áp suất cực đại, viết hàm từng khoảng thời gian và so sánh cực trị, bước nhảy cùng hai giá trị một phía.

\textbf{Bước 1 — Đọc đúng điều kiện biên.} Biên $u(0,t)=0$ dùng mở rộng lẻ; biên $u_x(0,t)=0$ dùng mở rộng chẵn. Với dữ liệu trên $x\ge0$, đặt $\widetilde f(-s)=-f(s)$ hoặc $f(s)$ tương ứng, làm tương tự cho $g$. Mở rộng lẻ là $C^2$ nếu $f(0)=f''(0)=0$, và mở rộng lẻ của $g$ là $C^1$ nếu $g(0)=0$. Mở rộng chẵn đòi $f'(0)=g'(0)=0$ cho mức trơn ấy.

\textbf{Bước 2 — Lập nguyên hàm sau khi phản xạ.} Đặt $H(s)=\int_0^s\widetilde g(r)dr$. Nếu $\widetilde g$ lẻ thì $H$ chẵn; nếu $\widetilde g$ chẵn thì $H$ lẻ. Chính nguyên hàm, chứ không phải vận tốc, xuất hiện trong


\[
u=\frac{\widetilde f(x-at)+\widetilde f(x+at)}2+\frac{H(x+at)-H(x-at)}{2a}.
\]


\textbf{Bước 3 — Tách miền trước và sau khi chạm biên.} Với $u(0,t)=h(t)$, ở $x\ge at$ dùng d'Alembert với dữ liệu gốc. Ở $0<x<at$:


\[
u=h(t-x/a)+\frac{f(x+at)-f(at-x)}2+\frac1{2a}\int_{at-x}^{at+x}g(s)ds.
\]


Với $u_x(0,t)=h(t)$, dùng nghiệm của các mở rộng chẵn rồi trừ $a\int_0^{t-x/a}h(s)ds$ trên miền $x<at$. Tính riêng đạo hàm theo $x$ của hạng này để kiểm tra biên.

\textbf{Bước 4 — Trên đoạn hữu hạn, phản xạ ở cả hai đầu.} Dirichlet ở hai đầu cho mở rộng lẻ tại $0,L$, chu kỳ không gian $2L$. Neumann ở hai đầu cho mở rộng chẵn tại $0,L$, cũng chu kỳ $2L$. Một đầu lẻ, một đầu chẵn cho $\widetilde f(s+2L)=-\widetilde f(s)$ và chu kỳ $4L$, tương tự $g$. Vì thế chu kỳ thời gian thường là $2L/a$ hoặc $4L/a$. Với Neumann hai đầu còn phải xét trung bình $g$: trung bình khác không sinh hạng tuyến tính theo $t$, nên nghiệm không tuần hoàn.

\textbf{Bước 5 — Kiểm tra góc và đường ghép.} Biên Dirichlet không thuần nhất cần $h(0)=f(0)$, $h'(0)=g(0)$, $h''(0)=a^2f''(0)$ để có nghiệm $C^2$ cho phương trình thuần nhất. So sánh hai vết trên $x=at$, gồm $u,u_x,u_t$ rồi đạo hàm cấp hai. Nếu không tương thích, nêu rõ mức trơn còn lại thay vì khẳng định nghiệm cổ điển trên toàn miền.
\subsection*{Lời giải áp dụng cho bài}
\textbf{Bước 1. Đọc dữ liệu và chọn phản xạ.} Đoạn $(0,\pi)$, $a=2$, $u_x(0,t)=u_x(\pi,t)=1$, $f=x$, $g=x$.

\textbf{Bước 2. Tính từng nhánh và nối nguyên hàm.} Lấy $b=x$. Phần $w=u-x$ có li độ $0$, vận tốc $x$ và biên Neumann không. Mở rộng chẵn vận tốc thành tam giác $q(s)=\pi-|\operatorname{mod}(s,2\pi)-\pi|$. Tích phân từng nhánh cho $H$ ở dưới, mỗi nửa chu kỳ góp $\pi^2/2$.
\[f_e(s)=0,\qquad H(s)=k\pi^2/2+J(r),\ r=s-k\pi,\ k=\lfloor s/\pi\rfloor;\ J(r)=r^2/2\ (k\text{ chẵn}),\ \pi r-r^2/2\ (k\text{ lẻ}).\]
Ở đây $\operatorname{clip}(s,p,q)=\min(q,\max(p,s))$; $\operatorname{mod}(s,L)\in[0,L)$ kể cả với $s<0$.

\textbf{Bước 3. Xác định hai sóng chạy và ghép nghiệm.} $A=f_e/2-H/(2a)$, $B=f_e/2+H/(2a)$. Thay hai cận:
\[u(x,t)=\frac{f_e(x-2t)+f_e(x+2t)}2+\frac{H(x+2t)-H(x-2t)}{4}+x.\]

\textbf{Bước 4. Trả lời các yêu cầu và kiểm tra.} Trung bình vận tốc là $\pi/2$, sinh hạng $\pi t/2$. Các mode còn lại có $B_n=2(( -1)^n-1)/(\pi n^2)$, chia $2n$ cho hệ số thời gian. Nghiệm không tuần hoàn; tương thích cấp hai ở góc không đúng vì $g'=1$ nhưng đạo hàm thời gian của biên bằng $0$.

Chọn $v=x$. $w=u-x$ có biên Neumann 0, dữ liệu $w=0,w_t=x$. Nghiệm $u=x+\pi t/2-(2/\pi)\sum_{n\text{ lẻ}}\sin(2nt)\cos(nx)/n^3$. Có mode tuyến tính nên không tuần hoàn.

Tại $t=0$, hiệu nguyên hàm bằng $0$ và đạo hàm thời gian cho $H'(x)=g(x)$ trong miền dữ liệu; half-sum cho đúng li độ sau khi cộng hàm nâng biên. Mỗi hạng là sóng của $x\pm at$, nên thỏa PDE ngoài các mặt gãy và theo nghĩa phân phối trên toàn miền. Đồ thị hai sóng ban đầu và các lát cắt thời gian được tính trực tiếp từ các nhánh trên.

\textbf{Chi tiết hệ số cos của dữ liệu $x$.} Trên $(0,\pi)$, trung bình là $A_0=\pi^{-1}\int_0^\pi x\,dx=\pi/2$. Với $n\ge1$,
\[\int_0^\pi x\cos nx\,dx=\left[\frac{x\sin nx}{n}\right]_0^\pi-\frac1n\int_0^\pi\sin nx\,dx
=\frac{(-1)^n-1}{n^2}.\]
Nhân $2/\pi$ cho $2((-1)^n-1)/(\pi n^2)$. Nếu dữ liệu này là li độ thì giữ hệ số với $\cos2nt$; nếu là vận tốc thì chia $2n$ và nhân $\sin2nt$. Vì $n$ lẻ có $(-1)^n-1=-2$, hệ số lần lượt là $-4/(\pi n^2)$ và $-2/(\pi n^3)$.
\begin{center}\includegraphics[width=.97\linewidth,height=11cm,keepaspectratio]{_xu_ly/graphs/BT_3_2_3_19.png}\end{center}
\textit{Đồ thị minh họa; với nghiệm số, xem phương pháp và tham số trong mã Python kèm theo.}

\clearpage\noindent\texttt{\detokenize{BT_3_2_3_20}}\par
\section*{Lời giải bổ sung — Bài tập 20 · Mục 3.2.3}

\subsection*{Kiến thức và trình tự giải}
\textbf{Bước 1 — Liệt kê các khoảng của dữ liệu.} Ghi đủ điểm đầu, điểm cuối hỗ trợ, điểm gãy và bước nhảy của $f,g$. Nếu có biên, phản xạ chẵn hoặc lẻ trước khi tích phân; trên đoạn hữu hạn tiếp tục phản xạ đến một chu kỳ hoàn chỉnh.

\textbf{Bước 2 — Tính nguyên hàm liên tục.} Lấy $H'=\widetilde g$ trên từng khoảng và chọn hằng số để hai vết của $H$ bằng nhau. Viết $[s]_p^q=\min(q,\max(p,s))$ hoặc $(s)_+=\max(s,0)$ khi cần gộp nhiều nhánh. Một nguyên hàm liên tục tránh tạo ra bước nhảy giả trong nghiệm.

\textbf{Bước 3 — Xác định sóng phải và sóng trái.} Đặt $A(s)=\widetilde f(s)/2-H(s)/(2a)$, $B(s)=\widetilde f(s)/2+H(s)/(2a)$; khi đó $u=A(x-at)+B(x+at)$. Với từng $t$ được yêu cầu, dịch đồ thị $A$ sang phải $at$ và $B$ sang trái $at$, rồi cộng tung độ. Các đường $x-at=b$ và $x+at=b$ chia mặt phẳng thành các miền có công thức đơn giản.

\textbf{Bước 4 — Xét liên tục, đạo hàm và giao của các mặt sóng.} Nếu $f$ có bước nhảy, $u$ thường có bước nhảy; nếu $f$ liên tục và $g$ bị chặn thì $u$ liên tục. Tại thời điểm cố định,


\[
u_x=\frac{\widetilde f'(x-at)+\widetilde f'(x+at)}2+\frac{\widetilde g(x+at)-\widetilde g(x-at)}{2a}.
\]


Tính giới hạn trái/phải của biểu thức này tại từng vị trí ứng viên. Có thể có triệt tiêu khi hai đường gặp nhau, nên không chỉ liệt kê đường rồi khẳng định mọi đạo hàm đều mất. Với tính $C^1$ theo cả $x,t$, còn kiểm tra $u_t$. Với $C^2$, xét tiếp các góc của $g$ và đạo hàm cấp hai của $f$.

\textbf{Bước 5 — Tính thời điểm đến, giới hạn và ngưỡng.} Sóng ở $(x,t)$ chỉ nhận dữ liệu trong $[x-at,x+at]$. Giải bất đẳng thức giao hỗ trợ để tìm thời điểm đến và đi qua. Nếu $H$ hằng ngoài đoạn hỗ trợ, dùng các giá trị hằng này để tính giới hạn khi $t\to\infty$. Nếu hỏi áp suất cực đại, viết hàm từng khoảng thời gian và so sánh cực trị, bước nhảy cùng hai giá trị một phía.

\textbf{Bước 1 — Đọc đúng điều kiện biên.} Biên $u(0,t)=0$ dùng mở rộng lẻ; biên $u_x(0,t)=0$ dùng mở rộng chẵn. Với dữ liệu trên $x\ge0$, đặt $\widetilde f(-s)=-f(s)$ hoặc $f(s)$ tương ứng, làm tương tự cho $g$. Mở rộng lẻ là $C^2$ nếu $f(0)=f''(0)=0$, và mở rộng lẻ của $g$ là $C^1$ nếu $g(0)=0$. Mở rộng chẵn đòi $f'(0)=g'(0)=0$ cho mức trơn ấy.

\textbf{Bước 2 — Lập nguyên hàm sau khi phản xạ.} Đặt $H(s)=\int_0^s\widetilde g(r)dr$. Nếu $\widetilde g$ lẻ thì $H$ chẵn; nếu $\widetilde g$ chẵn thì $H$ lẻ. Chính nguyên hàm, chứ không phải vận tốc, xuất hiện trong


\[
u=\frac{\widetilde f(x-at)+\widetilde f(x+at)}2+\frac{H(x+at)-H(x-at)}{2a}.
\]


\textbf{Bước 3 — Tách miền trước và sau khi chạm biên.} Với $u(0,t)=h(t)$, ở $x\ge at$ dùng d'Alembert với dữ liệu gốc. Ở $0<x<at$:


\[
u=h(t-x/a)+\frac{f(x+at)-f(at-x)}2+\frac1{2a}\int_{at-x}^{at+x}g(s)ds.
\]


Với $u_x(0,t)=h(t)$, dùng nghiệm của các mở rộng chẵn rồi trừ $a\int_0^{t-x/a}h(s)ds$ trên miền $x<at$. Tính riêng đạo hàm theo $x$ của hạng này để kiểm tra biên.

\textbf{Bước 4 — Trên đoạn hữu hạn, phản xạ ở cả hai đầu.} Dirichlet ở hai đầu cho mở rộng lẻ tại $0,L$, chu kỳ không gian $2L$. Neumann ở hai đầu cho mở rộng chẵn tại $0,L$, cũng chu kỳ $2L$. Một đầu lẻ, một đầu chẵn cho $\widetilde f(s+2L)=-\widetilde f(s)$ và chu kỳ $4L$, tương tự $g$. Vì thế chu kỳ thời gian thường là $2L/a$ hoặc $4L/a$. Với Neumann hai đầu còn phải xét trung bình $g$: trung bình khác không sinh hạng tuyến tính theo $t$, nên nghiệm không tuần hoàn.

\textbf{Bước 5 — Kiểm tra góc và đường ghép.} Biên Dirichlet không thuần nhất cần $h(0)=f(0)$, $h'(0)=g(0)$, $h''(0)=a^2f''(0)$ để có nghiệm $C^2$ cho phương trình thuần nhất. So sánh hai vết trên $x=at$, gồm $u,u_x,u_t$ rồi đạo hàm cấp hai. Nếu không tương thích, nêu rõ mức trơn còn lại thay vì khẳng định nghiệm cổ điển trên toàn miền.
\subsection*{Lời giải áp dụng cho bài}
\textbf{Bước 1. Đọc dữ liệu và chọn phản xạ.} Đoạn $(0,\pi)$, Dirichlet hai đầu, $a=2$, $f=0$, vận tốc $g$; minh họa $g=\min(x,\pi-x)$.

\textbf{Bước 2. Tính từng nhánh và nối nguyên hàm.} Với $g$ tổng quát, mở rộng lẻ chu kỳ $2\pi$. Với tam giác cụ thể, tích phân lần lượt trên $[0,\pi/2]$ và $[\pi/2,\pi]$, rồi phản xạ nguyên hàm chẵn và tuần hoàn, được $H$ dưới đây.
\[f_e(s)=0,\qquad H(s)=J(r),\quad r=\operatorname{mod}(s,2\pi).\]
\[J(r)=\begin{cases}r^2/2&0\le r\le\pi/2,\\\pi r-r^2/2-\pi^2/4&\pi/2<r\le\pi,\\J(2\pi-r)&\pi<r<2\pi.\end{cases}\]
Ở đây $\operatorname{clip}(s,p,q)=\min(q,\max(p,s))$; $\operatorname{mod}(s,L)\in[0,L)$ kể cả với $s<0$.

\textbf{Bước 3. Xác định hai sóng chạy và ghép nghiệm.} $A=f_e/2-H/(2a)$, $B=f_e/2+H/(2a)$. Thay hai cận:
\[u(x,t)=\frac{f_e(x-2t)+f_e(x+2t)}2+\frac{H(x+2t)-H(x-2t)}{4}.\]

\textbf{Bước 4. Trả lời các yêu cầu và kiểm tra.} Lấy đạo hàm ở biên: $u_x(0,t)=\widetilde g(2t)/2$. Đổi biến $s=2t$ cho $\int_0^{\pi/2}|u_x(0,t)|^2dt=\frac18\int_0^\pi g(s)^2ds$. Tích phân trên thời gian lớn hơn chỉ thêm số không âm, chứng minh yêu cầu với $T\ge\pi/2$. Hai tích phân của tam giác cho $B_n=4\sin(n\pi/2)/(\pi n^2)$, nên chia $2n$ cho chuỗi.

Mô hình biên cố định $u_{tt}=4u_{xx}$ trên $(0,\pi)$, $f=0,u_t=g$. $u_x(0,t)=\widetilde g(2t)/2$. Do đó $\int_0^{\pi/2}|u_x(0,t)|^2dt=\frac18\int_0^\pi g^2$, suy bất đẳng thức với $T\ge\pi/2$. Với $g=\min(x,\pi-x)$, $u=\sum_{n\ge1}2\sin(n\pi/2)\sin(2nt)\sin(nx)/(\pi n^3)$.

Tại $t=0$, hiệu nguyên hàm bằng $0$ và đạo hàm thời gian cho $H'(x)=g(x)$ trong miền dữ liệu; half-sum cho đúng li độ. Mỗi hạng là sóng của $x\pm at$, nên thỏa PDE ngoài các mặt gãy và theo nghĩa phân phối trên toàn miền. Đồ thị hai sóng ban đầu và các lát cắt thời gian được tính trực tiếp từ các nhánh trên.

\textbf{Chi tiết hệ số của tam giác.} Với $f(x)=\min(x,\pi-x)$, chia tại $b=\pi/2$. Đổi $y=\pi-x$ ở nửa sau cho
\[\int_0^\pi f(x)\sin nx\,dx=(1+(-1)^{n+1})\int_0^b x\sin nx\,dx.\]
Tích phân cuối bằng $[-x\cos nx/n+\sin nx/n^2]_0^b=-b\cos(nb)/n+\sin(nb)/n^2$. Nếu $n$ chẵn, hệ số trước tích phân bằng $0$; nếu $n$ lẻ, $\cos(nb)=0$. Nhân $2/\pi$ được $4\sin(n\pi/2)/(\pi n^2)$.
\begin{center}\includegraphics[width=.97\linewidth,height=11cm,keepaspectratio]{_xu_ly/graphs/BT_3_2_3_20.png}\end{center}
\textit{Đồ thị minh họa; với nghiệm số, xem phương pháp và tham số trong mã Python kèm theo.}

\clearpage\noindent\texttt{\detokenize{BT_3_2_3_21}}\par
\section*{Lời giải bổ sung — Bài tập 21 · Mục 3.2.3}

\subsection*{Kiến thức và trình tự giải}
\textbf{Bước 1 — Liệt kê các khoảng của dữ liệu.} Ghi đủ điểm đầu, điểm cuối hỗ trợ, điểm gãy và bước nhảy của $f,g$. Nếu có biên, phản xạ chẵn hoặc lẻ trước khi tích phân; trên đoạn hữu hạn tiếp tục phản xạ đến một chu kỳ hoàn chỉnh.

\textbf{Bước 2 — Tính nguyên hàm liên tục.} Lấy $H'=\widetilde g$ trên từng khoảng và chọn hằng số để hai vết của $H$ bằng nhau. Viết $[s]_p^q=\min(q,\max(p,s))$ hoặc $(s)_+=\max(s,0)$ khi cần gộp nhiều nhánh. Một nguyên hàm liên tục tránh tạo ra bước nhảy giả trong nghiệm.

\textbf{Bước 3 — Xác định sóng phải và sóng trái.} Đặt $A(s)=\widetilde f(s)/2-H(s)/(2a)$, $B(s)=\widetilde f(s)/2+H(s)/(2a)$; khi đó $u=A(x-at)+B(x+at)$. Với từng $t$ được yêu cầu, dịch đồ thị $A$ sang phải $at$ và $B$ sang trái $at$, rồi cộng tung độ. Các đường $x-at=b$ và $x+at=b$ chia mặt phẳng thành các miền có công thức đơn giản.

\textbf{Bước 4 — Xét liên tục, đạo hàm và giao của các mặt sóng.} Nếu $f$ có bước nhảy, $u$ thường có bước nhảy; nếu $f$ liên tục và $g$ bị chặn thì $u$ liên tục. Tại thời điểm cố định,


\[
u_x=\frac{\widetilde f'(x-at)+\widetilde f'(x+at)}2+\frac{\widetilde g(x+at)-\widetilde g(x-at)}{2a}.
\]


Tính giới hạn trái/phải của biểu thức này tại từng vị trí ứng viên. Có thể có triệt tiêu khi hai đường gặp nhau, nên không chỉ liệt kê đường rồi khẳng định mọi đạo hàm đều mất. Với tính $C^1$ theo cả $x,t$, còn kiểm tra $u_t$. Với $C^2$, xét tiếp các góc của $g$ và đạo hàm cấp hai của $f$.

\textbf{Bước 5 — Tính thời điểm đến, giới hạn và ngưỡng.} Sóng ở $(x,t)$ chỉ nhận dữ liệu trong $[x-at,x+at]$. Giải bất đẳng thức giao hỗ trợ để tìm thời điểm đến và đi qua. Nếu $H$ hằng ngoài đoạn hỗ trợ, dùng các giá trị hằng này để tính giới hạn khi $t\to\infty$. Nếu hỏi áp suất cực đại, viết hàm từng khoảng thời gian và so sánh cực trị, bước nhảy cùng hai giá trị một phía.

\textbf{Bước 1 — Đọc đúng điều kiện biên.} Biên $u(0,t)=0$ dùng mở rộng lẻ; biên $u_x(0,t)=0$ dùng mở rộng chẵn. Với dữ liệu trên $x\ge0$, đặt $\widetilde f(-s)=-f(s)$ hoặc $f(s)$ tương ứng, làm tương tự cho $g$. Mở rộng lẻ là $C^2$ nếu $f(0)=f''(0)=0$, và mở rộng lẻ của $g$ là $C^1$ nếu $g(0)=0$. Mở rộng chẵn đòi $f'(0)=g'(0)=0$ cho mức trơn ấy.

\textbf{Bước 2 — Lập nguyên hàm sau khi phản xạ.} Đặt $H(s)=\int_0^s\widetilde g(r)dr$. Nếu $\widetilde g$ lẻ thì $H$ chẵn; nếu $\widetilde g$ chẵn thì $H$ lẻ. Chính nguyên hàm, chứ không phải vận tốc, xuất hiện trong


\[
u=\frac{\widetilde f(x-at)+\widetilde f(x+at)}2+\frac{H(x+at)-H(x-at)}{2a}.
\]


\textbf{Bước 3 — Tách miền trước và sau khi chạm biên.} Với $u(0,t)=h(t)$, ở $x\ge at$ dùng d'Alembert với dữ liệu gốc. Ở $0<x<at$:


\[
u=h(t-x/a)+\frac{f(x+at)-f(at-x)}2+\frac1{2a}\int_{at-x}^{at+x}g(s)ds.
\]


Với $u_x(0,t)=h(t)$, dùng nghiệm của các mở rộng chẵn rồi trừ $a\int_0^{t-x/a}h(s)ds$ trên miền $x<at$. Tính riêng đạo hàm theo $x$ của hạng này để kiểm tra biên.

\textbf{Bước 4 — Trên đoạn hữu hạn, phản xạ ở cả hai đầu.} Dirichlet ở hai đầu cho mở rộng lẻ tại $0,L$, chu kỳ không gian $2L$. Neumann ở hai đầu cho mở rộng chẵn tại $0,L$, cũng chu kỳ $2L$. Một đầu lẻ, một đầu chẵn cho $\widetilde f(s+2L)=-\widetilde f(s)$ và chu kỳ $4L$, tương tự $g$. Vì thế chu kỳ thời gian thường là $2L/a$ hoặc $4L/a$. Với Neumann hai đầu còn phải xét trung bình $g$: trung bình khác không sinh hạng tuyến tính theo $t$, nên nghiệm không tuần hoàn.

\textbf{Bước 5 — Kiểm tra góc và đường ghép.} Biên Dirichlet không thuần nhất cần $h(0)=f(0)$, $h'(0)=g(0)$, $h''(0)=a^2f''(0)$ để có nghiệm $C^2$ cho phương trình thuần nhất. So sánh hai vết trên $x=at$, gồm $u,u_x,u_t$ rồi đạo hàm cấp hai. Nếu không tương thích, nêu rõ mức trơn còn lại thay vì khẳng định nghiệm cổ điển trên toàn miền.
\subsection*{Lời giải áp dụng cho bài}
\textbf{Bước 1. Đọc dữ liệu và chọn phản xạ.} Đoạn $(0,1)$, Dirichlet tại $0$, Neumann tại $1$, $a=1$, $f=0$, vận tốc $g$; minh họa $g=1$.

\textbf{Bước 2. Tính từng nhánh và nối nguyên hàm.} Mở rộng lẻ tại $0$ và chẵn tại $1$ cho chu kỳ $4$, đổi dấu sau $2$. Với $g=1$, vận tốc mở rộng là $1$ trên $(0,2)$, $-1$ trên $(2,4)$. Nguyên hàm có nhánh $s$ rồi $4-s$; thêm hằng số không đổi nghiệm.
\[f_e(s)=0,\qquad H(s)=1-|\operatorname{mod}(s,4)-2|.\]
Ở đây $\operatorname{clip}(s,p,q)=\min(q,\max(p,s))$; $\operatorname{mod}(s,L)\in[0,L)$ kể cả với $s<0$.

\textbf{Bước 3. Xác định hai sóng chạy và ghép nghiệm.} $A=f_e/2-H/(2a)$, $B=f_e/2+H/(2a)$. Thay hai cận:
\[u(x,t)=\frac{f_e(x-1t)+f_e(x+1t)}2+\frac{H(x+1t)-H(x-1t)}{2}.\]

\textbf{Bước 4. Trả lời các yêu cầu và kiểm tra.} Với $g$ tổng quát, $u_x(0,t)=\widetilde g(t)=g(t)$ khi $0<t<1$. Do đó $\int_0^1g^2=\int_0^1|u_x(0,t)|^2dt\le\int_0^T|u_x(0,t)|^2dt$ với $T\ge1$. Với $g=1$, công thức dưới có chu kỳ $4$ và đổi dấu sau $2$; không có chu kỳ $2$ cho nghiệm khác không.

Tại $t=0$, hiệu nguyên hàm bằng $0$ và đạo hàm thời gian cho $H'(x)=g(x)$ trong miền dữ liệu; half-sum cho đúng li độ. Mỗi hạng là sóng của $x\pm at$, nên thỏa PDE ngoài các mặt gãy và theo nghĩa phân phối trên toàn miền. Đồ thị hai sóng ban đầu và các lát cắt thời gian được tính trực tiếp từ các nhánh trên.
\begin{center}\includegraphics[width=.97\linewidth,height=11cm,keepaspectratio]{_xu_ly/graphs/BT_3_2_3_21.png}\end{center}
\textit{Đồ thị minh họa; với nghiệm số, xem phương pháp và tham số trong mã Python kèm theo.}

\clearpage\noindent\texttt{\detokenize{BT_3_2_3_22}}\par
\section*{Lời giải bổ sung — Bài tập 22 · Mục 3.2.3}

\subsection*{Kiến thức và trình tự giải}
\textbf{Bước 1 — Liệt kê các khoảng của dữ liệu.} Ghi đủ điểm đầu, điểm cuối hỗ trợ, điểm gãy và bước nhảy của $f,g$. Nếu có biên, phản xạ chẵn hoặc lẻ trước khi tích phân; trên đoạn hữu hạn tiếp tục phản xạ đến một chu kỳ hoàn chỉnh.

\textbf{Bước 2 — Tính nguyên hàm liên tục.} Lấy $H'=\widetilde g$ trên từng khoảng và chọn hằng số để hai vết của $H$ bằng nhau. Viết $[s]_p^q=\min(q,\max(p,s))$ hoặc $(s)_+=\max(s,0)$ khi cần gộp nhiều nhánh. Một nguyên hàm liên tục tránh tạo ra bước nhảy giả trong nghiệm.

\textbf{Bước 3 — Xác định sóng phải và sóng trái.} Đặt $A(s)=\widetilde f(s)/2-H(s)/(2a)$, $B(s)=\widetilde f(s)/2+H(s)/(2a)$; khi đó $u=A(x-at)+B(x+at)$. Với từng $t$ được yêu cầu, dịch đồ thị $A$ sang phải $at$ và $B$ sang trái $at$, rồi cộng tung độ. Các đường $x-at=b$ và $x+at=b$ chia mặt phẳng thành các miền có công thức đơn giản.

\textbf{Bước 4 — Xét liên tục, đạo hàm và giao của các mặt sóng.} Nếu $f$ có bước nhảy, $u$ thường có bước nhảy; nếu $f$ liên tục và $g$ bị chặn thì $u$ liên tục. Tại thời điểm cố định,


\[
u_x=\frac{\widetilde f'(x-at)+\widetilde f'(x+at)}2+\frac{\widetilde g(x+at)-\widetilde g(x-at)}{2a}.
\]


Tính giới hạn trái/phải của biểu thức này tại từng vị trí ứng viên. Có thể có triệt tiêu khi hai đường gặp nhau, nên không chỉ liệt kê đường rồi khẳng định mọi đạo hàm đều mất. Với tính $C^1$ theo cả $x,t$, còn kiểm tra $u_t$. Với $C^2$, xét tiếp các góc của $g$ và đạo hàm cấp hai của $f$.

\textbf{Bước 5 — Tính thời điểm đến, giới hạn và ngưỡng.} Sóng ở $(x,t)$ chỉ nhận dữ liệu trong $[x-at,x+at]$. Giải bất đẳng thức giao hỗ trợ để tìm thời điểm đến và đi qua. Nếu $H$ hằng ngoài đoạn hỗ trợ, dùng các giá trị hằng này để tính giới hạn khi $t\to\infty$. Nếu hỏi áp suất cực đại, viết hàm từng khoảng thời gian và so sánh cực trị, bước nhảy cùng hai giá trị một phía.

\textbf{Bước 1 — Đọc đúng điều kiện biên.} Biên $u(0,t)=0$ dùng mở rộng lẻ; biên $u_x(0,t)=0$ dùng mở rộng chẵn. Với dữ liệu trên $x\ge0$, đặt $\widetilde f(-s)=-f(s)$ hoặc $f(s)$ tương ứng, làm tương tự cho $g$. Mở rộng lẻ là $C^2$ nếu $f(0)=f''(0)=0$, và mở rộng lẻ của $g$ là $C^1$ nếu $g(0)=0$. Mở rộng chẵn đòi $f'(0)=g'(0)=0$ cho mức trơn ấy.

\textbf{Bước 2 — Lập nguyên hàm sau khi phản xạ.} Đặt $H(s)=\int_0^s\widetilde g(r)dr$. Nếu $\widetilde g$ lẻ thì $H$ chẵn; nếu $\widetilde g$ chẵn thì $H$ lẻ. Chính nguyên hàm, chứ không phải vận tốc, xuất hiện trong


\[
u=\frac{\widetilde f(x-at)+\widetilde f(x+at)}2+\frac{H(x+at)-H(x-at)}{2a}.
\]


\textbf{Bước 3 — Tách miền trước và sau khi chạm biên.} Với $u(0,t)=h(t)$, ở $x\ge at$ dùng d'Alembert với dữ liệu gốc. Ở $0<x<at$:


\[
u=h(t-x/a)+\frac{f(x+at)-f(at-x)}2+\frac1{2a}\int_{at-x}^{at+x}g(s)ds.
\]


Với $u_x(0,t)=h(t)$, dùng nghiệm của các mở rộng chẵn rồi trừ $a\int_0^{t-x/a}h(s)ds$ trên miền $x<at$. Tính riêng đạo hàm theo $x$ của hạng này để kiểm tra biên.

\textbf{Bước 4 — Trên đoạn hữu hạn, phản xạ ở cả hai đầu.} Dirichlet ở hai đầu cho mở rộng lẻ tại $0,L$, chu kỳ không gian $2L$. Neumann ở hai đầu cho mở rộng chẵn tại $0,L$, cũng chu kỳ $2L$. Một đầu lẻ, một đầu chẵn cho $\widetilde f(s+2L)=-\widetilde f(s)$ và chu kỳ $4L$, tương tự $g$. Vì thế chu kỳ thời gian thường là $2L/a$ hoặc $4L/a$. Với Neumann hai đầu còn phải xét trung bình $g$: trung bình khác không sinh hạng tuyến tính theo $t$, nên nghiệm không tuần hoàn.

\textbf{Bước 5 — Kiểm tra góc và đường ghép.} Biên Dirichlet không thuần nhất cần $h(0)=f(0)$, $h'(0)=g(0)$, $h''(0)=a^2f''(0)$ để có nghiệm $C^2$ cho phương trình thuần nhất. So sánh hai vết trên $x=at$, gồm $u,u_x,u_t$ rồi đạo hàm cấp hai. Nếu không tương thích, nêu rõ mức trơn còn lại thay vì khẳng định nghiệm cổ điển trên toàn miền.
\subsection*{Lời giải áp dụng cho bài}
\textbf{Bước 1. Đọc dữ liệu và chọn phản xạ.} Đoạn $(0,1)$, Dirichlet hai đầu, $a=1$, $g=0$, $f=4x$ trên $[0,1/4]$, $2-4x$ trên $[1/4,1/2]$, $0$ trên $[1/2,1]$.

\textbf{Bước 2. Tính từng nhánh và nối nguyên hàm.} Mở rộng lẻ tại hai đầu và chu kỳ $2$. Vì vận tốc bằng không, hai sóng chạy đều bằng một nửa $f_e$; không cần tích phân. Công thức này mô tả cả phản xạ tại các thời gian lớn hơn.
\[f_e(s)=\widetilde f\text{ là thác triển lẻ chu kỳ }2\text{ của tam giác trong đề},\qquad H(s)=0.\]
Ở đây $\operatorname{clip}(s,p,q)=\min(q,\max(p,s))$; $\operatorname{mod}(s,L)\in[0,L)$ kể cả với $s<0$.

\textbf{Bước 3. Xác định hai sóng chạy và ghép nghiệm.} $A=f_e/2-H/(2a)$, $B=f_e/2+H/(2a)$. Thay hai cận:
\[u(x,t)=\frac{f_e(x-1t)+f_e(x+1t)}2+\frac{H(x+1t)-H(x-1t)}{2}.\]

\textbf{Bước 4. Trả lời các yêu cầu và kiểm tra.} Với $1/2<x<1$, trước $t=x-1/2$, cả hai cận đều ở vùng dữ liệu không hoặc chưa có sóng phản xạ, nên $u=0$. Ngay sau thời điểm ấy, cận trái vào phần tam giác dương, nên $u\ne0$. Đó là thời điểm đến (giá trị đúng tại thời điểm đến vẫn $0$). Suốt $0\le t\le1/4$, các điểm cân bằng là $[3/4,1]$ cùng đầu cố định $0$. Với tốc độ $c$, thời điểm đến là $(x-1/2)/c$; vùng cân bằng suốt khoảng thời gian này là $[1/2+c/4,1]$ nếu cận không vượt $1$, cộng hai đầu cố định. Nếu $c\ge2$ chỉ còn hai đầu. Nếu đồng thời đổi khoảng thời gian thành $[0,1/(4c)]$, vùng giữ nguyên là $[3/4,1]$ cùng $0$.

Tại $t=0$, hiệu nguyên hàm bằng $0$ và đạo hàm thời gian cho $H'(x)=g(x)$ trong miền dữ liệu; half-sum cho đúng li độ. Mỗi hạng là sóng của $x\pm at$, nên thỏa PDE ngoài các mặt gãy và theo nghĩa phân phối trên toàn miền. Đồ thị hai sóng ban đầu và các lát cắt thời gian được tính trực tiếp từ các nhánh trên.
\begin{center}\includegraphics[width=.97\linewidth,height=11cm,keepaspectratio]{_xu_ly/graphs/BT_3_2_3_22.png}\end{center}
\textit{Đồ thị minh họa; với nghiệm số, xem phương pháp và tham số trong mã Python kèm theo.}

\clearpage\noindent\texttt{\detokenize{VD_3_3_2_01}}\par
\section*{Lời giải bổ sung — Ví dụ 1 · Mục 3.3.2}

\subsection*{Kiến thức và trình tự giải}
\textbf{Bước 1 — Tách biến và chọn bài toán trị riêng theo biên.} Đặt $u=X(x)T(t)$ trong $u_{tt}=a^2u_{xx}$, suy $X''+\lambda X=0$, $T''+a^2\lambda T=0$. Tích phân từng phần cho $\lambda\int X^2=\int (X')^2$ khi hạng biên bằng không; do đó $\lambda\ge0$. Giải cả trường hợp $\lambda=0$ trước khi xét $\lambda=k^2>0$.

\textbf{Bước 2 — Tìm hệ hàm riêng.} Trên $(0,L)$: Dirichlet hai đầu dùng $\sin(n\pi x/L)$, $n\ge1$; Neumann hai đầu dùng $1$ và $\cos(n\pi x/L)$; Dirichlet–Neumann dùng $\sin((n+1/2)\pi x/L)$, $n\ge0$; Neumann–Dirichlet dùng cos của các tần số nửa nguyên ấy. Chuẩn bình phương của mode khác không là $L/2$, của mode hằng là $L$.

\textbf{Bước 3 — Chiếu cả hai dữ liệu.} Với mode khác không đặt


\[
A_n=\frac2L\int_0^L fX_n\,dx,\quad B_n=\frac2L\int_0^L gX_n\,dx,\quad \omega_n=ak_n.
\]


Tích phân từng khoảng nếu dữ liệu chia đoạn; dùng tích phân từng phần hoặc công thức tích–tổng khi là đa thức/lượng giác. Hệ số $B_n$ là hệ số của vận tốc, nên hệ số đứng trước $\sin\omega_nt$ phải là $B_n/\omega_n$.

\textbf{Bước 4 — Ghép nghiệm và xét mode hằng.}


\[
u=\sum_n\left(A_n\cos\omega_nt+\frac{B_n}{\omega_n}\sin\omega_nt\right)X_n.
\]


Với Neumann hai đầu, thêm $A_0+B_0t$, trong đó $A_0=L^{-1}\int f$, $B_0=L^{-1}\int g$. Với tần số bị thay đổi bởi hạng $u$ hoặc cản $u_t$, giải lại ODE của từng mode; không dùng nguyên tần số $ak_n$.

\textbf{Bước 5 — Kiểm tra biên, dữ liệu và hội tụ.} Mỗi mode thỏa biên; tại $t=0$ hai chuỗi cho khai triển của $f,g$. Nếu dữ liệu không tương thích ở góc, giải theo nghĩa năng lượng hoặc từng miền và nêu đường mất tính trơn; không đạo hàm một chuỗi yếu tùy ý. Chu kỳ được suy từ tất cả mode có hệ số khác không, kể cả mode hằng; hai tần số có tỷ số vô tỷ ngăn một chu kỳ chung.

\textbf{Chi tiết tính tích phân hệ số.} Với $k=n\pi/L$, hai nguyên hàm thường dùng là


\[
\int x\sin kx\,dx=-x\cos kx/k+\sin kx/k^2,\qquad
\int x\cos kx\,dx=x\sin kx/k+\cos kx/k^2.
\]


Với đa thức $P$, ghi cả hai lần tích phân:


\[
\int_0^L P\sin kx\,dx
=\left[-P\cos kx/k\right]_0^L+
\left[P'\sin kx/k^2\right]_0^L-
\frac1{k^2}\int_0^L P''\sin kx\,dx.
\]


Chỉ bỏ hạng biên sau khi thế $0,L$ và điều kiện thực sự của $P$. Nếu dữ liệu chia đoạn, áp dụng công thức trên từng đoạn, rồi cộng các hạng tại điểm chia. Cuối cùng nhân nghịch đảo chuẩn của hàm riêng; nếu hệ số là của vận tốc, chia thêm tần số khi ghép nghiệm. Với mẫu số bằng $0$ trong tích–tổng, quay lại tích phân gốc để tính riêng.
\subsection*{Lời giải áp dụng cho bài}
\textbf{Bước 1. Chọn sine theo biên cố định.} Trên $(0,L)$, $X_n=\sin(n\pi x/L)$, $\omega_n=an\pi/L$. Hệ số của li độ là $A_n$, của vận tốc là $B_n$; nghiệm mode là $(A_n\cos\omega_nt+B_n\sin\omega_nt/\omega_n)X_n$.

\textbf{Bước 2. Câu (i).} $a=L=1$, $f=\sin\pi x$, $g=\sin^3\pi x=(3\sin\pi x-\sin3\pi x)/4$. Trực giao cho $A_1=1$, $B_1=3/4$, $B_3=-1/4$, các hệ số khác bằng $0$. Do đó
\[u=\sin\pi x\cos\pi t+\frac3{4\pi}\sin\pi x\sin\pi t-\frac1{12\pi}\sin3\pi x\sin3\pi t.\]
Đạo hàm thời gian tại $0$ khôi phục $g$; từng mode thỏa PDE và biên. Nghiệm $C^\infty$, chu kỳ cơ bản $2$.

\textbf{Bước 3. Câu (ii).} $a=2$, $L=\pi$, dữ liệu trong sách là $f(x)=\sin(\pi x)$, không phải $\sin x$. Dùng $2\sin\pi x\sin nx=\cos((\pi-n)x)-\cos((\pi+n)x)$:
\[A_n=\frac2\pi\int_0^\pi\sin(\pi x)\sin nx\,dx
=\frac{2(-1)^nn\sin(\pi^2)}{\pi(\pi^2-n^2)},\qquad B_n=0.\]
Không có mẫu số bằng không vì $n$ nguyên. Nghiệm là $u=\sum_{n\ge1}A_n\cos2nt\sin nx$, chu kỳ $\pi$.

\textbf{Bước 4. Nêu mức trơn đúng.} $f(\pi)=\sin(\pi^2)\ne0$, trái biên $u(\pi,t)=0$. Do đó câu (ii) không có nghiệm liên tục đến góc $(\pi,0)$ với mọi dữ liệu nguyên văn. Chuỗi xác định nghiệm theo nghĩa $L^2$/ từng miền; bước nhảy của mở rộng lẻ phát sinh mặt sóng từ đầu phải. Câu (i) không có vấn đề tương thích ấy.

\textbf{Chi tiết tích phân câu (ii).} Từ tích–tổng,
\[A_n=\frac1\pi\left[\frac{\sin((\pi-n)x)}{\pi-n}-\frac{\sin((\pi+n)x)}{\pi+n}\right]_0^\pi.\]
Vì $\sin(\pi^2-n\pi)=\sin(\pi^2+n\pi)=(-1)^n\sin\pi^2$,
\[A_n=\frac{(-1)^n\sin\pi^2}{\pi}\left(\frac1{\pi-n}-\frac1{\pi+n}\right)
=\frac{2(-1)^nn\sin\pi^2}{\pi(\pi^2-n^2)}.\]

\clearpage\noindent\texttt{\detokenize{VD_3_3_2_02}}\par
\section*{Lời giải bổ sung — Ví dụ 2 · Mục 3.3.2}

\subsection*{Kiến thức và trình tự giải}
\textbf{Bước 1 — Tách biến và chọn bài toán trị riêng theo biên.} Đặt $u=X(x)T(t)$ trong $u_{tt}=a^2u_{xx}$, suy $X''+\lambda X=0$, $T''+a^2\lambda T=0$. Tích phân từng phần cho $\lambda\int X^2=\int (X')^2$ khi hạng biên bằng không; do đó $\lambda\ge0$. Giải cả trường hợp $\lambda=0$ trước khi xét $\lambda=k^2>0$.

\textbf{Bước 2 — Tìm hệ hàm riêng.} Trên $(0,L)$: Dirichlet hai đầu dùng $\sin(n\pi x/L)$, $n\ge1$; Neumann hai đầu dùng $1$ và $\cos(n\pi x/L)$; Dirichlet–Neumann dùng $\sin((n+1/2)\pi x/L)$, $n\ge0$; Neumann–Dirichlet dùng cos của các tần số nửa nguyên ấy. Chuẩn bình phương của mode khác không là $L/2$, của mode hằng là $L$.

\textbf{Bước 3 — Chiếu cả hai dữ liệu.} Với mode khác không đặt


\[
A_n=\frac2L\int_0^L fX_n\,dx,\quad B_n=\frac2L\int_0^L gX_n\,dx,\quad \omega_n=ak_n.
\]


Tích phân từng khoảng nếu dữ liệu chia đoạn; dùng tích phân từng phần hoặc công thức tích–tổng khi là đa thức/lượng giác. Hệ số $B_n$ là hệ số của vận tốc, nên hệ số đứng trước $\sin\omega_nt$ phải là $B_n/\omega_n$.

\textbf{Bước 4 — Ghép nghiệm và xét mode hằng.}


\[
u=\sum_n\left(A_n\cos\omega_nt+\frac{B_n}{\omega_n}\sin\omega_nt\right)X_n.
\]


Với Neumann hai đầu, thêm $A_0+B_0t$, trong đó $A_0=L^{-1}\int f$, $B_0=L^{-1}\int g$. Với tần số bị thay đổi bởi hạng $u$ hoặc cản $u_t$, giải lại ODE của từng mode; không dùng nguyên tần số $ak_n$.

\textbf{Bước 5 — Kiểm tra biên, dữ liệu và hội tụ.} Mỗi mode thỏa biên; tại $t=0$ hai chuỗi cho khai triển của $f,g$. Nếu dữ liệu không tương thích ở góc, giải theo nghĩa năng lượng hoặc từng miền và nêu đường mất tính trơn; không đạo hàm một chuỗi yếu tùy ý. Chu kỳ được suy từ tất cả mode có hệ số khác không, kể cả mode hằng; hai tần số có tỷ số vô tỷ ngăn một chu kỳ chung.

\textbf{Chi tiết tính tích phân hệ số.} Với $k=n\pi/L$, hai nguyên hàm thường dùng là


\[
\int x\sin kx\,dx=-x\cos kx/k+\sin kx/k^2,\qquad
\int x\cos kx\,dx=x\sin kx/k+\cos kx/k^2.
\]


Với đa thức $P$, ghi cả hai lần tích phân:


\[
\int_0^L P\sin kx\,dx
=\left[-P\cos kx/k\right]_0^L+
\left[P'\sin kx/k^2\right]_0^L-
\frac1{k^2}\int_0^L P''\sin kx\,dx.
\]


Chỉ bỏ hạng biên sau khi thế $0,L$ và điều kiện thực sự của $P$. Nếu dữ liệu chia đoạn, áp dụng công thức trên từng đoạn, rồi cộng các hạng tại điểm chia. Cuối cùng nhân nghịch đảo chuẩn của hàm riêng; nếu hệ số là của vận tốc, chia thêm tần số khi ghép nghiệm. Với mẫu số bằng $0$ trong tích–tổng, quay lại tích phân gốc để tính riêng.
\subsection*{Lời giải áp dụng cho bài}
\textbf{Bước 1. Giữ mode hằng của Neumann.} Hệ hàm riêng là $X_0=1$, $X_n=\cos(n\pi x/L)$, $n\ge1$. Mode $0$ bằng $A_0+B_0t$ với hệ số là trung bình của dữ liệu.

\textbf{Bước 2. Câu (i).} $a=L=1$, $f=0$, $g=\cos^2\pi x=(1+\cos2\pi x)/2$. Vì vậy $B_0=1/2$, $B_2=1/2$, còn lại bằng $0$:
\[u=t/2+\cos2\pi x\sin2\pi t/(4\pi).\]
Đạo hàm không gian là bội của $\sin2\pi x$, bằng không ở hai đầu. Đạo hàm thời gian tại $0$ cho đúng $g$. Hạng tuyến tính khiến nghiệm không tuần hoàn dù phần dao động có chu kỳ $1$.

\textbf{Bước 3. Câu (ii).} $a=2$, $L=\pi$, $f=\sin x$, $g=0$. $A_0=2/\pi$. Tích–tổng cho
\[A_n=\frac2\pi\int_0^\pi\sin x\cos nx\,dx
=\frac{2(1+(-1)^n)}{\pi(1-n^2)}\quad(n\ne1),\qquad A_1=0.\]
Tại $n=1$ phải tích phân riêng, không dùng mẫu bằng $0$. Chỉ còn các mode chẵn:
\[u=\frac2\pi+\sum_{m\ge1}\frac4{\pi(1-4m^2)}\cos4mt\cos2mx.\]
\textbf{Bước 4. Chu kỳ và trơn.} Mode $m=1$ khác không cho chu kỳ cơ bản $\pi/2$. Dữ liệu $f'(0)=1$, $f'(\pi)=-1$ không thỏa biên Neumann; mở rộng chẵn có góc, nên không có nghiệm $C^1$ trên toàn miền qua các mặt phản xạ. Chuỗi và công thức phản xạ xác định cùng nghiệm liên tục từng miền.

\textbf{Chi tiết tích–tổng và ngoại lệ $n=1$.} Với $n\ne1$,
\[\begin{aligned}
\int_0^\pi\sin x\cos nx\,dx
&=\frac12\int_0^\pi[\sin((1+n)x)+\sin((1-n)x)]dx\\
&=\frac{1+(-1)^n}{2}\left(\frac1{1+n}+\frac1{1-n}\right)
=\frac{1+(-1)^n}{1-n^2}.
\end{aligned}\]
Nhân $2/\pi$ được $A_n$. Với $n=1$, tích phân riêng $\int_0^\pi\sin x\cos x\,dx=[\sin^2x/2]_0^\pi=0$. Mode hằng dùng chuẩn $\pi$, nên $A_0=\pi^{-1}[-\cos x]_0^\pi=2/\pi$, không dùng chuẩn $\pi/2$ của các mode khác.

\clearpage\noindent\texttt{\detokenize{VD_3_3_2_03}}\par
\section*{Lời giải bổ sung — Ví dụ 3 · Mục 3.3.2}

\subsection*{Kiến thức và trình tự giải}
\textbf{Bước 1 — Tách biến và chọn bài toán trị riêng theo biên.} Đặt $u=X(x)T(t)$ trong $u_{tt}=a^2u_{xx}$, suy $X''+\lambda X=0$, $T''+a^2\lambda T=0$. Tích phân từng phần cho $\lambda\int X^2=\int (X')^2$ khi hạng biên bằng không; do đó $\lambda\ge0$. Giải cả trường hợp $\lambda=0$ trước khi xét $\lambda=k^2>0$.

\textbf{Bước 2 — Tìm hệ hàm riêng.} Trên $(0,L)$: Dirichlet hai đầu dùng $\sin(n\pi x/L)$, $n\ge1$; Neumann hai đầu dùng $1$ và $\cos(n\pi x/L)$; Dirichlet–Neumann dùng $\sin((n+1/2)\pi x/L)$, $n\ge0$; Neumann–Dirichlet dùng cos của các tần số nửa nguyên ấy. Chuẩn bình phương của mode khác không là $L/2$, của mode hằng là $L$.

\textbf{Bước 3 — Chiếu cả hai dữ liệu.} Với mode khác không đặt


\[
A_n=\frac2L\int_0^L fX_n\,dx,\quad B_n=\frac2L\int_0^L gX_n\,dx,\quad \omega_n=ak_n.
\]


Tích phân từng khoảng nếu dữ liệu chia đoạn; dùng tích phân từng phần hoặc công thức tích–tổng khi là đa thức/lượng giác. Hệ số $B_n$ là hệ số của vận tốc, nên hệ số đứng trước $\sin\omega_nt$ phải là $B_n/\omega_n$.

\textbf{Bước 4 — Ghép nghiệm và xét mode hằng.}


\[
u=\sum_n\left(A_n\cos\omega_nt+\frac{B_n}{\omega_n}\sin\omega_nt\right)X_n.
\]


Với Neumann hai đầu, thêm $A_0+B_0t$, trong đó $A_0=L^{-1}\int f$, $B_0=L^{-1}\int g$. Với tần số bị thay đổi bởi hạng $u$ hoặc cản $u_t$, giải lại ODE của từng mode; không dùng nguyên tần số $ak_n$.

\textbf{Bước 5 — Kiểm tra biên, dữ liệu và hội tụ.} Mỗi mode thỏa biên; tại $t=0$ hai chuỗi cho khai triển của $f,g$. Nếu dữ liệu không tương thích ở góc, giải theo nghĩa năng lượng hoặc từng miền và nêu đường mất tính trơn; không đạo hàm một chuỗi yếu tùy ý. Chu kỳ được suy từ tất cả mode có hệ số khác không, kể cả mode hằng; hai tần số có tỷ số vô tỷ ngăn một chu kỳ chung.

\textbf{Chi tiết tính tích phân hệ số.} Với $k=n\pi/L$, hai nguyên hàm thường dùng là


\[
\int x\sin kx\,dx=-x\cos kx/k+\sin kx/k^2,\qquad
\int x\cos kx\,dx=x\sin kx/k+\cos kx/k^2.
\]


Với đa thức $P$, ghi cả hai lần tích phân:


\[
\int_0^L P\sin kx\,dx
=\left[-P\cos kx/k\right]_0^L+
\left[P'\sin kx/k^2\right]_0^L-
\frac1{k^2}\int_0^L P''\sin kx\,dx.
\]


Chỉ bỏ hạng biên sau khi thế $0,L$ và điều kiện thực sự của $P$. Nếu dữ liệu chia đoạn, áp dụng công thức trên từng đoạn, rồi cộng các hạng tại điểm chia. Cuối cùng nhân nghịch đảo chuẩn của hàm riêng; nếu hệ số là của vận tốc, chia thêm tần số khi ghép nghiệm. Với mẫu số bằng $0$ trong tích–tổng, quay lại tích phân gốc để tính riêng.
\subsection*{Lời giải áp dụng cho bài}
\textbf{Bước 1. Trị riêng Dirichlet–Neumann.} $X(0)=0$, $X'(L)=0$. Tích phân từng phần cho $\lambda\int X^2=\int(X')^2>0$ vì nghiệm hằng thỏa $X(0)=0$ chỉ là $0$. Với $\lambda=k^2$, $X=C\sin kx$, $\cos kL=0$, nên $k_n=(n+1/2)\pi/L$, $n\ge0$.

\textbf{Bước 2. Trực giao và hệ số.} Với hai trị riêng, trừ $X_m''X_n-X_n''X_m$ và tích phân. Hạng biên $[X_m'X_n-X_n'X_m]_0^L=0$, nên tích vô hướng bằng $0$ nếu khác trị riêng. Chuẩn là $L/2$, suy $A_n=(2/L)\int f\sin k_nx$, $B_n=(2/L)\int g\sin k_nx$.

\textbf{Bước 3. Ghép nghiệm.} $u=\sum_{n\ge0}(A_n\cos ak_nt+B_n\sin ak_nt/(ak_n))\sin k_nx$. Mỗi mode đổi dấu sau $2L/a$ và có chu kỳ $4L/a$. Tính duy nhất theo năng lượng: thông lượng ở $0$ hết vì $u_t=0$, ở $L$ hết vì $u_x=0$.

\textbf{Bước 4. Dữ liệu cụ thể.} $a=1$, $L=\pi$, $f=x$, $g=\sin^3(x/2)$. Với $k_n=n+1/2$, tích phân từng phần cho $\int_0^\pi x\sin k_nx=(-1)^n/k_n^2$, nên $A_n=2(-1)^n/(\pi k_n^2)$. Từ công thức sin lập phương, $B_0=3/4$, $B_1=-1/4$:
\[u=\sum_{n\ge0}\frac{2(-1)^n}{\pi k_n^2}\cos k_nt\sin k_nx
+\frac32\sin(t/2)\sin(x/2)-\frac16\sin(3t/2)\sin(3x/2).\]
Chu kỳ cơ bản $4\pi$. $f'(\pi)=1$ trái biên Neumann, nên nghiệm không $C^1$ qua mặt sóng từ góc phải; đây là nghiệm năng lượng/từng miền, không nghiệm cổ điển toàn bao đóng.

\textbf{Chi tiết điều kiện biên, chuẩn và tích phân dữ liệu.} Từ $X''+k^2X=0$, $X=A\cos kx+B\sin kx$. $X(0)=0$ cho $A=0$; $X'(L)=Bk\cos kL=0$ với $Bk\ne0$ cho $\cos kL=0$. Chuẩn
\[\int_0^L\sin^2k_nx\,dx=\left[x/2-\sin(2k_nx)/(4k_n)\right]_0^L=L/2.\]
Với $L=\pi$,
\[\int_0^\pi x\sin k_nx\,dx=
\left[-x\cos(k_nx)/k_n+\sin(k_nx)/k_n^2\right]_0^\pi
=(-1)^n/k_n^2.\]
Đã dùng $\cos(k_n\pi)=0$, $\sin(k_n\pi)=(-1)^n$. Khai triển $g=3\sin(x/2)/4-\sin(3x/2)/4$ rồi chiếu lên hệ sine trực giao cho $B_0=3/4,B_1=-1/4$.

\textbf{Chi tiết trực giao và dấu trị riêng.} Nhân $-X''=\lambda X$ bởi $X$, tích phân:
\[\lambda\int_0^LX^2dx=-[X'X]_0^L+\int_0^L(X')^2dx=\int_0^L(X')^2dx.\]
Hạng biên hết vì mỗi đầu có $X=0$ hoặc $X'=0$. Nếu $\lambda=0$, $\int(X')^2=0$ cho $X$ hằng; khi có một đầu Dirichlet, hằng đó bằng $0$. Với hai trị riêng,
\[(\lambda_m-\lambda_n)\int_0^LX_mX_n\,dx
=\int_0^L(-X_m''X_n+X_n''X_m)dx
=[-X_m'X_n+X_n'X_m]_0^L=0.\]
Nếu $\lambda_m\ne\lambda_n$, chia hiệu trị riêng cho tích vô hướng bằng $0$.

\clearpage\noindent\texttt{\detokenize{BT_3_3_3_01}}\par
\section*{Lời giải bổ sung — Bài tập 1 · Mục 3.3.3}

\subsection*{Kiến thức và trình tự giải}
\textbf{Bước 1 — Tách biến và chọn bài toán trị riêng theo biên.} Đặt $u=X(x)T(t)$ trong $u_{tt}=a^2u_{xx}$, suy $X''+\lambda X=0$, $T''+a^2\lambda T=0$. Tích phân từng phần cho $\lambda\int X^2=\int (X')^2$ khi hạng biên bằng không; do đó $\lambda\ge0$. Giải cả trường hợp $\lambda=0$ trước khi xét $\lambda=k^2>0$.

\textbf{Bước 2 — Tìm hệ hàm riêng.} Trên $(0,L)$: Dirichlet hai đầu dùng $\sin(n\pi x/L)$, $n\ge1$; Neumann hai đầu dùng $1$ và $\cos(n\pi x/L)$; Dirichlet–Neumann dùng $\sin((n+1/2)\pi x/L)$, $n\ge0$; Neumann–Dirichlet dùng cos của các tần số nửa nguyên ấy. Chuẩn bình phương của mode khác không là $L/2$, của mode hằng là $L$.

\textbf{Bước 3 — Chiếu cả hai dữ liệu.} Với mode khác không đặt


\[
A_n=\frac2L\int_0^L fX_n\,dx,\quad B_n=\frac2L\int_0^L gX_n\,dx,\quad \omega_n=ak_n.
\]


Tích phân từng khoảng nếu dữ liệu chia đoạn; dùng tích phân từng phần hoặc công thức tích–tổng khi là đa thức/lượng giác. Hệ số $B_n$ là hệ số của vận tốc, nên hệ số đứng trước $\sin\omega_nt$ phải là $B_n/\omega_n$.

\textbf{Bước 4 — Ghép nghiệm và xét mode hằng.}


\[
u=\sum_n\left(A_n\cos\omega_nt+\frac{B_n}{\omega_n}\sin\omega_nt\right)X_n.
\]


Với Neumann hai đầu, thêm $A_0+B_0t$, trong đó $A_0=L^{-1}\int f$, $B_0=L^{-1}\int g$. Với tần số bị thay đổi bởi hạng $u$ hoặc cản $u_t$, giải lại ODE của từng mode; không dùng nguyên tần số $ak_n$.

\textbf{Bước 5 — Kiểm tra biên, dữ liệu và hội tụ.} Mỗi mode thỏa biên; tại $t=0$ hai chuỗi cho khai triển của $f,g$. Nếu dữ liệu không tương thích ở góc, giải theo nghĩa năng lượng hoặc từng miền và nêu đường mất tính trơn; không đạo hàm một chuỗi yếu tùy ý. Chu kỳ được suy từ tất cả mode có hệ số khác không, kể cả mode hằng; hai tần số có tỷ số vô tỷ ngăn một chu kỳ chung.

\textbf{Chi tiết tính tích phân hệ số.} Với $k=n\pi/L$, hai nguyên hàm thường dùng là


\[
\int x\sin kx\,dx=-x\cos kx/k+\sin kx/k^2,\qquad
\int x\cos kx\,dx=x\sin kx/k+\cos kx/k^2.
\]


Với đa thức $P$, ghi cả hai lần tích phân:


\[
\int_0^L P\sin kx\,dx
=\left[-P\cos kx/k\right]_0^L+
\left[P'\sin kx/k^2\right]_0^L-
\frac1{k^2}\int_0^L P''\sin kx\,dx.
\]


Chỉ bỏ hạng biên sau khi thế $0,L$ và điều kiện thực sự của $P$. Nếu dữ liệu chia đoạn, áp dụng công thức trên từng đoạn, rồi cộng các hạng tại điểm chia. Cuối cùng nhân nghịch đảo chuẩn của hàm riêng; nếu hệ số là của vận tốc, chia thêm tần số khi ghép nghiệm. Với mẫu số bằng $0$ trong tích–tổng, quay lại tích phân gốc để tính riêng.
\subsection*{Lời giải áp dụng cho bài}
\textbf{Bước 1. Công thức chung.} Hai biên cố định cho $X_n=\sin(n\pi x/L)$, $\omega_n=an\pi/L$ và hệ số $A_n=(2/L)\int fX_n$, $B_n=(2/L)\int gX_n$. Dùng $u=\sum(A_n\cos\omega_nt+B_n\sin\omega_nt/\omega_n)X_n$.

\textbf{Bước 2. Tính hệ số từng dữ liệu.}

(i) $a=L=1$. $A_1=1$, còn lại $0$. Hai lần tích phân từng phần cho
\[\int_0^1x(1-x)\sin(n\pi x)dx=\frac{2(1-(-1)^n)}{n^3\pi^3}.\]
Do đó $B_n=4(1-(-1)^n)/(n^3\pi^3)$, và
\[u=\sin\pi x\cos\pi t+\sum_{n\text{ lẻ}}\frac8{n^4\pi^4}\sin(n\pi t)\sin(n\pi x).\]
Chu kỳ $2$; dữ liệu tương thích cho $C^2$.

(ii) $a=2$, $L=\pi$, $f=\sin x$, $g=1$. $A_1=1$; $B_n=(2/\pi)\int_0^\pi\sin nx=2(1-(-1)^n)/(\pi n)$. Suy $u=\sin x\cos2t+\sum_{n\text{ lẻ}}2\sin2nt\sin nx/(\pi n^2)$, chu kỳ $\pi$. Vận tốc không tương thích ở hai đầu, nên có gãy cấp một.

(iii) $a=1$, $L=\pi$, $f=1$, $g=0$. Cùng tích phân trên cho $A_n=2(1-(-1)^n)/(\pi n)$, nên $u=\sum_{n\text{ lẻ}}4\cos nt\sin nx/(\pi n)$. Chu kỳ $2\pi$; li độ không tương thích biên tạo bước nhảy.

(iv) $a=L=1$, $f=x(1-x)/4$, $g=0$. Chia hệ số của đa thức ở (i) cho $4$: $A_n=(1-(-1)^n)/(n^3\pi^3)$. Vậy $u=\sum_{n\text{ lẻ}}2\cos(n\pi t)\sin(n\pi x)/(n^3\pi^3)$, chu kỳ $2$. $f''(0)=f''(1)=-1/2$ không khớp biên cấp hai, nên không $C^2$ qua mặt phản xạ, dù nghiệm $C^1$.

(v) Cùng $f$, $g=-1$: giữ $A_n$ trên, thêm $B_n=-2(1-(-1)^n)/(n\pi)$. Nghiệm theo công thức chung, chu kỳ $2$, không $C^1$ qua mặt sóng từ góc vì vận tốc biên khác không.

(vi) $f=0$, $g=-1/2$ trên $(0,1/2)$, $1/2$ trên $(1/2,1)$. Tách tích phân tại $1/2$:
\[B_n=-\int_0^{1/2}\sin(n\pi x)dx+\int_{1/2}^1\sin(n\pi x)dx
=\frac{2\cos(n\pi/2)-1-(-1)^n}{n\pi}.\]
Chỉ $n\equiv2\pmod4$ khác không, với $B_n=-4/(n\pi)$. Nghiệm $u=-\sum_{n\equiv2\ (4)}4\sin(n\pi t)\sin(n\pi x)/(n^2\pi^2)$ có chu kỳ cơ bản $1$. Vận tốc bước nhảy cho nghiệm liên tục nhưng gãy cấp một.

\textbf{Bước 3. Kiểm tra.} Mỗi mode thỏa PDE và biên, hai chuỗi tại $t=0$ đúng các hệ số đã tính. Chỉ đạo hàm chuỗi ở nơi hội tụ đủ; với dữ liệu không tương thích dùng phản xạ để xác định các mặt gãy, không khẳng định cổ điển toàn miền.

\textbf{Chi tiết hai lần tích phân từng phần ở (i).} Đặt $k=n\pi$, $P=x(1-x)$, $P'=1-2x$, $P''=-2$.
\[\begin{aligned}
I&=\int_0^1P\sin kx\,dx
=\left[-P\frac{\cos kx}{k}\right]_0^1+\frac1k\int_0^1P'\cos kx\,dx\\
&=\frac1k\left[\frac{P'\sin kx}{k}\right]_0^1-\frac1{k^2}\int_0^1P''\sin kx\,dx\\
&=\frac2{k^2}\left[-\frac{\cos kx}{k}\right]_0^1=\frac{2(1-(-1)^n)}{k^3}.
\end{aligned}\]
Hạng biên lần đầu hết vì $P(0)=P(1)=0$; lần sau hết vì $\sin0=\sin n\pi=0$. Hệ số vận tốc là $B_n=2I$; khi ghép nghiệm còn chia $\omega_n=k$. Với $n$ lẻ, $2I/k=8/k^4$, không bỏ hệ số chia tần số này.

Ở (ii),(iii),(v), dùng $\int_0^L\sin(n\pi x/L)dx=L(1-(-1)^n)/(n\pi)$, rồi nhân $2/L$ và giá trị hằng của dữ liệu. Ở (iv), $\int fX_n=I/4$ nên $A_n=I/2=(1-(-1)^n)/k^3$. Ở (vi),
\[-\int_0^{1/2}\sin kx\,dx=\frac{\cos(k/2)-1}{k},\quad
\int_{1/2}^1\sin kx\,dx=\frac{\cos(k/2)-\cos k}{k}.\]
Cộng rồi xét $n$ lẻ, $n\equiv0\pmod4$, $n\equiv2\pmod4$ cho lần lượt $0,0,-4/k$.

\clearpage\noindent\texttt{\detokenize{BT_3_3_3_02}}\par
\section*{Lời giải bổ sung — Bài tập 2 · Mục 3.3.3}

\subsection*{Kiến thức và trình tự giải}
\textbf{Bước 1 — Tách biến và chọn bài toán trị riêng theo biên.} Đặt $u=X(x)T(t)$ trong $u_{tt}=a^2u_{xx}$, suy $X''+\lambda X=0$, $T''+a^2\lambda T=0$. Tích phân từng phần cho $\lambda\int X^2=\int (X')^2$ khi hạng biên bằng không; do đó $\lambda\ge0$. Giải cả trường hợp $\lambda=0$ trước khi xét $\lambda=k^2>0$.

\textbf{Bước 2 — Tìm hệ hàm riêng.} Trên $(0,L)$: Dirichlet hai đầu dùng $\sin(n\pi x/L)$, $n\ge1$; Neumann hai đầu dùng $1$ và $\cos(n\pi x/L)$; Dirichlet–Neumann dùng $\sin((n+1/2)\pi x/L)$, $n\ge0$; Neumann–Dirichlet dùng cos của các tần số nửa nguyên ấy. Chuẩn bình phương của mode khác không là $L/2$, của mode hằng là $L$.

\textbf{Bước 3 — Chiếu cả hai dữ liệu.} Với mode khác không đặt


\[
A_n=\frac2L\int_0^L fX_n\,dx,\quad B_n=\frac2L\int_0^L gX_n\,dx,\quad \omega_n=ak_n.
\]


Tích phân từng khoảng nếu dữ liệu chia đoạn; dùng tích phân từng phần hoặc công thức tích–tổng khi là đa thức/lượng giác. Hệ số $B_n$ là hệ số của vận tốc, nên hệ số đứng trước $\sin\omega_nt$ phải là $B_n/\omega_n$.

\textbf{Bước 4 — Ghép nghiệm và xét mode hằng.}


\[
u=\sum_n\left(A_n\cos\omega_nt+\frac{B_n}{\omega_n}\sin\omega_nt\right)X_n.
\]


Với Neumann hai đầu, thêm $A_0+B_0t$, trong đó $A_0=L^{-1}\int f$, $B_0=L^{-1}\int g$. Với tần số bị thay đổi bởi hạng $u$ hoặc cản $u_t$, giải lại ODE của từng mode; không dùng nguyên tần số $ak_n$.

\textbf{Bước 5 — Kiểm tra biên, dữ liệu và hội tụ.} Mỗi mode thỏa biên; tại $t=0$ hai chuỗi cho khai triển của $f,g$. Nếu dữ liệu không tương thích ở góc, giải theo nghĩa năng lượng hoặc từng miền và nêu đường mất tính trơn; không đạo hàm một chuỗi yếu tùy ý. Chu kỳ được suy từ tất cả mode có hệ số khác không, kể cả mode hằng; hai tần số có tỷ số vô tỷ ngăn một chu kỳ chung.

\textbf{Chi tiết tính tích phân hệ số.} Với $k=n\pi/L$, hai nguyên hàm thường dùng là


\[
\int x\sin kx\,dx=-x\cos kx/k+\sin kx/k^2,\qquad
\int x\cos kx\,dx=x\sin kx/k+\cos kx/k^2.
\]


Với đa thức $P$, ghi cả hai lần tích phân:


\[
\int_0^L P\sin kx\,dx
=\left[-P\cos kx/k\right]_0^L+
\left[P'\sin kx/k^2\right]_0^L-
\frac1{k^2}\int_0^L P''\sin kx\,dx.
\]


Chỉ bỏ hạng biên sau khi thế $0,L$ và điều kiện thực sự của $P$. Nếu dữ liệu chia đoạn, áp dụng công thức trên từng đoạn, rồi cộng các hạng tại điểm chia. Cuối cùng nhân nghịch đảo chuẩn của hàm riêng; nếu hệ số là của vận tốc, chia thêm tần số khi ghép nghiệm. Với mẫu số bằng $0$ trong tích–tổng, quay lại tích phân gốc để tính riêng.
\subsection*{Lời giải áp dụng cho bài}
\textbf{Bước 1. Chọn hệ sine.} Trên $(0,\pi)$, $f=\min(x,\pi-x)$, $g=0$, $a=0.0002$. Tách tích phân tại $\pi/2$ và đổi biến $y=\pi-x$ ở nửa sau:
\[A_n=\frac2\pi(1+(-1)^{n+1})\int_0^{\pi/2}x\sin nx\,dx.\]
Với $n$ chẵn bằng $0$. Với $n$ lẻ, $\cos(n\pi/2)=0$, tích phân bằng $\sin(n\pi/2)/n^2$, nên chung $A_n=4\sin(n\pi/2)/(\pi n^2)$.

\textbf{Bước 2. Giải ODE.} $T_n''+a^2n^2T_n=0$, $T_n(0)=A_n$, $T_n'(0)=0$, cho
\[u=\sum_{n\ge1}\frac{4\sin(n\pi/2)}{\pi n^2}\cos(0.0002nt)\sin nx.\]
\textbf{Bước 3. Chu kỳ và tính trơn.} Mode $1$ khác không, mọi tần số là bội nguyên của $a$, nên chu kỳ cơ bản $2\pi/a=10000\pi$. Li độ có góc tại $\pi/2$; hai góc chạy theo $x=\pi/2\pm at$ rồi phản xạ ở hai đầu. Nghiệm liên tục/ năng lượng, không $C^1$ trên các nhánh truyền góc. Thế $t=0$ thu lại khai triển tam giác; biên đúng theo từng mode.

\textbf{Chi tiết hai tích phân của tam giác.} Đặt $k=n$, $b=\pi/2$. Đổi $y=\pi-x$ ở nửa sau:
\[\int_b^\pi(\pi-x)\sin nx\,dx=(-1)^{n+1}\int_0^b y\sin ny\,dy.\]
Tích phân từng phần còn lại là
\[\int_0^b x\sin nx\,dx=\left[-x\cos nx/n+\sin nx/n^2\right]_0^b
=-b\cos(nb)/n+\sin(nb)/n^2.\]
Với $n$ chẵn, hai nửa triệt tiêu; với $n$ lẻ, $\cos(nb)=0$ và hệ số trước tích phân là $4/\pi$. Suy $A_n=4\sin(n\pi/2)/(\pi n^2)$ cho mọi $n$.

\clearpage\noindent\texttt{\detokenize{BT_3_3_3_03}}\par
\section*{Lời giải bổ sung — Bài tập 3 · Mục 3.3.3}

\subsection*{Kiến thức và trình tự giải}
\textbf{Bước 1 — Tách biến và chọn bài toán trị riêng theo biên.} Đặt $u=X(x)T(t)$ trong $u_{tt}=a^2u_{xx}$, suy $X''+\lambda X=0$, $T''+a^2\lambda T=0$. Tích phân từng phần cho $\lambda\int X^2=\int (X')^2$ khi hạng biên bằng không; do đó $\lambda\ge0$. Giải cả trường hợp $\lambda=0$ trước khi xét $\lambda=k^2>0$.

\textbf{Bước 2 — Tìm hệ hàm riêng.} Trên $(0,L)$: Dirichlet hai đầu dùng $\sin(n\pi x/L)$, $n\ge1$; Neumann hai đầu dùng $1$ và $\cos(n\pi x/L)$; Dirichlet–Neumann dùng $\sin((n+1/2)\pi x/L)$, $n\ge0$; Neumann–Dirichlet dùng cos của các tần số nửa nguyên ấy. Chuẩn bình phương của mode khác không là $L/2$, của mode hằng là $L$.

\textbf{Bước 3 — Chiếu cả hai dữ liệu.} Với mode khác không đặt


\[
A_n=\frac2L\int_0^L fX_n\,dx,\quad B_n=\frac2L\int_0^L gX_n\,dx,\quad \omega_n=ak_n.
\]


Tích phân từng khoảng nếu dữ liệu chia đoạn; dùng tích phân từng phần hoặc công thức tích–tổng khi là đa thức/lượng giác. Hệ số $B_n$ là hệ số của vận tốc, nên hệ số đứng trước $\sin\omega_nt$ phải là $B_n/\omega_n$.

\textbf{Bước 4 — Ghép nghiệm và xét mode hằng.}


\[
u=\sum_n\left(A_n\cos\omega_nt+\frac{B_n}{\omega_n}\sin\omega_nt\right)X_n.
\]


Với Neumann hai đầu, thêm $A_0+B_0t$, trong đó $A_0=L^{-1}\int f$, $B_0=L^{-1}\int g$. Với tần số bị thay đổi bởi hạng $u$ hoặc cản $u_t$, giải lại ODE của từng mode; không dùng nguyên tần số $ak_n$.

\textbf{Bước 5 — Kiểm tra biên, dữ liệu và hội tụ.} Mỗi mode thỏa biên; tại $t=0$ hai chuỗi cho khai triển của $f,g$. Nếu dữ liệu không tương thích ở góc, giải theo nghĩa năng lượng hoặc từng miền và nêu đường mất tính trơn; không đạo hàm một chuỗi yếu tùy ý. Chu kỳ được suy từ tất cả mode có hệ số khác không, kể cả mode hằng; hai tần số có tỷ số vô tỷ ngăn một chu kỳ chung.

\textbf{Chi tiết tính tích phân hệ số.} Với $k=n\pi/L$, hai nguyên hàm thường dùng là


\[
\int x\sin kx\,dx=-x\cos kx/k+\sin kx/k^2,\qquad
\int x\cos kx\,dx=x\sin kx/k+\cos kx/k^2.
\]


Với đa thức $P$, ghi cả hai lần tích phân:


\[
\int_0^L P\sin kx\,dx
=\left[-P\cos kx/k\right]_0^L+
\left[P'\sin kx/k^2\right]_0^L-
\frac1{k^2}\int_0^L P''\sin kx\,dx.
\]


Chỉ bỏ hạng biên sau khi thế $0,L$ và điều kiện thực sự của $P$. Nếu dữ liệu chia đoạn, áp dụng công thức trên từng đoạn, rồi cộng các hạng tại điểm chia. Cuối cùng nhân nghịch đảo chuẩn của hàm riêng; nếu hệ số là của vận tốc, chia thêm tần số khi ghép nghiệm. Với mẫu số bằng $0$ trong tích–tổng, quay lại tích phân gốc để tính riêng.
\subsection*{Lời giải áp dụng cho bài}
\textbf{Bước 1. Biên Dirichlet–Neumann.} Tách biến cho $k_n=(n+1/2)\pi/L$. Khi dịch $2L/a$, pha tăng $(2n+1)\pi$; cos và sin đều đổi dấu, nên $u(t+2L/a)=-u(t)$. Dịch thêm lần nữa cho chu kỳ $4L/a$.

\textbf{Bước 2. Kiểm tra khẳng định chu kỳ ngắn.} Nếu còn $u(t+2L/a)=u(t)$ thì $u=-u$, suy $u=0$. Lấy $f=\sin(\pi x/(2L))$, $g=0$ cho $u=f\cos(a\pi t/(2L))$ là phản ví dụ rõ ràng với chu kỳ cơ bản $4L/a$. Vì vậy chu kỳ $2L/a$ trong câu (a) không đúng cho nghiệm khác không.

\textbf{Bước 3. Biên Neumann hai đầu.} Tích phân PDE theo $x$ cho trung bình $\bar u''=(a^2/L)[u_x]_0^L=0$, nên $\bar u=\bar f+t\bar g$. Phần còn lại gồm các tần số $an\pi/L$, tất cả có chu kỳ $2L/a$. Nếu $\bar g=0$ thì nghiệm tuần hoàn; nếu $\bar g\ne0$, trung bình tăng tuyến tính và không thể tuần hoàn. Điều kiện cần và đủ là $\int_0^Lg=0$. Nghiệm hằng theo thời gian có mọi chu kỳ dương.

\clearpage\noindent\texttt{\detokenize{BT_3_3_3_04}}\par
\section*{Lời giải bổ sung — Bài tập 4 · Mục 3.3.3}

\subsection*{Kiến thức và trình tự giải}
\textbf{Bước 1 — Tách biến và chọn bài toán trị riêng theo biên.} Đặt $u=X(x)T(t)$ trong $u_{tt}=a^2u_{xx}$, suy $X''+\lambda X=0$, $T''+a^2\lambda T=0$. Tích phân từng phần cho $\lambda\int X^2=\int (X')^2$ khi hạng biên bằng không; do đó $\lambda\ge0$. Giải cả trường hợp $\lambda=0$ trước khi xét $\lambda=k^2>0$.

\textbf{Bước 2 — Tìm hệ hàm riêng.} Trên $(0,L)$: Dirichlet hai đầu dùng $\sin(n\pi x/L)$, $n\ge1$; Neumann hai đầu dùng $1$ và $\cos(n\pi x/L)$; Dirichlet–Neumann dùng $\sin((n+1/2)\pi x/L)$, $n\ge0$; Neumann–Dirichlet dùng cos của các tần số nửa nguyên ấy. Chuẩn bình phương của mode khác không là $L/2$, của mode hằng là $L$.

\textbf{Bước 3 — Chiếu cả hai dữ liệu.} Với mode khác không đặt


\[
A_n=\frac2L\int_0^L fX_n\,dx,\quad B_n=\frac2L\int_0^L gX_n\,dx,\quad \omega_n=ak_n.
\]


Tích phân từng khoảng nếu dữ liệu chia đoạn; dùng tích phân từng phần hoặc công thức tích–tổng khi là đa thức/lượng giác. Hệ số $B_n$ là hệ số của vận tốc, nên hệ số đứng trước $\sin\omega_nt$ phải là $B_n/\omega_n$.

\textbf{Bước 4 — Ghép nghiệm và xét mode hằng.}


\[
u=\sum_n\left(A_n\cos\omega_nt+\frac{B_n}{\omega_n}\sin\omega_nt\right)X_n.
\]


Với Neumann hai đầu, thêm $A_0+B_0t$, trong đó $A_0=L^{-1}\int f$, $B_0=L^{-1}\int g$. Với tần số bị thay đổi bởi hạng $u$ hoặc cản $u_t$, giải lại ODE của từng mode; không dùng nguyên tần số $ak_n$.

\textbf{Bước 5 — Kiểm tra biên, dữ liệu và hội tụ.} Mỗi mode thỏa biên; tại $t=0$ hai chuỗi cho khai triển của $f,g$. Nếu dữ liệu không tương thích ở góc, giải theo nghĩa năng lượng hoặc từng miền và nêu đường mất tính trơn; không đạo hàm một chuỗi yếu tùy ý. Chu kỳ được suy từ tất cả mode có hệ số khác không, kể cả mode hằng; hai tần số có tỷ số vô tỷ ngăn một chu kỳ chung.

\textbf{Chi tiết tính tích phân hệ số.} Với $k=n\pi/L$, hai nguyên hàm thường dùng là


\[
\int x\sin kx\,dx=-x\cos kx/k+\sin kx/k^2,\qquad
\int x\cos kx\,dx=x\sin kx/k+\cos kx/k^2.
\]


Với đa thức $P$, ghi cả hai lần tích phân:


\[
\int_0^L P\sin kx\,dx
=\left[-P\cos kx/k\right]_0^L+
\left[P'\sin kx/k^2\right]_0^L-
\frac1{k^2}\int_0^L P''\sin kx\,dx.
\]


Chỉ bỏ hạng biên sau khi thế $0,L$ và điều kiện thực sự của $P$. Nếu dữ liệu chia đoạn, áp dụng công thức trên từng đoạn, rồi cộng các hạng tại điểm chia. Cuối cùng nhân nghịch đảo chuẩn của hàm riêng; nếu hệ số là của vận tốc, chia thêm tần số khi ghép nghiệm. Với mẫu số bằng $0$ trong tích–tổng, quay lại tích phân gốc để tính riêng.
\subsection*{Lời giải áp dụng cho bài}
\textbf{Bước 1. Hệ cos và mode hằng.} $A_0,B_0$ là trung bình dữ liệu; mode khác không có chuẩn $L/2$. Không được dùng công thức sine của biên cố định.

\textbf{Bước 2. Câu (i).} $a=L=1$, $f=\cos2\pi x$, $g=x-1/2$. $A_2=1$, $B_0=0$. Tích phân từng phần cho
\[B_n=2\int_0^1(x-1/2)\cos(n\pi x)dx=\frac{2((-1)^n-1)}{n^2\pi^2}.\]
Do đó $u=\cos2\pi x\cos2\pi t-\sum_{n\text{ lẻ}}4\sin(n\pi t)\cos(n\pi x)/(n^3\pi^3)$. Mode $1$ khác không cho chu kỳ cơ bản $2$. $g'=1$ ở hai đầu trái đạo hàm thời gian của biên Neumann, nên không $C^2$ qua các mặt sóng góc.

\textbf{Bước 3. Câu (ii).} $a=2$, $L=\pi$, $f=\sin^2x=(1-\cos2x)/2$, $g=\sin x$. $A_0=1/2$, $A_2=-1/2$, $B_0=2/\pi$. Tích–tổng cho $B_n=2(1+(-1)^n)/(\pi(1-n^2))$ nếu $n\ne1$, và $B_1=0$ tính riêng. Suy
\[u=\tfrac12-\tfrac12\cos2x\cos4t+2t/\pi+\sum_{n\ge1}\frac{B_n}{2n}\sin2nt\cos nx.\]
\textbf{Bước 4. Kiểm tra.} Mode hằng đảm bảo vận tốc trung bình đúng; bỏ nó sẽ làm sai $g$. Nó cũng khiến câu (ii) không tuần hoàn. Biên đúng theo mỗi cos; tại $t=0$ hai chuỗi khôi phục dữ liệu. $g'(0),g'(\pi)$ khác không nên vẫn có vấn đề tương thích cấp hai.

\textbf{Chi tiết tích phân câu (i).} Đặt $k=n\pi$:
\[\begin{aligned}
B_n&=2\left[\frac{(x-1/2)\sin kx}{k}\right]_0^1-\frac2k\int_0^1\sin kx\,dx\\
&=-\frac2k\left[-\frac{\cos kx}{k}\right]_0^1
=\frac{2((-1)^n-1)}{k^2}.
\end{aligned}\]
Vì $\sin0=\sin k=0$, hạng biên đầu hết. $B_0=\int_0^1(x-1/2)dx=0$. Ở (ii), trung bình $B_0=\pi^{-1}\int_0^\pi\sin x\,dx=2/\pi$; các mode khác dùng $2\sin x\cos nx=\sin((1+n)x)+\sin((1-n)x)$. Tích phân cho $B_n=2(1+(-1)^n)/[\pi(1-n^2)]$ khi $n\ne1$, còn $B_1=\pi^{-1}\int_0^\pi\sin2x\,dx=0$.

\clearpage\noindent\texttt{\detokenize{BT_3_3_3_05}}\par
\section*{Lời giải bổ sung — Bài tập 5 · Mục 3.3.3}

\subsection*{Kiến thức và trình tự giải}
\textbf{Bước 1 — Tách biến và chọn bài toán trị riêng theo biên.} Đặt $u=X(x)T(t)$ trong $u_{tt}=a^2u_{xx}$, suy $X''+\lambda X=0$, $T''+a^2\lambda T=0$. Tích phân từng phần cho $\lambda\int X^2=\int (X')^2$ khi hạng biên bằng không; do đó $\lambda\ge0$. Giải cả trường hợp $\lambda=0$ trước khi xét $\lambda=k^2>0$.

\textbf{Bước 2 — Tìm hệ hàm riêng.} Trên $(0,L)$: Dirichlet hai đầu dùng $\sin(n\pi x/L)$, $n\ge1$; Neumann hai đầu dùng $1$ và $\cos(n\pi x/L)$; Dirichlet–Neumann dùng $\sin((n+1/2)\pi x/L)$, $n\ge0$; Neumann–Dirichlet dùng cos của các tần số nửa nguyên ấy. Chuẩn bình phương của mode khác không là $L/2$, của mode hằng là $L$.

\textbf{Bước 3 — Chiếu cả hai dữ liệu.} Với mode khác không đặt


\[
A_n=\frac2L\int_0^L fX_n\,dx,\quad B_n=\frac2L\int_0^L gX_n\,dx,\quad \omega_n=ak_n.
\]


Tích phân từng khoảng nếu dữ liệu chia đoạn; dùng tích phân từng phần hoặc công thức tích–tổng khi là đa thức/lượng giác. Hệ số $B_n$ là hệ số của vận tốc, nên hệ số đứng trước $\sin\omega_nt$ phải là $B_n/\omega_n$.

\textbf{Bước 4 — Ghép nghiệm và xét mode hằng.}


\[
u=\sum_n\left(A_n\cos\omega_nt+\frac{B_n}{\omega_n}\sin\omega_nt\right)X_n.
\]


Với Neumann hai đầu, thêm $A_0+B_0t$, trong đó $A_0=L^{-1}\int f$, $B_0=L^{-1}\int g$. Với tần số bị thay đổi bởi hạng $u$ hoặc cản $u_t$, giải lại ODE của từng mode; không dùng nguyên tần số $ak_n$.

\textbf{Bước 5 — Kiểm tra biên, dữ liệu và hội tụ.} Mỗi mode thỏa biên; tại $t=0$ hai chuỗi cho khai triển của $f,g$. Nếu dữ liệu không tương thích ở góc, giải theo nghĩa năng lượng hoặc từng miền và nêu đường mất tính trơn; không đạo hàm một chuỗi yếu tùy ý. Chu kỳ được suy từ tất cả mode có hệ số khác không, kể cả mode hằng; hai tần số có tỷ số vô tỷ ngăn một chu kỳ chung.

\textbf{Chi tiết tính tích phân hệ số.} Với $k=n\pi/L$, hai nguyên hàm thường dùng là


\[
\int x\sin kx\,dx=-x\cos kx/k+\sin kx/k^2,\qquad
\int x\cos kx\,dx=x\sin kx/k+\cos kx/k^2.
\]


Với đa thức $P$, ghi cả hai lần tích phân:


\[
\int_0^L P\sin kx\,dx
=\left[-P\cos kx/k\right]_0^L+
\left[P'\sin kx/k^2\right]_0^L-
\frac1{k^2}\int_0^L P''\sin kx\,dx.
\]


Chỉ bỏ hạng biên sau khi thế $0,L$ và điều kiện thực sự của $P$. Nếu dữ liệu chia đoạn, áp dụng công thức trên từng đoạn, rồi cộng các hạng tại điểm chia. Cuối cùng nhân nghịch đảo chuẩn của hàm riêng; nếu hệ số là của vận tốc, chia thêm tần số khi ghép nghiệm. Với mẫu số bằng $0$ trong tích–tổng, quay lại tích phân gốc để tính riêng.
\subsection*{Lời giải áp dụng cho bài}
\textbf{Bước 1. Trị riêng và trực giao.} $X'(0)=0$, $X(L)=0$. Tích phân từng phần cho $\lambda\int X^2=\int X'^2$. Nếu $\lambda=0$, $X$ hằng rồi $X(L)=0$ nên chỉ có nghiệm không; vì vậy $\lambda>0$. Với $\lambda=k^2$, điều kiện đầu cho $X=C\cos kx$, điều kiện cuối cho $\cos kL=0$. Suy $k_n=(n+1/2)\pi/L$. Trừ hai phương trình của $X_m,X_n$ và tích phân; các hạng biên bằng không, chứng minh trực giao.

\textbf{Bước 2. Chiếu dữ liệu và ghép.} $A_n=(2/L)\int f\cos k_nx$, $B_n=(2/L)\int g\cos k_nx$; $u=\sum(A_n\cos ak_nt+B_n\sin ak_nt/(ak_n))\cos k_nx$. Mọi mode có chu kỳ $4L/a$ và đổi dấu sau $2L/a$.

\textbf{Bước 3. Trường hợp cụ thể.} $a=1$, $L=\pi$, $f=\pi-x$, $g=\cos^3(x/2)$. Tích phân từng phần cho $\int_0^\pi(\pi-x)\cos k_nx=1/k_n^2$, nên $A_n=2/(\pi k_n^2)$. Công thức cos lập phương cho $B_0=3/4$, $B_1=1/4$. Vì vậy
\[u=\sum_{n\ge0}\frac2{\pi k_n^2}\cos k_nt\cos k_nx
+\tfrac32\sin(t/2)\cos(x/2)+\tfrac16\sin(3t/2)\cos(3x/2).\]
\textbf{Bước 4. Kết luận.} Mode đầu khác không nên chu kỳ cơ bản $4\pi$. $f'(0)=-1$ không tương thích Neumann, nên nghiệm không $C^1$ qua mặt phản xạ từ góc trái. Hai chuỗi vẫn đúng dữ liệu trong miền và từng mode thỏa PDE/biên; tính duy nhất trong lớp năng lượng theo năng lượng của hiệu.

\textbf{Chi tiết điều kiện biên và tích phân dữ liệu.} $X=A\cos kx+B\sin kx$, $X'(0)=Bk=0$ cho $B=0$ khi $k>0$; $X(L)=A\cos kL=0$ cho $\cos kL=0$. Chuẩn $\int_0^L\cos^2k_nx\,dx=L/2$ vì $\sin(2k_nL)=0$. Với $L=\pi$,
\[\begin{aligned}
\int_0^\pi(\pi-x)\cos k_nx\,dx
&=\left[\frac{(\pi-x)\sin k_nx}{k_n}\right]_0^\pi
+\frac1{k_n}\int_0^\pi\sin k_nx\,dx\\
&=\frac{1-\cos(k_n\pi)}{k_n^2}=1/k_n^2.
\end{aligned}\]
Chiếu $g=(3\cos(x/2)+\cos(3x/2))/4$ lên hệ cos trực giao cho $B_0=3/4,B_1=1/4$, rồi chia tần số $k_0=1/2,k_1=3/2$ khi ghép nghiệm.

\textbf{Chi tiết trực giao và dấu trị riêng.} Nhân $-X''=\lambda X$ bởi $X$, tích phân:
\[\lambda\int_0^LX^2dx=-[X'X]_0^L+\int_0^L(X')^2dx=\int_0^L(X')^2dx.\]
Hạng biên hết vì mỗi đầu có $X=0$ hoặc $X'=0$. Nếu $\lambda=0$, $\int(X')^2=0$ cho $X$ hằng; khi có một đầu Dirichlet, hằng đó bằng $0$. Với hai trị riêng,
\[(\lambda_m-\lambda_n)\int_0^LX_mX_n\,dx
=\int_0^L(-X_m''X_n+X_n''X_m)dx
=[-X_m'X_n+X_n'X_m]_0^L=0.\]
Nếu $\lambda_m\ne\lambda_n$, chia hiệu trị riêng cho tích vô hướng bằng $0$.

\clearpage\noindent\texttt{\detokenize{VD_3_4_2_01}}\par
\section*{Lời giải bổ sung — Ví dụ 1 · Mục 3.4.2}

\subsection*{Kiến thức và trình tự giải}
\textbf{Bước 1 — Đưa điều kiện biên về thuần nhất.} Với $u(0,t)=\mu(t)$, $u(L,t)=\nu(t)$, đặt $b=(1-x/L)\mu+(x/L)\nu$. Với Neumann hai đầu đặt $b=\mu x+(\nu-\mu)x^2/(2L)$. Với Dirichlet–Neumann đặt $b=\mu+x\nu$. Viết $u=b+w$; khi đó $w$ có biên thuần nhất, nguồn $Q=F+a^2b_{xx}-b_{tt}$ và dữ liệu $f-b(x,0)$, $g-b_t(x,0)$.

\textbf{Bước 2 — Chiếu nguồn và dữ liệu lên cùng hệ hàm riêng.} Tính $Q_n(t)=\int QX_n/\int X_n^2$, $A_n$, $B_n$. Không bỏ các dữ liệu mới phát sinh từ hàm nâng biên; chúng có thể khác không ngay khi dữ liệu gốc bằng không.

\textbf{Bước 3 — Giải ODE có nguồn cho từng mode.}


\[
y_n''+\omega_n^2y_n=Q_n(t),\qquad
y_n=A_n\cos\omega_nt+\frac{B_n}{\omega_n}\sin\omega_nt+\frac1{\omega_n}\int_0^t\sin(\omega_n(t-s))Q_n(s)ds.
\]


Với mode hằng Neumann, thay công thức bằng $y_0=A_0+B_0t+\int_0^t(t-s)Q_0(s)ds$.

\textbf{Bước 4 — Xử lý cộng hưởng trước khi chia.} Nguồn $\cos\alpha t$ có đáp ứng dữ liệu không là $(\cos\alpha t-\cos\omega t)/(\omega^2-\alpha^2)$; tại $\alpha=\omega$ lấy $t\sin\omega t/(2\omega)$. Nguồn $\sin\alpha t$ có đáp ứng $(\sin\alpha t-(\alpha/\omega)\sin\omega t)/(\omega^2-\alpha^2)$; tại cộng hưởng là $\sin\omega t/(2\omega^2)-t\cos\omega t/(2\omega)$. Không thế vào phân số $0/0$.

\textbf{Bước 5 — Cộng lại hàm nâng và kiểm tra duy nhất.} Thế $u=b+\sum y_nX_n$ vào PDE và mọi dữ liệu. Với hiệu hai nghiệm, nguồn và biên bằng không; nhân phương trình với $w_t$, tích phân từng phần để có $E'=0$, $E(0)=0$, suy $w=0$. Các góc không tương thích phải được ghi riêng.

\textbf{Chi tiết xác nhận tích phân nguồn.} Với $D(t)=\omega^{-1}\int_0^t\sin(\omega(t-s))Q(s)ds$, lần đạo hàm đầu cho $D'=\int_0^t\cos(\omega(t-s))Q(s)ds$, vì $\sin0=0$. Lần sau cho $D''=Q(t)-\omega\int_0^t\sin(\omega(t-s))Q(s)ds=Q-\omega^2D$. Do $D(0)=D'(0)=0$, tích phân này thêm đúng nguồn mà giữ nguyên hai dữ liệu đầu.
\subsection*{Lời giải áp dụng cho bài}
\textbf{Bước 1. Hệ sine và đáp ứng nguồn cos.} $a=L=1$, $X_n=\sin(n\pi x)$, $\omega_n=n\pi$. Khai triển $x$ có hệ số $c_n=2\int_0^1x\sin(n\pi x)dx=2(-1)^{n+1}/(n\pi)$. Đặt
\[C_\omega(\alpha,t)=\frac{\cos\alpha t-\cos\omega t}{\omega^2-\alpha^2},\qquad
C_\omega(\omega,t)=\frac{t\sin\omega t}{2\omega}.\]
Đây là nghiệm $C''+\omega^2C=\cos\alpha t$, $C(0)=C'(0)=0$; nhánh cộng hưởng lấy giới hạn hoặc tích phân Duhamel.

\textbf{Bước 2. Câu (i), biên không và nguồn $x\cos\alpha t$.} Dữ liệu đầu không, nên $y_n=c_nC_{\omega_n}$ và $u=\sum_{n\ge1}c_nC_{n\pi}\sin(n\pi x)$. Nếu $\alpha=m\pi$, chỉ mode $m$ dùng nhánh cộng hưởng, không chia $0$.

\textbf{Bước 3. Câu (ii), nguồn không và biên phải $\cos\alpha t$.} Lấy $b=x\cos\alpha t$, $w=u-b$. Khi đó $Q=\alpha^2x\cos\alpha t$, $w(x,0)=-x$, $w_t(x,0)=0$. Suy
\[u=x\cos\alpha t+\sum_{n\ge1}c_n[-\cos(n\pi t)+\alpha^2C_{n\pi}(\alpha,t)]\sin(n\pi x).\]
\textbf{Bước 4. Kiểm tra và góc.} Hàm nâng cho đúng biên, ODE cho nguồn của $w$ rồi triệt tiêu với $b_{tt}$. Tại $t=0$ khai triển $-x$ bù li độ của $b$ trong miền. Nhưng ở $(1,0)$, biên phải bằng $1$ trái li độ đầu $0$, nên câu (ii) không có nghiệm liên tục trên toàn bao đóng. Câu (i) có nguồn không khớp gia tốc biên tại góc phải nếu đòi $C^2$ đến góc; chuỗi vẫn giải trong miền với cách hiểu thích hợp.

\textbf{Chi tiết hệ số nguồn và đáp ứng.} Tích phân từng phần:
\[c_n=2\left[-\frac{x\cos(n\pi x)}{n\pi}+
\frac{\sin(n\pi x)}{n^2\pi^2}\right]_0^1
=-\frac{2(-1)^n}{n\pi}.\]
Với $\alpha^2\ne\omega^2$, thử $C_p=A\cos\alpha t$: $(\omega^2-\alpha^2)A=1$. Tổng $C=A\cos\alpha t+B\cos\omega t+D\sin\omega t$ có dữ liệu không khi $B=-A,D=0$. Với $\alpha=\omega$, thử $C=bt\sin\omega t$; đạo hàm hai lần cho $C''+\omega^2C=2b\omega\cos\omega t$, nên $b=1/(2\omega)$. Trong câu (ii), $b_{xx}=0$, $b_{tt}=-\alpha^2x\cos\alpha t$; vì $u_{tt}-u_{xx}=0$, phần $w$ có nguồn $-b_{tt}=\alpha^2x\cos\alpha t$ và dữ liệu $w_0=-x,w_{t,0}=0$.

\clearpage\noindent\texttt{\detokenize{VD_3_4_2_02}}\par
\section*{Lời giải bổ sung — Ví dụ 2 · Mục 3.4.2}

\subsection*{Kiến thức và trình tự giải}
\textbf{Bước 1 — Đưa điều kiện biên về thuần nhất.} Với $u(0,t)=\mu(t)$, $u(L,t)=\nu(t)$, đặt $b=(1-x/L)\mu+(x/L)\nu$. Với Neumann hai đầu đặt $b=\mu x+(\nu-\mu)x^2/(2L)$. Với Dirichlet–Neumann đặt $b=\mu+x\nu$. Viết $u=b+w$; khi đó $w$ có biên thuần nhất, nguồn $Q=F+a^2b_{xx}-b_{tt}$ và dữ liệu $f-b(x,0)$, $g-b_t(x,0)$.

\textbf{Bước 2 — Chiếu nguồn và dữ liệu lên cùng hệ hàm riêng.} Tính $Q_n(t)=\int QX_n/\int X_n^2$, $A_n$, $B_n$. Không bỏ các dữ liệu mới phát sinh từ hàm nâng biên; chúng có thể khác không ngay khi dữ liệu gốc bằng không.

\textbf{Bước 3 — Giải ODE có nguồn cho từng mode.}


\[
y_n''+\omega_n^2y_n=Q_n(t),\qquad
y_n=A_n\cos\omega_nt+\frac{B_n}{\omega_n}\sin\omega_nt+\frac1{\omega_n}\int_0^t\sin(\omega_n(t-s))Q_n(s)ds.
\]


Với mode hằng Neumann, thay công thức bằng $y_0=A_0+B_0t+\int_0^t(t-s)Q_0(s)ds$.

\textbf{Bước 4 — Xử lý cộng hưởng trước khi chia.} Nguồn $\cos\alpha t$ có đáp ứng dữ liệu không là $(\cos\alpha t-\cos\omega t)/(\omega^2-\alpha^2)$; tại $\alpha=\omega$ lấy $t\sin\omega t/(2\omega)$. Nguồn $\sin\alpha t$ có đáp ứng $(\sin\alpha t-(\alpha/\omega)\sin\omega t)/(\omega^2-\alpha^2)$; tại cộng hưởng là $\sin\omega t/(2\omega^2)-t\cos\omega t/(2\omega)$. Không thế vào phân số $0/0$.

\textbf{Bước 5 — Cộng lại hàm nâng và kiểm tra duy nhất.} Thế $u=b+\sum y_nX_n$ vào PDE và mọi dữ liệu. Với hiệu hai nghiệm, nguồn và biên bằng không; nhân phương trình với $w_t$, tích phân từng phần để có $E'=0$, $E(0)=0$, suy $w=0$. Các góc không tương thích phải được ghi riêng.

\textbf{Chi tiết xác nhận tích phân nguồn.} Với $D(t)=\omega^{-1}\int_0^t\sin(\omega(t-s))Q(s)ds$, lần đạo hàm đầu cho $D'=\int_0^t\cos(\omega(t-s))Q(s)ds$, vì $\sin0=0$. Lần sau cho $D''=Q(t)-\omega\int_0^t\sin(\omega(t-s))Q(s)ds=Q-\omega^2D$. Do $D(0)=D'(0)=0$, tích phân này thêm đúng nguồn mà giữ nguyên hai dữ liệu đầu.
\subsection*{Lời giải áp dụng cho bài}
\textbf{Bước 1. Tính duy nhất.} Với hiệu $v$ của hai nghiệm cùng dữ liệu và biên Neumann, $v_{tt}=a^2v_{xx}$, $v_x(0,t)=v_x(L,t)=0$. Đạo hàm $E=\frac12\int(v_t^2+a^2v_x^2)$ cho $E'=a^2[v_tv_x]_0^L=0$. $E(0)=0$ suy $v_t=v_x=0$; li độ đầu bằng $0$ xác định hằng số là $0$.

\textbf{Bước 2. Nâng biên tổng quát.} Với $u_x(0,t)=\mu(t)$, $u_x(L,t)=\nu(t)$, chọn $b=\mu x+(\nu-\mu)x^2/(2L)$, vì $b_x(0)=\mu$, $b_x(L)=\nu$. Đặt $w=u-b$; nguồn mới là $Q=F+a^2(\nu-\mu)/L-b_{tt}$, hai dữ liệu mới là $f-b_0$, $g-b_{t,0}$.

\textbf{Bước 3. Tách mode.} Với $X_0=1$, $X_n=\cos(n\pi x/L)$, lấy hệ số $A_0,B_0,Q_0$ là trung bình và các hệ số $n\ge1$ là $2/L$ lần tích phân với $X_n$. Mode hằng giải $y_0''=Q_0$, nên $y_0=A_0+B_0t+\int_0^t(t-s)Q_0(s)ds$. Mode còn lại giải
\[y_n=A_n\cos\omega_nt+\frac{B_n}{\omega_n}\sin\omega_nt+\frac1{\omega_n}\int_0^t\sin(\omega_n(t-s))Q_n(s)ds.\]
\textbf{Bước 4. Trở lại nghiệm.} $u=b+y_0+\sum_{n\ge1}y_nX_n$. Kiểm tra $t=0$ sau khi cộng hàm nâng; kiểm tra từng đạo hàm biên. Công thức mode hằng khác các mode dao động, nên không áp dụng mẫu $1/\omega_0$. Nếu dữ liệu không tương thích $f'(0)=\mu(0)$, $g'(0)=\mu'(0)$ (và tương tự đầu phải), nêu rõ mức trơn và xét nghiệm yếu thay vì khẳng định cổ điển đến góc.

\textbf{Chi tiết tích phân từng phần trong đạo hàm năng lượng.} Đặt $E=\frac12\int_0^L(u_t^2+a^2u_x^2+\kappa u^2)dx$. Với $u_{tt}-a^2u_{xx}+\gamma u_t+\kappa u=F$,
\[\begin{aligned}
E'&=\int_0^L(u_tu_{tt}+a^2u_xu_{xt}+\kappa uu_t)dx\\
&=\int_0^Lu_t(a^2u_{xx}-\gamma u_t-\kappa u+F)dx
+a^2\int_0^Lu_xu_{xt}dx+\kappa\int_0^Luu_tdx\\
&=a^2[u_tu_x]_0^L-a^2\int_0^Lu_{tx}u_xdx
+a^2\int_0^Lu_xu_{xt}dx-\gamma\int_0^Lu_t^2dx+\int_0^LFu_tdx\\
&=a^2[u_tu_x]_0^L-\gamma\int_0^Lu_t^2dx+\int_0^LFu_tdx.
\end{aligned}\]
Hai hạng $\kappa uu_t$ triệt tiêu ngay khi thế PDE; hai tích phân $u_xu_{xt}$ triệt tiêu sau tích phân từng phần. Với biên Dirichlet thuần nhất, đạo hàm vết theo $t$ cho $u_t=0$; với biên Neumann thuần nhất, $u_x=0$. Nếu định nghĩa năng lượng không có $1/2$, nhân cả đẳng thức đạo hàm với $2$.

\textbf{Chi tiết đạo hàm đáp ứng mode nguồn.} Đặt $D_n(t)=\omega_n^{-1}\int_0^t\sin(\omega_n(t-s))Q_n(s)ds$. Vì $\sin0=0$,
\[D_n'=\int_0^t\cos(\omega_n(t-s))Q_n(s)ds,\quad
D_n''=Q_n(t)-\omega_n\int_0^t\sin(\omega_n(t-s))Q_n(s)ds=Q_n-\omega_n^2D_n.\]
Tại $0$, $D_n=D_n'=0$. Vì thế tích phân vừa cho đúng nguồn, vừa giữ nguyên cả hai dữ liệu ban đầu của các hạng thuần nhất. Với mode hằng, $D_0=\int_0^t(t-s)Q_0(s)ds$, $D_0'=\int_0^tQ_0(s)ds$, $D_0''=Q_0(t)$, không chia cho $\omega_0=0$.

\clearpage\noindent\texttt{\detokenize{BT_3_4_3_01}}\par
\section*{Lời giải bổ sung — Bài tập 1 · Mục 3.4.3}

\subsection*{Kiến thức và trình tự giải}
\textbf{Bước 1 — Đưa điều kiện biên về thuần nhất.} Với $u(0,t)=\mu(t)$, $u(L,t)=\nu(t)$, đặt $b=(1-x/L)\mu+(x/L)\nu$. Với Neumann hai đầu đặt $b=\mu x+(\nu-\mu)x^2/(2L)$. Với Dirichlet–Neumann đặt $b=\mu+x\nu$. Viết $u=b+w$; khi đó $w$ có biên thuần nhất, nguồn $Q=F+a^2b_{xx}-b_{tt}$ và dữ liệu $f-b(x,0)$, $g-b_t(x,0)$.

\textbf{Bước 2 — Chiếu nguồn và dữ liệu lên cùng hệ hàm riêng.} Tính $Q_n(t)=\int QX_n/\int X_n^2$, $A_n$, $B_n$. Không bỏ các dữ liệu mới phát sinh từ hàm nâng biên; chúng có thể khác không ngay khi dữ liệu gốc bằng không.

\textbf{Bước 3 — Giải ODE có nguồn cho từng mode.}


\[
y_n''+\omega_n^2y_n=Q_n(t),\qquad
y_n=A_n\cos\omega_nt+\frac{B_n}{\omega_n}\sin\omega_nt+\frac1{\omega_n}\int_0^t\sin(\omega_n(t-s))Q_n(s)ds.
\]


Với mode hằng Neumann, thay công thức bằng $y_0=A_0+B_0t+\int_0^t(t-s)Q_0(s)ds$.

\textbf{Bước 4 — Xử lý cộng hưởng trước khi chia.} Nguồn $\cos\alpha t$ có đáp ứng dữ liệu không là $(\cos\alpha t-\cos\omega t)/(\omega^2-\alpha^2)$; tại $\alpha=\omega$ lấy $t\sin\omega t/(2\omega)$. Nguồn $\sin\alpha t$ có đáp ứng $(\sin\alpha t-(\alpha/\omega)\sin\omega t)/(\omega^2-\alpha^2)$; tại cộng hưởng là $\sin\omega t/(2\omega^2)-t\cos\omega t/(2\omega)$. Không thế vào phân số $0/0$.

\textbf{Bước 5 — Cộng lại hàm nâng và kiểm tra duy nhất.} Thế $u=b+\sum y_nX_n$ vào PDE và mọi dữ liệu. Với hiệu hai nghiệm, nguồn và biên bằng không; nhân phương trình với $w_t$, tích phân từng phần để có $E'=0$, $E(0)=0$, suy $w=0$. Các góc không tương thích phải được ghi riêng.

\textbf{Chi tiết xác nhận tích phân nguồn.} Với $D(t)=\omega^{-1}\int_0^t\sin(\omega(t-s))Q(s)ds$, lần đạo hàm đầu cho $D'=\int_0^t\cos(\omega(t-s))Q(s)ds$, vì $\sin0=0$. Lần sau cho $D''=Q(t)-\omega\int_0^t\sin(\omega(t-s))Q(s)ds=Q-\omega^2D$. Do $D(0)=D'(0)=0$, tích phân này thêm đúng nguồn mà giữ nguyên hai dữ liệu đầu.
\subsection*{Lời giải áp dụng cho bài}
\textbf{Bước 1. Tính hệ số và nhân đáp ứng.} Trên $(0,\pi)$, khai triển $1$ theo sine có $c_n=2(1-(-1)^n)/(\pi n)$. Đặt
\[R_\omega(\alpha,t)=\frac{\sin\alpha t-(\alpha/\omega)\sin\omega t}{\omega^2-\alpha^2}.\]
Nó thỏa $R''+\omega^2R=\sin\alpha t$, $R(0)=R'(0)=0$. Khi $\alpha=\omega$, dùng $R=\sin\omega t/(2\omega^2)-t\cos\omega t/(2\omega)$.

\textbf{Bước 2. Câu (i).} $a=1$, biên và dữ liệu đầu không, nguồn $\sin\alpha t$. Chiếu nguồn cho $y_n''+n^2y_n=c_n\sin\alpha t$, nên $u=\sum c_nR_n\sin nx$. Tại cộng hưởng $\alpha=m$, thay đúng mode $m$ bằng nhánh giới hạn; nếu $c_m=0$, mode ấy không được kích thích.

\textbf{Bước 3. Câu (ii).} $a=2$, nguồn không, cả hai biên bằng $\sin\alpha t$. Lấy $b=\sin\alpha t$, $w=u-b$, suy $Q=\alpha^2\sin\alpha t$, $w_0=0$, $w_{t,0}=-\alpha$. Hệ số vận tốc là $-\alpha c_n$, vì vậy
\[u=\sin\alpha t+\sum_{n\ge1}c_n\left[-\frac\alpha{2n}\sin2nt+\alpha^2R_{2n}(\alpha,t)\right]\sin nx.\]
Đạo hàm ban đầu của hàm nâng và chuỗi triệt tiêu trong miền. Tại góc, vận tốc biên $\alpha$ trái $g=0$, nên không có nghiệm $C^1$ trên toàn bao đóng khi $\alpha\ne0$.

\textbf{Bước 4. Câu (iii).} $a=1$, nguồn $1$, biên trái $0$, biên phải $-\pi^2/2$, dữ liệu đầu không. Chọn $P=-x^2/2$: $P_{tt}-P_{xx}=1$ và đúng biên. $w=u-P$ thuần nhất, li độ $x^2/2$, vận tốc $0$. Tích phân từng phần hai lần cho
\[A_n=\frac1\pi\int_0^\pi x^2\sin nx\,dx
=-\frac{\pi(-1)^n}{n}+\frac{2((-1)^n-1)}{\pi n^3}.\]
Suy $u=-x^2/2+\sum A_n\cos nt\sin nx$. Li độ đầu và biên phải mâu thuẫn ở góc; công thức là nghiệm từng miền/ yếu. Mỗi câu đã kiểm tra nguồn qua ODE và dữ liệu qua hệ số, không chỉ ghi chuỗi hình thức.

\textbf{Chi tiết đáp ứng sine và tích phân đa thức.} Không cộng hưởng, thử $R_p=A\sin\alpha t$, có $A=1/(\omega^2-\alpha^2)$. Viết $R=R_p+B\cos\omega t+D\sin\omega t$; dữ liệu $R(0)=0$ cho $B=0$, $R'(0)=0$ cho $D=-\alpha A/\omega$. Ở cộng hưởng, $-t\cos\omega t/(2\omega)$ có nguồn $\sin\omega t$ nhưng vận tốc đầu $-1/(2\omega)$; cộng $\sin\omega t/(2\omega^2)$ để bù.

Trong (iii), đặt $I=\int_0^\pi x^2\sin nx\,dx$. Hai lần tích phân:
\[\begin{aligned}
I&=\left[-x^2\cos nx/n\right]_0^\pi+\frac2n\int_0^\pi x\cos nx\,dx,\\
\int_0^\pi x\cos nx\,dx
&=\left[x\sin nx/n+\cos nx/n^2\right]_0^\pi
=\frac{(-1)^n-1}{n^2}.
\end{aligned}\]
Vậy $I=-\pi^2(-1)^n/n+2((-1)^n-1)/n^3$, rồi chia $\pi$ cho $A_n$.

\clearpage\noindent\texttt{\detokenize{BT_3_4_3_02}}\par
\section*{Lời giải bổ sung — Bài tập 2 · Mục 3.4.3}

\subsection*{Kiến thức và trình tự giải}
\textbf{Bước 1 — Đưa điều kiện biên về thuần nhất.} Với $u(0,t)=\mu(t)$, $u(L,t)=\nu(t)$, đặt $b=(1-x/L)\mu+(x/L)\nu$. Với Neumann hai đầu đặt $b=\mu x+(\nu-\mu)x^2/(2L)$. Với Dirichlet–Neumann đặt $b=\mu+x\nu$. Viết $u=b+w$; khi đó $w$ có biên thuần nhất, nguồn $Q=F+a^2b_{xx}-b_{tt}$ và dữ liệu $f-b(x,0)$, $g-b_t(x,0)$.

\textbf{Bước 2 — Chiếu nguồn và dữ liệu lên cùng hệ hàm riêng.} Tính $Q_n(t)=\int QX_n/\int X_n^2$, $A_n$, $B_n$. Không bỏ các dữ liệu mới phát sinh từ hàm nâng biên; chúng có thể khác không ngay khi dữ liệu gốc bằng không.

\textbf{Bước 3 — Giải ODE có nguồn cho từng mode.}


\[
y_n''+\omega_n^2y_n=Q_n(t),\qquad
y_n=A_n\cos\omega_nt+\frac{B_n}{\omega_n}\sin\omega_nt+\frac1{\omega_n}\int_0^t\sin(\omega_n(t-s))Q_n(s)ds.
\]


Với mode hằng Neumann, thay công thức bằng $y_0=A_0+B_0t+\int_0^t(t-s)Q_0(s)ds$.

\textbf{Bước 4 — Xử lý cộng hưởng trước khi chia.} Nguồn $\cos\alpha t$ có đáp ứng dữ liệu không là $(\cos\alpha t-\cos\omega t)/(\omega^2-\alpha^2)$; tại $\alpha=\omega$ lấy $t\sin\omega t/(2\omega)$. Nguồn $\sin\alpha t$ có đáp ứng $(\sin\alpha t-(\alpha/\omega)\sin\omega t)/(\omega^2-\alpha^2)$; tại cộng hưởng là $\sin\omega t/(2\omega^2)-t\cos\omega t/(2\omega)$. Không thế vào phân số $0/0$.

\textbf{Bước 5 — Cộng lại hàm nâng và kiểm tra duy nhất.} Thế $u=b+\sum y_nX_n$ vào PDE và mọi dữ liệu. Với hiệu hai nghiệm, nguồn và biên bằng không; nhân phương trình với $w_t$, tích phân từng phần để có $E'=0$, $E(0)=0$, suy $w=0$. Các góc không tương thích phải được ghi riêng.

\textbf{Chi tiết xác nhận tích phân nguồn.} Với $D(t)=\omega^{-1}\int_0^t\sin(\omega(t-s))Q(s)ds$, lần đạo hàm đầu cho $D'=\int_0^t\cos(\omega(t-s))Q(s)ds$, vì $\sin0=0$. Lần sau cho $D''=Q(t)-\omega\int_0^t\sin(\omega(t-s))Q(s)ds=Q-\omega^2D$. Do $D(0)=D'(0)=0$, tích phân này thêm đúng nguồn mà giữ nguyên hai dữ liệu đầu.
\subsection*{Lời giải áp dụng cho bài}
\textbf{Bước 1. Giữ mode hằng Neumann.} Hệ hàm riêng trên $(0,\pi)$ là $1,\cos nx$. Mode hằng dùng $y_0''=Q_0$; các mode $n\ge1$ dùng ODE dao động và $R_\omega$ như Bài 1.

\textbf{Bước 2. Câu (i).} Biên Neumann không, $a=1$, nguồn $\sin\alpha t$ không phụ thuộc $x$. Chỉ mode $0$ có nguồn khác không. Với dữ liệu không,
\[u=\int_0^t(t-s)\sin\alpha s\,ds=t/\alpha-\sin\alpha t/\alpha^2.\]
Thế lại cho $u_{tt}=\sin\alpha t$, $u_{xx}=0$, dữ liệu không. Nếu $\alpha=0$, lấy giới hạn $u=0$.

\textbf{Bước 3. Câu (ii).} $a=2$, $\mu=\nu=\sin\alpha t$, nguồn không. Lấy $b=x\sin\alpha t$, $w=u-b$, suy $Q=\alpha^2x\sin\alpha t$, $w_0=0$, $w_{t,0}=-\alpha x$. Khai triển $x$ theo cos có $d_0=\pi/2$, $d_n=2((-1)^n-1)/(\pi n^2)$, do tích phân từng phần.

\textbf{Bước 4. Giải từng ODE và ghép.} Mode hằng là $y_0=-d_0\sin\alpha t$ (đạo hàm đầu $-\alpha d_0$ và $y_0''=\alpha^2d_0\sin\alpha t$). Các mode khác là $y_n=d_n[-\alpha\sin2nt/(2n)+\alpha^2R_{2n}]$. Do đó
\[u=(x-\pi/2)\sin\alpha t+\sum_{n\ge1}d_n[-\alpha\sin2nt/(2n)+\alpha^2R_{2n}(\alpha,t)]\cos nx.\]
Hàm nâng cho đúng đạo hàm biên; các ODE cho đúng PDE. Tại góc, $g'=0$ trái $\mu'(0)=\nu'(0)=\alpha$, nên không $C^2$ đến góc nếu $\alpha\ne0$. Tại $\alpha=2n$ phải dùng nhánh cộng hưởng, không chia mẫu bằng không.

\textbf{Chi tiết tích phân nguồn không gian hằng và hệ số cos.} Với $\alpha\ne0$,
\[\int_0^t(t-s)\sin\alpha s\,ds
=\left[-(t-s)\cos\alpha s/\alpha\right]_0^t-\frac1\alpha\int_0^t\cos\alpha s\,ds
=t/\alpha-\sin\alpha t/\alpha^2.\]
Với dữ liệu $x$, $d_0=\pi^{-1}[x^2/2]_0^\pi=\pi/2$ và
\[d_n=\frac2\pi\left[x\sin nx/n+\cos nx/n^2\right]_0^\pi
=\frac{2((-1)^n-1)}{\pi n^2}.\]
Mode hằng giải $y_0''=\alpha^2d_0\sin\alpha t$ với $y_0'(0)=-\alpha d_0$: tích phân cho $y_0'=-\alpha d_0\cos\alpha t$, $y_0=-d_0\sin\alpha t$. Đây là phép bù tạo ra $(x-\pi/2)\sin\alpha t$.

\clearpage\noindent\texttt{\detokenize{BT_3_4_3_03}}\par
\section*{Lời giải bổ sung — Bài tập 3 · Mục 3.4.3}

\subsection*{Kiến thức và trình tự giải}
\textbf{Bước 1 — Đưa điều kiện biên về thuần nhất.} Với $u(0,t)=\mu(t)$, $u(L,t)=\nu(t)$, đặt $b=(1-x/L)\mu+(x/L)\nu$. Với Neumann hai đầu đặt $b=\mu x+(\nu-\mu)x^2/(2L)$. Với Dirichlet–Neumann đặt $b=\mu+x\nu$. Viết $u=b+w$; khi đó $w$ có biên thuần nhất, nguồn $Q=F+a^2b_{xx}-b_{tt}$ và dữ liệu $f-b(x,0)$, $g-b_t(x,0)$.

\textbf{Bước 2 — Chiếu nguồn và dữ liệu lên cùng hệ hàm riêng.} Tính $Q_n(t)=\int QX_n/\int X_n^2$, $A_n$, $B_n$. Không bỏ các dữ liệu mới phát sinh từ hàm nâng biên; chúng có thể khác không ngay khi dữ liệu gốc bằng không.

\textbf{Bước 3 — Giải ODE có nguồn cho từng mode.}


\[
y_n''+\omega_n^2y_n=Q_n(t),\qquad
y_n=A_n\cos\omega_nt+\frac{B_n}{\omega_n}\sin\omega_nt+\frac1{\omega_n}\int_0^t\sin(\omega_n(t-s))Q_n(s)ds.
\]


Với mode hằng Neumann, thay công thức bằng $y_0=A_0+B_0t+\int_0^t(t-s)Q_0(s)ds$.

\textbf{Bước 4 — Xử lý cộng hưởng trước khi chia.} Nguồn $\cos\alpha t$ có đáp ứng dữ liệu không là $(\cos\alpha t-\cos\omega t)/(\omega^2-\alpha^2)$; tại $\alpha=\omega$ lấy $t\sin\omega t/(2\omega)$. Nguồn $\sin\alpha t$ có đáp ứng $(\sin\alpha t-(\alpha/\omega)\sin\omega t)/(\omega^2-\alpha^2)$; tại cộng hưởng là $\sin\omega t/(2\omega^2)-t\cos\omega t/(2\omega)$. Không thế vào phân số $0/0$.

\textbf{Bước 5 — Cộng lại hàm nâng và kiểm tra duy nhất.} Thế $u=b+\sum y_nX_n$ vào PDE và mọi dữ liệu. Với hiệu hai nghiệm, nguồn và biên bằng không; nhân phương trình với $w_t$, tích phân từng phần để có $E'=0$, $E(0)=0$, suy $w=0$. Các góc không tương thích phải được ghi riêng.

\textbf{Chi tiết xác nhận tích phân nguồn.} Với $D(t)=\omega^{-1}\int_0^t\sin(\omega(t-s))Q(s)ds$, lần đạo hàm đầu cho $D'=\int_0^t\cos(\omega(t-s))Q(s)ds$, vì $\sin0=0$. Lần sau cho $D''=Q(t)-\omega\int_0^t\sin(\omega(t-s))Q(s)ds=Q-\omega^2D$. Do $D(0)=D'(0)=0$, tích phân này thêm đúng nguồn mà giữ nguyên hai dữ liệu đầu.
\subsection*{Lời giải áp dụng cho bài}
\textbf{Bước 1. Chứng minh duy nhất.} Hiệu $v$ có nguồn và dữ liệu không, $v(0,t)=v_x(L,t)=0$. Đạo hàm năng lượng $E'=a^2[v_tv_x]_0^L$: ở đầu trái $v_t=0$ vì vết $v=0$, đầu phải $v_x=0$. Suy $E=E(0)=0$, $v$ hằng cả không gian và thời gian, rồi dữ liệu đầu buộc $v=0$.

\textbf{Bước 2. Nâng biên.} Lấy $b=\mu(t)+x\nu(t)$ để $b(0,t)=\mu$, $b_x(L,t)=\nu$. $w=u-b$ có biên hỗn hợp không, nguồn $Q=F-\mu''-x\nu''$, li độ $f-\mu(0)-x\nu(0)$, vận tốc $g-\mu'(0)-x\nu'(0)$.

\textbf{Bước 3. Chiếu nguồn cùng dữ liệu.} $k_n=(n+1/2)\pi/L$, $X_n=\sin k_nx$, chuẩn $L/2$. Tính $A_n=(2/L)\int w_0X_n$, $B_n=(2/L)\int w_{t,0}X_n$, $Q_n(s)=(2/L)\int Q(x,s)X_n(x)dx$.

\textbf{Bước 4. Giải và trả về biến gốc.}
\[u=b+\sum_{n\ge0}\left[A_n\cos\omega_nt+\frac{B_n}{\omega_n}\sin\omega_nt+\int_0^t\frac{\sin(\omega_n(t-s))}{\omega_n}Q_n(s)ds\right]X_n,\quad\omega_n=ak_n.\]
Với $\mu=\nu=0$, $b=0$ và dữ liệu/nguồn chiếu chính là dữ liệu gốc. Đạo hàm hai lần của tích phân mode cho $Q_n$, còn hai dữ liệu ban đầu của nó bằng không. Cộng hàm nâng xác nhận mọi biên. Cần điều kiện tương thích ở góc nếu muốn nghiệm cổ điển, chẳng hạn $f(0)=\mu(0)$, $g(0)=\mu'(0)$, $f'(L)=\nu(0)$, $g'(L)=\nu'(0)$.

\textbf{Chi tiết tích phân từng phần trong đạo hàm năng lượng.} Đặt $E=\frac12\int_0^L(u_t^2+a^2u_x^2+\kappa u^2)dx$. Với $u_{tt}-a^2u_{xx}+\gamma u_t+\kappa u=F$,
\[\begin{aligned}
E'&=\int_0^L(u_tu_{tt}+a^2u_xu_{xt}+\kappa uu_t)dx\\
&=\int_0^Lu_t(a^2u_{xx}-\gamma u_t-\kappa u+F)dx
+a^2\int_0^Lu_xu_{xt}dx+\kappa\int_0^Luu_tdx\\
&=a^2[u_tu_x]_0^L-a^2\int_0^Lu_{tx}u_xdx
+a^2\int_0^Lu_xu_{xt}dx-\gamma\int_0^Lu_t^2dx+\int_0^LFu_tdx\\
&=a^2[u_tu_x]_0^L-\gamma\int_0^Lu_t^2dx+\int_0^LFu_tdx.
\end{aligned}\]
Hai hạng $\kappa uu_t$ triệt tiêu ngay khi thế PDE; hai tích phân $u_xu_{xt}$ triệt tiêu sau tích phân từng phần. Với biên Dirichlet thuần nhất, đạo hàm vết theo $t$ cho $u_t=0$; với biên Neumann thuần nhất, $u_x=0$. Nếu định nghĩa năng lượng không có $1/2$, nhân cả đẳng thức đạo hàm với $2$.

\textbf{Chi tiết đạo hàm đáp ứng mode nguồn.} Đặt $D_n(t)=\omega_n^{-1}\int_0^t\sin(\omega_n(t-s))Q_n(s)ds$. Vì $\sin0=0$,
\[D_n'=\int_0^t\cos(\omega_n(t-s))Q_n(s)ds,\quad
D_n''=Q_n(t)-\omega_n\int_0^t\sin(\omega_n(t-s))Q_n(s)ds=Q_n-\omega_n^2D_n.\]
Tại $0$, $D_n=D_n'=0$. Vì thế tích phân vừa cho đúng nguồn, vừa giữ nguyên cả hai dữ liệu ban đầu của các hạng thuần nhất. Với mode hằng, $D_0=\int_0^t(t-s)Q_0(s)ds$, $D_0'=\int_0^tQ_0(s)ds$, $D_0''=Q_0(t)$, không chia cho $\omega_0=0$.

\clearpage\noindent\texttt{\detokenize{BT_3_4_3_04}}\par
\section*{Lời giải bổ sung — Bài tập 4 · Mục 3.4.3}

\subsection*{Kiến thức và trình tự giải}
\textbf{Bước 1 — Đưa điều kiện biên về thuần nhất.} Với $u(0,t)=\mu(t)$, $u(L,t)=\nu(t)$, đặt $b=(1-x/L)\mu+(x/L)\nu$. Với Neumann hai đầu đặt $b=\mu x+(\nu-\mu)x^2/(2L)$. Với Dirichlet–Neumann đặt $b=\mu+x\nu$. Viết $u=b+w$; khi đó $w$ có biên thuần nhất, nguồn $Q=F+a^2b_{xx}-b_{tt}$ và dữ liệu $f-b(x,0)$, $g-b_t(x,0)$.

\textbf{Bước 2 — Chiếu nguồn và dữ liệu lên cùng hệ hàm riêng.} Tính $Q_n(t)=\int QX_n/\int X_n^2$, $A_n$, $B_n$. Không bỏ các dữ liệu mới phát sinh từ hàm nâng biên; chúng có thể khác không ngay khi dữ liệu gốc bằng không.

\textbf{Bước 3 — Giải ODE có nguồn cho từng mode.}


\[
y_n''+\omega_n^2y_n=Q_n(t),\qquad
y_n=A_n\cos\omega_nt+\frac{B_n}{\omega_n}\sin\omega_nt+\frac1{\omega_n}\int_0^t\sin(\omega_n(t-s))Q_n(s)ds.
\]


Với mode hằng Neumann, thay công thức bằng $y_0=A_0+B_0t+\int_0^t(t-s)Q_0(s)ds$.

\textbf{Bước 4 — Xử lý cộng hưởng trước khi chia.} Nguồn $\cos\alpha t$ có đáp ứng dữ liệu không là $(\cos\alpha t-\cos\omega t)/(\omega^2-\alpha^2)$; tại $\alpha=\omega$ lấy $t\sin\omega t/(2\omega)$. Nguồn $\sin\alpha t$ có đáp ứng $(\sin\alpha t-(\alpha/\omega)\sin\omega t)/(\omega^2-\alpha^2)$; tại cộng hưởng là $\sin\omega t/(2\omega^2)-t\cos\omega t/(2\omega)$. Không thế vào phân số $0/0$.

\textbf{Bước 5 — Cộng lại hàm nâng và kiểm tra duy nhất.} Thế $u=b+\sum y_nX_n$ vào PDE và mọi dữ liệu. Với hiệu hai nghiệm, nguồn và biên bằng không; nhân phương trình với $w_t$, tích phân từng phần để có $E'=0$, $E(0)=0$, suy $w=0$. Các góc không tương thích phải được ghi riêng.

\textbf{Chi tiết xác nhận tích phân nguồn.} Với $D(t)=\omega^{-1}\int_0^t\sin(\omega(t-s))Q(s)ds$, lần đạo hàm đầu cho $D'=\int_0^t\cos(\omega(t-s))Q(s)ds$, vì $\sin0=0$. Lần sau cho $D''=Q(t)-\omega\int_0^t\sin(\omega(t-s))Q(s)ds=Q-\omega^2D$. Do $D(0)=D'(0)=0$, tích phân này thêm đúng nguồn mà giữ nguyên hai dữ liệu đầu.
\subsection*{Lời giải áp dụng cho bài}
\textbf{Bước 1. Nhận dạng một mode.} Biên cố định trên $(0,1)$ và nguồn $\sin\pi x\sin t$ gợi $u=\sin\pi x\,T(t)$. Đạo hàm $u_{xx}=-\pi^2\sin\pi x\,T$; vì tốc độ bằng $1$, PDE trở thành $T''+\pi^2T=\sin t$, dữ liệu $T(0)=T'(0)=0$.

\textbf{Bước 2. Nghiệm riêng và thuần nhất.} Nguồn tần số $1$ khác $\pi$, nên thử $T_p=C\sin t$, được $C=1/(\pi^2-1)$. Nghiệm tổng $T=C\sin t+A\cos\pi t+B\sin\pi t$. Dữ liệu đầu cho $A=0$, $C+\pi B=0$, nên $B=-1/(\pi(\pi^2-1))$.

\textbf{Bước 3. Nghiệm.}
\[u(x,t)=\frac{\sin t-\sin\pi t/\pi}{\pi^2-1}\sin\pi x.\]
\textbf{Bước 4. Kiểm tra.} Hai biên bằng không vì $\sin\pi x=0$; $T(0)=0$ và $T'(0)=(1-1)/(\pi^2-1)=0$. Thế $T''+\pi^2T$ cho $\sin t$, nên đúng nguồn. Tính duy nhất theo năng lượng của hiệu hai nghiệm biên cố định.

\clearpage\noindent\texttt{\detokenize{VD_3_5_2_01}}\par
\section*{Lời giải bổ sung — Ví dụ 1 · Mục 3.5.2}

\subsection*{Kiến thức và trình tự giải}
\textbf{Bước 1 — Chọn năng lượng theo toán tử.} Với sóng dùng $E=\frac12\int(u_t^2+a^2u_x^2)dx$. Với Klein–Gordon có thêm $u$, thêm $\frac12\int u^2dx$. Với dầm $u_{tt}=-a^2u_{xxxx}$ dùng $E=\frac12\int(u_t^2+a^2u_{xx}^2)dx$, không thay $u_{xx}$ bằng $u_x$.

\textbf{Bước 2 — Đạo hàm và tích phân từng phần.} Với $u_{tt}-a^2u_{xx}+\gamma u_t+\kappa u=F$, năng lượng có hạng $\kappa u^2$ cho


\[
E'=a^2[u_tu_x]_0^L-\gamma\int u_t^2dx+\int Fu_tdx.
\]


Với dầm, tích phân hai lần cho $E'=a^2[u_{tx}u_{xx}-u_tu_{xxx}]_0^L$. Viết riêng từng hạng biên rồi dùng điều kiện biên và đạo hàm theo thời gian của điều kiện đó để cho chúng bằng không.

\textbf{Bước 3 — Suy hệ quả.} Không cản, không nguồn thì $E$ bảo toàn. Có cản $\gamma>0$ thì $E'\le0$; nếu $E$ không có hệ số $1/2$, $E'=-2\gamma\int u_t^2\ge-2\gamma E$, suy $e^{-2\gamma t}E(0)\le E(t)\le E(0)$. Với hiệu hai nghiệm cùng dữ liệu, $E(0)=0$ nên có duy nhất.

\textbf{Bước 4 — Tách biến nếu phải tìm nghiệm cụ thể.} Sóng có cản $\gamma=1$ cho $T''+T'+a^2\lambda T=0$. Đặt $T=e^{-t/2}Y$, thu $Y''+(a^2\lambda-1/4)Y=0$, $Y(0)=A$, $Y'(0)=B+A/2$. Phân biệt dao động, nghiệm kép và hàm hyperbolic. Dầm cho $X''''=k^4X$, $T''+a^2k^4T=0$; dùng đủ bốn điều kiện biên để tìm phương trình của $kL$.

\textbf{Bước 5 — Kiểm tra hệ số và khẳng định của đề.} Năng lượng đầu chứa cả $g^2$ và các đạo hàm của $f$. Thử $f=0$, $g\ne0$ để phát hiện công thức bỏ vận tốc. Với dầm, nghiệm $k=0$ phải giải trực tiếp bằng đa thức bậc ba; nghiệm $0$ của phương trình lượng giác–hyperbolic có thể chỉ là nghiệm giả do cơ sở suy biến.

\textbf{Chi tiết phép triệt tiêu tích phân.} Đạo hàm động năng cho $\int u_tu_{tt}$; dùng PDE thay $u_{tt}$. Với sóng, $\int u_tu_{xx}=[u_tu_x]-\int u_{tx}u_x$, và tích phân cuối triệt tiêu đạo hàm thế năng $\int u_xu_{xt}$. Với dầm,


\[
\int u_tu_{xxxx}=[u_tu_{xxx}-u_{tx}u_{xx}]+\int u_{txx}u_{xx}.
\]


Tích phân cuối triệt tiêu đạo hàm $\frac12\int u_{xx}^2$ sau khi nhân hệ số và giữ đúng dấu PDE. Viết hạng biên trước, sau đó mới dùng điều kiện của từng đầu.

\textbf{Bước 1 — Tách biến và chọn bài toán trị riêng theo biên.} Đặt $u=X(x)T(t)$ trong $u_{tt}=a^2u_{xx}$, suy $X''+\lambda X=0$, $T''+a^2\lambda T=0$. Tích phân từng phần cho $\lambda\int X^2=\int (X')^2$ khi hạng biên bằng không; do đó $\lambda\ge0$. Giải cả trường hợp $\lambda=0$ trước khi xét $\lambda=k^2>0$.

\textbf{Bước 2 — Tìm hệ hàm riêng.} Trên $(0,L)$: Dirichlet hai đầu dùng $\sin(n\pi x/L)$, $n\ge1$; Neumann hai đầu dùng $1$ và $\cos(n\pi x/L)$; Dirichlet–Neumann dùng $\sin((n+1/2)\pi x/L)$, $n\ge0$; Neumann–Dirichlet dùng cos của các tần số nửa nguyên ấy. Chuẩn bình phương của mode khác không là $L/2$, của mode hằng là $L$.

\textbf{Bước 3 — Chiếu cả hai dữ liệu.} Với mode khác không đặt


\[
A_n=\frac2L\int_0^L fX_n\,dx,\quad B_n=\frac2L\int_0^L gX_n\,dx,\quad \omega_n=ak_n.
\]


Tích phân từng khoảng nếu dữ liệu chia đoạn; dùng tích phân từng phần hoặc công thức tích–tổng khi là đa thức/lượng giác. Hệ số $B_n$ là hệ số của vận tốc, nên hệ số đứng trước $\sin\omega_nt$ phải là $B_n/\omega_n$.

\textbf{Bước 4 — Ghép nghiệm và xét mode hằng.}


\[
u=\sum_n\left(A_n\cos\omega_nt+\frac{B_n}{\omega_n}\sin\omega_nt\right)X_n.
\]


Với Neumann hai đầu, thêm $A_0+B_0t$, trong đó $A_0=L^{-1}\int f$, $B_0=L^{-1}\int g$. Với tần số bị thay đổi bởi hạng $u$ hoặc cản $u_t$, giải lại ODE của từng mode; không dùng nguyên tần số $ak_n$.

\textbf{Bước 5 — Kiểm tra biên, dữ liệu và hội tụ.} Mỗi mode thỏa biên; tại $t=0$ hai chuỗi cho khai triển của $f,g$. Nếu dữ liệu không tương thích ở góc, giải theo nghĩa năng lượng hoặc từng miền và nêu đường mất tính trơn; không đạo hàm một chuỗi yếu tùy ý. Chu kỳ được suy từ tất cả mode có hệ số khác không, kể cả mode hằng; hai tần số có tỷ số vô tỷ ngăn một chu kỳ chung.

\textbf{Chi tiết tính tích phân hệ số.} Với $k=n\pi/L$, hai nguyên hàm thường dùng là


\[
\int x\sin kx\,dx=-x\cos kx/k+\sin kx/k^2,\qquad
\int x\cos kx\,dx=x\sin kx/k+\cos kx/k^2.
\]


Với đa thức $P$, ghi cả hai lần tích phân:


\[
\int_0^L P\sin kx\,dx
=\left[-P\cos kx/k\right]_0^L+
\left[P'\sin kx/k^2\right]_0^L-
\frac1{k^2}\int_0^L P''\sin kx\,dx.
\]


Chỉ bỏ hạng biên sau khi thế $0,L$ và điều kiện thực sự của $P$. Nếu dữ liệu chia đoạn, áp dụng công thức trên từng đoạn, rồi cộng các hạng tại điểm chia. Cuối cùng nhân nghịch đảo chuẩn của hàm riêng; nếu hệ số là của vận tốc, chia thêm tần số khi ghép nghiệm. Với mẫu số bằng $0$ trong tích–tổng, quay lại tích phân gốc để tính riêng.
\subsection*{Lời giải áp dụng cho bài}
\textbf{Bước 1. Năng lượng của sóng có cản.} PDE là $u_{tt}+u_t=4u_{xx}$, biên cố định trên $(0,\pi)$. Với $E=\frac12\int(u_t^2+4u_x^2)$, tích phân từng phần cho $E'=4[u_tu_x]_0^\pi-\int u_t^2=-\int u_t^2\le0$. Do vậy năng lượng giảm; hiệu hai nghiệm cùng dữ liệu có năng lượng đầu $0$, chứng minh duy nhất.

\textbf{Bước 2. Tách biến và bỏ cản bằng đổi ẩn.} Mode $\sin nx$ có $T_n''+T_n'+4n^2T_n=0$. Đặt $T_n=e^{-t/2}Y_n$, được $Y_n''+(4n^2-1/4)Y_n=0$, $Y_n(0)=A_n$, $Y_n'(0)=B_n+A_n/2$. Đặt $\omega_n=\sqrt{4n^2-1/4}$:
\[T_n=e^{-t/2}[A_n\cos\omega_nt+(B_n+A_n/2)\sin\omega_nt/\omega_n].\]
Hạng $A_n/2$ cần để bù đạo hàm của $e^{-t/2}$.

\textbf{Bước 3. Dữ liệu cụ thể.} $f=\sin^3x$ cho $A_1=3/4$, $A_3=-1/4$, các hệ số khác $0$. $g=1$ cho $B_n=2(1-(-1)^n)/(\pi n)$. Nghiệm $u=\sum_{n\ge1}T_n\sin nx$ với các hệ số này.

\textbf{Bước 4. Kiểm tra.} Tại $t=0$, $T_n=A_n$, $T_n'=B_n$, đúng cả hai dữ liệu. Biên đúng theo sine và ODE chứng minh PDE. Tuy nhiên $g(0)=g(\pi)=1$ trái biên cố định, nên công thức là nghiệm năng lượng, có mặt gãy từ góc; không có nghiệm $C^1$ trên toàn bao đóng. Đẳng thức năng lượng vẫn đúng trong lớp năng lượng bằng xấp xỉ dữ liệu.

\textbf{Chi tiết tích phân của dữ liệu hằng.} Với $k_n=n\pi/L$,
\[\int_0^L\sin k_nx\,dx=\left[-\frac{\cos k_nx}{k_n}\right]_0^L=\frac{1-(-1)^n}{k_n}.\]
Nhân $2/L$ cho hệ số của hằng $1$ là $2(1-(-1)^n)/(n\pi)$; hằng $-1$ đổi dấu. Khi chiếu vận tốc, còn chia $\omega_n$ trong hệ số $\sin\omega_nt$; khi chiếu li độ, hệ số $\cos\omega_nt$ giữ nguyên.

\textbf{Chi tiết tích phân từng phần trong đạo hàm năng lượng.} Đặt $E=\frac12\int_0^L(u_t^2+a^2u_x^2+\kappa u^2)dx$. Với $u_{tt}-a^2u_{xx}+\gamma u_t+\kappa u=F$,
\[\begin{aligned}
E'&=\int_0^L(u_tu_{tt}+a^2u_xu_{xt}+\kappa uu_t)dx\\
&=\int_0^Lu_t(a^2u_{xx}-\gamma u_t-\kappa u+F)dx
+a^2\int_0^Lu_xu_{xt}dx+\kappa\int_0^Luu_tdx\\
&=a^2[u_tu_x]_0^L-a^2\int_0^Lu_{tx}u_xdx
+a^2\int_0^Lu_xu_{xt}dx-\gamma\int_0^Lu_t^2dx+\int_0^LFu_tdx\\
&=a^2[u_tu_x]_0^L-\gamma\int_0^Lu_t^2dx+\int_0^LFu_tdx.
\end{aligned}\]
Hai hạng $\kappa uu_t$ triệt tiêu ngay khi thế PDE; hai tích phân $u_xu_{xt}$ triệt tiêu sau tích phân từng phần. Với biên Dirichlet thuần nhất, đạo hàm vết theo $t$ cho $u_t=0$; với biên Neumann thuần nhất, $u_x=0$. Nếu định nghĩa năng lượng không có $1/2$, nhân cả đẳng thức đạo hàm với $2$.

\textbf{Chi tiết đổi ẩn trong ODE có cản.} Với $T=e^{-t/2}Y$,
\[T'=e^{-t/2}(Y'-Y/2),\quad T''=e^{-t/2}(Y''-Y'+Y/4).\]
Do đó $T''+T'+\Omega^2T=e^{-t/2}[Y''+(\Omega^2-1/4)Y]$. Dữ liệu $T(0)=A,T'(0)=B$ trở thành $Y(0)=A,Y'(0)=B+A/2$. Nếu $\Omega=0$, giải trực tiếp $T''+T'=0$: đặt $Z=T'$, $Z=Be^{-t}$, rồi $T=A+B(1-e^{-t})$.

\clearpage\noindent\texttt{\detokenize{VD_3_5_2_02}}\par
\section*{Lời giải bổ sung — Ví dụ 2 · Mục 3.5.2}

\subsection*{Kiến thức và trình tự giải}
\textbf{Bước 1 — Chọn năng lượng theo toán tử.} Với sóng dùng $E=\frac12\int(u_t^2+a^2u_x^2)dx$. Với Klein–Gordon có thêm $u$, thêm $\frac12\int u^2dx$. Với dầm $u_{tt}=-a^2u_{xxxx}$ dùng $E=\frac12\int(u_t^2+a^2u_{xx}^2)dx$, không thay $u_{xx}$ bằng $u_x$.

\textbf{Bước 2 — Đạo hàm và tích phân từng phần.} Với $u_{tt}-a^2u_{xx}+\gamma u_t+\kappa u=F$, năng lượng có hạng $\kappa u^2$ cho


\[
E'=a^2[u_tu_x]_0^L-\gamma\int u_t^2dx+\int Fu_tdx.
\]


Với dầm, tích phân hai lần cho $E'=a^2[u_{tx}u_{xx}-u_tu_{xxx}]_0^L$. Viết riêng từng hạng biên rồi dùng điều kiện biên và đạo hàm theo thời gian của điều kiện đó để cho chúng bằng không.

\textbf{Bước 3 — Suy hệ quả.} Không cản, không nguồn thì $E$ bảo toàn. Có cản $\gamma>0$ thì $E'\le0$; nếu $E$ không có hệ số $1/2$, $E'=-2\gamma\int u_t^2\ge-2\gamma E$, suy $e^{-2\gamma t}E(0)\le E(t)\le E(0)$. Với hiệu hai nghiệm cùng dữ liệu, $E(0)=0$ nên có duy nhất.

\textbf{Bước 4 — Tách biến nếu phải tìm nghiệm cụ thể.} Sóng có cản $\gamma=1$ cho $T''+T'+a^2\lambda T=0$. Đặt $T=e^{-t/2}Y$, thu $Y''+(a^2\lambda-1/4)Y=0$, $Y(0)=A$, $Y'(0)=B+A/2$. Phân biệt dao động, nghiệm kép và hàm hyperbolic. Dầm cho $X''''=k^4X$, $T''+a^2k^4T=0$; dùng đủ bốn điều kiện biên để tìm phương trình của $kL$.

\textbf{Bước 5 — Kiểm tra hệ số và khẳng định của đề.} Năng lượng đầu chứa cả $g^2$ và các đạo hàm của $f$. Thử $f=0$, $g\ne0$ để phát hiện công thức bỏ vận tốc. Với dầm, nghiệm $k=0$ phải giải trực tiếp bằng đa thức bậc ba; nghiệm $0$ của phương trình lượng giác–hyperbolic có thể chỉ là nghiệm giả do cơ sở suy biến.

\textbf{Chi tiết phép triệt tiêu tích phân.} Đạo hàm động năng cho $\int u_tu_{tt}$; dùng PDE thay $u_{tt}$. Với sóng, $\int u_tu_{xx}=[u_tu_x]-\int u_{tx}u_x$, và tích phân cuối triệt tiêu đạo hàm thế năng $\int u_xu_{xt}$. Với dầm,


\[
\int u_tu_{xxxx}=[u_tu_{xxx}-u_{tx}u_{xx}]+\int u_{txx}u_{xx}.
\]


Tích phân cuối triệt tiêu đạo hàm $\frac12\int u_{xx}^2$ sau khi nhân hệ số và giữ đúng dấu PDE. Viết hạng biên trước, sau đó mới dùng điều kiện của từng đầu.

\textbf{Bước 1 — Tách biến và chọn bài toán trị riêng theo biên.} Đặt $u=X(x)T(t)$ trong $u_{tt}=a^2u_{xx}$, suy $X''+\lambda X=0$, $T''+a^2\lambda T=0$. Tích phân từng phần cho $\lambda\int X^2=\int (X')^2$ khi hạng biên bằng không; do đó $\lambda\ge0$. Giải cả trường hợp $\lambda=0$ trước khi xét $\lambda=k^2>0$.

\textbf{Bước 2 — Tìm hệ hàm riêng.} Trên $(0,L)$: Dirichlet hai đầu dùng $\sin(n\pi x/L)$, $n\ge1$; Neumann hai đầu dùng $1$ và $\cos(n\pi x/L)$; Dirichlet–Neumann dùng $\sin((n+1/2)\pi x/L)$, $n\ge0$; Neumann–Dirichlet dùng cos của các tần số nửa nguyên ấy. Chuẩn bình phương của mode khác không là $L/2$, của mode hằng là $L$.

\textbf{Bước 3 — Chiếu cả hai dữ liệu.} Với mode khác không đặt


\[
A_n=\frac2L\int_0^L fX_n\,dx,\quad B_n=\frac2L\int_0^L gX_n\,dx,\quad \omega_n=ak_n.
\]


Tích phân từng khoảng nếu dữ liệu chia đoạn; dùng tích phân từng phần hoặc công thức tích–tổng khi là đa thức/lượng giác. Hệ số $B_n$ là hệ số của vận tốc, nên hệ số đứng trước $\sin\omega_nt$ phải là $B_n/\omega_n$.

\textbf{Bước 4 — Ghép nghiệm và xét mode hằng.}


\[
u=\sum_n\left(A_n\cos\omega_nt+\frac{B_n}{\omega_n}\sin\omega_nt\right)X_n.
\]


Với Neumann hai đầu, thêm $A_0+B_0t$, trong đó $A_0=L^{-1}\int f$, $B_0=L^{-1}\int g$. Với tần số bị thay đổi bởi hạng $u$ hoặc cản $u_t$, giải lại ODE của từng mode; không dùng nguyên tần số $ak_n$.

\textbf{Bước 5 — Kiểm tra biên, dữ liệu và hội tụ.} Mỗi mode thỏa biên; tại $t=0$ hai chuỗi cho khai triển của $f,g$. Nếu dữ liệu không tương thích ở góc, giải theo nghĩa năng lượng hoặc từng miền và nêu đường mất tính trơn; không đạo hàm một chuỗi yếu tùy ý. Chu kỳ được suy từ tất cả mode có hệ số khác không, kể cả mode hằng; hai tần số có tỷ số vô tỷ ngăn một chu kỳ chung.

\textbf{Chi tiết tính tích phân hệ số.} Với $k=n\pi/L$, hai nguyên hàm thường dùng là


\[
\int x\sin kx\,dx=-x\cos kx/k+\sin kx/k^2,\qquad
\int x\cos kx\,dx=x\sin kx/k+\cos kx/k^2.
\]


Với đa thức $P$, ghi cả hai lần tích phân:


\[
\int_0^L P\sin kx\,dx
=\left[-P\cos kx/k\right]_0^L+
\left[P'\sin kx/k^2\right]_0^L-
\frac1{k^2}\int_0^L P''\sin kx\,dx.
\]


Chỉ bỏ hạng biên sau khi thế $0,L$ và điều kiện thực sự của $P$. Nếu dữ liệu chia đoạn, áp dụng công thức trên từng đoạn, rồi cộng các hạng tại điểm chia. Cuối cùng nhân nghịch đảo chuẩn của hàm riêng; nếu hệ số là của vận tốc, chia thêm tần số khi ghép nghiệm. Với mẫu số bằng $0$ trong tích–tổng, quay lại tích phân gốc để tính riêng.
\subsection*{Lời giải áp dụng cho bài}
\textbf{Bước 1. Năng lượng đúng của dầm.} $u_{tt}=-a^2u_{xxxx}$, hai đầu ngàm $u=u_x=0$. Nhân với $u_t$ và tích phân hai lần:
\[E=\tfrac12\int_0^L(u_t^2+a^2u_{xx}^2)dx,\qquad E'=a^2[u_{tx}u_{xx}-u_tu_{xxx}]_0^L=0.\]
Do $u_t=u_{tx}=0$ ở hai đầu. Năng lượng có $u_x^2$ trong đề không phù hợp toán tử cấp bốn; cần sửa thành $u_{xx}^2$. Với hiệu hai nghiệm, $E=0$ suy $u_t=u_{xx}=0$; hàm affine còn lại bị biên ngàm và dữ liệu đầu buộc bằng $0$, nên duy nhất.

\textbf{Bước 2. Trị riêng dương và trực giao.} $X''''=\lambda X$, $X=X'=0$ ở hai đầu. Tích phân hai lần cho $\lambda\int X^2=\int(X'')^2$. Nếu $\lambda=0$, $X$ affine rồi biên loại hết; vậy $\lambda>0$. Trừ tích phân hai phương trình cho $(\lambda_m-\lambda_n)\int X_mX_n=0$ vì hạng biên $[X_m'''X_n-X_m''X_n'-X_n'''X_m+X_n''X_m']$ bằng $0$.

\textbf{Bước 3. Giải đủ bốn điều kiện.} Đặt $\lambda=k^4$, $\beta=kL$. Hai điều kiện trái cho $X=A(\cosh kx-\cos kx)+B(\sinh kx-\sin kx)$. Hai điều kiện phải tạo định thức $2(1-\cosh\beta\cos\beta)=0$, nên $\cosh\beta\cos\beta=1$. Chỉ lấy nghiệm dương; $\beta=0$ là nghiệm giả của cơ sở suy biến. Nghiệm đầu $\beta_1\approx4.730040745$. Chọn
\[X_n=\cosh k_nx-\cos k_nx-\sigma_n(\sinh k_nx-\sin k_nx),\quad
\sigma_n=\frac{\cosh\beta_n-\cos\beta_n}{\sinh\beta_n-\sin\beta_n}.\]
Phương trình đặc trưng bảo đảm cả $X_n(L)=X_n'(L)=0$.

\textbf{Bước 4. Chiếu dữ liệu và ghép.} $A_n=\int fX_n/\int X_n^2$, $B_n=\int gX_n/\int X_n^2$, tần số $\omega_n=ak_n^2$. Nghiệm $u=\sum(A_n\cos\omega_nt+B_n\sin\omega_nt/\omega_n)X_n$. Kiểm tra bốn biên bằng hàm riêng, hai dữ liệu bằng hệ số; không dùng tần số $ak_n$ của phương trình sóng cấp hai. Với dữ liệu đủ trơn và tương thích thu nghiệm cổ điển; với dữ liệu năng lượng hiểu chuỗi trong không gian thích hợp.

\textbf{Chi tiết hai tích phân năng lượng và định thức.} Nhân PDE với $u_t$:
\[\begin{aligned}
\int_0^Lu_tu_{xxxx}dx
&=[u_tu_{xxx}]_0^L-\int_0^Lu_{tx}u_{xxx}dx\\
&=[u_tu_{xxx}-u_{tx}u_{xx}]_0^L+\int_0^Lu_{txx}u_{xx}dx.
\end{aligned}\]
Hạng tích phân cuối bằng $\frac12(d/dt)\int u_{xx}^2$; chuyển sang cùng động năng cho $E'$ đã viết.

Đặt $C=\cosh\beta,c=\cos\beta,S=\sinh\beta,s=\sin\beta$. Hai điều kiện phải, sau chia đạo hàm cho $k$, là
\[\begin{pmatrix}C-c&S-s\\S+s&C-c\end{pmatrix}\binom AB=\binom00.\]
Định thức $(C-c)^2-(S-s)(S+s)=C^2-2Cc+c^2-S^2+s^2=2-2Cc$, vì $C^2-S^2=c^2+s^2=1$. Dòng đầu cho $B/A=-(C-c)/(S-s)$; dòng hai đúng khi định thức bằng $0$. Đây là nguồn gốc của $\sigma_n$.

\clearpage\noindent\texttt{\detokenize{BT_3_5_3_01}}\par
\section*{Lời giải bổ sung — Bài tập 1 · Mục 3.5.3}

\subsection*{Kiến thức và trình tự giải}
\textbf{Bước 1 — Chọn năng lượng theo toán tử.} Với sóng dùng $E=\frac12\int(u_t^2+a^2u_x^2)dx$. Với Klein–Gordon có thêm $u$, thêm $\frac12\int u^2dx$. Với dầm $u_{tt}=-a^2u_{xxxx}$ dùng $E=\frac12\int(u_t^2+a^2u_{xx}^2)dx$, không thay $u_{xx}$ bằng $u_x$.

\textbf{Bước 2 — Đạo hàm và tích phân từng phần.} Với $u_{tt}-a^2u_{xx}+\gamma u_t+\kappa u=F$, năng lượng có hạng $\kappa u^2$ cho


\[
E'=a^2[u_tu_x]_0^L-\gamma\int u_t^2dx+\int Fu_tdx.
\]


Với dầm, tích phân hai lần cho $E'=a^2[u_{tx}u_{xx}-u_tu_{xxx}]_0^L$. Viết riêng từng hạng biên rồi dùng điều kiện biên và đạo hàm theo thời gian của điều kiện đó để cho chúng bằng không.

\textbf{Bước 3 — Suy hệ quả.} Không cản, không nguồn thì $E$ bảo toàn. Có cản $\gamma>0$ thì $E'\le0$; nếu $E$ không có hệ số $1/2$, $E'=-2\gamma\int u_t^2\ge-2\gamma E$, suy $e^{-2\gamma t}E(0)\le E(t)\le E(0)$. Với hiệu hai nghiệm cùng dữ liệu, $E(0)=0$ nên có duy nhất.

\textbf{Bước 4 — Tách biến nếu phải tìm nghiệm cụ thể.} Sóng có cản $\gamma=1$ cho $T''+T'+a^2\lambda T=0$. Đặt $T=e^{-t/2}Y$, thu $Y''+(a^2\lambda-1/4)Y=0$, $Y(0)=A$, $Y'(0)=B+A/2$. Phân biệt dao động, nghiệm kép và hàm hyperbolic. Dầm cho $X''''=k^4X$, $T''+a^2k^4T=0$; dùng đủ bốn điều kiện biên để tìm phương trình của $kL$.

\textbf{Bước 5 — Kiểm tra hệ số và khẳng định của đề.} Năng lượng đầu chứa cả $g^2$ và các đạo hàm của $f$. Thử $f=0$, $g\ne0$ để phát hiện công thức bỏ vận tốc. Với dầm, nghiệm $k=0$ phải giải trực tiếp bằng đa thức bậc ba; nghiệm $0$ của phương trình lượng giác–hyperbolic có thể chỉ là nghiệm giả do cơ sở suy biến.

\textbf{Chi tiết phép triệt tiêu tích phân.} Đạo hàm động năng cho $\int u_tu_{tt}$; dùng PDE thay $u_{tt}$. Với sóng, $\int u_tu_{xx}=[u_tu_x]-\int u_{tx}u_x$, và tích phân cuối triệt tiêu đạo hàm thế năng $\int u_xu_{xt}$. Với dầm,


\[
\int u_tu_{xxxx}=[u_tu_{xxx}-u_{tx}u_{xx}]+\int u_{txx}u_{xx}.
\]


Tích phân cuối triệt tiêu đạo hàm $\frac12\int u_{xx}^2$ sau khi nhân hệ số và giữ đúng dấu PDE. Viết hạng biên trước, sau đó mới dùng điều kiện của từng đầu.

\textbf{Bước 1 — Tách biến và chọn bài toán trị riêng theo biên.} Đặt $u=X(x)T(t)$ trong $u_{tt}=a^2u_{xx}$, suy $X''+\lambda X=0$, $T''+a^2\lambda T=0$. Tích phân từng phần cho $\lambda\int X^2=\int (X')^2$ khi hạng biên bằng không; do đó $\lambda\ge0$. Giải cả trường hợp $\lambda=0$ trước khi xét $\lambda=k^2>0$.

\textbf{Bước 2 — Tìm hệ hàm riêng.} Trên $(0,L)$: Dirichlet hai đầu dùng $\sin(n\pi x/L)$, $n\ge1$; Neumann hai đầu dùng $1$ và $\cos(n\pi x/L)$; Dirichlet–Neumann dùng $\sin((n+1/2)\pi x/L)$, $n\ge0$; Neumann–Dirichlet dùng cos của các tần số nửa nguyên ấy. Chuẩn bình phương của mode khác không là $L/2$, của mode hằng là $L$.

\textbf{Bước 3 — Chiếu cả hai dữ liệu.} Với mode khác không đặt


\[
A_n=\frac2L\int_0^L fX_n\,dx,\quad B_n=\frac2L\int_0^L gX_n\,dx,\quad \omega_n=ak_n.
\]


Tích phân từng khoảng nếu dữ liệu chia đoạn; dùng tích phân từng phần hoặc công thức tích–tổng khi là đa thức/lượng giác. Hệ số $B_n$ là hệ số của vận tốc, nên hệ số đứng trước $\sin\omega_nt$ phải là $B_n/\omega_n$.

\textbf{Bước 4 — Ghép nghiệm và xét mode hằng.}


\[
u=\sum_n\left(A_n\cos\omega_nt+\frac{B_n}{\omega_n}\sin\omega_nt\right)X_n.
\]


Với Neumann hai đầu, thêm $A_0+B_0t$, trong đó $A_0=L^{-1}\int f$, $B_0=L^{-1}\int g$. Với tần số bị thay đổi bởi hạng $u$ hoặc cản $u_t$, giải lại ODE của từng mode; không dùng nguyên tần số $ak_n$.

\textbf{Bước 5 — Kiểm tra biên, dữ liệu và hội tụ.} Mỗi mode thỏa biên; tại $t=0$ hai chuỗi cho khai triển của $f,g$. Nếu dữ liệu không tương thích ở góc, giải theo nghĩa năng lượng hoặc từng miền và nêu đường mất tính trơn; không đạo hàm một chuỗi yếu tùy ý. Chu kỳ được suy từ tất cả mode có hệ số khác không, kể cả mode hằng; hai tần số có tỷ số vô tỷ ngăn một chu kỳ chung.

\textbf{Chi tiết tính tích phân hệ số.} Với $k=n\pi/L$, hai nguyên hàm thường dùng là


\[
\int x\sin kx\,dx=-x\cos kx/k+\sin kx/k^2,\qquad
\int x\cos kx\,dx=x\sin kx/k+\cos kx/k^2.
\]


Với đa thức $P$, ghi cả hai lần tích phân:


\[
\int_0^L P\sin kx\,dx
=\left[-P\cos kx/k\right]_0^L+
\left[P'\sin kx/k^2\right]_0^L-
\frac1{k^2}\int_0^L P''\sin kx\,dx.
\]


Chỉ bỏ hạng biên sau khi thế $0,L$ và điều kiện thực sự của $P$. Nếu dữ liệu chia đoạn, áp dụng công thức trên từng đoạn, rồi cộng các hạng tại điểm chia. Cuối cùng nhân nghịch đảo chuẩn của hàm riêng; nếu hệ số là của vận tốc, chia thêm tần số khi ghép nghiệm. Với mẫu số bằng $0$ trong tích–tổng, quay lại tích phân gốc để tính riêng.
\subsection*{Lời giải áp dụng cho bài}
\textbf{Bước 1. Sửa khẳng định bảo toàn và chứng minh duy nhất.} Phương trình (3.5.46) là $u_{tt}+u_t=a^2u_{xx}$. Với $I=\frac12\int(u_t^2+a^2u_x^2)$, $I'=a^2[u_tu_x]_0^L-\int u_t^2$. Ở biên Dirichlet, $u_t=0$; ở biên Neumann, $u_x=0$. Cả ba kiểu biên đều cho $I'=-\int u_t^2$, không bảo toàn nói chung. Hiệu hai nghiệm có $I(0)=0$, suy $u_t=u_x=0$, rồi dữ liệu đầu buộc hằng số bằng $0$.

\textbf{Bước 2. Ba bài trị riêng.} Tách $X''+\lambda X=0$. Tích phân cho $\lambda\int X^2=\int(X')^2$; hai kiểu hỗn hợp loại mode hằng, còn Neumann hai đầu giữ nó. Trừ hai phương trình chứng minh trực giao như mục 3.3. Cụ thể:

(a) $X(0)=X'(L)=0$: $X_n=\sin k_nx$, $k_n=(n+1/2)\pi/L$, $n\ge0$.

(b) $X'(0)=X(L)=0$: cùng $k_n$, $X_n=\cos k_nx$.

(c) $X'(0)=X'(L)=0$: $X_0=1$, $X_n=\cos(n\pi x/L)$, $n\ge1$. Chuẩn là $L$ cho mode $0$, $L/2$ cho mode khác không.

\textbf{Bước 3. Giải ODE và chiếu cả hai dữ liệu.} Đặt $A_n=\int fX_n/\int X_n^2$, $B_n=\int gX_n/\int X_n^2$. $T_n''+T_n'+a^2\lambda_nT_n=0$. Đổi $T_n=e^{-t/2}Y_n$ cho $Y_n''+(a^2\lambda_n-1/4)Y_n=0$. Nếu $\omega_n^2=a^2\lambda_n-1/4>0$,
\[T_n=e^{-t/2}[A_n\cos\omega_nt+(B_n+A_n/2)\sin\omega_nt/\omega_n].\]
Nếu $\omega_n=0$, dùng $e^{-t/2}[A_n+(B_n+A_n/2)t]$. Nếu $\omega_n^2<0$, thay cos/sin bằng cosh/sinh với $\sqrt{-\omega_n^2}$. Mode hằng là $A_0+B_0(1-e^{-t})$, giải trực tiếp $T_0''+T_0'=0$.

\textbf{Bước 4. Nghiệm và ví dụ cho cả ba kiểu.} $u=\sum T_nX_n$. Với $a=L=1$, lấy lần lượt $f=\sin(\pi x/2)$, $f=\cos(\pi x/2)$, $f=\cos\pi x$, tất cả $g=0$. Chỉ có một mode, $A=1$, $B=0$, nên nghiệm là $X(x)e^{-t/2}(\cos\omega t+\sin\omega t/(2\omega))$, với $\omega=\sqrt{\pi^2/4-1/4}$ cho (a),(b), $\sqrt{\pi^2-1/4}$ cho (c). Đồ thị thể hiện giảm biên độ; ví dụ mode hằng $f=1,g=0$ ở (c) cho nghiệm hằng $1$, phù hợp vì năng lượng đạo hàm bằng $0$.

\textbf{Chi tiết tích phân từng phần trong đạo hàm năng lượng.} Đặt $E=\frac12\int_0^L(u_t^2+a^2u_x^2+\kappa u^2)dx$. Với $u_{tt}-a^2u_{xx}+\gamma u_t+\kappa u=F$,
\[\begin{aligned}
E'&=\int_0^L(u_tu_{tt}+a^2u_xu_{xt}+\kappa uu_t)dx\\
&=\int_0^Lu_t(a^2u_{xx}-\gamma u_t-\kappa u+F)dx
+a^2\int_0^Lu_xu_{xt}dx+\kappa\int_0^Luu_tdx\\
&=a^2[u_tu_x]_0^L-a^2\int_0^Lu_{tx}u_xdx
+a^2\int_0^Lu_xu_{xt}dx-\gamma\int_0^Lu_t^2dx+\int_0^LFu_tdx\\
&=a^2[u_tu_x]_0^L-\gamma\int_0^Lu_t^2dx+\int_0^LFu_tdx.
\end{aligned}\]
Hai hạng $\kappa uu_t$ triệt tiêu ngay khi thế PDE; hai tích phân $u_xu_{xt}$ triệt tiêu sau tích phân từng phần. Với biên Dirichlet thuần nhất, đạo hàm vết theo $t$ cho $u_t=0$; với biên Neumann thuần nhất, $u_x=0$. Nếu định nghĩa năng lượng không có $1/2$, nhân cả đẳng thức đạo hàm với $2$.

\textbf{Chi tiết trực giao và dấu trị riêng.} Nhân $-X''=\lambda X$ bởi $X$, tích phân:
\[\lambda\int_0^LX^2dx=-[X'X]_0^L+\int_0^L(X')^2dx=\int_0^L(X')^2dx.\]
Hạng biên hết vì mỗi đầu có $X=0$ hoặc $X'=0$. Nếu $\lambda=0$, $\int(X')^2=0$ cho $X$ hằng; khi có một đầu Dirichlet, hằng đó bằng $0$. Với hai trị riêng,
\[(\lambda_m-\lambda_n)\int_0^LX_mX_n\,dx
=\int_0^L(-X_m''X_n+X_n''X_m)dx
=[-X_m'X_n+X_n'X_m]_0^L=0.\]
Nếu $\lambda_m\ne\lambda_n$, chia hiệu trị riêng cho tích vô hướng bằng $0$.

\textbf{Chi tiết đổi ẩn trong ODE có cản.} Với $T=e^{-t/2}Y$,
\[T'=e^{-t/2}(Y'-Y/2),\quad T''=e^{-t/2}(Y''-Y'+Y/4).\]
Do đó $T''+T'+\Omega^2T=e^{-t/2}[Y''+(\Omega^2-1/4)Y]$. Dữ liệu $T(0)=A,T'(0)=B$ trở thành $Y(0)=A,Y'(0)=B+A/2$. Nếu $\Omega=0$, giải trực tiếp $T''+T'=0$: đặt $Z=T'$, $Z=Be^{-t}$, rồi $T=A+B(1-e^{-t})$.
\begin{center}\includegraphics[width=.97\linewidth,height=11cm,keepaspectratio]{_xu_ly/graphs/BT_3_5_3_01.png}\end{center}
\textit{Đồ thị minh họa; với nghiệm số, xem phương pháp và tham số trong mã Python kèm theo.}

\clearpage\noindent\texttt{\detokenize{BT_3_5_3_02}}\par
\section*{Lời giải bổ sung — Bài tập 2 · Mục 3.5.3}

\subsection*{Kiến thức và trình tự giải}
\textbf{Bước 1 — Chọn năng lượng theo toán tử.} Với sóng dùng $E=\frac12\int(u_t^2+a^2u_x^2)dx$. Với Klein–Gordon có thêm $u$, thêm $\frac12\int u^2dx$. Với dầm $u_{tt}=-a^2u_{xxxx}$ dùng $E=\frac12\int(u_t^2+a^2u_{xx}^2)dx$, không thay $u_{xx}$ bằng $u_x$.

\textbf{Bước 2 — Đạo hàm và tích phân từng phần.} Với $u_{tt}-a^2u_{xx}+\gamma u_t+\kappa u=F$, năng lượng có hạng $\kappa u^2$ cho


\[
E'=a^2[u_tu_x]_0^L-\gamma\int u_t^2dx+\int Fu_tdx.
\]


Với dầm, tích phân hai lần cho $E'=a^2[u_{tx}u_{xx}-u_tu_{xxx}]_0^L$. Viết riêng từng hạng biên rồi dùng điều kiện biên và đạo hàm theo thời gian của điều kiện đó để cho chúng bằng không.

\textbf{Bước 3 — Suy hệ quả.} Không cản, không nguồn thì $E$ bảo toàn. Có cản $\gamma>0$ thì $E'\le0$; nếu $E$ không có hệ số $1/2$, $E'=-2\gamma\int u_t^2\ge-2\gamma E$, suy $e^{-2\gamma t}E(0)\le E(t)\le E(0)$. Với hiệu hai nghiệm cùng dữ liệu, $E(0)=0$ nên có duy nhất.

\textbf{Bước 4 — Tách biến nếu phải tìm nghiệm cụ thể.} Sóng có cản $\gamma=1$ cho $T''+T'+a^2\lambda T=0$. Đặt $T=e^{-t/2}Y$, thu $Y''+(a^2\lambda-1/4)Y=0$, $Y(0)=A$, $Y'(0)=B+A/2$. Phân biệt dao động, nghiệm kép và hàm hyperbolic. Dầm cho $X''''=k^4X$, $T''+a^2k^4T=0$; dùng đủ bốn điều kiện biên để tìm phương trình của $kL$.

\textbf{Bước 5 — Kiểm tra hệ số và khẳng định của đề.} Năng lượng đầu chứa cả $g^2$ và các đạo hàm của $f$. Thử $f=0$, $g\ne0$ để phát hiện công thức bỏ vận tốc. Với dầm, nghiệm $k=0$ phải giải trực tiếp bằng đa thức bậc ba; nghiệm $0$ của phương trình lượng giác–hyperbolic có thể chỉ là nghiệm giả do cơ sở suy biến.

\textbf{Chi tiết phép triệt tiêu tích phân.} Đạo hàm động năng cho $\int u_tu_{tt}$; dùng PDE thay $u_{tt}$. Với sóng, $\int u_tu_{xx}=[u_tu_x]-\int u_{tx}u_x$, và tích phân cuối triệt tiêu đạo hàm thế năng $\int u_xu_{xt}$. Với dầm,


\[
\int u_tu_{xxxx}=[u_tu_{xxx}-u_{tx}u_{xx}]+\int u_{txx}u_{xx}.
\]


Tích phân cuối triệt tiêu đạo hàm $\frac12\int u_{xx}^2$ sau khi nhân hệ số và giữ đúng dấu PDE. Viết hạng biên trước, sau đó mới dùng điều kiện của từng đầu.

\textbf{Bước 1 — Tách biến và chọn bài toán trị riêng theo biên.} Đặt $u=X(x)T(t)$ trong $u_{tt}=a^2u_{xx}$, suy $X''+\lambda X=0$, $T''+a^2\lambda T=0$. Tích phân từng phần cho $\lambda\int X^2=\int (X')^2$ khi hạng biên bằng không; do đó $\lambda\ge0$. Giải cả trường hợp $\lambda=0$ trước khi xét $\lambda=k^2>0$.

\textbf{Bước 2 — Tìm hệ hàm riêng.} Trên $(0,L)$: Dirichlet hai đầu dùng $\sin(n\pi x/L)$, $n\ge1$; Neumann hai đầu dùng $1$ và $\cos(n\pi x/L)$; Dirichlet–Neumann dùng $\sin((n+1/2)\pi x/L)$, $n\ge0$; Neumann–Dirichlet dùng cos của các tần số nửa nguyên ấy. Chuẩn bình phương của mode khác không là $L/2$, của mode hằng là $L$.

\textbf{Bước 3 — Chiếu cả hai dữ liệu.} Với mode khác không đặt


\[
A_n=\frac2L\int_0^L fX_n\,dx,\quad B_n=\frac2L\int_0^L gX_n\,dx,\quad \omega_n=ak_n.
\]


Tích phân từng khoảng nếu dữ liệu chia đoạn; dùng tích phân từng phần hoặc công thức tích–tổng khi là đa thức/lượng giác. Hệ số $B_n$ là hệ số của vận tốc, nên hệ số đứng trước $\sin\omega_nt$ phải là $B_n/\omega_n$.

\textbf{Bước 4 — Ghép nghiệm và xét mode hằng.}


\[
u=\sum_n\left(A_n\cos\omega_nt+\frac{B_n}{\omega_n}\sin\omega_nt\right)X_n.
\]


Với Neumann hai đầu, thêm $A_0+B_0t$, trong đó $A_0=L^{-1}\int f$, $B_0=L^{-1}\int g$. Với tần số bị thay đổi bởi hạng $u$ hoặc cản $u_t$, giải lại ODE của từng mode; không dùng nguyên tần số $ak_n$.

\textbf{Bước 5 — Kiểm tra biên, dữ liệu và hội tụ.} Mỗi mode thỏa biên; tại $t=0$ hai chuỗi cho khai triển của $f,g$. Nếu dữ liệu không tương thích ở góc, giải theo nghĩa năng lượng hoặc từng miền và nêu đường mất tính trơn; không đạo hàm một chuỗi yếu tùy ý. Chu kỳ được suy từ tất cả mode có hệ số khác không, kể cả mode hằng; hai tần số có tỷ số vô tỷ ngăn một chu kỳ chung.

\textbf{Chi tiết tính tích phân hệ số.} Với $k=n\pi/L$, hai nguyên hàm thường dùng là


\[
\int x\sin kx\,dx=-x\cos kx/k+\sin kx/k^2,\qquad
\int x\cos kx\,dx=x\sin kx/k+\cos kx/k^2.
\]


Với đa thức $P$, ghi cả hai lần tích phân:


\[
\int_0^L P\sin kx\,dx
=\left[-P\cos kx/k\right]_0^L+
\left[P'\sin kx/k^2\right]_0^L-
\frac1{k^2}\int_0^L P''\sin kx\,dx.
\]


Chỉ bỏ hạng biên sau khi thế $0,L$ và điều kiện thực sự của $P$. Nếu dữ liệu chia đoạn, áp dụng công thức trên từng đoạn, rồi cộng các hạng tại điểm chia. Cuối cùng nhân nghịch đảo chuẩn của hàm riêng; nếu hệ số là của vận tốc, chia thêm tần số khi ghép nghiệm. Với mẫu số bằng $0$ trong tích–tổng, quay lại tích phân gốc để tính riêng.
\subsection*{Lời giải áp dụng cho bài}
\textbf{Bước 1. Năng lượng và tính duy nhất.} Dầm $u_{tt}=-a^2u_{xxxx}$ dùng $E=\frac12\int(u_t^2+a^2u_{xx}^2)$, không phải năng lượng chứa $u_x^2$. Tích phân hai lần cho $E'=a^2[u_{tx}u_{xx}-u_tu_{xxx}]_0^L$. Mỗi cặp điều kiện trong đề làm cả hai tích biên bằng $0$: ngàm có $u_t=u_{tx}=0$; tựa có $u_t=u_{xx}=0$; đầu trượt có $u_{tx}=u_{xxx}=0$. Do đó $E$ bảo toàn. Hiệu hai nghiệm có $E=0$, $u_t=u_{xx}=0$, rồi dữ liệu đầu loại hàm affine còn lại, chứng minh duy nhất.

\textbf{Bước 2. Đọc đúng năm kiểu biên của bản scan.} (a) ngàm trái, tựa phải: $X(0)=X'(0)=X(L)=X''(L)=0$. (b) ngàm trái, trượt phải: $X(0)=X'(0)=X'(L)=X'''(L)=0$. Hai kiểu này khác hai đầu ngàm/ngàm–tự do của lời giải cũ.

\textbf{Bước 3. Giải phương trình đặc trưng.} Đặt $X''''=k^4X$, $\beta=kL$. Với (a),(b), hai điều kiện trái cho $X=A(\cosh kx-\cos kx)+B(\sinh kx-\sin kx)$.

(a) Từ $X(L)=X''(L)=0$, cộng/trừ suy $A\cosh\beta+B\sinh\beta=0$, $A\cos\beta+B\sin\beta=0$. Định thức cho $\cosh\beta\sin\beta-\sinh\beta\cos\beta=0$, hay $\tan\beta=\tanh\beta$. Chọn
\[X=\cosh kx-\cos kx-\coth\beta(\sinh kx-\sin kx).\]
Nghiệm dương đầu $\beta_1\approx3.926602312$.

(b) Từ $X'(L)=X'''(L)=0$, suy $A\sinh\beta+B\cosh\beta=0$, $A\sin\beta-B\cos\beta=0$. Định thức cho $\sinh\beta\cos\beta+\cosh\beta\sin\beta=0$, hay $\tan\beta=-\tanh\beta$. Chọn
\[X=\cosh kx-\cos kx-\tanh\beta(\sinh kx-\sin kx),\]
nghiệm dương đầu $\beta_1\approx2.365020372$.

(c) Tựa hai đầu $X=X''=0$: điều kiện trái loại cos/cosh, điều kiện phải loại sinh, nên $k_n=n\pi/L$, $X_n=\sin k_nx$, $n\ge1$.

(d) Tựa trái, trượt phải $X(0)=X''(0)=X'(L)=X'''(L)=0$: còn sine và sinh; cộng/trừ hai điều kiện phải loại sinh rồi cho $\cos kL=0$. $k_n=(n+1/2)\pi/L$, $X_n=\sin k_nx$, $n\ge0$.

(e) $X'=X'''=0$ hai đầu: $X_n=\cos(n\pi x/L)$, $n\ge1$, cộng mode $X_0=1$, $\lambda_0=0$.

\textbf{Bước 4. Dương, trực giao và nghiệm.} Tích phân hai lần cho $\lambda\int X^2=\int(X'')^2\ge0$ dưới cả năm kiểu biên. Nếu $\lambda=0$, $X$ affine; chỉ (e) giữ hằng. Các nghiệm $\beta=0$ trong (a),(b) phải loại bằng kiểm tra đa thức, không coi là mode. Trừ hai phương trình tích phân chứng minh trực giao với các trị riêng khác nhau. Lấy $A_n=\int fX_n/\int X_n^2$, $B_n=\int gX_n/\int X_n^2$:
\[u=\sum_n[A_n\cos(ak_n^2t)+B_n\sin(ak_n^2t)/(ak_n^2)]X_n.\]
Ở (e) thêm $\bar f+t\bar g$. Với $a=L=1$, mỗi kiểu lấy $f$ là hàm riêng dương đầu vừa tìm, $g=0$, cho ví dụ $u=X_1\cos(k_1^2t)$; đồ thị kèm đủ năm kiểu biên và không dùng nhầm các phương trình $\cos\beta\cosh\beta=\pm1$.

\textbf{Chi tiết giải đủ các điều kiện của (a),(b).} Viết $C=\cosh\beta,c=\cos\beta,S=\sinh\beta,s=\sin\beta$.
Ở (a), $X(L)=A(C-c)+B(S-s)=0$, $X''(L)/k^2=A(C+c)+B(S+s)=0$. Cộng hai dòng rồi chia $2$ cho $AC+BS=0$; lấy dòng sau trừ dòng trước rồi chia $2$ cho $Ac+Bs=0$. Định thức của hệ này là $Cs-Sc=0$, và $B/A=-C/S=-\coth\beta$.

Ở (b), $X'(L)/k=A(S+s)+B(C-c)=0$, $X'''(L)/k^3=A(S-s)+B(C+c)=0$. Cộng cho $AS+BC=0$; lấy dòng đầu trừ dòng sau cho $As-Bc=0$. Định thức là $-Sc-Cs=0$, và $B/A=-S/C=-\tanh\beta$. Với $\beta>0$, $\sinh\beta,\cosh\beta$ khác $0$; $\cos\beta=0$ không thỏa các định thức, nên các phép chia để viết tan hợp lệ tại nghiệm.

\textbf{Chi tiết ba kiểu còn lại.} Từ $X=A\cosh kx+B\sinh kx+C\cos kx+D\sin kx$, ở (c),(d), $X(0)=A+C=0$, $X''(0)/k^2=A-C=0$ cho $A=C=0$. (c) Hai dòng phải $B\sinh\beta+D\sin\beta=0$, $B\sinh\beta-D\sin\beta=0$ cho $B=0$, $\sin\beta=0$. (d) Hai dòng phải $B\cosh\beta+D\cos\beta=0$, $B\cosh\beta-D\cos\beta=0$ cho $B=0$, $\cos\beta=0$.

Ở (e), $X'(0)/k=B+D=0$, $X'''(0)/k^3=B-D=0$ cho $B=D=0$. Hai dòng phải $A\sinh\beta-C\sin\beta=0$, $A\sinh\beta+C\sin\beta=0$ cho $A=0$, $\sin\beta=0$. Riêng $\lambda=0$, $X''''=0$ cho $X=A+Bx+Cx^2+Dx^3$; $X'''=0$ cho $D=0$, $X'(0)=0$ cho $B=0$, $X'(L)=0$ cho $C=0$, còn hằng $A$.

\textbf{Chi tiết tích phân chứng minh trực giao dầm.} Tích phân từng phần hai lần:
\[\int_0^L(X_m''''X_n-X_n''''X_m)dx=
[X_m'''X_n-X_m''X_n'-X_n'''X_m+X_n''X_m']_0^L.\]
Hai tích phân còn lại là $\int X_m''X_n''$ và $\int X_n''X_m''$, triệt tiêu khi trừ. Tại đầu ngàm $X=X'=0$, đầu tựa $X=X''=0$, đầu trượt $X'=X'''=0$, mỗi tích biên mất một thừa số. Thay phương trình hàm riêng cho $(\lambda_m-\lambda_n)\int X_mX_n=0$, rồi chia hiệu trị riêng.

\textbf{Chi tiết hai lần tích phân từng phần của dầm.}
\[\int_0^L u_tu_{xxxx}dx
=[u_tu_{xxx}]_0^L-\int_0^Lu_{tx}u_{xxx}dx
=[u_tu_{xxx}-u_{tx}u_{xx}]_0^L+\int_0^Lu_{txx}u_{xx}dx.\]
Thế $u_{tt}=-a^2u_{xxxx}$ vào đạo hàm động năng. Tích phân cuối cùng triệt tiêu với đạo hàm $\frac{a^2}{2}\int u_{xx}^2$, còn $E'=a^2[u_{tx}u_{xx}-u_tu_{xxx}]_0^L$.
\begin{center}\includegraphics[width=.97\linewidth,height=11cm,keepaspectratio]{_xu_ly/graphs/BT_3_5_3_02.png}\end{center}
\textit{Đồ thị minh họa; với nghiệm số, xem phương pháp và tham số trong mã Python kèm theo.}

\clearpage\noindent\texttt{\detokenize{BT_3_5_3_03}}\par
\section*{Lời giải bổ sung — Bài tập 3 · Mục 3.5.3}

\subsection*{Kiến thức và trình tự giải}
\textbf{Bước 1 — Chọn năng lượng theo toán tử.} Với sóng dùng $E=\frac12\int(u_t^2+a^2u_x^2)dx$. Với Klein–Gordon có thêm $u$, thêm $\frac12\int u^2dx$. Với dầm $u_{tt}=-a^2u_{xxxx}$ dùng $E=\frac12\int(u_t^2+a^2u_{xx}^2)dx$, không thay $u_{xx}$ bằng $u_x$.

\textbf{Bước 2 — Đạo hàm và tích phân từng phần.} Với $u_{tt}-a^2u_{xx}+\gamma u_t+\kappa u=F$, năng lượng có hạng $\kappa u^2$ cho


\[
E'=a^2[u_tu_x]_0^L-\gamma\int u_t^2dx+\int Fu_tdx.
\]


Với dầm, tích phân hai lần cho $E'=a^2[u_{tx}u_{xx}-u_tu_{xxx}]_0^L$. Viết riêng từng hạng biên rồi dùng điều kiện biên và đạo hàm theo thời gian của điều kiện đó để cho chúng bằng không.

\textbf{Bước 3 — Suy hệ quả.} Không cản, không nguồn thì $E$ bảo toàn. Có cản $\gamma>0$ thì $E'\le0$; nếu $E$ không có hệ số $1/2$, $E'=-2\gamma\int u_t^2\ge-2\gamma E$, suy $e^{-2\gamma t}E(0)\le E(t)\le E(0)$. Với hiệu hai nghiệm cùng dữ liệu, $E(0)=0$ nên có duy nhất.

\textbf{Bước 4 — Tách biến nếu phải tìm nghiệm cụ thể.} Sóng có cản $\gamma=1$ cho $T''+T'+a^2\lambda T=0$. Đặt $T=e^{-t/2}Y$, thu $Y''+(a^2\lambda-1/4)Y=0$, $Y(0)=A$, $Y'(0)=B+A/2$. Phân biệt dao động, nghiệm kép và hàm hyperbolic. Dầm cho $X''''=k^4X$, $T''+a^2k^4T=0$; dùng đủ bốn điều kiện biên để tìm phương trình của $kL$.

\textbf{Bước 5 — Kiểm tra hệ số và khẳng định của đề.} Năng lượng đầu chứa cả $g^2$ và các đạo hàm của $f$. Thử $f=0$, $g\ne0$ để phát hiện công thức bỏ vận tốc. Với dầm, nghiệm $k=0$ phải giải trực tiếp bằng đa thức bậc ba; nghiệm $0$ của phương trình lượng giác–hyperbolic có thể chỉ là nghiệm giả do cơ sở suy biến.

\textbf{Chi tiết phép triệt tiêu tích phân.} Đạo hàm động năng cho $\int u_tu_{tt}$; dùng PDE thay $u_{tt}$. Với sóng, $\int u_tu_{xx}=[u_tu_x]-\int u_{tx}u_x$, và tích phân cuối triệt tiêu đạo hàm thế năng $\int u_xu_{xt}$. Với dầm,


\[
\int u_tu_{xxxx}=[u_tu_{xxx}-u_{tx}u_{xx}]+\int u_{txx}u_{xx}.
\]


Tích phân cuối triệt tiêu đạo hàm $\frac12\int u_{xx}^2$ sau khi nhân hệ số và giữ đúng dấu PDE. Viết hạng biên trước, sau đó mới dùng điều kiện của từng đầu.

\textbf{Bước 1 — Tách biến và chọn bài toán trị riêng theo biên.} Đặt $u=X(x)T(t)$ trong $u_{tt}=a^2u_{xx}$, suy $X''+\lambda X=0$, $T''+a^2\lambda T=0$. Tích phân từng phần cho $\lambda\int X^2=\int (X')^2$ khi hạng biên bằng không; do đó $\lambda\ge0$. Giải cả trường hợp $\lambda=0$ trước khi xét $\lambda=k^2>0$.

\textbf{Bước 2 — Tìm hệ hàm riêng.} Trên $(0,L)$: Dirichlet hai đầu dùng $\sin(n\pi x/L)$, $n\ge1$; Neumann hai đầu dùng $1$ và $\cos(n\pi x/L)$; Dirichlet–Neumann dùng $\sin((n+1/2)\pi x/L)$, $n\ge0$; Neumann–Dirichlet dùng cos của các tần số nửa nguyên ấy. Chuẩn bình phương của mode khác không là $L/2$, của mode hằng là $L$.

\textbf{Bước 3 — Chiếu cả hai dữ liệu.} Với mode khác không đặt


\[
A_n=\frac2L\int_0^L fX_n\,dx,\quad B_n=\frac2L\int_0^L gX_n\,dx,\quad \omega_n=ak_n.
\]


Tích phân từng khoảng nếu dữ liệu chia đoạn; dùng tích phân từng phần hoặc công thức tích–tổng khi là đa thức/lượng giác. Hệ số $B_n$ là hệ số của vận tốc, nên hệ số đứng trước $\sin\omega_nt$ phải là $B_n/\omega_n$.

\textbf{Bước 4 — Ghép nghiệm và xét mode hằng.}


\[
u=\sum_n\left(A_n\cos\omega_nt+\frac{B_n}{\omega_n}\sin\omega_nt\right)X_n.
\]


Với Neumann hai đầu, thêm $A_0+B_0t$, trong đó $A_0=L^{-1}\int f$, $B_0=L^{-1}\int g$. Với tần số bị thay đổi bởi hạng $u$ hoặc cản $u_t$, giải lại ODE của từng mode; không dùng nguyên tần số $ak_n$.

\textbf{Bước 5 — Kiểm tra biên, dữ liệu và hội tụ.} Mỗi mode thỏa biên; tại $t=0$ hai chuỗi cho khai triển của $f,g$. Nếu dữ liệu không tương thích ở góc, giải theo nghĩa năng lượng hoặc từng miền và nêu đường mất tính trơn; không đạo hàm một chuỗi yếu tùy ý. Chu kỳ được suy từ tất cả mode có hệ số khác không, kể cả mode hằng; hai tần số có tỷ số vô tỷ ngăn một chu kỳ chung.

\textbf{Chi tiết tính tích phân hệ số.} Với $k=n\pi/L$, hai nguyên hàm thường dùng là


\[
\int x\sin kx\,dx=-x\cos kx/k+\sin kx/k^2,\qquad
\int x\cos kx\,dx=x\sin kx/k+\cos kx/k^2.
\]


Với đa thức $P$, ghi cả hai lần tích phân:


\[
\int_0^L P\sin kx\,dx
=\left[-P\cos kx/k\right]_0^L+
\left[P'\sin kx/k^2\right]_0^L-
\frac1{k^2}\int_0^L P''\sin kx\,dx.
\]


Chỉ bỏ hạng biên sau khi thế $0,L$ và điều kiện thực sự của $P$. Nếu dữ liệu chia đoạn, áp dụng công thức trên từng đoạn, rồi cộng các hạng tại điểm chia. Cuối cùng nhân nghịch đảo chuẩn của hàm riêng; nếu hệ số là của vận tốc, chia thêm tần số khi ghép nghiệm. Với mẫu số bằng $0$ trong tích–tổng, quay lại tích phân gốc để tính riêng.
\subsection*{Lời giải áp dụng cho bài}
\textbf{Bước 1. Hệ sine và dữ liệu tam giác.} PDE $u_{tt}=u_{xx}-u$ trên $(0,\pi)$, biên không. $f=\min(x,\pi-x)$, $g=0$ có $A_n=4\sin(n\pi/2)/(\pi n^2)$, tính bằng tách tích phân tại $\pi/2$ như Bài 3.3.3.2.

\textbf{Bước 2. Giải đúng tần số mới.} $T_n''+(n^2+1)T_n=0$, $T_n(0)=A_n$, $T_n'(0)=0$. Do đó
\[u=\sum_{n\ge1}\frac{4\sin(n\pi/2)}{\pi n^2}\cos(\sqrt{n^2+1}\,t)\sin nx.\]
Hạng $+1$ dưới căn đến từ $-u$ trong PDE, không thể dùng tần số $n$.

\textbf{Bước 3. Không có chu kỳ chung.} Hệ số mode $1,3$ đều khác không. Nếu chu kỳ $T>0$ tồn tại, chiếu lên hai mode cho $\sqrt2T$ và $\sqrt{10}T$ đều là bội của $2\pi$, buộc tỷ số $\sqrt5$ hữu tỷ, mâu thuẫn. Vậy nghiệm này không tuần hoàn.

\textbf{Bước 4. Kiểm tra và mức trơn.} Mỗi mode thỏa PDE/biên và đúng dữ liệu. Tam giác có góc nên nghiệm thuộc lớp năng lượng, không cổ điển qua mặt truyền góc. Năng lượng bảo toàn phải chứa $u^2$: $\frac12\int(u_t^2+u_x^2+u^2)$.

\textbf{Chi tiết hệ số của tam giác.} Với $f(x)=\min(x,\pi-x)$, chia tại $b=\pi/2$. Đổi $y=\pi-x$ ở nửa sau cho
\[\int_0^\pi f(x)\sin nx\,dx=(1+(-1)^{n+1})\int_0^b x\sin nx\,dx.\]
Tích phân cuối bằng $[-x\cos nx/n+\sin nx/n^2]_0^b=-b\cos(nb)/n+\sin(nb)/n^2$. Nếu $n$ chẵn, hệ số trước tích phân bằng $0$; nếu $n$ lẻ, $\cos(nb)=0$. Nhân $2/\pi$ được $4\sin(n\pi/2)/(\pi n^2)$.

\clearpage\noindent\texttt{\detokenize{BT_3_5_3_04}}\par
\section*{Lời giải bổ sung — Bài tập 4 · Mục 3.5.3}

\subsection*{Kiến thức và trình tự giải}
\textbf{Bước 1 — Chọn năng lượng theo toán tử.} Với sóng dùng $E=\frac12\int(u_t^2+a^2u_x^2)dx$. Với Klein–Gordon có thêm $u$, thêm $\frac12\int u^2dx$. Với dầm $u_{tt}=-a^2u_{xxxx}$ dùng $E=\frac12\int(u_t^2+a^2u_{xx}^2)dx$, không thay $u_{xx}$ bằng $u_x$.

\textbf{Bước 2 — Đạo hàm và tích phân từng phần.} Với $u_{tt}-a^2u_{xx}+\gamma u_t+\kappa u=F$, năng lượng có hạng $\kappa u^2$ cho


\[
E'=a^2[u_tu_x]_0^L-\gamma\int u_t^2dx+\int Fu_tdx.
\]


Với dầm, tích phân hai lần cho $E'=a^2[u_{tx}u_{xx}-u_tu_{xxx}]_0^L$. Viết riêng từng hạng biên rồi dùng điều kiện biên và đạo hàm theo thời gian của điều kiện đó để cho chúng bằng không.

\textbf{Bước 3 — Suy hệ quả.} Không cản, không nguồn thì $E$ bảo toàn. Có cản $\gamma>0$ thì $E'\le0$; nếu $E$ không có hệ số $1/2$, $E'=-2\gamma\int u_t^2\ge-2\gamma E$, suy $e^{-2\gamma t}E(0)\le E(t)\le E(0)$. Với hiệu hai nghiệm cùng dữ liệu, $E(0)=0$ nên có duy nhất.

\textbf{Bước 4 — Tách biến nếu phải tìm nghiệm cụ thể.} Sóng có cản $\gamma=1$ cho $T''+T'+a^2\lambda T=0$. Đặt $T=e^{-t/2}Y$, thu $Y''+(a^2\lambda-1/4)Y=0$, $Y(0)=A$, $Y'(0)=B+A/2$. Phân biệt dao động, nghiệm kép và hàm hyperbolic. Dầm cho $X''''=k^4X$, $T''+a^2k^4T=0$; dùng đủ bốn điều kiện biên để tìm phương trình của $kL$.

\textbf{Bước 5 — Kiểm tra hệ số và khẳng định của đề.} Năng lượng đầu chứa cả $g^2$ và các đạo hàm của $f$. Thử $f=0$, $g\ne0$ để phát hiện công thức bỏ vận tốc. Với dầm, nghiệm $k=0$ phải giải trực tiếp bằng đa thức bậc ba; nghiệm $0$ của phương trình lượng giác–hyperbolic có thể chỉ là nghiệm giả do cơ sở suy biến.

\textbf{Chi tiết phép triệt tiêu tích phân.} Đạo hàm động năng cho $\int u_tu_{tt}$; dùng PDE thay $u_{tt}$. Với sóng, $\int u_tu_{xx}=[u_tu_x]-\int u_{tx}u_x$, và tích phân cuối triệt tiêu đạo hàm thế năng $\int u_xu_{xt}$. Với dầm,


\[
\int u_tu_{xxxx}=[u_tu_{xxx}-u_{tx}u_{xx}]+\int u_{txx}u_{xx}.
\]


Tích phân cuối triệt tiêu đạo hàm $\frac12\int u_{xx}^2$ sau khi nhân hệ số và giữ đúng dấu PDE. Viết hạng biên trước, sau đó mới dùng điều kiện của từng đầu.

\textbf{Bước 1 — Tách biến và chọn bài toán trị riêng theo biên.} Đặt $u=X(x)T(t)$ trong $u_{tt}=a^2u_{xx}$, suy $X''+\lambda X=0$, $T''+a^2\lambda T=0$. Tích phân từng phần cho $\lambda\int X^2=\int (X')^2$ khi hạng biên bằng không; do đó $\lambda\ge0$. Giải cả trường hợp $\lambda=0$ trước khi xét $\lambda=k^2>0$.

\textbf{Bước 2 — Tìm hệ hàm riêng.} Trên $(0,L)$: Dirichlet hai đầu dùng $\sin(n\pi x/L)$, $n\ge1$; Neumann hai đầu dùng $1$ và $\cos(n\pi x/L)$; Dirichlet–Neumann dùng $\sin((n+1/2)\pi x/L)$, $n\ge0$; Neumann–Dirichlet dùng cos của các tần số nửa nguyên ấy. Chuẩn bình phương của mode khác không là $L/2$, của mode hằng là $L$.

\textbf{Bước 3 — Chiếu cả hai dữ liệu.} Với mode khác không đặt


\[
A_n=\frac2L\int_0^L fX_n\,dx,\quad B_n=\frac2L\int_0^L gX_n\,dx,\quad \omega_n=ak_n.
\]


Tích phân từng khoảng nếu dữ liệu chia đoạn; dùng tích phân từng phần hoặc công thức tích–tổng khi là đa thức/lượng giác. Hệ số $B_n$ là hệ số của vận tốc, nên hệ số đứng trước $\sin\omega_nt$ phải là $B_n/\omega_n$.

\textbf{Bước 4 — Ghép nghiệm và xét mode hằng.}


\[
u=\sum_n\left(A_n\cos\omega_nt+\frac{B_n}{\omega_n}\sin\omega_nt\right)X_n.
\]


Với Neumann hai đầu, thêm $A_0+B_0t$, trong đó $A_0=L^{-1}\int f$, $B_0=L^{-1}\int g$. Với tần số bị thay đổi bởi hạng $u$ hoặc cản $u_t$, giải lại ODE của từng mode; không dùng nguyên tần số $ak_n$.

\textbf{Bước 5 — Kiểm tra biên, dữ liệu và hội tụ.} Mỗi mode thỏa biên; tại $t=0$ hai chuỗi cho khai triển của $f,g$. Nếu dữ liệu không tương thích ở góc, giải theo nghĩa năng lượng hoặc từng miền và nêu đường mất tính trơn; không đạo hàm một chuỗi yếu tùy ý. Chu kỳ được suy từ tất cả mode có hệ số khác không, kể cả mode hằng; hai tần số có tỷ số vô tỷ ngăn một chu kỳ chung.

\textbf{Chi tiết tính tích phân hệ số.} Với $k=n\pi/L$, hai nguyên hàm thường dùng là


\[
\int x\sin kx\,dx=-x\cos kx/k+\sin kx/k^2,\qquad
\int x\cos kx\,dx=x\sin kx/k+\cos kx/k^2.
\]


Với đa thức $P$, ghi cả hai lần tích phân:


\[
\int_0^L P\sin kx\,dx
=\left[-P\cos kx/k\right]_0^L+
\left[P'\sin kx/k^2\right]_0^L-
\frac1{k^2}\int_0^L P''\sin kx\,dx.
\]


Chỉ bỏ hạng biên sau khi thế $0,L$ và điều kiện thực sự của $P$. Nếu dữ liệu chia đoạn, áp dụng công thức trên từng đoạn, rồi cộng các hạng tại điểm chia. Cuối cùng nhân nghịch đảo chuẩn của hàm riêng; nếu hệ số là của vận tốc, chia thêm tần số khi ghép nghiệm. Với mẫu số bằng $0$ trong tích–tổng, quay lại tích phân gốc để tính riêng.
\subsection*{Lời giải áp dụng cho bài}
\textbf{Bước 1. Tách phương trình điện báo.} Với $u_{tt}+u_t+u=u_{xx}$, hai biên cố định trên $(0,1)$, đặt $u=\sum T_n\sin(n\pi x)$. ODE là $T_n''+T_n'+(n^2\pi^2+1)T_n=0$.

\textbf{Bước 2. Đổi ẩn và dữ liệu.} $A_n=2\int_0^1f(x)\sin(n\pi x)dx$, vận tốc đầu bằng $0$. Đặt $T_n=e^{-t/2}Y_n$ cho $Y_n''+(n^2\pi^2+3/4)Y_n=0$, $Y_n(0)=A_n$, $Y_n'(0)=A_n/2$.

\textbf{Bước 3. Nghiệm.}
\[u=e^{-t/2}\sum_{n\ge1}A_n\left[\cos\omega_nt+\frac{\sin\omega_nt}{2\omega_n}\right]\sin(n\pi x),\qquad\omega_n=\sqrt{n^2\pi^2+3/4}.\]
\textbf{Bước 4. Kiểm tra.} Tại $0$, đạo hàm của ngoặc là $A_n/2$, triệt tiêu đạo hàm $-A_n/2$ của hệ số mũ; vì thế vận tốc đầu đúng bằng $0$. Biên theo sine và PDE theo ODE. Năng lượng có thêm $u^2$, đạo hàm bằng $-\int u_t^2$, chứng minh duy nhất trong lớp thích hợp. Cần $f$ tương thích biên nếu đòi nghiệm trơn đến góc; công thức hệ số áp dụng cho dữ liệu năng lượng.

\textbf{Chi tiết tích phân từng phần trong đạo hàm năng lượng.} Đặt $E=\frac12\int_0^L(u_t^2+a^2u_x^2+\kappa u^2)dx$. Với $u_{tt}-a^2u_{xx}+\gamma u_t+\kappa u=F$,
\[\begin{aligned}
E'&=\int_0^L(u_tu_{tt}+a^2u_xu_{xt}+\kappa uu_t)dx\\
&=\int_0^Lu_t(a^2u_{xx}-\gamma u_t-\kappa u+F)dx
+a^2\int_0^Lu_xu_{xt}dx+\kappa\int_0^Luu_tdx\\
&=a^2[u_tu_x]_0^L-a^2\int_0^Lu_{tx}u_xdx
+a^2\int_0^Lu_xu_{xt}dx-\gamma\int_0^Lu_t^2dx+\int_0^LFu_tdx\\
&=a^2[u_tu_x]_0^L-\gamma\int_0^Lu_t^2dx+\int_0^LFu_tdx.
\end{aligned}\]
Hai hạng $\kappa uu_t$ triệt tiêu ngay khi thế PDE; hai tích phân $u_xu_{xt}$ triệt tiêu sau tích phân từng phần. Với biên Dirichlet thuần nhất, đạo hàm vết theo $t$ cho $u_t=0$; với biên Neumann thuần nhất, $u_x=0$. Nếu định nghĩa năng lượng không có $1/2$, nhân cả đẳng thức đạo hàm với $2$.

\textbf{Chi tiết đổi ẩn trong ODE có cản.} Với $T=e^{-t/2}Y$,
\[T'=e^{-t/2}(Y'-Y/2),\quad T''=e^{-t/2}(Y''-Y'+Y/4).\]
Do đó $T''+T'+\Omega^2T=e^{-t/2}[Y''+(\Omega^2-1/4)Y]$. Dữ liệu $T(0)=A,T'(0)=B$ trở thành $Y(0)=A,Y'(0)=B+A/2$. Nếu $\Omega=0$, giải trực tiếp $T''+T'=0$: đặt $Z=T'$, $Z=Be^{-t}$, rồi $T=A+B(1-e^{-t})$.

\clearpage\noindent\texttt{\detokenize{BT_3_5_3_05}}\par
\section*{Lời giải bổ sung — Bài tập 5 · Mục 3.5.3}

\subsection*{Kiến thức và trình tự giải}
\textbf{Bước 1 — Chọn năng lượng theo toán tử.} Với sóng dùng $E=\frac12\int(u_t^2+a^2u_x^2)dx$. Với Klein–Gordon có thêm $u$, thêm $\frac12\int u^2dx$. Với dầm $u_{tt}=-a^2u_{xxxx}$ dùng $E=\frac12\int(u_t^2+a^2u_{xx}^2)dx$, không thay $u_{xx}$ bằng $u_x$.

\textbf{Bước 2 — Đạo hàm và tích phân từng phần.} Với $u_{tt}-a^2u_{xx}+\gamma u_t+\kappa u=F$, năng lượng có hạng $\kappa u^2$ cho


\[
E'=a^2[u_tu_x]_0^L-\gamma\int u_t^2dx+\int Fu_tdx.
\]


Với dầm, tích phân hai lần cho $E'=a^2[u_{tx}u_{xx}-u_tu_{xxx}]_0^L$. Viết riêng từng hạng biên rồi dùng điều kiện biên và đạo hàm theo thời gian của điều kiện đó để cho chúng bằng không.

\textbf{Bước 3 — Suy hệ quả.} Không cản, không nguồn thì $E$ bảo toàn. Có cản $\gamma>0$ thì $E'\le0$; nếu $E$ không có hệ số $1/2$, $E'=-2\gamma\int u_t^2\ge-2\gamma E$, suy $e^{-2\gamma t}E(0)\le E(t)\le E(0)$. Với hiệu hai nghiệm cùng dữ liệu, $E(0)=0$ nên có duy nhất.

\textbf{Bước 4 — Tách biến nếu phải tìm nghiệm cụ thể.} Sóng có cản $\gamma=1$ cho $T''+T'+a^2\lambda T=0$. Đặt $T=e^{-t/2}Y$, thu $Y''+(a^2\lambda-1/4)Y=0$, $Y(0)=A$, $Y'(0)=B+A/2$. Phân biệt dao động, nghiệm kép và hàm hyperbolic. Dầm cho $X''''=k^4X$, $T''+a^2k^4T=0$; dùng đủ bốn điều kiện biên để tìm phương trình của $kL$.

\textbf{Bước 5 — Kiểm tra hệ số và khẳng định của đề.} Năng lượng đầu chứa cả $g^2$ và các đạo hàm của $f$. Thử $f=0$, $g\ne0$ để phát hiện công thức bỏ vận tốc. Với dầm, nghiệm $k=0$ phải giải trực tiếp bằng đa thức bậc ba; nghiệm $0$ của phương trình lượng giác–hyperbolic có thể chỉ là nghiệm giả do cơ sở suy biến.

\textbf{Chi tiết phép triệt tiêu tích phân.} Đạo hàm động năng cho $\int u_tu_{tt}$; dùng PDE thay $u_{tt}$. Với sóng, $\int u_tu_{xx}=[u_tu_x]-\int u_{tx}u_x$, và tích phân cuối triệt tiêu đạo hàm thế năng $\int u_xu_{xt}$. Với dầm,


\[
\int u_tu_{xxxx}=[u_tu_{xxx}-u_{tx}u_{xx}]+\int u_{txx}u_{xx}.
\]


Tích phân cuối triệt tiêu đạo hàm $\frac12\int u_{xx}^2$ sau khi nhân hệ số và giữ đúng dấu PDE. Viết hạng biên trước, sau đó mới dùng điều kiện của từng đầu.

\textbf{Bước 1 — Tách biến và chọn bài toán trị riêng theo biên.} Đặt $u=X(x)T(t)$ trong $u_{tt}=a^2u_{xx}$, suy $X''+\lambda X=0$, $T''+a^2\lambda T=0$. Tích phân từng phần cho $\lambda\int X^2=\int (X')^2$ khi hạng biên bằng không; do đó $\lambda\ge0$. Giải cả trường hợp $\lambda=0$ trước khi xét $\lambda=k^2>0$.

\textbf{Bước 2 — Tìm hệ hàm riêng.} Trên $(0,L)$: Dirichlet hai đầu dùng $\sin(n\pi x/L)$, $n\ge1$; Neumann hai đầu dùng $1$ và $\cos(n\pi x/L)$; Dirichlet–Neumann dùng $\sin((n+1/2)\pi x/L)$, $n\ge0$; Neumann–Dirichlet dùng cos của các tần số nửa nguyên ấy. Chuẩn bình phương của mode khác không là $L/2$, của mode hằng là $L$.

\textbf{Bước 3 — Chiếu cả hai dữ liệu.} Với mode khác không đặt


\[
A_n=\frac2L\int_0^L fX_n\,dx,\quad B_n=\frac2L\int_0^L gX_n\,dx,\quad \omega_n=ak_n.
\]


Tích phân từng khoảng nếu dữ liệu chia đoạn; dùng tích phân từng phần hoặc công thức tích–tổng khi là đa thức/lượng giác. Hệ số $B_n$ là hệ số của vận tốc, nên hệ số đứng trước $\sin\omega_nt$ phải là $B_n/\omega_n$.

\textbf{Bước 4 — Ghép nghiệm và xét mode hằng.}


\[
u=\sum_n\left(A_n\cos\omega_nt+\frac{B_n}{\omega_n}\sin\omega_nt\right)X_n.
\]


Với Neumann hai đầu, thêm $A_0+B_0t$, trong đó $A_0=L^{-1}\int f$, $B_0=L^{-1}\int g$. Với tần số bị thay đổi bởi hạng $u$ hoặc cản $u_t$, giải lại ODE của từng mode; không dùng nguyên tần số $ak_n$.

\textbf{Bước 5 — Kiểm tra biên, dữ liệu và hội tụ.} Mỗi mode thỏa biên; tại $t=0$ hai chuỗi cho khai triển của $f,g$. Nếu dữ liệu không tương thích ở góc, giải theo nghĩa năng lượng hoặc từng miền và nêu đường mất tính trơn; không đạo hàm một chuỗi yếu tùy ý. Chu kỳ được suy từ tất cả mode có hệ số khác không, kể cả mode hằng; hai tần số có tỷ số vô tỷ ngăn một chu kỳ chung.

\textbf{Chi tiết tính tích phân hệ số.} Với $k=n\pi/L$, hai nguyên hàm thường dùng là


\[
\int x\sin kx\,dx=-x\cos kx/k+\sin kx/k^2,\qquad
\int x\cos kx\,dx=x\sin kx/k+\cos kx/k^2.
\]


Với đa thức $P$, ghi cả hai lần tích phân:


\[
\int_0^L P\sin kx\,dx
=\left[-P\cos kx/k\right]_0^L+
\left[P'\sin kx/k^2\right]_0^L-
\frac1{k^2}\int_0^L P''\sin kx\,dx.
\]


Chỉ bỏ hạng biên sau khi thế $0,L$ và điều kiện thực sự của $P$. Nếu dữ liệu chia đoạn, áp dụng công thức trên từng đoạn, rồi cộng các hạng tại điểm chia. Cuối cùng nhân nghịch đảo chuẩn của hàm riêng; nếu hệ số là của vận tốc, chia thêm tần số khi ghép nghiệm. Với mẫu số bằng $0$ trong tích–tổng, quay lại tích phân gốc để tính riêng.
\subsection*{Lời giải áp dụng cho bài}
\textbf{Câu (i): tìm một hàm riêng của toán tử có đạo hàm cấp một.} Đặt $f=e^{-x}(\pi\cos\pi x+\sin\pi x)$. Đạo hàm trực tiếp cho $f'=-(\pi^2+1)e^{-x}\sin\pi x$, nên biên $f'(0)=f'(1)=0$. Đạo hàm thêm cho $f''+2f'=-(\pi^2+1)f$. Vì vậy $u=fT$ dẫn tới $T''+(\pi^2+1)T=0$, $T(0)=1,T'(0)=0$:
\[u=f(x)\cos(\sqrt{\pi^2+1}\,t).\]
Các tính toán trên kiểm tra đồng thời PDE, biên Neumann và dữ liệu đầu.

\textbf{Câu (ii), bước 1: đổi tọa độ và khử đạo hàm cấp một.} $z=\log(1+x)$, $\ell=\log2$. Vì $\partial_x=e^{-z}\partial_z$, $\partial_x((1+x)^2u_x)=u_{zz}+u_z$. Đặt $u=e^{-z/2}v$; tính $u_{zz}+u_z=e^{-z/2}(v_{zz}-v/4)$. Nguồn $t/\sqrt{1+x}=te^{-z/2}$, nên $v_{tt}=v_{zz}-v/4+t$, biên Dirichlet ở $z=0,\ell$.

\textbf{Câu (ii), bước 2: chiếu và giải nguồn tuyến tính.} $v_0=e^{z/2}(e^z-1)(2-e^z)$, $v_{t,0}=0$. Với $k_n=n\pi/\ell$, $\omega_n=\sqrt{k_n^2+1/4}$,
\[A_n=\frac2\ell\int_0^\ell e^{z/2}(e^z-1)(2-e^z)\sin k_nz\,dz,\quad c_n=\frac{2(1-(-1)^n)}{n\pi}.\]
ODE $T_n''+\omega_n^2T_n=c_nt$ cho nghiệm riêng dữ liệu không $c_n(t/\omega_n^2-\sin\omega_nt/\omega_n^3)$. Vì vậy
\[u=e^{-z/2}\sum_{n\ge1}\left[A_n\cos\omega_nt+c_n\left(\frac t{\omega_n^2}-\frac{\sin\omega_nt}{\omega_n^3}\right)\right]\sin k_nz.\]
Tích phân của $A_n$ là tích phân sơ cấp: khai triển $v_0=-e^{5z/2}+3e^{3z/2}-2e^{z/2}$ rồi dùng
\[\int_0^\ell e^{pz}\sin k_nz\,dz=\frac{k_n(1-(-1)^ne^{p\ell})}{p^2+k_n^2}.\]
Thay $p=5/2,3/2,1/2$ cùng các hệ số $-1,3,-2$ cho hệ số tường minh.

\textbf{Câu (iii), bước 1: chiếu nguồn và vận tốc.} PDE $u_{tt}+2u_t=u_{xx}+\sin x$, biên không trên $(0,1)$, $f=0,g=1$. Với $k_n=n\pi$, $\omega_n=\sqrt{k_n^2-1}$, vận tốc có $c_n=2(1-(-1)^n)/(n\pi)$. Tích–tổng cho nguồn
\[F_n=2\int_0^1\sin x\sin k_nx\,dx=\frac{2(-1)^{n+1}k_n\sin1}{k_n^2-1}.\]
\textbf{Câu (iii), bước 2: giải ODE có cản.} $T_n''+2T_n'+k_n^2T_n=F_n$, $T_n(0)=0$, $T_n'(0)=c_n$. Nghiệm riêng tĩnh $F_n/k_n^2$ được bù hai dữ liệu đầu bởi phần thuần nhất:
\[u=\sum_{n\ge1}\left[c_ne^{-t}\frac{\sin\omega_nt}{\omega_n}+\frac{F_n}{k_n^2}\left\{1-e^{-t}\left(\cos\omega_nt+\frac{\sin\omega_nt}{\omega_n}\right)\right\}\right]\sin k_nx.\]
Thế $t=0$ kiểm tra hai dữ liệu; từng ODE kiểm tra PDE. $g=1$ không tương thích biên cố định ở góc, nên câu (iii) không có nghiệm cổ điển toàn bao đóng; nguồn cũng có vết không phù hợp nếu đòi cấp hai đến biên.

\textbf{Chi tiết đạo hàm ở (i).} Với $f=e^{-x}(\pi\cos\pi x+\sin\pi x)$, quy tắc tích cho
\[f'=e^{-x}[-\pi\cos\pi x-\sin\pi x-\pi^2\sin\pi x+\pi\cos\pi x]
=-(\pi^2+1)e^{-x}\sin\pi x.\]
Đạo hàm tiếp: $f''=(\pi^2+1)e^{-x}(\sin\pi x-\pi\cos\pi x)$. Vì $2f'=-2(\pi^2+1)e^{-x}\sin\pi x$, cộng cho $f''+2f'=-(\pi^2+1)e^{-x}(\sin\pi x+\pi\cos\pi x)=-(\pi^2+1)f$.

\textbf{Chi tiết đổi ẩn và tính hệ số ở (ii).} $u=e^{-z/2}v$ cho $u_z=e^{-z/2}(v_z-v/2)$, $u_{zz}=e^{-z/2}(v_{zz}-v_z+v/4)$. Cộng $u_z$ được $e^{-z/2}(v_{zz}-v/4)$. Vì $(e^z-1)(2-e^z)=-e^{2z}+3e^z-2$, dữ liệu mới là $-e^{5z/2}+3e^{3z/2}-2e^{z/2}$. Nguyên hàm
\[\int e^{pz}\sin k_nz\,dz=
\frac{e^{pz}(p\sin k_nz-k_n\cos k_nz)}{p^2+k_n^2}\]
được kiểm tra bằng đạo hàm tích. Tại $0,\ell$, $\sin k_nz=0$, $\cos k_n\ell=(-1)^n$, nên
\[A_n=\frac{2k_n}{\ell}\left[
-\frac{1-(-1)^n2^{5/2}}{k_n^2+25/4}
+3\frac{1-(-1)^n2^{3/2}}{k_n^2+9/4}
-2\frac{1-(-1)^n2^{1/2}}{k_n^2+1/4}\right].\]
Với nguồn $c_nt$, nghiệm riêng $c_nt/\omega_n^2$ có vận tốc đầu $c_n/\omega_n^2$; trừ $c_n\sin\omega_nt/\omega_n^3$ để bù, cho đáp ứng dữ liệu không.

\textbf{Chi tiết nguồn và hằng số của (iii).} Tích–tổng cho
\[F_n=\left[\frac{\sin((1-k_n)x)}{1-k_n}
-\frac{\sin((1+k_n)x)}{1+k_n}\right]_0^1
=(-1)^n\sin1\left(\frac1{1-k_n}-\frac1{1+k_n}\right).\]
Rút gọn cho $F_n=2(-1)^{n+1}k_n\sin1/(k_n^2-1)$. Đặt $H_n=F_n/k_n^2$; tổng nghiệm ODE là $T_n=H_n+e^{-t}(A\cos\omega_nt+B\sin\omega_nt)$. $T_n(0)=0$ cho $A=-H_n$; $T_n'(0)=c_n$ cho $-A+\omega_nB=c_n$, nên $B=(c_n-H_n)/\omega_n$. Thế lại và gom hai nhóm cho công thức ở bước 2.

\clearpage\noindent\texttt{\detokenize{BT_3_5_3_06}}\par
\section*{Lời giải bổ sung — Bài tập 6 · Mục 3.5.3}

\subsection*{Kiến thức và trình tự giải}
\textbf{Bước 1 — Chọn năng lượng theo toán tử.} Với sóng dùng $E=\frac12\int(u_t^2+a^2u_x^2)dx$. Với Klein–Gordon có thêm $u$, thêm $\frac12\int u^2dx$. Với dầm $u_{tt}=-a^2u_{xxxx}$ dùng $E=\frac12\int(u_t^2+a^2u_{xx}^2)dx$, không thay $u_{xx}$ bằng $u_x$.

\textbf{Bước 2 — Đạo hàm và tích phân từng phần.} Với $u_{tt}-a^2u_{xx}+\gamma u_t+\kappa u=F$, năng lượng có hạng $\kappa u^2$ cho


\[
E'=a^2[u_tu_x]_0^L-\gamma\int u_t^2dx+\int Fu_tdx.
\]


Với dầm, tích phân hai lần cho $E'=a^2[u_{tx}u_{xx}-u_tu_{xxx}]_0^L$. Viết riêng từng hạng biên rồi dùng điều kiện biên và đạo hàm theo thời gian của điều kiện đó để cho chúng bằng không.

\textbf{Bước 3 — Suy hệ quả.} Không cản, không nguồn thì $E$ bảo toàn. Có cản $\gamma>0$ thì $E'\le0$; nếu $E$ không có hệ số $1/2$, $E'=-2\gamma\int u_t^2\ge-2\gamma E$, suy $e^{-2\gamma t}E(0)\le E(t)\le E(0)$. Với hiệu hai nghiệm cùng dữ liệu, $E(0)=0$ nên có duy nhất.

\textbf{Bước 4 — Tách biến nếu phải tìm nghiệm cụ thể.} Sóng có cản $\gamma=1$ cho $T''+T'+a^2\lambda T=0$. Đặt $T=e^{-t/2}Y$, thu $Y''+(a^2\lambda-1/4)Y=0$, $Y(0)=A$, $Y'(0)=B+A/2$. Phân biệt dao động, nghiệm kép và hàm hyperbolic. Dầm cho $X''''=k^4X$, $T''+a^2k^4T=0$; dùng đủ bốn điều kiện biên để tìm phương trình của $kL$.

\textbf{Bước 5 — Kiểm tra hệ số và khẳng định của đề.} Năng lượng đầu chứa cả $g^2$ và các đạo hàm của $f$. Thử $f=0$, $g\ne0$ để phát hiện công thức bỏ vận tốc. Với dầm, nghiệm $k=0$ phải giải trực tiếp bằng đa thức bậc ba; nghiệm $0$ của phương trình lượng giác–hyperbolic có thể chỉ là nghiệm giả do cơ sở suy biến.

\textbf{Chi tiết phép triệt tiêu tích phân.} Đạo hàm động năng cho $\int u_tu_{tt}$; dùng PDE thay $u_{tt}$. Với sóng, $\int u_tu_{xx}=[u_tu_x]-\int u_{tx}u_x$, và tích phân cuối triệt tiêu đạo hàm thế năng $\int u_xu_{xt}$. Với dầm,


\[
\int u_tu_{xxxx}=[u_tu_{xxx}-u_{tx}u_{xx}]+\int u_{txx}u_{xx}.
\]


Tích phân cuối triệt tiêu đạo hàm $\frac12\int u_{xx}^2$ sau khi nhân hệ số và giữ đúng dấu PDE. Viết hạng biên trước, sau đó mới dùng điều kiện của từng đầu.

\textbf{Bước 1 — Tách biến và chọn bài toán trị riêng theo biên.} Đặt $u=X(x)T(t)$ trong $u_{tt}=a^2u_{xx}$, suy $X''+\lambda X=0$, $T''+a^2\lambda T=0$. Tích phân từng phần cho $\lambda\int X^2=\int (X')^2$ khi hạng biên bằng không; do đó $\lambda\ge0$. Giải cả trường hợp $\lambda=0$ trước khi xét $\lambda=k^2>0$.

\textbf{Bước 2 — Tìm hệ hàm riêng.} Trên $(0,L)$: Dirichlet hai đầu dùng $\sin(n\pi x/L)$, $n\ge1$; Neumann hai đầu dùng $1$ và $\cos(n\pi x/L)$; Dirichlet–Neumann dùng $\sin((n+1/2)\pi x/L)$, $n\ge0$; Neumann–Dirichlet dùng cos của các tần số nửa nguyên ấy. Chuẩn bình phương của mode khác không là $L/2$, của mode hằng là $L$.

\textbf{Bước 3 — Chiếu cả hai dữ liệu.} Với mode khác không đặt


\[
A_n=\frac2L\int_0^L fX_n\,dx,\quad B_n=\frac2L\int_0^L gX_n\,dx,\quad \omega_n=ak_n.
\]


Tích phân từng khoảng nếu dữ liệu chia đoạn; dùng tích phân từng phần hoặc công thức tích–tổng khi là đa thức/lượng giác. Hệ số $B_n$ là hệ số của vận tốc, nên hệ số đứng trước $\sin\omega_nt$ phải là $B_n/\omega_n$.

\textbf{Bước 4 — Ghép nghiệm và xét mode hằng.}


\[
u=\sum_n\left(A_n\cos\omega_nt+\frac{B_n}{\omega_n}\sin\omega_nt\right)X_n.
\]


Với Neumann hai đầu, thêm $A_0+B_0t$, trong đó $A_0=L^{-1}\int f$, $B_0=L^{-1}\int g$. Với tần số bị thay đổi bởi hạng $u$ hoặc cản $u_t$, giải lại ODE của từng mode; không dùng nguyên tần số $ak_n$.

\textbf{Bước 5 — Kiểm tra biên, dữ liệu và hội tụ.} Mỗi mode thỏa biên; tại $t=0$ hai chuỗi cho khai triển của $f,g$. Nếu dữ liệu không tương thích ở góc, giải theo nghĩa năng lượng hoặc từng miền và nêu đường mất tính trơn; không đạo hàm một chuỗi yếu tùy ý. Chu kỳ được suy từ tất cả mode có hệ số khác không, kể cả mode hằng; hai tần số có tỷ số vô tỷ ngăn một chu kỳ chung.

\textbf{Chi tiết tính tích phân hệ số.} Với $k=n\pi/L$, hai nguyên hàm thường dùng là


\[
\int x\sin kx\,dx=-x\cos kx/k+\sin kx/k^2,\qquad
\int x\cos kx\,dx=x\sin kx/k+\cos kx/k^2.
\]


Với đa thức $P$, ghi cả hai lần tích phân:


\[
\int_0^L P\sin kx\,dx
=\left[-P\cos kx/k\right]_0^L+
\left[P'\sin kx/k^2\right]_0^L-
\frac1{k^2}\int_0^L P''\sin kx\,dx.
\]


Chỉ bỏ hạng biên sau khi thế $0,L$ và điều kiện thực sự của $P$. Nếu dữ liệu chia đoạn, áp dụng công thức trên từng đoạn, rồi cộng các hạng tại điểm chia. Cuối cùng nhân nghịch đảo chuẩn của hàm riêng; nếu hệ số là của vận tốc, chia thêm tần số khi ghép nghiệm. Với mẫu số bằng $0$ trong tích–tổng, quay lại tích phân gốc để tính riêng.
\subsection*{Lời giải áp dụng cho bài}
\textbf{Bước 1. Đạo hàm năng lượng Klein–Gordon.} Với $u_{tt}-u_{xx}+u=0$, đặt $E=\int_0^L(u^2+u_t^2+u_x^2)dx$. Tích phân từng phần:
\[E'=2\int_0^Lu_t(u_{tt}-u_{xx}+u)dx+2[u_tu_x]_0^L=0.\]
Ở biên cố định, $u_t=0$, nên thông lượng hết.

\textbf{Bước 2. Thế dữ liệu đầu và sửa câu (i).}
\[E(t)=E(0)=\int_0^L(f^2+g^2+(f')^2)dx.\]
Công thức in thiếu $g^2$. Nó chỉ đúng nếu thêm $g=0$. Lấy $f=0$, $g=\sin(\pi x/L)$: tại $t=0$ vế trái là $L/2$, còn vế phải in bằng $0$, nên yêu cầu nguyên văn sai.

\textbf{Bước 3. Câu (ii).} Khai triển và tích phân hạng chéo:
\[\int_0^L|f+f'|^2dx=\int_0^L(f^2+(f')^2)dx+[f^2]_0^L.\]
Nếu $f(0)=f(L)$ thì hạng biên hết; để bằng năng lượng còn cần $g=0$. Dưới biên cố định tương thích, hai giá trị $f$ đều $0$, nhưng điều đó không xóa năng lượng vận tốc. Công thức đúng tổng quát là $E=\int|f+f'|^2+\int g^2$ khi hai giá trị đầu bằng nhau. Đoạn không gian ghi $(0,1)$ cùng cận $L$ cần thống nhất; chứng minh đúng trên mọi $(0,L)$.

\textbf{Chi tiết tích phân từng phần trong đạo hàm năng lượng.} Đặt $E=\frac12\int_0^L(u_t^2+a^2u_x^2+\kappa u^2)dx$. Với $u_{tt}-a^2u_{xx}+\gamma u_t+\kappa u=F$,
\[\begin{aligned}
E'&=\int_0^L(u_tu_{tt}+a^2u_xu_{xt}+\kappa uu_t)dx\\
&=\int_0^Lu_t(a^2u_{xx}-\gamma u_t-\kappa u+F)dx
+a^2\int_0^Lu_xu_{xt}dx+\kappa\int_0^Luu_tdx\\
&=a^2[u_tu_x]_0^L-a^2\int_0^Lu_{tx}u_xdx
+a^2\int_0^Lu_xu_{xt}dx-\gamma\int_0^Lu_t^2dx+\int_0^LFu_tdx\\
&=a^2[u_tu_x]_0^L-\gamma\int_0^Lu_t^2dx+\int_0^LFu_tdx.
\end{aligned}\]
Hai hạng $\kappa uu_t$ triệt tiêu ngay khi thế PDE; hai tích phân $u_xu_{xt}$ triệt tiêu sau tích phân từng phần. Với biên Dirichlet thuần nhất, đạo hàm vết theo $t$ cho $u_t=0$; với biên Neumann thuần nhất, $u_x=0$. Nếu định nghĩa năng lượng không có $1/2$, nhân cả đẳng thức đạo hàm với $2$.

\clearpage\noindent\texttt{\detokenize{BT_3_5_3_07}}\par
\section*{Lời giải bổ sung — Bài tập 7 · Mục 3.5.3}

\subsection*{Kiến thức và trình tự giải}
\textbf{Bước 1 — Chọn năng lượng theo toán tử.} Với sóng dùng $E=\frac12\int(u_t^2+a^2u_x^2)dx$. Với Klein–Gordon có thêm $u$, thêm $\frac12\int u^2dx$. Với dầm $u_{tt}=-a^2u_{xxxx}$ dùng $E=\frac12\int(u_t^2+a^2u_{xx}^2)dx$, không thay $u_{xx}$ bằng $u_x$.

\textbf{Bước 2 — Đạo hàm và tích phân từng phần.} Với $u_{tt}-a^2u_{xx}+\gamma u_t+\kappa u=F$, năng lượng có hạng $\kappa u^2$ cho


\[
E'=a^2[u_tu_x]_0^L-\gamma\int u_t^2dx+\int Fu_tdx.
\]


Với dầm, tích phân hai lần cho $E'=a^2[u_{tx}u_{xx}-u_tu_{xxx}]_0^L$. Viết riêng từng hạng biên rồi dùng điều kiện biên và đạo hàm theo thời gian của điều kiện đó để cho chúng bằng không.

\textbf{Bước 3 — Suy hệ quả.} Không cản, không nguồn thì $E$ bảo toàn. Có cản $\gamma>0$ thì $E'\le0$; nếu $E$ không có hệ số $1/2$, $E'=-2\gamma\int u_t^2\ge-2\gamma E$, suy $e^{-2\gamma t}E(0)\le E(t)\le E(0)$. Với hiệu hai nghiệm cùng dữ liệu, $E(0)=0$ nên có duy nhất.

\textbf{Bước 4 — Tách biến nếu phải tìm nghiệm cụ thể.} Sóng có cản $\gamma=1$ cho $T''+T'+a^2\lambda T=0$. Đặt $T=e^{-t/2}Y$, thu $Y''+(a^2\lambda-1/4)Y=0$, $Y(0)=A$, $Y'(0)=B+A/2$. Phân biệt dao động, nghiệm kép và hàm hyperbolic. Dầm cho $X''''=k^4X$, $T''+a^2k^4T=0$; dùng đủ bốn điều kiện biên để tìm phương trình của $kL$.

\textbf{Bước 5 — Kiểm tra hệ số và khẳng định của đề.} Năng lượng đầu chứa cả $g^2$ và các đạo hàm của $f$. Thử $f=0$, $g\ne0$ để phát hiện công thức bỏ vận tốc. Với dầm, nghiệm $k=0$ phải giải trực tiếp bằng đa thức bậc ba; nghiệm $0$ của phương trình lượng giác–hyperbolic có thể chỉ là nghiệm giả do cơ sở suy biến.

\textbf{Chi tiết phép triệt tiêu tích phân.} Đạo hàm động năng cho $\int u_tu_{tt}$; dùng PDE thay $u_{tt}$. Với sóng, $\int u_tu_{xx}=[u_tu_x]-\int u_{tx}u_x$, và tích phân cuối triệt tiêu đạo hàm thế năng $\int u_xu_{xt}$. Với dầm,


\[
\int u_tu_{xxxx}=[u_tu_{xxx}-u_{tx}u_{xx}]+\int u_{txx}u_{xx}.
\]


Tích phân cuối triệt tiêu đạo hàm $\frac12\int u_{xx}^2$ sau khi nhân hệ số và giữ đúng dấu PDE. Viết hạng biên trước, sau đó mới dùng điều kiện của từng đầu.

\textbf{Bước 1 — Tách biến và chọn bài toán trị riêng theo biên.} Đặt $u=X(x)T(t)$ trong $u_{tt}=a^2u_{xx}$, suy $X''+\lambda X=0$, $T''+a^2\lambda T=0$. Tích phân từng phần cho $\lambda\int X^2=\int (X')^2$ khi hạng biên bằng không; do đó $\lambda\ge0$. Giải cả trường hợp $\lambda=0$ trước khi xét $\lambda=k^2>0$.

\textbf{Bước 2 — Tìm hệ hàm riêng.} Trên $(0,L)$: Dirichlet hai đầu dùng $\sin(n\pi x/L)$, $n\ge1$; Neumann hai đầu dùng $1$ và $\cos(n\pi x/L)$; Dirichlet–Neumann dùng $\sin((n+1/2)\pi x/L)$, $n\ge0$; Neumann–Dirichlet dùng cos của các tần số nửa nguyên ấy. Chuẩn bình phương của mode khác không là $L/2$, của mode hằng là $L$.

\textbf{Bước 3 — Chiếu cả hai dữ liệu.} Với mode khác không đặt


\[
A_n=\frac2L\int_0^L fX_n\,dx,\quad B_n=\frac2L\int_0^L gX_n\,dx,\quad \omega_n=ak_n.
\]


Tích phân từng khoảng nếu dữ liệu chia đoạn; dùng tích phân từng phần hoặc công thức tích–tổng khi là đa thức/lượng giác. Hệ số $B_n$ là hệ số của vận tốc, nên hệ số đứng trước $\sin\omega_nt$ phải là $B_n/\omega_n$.

\textbf{Bước 4 — Ghép nghiệm và xét mode hằng.}


\[
u=\sum_n\left(A_n\cos\omega_nt+\frac{B_n}{\omega_n}\sin\omega_nt\right)X_n.
\]


Với Neumann hai đầu, thêm $A_0+B_0t$, trong đó $A_0=L^{-1}\int f$, $B_0=L^{-1}\int g$. Với tần số bị thay đổi bởi hạng $u$ hoặc cản $u_t$, giải lại ODE của từng mode; không dùng nguyên tần số $ak_n$.

\textbf{Bước 5 — Kiểm tra biên, dữ liệu và hội tụ.} Mỗi mode thỏa biên; tại $t=0$ hai chuỗi cho khai triển của $f,g$. Nếu dữ liệu không tương thích ở góc, giải theo nghĩa năng lượng hoặc từng miền và nêu đường mất tính trơn; không đạo hàm một chuỗi yếu tùy ý. Chu kỳ được suy từ tất cả mode có hệ số khác không, kể cả mode hằng; hai tần số có tỷ số vô tỷ ngăn một chu kỳ chung.

\textbf{Chi tiết tính tích phân hệ số.} Với $k=n\pi/L$, hai nguyên hàm thường dùng là


\[
\int x\sin kx\,dx=-x\cos kx/k+\sin kx/k^2,\qquad
\int x\cos kx\,dx=x\sin kx/k+\cos kx/k^2.
\]


Với đa thức $P$, ghi cả hai lần tích phân:


\[
\int_0^L P\sin kx\,dx
=\left[-P\cos kx/k\right]_0^L+
\left[P'\sin kx/k^2\right]_0^L-
\frac1{k^2}\int_0^L P''\sin kx\,dx.
\]


Chỉ bỏ hạng biên sau khi thế $0,L$ và điều kiện thực sự của $P$. Nếu dữ liệu chia đoạn, áp dụng công thức trên từng đoạn, rồi cộng các hạng tại điểm chia. Cuối cùng nhân nghịch đảo chuẩn của hàm riêng; nếu hệ số là của vận tốc, chia thêm tần số khi ghép nghiệm. Với mẫu số bằng $0$ trong tích–tổng, quay lại tích phân gốc để tính riêng.
\subsection*{Lời giải áp dụng cho bài}
\textbf{Bước 1. Đạo hàm năng lượng.} Với $u_{tt}+\xi u_t=u_{xx}$, $\xi\ge0$, biên cố định, $E=\int_0^L(u_t^2+u_x^2)dx$. Nhân PDE với $2u_t$ và tích phân từng phần:
\[E'=2[u_tu_x]_0^L-2\xi\int_0^Lu_t^2dx=-2\xi\int_0^Lu_t^2dx.\]
Hạng biên hết vì đạo hàm thời gian của vết Dirichlet bằng $0$.

\textbf{Bước 2. Ước lượng trên.} Tích phân bình phương không âm, nên $E'\le0$ và $E(t)\le E(0)$.

\textbf{Bước 3. Ước lượng dưới.} $\int u_t^2\le E$, nên $E'\ge-2\xi E$. Do đó $(e^{2\xi t}E)'=e^{2\xi t}(E'+2\xi E)\ge0$. So với giá trị đầu:
\[e^{-2\xi t}E(0)\le E(t)\le E(0).\]
\textbf{Bước 4. Kiểm tra hệ số.} Vì định nghĩa $E$ không có $1/2$, hệ số đạo hàm là $2\xi$. Với $\xi=0$, hai cận bằng nhau và thu bảo toàn năng lượng. Với $E(0)=0$ suy nghiệm duy nhất bằng $0$ khi li độ đầu cũng bằng $0$. Các công thức đúng trên đoạn $(0,L)$ sau khi thống nhất ký hiệu cận phải của đề.

\clearpage\noindent\texttt{\detokenize{VD_3_6_2_01}}\par
\section*{Lời giải bổ sung — Ví dụ 1 · Mục 3.6.2}

\subsection*{Kiến thức và trình tự giải}
\textbf{Bước 1 — Chọn Kirchhoff (3D) hoặc Poisson (2D).} Với li độ đầu bằng không, vận tốc $\psi$, tốc độ $c$, đặt $\rho=ct$. Trong ba chiều,


\[
u(P,t)=\frac t{4\pi}\int_{S^2}\psi(P+\rho\omega)d\omega.
\]


Trong hai chiều,


\[
u(P,t)=\frac1{2\pi c}\int_{|Q-P|<ct}\frac{\psi(Q)}{\sqrt{c^2t^2-|Q-P|^2}}dQ.
\]


Với li độ $f$ khác không, thêm đạo hàm theo $t$ của chính biểu thức tương ứng khi thay $\psi$ bằng $f$. Với nguồn và dữ liệu đầu không, tích phân theo thời gian của nghiệm có vận tốc là nguồn tại từng thời điểm (Duhamel).

\textbf{Bước 2 — Tận dụng hình học và đối xứng.} Chỉ thị miền $D$ biến Kirchhoff thành $t$ nhân tỷ lệ diện tích mặt cầu nằm trong $D$. Chọn trục qua tâm miền và $P$, đặt $\mu=\cos\theta$, $d\omega=d\mu\,d\phi$. Nếu điều kiện không phụ thuộc $\phi$, tỷ lệ là $|E|/2$ với $E\subset[-1,1]$. Mỗi nửa không gian cho một bất đẳng thức tuyến tính theo $\mu$; quả cầu cho một ngưỡng chỏm. Miền hợp/rút bỏ cho tổng/hiệu nếu các phần không chồng diện tích.

\textbf{Bước 3 — Viết chính xác các khoảng hợp lệ.} Giải bất đẳng thức gốc, sau đó mới bình phương khi có căn. Chia $[-1,1]$ ở nghiệm phương trình biên, kiểm tra dấu ở trung điểm từng khoảng, loại nghiệm ngoại lai. Nếu còn góc phương vị, dùng độ dài tập $\cos\phi<v$ rồi còn một tích phân theo $\mu$. Với Poisson, tích phân theo một tọa độ cho hiệu arcsin; cắt các cận vào đĩa trước khi tính.

\textbf{Bước 4 — Tính nguồn bằng nguyên hàm thời gian.} Với quả cầu bán kính $R$ cách $P$ khoảng $r>0$, đặt


\[
K_R(r,t)=t\left[\frac{R^2-(r-ct)^2}{4rct}\right]_0^1,\quad J_R(r,t)=\int_0^tK_R(r,s)ds.
\]


Tách ba giai đoạn: cả mặt cầu nằm trong miền, giao một chỏm, không giao. Nguồn bật trên $[0,T]$ cho $J_R(r,t)-J_R(r,(t-T)_+)$. Trường hợp $r=0$ tính trực tiếp, tránh chia cho $r$.

\textbf{Bước 5 — Kiểm tra chuẩn hóa, mặt sóng và cực đại.} Nếu $\psi=1$ trên toàn không gian thì $u=t$, dùng để kiểm tra hệ số. Với dữ liệu bước nhảy, chỉ kết luận cổ điển ngoài các mặt sóng. Cực đại phải xét cả điểm chuyển nhánh và giá trị một phía; nêu rõ supremum nếu chỉ đạt theo giới hạn. Nếu dùng tích phân số, ghi tham số và cách kiểm tra hội tụ; đồ thị không tự chứng minh một cực đại toàn cục.
\subsection*{Lời giải áp dụng cho bài}
\textbf{Bước 1. Đưa Kirchhoff về diện tích chỏm.} Tốc độ $c=2$, li độ đầu không, vận tốc là chỉ thị $z>0$. Tại $P=(x,y,z)$, điểm trên mặt cầu là $P+2t\omega$ và thuộc miền vận tốc khi $z+2t\omega_z>0$. Đặt $\mu=\omega_z$; diện tích hướng có $\mu$ trong một khoảng bằng $2\pi$ nhân độ dài khoảng ấy.

\textbf{Bước 2. Cắt ngưỡng vào $[-1,1]$.} Điều kiện là $\mu>-z/(2t)$. Nếu $z\ge2t$, cả mặt cầu nằm phía trên; nếu $z\le-2t$, không có phần nào; nếu $|z|<2t$, độ dài khoảng là $1+z/(2t)$. Vì $u=t|E|/2$,
\[u(P,t)=\begin{cases}t&z\ge2t,\\t/2+z/4&|z|<2t,\\0&z\le-2t.\end{cases}\]
\textbf{Bước 3. Kiểm tra.} Ngoài các mặt $z=\pm2t$, mỗi nhánh affine nên $u_{tt}=4\Delta u=0$. $u(P,0)=0$; với $z\ne0$, $u_t(P,0)$ là $1$ hoặc $0$ đúng dữ liệu. Nghiệm liên tục nhưng gãy cấp một trên các mặt sóng. Tại $z=0$ dữ liệu chỉ thị gián đoạn nên không có vết cổ điển thống nhất; công thức cho $u=t/2$, là cách hiểu từ trung bình mặt cầu.

\clearpage\noindent\texttt{\detokenize{VD_3_6_2_02}}\par
\section*{Lời giải bổ sung — Ví dụ 2 · Mục 3.6.2}

\subsection*{Kiến thức và trình tự giải}
\textbf{Bước 1 — Chọn Kirchhoff (3D) hoặc Poisson (2D).} Với li độ đầu bằng không, vận tốc $\psi$, tốc độ $c$, đặt $\rho=ct$. Trong ba chiều,


\[
u(P,t)=\frac t{4\pi}\int_{S^2}\psi(P+\rho\omega)d\omega.
\]


Trong hai chiều,


\[
u(P,t)=\frac1{2\pi c}\int_{|Q-P|<ct}\frac{\psi(Q)}{\sqrt{c^2t^2-|Q-P|^2}}dQ.
\]


Với li độ $f$ khác không, thêm đạo hàm theo $t$ của chính biểu thức tương ứng khi thay $\psi$ bằng $f$. Với nguồn và dữ liệu đầu không, tích phân theo thời gian của nghiệm có vận tốc là nguồn tại từng thời điểm (Duhamel).

\textbf{Bước 2 — Tận dụng hình học và đối xứng.} Chỉ thị miền $D$ biến Kirchhoff thành $t$ nhân tỷ lệ diện tích mặt cầu nằm trong $D$. Chọn trục qua tâm miền và $P$, đặt $\mu=\cos\theta$, $d\omega=d\mu\,d\phi$. Nếu điều kiện không phụ thuộc $\phi$, tỷ lệ là $|E|/2$ với $E\subset[-1,1]$. Mỗi nửa không gian cho một bất đẳng thức tuyến tính theo $\mu$; quả cầu cho một ngưỡng chỏm. Miền hợp/rút bỏ cho tổng/hiệu nếu các phần không chồng diện tích.

\textbf{Bước 3 — Viết chính xác các khoảng hợp lệ.} Giải bất đẳng thức gốc, sau đó mới bình phương khi có căn. Chia $[-1,1]$ ở nghiệm phương trình biên, kiểm tra dấu ở trung điểm từng khoảng, loại nghiệm ngoại lai. Nếu còn góc phương vị, dùng độ dài tập $\cos\phi<v$ rồi còn một tích phân theo $\mu$. Với Poisson, tích phân theo một tọa độ cho hiệu arcsin; cắt các cận vào đĩa trước khi tính.

\textbf{Bước 4 — Tính nguồn bằng nguyên hàm thời gian.} Với quả cầu bán kính $R$ cách $P$ khoảng $r>0$, đặt


\[
K_R(r,t)=t\left[\frac{R^2-(r-ct)^2}{4rct}\right]_0^1,\quad J_R(r,t)=\int_0^tK_R(r,s)ds.
\]


Tách ba giai đoạn: cả mặt cầu nằm trong miền, giao một chỏm, không giao. Nguồn bật trên $[0,T]$ cho $J_R(r,t)-J_R(r,(t-T)_+)$. Trường hợp $r=0$ tính trực tiếp, tránh chia cho $r$.

\textbf{Bước 5 — Kiểm tra chuẩn hóa, mặt sóng và cực đại.} Nếu $\psi=1$ trên toàn không gian thì $u=t$, dùng để kiểm tra hệ số. Với dữ liệu bước nhảy, chỉ kết luận cổ điển ngoài các mặt sóng. Cực đại phải xét cả điểm chuyển nhánh và giá trị một phía; nêu rõ supremum nếu chỉ đạt theo giới hạn. Nếu dùng tích phân số, ghi tham số và cách kiểm tra hội tụ; đồ thị không tự chứng minh một cực đại toàn cục.
\subsection*{Lời giải áp dụng cho bài}
\textbf{Bước 1. Chọn trục theo điểm quan sát.} $c=2$, $\psi=\chi_{B_R(0)}$, $r=|P|$, $\rho=2t$. Nếu $r>0$, lấy trục của $\mu$ hướng về tâm từ $P$. Khi đó $|P+\rho\omega|^2=r^2+\rho^2-2r\rho\mu$.

\textbf{Bước 2. Bất đẳng thức quả cầu.} Điều kiện nằm trong $B_R$ tương đương $\mu>(r^2+\rho^2-R^2)/(2r\rho)$. Tỷ lệ diện tích chỏm là ngưỡng đã cắt:
\[K_R(r,t)=t\left[\frac{R^2-(r-2t)^2}{8rt}\right]_0^1,\qquad [s]_0^1=\min(1,\max(0,s)).\]
\textbf{Bước 3. Giải thích các nhánh.} Nếu $r<R$ và $2t<R-r$, cả mặt cầu ở trong, $u=t$. Nếu $|R-r|<2t<R+r$, chỉ một chỏm giao, $u=[R^2-(r-2t)^2]/(8r)$. Nếu không giao, $u=0$. Với $r=0$, tính riêng: $u=t$ khi $2t<R$, $0$ khi $2t>R$; giá trị tại $2t=R$ phụ thuộc quy ước chỉ thị trên toàn mặt cầu, với quả cầu mở là $0$.

\textbf{Bước 4. Kiểm tra.} Khi điểm ban đầu nằm trong quả cầu, nhánh đầu cho vận tốc $1$; ở ngoài cho $0$. Thời gian đến/đi là khoảng cách tới mặt cầu chia tốc độ. Dữ liệu gián đoạn cho nghiệm không cổ điển trên các mặt sóng; trường hợp tâm quan sát có thể có bước nhảy khi toàn mặt cầu đồng thời đi qua biên.

\clearpage\noindent\texttt{\detokenize{VD_3_6_2_03}}\par
\section*{Lời giải bổ sung — Ví dụ 3 · Mục 3.6.2}

\subsection*{Kiến thức và trình tự giải}
\textbf{Bước 1 — Chọn Kirchhoff (3D) hoặc Poisson (2D).} Với li độ đầu bằng không, vận tốc $\psi$, tốc độ $c$, đặt $\rho=ct$. Trong ba chiều,


\[
u(P,t)=\frac t{4\pi}\int_{S^2}\psi(P+\rho\omega)d\omega.
\]


Trong hai chiều,


\[
u(P,t)=\frac1{2\pi c}\int_{|Q-P|<ct}\frac{\psi(Q)}{\sqrt{c^2t^2-|Q-P|^2}}dQ.
\]


Với li độ $f$ khác không, thêm đạo hàm theo $t$ của chính biểu thức tương ứng khi thay $\psi$ bằng $f$. Với nguồn và dữ liệu đầu không, tích phân theo thời gian của nghiệm có vận tốc là nguồn tại từng thời điểm (Duhamel).

\textbf{Bước 2 — Tận dụng hình học và đối xứng.} Chỉ thị miền $D$ biến Kirchhoff thành $t$ nhân tỷ lệ diện tích mặt cầu nằm trong $D$. Chọn trục qua tâm miền và $P$, đặt $\mu=\cos\theta$, $d\omega=d\mu\,d\phi$. Nếu điều kiện không phụ thuộc $\phi$, tỷ lệ là $|E|/2$ với $E\subset[-1,1]$. Mỗi nửa không gian cho một bất đẳng thức tuyến tính theo $\mu$; quả cầu cho một ngưỡng chỏm. Miền hợp/rút bỏ cho tổng/hiệu nếu các phần không chồng diện tích.

\textbf{Bước 3 — Viết chính xác các khoảng hợp lệ.} Giải bất đẳng thức gốc, sau đó mới bình phương khi có căn. Chia $[-1,1]$ ở nghiệm phương trình biên, kiểm tra dấu ở trung điểm từng khoảng, loại nghiệm ngoại lai. Nếu còn góc phương vị, dùng độ dài tập $\cos\phi<v$ rồi còn một tích phân theo $\mu$. Với Poisson, tích phân theo một tọa độ cho hiệu arcsin; cắt các cận vào đĩa trước khi tính.

\textbf{Bước 4 — Tính nguồn bằng nguyên hàm thời gian.} Với quả cầu bán kính $R$ cách $P$ khoảng $r>0$, đặt


\[
K_R(r,t)=t\left[\frac{R^2-(r-ct)^2}{4rct}\right]_0^1,\quad J_R(r,t)=\int_0^tK_R(r,s)ds.
\]


Tách ba giai đoạn: cả mặt cầu nằm trong miền, giao một chỏm, không giao. Nguồn bật trên $[0,T]$ cho $J_R(r,t)-J_R(r,(t-T)_+)$. Trường hợp $r=0$ tính trực tiếp, tránh chia cho $r$.

\textbf{Bước 5 — Kiểm tra chuẩn hóa, mặt sóng và cực đại.} Nếu $\psi=1$ trên toàn không gian thì $u=t$, dùng để kiểm tra hệ số. Với dữ liệu bước nhảy, chỉ kết luận cổ điển ngoài các mặt sóng. Cực đại phải xét cả điểm chuyển nhánh và giá trị một phía; nêu rõ supremum nếu chỉ đạt theo giới hạn. Nếu dùng tích phân số, ghi tham số và cách kiểm tra hội tụ; đồ thị không tự chứng minh một cực đại toàn cục.
\subsection*{Lời giải áp dụng cho bài}
\textbf{Bước 1. Dùng phương pháp hạ chiều.} Dữ liệu hai chiều $u_0=0$, $u_{t,0}=f(x)$ không phụ thuộc $y$. Mở rộng sang ba chiều cũng không phụ thuộc $y,z$. Kirchhoff cho $u=t(4\pi)^{-1}\int_{S^2}f(x+ct\omega_x)d\omega$.

\textbf{Bước 2. Tích phân góc.} Chọn trục cực là trục $x$, $\mu=\omega_x$, $d\omega=d\phi\,d\mu$. Tích phân theo $\phi$ cho $2\pi$, nên $u=(t/2)\int_{-1}^1f(x+ct\mu)d\mu$. Đổi $s=x+ct\mu$, $ds=ct\,d\mu$:
\[u(x,y,t)=\frac1{2c}\int_{x-ct}^{x+ct}f(s)ds.\]
\textbf{Bước 3. Kiểm tra bằng d'Alembert.} Nghiệm không phụ thuộc $y$, nên $\Delta u=u_{xx}$. Công thức vừa tìm đúng d'Alembert một chiều với li độ không và vận tốc $f$. Tại $0$ tích phân có độ dài $0$, còn đạo hàm thời gian cho $f(x)$. Nếu $f$ không trơn, công thức vẫn có nghĩa và mức trơn được xét tại các cận chạm điểm gãy của $f$.

\clearpage\noindent\texttt{\detokenize{VD_3_6_2_04}}\par
\section*{Lời giải bổ sung — Ví dụ 4 · Mục 3.6.2}

\subsection*{Kiến thức và trình tự giải}
\textbf{Bước 1 — Chọn Kirchhoff (3D) hoặc Poisson (2D).} Với li độ đầu bằng không, vận tốc $\psi$, tốc độ $c$, đặt $\rho=ct$. Trong ba chiều,


\[
u(P,t)=\frac t{4\pi}\int_{S^2}\psi(P+\rho\omega)d\omega.
\]


Trong hai chiều,


\[
u(P,t)=\frac1{2\pi c}\int_{|Q-P|<ct}\frac{\psi(Q)}{\sqrt{c^2t^2-|Q-P|^2}}dQ.
\]


Với li độ $f$ khác không, thêm đạo hàm theo $t$ của chính biểu thức tương ứng khi thay $\psi$ bằng $f$. Với nguồn và dữ liệu đầu không, tích phân theo thời gian của nghiệm có vận tốc là nguồn tại từng thời điểm (Duhamel).

\textbf{Bước 2 — Tận dụng hình học và đối xứng.} Chỉ thị miền $D$ biến Kirchhoff thành $t$ nhân tỷ lệ diện tích mặt cầu nằm trong $D$. Chọn trục qua tâm miền và $P$, đặt $\mu=\cos\theta$, $d\omega=d\mu\,d\phi$. Nếu điều kiện không phụ thuộc $\phi$, tỷ lệ là $|E|/2$ với $E\subset[-1,1]$. Mỗi nửa không gian cho một bất đẳng thức tuyến tính theo $\mu$; quả cầu cho một ngưỡng chỏm. Miền hợp/rút bỏ cho tổng/hiệu nếu các phần không chồng diện tích.

\textbf{Bước 3 — Viết chính xác các khoảng hợp lệ.} Giải bất đẳng thức gốc, sau đó mới bình phương khi có căn. Chia $[-1,1]$ ở nghiệm phương trình biên, kiểm tra dấu ở trung điểm từng khoảng, loại nghiệm ngoại lai. Nếu còn góc phương vị, dùng độ dài tập $\cos\phi<v$ rồi còn một tích phân theo $\mu$. Với Poisson, tích phân theo một tọa độ cho hiệu arcsin; cắt các cận vào đĩa trước khi tính.

\textbf{Bước 4 — Tính nguồn bằng nguyên hàm thời gian.} Với quả cầu bán kính $R$ cách $P$ khoảng $r>0$, đặt


\[
K_R(r,t)=t\left[\frac{R^2-(r-ct)^2}{4rct}\right]_0^1,\quad J_R(r,t)=\int_0^tK_R(r,s)ds.
\]


Tách ba giai đoạn: cả mặt cầu nằm trong miền, giao một chỏm, không giao. Nguồn bật trên $[0,T]$ cho $J_R(r,t)-J_R(r,(t-T)_+)$. Trường hợp $r=0$ tính trực tiếp, tránh chia cho $r$.

\textbf{Bước 5 — Kiểm tra chuẩn hóa, mặt sóng và cực đại.} Nếu $\psi=1$ trên toàn không gian thì $u=t$, dùng để kiểm tra hệ số. Với dữ liệu bước nhảy, chỉ kết luận cổ điển ngoài các mặt sóng. Cực đại phải xét cả điểm chuyển nhánh và giá trị một phía; nêu rõ supremum nếu chỉ đạt theo giới hạn. Nếu dùng tích phân số, ghi tham số và cách kiểm tra hội tụ; đồ thị không tự chứng minh một cực đại toàn cục.
\subsection*{Lời giải áp dụng cho bài}
\textbf{Bước 1. Tách hai nguồn vận tốc.} $c=2$, $u_0=0$, $\psi=\chi_{[0,1]}(x)-\chi_{[0,1]}(y)$. Tuyến tính cho hiệu của một nghiệm chỉ phụ thuộc $x$ và một nghiệm chỉ phụ thuộc $y$.

\textbf{Bước 2. Tính độ dài giao.} Đặt $h(s)=[s]_0^1$, nguyên hàm của $\chi_{[0,1]}$. Theo Ví dụ 3,
\[u=\frac14\{h(x+2t)-h(x-2t)-h(y+2t)+h(y-2t)\}.\]
Mỗi hiệu $h$ bằng độ dài giao $[q-2t,q+2t]\cap[0,1]$ với $q=x$ hoặc $y$, nên có thể đọc mọi nhánh từ min/max.

\textbf{Bước 3. Kiểm tra và đối xứng.} Tại $t=0$, $u=0$; đạo hàm thời gian cho $\chi_{[0,1]}(x)-\chi_{[0,1]}(y)$. Hai phần thỏa PDE một chiều, nên tổng thỏa PDE hai chiều. Đổi $x,y$ đổi dấu, do đó $u(x,x,t)=0$. Khi cả hai khoảng ảnh hưởng đã bao trọn $[0,1]$, hai phần cùng bằng $1/4$ và triệt tiêu. Các mặt gãy là $x\pm2t=0,1$ và $y\pm2t=0,1$.

\clearpage\noindent\texttt{\detokenize{VD_3_6_2_05}}\par
\section*{Lời giải bổ sung — Ví dụ 5 · Mục 3.6.2}

\subsection*{Kiến thức và trình tự giải}
\textbf{Bước 1 — Chọn Kirchhoff (3D) hoặc Poisson (2D).} Với li độ đầu bằng không, vận tốc $\psi$, tốc độ $c$, đặt $\rho=ct$. Trong ba chiều,


\[
u(P,t)=\frac t{4\pi}\int_{S^2}\psi(P+\rho\omega)d\omega.
\]


Trong hai chiều,


\[
u(P,t)=\frac1{2\pi c}\int_{|Q-P|<ct}\frac{\psi(Q)}{\sqrt{c^2t^2-|Q-P|^2}}dQ.
\]


Với li độ $f$ khác không, thêm đạo hàm theo $t$ của chính biểu thức tương ứng khi thay $\psi$ bằng $f$. Với nguồn và dữ liệu đầu không, tích phân theo thời gian của nghiệm có vận tốc là nguồn tại từng thời điểm (Duhamel).

\textbf{Bước 2 — Tận dụng hình học và đối xứng.} Chỉ thị miền $D$ biến Kirchhoff thành $t$ nhân tỷ lệ diện tích mặt cầu nằm trong $D$. Chọn trục qua tâm miền và $P$, đặt $\mu=\cos\theta$, $d\omega=d\mu\,d\phi$. Nếu điều kiện không phụ thuộc $\phi$, tỷ lệ là $|E|/2$ với $E\subset[-1,1]$. Mỗi nửa không gian cho một bất đẳng thức tuyến tính theo $\mu$; quả cầu cho một ngưỡng chỏm. Miền hợp/rút bỏ cho tổng/hiệu nếu các phần không chồng diện tích.

\textbf{Bước 3 — Viết chính xác các khoảng hợp lệ.} Giải bất đẳng thức gốc, sau đó mới bình phương khi có căn. Chia $[-1,1]$ ở nghiệm phương trình biên, kiểm tra dấu ở trung điểm từng khoảng, loại nghiệm ngoại lai. Nếu còn góc phương vị, dùng độ dài tập $\cos\phi<v$ rồi còn một tích phân theo $\mu$. Với Poisson, tích phân theo một tọa độ cho hiệu arcsin; cắt các cận vào đĩa trước khi tính.

\textbf{Bước 4 — Tính nguồn bằng nguyên hàm thời gian.} Với quả cầu bán kính $R$ cách $P$ khoảng $r>0$, đặt


\[
K_R(r,t)=t\left[\frac{R^2-(r-ct)^2}{4rct}\right]_0^1,\quad J_R(r,t)=\int_0^tK_R(r,s)ds.
\]


Tách ba giai đoạn: cả mặt cầu nằm trong miền, giao một chỏm, không giao. Nguồn bật trên $[0,T]$ cho $J_R(r,t)-J_R(r,(t-T)_+)$. Trường hợp $r=0$ tính trực tiếp, tránh chia cho $r$.

\textbf{Bước 5 — Kiểm tra chuẩn hóa, mặt sóng và cực đại.} Nếu $\psi=1$ trên toàn không gian thì $u=t$, dùng để kiểm tra hệ số. Với dữ liệu bước nhảy, chỉ kết luận cổ điển ngoài các mặt sóng. Cực đại phải xét cả điểm chuyển nhánh và giá trị một phía; nêu rõ supremum nếu chỉ đạt theo giới hạn. Nếu dùng tích phân số, ghi tham số và cách kiểm tra hội tụ; đồ thị không tự chứng minh một cực đại toàn cục.
\subsection*{Lời giải áp dụng cho bài}
\textbf{Bước 1. Duhamel cho nguồn bật trong một đơn vị thời gian.} $c=2$, dữ liệu đầu không, $F(P,s)=\chi_{[0,1]}(s)\chi_{B_2(0)}(P)$. Xung tại thời điểm $s$ có nghiệm $K_2(r,t-s)$, $r=|P|$, nên
\[u=\int_0^{\min(t,1)}K_2(r,t-s)ds=J_2(r,t)-J_2(r,(t-1)_+).\]
\textbf{Bước 2. Tính nguyên hàm chỏm với tốc độ tổng quát.} Với $R,c>0$, đặt $b=|R-r|/c$, $d=(R+r)/c$, $q=[t]_b^d$ và
\[A(s)=\frac{(R^2-r^2)s+rcs^2-c^2s^3/3}{4rc}.\]
Ở phần giao chỏm, $K_R=[R^2-(r-cs)^2]/(4rc)=A'(s)$. Nếu $r<R$, trước $b$ có $K_R=s$; nếu $r\ge R$, trước $b$ bằng $0$. Do đó, với $r>0$,
\[J_R(r,t)=\tfrac12\min(t,(R-r)_+/c)^2+A(q)-A(b).\]
Với $r=0$, $J_R(0,t)=\frac12\min(t,R/c)^2$.

\textbf{Bước 3. Áp dụng $R=c=2$ và kiểm tra.} Thay hai bộ thời gian $t,(t-1)_+$ vào công thức; mọi nhánh đều là đa thức và min/max. Duhamel đảm bảo hai dữ liệu đầu không và PDE có đúng nguồn. Khi $t\ge1+(r+2)/2$, cả hai $J$ đã đạt cùng hằng số, nên $u=0$ tại điểm ấy: xung nguồn hữu hạn đi qua hoàn toàn trong ba chiều.

\clearpage\noindent\texttt{\detokenize{BT_3_6_3_01}}\par
\section*{Lời giải bổ sung — Bài tập 1 · Mục 3.6.3}

\subsection*{Kiến thức và trình tự giải}
\textbf{Bước 1 — Chọn Kirchhoff (3D) hoặc Poisson (2D).} Với li độ đầu bằng không, vận tốc $\psi$, tốc độ $c$, đặt $\rho=ct$. Trong ba chiều,


\[
u(P,t)=\frac t{4\pi}\int_{S^2}\psi(P+\rho\omega)d\omega.
\]


Trong hai chiều,


\[
u(P,t)=\frac1{2\pi c}\int_{|Q-P|<ct}\frac{\psi(Q)}{\sqrt{c^2t^2-|Q-P|^2}}dQ.
\]


Với li độ $f$ khác không, thêm đạo hàm theo $t$ của chính biểu thức tương ứng khi thay $\psi$ bằng $f$. Với nguồn và dữ liệu đầu không, tích phân theo thời gian của nghiệm có vận tốc là nguồn tại từng thời điểm (Duhamel).

\textbf{Bước 2 — Tận dụng hình học và đối xứng.} Chỉ thị miền $D$ biến Kirchhoff thành $t$ nhân tỷ lệ diện tích mặt cầu nằm trong $D$. Chọn trục qua tâm miền và $P$, đặt $\mu=\cos\theta$, $d\omega=d\mu\,d\phi$. Nếu điều kiện không phụ thuộc $\phi$, tỷ lệ là $|E|/2$ với $E\subset[-1,1]$. Mỗi nửa không gian cho một bất đẳng thức tuyến tính theo $\mu$; quả cầu cho một ngưỡng chỏm. Miền hợp/rút bỏ cho tổng/hiệu nếu các phần không chồng diện tích.

\textbf{Bước 3 — Viết chính xác các khoảng hợp lệ.} Giải bất đẳng thức gốc, sau đó mới bình phương khi có căn. Chia $[-1,1]$ ở nghiệm phương trình biên, kiểm tra dấu ở trung điểm từng khoảng, loại nghiệm ngoại lai. Nếu còn góc phương vị, dùng độ dài tập $\cos\phi<v$ rồi còn một tích phân theo $\mu$. Với Poisson, tích phân theo một tọa độ cho hiệu arcsin; cắt các cận vào đĩa trước khi tính.

\textbf{Bước 4 — Tính nguồn bằng nguyên hàm thời gian.} Với quả cầu bán kính $R$ cách $P$ khoảng $r>0$, đặt


\[
K_R(r,t)=t\left[\frac{R^2-(r-ct)^2}{4rct}\right]_0^1,\quad J_R(r,t)=\int_0^tK_R(r,s)ds.
\]


Tách ba giai đoạn: cả mặt cầu nằm trong miền, giao một chỏm, không giao. Nguồn bật trên $[0,T]$ cho $J_R(r,t)-J_R(r,(t-T)_+)$. Trường hợp $r=0$ tính trực tiếp, tránh chia cho $r$.

\textbf{Bước 5 — Kiểm tra chuẩn hóa, mặt sóng và cực đại.} Nếu $\psi=1$ trên toàn không gian thì $u=t$, dùng để kiểm tra hệ số. Với dữ liệu bước nhảy, chỉ kết luận cổ điển ngoài các mặt sóng. Cực đại phải xét cả điểm chuyển nhánh và giá trị một phía; nêu rõ supremum nếu chỉ đạt theo giới hạn. Nếu dùng tích phân số, ghi tham số và cách kiểm tra hội tụ; đồ thị không tự chứng minh một cực đại toàn cục.
\subsection*{Lời giải áp dụng cho bài}
\textbf{Bước 1. Quy các miền về cầu, nửa không gian và lớp.} Với vận tốc chỉ thị, Kirchhoff bằng $t$ nhân tỷ lệ diện tích mặt cầu. Các quả cầu bị bỏ ở (a),(b),(c) nằm trong miền lớn và đôi cầu ở (a) không chồng thể tích, nên được lấy hiệu các chỉ thị. Các quả cầu ở (d) rời nhau, nên cộng ba đóng góp. Đặt $\mu$ hướng về tâm cầu; điều kiện $r^2+c^2t^2-2rct\mu<R^2$ cho ngưỡng chỏm.

\textbf{Bước 2. Tính các khoảng cách đúng điểm quan sát.} (a) $P=(0,100,0)$ cách tâm lớn $100$, cách hai tâm $(\pm1,0,0)$ cùng $\sqrt{10001}$. (b) $P_z=100$, cầu bị bỏ có tâm $z=-2$, cách $102$; nửa không gian $z<0$ cho $\mu<-100/(2t)$, tỷ lệ $(1-50/t)_+/2$. (c) lớp $0<z<4$ cho $-100/(3t)<\mu<-96/(3t)$, còn cầu tâm $z=2$ cách $98$. (d) khoảng cách là $98$ tới $(0,2,0)$ và $\sqrt{10004}$ tới hai tâm $(\pm2,0,0)$. (e) hiệu hai góc phần tư là $(\operatorname{sgn}x+\operatorname{sgn}y)/2$; mỗi tọa độ của hướng cầu đều phân bố đều trên $[-1,1]$.

\textbf{Bước 3. Thay tỷ lệ diện tích vào Kirchhoff.}

Kirchhoff với li độ đầu 0: $u(P,t)=t(4\pi)^{-1}\int_{S^2}\psi(P+ct\omega)d\omega$. Đối với một quả cầu bán kính $R$ có tâm cách $P$ một khoảng $d>0$, diện tích chỏm cho
\[K_R(d,t)=t\left[\frac{R^2-(d-ct)^2}{4dct}\right]_0^1,\quad [s]_0^1=\max(0,\min(1,s)).\]
Nếu $d=0$, $K_R(0,t)=t\chi_{ct<R}$. Giá trị trên mặt sóng hiểu theo giới hạn; dữ liệu bước nhảy không cho nghiệm cổ điển trên mặt ấy.

(a) $u=K_2(100,t)-2K_1(\sqrt{10001},t)$, $c=2$.
(b) $u=(t-50)_+/2-K_1(102,t)$, $c=2$; hạng đầu là tỷ lệ mặt cầu ở nửa không gian $z<0$.
(c) Với $c=3$, tỷ lệ của lớp $0<z<4$ cho
\[u=\frac t2\left[\left[-\frac{96}{3t}\right]_{-1}^1-\left[-\frac{100}{3t}\right]_{-1}^1\right]-K_1(98,t).\]
Ở đây $[s]_{-1}^1=\max(-1,\min(1,s))$.
(d) $u=K_1(98,t)+2K_1(\sqrt{10004},t)$, $c=3$.
(e) $\psi=(\operatorname{sgn}x+\operatorname{sgn}y)/2$. Tọa độ mỗi hướng cầu phân bố đều trên $[-1,1]$; ở $P=(100,100,0)$ trung bình sign là $\min(1,100/(3t))$. Do đó $u=\min(t,100/3)$.

\textbf{Bước 4. Đối chiếu và kết luận.} Ở (a),(c),(d), sau khi mặt cầu ra ngoài toàn bộ các miền bị chặn, đóng góp bằng $0$; ở (b),(e), miền vận tốc không bị chặn nên không suy như vậy. Tại các điểm dữ liệu liên tục, đạo hàm đầu đúng bằng chỉ thị; chuẩn hóa $\psi=1$ cho $u=t$. Các mặt tiếp xúc với biên miền là nơi cần xét vết một phía.

\clearpage\noindent\texttt{\detokenize{BT_3_6_3_02}}\par
\section*{Lời giải bổ sung — Bài tập 2 · Mục 3.6.3}

\subsection*{Kiến thức và trình tự giải}
\textbf{Bước 1 — Chọn Kirchhoff (3D) hoặc Poisson (2D).} Với li độ đầu bằng không, vận tốc $\psi$, tốc độ $c$, đặt $\rho=ct$. Trong ba chiều,


\[
u(P,t)=\frac t{4\pi}\int_{S^2}\psi(P+\rho\omega)d\omega.
\]


Trong hai chiều,


\[
u(P,t)=\frac1{2\pi c}\int_{|Q-P|<ct}\frac{\psi(Q)}{\sqrt{c^2t^2-|Q-P|^2}}dQ.
\]


Với li độ $f$ khác không, thêm đạo hàm theo $t$ của chính biểu thức tương ứng khi thay $\psi$ bằng $f$. Với nguồn và dữ liệu đầu không, tích phân theo thời gian của nghiệm có vận tốc là nguồn tại từng thời điểm (Duhamel).

\textbf{Bước 2 — Tận dụng hình học và đối xứng.} Chỉ thị miền $D$ biến Kirchhoff thành $t$ nhân tỷ lệ diện tích mặt cầu nằm trong $D$. Chọn trục qua tâm miền và $P$, đặt $\mu=\cos\theta$, $d\omega=d\mu\,d\phi$. Nếu điều kiện không phụ thuộc $\phi$, tỷ lệ là $|E|/2$ với $E\subset[-1,1]$. Mỗi nửa không gian cho một bất đẳng thức tuyến tính theo $\mu$; quả cầu cho một ngưỡng chỏm. Miền hợp/rút bỏ cho tổng/hiệu nếu các phần không chồng diện tích.

\textbf{Bước 3 — Viết chính xác các khoảng hợp lệ.} Giải bất đẳng thức gốc, sau đó mới bình phương khi có căn. Chia $[-1,1]$ ở nghiệm phương trình biên, kiểm tra dấu ở trung điểm từng khoảng, loại nghiệm ngoại lai. Nếu còn góc phương vị, dùng độ dài tập $\cos\phi<v$ rồi còn một tích phân theo $\mu$. Với Poisson, tích phân theo một tọa độ cho hiệu arcsin; cắt các cận vào đĩa trước khi tính.

\textbf{Bước 4 — Tính nguồn bằng nguyên hàm thời gian.} Với quả cầu bán kính $R$ cách $P$ khoảng $r>0$, đặt


\[
K_R(r,t)=t\left[\frac{R^2-(r-ct)^2}{4rct}\right]_0^1,\quad J_R(r,t)=\int_0^tK_R(r,s)ds.
\]


Tách ba giai đoạn: cả mặt cầu nằm trong miền, giao một chỏm, không giao. Nguồn bật trên $[0,T]$ cho $J_R(r,t)-J_R(r,(t-T)_+)$. Trường hợp $r=0$ tính trực tiếp, tránh chia cho $r$.

\textbf{Bước 5 — Kiểm tra chuẩn hóa, mặt sóng và cực đại.} Nếu $\psi=1$ trên toàn không gian thì $u=t$, dùng để kiểm tra hệ số. Với dữ liệu bước nhảy, chỉ kết luận cổ điển ngoài các mặt sóng. Cực đại phải xét cả điểm chuyển nhánh và giá trị một phía; nêu rõ supremum nếu chỉ đạt theo giới hạn. Nếu dùng tích phân số, ghi tham số và cách kiểm tra hội tụ; đồ thị không tự chứng minh một cực đại toàn cục.
\subsection*{Lời giải áp dụng cho bài}
\textbf{Bước 1. Tham số hóa mặt cầu và giữ dấu của nón.} Tốc độ ký hiệu $c$, $\rho=ct$. Trên trục $z$, điểm cầu có chiều cao $z+\rho\mu$, bán kính trụ $\rho\sqrt{1-\mu^2}$. Miền $D$ là giao của $z'>\alpha r'$ với quả cầu $|Q|<R$; điều kiện $z'>0$ không được bỏ khi bình phương. Tích phân phương vị cho $2\pi$, nên $u=t|E|/2$.

\textbf{Bước 2. Giải từng ngưỡng; ở mặt phẳng $z=0$ khử góc phương vị.} Nhờ đối xứng chỉ cần $P=(r,0,0)$. Đặt $\omega=(\sqrt{1-\mu^2}\cos\phi,\sqrt{1-\mu^2}\sin\phi,\mu)$. Nón đòi $\mu>0$ và $r'^2<\rho^2\mu^2/\alpha^2$; cầu đòi $r'^2+\rho^2\mu^2<R^2$. Viết $r'^2=r^2+\rho^2(1-\mu^2)+2r\rho\sqrt{1-\mu^2}\cos\phi$, được hai cận trên cho $\cos\phi$. Độ dài góc đúng là $A(\min(B,C))$, không nhân hai tỷ lệ chỏm độc lập.

\textbf{Bước 3. Công thức tính $u$ và cực trị.}

(a) Đặt $\rho=ct$, $\mu=\cos\theta$. Tại $P=(0,0,z)$, miền hướng hợp lệ là
\[E=\{\mu\in[-1,1]:z+\rho\mu>\alpha\rho\sqrt{1-\mu^2},\quad z^2+\rho^2+2z\rho\mu<R^2\}.\]
Kirchhoff cho kết quả chính xác $u(0,0,z,t)=t|E|/2$. Để tính $|E|$ bằng hữu hạn phép toán: chia $[-1,1]$ ở các nghiệm thực của
\[(1+\alpha^2)\rho^2\mu^2+2z\rho\mu+z^2-\alpha^2\rho^2=0,\]
và điểm $(R^2-z^2-\rho^2)/(2z\rho)$ nếu $z\ne0$, cùng điểm $-z/\rho$. Trên từng khoảng kiểm tra hai bất đẳng thức gốc tại trung điểm rồi cộng độ dài các khoảng đúng. Kiểm tra dấu $z+\rho\mu$ loại nghiệm ngoại lai do bình phương. Khi $z=0$, $u=t(1-\alpha/\sqrt{1+\alpha^2})/2$ nếu $ct<R$, và 0 nếu $ct>R$.

(b) Do đối xứng quay, chỉ cần $r=\sqrt{x^2+y^2}$. Với $r>0$, đặt $D=2r\rho\sqrt{1-\mu^2}$,
\[B=\frac{R^2-r^2-\rho^2}{D},\qquad C=\frac{\rho^2\mu^2/\alpha^2-r^2-\rho^2(1-\mu^2)}D.\]
Gọi $A(v)=0$ nếu $v\le-1$, $A(v)=2\pi-2\arccos v$ nếu $|v|<1$, $A(v)=2\pi$ nếu $v\ge1$. Khi đó
\[U_r(t)=\frac t{4\pi}\int_0^1A(\min(B,C))d\mu.\]
Đây là tích phân một chiều đã loại hoàn toàn tích phân mặt. $U=0$ khi $ct>r+R$; để tìm giá trị lớn nhất, tính $U$ ở các điểm chuyển nhánh và các nghiệm $U'(t)=0$ trong $0<t<(r+R)/c$, lấy giá trị lớn nhất. Kèm mã tính tích phân và tối ưu cho $r,R,\alpha,c$ bất kỳ trong file Python chung. Không có một con số duy nhất vì đề để các tham số tùy ý. Với $r=0$, supremum là $R(1-\alpha/\sqrt{1+\alpha^2})/(2c)$; dữ liệu chỉ thị gián đoạn ở mặt cầu có thể khiến supremum chỉ đạt theo giới hạn trái.

\textbf{Bước 4. Đối chiếu và kết luận.} Với $r>0$, công thức tích phân và bài toán cực trị một biến là biểu diễn chính xác theo tham số; các con số do mã tích phân/tối ưu trả về là xấp xỉ. Mã chia lưới thời gian, tinh chỉnh từng cực đại và cho phép tăng độ phân giải để kiểm tra hội tụ, nên không dùng một lần tối ưu cục bộ làm bằng chứng cực đại toàn cục. Ở $r=0$ công thức supremum nêu trên là tường minh, suy bằng lấy $t\uparrow R/c$. Điều kiện nón strict và biên cầu phải được giữ khi xét đạt cực đại.

\textbf{Chi tiết nghiệm biên nón và độ dài phương vị.} Phương trình bậc hai ở (a) có biệt thức
\[4z^2\rho^2-4(1+\alpha^2)\rho^2(z^2-\alpha^2\rho^2)
=4\alpha^2\rho^2[(1+\alpha^2)\rho^2-z^2].\]
Nếu biểu thức dưới căn không âm và $\rho>0$, hai nghiệm là
\[\mu_\pm=\frac{-z\pm\alpha\sqrt{(1+\alpha^2)\rho^2-z^2}}{(1+\alpha^2)\rho}.\]
Giữ những nghiệm thuộc $[-1,1]$, thêm các điểm cắt cầu và dấu đã nêu, rồi kiểm tra bất đẳng thức gốc trên từng khoảng. Ở (b), $\cos\phi<v$ với $-1<v<1$ cho $\arccos v<\phi<2\pi-\arccos v$; độ dài là $2\pi-2\arccos v$. Hai điều kiện cùng là cận trên nên lấy $\min(B,C)$. Khi $r=0$, nón cho $\mu>\alpha/\sqrt{1+\alpha^2}$ và cầu cho $\rho<R$, nên tỷ lệ diện tích bằng $(1-\alpha/\sqrt{1+\alpha^2})/2$.
\begin{center}\includegraphics[width=.97\linewidth]{_xu_ly/graphs/cone_torus.png}\end{center}

\clearpage\noindent\texttt{\detokenize{BT_3_6_3_03}}\par
\section*{Lời giải bổ sung — Bài tập 3 · Mục 3.6.3}

\subsection*{Kiến thức và trình tự giải}
\textbf{Bước 1 — Chọn Kirchhoff (3D) hoặc Poisson (2D).} Với li độ đầu bằng không, vận tốc $\psi$, tốc độ $c$, đặt $\rho=ct$. Trong ba chiều,


\[
u(P,t)=\frac t{4\pi}\int_{S^2}\psi(P+\rho\omega)d\omega.
\]


Trong hai chiều,


\[
u(P,t)=\frac1{2\pi c}\int_{|Q-P|<ct}\frac{\psi(Q)}{\sqrt{c^2t^2-|Q-P|^2}}dQ.
\]


Với li độ $f$ khác không, thêm đạo hàm theo $t$ của chính biểu thức tương ứng khi thay $\psi$ bằng $f$. Với nguồn và dữ liệu đầu không, tích phân theo thời gian của nghiệm có vận tốc là nguồn tại từng thời điểm (Duhamel).

\textbf{Bước 2 — Tận dụng hình học và đối xứng.} Chỉ thị miền $D$ biến Kirchhoff thành $t$ nhân tỷ lệ diện tích mặt cầu nằm trong $D$. Chọn trục qua tâm miền và $P$, đặt $\mu=\cos\theta$, $d\omega=d\mu\,d\phi$. Nếu điều kiện không phụ thuộc $\phi$, tỷ lệ là $|E|/2$ với $E\subset[-1,1]$. Mỗi nửa không gian cho một bất đẳng thức tuyến tính theo $\mu$; quả cầu cho một ngưỡng chỏm. Miền hợp/rút bỏ cho tổng/hiệu nếu các phần không chồng diện tích.

\textbf{Bước 3 — Viết chính xác các khoảng hợp lệ.} Giải bất đẳng thức gốc, sau đó mới bình phương khi có căn. Chia $[-1,1]$ ở nghiệm phương trình biên, kiểm tra dấu ở trung điểm từng khoảng, loại nghiệm ngoại lai. Nếu còn góc phương vị, dùng độ dài tập $\cos\phi<v$ rồi còn một tích phân theo $\mu$. Với Poisson, tích phân theo một tọa độ cho hiệu arcsin; cắt các cận vào đĩa trước khi tính.

\textbf{Bước 4 — Tính nguồn bằng nguyên hàm thời gian.} Với quả cầu bán kính $R$ cách $P$ khoảng $r>0$, đặt


\[
K_R(r,t)=t\left[\frac{R^2-(r-ct)^2}{4rct}\right]_0^1,\quad J_R(r,t)=\int_0^tK_R(r,s)ds.
\]


Tách ba giai đoạn: cả mặt cầu nằm trong miền, giao một chỏm, không giao. Nguồn bật trên $[0,T]$ cho $J_R(r,t)-J_R(r,(t-T)_+)$. Trường hợp $r=0$ tính trực tiếp, tránh chia cho $r$.

\textbf{Bước 5 — Kiểm tra chuẩn hóa, mặt sóng và cực đại.} Nếu $\psi=1$ trên toàn không gian thì $u=t$, dùng để kiểm tra hệ số. Với dữ liệu bước nhảy, chỉ kết luận cổ điển ngoài các mặt sóng. Cực đại phải xét cả điểm chuyển nhánh và giá trị một phía; nêu rõ supremum nếu chỉ đạt theo giới hạn. Nếu dùng tích phân số, ghi tham số và cách kiểm tra hội tụ; đồ thị không tự chứng minh một cực đại toàn cục.
\subsection*{Lời giải áp dụng cho bài}
\textbf{Bước 1. Đưa điều kiện xuyến về một biến góc.} Trên trục $z$, bán kính trụ của điểm cầu là $\rho\sqrt{1-\mu^2}$, chiều cao $z+\rho\mu$. Thay vào $(r'-R)^2+z'^2<a^2$, khai triển để được tập $E$ dưới đây. Bình phương phương trình biên có thể sinh nghiệm ngoại lai, nên chia các khoảng và kiểm tra bất đẳng thức gốc.

\textbf{Bước 2. Ở mặt phẳng $z=0$, tính góc phương vị.} Chọn $P=(r,0,0)$, đặt $h=\sqrt{a^2-\rho^2\mu^2}$. Điều kiện xuyến tương đương $R-h<r'<R+h$, với $|\rho\mu|<a$. Vì $R>a$, cận dưới dương nên có thể bình phương hai vế. Thay $r'^2=r^2+\rho^2(1-\mu^2)+2r\rho\sqrt{1-\mu^2}\cos\phi$ cho $C_-<\cos\phi<C_+$. Độ dài góc là hiệu $A(C_+)-A(C_-)$.

\textbf{Bước 3. Tích phân một chiều và cực trị.}

(a) Đặt $\rho=ct$. Tại $(0,0,z)$, điều kiện nằm trong xuyến là
\[E=\{\mu\in[-1,1]:\rho^2+R^2+z^2+2z\rho\mu-2R\rho\sqrt{1-\mu^2}<a^2\}.\]
Do đó $u=t|E|/2$. Tính các đầu mút bằng nghiệm của
\[(\rho^2+R^2+z^2-a^2+2z\rho\mu)^2=4R^2\rho^2(1-\mu^2),\]
kiểm tra bất đẳng thức gốc ở từng khoảng để loại nghiệm ngoại lai. Đây là phương trình bậc hai theo $\mu$, nên kết quả hữu hạn phép toán. Với $z=0$, miền đúng là $|\mu|<\sqrt{1-h^2}$ khi $h=(\rho^2+R^2-a^2)/(2R\rho)<1$, nên $u=t\sqrt{1-h^2}$; ngoài $R-a<\rho<R+a$, $u=0$.

(b) Với $r=\sqrt{x^2+y^2}>0$, $|\mu|<\min(1,a/\rho)$, đặt $h=\sqrt{a^2-\rho^2\mu^2}$, $D=2r\rho\sqrt{1-\mu^2}$,
\[C_\pm=\frac{(R\pm h)^2-r^2-\rho^2(1-\mu^2)}D.\]
Độ dài góc phương vị là $A(C_+)-A(C_-)$ với $A$ định nghĩa ở bài nón. Vậy
\[U_r(t)=\frac t{4\pi}\int_{-\min(1,a/\rho)}^{\min(1,a/\rho)}[A(C_+)-A(C_-)]d\mu.\]
Chỉ cần xét $0<t<(r+R+a)/c$; tính các cực trị $U'=0$ và điểm chuyển nhánh, lấy lớn nhất. Mã Python kèm theo thực hiện bước tích phân/tối ưu với tham số tùy chọn. Riêng $r=0$, từ câu (a), $U^2=[4R^2\rho^2-(\rho^2+R^2-a^2)^2]/(4R^2c^2)$; cực đại tại $\rho=\sqrt{R^2+a^2}$, cho $\max U=a/c$.

\textbf{Bước 4. Đối chiếu và kết luận.} Ở $r=0$, $U^2$ là một đa thức bậc hai theo $\rho^2$ với hệ số đầu âm. Đạo hàm bằng $0$ tại $\rho^2=R^2+a^2$, nằm giữa $(R-a)^2,(R+a)^2$, nên đây là cực đại toàn cục $a/c$. Ở $r>0$, biểu thức tích phân và cực trị một biến là chính xác theo tham số; kết quả từ mã là xấp xỉ, cần tăng cả độ chính xác tích phân và số khoảng thời gian để kiểm tra ổn định.

\textbf{Chi tiết nghiệm bậc hai tại điểm $(0,0,z)$.} Đặt $A=\rho^2+R^2+z^2-a^2$. Phương trình biên sau bình phương là
\[4\rho^2(z^2+R^2)\mu^2+4Az\rho\mu+A^2-4R^2\rho^2=0.\]
Với $\rho>0$, hai nghiệm thực nếu biểu thức dưới căn không âm:
\[\mu_\pm=\frac{-Az\pm R\sqrt{4\rho^2(z^2+R^2)-A^2}}{2\rho(z^2+R^2)}.\]
Giữ các nghiệm trong $[-1,1]$ và kiểm tra $A+2z\rho\mu=2R\rho\sqrt{1-\mu^2}$ để loại nghiệm ngoại lai. Chia khoảng ở các nghiệm còn lại rồi kiểm tra bất đẳng thức gốc trên từng khoảng để cộng độ dài; tại $t=0$ dùng dữ liệu $u=0$.

\textbf{Chi tiết khai triển biên và cực đại trên trục.} Ở (a), $(\rho\sqrt{1-\mu^2}-R)^2+(z+\rho\mu)^2$ khai triển thành $\rho^2+R^2+z^2+2z\rho\mu-2R\rho\sqrt{1-\mu^2}$: hai hạng $\rho^2\mu^2$ triệt tiêu. Với $z=r=0$, đặt $q=\rho^2$,
\[4R^2c^2U^2=4R^2q-(q+R^2-a^2)^2.\]
Đạo hàm theo $q$ bằng $4R^2-2(q+R^2-a^2)=2(R^2+a^2-q)$, đạo hàm bậc hai bằng $-2$. Cực đại duy nhất là $q=R^2+a^2$; thay vào được $4R^2c^2U^2=4R^2a^2$, tức $U=a/c$. $q-(R-a)^2=2Ra>0$ và $(R+a)^2-q=2Ra>0$, nên điểm cực đại nằm trong miền có sóng. Ở (b), $C_-<\cos\phi<C_+$ cho độ dài phương vị bằng hiệu hai độ dài $A(C_+)-A(C_-)$ sau khi cắt cận.
\begin{center}\includegraphics[width=.97\linewidth]{_xu_ly/graphs/cone_torus.png}\end{center}

\clearpage\noindent\texttt{\detokenize{BT_3_6_3_04}}\par
\section*{Lời giải bổ sung — Bài tập 4 · Mục 3.6.3}

\subsection*{Kiến thức và trình tự giải}
\textbf{Bước 1 — Chọn Kirchhoff (3D) hoặc Poisson (2D).} Với li độ đầu bằng không, vận tốc $\psi$, tốc độ $c$, đặt $\rho=ct$. Trong ba chiều,


\[
u(P,t)=\frac t{4\pi}\int_{S^2}\psi(P+\rho\omega)d\omega.
\]


Trong hai chiều,


\[
u(P,t)=\frac1{2\pi c}\int_{|Q-P|<ct}\frac{\psi(Q)}{\sqrt{c^2t^2-|Q-P|^2}}dQ.
\]


Với li độ $f$ khác không, thêm đạo hàm theo $t$ của chính biểu thức tương ứng khi thay $\psi$ bằng $f$. Với nguồn và dữ liệu đầu không, tích phân theo thời gian của nghiệm có vận tốc là nguồn tại từng thời điểm (Duhamel).

\textbf{Bước 2 — Tận dụng hình học và đối xứng.} Chỉ thị miền $D$ biến Kirchhoff thành $t$ nhân tỷ lệ diện tích mặt cầu nằm trong $D$. Chọn trục qua tâm miền và $P$, đặt $\mu=\cos\theta$, $d\omega=d\mu\,d\phi$. Nếu điều kiện không phụ thuộc $\phi$, tỷ lệ là $|E|/2$ với $E\subset[-1,1]$. Mỗi nửa không gian cho một bất đẳng thức tuyến tính theo $\mu$; quả cầu cho một ngưỡng chỏm. Miền hợp/rút bỏ cho tổng/hiệu nếu các phần không chồng diện tích.

\textbf{Bước 3 — Viết chính xác các khoảng hợp lệ.} Giải bất đẳng thức gốc, sau đó mới bình phương khi có căn. Chia $[-1,1]$ ở nghiệm phương trình biên, kiểm tra dấu ở trung điểm từng khoảng, loại nghiệm ngoại lai. Nếu còn góc phương vị, dùng độ dài tập $\cos\phi<v$ rồi còn một tích phân theo $\mu$. Với Poisson, tích phân theo một tọa độ cho hiệu arcsin; cắt các cận vào đĩa trước khi tính.

\textbf{Bước 4 — Tính nguồn bằng nguyên hàm thời gian.} Với quả cầu bán kính $R$ cách $P$ khoảng $r>0$, đặt


\[
K_R(r,t)=t\left[\frac{R^2-(r-ct)^2}{4rct}\right]_0^1,\quad J_R(r,t)=\int_0^tK_R(r,s)ds.
\]


Tách ba giai đoạn: cả mặt cầu nằm trong miền, giao một chỏm, không giao. Nguồn bật trên $[0,T]$ cho $J_R(r,t)-J_R(r,(t-T)_+)$. Trường hợp $r=0$ tính trực tiếp, tránh chia cho $r$.

\textbf{Bước 5 — Kiểm tra chuẩn hóa, mặt sóng và cực đại.} Nếu $\psi=1$ trên toàn không gian thì $u=t$, dùng để kiểm tra hệ số. Với dữ liệu bước nhảy, chỉ kết luận cổ điển ngoài các mặt sóng. Cực đại phải xét cả điểm chuyển nhánh và giá trị một phía; nêu rõ supremum nếu chỉ đạt theo giới hạn. Nếu dùng tích phân số, ghi tham số và cách kiểm tra hội tụ; đồ thị không tự chứng minh một cực đại toàn cục.
\subsection*{Lời giải áp dụng cho bài}
\textbf{Bước 1. Kiểm tra hai dữ liệu trước khi dùng công thức in.} Tại $t=0$, công thức $\psi+t\varphi$ trong sách cho li độ $\psi$ và vận tốc $\varphi$, đảo với dữ liệu đã ghi. Vì thế cần đổi lại $\varphi+t\psi$. Đặt $I(t)=\int_0^t(t-s)g(s)ds$, ta có $I(0)=I'(0)=0$, $I''=g$.

\textbf{Bước 2. Trường hợp điều hòa.} Mọi Laplace của $f,\varphi,\psi$ đều bằng không, nên nghiệm không có đóng góp $c^2\Delta u$; chỉ cần giải ODE thời gian tại mỗi điểm. Trường hợp $k$-điều hòa dùng chuỗi Taylor chẵn/lẻ theo thời gian và chuỗi Duhamel, nhưng kết thúc vì $\Delta^k$ của từng dữ liệu bằng không.

\textbf{Bước 3. Nghiệm sửa đúng và chuỗi hữu hạn.}

(a) Với dữ liệu đầu $u(0)=\varphi,u_t(0)=\psi$, nghiệm đúng là
\[u=\varphi+t\psi+f\int_0^t(t-s)g(s)ds.\]
Do các hàm không gian điều hòa, $\Delta u=0$ và $u_{tt}=fg$. Dòng công thức trong đề đảo $\varphi,\psi$ nên không thỏa dữ liệu đã ghi; cần sửa như trên.
(b) Với $\Delta^kf=\Delta^k\varphi=\Delta^k\psi=0$, chuỗi hữu hạn:
\begin{align*}
u=\sum_{j=0}^{k-1}c^{2j}\bigg[&\frac{t^{2j}}{(2j)!}\Delta^j\varphi+\frac{t^{2j+1}}{(2j+1)!}\Delta^j\psi\\
&+\Delta^jf\int_0^t\frac{(t-s)^{2j+1}}{(2j+1)!}g(s)ds\bigg].
\end{align*}
Đạo hàm hai lần làm dịch chỉ số; hạng cuối hết vì $\Delta^k=0$, chứng minh phương trình và dữ liệu.

\textbf{Bước 4. Đối chiếu và kết luận.} Để kiểm tra chuỗi ở (b), đạo hàm hai lần hạng li độ chuyển $t^{2j}/(2j)!$ thành $t^{2j-2}/(2j-2)!$; Laplace hạng trước sinh cùng $\Delta^j\varphi$. Tương tự cho vận tốc và tích phân nguồn. Chỉ hạng nguồn $j=0$ sinh $fg$, các hạng cuối mất vì $\Delta^k=0$. Tại $0$ chỉ hai hạng $j=0$ có thể cho li độ/vận tốc, nên đúng $\varphi,\psi$.

\clearpage\noindent\texttt{\detokenize{BT_3_6_3_05}}\par
\section*{Lời giải bổ sung — Bài tập 5 · Mục 3.6.3}

\subsection*{Kiến thức và trình tự giải}
\textbf{Bước 1 — Chọn Kirchhoff (3D) hoặc Poisson (2D).} Với li độ đầu bằng không, vận tốc $\psi$, tốc độ $c$, đặt $\rho=ct$. Trong ba chiều,


\[
u(P,t)=\frac t{4\pi}\int_{S^2}\psi(P+\rho\omega)d\omega.
\]


Trong hai chiều,


\[
u(P,t)=\frac1{2\pi c}\int_{|Q-P|<ct}\frac{\psi(Q)}{\sqrt{c^2t^2-|Q-P|^2}}dQ.
\]


Với li độ $f$ khác không, thêm đạo hàm theo $t$ của chính biểu thức tương ứng khi thay $\psi$ bằng $f$. Với nguồn và dữ liệu đầu không, tích phân theo thời gian của nghiệm có vận tốc là nguồn tại từng thời điểm (Duhamel).

\textbf{Bước 2 — Tận dụng hình học và đối xứng.} Chỉ thị miền $D$ biến Kirchhoff thành $t$ nhân tỷ lệ diện tích mặt cầu nằm trong $D$. Chọn trục qua tâm miền và $P$, đặt $\mu=\cos\theta$, $d\omega=d\mu\,d\phi$. Nếu điều kiện không phụ thuộc $\phi$, tỷ lệ là $|E|/2$ với $E\subset[-1,1]$. Mỗi nửa không gian cho một bất đẳng thức tuyến tính theo $\mu$; quả cầu cho một ngưỡng chỏm. Miền hợp/rút bỏ cho tổng/hiệu nếu các phần không chồng diện tích.

\textbf{Bước 3 — Viết chính xác các khoảng hợp lệ.} Giải bất đẳng thức gốc, sau đó mới bình phương khi có căn. Chia $[-1,1]$ ở nghiệm phương trình biên, kiểm tra dấu ở trung điểm từng khoảng, loại nghiệm ngoại lai. Nếu còn góc phương vị, dùng độ dài tập $\cos\phi<v$ rồi còn một tích phân theo $\mu$. Với Poisson, tích phân theo một tọa độ cho hiệu arcsin; cắt các cận vào đĩa trước khi tính.

\textbf{Bước 4 — Tính nguồn bằng nguyên hàm thời gian.} Với quả cầu bán kính $R$ cách $P$ khoảng $r>0$, đặt


\[
K_R(r,t)=t\left[\frac{R^2-(r-ct)^2}{4rct}\right]_0^1,\quad J_R(r,t)=\int_0^tK_R(r,s)ds.
\]


Tách ba giai đoạn: cả mặt cầu nằm trong miền, giao một chỏm, không giao. Nguồn bật trên $[0,T]$ cho $J_R(r,t)-J_R(r,(t-T)_+)$. Trường hợp $r=0$ tính trực tiếp, tránh chia cho $r$.

\textbf{Bước 5 — Kiểm tra chuẩn hóa, mặt sóng và cực đại.} Nếu $\psi=1$ trên toàn không gian thì $u=t$, dùng để kiểm tra hệ số. Với dữ liệu bước nhảy, chỉ kết luận cổ điển ngoài các mặt sóng. Cực đại phải xét cả điểm chuyển nhánh và giá trị một phía; nêu rõ supremum nếu chỉ đạt theo giới hạn. Nếu dùng tích phân số, ghi tham số và cách kiểm tra hội tụ; đồ thị không tự chứng minh một cực đại toàn cục.
\subsection*{Lời giải áp dụng cho bài}
\textbf{Bước 1. Duhamel của nguồn cầu không đổi theo thời gian.} Li độ/vận tốc đầu bằng không, tốc độ $1$. Tại mỗi thời điểm nguồn $s$, vận tốc xung là $\chi_{B_1}$, nên đóng góp ở thời gian trễ $\tau=t-s$ là $K_1(r,\tau)$. Tích phân $s$ cho $u=J_1(r,t)$.

\textbf{Bước 2. Tính nguyên hàm ở từng giai đoạn.} Với cầu bán kính $R$, nếu $r<R$, trước $R-r$ có $K_R(s)=s$; nếu $r\ge R$, trước $r-R$ có $K_R=0$. Trong $|R-r|<s<R+r$, $K_R=(R^2-(r-s)^2)/(4r)$ là đa thức bậc hai. Sau $R+r$ đóng góp xung bằng không, nên $J_R$ giữ hằng số. Tích phân ba nhánh cho công thức dưới đây.

\textbf{Bước 3. Nghiệm và trạng thái dừng.}

Đặt $r=|P|$, $J_R(r,t)=\int_0^tK_R(r,s)ds$ với $c=1$. Duhamel cho $u=J_1(r,t)$. Khi $r>0$, đặt $b=|R-r|$, $q=\min(\max(t,b),R+r)$,
\[A(s)=\frac{(R^2-r^2)s+rs^2-s^3/3}{4r},\quad J_R=\frac12\min(t,(R-r)_+)^2+A(q)-A(b).\]
Với $r=0$, $J_R(0,t)=\min(t,R)^2/2$. Khi $t\ge r+R$, nghiệm bằng thế Newton của quả cầu: $R^3/(3r)$ nếu $r\ge R$, $(3R^2-r^2)/6$ nếu $r\le R$. Kiểm tra $\Delta u=-\chi_{r<R}$ cho trạng thái dừng.

\textbf{Bước 4. Đối chiếu và kết luận.} Với $R=1$, trạng thái dừng là $(3-r^2)/6$ bên trong và $1/(3r)$ bên ngoài. Công thức Laplace hướng tâm $v_{rr}+2v_r/r$ cho $-1$ bên trong và $0$ bên ngoài; giá trị cùng đạo hàm nối tại $r=1$, nên không sinh nguồn mặt cầu. Hai dữ liệu đầu của $J$ bằng $0$ vì $J(0)=0$ và $J_t=K$, $K(r,0)=0$.

\textbf{Chi tiết nguyên hàm và trạng thái dừng.} Trong nhánh chỏm,
\[K_R(r,s)=\frac{R^2-r^2+2rs-s^2}{4r},\quad
\int K_Rds=\frac{(R^2-r^2)s+rs^2-s^3/3}{4r}=A(s).\]
Nếu $r\ge R$, trước $b=r-R$ tích phân bằng $0$, và
\[J_\infty=A(r+R)-A(r-R)
=\frac{2R(R^2-r^2)+4r^2R-(2r^2R+2R^3/3)}{4r}
=\frac{R^3}{3r}.\]
Nếu $r<R$, trước $b=R-r$ tích phân bằng $(R-r)^2/2$. Tính hiệu:
\[A(R+r)-A(R-r)=Rr-\frac23r^2,\quad
J_\infty=\frac{(R-r)^2}{2}+Rr-\frac23r^2=\frac{3R^2-r^2}{6}.\]
Tại $r=R$, hai giá trị bằng $R^2/3$; đạo hàm trong $-r/3$ và ngoài $-R^3/(3r^2)$ cùng bằng $-R/3$.

\clearpage\noindent\texttt{\detokenize{BT_3_6_3_06}}\par
\section*{Lời giải bổ sung — Bài tập 6 · Mục 3.6.3}

\subsection*{Kiến thức và trình tự giải}
\textbf{Bước 1 — Chọn Kirchhoff (3D) hoặc Poisson (2D).} Với li độ đầu bằng không, vận tốc $\psi$, tốc độ $c$, đặt $\rho=ct$. Trong ba chiều,


\[
u(P,t)=\frac t{4\pi}\int_{S^2}\psi(P+\rho\omega)d\omega.
\]


Trong hai chiều,


\[
u(P,t)=\frac1{2\pi c}\int_{|Q-P|<ct}\frac{\psi(Q)}{\sqrt{c^2t^2-|Q-P|^2}}dQ.
\]


Với li độ $f$ khác không, thêm đạo hàm theo $t$ của chính biểu thức tương ứng khi thay $\psi$ bằng $f$. Với nguồn và dữ liệu đầu không, tích phân theo thời gian của nghiệm có vận tốc là nguồn tại từng thời điểm (Duhamel).

\textbf{Bước 2 — Tận dụng hình học và đối xứng.} Chỉ thị miền $D$ biến Kirchhoff thành $t$ nhân tỷ lệ diện tích mặt cầu nằm trong $D$. Chọn trục qua tâm miền và $P$, đặt $\mu=\cos\theta$, $d\omega=d\mu\,d\phi$. Nếu điều kiện không phụ thuộc $\phi$, tỷ lệ là $|E|/2$ với $E\subset[-1,1]$. Mỗi nửa không gian cho một bất đẳng thức tuyến tính theo $\mu$; quả cầu cho một ngưỡng chỏm. Miền hợp/rút bỏ cho tổng/hiệu nếu các phần không chồng diện tích.

\textbf{Bước 3 — Viết chính xác các khoảng hợp lệ.} Giải bất đẳng thức gốc, sau đó mới bình phương khi có căn. Chia $[-1,1]$ ở nghiệm phương trình biên, kiểm tra dấu ở trung điểm từng khoảng, loại nghiệm ngoại lai. Nếu còn góc phương vị, dùng độ dài tập $\cos\phi<v$ rồi còn một tích phân theo $\mu$. Với Poisson, tích phân theo một tọa độ cho hiệu arcsin; cắt các cận vào đĩa trước khi tính.

\textbf{Bước 4 — Tính nguồn bằng nguyên hàm thời gian.} Với quả cầu bán kính $R$ cách $P$ khoảng $r>0$, đặt


\[
K_R(r,t)=t\left[\frac{R^2-(r-ct)^2}{4rct}\right]_0^1,\quad J_R(r,t)=\int_0^tK_R(r,s)ds.
\]


Tách ba giai đoạn: cả mặt cầu nằm trong miền, giao một chỏm, không giao. Nguồn bật trên $[0,T]$ cho $J_R(r,t)-J_R(r,(t-T)_+)$. Trường hợp $r=0$ tính trực tiếp, tránh chia cho $r$.

\textbf{Bước 5 — Kiểm tra chuẩn hóa, mặt sóng và cực đại.} Nếu $\psi=1$ trên toàn không gian thì $u=t$, dùng để kiểm tra hệ số. Với dữ liệu bước nhảy, chỉ kết luận cổ điển ngoài các mặt sóng. Cực đại phải xét cả điểm chuyển nhánh và giá trị một phía; nêu rõ supremum nếu chỉ đạt theo giới hạn. Nếu dùng tích phân số, ghi tham số và cách kiểm tra hội tụ; đồ thị không tự chứng minh một cực đại toàn cục.
\subsection*{Lời giải áp dụng cho bài}
\textbf{Bước 1. Chồng nghiệm vận tốc và nguồn.} Hai cầu vận tốc có dấu trái nhau cho $K_1(d_+,t)-K_1(d_-,t)$. Nguồn chỉ bật khi $0<s<T$, nên phần nguồn là $\int_0^{\min(t,T)}K_1(r,t-s)ds$. Đổi $\tau=t-s$ cho hiệu hai nguyên hàm $J$.

\textbf{Bước 2. Dùng đối xứng và thời gian yêu cầu.} Ở $x=0$, khoảng cách tới hai tâm $(\pm2,0,0)$ bằng nhau, nên phần vận tốc triệt tiêu hoàn toàn; câu (a) chỉ còn nguồn với $T=1$. Câu (b) giữ hai khoảng cách khác nhau và thay $t=r$ vào biểu thức, không thay $d_\pm$ bằng $r$.

\textbf{Bước 3. Công thức áp dụng.}

Gọi $d_\pm=|P-(\pm2,0,0)|$, $r=|P|$. Chồng Kirchhoff và Duhamel cho
\[u=K_1(d_+,t)-K_1(d_-,t)+J_1(r,t)-J_1(r,(t-T)_+),\]
với $c=1$, $K,J$ có biểu thức chỏm và nguyên hàm ở hai bài trước.
(a) $P=(0,y,z)$ cho $d_+=d_-$, nên $u=J_1(\sqrt{y^2+z^2},t)-J_1(\sqrt{y^2+z^2},(t-1)_+)$.
(b) Thay $t=r$: $u=K_1(d_+,r)-K_1(d_-,r)+J_1(r,r)-J_1(r,(r-T)_+)$. Đây là biểu thức tường minh bằng min/max và đa thức, áp dụng mọi $P$ kể cả $r=0$ theo giới hạn.

\textbf{Bước 4. Đối chiếu và kết luận.} Tại $t=0$, cả $J$ và $K$ bằng $0$; đạo hàm của phần $K$ cho đúng hiệu hai chỉ thị cầu và phần $J$ có vận tốc đầu không. Với $t=r=0$, kết quả là $0$ và dùng nhánh tâm của $J$ để tránh chia $r$. Các min/max của $K,J$ nêu trong phần kiến thức và Ví dụ 3.6.2.5 làm biểu thức câu (b) hoàn toàn tường minh.

\textbf{Chi tiết các hàm dùng trong biểu thức tường minh.} Với tốc độ $1$, bán kính $1$, $d>0,t>0$,
\[K_1(d,t)=t\left[\frac{1-(d-t)^2}{4dt}\right]_0^1.\]
Với $r>0$, đặt $b=|1-r|$, $q=[t]_b^{r+1}$ và $A(s)=((1-r^2)s+rs^2-s^3/3)/(4r)$:
\[J_1(r,t)=\tfrac12\min(t,(1-r)_+)^2+A(q)-A(b).\]
Tại tâm, $K_1(0,t)=t$ khi $t<1$, bằng $0$ khi $t>1$, $J_1(0,t)=\min(t,1)^2/2$; tại $t=0$, $K_1=J_1=0$. Khi đổi $\tau=t-s$, $s$ chạy $0$ đến $\min(t,T)$ nên $\tau$ chạy ngược từ $t$ đến $(t-T)_+$:
\[\int_0^{\min(t,T)}K_1(r,t-s)ds=\int_{(t-T)_+}^{t}K_1(r,\tau)d\tau=J_1(r,t)-J_1(r,(t-T)_+).\]

\clearpage\noindent\texttt{\detokenize{BT_3_6_3_07}}\par
\section*{Lời giải bổ sung — Bài tập 7 · Mục 3.6.3}

\subsection*{Kiến thức và trình tự giải}
\textbf{Bước 1 — Chọn Kirchhoff (3D) hoặc Poisson (2D).} Với li độ đầu bằng không, vận tốc $\psi$, tốc độ $c$, đặt $\rho=ct$. Trong ba chiều,


\[
u(P,t)=\frac t{4\pi}\int_{S^2}\psi(P+\rho\omega)d\omega.
\]


Trong hai chiều,


\[
u(P,t)=\frac1{2\pi c}\int_{|Q-P|<ct}\frac{\psi(Q)}{\sqrt{c^2t^2-|Q-P|^2}}dQ.
\]


Với li độ $f$ khác không, thêm đạo hàm theo $t$ của chính biểu thức tương ứng khi thay $\psi$ bằng $f$. Với nguồn và dữ liệu đầu không, tích phân theo thời gian của nghiệm có vận tốc là nguồn tại từng thời điểm (Duhamel).

\textbf{Bước 2 — Tận dụng hình học và đối xứng.} Chỉ thị miền $D$ biến Kirchhoff thành $t$ nhân tỷ lệ diện tích mặt cầu nằm trong $D$. Chọn trục qua tâm miền và $P$, đặt $\mu=\cos\theta$, $d\omega=d\mu\,d\phi$. Nếu điều kiện không phụ thuộc $\phi$, tỷ lệ là $|E|/2$ với $E\subset[-1,1]$. Mỗi nửa không gian cho một bất đẳng thức tuyến tính theo $\mu$; quả cầu cho một ngưỡng chỏm. Miền hợp/rút bỏ cho tổng/hiệu nếu các phần không chồng diện tích.

\textbf{Bước 3 — Viết chính xác các khoảng hợp lệ.} Giải bất đẳng thức gốc, sau đó mới bình phương khi có căn. Chia $[-1,1]$ ở nghiệm phương trình biên, kiểm tra dấu ở trung điểm từng khoảng, loại nghiệm ngoại lai. Nếu còn góc phương vị, dùng độ dài tập $\cos\phi<v$ rồi còn một tích phân theo $\mu$. Với Poisson, tích phân theo một tọa độ cho hiệu arcsin; cắt các cận vào đĩa trước khi tính.

\textbf{Bước 4 — Tính nguồn bằng nguyên hàm thời gian.} Với quả cầu bán kính $R$ cách $P$ khoảng $r>0$, đặt


\[
K_R(r,t)=t\left[\frac{R^2-(r-ct)^2}{4rct}\right]_0^1,\quad J_R(r,t)=\int_0^tK_R(r,s)ds.
\]


Tách ba giai đoạn: cả mặt cầu nằm trong miền, giao một chỏm, không giao. Nguồn bật trên $[0,T]$ cho $J_R(r,t)-J_R(r,(t-T)_+)$. Trường hợp $r=0$ tính trực tiếp, tránh chia cho $r$.

\textbf{Bước 5 — Kiểm tra chuẩn hóa, mặt sóng và cực đại.} Nếu $\psi=1$ trên toàn không gian thì $u=t$, dùng để kiểm tra hệ số. Với dữ liệu bước nhảy, chỉ kết luận cổ điển ngoài các mặt sóng. Cực đại phải xét cả điểm chuyển nhánh và giá trị một phía; nêu rõ supremum nếu chỉ đạt theo giới hạn. Nếu dùng tích phân số, ghi tham số và cách kiểm tra hội tụ; đồ thị không tự chứng minh một cực đại toàn cục.
\subsection*{Lời giải áp dụng cho bài}
\textbf{Bước 1. Tích phân Poisson theo tung độ trước.} Với $c=2$, tại hoành độ $X$, đĩa ảnh hưởng có tung độ $P_y\pm h(X)$. Một dải $l<Y<v$ góp nguyên hàm $\arcsin((Y-P_y)/h)$, sau khi cắt hai cận vào $[-1,1]$. Các miền trong đề có cận tung độ tường minh theo $|X|$, nên chỉ còn một tích phân theo $X$.

\textbf{Bước 2. Hai đối xứng cần dùng.} Câu (a): trước $2t<98/\sqrt5$, cả đĩa quanh $(0,100)$ nằm trong miền $D_1$, vận tốc bằng $2$, nên $u=2t$. Câu (b): $\psi(Y,X)=-\psi(X,Y)$ và điểm quan sát $(100,100)$ bất biến khi đổi tọa độ. Đĩa và nhân Poisson cũng bất biến, nên tích phân đổi thành đối của chính nó: $u(100,100,t)=0$ với mọi $t$. Công thức tích phân dưới đây áp dụng cả cho các điểm khác khi vẽ đồ thị, nhưng tại điểm của câu (b) kết quả rút gọn là $0$.

\textbf{Bước 3. Tính các dải và dựng đồ thị.}

Poisson hai chiều: $u(P,t)=(2\pi c)^{-1}\int_{|Q-P|<ct}\psi(Q)/\sqrt{c^2t^2-|Q-P|^2}\,dQ$. Tích phân theo tung độ được thực hiện chính xác. Đặt $h(X)=\sqrt{c^2t^2-(X-P_x)^2}$,
\[B_X(l,v)=\arcsin\left[\frac{v-P_y}{h(X)}\right]_{-1}^1-\arcsin\left[\frac{l-P_y}{h(X)}\right]_{-1}^1.\]
(a) Với $P=(0,100),c=2$,
\[u=\frac1{4\pi}\int_{-2t}^{2t}\{2B_X(2+2|X|,\infty)-B_X(-\infty,\min(2-2|X|,0))+B_X(\max(2-2|X|,0),2+2|X|)\}dX.\]
(b) Với $P=(100,100),c=2$, tích phân $X\in[100-2t,100+2t]$: nếu $|X|>1$, tích phân tung độ là $B_X(-1,1)$; nếu $|X|<1$, là $-B_X(-\infty,-1)-B_X(1,\infty)$. Nhân $1/(4\pi)$. Đây là tích phân một chiều của hàm sơ cấp; không có kỳ dị căn ở tung độ nữa. Ký hiệu $D_1$ thứ hai trong đề là $D_2$.
(c) Đồ thị kèm theo được tính từ tích phân Poisson; mã cho phép đổi điểm/thời gian. Dữ liệu chỉ thị nên hiểu nghiệm yếu và sự gãy ở biên vùng ảnh hưởng.

\textbf{Bước 4. Đối chiếu và kết luận.} Với vận tốc hằng $1$, tích phân hai chiều của nhân bằng $2\pi ct$, cho $u=t$, kiểm tra hệ số $1/(4\pi)$ khi $c=2$. Đồ thị bề mặt dùng tích phân số, phải đọc tham số lưới trong mã; đường $x=y$ của câu (b) luôn bằng $0$ theo đối xứng, là một kiểm tra độc lập cho đồ thị.

\textbf{Chi tiết loại tích phân tung độ và kiểm tra hệ số.} Tại $X$ cố định, đặt $h^2=c^2t^2-(X-P_x)^2$ và $z=(Y-P_y)/h$. Khi $h>0$,
\[\int\frac{dY}{\sqrt{h^2-(Y-P_y)^2}}
=\int\frac{dz}{\sqrt{1-z^2}}=\arcsin z.\]
Chỉ lấy giao của dải với $P_y-h<Y<P_y+h$, tương đương cắt $z$ vào $[-1,1]$, cho $B_X(l,v)$. Với vận tốc hằng $1$, dùng tọa độ cực quanh $P$:
\[\int_{|Q-P|<ct}\frac{dQ}{\sqrt{c^2t^2-|Q-P|^2}}
=2\pi\int_0^{ct}\frac{r\,dr}{\sqrt{c^2t^2-r^2}}
=2\pi ct.\]
Nhân $1/(2\pi c)$ cho $u=t$. Ở (a), khoảng cách từ $(0,100)$ đến hai đường $Y=2\pm2X$ là $|100-2|/\sqrt{1+4}=98/\sqrt5$, nên trước $2t<98/\sqrt5$ cả đĩa có vận tốc $2$.
\begin{center}\includegraphics[width=.97\linewidth,height=11cm,keepaspectratio]{_xu_ly/graphs/BT_3_6_3_07.png}\end{center}
\textit{Đồ thị minh họa; với nghiệm số, xem phương pháp và tham số trong mã Python kèm theo.}

\clearpage\noindent\texttt{\detokenize{BT_3_6_3_08}}\par
\section*{Lời giải bổ sung — Bài tập 8 · Mục 3.6.3}

\subsection*{Kiến thức và trình tự giải}
\textbf{Bước 1 — Chọn Kirchhoff (3D) hoặc Poisson (2D).} Với li độ đầu bằng không, vận tốc $\psi$, tốc độ $c$, đặt $\rho=ct$. Trong ba chiều,


\[
u(P,t)=\frac t{4\pi}\int_{S^2}\psi(P+\rho\omega)d\omega.
\]


Trong hai chiều,


\[
u(P,t)=\frac1{2\pi c}\int_{|Q-P|<ct}\frac{\psi(Q)}{\sqrt{c^2t^2-|Q-P|^2}}dQ.
\]


Với li độ $f$ khác không, thêm đạo hàm theo $t$ của chính biểu thức tương ứng khi thay $\psi$ bằng $f$. Với nguồn và dữ liệu đầu không, tích phân theo thời gian của nghiệm có vận tốc là nguồn tại từng thời điểm (Duhamel).

\textbf{Bước 2 — Tận dụng hình học và đối xứng.} Chỉ thị miền $D$ biến Kirchhoff thành $t$ nhân tỷ lệ diện tích mặt cầu nằm trong $D$. Chọn trục qua tâm miền và $P$, đặt $\mu=\cos\theta$, $d\omega=d\mu\,d\phi$. Nếu điều kiện không phụ thuộc $\phi$, tỷ lệ là $|E|/2$ với $E\subset[-1,1]$. Mỗi nửa không gian cho một bất đẳng thức tuyến tính theo $\mu$; quả cầu cho một ngưỡng chỏm. Miền hợp/rút bỏ cho tổng/hiệu nếu các phần không chồng diện tích.

\textbf{Bước 3 — Viết chính xác các khoảng hợp lệ.} Giải bất đẳng thức gốc, sau đó mới bình phương khi có căn. Chia $[-1,1]$ ở nghiệm phương trình biên, kiểm tra dấu ở trung điểm từng khoảng, loại nghiệm ngoại lai. Nếu còn góc phương vị, dùng độ dài tập $\cos\phi<v$ rồi còn một tích phân theo $\mu$. Với Poisson, tích phân theo một tọa độ cho hiệu arcsin; cắt các cận vào đĩa trước khi tính.

\textbf{Bước 4 — Tính nguồn bằng nguyên hàm thời gian.} Với quả cầu bán kính $R$ cách $P$ khoảng $r>0$, đặt


\[
K_R(r,t)=t\left[\frac{R^2-(r-ct)^2}{4rct}\right]_0^1,\quad J_R(r,t)=\int_0^tK_R(r,s)ds.
\]


Tách ba giai đoạn: cả mặt cầu nằm trong miền, giao một chỏm, không giao. Nguồn bật trên $[0,T]$ cho $J_R(r,t)-J_R(r,(t-T)_+)$. Trường hợp $r=0$ tính trực tiếp, tránh chia cho $r$.

\textbf{Bước 5 — Kiểm tra chuẩn hóa, mặt sóng và cực đại.} Nếu $\psi=1$ trên toàn không gian thì $u=t$, dùng để kiểm tra hệ số. Với dữ liệu bước nhảy, chỉ kết luận cổ điển ngoài các mặt sóng. Cực đại phải xét cả điểm chuyển nhánh và giá trị một phía; nêu rõ supremum nếu chỉ đạt theo giới hạn. Nếu dùng tích phân số, ghi tham số và cách kiểm tra hội tụ; đồ thị không tự chứng minh một cực đại toàn cục.
\subsection*{Lời giải áp dụng cho bài}
\textbf{Bước 1. Hạ về hai bài một chiều có nguồn.} Vì $F=\chi_{[0,1]}(x)+\chi_{[0,1]}(y)$, đặt $u=U(x,t)+U(y,t)$. Mỗi $U$ có dữ liệu không, giải $U_{tt}-U_{xx}=\chi_{[0,1]}$.

\textbf{Bước 2. Tích phân hai lần hàm chỉ thị.} $H(s)=s_+-(s-1)_+$ là nguyên hàm liên tục của chỉ thị, và $Q(s)=(s_+^2-(s-1)_+^2)/2$ thỏa $Q''=\chi_{(0,1)}$. Duhamel cho $U=\frac12\int_0^t[H(x+\tau)-H(x-\tau)]d\tau=\frac12[Q(x+t)+Q(x-t)-2Q(x)]$. Khai triển $Q$ cho công thức dưới đây.

\textbf{Bước 3. Nghiệm và kiểm tra.}

Nguồn tách thành hai hàm chỉ phụ thuộc một tọa độ nên $u=U(x,t)+U(y,t)$. Duhamel một chiều với nguồn $\chi_{[0,1]}$ cho
\[U(x,t)=\frac14\{(x+t)_+^2-2x_+^2+(x-t)_+^2-(x-1+t)_+^2+2(x-1)_+^2-(x-1-t)_+^2\}.\]
Đạo hàm hai lần theo $t,x$ cho $U_{tt}-U_{xx}=\chi_{(0,1)}(x)$; $U(0)=U_t(0)=0$ theo thời gian. Do đó tổng đúng phương trình hai chiều.

\textbf{Bước 4. Đối chiếu và kết luận.} Từ $U=\frac12(Q(x+t)+Q(x-t)-2Q(x))$, có $U_{tt}-U_{xx}=Q''(x)=\chi_{(0,1)}(x)$. Tại $t=0$ các hạng triệt tiêu và $U_t=0$. Vì tổng không có đạo hàm chéo giữa $x,y$, PDE hai chiều có đúng hai nguồn. Công thức đúng theo nghĩa phân phối cả tại các đường cận chạm $0,1$.

\clearpage\noindent\texttt{\detokenize{BT_3_6_3_09}}\par
\section*{Lời giải bổ sung — Bài tập 9 · Mục 3.6.3}

\subsection*{Kiến thức và trình tự giải}
\textbf{Bước 1 — Giữ nhất quán dấu toán tử.} Đặt $q=x^2+y^2-t^2$, $\Box=\partial_{xx}+\partial_{yy}-\partial_{tt}$. Phương trình $u_{tt}-\Delta u=\Phi$ là $-\Box u=\Phi$. Dữ liệu trên nón $q=0$ gợi ý mỗi hạng của nghiệm chứa ít nhất một thừa số $q$.

\textbf{Bước 2 — Tính quy tắc đạo hàm tích.} Với $H$ thuần nhất bậc $m$, dùng Euler $xH_x+yH_y+tH_t=mH$ để có $\Box(q^jH)=q^j\Box H+2j(2m+2j+1)q^{j-1}H$. Tính trực tiếp $\Box q=6$ để kiểm tra hệ số.

\textbf{Bước 3 — Lập truy hồi hữu hạn.} Viết $u=\sum a_kq^{k+1}\Box^k\Phi$. Vì $\Phi$ bậc $d-2$, hệ số khi áp dụng $\Box$ là $C_k=2(k+1)(2d-2k-1)$. Chọn $a_0=-1/C_0$, $a_k=-a_{k-1}/C_k$ để triệt tiêu các hạng còn lại. Các đạo hàm lặp của đa thức hết sau hữu hạn bước.

\textbf{Bước 4 — Chứng minh duy nhất bằng năng lượng trên miền co.} Với hiệu $v$ của hai nghiệm, đặt $E(t)=\frac12\int_{x^2+y^2<t^2}(v_t^2+|\nabla v|^2)dxdy$ cho $t<0$. Đạo hàm cả tích phân lẫn biên chuyển động. Vết $v(r,-r)=0$ cho $v_t=v_r$ và đạo hàm tiếp tuyến góc bằng không, nên thông lượng năng lượng bằng không. Nếu nghiệm đủ trơn đến đỉnh, $E(t)\to0$ khi $t\uparrow0$; do đó $E=0$, suy $v=0$ trong nón.

\textbf{Bước 5 — Thế lại nghiệm cụ thể.} Kiểm tra cả dấu PDE và $u|_{q=0}=0$. Với đa thức nguồn cụ thể, tính $\Box\Phi$ trước: nếu bằng không thì chỉ còn hạng $a_0q\Phi$.
\subsection*{Lời giải áp dụng cho bài}
\textbf{Bước 1. Đạo hàm tích và Euler.} $q=x^2+y^2-t^2$, $\Box=\Delta-\partial_{tt}$, nên $\Box q=6$. Với $H$ thuần nhất bậc $m$, hạng đạo hàm chéo bằng $4j q^{j-1}(xH_x+yH_y+tH_t)=4jm q^{j-1}H$. Cộng $\Box q^j=2j(2j+1)q^{j-1}$ cho công thức tích dưới đây.

\textbf{Bước 2. Lập hệ số thay vì chép dấu in.} Đặt $u=\sum a_kq^{k+1}\Box^k\Phi$. Hệ số hạng hạ bậc là $C_k=2(k+1)(2d-2k-1)$. Vì PDE là $-\Box u=\Phi$, cần $a_0=-1/C_0$; các hạng sau triệt tiêu khi $a_k=-a_{k-1}/C_k$. Suy biểu thức giai thừa kép, có dấu âm phía ngoài và số mũ $q^{k+1}$, không phải $q^{2k+2}$.

\textbf{Bước 3. Công thức đúng, chứng minh duy nhất và nguồn cụ thể.}

Đặt $q=x^2+y^2-t^2$, $\Box=\Delta_{x,y}-\partial_{tt}$; phương trình là $-\Box u=\Phi$. Với $H$ thuần nhất bậc $m$,
\[\Box(q^jH)=q^j\Box H+2j(2m+2j+1)q^{j-1}H.\]
Lấy $H=\Box^k\Phi$ bậc $d-2-2k$, $j=k+1$: hệ số là $2(k+1)(2d-2k-1)$. Quy nạp triệt tiêu các hạng cho
\[u=-\sum_{k\ge0}\frac{(-1)^k(2d-2k-3)!!}{(2k+2)!!(2d-1)!!}\,q^{k+1}\Box^k\Phi.\]
Tổng hữu hạn vì $\Phi$ đa thức. Đề có dấu dương và chỉ số giai thừa kép không nhất quán; công thức trên thỏa $u_{tt}-\Delta u=\Phi$ và $u=0$ trên nón.

(ii-a,b) Với hiệu hai nghiệm nguồn và vết nón 0, đặt năng lượng trên $B_{-t}$, $t<0$. Đạo hàm miền di động cho $E'=\int_{r=-t}[u_tu_r-(u_t^2+|\nabla u|^2)/2]$. Trên nón, vết 0 buộc đạo hàm tiếp tuyến $u_t=u_r$ và đạo hàm góc 0, nên thông lượng bằng 0. Nếu nghiệm đủ trơn đến đỉnh, $E(t)\to0$ khi $t\uparrow0$, vì diện tích miền co về 0. Vậy $E=0$, $u_t=\nabla u=0$ trên mọi mặt $t=-h$, suy $u=0$ và duy nhất trong nón. Điều kiện chính quy ở đỉnh cần thiết cho lập luận.

(ii-c) $\Phi=xyt$ có bậc 3 nên $d=5$, $\Box\Phi=0$:
\[u=-\frac{(x^2+y^2-t^2)xyt}{18}.\]
Kiểm tra $u_{tt}-u_{xx}-u_{yy}=xyt$ và vết nón 0, nên đây là nghiệm duy nhất trong lớp chính quy đã nêu.

\textbf{Bước 4. Đối chiếu và kết luận.} Đạo hàm biên chuyển động trong năng lượng ở (ii) có tốc độ bán kính $-1$, nên hạng biên mang dấu âm. Vết $u(r,-r)=0$ cho $u_r-u_t=0$; thế vào thông lượng triệt tiêu bình phương. Với $\Phi=xyt$, tính trực tiếp $u=-[x^3yt+xy^3t-xyt^3]/18$: $u_{tt}=xyt/3$, $u_{xx}=-xyt/3$, $u_{yy}=-xyt/3$, suy hiệu bằng $xyt$. Hệ số $-1/18$ và dấu đã được kiểm tra độc lập.

\textbf{Chi tiết đạo hàm tích và truy hồi hệ số.} $\nabla_{x,y}q=(2x,2y)$, $q_t=-2t$, nên $|\nabla q|^2-q_t^2=4q$ và $\Box q=6$. Quy tắc dây chuyền:
\[\Box q^j=j q^{j-1}\Box q+j(j-1)q^{j-2}(|\nabla q|^2-q_t^2)
=[6j+4j(j-1)]q^{j-1}=2j(2j+1)q^{j-1}.\]
Hạng chéo trong $\Box(q^jH)$ là
\[2(\nabla q^j\cdot\nabla H-(q^j)_tH_t)
=4jq^{j-1}(xH_x+yH_y+tH_t)=4jmq^{j-1}H.\]
Gọi $H_k=\Box^k\Phi$, $C_k=2(k+1)(2d-2k-1)$. Thế chuỗi:
\[\Box u=a_0C_0\Phi+\sum_{k\ge1}(a_{k-1}+C_ka_k)q^kH_k,\]
vì hạng cuối $H_{k+1}$ hết khi bậc đa thức âm. Muốn $\Box u=-\Phi$ cần $a_0C_0=-1$, $a_{k-1}+C_ka_k=0$. Do đó
\[a_k=\frac{(-1)^{k+1}}{C_0\cdots C_k}
=\frac{(-1)^{k+1}(2d-2k-3)!!}{(2k+2)!!(2d-1)!!},\]
vì tích các thừa số chẵn là $(2k+2)!!$, tích lẻ là $(2d-1)!!/(2d-2k-3)!!$.

\clearpage\noindent\texttt{\detokenize{VD_4_1_2_01}}\par
\section*{Lời giải bổ sung — Ví dụ 1 · Mục 4.1.2}

\subsection*{Kiến thức và trình tự giải}
\textbf{Bước 1 — Đọc hệ số và điều kiện biên.} Viết phương trình dưới dạng $u_t=Ku_{xx}$, $K>0$. Nếu đề ký hiệu tốc độ $a$ và viết $u_t=a^2u_{xx}$ thì thay $K=a^2$; không nhầm $a$ với $K$. Kiểm tra điều kiện đầu $f$ có khớp giá trị biên tại $t=0$ hay không. Sự không tương thích tại góc không ngăn nghiệm trơn với $t>0$, nhưng ngăn tính liên tục đến góc khi biên Dirichlet khác vết của $f$.

\textbf{Bước 2 — Tách biến và giải bài toán không gian.} Thế $u=XT$ vào PDE, chia cho $KXT$ để có $X''/X=T'/(KT)=-\lambda$. Do đó $X''+\lambda X=0$, $T'+K\lambda T=0$. Với $\lambda=k^2>0$, $X=A\cos kx+B\sin kx$; thay từng điều kiện biên để tìm $k$. Hai biên Dirichlet cho $X_n=\sin(n\pi x/L)$, $n\ge1$; hai biên Neumann cho $1$ và $\cos(n\pi x/L)$; biên Dirichlet–Neumann cho sine, biên Neumann–Dirichlet cho cosine, với $k_n=(n+1/2)\pi/L$, $n\ge0$. Trường hợp $\lambda=0$ phải giải riêng $X=A+Bx$; chỉ giữ mode hằng nếu biên cho phép.

\textbf{Bước 3 — Tính hệ số, không chỉ ghi tích phân.} Đặt $N_n=\int_0^L X_n^2dx$, $c_n=N_n^{-1}\int_0^L fX_n\,dx$. Với các mode sine/cosine chuẩn, $N_n=L/2$; mode hằng có $N_0=L$. Nếu $f$ là đa thức, tích phân từng phần đến hết; nếu $f$ chia đoạn, tách đúng cận; nếu $f$ lượng giác, dùng tích–tổng. Chẳng hạn với $k=n\pi/L$, $f(0)=f(L)=0$,


\[
\int_0^Lf\sin kx\,dx=\frac1k\int_0^Lf'\cos kx\,dx
=-\frac1{k^2}\int_0^Lf''\sin kx\,dx.
\]


Các hạng biên phải được tính trước khi bỏ. Khi tích–tổng tạo mẫu $k-m=0$, tính trực tiếp tích phân ở tần số đó.

\textbf{Bước 4 — Giải ODE và ghép nghiệm.} Nhân $T_n'+K\lambda_nT_n=0$ với $e^{K\lambda_nt}$, được $(e^{K\lambda_nt}T_n)'=0$. Vì $T_n(0)=c_n$,


\[
u(x,t)=\sum_n c_ne^{-K\lambda_nt}X_n(x).
\]


Thêm $c_0=L^{-1}\int_0^Lf$ cho Neumann hai đầu. Kiểm tra từng mode: $X_n''=-\lambda_nX_n$, nên đạo hàm theo $t$ và $K$ lần đạo hàm theo $x$ cấp hai bằng nhau.

\textbf{Bước 5 — Kiểm tra dữ liệu, hội tụ và thời gian dài.} Với $t\ge\varepsilon>0$, thừa số $e^{-K\lambda_nt}$ cho phép lấy đạo hàm chuỗi. Với dữ liệu có bước nhảy, giới hạn ban đầu tại bước nhảy là trung bình hai vết. Nếu $f\in L^1$, $|c_n|\le2\|f\|_1/L$ cho hệ chuẩn, nên tổng các mode dương tiến về không đều theo $x$. Neumann giữ trung bình; Dirichlet và biên hỗn hợp có giới hạn không. Một nghiệm giảm theo hàm mũ và khác không không thể tuần hoàn theo thời gian.
\subsection*{Lời giải áp dụng cho bài}
\textbf{Bước 1. Tách biến theo hai biên Dirichlet.}
Với $u_t=Ku_{xx}$, $u(0,t)=u(L,t)=0$, đặt $u=XT$. Ta có $X''+\lambda X=0$, $T'+K\lambda T=0$. Điều kiện $X(0)=X(L)=0$ cho $X_n=\sin(n\pi x/L)$, $\lambda_n=(n\pi/L)^2$, chuẩn $L/2$. Vì vậy $u=\sum c_ne^{-K\lambda_nt}X_n$, $c_n=(2/L)\int_0^LfX_n$.

\textbf{Bước 2. Tính riêng từng dữ liệu.}
(i) $K=L=1$. Từ $\sin3z=3\sin z-4\sin^3z$ suy $\sin^3(\pi x)=3\sin(\pi x)/4-\sin(3\pi x)/4$. Trực giao cho $c_1=3/4,c_3=-1/4$, các hệ số khác bằng không. Do đó
\[u=\frac34e^{-\pi^2t}\sin\pi x-\frac14e^{-9\pi^2t}\sin3\pi x.\]
Tại $t=0$, công thức trả lại $f$; hai sine bằng không ở hai đầu. Đạo hàm mỗi hạng theo $t$ bằng đạo hàm theo $x$ cấp hai, nên PDE đúng. Đây là nghiệm trơn đến $t=0$.

(ii) $K=2,L=\pi$, nhưng $f=\sin(\pi x)$, không phải $\sin x$. Dùng $2\sin(\pi x)\sin(nx)=\cos((\pi-n)x)-\cos((\pi+n)x)$:
\begin{align*}
c_n&=\frac1\pi\left[\frac{\sin((\pi-n)x)}{\pi-n}-\frac{\sin((\pi+n)x)}{\pi+n}\right]_0^\pi\\
&=\frac{(-1)^n\sin(\pi^2)}{\pi}\left(\frac1{\pi-n}-\frac1{\pi+n}\right)
=\frac{2(-1)^nn\sin(\pi^2)}{\pi(\pi^2-n^2)}.
\end{align*}
Không có mẫu bằng không vì $\pi$ không nguyên. Nghiệm là
\[u=\frac{2\sin(\pi^2)}{\pi}\sum_{n=1}^\infty\frac{(-1)^nn}{\pi^2-n^2}e^{-2n^2t}\sin nx.\]
\textbf{Bước 3. Kiểm tra ý nghĩa nghiệm.} Chuỗi trơn với $t>0$, có biên $0$, và tiến về $f(x)$ khi $0<x<\pi$, $t\downarrow0$. Tuy nhiên $f(\pi)=\sin(\pi^2)\ne0$: không có nghiệm liên tục đến góc $(\pi,0)$ thỏa mọi dữ kiện. Khái niệm nghiệm thông thường nếu chỉ yêu cầu liên tục trên toàn bao đóng không được thỏa ở câu này.
\begin{center}\includegraphics[width=.97\linewidth,height=11cm,keepaspectratio]{_xu_ly/graphs/VD_4_1_2_01.png}\end{center}
\textit{Đồ thị minh họa; với nghiệm số, xem phương pháp và tham số trong mã Python kèm theo.}

\clearpage\noindent\texttt{\detokenize{VD_4_1_2_02}}\par
\section*{Lời giải bổ sung — Ví dụ 2 · Mục 4.1.2}

\subsection*{Kiến thức và trình tự giải}
\textbf{Bước 1 — Đọc hệ số và điều kiện biên.} Viết phương trình dưới dạng $u_t=Ku_{xx}$, $K>0$. Nếu đề ký hiệu tốc độ $a$ và viết $u_t=a^2u_{xx}$ thì thay $K=a^2$; không nhầm $a$ với $K$. Kiểm tra điều kiện đầu $f$ có khớp giá trị biên tại $t=0$ hay không. Sự không tương thích tại góc không ngăn nghiệm trơn với $t>0$, nhưng ngăn tính liên tục đến góc khi biên Dirichlet khác vết của $f$.

\textbf{Bước 2 — Tách biến và giải bài toán không gian.} Thế $u=XT$ vào PDE, chia cho $KXT$ để có $X''/X=T'/(KT)=-\lambda$. Do đó $X''+\lambda X=0$, $T'+K\lambda T=0$. Với $\lambda=k^2>0$, $X=A\cos kx+B\sin kx$; thay từng điều kiện biên để tìm $k$. Hai biên Dirichlet cho $X_n=\sin(n\pi x/L)$, $n\ge1$; hai biên Neumann cho $1$ và $\cos(n\pi x/L)$; biên Dirichlet–Neumann cho sine, biên Neumann–Dirichlet cho cosine, với $k_n=(n+1/2)\pi/L$, $n\ge0$. Trường hợp $\lambda=0$ phải giải riêng $X=A+Bx$; chỉ giữ mode hằng nếu biên cho phép.

\textbf{Bước 3 — Tính hệ số, không chỉ ghi tích phân.} Đặt $N_n=\int_0^L X_n^2dx$, $c_n=N_n^{-1}\int_0^L fX_n\,dx$. Với các mode sine/cosine chuẩn, $N_n=L/2$; mode hằng có $N_0=L$. Nếu $f$ là đa thức, tích phân từng phần đến hết; nếu $f$ chia đoạn, tách đúng cận; nếu $f$ lượng giác, dùng tích–tổng. Chẳng hạn với $k=n\pi/L$, $f(0)=f(L)=0$,


\[
\int_0^Lf\sin kx\,dx=\frac1k\int_0^Lf'\cos kx\,dx
=-\frac1{k^2}\int_0^Lf''\sin kx\,dx.
\]


Các hạng biên phải được tính trước khi bỏ. Khi tích–tổng tạo mẫu $k-m=0$, tính trực tiếp tích phân ở tần số đó.

\textbf{Bước 4 — Giải ODE và ghép nghiệm.} Nhân $T_n'+K\lambda_nT_n=0$ với $e^{K\lambda_nt}$, được $(e^{K\lambda_nt}T_n)'=0$. Vì $T_n(0)=c_n$,


\[
u(x,t)=\sum_n c_ne^{-K\lambda_nt}X_n(x).
\]


Thêm $c_0=L^{-1}\int_0^Lf$ cho Neumann hai đầu. Kiểm tra từng mode: $X_n''=-\lambda_nX_n$, nên đạo hàm theo $t$ và $K$ lần đạo hàm theo $x$ cấp hai bằng nhau.

\textbf{Bước 5 — Kiểm tra dữ liệu, hội tụ và thời gian dài.} Với $t\ge\varepsilon>0$, thừa số $e^{-K\lambda_nt}$ cho phép lấy đạo hàm chuỗi. Với dữ liệu có bước nhảy, giới hạn ban đầu tại bước nhảy là trung bình hai vết. Nếu $f\in L^1$, $|c_n|\le2\|f\|_1/L$ cho hệ chuẩn, nên tổng các mode dương tiến về không đều theo $x$. Neumann giữ trung bình; Dirichlet và biên hỗn hợp có giới hạn không. Một nghiệm giảm theo hàm mũ và khác không không thể tuần hoàn theo thời gian.
\subsection*{Lời giải áp dụng cho bài}
\textbf{Bước 1. Cơ sở Neumann.} $X'(0)=X'(L)=0$ cho mode hằng $1$ và $X_n=\cos(n\pi x/L)$, $n\ge1$. Chuẩn lần lượt $L,L/2$, nên
\[u=c_0+\sum_{n\ge1}c_ne^{-K(n\pi/L)^2t}\cos(n\pi x/L),\quad c_0=\frac1L\int f,\quad c_n=\frac2L\int f\cos(n\pi x/L).\]
\textbf{Bước 2. Dữ liệu (i).} $\cos^2(\pi x)=(1+\cos2\pi x)/2$, nên với $K=L=1$, $c_0=c_2=1/2$, các hệ số khác bằng không:
\[u=\frac12+\frac12e^{-4\pi^2t}\cos2\pi x.\]
Đạo hàm biên là $-\pi e^{-4\pi^2t}\sin2\pi x$, bằng không ở $0,1$. Thế $t=0$ trả lại $\cos^2\pi x$; đạo hàm $u_t,u_{xx}$ đều bằng $-2\pi^2e^{-4\pi^2t}\cos2\pi x$. Nghiệm trơn và tiến đều về $1/2$.

\textbf{Bước 3. Dữ liệu (ii).} $K=2,L=\pi,f=\sin x$, $c_0=2/\pi$. Với $n\ne1$,
\begin{align*}
\int_0^\pi\sin x\cos nx\,dx
&=\frac12\int_0^\pi[\sin((1+n)x)+\sin((1-n)x)]dx\\
&=\frac{1+(-1)^n}{1-n^2}.
\end{align*}
Khi $n=1$, $\int\sin x\cos x=[\sin^2x/2]_0^\pi=0$. Vì thế $c_n=0$ khi $n$ lẻ, $c_{2j}=4/[\pi(1-4j^2)]$. Suy
\[u=\frac2\pi+\frac4\pi\sum_{j\ge1}\frac{e^{-8j^2t}}{1-4j^2}\cos2jx.\]
\textbf{Bước 4. Kiểm tra tính thông thường.} Hệ số $O(j^{-2})$ cho hội tụ đều đến $t=0$, trả lại $f$ liên tục. Khi $t>0$, chuỗi đạo hàm hội tụ đều, đúng PDE và Neumann. Vì $f'(0)=1,f'(\pi)=-1$, đạo hàm không thể liên tục đến hai góc và đồng thời bằng $0$ trên biên. Nghiệm vẫn thuộc lớp $C^{2,1}$ với $t>0$ và liên tục trên bao đóng, là lớp nghiệm thông thường được dùng trong ví dụ; không khẳng định trơn đến góc. Ước lượng
\[\left|u-\frac2\pi\right|\le\frac4\pi e^{-8t}\sum_{j\ge1}\frac1{4j^2-1}\longrightarrow0\]
chứng minh giới hạn đều $2/\pi$.
\begin{center}\includegraphics[width=.97\linewidth,height=11cm,keepaspectratio]{_xu_ly/graphs/VD_4_1_2_02.png}\end{center}
\textit{Đồ thị minh họa; với nghiệm số, xem phương pháp và tham số trong mã Python kèm theo.}

\clearpage\noindent\texttt{\detokenize{VD_4_1_2_03}}\par
\section*{Lời giải bổ sung — Ví dụ 3 · Mục 4.1.2}

\subsection*{Kiến thức và trình tự giải}
\textbf{Bước 1 — Chọn đúng đại lượng.} Năng lượng bình phương là $I(t)=\frac12\int_0^Lu^2dx$. Để chứng minh duy nhất, áp dụng cho hiệu $w$ của hai nghiệm cùng dữ kiện; $w$ có dữ liệu đầu bằng không và biên thuần nhất. Tổng nhiệt $E=\int_0^Lu\,dx$ là đại lượng khác, dùng khi đề hỏi bảo toàn hoặc suy giảm tổng nhiệt.

\textbf{Bước 2 — Đạo hàm và tích phân từng phần.} Với $u_t=Ku_{xx}$,


\[
I'=\int_0^Luu_t\,dx
=K[uu_x]_0^L-K\int_0^Lu_x^2dx.
\]


Thay biên tại $L$ rồi trừ giá trị tại $0$. Dirichlet/Neumann thuần nhất làm hạng biên bằng không. Robin $u_x(0)=u(0)$ đóng góp $-u(0)^2$; Robin $u_x(L)=-u(L)$ đóng góp $-u(L)^2$. Biên tuần hoàn triệt tiêu hai đầu. Một điều kiện ghép khác phải được kiểm tra trực tiếp, vì hạng biên có thể dương.

\textbf{Bước 3 — Suy duy nhất và tính dương của trị riêng.} Nếu $I'\le0$ và $I(0)=0$ thì $0\le I(t)\le0$, nên $w=0$. Với $X''+\lambda X=0$, nhân với $X$ rồi tích phân:


\[
\lambda\int_0^LX^2dx=\int_0^L(X')^2dx-[XX']_0^L.
\]


Nếu vế phải không âm thì $\lambda\ge0$. Khi bằng không, buộc từng số hạng không âm bằng không và kiểm tra biên để quyết định có mode hằng hay chỉ nghiệm tầm thường.

\textbf{Bước 4 — Chứng minh trực giao.} Nhân hai phương trình riêng lần lượt với $X_m,X_n$, trừ rồi tích phân:


\[
(\lambda_n-\lambda_m)\int_0^LX_nX_mdx
=[X_nX_m'-X_n'X_m]_0^L.
\]


Biên tự liên hợp làm vế phải bằng không; $\lambda_n\ne\lambda_m$ suy trực giao. Với biên không tự liên hợp hoặc đạo hàm bậc ba, không tự động dùng công thức chiếu trực giao.

\textbf{Bước 5 — Khi có hạng phản ứng hoặc bậc đạo hàm khác.} Với Dirichlet và $u_t=Ku_{xx}+Au_x+Bu$, hạng $A\int uu_x=A[u^2]_0^L/2=0$, nên $I'=-K\int u_x^2+2BI\le2BI$; $e^{-2Bt}I$ không tăng, vẫn suy duy nhất. Với $u_t=Ku_{xxx}$, tích phân hai lần cho $I'=K[uu_{xx}-u_x^2/2]_0^L$. Với tổng nhiệt trong bài Neumann có $-ru$, $E'=K[u_x]_0^L-rE$; giải ODE này trước khi kết luận.
\subsection*{Lời giải áp dụng cho bài}
\textbf{Bước 1. Năng lượng và duy nhất.} Biên là $u_x(0)=u(0)$, $u_x(L)=0$. Tích phân từng phần:
\[I'=\int uu_t=K[uu_x]_0^L-K\int u_x^2
=-K u(0,t)^2-K\int_0^Lu_x^2dx\le0.\]
Áp dụng cho hiệu hai nghiệm có dữ liệu đầu $0$ cho $I(0)=0$, nên hiệu bằng không.

\textbf{Bước 2. Trị riêng và trực giao.} Thế $u=XT$ cho $X''+\lambda X=0$, $X'(0)=X(0)$, $X'(L)=0$, $T'=-K\lambda T$. Nhân ODE với $X$:
\[\lambda\int X^2=\int(X')^2+X(0)^2.\]
Nếu $\lambda=0$, $X'=0$ và $X(0)=0$, suy $X=0$; mọi trị riêng không tầm thường có $\lambda>0$. Trừ hai phương trình như đẳng thức Green cho $(\lambda_n-\lambda_m)\int X_nX_m=[X_nX_m'-X_n'X_m]_0^L=0$, nên trực giao.

\textbf{Bước 3. Tìm $k$ và chuẩn hóa.} Với $\lambda=k^2$, $X=A\cos kx+B\sin kx$. Biên trái cho $A=kB$; chọn $A=1$, được $X=\cos kx+\sin kx/k$. Biên phải cho $-k\sin kL+\cos kL=0$, tức $k\tan kL=1$. Trên mỗi khoảng $((n-1)\pi/L,(n-1/2)\pi/L)$, $k\tan kL$ tăng từ $0$ đến $\infty$, nên có đúng một nghiệm $k_n$, $n\ge1$. Chuẩn là
\[N(k)=\frac L2(1+k^{-2})+\frac{\sin2kL}{4k}(1-k^{-2})+\frac{\sin^2kL}{k^2}.\]
Do đó $u=\sum c_ne^{-Kk_n^2t}X_n$, $c_n=\int fX_n/N(k_n)$.

\textbf{Bước 4. Tính dữ liệu cụ thể $K=2,L=\pi,f=1+2\sin(x/2)$.}
Đặt $q=1/2$, và với $r\ne0$ đặt $D(r)=(1-\cos\pi r)/r$, $E(r)=\sin\pi r/r$, mở rộng $D(0)=0,E(0)=\pi$ theo giới hạn. Tích–tổng và tích phân từng hạng cho
\begin{align*}
J(k)=\int_0^\pi fX\,dx
={}&\frac{\sin\pi k}{k}+\frac{1-\cos\pi k}{k^2}\\
&+D(q+k)+D(q-k)+\frac{E(q-k)-E(q+k)}k.
\end{align*}
Vậy $c_n=J(k_n)/N(k_n)$ và
\[u(x,t)=\sum_{n\ge1}\frac{J(k_n)}{N(k_n)}e^{-2k_n^2t}
\left(\cos k_nx+\frac{\sin k_nx}{k_n}\right).\]
Không còn tích phân hệ số chưa tính. Với $L=\pi$, các nghiệm đầu là $0.3834486028,1.2187220137,2.1391968186$. Mỗi hạng thỏa hai biên và PDE. $f'(0)=1=f(0)$, $f'(\pi)=0$, nên dữ liệu khớp biên; chuỗi trả lại $f$ theo khai triển Sturm–Liouville và trơn khi $t>0$. Tất cả $k_n>0$, suy giới hạn đều bằng $0$. Không mặc định các đạo hàm cấp cao liên tục ở góc, vì điều đó cần thêm tương thích.
\begin{center}\includegraphics[width=.97\linewidth,height=11cm,keepaspectratio]{_xu_ly/graphs/VD_4_1_2_03.png}\end{center}
\textit{Đồ thị minh họa; với nghiệm số, xem phương pháp và tham số trong mã Python kèm theo.}

\clearpage\noindent\texttt{\detokenize{BT_4_1_3_01}}\par
\section*{Lời giải bổ sung — Bài tập 1 · Mục 4.1.3}

\subsection*{Kiến thức và trình tự giải}
\textbf{Bước 1 — Đọc hệ số và điều kiện biên.} Viết phương trình dưới dạng $u_t=Ku_{xx}$, $K>0$. Nếu đề ký hiệu tốc độ $a$ và viết $u_t=a^2u_{xx}$ thì thay $K=a^2$; không nhầm $a$ với $K$. Kiểm tra điều kiện đầu $f$ có khớp giá trị biên tại $t=0$ hay không. Sự không tương thích tại góc không ngăn nghiệm trơn với $t>0$, nhưng ngăn tính liên tục đến góc khi biên Dirichlet khác vết của $f$.

\textbf{Bước 2 — Tách biến và giải bài toán không gian.} Thế $u=XT$ vào PDE, chia cho $KXT$ để có $X''/X=T'/(KT)=-\lambda$. Do đó $X''+\lambda X=0$, $T'+K\lambda T=0$. Với $\lambda=k^2>0$, $X=A\cos kx+B\sin kx$; thay từng điều kiện biên để tìm $k$. Hai biên Dirichlet cho $X_n=\sin(n\pi x/L)$, $n\ge1$; hai biên Neumann cho $1$ và $\cos(n\pi x/L)$; biên Dirichlet–Neumann cho sine, biên Neumann–Dirichlet cho cosine, với $k_n=(n+1/2)\pi/L$, $n\ge0$. Trường hợp $\lambda=0$ phải giải riêng $X=A+Bx$; chỉ giữ mode hằng nếu biên cho phép.

\textbf{Bước 3 — Tính hệ số, không chỉ ghi tích phân.} Đặt $N_n=\int_0^L X_n^2dx$, $c_n=N_n^{-1}\int_0^L fX_n\,dx$. Với các mode sine/cosine chuẩn, $N_n=L/2$; mode hằng có $N_0=L$. Nếu $f$ là đa thức, tích phân từng phần đến hết; nếu $f$ chia đoạn, tách đúng cận; nếu $f$ lượng giác, dùng tích–tổng. Chẳng hạn với $k=n\pi/L$, $f(0)=f(L)=0$,


\[
\int_0^Lf\sin kx\,dx=\frac1k\int_0^Lf'\cos kx\,dx
=-\frac1{k^2}\int_0^Lf''\sin kx\,dx.
\]


Các hạng biên phải được tính trước khi bỏ. Khi tích–tổng tạo mẫu $k-m=0$, tính trực tiếp tích phân ở tần số đó.

\textbf{Bước 4 — Giải ODE và ghép nghiệm.} Nhân $T_n'+K\lambda_nT_n=0$ với $e^{K\lambda_nt}$, được $(e^{K\lambda_nt}T_n)'=0$. Vì $T_n(0)=c_n$,


\[
u(x,t)=\sum_n c_ne^{-K\lambda_nt}X_n(x).
\]


Thêm $c_0=L^{-1}\int_0^Lf$ cho Neumann hai đầu. Kiểm tra từng mode: $X_n''=-\lambda_nX_n$, nên đạo hàm theo $t$ và $K$ lần đạo hàm theo $x$ cấp hai bằng nhau.

\textbf{Bước 5 — Kiểm tra dữ liệu, hội tụ và thời gian dài.} Với $t\ge\varepsilon>0$, thừa số $e^{-K\lambda_nt}$ cho phép lấy đạo hàm chuỗi. Với dữ liệu có bước nhảy, giới hạn ban đầu tại bước nhảy là trung bình hai vết. Nếu $f\in L^1$, $|c_n|\le2\|f\|_1/L$ cho hệ chuẩn, nên tổng các mode dương tiến về không đều theo $x$. Neumann giữ trung bình; Dirichlet và biên hỗn hợp có giới hạn không. Một nghiệm giảm theo hàm mũ và khác không không thể tuần hoàn theo thời gian.
\subsection*{Lời giải áp dụng cho bài}
\textbf{Bước 1. Đọc hệ số trước khi tách biến.} Phương trình (4.1.1) của sách là $u_t=Ku_{xx}$, nhưng bài ghi $a$ mà không định nghĩa quan hệ giữa $a$ và $K$. Vì vậy giữ $K$ trong công thức tổng quát; ở đáp án số dưới đây đọc $a$ là lỗi ký hiệu thay cho $K$. Nếu một quy ước riêng là $K=a^2$, cần thay lại $K$ và ghi rõ, không tự suy từ đề. Hai biên Dirichlet cho $k_n=n\pi/L$, $X_n=\sin k_nx$ và
\[u=\sum_{n\ge1}c_ne^{-Kk_n^2t}\sin k_nx,\qquad c_n=\frac2L\int_0^Lf(x)\sin k_nx\,dx.\]
\textbf{Bước 2. Câu (i), đọc $K=L=1,f=x(1-x)$.} Đặt $k=n\pi$. Tích phân từng phần hai lần, dùng $f(0)=f(1)=0$, $f''=-2$:
\begin{align*}
\int_0^1f\sin kx\,dx
&=\frac1k\int_0^1f'\cos kx\,dx
=-\frac1{k^2}\int_0^1f''\sin kx\,dx\\
&=\frac2{k^2}\left[-\frac{\cos kx}k\right]_0^1
=\frac{2(1-(-1)^n)}{k^3}.
\end{align*}
Nhân với $2$ cho $c_n=4(1-(-1)^n)/(n^3\pi^3)$, tức $0$ khi $n$ chẵn, $8/(n^3\pi^3)$ khi $n$ lẻ:
\[u=\sum_{\substack{n\ge1\\n\ {\rm lẻ}}}\frac8{n^3\pi^3}e^{-n^2\pi^2t}\sin n\pi x.\]
\textbf{Bước 3. Câu (ii), $L=\pi,f=\sin x$.} Trực giao $\int_0^\pi\sin x\sin nx=\pi\delta_{n1}/2$ cho $c_1=1$, các hệ số khác $0$. Công thức không phụ thuộc cách hiểu ký hiệu là
\[u=e^{-Kt}\sin x.\]
Nếu $a=2$ là $K=2$, nghiệm là $e^{-2t}\sin x$. Nếu bổ sung quy ước $K=a^2=4$, nghiệm là $e^{-4t}\sin x$. Bản lời giải ngắn trước đây mặc định $K=a^2$; không có căn cứ cho sự mặc định đó trong (4.1.1).
\textbf{Bước 4. Kiểm tra và chu kỳ.} Hai nghiệm thỏa biên, PDE trên $t>0$ và trả lại $f$ liên tục lúc đầu. Câu (i) không trơn cấp hai đến góc theo mọi hướng vì $f''(0)=f''(1)=-2$, còn biên cố định có $u_t=0$; điều này không ngăn nghiệm thông thường trơn trong miền và liên tục trên bao đóng. Nếu có chu kỳ $T>0$, hệ số mode đầu sẽ thỏa $c_1e^{-Kk_1^2(t+T)}=c_1e^{-Kk_1^2t}$, suy $e^{-Kk_1^2T}=1$, vô lý với $K>0$. Cả hai nghiệm đều không tuần hoàn theo thời gian.
\begin{center}\includegraphics[width=.97\linewidth,height=11cm,keepaspectratio]{_xu_ly/graphs/BT_4_1_3_01.png}\end{center}
\textit{Đồ thị minh họa; với nghiệm số, xem phương pháp và tham số trong mã Python kèm theo.}

\clearpage\noindent\texttt{\detokenize{BT_4_1_3_02}}\par
\section*{Lời giải bổ sung — Bài tập 2 · Mục 4.1.3}

\subsection*{Kiến thức và trình tự giải}
\textbf{Bước 1 — Chọn đúng đại lượng.} Năng lượng bình phương là $I(t)=\frac12\int_0^Lu^2dx$. Để chứng minh duy nhất, áp dụng cho hiệu $w$ của hai nghiệm cùng dữ kiện; $w$ có dữ liệu đầu bằng không và biên thuần nhất. Tổng nhiệt $E=\int_0^Lu\,dx$ là đại lượng khác, dùng khi đề hỏi bảo toàn hoặc suy giảm tổng nhiệt.

\textbf{Bước 2 — Đạo hàm và tích phân từng phần.} Với $u_t=Ku_{xx}$,


\[
I'=\int_0^Luu_t\,dx
=K[uu_x]_0^L-K\int_0^Lu_x^2dx.
\]


Thay biên tại $L$ rồi trừ giá trị tại $0$. Dirichlet/Neumann thuần nhất làm hạng biên bằng không. Robin $u_x(0)=u(0)$ đóng góp $-u(0)^2$; Robin $u_x(L)=-u(L)$ đóng góp $-u(L)^2$. Biên tuần hoàn triệt tiêu hai đầu. Một điều kiện ghép khác phải được kiểm tra trực tiếp, vì hạng biên có thể dương.

\textbf{Bước 3 — Suy duy nhất và tính dương của trị riêng.} Nếu $I'\le0$ và $I(0)=0$ thì $0\le I(t)\le0$, nên $w=0$. Với $X''+\lambda X=0$, nhân với $X$ rồi tích phân:


\[
\lambda\int_0^LX^2dx=\int_0^L(X')^2dx-[XX']_0^L.
\]


Nếu vế phải không âm thì $\lambda\ge0$. Khi bằng không, buộc từng số hạng không âm bằng không và kiểm tra biên để quyết định có mode hằng hay chỉ nghiệm tầm thường.

\textbf{Bước 4 — Chứng minh trực giao.} Nhân hai phương trình riêng lần lượt với $X_m,X_n$, trừ rồi tích phân:


\[
(\lambda_n-\lambda_m)\int_0^LX_nX_mdx
=[X_nX_m'-X_n'X_m]_0^L.
\]


Biên tự liên hợp làm vế phải bằng không; $\lambda_n\ne\lambda_m$ suy trực giao. Với biên không tự liên hợp hoặc đạo hàm bậc ba, không tự động dùng công thức chiếu trực giao.

\textbf{Bước 5 — Khi có hạng phản ứng hoặc bậc đạo hàm khác.} Với Dirichlet và $u_t=Ku_{xx}+Au_x+Bu$, hạng $A\int uu_x=A[u^2]_0^L/2=0$, nên $I'=-K\int u_x^2+2BI\le2BI$; $e^{-2Bt}I$ không tăng, vẫn suy duy nhất. Với $u_t=Ku_{xxx}$, tích phân hai lần cho $I'=K[uu_{xx}-u_x^2/2]_0^L$. Với tổng nhiệt trong bài Neumann có $-ru$, $E'=K[u_x]_0^L-rE$; giải ODE này trước khi kết luận.
\subsection*{Lời giải áp dụng cho bài}
\textbf{(a) Duy nhất bằng năng lượng.} Lấy $w=u_1-u_2$. Khi đó $w_t=Kw_{xx}$, $w(x,0)=0$; hai biên Dirichlet hoặc Neumann thuần nhất cho $[ww_x]_0^L=0$. Vì
\[I'=\int_0^Lww_tdx=K[ww_x]_0^L-K\int_0^Lw_x^2dx\le0,\qquad I(0)=0,\]
nên $0\le I(t)\le0$, $w=0$. Nếu nghiệm chỉ trơn với $t>0$, tích phân từ $\varepsilon$ đến $t$ rồi cho $\varepsilon\downarrow0$ trong lớp có dữ liệu đầu liên tục/$L^2$ thích hợp.

\textbf{(b) Giới hạn đều với $f\in L^1$, không cần giả thiết $L^2$.}
Với Dirichlet, $|c_n|\le2\|f\|_1/L$. Đặt $\gamma=K\pi^2/L^2$, và dùng $n^2\ge1+3(n-1)$, $n\ge1$:
\[\sup_{x\in[0,L]}|u(x,t)|
\le\frac{2\|f\|_1}L\sum_{n\ge1}e^{-\gamma n^2t}
\le\frac{2\|f\|_1}L\frac{e^{-\gamma t}}{1-e^{-3\gamma t}}
\longrightarrow0.\]
Với Neumann, $c_0=L^{-1}\int_0^Lf$, và $|c_n|\le2\|f\|_1/L$ cho $n\ge1$. Cùng ước lượng cho
\[\sup_x|u(x,t)-c_0|\le\frac{2\|f\|_1}L
\frac{e^{-\gamma t}}{1-e^{-3\gamma t}}\longrightarrow0.\]
Đây là giới hạn đều được yêu cầu. Với dữ liệu $L^1$ bất kỳ, nghiệm lấy theo nhân nhiệt và trả lại dữ liệu trong $L^1$ khi $t\downarrow0$; không buộc phải liên tục đến mặt đầu.

\clearpage\noindent\texttt{\detokenize{BT_4_1_3_03}}\par
\section*{Lời giải bổ sung — Bài tập 3 · Mục 4.1.3}

\subsection*{Kiến thức và trình tự giải}
\textbf{Bước 1 — Đọc hệ số và điều kiện biên.} Viết phương trình dưới dạng $u_t=Ku_{xx}$, $K>0$. Nếu đề ký hiệu tốc độ $a$ và viết $u_t=a^2u_{xx}$ thì thay $K=a^2$; không nhầm $a$ với $K$. Kiểm tra điều kiện đầu $f$ có khớp giá trị biên tại $t=0$ hay không. Sự không tương thích tại góc không ngăn nghiệm trơn với $t>0$, nhưng ngăn tính liên tục đến góc khi biên Dirichlet khác vết của $f$.

\textbf{Bước 2 — Tách biến và giải bài toán không gian.} Thế $u=XT$ vào PDE, chia cho $KXT$ để có $X''/X=T'/(KT)=-\lambda$. Do đó $X''+\lambda X=0$, $T'+K\lambda T=0$. Với $\lambda=k^2>0$, $X=A\cos kx+B\sin kx$; thay từng điều kiện biên để tìm $k$. Hai biên Dirichlet cho $X_n=\sin(n\pi x/L)$, $n\ge1$; hai biên Neumann cho $1$ và $\cos(n\pi x/L)$; biên Dirichlet–Neumann cho sine, biên Neumann–Dirichlet cho cosine, với $k_n=(n+1/2)\pi/L$, $n\ge0$. Trường hợp $\lambda=0$ phải giải riêng $X=A+Bx$; chỉ giữ mode hằng nếu biên cho phép.

\textbf{Bước 3 — Tính hệ số, không chỉ ghi tích phân.} Đặt $N_n=\int_0^L X_n^2dx$, $c_n=N_n^{-1}\int_0^L fX_n\,dx$. Với các mode sine/cosine chuẩn, $N_n=L/2$; mode hằng có $N_0=L$. Nếu $f$ là đa thức, tích phân từng phần đến hết; nếu $f$ chia đoạn, tách đúng cận; nếu $f$ lượng giác, dùng tích–tổng. Chẳng hạn với $k=n\pi/L$, $f(0)=f(L)=0$,


\[
\int_0^Lf\sin kx\,dx=\frac1k\int_0^Lf'\cos kx\,dx
=-\frac1{k^2}\int_0^Lf''\sin kx\,dx.
\]


Các hạng biên phải được tính trước khi bỏ. Khi tích–tổng tạo mẫu $k-m=0$, tính trực tiếp tích phân ở tần số đó.

\textbf{Bước 4 — Giải ODE và ghép nghiệm.} Nhân $T_n'+K\lambda_nT_n=0$ với $e^{K\lambda_nt}$, được $(e^{K\lambda_nt}T_n)'=0$. Vì $T_n(0)=c_n$,


\[
u(x,t)=\sum_n c_ne^{-K\lambda_nt}X_n(x).
\]


Thêm $c_0=L^{-1}\int_0^Lf$ cho Neumann hai đầu. Kiểm tra từng mode: $X_n''=-\lambda_nX_n$, nên đạo hàm theo $t$ và $K$ lần đạo hàm theo $x$ cấp hai bằng nhau.

\textbf{Bước 5 — Kiểm tra dữ liệu, hội tụ và thời gian dài.} Với $t\ge\varepsilon>0$, thừa số $e^{-K\lambda_nt}$ cho phép lấy đạo hàm chuỗi. Với dữ liệu có bước nhảy, giới hạn ban đầu tại bước nhảy là trung bình hai vết. Nếu $f\in L^1$, $|c_n|\le2\|f\|_1/L$ cho hệ chuẩn, nên tổng các mode dương tiến về không đều theo $x$. Neumann giữ trung bình; Dirichlet và biên hỗn hợp có giới hạn không. Một nghiệm giảm theo hàm mũ và khác không không thể tuần hoàn theo thời gian.
\subsection*{Lời giải áp dụng cho bài}
\textbf{Bước 1. Cơ sở sine.} Đề ghi $a=1$ trong khi (4.1.1) dùng $K$; lời giải đọc $a$ thay $K$, tức $K=1$. Nếu dùng một quy ước $K=a^2$, ở $a=1$ vẫn cho cùng hệ số $K=1$. $L=\pi$ cho $X_n=\sin nx$, tốc độ giảm $n^2$, $c_n=2\int_0^\pi f\sin nx/\pi$.
\textbf{Bước 2. Câu (i).} $\cos^2x=(1+\cos2x)/2$, $\sin x\cos2x=(\sin3x-\sin x)/2$, nên
\[\sin x\cos^2x=\frac14(\sin x+\sin3x),\qquad
u=\frac14(e^{-t}\sin x+e^{-9t}\sin3x).\]
Thế $t=0$ và hai đầu mút kiểm tra trực tiếp; đạo hàm mỗi mode cho đúng PDE.
\textbf{Bước 3. Câu (ii).} Vì $f=1$ trên $(0,\pi/2]$ và $0$ trên $(\pi/2,\pi]$,
\[c_n=\frac2\pi\int_0^{\pi/2}\sin nx\,dx
=\frac2{\pi n}(1-\cos(n\pi/2)).\]
Theo $n$ modulo $4$: $c_n=0$ khi $n=4j$, $4/(\pi n)$ khi $n=4j+2$, $2/(\pi n)$ khi $n$ lẻ. Do đó
\[u=\sum_{n\ge1}\frac{2(1-\cos(n\pi/2))}{\pi n}e^{-n^2t}\sin nx.\]
\textbf{Bước 4. Dữ liệu gián đoạn.} Với $t>0$, chuỗi trơn và biên $0$. Khi $t\downarrow0$, nghiệm tiến về $1$ trên $0<x<\pi/2$, về $0$ trên $\pi/2<x<\pi$, và về $1/2$ tại $\pi/2$. Giá trị đầu in bằng $1$ ở điểm nhảy không tạo một nghiệm liên tục tại điểm đó; nghiệm được hiểu theo nghĩa $L^p$ hoặc tại các điểm liên tục. Mọi mode đều giảm nên giới hạn lâu dài bằng $0$.

\clearpage\noindent\texttt{\detokenize{BT_4_1_3_04}}\par
\section*{Lời giải bổ sung — Bài tập 4 · Mục 4.1.3}

\subsection*{Kiến thức và trình tự giải}
\textbf{Bước 1 — Đọc hệ số và điều kiện biên.} Viết phương trình dưới dạng $u_t=Ku_{xx}$, $K>0$. Nếu đề ký hiệu tốc độ $a$ và viết $u_t=a^2u_{xx}$ thì thay $K=a^2$; không nhầm $a$ với $K$. Kiểm tra điều kiện đầu $f$ có khớp giá trị biên tại $t=0$ hay không. Sự không tương thích tại góc không ngăn nghiệm trơn với $t>0$, nhưng ngăn tính liên tục đến góc khi biên Dirichlet khác vết của $f$.

\textbf{Bước 2 — Tách biến và giải bài toán không gian.} Thế $u=XT$ vào PDE, chia cho $KXT$ để có $X''/X=T'/(KT)=-\lambda$. Do đó $X''+\lambda X=0$, $T'+K\lambda T=0$. Với $\lambda=k^2>0$, $X=A\cos kx+B\sin kx$; thay từng điều kiện biên để tìm $k$. Hai biên Dirichlet cho $X_n=\sin(n\pi x/L)$, $n\ge1$; hai biên Neumann cho $1$ và $\cos(n\pi x/L)$; biên Dirichlet–Neumann cho sine, biên Neumann–Dirichlet cho cosine, với $k_n=(n+1/2)\pi/L$, $n\ge0$. Trường hợp $\lambda=0$ phải giải riêng $X=A+Bx$; chỉ giữ mode hằng nếu biên cho phép.

\textbf{Bước 3 — Tính hệ số, không chỉ ghi tích phân.} Đặt $N_n=\int_0^L X_n^2dx$, $c_n=N_n^{-1}\int_0^L fX_n\,dx$. Với các mode sine/cosine chuẩn, $N_n=L/2$; mode hằng có $N_0=L$. Nếu $f$ là đa thức, tích phân từng phần đến hết; nếu $f$ chia đoạn, tách đúng cận; nếu $f$ lượng giác, dùng tích–tổng. Chẳng hạn với $k=n\pi/L$, $f(0)=f(L)=0$,


\[
\int_0^Lf\sin kx\,dx=\frac1k\int_0^Lf'\cos kx\,dx
=-\frac1{k^2}\int_0^Lf''\sin kx\,dx.
\]


Các hạng biên phải được tính trước khi bỏ. Khi tích–tổng tạo mẫu $k-m=0$, tính trực tiếp tích phân ở tần số đó.

\textbf{Bước 4 — Giải ODE và ghép nghiệm.} Nhân $T_n'+K\lambda_nT_n=0$ với $e^{K\lambda_nt}$, được $(e^{K\lambda_nt}T_n)'=0$. Vì $T_n(0)=c_n$,


\[
u(x,t)=\sum_n c_ne^{-K\lambda_nt}X_n(x).
\]


Thêm $c_0=L^{-1}\int_0^Lf$ cho Neumann hai đầu. Kiểm tra từng mode: $X_n''=-\lambda_nX_n$, nên đạo hàm theo $t$ và $K$ lần đạo hàm theo $x$ cấp hai bằng nhau.

\textbf{Bước 5 — Kiểm tra dữ liệu, hội tụ và thời gian dài.} Với $t\ge\varepsilon>0$, thừa số $e^{-K\lambda_nt}$ cho phép lấy đạo hàm chuỗi. Với dữ liệu có bước nhảy, giới hạn ban đầu tại bước nhảy là trung bình hai vết. Nếu $f\in L^1$, $|c_n|\le2\|f\|_1/L$ cho hệ chuẩn, nên tổng các mode dương tiến về không đều theo $x$. Neumann giữ trung bình; Dirichlet và biên hỗn hợp có giới hạn không. Một nghiệm giảm theo hàm mũ và khác không không thể tuần hoàn theo thời gian.
\subsection*{Lời giải áp dụng cho bài}
\textbf{Bước 1. Khai triển Neumann.} Phương trình (4.1.1) dùng $K$, nhưng bài cho $a$ mà không định nghĩa quan hệ. Giữ $K$ trong công thức tổng quát và đọc $a$ thay $K$ ở đáp án số; một quy ước $K=a^2$ sẽ cho tốc độ khác ở câu (ii). Ta có $u=c_0+\sum_{n\ge1}c_ne^{-K(n\pi/L)^2t}\cos(n\pi x/L)$, $c_0=L^{-1}\int f$, $c_n=2L^{-1}\int f\cos(n\pi x/L)$.
\textbf{Bước 2. Câu (i).} Đọc $K=L=1$, $\sin^2(2\pi x)=(1-\cos4\pi x)/2$. Do đó
\[u=\frac12-\frac12e^{-16\pi^2t}\cos4\pi x.\]
Biên đạo hàm đúng vì sine bằng $0$ ở $0,1$; tại $t=0$ trả lại $f$, và $u_t=u_{xx}=8\pi^2e^{-16\pi^2t}\cos4\pi x$. Nghiệm trơn đến mặt đầu và tiến đều về $1/2$.
\textbf{Bước 3. Câu (ii).} $L=\pi$, $f=\sin3x$. Trung bình là
\[c_0=\frac1\pi[-\cos3x/3]_0^\pi=\frac2{3\pi}.\]
Với $n\ne3$, tích–tổng cho
\begin{align*}
\int_0^\pi\sin3x\cos nx\,dx
&=\frac12\left[\frac{1-\cos((3+n)\pi)}{3+n}
+\frac{1-\cos((3-n)\pi)}{3-n}\right]\\
&=\frac{3(1+(-1)^n)}{9-n^2}.
\end{align*}
Khi $n=3$, tích phân bằng $[\sin^23x/6]_0^\pi=0$. Vậy
\[u=\frac2{3\pi}+\sum_{\substack{n\ge2\\n\ {\rm chẵn}}}
\frac{12}{\pi(9-n^2)}e^{-Kn^2t}\cos nx.\]
Nếu đọc $a=2$ thay $K$, mũ là $-2n^2t$; nếu bổ sung $K=a^2$, mũ là $-4n^2t$. Các hệ số không gian và trung bình không đổi. Giữ cả hai cách đọc để không mặc nhiên sửa định nghĩa hệ số của đề.
\textbf{Bước 4. Kiểm tra tính trơn.} Hệ số $O(n^{-2})$ cho hội tụ đều và liên tục đến mặt đầu; chuỗi đạo hàm trơn với $t>0$. $f'(0)=3$, $f'(\pi)=-3$ không bằng đạo hàm biên Neumann $0$, nên không trơn đến góc. Trong lớp thông thường của sách, trơn với $t>0$ và liên tục trên bao đóng vẫn được thỏa. Giới hạn đều bằng trung bình $2/(3\pi)$.

\clearpage\noindent\texttt{\detokenize{BT_4_1_3_05}}\par
\section*{Lời giải bổ sung — Bài tập 5 · Mục 4.1.3}

\subsection*{Kiến thức và trình tự giải}
\textbf{Bước 1 — Chọn đúng đại lượng.} Năng lượng bình phương là $I(t)=\frac12\int_0^Lu^2dx$. Để chứng minh duy nhất, áp dụng cho hiệu $w$ của hai nghiệm cùng dữ kiện; $w$ có dữ liệu đầu bằng không và biên thuần nhất. Tổng nhiệt $E=\int_0^Lu\,dx$ là đại lượng khác, dùng khi đề hỏi bảo toàn hoặc suy giảm tổng nhiệt.

\textbf{Bước 2 — Đạo hàm và tích phân từng phần.} Với $u_t=Ku_{xx}$,


\[
I'=\int_0^Luu_t\,dx
=K[uu_x]_0^L-K\int_0^Lu_x^2dx.
\]


Thay biên tại $L$ rồi trừ giá trị tại $0$. Dirichlet/Neumann thuần nhất làm hạng biên bằng không. Robin $u_x(0)=u(0)$ đóng góp $-u(0)^2$; Robin $u_x(L)=-u(L)$ đóng góp $-u(L)^2$. Biên tuần hoàn triệt tiêu hai đầu. Một điều kiện ghép khác phải được kiểm tra trực tiếp, vì hạng biên có thể dương.

\textbf{Bước 3 — Suy duy nhất và tính dương của trị riêng.} Nếu $I'\le0$ và $I(0)=0$ thì $0\le I(t)\le0$, nên $w=0$. Với $X''+\lambda X=0$, nhân với $X$ rồi tích phân:


\[
\lambda\int_0^LX^2dx=\int_0^L(X')^2dx-[XX']_0^L.
\]


Nếu vế phải không âm thì $\lambda\ge0$. Khi bằng không, buộc từng số hạng không âm bằng không và kiểm tra biên để quyết định có mode hằng hay chỉ nghiệm tầm thường.

\textbf{Bước 4 — Chứng minh trực giao.} Nhân hai phương trình riêng lần lượt với $X_m,X_n$, trừ rồi tích phân:


\[
(\lambda_n-\lambda_m)\int_0^LX_nX_mdx
=[X_nX_m'-X_n'X_m]_0^L.
\]


Biên tự liên hợp làm vế phải bằng không; $\lambda_n\ne\lambda_m$ suy trực giao. Với biên không tự liên hợp hoặc đạo hàm bậc ba, không tự động dùng công thức chiếu trực giao.

\textbf{Bước 5 — Khi có hạng phản ứng hoặc bậc đạo hàm khác.} Với Dirichlet và $u_t=Ku_{xx}+Au_x+Bu$, hạng $A\int uu_x=A[u^2]_0^L/2=0$, nên $I'=-K\int u_x^2+2BI\le2BI$; $e^{-2Bt}I$ không tăng, vẫn suy duy nhất. Với $u_t=Ku_{xxx}$, tích phân hai lần cho $I'=K[uu_{xx}-u_x^2/2]_0^L$. Với tổng nhiệt trong bài Neumann có $-ru$, $E'=K[u_x]_0^L-rE$; giải ODE này trước khi kết luận.
\subsection*{Lời giải áp dụng cho bài}
\textbf{Bước 1. Thay từng biên vào năng lượng.}
\[I'=K\{u(L)u_x(L)-u(0)u_x(0)-\int_0^Lu_x^2dx\}.\]
Ở (a), Neumann trái và Dirichlet phải cho hạng biên $0$. Ở (b), $u_x(0)=u(0),u(L)=0$ cho $-u(0)^2$. Ở (c), $u_x(0)=0,u_x(L)=-u(L)$ cho $-u(L)^2$. Ở (d), tổng hai đầu là $-u(0)^2-u(L)^2$. Ở (e), hai giá trị và hai đạo hàm bằng nhau nên triệt tiêu. Do đó (a)--(e) có $I'\le0$.

\textbf{Bước 2. Trị riêng và trực giao cho (a)--(e).}
Đặt $u=XT$: $X''+\lambda X=0$, $T'=-K\lambda T$. Nhân ODE với $X$ cho $\lambda\int X^2=\int(X')^2-[XX']_0^L$. Các hạng biên là các số không âm vừa tính. Với $\lambda=0$, $X$ hằng; (a)--(d) buộc hằng bằng $0$, còn (e) cho $X=1$. Vì vậy dòng đề nói (d) có trị riêng $0$ phải sửa thành (e). Đẳng thức $(\lambda_n-\lambda_m)\int X_nX_m=[X_nX_m'-X_n'X_m]_0^L$ có vế phải $0$ ở từng biên (a)--(e); suy trực giao khi hai trị riêng khác nhau.

\textbf{Bước 3. Giải đầy đủ các điều kiện trên $A,B$.}
Với $X=A\cos kx+B\sin kx$, $k>0$:
\begin{align*}
(a)&:\ B=0,\ \cos kL=0;\quad k_n=(n+1/2)\pi/L,\ X_n=\cos k_nx,\ n\ge0;\\
(b)&:\ A=kB,\ X=\cos kx+\sin kx/k;\quad k\cos kL+\sin kL=0;\\
(c)&:\ B=0,\ -k\sin kL+\cos kL=0;\quad X=\cos kx;\\
(d)&:\ A=kB,\quad (1-k^2)\sin kL+2k\cos kL=0;\\
(e)&:\ k_n=2n\pi/L,\quad 1,\cos k_nx,\sin k_nx.
\end{align*}
Ở (b),(c),(d) chỉ lấy các nghiệm $k>0$; $k=0$ sinh ra sau nhân mẫu ở (d) là nghiệm ngoại lai. Với $X=\cos kx$, $N=L/2+\sin2kL/(4k)$. Với $X=\cos kx+\sin kx/k$,
\[N=\frac L2(1+k^{-2})+\frac{\sin2kL}{4k}(1-k^{-2})+\frac{\sin^2kL}{k^2}.\]
Nghiệm (a)--(d) là $\sum (\int fX_n/N_n)e^{-Kk_n^2t}X_n$. Nghiệm (e) có thêm trung bình và hai hệ số $2\int f\cos/L$, $2\int f\sin/L$. Với $f\in L^1$, nhân nhiệt làm trơn dữ liệu và các mode dương suy giảm đều; (a)--(d) tiến về $0$, (e) tiến về $L^{-1}\int f$.

\textbf{Bước 4. Điều kiện (f) không có tính chất năng lượng mà đề yêu cầu.}
Ở đây $u(0)=0$, $u_x(0)=-u_x(L)$, nên hạng biên là $u(L)u_x(L)$, không xác định dấu. Phản ví dụ cụ thể: đặt $y=x/L$ và
\[f(x)=y-3y(1-y)^{10}+y^{10}(y-1).\]
Ta có $f(0)=0$, $f(L)=1$, $f'(0)=-2/L$, $f'(L)=2/L$, đúng biên (f). Với $p(y)=y(1-y)^{10}$, $\int_0^1(p')^2dy=10/399$, và $q(y)=y^{10}(y-1)$ có cùng chuẩn đạo hàm, $\int p'=\int q'=0$. Bất đẳng thức tam giác trong $L^2$ cho
\[\int_0^L(f')^2dx\le\frac1L\left(1+16\frac{10}{399}\right)<\frac2L=f(L)f'(L).\]
Vì vậy $I'(0)>0$ cho nghiệm có dữ liệu này, theo vết đạo hàm tại $0+$; không thể chứng minh năng lượng không tăng.

\textbf{Bước 5. Khai triển đúng cho (f) phải có hàm liên kết.}
$X(0)=0$ cho $X=\sin kx$, biên đạo hàm cho $\cos kL=-1$, tức $k_n=(2n+1)\pi/L$, $n\ge0$. Nhưng các sine này chỉ biểu diễn phần đối xứng quanh $L/2$. Đặt $Y_n=x\cos k_nx$; cả $X_n,Y_n$ đều thỏa (f) và $Y_n''=-k_n^2Y_n-2k_nX_n$. Tách dữ liệu
\[d(x)=f(x)-f(L-x),\quad r(x)=f(x)-\frac xL d(x).\]
$d$ phản đối xứng, $r$ đối xứng. Khai triển cosine lẻ và sine lẻ cho
\[b_n=\frac2{L^2}\int_0^Ld(x)\cos k_nx\,dx,\qquad
a_n=\frac2L\int_0^Lr(x)\sin k_nx\,dx.\]
Khi đó $f=\sum(a_nX_n+b_nY_n)$ trong $L^2$, vì $\sum b_nY_n=xd(x)/L$. Nghiệm là
\[u=\sum_{n\ge0}e^{-Kk_n^2t}
\{(a_n-2Kk_nb_nt)\sin k_nx+b_nx\cos k_nx\}.\]
Đạo hàm trực tiếp dùng đẳng thức của $Y_n''$ kiểm tra PDE, biên và dữ liệu. Với $f\in L^1$, các hệ số $a_n,b_n$ bị chặn bởi hằng nhân $\|f\|_1$; thừa số $(1+k_nt)e^{-Kk_n^2t}$ suy giảm đều, nên giới hạn vẫn bằng $0$, dù năng lượng không nhất thiết giảm ở mọi thời điểm.

\textbf{Bước 6. Chọn dữ kiện cụ thể và tính cho đủ sáu biên.}
Lấy $K=L=1$ ở tất cả các ví dụ sau. Ở (a), $k=\pi/2$, chọn $f=\cos(\pi x/2)$; $\int fX=N=1/2$, nên hệ số duy nhất bằng $1$ và $u=e^{-\pi^2t/4}\cos(\pi x/2)$.
Ở (b), chọn nghiệm dương đầu $k=2.0287578381$ của $k\cos k+\sin k=0$, $f=\cos kx+\sin kx/k$. Vì $f=X_1$, $\int fX_1/N_1=1$ và mọi hệ số khác bằng $0$ do trực giao, nên $u=e^{-k^2t}(\cos kx+\sin kx/k)$. Ở (c), chọn $k=0.8603335890$ của $k\sin k=\cos k$, $f=\cos kx$; cùng phép chiếu cho $u=e^{-k^2t}\cos kx$. Ở (d), chọn $k=1.3065423742$ của $(1-k^2)\sin k+2k\cos k=0$, $f=\cos kx+\sin kx/k$, nên $u=e^{-k^2t}f(x)$. Các số thập phân chỉ là xấp xỉ; công thức chính xác dùng nghiệm của phương trình tương ứng.
Ở (e), chọn $f=1+\cos2\pi x$. Các tích phân cho $c_0=\int f=1$, $c_1=2\int f\cos2\pi x=1$, các hệ số khác bằng $0$. Vậy $u=1+e^{-4\pi^2t}\cos2\pi x$.
Ở (f), chọn $f=x\cos\pi x$. $d(x)=\cos\pi x$, $r(x)=0$, nên $b_0=2\int_0^1\cos^2\pi x\,dx=1$, $a_n=0$, các $b_n$ khác bằng $0$. Do đó $u=e^{-\pi^2t}(x\cos\pi x-2\pi t\sin\pi x)$. Đây là phép tính hệ số cụ thể với hàm liên kết. Thế từng $f$ vào cặp biên và từng $u$ vào PDE kiểm tra đúng sáu bài đã chọn.

\clearpage\noindent\texttt{\detokenize{VD_4_2_2_01}}\par
\section*{Lời giải bổ sung — Ví dụ 1 · Mục 4.2.2}

\subsection*{Kiến thức và trình tự giải}
\textbf{Bước 1 — Viết toán tử và biên parabol.} Đặt $P=\partial_t-a(x,t)\partial_{xx}-b(x,t)\partial_x-c(x,t)$ với $a\ge0$, $c\le0$. Biên parabol gồm mặt $t=0$ và các biên không gian; mặt $t=T$ vẫn thuộc miền được ước lượng. Với phương trình nhiệt, $a=K>0$, $b=c=0$.

\textbf{Bước 2 — Loại cực đại dương bằng dấu đạo hàm.} Tại cực đại dương nội miền, $w_x=0$, $w_{xx}\le0$, $w_t=0$; nếu cực đại trên mặt thời gian trên thì đạo hàm một phía cho $w_t\ge0$. Suy $Pw\ge0$, mâu thuẫn $Pw<0$. Với $Pw\le0$, xét $w_\varepsilon=w-\varepsilon e^{\gamma t}$, $\gamma>0$: $Pw_\varepsilon=Pw-\varepsilon(\gamma-c)e^{\gamma t}<0$. Chứng minh $w_\varepsilon\le0$ rồi cho $\varepsilon\downarrow0$.

\textbf{Bước 3 — So sánh hai hàm bằng hiệu.} Để chứng minh $v\le u\le V$, kiểm tra $Pv\le Pu\le PV$ và thứ tự của dữ liệu đầu, biên. Áp dụng bước 2 cho $v-u$ và $u-V$. Nếu có nguồn dương, hằng số không nhất thiết là nghiệm trên, nên phải tính $PV-Pu$; không dùng cực đại của dữ liệu để chặn một bài có nguồn mà chưa kiểm tra.

\textbf{Bước 4 — Xử lý Neumann, Robin và tuần hoàn.} Neumann thuần nhất có thể phản xạ chẵn qua đầu mút; điểm biên trở thành điểm trong. Robin viết theo pháp tuyến ngoài $\partial_\nu u+\sigma u=0$, $\sigma\ge0$. Tại cực đại dương không hằng, bổ đề Hopf cho $\partial_\nu u>0$, trái với $\partial_\nu u=-\sigma u\le0$. Với Robin phải thêm mức $0$ vào cận hằng nếu dữ liệu mang một dấu. Biên tuần hoàn ghép hai đầu thành đường tròn; không còn biên không gian.

\textbf{Bước 5 — Miền vô hạn và nguồn bị chặn.} Trên toàn trục/nửa trục, cần nghiệm bị chặn trên mỗi dải thời gian (hoặc điều kiện tăng trưởng đủ dùng). Dùng rào $\varepsilon((x-x_0)^2+2Kt)$ trên đoạn lớn, áp dụng cực đại trên miền hữu hạn rồi cho $\varepsilon\downarrow0$; cách này loại ảnh hưởng từ vô cực. Nếu $|Pu|\le N$ và $|u|\le M$ trên biên parabol, đặt $V=M+Nt$. Vì $PV=N-cV\ge N$, suy $|u|\le M+Nt\le M+NT$. $M$ ở đây là cận của $|u|$, không phải cực đại có dấu.
\subsection*{Lời giải áp dụng cho bài}
\textbf{Bước 1. Đọc đúng dữ liệu.} Đề là $u_t=u_{xx}$ trên $0<x<1$, biên $u(0,t)=u(1,t)=0$ và $f=\sin^2(\pi x)$. Ký hiệu trong bản scan là $\pi$, không phải tần số nguyên $n$; không cần đặt lại miền hay giới hạn $n$.
\textbf{Bước 2. Chọn hai nghiệm so sánh.} Đặt $v=0$, $V=e^{-\pi^2t}\sin\pi x$. Tính $V_t=-\pi^2V=V_{xx}$, $V(0,t)=V(1,t)=0$. Trên $[0,1]$, $0\le\sin\pi x\le1$, nên $0\le\sin^2\pi x\le\sin\pi x$.
\textbf{Bước 3. Áp dụng cực đại cho hai hiệu.} $v-u=-u$ giải phương trình nhiệt, dữ liệu đầu và biên không dương, nên $-u\le0$. Tương tự $u-V$ có dữ liệu đầu $\sin\pi x(\sin\pi x-1)\le0$ và biên $0$, nên $u-V\le0$. Kết luận
\[0\le u(x,t)\le e^{-\pi^2t}\sin\pi x,\quad 0\le x\le1,\ t\ge0.\]
Nếu cần chứng minh cực đại yếu, xét mỗi hiệu trừ $\varepsilon e^t$; toán tử $\partial_t-\partial_{xx}$ cho giá trị âm nghiêm ngặt, loại cực đại dương, rồi cho $\varepsilon\downarrow0$.

\clearpage\noindent\texttt{\detokenize{VD_4_2_2_02}}\par
\section*{Lời giải bổ sung — Ví dụ 2 · Mục 4.2.2}

\subsection*{Kiến thức và trình tự giải}
\textbf{Bước 1 — Viết toán tử và biên parabol.} Đặt $P=\partial_t-a(x,t)\partial_{xx}-b(x,t)\partial_x-c(x,t)$ với $a\ge0$, $c\le0$. Biên parabol gồm mặt $t=0$ và các biên không gian; mặt $t=T$ vẫn thuộc miền được ước lượng. Với phương trình nhiệt, $a=K>0$, $b=c=0$.

\textbf{Bước 2 — Loại cực đại dương bằng dấu đạo hàm.} Tại cực đại dương nội miền, $w_x=0$, $w_{xx}\le0$, $w_t=0$; nếu cực đại trên mặt thời gian trên thì đạo hàm một phía cho $w_t\ge0$. Suy $Pw\ge0$, mâu thuẫn $Pw<0$. Với $Pw\le0$, xét $w_\varepsilon=w-\varepsilon e^{\gamma t}$, $\gamma>0$: $Pw_\varepsilon=Pw-\varepsilon(\gamma-c)e^{\gamma t}<0$. Chứng minh $w_\varepsilon\le0$ rồi cho $\varepsilon\downarrow0$.

\textbf{Bước 3 — So sánh hai hàm bằng hiệu.} Để chứng minh $v\le u\le V$, kiểm tra $Pv\le Pu\le PV$ và thứ tự của dữ liệu đầu, biên. Áp dụng bước 2 cho $v-u$ và $u-V$. Nếu có nguồn dương, hằng số không nhất thiết là nghiệm trên, nên phải tính $PV-Pu$; không dùng cực đại của dữ liệu để chặn một bài có nguồn mà chưa kiểm tra.

\textbf{Bước 4 — Xử lý Neumann, Robin và tuần hoàn.} Neumann thuần nhất có thể phản xạ chẵn qua đầu mút; điểm biên trở thành điểm trong. Robin viết theo pháp tuyến ngoài $\partial_\nu u+\sigma u=0$, $\sigma\ge0$. Tại cực đại dương không hằng, bổ đề Hopf cho $\partial_\nu u>0$, trái với $\partial_\nu u=-\sigma u\le0$. Với Robin phải thêm mức $0$ vào cận hằng nếu dữ liệu mang một dấu. Biên tuần hoàn ghép hai đầu thành đường tròn; không còn biên không gian.

\textbf{Bước 5 — Miền vô hạn và nguồn bị chặn.} Trên toàn trục/nửa trục, cần nghiệm bị chặn trên mỗi dải thời gian (hoặc điều kiện tăng trưởng đủ dùng). Dùng rào $\varepsilon((x-x_0)^2+2Kt)$ trên đoạn lớn, áp dụng cực đại trên miền hữu hạn rồi cho $\varepsilon\downarrow0$; cách này loại ảnh hưởng từ vô cực. Nếu $|Pu|\le N$ và $|u|\le M$ trên biên parabol, đặt $V=M+Nt$. Vì $PV=N-cV\ge N$, suy $|u|\le M+Nt\le M+NT$. $M$ ở đây là cận của $|u|$, không phải cực đại có dấu.
\subsection*{Lời giải áp dụng cho bài}
\textbf{Bước 1. Cận dữ liệu và rào trên đoạn lớn.} Giả thiết $u$ bị chặn trên $\mathbb R\times[0,T]$. Đặt $M=\sup_{\mathbb R\times[0,T]}u$, $N=\sup_xu(x,0)$ và cố định $(x_0,t_0)$, $0<t_0\le T$. Với $\varepsilon>0$,
\[w_\varepsilon=u-\varepsilon\{(x-x_0)^2+2Kt\}.\]
Tính $(w_\varepsilon)_t-K(w_\varepsilon)_{xx}=0$, vì $2K-K\cdot2=0$.
\textbf{Bước 2. Kiểm tra toàn bộ biên.} Trên mặt $t=0$, $w_\varepsilon\le N$. Trên $x=x_0\pm R$, $w_\varepsilon\le M-\varepsilon R^2$. Chọn $R>|x_0|+1$ đủ lớn để $M-\varepsilon R^2\le N$; áp dụng cực đại trên $(x_0-R,x_0+R)\times(0,t_0]$ được $w_\varepsilon\le N$.
\textbf{Bước 3. Bỏ rào và chứng minh cận dưới.} Tại điểm cố định,
\[u(x_0,t_0)\le N+2K\varepsilon t_0.\]
Cho $\varepsilon\downarrow0$, suy $u\le N$. Áp dụng cùng lập luận cho $-u$ cho $u\ge\inf_xu(x,0)$. Do $(x_0,t_0)$ tùy ý,
\[\inf_yu(y,0)\le u(x,t)\le\sup_yu(y,0).\]
Giả thiết bị chặn chính là điều cho phép chọn $R$ bằng cận $M$ độc lập theo $x$; không bỏ bước này trên miền vô hạn.

\clearpage\noindent\texttt{\detokenize{VD_4_2_2_03}}\par
\section*{Lời giải bổ sung — Ví dụ 3 · Mục 4.2.2}

\subsection*{Kiến thức và trình tự giải}
\textbf{Bước 1 — Viết toán tử và biên parabol.} Đặt $P=\partial_t-a(x,t)\partial_{xx}-b(x,t)\partial_x-c(x,t)$ với $a\ge0$, $c\le0$. Biên parabol gồm mặt $t=0$ và các biên không gian; mặt $t=T$ vẫn thuộc miền được ước lượng. Với phương trình nhiệt, $a=K>0$, $b=c=0$.

\textbf{Bước 2 — Loại cực đại dương bằng dấu đạo hàm.} Tại cực đại dương nội miền, $w_x=0$, $w_{xx}\le0$, $w_t=0$; nếu cực đại trên mặt thời gian trên thì đạo hàm một phía cho $w_t\ge0$. Suy $Pw\ge0$, mâu thuẫn $Pw<0$. Với $Pw\le0$, xét $w_\varepsilon=w-\varepsilon e^{\gamma t}$, $\gamma>0$: $Pw_\varepsilon=Pw-\varepsilon(\gamma-c)e^{\gamma t}<0$. Chứng minh $w_\varepsilon\le0$ rồi cho $\varepsilon\downarrow0$.

\textbf{Bước 3 — So sánh hai hàm bằng hiệu.} Để chứng minh $v\le u\le V$, kiểm tra $Pv\le Pu\le PV$ và thứ tự của dữ liệu đầu, biên. Áp dụng bước 2 cho $v-u$ và $u-V$. Nếu có nguồn dương, hằng số không nhất thiết là nghiệm trên, nên phải tính $PV-Pu$; không dùng cực đại của dữ liệu để chặn một bài có nguồn mà chưa kiểm tra.

\textbf{Bước 4 — Xử lý Neumann, Robin và tuần hoàn.} Neumann thuần nhất có thể phản xạ chẵn qua đầu mút; điểm biên trở thành điểm trong. Robin viết theo pháp tuyến ngoài $\partial_\nu u+\sigma u=0$, $\sigma\ge0$. Tại cực đại dương không hằng, bổ đề Hopf cho $\partial_\nu u>0$, trái với $\partial_\nu u=-\sigma u\le0$. Với Robin phải thêm mức $0$ vào cận hằng nếu dữ liệu mang một dấu. Biên tuần hoàn ghép hai đầu thành đường tròn; không còn biên không gian.

\textbf{Bước 5 — Miền vô hạn và nguồn bị chặn.} Trên toàn trục/nửa trục, cần nghiệm bị chặn trên mỗi dải thời gian (hoặc điều kiện tăng trưởng đủ dùng). Dùng rào $\varepsilon((x-x_0)^2+2Kt)$ trên đoạn lớn, áp dụng cực đại trên miền hữu hạn rồi cho $\varepsilon\downarrow0$; cách này loại ảnh hưởng từ vô cực. Nếu $|Pu|\le N$ và $|u|\le M$ trên biên parabol, đặt $V=M+Nt$. Vì $PV=N-cV\ge N$, suy $|u|\le M+Nt\le M+NT$. $M$ ở đây là cận của $|u|$, không phải cực đại có dấu.
\subsection*{Lời giải áp dụng cho bài}
\textbf{Bước 1. Đổi biên Neumann thành điểm trong bằng phản xạ.} Mở rộng $u$ chẵn qua $0$, rồi chẵn qua $L$, và tuần hoàn chu kỳ $2L$. $u_x(0,t)=u_x(L,t)=0$ làm đạo hàm cấp một ghép liên tục; đạo hàm cấp hai và $u_t$ ghép theo tính chẵn và PDE, nên hàm mở rộng giải phương trình nhiệt với $t>0$.
\textbf{Bước 2. Áp dụng cực đại trên đường tròn.} Đặt $m=\min_{[0,L]}f$, $M=\max_{[0,L]}f$. Hai hằng $m,M$ đều giải PDE và có biên Neumann. Với $w=u-M$, xét $w_\varepsilon=w-\varepsilon e^t$ trên miền tuần hoàn. Vì $(\partial_t-K\partial_{xx})w_\varepsilon=-\varepsilon e^t<0$, cực đại dương không thể ở $t>0$; ở $t=0$ cũng không dương. Do đó $u\le M+\varepsilon e^t$, cho $\varepsilon\downarrow0$ suy $u\le M$. Thay $u$ bằng $-u$ suy $u\ge m$.
\textbf{Bước 3. Kết luận.}
\[\min_{y\in[0,L]}u(y,0)\le u(x,t)\le\max_{y\in[0,L]}u(y,0).\]
Nếu dữ liệu không khớp đạo hàm tại góc, phản xạ vẫn dùng cho $t>0$; tính liên tục của nghiệm đến mặt đầu đủ để áp dụng cực đại.

\clearpage\noindent\texttt{\detokenize{BT_4_2_3_01}}\par
\section*{Lời giải bổ sung — Bài tập 1 · Mục 4.2.3}

\subsection*{Kiến thức và trình tự giải}
\textbf{Bước 1 — Viết toán tử và biên parabol.} Đặt $P=\partial_t-a(x,t)\partial_{xx}-b(x,t)\partial_x-c(x,t)$ với $a\ge0$, $c\le0$. Biên parabol gồm mặt $t=0$ và các biên không gian; mặt $t=T$ vẫn thuộc miền được ước lượng. Với phương trình nhiệt, $a=K>0$, $b=c=0$.

\textbf{Bước 2 — Loại cực đại dương bằng dấu đạo hàm.} Tại cực đại dương nội miền, $w_x=0$, $w_{xx}\le0$, $w_t=0$; nếu cực đại trên mặt thời gian trên thì đạo hàm một phía cho $w_t\ge0$. Suy $Pw\ge0$, mâu thuẫn $Pw<0$. Với $Pw\le0$, xét $w_\varepsilon=w-\varepsilon e^{\gamma t}$, $\gamma>0$: $Pw_\varepsilon=Pw-\varepsilon(\gamma-c)e^{\gamma t}<0$. Chứng minh $w_\varepsilon\le0$ rồi cho $\varepsilon\downarrow0$.

\textbf{Bước 3 — So sánh hai hàm bằng hiệu.} Để chứng minh $v\le u\le V$, kiểm tra $Pv\le Pu\le PV$ và thứ tự của dữ liệu đầu, biên. Áp dụng bước 2 cho $v-u$ và $u-V$. Nếu có nguồn dương, hằng số không nhất thiết là nghiệm trên, nên phải tính $PV-Pu$; không dùng cực đại của dữ liệu để chặn một bài có nguồn mà chưa kiểm tra.

\textbf{Bước 4 — Xử lý Neumann, Robin và tuần hoàn.} Neumann thuần nhất có thể phản xạ chẵn qua đầu mút; điểm biên trở thành điểm trong. Robin viết theo pháp tuyến ngoài $\partial_\nu u+\sigma u=0$, $\sigma\ge0$. Tại cực đại dương không hằng, bổ đề Hopf cho $\partial_\nu u>0$, trái với $\partial_\nu u=-\sigma u\le0$. Với Robin phải thêm mức $0$ vào cận hằng nếu dữ liệu mang một dấu. Biên tuần hoàn ghép hai đầu thành đường tròn; không còn biên không gian.

\textbf{Bước 5 — Miền vô hạn và nguồn bị chặn.} Trên toàn trục/nửa trục, cần nghiệm bị chặn trên mỗi dải thời gian (hoặc điều kiện tăng trưởng đủ dùng). Dùng rào $\varepsilon((x-x_0)^2+2Kt)$ trên đoạn lớn, áp dụng cực đại trên miền hữu hạn rồi cho $\varepsilon\downarrow0$; cách này loại ảnh hưởng từ vô cực. Nếu $|Pu|\le N$ và $|u|\le M$ trên biên parabol, đặt $V=M+Nt$. Vì $PV=N-cV\ge N$, suy $|u|\le M+Nt\le M+NT$. $M$ ở đây là cận của $|u|$, không phải cực đại có dấu.
\subsection*{Lời giải áp dụng cho bài}
\textbf{Bước 1. Phát biểu cận riêng cho từng biên.} Với thời gian $[0,T]$, đặt $m_f=\min f$, $M_f=\max f$. Ở (a),
\[\min(m_f,\inf_{0\le s\le T}\mu(s))\le u\le
\max(M_f,\sup_{0\le s\le T}\mu(s)).\]
Ở (b), thêm mức $0$ do biên Robin:
\[\min(0,m_f,\inf_s\nu(s))\le u\le\max(0,M_f,\sup_s\nu(s)).\]
Ở (c),(d), $\min(0,m_f)\le u\le\max(0,M_f)$. Ở (e), $m_f\le u\le M_f$.

\textbf{Bước 2. Kiểm tra các nghiệm trên hằng.} Gọi cận trên là $C$, $w=u-C$. Ở (a), $w_x(0)=0$, $w(L)=\mu-C\le0$. Ở (b), pháp tuyến trái là $-\partial_x$, nên $(\partial_\nu+1)w=-C\le0$ tại $0$, và $w(L)=\nu-C\le0$. Ở (c), $w_x(0)=0$, $(\partial_\nu+1)w=-C\le0$ tại $L$. Ở (d), biểu thức Robin này bằng $-C\le0$ ở cả hai đầu. Lý do chọn $C\ge0$ ở Robin nằm chính ở dấu $-C$.

\textbf{Bước 3. Loại cực đại dương.} Trong miền, dùng $w-\varepsilon e^{\gamma t}$ như chứng minh cực đại yếu. Biên Dirichlet không dương. Neumann phản xạ chẵn để trở thành điểm trong, hoặc dùng Hopf; Robin ở cực đại dương không hằng có $\partial_\nu w>0$, nên $\partial_\nu w+w>0$, trái với biểu thức không dương. Trường hợp hằng dương cũng bị dữ liệu đầu và/hoặc Robin loại. Biên (e) ghép thành đường tròn vì cả giá trị và đạo hàm bằng nhau; cực đại tại đường ghép vẫn là cực đại trong. Suy $w\le0$.

\textbf{Bước 4. Cận dưới.} Áp dụng cận trên cho $-u$, với dữ liệu $-f,-\mu,-\nu$, rồi đổi dấu. Điều này trả lại đúng năm cận đã phát biểu, gồm các mức $0$ cần thiết ở Robin.

\clearpage\noindent\texttt{\detokenize{BT_4_2_3_02}}\par
\section*{Lời giải bổ sung — Bài tập 2 · Mục 4.2.3}

\subsection*{Kiến thức và trình tự giải}
\textbf{Bước 1 — Viết toán tử và biên parabol.} Đặt $P=\partial_t-a(x,t)\partial_{xx}-b(x,t)\partial_x-c(x,t)$ với $a\ge0$, $c\le0$. Biên parabol gồm mặt $t=0$ và các biên không gian; mặt $t=T$ vẫn thuộc miền được ước lượng. Với phương trình nhiệt, $a=K>0$, $b=c=0$.

\textbf{Bước 2 — Loại cực đại dương bằng dấu đạo hàm.} Tại cực đại dương nội miền, $w_x=0$, $w_{xx}\le0$, $w_t=0$; nếu cực đại trên mặt thời gian trên thì đạo hàm một phía cho $w_t\ge0$. Suy $Pw\ge0$, mâu thuẫn $Pw<0$. Với $Pw\le0$, xét $w_\varepsilon=w-\varepsilon e^{\gamma t}$, $\gamma>0$: $Pw_\varepsilon=Pw-\varepsilon(\gamma-c)e^{\gamma t}<0$. Chứng minh $w_\varepsilon\le0$ rồi cho $\varepsilon\downarrow0$.

\textbf{Bước 3 — So sánh hai hàm bằng hiệu.} Để chứng minh $v\le u\le V$, kiểm tra $Pv\le Pu\le PV$ và thứ tự của dữ liệu đầu, biên. Áp dụng bước 2 cho $v-u$ và $u-V$. Nếu có nguồn dương, hằng số không nhất thiết là nghiệm trên, nên phải tính $PV-Pu$; không dùng cực đại của dữ liệu để chặn một bài có nguồn mà chưa kiểm tra.

\textbf{Bước 4 — Xử lý Neumann, Robin và tuần hoàn.} Neumann thuần nhất có thể phản xạ chẵn qua đầu mút; điểm biên trở thành điểm trong. Robin viết theo pháp tuyến ngoài $\partial_\nu u+\sigma u=0$, $\sigma\ge0$. Tại cực đại dương không hằng, bổ đề Hopf cho $\partial_\nu u>0$, trái với $\partial_\nu u=-\sigma u\le0$. Với Robin phải thêm mức $0$ vào cận hằng nếu dữ liệu mang một dấu. Biên tuần hoàn ghép hai đầu thành đường tròn; không còn biên không gian.

\textbf{Bước 5 — Miền vô hạn và nguồn bị chặn.} Trên toàn trục/nửa trục, cần nghiệm bị chặn trên mỗi dải thời gian (hoặc điều kiện tăng trưởng đủ dùng). Dùng rào $\varepsilon((x-x_0)^2+2Kt)$ trên đoạn lớn, áp dụng cực đại trên miền hữu hạn rồi cho $\varepsilon\downarrow0$; cách này loại ảnh hưởng từ vô cực. Nếu $|Pu|\le N$ và $|u|\le M$ trên biên parabol, đặt $V=M+Nt$. Vì $PV=N-cV\ge N$, suy $|u|\le M+Nt\le M+NT$. $M$ ở đây là cận của $|u|$, không phải cực đại có dấu.
\subsection*{Lời giải áp dụng cho bài}
\textbf{Bước 1. Phát biểu trên nửa trục.} Với nghiệm bị chặn trên $[0,\infty)\times[0,T]$ và $f=u(\cdot,0)$:
\begin{align*}
(a)&:\quad \min(\inf_{x\ge0}f,\inf_{0\le s\le T}\mu(s))
\le u\le\max(\sup_{x\ge0}f,\sup_s\mu(s));\\
(b)&:\quad \inf_{x\ge0}f\le u\le\sup_{x\ge0}f.
\end{align*}
\textbf{Bước 2. Dirichlet và vô cực.} Gọi cận trên (a) là $C$, đặt $w=u-C\le0$ trên mặt đầu và tại $x=0$. Xét $w_\varepsilon=w-\varepsilon(x^2+2Kt)$: vẫn giải PDE, không dương trên hai phần biên này. Vì $w$ bị chặn, chọn $R$ lớn làm $w_\varepsilon(R,t)\le0$ cho $0\le t\le T$. Cực đại trên $(0,R)$ cho $w_\varepsilon\le0$. Với $(x,t)$ cố định, cho $\varepsilon\downarrow0$ được $w\le0$. Áp dụng cho $-u$ được cận dưới.
\textbf{Bước 3. Neumann.} Mở rộng chẵn $\widetilde u(x,t)=u(|x|,t)$. Biên $u_x(0,t)=0$ làm nghiệm mở rộng trơn qua $0$ với $t>0$, bị chặn trên toàn trục. Cận toàn trục bằng rào bậc hai cho $\inf f\le\widetilde u\le\sup f$, hạn chế lại $x\ge0$ là kết quả (b).
Miền in $\mathbb R$ ở dòng PDE không khớp phần mở đầu nửa trục; chứng minh dùng miền $x>0$ và biên được nêu trong bài.

\clearpage\noindent\texttt{\detokenize{BT_4_2_3_03}}\par
\section*{Lời giải bổ sung — Bài tập 3 · Mục 4.2.3}

\subsection*{Kiến thức và trình tự giải}
\textbf{Bước 1 — Viết toán tử và biên parabol.} Đặt $P=\partial_t-a(x,t)\partial_{xx}-b(x,t)\partial_x-c(x,t)$ với $a\ge0$, $c\le0$. Biên parabol gồm mặt $t=0$ và các biên không gian; mặt $t=T$ vẫn thuộc miền được ước lượng. Với phương trình nhiệt, $a=K>0$, $b=c=0$.

\textbf{Bước 2 — Loại cực đại dương bằng dấu đạo hàm.} Tại cực đại dương nội miền, $w_x=0$, $w_{xx}\le0$, $w_t=0$; nếu cực đại trên mặt thời gian trên thì đạo hàm một phía cho $w_t\ge0$. Suy $Pw\ge0$, mâu thuẫn $Pw<0$. Với $Pw\le0$, xét $w_\varepsilon=w-\varepsilon e^{\gamma t}$, $\gamma>0$: $Pw_\varepsilon=Pw-\varepsilon(\gamma-c)e^{\gamma t}<0$. Chứng minh $w_\varepsilon\le0$ rồi cho $\varepsilon\downarrow0$.

\textbf{Bước 3 — So sánh hai hàm bằng hiệu.} Để chứng minh $v\le u\le V$, kiểm tra $Pv\le Pu\le PV$ và thứ tự của dữ liệu đầu, biên. Áp dụng bước 2 cho $v-u$ và $u-V$. Nếu có nguồn dương, hằng số không nhất thiết là nghiệm trên, nên phải tính $PV-Pu$; không dùng cực đại của dữ liệu để chặn một bài có nguồn mà chưa kiểm tra.

\textbf{Bước 4 — Xử lý Neumann, Robin và tuần hoàn.} Neumann thuần nhất có thể phản xạ chẵn qua đầu mút; điểm biên trở thành điểm trong. Robin viết theo pháp tuyến ngoài $\partial_\nu u+\sigma u=0$, $\sigma\ge0$. Tại cực đại dương không hằng, bổ đề Hopf cho $\partial_\nu u>0$, trái với $\partial_\nu u=-\sigma u\le0$. Với Robin phải thêm mức $0$ vào cận hằng nếu dữ liệu mang một dấu. Biên tuần hoàn ghép hai đầu thành đường tròn; không còn biên không gian.

\textbf{Bước 5 — Miền vô hạn và nguồn bị chặn.} Trên toàn trục/nửa trục, cần nghiệm bị chặn trên mỗi dải thời gian (hoặc điều kiện tăng trưởng đủ dùng). Dùng rào $\varepsilon((x-x_0)^2+2Kt)$ trên đoạn lớn, áp dụng cực đại trên miền hữu hạn rồi cho $\varepsilon\downarrow0$; cách này loại ảnh hưởng từ vô cực. Nếu $|Pu|\le N$ và $|u|\le M$ trên biên parabol, đặt $V=M+Nt$. Vì $PV=N-cV\ge N$, suy $|u|\le M+Nt\le M+NT$. $M$ ở đây là cận của $|u|$, không phải cực đại có dấu.
\subsection*{Lời giải áp dụng cho bài}
\textbf{Bước 1. Đặt $P=\partial_t-a\partial_{xx}-b\partial_x-c$, $a\ge0,c\le0$.}
(a) Nếu $w$ có cực đại địa phương dương ở một điểm trong, thì $w_x=0,w_{xx}\le0,w_t=0$, nên $Pw=-aw_{xx}-cw\ge0$, trái với $Pw<0$. Nếu điểm ở mặt thời gian trên thì $w_t\ge0$ một phía vẫn cho cùng mâu thuẫn.

(b) Đặt $v=u-\varepsilon e^t$. $Pv=Pu-\varepsilon(1-c)e^t<0$ và $v<0$ trên biên parabol. Nếu $v$ dương ở đâu, cực đại dương trên bao đóng phải ở trong hoặc mặt thời gian trên, trái với (a). Vì thế $u\le\varepsilon e^t$; cho $\varepsilon\downarrow0$ được $u\le0$.

(c) Đặt $C=\max(0,\max_{\partial_pQ_T}u)$. Vì $P(u-C)=cC\le0$ và $u-C\le0$ trên biên, (b) cho $u\le C$. Nếu cực đại toàn miền dương, cận này buộc cực đại bằng cực đại biên. Còn trường hợp cực đại bằng $0$, cần loại khả năng mọi giá trị biên đều âm.
Nếu $\max_{\partial_pQ_T}u=-\eta<0$, tính liên tục trên bao đóng cho một dải gần mặt đầu và hai biên không gian mà $u<-\eta/2$. Cắt bỏ các dải này để được hình trụ nhỏ hơn chứa điểm có $u=0$, với biên parabol $u\le-\eta/2$. Trên hình trụ nhỏ, $c$ có cận $|c|\le C_0$ vì liên tục trên tập compact nằm trong miền. Lấy $t_*$ là thời gian đầu của hình trụ đó và $V=-(\eta/2)e^{-C_0(t-t_*)}$. Khi đó $PV=(-C_0-c)V\ge0$, $u\le V$ trên biên. Áp dụng (b) cho $u-V$ được $u\le V<0$, mâu thuẫn điểm có $u=0$. Vậy cực đại không âm cũng đạt trên biên parabol. Đổi $u$ thành $-u$ cho phát biểu và chứng minh tương tự: cực tiểu không dương đạt trên biên.

(d) Sửa $M$ thành $M=\max_{\partial_pQ_T}|u|$. Đặt $V=M+Nt\ge0$. Khi đó $PV=N-cV\ge N\ge F$. Do $P(u-V)=F-PV\le0$ và $u-V\le0$ trên biên, suy $u\le V$. Với $-u$, nguồn là $-F\le N$, suy $-u\le V$. Vậy
\[|u(x,t)|\le M+Nt\le M+NT.\]
Nếu giữ $M=\max_{\partial_pQ_T}u$ có dấu như bản in, kết luận sai: chọn $a=1,b=c=0,F=0,u=-1$ thì $M=-1,N=0$, không thể có $1\le-1$.

\clearpage\noindent\texttt{\detokenize{BT_4_2_3_04}}\par
\section*{Lời giải bổ sung — Bài tập 4 · Mục 4.2.3}

\subsection*{Kiến thức và trình tự giải}
\textbf{Bước 1 — Viết toán tử và biên parabol.} Đặt $P=\partial_t-a(x,t)\partial_{xx}-b(x,t)\partial_x-c(x,t)$ với $a\ge0$, $c\le0$. Biên parabol gồm mặt $t=0$ và các biên không gian; mặt $t=T$ vẫn thuộc miền được ước lượng. Với phương trình nhiệt, $a=K>0$, $b=c=0$.

\textbf{Bước 2 — Loại cực đại dương bằng dấu đạo hàm.} Tại cực đại dương nội miền, $w_x=0$, $w_{xx}\le0$, $w_t=0$; nếu cực đại trên mặt thời gian trên thì đạo hàm một phía cho $w_t\ge0$. Suy $Pw\ge0$, mâu thuẫn $Pw<0$. Với $Pw\le0$, xét $w_\varepsilon=w-\varepsilon e^{\gamma t}$, $\gamma>0$: $Pw_\varepsilon=Pw-\varepsilon(\gamma-c)e^{\gamma t}<0$. Chứng minh $w_\varepsilon\le0$ rồi cho $\varepsilon\downarrow0$.

\textbf{Bước 3 — So sánh hai hàm bằng hiệu.} Để chứng minh $v\le u\le V$, kiểm tra $Pv\le Pu\le PV$ và thứ tự của dữ liệu đầu, biên. Áp dụng bước 2 cho $v-u$ và $u-V$. Nếu có nguồn dương, hằng số không nhất thiết là nghiệm trên, nên phải tính $PV-Pu$; không dùng cực đại của dữ liệu để chặn một bài có nguồn mà chưa kiểm tra.

\textbf{Bước 4 — Xử lý Neumann, Robin và tuần hoàn.} Neumann thuần nhất có thể phản xạ chẵn qua đầu mút; điểm biên trở thành điểm trong. Robin viết theo pháp tuyến ngoài $\partial_\nu u+\sigma u=0$, $\sigma\ge0$. Tại cực đại dương không hằng, bổ đề Hopf cho $\partial_\nu u>0$, trái với $\partial_\nu u=-\sigma u\le0$. Với Robin phải thêm mức $0$ vào cận hằng nếu dữ liệu mang một dấu. Biên tuần hoàn ghép hai đầu thành đường tròn; không còn biên không gian.

\textbf{Bước 5 — Miền vô hạn và nguồn bị chặn.} Trên toàn trục/nửa trục, cần nghiệm bị chặn trên mỗi dải thời gian (hoặc điều kiện tăng trưởng đủ dùng). Dùng rào $\varepsilon((x-x_0)^2+2Kt)$ trên đoạn lớn, áp dụng cực đại trên miền hữu hạn rồi cho $\varepsilon\downarrow0$; cách này loại ảnh hưởng từ vô cực. Nếu $|Pu|\le N$ và $|u|\le M$ trên biên parabol, đặt $V=M+Nt$. Vì $PV=N-cV\ge N$, suy $|u|\le M+Nt\le M+NT$. $M$ ở đây là cận của $|u|$, không phải cực đại có dấu.
\subsection*{Lời giải áp dụng cho bài}
\textbf{(a) Tìm nghiệm dừng.} Đặt $u=p(x)$, PDE thành $-p''=1$. Tích phân hai lần: $p=-x^2/2+Ax+B$. Từ $p(0)=p(1)=0$, $B=0,A=1/2$, nên $p=x(1-x)/2$.
\textbf{(b) Chọn và kiểm tra hai rào.} Đặt $v=p(1-e^{-8t})$. Vì $p''=-1$,
\[v_t=8pe^{-8t},\quad v_{xx}=-(1-e^{-8t}),\quad
v_t-v_{xx}=1+e^{-8t}(8p-1).\]
Trên $[0,1]$, $0\le p\le1/8$, nên $8p-1\le0$, suy $Pv\le1=Pu= Pp$, với $P=\partial_t-\partial_{xx}$. Hai rào có biên $0$, và $v(x,0)=0=u(x,0)\le p(x)$.
\textbf{Áp dụng so sánh.} $P(v-u)\le0$, dữ liệu của $v-u$ không dương, nên $v\le u$. Tương tự $P(u-p)=0$, dữ liệu đầu $-p\le0$, nên $u\le p$. Kết luận
\[\frac{x(1-x)}2(1-e^{-8t})\le u(x,t)\le\frac{x(1-x)}2.\]
Hai rào cùng tiến về $p$ đều, nên còn có $u\to p$ đều theo $x$.

\clearpage\noindent\texttt{\detokenize{BT_4_2_3_05}}\par
\section*{Lời giải bổ sung — Bài tập 5 · Mục 4.2.3}

\subsection*{Kiến thức và trình tự giải}
\textbf{Bước 1 — Viết toán tử và biên parabol.} Đặt $P=\partial_t-a(x,t)\partial_{xx}-b(x,t)\partial_x-c(x,t)$ với $a\ge0$, $c\le0$. Biên parabol gồm mặt $t=0$ và các biên không gian; mặt $t=T$ vẫn thuộc miền được ước lượng. Với phương trình nhiệt, $a=K>0$, $b=c=0$.

\textbf{Bước 2 — Loại cực đại dương bằng dấu đạo hàm.} Tại cực đại dương nội miền, $w_x=0$, $w_{xx}\le0$, $w_t=0$; nếu cực đại trên mặt thời gian trên thì đạo hàm một phía cho $w_t\ge0$. Suy $Pw\ge0$, mâu thuẫn $Pw<0$. Với $Pw\le0$, xét $w_\varepsilon=w-\varepsilon e^{\gamma t}$, $\gamma>0$: $Pw_\varepsilon=Pw-\varepsilon(\gamma-c)e^{\gamma t}<0$. Chứng minh $w_\varepsilon\le0$ rồi cho $\varepsilon\downarrow0$.

\textbf{Bước 3 — So sánh hai hàm bằng hiệu.} Để chứng minh $v\le u\le V$, kiểm tra $Pv\le Pu\le PV$ và thứ tự của dữ liệu đầu, biên. Áp dụng bước 2 cho $v-u$ và $u-V$. Nếu có nguồn dương, hằng số không nhất thiết là nghiệm trên, nên phải tính $PV-Pu$; không dùng cực đại của dữ liệu để chặn một bài có nguồn mà chưa kiểm tra.

\textbf{Bước 4 — Xử lý Neumann, Robin và tuần hoàn.} Neumann thuần nhất có thể phản xạ chẵn qua đầu mút; điểm biên trở thành điểm trong. Robin viết theo pháp tuyến ngoài $\partial_\nu u+\sigma u=0$, $\sigma\ge0$. Tại cực đại dương không hằng, bổ đề Hopf cho $\partial_\nu u>0$, trái với $\partial_\nu u=-\sigma u\le0$. Với Robin phải thêm mức $0$ vào cận hằng nếu dữ liệu mang một dấu. Biên tuần hoàn ghép hai đầu thành đường tròn; không còn biên không gian.

\textbf{Bước 5 — Miền vô hạn và nguồn bị chặn.} Trên toàn trục/nửa trục, cần nghiệm bị chặn trên mỗi dải thời gian (hoặc điều kiện tăng trưởng đủ dùng). Dùng rào $\varepsilon((x-x_0)^2+2Kt)$ trên đoạn lớn, áp dụng cực đại trên miền hữu hạn rồi cho $\varepsilon\downarrow0$; cách này loại ảnh hưởng từ vô cực. Nếu $|Pu|\le N$ và $|u|\le M$ trên biên parabol, đặt $V=M+Nt$. Vì $PV=N-cV\ge N$, suy $|u|\le M+Nt\le M+NT$. $M$ ở đây là cận của $|u|$, không phải cực đại có dấu.
\subsection*{Lời giải áp dụng cho bài}
\textbf{Bước 1. So dữ liệu đầu.} Trên $0\le x\le\pi$, đặt $s=\sin x\in[0,1]$. Vì $s^2\le s$, $0\le f=\sin^2x\le\sin x$.
\textbf{Bước 2. Kiểm tra nghiệm so sánh.} $V=e^{-t}\sin x$ có $V_t=-e^{-t}\sin x=V_{xx}$, $V(0,t)=V(\pi,t)=0$, $V(x,0)=\sin x$. Hàm $0$ cũng giải PDE và biên.
\textbf{Bước 3. Dùng nguyên lý cực đại cho hiệu.} $P(-u)=0$, biên $0$, dữ liệu đầu $-\sin^2x\le0$, nên $u\ge0$. $P(u-V)=0$, biên $0$, dữ liệu đầu $\sin x(\sin x-1)\le0$, nên $u\le V$. Vì thế
\[0\le u(x,t)\le e^{-t}\sin x.\]
Kết quả đúng trên mọi dải thời gian hữu hạn, nên đúng với mọi $t\ge0$; cận trên còn chứng minh suy giảm đều về $0$.

\clearpage\noindent\texttt{\detokenize{BT_4_2_3_06}}\par
\section*{Lời giải bổ sung — Bài tập 6 · Mục 4.2.3}

\subsection*{Kiến thức và trình tự giải}
\textbf{Bước 1 — Viết toán tử và biên parabol.} Đặt $P=\partial_t-a(x,t)\partial_{xx}-b(x,t)\partial_x-c(x,t)$ với $a\ge0$, $c\le0$. Biên parabol gồm mặt $t=0$ và các biên không gian; mặt $t=T$ vẫn thuộc miền được ước lượng. Với phương trình nhiệt, $a=K>0$, $b=c=0$.

\textbf{Bước 2 — Loại cực đại dương bằng dấu đạo hàm.} Tại cực đại dương nội miền, $w_x=0$, $w_{xx}\le0$, $w_t=0$; nếu cực đại trên mặt thời gian trên thì đạo hàm một phía cho $w_t\ge0$. Suy $Pw\ge0$, mâu thuẫn $Pw<0$. Với $Pw\le0$, xét $w_\varepsilon=w-\varepsilon e^{\gamma t}$, $\gamma>0$: $Pw_\varepsilon=Pw-\varepsilon(\gamma-c)e^{\gamma t}<0$. Chứng minh $w_\varepsilon\le0$ rồi cho $\varepsilon\downarrow0$.

\textbf{Bước 3 — So sánh hai hàm bằng hiệu.} Để chứng minh $v\le u\le V$, kiểm tra $Pv\le Pu\le PV$ và thứ tự của dữ liệu đầu, biên. Áp dụng bước 2 cho $v-u$ và $u-V$. Nếu có nguồn dương, hằng số không nhất thiết là nghiệm trên, nên phải tính $PV-Pu$; không dùng cực đại của dữ liệu để chặn một bài có nguồn mà chưa kiểm tra.

\textbf{Bước 4 — Xử lý Neumann, Robin và tuần hoàn.} Neumann thuần nhất có thể phản xạ chẵn qua đầu mút; điểm biên trở thành điểm trong. Robin viết theo pháp tuyến ngoài $\partial_\nu u+\sigma u=0$, $\sigma\ge0$. Tại cực đại dương không hằng, bổ đề Hopf cho $\partial_\nu u>0$, trái với $\partial_\nu u=-\sigma u\le0$. Với Robin phải thêm mức $0$ vào cận hằng nếu dữ liệu mang một dấu. Biên tuần hoàn ghép hai đầu thành đường tròn; không còn biên không gian.

\textbf{Bước 5 — Miền vô hạn và nguồn bị chặn.} Trên toàn trục/nửa trục, cần nghiệm bị chặn trên mỗi dải thời gian (hoặc điều kiện tăng trưởng đủ dùng). Dùng rào $\varepsilon((x-x_0)^2+2Kt)$ trên đoạn lớn, áp dụng cực đại trên miền hữu hạn rồi cho $\varepsilon\downarrow0$; cách này loại ảnh hưởng từ vô cực. Nếu $|Pu|\le N$ và $|u|\le M$ trên biên parabol, đặt $V=M+Nt$. Vì $PV=N-cV\ge N$, suy $|u|\le M+Nt\le M+NT$. $M$ ở đây là cận của $|u|$, không phải cực đại có dấu.
\subsection*{Lời giải áp dụng cho bài}
\textbf{Bước 1. Đúng điều kiện biên.} Bài có $u(0,t)=0$, $u_x(\pi,t)=0$, là Dirichlet–Neumann; không phải Neumann ở cả hai đầu. $f=\sin x+\frac12\sin2x=\sin x(1+\cos x)\ge0$ trên $[0,\pi]$.
\textbf{Bước 2. Tính cực đại của $f$.}
\[f'=\cos x+\cos2x=2\cos^2x+\cos x-1
=(2\cos x-1)(\cos x+1).\]
Trong $(0,\pi)$, điểm tới hạn là $\cos x=1/2$, tức $x=\pi/3$. Hai đầu có $f=0$, còn
\[f(\pi/3)=\frac{\sqrt3}2\left(1+\frac12\right)=\frac{3\sqrt3}4.\]
Vậy $\min f=0$, $\max f=3\sqrt3/4$.
\textbf{Bước 3. So sánh hằng.} Hằng $0$ và $C=3\sqrt3/4$ giải PDE, có đạo hàm bằng $0$ tại $\pi$. Tại biên Dirichlet $0$, $0=u(0,t)\le C$; mặt đầu có $0\le f\le C$. Phản xạ chẵn tại $\pi$ hoặc dùng nguyên lý cực đại cho biên hỗn hợp suy
\[0\le u(x,t)\le\frac{3\sqrt3}4.\]
Không bảo toàn trung bình: nhiệt thoát qua đầu Dirichlet. Nếu xét thời gian dài, cơ sở sine nửa nguyên $k_n=n+1/2$ chỉ có trị riêng dương, nên nghiệm tiến về $0$, không phải $2/\pi$.

\clearpage\noindent\texttt{\detokenize{BT_4_2_3_07}}\par
\section*{Lời giải bổ sung — Bài tập 7 · Mục 4.2.3}

\subsection*{Kiến thức và trình tự giải}
\textbf{Bước 1 — Viết toán tử và biên parabol.} Đặt $P=\partial_t-a(x,t)\partial_{xx}-b(x,t)\partial_x-c(x,t)$ với $a\ge0$, $c\le0$. Biên parabol gồm mặt $t=0$ và các biên không gian; mặt $t=T$ vẫn thuộc miền được ước lượng. Với phương trình nhiệt, $a=K>0$, $b=c=0$.

\textbf{Bước 2 — Loại cực đại dương bằng dấu đạo hàm.} Tại cực đại dương nội miền, $w_x=0$, $w_{xx}\le0$, $w_t=0$; nếu cực đại trên mặt thời gian trên thì đạo hàm một phía cho $w_t\ge0$. Suy $Pw\ge0$, mâu thuẫn $Pw<0$. Với $Pw\le0$, xét $w_\varepsilon=w-\varepsilon e^{\gamma t}$, $\gamma>0$: $Pw_\varepsilon=Pw-\varepsilon(\gamma-c)e^{\gamma t}<0$. Chứng minh $w_\varepsilon\le0$ rồi cho $\varepsilon\downarrow0$.

\textbf{Bước 3 — So sánh hai hàm bằng hiệu.} Để chứng minh $v\le u\le V$, kiểm tra $Pv\le Pu\le PV$ và thứ tự của dữ liệu đầu, biên. Áp dụng bước 2 cho $v-u$ và $u-V$. Nếu có nguồn dương, hằng số không nhất thiết là nghiệm trên, nên phải tính $PV-Pu$; không dùng cực đại của dữ liệu để chặn một bài có nguồn mà chưa kiểm tra.

\textbf{Bước 4 — Xử lý Neumann, Robin và tuần hoàn.} Neumann thuần nhất có thể phản xạ chẵn qua đầu mút; điểm biên trở thành điểm trong. Robin viết theo pháp tuyến ngoài $\partial_\nu u+\sigma u=0$, $\sigma\ge0$. Tại cực đại dương không hằng, bổ đề Hopf cho $\partial_\nu u>0$, trái với $\partial_\nu u=-\sigma u\le0$. Với Robin phải thêm mức $0$ vào cận hằng nếu dữ liệu mang một dấu. Biên tuần hoàn ghép hai đầu thành đường tròn; không còn biên không gian.

\textbf{Bước 5 — Miền vô hạn và nguồn bị chặn.} Trên toàn trục/nửa trục, cần nghiệm bị chặn trên mỗi dải thời gian (hoặc điều kiện tăng trưởng đủ dùng). Dùng rào $\varepsilon((x-x_0)^2+2Kt)$ trên đoạn lớn, áp dụng cực đại trên miền hữu hạn rồi cho $\varepsilon\downarrow0$; cách này loại ảnh hưởng từ vô cực. Nếu $|Pu|\le N$ và $|u|\le M$ trên biên parabol, đặt $V=M+Nt$. Vì $PV=N-cV\ge N$, suy $|u|\le M+Nt\le M+NT$. $M$ ở đây là cận của $|u|$, không phải cực đại có dấu.
\subsection*{Lời giải áp dụng cho bài}
\textbf{Bước 1. Tìm cận dữ liệu biên.} Đặt $h(t)=te^{-t}$. $h'(t)=e^{-t}(1-t)$, nên tăng trên $[0,1]$, giảm sau $1$, $h(1)=e^{-1}$. Vì vậy $0\le h(t)\le e^{-1}$.
\textbf{Bước 2. Kiểm tra hiệu trên biên parabol.} $P=\partial_t-\partial_{xx}$; $Pu=0$. Với $w=u-e^{-1}$, $Pw=0$, mặt đầu $w=-e^{-1}$, tại $x=0$ có $w=te^{-t}-e^{-1}\le0$, tại $x=1$ có $w=-e^{-1}$. Cực đại yếu cho $u\le e^{-1}$. Với $-u$, cả dữ liệu đầu lẫn hai biên không dương, nên $u\ge0$.
\textbf{Bước 3. Kết luận.} Trên từng dải $[0,T]$,
\[0\le u(x,t)\le e^{-1}.\]
Do $T$ tùy ý, cận đúng với mọi $t>0$. Dữ liệu đầu tại $x=0$ khớp $h(0)=0$, nên không có mâu thuẫn góc về giá trị.

\clearpage\noindent\texttt{\detokenize{BT_4_2_3_08}}\par
\section*{Lời giải bổ sung — Bài tập 8 · Mục 4.2.3}

\subsection*{Kiến thức và trình tự giải}
\textbf{Bước 1 — Viết toán tử và biên parabol.} Đặt $P=\partial_t-a(x,t)\partial_{xx}-b(x,t)\partial_x-c(x,t)$ với $a\ge0$, $c\le0$. Biên parabol gồm mặt $t=0$ và các biên không gian; mặt $t=T$ vẫn thuộc miền được ước lượng. Với phương trình nhiệt, $a=K>0$, $b=c=0$.

\textbf{Bước 2 — Loại cực đại dương bằng dấu đạo hàm.} Tại cực đại dương nội miền, $w_x=0$, $w_{xx}\le0$, $w_t=0$; nếu cực đại trên mặt thời gian trên thì đạo hàm một phía cho $w_t\ge0$. Suy $Pw\ge0$, mâu thuẫn $Pw<0$. Với $Pw\le0$, xét $w_\varepsilon=w-\varepsilon e^{\gamma t}$, $\gamma>0$: $Pw_\varepsilon=Pw-\varepsilon(\gamma-c)e^{\gamma t}<0$. Chứng minh $w_\varepsilon\le0$ rồi cho $\varepsilon\downarrow0$.

\textbf{Bước 3 — So sánh hai hàm bằng hiệu.} Để chứng minh $v\le u\le V$, kiểm tra $Pv\le Pu\le PV$ và thứ tự của dữ liệu đầu, biên. Áp dụng bước 2 cho $v-u$ và $u-V$. Nếu có nguồn dương, hằng số không nhất thiết là nghiệm trên, nên phải tính $PV-Pu$; không dùng cực đại của dữ liệu để chặn một bài có nguồn mà chưa kiểm tra.

\textbf{Bước 4 — Xử lý Neumann, Robin và tuần hoàn.} Neumann thuần nhất có thể phản xạ chẵn qua đầu mút; điểm biên trở thành điểm trong. Robin viết theo pháp tuyến ngoài $\partial_\nu u+\sigma u=0$, $\sigma\ge0$. Tại cực đại dương không hằng, bổ đề Hopf cho $\partial_\nu u>0$, trái với $\partial_\nu u=-\sigma u\le0$. Với Robin phải thêm mức $0$ vào cận hằng nếu dữ liệu mang một dấu. Biên tuần hoàn ghép hai đầu thành đường tròn; không còn biên không gian.

\textbf{Bước 5 — Miền vô hạn và nguồn bị chặn.} Trên toàn trục/nửa trục, cần nghiệm bị chặn trên mỗi dải thời gian (hoặc điều kiện tăng trưởng đủ dùng). Dùng rào $\varepsilon((x-x_0)^2+2Kt)$ trên đoạn lớn, áp dụng cực đại trên miền hữu hạn rồi cho $\varepsilon\downarrow0$; cách này loại ảnh hưởng từ vô cực. Nếu $|Pu|\le N$ và $|u|\le M$ trên biên parabol, đặt $V=M+Nt$. Vì $PV=N-cV\ge N$, suy $|u|\le M+Nt\le M+NT$. $M$ ở đây là cận của $|u|$, không phải cực đại có dấu.
\subsection*{Lời giải áp dụng cho bài}
\textbf{(a) Thế trực tiếp nghiệm.} Với $u=(t+1)x(1-x)$,
\[u_t=x(1-x),\quad u_x=(t+1)(1-2x),\quad u_{xx}=-2(t+1).\]
Suy $u_t-u_{xx}=x(1-x)+2(t+1)$, đúng nguồn. Tại $x=0,1$, $u=0$; tại $t=0$, $u=x(1-x)$.
\textbf{(b) Cực trị dữ liệu.} $f'=1-2x=0$ cho $x=1/2$, $f(1/2)=1/4$; hai đầu bằng $0$ và $f\ge0$. Biên cũng bằng $0$, nên $m=0$, $M=1/4$.
\textbf{(c) Nhiệt độ vượt $M$.} Tại $x=1/2$,
\[u(1/2,t)=\frac{t+1}4>\frac14=M\quad\text{với mọi }t>0.\]
Ví dụ $t=1$ cho $u=1/2$. Nguyên lý cực đại của phương trình thuần nhất không áp dụng trực tiếp vì $Pu=2(t+1)+x(1-x)>0$; nguồn dương làm nhiệt tăng.
\begin{center}\includegraphics[width=.97\linewidth,height=11cm,keepaspectratio]{_xu_ly/graphs/BT_4_2_3_08.png}\end{center}
\textit{Đồ thị minh họa; với nghiệm số, xem phương pháp và tham số trong mã Python kèm theo.}

\clearpage\noindent\texttt{\detokenize{VD_4_3_2_01}}\par
\section*{Lời giải bổ sung — Ví dụ 1 · Mục 4.3.2}

\subsection*{Kiến thức và trình tự giải}
\textbf{Bước 1 — Tìm cơ sở và chuẩn.} Giải $-X_n''=\lambda_nX_n$ với đúng biên thuần nhất. Đặt $N_n=\int_0^LX_n^2dx$; các hàm riêng ở biên Robin không có chuẩn mặc định $L/2$. Từ $c_n=N_n^{-1}\int fX_n$ trong chuỗi nghiệm, đổi thứ tự tổng–tích phân với $t>0$ để nhận ra


\[
G(x,y,t)=\sum_n\frac{X_n(x)X_n(y)}{N_n}e^{-K\lambda_nt},\qquad
u(x,t)=\int_0^LG(x,y,t)f(y)dy.
\]


\textbf{Bước 2 — Tính chuẩn Robin đến biểu thức sơ cấp.} Nếu $X=\cos kx+\sin kx/k$, dùng


\[
N=\frac L2(1+k^{-2})+\frac{\sin(2kL)}{4k}(1-k^{-2})+\frac{\sin^2(kL)}{k^2}.
\]


Nếu $X=\cos kx$, $N=L/2+\sin(2kL)/(4k)$. Hai biên tuần hoàn có mode hằng $1/L$ trong $G$ và cặp sine/cosine tần số $2n\pi/L$.

\textbf{Bước 3 — Kiểm tra PDE, biên và đối xứng.} Đạo hàm $e^{-K\lambda_nt}X_n(x)$ cho $G_t=KG_{xx}$; đạo hàm theo $y$ tương tự. Biên được kiểm tra trên từng $X_n$. Mỗi tích $X_n(x)X_n(y)$ đối xứng; hội tụ đều của chuỗi đạo hàm trên $t\ge\varepsilon$ biện minh các thao tác. Dữ liệu delta ban đầu hiểu qua $\int Gf\to f$ trong miền hoặc trong $L^2$, không bằng cách thế $t=0$ trực tiếp vào một chuỗi không hội tụ.

\textbf{Bước 4 — Chứng minh dương và khối lượng.} Nguyên lý cực đại cho $\int Gf\ge0$ với $f\ge0$, suy $G\ge0$ bằng cách lấy $f$ tập trung gần một điểm. Kết hợp tính nửa nhóm và nguyên lý mạnh để suy $G>0$ trong miền. $m(x,t)=\int G(x,y,t)dy$ là nghiệm có dữ liệu đầu $1$. Neumann/tuần hoàn giữ $m=1$; Dirichlet hoặc Robin hấp thụ cho $0<m<1$.

\textbf{Bước 5 — Nửa nhóm và thời gian dài.} Trực giao cho $\int G(x,z,s)G(z,y,t)dz=G(x,y,s+t)$. Với $H(t)=\max_xm(x,t)$ ở Dirichlet, $m(x,t+s)\le H(t)m(x,s)$, nên $H(t+s)\le H(t)H(s)<H(t)$. Để chứng minh $H(t)\to1$ lúc đầu, giữ một $x_0$ trong miền, dùng $m(x_0,t)\to1$ và $m(x_0,t)\le H(t)\le1$.
\subsection*{Lời giải áp dụng cho bài}
\textbf{Bước 1. Nhận diện hàm Green từ chuỗi Neumann.}
Với $u_t=Ku_{xx}$, $u_x(0,t)=u_x(L,t)=0$, hệ riêng là $X_0=1$, $\lambda_0=0$, $N_0=L$ và $X_n=\cos(n\pi x/L)$, $\lambda_n=(n\pi/L)^2$, $N_n=L/2$. Chiếu dữ liệu rồi đổi tổng–tích phân với $t>0$:
\begin{align*}
u(x,t)&=\frac1L\int_0^Lf(y)dy+
\sum_{n\ge1}\frac2L e^{-K(n\pi/L)^2t}\cos(n\pi x/L)\int_0^Lf(y)\cos(n\pi y/L)dy\\
&=\int_0^LG_N(x,y,t)f(y)dy,\\
G_N(x,y,t)&=\frac1L+\frac2L\sum_{n\ge1}e^{-K(n\pi/L)^2t}
\cos(n\pi x/L)\cos(n\pi y/L).
\end{align*}
\textbf{Bước 2. Kiểm tra đạo hàm, biên và đối xứng.}
Mỗi đạo hàm cấp hữu hạn bị chặn bởi hằng nhân $n^m e^{-K(n\pi/L)^2\varepsilon}$ khi $t\ge\varepsilon>0$; tổng này hội tụ. Do đó được đạo hàm từng số hạng. Đạo hàm theo $t$ nhân $-K(n\pi/L)^2$, đạo hàm theo $x$ hai lần nhân $-(n\pi/L)^2$, nên $G_{Nt}=KG_{Nxx}=KG_{Nyy}$. Tích hai cosine đối xứng theo $x,y$ và tuần hoàn $2L$ theo từng biến; đạo hàm theo $x$ chứa sine, bằng $0$ tại $0,L$.
\textbf{Bước 3. Tính khối lượng và chứng minh dương.} Vì $\int_0^L\cos(n\pi y/L)dy=0$,
\[\int_0^LG_N(x,y,t)dy=1.\]
Đặt $g_K(z,t)=(4\pi Kt)^{-1/2}e^{-z^2/(4Kt)}$. Phản xạ chẵn và tuần hoàn cho
\[G_N(x,y,t)=\sum_{j\in\mathbb Z}\{g_K(x-y+2jL,t)+g_K(x+y+2jL,t)\}.\]
Có thể kiểm tra đẳng thức bằng hệ số Fourier: nhân tuần hoàn $\sum_jg_K(z+2jL,t)$ có hệ số $(2L)^{-1}e^{-K(n\pi/L)^2t}$; cộng hai ảnh và dùng $\cos(a-b)+\cos(a+b)=2\cos a\cos b$ trả lại chuỗi trên. Mỗi Gauss dương nên $G_N>0$.
\textbf{Bước 4. Dữ liệu đầu, nửa nhóm và giới hạn.} Với $f$ liên tục, xấp xỉ đơn vị của nhân nhiệt trên hàm mở rộng chẵn cho $u\to f$ đều. Trực giao cosine cho $\int_0^LG_N(x,z,s)G_N(z,y,t)dz=G_N(x,y,s+t)$. Vì $G_N>0$ và tổng trọng số $1$, $\min f\le u\le\max f$; tính đối xứng cho $\int_0^Lu(x,t)dx=\int f$. Các mode dương giảm theo mũ, nên $u\to L^{-1}\int f$ đều. Không đòi $\int G_N$ bằng $1$ bằng một phép tích phân trên toàn trục: tích phân chuẩn hóa ở đây là trên một đoạn $[0,L]$.

\clearpage\noindent\texttt{\detokenize{BT_4_3_3_01}}\par
\section*{Lời giải bổ sung — Bài tập 1 · Mục 4.3.3}

\subsection*{Kiến thức và trình tự giải}
\textbf{Bước 1 — Tìm cơ sở và chuẩn.} Giải $-X_n''=\lambda_nX_n$ với đúng biên thuần nhất. Đặt $N_n=\int_0^LX_n^2dx$; các hàm riêng ở biên Robin không có chuẩn mặc định $L/2$. Từ $c_n=N_n^{-1}\int fX_n$ trong chuỗi nghiệm, đổi thứ tự tổng–tích phân với $t>0$ để nhận ra


\[
G(x,y,t)=\sum_n\frac{X_n(x)X_n(y)}{N_n}e^{-K\lambda_nt},\qquad
u(x,t)=\int_0^LG(x,y,t)f(y)dy.
\]


\textbf{Bước 2 — Tính chuẩn Robin đến biểu thức sơ cấp.} Nếu $X=\cos kx+\sin kx/k$, dùng


\[
N=\frac L2(1+k^{-2})+\frac{\sin(2kL)}{4k}(1-k^{-2})+\frac{\sin^2(kL)}{k^2}.
\]


Nếu $X=\cos kx$, $N=L/2+\sin(2kL)/(4k)$. Hai biên tuần hoàn có mode hằng $1/L$ trong $G$ và cặp sine/cosine tần số $2n\pi/L$.

\textbf{Bước 3 — Kiểm tra PDE, biên và đối xứng.} Đạo hàm $e^{-K\lambda_nt}X_n(x)$ cho $G_t=KG_{xx}$; đạo hàm theo $y$ tương tự. Biên được kiểm tra trên từng $X_n$. Mỗi tích $X_n(x)X_n(y)$ đối xứng; hội tụ đều của chuỗi đạo hàm trên $t\ge\varepsilon$ biện minh các thao tác. Dữ liệu delta ban đầu hiểu qua $\int Gf\to f$ trong miền hoặc trong $L^2$, không bằng cách thế $t=0$ trực tiếp vào một chuỗi không hội tụ.

\textbf{Bước 4 — Chứng minh dương và khối lượng.} Nguyên lý cực đại cho $\int Gf\ge0$ với $f\ge0$, suy $G\ge0$ bằng cách lấy $f$ tập trung gần một điểm. Kết hợp tính nửa nhóm và nguyên lý mạnh để suy $G>0$ trong miền. $m(x,t)=\int G(x,y,t)dy$ là nghiệm có dữ liệu đầu $1$. Neumann/tuần hoàn giữ $m=1$; Dirichlet hoặc Robin hấp thụ cho $0<m<1$.

\textbf{Bước 5 — Nửa nhóm và thời gian dài.} Trực giao cho $\int G(x,z,s)G(z,y,t)dz=G(x,y,s+t)$. Với $H(t)=\max_xm(x,t)$ ở Dirichlet, $m(x,t+s)\le H(t)m(x,s)$, nên $H(t+s)\le H(t)H(s)<H(t)$. Để chứng minh $H(t)\to1$ lúc đầu, giữ một $x_0$ trong miền, dùng $m(x_0,t)\to1$ và $m(x_0,t)\le H(t)\le1$.
\subsection*{Lời giải áp dụng cho bài}
\textbf{Bước 1. Nhân Dirichlet và hội tụ.}
\[G(x,y,t)=\frac2L\sum_{n\ge1}e^{-K(n\pi/L)^2t}
\sin(n\pi x/L)\sin(n\pi y/L).\]
Với $t\ge\varepsilon>0$, chuỗi và các đạo hàm hội tụ đều. Với dữ liệu $f$, $u=\int Gf$ giải PDE, biên $0$ và trả lại $f$ tại các điểm liên tục.

\textbf{(a) Dương yếu rồi dương nghiêm ngặt.}
Nếu $f\ge0$, cực đại cho $u\ge0$. Nếu $f\not\equiv0$, nguyên lý cực đại mạnh cho $u(x,t)>0$ khi $0<x<L,t>0$. Để suy $G\ge0$, nếu $G(x_0,y_0,t_0)<0$, tính liên tục theo $y$ cho một khoảng mà $G<0$; chọn $f\ge0$, không bằng $0$, hỗ trợ trong khoảng này, sẽ có $\int G(x_0,y,t_0)f(y)dy<0$, mâu thuẫn. Để chứng minh $G>0$ tại mọi điểm, chú ý
\[G(y,y,s)=\frac2L\sum_{n\ge1}e^{-K(n\pi/L)^2s}\sin^2(n\pi y/L)>0\quad(0<y<L).\]
Với $s=t/2$, hàm $G(\cdot,y,s)\ge0$ và không bằng không. Nguyên lý mạnh cho nghiệm tiến thêm $s$; tính nửa nhóm suy $G(x,y,t)>0$.

\textbf{(b) Cận nhỏ hơn $1$ và khối lượng.}
Với $f\le1$, hàm $1-u$ giải PDE, dữ liệu đầu không âm và biên bằng $1$, nên $1-u>0$ trong miền, tức $\int Gf<1$. Để chuyển từ $f$ vanishing ở hai đầu sang dữ liệu $1$, chọn $f_j(y)=\min(1,jy,j(L-y))$: $0\le f_j\uparrow1$ trong miền, $f_j(0)=f_j(L)=0$. Vì $G\ge0$, hội tụ đơn điệu cho
\[m(x,t)=\int_0^LG(x,y,t)dy=\lim_j\int Gf_j\le1.\]
Hàm $m$ là nghiệm với dữ liệu đầu $1$ trong miền và biên $0$ cho $t>0$. Nguyên lý mạnh cho $m>0$ và $1-m>0$ trong miền; vì thế $0<m<1$, $m(0,t)=m(L,t)=0$. Tính liên tục trên đoạn với $t>0$ cho $H(t)=\max_xm(x,t)$ thỏa $0<H(t)<1$.

\textbf{(c) Tính giảm nghiêm ngặt và giới hạn đầu.}
Trực giao sine cho đẳng thức nửa nhóm; tích phân theo $y$ rồi dùng $m(z,t)\le H(t)$:
\[m(x,t+s)=\int_0^LG(x,z,s)m(z,t)dz\le H(t)m(x,s).\]
Lấy maximum theo $x$: $H(t+s)\le H(t)H(s)<H(t)$ với mọi $s,t>0$, nên $H$ giảm chặt. Giữ $x_0=L/2$: dữ liệu $1$ liên tục gần $x_0$ cho $m(x_0,t)\to1$ khi $t\downarrow0$. Vì $m(x_0,t)\le H(t)\le1$, định lý kẹp cho $\lim_{t\downarrow0}H(t)=1$. Không yêu cầu hội tụ đều của $m$ đến $1$ ở hai đầu, vì $m$ luôn bằng $0$ tại đó.

\clearpage\noindent\texttt{\detokenize{BT_4_3_3_02}}\par
\section*{Lời giải bổ sung — Bài tập 2 · Mục 4.3.3}

\subsection*{Kiến thức và trình tự giải}
\textbf{Bước 1 — Tìm cơ sở và chuẩn.} Giải $-X_n''=\lambda_nX_n$ với đúng biên thuần nhất. Đặt $N_n=\int_0^LX_n^2dx$; các hàm riêng ở biên Robin không có chuẩn mặc định $L/2$. Từ $c_n=N_n^{-1}\int fX_n$ trong chuỗi nghiệm, đổi thứ tự tổng–tích phân với $t>0$ để nhận ra


\[
G(x,y,t)=\sum_n\frac{X_n(x)X_n(y)}{N_n}e^{-K\lambda_nt},\qquad
u(x,t)=\int_0^LG(x,y,t)f(y)dy.
\]


\textbf{Bước 2 — Tính chuẩn Robin đến biểu thức sơ cấp.} Nếu $X=\cos kx+\sin kx/k$, dùng


\[
N=\frac L2(1+k^{-2})+\frac{\sin(2kL)}{4k}(1-k^{-2})+\frac{\sin^2(kL)}{k^2}.
\]


Nếu $X=\cos kx$, $N=L/2+\sin(2kL)/(4k)$. Hai biên tuần hoàn có mode hằng $1/L$ trong $G$ và cặp sine/cosine tần số $2n\pi/L$.

\textbf{Bước 3 — Kiểm tra PDE, biên và đối xứng.} Đạo hàm $e^{-K\lambda_nt}X_n(x)$ cho $G_t=KG_{xx}$; đạo hàm theo $y$ tương tự. Biên được kiểm tra trên từng $X_n$. Mỗi tích $X_n(x)X_n(y)$ đối xứng; hội tụ đều của chuỗi đạo hàm trên $t\ge\varepsilon$ biện minh các thao tác. Dữ liệu delta ban đầu hiểu qua $\int Gf\to f$ trong miền hoặc trong $L^2$, không bằng cách thế $t=0$ trực tiếp vào một chuỗi không hội tụ.

\textbf{Bước 4 — Chứng minh dương và khối lượng.} Nguyên lý cực đại cho $\int Gf\ge0$ với $f\ge0$, suy $G\ge0$ bằng cách lấy $f$ tập trung gần một điểm. Kết hợp tính nửa nhóm và nguyên lý mạnh để suy $G>0$ trong miền. $m(x,t)=\int G(x,y,t)dy$ là nghiệm có dữ liệu đầu $1$. Neumann/tuần hoàn giữ $m=1$; Dirichlet hoặc Robin hấp thụ cho $0<m<1$.

\textbf{Bước 5 — Nửa nhóm và thời gian dài.} Trực giao cho $\int G(x,z,s)G(z,y,t)dz=G(x,y,s+t)$. Với $H(t)=\max_xm(x,t)$ ở Dirichlet, $m(x,t+s)\le H(t)m(x,s)$, nên $H(t+s)\le H(t)H(s)<H(t)$. Để chứng minh $H(t)\to1$ lúc đầu, giữ một $x_0$ trong miền, dùng $m(x_0,t)\to1$ và $m(x_0,t)\le H(t)\le1$.
\subsection*{Lời giải áp dụng cho bài}
\textbf{Bước 1. Tách biến và giải từng điều kiện biên.}
Đặt $X=A\cos kx+B\sin kx$, $-X''=k^2X$.
(a) $X'(0)=0$ cho $B=0$, $X(L)=0$ cho $\cos kL=0$. Vậy $k_n=(n+1/2)\pi/L$, $X_n=\cos k_nx$, $N_n=L/2$, $n\ge0$:
\[G_a=\frac2L\sum_{n\ge0}e^{-Kk_n^2t}\cos k_nx\cos k_ny.\]
(b) $X(0)=X'(0)$ cho $A=kB$, chọn $X=\cos kx+\sin kx/k$. Biên $X(L)=0$ cho $k\cos kL+\sin kL=0$.
(c) $X'(0)=0$ cho $X=\cos kx$, $X(L)+X'(L)=0$ cho $\cos kL-k\sin kL=0$.
(d) $X'(0)=X(0)$ cho $X=\cos kx+\sin kx/k$, và $X(L)+X'(L)=0$ cho $(1-k^2)\sin kL+2k\cos kL=0$.

\textbf{Bước 2. Chuẩn hóa các biên Robin đến công thức rõ.}
Ở (b),(d), dùng $\cos^2kx=(1+\cos2kx)/2$, $\sin^2kx=(1-\cos2kx)/2$ và $2\sin kx\cos kx=\sin2kx$:
\[N(k)=\int_0^LX^2dx
=\frac L2(1+k^{-2})+\frac{\sin2kL}{4k}(1-k^{-2})+\frac{\sin^2kL}{k^2}.\]
Ở (c), $N(k)=L/2+\sin2kL/(4k)$. Vì vậy trong mỗi trường hợp (b),(c),(d),
\[G_j(x,y,t)=\sum_{k_n>0}\frac{X_n(x)X_n(y)}{N(k_n)}e^{-Kk_n^2t},\quad j=b,c,d,\]
với $k_n,X_n$ vừa xác định. Đây là các chuẩn đã tính, không bỏ lại $\int X_n^2$ chưa thực hiện. Trị riêng $0$ không có: $X=A+Bx$ thay vào từng cặp biên buộc $A=B=0$. Bỏ nghiệm $k=0$ ngoại lai của phương trình ở (d).

\textbf{Bước 3. Biên tuần hoàn (e).}
Các mode $1,\cos(2n\pi x/L),\sin(2n\pi x/L)$ có chuẩn $L,L/2,L/2$. Cộng hai tích bằng $\cos(a-b)$:
\[G_e=\frac1L+\frac2L\sum_{n\ge1}e^{-K(2n\pi/L)^2t}\cos\frac{2n\pi(x-y)}L.\]
Biểu diễn ảnh theo một chu kỳ là $G_e=\sum_{j\in\mathbb Z}g_K(x-y+jL,t)$; mỗi hạng dương và $\int_0^LG_e\,dy=\int_{\mathbb R}g_K\,dy=1$.

\textbf{Bước 4. Chứng minh các tính chất cho từng nhân.}
Trên $t\ge\varepsilon$, các hệ có $k_n$ tăng tuyến tính và chuẩn không suy biến, nên chuỗi đạo hàm hội tụ đều. Từ $X_n''=-k_n^2X_n$ suy $G_t=KG_{xx}=KG_{yy}$; từng $X_n$ cho đúng cặp biên (a)--(e); tích $X_n(x)X_n(y)$ cho đối xứng. Định lý khai triển Sturm–Liouville tự liên hợp ở (a)--(d), hoặc Fourier tuần hoàn ở (e), cho $\int G_jf\to f$ trong $L^2$ và tại điểm liên tục thích hợp. Trực giao cho tính nửa nhóm.

Nguyên lý cực đại ở từng biên cho $\int G_jf\ge0$ nếu $f\ge0$, suy $G_j\ge0$ bằng dữ liệu tập trung. Cùng lập luận nửa nhóm và nguyên lý mạnh cho $G_j>0$ trong miền. $m_j=\int G_jdy$ có dữ liệu đầu $1$: ở (a)--(d), biên Dirichlet hoặc Robin hấp thụ làm $0<m_j<1$ trong miền, còn (e) giữ $m_e=1$. Tất cả mode ở (a)--(d) dương, nên $G_j\to0$; ở (e), $G_e\to1/L$. Do đó nghiệm lần lượt tiến về $0$ hoặc trung bình dữ liệu. Các kết luận về dương, khối lượng và giới hạn được chứng minh bằng đúng biên của từng trường hợp.

\clearpage\noindent\texttt{\detokenize{VD_4_4_2_01}}\par
\section*{Lời giải bổ sung — Ví dụ 1 · Mục 4.4.2}

\subsection*{Kiến thức và trình tự giải}
\textbf{Bước 1 — Khử điều kiện biên không thuần nhất.} Viết $u=b+w$ với $b$ thỏa đúng biên. Dirichlet hai đầu: $b=\mu+(\nu-\mu)x/L$. Neumann hai đầu: $b=\mu x+(\nu-\mu)x^2/(2L)$. Neumann–Dirichlet: $b=\nu+\mu(x-L)$; Dirichlet–Neumann: $b=\mu+x\nu$. Robin trái–Dirichlet phải: $b=\mu+d+dx$, $d=(\nu-\mu)/(L+1)$. Neumann trái–Robin phải: $b=\nu-\mu(L+1)+\mu x$. Hai Robin: $b=\mu+d+dx$, $d=(\nu-\mu)/(L+2)$.

\textbf{Bước 2 — Tính lại nguồn và dữ liệu đầu.} Từ $u_t=Ku_{xx}+F$ suy $w_t=Kw_{xx}+Q$, $Q=F+Kb_{xx}-b_t$, $w(x,0)=f(x)-b(x,0)$. Tính $b_t,b_{xx}$ trước khi chiếu; với hàm nâng bậc hai Neumann, $Kb_{xx}$ không bằng không. Nếu dùng trạng thái dừng $P$, giải $KP''+F=0$ và đúng biên rồi đặt $w=u-P$.

\textbf{Bước 3 — Chiếu và giải ODE bằng nhân tử.} Với $N_n=\int X_n^2$, đặt $d_n=N_n^{-1}\int(f-b_0)X_n$, $Q_n(t)=N_n^{-1}\int QX_n$. Phương trình mode là $T_n'+K\lambda_nT_n=Q_n$. Nhân với $e^{K\lambda_nt}$, tích phân từ $0$ đến $t$, rồi chia:


\[
T_n=d_ne^{-K\lambda_nt}+\int_0^te^{-K\lambda_n(t-s)}Q_n(s)ds,\qquad
u=b+\sum_nT_nX_n.
\]


Đây cũng là công thức Green $u=b+\int G(f-b_0)+\int_0^t\int GQ$ cho biên đã được khử.

\textbf{Bước 4 — Thực hiện tích phân thời gian.} Với $Q_n(s)=C_ne^{-\rho s}$, đặt $d=K\lambda_n$:


\[
\int_0^te^{-d(t-s)}C_ne^{-\rho s}ds
=C_ne^{-dt}\int_0^te^{(d-\rho)s}ds
=C_n\frac{e^{-\rho t}-e^{-dt}}{d-\rho}.
\]


Nếu $d=\rho$, kết quả là $C_nte^{-dt}$. Với nguồn hằng, kết quả $C_n(1-e^{-dt})/d$ khi $d>0$; mode $d=0$ phải dùng $\int_0^tQ_0(s)ds$. Nguồn đa thức thời gian tích phân từng phần hoặc giải ODE với nghiệm riêng đa thức.

\textbf{Bước 5 — Kiểm tra và tìm giới hạn.} Mỗi mode đúng PDE, biên; $T_n(0)=d_n$ trả lại dữ liệu đầu sau khi cộng $b_0$. Tách mode hằng Neumann: $\frac d{dt}\int u=K(\nu-\mu)+\int F$. Trạng thái dừng chỉ tồn tại nếu tổng nguồn và thông lượng cân bằng. Với các mode dương, kiểm tra hội tụ đủ để lấy giới hạn đều; không bỏ một mode bằng không hoặc mode tăng.
\subsection*{Lời giải áp dụng cho bài}
\textbf{Bước 1. Công thức mode có nguồn.}
Với Dirichlet thuần nhất trên $(0,\pi)$, đặt $u=\sum T_n\sin nx$. Chiếu PDE cho $T_n'+n^2T_n=F_n$, $T_n(0)=f_n$. Nhân với $e^{n^2t}$ rồi tích phân:
\[T_n=f_ne^{-n^2t}+\int_0^te^{-n^2(t-s)}F_n(s)ds.\]
\textbf{Bước 2. Ví dụ 1(i), $F=1,f=0$.}
\[F_n=\frac2\pi[-\cos ny/n]_0^\pi=\frac{2(1-(-1)^n)}{\pi n}.\]
Vì $\int_0^te^{-n^2(t-s)}ds=(1-e^{-n^2t})/n^2$,
\[u=\sum_{n\ge1}\frac{2(1-(-1)^n)}{\pi n^3}(1-e^{-n^2t})\sin nx.\]
Trạng thái dừng giải $P''=-1$, $P(0)=P(\pi)=0$, nên $P=x(\pi-x)/2$. Hệ số sine của $P$ chính là $F_n/n^2$; do đó $u\to P$, không tiến về hằng $1$. Đồ thị trong sách cần được đối chiếu với nghiệm này.

\textbf{Bước 3. Ví dụ 1(ii), biên $T_1,T_2$, nguồn và dữ liệu đầu bằng không.}
Đặt $b=T_1+(T_2-T_1)x/\pi$ và $w=u-b$. $b_t=b_{xx}=0$, nên $w_t=w_{xx}$, biên $0$, dữ liệu đầu $-b$. Tích phân từng phần:
\[\int_0^\pi b\sin nx\,dx
=\left[-\frac{b\cos nx}n\right]_0^\pi+\frac{b'}n\int_0^\pi\cos nx\,dx
=\frac{T_1-(-1)^nT_2}n.\]
Vì vậy
\[u=b-\sum_{n\ge1}\frac{2(T_1-(-1)^nT_2)}{\pi n}e^{-n^2t}\sin nx.\]
Biên được bảo đảm bởi $b$, mặt đầu trong miền bằng $0$ theo khai triển sine của $b$, và $u\to b$ đều lúc thời gian dài. Nếu $T_1$ hoặc $T_2$ khác $0$, không có liên tục tại góc tương ứng với dữ liệu đầu $0$.

\textbf{Bước 4. Ví dụ 2 nằm tiếp trong cùng phần scan: Neumann có nguồn.}
Với $u_x(0)=\mu,u_x(L)=\nu$, chọn $b=\mu x+(\nu-\mu)x^2/(2L)$. Khi đó $w=u-b$ có biên Neumann $0$, nguồn $Q=F+K(\nu-\mu)/L-b_t$, dữ liệu $f-b_0$. Nhân Green Neumann cho $w=\int G_N(f-b_0)+\int_0^t\int G_NQ$. Cũng có công thức trực tiếp
\[u=\int_0^LG_N(x,y,t)f(y)dy+
\int_0^t\!\int_0^LG_N(x,y,t-s)F(y,s)dy\,ds+
K\int_0^t\{G_N(x,L,t-s)\nu(s)-G_N(x,0,t-s)\mu(s)\}ds.\]
Dấu hai đầu được suy từ $K[G_Nu_y-(G_N)_yu]_0^L$.
Với $K=2,L=1,\mu=\nu=f=0,F=e^{-t}\cos^2\pi x$, dùng $\cos^2\pi x=(1+\cos2\pi x)/2$. Chỉ có mode $0,2$:
\[T_0'=\tfrac12e^{-t},\quad T_2'+8\pi^2T_2=\tfrac12e^{-t},\quad T_0(0)=T_2(0)=0.\]
Tích phân cho
\[u=\frac{1-e^{-t}}2+\frac{e^{-t}-e^{-8\pi^2t}}{2(8\pi^2-1)}\cos2\pi x.\]
Đúng biên, nguồn, mặt đầu; giới hạn bằng $1/2$. Hai ví dụ có đánh số 1 và 2 được giữ chung ở đây vì cùng nằm trong ảnh đề của mục này.
\begin{center}\includegraphics[width=.97\linewidth,height=11cm,keepaspectratio]{_xu_ly/graphs/VD_4_4_2_01.png}\end{center}
\textit{Đồ thị minh họa; với nghiệm số, xem phương pháp và tham số trong mã Python kèm theo.}

\clearpage\noindent\texttt{\detokenize{BT_4_4_3_01}}\par
\section*{Lời giải bổ sung — Bài tập 1 · Mục 4.4.3}

\subsection*{Kiến thức và trình tự giải}
\textbf{Bước 1 — Khử điều kiện biên không thuần nhất.} Viết $u=b+w$ với $b$ thỏa đúng biên. Dirichlet hai đầu: $b=\mu+(\nu-\mu)x/L$. Neumann hai đầu: $b=\mu x+(\nu-\mu)x^2/(2L)$. Neumann–Dirichlet: $b=\nu+\mu(x-L)$; Dirichlet–Neumann: $b=\mu+x\nu$. Robin trái–Dirichlet phải: $b=\mu+d+dx$, $d=(\nu-\mu)/(L+1)$. Neumann trái–Robin phải: $b=\nu-\mu(L+1)+\mu x$. Hai Robin: $b=\mu+d+dx$, $d=(\nu-\mu)/(L+2)$.

\textbf{Bước 2 — Tính lại nguồn và dữ liệu đầu.} Từ $u_t=Ku_{xx}+F$ suy $w_t=Kw_{xx}+Q$, $Q=F+Kb_{xx}-b_t$, $w(x,0)=f(x)-b(x,0)$. Tính $b_t,b_{xx}$ trước khi chiếu; với hàm nâng bậc hai Neumann, $Kb_{xx}$ không bằng không. Nếu dùng trạng thái dừng $P$, giải $KP''+F=0$ và đúng biên rồi đặt $w=u-P$.

\textbf{Bước 3 — Chiếu và giải ODE bằng nhân tử.} Với $N_n=\int X_n^2$, đặt $d_n=N_n^{-1}\int(f-b_0)X_n$, $Q_n(t)=N_n^{-1}\int QX_n$. Phương trình mode là $T_n'+K\lambda_nT_n=Q_n$. Nhân với $e^{K\lambda_nt}$, tích phân từ $0$ đến $t$, rồi chia:


\[
T_n=d_ne^{-K\lambda_nt}+\int_0^te^{-K\lambda_n(t-s)}Q_n(s)ds,\qquad
u=b+\sum_nT_nX_n.
\]


Đây cũng là công thức Green $u=b+\int G(f-b_0)+\int_0^t\int GQ$ cho biên đã được khử.

\textbf{Bước 4 — Thực hiện tích phân thời gian.} Với $Q_n(s)=C_ne^{-\rho s}$, đặt $d=K\lambda_n$:


\[
\int_0^te^{-d(t-s)}C_ne^{-\rho s}ds
=C_ne^{-dt}\int_0^te^{(d-\rho)s}ds
=C_n\frac{e^{-\rho t}-e^{-dt}}{d-\rho}.
\]


Nếu $d=\rho$, kết quả là $C_nte^{-dt}$. Với nguồn hằng, kết quả $C_n(1-e^{-dt})/d$ khi $d>0$; mode $d=0$ phải dùng $\int_0^tQ_0(s)ds$. Nguồn đa thức thời gian tích phân từng phần hoặc giải ODE với nghiệm riêng đa thức.

\textbf{Bước 5 — Kiểm tra và tìm giới hạn.} Mỗi mode đúng PDE, biên; $T_n(0)=d_n$ trả lại dữ liệu đầu sau khi cộng $b_0$. Tách mode hằng Neumann: $\frac d{dt}\int u=K(\nu-\mu)+\int F$. Trạng thái dừng chỉ tồn tại nếu tổng nguồn và thông lượng cân bằng. Với các mode dương, kiểm tra hội tụ đủ để lấy giới hạn đều; không bỏ một mode bằng không hoặc mode tăng.
\subsection*{Lời giải áp dụng cho bài}
\textbf{Bước 1. Khử biên Dirichlet và dùng Duhamel.}
Đặt $b=\mu+(\nu-\mu)x/L$, $w=u-b$. Khi đó $w_t=Kw_{xx}+F-b_t$, biên $0$, mặt đầu $f-b_0$. Với $k_n=n\pi/L$ và hệ số sine chuẩn $q_n=(2/L)\int q\sin k_nx$,
\[u=b+\sum_{n\ge1}\left[(f-b_0)_ne^{-Kk_n^2t}
+\int_0^te^{-Kk_n^2(t-s)}(F-b_t)_n(s)ds\right]\sin k_nx.\]

\textbf{Bước 2. Câu (i).} $K=1,L=2,b=1+x$, $q=f-b=x^2-2x-2$, $k=n\pi/2$. Tích phân hai lần cho
\begin{align*}
d_n=\int_0^2q\sin kx\,dx
&=\left[-q\cos kx/k\right]_0^2-\frac1{k^2}\int_0^2q''\sin kx\,dx\\
&=\frac{2((-1)^n-1)}k-\frac{2(1-(-1)^n)}{k^3}
=-2(1-(-1)^n)(k^{-1}+k^{-3}).
\end{align*}
Nguồn bằng $0$, nên
\[u=1+x-\sum_{n\ge1}2(1-(-1)^n)(k_n^{-1}+k_n^{-3})e^{-k_n^2t}\sin k_nx.\]
Thời gian dài các mode mất đi đều: $u\to1+x$. $f(0)=-1\ne1$, $f(2)=1\ne3$, nên không liên tục tại hai góc; chuỗi trả lại $f$ trong miền và biên được giữ bởi $b$ khi $t>0$.

\textbf{Bước 3. Câu (ii).} $b=0,f=2x-x^2,F=1-|x-1|$, $k=n\pi/2$. Vì $f(0)=f(2)=0,f''=-2$,
\[f_n=-\frac1{k^2}\int_0^2f''\sin kx\,dx=\frac{2(1-(-1)^n)}{k^3}.\]
Tách nguồn tại $1$, dùng nguyên hàm của $x\sin kx$:
\begin{align*}
F_n&=\int_0^1x\sin kx\,dx+\int_1^2(2-x)\sin kx\,dx\\
&=\left[-\frac{x\cos kx}k+\frac{\sin kx}{k^2}\right]_0^1
+\left[-\frac{(2-x)\cos kx}k-\frac{\sin kx}{k^2}\right]_1^2
=\frac{2\sin k}{k^2}.
\end{align*}
Do nguồn hằng theo thời gian,
\[u=\sum_{n\ge1}\left[\frac{2(1-(-1)^n)}{k_n^3}e^{-k_n^2t}
+\frac{2\sin k_n}{k_n^4}(1-e^{-k_n^2t})\right]\sin k_nx.\]
Trạng thái dừng giải $-P''=F$, biên $0$. Đối xứng quanh $1$ cho
\[P(x)=\begin{cases}x/2-x^3/6,&0\le x\le1,\\(2-x)/2-(2-x)^3/6,&1\le x\le2.\end{cases}\]
Hai vết $P,P'$ tại $1$ bằng nhau; $P''=-F$, $P(0)=P(2)=0$. Các hệ số dừng bằng $F_n/k_n^2$, nên $u\to P$ đều.

\clearpage\noindent\texttt{\detokenize{BT_4_4_3_02}}\par
\section*{Lời giải bổ sung — Bài tập 2 · Mục 4.4.3}

\subsection*{Kiến thức và trình tự giải}
\textbf{Bước 1 — Khử điều kiện biên không thuần nhất.} Viết $u=b+w$ với $b$ thỏa đúng biên. Dirichlet hai đầu: $b=\mu+(\nu-\mu)x/L$. Neumann hai đầu: $b=\mu x+(\nu-\mu)x^2/(2L)$. Neumann–Dirichlet: $b=\nu+\mu(x-L)$; Dirichlet–Neumann: $b=\mu+x\nu$. Robin trái–Dirichlet phải: $b=\mu+d+dx$, $d=(\nu-\mu)/(L+1)$. Neumann trái–Robin phải: $b=\nu-\mu(L+1)+\mu x$. Hai Robin: $b=\mu+d+dx$, $d=(\nu-\mu)/(L+2)$.

\textbf{Bước 2 — Tính lại nguồn và dữ liệu đầu.} Từ $u_t=Ku_{xx}+F$ suy $w_t=Kw_{xx}+Q$, $Q=F+Kb_{xx}-b_t$, $w(x,0)=f(x)-b(x,0)$. Tính $b_t,b_{xx}$ trước khi chiếu; với hàm nâng bậc hai Neumann, $Kb_{xx}$ không bằng không. Nếu dùng trạng thái dừng $P$, giải $KP''+F=0$ và đúng biên rồi đặt $w=u-P$.

\textbf{Bước 3 — Chiếu và giải ODE bằng nhân tử.} Với $N_n=\int X_n^2$, đặt $d_n=N_n^{-1}\int(f-b_0)X_n$, $Q_n(t)=N_n^{-1}\int QX_n$. Phương trình mode là $T_n'+K\lambda_nT_n=Q_n$. Nhân với $e^{K\lambda_nt}$, tích phân từ $0$ đến $t$, rồi chia:


\[
T_n=d_ne^{-K\lambda_nt}+\int_0^te^{-K\lambda_n(t-s)}Q_n(s)ds,\qquad
u=b+\sum_nT_nX_n.
\]


Đây cũng là công thức Green $u=b+\int G(f-b_0)+\int_0^t\int GQ$ cho biên đã được khử.

\textbf{Bước 4 — Thực hiện tích phân thời gian.} Với $Q_n(s)=C_ne^{-\rho s}$, đặt $d=K\lambda_n$:


\[
\int_0^te^{-d(t-s)}C_ne^{-\rho s}ds
=C_ne^{-dt}\int_0^te^{(d-\rho)s}ds
=C_n\frac{e^{-\rho t}-e^{-dt}}{d-\rho}.
\]


Nếu $d=\rho$, kết quả là $C_nte^{-dt}$. Với nguồn hằng, kết quả $C_n(1-e^{-dt})/d$ khi $d>0$; mode $d=0$ phải dùng $\int_0^tQ_0(s)ds$. Nguồn đa thức thời gian tích phân từng phần hoặc giải ODE với nghiệm riêng đa thức.

\textbf{Bước 5 — Kiểm tra và tìm giới hạn.} Mỗi mode đúng PDE, biên; $T_n(0)=d_n$ trả lại dữ liệu đầu sau khi cộng $b_0$. Tách mode hằng Neumann: $\frac d{dt}\int u=K(\nu-\mu)+\int F$. Trạng thái dừng chỉ tồn tại nếu tổng nguồn và thông lượng cân bằng. Với các mode dương, kiểm tra hội tụ đủ để lấy giới hạn đều; không bỏ một mode bằng không hoặc mode tăng.
\subsection*{Lời giải áp dụng cho bài}
\textbf{Bước 1. Công thức Neumann có nguồn và biên không thuần nhất.}
Chọn $b=\mu x+(\nu-\mu)x^2/(2L)$, $w=u-b$, $Q=F+K(\nu-\mu)/L-b_t$, $q=f-b_0$. Đặt $q_0=L^{-1}\int q$, $q_n=(2/L)\int q\cos(n\pi x/L)$, tương tự $Q_n$. Khi đó
\[u=b+q_0+\int_0^tQ_0(s)ds+
\sum_{n\ge1}\left[q_ne^{-Kk_n^2t}+
\int_0^te^{-Kk_n^2(t-s)}Q_n(s)ds\right]\cos k_nx.\]
Mode $0$ không được bỏ.

\textbf{Bước 2. Câu (i), $K=1,L=\pi,F=\cos x,f=\cos x+\cos2x$, biên $0$.}
Trực giao cosine cho $f_1=f_2=1$, $F_1=1$, các hệ số khác bằng $0$. Hai ODE là $T_1'+T_1=1,T_1(0)=1$ và $T_2'+4T_2=0,T_2(0)=1$. Suy $T_1=1,T_2=e^{-4t}$:
\[u=\cos x+e^{-4t}\cos2x,\qquad u\to\cos x.\]
Kiểm tra $u_t=-4e^{-4t}\cos2x=u_{xx}+\cos x$, $u_x(0)=u_x(\pi)=0$, mặt đầu đúng.
Nếu thay đạo hàm trái bằng $1$, phải đọc đầu phải là $\pi$ theo miền; dòng $2\pi$ trong đề không thuộc đoạn. Khi đó
\[\frac d{dt}\int_0^\pi u\,dx=[u_x]_0^\pi+\int_0^\pi\cos x\,dx=-1.\]
Tổng nhiệt giảm tuyến tính, không có giới hạn dừng hữu hạn. Tương đương, bài dừng $P''+\cos x=0$ buộc $P'(\pi)-P'(0)=0$, trái với $0-1=-1$.

\textbf{Bước 3. Câu (ii), $\mu=\nu=A,F=0$.}
Đặt $b=Ax$. $b_t=b_{xx}=0$, nên $w=u-Ax$ giải Neumann thuần nhất với $q=f-Ax$. Giới hạn đều của $w$ là $q_0=L^{-1}\int_0^L(f(y)-Ay)dy$. Do đó
\[v(x)=Ax+\frac1L\int_0^Lf(y)dy-\frac{AL}2.\]
Tính $v''=0$, $v'(0)=v'(L)=A$, đúng bài dừng. Hằng được xác định bởi bảo toàn $\int u=\int f$; các nghiệm dừng khác $Ax+C$ có tổng nhiệt khác và không phải giới hạn của cùng dữ liệu.

\clearpage\noindent\texttt{\detokenize{BT_4_4_3_03}}\par
\section*{Lời giải bổ sung — Bài tập 3 · Mục 4.4.3}

\subsection*{Kiến thức và trình tự giải}
\textbf{Bước 1 — Khử điều kiện biên không thuần nhất.} Viết $u=b+w$ với $b$ thỏa đúng biên. Dirichlet hai đầu: $b=\mu+(\nu-\mu)x/L$. Neumann hai đầu: $b=\mu x+(\nu-\mu)x^2/(2L)$. Neumann–Dirichlet: $b=\nu+\mu(x-L)$; Dirichlet–Neumann: $b=\mu+x\nu$. Robin trái–Dirichlet phải: $b=\mu+d+dx$, $d=(\nu-\mu)/(L+1)$. Neumann trái–Robin phải: $b=\nu-\mu(L+1)+\mu x$. Hai Robin: $b=\mu+d+dx$, $d=(\nu-\mu)/(L+2)$.

\textbf{Bước 2 — Tính lại nguồn và dữ liệu đầu.} Từ $u_t=Ku_{xx}+F$ suy $w_t=Kw_{xx}+Q$, $Q=F+Kb_{xx}-b_t$, $w(x,0)=f(x)-b(x,0)$. Tính $b_t,b_{xx}$ trước khi chiếu; với hàm nâng bậc hai Neumann, $Kb_{xx}$ không bằng không. Nếu dùng trạng thái dừng $P$, giải $KP''+F=0$ và đúng biên rồi đặt $w=u-P$.

\textbf{Bước 3 — Chiếu và giải ODE bằng nhân tử.} Với $N_n=\int X_n^2$, đặt $d_n=N_n^{-1}\int(f-b_0)X_n$, $Q_n(t)=N_n^{-1}\int QX_n$. Phương trình mode là $T_n'+K\lambda_nT_n=Q_n$. Nhân với $e^{K\lambda_nt}$, tích phân từ $0$ đến $t$, rồi chia:


\[
T_n=d_ne^{-K\lambda_nt}+\int_0^te^{-K\lambda_n(t-s)}Q_n(s)ds,\qquad
u=b+\sum_nT_nX_n.
\]


Đây cũng là công thức Green $u=b+\int G(f-b_0)+\int_0^t\int GQ$ cho biên đã được khử.

\textbf{Bước 4 — Thực hiện tích phân thời gian.} Với $Q_n(s)=C_ne^{-\rho s}$, đặt $d=K\lambda_n$:


\[
\int_0^te^{-d(t-s)}C_ne^{-\rho s}ds
=C_ne^{-dt}\int_0^te^{(d-\rho)s}ds
=C_n\frac{e^{-\rho t}-e^{-dt}}{d-\rho}.
\]


Nếu $d=\rho$, kết quả là $C_nte^{-dt}$. Với nguồn hằng, kết quả $C_n(1-e^{-dt})/d$ khi $d>0$; mode $d=0$ phải dùng $\int_0^tQ_0(s)ds$. Nguồn đa thức thời gian tích phân từng phần hoặc giải ODE với nghiệm riêng đa thức.

\textbf{Bước 5 — Kiểm tra và tìm giới hạn.} Mỗi mode đúng PDE, biên; $T_n(0)=d_n$ trả lại dữ liệu đầu sau khi cộng $b_0$. Tách mode hằng Neumann: $\frac d{dt}\int u=K(\nu-\mu)+\int F$. Trạng thái dừng chỉ tồn tại nếu tổng nguồn và thông lượng cân bằng. Với các mode dương, kiểm tra hội tụ đủ để lấy giới hạn đều; không bỏ một mode bằng không hoặc mode tăng.
\subsection*{Lời giải áp dụng cho bài}
\textbf{Bước 1. Biên Neumann–Dirichlet và Green.}
Chọn $b=\nu+\mu(x-L)$, có $b_x(0)=\mu,b(L)=\nu$, $b_{xx}=0$. Khi $w=u-b$, nguồn $Q=F-b_t$, dữ liệu $q=f-b_0$, cơ sở $X_n=\cos k_nx$, $k_n=(n+1/2)\pi/L$, $N_n=L/2$. Nhân $G=\sum X_n(x)X_n(y)e^{-Kk_n^2t}/N_n$ cho
\[u=b+\int_0^LGq+\int_0^t\!\int_0^LG(x,y,t-s)Q(y,s)dy\,ds.\]
Công thức trực tiếp, suy bằng $K[Gu_y-G_yu]_0^L$, có thêm các biên $-K\int_0^t\{G(x,0,t-s)\mu(s)+G_y(x,L,t-s)\nu(s)\}ds$ vào phần dữ liệu đầu và nguồn của $u$.

\textbf{Bước 2. Tính hệ số câu cụ thể.}
$K=2,L=1,\mu=\nu=f=0,F=e^{-t}\cos^2\pi x$, nên $b=0$. Với $k=(n+1/2)\pi$, dùng tích–tổng:
\begin{align*}
C_n=2\int_0^1\cos^2\pi y\cos ky\,dy
&=\frac{\sin k}k+\frac12\left(\frac{\sin(k+2\pi)}{k+2\pi}
+\frac{\sin(k-2\pi)}{k-2\pi}\right)\\
&=\frac{2\sin k(k^2-2\pi^2)}{k(k^2-4\pi^2)}
=\frac{2(-1)^n(k^2-2\pi^2)}{k(k^2-4\pi^2)}.
\end{align*}
\textbf{Bước 3. Thực hiện tích phân thời gian.}
$T_n'+2k_n^2T_n=C_ne^{-t},T_n(0)=0$, nên
\[T_n=C_ne^{-2k_n^2t}\int_0^te^{(2k_n^2-1)s}ds
=C_n\frac{e^{-t}-e^{-2k_n^2t}}{2k_n^2-1}.\]
Vậy $u=\sum_{n\ge0}T_n\cos k_nx$. Mỗi cosine có $X_n'(0)=0,X_n(1)=0$; ODE chứng minh đúng nguồn và $T_n(0)=0$ trả lại mặt đầu. Hai số mũ đều âm, nên $u\to0$ đều.

\clearpage\noindent\texttt{\detokenize{BT_4_4_3_04}}\par
\section*{Lời giải bổ sung — Bài tập 4 · Mục 4.4.3}

\subsection*{Kiến thức và trình tự giải}
\textbf{Bước 1 — Khử điều kiện biên không thuần nhất.} Viết $u=b+w$ với $b$ thỏa đúng biên. Dirichlet hai đầu: $b=\mu+(\nu-\mu)x/L$. Neumann hai đầu: $b=\mu x+(\nu-\mu)x^2/(2L)$. Neumann–Dirichlet: $b=\nu+\mu(x-L)$; Dirichlet–Neumann: $b=\mu+x\nu$. Robin trái–Dirichlet phải: $b=\mu+d+dx$, $d=(\nu-\mu)/(L+1)$. Neumann trái–Robin phải: $b=\nu-\mu(L+1)+\mu x$. Hai Robin: $b=\mu+d+dx$, $d=(\nu-\mu)/(L+2)$.

\textbf{Bước 2 — Tính lại nguồn và dữ liệu đầu.} Từ $u_t=Ku_{xx}+F$ suy $w_t=Kw_{xx}+Q$, $Q=F+Kb_{xx}-b_t$, $w(x,0)=f(x)-b(x,0)$. Tính $b_t,b_{xx}$ trước khi chiếu; với hàm nâng bậc hai Neumann, $Kb_{xx}$ không bằng không. Nếu dùng trạng thái dừng $P$, giải $KP''+F=0$ và đúng biên rồi đặt $w=u-P$.

\textbf{Bước 3 — Chiếu và giải ODE bằng nhân tử.} Với $N_n=\int X_n^2$, đặt $d_n=N_n^{-1}\int(f-b_0)X_n$, $Q_n(t)=N_n^{-1}\int QX_n$. Phương trình mode là $T_n'+K\lambda_nT_n=Q_n$. Nhân với $e^{K\lambda_nt}$, tích phân từ $0$ đến $t$, rồi chia:


\[
T_n=d_ne^{-K\lambda_nt}+\int_0^te^{-K\lambda_n(t-s)}Q_n(s)ds,\qquad
u=b+\sum_nT_nX_n.
\]


Đây cũng là công thức Green $u=b+\int G(f-b_0)+\int_0^t\int GQ$ cho biên đã được khử.

\textbf{Bước 4 — Thực hiện tích phân thời gian.} Với $Q_n(s)=C_ne^{-\rho s}$, đặt $d=K\lambda_n$:


\[
\int_0^te^{-d(t-s)}C_ne^{-\rho s}ds
=C_ne^{-dt}\int_0^te^{(d-\rho)s}ds
=C_n\frac{e^{-\rho t}-e^{-dt}}{d-\rho}.
\]


Nếu $d=\rho$, kết quả là $C_nte^{-dt}$. Với nguồn hằng, kết quả $C_n(1-e^{-dt})/d$ khi $d>0$; mode $d=0$ phải dùng $\int_0^tQ_0(s)ds$. Nguồn đa thức thời gian tích phân từng phần hoặc giải ODE với nghiệm riêng đa thức.

\textbf{Bước 5 — Kiểm tra và tìm giới hạn.} Mỗi mode đúng PDE, biên; $T_n(0)=d_n$ trả lại dữ liệu đầu sau khi cộng $b_0$. Tách mode hằng Neumann: $\frac d{dt}\int u=K(\nu-\mu)+\int F$. Trạng thái dừng chỉ tồn tại nếu tổng nguồn và thông lượng cân bằng. Với các mode dương, kiểm tra hội tụ đủ để lấy giới hạn đều; không bỏ một mode bằng không hoặc mode tăng.
\subsection*{Lời giải áp dụng cho bài}
\textbf{Bước 1. Khử Robin trái–Dirichlet phải.}
Với $b=a+dx$, $b(0)-b_x(0)=a-d=\mu$, $b(L)=a+dL=\nu$, nên $d=(\nu-\mu)/(L+1)$, $a=\mu+d$. $w=u-b$ có nguồn $Q=F-b_t$, dữ liệu $q=f-b_0$.
$X=\cos kx+\sin kx/k$ thỏa $X(0)-X'(0)=0$; $X(L)=0$ cho $k\cos kL+\sin kL=0$. Chuẩn
\[N(k)=\frac L2(1+k^{-2})+\frac{\sin2kL}{4k}(1-k^{-2})+\frac{\sin^2kL}{k^2}.\]
Nhân Green $G=\sum XXe^{-Kk^2t}/N$ cho $u=b+\int Gq+\int_0^t\int GQ$. Công thức trực tiếp cho $u$ có hạng biên $K\int_0^t[G(x,0,t-s)\mu(s)-G_y(x,L,t-s)\nu(s)]ds$, vì $G_y(0)=G(0),G(L)=0$.

\textbf{Bước 2. Tính hệ số cụ thể $K=2,L=1$, các biên và mặt đầu bằng $0$.}
Đặt
\[A(k)=\int_0^1\cos^2\pi y\cos ky\,dy
=\frac{\sin k(k^2-2\pi^2)}{k(k^2-4\pi^2)},\quad
B(k)=\int_0^1\cos^2\pi y\sin ky\,dy
=\frac{(1-\cos k)(k^2-2\pi^2)}{k(k^2-4\pi^2)}.\]
Các đẳng thức có được bằng $\cos^2\pi y=(1+\cos2\pi y)/2$ và tích–tổng, rồi tích phân $\cos ry,\sin ry$ trên $0,1$. Các nghiệm riêng không bằng $0,2\pi$, nên các mẫu ở đây không bằng không. Hệ số nguồn là
\[C_n=\frac{A(k_n)+B(k_n)/k_n}{N(k_n)},\qquad k_n\cos k_n+\sin k_n=0.\]
\textbf{Bước 3. Giải ODE và ghép.}
\[u=\sum_{n\ge1}\frac{A(k_n)+B(k_n)/k_n}{N(k_n)}
\frac{e^{-t}-e^{-2k_n^2t}}{2k_n^2-1}
\left(\cos k_nx+\frac{\sin k_nx}{k_n}\right).\]
Thương thời gian xuất phát từ $e^{-2k_n^2t}\int_0^te^{(2k_n^2-1)s}ds$; nếu trùng tốc độ thì dùng $te^{-t}$. Mỗi mode thỏa hai biên; $T_n'+2k_n^2T_n=C_ne^{-t}$ cho đúng nguồn, $T_n(0)=0$ cho dữ liệu đầu. Tất cả tốc độ giảm dương, nên giới hạn bằng $0$.

\clearpage\noindent\texttt{\detokenize{BT_4_4_3_05}}\par
\section*{Lời giải bổ sung — Bài tập 5 · Mục 4.4.3}

\subsection*{Kiến thức và trình tự giải}
\textbf{Bước 1 — Khử điều kiện biên không thuần nhất.} Viết $u=b+w$ với $b$ thỏa đúng biên. Dirichlet hai đầu: $b=\mu+(\nu-\mu)x/L$. Neumann hai đầu: $b=\mu x+(\nu-\mu)x^2/(2L)$. Neumann–Dirichlet: $b=\nu+\mu(x-L)$; Dirichlet–Neumann: $b=\mu+x\nu$. Robin trái–Dirichlet phải: $b=\mu+d+dx$, $d=(\nu-\mu)/(L+1)$. Neumann trái–Robin phải: $b=\nu-\mu(L+1)+\mu x$. Hai Robin: $b=\mu+d+dx$, $d=(\nu-\mu)/(L+2)$.

\textbf{Bước 2 — Tính lại nguồn và dữ liệu đầu.} Từ $u_t=Ku_{xx}+F$ suy $w_t=Kw_{xx}+Q$, $Q=F+Kb_{xx}-b_t$, $w(x,0)=f(x)-b(x,0)$. Tính $b_t,b_{xx}$ trước khi chiếu; với hàm nâng bậc hai Neumann, $Kb_{xx}$ không bằng không. Nếu dùng trạng thái dừng $P$, giải $KP''+F=0$ và đúng biên rồi đặt $w=u-P$.

\textbf{Bước 3 — Chiếu và giải ODE bằng nhân tử.} Với $N_n=\int X_n^2$, đặt $d_n=N_n^{-1}\int(f-b_0)X_n$, $Q_n(t)=N_n^{-1}\int QX_n$. Phương trình mode là $T_n'+K\lambda_nT_n=Q_n$. Nhân với $e^{K\lambda_nt}$, tích phân từ $0$ đến $t$, rồi chia:


\[
T_n=d_ne^{-K\lambda_nt}+\int_0^te^{-K\lambda_n(t-s)}Q_n(s)ds,\qquad
u=b+\sum_nT_nX_n.
\]


Đây cũng là công thức Green $u=b+\int G(f-b_0)+\int_0^t\int GQ$ cho biên đã được khử.

\textbf{Bước 4 — Thực hiện tích phân thời gian.} Với $Q_n(s)=C_ne^{-\rho s}$, đặt $d=K\lambda_n$:


\[
\int_0^te^{-d(t-s)}C_ne^{-\rho s}ds
=C_ne^{-dt}\int_0^te^{(d-\rho)s}ds
=C_n\frac{e^{-\rho t}-e^{-dt}}{d-\rho}.
\]


Nếu $d=\rho$, kết quả là $C_nte^{-dt}$. Với nguồn hằng, kết quả $C_n(1-e^{-dt})/d$ khi $d>0$; mode $d=0$ phải dùng $\int_0^tQ_0(s)ds$. Nguồn đa thức thời gian tích phân từng phần hoặc giải ODE với nghiệm riêng đa thức.

\textbf{Bước 5 — Kiểm tra và tìm giới hạn.} Mỗi mode đúng PDE, biên; $T_n(0)=d_n$ trả lại dữ liệu đầu sau khi cộng $b_0$. Tách mode hằng Neumann: $\frac d{dt}\int u=K(\nu-\mu)+\int F$. Trạng thái dừng chỉ tồn tại nếu tổng nguồn và thông lượng cân bằng. Với các mode dương, kiểm tra hội tụ đủ để lấy giới hạn đều; không bỏ một mode bằng không hoặc mode tăng.
\subsection*{Lời giải áp dụng cho bài}
\textbf{Bước 1. Khử Neumann trái–Robin phải.}
Lấy $b=\nu-\mu(L+1)+\mu x$. Khi đó $b_x(0)=\mu$, $b(L)+b_x(L)=\nu$, $b_{xx}=0$; $w=u-b$ có nguồn $Q=F-b_t$, mặt đầu $q=f-b_0$. $X'(0)=0$ cho $X=\cos kx$, biên phải cho $\cos kL-k\sin kL=0$, chuẩn $N(k)=L/2+\sin2kL/(4k)$. Đặt $G=\sum X(x)X(y)e^{-Kk^2t}/N$. Công thức Green là $u=b+\int Gq+\int_0^t\int GQ$; nếu dùng dữ liệu gốc, hạng biên là $K\int_0^t[G(x,L,t-s)\nu(s)-G(x,0,t-s)\mu(s)]ds$.

\textbf{Bước 2. Tính nguồn cụ thể.}
$K=2,L=1,\mu=\nu=f=0$. Đặt $k_n>0$ là nghiệm $k\sin k=\cos k$. Từ tích–tổng,
\begin{align*}
A(k)&=\int_0^1\cos^2\pi y\cos ky\,dy\\
&=\frac{\sin k}{2k}+\frac14\left(\frac{\sin(k+2\pi)}{k+2\pi}+\frac{\sin(k-2\pi)}{k-2\pi}\right)
=\frac{\sin k(k^2-2\pi^2)}{k(k^2-4\pi^2)}.
\end{align*}
Vậy $C_n=A(k_n)/N(k_n)$, $N(k_n)=1/2+\sin2k_n/(4k_n)$.
\textbf{Bước 3. Tích phân thời gian và kiểm tra.}
\[u=\sum_{n\ge1}\frac{A(k_n)}{N(k_n)}
\frac{e^{-t}-e^{-2k_n^2t}}{2k_n^2-1}\cos k_nx.\]
Thế $T_n$ vào $T_n'+2k_n^2T_n=C_ne^{-t}$ trả lại nguồn; $T_n(0)=0$. $X_n'(0)=0$ và $X_n(1)+X_n'(1)=0$ trả lại hai biên. Nếu một mẫu thời gian bằng $0$, biểu thức được hiểu bằng $te^{-t}$, không bỏ mode đó. Vì các trị riêng đều dương, nghiệm tiến đều về $0$.

\clearpage\noindent\texttt{\detokenize{BT_4_4_3_06}}\par
\section*{Lời giải bổ sung — Bài tập 6 · Mục 4.4.3}

\subsection*{Kiến thức và trình tự giải}
\textbf{Bước 1 — Khử điều kiện biên không thuần nhất.} Viết $u=b+w$ với $b$ thỏa đúng biên. Dirichlet hai đầu: $b=\mu+(\nu-\mu)x/L$. Neumann hai đầu: $b=\mu x+(\nu-\mu)x^2/(2L)$. Neumann–Dirichlet: $b=\nu+\mu(x-L)$; Dirichlet–Neumann: $b=\mu+x\nu$. Robin trái–Dirichlet phải: $b=\mu+d+dx$, $d=(\nu-\mu)/(L+1)$. Neumann trái–Robin phải: $b=\nu-\mu(L+1)+\mu x$. Hai Robin: $b=\mu+d+dx$, $d=(\nu-\mu)/(L+2)$.

\textbf{Bước 2 — Tính lại nguồn và dữ liệu đầu.} Từ $u_t=Ku_{xx}+F$ suy $w_t=Kw_{xx}+Q$, $Q=F+Kb_{xx}-b_t$, $w(x,0)=f(x)-b(x,0)$. Tính $b_t,b_{xx}$ trước khi chiếu; với hàm nâng bậc hai Neumann, $Kb_{xx}$ không bằng không. Nếu dùng trạng thái dừng $P$, giải $KP''+F=0$ và đúng biên rồi đặt $w=u-P$.

\textbf{Bước 3 — Chiếu và giải ODE bằng nhân tử.} Với $N_n=\int X_n^2$, đặt $d_n=N_n^{-1}\int(f-b_0)X_n$, $Q_n(t)=N_n^{-1}\int QX_n$. Phương trình mode là $T_n'+K\lambda_nT_n=Q_n$. Nhân với $e^{K\lambda_nt}$, tích phân từ $0$ đến $t$, rồi chia:


\[
T_n=d_ne^{-K\lambda_nt}+\int_0^te^{-K\lambda_n(t-s)}Q_n(s)ds,\qquad
u=b+\sum_nT_nX_n.
\]


Đây cũng là công thức Green $u=b+\int G(f-b_0)+\int_0^t\int GQ$ cho biên đã được khử.

\textbf{Bước 4 — Thực hiện tích phân thời gian.} Với $Q_n(s)=C_ne^{-\rho s}$, đặt $d=K\lambda_n$:


\[
\int_0^te^{-d(t-s)}C_ne^{-\rho s}ds
=C_ne^{-dt}\int_0^te^{(d-\rho)s}ds
=C_n\frac{e^{-\rho t}-e^{-dt}}{d-\rho}.
\]


Nếu $d=\rho$, kết quả là $C_nte^{-dt}$. Với nguồn hằng, kết quả $C_n(1-e^{-dt})/d$ khi $d>0$; mode $d=0$ phải dùng $\int_0^tQ_0(s)ds$. Nguồn đa thức thời gian tích phân từng phần hoặc giải ODE với nghiệm riêng đa thức.

\textbf{Bước 5 — Kiểm tra và tìm giới hạn.} Mỗi mode đúng PDE, biên; $T_n(0)=d_n$ trả lại dữ liệu đầu sau khi cộng $b_0$. Tách mode hằng Neumann: $\frac d{dt}\int u=K(\nu-\mu)+\int F$. Trạng thái dừng chỉ tồn tại nếu tổng nguồn và thông lượng cân bằng. Với các mode dương, kiểm tra hội tụ đủ để lấy giới hạn đều; không bỏ một mode bằng không hoặc mode tăng.
\subsection*{Lời giải áp dụng cho bài}
\textbf{Bước 1. Giải hệ nâng hai biên Robin.}
Với $b=a+dx$, $a-d=\mu$, $a+dL+d=\nu$. Trừ hai phương trình được $d=(\nu-\mu)/(L+2)$, $a=\mu+d$. Đặt $w=u-b$, $Q=F-b_t$, $q=f-b_0$. Hàm riêng là $X=\cos kx+\sin kx/k$; biên trái đúng sẵn, biên phải cho
\[(1-k^2)\sin kL+2k\cos kL=0.\]
Lấy $k>0$. Ở $k=0$, $X=A+Bx$ cho $A=B$ và $A+B(L+1)=0$, suy $A=B=0$, nên không có trị riêng $0$.
\textbf{Bước 2. Chuẩn và Green.}
\[N(k)=\frac L2(1+k^{-2})+\frac{\sin2kL}{4k}(1-k^{-2})+\frac{\sin^2kL}{k^2},\quad
G=\sum\frac{X_n(x)X_n(y)}{N(k_n)}e^{-Kk_n^2t}.\]
Vậy $u=b+\int Gq+\int_0^t\int GQ$. Công thức trên dữ liệu gốc có hai hạng biên cùng dấu dương:
\[K\int_0^t\{G(x,0,t-s)\mu(s)+G(x,L,t-s)\nu(s)\}ds.\]
Thật vậy $G_y(0)=G(0)$ và $G_y(L)=-G(L)$ trong $K[Gu_y-G_yu]_0^L$.
\textbf{Bước 3. Tính các tích phân của câu cụ thể $K=2,L=1$, dữ liệu đầu và biên $0$.}
Tích–tổng trên $[0,1]$ cho
\[A(k)=\frac{\sin k(k^2-2\pi^2)}{k(k^2-4\pi^2)},\qquad
B(k)=\frac{(1-\cos k)(k^2-2\pi^2)}{k(k^2-4\pi^2)}.\]
Chúng lần lượt là tích phân của $\cos^2\pi y$ nhân với $\cos ky$ và $\sin ky$. Do đó $C_n=(A(k_n)+B(k_n)/k_n)/N(k_n)$, không còn tích phân không gian chưa tính.
\textbf{Bước 4. Mode nguồn.}
\[T_n=C_ne^{-2k_n^2t}\int_0^te^{(2k_n^2-1)s}ds
=C_n\frac{e^{-t}-e^{-2k_n^2t}}{2k_n^2-1},\qquad
u=\sum_{n\ge1}T_nX_n.\]
Thương bằng $te^{-t}$ khi $2k_n^2=1$. ODE và từng hàm riêng kiểm tra đầy đủ PDE, biên, mặt đầu; các mode giảm cho giới hạn $0$.

\clearpage\noindent\texttt{\detokenize{BT_4_4_3_07}}\par
\section*{Lời giải bổ sung — Bài tập 7 · Mục 4.4.3}

\subsection*{Kiến thức và trình tự giải}
\textbf{Bước 1 — Khử điều kiện biên không thuần nhất.} Viết $u=b+w$ với $b$ thỏa đúng biên. Dirichlet hai đầu: $b=\mu+(\nu-\mu)x/L$. Neumann hai đầu: $b=\mu x+(\nu-\mu)x^2/(2L)$. Neumann–Dirichlet: $b=\nu+\mu(x-L)$; Dirichlet–Neumann: $b=\mu+x\nu$. Robin trái–Dirichlet phải: $b=\mu+d+dx$, $d=(\nu-\mu)/(L+1)$. Neumann trái–Robin phải: $b=\nu-\mu(L+1)+\mu x$. Hai Robin: $b=\mu+d+dx$, $d=(\nu-\mu)/(L+2)$.

\textbf{Bước 2 — Tính lại nguồn và dữ liệu đầu.} Từ $u_t=Ku_{xx}+F$ suy $w_t=Kw_{xx}+Q$, $Q=F+Kb_{xx}-b_t$, $w(x,0)=f(x)-b(x,0)$. Tính $b_t,b_{xx}$ trước khi chiếu; với hàm nâng bậc hai Neumann, $Kb_{xx}$ không bằng không. Nếu dùng trạng thái dừng $P$, giải $KP''+F=0$ và đúng biên rồi đặt $w=u-P$.

\textbf{Bước 3 — Chiếu và giải ODE bằng nhân tử.} Với $N_n=\int X_n^2$, đặt $d_n=N_n^{-1}\int(f-b_0)X_n$, $Q_n(t)=N_n^{-1}\int QX_n$. Phương trình mode là $T_n'+K\lambda_nT_n=Q_n$. Nhân với $e^{K\lambda_nt}$, tích phân từ $0$ đến $t$, rồi chia:


\[
T_n=d_ne^{-K\lambda_nt}+\int_0^te^{-K\lambda_n(t-s)}Q_n(s)ds,\qquad
u=b+\sum_nT_nX_n.
\]


Đây cũng là công thức Green $u=b+\int G(f-b_0)+\int_0^t\int GQ$ cho biên đã được khử.

\textbf{Bước 4 — Thực hiện tích phân thời gian.} Với $Q_n(s)=C_ne^{-\rho s}$, đặt $d=K\lambda_n$:


\[
\int_0^te^{-d(t-s)}C_ne^{-\rho s}ds
=C_ne^{-dt}\int_0^te^{(d-\rho)s}ds
=C_n\frac{e^{-\rho t}-e^{-dt}}{d-\rho}.
\]


Nếu $d=\rho$, kết quả là $C_nte^{-dt}$. Với nguồn hằng, kết quả $C_n(1-e^{-dt})/d$ khi $d>0$; mode $d=0$ phải dùng $\int_0^tQ_0(s)ds$. Nguồn đa thức thời gian tích phân từng phần hoặc giải ODE với nghiệm riêng đa thức.

\textbf{Bước 5 — Kiểm tra và tìm giới hạn.} Mỗi mode đúng PDE, biên; $T_n(0)=d_n$ trả lại dữ liệu đầu sau khi cộng $b_0$. Tách mode hằng Neumann: $\frac d{dt}\int u=K(\nu-\mu)+\int F$. Trạng thái dừng chỉ tồn tại nếu tổng nguồn và thông lượng cân bằng. Với các mode dương, kiểm tra hội tụ đủ để lấy giới hạn đều; không bỏ một mode bằng không hoặc mode tăng.
\subsection*{Lời giải áp dụng cho bài}
\textbf{Bước 1. Biên Dirichlet–Neumann và công thức Green.}
Lấy $b=\mu+x\nu$, $w=u-b$, $Q=F-b_t$, $q=f-b_0$. Biên thuần nhất có $X_n=\sin k_nx$, $k_n=(n+1/2)\pi/L$, $N_n=L/2$. Nhân $G=\sum XXe^{-Kk^2t}/N$ cho $u=b+\int Gq+\int_0^t\int GQ$. Công thức trực tiếp có hạng biên
\[K\int_0^t\{G_y(x,0,t-s)\mu(s)+G(x,L,t-s)\nu(s)\}ds,\]
do $G(0)=0,G_y(L)=0$. Dòng tham chiếu (4.4.34) trong đề phải đọc theo biên Dirichlet–Neumann (4.4.33) được nêu ngay trước.

\textbf{Bước 2. Dữ kiện cụ thể.}
$K=L=1,\mu=1-e^{-t},\nu=e^{-t}-1$, $f=F=0$. Khi đó $b=(1-x)(1-e^{-t})$, $b_t=(1-x)e^{-t}$, $q=0$, $Q=-(1-x)e^{-t}$.
Với $k=(n+1/2)\pi$, $\cos k=0,\sin k=(-1)^n$,
\begin{align*}
c_n&=2\int_0^1(1-y)\sin ky\,dy\\
&=2\left[\frac{-(1-y)\cos ky}k-\frac{\sin ky}{k^2}\right]_0^1
=2\left(\frac1k-\frac{(-1)^n}{k^2}\right).
\end{align*}
\textbf{Bước 3. Tích phân và ghép.}
$T_n'+k_n^2T_n=-c_ne^{-t}$, $T_n(0)=0$, nên
\[u=(1-x)(1-e^{-t})
-\sum_{n\ge0}c_n\frac{e^{-t}-e^{-k_n^2t}}{k_n^2-1}\sin k_nx.\]
$k_n$ nửa nguyên nhân $\pi$, không bằng $1$. Thế biên: tại $0$, chuỗi bằng $0$, $u=1-e^{-t}$; đạo hàm tại $1$ của chuỗi bằng $0$, $u_x=e^{-t}-1$. Tại $t=0$ mọi hạng bằng $0$. Nguồn phụ $-b_t$ bị triệt tiêu khi cộng lại $b_t$, cho PDE gốc không nguồn.
\textbf{Bước 4. Giới hạn.} $b\to1-x$ đều. Các hệ số quá độ $O(k_n^{-3})$ và hai mũ giảm cho chuỗi tiến đều về $0$, nên $\lim_{t\to\infty}u=1-x$.

\clearpage\noindent\texttt{\detokenize{BT_4_4_3_08}}\par
\section*{Lời giải bổ sung — Bài tập 8 · Mục 4.4.3}

\subsection*{Kiến thức và trình tự giải}
\textbf{Bước 1 — Khử điều kiện biên không thuần nhất.} Viết $u=b+w$ với $b$ thỏa đúng biên. Dirichlet hai đầu: $b=\mu+(\nu-\mu)x/L$. Neumann hai đầu: $b=\mu x+(\nu-\mu)x^2/(2L)$. Neumann–Dirichlet: $b=\nu+\mu(x-L)$; Dirichlet–Neumann: $b=\mu+x\nu$. Robin trái–Dirichlet phải: $b=\mu+d+dx$, $d=(\nu-\mu)/(L+1)$. Neumann trái–Robin phải: $b=\nu-\mu(L+1)+\mu x$. Hai Robin: $b=\mu+d+dx$, $d=(\nu-\mu)/(L+2)$.

\textbf{Bước 2 — Tính lại nguồn và dữ liệu đầu.} Từ $u_t=Ku_{xx}+F$ suy $w_t=Kw_{xx}+Q$, $Q=F+Kb_{xx}-b_t$, $w(x,0)=f(x)-b(x,0)$. Tính $b_t,b_{xx}$ trước khi chiếu; với hàm nâng bậc hai Neumann, $Kb_{xx}$ không bằng không. Nếu dùng trạng thái dừng $P$, giải $KP''+F=0$ và đúng biên rồi đặt $w=u-P$.

\textbf{Bước 3 — Chiếu và giải ODE bằng nhân tử.} Với $N_n=\int X_n^2$, đặt $d_n=N_n^{-1}\int(f-b_0)X_n$, $Q_n(t)=N_n^{-1}\int QX_n$. Phương trình mode là $T_n'+K\lambda_nT_n=Q_n$. Nhân với $e^{K\lambda_nt}$, tích phân từ $0$ đến $t$, rồi chia:


\[
T_n=d_ne^{-K\lambda_nt}+\int_0^te^{-K\lambda_n(t-s)}Q_n(s)ds,\qquad
u=b+\sum_nT_nX_n.
\]


Đây cũng là công thức Green $u=b+\int G(f-b_0)+\int_0^t\int GQ$ cho biên đã được khử.

\textbf{Bước 4 — Thực hiện tích phân thời gian.} Với $Q_n(s)=C_ne^{-\rho s}$, đặt $d=K\lambda_n$:


\[
\int_0^te^{-d(t-s)}C_ne^{-\rho s}ds
=C_ne^{-dt}\int_0^te^{(d-\rho)s}ds
=C_n\frac{e^{-\rho t}-e^{-dt}}{d-\rho}.
\]


Nếu $d=\rho$, kết quả là $C_nte^{-dt}$. Với nguồn hằng, kết quả $C_n(1-e^{-dt})/d$ khi $d>0$; mode $d=0$ phải dùng $\int_0^tQ_0(s)ds$. Nguồn đa thức thời gian tích phân từng phần hoặc giải ODE với nghiệm riêng đa thức.

\textbf{Bước 5 — Kiểm tra và tìm giới hạn.} Mỗi mode đúng PDE, biên; $T_n(0)=d_n$ trả lại dữ liệu đầu sau khi cộng $b_0$. Tách mode hằng Neumann: $\frac d{dt}\int u=K(\nu-\mu)+\int F$. Trạng thái dừng chỉ tồn tại nếu tổng nguồn và thông lượng cân bằng. Với các mode dương, kiểm tra hội tụ đủ để lấy giới hạn đều; không bỏ một mode bằng không hoặc mode tăng.
\subsection*{Lời giải áp dụng cho bài}
\textbf{Bước 1. Cơ sở cosine và mode trung bình.}
Đặt $f_0=\int_0^1f$, $g_0=\int_0^1g$, $f_n=2\int_0^1f(y)\cos n\pi y\,dy$, $g_n=2\int_0^1g(y)\cos n\pi y\,dy$, $n\ge1$. Hai biên Neumann $0$ cho
\[u=T_0+\sum_{n\ge1}T_n\cos n\pi x.\]
Chiếu PDE: $T_0'=g_0e^{-2t},T_0(0)=f_0$; $T_n'+n^2\pi^2T_n=g_ne^{-2t},T_n(0)=f_n$.
\textbf{Bước 2. Giải từng ODE, kể cả mode $0$.}
\[T_0=f_0+\frac{g_0}2(1-e^{-2t}),\quad
T_n=f_ne^{-n^2\pi^2t}+g_ne^{-n^2\pi^2t}
\int_0^te^{(n^2\pi^2-2)s}ds.\]
Vì $n^2\pi^2>2$, tích phân cuối bằng $(e^{(n^2\pi^2-2)t}-1)/(n^2\pi^2-2)$:
\[u=f_0+\frac{g_0}2(1-e^{-2t})+
\sum_{n\ge1}\left[f_ne^{-n^2\pi^2t}
+g_n\frac{e^{-2t}-e^{-n^2\pi^2t}}{n^2\pi^2-2}\right]\cos n\pi x.\]
\textbf{Bước 3. Kiểm tra và giới hạn đều.}
Đạo hàm cosine bằng $0$ ở hai đầu; ODE kiểm tra PDE; $T_n(0)=f_n$ trả lại mặt đầu. Nếu $f,g\in L^1$, $|f_n|\le2\|f\|_1$, $|g_n|\le2\|g\|_1$. Phần đầu của chuỗi bị chặn bởi $2\|f\|_1\sum e^{-n^2\pi^2t}\to0$. Phần nguồn bị chặn bởi $2\|g\|_1(e^{-2t}+e^{-\pi^2t})\sum(n^2\pi^2-2)^{-1}\to0$. Suy
\[\lim_{t\to\infty}u(x,t)=\int_0^1f(y)dy+\frac12\int_0^1g(y)dy\]
đều theo $x$. Tích phân PDE cho tổng nhiệt $E'=g_0e^{-2t}$ cũng kiểm tra độc lập mode trung bình và giới hạn này.

\clearpage\noindent\texttt{\detokenize{VD_4_5_2_01}}\par
\section*{Lời giải bổ sung — Ví dụ 1 · Mục 4.5.2}

\subsection*{Kiến thức và trình tự giải}
\textbf{Bước 1 — Loại đạo hàm cấp một bằng tính đạo hàm đầy đủ.} Với $u_t=Ku_{xx}+Au_x+Bu$, đặt $u=e^{-hx}v$, $h=A/(2K)$. Khi đó


\[
u_t=e^{-hx}v_t,\quad u_x=e^{-hx}(v_x-hv),\quad
u_{xx}=e^{-hx}(v_{xx}-2hv_x+h^2v).
\]


Thay vào PDE, hệ số $v_x$ là $A-2Kh=0$, hệ số $v$ là $Kh^2-Ah+B=B-A^2/(4K)$. Vì thế $v_t=Kv_{xx}+Cv$, $C=B-A^2/(4K)$; có thể tiếp tục đặt $v=e^{Ct}z$ để $z_t=Kz_{xx}$.

\textbf{Bước 2 — Đổi dữ liệu và biên cùng phép đổi ẩn.} $v(x,0)=e^{hx}f(x)$. Dirichlet thuần nhất giữ nguyên; Neumann biến thành $v_x-hv=0$, nên không giữ nguyên một cơ sở cosine nếu $A\ne0$. Với Dirichlet:


\[
u=e^{-hx}\sum_{n\ge1}\left(\frac2L\int_0^Le^{hy}f(y)\sin\frac{n\pi y}Ldy\right)
e^{[C-K(n\pi/L)^2]t}\sin\frac{n\pi x}L.
\]


\textbf{Bước 3 — Tính hệ số mũ–đa thức.} Dùng nguyên hàm


\[
\int e^{hy}\sin ky\,dy
=\frac{e^{hy}(h\sin ky-k\cos ky)}{h^2+k^2}.
\]


Với $y e^{hy}$ hoặc $y(1-y)e^{hy}$, lấy đạo hàm theo $h$ của tích phân đã tính trên cận cố định; trình bày đạo hàm để tìm hệ số thực. Cách này tránh để lại một tích phân sơ cấp chưa thực hiện.

\textbf{Bước 4 — Hệ số không gian biến thiên.} Với $x^2u_{xx}$ trên $x>0$, đặt $z=\ln x$: nếu $U(z,t)=u(e^z,t)$ thì $u_x=U_z/x$, $x^2u_{xx}=U_{zz}-U_z$. Đổi miền và biên sang biến $z$ rồi dùng bước 1. Biên không thuần nhất nên tìm trạng thái dừng trước khi khai triển.

\textbf{Bước 5 — Duy nhất và ổn định.} Với hiệu nghiệm Dirichlet, $I'=-K\int w_x^2+2BI\le2BI$ suy $w=0$ khi $I(0)=0$. Giới hạn phụ thuộc các số $C-K\lambda_n$ có hệ số khác không: âm thì giảm, bằng không thì tồn tại mode dừng, dương thì tăng. Phải kiểm tra tính thống nhất giữa miền PDE và vị trí biên trước khi dùng công thức.
\subsection*{Lời giải áp dụng cho bài}
\textbf{Bước 1. Tính phép đổi ẩn.}
Đặt $u=e^{-hx}v$, $h=A/(2K)$. Khi đó $u_x=e^{-hx}(v_x-hv)$, $u_{xx}=e^{-hx}(v_{xx}-2hv_x+h^2v)$, $u_t=e^{-hx}v_t$. Thay vào $u_t=Ku_{xx}+Au_x+Bu$ cho
\[v_t=Kv_{xx}+(B-A^2/(4K))v,\quad v(x,0)=e^{hx}f(x).\]
Hai biên Dirichlet vẫn bằng $0$, nên $u=e^{-hx}\sum c_ne^{[B-A^2/(4K)-Kn^2]t}\sin nx$ khi $L=\pi$.

\textbf{Bước 2. Câu (i), $K=1,A=0,B=-1,f=x$.}
\[c_n=\frac2\pi\int_0^\pi x\sin nx\,dx
=\frac2\pi\left[-\frac{x\cos nx}n+\frac{\sin nx}{n^2}\right]_0^\pi
=\frac{2(-1)^{n+1}}n.\]
Vì vậy $u=\sum_{n\ge1}2(-1)^{n+1}e^{-(n^2+1)t}\sin nx/n$.

\textbf{Bước 3. Câu (ii), $K=1,A=-1,B=0,f=1$.}
$h=-1/2$, nên dữ liệu của $v$ là $e^{-x/2}$. Nguyên hàm $\int e^{hx}\sin nx=e^{hx}(h\sin nx-n\cos nx)/(h^2+n^2)$ cho
\[c_n=\frac2\pi\left[\frac{e^{-x/2}(-\sin nx/2-n\cos nx)}{n^2+1/4}\right]_0^\pi
=\frac{2n(1-(-1)^ne^{-\pi/2})}{\pi(n^2+1/4)}.\]
Suy $u=e^{x/2}\sum c_ne^{-(n^2+1/4)t}\sin nx$.

\textbf{Bước 4. Duy nhất và kiểm tra.}
Với hiệu $w$ có dữ liệu đầu $0$, biên $0$,
\begin{align*}
I'&=K[ww_x]_0^\pi-K\int w_x^2+\frac A2[w^2]_0^\pi+B\int w^2\\
&=-K\int w_x^2+2BI\le0
\end{align*}
trong cả hai câu vì $B\le0$. $I(0)=0$ suy $w=0$. Phép đổi ẩn và các ODE mode kiểm tra PDE; sine kiểm tra biên; Fourier trả lại dữ liệu trong miền. $f(\pi)=\pi$ ở (i), và $f(0)=f(\pi)=1$ ở (ii) không khớp biên $0$, nên không liên tục tại các góc này. Hai nghiệm trơn khi $t>0$ và đều tiến về $0$.

\clearpage\noindent\texttt{\detokenize{VD_4_5_2_02}}\par
\section*{Lời giải bổ sung — Ví dụ 2 · Mục 4.5.2}

\subsection*{Kiến thức và trình tự giải}
\textbf{Bước 1 — Phân biệt đạo hàm cấp hai, ba và bốn.} Quy ước $\widehat f(\xi)=\int e^{-i\xi x}f(x)dx$, biến đổi ngược có hệ số $1/(2\pi)$. Khi đó $\widehat{f^{(m)}}=(i\xi)^m\widehat f$, $\widehat{xf}=i\partial_\xi\widehat f$, $\widehat{xf_x}=-\widehat f-\xi\partial_\xi\widehat f$.

\textbf{Bước 2 — Giải phương trình của biến đổi.} $u_t=u_{xx}$ cho $\widehat u=e^{-t\xi^2}\widehat f$. $u_t=u_{xxx}$ cho $\widehat u=e^{-it\xi^3}\widehat f$. $u_t=-u_{xxxx}$ cho $\widehat u=e^{-t\xi^4}\widehat f$. Với $xu_x$, giải PDE cấp một theo $\xi$ bằng đường đặc trưng, gồm cả hệ số biên độ $-\widehat u$.

\textbf{Bước 3 — Biến đổi ngược và ý nghĩa tích phân.} Nhân bậc ba là $A_t(x)=(3t)^{-1/3}\operatorname{Ai}(-x/(3t)^{1/3})$, tích phân dao động. Nhân bậc bốn là $B_t(x)=(2\pi)^{-1}\int e^{i\xi x-t\xi^4}d\xi$, hội tụ tuyệt đối theo $\xi$ và có thể đổi dấu. Không áp dụng tính dương của Gauss hay nguyên lý cực đại cấp hai cho hai nhân này.

\textbf{Bước 4 — Miền hữu hạn cần kiểm tra biên và toán tử liên hợp.} Với $u_t=Ku_{xxx}$, $u(0)=u_x(0)=u(L)=0$, năng lượng cho $I'=-K u_x(L)^2/2$. Toán tử liên hợp là $-KD^3$ với $Y(0)=Y(L)=Y'(L)=0$. Đây không phải bài Sturm–Liouville tự liên hợp cấp hai; không suy mọi trị riêng đều thực và không tự dùng hệ số $\int fX/\int X^2$. Có thể giải bằng biến đổi Laplace theo thời gian và nghiệm cơ bản của ODE bậc ba, tránh khẳng định cơ sở chưa chứng minh.

\textbf{Bước 5 — Nêu lớp dữ liệu đủ cho công thức.} Với bậc ba trên toàn trục, công thức Fourier hợp lệ trước hết cho dữ liệu Schwartz hoặc $L^2$ theo nghĩa biến đổi Fourier. Dữ liệu liên tục bị chặn là phân phối tempered; phép nhân $e^{-it\xi^3}\widehat f$ vẫn hợp lệ vì các đạo hàm của thừa số tăng đa thức, nên xác định nghiệm trong $\mathcal S\prime$. Tuy nhiên điều đó không bảo đảm tích phân Airy hội tụ tuyệt đối hay nghiệm tự động là hàm cổ điển. Ghi đúng ý nghĩa của công thức theo lớp dữ liệu.
\subsection*{Lời giải áp dụng cho bài}
\textbf{Bước 1. Năng lượng bậc ba và tính duy nhất.}
Với $K>0$, $u_t=Ku_{xxx}$, $u(0)=u_x(0)=u(L)=0$, hai lần tích phân từng phần cho
\[I'=K\int uu_{xxx}
=K[uu_{xx}]_0^L-\frac K2[u_x^2]_0^L
=-\frac K2u_x(L,t)^2\le0.\]
Bản in bỏ hệ số $1/2$ ở đẳng thức cuối; dấu suy giảm vẫn đúng. Áp dụng cho hiệu hai nghiệm có mặt đầu $0$, suy $I=0$, tức duy nhất trong lớp nghiệm đủ trơn/năng lượng thích hợp.

\textbf{Bước 2. Vì sao không dùng cơ sở trực giao cấp hai.}
Tách biến $u=XT$ cho $X'''=\lambda X$, $T'=K\lambda T$ và ba biên $X(0)=X'(0)=X(L)=0$. Với hàm riêng phức, lấy phần thực của tích phân với $\overline X$:
\[\operatorname{Re}\lambda\int_0^L|X|^2dx=-\frac12|X'(L)|^2\le0.\]
Điều này không chứng minh $\lambda$ thực. Toán tử $D^3$ không tự liên hợp; liên hợp là $-D^3$ với biên $Y(0)=Y(L)=Y'(L)=0$, suy từ hạng $[X''\overline Y-X'\overline{Y'}+X\overline{Y''}]_0^L$. Vì vậy chưa được khẳng định dãy hàm riêng thực trong sách là cơ sở đầy đủ.
Nếu chỉ xét nhánh thực $\lambda=-k^3$, nghiệm thỏa hai biên trái là
\[X=e^{-kx}-e^{kx/2}\cos(\sqrt3kx/2)+\sqrt3e^{kx/2}\sin(\sqrt3kx/2)
=e^{-kx}+2e^{kx/2}\sin(\sqrt3kx/2-\pi/6).\]
Biên phải cho $e^{-3kL/2}=-2\sin(\sqrt3kL/2-\pi/6)$. Dấu trước sine trong $X$ phải là $+$; dấu $-$ ở công thức hàm riêng của bản in làm $X(0)=2$, sai biên.

\textbf{Bước 3. Giải đầy đủ bằng Laplace, không giả định cơ sở chưa chứng minh.}
Đặt $U(x,p)=\int_0^\infty e^{-pt}u(x,t)dt$, $\operatorname{Re}p>0$. Từ $pU-f=KU'''$:
\[U'''-\frac pK U=-\frac fK,\quad U(0,p)=U_x(0,p)=U(L,p)=0.\]
Đặt nghiệm cơ bản
\[S_p(x)=\sum_{m=0}^\infty\frac{(p/K)^m x^{3m+2}}{(3m+2)!}.\]
Đạo hàm chuỗi cho $S_p'''=(p/K)S_p$, $S_p(0)=S_p'(0)=0,S_p''(0)=1$. Biến thiên hằng số cho
\[U(x,p)=C(p)S_p(x)-\frac1K\int_0^xS_p(x-y)f(y)dy.\]
Đạo hàm tích phân ba lần tạo đúng $f(x)$ nhờ $S_p''(0)=1$. Biên $U(L)=0$ xác định
\[C(p)=\frac{\int_0^LS_p(L-y)f(y)dy}{KS_p(L)}.\]
Mẫu không bằng $0$ khi $\operatorname{Re}p>0$: nếu $S_p(L)=0$, $S_p$ là hàm riêng không tầm thường với trị riêng $p/K$ có phần thực dương, trái đẳng thức năng lượng.

\textbf{Bước 4. Công thức nghiệm và kiểm tra.}
\[
u(x,t)=\mathcal L^{-1}_{p\to t}
\left\{\frac{S_p(x)}{KS_p(L)}\int_0^LS_p(L-y)f(y)dy
-\frac1K\int_0^xS_p(x-y)f(y)dy\right\}.
\]
Ở đây $\mathcal L^{-1}H=(2\pi i)^{-1}\int_{\gamma-i\infty}^{\gamma+i\infty}e^{pt}H(p)dp$, $\gamma>0$, hiểu theo phép đảo Laplace trong lớp năng lượng; với dữ liệu trơn tương thích có nghiệm cổ điển. Công thức $U$ thỏa ODE, cả ba biên và hạng mặt đầu $f$ trong $pU-f=KU'''$, nên phép đảo trả lại đúng bài toán. Với $f\in L^2$, nghiệm được hiểu là nửa nhóm co sinh bởi $KD^3$ trên miền $H^3$ có ba biên trên; tiêu tán vừa chứng minh và nghiệm resolvent cho mọi $p>0$ bảo đảm tồn tại và duy nhất. Đây là công thức giải cho dữ liệu tổng quát mà không bỏ các thành phần phổ phức hoặc giả định các hệ số trực giao.

\clearpage\noindent\texttt{\detokenize{BT_4_5_3_01}}\par
\section*{Lời giải bổ sung — Bài tập 1 · Mục 4.5.3}

\subsection*{Kiến thức và trình tự giải}
\textbf{Bước 1 — Loại đạo hàm cấp một bằng tính đạo hàm đầy đủ.} Với $u_t=Ku_{xx}+Au_x+Bu$, đặt $u=e^{-hx}v$, $h=A/(2K)$. Khi đó


\[
u_t=e^{-hx}v_t,\quad u_x=e^{-hx}(v_x-hv),\quad
u_{xx}=e^{-hx}(v_{xx}-2hv_x+h^2v).
\]


Thay vào PDE, hệ số $v_x$ là $A-2Kh=0$, hệ số $v$ là $Kh^2-Ah+B=B-A^2/(4K)$. Vì thế $v_t=Kv_{xx}+Cv$, $C=B-A^2/(4K)$; có thể tiếp tục đặt $v=e^{Ct}z$ để $z_t=Kz_{xx}$.

\textbf{Bước 2 — Đổi dữ liệu và biên cùng phép đổi ẩn.} $v(x,0)=e^{hx}f(x)$. Dirichlet thuần nhất giữ nguyên; Neumann biến thành $v_x-hv=0$, nên không giữ nguyên một cơ sở cosine nếu $A\ne0$. Với Dirichlet:


\[
u=e^{-hx}\sum_{n\ge1}\left(\frac2L\int_0^Le^{hy}f(y)\sin\frac{n\pi y}Ldy\right)
e^{[C-K(n\pi/L)^2]t}\sin\frac{n\pi x}L.
\]


\textbf{Bước 3 — Tính hệ số mũ–đa thức.} Dùng nguyên hàm


\[
\int e^{hy}\sin ky\,dy
=\frac{e^{hy}(h\sin ky-k\cos ky)}{h^2+k^2}.
\]


Với $y e^{hy}$ hoặc $y(1-y)e^{hy}$, lấy đạo hàm theo $h$ của tích phân đã tính trên cận cố định; trình bày đạo hàm để tìm hệ số thực. Cách này tránh để lại một tích phân sơ cấp chưa thực hiện.

\textbf{Bước 4 — Hệ số không gian biến thiên.} Với $x^2u_{xx}$ trên $x>0$, đặt $z=\ln x$: nếu $U(z,t)=u(e^z,t)$ thì $u_x=U_z/x$, $x^2u_{xx}=U_{zz}-U_z$. Đổi miền và biên sang biến $z$ rồi dùng bước 1. Biên không thuần nhất nên tìm trạng thái dừng trước khi khai triển.

\textbf{Bước 5 — Duy nhất và ổn định.} Với hiệu nghiệm Dirichlet, $I'=-K\int w_x^2+2BI\le2BI$ suy $w=0$ khi $I(0)=0$. Giới hạn phụ thuộc các số $C-K\lambda_n$ có hệ số khác không: âm thì giảm, bằng không thì tồn tại mode dừng, dương thì tăng. Phải kiểm tra tính thống nhất giữa miền PDE và vị trí biên trước khi dùng công thức.
\subsection*{Lời giải áp dụng cho bài}
\textbf{Bước 1. Phép đổi ẩn cho hai biên Dirichlet.}
$u=e^{-Ax/(2K)}v$ cho $v_t=Kv_{xx}+(B-A^2/(4K))v$, dữ liệu $e^{Ax/(2K)}f$, biên $0$. Cơ sở là $\sin nx$ vì $L=\pi$.
\textbf{Bước 2. Câu (i), $A=-1,B=0,K=1,f=1$.}
Đặt $h=-1/2$, dùng nguyên hàm $e^{hx}(h\sin nx-n\cos nx)/(h^2+n^2)$:
\[c_n=\frac2\pi\int_0^\pi e^{-y/2}\sin ny\,dy
=\frac{2n(1-(-1)^ne^{-\pi/2})}{\pi(n^2+1/4)}.\]
Suy
\[u=e^{x/2}\sum_{n\ge1}c_ne^{-(n^2+1/4)t}\sin nx.\]
\textbf{Bước 3. Câu (ii), $A=0,B=-1,K=1,f=x$.}
Tích phân từng phần:
\[c_n=\frac2\pi[-y\cos ny/n+\sin ny/n^2]_0^\pi
=\frac{2(-1)^{n+1}}n,\qquad
u=\sum_{n\ge1}\frac{2(-1)^{n+1}}n e^{-(n^2+1)t}\sin nx.\]
\textbf{Bước 4. Chứng minh duy nhất bằng năng lượng.}
Hiệu $w$ có dữ liệu $0$ và biên $0$. Tính
\[I'=K[ww_x]_0^\pi+\frac A2[w^2]_0^\pi-K\int w_x^2+B\int w^2
=-K\int w_x^2+2BI\le0.\]
Vì $I(0)=0$, $w=0$. Phép đổi ẩn kiểm tra PDE, các sine kiểm tra biên, Fourier kiểm tra dữ liệu đầu trong miền. Dữ liệu không khớp Dirichlet tại các góc như đã thấy từ $f=1$ hoặc $f(\pi)=\pi$; nghiệm trơn khi $t>0$ và tiến đều về $0$.

\clearpage\noindent\texttt{\detokenize{BT_4_5_3_02}}\par
\section*{Lời giải bổ sung — Bài tập 2 · Mục 4.5.3}

\subsection*{Kiến thức và trình tự giải}
\textbf{Bước 1 — Loại đạo hàm cấp một bằng tính đạo hàm đầy đủ.} Với $u_t=Ku_{xx}+Au_x+Bu$, đặt $u=e^{-hx}v$, $h=A/(2K)$. Khi đó


\[
u_t=e^{-hx}v_t,\quad u_x=e^{-hx}(v_x-hv),\quad
u_{xx}=e^{-hx}(v_{xx}-2hv_x+h^2v).
\]


Thay vào PDE, hệ số $v_x$ là $A-2Kh=0$, hệ số $v$ là $Kh^2-Ah+B=B-A^2/(4K)$. Vì thế $v_t=Kv_{xx}+Cv$, $C=B-A^2/(4K)$; có thể tiếp tục đặt $v=e^{Ct}z$ để $z_t=Kz_{xx}$.

\textbf{Bước 2 — Đổi dữ liệu và biên cùng phép đổi ẩn.} $v(x,0)=e^{hx}f(x)$. Dirichlet thuần nhất giữ nguyên; Neumann biến thành $v_x-hv=0$, nên không giữ nguyên một cơ sở cosine nếu $A\ne0$. Với Dirichlet:


\[
u=e^{-hx}\sum_{n\ge1}\left(\frac2L\int_0^Le^{hy}f(y)\sin\frac{n\pi y}Ldy\right)
e^{[C-K(n\pi/L)^2]t}\sin\frac{n\pi x}L.
\]


\textbf{Bước 3 — Tính hệ số mũ–đa thức.} Dùng nguyên hàm


\[
\int e^{hy}\sin ky\,dy
=\frac{e^{hy}(h\sin ky-k\cos ky)}{h^2+k^2}.
\]


Với $y e^{hy}$ hoặc $y(1-y)e^{hy}$, lấy đạo hàm theo $h$ của tích phân đã tính trên cận cố định; trình bày đạo hàm để tìm hệ số thực. Cách này tránh để lại một tích phân sơ cấp chưa thực hiện.

\textbf{Bước 4 — Hệ số không gian biến thiên.} Với $x^2u_{xx}$ trên $x>0$, đặt $z=\ln x$: nếu $U(z,t)=u(e^z,t)$ thì $u_x=U_z/x$, $x^2u_{xx}=U_{zz}-U_z$. Đổi miền và biên sang biến $z$ rồi dùng bước 1. Biên không thuần nhất nên tìm trạng thái dừng trước khi khai triển.

\textbf{Bước 5 — Duy nhất và ổn định.} Với hiệu nghiệm Dirichlet, $I'=-K\int w_x^2+2BI\le2BI$ suy $w=0$ khi $I(0)=0$. Giới hạn phụ thuộc các số $C-K\lambda_n$ có hệ số khác không: âm thì giảm, bằng không thì tồn tại mode dừng, dương thì tăng. Phải kiểm tra tính thống nhất giữa miền PDE và vị trí biên trước khi dùng công thức.
\subsection*{Lời giải áp dụng cho bài}
\textbf{Bước 1. Câu (i), phản ứng $\alpha u$.}
$K=L=1,A=0,B=\alpha$. Với $f=x(1-x)$, tích phân từng phần hai lần, $f''=-2$, cho
\[c_n=2\int_0^1f\sin n\pi x\,dx
=\frac{4(1-(-1)^n)}{n^3\pi^3}.\]
$T_n'+(n^2\pi^2-\alpha)T_n=0$, nên
\[u=\sum_{\substack{n\ge1\\n\ {\rm lẻ}}}\frac8{n^3\pi^3}
e^{(\alpha-n^2\pi^2)t}\sin n\pi x.\]
Mode đầu có hệ số $8/\pi^3>0$. Nếu $\alpha<\pi^2$, mọi mode giảm và giới hạn đều bằng $0$; nếu $\alpha=\pi^2$, chỉ mode đầu đứng yên và giới hạn là $8\sin\pi x/\pi^3$; nếu $\alpha>\pi^2$, mode đầu tăng, nghiệm không bị chặn theo thời gian.

\textbf{Bước 2. Câu (ii), vận chuyển $\alpha u_x$.}
Đặt $h=\alpha/2$, $u=e^{-hx}v$. Đạo hàm cho $v_t=v_{xx}-h^2v$ và $v(x,0)=e^{hx}x(1-x)$. Với $k=n\pi$, $s_n=(-1)^n$, đặt $D=h^2+k^2$, $E=s_ne^h$.
\textbf{Bước 3. Tính hệ số không để lại tích phân.}
\[J(h)=\int_0^1e^{hy}\sin ky\,dy=\frac{k(1-E)}D.\]
Do $\partial_h(e^{hy})=ye^{hy}$, $\int e^{hy}y(1-y)\sin ky=J'-J''$. Tính
\[J'=-\frac{kE}D-\frac{2kh(1-E)}{D^2},\quad
J''=-\frac{kE}D+\frac{4khE-2k(1-E)}{D^2}
+\frac{8kh^2(1-E)}{D^3}.\]
Vì vậy hệ số $c_n=2(J'-J'')$ là
\[c_n=\frac{4k\{(1-h)-(1+h)E\}}{D^2}
-\frac{16kh^2(1-E)}{D^3}.\]
Nghiệm:
\[u=e^{-hx}\sum_{n\ge1}c_ne^{-(n^2\pi^2+h^2)t}\sin n\pi x.\]
Mọi tốc độ giảm đều dương với $\alpha$ thực, nên giới hạn đều bằng $0$. Khi $\alpha=0$, công thức hệ số trở về $4(1-(-1)^n)/(n^3\pi^3)$, kiểm tra khớp câu (i) với $\alpha=0$. Các mode và phép đổi ẩn kiểm tra PDE, biên, dữ liệu của từng câu.

\clearpage\noindent\texttt{\detokenize{BT_4_5_3_03}}\par
\section*{Lời giải bổ sung — Bài tập 3 · Mục 4.5.3}

\subsection*{Kiến thức và trình tự giải}
\textbf{Bước 1 — Loại đạo hàm cấp một bằng tính đạo hàm đầy đủ.} Với $u_t=Ku_{xx}+Au_x+Bu$, đặt $u=e^{-hx}v$, $h=A/(2K)$. Khi đó


\[
u_t=e^{-hx}v_t,\quad u_x=e^{-hx}(v_x-hv),\quad
u_{xx}=e^{-hx}(v_{xx}-2hv_x+h^2v).
\]


Thay vào PDE, hệ số $v_x$ là $A-2Kh=0$, hệ số $v$ là $Kh^2-Ah+B=B-A^2/(4K)$. Vì thế $v_t=Kv_{xx}+Cv$, $C=B-A^2/(4K)$; có thể tiếp tục đặt $v=e^{Ct}z$ để $z_t=Kz_{xx}$.

\textbf{Bước 2 — Đổi dữ liệu và biên cùng phép đổi ẩn.} $v(x,0)=e^{hx}f(x)$. Dirichlet thuần nhất giữ nguyên; Neumann biến thành $v_x-hv=0$, nên không giữ nguyên một cơ sở cosine nếu $A\ne0$. Với Dirichlet:


\[
u=e^{-hx}\sum_{n\ge1}\left(\frac2L\int_0^Le^{hy}f(y)\sin\frac{n\pi y}Ldy\right)
e^{[C-K(n\pi/L)^2]t}\sin\frac{n\pi x}L.
\]


\textbf{Bước 3 — Tính hệ số mũ–đa thức.} Dùng nguyên hàm


\[
\int e^{hy}\sin ky\,dy
=\frac{e^{hy}(h\sin ky-k\cos ky)}{h^2+k^2}.
\]


Với $y e^{hy}$ hoặc $y(1-y)e^{hy}$, lấy đạo hàm theo $h$ của tích phân đã tính trên cận cố định; trình bày đạo hàm để tìm hệ số thực. Cách này tránh để lại một tích phân sơ cấp chưa thực hiện.

\textbf{Bước 4 — Hệ số không gian biến thiên.} Với $x^2u_{xx}$ trên $x>0$, đặt $z=\ln x$: nếu $U(z,t)=u(e^z,t)$ thì $u_x=U_z/x$, $x^2u_{xx}=U_{zz}-U_z$. Đổi miền và biên sang biến $z$ rồi dùng bước 1. Biên không thuần nhất nên tìm trạng thái dừng trước khi khai triển.

\textbf{Bước 5 — Duy nhất và ổn định.} Với hiệu nghiệm Dirichlet, $I'=-K\int w_x^2+2BI\le2BI$ suy $w=0$ khi $I(0)=0$. Giới hạn phụ thuộc các số $C-K\lambda_n$ có hệ số khác không: âm thì giảm, bằng không thì tồn tại mode dừng, dương thì tăng. Phải kiểm tra tính thống nhất giữa miền PDE và vị trí biên trước khi dùng công thức.
\subsection*{Lời giải áp dụng cho bài}
\textbf{Bước 1. Chỉ ra chỗ thiếu của đề gốc.}
PDE và dữ liệu đầu ghi $0<x<2$, nhưng hai biên đặt tại $0$ và $1$, không có biên tại $2$. Không thể dùng sine $n\pi x$ như một cơ sở Dirichlet trên toàn $(0,2)$ rồi tự coi $1$ là đầu phải. Đây không phải bài Dirichlet trên đoạn được ghi. Trong lớp nghiệm cổ điển với $t>0$ và liên tục đến dữ liệu đầu trong miền, còn có mâu thuẫn: nghiệm trên $(0,1)$ được xác định duy nhất bởi hai biên $0$. Nghiệm nhiệt hệ số hằng giải tích theo $x$ ở thời gian dương, nên thác triển qua điểm trong $1$ buộc trùng phản xạ lẻ của nghiệm sine. Khi $1<x<2$, mặt đầu thác triển là $-f(2-x)=x^2-3x+2$, khác $f=x-x^2$ đã in, với hiệu $-2(x-1)^2$. Vì vậy cần sửa miền hoặc biên/dữ liệu để có nghiệm trong lớp này.
\textbf{Bước 2. Nếu miền dự kiến là $(0,1)$.}
Đây là giả thiết sửa đề, cần ghi riêng. Khi ấy $X_n=\sin n\pi x$, $\lambda_n=n^2\pi^2$. Với $f=x-x^2$,
\[c_n=2\int_0^1f\sin kx\,dx
=-\frac2{k^2}\int_0^1(-2)\sin kx\,dx
=\frac{4(1-(-1)^n)}{k^3},\quad k=n\pi.\]
Mode thỏa $T_n'=(\alpha-k^2)T_n$:
\[u=\sum_{\substack{n\ge1\\n\ {\rm lẻ}}}\frac8{n^3\pi^3}
e^{(\alpha-n^2\pi^2)t}\sin n\pi x,\quad 0<x<1.\]
Sine bằng $0$ ở $0,1$, chuỗi đầu bằng $x-x^2$, và $u_t-u_{xx}=\alpha u$ theo từng mode. Nếu $\alpha<\pi^2$, giới hạn $0$; nếu bằng $\pi^2$, giới hạn $8\sin\pi x/\pi^3$; nếu lớn hơn, mode đầu tăng. Những kết quả này giải bài đã sửa miền, không tự bổ sung điều kiện cho phần $(1,2)$ của đề gốc.

\clearpage\noindent\texttt{\detokenize{BT_4_5_3_04}}\par
\section*{Lời giải bổ sung — Bài tập 4 · Mục 4.5.3}

\subsection*{Kiến thức và trình tự giải}
\textbf{Bước 1 — Loại đạo hàm cấp một bằng tính đạo hàm đầy đủ.} Với $u_t=Ku_{xx}+Au_x+Bu$, đặt $u=e^{-hx}v$, $h=A/(2K)$. Khi đó


\[
u_t=e^{-hx}v_t,\quad u_x=e^{-hx}(v_x-hv),\quad
u_{xx}=e^{-hx}(v_{xx}-2hv_x+h^2v).
\]


Thay vào PDE, hệ số $v_x$ là $A-2Kh=0$, hệ số $v$ là $Kh^2-Ah+B=B-A^2/(4K)$. Vì thế $v_t=Kv_{xx}+Cv$, $C=B-A^2/(4K)$; có thể tiếp tục đặt $v=e^{Ct}z$ để $z_t=Kz_{xx}$.

\textbf{Bước 2 — Đổi dữ liệu và biên cùng phép đổi ẩn.} $v(x,0)=e^{hx}f(x)$. Dirichlet thuần nhất giữ nguyên; Neumann biến thành $v_x-hv=0$, nên không giữ nguyên một cơ sở cosine nếu $A\ne0$. Với Dirichlet:


\[
u=e^{-hx}\sum_{n\ge1}\left(\frac2L\int_0^Le^{hy}f(y)\sin\frac{n\pi y}Ldy\right)
e^{[C-K(n\pi/L)^2]t}\sin\frac{n\pi x}L.
\]


\textbf{Bước 3 — Tính hệ số mũ–đa thức.} Dùng nguyên hàm


\[
\int e^{hy}\sin ky\,dy
=\frac{e^{hy}(h\sin ky-k\cos ky)}{h^2+k^2}.
\]


Với $y e^{hy}$ hoặc $y(1-y)e^{hy}$, lấy đạo hàm theo $h$ của tích phân đã tính trên cận cố định; trình bày đạo hàm để tìm hệ số thực. Cách này tránh để lại một tích phân sơ cấp chưa thực hiện.

\textbf{Bước 4 — Hệ số không gian biến thiên.} Với $x^2u_{xx}$ trên $x>0$, đặt $z=\ln x$: nếu $U(z,t)=u(e^z,t)$ thì $u_x=U_z/x$, $x^2u_{xx}=U_{zz}-U_z$. Đổi miền và biên sang biến $z$ rồi dùng bước 1. Biên không thuần nhất nên tìm trạng thái dừng trước khi khai triển.

\textbf{Bước 5 — Duy nhất và ổn định.} Với hiệu nghiệm Dirichlet, $I'=-K\int w_x^2+2BI\le2BI$ suy $w=0$ khi $I(0)=0$. Giới hạn phụ thuộc các số $C-K\lambda_n$ có hệ số khác không: âm thì giảm, bằng không thì tồn tại mode dừng, dương thì tăng. Phải kiểm tra tính thống nhất giữa miền PDE và vị trí biên trước khi dùng công thức.
\subsection*{Lời giải áp dụng cho bài}
\textbf{Bước 1. Kiểm tra miền.}
Đề ghi PDE trên $(0,2)$, hai biên tại $0,1$ và dữ liệu đầu chỉ trên $(0,1)$. Vì $1$ là điểm trong của miền PDE, đây không phải bài Dirichlet chuẩn trên đoạn $(0,2)$. Cần làm rõ miền dự kiến trước khi đưa ra đáp án theo tách biến trên đoạn.
\textbf{Bước 2. Giải bài nếu miền dự kiến là $(0,1)$.}
Đặt $h=\beta/2$, $u=e^{-hx}v$. Tính $u_x=e^{-hx}(v_x-hv)$, $u_{xx}=e^{-hx}(v_{xx}-2hv_x+h^2v)$; suy $v_t=v_{xx}-h^2v$ và $v(x,0)=e^{hx}x(1-x)$.
\textbf{Bước 3. Thực hiện tích phân hệ số.}
Với $k=n\pi$, $D=h^2+k^2,E=(-1)^ne^h$, nguyên hàm cho $J=\int_0^1e^{hy}\sin ky=k(1-E)/D$. Hệ số $2\int e^{hy}y(1-y)\sin ky=2(J'-J'')$:
\[c_n=\frac{4k\{(1-h)-(1+h)E\}}{D^2}
-\frac{16kh^2(1-E)}{D^3}.\]
Đạo hàm được tính từ $E'=E,D'=2h$: $J'=-kE/D-2kh(1-E)/D^2$ và $J''=-kE/D+(4khE-2k(1-E))/D^2+8kh^2(1-E)/D^3$.
\textbf{Bước 4. Ghép và kiểm tra.}
\[u=e^{-hx}\sum_{n\ge1}c_ne^{-(n^2\pi^2+h^2)t}\sin n\pi x,\quad0<x<1.\]
Mỗi mode thỏa biên, ODE và dữ liệu biến đổi; cộng lại trả về PDE gốc. Mọi mũ âm, nên giới hạn đều $0$ với mọi $\beta$ thực. Công thức trên giải miền sửa $(0,1)$. Nếu chủ ý giữ miền PDE $(0,2)$ và yêu cầu nghiệm cổ điển qua $1$ với $t>0$, tính giải tích theo không gian buộc thác triển khác: với $v=e^{hx}u$, phản xạ lẻ qua $1$ cho $u(x,t)=-e^{2h(1-x)}u(2-x,t)$ trên $(1,2)$. Nó xác định dữ liệu và biên phía phải không được in; không được tự coi sine $n\pi x$ là cơ sở Dirichlet tại $2$.

\clearpage\noindent\texttt{\detokenize{BT_4_5_3_05}}\par
\section*{Lời giải bổ sung — Bài tập 5 · Mục 4.5.3}

\subsection*{Kiến thức và trình tự giải}
\textbf{Bước 1 — Chọn đúng đại lượng.} Năng lượng bình phương là $I(t)=\frac12\int_0^Lu^2dx$. Để chứng minh duy nhất, áp dụng cho hiệu $w$ của hai nghiệm cùng dữ kiện; $w$ có dữ liệu đầu bằng không và biên thuần nhất. Tổng nhiệt $E=\int_0^Lu\,dx$ là đại lượng khác, dùng khi đề hỏi bảo toàn hoặc suy giảm tổng nhiệt.

\textbf{Bước 2 — Đạo hàm và tích phân từng phần.} Với $u_t=Ku_{xx}$,


\[
I'=\int_0^Luu_t\,dx
=K[uu_x]_0^L-K\int_0^Lu_x^2dx.
\]


Thay biên tại $L$ rồi trừ giá trị tại $0$. Dirichlet/Neumann thuần nhất làm hạng biên bằng không. Robin $u_x(0)=u(0)$ đóng góp $-u(0)^2$; Robin $u_x(L)=-u(L)$ đóng góp $-u(L)^2$. Biên tuần hoàn triệt tiêu hai đầu. Một điều kiện ghép khác phải được kiểm tra trực tiếp, vì hạng biên có thể dương.

\textbf{Bước 3 — Suy duy nhất và tính dương của trị riêng.} Nếu $I'\le0$ và $I(0)=0$ thì $0\le I(t)\le0$, nên $w=0$. Với $X''+\lambda X=0$, nhân với $X$ rồi tích phân:


\[
\lambda\int_0^LX^2dx=\int_0^L(X')^2dx-[XX']_0^L.
\]


Nếu vế phải không âm thì $\lambda\ge0$. Khi bằng không, buộc từng số hạng không âm bằng không và kiểm tra biên để quyết định có mode hằng hay chỉ nghiệm tầm thường.

\textbf{Bước 4 — Chứng minh trực giao.} Nhân hai phương trình riêng lần lượt với $X_m,X_n$, trừ rồi tích phân:


\[
(\lambda_n-\lambda_m)\int_0^LX_nX_mdx
=[X_nX_m'-X_n'X_m]_0^L.
\]


Biên tự liên hợp làm vế phải bằng không; $\lambda_n\ne\lambda_m$ suy trực giao. Với biên không tự liên hợp hoặc đạo hàm bậc ba, không tự động dùng công thức chiếu trực giao.

\textbf{Bước 5 — Khi có hạng phản ứng hoặc bậc đạo hàm khác.} Với Dirichlet và $u_t=Ku_{xx}+Au_x+Bu$, hạng $A\int uu_x=A[u^2]_0^L/2=0$, nên $I'=-K\int u_x^2+2BI\le2BI$; $e^{-2Bt}I$ không tăng, vẫn suy duy nhất. Với $u_t=Ku_{xxx}$, tích phân hai lần cho $I'=K[uu_{xx}-u_x^2/2]_0^L$. Với tổng nhiệt trong bài Neumann có $-ru$, $E'=K[u_x]_0^L-rE$; giải ODE này trước khi kết luận.
\subsection*{Lời giải áp dụng cho bài}
\textbf{Bước 1. Phân biệt tổng nhiệt với năng lượng bình phương.}
Đề đặt $E(t)=\int_0^Lu(x,t)dx$, là tổng nhiệt có dấu. Với $u_t=ku_{xx}-ru$ và hai biên Neumann $0$,
\begin{align*}
E'(t)&=\int_0^Lu_tdx=k\int_0^Lu_{xx}dx-r\int_0^Lu\,dx\\
&=k(u_x(L,t)-u_x(0,t))-rE(t)=-rE(t).
\end{align*}
\textbf{Bước 2. Giải ODE.}
Nhân với $e^{rt}$ cho $(e^{rt}E)'=0$, nên $e^{rt}E(t)=E(0)=\int_0^Lf$. Do đó
\[E(t)=e^{-rt}\int_0^Lf(x)dx.\]
Kết quả đúng cả khi $r=0$ (bảo toàn), $r>0$ (giảm độ lớn) và $r<0$ (tăng độ lớn). Với dữ liệu không trơn, thực hiện từ $\varepsilon>0$ rồi lấy giới hạn nếu $u(\cdot,t)\to f$ trong $L^1$.
\textbf{Bước 3. Kiểm tra bằng mode hằng.}
Nhân cosine có mode trung bình $c_0(t)=E(t)/L$; PDE chiếu lên $1$ cho $c_0'=-rc_0$. Các cosine khác có tích phân bằng $0$, nên không đóng góp tổng nhiệt. Không thay $E$ bằng $\frac12\int u^2$, vì đại lượng sau còn có hạng $-k\int u_x^2$ trong đạo hàm.

\clearpage\noindent\texttt{\detokenize{BT_4_5_3_06}}\par
\section*{Lời giải bổ sung — Bài tập 6 · Mục 4.5.3}

\subsection*{Kiến thức và trình tự giải}
\textbf{Bước 1 — Loại đạo hàm cấp một bằng tính đạo hàm đầy đủ.} Với $u_t=Ku_{xx}+Au_x+Bu$, đặt $u=e^{-hx}v$, $h=A/(2K)$. Khi đó


\[
u_t=e^{-hx}v_t,\quad u_x=e^{-hx}(v_x-hv),\quad
u_{xx}=e^{-hx}(v_{xx}-2hv_x+h^2v).
\]


Thay vào PDE, hệ số $v_x$ là $A-2Kh=0$, hệ số $v$ là $Kh^2-Ah+B=B-A^2/(4K)$. Vì thế $v_t=Kv_{xx}+Cv$, $C=B-A^2/(4K)$; có thể tiếp tục đặt $v=e^{Ct}z$ để $z_t=Kz_{xx}$.

\textbf{Bước 2 — Đổi dữ liệu và biên cùng phép đổi ẩn.} $v(x,0)=e^{hx}f(x)$. Dirichlet thuần nhất giữ nguyên; Neumann biến thành $v_x-hv=0$, nên không giữ nguyên một cơ sở cosine nếu $A\ne0$. Với Dirichlet:


\[
u=e^{-hx}\sum_{n\ge1}\left(\frac2L\int_0^Le^{hy}f(y)\sin\frac{n\pi y}Ldy\right)
e^{[C-K(n\pi/L)^2]t}\sin\frac{n\pi x}L.
\]


\textbf{Bước 3 — Tính hệ số mũ–đa thức.} Dùng nguyên hàm


\[
\int e^{hy}\sin ky\,dy
=\frac{e^{hy}(h\sin ky-k\cos ky)}{h^2+k^2}.
\]


Với $y e^{hy}$ hoặc $y(1-y)e^{hy}$, lấy đạo hàm theo $h$ của tích phân đã tính trên cận cố định; trình bày đạo hàm để tìm hệ số thực. Cách này tránh để lại một tích phân sơ cấp chưa thực hiện.

\textbf{Bước 4 — Hệ số không gian biến thiên.} Với $x^2u_{xx}$ trên $x>0$, đặt $z=\ln x$: nếu $U(z,t)=u(e^z,t)$ thì $u_x=U_z/x$, $x^2u_{xx}=U_{zz}-U_z$. Đổi miền và biên sang biến $z$ rồi dùng bước 1. Biên không thuần nhất nên tìm trạng thái dừng trước khi khai triển.

\textbf{Bước 5 — Duy nhất và ổn định.} Với hiệu nghiệm Dirichlet, $I'=-K\int w_x^2+2BI\le2BI$ suy $w=0$ khi $I(0)=0$. Giới hạn phụ thuộc các số $C-K\lambda_n$ có hệ số khác không: âm thì giảm, bằng không thì tồn tại mode dừng, dương thì tăng. Phải kiểm tra tính thống nhất giữa miền PDE và vị trí biên trước khi dùng công thức.
\subsection*{Lời giải áp dụng cho bài}
\textbf{(i) Dữ liệu đầu và biên đều bằng $0$.}
Với $u_t=u_{xx}+2u_x$ trên $(0,\pi)$, $u=0$ thỏa mọi dữ kiện. Hiệu nghiệm có $I'=-\int u_x^2+[u^2]_0^\pi=-\int u_x^2\le0$ và $I(0)=0$, nên đây là nghiệm duy nhất.

\textbf{(ii) Hệ số $x^2$, biên tại $1,e$.}
Đặt $z=\ln x$, $U(z,t)=u(e^z,t)$. Tính $u_x=U_z/x$, $x^2u_{xx}=U_{zz}-U_z$, nên $U_t=U_{zz}-U_z+4U$ trên $0<z<1$. Đặt $U=e^{z/2}v$:
\[U_{zz}-U_z+4U=e^{z/2}(v_{zz}+15v/4),\quad v_t=v_{zz}+15v/4.\]
Biên $v(0)=-1$, $v(1)=e^{-1/2}$, dữ liệu đầu $v(z,0)=e^{z/2}$. Đặt $b=\sqrt{15}/2$, $D=(e^{-1/2}+\cos b)/\sin b$. Nghiệm dừng $S=-\cos bz+D\sin bz$ giải $S''+b^2S=0$, đúng hai biên. Vì $v-S$ có biên $0$,
\[u=P(x)+\sqrt x\sum_{n\ge1}A_ne^{(15/4-n^2\pi^2)t}\sin(n\pi\ln x),\quad
P(x)=\sqrt x\,S(\ln x).\]
Để tính $A_n$, đặt $k=n\pi,s=(-1)^n$:
\begin{align*}
\int_0^1e^{z/2}\sin kz\,dz&=\frac{k(1-se^{1/2})}{k^2+1/4},\\
\int_0^1\cos bz\sin kz\,dz&=\frac{k(1-s\cos b)}{k^2-b^2},\\
\int_0^1\sin bz\sin kz\,dz&=-\frac{sk\sin b}{k^2-b^2}.
\end{align*}
Các đẳng thức lần lượt dùng nguyên hàm mũ–sine và tích–tổng. Cộng theo dữ liệu $e^{z/2}+\cos bz-D\sin bz$ cho
\[A_n=2k\left\{\frac{1-se^{1/2}}{k^2+1/4}
+\frac{1+se^{-1/2}}{k^2-15/4}\right\}.\]
Mọi $15/4-n^2\pi^2<0$, nên $u\to P$ đều. $P(1)=-1,P(e)=1$, $x^2P''+4P=0$; dữ liệu đầu $x$ không khớp giá trị biên tại góc, nên chỉ có liên tục đến mặt đầu trong miền.

\textbf{(iii) Vận chuyển với dữ liệu $x$ trên $(0,1)$.}
Đặt $u=e^{-x}v$ cho $v_t=v_{xx}-v$, $v(x,0)=xe^x$, biên $0$. Với $k=n\pi,s=(-1)^n$, nguyên hàm cho $J(h)=\int_0^1e^{hy}\sin ky=k(1-se^h)/(h^2+k^2)$. Do đó
\[A_n=2J'(1)=-\frac{2kse}{1+k^2}-\frac{4k(1-se)}{(1+k^2)^2},\qquad
u=e^{-x}\sum_{n\ge1}A_ne^{-(n^2\pi^2+1)t}\sin n\pi x.\]
Tại $t=0$, sine trả lại $xe^x$, rồi nhân $e^{-x}$ trả lại $x$ trong miền. Sine kiểm tra hai biên; phép đổi ẩn kiểm tra PDE. Giá trị $f(1)=1$ không khớp biên $0$ ở góc.

\textbf{(iv) Nguồn $-x$, Dirichlet–Neumann.}
Nghiệm dừng thỏa $P''=x,P(0)=-1,P'(\pi)=0$: tích phân cho $P=x^3/6-\pi^2x/2-1$. Khi $w=u-P$, nguồn bằng $0$, biên $w(0)=w_x(\pi)=0$, mặt đầu $q=1+\pi^2x/2$. Cơ sở $X_n=\sin k_nx$, $k_n=n+1/2$, $N_n=\pi/2$.
\[\int_0^\pi\sin k_ny\,dy=1/k_n,\quad
\int_0^\pi y\sin k_ny\,dy=(-1)^n/k_n^2.\]
Vì vậy
\[A_n=\frac2{\pi k_n}+\frac{\pi(-1)^n}{k_n^2},\qquad
u=P+\sum_{n\ge0}A_ne^{-k_n^2t}\sin k_nx.\]
Biên đúng, mặt đầu trong miền bằng $x^3/6$; $P''-x=0$ và từng mode cho PDE. Giới hạn đều $P$, nhưng mặt đầu không khớp $u(0,0)=-1$.

\textbf{(v) Nguồn và đạo hàm biên phụ thuộc thời gian.}
Đặt $u=xt+w$, vì $(xt)(0,t)=0$, $(xt)_x(\pi,t)=t$. Từ PDE:
\[w_t=w_{xx}+w+\sin x+x(t-1+1/\pi),\quad
w(0,t)=w_x(\pi,t)=0,\quad w(x,0)=0.\]
Với $k_n=n+1/2,s_n=(-1)^n$, tích–tổng và tích phân từng phần cho
\[S_n=\frac2\pi\int_0^\pi\sin y\sin k_ny\,dy
=\frac{2s_n}{\pi(1-k_n^2)},\quad
R_n=\frac2\pi\int_0^\pi y\sin k_ny\,dy=\frac{2s_n}{\pi k_n^2}.\]
Đặt $d_n=1-k_n^2\ne0$, $C_n=S_n+R_n(-1+1/\pi)$. ODE là $T_n'=d_nT_n+R_nt+C_n$, $T_n(0)=0$. Hai tích phân thời gian thực hiện đến hết:
\[\int_0^te^{d(t-s)}ds=\frac{e^{dt}-1}d,\quad
\int_0^te^{d(t-s)}s\,ds=\frac{e^{dt}-1-dt}{d^2}.\]
Do đó
\[u=xt+\sum_{n\ge0}\left\{
R_n\frac{e^{d_nt}-1-d_nt}{d_n^2}
+C_n\frac{e^{d_nt}-1}{d_n}\right\}\sin k_nx.\]
Mỗi mode đúng ODE và biên thuần nhất; cộng $xt$ trả lại nguồn và đạo hàm biên. Các hệ số $O(k_n^{-4})$ ở thời gian cố định cho hội tụ của nghiệm. Mode $n=0$ có $d_0=3/4>0$; không được kết luận mọi mode đều giảm.

\textbf{(vi) Biên tuần hoàn chu kỳ $1$.}
$f=1+2\cos2\pi x-2\sin3\pi x$ đối xứng qua $1/2$, nên các hệ số sine chu kỳ $1$ bằng $0$ (đổi biến $x\mapsto1-x$ trong tích phân). Trung bình:
\[c_0=\int_0^1f\,dx=1-\frac4{3\pi}.\]
Với $n\ge1$, $\int_0^1\sin3\pi x\cos2n\pi x\,dx=6/[\pi(9-4n^2)]$ từ tích–tổng. Vì vậy
\[c_n=2\delta_{n1}-\frac{24}{\pi(9-4n^2)},\qquad
u=1-\frac4{3\pi}+\sum_{n\ge1}c_ne^{-4n^2\pi^2t}\cos2n\pi x.\]
Các cosine có cùng giá trị và đạo hàm tại $0,1$; ODE kiểm tra PDE, chuỗi ban đầu trả lại $f$ liên tục. $f'(0)=-6\pi$, $f'(1)=6\pi$ không khớp đạo hàm tuần hoàn ở góc, nên không trơn đến góc. Giới hạn đều là $1-4/(3\pi)$.

\clearpage\noindent\texttt{\detokenize{VD_4_6_2_01}}\par
\section*{Lời giải bổ sung — Ví dụ 1 · Mục 4.6.2}

\subsection*{Kiến thức và trình tự giải}
\textbf{Bước 1 — Chọn nhân nhiệt theo miền và hệ số.} Trên $\mathbb R$, $g_K(z,t)=(4\pi Kt)^{-1/2}e^{-z^2/(4Kt)}$, $u=g_K*f$. Trên nửa trục, thác triển lẻ cho Dirichlet $0$, chẵn cho Neumann $0$; sau đổi $y\mapsto-y$ ở nửa âm:


\[
H_D(x,y,t)=g_K(x-y,t)-g_K(x+y,t),\qquad
H_N(x,y,t)=g_K(x-y,t)+g_K(x+y,t).
\]


Kiểm tra $H_D(0,y,t)=0$, $\partial_xH_N(0,y,t)=0$. Nửa trục phải có một điều kiện biên; dữ liệu đầu một mình không xác định duy nhất.

\textbf{Bước 2 — Tính các tích phân không gian.} Mode $\sin ky,\cos ky$ biến thành $e^{-Kk^2t}\sin kx,e^{-Kk^2t}\cos kx$. Với Gauss $e^{-\alpha(y-m)^2}$, $\alpha>0$, hoàn thành bình phương:


\[
-\frac{(x-y)^2}{4Kt}-\alpha(y-m)^2
=-\frac{1+4\alpha Kt}{4Kt}\left(y-\frac{x+4\alpha Ktm}{1+4\alpha Kt}\right)^2
-\frac{\alpha(x-m)^2}{1+4\alpha Kt}.
\]


Đổi biến để $\int_{\mathbb R}e^{-z^2}dz=\sqrt\pi$, suy $(g_K*f)(x)=e^{-\alpha(x-m)^2/(1+4\alpha Kt)}/\sqrt{1+4\alpha Kt}$. Với dữ liệu chỉ thị đoạn, cận hữu hạn tạo $\operatorname{erf}$; với $e^{-y}$ trên nửa trục, cận $0,\infty$ tạo $\operatorname{erfc}$, trong đó $\operatorname{erf}z=2\int_0^ze^{-r^2}dr/\sqrt\pi$, $\operatorname{erfc}z=1-\operatorname{erf}z$.

\textbf{Bước 3 — Thêm nguồn và biên.} Nguồn cho $u=g_K*f+\int_0^tg_K(\cdot,t-s)*F(\cdot,s)\,ds$, hoặc thay $g_K$ bằng $H_D,H_N$. Biên Dirichlet $\mu(t)$ trên nửa trục thêm


\[
\int_0^t\frac{x\,e^{-x^2/[4K(t-s)]}}{\sqrt{4\pi K}(t-s)^{3/2}}\mu(s)ds.
\]


Nguồn một mode giảm cùng tốc độ khuếch tán cho tích phân thời gian hằng sau khi gom mũ, nên xuất hiện $te^{-Kk^2t}$.

\textbf{Bước 4 — Đổi biến khi hệ số phụ thuộc thời gian hoặc có vận chuyển.} Với $u_t=a(t)u_{xx}+b(t)u_x+c(t)u$, đặt $A=\int_0^ta$, $B=\int_0^tb$, $C=\int_0^tc$, $u=e^Cv(x+B,A)$. Tính $u_t=e^C(cv+bv_z+av_\tau)$, $u_x=e^Cv_z$, $u_{xx}=e^Cv_{zz}$, suy $v_\tau=v_{zz}$. Dấu dịch là $x+B$, không phải $x-B$. Với hạng $xu_x$, dùng $z=e^{\pm t}x$ và chọn đồng hồ $\tau'(t)=e^{\pm2t}$.

\textbf{Bước 5 — Nhiều chiều và kiểm tra.} Nhân $n$ chiều là tích các nhân một chiều. Dữ liệu dạng tích tách được tích phân không gian bằng Fubini; nguồn dạng tích vẫn phải tích phân theo một thời gian chung $s$. Kiểm tra PDE bằng đạo hàm Duhamel: đạo hàm cận trên trả lại $F(x,t)$, phần còn lại bằng $K\Delta u$. Nhân Gauss là xấp xỉ đơn vị, nên trả lại dữ liệu tại điểm liên tục; dữ liệu nhảy không cho nghiệm liên tục đến mọi điểm lúc đầu.
\subsection*{Lời giải áp dụng cho bài}
\textbf{Bước 1. Nhân nhiệt trên toàn trục.}
Đặt $g_K(z,t)=(4\pi Kt)^{-1/2}e^{-z^2/(4Kt)}$. Nghiệm $u=\int_{\mathbb R}g_K(x-y,t)f(y)dy$. Tích phân của $g_K$ bằng $1$, và biến đổi Fourier của nó là $e^{-Kt\xi^2}$, nên một mode $\cos ky$ biến thành $e^{-Kk^2t}\cos kx$.
\textbf{Bước 2. Câu (i), $K=1,f=\cos^23x$.}
$\cos^23x=(1+\cos6x)/2$, suy
\[u=\frac12+\frac12e^{-36t}\cos6x.\]
Kiểm tra $u_t=-18e^{-36t}\cos6x=u_{xx}$, $u(x,0)=\cos^23x$.
\textbf{Bước 3. Câu (ii), $K=2,f=e^{-x^2}$.}
\[u=\frac1{\sqrt{8\pi t}}\int_{\mathbb R}
\exp\left(-\frac{(x-y)^2}{8t}-y^2\right)dy.\]
Hoàn thành bình phương:
\[-\frac{(x-y)^2}{8t}-y^2
=-\frac{1+8t}{8t}\left(y-\frac{x}{1+8t}\right)^2-\frac{x^2}{1+8t}.\]
Đặt $z=\sqrt{(1+8t)/(8t)}(y-x/(1+8t))$, $dy=\sqrt{8t/(1+8t)}dz$. Tích phân $\int e^{-z^2}dz=\sqrt\pi$ cho
\[u=\frac1{\sqrt{1+8t}}\exp\left(-\frac{x^2}{1+8t}\right).\]
\textbf{Bước 4. Kiểm tra đạo hàm.} Đặt $D=1+8t$. $u_t/u=-4/D+8x^2/D^2$, $u_{xx}/u=-2/D+4x^2/D^2$, nên $u_t=2u_{xx}$. Khi $t\downarrow0$, $D\to1$ trả lại $f$.
\begin{center}\includegraphics[width=.97\linewidth,height=11cm,keepaspectratio]{_xu_ly/graphs/VD_4_6_2_01.png}\end{center}
\textit{Đồ thị minh họa; với nghiệm số, xem phương pháp và tham số trong mã Python kèm theo.}

\clearpage\noindent\texttt{\detokenize{VD_4_6_2_02}}\par
\section*{Lời giải bổ sung — Ví dụ 2 · Mục 4.6.2}

\subsection*{Kiến thức và trình tự giải}
\textbf{Bước 1 — Chọn nhân nhiệt theo miền và hệ số.} Trên $\mathbb R$, $g_K(z,t)=(4\pi Kt)^{-1/2}e^{-z^2/(4Kt)}$, $u=g_K*f$. Trên nửa trục, thác triển lẻ cho Dirichlet $0$, chẵn cho Neumann $0$; sau đổi $y\mapsto-y$ ở nửa âm:


\[
H_D(x,y,t)=g_K(x-y,t)-g_K(x+y,t),\qquad
H_N(x,y,t)=g_K(x-y,t)+g_K(x+y,t).
\]


Kiểm tra $H_D(0,y,t)=0$, $\partial_xH_N(0,y,t)=0$. Nửa trục phải có một điều kiện biên; dữ liệu đầu một mình không xác định duy nhất.

\textbf{Bước 2 — Tính các tích phân không gian.} Mode $\sin ky,\cos ky$ biến thành $e^{-Kk^2t}\sin kx,e^{-Kk^2t}\cos kx$. Với Gauss $e^{-\alpha(y-m)^2}$, $\alpha>0$, hoàn thành bình phương:


\[
-\frac{(x-y)^2}{4Kt}-\alpha(y-m)^2
=-\frac{1+4\alpha Kt}{4Kt}\left(y-\frac{x+4\alpha Ktm}{1+4\alpha Kt}\right)^2
-\frac{\alpha(x-m)^2}{1+4\alpha Kt}.
\]


Đổi biến để $\int_{\mathbb R}e^{-z^2}dz=\sqrt\pi$, suy $(g_K*f)(x)=e^{-\alpha(x-m)^2/(1+4\alpha Kt)}/\sqrt{1+4\alpha Kt}$. Với dữ liệu chỉ thị đoạn, cận hữu hạn tạo $\operatorname{erf}$; với $e^{-y}$ trên nửa trục, cận $0,\infty$ tạo $\operatorname{erfc}$, trong đó $\operatorname{erf}z=2\int_0^ze^{-r^2}dr/\sqrt\pi$, $\operatorname{erfc}z=1-\operatorname{erf}z$.

\textbf{Bước 3 — Thêm nguồn và biên.} Nguồn cho $u=g_K*f+\int_0^tg_K(\cdot,t-s)*F(\cdot,s)\,ds$, hoặc thay $g_K$ bằng $H_D,H_N$. Biên Dirichlet $\mu(t)$ trên nửa trục thêm


\[
\int_0^t\frac{x\,e^{-x^2/[4K(t-s)]}}{\sqrt{4\pi K}(t-s)^{3/2}}\mu(s)ds.
\]


Nguồn một mode giảm cùng tốc độ khuếch tán cho tích phân thời gian hằng sau khi gom mũ, nên xuất hiện $te^{-Kk^2t}$.

\textbf{Bước 4 — Đổi biến khi hệ số phụ thuộc thời gian hoặc có vận chuyển.} Với $u_t=a(t)u_{xx}+b(t)u_x+c(t)u$, đặt $A=\int_0^ta$, $B=\int_0^tb$, $C=\int_0^tc$, $u=e^Cv(x+B,A)$. Tính $u_t=e^C(cv+bv_z+av_\tau)$, $u_x=e^Cv_z$, $u_{xx}=e^Cv_{zz}$, suy $v_\tau=v_{zz}$. Dấu dịch là $x+B$, không phải $x-B$. Với hạng $xu_x$, dùng $z=e^{\pm t}x$ và chọn đồng hồ $\tau'(t)=e^{\pm2t}$.

\textbf{Bước 5 — Nhiều chiều và kiểm tra.} Nhân $n$ chiều là tích các nhân một chiều. Dữ liệu dạng tích tách được tích phân không gian bằng Fubini; nguồn dạng tích vẫn phải tích phân theo một thời gian chung $s$. Kiểm tra PDE bằng đạo hàm Duhamel: đạo hàm cận trên trả lại $F(x,t)$, phần còn lại bằng $K\Delta u$. Nhân Gauss là xấp xỉ đơn vị, nên trả lại dữ liệu tại điểm liên tục; dữ liệu nhảy không cho nghiệm liên tục đến mọi điểm lúc đầu.
\subsection*{Lời giải áp dụng cho bài}
\textbf{Bước 1. Thác triển lẻ.}
Với $u(0,t)=0$, đặt $\widetilde f(y)=f(y)$ khi $y>0$, $\widetilde f(y)=-f(-y)$ khi $y<0$. Vì Gauss chẵn, nghiệm toàn trục của dữ liệu lẻ vẫn lẻ, nên bằng $0$ tại $x=0$. Tách tích phân và đổi $y\mapsto-y$ trên nửa âm:
\[u=\int_0^\infty[g_K(x-y,t)-g_K(x+y,t)]f(y)dy.\]
\textbf{Bước 2. Câu cụ thể $K=2,f=e^{-y}$.}
Đặt $I_\pm=\int_0^\infty g_2(x\mp y,t)e^{-y}dy$, trong đó $I_+$ dùng $x-y$, $I_-$ dùng $x+y$. Hoàn thành bình phương:
\begin{align*}
-\frac{(x-y)^2}{8t}-y&=-\frac{(y+4t-x)^2}{8t}+2t-x,\\
-\frac{(x+y)^2}{8t}-y&=-\frac{(y+4t+x)^2}{8t}+2t+x.
\end{align*}
Đổi biến $z=(y+4t\mp x)/\sqrt{8t}$; cận $y=0$ trở thành $(4t\mp x)/\sqrt{8t}$, và dùng $\int_a^\infty e^{-z^2}dz=\sqrt\pi\,\operatorname{erfc}(a)/2$:
\[I_+=\tfrac12e^{2t-x}\operatorname{erfc}\frac{4t-x}{\sqrt{8t}},\quad
I_-=\tfrac12e^{2t+x}\operatorname{erfc}\frac{4t+x}{\sqrt{8t}}.\]
Vậy
\[u=\frac12e^{2t-x}\operatorname{erfc}\frac{4t-x}{\sqrt{8t}}
-\frac12e^{2t+x}\operatorname{erfc}\frac{4t+x}{\sqrt{8t}}.\]
\textbf{Bước 3. Kiểm tra.} Tại $x=0$ hai hạng triệt tiêu. Với $x>0,t\downarrow0$, đối số thứ nhất tiến về $-\infty$ nên erfc tiến về $2$, đối số thứ hai về $\infty$ nên erfc về $0$; suy $u\to e^{-x}$. Nhân ảnh giải PDE nên $u_t=2u_{xx}$. $f(0)=1$ không khớp biên $0$, nên không liên tục ở góc $(0,0)$; nghiệm trơn với $x\ge0,t>0$.

\clearpage\noindent\texttt{\detokenize{VD_4_6_2_03}}\par
\section*{Lời giải bổ sung — Ví dụ 3 · Mục 4.6.2}

\subsection*{Kiến thức và trình tự giải}
\textbf{Bước 1 — Chọn nhân nhiệt theo miền và hệ số.} Trên $\mathbb R$, $g_K(z,t)=(4\pi Kt)^{-1/2}e^{-z^2/(4Kt)}$, $u=g_K*f$. Trên nửa trục, thác triển lẻ cho Dirichlet $0$, chẵn cho Neumann $0$; sau đổi $y\mapsto-y$ ở nửa âm:


\[
H_D(x,y,t)=g_K(x-y,t)-g_K(x+y,t),\qquad
H_N(x,y,t)=g_K(x-y,t)+g_K(x+y,t).
\]


Kiểm tra $H_D(0,y,t)=0$, $\partial_xH_N(0,y,t)=0$. Nửa trục phải có một điều kiện biên; dữ liệu đầu một mình không xác định duy nhất.

\textbf{Bước 2 — Tính các tích phân không gian.} Mode $\sin ky,\cos ky$ biến thành $e^{-Kk^2t}\sin kx,e^{-Kk^2t}\cos kx$. Với Gauss $e^{-\alpha(y-m)^2}$, $\alpha>0$, hoàn thành bình phương:


\[
-\frac{(x-y)^2}{4Kt}-\alpha(y-m)^2
=-\frac{1+4\alpha Kt}{4Kt}\left(y-\frac{x+4\alpha Ktm}{1+4\alpha Kt}\right)^2
-\frac{\alpha(x-m)^2}{1+4\alpha Kt}.
\]


Đổi biến để $\int_{\mathbb R}e^{-z^2}dz=\sqrt\pi$, suy $(g_K*f)(x)=e^{-\alpha(x-m)^2/(1+4\alpha Kt)}/\sqrt{1+4\alpha Kt}$. Với dữ liệu chỉ thị đoạn, cận hữu hạn tạo $\operatorname{erf}$; với $e^{-y}$ trên nửa trục, cận $0,\infty$ tạo $\operatorname{erfc}$, trong đó $\operatorname{erf}z=2\int_0^ze^{-r^2}dr/\sqrt\pi$, $\operatorname{erfc}z=1-\operatorname{erf}z$.

\textbf{Bước 3 — Thêm nguồn và biên.} Nguồn cho $u=g_K*f+\int_0^tg_K(\cdot,t-s)*F(\cdot,s)\,ds$, hoặc thay $g_K$ bằng $H_D,H_N$. Biên Dirichlet $\mu(t)$ trên nửa trục thêm


\[
\int_0^t\frac{x\,e^{-x^2/[4K(t-s)]}}{\sqrt{4\pi K}(t-s)^{3/2}}\mu(s)ds.
\]


Nguồn một mode giảm cùng tốc độ khuếch tán cho tích phân thời gian hằng sau khi gom mũ, nên xuất hiện $te^{-Kk^2t}$.

\textbf{Bước 4 — Đổi biến khi hệ số phụ thuộc thời gian hoặc có vận chuyển.} Với $u_t=a(t)u_{xx}+b(t)u_x+c(t)u$, đặt $A=\int_0^ta$, $B=\int_0^tb$, $C=\int_0^tc$, $u=e^Cv(x+B,A)$. Tính $u_t=e^C(cv+bv_z+av_\tau)$, $u_x=e^Cv_z$, $u_{xx}=e^Cv_{zz}$, suy $v_\tau=v_{zz}$. Dấu dịch là $x+B$, không phải $x-B$. Với hạng $xu_x$, dùng $z=e^{\pm t}x$ và chọn đồng hồ $\tau'(t)=e^{\pm2t}$.

\textbf{Bước 5 — Nhiều chiều và kiểm tra.} Nhân $n$ chiều là tích các nhân một chiều. Dữ liệu dạng tích tách được tích phân không gian bằng Fubini; nguồn dạng tích vẫn phải tích phân theo một thời gian chung $s$. Kiểm tra PDE bằng đạo hàm Duhamel: đạo hàm cận trên trả lại $F(x,t)$, phần còn lại bằng $K\Delta u$. Nhân Gauss là xấp xỉ đơn vị, nên trả lại dữ liệu tại điểm liên tục; dữ liệu nhảy không cho nghiệm liên tục đến mọi điểm lúc đầu.
\subsection*{Lời giải áp dụng cho bài}
\textbf{Bước 1. Chuyển bài $u_t=u_{xx}+xu$ về nhiệt.}
Đặt $u=e^{tx+t^3/3}v(z,t)$, $z=x+t^2$. Tính
\begin{align*}
u_t&=e^{tx+t^3/3}\{(x+t^2)v+2tv_z+v_t\},\\
u_{xx}+xu&=e^{tx+t^3/3}\{(t^2+x)v+2tv_z+v_{zz}\}.
\end{align*}
Do đó $v_t=v_{zz}$, $v(z,0)=f(z)$. Phép đổi biến này biện minh công thức cho $f$ liên tục bị chặn mà không cần đánh giá $\widehat f$ tại tần số phức.
\textbf{Bước 2. Fourier cho bài đã đổi.} Với quy ước $\widehat v=\int e^{-i\xi z}v\,dz$, $\widehat v_t=-\xi^2\widehat v$, suy $\widehat v=e^{-t\xi^2}\widehat f$. Biến đổi ngược cho $v=g_1*f$. Vì vậy
\[u=e^{tx+t^3/3}\int_{\mathbb R}g_1(x+t^2-y,t)f(y)dy.\]
Gom mũ:
\[tx+t^3/3-\frac{(x+t^2-y)^2}{4t}
=-\frac{(x-y)^2}{4t}+\frac{t(x+y)}2+\frac{t^3}{12}.\]
Vậy
\[u=\int_{\mathbb R}G(x,y,t)f(y)dy,\qquad
G=\frac1{\sqrt{4\pi t}}
\exp\left[-\frac{(x-y)^2}{4t}+\frac{t(x+y)}2+\frac{t^3}{12}\right].\]
\textbf{Bước 3. Kiểm tra Green.} $G>0$, đối xứng theo $x,y$, trơn với $t>0$. Với $\ell=\log G$,
\[\ell_x=-\frac{x-y}{2t}+\frac t2,\quad \ell_{xx}=-\frac1{2t},\quad
\ell_t=-\frac1{2t}+\frac{(x-y)^2}{4t^2}+\frac{x+y}2+\frac{t^2}4.\]
Ta có $\ell_t=\ell_{xx}+\ell_x^2+x$, suy $G_t=G_{xx}+xG$. Từ biểu diễn Gauss đã dịch,
\[\int_{\mathbb R}G(x,y,t)dy=e^{tx+t^3/3}.\]
Nhân không có tổng bằng $1$ do hạng phản ứng $xu$.
\textbf{Bước 4. Dữ liệu đầu và sự tập trung.} Với $x$ trong tập compact, hệ số $e^{tx+t^3/3}\to1$ và tâm $x+t^2\to x$. Gauss là xấp xỉ đơn vị, nên $u(x,t)\to f(x)$ tại điểm liên tục, cũng liên tục khi $(x,t)\to(x_0,0)$. Với $\delta>0$, tích phân nhân trên $|y-x|>\delta$ tiến về $0$. Điều kiện này phải là khoảng cách $|y-x|$, không phải $|x|$ như một dòng sai biến trong bản in. Công thức không bảo đảm nghiệm bị chặn theo $x$ khi $t>0$, vì phản ứng tuyến tính theo $x$ có thể làm nghiệm tăng theo mũ.

\clearpage\noindent\texttt{\detokenize{VD_4_6_2_04}}\par
\section*{Lời giải bổ sung — Ví dụ 4 · Mục 4.6.2}

\subsection*{Kiến thức và trình tự giải}
\textbf{Bước 1 — Chọn nhân nhiệt theo miền và hệ số.} Trên $\mathbb R$, $g_K(z,t)=(4\pi Kt)^{-1/2}e^{-z^2/(4Kt)}$, $u=g_K*f$. Trên nửa trục, thác triển lẻ cho Dirichlet $0$, chẵn cho Neumann $0$; sau đổi $y\mapsto-y$ ở nửa âm:


\[
H_D(x,y,t)=g_K(x-y,t)-g_K(x+y,t),\qquad
H_N(x,y,t)=g_K(x-y,t)+g_K(x+y,t).
\]


Kiểm tra $H_D(0,y,t)=0$, $\partial_xH_N(0,y,t)=0$. Nửa trục phải có một điều kiện biên; dữ liệu đầu một mình không xác định duy nhất.

\textbf{Bước 2 — Tính các tích phân không gian.} Mode $\sin ky,\cos ky$ biến thành $e^{-Kk^2t}\sin kx,e^{-Kk^2t}\cos kx$. Với Gauss $e^{-\alpha(y-m)^2}$, $\alpha>0$, hoàn thành bình phương:


\[
-\frac{(x-y)^2}{4Kt}-\alpha(y-m)^2
=-\frac{1+4\alpha Kt}{4Kt}\left(y-\frac{x+4\alpha Ktm}{1+4\alpha Kt}\right)^2
-\frac{\alpha(x-m)^2}{1+4\alpha Kt}.
\]


Đổi biến để $\int_{\mathbb R}e^{-z^2}dz=\sqrt\pi$, suy $(g_K*f)(x)=e^{-\alpha(x-m)^2/(1+4\alpha Kt)}/\sqrt{1+4\alpha Kt}$. Với dữ liệu chỉ thị đoạn, cận hữu hạn tạo $\operatorname{erf}$; với $e^{-y}$ trên nửa trục, cận $0,\infty$ tạo $\operatorname{erfc}$, trong đó $\operatorname{erf}z=2\int_0^ze^{-r^2}dr/\sqrt\pi$, $\operatorname{erfc}z=1-\operatorname{erf}z$.

\textbf{Bước 3 — Thêm nguồn và biên.} Nguồn cho $u=g_K*f+\int_0^tg_K(\cdot,t-s)*F(\cdot,s)\,ds$, hoặc thay $g_K$ bằng $H_D,H_N$. Biên Dirichlet $\mu(t)$ trên nửa trục thêm


\[
\int_0^t\frac{x\,e^{-x^2/[4K(t-s)]}}{\sqrt{4\pi K}(t-s)^{3/2}}\mu(s)ds.
\]


Nguồn một mode giảm cùng tốc độ khuếch tán cho tích phân thời gian hằng sau khi gom mũ, nên xuất hiện $te^{-Kk^2t}$.

\textbf{Bước 4 — Đổi biến khi hệ số phụ thuộc thời gian hoặc có vận chuyển.} Với $u_t=a(t)u_{xx}+b(t)u_x+c(t)u$, đặt $A=\int_0^ta$, $B=\int_0^tb$, $C=\int_0^tc$, $u=e^Cv(x+B,A)$. Tính $u_t=e^C(cv+bv_z+av_\tau)$, $u_x=e^Cv_z$, $u_{xx}=e^Cv_{zz}$, suy $v_\tau=v_{zz}$. Dấu dịch là $x+B$, không phải $x-B$. Với hạng $xu_x$, dùng $z=e^{\pm t}x$ và chọn đồng hồ $\tau'(t)=e^{\pm2t}$.

\textbf{Bước 5 — Nhiều chiều và kiểm tra.} Nhân $n$ chiều là tích các nhân một chiều. Dữ liệu dạng tích tách được tích phân không gian bằng Fubini; nguồn dạng tích vẫn phải tích phân theo một thời gian chung $s$. Kiểm tra PDE bằng đạo hàm Duhamel: đạo hàm cận trên trả lại $F(x,t)$, phần còn lại bằng $K\Delta u$. Nhân Gauss là xấp xỉ đơn vị, nên trả lại dữ liệu tại điểm liên tục; dữ liệu nhảy không cho nghiệm liên tục đến mọi điểm lúc đầu.
\subsection*{Lời giải áp dụng cho bài}
\textbf{Bước 1. Nhận dạng một mode cộng hưởng.}
Nguồn $e^{-t}\cos x$ và mặt đầu $0$ cho phép đặt $u=T(t)\cos x$. Khi đó $u_{xx}=-T\cos x$, PDE thành $T'+T=e^{-t}$, $T(0)=0$.
\textbf{Bước 2. Thực hiện tích phân.} Nhân với $e^t$:
\[(e^tT)'=1,\quad e^tT=t,\quad T=te^{-t}.\]
Tương đương Duhamel: $u=\int_0^te^{-(t-s)}e^{-s}\cos x\,ds=e^{-t}\cos x\int_0^tds$.
\textbf{Bước 3. Kiểm tra.}
\[u=te^{-t}\cos x,\quad u_t=(1-t)e^{-t}\cos x,\quad
u_{xx}=-te^{-t}\cos x.\]
Suy $u_t-u_{xx}=e^{-t}\cos x$ và $u(x,0)=0$. Thừa số $t$ xuất hiện vì tốc độ giảm của nguồn bằng tốc độ giảm của mode nhiệt.

\clearpage\noindent\texttt{\detokenize{VD_4_6_2_05}}\par
\section*{Lời giải bổ sung — Ví dụ 5 · Mục 4.6.2}

\subsection*{Kiến thức và trình tự giải}
\textbf{Bước 1 — Chọn nhân nhiệt theo miền và hệ số.} Trên $\mathbb R$, $g_K(z,t)=(4\pi Kt)^{-1/2}e^{-z^2/(4Kt)}$, $u=g_K*f$. Trên nửa trục, thác triển lẻ cho Dirichlet $0$, chẵn cho Neumann $0$; sau đổi $y\mapsto-y$ ở nửa âm:


\[
H_D(x,y,t)=g_K(x-y,t)-g_K(x+y,t),\qquad
H_N(x,y,t)=g_K(x-y,t)+g_K(x+y,t).
\]


Kiểm tra $H_D(0,y,t)=0$, $\partial_xH_N(0,y,t)=0$. Nửa trục phải có một điều kiện biên; dữ liệu đầu một mình không xác định duy nhất.

\textbf{Bước 2 — Tính các tích phân không gian.} Mode $\sin ky,\cos ky$ biến thành $e^{-Kk^2t}\sin kx,e^{-Kk^2t}\cos kx$. Với Gauss $e^{-\alpha(y-m)^2}$, $\alpha>0$, hoàn thành bình phương:


\[
-\frac{(x-y)^2}{4Kt}-\alpha(y-m)^2
=-\frac{1+4\alpha Kt}{4Kt}\left(y-\frac{x+4\alpha Ktm}{1+4\alpha Kt}\right)^2
-\frac{\alpha(x-m)^2}{1+4\alpha Kt}.
\]


Đổi biến để $\int_{\mathbb R}e^{-z^2}dz=\sqrt\pi$, suy $(g_K*f)(x)=e^{-\alpha(x-m)^2/(1+4\alpha Kt)}/\sqrt{1+4\alpha Kt}$. Với dữ liệu chỉ thị đoạn, cận hữu hạn tạo $\operatorname{erf}$; với $e^{-y}$ trên nửa trục, cận $0,\infty$ tạo $\operatorname{erfc}$, trong đó $\operatorname{erf}z=2\int_0^ze^{-r^2}dr/\sqrt\pi$, $\operatorname{erfc}z=1-\operatorname{erf}z$.

\textbf{Bước 3 — Thêm nguồn và biên.} Nguồn cho $u=g_K*f+\int_0^tg_K(\cdot,t-s)*F(\cdot,s)\,ds$, hoặc thay $g_K$ bằng $H_D,H_N$. Biên Dirichlet $\mu(t)$ trên nửa trục thêm


\[
\int_0^t\frac{x\,e^{-x^2/[4K(t-s)]}}{\sqrt{4\pi K}(t-s)^{3/2}}\mu(s)ds.
\]


Nguồn một mode giảm cùng tốc độ khuếch tán cho tích phân thời gian hằng sau khi gom mũ, nên xuất hiện $te^{-Kk^2t}$.

\textbf{Bước 4 — Đổi biến khi hệ số phụ thuộc thời gian hoặc có vận chuyển.} Với $u_t=a(t)u_{xx}+b(t)u_x+c(t)u$, đặt $A=\int_0^ta$, $B=\int_0^tb$, $C=\int_0^tc$, $u=e^Cv(x+B,A)$. Tính $u_t=e^C(cv+bv_z+av_\tau)$, $u_x=e^Cv_z$, $u_{xx}=e^Cv_{zz}$, suy $v_\tau=v_{zz}$. Dấu dịch là $x+B$, không phải $x-B$. Với hạng $xu_x$, dùng $z=e^{\pm t}x$ và chọn đồng hồ $\tau'(t)=e^{\pm2t}$.

\textbf{Bước 5 — Nhiều chiều và kiểm tra.} Nhân $n$ chiều là tích các nhân một chiều. Dữ liệu dạng tích tách được tích phân không gian bằng Fubini; nguồn dạng tích vẫn phải tích phân theo một thời gian chung $s$. Kiểm tra PDE bằng đạo hàm Duhamel: đạo hàm cận trên trả lại $F(x,t)$, phần còn lại bằng $K\Delta u$. Nhân Gauss là xấp xỉ đơn vị, nên trả lại dữ liệu tại điểm liên tục; dữ liệu nhảy không cho nghiệm liên tục đến mọi điểm lúc đầu.
\subsection*{Lời giải áp dụng cho bài}
\textbf{(a) Nhân nhiều chiều tách tích.}
Với $u_t=K\Delta u$ trên $\mathbb R^n$, nhân là
\[g_K^{(n)}(x-y,t)=\prod_{j=1}^n g_K(x_j-y_j,t).\]
Dữ liệu $f(y)=\prod f_j(y_j)$ cho, nhờ Fubini,
\[u(x,t)=\int_{\mathbb R^n}\prod_{j=1}^n[g_K(x_j-y_j,t)f_j(y_j)]dy
=\prod_{j=1}^nu_j(x_j,t).\]
Kiểm tra trực tiếp cũng được: $u_t=\sum_j(u_j)_t\prod_{i\ne j}u_i$, còn $K\Delta u=\sum_jK(u_j)_{x_jx_j}\prod_{i\ne j}u_i$, bằng nhau vì mỗi $u_j$ giải nhiệt một chiều. Tại $t=0$, tích trả lại $f$.
\textbf{(b) $n=2,K=1,f_1=\cos x_1,f_2=\sin x_2$.}
Mode một chiều cho $u_1=e^{-t}\cos x_1$, $u_2=e^{-t}\sin x_2$. Vì vậy
\[u=e^{-2t}\cos x_1\sin x_2.\]
$u_t=-2u$, $\Delta u=-u-u=-2u$, và mặt đầu đúng. Tốc độ giảm tổng là tổng hai trị riêng $1+1=2$.

\clearpage\noindent\texttt{\detokenize{BT_4_6_3_01}}\par
\section*{Lời giải bổ sung — Bài tập 1 · Mục 4.6.3}

\subsection*{Kiến thức và trình tự giải}
\textbf{Bước 1 — Chọn nhân nhiệt theo miền và hệ số.} Trên $\mathbb R$, $g_K(z,t)=(4\pi Kt)^{-1/2}e^{-z^2/(4Kt)}$, $u=g_K*f$. Trên nửa trục, thác triển lẻ cho Dirichlet $0$, chẵn cho Neumann $0$; sau đổi $y\mapsto-y$ ở nửa âm:


\[
H_D(x,y,t)=g_K(x-y,t)-g_K(x+y,t),\qquad
H_N(x,y,t)=g_K(x-y,t)+g_K(x+y,t).
\]


Kiểm tra $H_D(0,y,t)=0$, $\partial_xH_N(0,y,t)=0$. Nửa trục phải có một điều kiện biên; dữ liệu đầu một mình không xác định duy nhất.

\textbf{Bước 2 — Tính các tích phân không gian.} Mode $\sin ky,\cos ky$ biến thành $e^{-Kk^2t}\sin kx,e^{-Kk^2t}\cos kx$. Với Gauss $e^{-\alpha(y-m)^2}$, $\alpha>0$, hoàn thành bình phương:


\[
-\frac{(x-y)^2}{4Kt}-\alpha(y-m)^2
=-\frac{1+4\alpha Kt}{4Kt}\left(y-\frac{x+4\alpha Ktm}{1+4\alpha Kt}\right)^2
-\frac{\alpha(x-m)^2}{1+4\alpha Kt}.
\]


Đổi biến để $\int_{\mathbb R}e^{-z^2}dz=\sqrt\pi$, suy $(g_K*f)(x)=e^{-\alpha(x-m)^2/(1+4\alpha Kt)}/\sqrt{1+4\alpha Kt}$. Với dữ liệu chỉ thị đoạn, cận hữu hạn tạo $\operatorname{erf}$; với $e^{-y}$ trên nửa trục, cận $0,\infty$ tạo $\operatorname{erfc}$, trong đó $\operatorname{erf}z=2\int_0^ze^{-r^2}dr/\sqrt\pi$, $\operatorname{erfc}z=1-\operatorname{erf}z$.

\textbf{Bước 3 — Thêm nguồn và biên.} Nguồn cho $u=g_K*f+\int_0^tg_K(\cdot,t-s)*F(\cdot,s)\,ds$, hoặc thay $g_K$ bằng $H_D,H_N$. Biên Dirichlet $\mu(t)$ trên nửa trục thêm


\[
\int_0^t\frac{x\,e^{-x^2/[4K(t-s)]}}{\sqrt{4\pi K}(t-s)^{3/2}}\mu(s)ds.
\]


Nguồn một mode giảm cùng tốc độ khuếch tán cho tích phân thời gian hằng sau khi gom mũ, nên xuất hiện $te^{-Kk^2t}$.

\textbf{Bước 4 — Đổi biến khi hệ số phụ thuộc thời gian hoặc có vận chuyển.} Với $u_t=a(t)u_{xx}+b(t)u_x+c(t)u$, đặt $A=\int_0^ta$, $B=\int_0^tb$, $C=\int_0^tc$, $u=e^Cv(x+B,A)$. Tính $u_t=e^C(cv+bv_z+av_\tau)$, $u_x=e^Cv_z$, $u_{xx}=e^Cv_{zz}$, suy $v_\tau=v_{zz}$. Dấu dịch là $x+B$, không phải $x-B$. Với hạng $xu_x$, dùng $z=e^{\pm t}x$ và chọn đồng hồ $\tau'(t)=e^{\pm2t}$.

\textbf{Bước 5 — Nhiều chiều và kiểm tra.} Nhân $n$ chiều là tích các nhân một chiều. Dữ liệu dạng tích tách được tích phân không gian bằng Fubini; nguồn dạng tích vẫn phải tích phân theo một thời gian chung $s$. Kiểm tra PDE bằng đạo hàm Duhamel: đạo hàm cận trên trả lại $F(x,t)$, phần còn lại bằng $K\Delta u$. Nhân Gauss là xấp xỉ đơn vị, nên trả lại dữ liệu tại điểm liên tục; dữ liệu nhảy không cho nghiệm liên tục đến mọi điểm lúc đầu.
\subsection*{Lời giải áp dụng cho bài}
\textbf{Bước 1. Dùng tích chập Gauss.} $u=g_K*f$, $g_K=(4\pi Kt)^{-1/2}e^{-z^2/(4Kt)}$.
\textbf{Bước 2. Câu (i), $K=4,f=\sin^3x$.}
Từ $\sin3x=3\sin x-4\sin^3x$, $f=(3\sin x-\sin3x)/4$. Nhân nhiệt làm mode $k$ giảm với $Kk^2$, nên
\[u=\frac34e^{-4t}\sin x-\frac14e^{-36t}\sin3x.\]
Thế $t=0$ trả lại $f$; từng hạng có $u_t=4u_{xx}$.
\textbf{Bước 3. Câu (ii), $K=2,f=e^{-x^2+x}$.}
Hoàn thành bình phương ban đầu: $-y^2+y=1/4-(y-1/2)^2$. Trong tích chập,
\[-\frac{(x-y)^2}{8t}-(y-1/2)^2
=-\frac{1+8t}{8t}\left(y-\frac{x+4t}{1+8t}\right)^2
-\frac{(x-1/2)^2}{1+8t}.\]
Đổi biến $z=\sqrt{(1+8t)/(8t)}(y-(x+4t)/(1+8t))$ và tích phân Gauss cho
\[u=\frac{e^{1/4}}{\sqrt{1+8t}}
\exp\left[-\frac{(x-1/2)^2}{1+8t}\right].\]
Đặt $D=1+8t$, đạo hàm $u_t/u=-4/D+8(x-1/2)^2/D^2$, $u_{xx}/u=-2/D+4(x-1/2)^2/D^2$ kiểm tra $u_t=2u_{xx}$. $D\to1$ trả lại $e^{-x^2+x}$.

\clearpage\noindent\texttt{\detokenize{BT_4_6_3_02}}\par
\section*{Lời giải bổ sung — Bài tập 2 · Mục 4.6.3}

\subsection*{Kiến thức và trình tự giải}
\textbf{Bước 1 — Chọn nhân nhiệt theo miền và hệ số.} Trên $\mathbb R$, $g_K(z,t)=(4\pi Kt)^{-1/2}e^{-z^2/(4Kt)}$, $u=g_K*f$. Trên nửa trục, thác triển lẻ cho Dirichlet $0$, chẵn cho Neumann $0$; sau đổi $y\mapsto-y$ ở nửa âm:


\[
H_D(x,y,t)=g_K(x-y,t)-g_K(x+y,t),\qquad
H_N(x,y,t)=g_K(x-y,t)+g_K(x+y,t).
\]


Kiểm tra $H_D(0,y,t)=0$, $\partial_xH_N(0,y,t)=0$. Nửa trục phải có một điều kiện biên; dữ liệu đầu một mình không xác định duy nhất.

\textbf{Bước 2 — Tính các tích phân không gian.} Mode $\sin ky,\cos ky$ biến thành $e^{-Kk^2t}\sin kx,e^{-Kk^2t}\cos kx$. Với Gauss $e^{-\alpha(y-m)^2}$, $\alpha>0$, hoàn thành bình phương:


\[
-\frac{(x-y)^2}{4Kt}-\alpha(y-m)^2
=-\frac{1+4\alpha Kt}{4Kt}\left(y-\frac{x+4\alpha Ktm}{1+4\alpha Kt}\right)^2
-\frac{\alpha(x-m)^2}{1+4\alpha Kt}.
\]


Đổi biến để $\int_{\mathbb R}e^{-z^2}dz=\sqrt\pi$, suy $(g_K*f)(x)=e^{-\alpha(x-m)^2/(1+4\alpha Kt)}/\sqrt{1+4\alpha Kt}$. Với dữ liệu chỉ thị đoạn, cận hữu hạn tạo $\operatorname{erf}$; với $e^{-y}$ trên nửa trục, cận $0,\infty$ tạo $\operatorname{erfc}$, trong đó $\operatorname{erf}z=2\int_0^ze^{-r^2}dr/\sqrt\pi$, $\operatorname{erfc}z=1-\operatorname{erf}z$.

\textbf{Bước 3 — Thêm nguồn và biên.} Nguồn cho $u=g_K*f+\int_0^tg_K(\cdot,t-s)*F(\cdot,s)\,ds$, hoặc thay $g_K$ bằng $H_D,H_N$. Biên Dirichlet $\mu(t)$ trên nửa trục thêm


\[
\int_0^t\frac{x\,e^{-x^2/[4K(t-s)]}}{\sqrt{4\pi K}(t-s)^{3/2}}\mu(s)ds.
\]


Nguồn một mode giảm cùng tốc độ khuếch tán cho tích phân thời gian hằng sau khi gom mũ, nên xuất hiện $te^{-Kk^2t}$.

\textbf{Bước 4 — Đổi biến khi hệ số phụ thuộc thời gian hoặc có vận chuyển.} Với $u_t=a(t)u_{xx}+b(t)u_x+c(t)u$, đặt $A=\int_0^ta$, $B=\int_0^tb$, $C=\int_0^tc$, $u=e^Cv(x+B,A)$. Tính $u_t=e^C(cv+bv_z+av_\tau)$, $u_x=e^Cv_z$, $u_{xx}=e^Cv_{zz}$, suy $v_\tau=v_{zz}$. Dấu dịch là $x+B$, không phải $x-B$. Với hạng $xu_x$, dùng $z=e^{\pm t}x$ và chọn đồng hồ $\tau'(t)=e^{\pm2t}$.

\textbf{Bước 5 — Nhiều chiều và kiểm tra.} Nhân $n$ chiều là tích các nhân một chiều. Dữ liệu dạng tích tách được tích phân không gian bằng Fubini; nguồn dạng tích vẫn phải tích phân theo một thời gian chung $s$. Kiểm tra PDE bằng đạo hàm Duhamel: đạo hàm cận trên trả lại $F(x,t)$, phần còn lại bằng $K\Delta u$. Nhân Gauss là xấp xỉ đơn vị, nên trả lại dữ liệu tại điểm liên tục; dữ liệu nhảy không cho nghiệm liên tục đến mọi điểm lúc đầu.
\subsection*{Lời giải áp dụng cho bài}
\textbf{Bước 1. Thác triển chẵn đúng biên Neumann.}
Đặt $\widetilde f(y)=f(|y|)$. Gauss chẵn bảo đảm nghiệm toàn trục chẵn theo $x$, nên $u_x(0,t)=0$. Tách hai nửa của tích chập rồi đổi $y\mapsto-y$:
\[u=\int_0^\infty[g_K(x-y,t)+g_K(x+y,t)]f(y)dy.\]
\textbf{Bước 2. Tính câu $K=2,f=e^{-y}$.}
Hai số mũ là
\[-(x-y)^2/(8t)-y=-(y+4t-x)^2/(8t)+2t-x,\]
\[-(x+y)^2/(8t)-y=-(y+4t+x)^2/(8t)+2t+x.\]
Đổi biến $z=(y+4t\mp x)/\sqrt{8t}$, các cận thành $(4t\mp x)/\sqrt{8t},\infty$. Dùng $\int_a^\infty e^{-z^2}dz=\sqrt\pi\,\operatorname{erfc}(a)/2$ cho
\[u=\frac12e^{2t-x}\operatorname{erfc}\frac{4t-x}{\sqrt{8t}}
+\frac12e^{2t+x}\operatorname{erfc}\frac{4t+x}{\sqrt{8t}}.\]
\textbf{Bước 3. Kiểm tra.}
Hai hạng đổi chỗ khi $x\mapsto-x$, nên $u$ chẵn, $u_x(0,t)=0$. Với $x>0,t\downarrow0$, hai erfc lần lượt về $2,0$, nên $u\to e^{-x}$. Nhân ảnh cho $u_t=2u_{xx}$. $f'(0)=-1$ khác Neumann $0$, nên không trơn đạo hàm đến góc, dù giá trị nghiệm liên tục đến mặt đầu.
\begin{center}\includegraphics[width=.97\linewidth,height=11cm,keepaspectratio]{_xu_ly/graphs/BT_4_6_3_02.png}\end{center}
\textit{Đồ thị minh họa; với nghiệm số, xem phương pháp và tham số trong mã Python kèm theo.}

\clearpage\noindent\texttt{\detokenize{BT_4_6_3_03}}\par
\section*{Lời giải bổ sung — Bài tập 3 · Mục 4.6.3}

\subsection*{Kiến thức và trình tự giải}
\textbf{Bước 1 — Chọn nhân nhiệt theo miền và hệ số.} Trên $\mathbb R$, $g_K(z,t)=(4\pi Kt)^{-1/2}e^{-z^2/(4Kt)}$, $u=g_K*f$. Trên nửa trục, thác triển lẻ cho Dirichlet $0$, chẵn cho Neumann $0$; sau đổi $y\mapsto-y$ ở nửa âm:


\[
H_D(x,y,t)=g_K(x-y,t)-g_K(x+y,t),\qquad
H_N(x,y,t)=g_K(x-y,t)+g_K(x+y,t).
\]


Kiểm tra $H_D(0,y,t)=0$, $\partial_xH_N(0,y,t)=0$. Nửa trục phải có một điều kiện biên; dữ liệu đầu một mình không xác định duy nhất.

\textbf{Bước 2 — Tính các tích phân không gian.} Mode $\sin ky,\cos ky$ biến thành $e^{-Kk^2t}\sin kx,e^{-Kk^2t}\cos kx$. Với Gauss $e^{-\alpha(y-m)^2}$, $\alpha>0$, hoàn thành bình phương:


\[
-\frac{(x-y)^2}{4Kt}-\alpha(y-m)^2
=-\frac{1+4\alpha Kt}{4Kt}\left(y-\frac{x+4\alpha Ktm}{1+4\alpha Kt}\right)^2
-\frac{\alpha(x-m)^2}{1+4\alpha Kt}.
\]


Đổi biến để $\int_{\mathbb R}e^{-z^2}dz=\sqrt\pi$, suy $(g_K*f)(x)=e^{-\alpha(x-m)^2/(1+4\alpha Kt)}/\sqrt{1+4\alpha Kt}$. Với dữ liệu chỉ thị đoạn, cận hữu hạn tạo $\operatorname{erf}$; với $e^{-y}$ trên nửa trục, cận $0,\infty$ tạo $\operatorname{erfc}$, trong đó $\operatorname{erf}z=2\int_0^ze^{-r^2}dr/\sqrt\pi$, $\operatorname{erfc}z=1-\operatorname{erf}z$.

\textbf{Bước 3 — Thêm nguồn và biên.} Nguồn cho $u=g_K*f+\int_0^tg_K(\cdot,t-s)*F(\cdot,s)\,ds$, hoặc thay $g_K$ bằng $H_D,H_N$. Biên Dirichlet $\mu(t)$ trên nửa trục thêm


\[
\int_0^t\frac{x\,e^{-x^2/[4K(t-s)]}}{\sqrt{4\pi K}(t-s)^{3/2}}\mu(s)ds.
\]


Nguồn một mode giảm cùng tốc độ khuếch tán cho tích phân thời gian hằng sau khi gom mũ, nên xuất hiện $te^{-Kk^2t}$.

\textbf{Bước 4 — Đổi biến khi hệ số phụ thuộc thời gian hoặc có vận chuyển.} Với $u_t=a(t)u_{xx}+b(t)u_x+c(t)u$, đặt $A=\int_0^ta$, $B=\int_0^tb$, $C=\int_0^tc$, $u=e^Cv(x+B,A)$. Tính $u_t=e^C(cv+bv_z+av_\tau)$, $u_x=e^Cv_z$, $u_{xx}=e^Cv_{zz}$, suy $v_\tau=v_{zz}$. Dấu dịch là $x+B$, không phải $x-B$. Với hạng $xu_x$, dùng $z=e^{\pm t}x$ và chọn đồng hồ $\tau'(t)=e^{\pm2t}$.

\textbf{Bước 5 — Nhiều chiều và kiểm tra.} Nhân $n$ chiều là tích các nhân một chiều. Dữ liệu dạng tích tách được tích phân không gian bằng Fubini; nguồn dạng tích vẫn phải tích phân theo một thời gian chung $s$. Kiểm tra PDE bằng đạo hàm Duhamel: đạo hàm cận trên trả lại $F(x,t)$, phần còn lại bằng $K\Delta u$. Nhân Gauss là xấp xỉ đơn vị, nên trả lại dữ liệu tại điểm liên tục; dữ liệu nhảy không cho nghiệm liên tục đến mọi điểm lúc đầu.
\subsection*{Lời giải áp dụng cho bài}
\textbf{Bước 1. Tích phân ba hệ số.}
Đặt $A(t)=\int_0^ta(s)ds$, $B(t)=\int_0^tb(s)ds$, $C(t)=\int_0^tc(s)ds$. Đề giả thiết $a(t)\ge a_0>0$, nên $A(t)\ge a_0t>0$ khi $t>0$.
\textbf{Bước 2. Tính đầy đủ phép đổi biến.}
Đặt $z=x+B(t)$, $\tau=A(t)$ và $u=e^{C(t)}v(z,\tau)$. Quy tắc dây chuyền cho
\[u_t=e^C(cv+bv_z+av_\tau),\quad
u_x=e^Cv_z,\quad u_{xx}=e^Cv_{zz}.\]
Thay vào $u_t=au_{xx}+bu_x+cu$ rồi triệt tiêu $bv_z,cv$ được $a(v_\tau-v_{zz})=0$, suy $v_\tau=v_{zz}$. Tại $t=0$, $A=B=C=0$, nên $v(z,0)=f(z)$.
\textbf{Bước 3. Fourier và biến đổi ngược.}
$\widehat v_\tau=-\xi^2\widehat v$ cho $\widehat v=e^{-\tau\xi^2}\widehat f$. Biến đổi ngược thành tích chập Gauss:
\[u(x,t)=\frac{e^{C(t)}}{\sqrt{4\pi A(t)}}\int_{\mathbb R}
\exp\left[-\frac{(x+B(t)-y)^2}{4A(t)}\right]f(y)dy.\]
\textbf{Bước 4. Kiểm tra.}
Phép tính bước 2 kiểm tra PDE; xấp xỉ đơn vị của Gauss và $B,C\to0$ trả lại dữ liệu đầu tại điểm liên tục. Dấu dịch $x+B$ đúng vì khi $a=c=0$, nghiệm của $u_t=b(t)u_x$ là $f(x+B(t))$. Do nhân dương, $|u|\le e^{C(t)}\|f\|_\infty$ trên mỗi dải thời gian hữu hạn.

\clearpage\noindent\texttt{\detokenize{BT_4_6_3_04}}\par
\section*{Lời giải bổ sung — Bài tập 4 · Mục 4.6.3}

\subsection*{Kiến thức và trình tự giải}
\textbf{Bước 1 — Phân biệt đạo hàm cấp hai, ba và bốn.} Quy ước $\widehat f(\xi)=\int e^{-i\xi x}f(x)dx$, biến đổi ngược có hệ số $1/(2\pi)$. Khi đó $\widehat{f^{(m)}}=(i\xi)^m\widehat f$, $\widehat{xf}=i\partial_\xi\widehat f$, $\widehat{xf_x}=-\widehat f-\xi\partial_\xi\widehat f$.

\textbf{Bước 2 — Giải phương trình của biến đổi.} $u_t=u_{xx}$ cho $\widehat u=e^{-t\xi^2}\widehat f$. $u_t=u_{xxx}$ cho $\widehat u=e^{-it\xi^3}\widehat f$. $u_t=-u_{xxxx}$ cho $\widehat u=e^{-t\xi^4}\widehat f$. Với $xu_x$, giải PDE cấp một theo $\xi$ bằng đường đặc trưng, gồm cả hệ số biên độ $-\widehat u$.

\textbf{Bước 3 — Biến đổi ngược và ý nghĩa tích phân.} Nhân bậc ba là $A_t(x)=(3t)^{-1/3}\operatorname{Ai}(-x/(3t)^{1/3})$, tích phân dao động. Nhân bậc bốn là $B_t(x)=(2\pi)^{-1}\int e^{i\xi x-t\xi^4}d\xi$, hội tụ tuyệt đối theo $\xi$ và có thể đổi dấu. Không áp dụng tính dương của Gauss hay nguyên lý cực đại cấp hai cho hai nhân này.

\textbf{Bước 4 — Miền hữu hạn cần kiểm tra biên và toán tử liên hợp.} Với $u_t=Ku_{xxx}$, $u(0)=u_x(0)=u(L)=0$, năng lượng cho $I'=-K u_x(L)^2/2$. Toán tử liên hợp là $-KD^3$ với $Y(0)=Y(L)=Y'(L)=0$. Đây không phải bài Sturm–Liouville tự liên hợp cấp hai; không suy mọi trị riêng đều thực và không tự dùng hệ số $\int fX/\int X^2$. Có thể giải bằng biến đổi Laplace theo thời gian và nghiệm cơ bản của ODE bậc ba, tránh khẳng định cơ sở chưa chứng minh.

\textbf{Bước 5 — Nêu lớp dữ liệu đủ cho công thức.} Với bậc ba trên toàn trục, công thức Fourier hợp lệ trước hết cho dữ liệu Schwartz hoặc $L^2$ theo nghĩa biến đổi Fourier. Dữ liệu liên tục bị chặn là phân phối tempered; phép nhân $e^{-it\xi^3}\widehat f$ vẫn hợp lệ vì các đạo hàm của thừa số tăng đa thức, nên xác định nghiệm trong $\mathcal S\prime$. Tuy nhiên điều đó không bảo đảm tích phân Airy hội tụ tuyệt đối hay nghiệm tự động là hàm cổ điển. Ghi đúng ý nghĩa của công thức theo lớp dữ liệu.
\subsection*{Lời giải áp dụng cho bài}
\textbf{Bước 1. Quy ước Fourier và các công thức dùng chung.}
$\widehat f(\xi)=\int e^{-i\xi x}f(x)dx$, biến đổi ngược $(2\pi)^{-1}\int e^{i\xi x}\widehat f(\xi)d\xi$. Tích phân từng phần cho $\widehat{u_{xx}}=-\xi^2\widehat u$, $\widehat{u_{xxx}}=-i\xi^3\widehat u$, $\widehat{u_{xxxx}}=\xi^4\widehat u$. Đạo hàm theo $\xi$ cho $\widehat{xu_x}=-\widehat u-\xi\widehat u_\xi$. Ký hiệu $g_K(z,t)=(4\pi Kt)^{-1/2}e^{-z^2/(4Kt)}$.

\textbf{(i) $u_t=u_{xx}+xu_x$.}
Đặt $W=\widehat u$: $W_t+\xi W_\xi=-(1+\xi^2)W$. Đặc trưng $\xi(s)=\xi_0e^s$ cho
\[W(\xi_0e^t,t)=e^{-t}\exp[-\xi_0^2(e^{2t}-1)/2]\widehat f(\xi_0).\]
Vì $\xi_0=e^{-t}\xi$,
\[W=e^{-t}e^{-\xi^2(1-e^{-2t})/2}\widehat f(e^{-t}\xi).\]
Đổi biến $\eta=e^{-t}\xi$ trong biến đổi ngược, đặt $z=e^tx,\tau=(e^{2t}-1)/2$, được
\[u(x,t)=(g_1(\cdot,\tau)*f)(e^tx).\]
Kiểm tra: $u_t=zv_z+e^{2t}v_\tau$, $u_{xx}=e^{2t}v_{zz}$, $xu_x=zv_z$, nên $v_\tau=v_{zz}$.

\textbf{(ii) $u_t=u_{xxx}$.}
$W_t=-i\xi^3W$ cho $W=e^{-it\xi^3}\widehat f$. Do đó
\[u=\mathcal F^{-1}(e^{-it\xi^3}\widehat f)=A_t*f,\qquad
A_t(x)=(3t)^{-1/3}\operatorname{Ai}(-x/(3t)^{1/3}).\]
Để thấy dấu Airy, đổi $\eta=(3t)^{1/3}\xi$ trong $(2\pi)^{-1}\int e^{i\xi x-it\xi^3}d\xi$ rồi đổi $\eta\mapsto-\eta$, được tích phân Airy có pha $\eta^3/3-\eta x/(3t)^{1/3}$. Với dữ liệu Schwartz, công thức cho nghiệm cổ điển; với $L^2$ là nghiệm Fourier trong $L^2$. Với dữ liệu chỉ liên tục bị chặn như đề, $f$ là phân phối tempered và phép nhân $e^{-it\xi^3}\widehat f$ hợp lệ vì các đạo hàm của thừa số có tăng trưởng đa thức; kết quả được hiểu trong $\mathcal S'$. Không khẳng định tích phân $A_t*f$ hội tụ tuyệt đối hoặc nghiệm tự động là hàm cổ điển.

\textbf{(iii) $u_t=-u_{xxxx}+xu_x$.}
$W_t+\xi W_\xi=-(1+\xi^4)W$. Cùng đặc trưng $\xi=\xi_0e^t$ cho
\[W=e^{-t}e^{-\xi^4(1-e^{-4t})/4}\widehat f(e^{-t}\xi).\]
Đặt $z=e^tx$, $\tau=(e^{4t}-1)/4$ và
\[B_\tau(z)=\frac1{2\pi}\int_{\mathbb R}e^{i\xi z-\tau\xi^4}d\xi.\]
Đổi biến Fourier ngược cho $u=(B_\tau*f)(e^tx)$. Kiểm tra $u_t=zv_z+e^{4t}v_\tau$, $u_{xxxx}=e^{4t}v_{zzzz}$, nên $v_\tau=-v_{zzzz}$. $B_\tau$ là hàm Schwartz với $\tau>0$, nhưng không phải nhân Gauss dương.

\textbf{(iv) $u_t=u_{xx}-xu_x$.}
$W_t-\xi W_\xi=(1-\xi^2)W$. Đặc trưng $\xi=\xi_0e^{-t}$ cho $W=e^t e^{-\xi^2(e^{2t}-1)/2}\widehat f(e^t\xi)$. Đổi biến trong Fourier ngược, đặt $z=e^{-t}x$, $\tau=(1-e^{-2t})/2$:
\[u(x,t)=(g_1(\cdot,\tau)*f)(e^{-t}x).\]
Kiểm tra $u_t=-zv_z+e^{-2t}v_\tau$, $u_{xx}=e^{-2t}v_{zz}$, $-xu_x=-zv_z$, nên đúng PDE. Ở (i),(iii),(iv), khi $t\downarrow0$, $z\to x,\tau\to0$, các nhân trả lại $f$.

\textbf{(v) $u_t=u_{xx}+tx$ trên nửa trục.}
Bản in lặp lại $u(x,0)=0$ hai lần, không cho biên tại $x=0$. Một nghiệm riêng là $u_p=xt^2/2$: $(u_p)_t=tx$, $(u_p)_{xx}=0$, $u_p(x,0)=0$. Nếu bổ sung $u(0,t)=0$ và lớp tăng trưởng phù hợp, đây là nghiệm của bài bổ sung. Với bản in, còn họ nghiệm
\[u=\frac{xt^2}2+\int_0^t
\frac{x e^{-x^2/[4(t-s)]}}{2\sqrt\pi(t-s)^{3/2}}\mu(s)ds,\]
với $\mu$ trơn tùy ý vanishing gần $0$. Hạng thêm giải nhiệt, dữ liệu đầu $0$ và có biên $\mu$, nên chứng minh không duy nhất khi thiếu biên. Dòng rời về $f=\chi_{(0,1)}$ không thống nhất với dữ liệu đầu $0$ đã in; không tự đưa nó vào nghiệm của hệ này.

\textbf{(vi) Dirichlet nửa trục, hệ số $k>0$.}
Thác triển lẻ $f(y)=1/(1+y^2)$ cho
\[u=\int_0^\infty
\frac{e^{-(x-y)^2/(4kt)}-e^{-(x+y)^2/(4kt)}}{\sqrt{4\pi kt}}
\frac{dy}{1+y^2}.\]
Hệ số $k$ phải có trong cả mẫu và số mũ. Hai ảnh triệt tiêu tại $x=0$; Gauss xấp xỉ đơn vị trả lại $f$ khi $x>0,t\downarrow0$; đạo hàm nhân cho $u_t=ku_{xx}$. Giá trị đầu $f(0)=1$ không khớp biên $0$ tại góc.

\textbf{(vii),(viii) Nguồn trên nửa trục nhưng thiếu biên.}
Đặt $H(x,y,t)=g_1(x-y,t)-g_1(x+y,t)$ và $P(x,t)=xe^{-x^2/(4t)}/(2\sqrt\pi\,t^{3/2})$. Với một biên Dirichlet được \emph{bổ sung} là $\mu(t)$, công thức Duhamel cho
\[u=\int_0^\infty H(x,y,t)f(y)dy+
\int_0^t\!\int_0^\infty H(x,y,t-s)F(y)dy\,ds+
\int_0^tP(x,t-s)\mu(s)ds.\]
Ở (vii), $f(y)=1/(1+y^2)$, $F(y)=e^{-y}$; ở (viii), $f=0$, $F(y)=e^{-y}/(1+y^2)$. Để kiểm tra hạng biên, đổi $r=t-s$:
\[\int_0^tP(x,r)dr=\operatorname{erfc}\frac{x}{2\sqrt t}\longrightarrow1
\quad(x\downarrow0).\]
Khối lượng của $P(x,r)$ tập trung tại $r=0$ khi $x\downarrow0$, nên hạng này tiến về $\mu(t)$ nếu $\mu$ liên tục. Nó giải nhiệt và có mặt đầu $0$ với $x>0$. Chọn hai hàm $\mu$ khác nhau vanishing gần $0$ cho hai nghiệm cùng PDE và dữ liệu đầu, khác biên; do đó đề gốc không xác định duy nhất. Các tích phân không gian với $1/(1+y^2)$ là biểu diễn chính xác; chúng không phải tích phân sơ cấp cần giả tạo một nguyên hàm.

\textbf{(ix) Neumann, dữ liệu $e^{-x}$, $K=1$.}
Nhân ảnh chẵn cho $u=\int_0^\infty[g_1(x-y,t)+g_1(x+y,t)]e^{-y}dy$. Hoàn thành bình phương:
\[-(x\mp y)^2/(4t)-y=-(y+2t\mp x)^2/(4t)+t\mp x.\]
Đổi biến $z=(y+2t\mp x)/\sqrt{4t}$ rồi dùng erfc:
\[u=\frac12e^{t-x}\operatorname{erfc}\frac{2t-x}{\sqrt{4t}}
+\frac12e^{t+x}\operatorname{erfc}\frac{2t+x}{\sqrt{4t}}.\]
Hàm chẵn theo $x$ cho $u_x(0,t)=0$, hai giới hạn erfc cho $u(x,0)=e^{-x}$ với $x>0$; nhân ảnh kiểm tra PDE.

\textbf{(x) Dirichlet bằng $1$.}
Đặt $w=u-1$. Khi đó $w_t=w_{xx}$, $w(0,t)=0$, $w(x,0)=1/(1+x^2)-1$. Thác triển lẻ cho
\[u=1+\int_0^\infty H(x,y,t)\left(\frac1{1+y^2}-1\right)dy.\]
Tại $x=0$ nhân $H$ bằng $0$, nên $u=1$. Mặt đầu trả lại $f$, và phép cộng hằng không đổi PDE. Dữ liệu góc về giá trị khớp vì $f(0)=1$.

\clearpage\noindent\texttt{\detokenize{BT_4_6_3_05}}\par
\section*{Lời giải bổ sung — Bài tập 5 · Mục 4.6.3}

\subsection*{Kiến thức và trình tự giải}
\textbf{Bước 1 — Chọn nhân nhiệt theo miền và hệ số.} Trên $\mathbb R$, $g_K(z,t)=(4\pi Kt)^{-1/2}e^{-z^2/(4Kt)}$, $u=g_K*f$. Trên nửa trục, thác triển lẻ cho Dirichlet $0$, chẵn cho Neumann $0$; sau đổi $y\mapsto-y$ ở nửa âm:


\[
H_D(x,y,t)=g_K(x-y,t)-g_K(x+y,t),\qquad
H_N(x,y,t)=g_K(x-y,t)+g_K(x+y,t).
\]


Kiểm tra $H_D(0,y,t)=0$, $\partial_xH_N(0,y,t)=0$. Nửa trục phải có một điều kiện biên; dữ liệu đầu một mình không xác định duy nhất.

\textbf{Bước 2 — Tính các tích phân không gian.} Mode $\sin ky,\cos ky$ biến thành $e^{-Kk^2t}\sin kx,e^{-Kk^2t}\cos kx$. Với Gauss $e^{-\alpha(y-m)^2}$, $\alpha>0$, hoàn thành bình phương:


\[
-\frac{(x-y)^2}{4Kt}-\alpha(y-m)^2
=-\frac{1+4\alpha Kt}{4Kt}\left(y-\frac{x+4\alpha Ktm}{1+4\alpha Kt}\right)^2
-\frac{\alpha(x-m)^2}{1+4\alpha Kt}.
\]


Đổi biến để $\int_{\mathbb R}e^{-z^2}dz=\sqrt\pi$, suy $(g_K*f)(x)=e^{-\alpha(x-m)^2/(1+4\alpha Kt)}/\sqrt{1+4\alpha Kt}$. Với dữ liệu chỉ thị đoạn, cận hữu hạn tạo $\operatorname{erf}$; với $e^{-y}$ trên nửa trục, cận $0,\infty$ tạo $\operatorname{erfc}$, trong đó $\operatorname{erf}z=2\int_0^ze^{-r^2}dr/\sqrt\pi$, $\operatorname{erfc}z=1-\operatorname{erf}z$.

\textbf{Bước 3 — Thêm nguồn và biên.} Nguồn cho $u=g_K*f+\int_0^tg_K(\cdot,t-s)*F(\cdot,s)\,ds$, hoặc thay $g_K$ bằng $H_D,H_N$. Biên Dirichlet $\mu(t)$ trên nửa trục thêm


\[
\int_0^t\frac{x\,e^{-x^2/[4K(t-s)]}}{\sqrt{4\pi K}(t-s)^{3/2}}\mu(s)ds.
\]


Nguồn một mode giảm cùng tốc độ khuếch tán cho tích phân thời gian hằng sau khi gom mũ, nên xuất hiện $te^{-Kk^2t}$.

\textbf{Bước 4 — Đổi biến khi hệ số phụ thuộc thời gian hoặc có vận chuyển.} Với $u_t=a(t)u_{xx}+b(t)u_x+c(t)u$, đặt $A=\int_0^ta$, $B=\int_0^tb$, $C=\int_0^tc$, $u=e^Cv(x+B,A)$. Tính $u_t=e^C(cv+bv_z+av_\tau)$, $u_x=e^Cv_z$, $u_{xx}=e^Cv_{zz}$, suy $v_\tau=v_{zz}$. Dấu dịch là $x+B$, không phải $x-B$. Với hạng $xu_x$, dùng $z=e^{\pm t}x$ và chọn đồng hồ $\tau'(t)=e^{\pm2t}$.

\textbf{Bước 5 — Nhiều chiều và kiểm tra.} Nhân $n$ chiều là tích các nhân một chiều. Dữ liệu dạng tích tách được tích phân không gian bằng Fubini; nguồn dạng tích vẫn phải tích phân theo một thời gian chung $s$. Kiểm tra PDE bằng đạo hàm Duhamel: đạo hàm cận trên trả lại $F(x,t)$, phần còn lại bằng $K\Delta u$. Nhân Gauss là xấp xỉ đơn vị, nên trả lại dữ liệu tại điểm liên tục; dữ liệu nhảy không cho nghiệm liên tục đến mọi điểm lúc đầu.
\subsection*{Lời giải áp dụng cho bài}
\textbf{Bước 1. Tách dữ liệu đầu và nguồn.}
Với $K=2$, $f=\chi_{[-2,2]}$, $F=e^{-2t}\sin x$, Duhamel cho $u=u_0+u_F$:
\[u_0=\frac1{\sqrt{8\pi t}}\int_{-2}^2e^{-(x-y)^2/(8t)}dy,\qquad
u_F=\int_0^te^{-2(t-s)}e^{-2s}\sin x\,ds.\]
\textbf{Bước 2. Thực hiện từng tích phân.}
Đặt $z=(y-x)/\sqrt{8t}$. Các cận là $(-2-x)/\sqrt{8t}$ và $(2-x)/\sqrt{8t}$; dùng erf lẻ:
\[u_0=\frac12\left[\operatorname{erf}\frac{x+2}{\sqrt{8t}}
-\operatorname{erf}\frac{x-2}{\sqrt{8t}}\right].\]
Trong $u_F$, hai mũ gom thành $e^{-2t}$ độc lập $s$, nên $u_F=te^{-2t}\sin x$. Vậy
\[u=\frac12\left[\operatorname{erf}\frac{x+2}{\sqrt{8t}}
-\operatorname{erf}\frac{x-2}{\sqrt{8t}}\right]+te^{-2t}\sin x.\]
\textbf{Bước 3. Kiểm tra.}
$u_0$ giải nhiệt không nguồn. $(u_F)_t-2(u_F)_{xx}=(1-2t)e^{-2t}\sin x+2te^{-2t}\sin x=e^{-2t}\sin x$. Khi $t\downarrow0$, $u_F\to0$, $u_0\to1$ nếu $|x|<2$, về $0$ nếu $|x|>2$, về $1/2$ tại $x=\pm2$. Giá trị tại hai điểm nhảy không bằng giá trị $1$ của chỉ thị đoạn đóng; điều kiện đầu phải hiểu theo $L^p$ hoặc ở điểm liên tục, không có nghiệm liên tục đến các điểm nhảy đó.

\clearpage\noindent\texttt{\detokenize{BT_4_6_3_06}}\par
\section*{Lời giải bổ sung — Bài tập 6 · Mục 4.6.3}

\subsection*{Kiến thức và trình tự giải}
\textbf{Bước 1 — Chọn nhân nhiệt theo miền và hệ số.} Trên $\mathbb R$, $g_K(z,t)=(4\pi Kt)^{-1/2}e^{-z^2/(4Kt)}$, $u=g_K*f$. Trên nửa trục, thác triển lẻ cho Dirichlet $0$, chẵn cho Neumann $0$; sau đổi $y\mapsto-y$ ở nửa âm:


\[
H_D(x,y,t)=g_K(x-y,t)-g_K(x+y,t),\qquad
H_N(x,y,t)=g_K(x-y,t)+g_K(x+y,t).
\]


Kiểm tra $H_D(0,y,t)=0$, $\partial_xH_N(0,y,t)=0$. Nửa trục phải có một điều kiện biên; dữ liệu đầu một mình không xác định duy nhất.

\textbf{Bước 2 — Tính các tích phân không gian.} Mode $\sin ky,\cos ky$ biến thành $e^{-Kk^2t}\sin kx,e^{-Kk^2t}\cos kx$. Với Gauss $e^{-\alpha(y-m)^2}$, $\alpha>0$, hoàn thành bình phương:


\[
-\frac{(x-y)^2}{4Kt}-\alpha(y-m)^2
=-\frac{1+4\alpha Kt}{4Kt}\left(y-\frac{x+4\alpha Ktm}{1+4\alpha Kt}\right)^2
-\frac{\alpha(x-m)^2}{1+4\alpha Kt}.
\]


Đổi biến để $\int_{\mathbb R}e^{-z^2}dz=\sqrt\pi$, suy $(g_K*f)(x)=e^{-\alpha(x-m)^2/(1+4\alpha Kt)}/\sqrt{1+4\alpha Kt}$. Với dữ liệu chỉ thị đoạn, cận hữu hạn tạo $\operatorname{erf}$; với $e^{-y}$ trên nửa trục, cận $0,\infty$ tạo $\operatorname{erfc}$, trong đó $\operatorname{erf}z=2\int_0^ze^{-r^2}dr/\sqrt\pi$, $\operatorname{erfc}z=1-\operatorname{erf}z$.

\textbf{Bước 3 — Thêm nguồn và biên.} Nguồn cho $u=g_K*f+\int_0^tg_K(\cdot,t-s)*F(\cdot,s)\,ds$, hoặc thay $g_K$ bằng $H_D,H_N$. Biên Dirichlet $\mu(t)$ trên nửa trục thêm


\[
\int_0^t\frac{x\,e^{-x^2/[4K(t-s)]}}{\sqrt{4\pi K}(t-s)^{3/2}}\mu(s)ds.
\]


Nguồn một mode giảm cùng tốc độ khuếch tán cho tích phân thời gian hằng sau khi gom mũ, nên xuất hiện $te^{-Kk^2t}$.

\textbf{Bước 4 — Đổi biến khi hệ số phụ thuộc thời gian hoặc có vận chuyển.} Với $u_t=a(t)u_{xx}+b(t)u_x+c(t)u$, đặt $A=\int_0^ta$, $B=\int_0^tb$, $C=\int_0^tc$, $u=e^Cv(x+B,A)$. Tính $u_t=e^C(cv+bv_z+av_\tau)$, $u_x=e^Cv_z$, $u_{xx}=e^Cv_{zz}$, suy $v_\tau=v_{zz}$. Dấu dịch là $x+B$, không phải $x-B$. Với hạng $xu_x$, dùng $z=e^{\pm t}x$ và chọn đồng hồ $\tau'(t)=e^{\pm2t}$.

\textbf{Bước 5 — Nhiều chiều và kiểm tra.} Nhân $n$ chiều là tích các nhân một chiều. Dữ liệu dạng tích tách được tích phân không gian bằng Fubini; nguồn dạng tích vẫn phải tích phân theo một thời gian chung $s$. Kiểm tra PDE bằng đạo hàm Duhamel: đạo hàm cận trên trả lại $F(x,t)$, phần còn lại bằng $K\Delta u$. Nhân Gauss là xấp xỉ đơn vị, nên trả lại dữ liệu tại điểm liên tục; dữ liệu nhảy không cho nghiệm liên tục đến mọi điểm lúc đầu.
\subsection*{Lời giải áp dụng cho bài}
\textbf{(a) Duhamel nhiều chiều và tách tích phân không gian.}
Với $f=0,F(x,t)=g(t)\prod_{j=1}^nF_j(x_j)$, nhân nhiệt $n$ chiều là $\prod_jg_K(x_j-y_j,t-s)$. Vì vậy
\begin{align*}
u(x,t)&=\int_0^t\!\int_{\mathbb R^n}
g(s)\prod_{j=1}^n[g_K(x_j-y_j,t-s)F_j(y_j)]dy\,ds\\
&=\int_0^tg(s)\prod_{j=1}^n
\left(\int_{\mathbb R}g_K(x_j-y_j,t-s)F_j(y_j)dy_j\right)ds.
\end{align*}
Fubini tách các tích phân không gian; không tách tích phân thời gian thành tích của $n$ tích phân, vì tất cả thừa số dùng cùng thời gian $s$.
\textbf{(b) $n=2,K=1,g(s)=\sin s,F_1=\cos x,F_2=e^{-2y^2}$.}
Đặt $\tau=t-s$. Mode cosine cho $g_1(\cdot,\tau)*F_1=e^{-\tau}\cos x$. Với Gauss, hoàn thành bình phương
\[-(y-z)^2/(4\tau)-2z^2
=-\frac{1+8\tau}{4\tau}\left(z-\frac y{1+8\tau}\right)^2
-\frac{2y^2}{1+8\tau},\]
nên $g_1*F_2=e^{-2y^2/(1+8\tau)}/\sqrt{1+8\tau}$. Do đó
\[u(x,y,t)=\cos x\int_0^t
\frac{\sin s\,e^{-(t-s)}}{\sqrt{1+8(t-s)}}
\exp\left[-\frac{2y^2}{1+8(t-s)}\right]ds.\]
Tích phân thời gian hữu hạn trên hàm đã biết là biểu diễn tường minh; không có nguyên hàm sơ cấp thông thường cần tính tiếp.
\textbf{Kiểm tra bằng đạo hàm Duhamel.}
Gọi $V(x,y,\tau)=e^{-\tau}\cos x\,e^{-2y^2/(1+8\tau)}/\sqrt{1+8\tau}$, thì $V_\tau=\Delta V$, $V(x,y,0)=\cos x\,e^{-2y^2}$. Do đó
\[u_t=\sin t\,V(x,y,0)+\int_0^t\sin s\,V_\tau(x,y,t-s)ds
=\sin t\cos x\,e^{-2y^2}+\Delta u.\]
Tại $t=0$ tích phân bằng $0$, đúng mặt đầu.

\clearpage\noindent\texttt{\detokenize{BT_5_1_3_01}}\par
\section*{Lời giải bổ sung — Bài tập 1 · Mục 5.1.3}
\textit{Áp dụng kiến thức và phương pháp trong tài liệu gốc.}\par
Đặt $S_{D,n}(q)=2D^{-1}\int_0^Dq(s)\sin(n\pi s/D)ds$, $k_n=n\pi/L,\ell_n=n\pi/M$. Chồng chất bốn bài có một cạnh khác 0 cho công thức đã tính hệ số:
\begin{align*}
u={}&\sum_{n\ge1}\sin k_nx\frac{S_{L,n}(\varphi_0)\sinh k_n(M-y)+S_{L,n}(\varphi_1)\sinh k_ny}{\sinh k_nM}\\
&+\sum_{n\ge1}\sin\ell_ny\frac{S_{M,n}(\psi_0)\sinh\ell_n(L-x)+S_{M,n}(\psi_1)\sinh\ell_nx}{\sinh\ell_nL}.
\end{align*}
Để áp dụng từng câu, các hệ số không bằng 0 là:

(i) $L=1,M=\pi$, $S(\varphi_1)=(3\delta_{n1}-\delta_{n3})/4$, $S(\psi_1)=2(1-(-1)^n)/(\pi n)$.
(ii) $L=M=1$, $S(\varphi_0)=S(\psi_0)=S(\psi_1)=2(1-(-1)^n)/(\pi n)$; cũng có thể viết $u=1-$chuỗi một cạnh trên với dữ liệu 1.
(iii) $u=\sin x\sinh y/\sinh\pi$.
(iv) $S(\varphi_1)=4(1-(-1)^n)/(n^3\pi^3)$.
(v) $S(\varphi_1)=4\sin(n\pi/2)/(n^2\pi^2)$.
(vi) $L=2,M=1$, $S(\varphi_1)=2[1+(-1)^n-2\cos(n\pi/2)]/(n\pi)$.
(vii) $u=\sin2\pi x\sinh(2\pi(1-y))/\sinh2\pi$.
(viii) $L=1,M=2$, $S(\varphi_0)=8/[\pi n(4-n^2)]$ nếu $n$ lẻ, bằng 0 nếu $n$ chẵn.
(ix) $S(\varphi_0)=\delta_{n2}$, $S(\varphi_1)=-2\cos(n\pi/2)/(n\pi)$.

Thay các hệ số trên vào hai chuỗi là nghiệm của từng trường hợp. Đạo hàm trong miền cho $\Delta u=0$; trên cạnh hội tụ về dữ liệu tại điểm liên tục, và về trung bình ở điểm gián đoạn. Nếu dữ liệu ở góc không khớp, không có nghiệm liên tục trên cả hình chữ nhật đóng; chuỗi vẫn giải bài trong miền và trên các cạnh mở.

\clearpage\noindent\texttt{\detokenize{BT_5_1_3_02}}\par
\section*{Lời giải bổ sung — Bài tập 2 · Mục 5.1.3}
\textit{Áp dụng kiến thức và phương pháp trong tài liệu gốc.}\par
Điều kiện (5.1.14) là $u_x(0,y)=\psi_0(y),u_x(L,y)=\psi_1(y)$, còn biên dưới/trên là Dirichlet. Đặt hệ số cosine $c_0(q)=L^{-1}\int q$, $c_n(q)=2L^{-1}\int q\cos(n\pi x/L)$; hệ số sine $s_n(q)=2M^{-1}\int q\sin(n\pi y/M)$.
\begin{align*}
u={}&c_0(\varphi_0)(1-y/M)+c_0(\varphi_1)y/M\\
&+\sum_{n\ge1}\cos k_nx\frac{c_n(\varphi_0)\sinh k_n(M-y)+c_n(\varphi_1)\sinh k_ny}{\sinh k_nM}\\
&+\sum_{n\ge1}\sin\ell_ny\frac{s_n(\psi_1)\cosh\ell_nx-s_n(\psi_0)\cosh\ell_n(L-x)}{\ell_n\sinh\ell_nL}.
\end{align*}
(i) $c_0(\varphi_1)=4/(3\pi)$, $c_n(\varphi_1)=24/[\pi(1-n^2)(9-n^2)]$ nếu $n$ chẵn, bằng 0 nếu $n$ lẻ; $s_n(\psi_1)=2(1-(-1)^n)/(\pi n)$. Các hệ số còn lại 0.
(ii) $c_0(\varphi_0)=1$, các cosine khác 0; $s_n(\psi_0)=s_n(\psi_1)=2(1-(-1)^n)/(\pi n)$. Thế vào công thức.
(iii) $u=\cos x\sinh y/\sinh\pi$. Các công thức thỏa đạo hàm biên trực tiếp; năng lượng cho tính duy nhất.

\clearpage\noindent\texttt{\detokenize{BT_5_1_3_03}}\par
\section*{Lời giải bổ sung — Bài tập 3 · Mục 5.1.3}
\textit{Áp dụng kiến thức và phương pháp trong tài liệu gốc.}\par
Với hiệu $w$ của hai nghiệm, $\int_\Omega|\nabla w|^2=\int_{\partial\Omega}w\partial_\nu w=0$, nên $w$ hằng; vì vậy nếu có một nghiệm thì mọi nghiệm là $u+C$.

Điều kiện tương thích (dùng đạo hàm tọa độ trong đề) là
\[\int_0^L(\varphi_1-\varphi_0)dx+\int_0^M(\psi_1-\psi_0)dy=0.\]
Đủ tồn tại nghiệm yếu Neumann, duy nhất sau khi ấn định trung bình; để có nghiệm cổ điển đến góc còn cần độ trơn và tương thích góc.

(iii) Dữ liệu cho $\int(\varphi_1-\varphi_0)=-1/2$, hai cạnh kia 0. Do đó \textbf{không có nghiệm}; không được triển khai Fourier bỏ qua mode 0.

\clearpage\noindent\texttt{\detokenize{BT_5_1_3_04}}\par
\section*{Lời giải bổ sung — Bài tập 4 · Mục 5.1.3}
\textit{Áp dụng kiến thức và phương pháp trong tài liệu gốc.}\par
Biên dưới $u_y=u$, trên $u_y=-u$ tương đương $\partial_\nu u=-u$ ở cả hai cạnh; hai cạnh đứng Neumann. Với hiệu $w$,
\[\int_\Omega|\nabla w|^2=-\int_{y=0,M}w^2\le0.\]
Do đó $w=0$. Kết luận ``vô số nghiệm sai khác hằng số'' của đề là sai: cộng hằng khác 0 phá điều kiện Robin.

Câu (ii), $L=M=1,u_x(0,y)=1,u_x(1,y)=0$. Gọi $Y_j=\cos k_jy+\sin k_jy/k_j$, với $(1-k_j^2)\sin k_j+2k_j\cos k_j=0$, $k_j>0$, $N_j=\int_0^1Y_j^2$, $d_j=\int_0^1Y_j/N_j$. Tách biến theo $y$ cho
\[u(x,y)=-\sum_j\frac{d_j\cosh(k_j(1-x))}{k_j\sinh k_j}Y_j(y).\]
Đạo hàm $x$ cho đúng dữ liệu hai cạnh; $Y_j$ thỏa Robin. Đây là nghiệm duy nhất.

\clearpage\noindent\texttt{\detokenize{BT_5_1_3_05}}\par
\section*{Lời giải bổ sung — Bài tập 5 · Mục 5.1.3}
\textit{Áp dụng kiến thức và phương pháp trong tài liệu gốc.}\par
Tìm đa thức bậc hai: $u=x^2/2-ax-y^2/2+by+C$. Khi đó $\Delta u=1-1=0$, $u_x=x-a$, $u_y=b-y$, đúng bốn điều kiện. Tất cả nghiệm sai khác nhau một hằng vì đây là Neumann thuần. Tổng thông lượng ngoài $ab-ab=0$, nên tương thích với mọi $a,b>0$.

\clearpage\noindent\texttt{\detokenize{BT_5_1_3_06}}\par
\section*{Lời giải bổ sung — Bài tập 6 · Mục 5.1.3}
\textit{Áp dụng kiến thức và phương pháp trong tài liệu gốc.}\par
Đặt $b=x^2y^2/2$, $w=u-b$. Khi đó $w_x(0)=w_x(1)=0$, $\Delta w=-x^2-y^2$, $w(x,0)=x$, $w(x,1)=-x^2/2$. Dùng cơ sở cosine $X_n=\cos n\pi x$, $q_0=\int_0^1q,q_n=2\int_0^1qX_n$:
\[w=\sum_{n\ge0}W_n(y)X_n(x).\]
Đặt $a_n=(x)_n$, $b_n=(-x^2/2)_n$, $F_n(y)=(-x^2-y^2)_n$ (đều là tích phân đa thức hữu hạn). Với $k=n\pi>0$,
\[W_n=\frac{a_n\sinh k(1-y)+b_n\sinh ky}{\sinh k}+\int_0^1\frac{-\sinh(k\min(y,s))\sinh(k[1-\max(y,s)])}{k\sinh k}F_n(s)ds.\]
Mode 0: $a_0=1/2,b_0=-1/6,F_0=-1/3-y^2$, nên
\[W_0=1/2-5y/12-y^2/6-y^4/12.\]
Thế vào $u=b+w$ là nghiệm; công thức Green đã giải mọi ODE hệ số.

\clearpage\noindent\texttt{\detokenize{BT_5_1_3_07}}\par
\section*{Lời giải bổ sung — Bài tập 7 · Mục 5.1.3}
\textit{Áp dụng kiến thức và phương pháp trong tài liệu gốc.}\par
Sine Fourier $h_n=2\pi^{-1}\int_0^\pi h(x)\sin nx\,dx$. Điều kiện giảm ở vô cực loại $e^{ny}$:
\[u=\sum_{n\ge1}h_ne^{-ny}\sin nx.\]
Nếu bỏ điều kiện ở vô cực, có thể cộng tùy ý $c_n\sinh(ny)\sin nx$, ít nhất các tổng hữu hạn, nên mất tính duy nhất. Trong lớp nghiệm bị chặn, điều kiện giảm vẫn suy ra và nghiệm trên duy nhất.

\clearpage\noindent\texttt{\detokenize{BT_5_1_3_08}}\par
\section*{Lời giải bổ sung — Bài tập 8 · Mục 5.1.3}
\textit{Áp dụng kiến thức và phương pháp trong tài liệu gốc.}\par
(i) Theo $y$, dùng $Y_j=\cos k_jy-\sin k_jy/k_j$, vì $Y'(0)=-Y(0),Y'(1)=0$. Các mode gồm cả khả năng trị riêng âm: viết chung $-Y_j''=\lambda_jY_j$, $d_j=\int_0^1Y_j/\int_0^1Y_j^2$. Mỗi hệ số $V_j$ giải $V_j''-\lambda_jV_j=0$, $V_j(0)=-d_j,V_j(1)=d_j$. Với $\lambda_j>0$,
\[V_j=d_j\frac{\sinh(\sqrt{\lambda_j}x)-\sinh(\sqrt{\lambda_j}(1-x))}{\sinh\sqrt{\lambda_j}}.\]
Với $\lambda_j=-\kappa_j^2<0$, thay sinh bởi sin. Ở đây trị riêng âm duy nhất thỏa $\kappa\tanh\kappa=1$ và $\kappa<\pi$, nên không có cộng hưởng với hai biên Dirichlet theo $x$. $u=\sum V_jY_j$.

(ii) $u=-\sum_{n\text{ lẻ}}4\sin nx\sinh(n(\pi-y))/[\pi n\sinh n\pi]$.

(iii) Miền ghi $(0,\pi)\times(0,1)$ nhưng biên ghi $x=1,y=\pi$: không khép miền, không xác định duy nhất như bản in. Nếu miền dự kiến là $(0,1)\times(0,\pi)$, dùng công thức bốn cạnh Dirichlet với cạnh dưới và phải bằng 1, hai cạnh còn lại 0.

(iv) Đặt $X=x-\pi$. Nghiệm là $1-2v$, với
\[v=\sum_{n\text{ lẻ}}\frac4{\pi n}\sin nX\frac{\sinh ny+\sinh(n(\pi-y))}{\sinh n\pi}.\]

(v) $u=1+w$, $w=0$ ở $x=\pi,2\pi$, $w_y(x,0)=-1,w_y(x,\pi)=0$. Do đó
\[w=\sum_{n\text{ lẻ}}\frac4{\pi n^2}\frac{\cosh(n(\pi-y))}{\sinh n\pi}\sin(n(x-\pi)).\]

\clearpage\noindent\texttt{\detokenize{BT_5_1_3_09}}\par
\section*{Lời giải bổ sung — Bài tập 9 · Mục 5.1.3}
\textit{Áp dụng kiến thức và phương pháp trong tài liệu gốc.}\par
Hai cạnh khác 0 có hệ số sine của $x$ hoặc $y$ bằng $2(-1)^{n+1}/(n\pi)$. Do đó
\[u=\sum_{n\ge1}\frac{2(-1)^{n+1}}{n\pi\sinh n\pi}[\sin n\pi x\sinh n\pi y+\sin n\pi y\sinh n\pi x].\]
Mặt khác $u=xy$ có $u_{xx}=u_{yy}=0$ và đúng bốn biên. Tính duy nhất Dirichlet suy chuỗi trên bằng $xy$ trong hình vuông.

\clearpage\noindent\texttt{\detokenize{VD_5_2_2_01}}\par
\section*{Lời giải bổ sung — Ví dụ 1 · Mục 5.2.2}
\textit{Áp dụng kiến thức và phương pháp trong tài liệu gốc.}\par
Ví dụ chưa có lời giải trong bản in. Với đĩa bán kính $R$, dùng phép kéo dài điều hòa $\mathcal P_R[\varphi]$:
\[\mathcal P_R[\varphi](r,\theta)=\frac{R^2-r^2}{2\pi}\int_0^{2\pi}\frac{\varphi(R\cos s,R\sin s)}{R^2+r^2-2Rr\cos(\theta-s)}ds.\]
(i) $u=(R^2+x^2-y^2)/2-2y$.
(ii) Biên $\min(x,0)=(x-|x|)/2$ cho
\[u=\frac x2-\frac R\pi-\frac{2R}\pi\sum_{n\ge1}\frac{(-1)^{n+1}(r/R)^{2n}}{4n^2-1}\cos2n\theta.\]
(iii) $u=1/2+(2/\pi)\sum_{n\text{ lẻ}}(r/R)^n\sin n\theta/n$. Các công thức đến từ Fourier của dữ liệu trên đường tròn; mọi nghiệm đều bị chặn tại tâm.

\clearpage\noindent\texttt{\detokenize{VD_5_2_2_02}}\par
\section*{Lời giải bổ sung — Ví dụ 2 · Mục 5.2.2}
\textit{Áp dụng kiến thức và phương pháp trong tài liệu gốc.}\par
(i) $u=x(x^2+y^2-R^2)/8$, vì $\Delta[x(x^2+y^2)]=8x$, và biên 0.
(ii) $u=(r^2-R^2)/4+\mathcal P_R[\min(x,0)]$, với $\mathcal P_R$ có chuỗi ở ví dụ trước. Hạng đầu có Laplace bằng 1 và biên 0, hạng sau điều hòa và đúng dữ liệu biên.

\clearpage\noindent\texttt{\detokenize{VD_5_2_2_03}}\par
\section*{Lời giải bổ sung — Ví dụ 3 · Mục 5.2.2}
\textit{Áp dụng kiến thức và phương pháp trong tài liệu gốc.}\par
(i) Pháp tuyến ngoài trên đường tròn là $e_r$, nên biên $u_r(R,\theta)=\varphi(R\cos\theta,R\sin\theta)$.
(ii) Hiệu $w$ có $\Delta w=0,\partial_\nu w=0$; Green cho $\int|\nabla w|^2=0$, nên mọi nghiệm sai khác hằng.
(iii) Với $\Delta u=F$, điều kiện cần và đủ Neumann là $\int_{\partial D}\varphi=\int_DF$. Có thể giải bằng một nghiệm riêng $P$ của $\Delta P=F$, rồi khai triển biên $h=\varphi-P_r$ có trung bình 0:
\[u=P+C+\sum_{n\ge1}\frac{r^n}{nR^{n-1}}(h_n^c\cos n\theta+h_n^s\sin n\theta).\]
Nếu $F$ chưa được chỉ định trong đề thì kết quả phải phụ thuộc $F$ như công thức này.
(iv) Khi $F=0,R=1$, trung bình $x^2-2xy+k$ là $1/2+k$, nên $k=-1/2$; nghiệm $u=(x^2-y^2)/4-xy+C$. Nếu $F$ khác 0, phải thay điều kiện $k$ bằng $\int_DF/(2\pi)-1/2$ và thêm nghiệm riêng; bản in không ghi $F$ cho câu này.

\clearpage\noindent\texttt{\detokenize{BT_5_2_3_01}}\par
\section*{Lời giải bổ sung — Bài tập 1 · Mục 5.2.3}
\textit{Áp dụng kiến thức và phương pháp trong tài liệu gốc.}\par
Đặt $q=r/R$. (i) $u=1/2+(2/\pi)\sum_{n\text{ lẻ}}q^n\sin n\theta/n$.
(ii) Biên $2y\chi_{y>0}=y+|y|$ cho
\[u=y+\frac{2R}\pi-\frac{4R}\pi\sum_{n\ge1}\frac{q^{2n}}{4n^2-1}\cos2n\theta.\]
(iii) $u=1-x$.
(iv) $u=(R^2+x^2-y^2)/2$.
(v) $u=(R^2-x^2+y^2)/2$.
(vi) $u=xy$.
(vii) $u=1+3y/\pi$.
Các đa thức được chọn điều hòa và khớp biên; hai trường hợp phân đoạn lấy hệ số Fourier trên đường tròn rồi nhân $(r/R)^n$.

\clearpage\noindent\texttt{\detokenize{BT_5_2_3_02}}\par
\section*{Lời giải bổ sung — Bài tập 2 · Mục 5.2.3}
\textit{Áp dụng kiến thức và phương pháp trong tài liệu gốc.}\par
Trong bản PDF, bài số 2 chỉ có số thứ tự, không có nội dung hoặc dữ kiện trước tiêu đề mục 5.3. Vì vậy không có đề toán để lập lời giải. File giữ dấu bài 2 để không bỏ sót mục được đánh số; không tự tạo một đề khác.

\clearpage\noindent\texttt{\detokenize{VD_5_3_2_01}}\par
\section*{Lời giải bổ sung — Ví dụ 1 · Mục 5.3.2}
\textit{Áp dụng kiến thức và phương pháp trong tài liệu gốc.}\par
Với hình quạt $0<r<R,0<\theta<\alpha$, đặt $b=(1-\theta/\alpha)g_0(r)+(\theta/\alpha)g_1(r)$, $w=u-b$, $H=F-\Delta b$, $h(\theta)=\varphi(\theta)-b(R,\theta)$. Đặt $k_n=n\pi/\alpha$, $H_n=2\alpha^{-1}\int_0^\alpha H\sin k_n\theta$, $h_n$ tương tự. Khi đó
\[u=b+\sum_{n\ge1}\left[h_n(r/R)^{k_n}+\int_0^R\mathcal G_n(r,s)H_n(s)s\,ds\right]\sin k_n\theta,\]
\[\mathcal G_n=-\frac1{2k_n}\left[(r_</r_>)^{k_n}-(r_<r_>/R^2)^{k_n}\right].\]
Đây là Green của ODE xuyên tâm, bị chặn tại 0 và bằng 0 tại $R$. Năng lượng Dirichlet cho duy nhất. Để nghiệm liên tục tại đỉnh, $g_0(0)=g_1(0)$ và dữ liệu phải tương thích.
Câu cụ thể (i) $\alpha=\pi/2,F=0,g_0=g_1=0,\varphi=xy$: $u=xy$ vì điều hòa và đúng cả ba biên. (ii) $\alpha=\pi,F=1$, mọi biên 0: đặt $k_n=n$, $H_n=2(1-(-1)^n)/(\pi n)$. Với $n$ lẻ,
\[u=\sum_{n\text{ lẻ}}\frac{H_n}{4-n^2}[r^2-R^2(r/R)^n]\sin n\theta.\]
Các mode chẵn bằng 0; mode cộng hưởng $n=2$ không có nguồn nên không xuất hiện. Thế Laplace cho $1$.

\clearpage\noindent\texttt{\detokenize{VD_5_3_2_02}}\par
\section*{Lời giải bổ sung — Ví dụ 2 · Mục 5.3.2}
\textit{Áp dụng kiến thức và phương pháp trong tài liệu gốc.}\par
Trong bản in, biên cạnh là $-u_\theta(r,0)=g_0(r),u_\theta(r,\alpha)=g_1(r)$; thông lượng pháp tuyến thực tế lần lượt là $g_0/r,g_1/r$, vì đạo hàm theo góc có thừa số $1/r$. Biên cung có thông lượng $\varphi$. Green cho duy nhất modulo hằng. Điều kiện tương thích:
\[\int_0^R\frac{g_0+g_1}{r}dr+R\int_0^\alpha\varphi(\theta)d\theta=\int_0^R\int_0^\alpha F(r,\theta)r\,d\theta dr.\]
Có thể khử biên cạnh bằng $b$ rồi khai triển cosine $\cos(n\pi\theta/\alpha)$; mỗi hệ số giải $(rW_n')'-k_n^2W_n/r=r(F-\Delta b)_n$, với đạo hàm đã cho tại $R$. Mode 0 chính là điều kiện tương thích trên, hằng tích phân còn lại tùy ý.

Câu (iii) của ví dụ chỉ ghi $R=1,\varphi=x^2-2xy+k$; chưa ghi $F,\alpha,g_0,g_1$. Vì vậy không đủ xác định một nghiệm cụ thể trong hình quạt. Khi các dữ liệu đó được chọn, phương trình trên xác định $k$ và các hệ số. Nếu ý định là đĩa toàn phần, $F=0$, thì $k=-1/2,u=(x^2-y^2)/4-xy+C$ như ví dụ 5.2.2.3; đây là một bài đã sửa, không phải dữ kiện gốc.

\clearpage\noindent\texttt{\detokenize{BT_5_3_3_01}}\par
\section*{Lời giải bổ sung — Bài tập 1 · Mục 5.3.3}
\textit{Áp dụng kiến thức và phương pháp trong tài liệu gốc.}\par
Tổng thông lượng trên hai cung là $b\int f_2-a\int f_1$, hai cạnh góc thông lượng 0. Vì $\Delta u=0$, điều kiện cần là $a\int f_1=b\int f_2$. Đặt $z=\ln(r/a)$, $H=\ln(b/a)$, $k_n=2n$, $c_{in}=4\pi^{-1}\int_0^{\pi/2}f_i\cos2n\theta$, $c_{i0}=2\pi^{-1}\int f_i$. Nghiệm là
\[u=C+ac_{10}\ln(r/a)+\sum_{n\ge1}\frac{bc_{2n}\cosh(k_nz)-ac_{1n}\cosh(k_n(H-z))}{k_n\sinh k_nH}\cos2n\theta.\]
Đạo hàm theo $r$ cho đúng $f_i$; mode 0 cần $ac_{10}=bc_{20}$, chính là điều kiện đã cho. Vì thế điều kiện đủ; năng lượng cho duy nhất modulo hằng, nên không duy nhất nếu chưa chuẩn hóa.

\clearpage\noindent\texttt{\detokenize{BT_5_3_3_02}}\par
\section*{Lời giải bổ sung — Bài tập 2 · Mục 5.3.3}
\textit{Áp dụng kiến thức và phương pháp trong tài liệu gốc.}\par
Đặt $z=\ln r,H=\ln2$. (i) Khai triển $1=\sum_{n\text{ lẻ}}4\sin n\theta/(\pi n)$:
\[u=\sum_{n\text{ lẻ}}\frac4{\pi n}\frac{\sinh(n(H-z))+\sinh nz}{\sinh nH}\sin n\theta.\]
(ii) $u=(4/r-r)\sin\theta/3$.
(iii) Trên vành khuyên toàn phần, $u=(4/r-r)\cos\theta/3+2(r-1/r)\sin\theta/3$. Các biểu thức đều là tổ hợp $r^n,r^{-n}$ nên điều hòa và thỏa đúng hai cung.

\clearpage\noindent\texttt{\detokenize{BT_5_3_3_03}}\par
\section*{Lời giải bổ sung — Bài tập 3 · Mục 5.3.3}
\textit{Áp dụng kiến thức và phương pháp trong tài liệu gốc.}\par
Đặt $z=\ln(r/r_1),H=\ln(r_2/r_1)$. Gọi $a_{i0},a_{in},b_{in}$ là các hệ số Fourier của $f_i$ (mode hằng là trung bình). Khi đó
\begin{align*}
u={}&a_{10}(1-z/H)+a_{20}z/H\\
&+\sum_{n\ge1}\frac{[a_{1n}\sinh n(H-z)+a_{2n}\sinh nz]\cos n\theta}{\sinh nH}\\
&+\sum_{n\ge1}\frac{[b_{1n}\sinh n(H-z)+b_{2n}\sinh nz]\sin n\theta}{\sinh nH}.
\end{align*}
Với $r_1=3,r_2=6$, $f_1=5\cos\theta,f_2=10\sin\theta$,
\[u=\left(\frac{20}r-\frac{5r}9\right)\cos\theta+\left(\frac{20r}9-\frac{20}r\right)\sin\theta.\]

\clearpage\noindent\texttt{\detokenize{BT_5_3_3_04}}\par
\section*{Lời giải bổ sung — Bài tập 4 · Mục 5.3.3}
\textit{Áp dụng kiến thức và phương pháp trong tài liệu gốc.}\par
Nghiệm điều hòa trơn tại tâm có dạng $a_0+\sum r^n(a_n\cos n\theta+b_n\sin n\theta)$. Biên $u_r-u=0$ tại $r=1$ cho $a_0=0$, $(n-1)a_n=(n-1)b_n=0$. Chỉ mode $n=1$ còn lại. Tất cả nghiệm là $u=Ax+By$; đây là Robin với dấu âm, nên khác bài Robin có dấu dương cho duy nhất.

\clearpage\noindent\texttt{\detokenize{BT_5_4_3_01}}\par
\section*{Lời giải bổ sung — Bài tập 1 · Mục 5.4.3}
\textit{Áp dụng kiến thức và phương pháp trong tài liệu gốc.}\par
(a)--(c) Với hàm xuyên tâm $r^p$, $\Delta r^p=p(p+n-2)r^{p-2}$. Vì $r>0$, Laplace bằng 0 đúng khi $p=0$ hoặc $p=2-n$. Trong $n=3,p=-1$ cho $1/r$ điều hòa ngoài gốc.

(d) Đặt $v=x-\xi$, $|\xi|=1$, $1-|x|^2=-|v|^2-2v\cdot\xi$. Khi đó nhân Poisson bằng $-|v|^{2-n}-2(v\cdot\xi)|v|^{-n}$. Với $n\ne2$, hạng sau là đạo hàm theo hướng của $|v|^{2-n}$ chia $2-n$; do đó điều hòa. Với $n=2$, dùng $(v\cdot\xi)|v|^{-2}=\partial_\xi\ln|v|$, cũng điều hòa. Miền đúng là $\mathbb R^n\setminus\{\xi\}$; ký hiệu $\mathbb R^3$ ở dòng cuối đề chỉ phù hợp khi $n=3$.

\clearpage\noindent\texttt{\detokenize{BT_5_4_3_02}}\par
\section*{Lời giải bổ sung — Bài tập 2 · Mục 5.4.3}
\textit{Áp dụng kiến thức và phương pháp trong tài liệu gốc.}\par
(a) Quy tắc tích cho $\Delta(u^2)=2u\Delta u+2|\nabla u|^2=2|\nabla u|^2\ge0$.
(b) Tính giá trị trung bình và Cauchy--Schwarz trên quả cầu bán kính $R$ cho $|u(x)|^2\le |B_R|^{-1}\int_{B_R(x)}u^2\le |B_R|^{-1}\|u\|_{L^2(\mathbb R^n)}^2$. Cho $R\to\infty$, vế phải về 0, nên $u=0$.

\clearpage\noindent\texttt{\detokenize{BT_5_4_3_03}}\par
\section*{Lời giải bổ sung — Bài tập 3 · Mục 5.4.3}
\textit{Áp dụng kiến thức và phương pháp trong tài liệu gốc.}\par
Nếu $\Omega$ là một khoảng liên thông, $u''=0$ suy $u'=m$, $u=mx+b$; chiều ngược hiển nhiên. Nếu $\Omega$ chỉ là tập mở như đề ghi, khẳng định một cặp $m,b$ chung là sai: trên $(0,1)\cup(2,3)$ lấy $u=0$ ở khoảng đầu và $u=1$ ở khoảng sau. Kết luận đúng là affine riêng trên từng thành phần liên thông; cần thêm liên thông để có một cặp toàn miền.

\clearpage\noindent\texttt{\detokenize{BT_5_4_3_04}}\par
\section*{Lời giải bổ sung — Bài tập 4 · Mục 5.4.3}
\textit{Áp dụng kiến thức và phương pháp trong tài liệu gốc.}\par
$u=6x^2yz-y^3z-yz^3$ cho $u_{xx}=12yz$, $u_{yy}=-6yz$, $u_{zz}=-6yz$. Tổng bằng 0, nên điều hòa.

\clearpage\noindent\texttt{\detokenize{BT_5_4_3_05}}\par
\section*{Lời giải bổ sung — Bài tập 5 · Mục 5.4.3}
\textit{Áp dụng kiến thức và phương pháp trong tài liệu gốc.}\par
Với đa thức in trong đề,
\begin{align*}
u_{xx}&=138xy^2z-40x^3z-6xz^3,\\
u_{yy}&=46x^3z-96xy^2z-14xz^3,\\
u_{zz}&=-42xy^2z-6x^3z+20xz^3.
\end{align*}
Các hệ số của $xy^2z,x^3z,xz^3$ lần lượt triệt tiêu, nên $\Delta u=0$.

\clearpage\noindent\texttt{\detokenize{BT_5_4_3_06}}\par
\section*{Lời giải bổ sung — Bài tập 6 · Mục 5.4.3}
\textit{Áp dụng kiến thức và phương pháp trong tài liệu gốc.}\par
$u_r=(1-r^{-2})\cos\theta,u_{rr}=2r^{-3}\cos\theta,u_{\theta\theta}=-(r+r^{-1})\cos\theta$. Thế vào $u_{rr}+u_r/r+u_{\theta\theta}/r^2$ được 0 với $r>0$.

\clearpage\noindent\texttt{\detokenize{BT_5_4_3_07}}\par
\section*{Lời giải bổ sung — Bài tập 7 · Mục 5.4.3}
\textit{Áp dụng kiến thức và phương pháp trong tài liệu gốc.}\par
Đặt $y=x/|x|^2,r=|x|$. Đạo hàm $\partial_i y_j=r^{-2}(\delta_{ij}-2x_ix_j/r^2)$, ma trận trong ngoặc trực giao. Quy tắc dây chuyền, cộng các đạo hàm cấp hai và hạng từ $r^{2-n}$ cho đồng nhất thức Kelvin:
\[\Delta_x[r^{2-n}u(y)]=r^{-n-2}(\Delta_yu)(y).\]
Vì vậy công thức đúng là $r^{-n-2}f(x/r^2)$; số mũ $-n-2$ cần giữ đúng. Nửa không gian $x_2>0$ bất biến qua phép đảo, nên miền hợp lệ. Với $f=0$, $v$ điều hòa.

\clearpage\noindent\texttt{\detokenize{BT_5_4_3_08}}\par
\section*{Lời giải bổ sung — Bài tập 8 · Mục 5.4.3}
\textit{Áp dụng kiến thức và phương pháp trong tài liệu gốc.}\par
$\Delta(uv)=u\Delta v+v\Delta u+2\nabla u\cdot\nabla v=2\nabla u\cdot\nabla v$. Do đó $uv$ điều hòa khi và chỉ khi tích vô hướng bằng 0.

\clearpage\noindent\texttt{\detokenize{BT_5_4_3_09}}\par
\section*{Lời giải bổ sung — Bài tập 9 · Mục 5.4.3}
\textit{Áp dụng kiến thức và phương pháp trong tài liệu gốc.}\par
$\Delta(u^2)=2|\nabla u|^2$. Nếu $u^2$ điều hòa, $\nabla u=0$; liên thông suy $u$ hằng. Chiều ngược cũng đúng.

\clearpage\noindent\texttt{\detokenize{BT_5_4_3_10}}\par
\section*{Lời giải bổ sung — Bài tập 10 · Mục 5.4.3}
\textit{Áp dụng kiến thức và phương pháp trong tài liệu gốc.}\par
Viết ma trận $T=A$. Với mọi ma trận đối xứng có vết 0, đa thức $u(y)=y^THy$ điều hòa. Giả thiết cho $0=\Delta(u\circ A)=2\operatorname{tr}(HAA^T)$. Trực giao với mọi $H$ vết 0 buộc $AA^T=c^2I$, $c\ge0$. Nếu $c=0$ thì $A=0$; nếu $c>0$, $A=cQ$ với $Q$ trực giao. Đây là kết luận, kể cả phép biến đổi 0.

\clearpage\noindent\texttt{\detokenize{BT_5_4_3_11}}\par
\section*{Lời giải bổ sung — Bài tập 11 · Mục 5.4.3}
\textit{Áp dụng kiến thức và phương pháp trong tài liệu gốc.}\par
Với $u=f(r)$, $(r^{n-1}f')'=0$. Nghiệm trên miền thủng là $C+Dr^{2-n}$ ($n\ne2$) hoặc $C+D\ln r$ ($n=2$). Trơn tại tâm quả cầu buộc $D=0$, nên $u$ hằng. Trong $n=1$, tính chẵn và affine cũng buộc hằng.

\clearpage\noindent\texttt{\detokenize{BT_5_4_3_12}}\par
\section*{Lời giải bổ sung — Bài tập 12 · Mục 5.4.3}
\textit{Áp dụng kiến thức và phương pháp trong tài liệu gốc.}\par
Đa thức thuần nhất $p$ bậc $m$ có $p(x/r^2)=r^{-2m}p(x)$. Kelvin của $p$ là $r^{2-n-2m}p(x)$; vì Kelvin giữ điều hòa, $p(x)/r^{2m+n-2}$ điều hòa ngoài gốc.

\clearpage\noindent\texttt{\detokenize{BT_5_4_3_13}}\par
\section*{Lời giải bổ sung — Bài tập 13 · Mục 5.4.3}
\textit{Áp dụng kiến thức và phương pháp trong tài liệu gốc.}\par
Ký hiệu $B$ trong đề dùng $|x|\le1$: nếu $u$ liên tục trên quả cầu đóng thì Heine--Cantor buộc liên tục đều, nên yêu cầu như vậy không thể thỏa. Với ý định là quả cầu mở, trong đĩa lấy
\[u(x,y)=\operatorname{Re}\exp\left[-\frac{1+x+iy}{1-x-iy}\right].\]
Hàm giải tích vì cực $z=1$ nằm ngoài miền; $\operatorname{Re}(1+z)/(1-z)>0$ nên $|u|\le1$. Khi $z\to1$ theo bán kính, $u\to0$; theo các đường có phần thực thương tiến về một hằng và phần ảo thay đổi, $u$ có nhiều giá trị giới hạn. Do đó không liên tục đều. Có thể mở rộng ví dụ sang $n\ge2$ bằng cách không phụ thuộc các tọa độ còn lại.

\clearpage\noindent\texttt{\detokenize{BT_5_4_3_14}}\par
\section*{Lời giải bổ sung — Bài tập 14 · Mục 5.4.3}
\textit{Áp dụng kiến thức và phương pháp trong tài liệu gốc.}\par
Đặt $v(x,y)=u(x,y)+u(x,-y)$. Đây là hàm điều hòa chẵn theo $y$, và $v(x,0)=0$, $v_y(x,0)=0$. Từ $v_{yy}=-\Delta_xv$, tất cả đạo hàm theo $y$ trên siêu phẳng đều 0 theo quy nạp. Hàm điều hòa giải tích, nên $v=0$ gần siêu phẳng, rồi trên toàn miền nhờ tính duy nhất giải tích. Suy $u(x,-y)=-u(x,y)$.

\clearpage\noindent\texttt{\detokenize{BT_5_4_3_15}}\par
\section*{Lời giải bổ sung — Bài tập 15 · Mục 5.4.3}
\textit{Áp dụng kiến thức và phương pháp trong tài liệu gốc.}\par
Ước lượng đạo hàm nội suy từ công thức Poisson: $|D^ku(x_0)|\le C_kR^{-k}\sup_{B_R(x_0)}|u|\le C'_kR^{-k}(1+(R+|x_0|)^p)$. Chọn số nguyên $k>p$, cho $R\to\infty$ được $D^ku=0$ với mọi $x_0$. Taylor suy $u$ là đa thức bậc không quá $\lfloor p\rfloor$.

\clearpage\noindent\texttt{\detokenize{BT_5_4_3_16}}\par
\section*{Lời giải bổ sung — Bài tập 16 · Mục 5.4.3}
\textit{Áp dụng kiến thức và phương pháp trong tài liệu gốc.}\par
Đặt $d=1-|x|$, dùng quả cầu $B_{d/2}(x)\subset B$. Công thức Poisson đạo hàm cho $|\nabla u(x)|\le C_nd^{-1}\|u\|_\infty$. Suy $\sup_B(1-|x|)|\nabla u(x)|\le C_n\|u\|_\infty<\infty$.

\clearpage\noindent\texttt{\detokenize{BT_5_4_3_17}}\par
\section*{Lời giải bổ sung — Bài tập 17 · Mục 5.4.3}
\textit{Áp dụng kiến thức và phương pháp trong tài liệu gốc.}\par
Chuẩn $\|u\|=|u(0)|+\sup_B(1-|x|)|\nabla u|$ có tính đồng nhất và tam giác; chuẩn 0 buộc $\nabla u=0,u(0)=0$, nên $u=0$. Với dãy Cauchy $u_j$, tích phân trên đoạn từ 0 đến $x$ cho
\[|u_j(x)-u_k(x)|\le\|u_j-u_k\|\left[1+\log\frac1{1-|x|}\right].\]
Do đó hội tụ đều trên compact đến $u$. Tính giá trị trung bình hoặc ước lượng đạo hàm cho $u$ điều hòa và $\nabla u_j\to\nabla u$ trên compact. Cho $k\to\infty$ trong $\sup(1-|x|)|\nabla(u_j-u_k)|\le\varepsilon$ được cùng cận cho $u_j-u$; giá trị tại 0 cũng hội tụ. Vậy $u_j\to u$ theo chuẩn, chứng minh Banach.

\clearpage\noindent\texttt{\detokenize{BT_5_4_3_18}}\par
\section*{Lời giải bổ sung — Bài tập 18 · Mục 5.4.3}
\textit{Áp dụng kiến thức và phương pháp trong tài liệu gốc.}\par
Trong đĩa, $u(z)=\operatorname{Re}\log(1-z)=\log|1-z|$ điều hòa và không bị chặn khi $z\to1$. Vì $|\nabla u|=1/|1-z|\le1/(1-|z|)$, điều kiện Bloch thỏa. Trong $n\ge2$, lấy hàm chỉ phụ thuộc $x_1,x_2$; $|1-(x_1+ix_2)|\ge1-|x|$ cho cùng cận.

\clearpage\noindent\texttt{\detokenize{BT_5_4_3_19}}\par
\section*{Lời giải bổ sung — Bài tập 19 · Mục 5.4.3}
\textit{Áp dụng kiến thức và phương pháp trong tài liệu gốc.}\par
Đặt $d=1-|x|$, $B_{d/2}(x)\subset B$. Trên quả cầu này, $|\nabla u|\le2M/d$ với $M=\sup(1-|x|)|\nabla u|$. Mỗi $u_{x_i}$ điều hòa. Ước lượng đạo hàm bậc $|\alpha|-1$ của $u_{x_i}$ cho $|D^\alpha u(x)|\le C_\alpha M d^{-|\alpha|}$. Nhân $d^{|\alpha|}$ rồi lấy supremum được kết luận.

\clearpage\noindent\texttt{\detokenize{BT_5_5_3_01}}\par
\section*{Lời giải bổ sung — Bài tập 1 · Mục 5.5.3}
\textit{Áp dụng kiến thức và phương pháp trong tài liệu gốc.}\par
Giá trị trung bình trên đường tròn: $u(0)=\frac1{2\pi}\int_0^{2\pi}\cos^2t\sin^2t\,dt=1/8$, vì $\cos^2t\sin^2t=(1-\cos4t)/8$.

\clearpage\noindent\texttt{\detokenize{BT_5_5_3_02}}\par
\section*{Lời giải bổ sung — Bài tập 2 · Mục 5.5.3}
\textit{Áp dụng kiến thức và phương pháp trong tài liệu gốc.}\par
Biên có cực đại 4 khi $\sin2\theta=1$ và cực tiểu $-2$ khi $\sin2\theta=-1$. Cực đại trên đĩa đóng bằng 4, đạt trên biên tại $\theta=\pi/4,5\pi/4$. Giá trị tâm là trung bình bằng 1. Nghiệm cụ thể $u=1+3r^2\sin2\theta/4=1+3xy/2$.

\clearpage\noindent\texttt{\detokenize{BT_5_5_3_03}}\par
\section*{Lời giải bổ sung — Bài tập 3 · Mục 5.5.3}
\textit{Áp dụng kiến thức và phương pháp trong tài liệu gốc.}\par
Hàm là nhân Poisson với cực $(1,0)$. Nó điều hòa trong đĩa mở, dương, nhưng không liên tục và không bị chặn ở cực: $u(r,0)=(1+r)/(1-r)\to\infty$. Ở các điểm biên khác cực, giới hạn bằng 0. Nguyên lý cực đại đòi liên tục trên miền đóng nên không áp dụng như đề phát biểu cho quả cầu đóng chứa cực. Không có mâu thuẫn: hàm không đạt cực đại/cực tiểu trong miền mở.

\clearpage\noindent\texttt{\detokenize{BT_5_5_3_04}}\par
\section*{Lời giải bổ sung — Bài tập 4 · Mục 5.5.3}
\textit{Áp dụng kiến thức và phương pháp trong tài liệu gốc.}\par
Mỗi hàm có Laplace 0 bằng đạo hàm trực tiếp (ở (e) ngoài gốc). Xét các cạnh/đường tròn theo cực đại:
(a) max $4$ tại $(2,0)$, min $-1$ tại $(0,1)$.
(b) max $e$ tại $(0,1)$, min 0 trên $x=\pm\pi/2$.
(c) max $5/4$ tại $(1,1/2)$, min $-1$ tại $(0,1)$.
(d) max $e$ tại $(1,0)$, min $-e$ tại $(1,\pm\pi)$.
(e) $u=\cos\theta/r$: max 1 tại $(1,0)$, min $-1$ tại $(-1,0)$.
(f) $u=(r^2/2)\sin2\theta$: max 2 tại $(\sqrt2,\sqrt2),(-\sqrt2,-\sqrt2)$; min $-2$ tại $(\sqrt2,-\sqrt2),(-\sqrt2,\sqrt2)$.

\clearpage\noindent\texttt{\detokenize{BT_5_5_3_05}}\par
\section*{Lời giải bổ sung — Bài tập 5 · Mục 5.5.3}
\textit{Áp dụng kiến thức và phương pháp trong tài liệu gốc.}\par
$u=2\log r$ có $u_{rr}+u_r/r=0$ với $r>1$, và $u_r=2/r>0$. Vì vậy không có cực trị địa phương trong $D$. Infimum 0 khi $r\downarrow1$ không đạt; supremum $+\infty$ khi $r\to\infty$.

\clearpage\noindent\texttt{\detokenize{BT_5_5_3_06}}\par
\section*{Lời giải bổ sung — Bài tập 6 · Mục 5.5.3}
\textit{Áp dụng kiến thức và phương pháp trong tài liệu gốc.}\par
Dùng trung bình điều hòa trên đường tròn tâm của đối số:
(a) $e^x\cos y$ tại $(0,0)$ cho $I_1=1$.
(b) $3+x(1+y)$ điều hòa; tại $(0,0)$ cho $I_2=3$.
(c) $\cos x\cosh y$ điều hòa; tại $(1,2)$ cho $I_3=\cos1\cosh2$.
(d) $\sin x\sinh y$ điều hòa; tại $(1,2)$ cho $I_4=\sin1\sinh2$.

\clearpage\noindent\texttt{\detokenize{BT_5_5_3_07}}\par
\section*{Lời giải bổ sung — Bài tập 7 · Mục 5.5.3}
\textit{Áp dụng kiến thức và phương pháp trong tài liệu gốc.}\par
Trên biên $x^4+y^4=1$, đặt $a=x^2,b=y^2\ge0$ thì $a^2+b^2=1$, nên $1\le a+b\le\sqrt2$. Hai dấu bằng đạt tại các trục và tại $a=b=1/\sqrt2$. Dữ liệu biên không hằng, nên nguyên lý cực đại mạnh cho $1<u<\sqrt2$ ở mọi điểm trong miền.

\clearpage\noindent\texttt{\detokenize{BT_5_5_3_08}}\par
\section*{Lời giải bổ sung — Bài tập 8 · Mục 5.5.3}
\textit{Áp dụng kiến thức và phương pháp trong tài liệu gốc.}\par
Nhân Poisson dương nên $h\ge0$ suy $u\ge0$. Nếu $h$ không đồng nhất 0 và liên tục, tích phân nhân dương cho $u>0$. Nếu $h=0$, duy nhất Dirichlet cho $u=0$.
Với $r=1/2$, $P_r(\theta)=(1-r^2)/(1-2r\cos\theta+r^2)$ nằm giữa $(1-r)/(1+r)=1/3$ và $(1+r)/(1-r)=3$. Vì trung bình $h$ bằng $u(0)=1$, tích phân cho $1/3\le u\le3$.

\clearpage\noindent\texttt{\detokenize{BT_5_5_3_09}}\par
\section*{Lời giải bổ sung — Bài tập 9 · Mục 5.5.3}
\textit{Áp dụng kiến thức và phương pháp trong tài liệu gốc.}\par
Bochner: $\Delta|\nabla u|^2=2\sum_{i,j}u_{ij}^2+2\nabla u\cdot\nabla\Delta u=2\sum u_{ij}^2\ge0$. Nguyên lý cực đại cho hàm dưới điều hòa suy supremum ở biên. Nếu giá trị cực đại đạt trong miền, nguyên lý mạnh buộc $|\nabla u|^2$ hằng và $D^2u=0$, tức $u$ affine; khi ấy cực đại cũng đạt trong miền. Vì thế chữ ``chỉ'' trong đề cần ngoại lệ affine/hằng.

\clearpage\noindent\texttt{\detokenize{BT_5_5_3_10}}\par
\section*{Lời giải bổ sung — Bài tập 10 · Mục 5.5.3}
\textit{Áp dụng kiến thức và phương pháp trong tài liệu gốc.}\par
Đặt $v=K(1-|x|^2)/(2n)$, $\Delta v=-K$. Do $|\Delta u|\le K$, $\Delta(u-v)\ge0$ và $\Delta(u+v)\le0$. Cả hai bằng 0 trên biên. Cực đại/cực tiểu cho $-v\le u\le v$, mạnh hơn cận $-K/(2n)\le u\le K/(2n)$ cần chứng minh.

\clearpage\noindent\texttt{\detokenize{BT_5_5_3_11}}\par
\section*{Lời giải bổ sung — Bài tập 11 · Mục 5.5.3}
\textit{Áp dụng kiến thức và phương pháp trong tài liệu gốc.}\par
Nếu $u$ hằng trên một đĩa mở, tính giải tích và liên thông suy hằng trên toàn $D(1)$, trái giả thiết. Với $0<r_1<r_2$, cực đại trên $\overline D(r_1)$ nhỏ nghiêm ngặt cực đại trên $\overline D(r_2)$: nếu bằng nhau, $u$ đạt cực đại ở điểm nội của đĩa lớn nên hằng. Với $r_1=0$, dùng cực đại mạnh tương tự. Vì vậy $\phi$ tăng nghiêm ngặt, kéo theo không giảm.

\clearpage\noindent\texttt{\detokenize{BT_5_5_3_12}}\par
\section*{Lời giải bổ sung — Bài tập 12 · Mục 5.5.3}
\textit{Áp dụng kiến thức và phương pháp trong tài liệu gốc.}\par
Không được suy trực tiếp $u_t\to0$ từ hội tụ $u\to A$. Lấy $\eta\in C_c^\infty(B_R)$, $T>0$:
\[\int[u(x,t+T)-u(x,t)]\eta dx=\int_t^{t+T}\int u(x,s)\Delta\eta(x)dx\,ds.\]
Hội tụ đều cho vế trái về 0, vế phải về $T\int A\Delta\eta$. Do đó $\Delta A=0$ theo phân phối; tính trơn nội của nghiệm điều hòa suy $A$ là hàm điều hòa cổ điển.

\clearpage\noindent\texttt{\detokenize{BT_5_5_3_13}}\par
\section*{Lời giải bổ sung — Bài tập 13 · Mục 5.5.3}
\textit{Áp dụng kiến thức và phương pháp trong tài liệu gốc.}\par
(a) Phát biểu thiếu giả thiết $\Delta u=0$: với $u$ chỉ $C^2$ tùy ý, công thức đúng là $\int_{\partial\Omega}u\partial_\nu u=\int_\Omega(|\nabla u|^2+u\Delta u)$, không luôn không âm. Phản ví dụ trên quả cầu đơn vị: $u=2-|x|^2$, tích phân biên $=-2|\partial B|<0$. Nếu thêm điều hòa, tích phân bằng $\int|\nabla u|^2\ge0$, bằng 0 đúng khi $u$ hằng.
(b) Hiệu $w$ của hai nghiệm cùng $f,\beta$ là điều hòa và $\partial_\nu w=-\alpha w$. Vì $\alpha>0$ hầu khắp, $\int|\nabla w|^2=-\int\alpha w^2\le0$ suy $w=0$. Vậy bài Robin có nhiều nhất một nghiệm.

\clearpage\noindent\texttt{\detokenize{BT_5_6_3_01}}\par
\section*{Lời giải bổ sung — Bài tập 1 · Mục 5.6.3}
\textit{Áp dụng kiến thức và phương pháp trong tài liệu gốc.}\par
Cho hàm riêng phức $\phi\ne0$, nhân $-\Delta\phi=\lambda\phi$ với $\overline\phi$ rồi tích phân:
\[\lambda\int_\Omega|\phi|^2=\int_\Omega|\nabla\phi|^2+\frac\alpha\beta\int_{\partial\Omega}|\phi|^2.\]
Vì $\alpha\beta>0$, $\alpha/\beta>0$, vế phải thực và dương: nếu bằng 0 thì $\phi$ hằng và biên buộc 0. Do đó mọi trị riêng thực dương.

\clearpage\noindent\texttt{\detokenize{BT_5_6_3_02}}\par
\section*{Lời giải bổ sung — Bài tập 2 · Mục 5.6.3}
\textit{Áp dụng kiến thức và phương pháp trong tài liệu gốc.}\par
Với $z=\ln(x+1)$, đặt $X=(x+1)^{-1/2}\sin(n\pi z)$. Tính $\partial_x[(x+1)^2X']=-(n^2\pi^2+1/4)X$. Tương tự theo $y$. Do đó
\[L\phi_{nm}=-\left[\pi^2(n^2+m^2)+\frac12\right]\phi_{nm}.\]
Trị riêng theo quy ước $L\phi=-\lambda\phi$ là $\lambda=\pi^2(n^2+m^2)+1/2$. Các hàm này không triệt tiêu ở $x=1,y=1$ nói chung; đề chỉ hỏi hàm riêng của toán tử, không được thêm kết luận biên Dirichlet. Nếu cần Dirichlet trên $(0,1)^2$, thay $n\pi\ln(x+1)$ bởi $n\pi\ln(x+1)/\ln2$ và sửa trị riêng tương ứng.

\clearpage\noindent\texttt{\detokenize{BT_5_6_3_03}}\par
\section*{Lời giải bổ sung — Bài tập 3 · Mục 5.6.3}
\textit{Áp dụng kiến thức và phương pháp trong tài liệu gốc.}\par
(a) Neumann hình chữ nhật: $\phi_{nm}=\cos(n\pi x/a)\cos(m\pi y/b)$, $\lambda_{nm}=(n\pi/a)^2+(m\pi/b)^2$, $n,m\ge0$, có $\lambda_{00}=0$.
(b) Dirichlet trái/dưới, Neumann phải/trên: $\phi_{nm}=\sin((n+1/2)\pi x/a)\sin((m+1/2)\pi y/b)$, $\lambda$ là tổng hai bình phương, $n,m\ge0$. Dòng $\partial_x\phi(x,b)$ của đề phải sửa thành $\partial_y\phi(x,b)$ để là Neumann ở cạnh trên. Giữ $\partial_x$ nguyên văn thì đây là đạo hàm tiếp tuyến, suy $\phi$ hằng trên cạnh trên; cùng Dirichlet ở góc trái buộc cạnh trên bằng 0, tức bài Dirichlet trên, có mode $\sin(m\pi y/b)$, $m\ge1$.
(c) Tách $\phi=XY$: giải hai bài một chiều
\[aX''+cX'=-\lambda_xX,\ X(0)=X(a_{\rm dom})=0;\quad bY''+dY'=-\lambda_yY,\ Y(0)=Y(b_{\rm dom})=0.\]
Trị riêng $\lambda=\lambda_x+\lambda_y$, hàm riêng $XY$. Khi $a(x),b(y)>0$, trọng số $\rho_x=a^{-1}e^{\int c/a}$, $\rho_y=b^{-1}e^{\int d/b}$ đưa về Sturm--Liouville; dữ kiện hàm tùy ý chỉ cho kết quả qua hai ODE này, không có danh sách số cụ thể.
(d) Đề có miền $0<x<e$, hệ số $x^2$ suy biến ở 0. Đặt $z=\ln x\in(-\infty,1)$; toán tử thành $\partial_{zz}+\partial_{yy}$. Không có phổ rời rạc như hình chữ nhật hữu hạn. Trong lớp trơn và liên tục bằng 0 ở $x=0,e$, không có hàm riêng tách biến không tầm thường: mode tăng/giảm mũ theo $z$ không thể đồng thời giảm ở $-\infty$ và triệt tiêu ở 1, mode dao động không có giới hạn 0 tại $-\infty$. Với không gian trọng số $dx\,dy/x$, phổ liên tục có hàm riêng suy rộng $\sin(k(1-\ln x))\sin(m\pi y)$, $\lambda=k^2+m^2\pi^2$, $k\ge0,m\ge1$. Nếu miền dự định là $(1,e)$, phổ rời rạc là $\sin(n\pi\ln x)\sin(m\pi y)$ với $\lambda=(n^2+m^2)\pi^2$.

\clearpage\noindent\texttt{\detokenize{BT_5_6_3_04}}\par
\section*{Lời giải bổ sung — Bài tập 4 · Mục 5.6.3}
\textit{Áp dụng kiến thức và phương pháp trong tài liệu gốc.}\par
(a) Đặt $\phi=e^{y-x}v$. Khi đó $(\Delta+2\partial_x-2\partial_y)\phi=e^{y-x}(\Delta v-2v)$. Do đó $\phi_{nm}=e^{y-x}\sin n\pi x\sin m\pi y$, $\lambda_{nm}=\pi^2(n^2+m^2)+2$, $n,m\ge1$.
(b) Chọn $\sigma=e^{2x-2y}$, tích có trọng số biến thành tích các sine, nên trực giao.
(c) $u=3+e^{y-x}\sum_{n,m\ge1}c_{nm}e^{-\lambda_{nm}t}\sin n\pi x\sin m\pi y$, với $c_{nm}=4(-1)^{n+m}/(nm\pi^2)$ là hệ số sine của $xy$.
(d) Mọi $\lambda_{nm}>0$, nên giới hạn bằng 3. Dữ liệu đầu không khớp biên ở cạnh phải/trên; nghiệm hiểu theo dữ liệu nội và biên cho $t>0$.

\clearpage\noindent\texttt{\detokenize{BT_5_6_3_05}}\par
\section*{Lời giải bổ sung — Bài tập 5 · Mục 5.6.3}
\textit{Áp dụng kiến thức và phương pháp trong tài liệu gốc.}\par
(a) Tích phân trên chu kỳ cho $\lambda\int|\phi|^2=\int|\nabla\phi|^2\ge0$.
(b) Cơ sở đầy đủ là $e^{i(nx+my)}$, $n,m\in\mathbb Z$, $\lambda=n^2+m^2$; phần thực là cos/sin của $nx+my$. Chỉ lấy $n,m\ge0$ như đề sẽ bỏ các mode $nx-my$.
(c) Tích phân trực giao cần tích hai hàm (hoặc liên hợp) và hệ số chuẩn hóa, không phải $\int\phi_{nm}$ đơn lẻ như dòng đề. Cơ sở chuẩn hóa $(2\pi)^{-1}e^{i(nx+my)}$ có $\int_{(-\pi,\pi)^2}\phi_{nm}\overline\phi_{n'm'}=\delta_{nn'}\delta_{mm'}$. Trên $(0,\pi)^2$, các sóng chéo đã in không thỏa công thức trực giao ấy nói chung.
(d) Fourier của $x$ trên $(-\pi,\pi)$ là $2\sum_{n\ge1}(-1)^{n+1}\sin nx/n$. Vì $xy$ là tích hai chuỗi,
\[u=4\sum_{n,m\ge1}\frac{(-1)^{n+m}}{nm}e^{-(n^2+m^2)t}\sin nx\sin my.\]
Giới hạn 0, nghiệm tuần hoàn trơn khi $t>0$; dữ liệu ban đầu được tuần hoàn hóa và có gián đoạn ở mép.

\clearpage\noindent\texttt{\detokenize{BT_5_6_3_06}}\par
\section*{Lời giải bổ sung — Bài tập 6 · Mục 5.6.3}
\textit{Áp dụng kiến thức và phương pháp trong tài liệu gốc.}\par
(a) Với $k_n=(n+1/2)\pi$, $\ell_m=(m+1/2)\pi$, $X_n=\sin k_nx,Y_m=\cos\ell_my$, $\lambda=k_n^2+\ell_m^2$, đặt
\[f_{nm}=4\int_0^1\int_0^1xX_nY_m\,dxdy,\quad F_{nm}=4\int_0^1\int_0^1xyX_nY_m\,dxdy.\]
Nghiệm $u=\sum[f_{nm}e^{-\lambda t}+F_{nm}(1-e^{-\lambda t})/\lambda]X_nY_m$.

(b) Đặt $u=e^{-x}v$. Toán tử thành $v_{xx}+v_{yy}-v$, với $v(0)=0,v_x(1)=v(1)$, $v_y(y=0)=0,v(y=1)=0$. Cơ sở $X_0=x$, $X_n=\sin k_nx$ với các nghiệm dương $\tan k_n=k_n$, và $Y_m=\cos\ell_my$. Đặt $\lambda_{nm}=k_n^2+1+\ell_m^2$ ($k_0=0$), $q_{nm}=\iint qX_nY_m/\iint X_n^2Y_m^2$. Dữ liệu của $v$ là $xe^x$, nguồn $txy e^x$:
\[u=e^{-x}\sum X_nY_m\left[(xe^x)_{nm}e^{-\lambda t}+(xy e^x)_{nm}\left(\frac t\lambda-\frac{1-e^{-\lambda t}}{\lambda^2}\right)\right].\]

(c) Dòng biến ở biên bị đảo thứ tự trong bản in; dùng cùng điều kiện hình học như (b). Sau $u=e^{-x}v$, hạng $+u$ khử $-v$, nguồn $ye^x$. Với $\Lambda=k_n^2+\ell_m^2>0$,
\[u=e^{-x}\sum\frac{(ye^x)_{nm}}\Lambda(1-e^{-\Lambda t})X_nY_m.\]
Các hệ số trên là tích phân một chiều của hàm sơ cấp, không chứa ẩn số chưa giải.

\clearpage\noindent\texttt{\detokenize{BT_5_6_3_07}}\par
\section*{Lời giải bổ sung — Bài tập 7 · Mục 5.6.3}
\textit{Áp dụng kiến thức và phương pháp trong tài liệu gốc.}\par
Đặt $\mathcal H$ là phép kéo dài điều hòa bốn cạnh trong hình chữ nhật theo công thức của bài 5.1.3.1, và $\mathcal H_N$ là phép kéo dài với hai cạnh đứng Neumann 0 (chuỗi cosine theo $x$, sinh theo $y$). Một nghiệm riêng $P$ cho nguồn giúp giải ngay $u=P+$phép kéo dài của dữ liệu biên trừ $P$.

(a) $P=-\sin2\pi x/(4\pi^2)+\sin3\pi x/(3\pi^2)$; $u=P+\mathcal H[-P,-P,0,0]$ trên hình vuông đơn vị.
(b) $P=-\sin\pi y/\pi^2$; $u=P+\mathcal H[-1,1,\sin\pi y/\pi^2,\sin\pi y/\pi^2]$.
(c) $P=-\sin2x/4-\sin3y/9$; $u=P+\mathcal H[-1+\sin2x/4,1+\sin2x/4,-1+\sin3y/9,1+\sin3y/9]$ trên $(0,\pi)^2$.
(d) $P=x^3y/6$; $u=P+\mathcal H[0,1-\pi x^3/6,0,-\pi^3y/6]$.
(e) Trên $(0,\pi)^2$, lấy cơ sở $X_n=\cos nx$, biên dưới 0, trên 1. Nguồn $F=xy$ có $F_0(y)=\pi y/2$, $F_n(y)=2[(-1)^n-1]y/(\pi n^2)$, $n\ge1$. Mode 0 $U_0=y/\pi+\pi(y^3-\pi^2y)/12$. Các mode $n\ge1$:
\[U_n(y)=\int_0^\pi\frac{-\sinh(n\min(y,s))\sinh(n[\pi-\max(y,s)])}{n\sinh n\pi}F_n(s)ds.\]
$u=U_0+\sum U_n\cos nx$. Biên Neumann theo $x$ thỏa tự động.

\clearpage\noindent\texttt{\detokenize{BT_5_6_3_08}}\par
\section*{Lời giải bổ sung — Bài tập 8 · Mục 5.6.3}
\textit{Áp dụng kiến thức và phương pháp trong tài liệu gốc.}\par
Đặt $S_{nm}=\sin n\pi x\sin m\pi y$, $\lambda_{nm}=\pi^2(n^2+m^2)$, $q_{nm}=4\iint qS_{nm}$. Các hệ số cần dùng: $(xy)_{nm}=4(-1)^{n+m}/(nm\pi^2)$; $(1)_{nm}=4(1-(-1)^n)(1-(-1)^m)/(nm\pi^2)$; $(x)_{nm}=4(-1)^{n+1}(1-(-1)^m)/(nm\pi^2)$.

(i) $u=\sum(xy)_{nm}e^{-\lambda t}S_{nm}$.
(ii) $u=\sum(1)_{nm}[e^{-2\lambda t}+t/(2\lambda)-(1-e^{-2\lambda t})/(4\lambda^2)]S_{nm}$.
(iii) Đổi thời gian $\tau=t^2$: $u=\sum(1)_{nm}e^{-\lambda t^2}S_{nm}$.
(iv) Biên Neumann và dữ liệu 1: chỉ phụ thuộc $x$. Fourier cosine của $x$ là $1/2+\sum 2[(-1)^n-1]\cos n\pi x/(n^2\pi^2)$, nên
\[u=1+t/2+\sum_{n\ge1}\frac{2[(-1)^n-1]}{n^4\pi^4}(1-e^{-n^2\pi^2t})\cos n\pi x.\]

(v) Bản in ghi $u_t=\Delta u$ và còn cho $u_t(x,y,0)=0$. Trong miền mở, dữ liệu $xy$ có $\Delta(xy)=0$, nên điều kiện thêm tương thích và nghiệm chính là (i); không phải phương trình sóng nếu giữ dấu $u_t$.
(vi) Bản in ghi phương trình nhiệt $u_t=2\Delta u+tx$ với $u(0)=0$ nhưng $u_t(0)=1$. Thế $t=0$ vào phương trình trong miền buộc $u_t(0)=0$, nên không có nghiệm cổ điển thỏa cả hai điều kiện. Bỏ điều kiện dư sai cho
\[u=\sum(x)_{nm}\left[\frac t{2\lambda}-\frac{1-e^{-2\lambda t}}{4\lambda^2}\right]S_{nm}.\]
Nếu ý định là sóng $u_{tt}=2\Delta u+tx$, nghiệm là $\sum[(1)_{nm}\sin(\sqrt{2\lambda}t)/\sqrt{2\lambda}+(x)_{nm}(t/(2\lambda)-\sin(\sqrt{2\lambda}t)/(2\lambda)^{3/2})]S_{nm}$.

(vii, lần đánh số vi thứ hai) Đặt $u=P+w$, với $P$ điều hòa, $P_x(0)=P_y(y=0)=0,P(1,y)=1,P(x,1)=x$. Tính $P$ bằng chồng hai chuỗi cos nửa nguyên: $X_n=\cos((n+1/2)\pi x)$; phép kéo dài cạnh trên $x$ có hệ số $2\int_0^1xX_n dx$, nhân $\cosh(k_ny)/\cosh k_n$; phép kéo dài cạnh phải 1 tương tự theo $y$. Sau đó
\[u=P+\sum c_{nm}e^{-(k_n^2+k_m^2)t}\cos k_nx\cos k_my,\quad c_{nm}=4\iint(xy-P)\cos k_nx\cos k_my.\]

\clearpage\noindent\texttt{\detokenize{BT_5_6_3_09}}\par
\section*{Lời giải bổ sung — Bài tập 9 · Mục 5.6.3}
\textit{Áp dụng kiến thức và phương pháp trong tài liệu gốc.}\par
Đặt $u=2x-1+w$. Khi đó $w(0,y)=w(1,y)=0$, $w_y(x,0)=w_y(x,1)=0$, $\Delta w=y$. Cosine của $y$ trên $(0,1)$ có $a_0=1/2$, $a_n=2[(-1)^n-1]/(n^2\pi^2)$. Mode 0 $w_0=x(x-1)/4$. Với $k_n=n\pi$,
\[w_n=\frac{a_n}{k_n^2}\left[\frac{\cosh(k_n(x-1/2))}{\cosh(k_n/2)}-1\right].\]
Vậy $u=2x-1+x(x-1)/4+\sum_{n\ge1}w_n(x)\cos n\pi y$. Mỗi mode giải $w_n''-k_n^2w_n=a_n$, đúng biên, chứng minh công thức.

\clearpage\noindent\texttt{\detokenize{BT_6_1_1_2_01}}\par
\section*{Lời giải bổ sung — Bài tập 1 · Mục 6.1.1.2}
\textit{Áp dụng kiến thức và phương pháp trong tài liệu gốc.}\par
Dùng quy ước $f=a_0+\sum_{n\ge1}(a_n\cos k_nx+b_n\sin k_nx)$, $k_n=n\pi/L$, $a_0=(2L)^{-1}\int_{-L}^Lf$, $a_n=L^{-1}\int f\cos$, $b_n=L^{-1}\int f\sin$.
(i) $a_0=\sin(\alpha L)/(\alpha L)$, $a_n=2(-1)^n\alpha\sin(\alpha L)/[L(\alpha^2-k_n^2)]$, $b_n=0$. Ở $\alpha=0$ hoặc cộng hưởng, lấy giới hạn (hàm hằng hoặc một mode cos). Chuỗi hội tụ đều đến phép kéo dài liên tục chu kỳ $2L$.
(ii) $b_n=2(1-(-1)^n)/(n\pi)$, các $a_n=0$. Chuỗi hội tụ đến sign trong miền trừ điểm 0, về 0 tại 0 và các điểm nối $\pm L$; không hội tụ đều qua bước nhảy.
(iii) $a_0=\sinh(\alpha L)/(\alpha L)$,
\[a_n=\frac{2\alpha(-1)^n\sinh(\alpha L)}{L(\alpha^2+k_n^2)},\quad b_n=\frac{2k_n(-1)^{n+1}\sinh(\alpha L)}{L(\alpha^2+k_n^2)}.\]
Với $\alpha=0$, $f=1$. Với $\alpha\ne0$, chuỗi về $e^{\alpha x}$ trong $(-L,L)$, và $\cosh\alpha L$ tại mép; hội tụ đều trên compact nội, không trên toàn chu kỳ.
(iv) Phần ảo của $\log(1+re^{ix})$ cho $f=\sum_{n\ge1}(-1)^{n+1}r^n\sin nx/n$.
(v) Hai lần phần thực cho $f=2\sum_{n\ge1}(-1)^{n+1}r^n\cos nx/n$.
(vi) Nhân Poisson cho $f=(1-r^2)^{-1}[1+2\sum_{n\ge1}r^n\cos nx]$. Ở (iv)--(vi), $0<r<1$, chuỗi hội tụ tuyệt đối và đều theo tiêu chuẩn Weierstrass.

\clearpage\noindent\texttt{\detokenize{BT_6_1_1_2_02}}\par
\section*{Lời giải bổ sung — Bài tập 2 · Mục 6.1.1.2}
\textit{Áp dụng kiến thức và phương pháp trong tài liệu gốc.}\par
(i) Khai triển Fourier $\cos(\alpha x)$ trên $[-\pi,\pi]$ và thế $x=\pi$ cho
\[\sum_{n\in\mathbb Z}\frac1{\alpha^2-n^2}=\frac\pi\alpha\cot\pi\alpha,\quad0<|\alpha|<1.\]
Điều kiện $|\alpha|<1$ trong đề còn phải loại $\alpha=0$ vì số hạng $n=0$ không xác định.
(ii) Khai triển $\cosh(\alpha x)$ (tích phân thực tương tự) và thế $x=\pi$ cho
\[\sum_{n\in\mathbb Z}\frac1{\alpha^2+n^2}=\frac\pi\alpha\coth\pi\alpha,\quad\alpha\ne0.\]
Các chuỗi hội tụ tuyệt đối vì số hạng $O(n^{-2})$, nên phép gom hai phía hợp lệ.

\clearpage\noindent\texttt{\detokenize{BT_6_1_1_2_03}}\par
\section*{Lời giải bổ sung — Bài tập 3 · Mục 6.1.1.2}
\textit{Áp dụng kiến thức và phương pháp trong tài liệu gốc.}\par
(i) $\sum\sin nx/n=(\pi-x)/2$ với $0<x<2\pi$, kéo dài chu kỳ $2\pi$; bằng 0 tại $x\in2\pi\mathbb Z$. Suy từ Fourier của hàm răng cưa hoặc phần ảo $-\log(1-re^{ix})$ rồi $r\uparrow1$.
(ii) Tích phân chuỗi (i), ấn định hằng bằng trung bình 0, cho $\sum\cos nx/n^2=\pi^2/6-\pi x/2+x^2/4$ trên $[0,2\pi]$, kéo dài tuần hoàn. Chuỗi hội tụ đều và tuyệt đối.
(iii) Tổng chuỗi hình học với $z=-e^{ix}/2$:
\[\sum_{n\ge1}(-1)^n2^{-n}\sin nx=\operatorname{Im}\frac z{1-z}=-\frac{2\sin x}{5+4\cos x}.\]
Chuỗi hội tụ tuyệt đối và đều.

\clearpage\noindent\texttt{\detokenize{BT_6_1_1_2_04}}\par
\section*{Lời giải bổ sung — Bài tập 4 · Mục 6.1.1.2}
\textit{Áp dụng kiến thức và phương pháp trong tài liệu gốc.}\par
(i) Tích phân từng phần $k$ lần, các hạng biên mất vì tuần hoàn; hệ số phức $\widehat f(n)=(in)^{-k}\widehat{f^{(k)}}(n)$. Bổ đề Riemann--Lebesgue cho $\widehat{f^{(k)}}(n)\to0$, nên $n^k(|a_n|+|b_n|)\to0$.
(ii) Giả thiết cho các chuỗi đạo hàm đến cấp $k$ được chặn bởi $M\sum n^{-2}$. Vì vậy tổng Fourier có đại diện $g\in C^k$, và $f=g$ hầu khắp nơi do tính duy nhất hệ số Fourier. Đề chỉ giả thiết $f$ khả tích nên không thể kết luận $f\in C^k$ tại mọi điểm: sửa một giá trị của $g$ không đổi bất kỳ hệ số nào nhưng có thể phá liên tục. Kết luận nguyên văn cần thêm giả thiết $f$ liên tục, hoặc hiểu $f$ như lớp tương đương hầu khắp nơi.

\clearpage\noindent\texttt{\detokenize{BT_6_1_1_2_05}}\par
\section*{Lời giải bổ sung — Bài tập 5 · Mục 6.1.1.2}
\textit{Áp dụng kiến thức và phương pháp trong tài liệu gốc.}\par
Tích thành tổng cho tích phân sine bằng
\[\frac{\sin((n-m)\pi/2)}{n-m}-\frac{\sin((n+m)\pi/2)}{n+m},\]
và cosine thay dấu trừ thành cộng. Vì $n-m$ chẵn nên $n+m$ cũng chẵn, hai tử bằng 0. Tuy nhiên đề cho $m,n\in\mathbb Z$ và chỉ loại $n=m$, chưa loại $n=-m$. Nếu $n=-m\ne0$, hai hàm cùng tần số và tích phân sine bằng $-\pi/2$, cosine bằng $\pi/2$, không bằng 0. Kết luận đúng với $n\ne\pm m$ (hoặc $m,n$ không âm khác nhau). Trường hợp chỉ số 0 xử lý trực tiếp; sine 0 đồng nhất 0.

\clearpage\noindent\texttt{\detokenize{BT_6_1_1_2_06}}\par
\section*{Lời giải bổ sung — Bài tập 6 · Mục 6.1.1.2}
\textit{Áp dụng kiến thức và phương pháp trong tài liệu gốc.}\par
Dùng $f=\sum b_n\sin k_nx$ và $f=a_0+\sum a_n\cos k_nx$, $k_n=n\pi/L$. Hệ số lần lượt $b_n=2L^{-1}\int_0^Lf\sin$, $a_n=2L^{-1}\int_0^Lf\cos$, $a_0=L^{-1}\int_0^Lf$.
(a) $b_n=2(1-(-1)^n)/(n\pi)$; cosine chỉ có $a_0=1$.
(b) $a_0=(1-\cos\alpha L)/(\alpha L)$,
\[b_n=\frac{2(-1)^nk_n\sin\alpha L}{L(\alpha^2-k_n^2)},\quad a_n=\frac{2\alpha(1-(-1)^n\cos\alpha L)}{L(\alpha^2-k_n^2)}.\]
Ở cộng hưởng lấy giới hạn; ví dụ $\alpha=k_N$ cho $b_N=1$, $b_n=0$ với $n\ne N$.
(c) $a_0=(e^{\alpha L}-1)/(\alpha L)$,
\[b_n=\frac{2k_n(1-(-1)^ne^{\alpha L})}{L(\alpha^2+k_n^2)},\quad a_n=\frac{2\alpha((-1)^ne^{\alpha L}-1)}{L(\alpha^2+k_n^2)}.\]
(d) $L=\pi$: $b_n=4(1-(-1)^n)/(\pi n^3)$, $a_0=\pi^2/6$, $a_n=-2(1+(-1)^n)/n^2$.
(e) $L=1$: $b_n=2[1+(-1)^n-2\cos(n\pi/2)]/(n\pi)$, $a_0=0$, $a_n=4\sin(n\pi/2)/(n\pi)$.

Sine là Fourier của thác triển lẻ chu kỳ $2L$, cosine là của thác triển chẵn. Tại điểm nội liên tục, cả hai về $f$; tại bước nhảy về trung bình hai phía. Sine luôn về 0 ở $0,L$, nên nếu $f$ không bằng 0 tại đó không có hội tụ đến giá trị biên. Cosine về giá trị biên một phía. Các cosine của (a)--(d) và sine của (b),(d) khi dữ liệu vanishes ở hai đầu có hội tụ đều nếu thác triển liên tục và các hệ số suy đủ nhanh; trường hợp bước nhảy (e) không đều. Các công thức hệ số chính xác ở trên cũng cho kiểm tra từng điều kiện biên.

\clearpage\noindent\texttt{\detokenize{BT_3_2_3_LT_01}}\par
\section*{Lời giải bổ sung — Bài tập 1 · Mục 3.2.3.LT}

\subsection*{Kiến thức và trình tự giải}
\textbf{Bước 1 — Đọc đúng điều kiện biên.} Biên $u(0,t)=0$ dùng mở rộng lẻ; biên $u_x(0,t)=0$ dùng mở rộng chẵn. Với dữ liệu trên $x\ge0$, đặt $\widetilde f(-s)=-f(s)$ hoặc $f(s)$ tương ứng, làm tương tự cho $g$. Mở rộng lẻ là $C^2$ nếu $f(0)=f''(0)=0$, và mở rộng lẻ của $g$ là $C^1$ nếu $g(0)=0$. Mở rộng chẵn đòi $f'(0)=g'(0)=0$ cho mức trơn ấy.

\textbf{Bước 2 — Lập nguyên hàm sau khi phản xạ.} Đặt $H(s)=\int_0^s\widetilde g(r)dr$. Nếu $\widetilde g$ lẻ thì $H$ chẵn; nếu $\widetilde g$ chẵn thì $H$ lẻ. Chính nguyên hàm, chứ không phải vận tốc, xuất hiện trong


\[
u=\frac{\widetilde f(x-at)+\widetilde f(x+at)}2+\frac{H(x+at)-H(x-at)}{2a}.
\]


\textbf{Bước 3 — Tách miền trước và sau khi chạm biên.} Với $u(0,t)=h(t)$, ở $x\ge at$ dùng d'Alembert với dữ liệu gốc. Ở $0<x<at$:


\[
u=h(t-x/a)+\frac{f(x+at)-f(at-x)}2+\frac1{2a}\int_{at-x}^{at+x}g(s)ds.
\]


Với $u_x(0,t)=h(t)$, dùng nghiệm của các mở rộng chẵn rồi trừ $a\int_0^{t-x/a}h(s)ds$ trên miền $x<at$. Tính riêng đạo hàm theo $x$ của hạng này để kiểm tra biên.

\textbf{Bước 4 — Trên đoạn hữu hạn, phản xạ ở cả hai đầu.} Dirichlet ở hai đầu cho mở rộng lẻ tại $0,L$, chu kỳ không gian $2L$. Neumann ở hai đầu cho mở rộng chẵn tại $0,L$, cũng chu kỳ $2L$. Một đầu lẻ, một đầu chẵn cho $\widetilde f(s+2L)=-\widetilde f(s)$ và chu kỳ $4L$, tương tự $g$. Vì thế chu kỳ thời gian thường là $2L/a$ hoặc $4L/a$. Với Neumann hai đầu còn phải xét trung bình $g$: trung bình khác không sinh hạng tuyến tính theo $t$, nên nghiệm không tuần hoàn.

\textbf{Bước 5 — Kiểm tra góc và đường ghép.} Biên Dirichlet không thuần nhất cần $h(0)=f(0)$, $h'(0)=g(0)$, $h''(0)=a^2f''(0)$ để có nghiệm $C^2$ cho phương trình thuần nhất. So sánh hai vết trên $x=at$, gồm $u,u_x,u_t$ rồi đạo hàm cấp hai. Nếu không tương thích, nêu rõ mức trơn còn lại thay vì khẳng định nghiệm cổ điển trên toàn miền.
\subsection*{Lời giải áp dụng cho bài}
\textbf{Bước 1. Tính liên tục tuyệt đối.} Với $F(x)=\int_0^xf(s)ds$, trên đoạn hữu hạn $[a,b]$ hàm Riemann khả tích $f$ bị chặn bởi $M$. Với các khoảng rời nhau $(a_j,b_j)$,
\[\sum_j|F(b_j)-F(a_j)|\le M\sum_j(b_j-a_j).\]
Chọn tổng độ dài nhỏ hơn $\varepsilon/M$ chứng minh liên tục tuyệt đối (trường hợp $M=0$ hiển nhiên), và $F'=f$ tại mọi điểm liên tục của $f$.

\textbf{Bước 2. Tính chẵn/lẻ.} Đổi biến $s=-r$ cho $F(-x)=-\int_0^xf(-r)dr$. Nếu $f$ lẻ, kết quả là $F(x)$; nếu $f$ chẵn, kết quả là $-F(x)$. Vì vậy nguyên hàm lấy gốc tại $0$ đổi tính chẵn/lẻ.

\textbf{Bước 3. Điều kiện tuần hoàn.} Nếu $f$ có chu kỳ $L$ thì
\[F(x+L)-F(x)=\int_x^{x+L}f(s)ds=\int_0^Lf(s)ds.\]
Do đó $F$ có chu kỳ $L$ khi và chỉ khi trung bình của $f$ bằng $0$. Nếu $f$ lẻ và tuần hoàn, tích phân trên chu kỳ đối xứng $[-L/2,L/2]$ bằng $0$, nên điều kiện đúng. Ví dụ $f=1$ tuần hoàn nhưng $F=x$ không tuần hoàn, cho thấy không được bỏ điều kiện trung bình.

\clearpage\noindent\texttt{\detokenize{BT_3_2_3_LT_02}}\par
\section*{Lời giải bổ sung — Bài tập 2 · Mục 3.2.3.LT}

\subsection*{Kiến thức và trình tự giải}
\textbf{Bước 1 — Đọc đúng điều kiện biên.} Biên $u(0,t)=0$ dùng mở rộng lẻ; biên $u_x(0,t)=0$ dùng mở rộng chẵn. Với dữ liệu trên $x\ge0$, đặt $\widetilde f(-s)=-f(s)$ hoặc $f(s)$ tương ứng, làm tương tự cho $g$. Mở rộng lẻ là $C^2$ nếu $f(0)=f''(0)=0$, và mở rộng lẻ của $g$ là $C^1$ nếu $g(0)=0$. Mở rộng chẵn đòi $f'(0)=g'(0)=0$ cho mức trơn ấy.

\textbf{Bước 2 — Lập nguyên hàm sau khi phản xạ.} Đặt $H(s)=\int_0^s\widetilde g(r)dr$. Nếu $\widetilde g$ lẻ thì $H$ chẵn; nếu $\widetilde g$ chẵn thì $H$ lẻ. Chính nguyên hàm, chứ không phải vận tốc, xuất hiện trong


\[
u=\frac{\widetilde f(x-at)+\widetilde f(x+at)}2+\frac{H(x+at)-H(x-at)}{2a}.
\]


\textbf{Bước 3 — Tách miền trước và sau khi chạm biên.} Với $u(0,t)=h(t)$, ở $x\ge at$ dùng d'Alembert với dữ liệu gốc. Ở $0<x<at$:


\[
u=h(t-x/a)+\frac{f(x+at)-f(at-x)}2+\frac1{2a}\int_{at-x}^{at+x}g(s)ds.
\]


Với $u_x(0,t)=h(t)$, dùng nghiệm của các mở rộng chẵn rồi trừ $a\int_0^{t-x/a}h(s)ds$ trên miền $x<at$. Tính riêng đạo hàm theo $x$ của hạng này để kiểm tra biên.

\textbf{Bước 4 — Trên đoạn hữu hạn, phản xạ ở cả hai đầu.} Dirichlet ở hai đầu cho mở rộng lẻ tại $0,L$, chu kỳ không gian $2L$. Neumann ở hai đầu cho mở rộng chẵn tại $0,L$, cũng chu kỳ $2L$. Một đầu lẻ, một đầu chẵn cho $\widetilde f(s+2L)=-\widetilde f(s)$ và chu kỳ $4L$, tương tự $g$. Vì thế chu kỳ thời gian thường là $2L/a$ hoặc $4L/a$. Với Neumann hai đầu còn phải xét trung bình $g$: trung bình khác không sinh hạng tuyến tính theo $t$, nên nghiệm không tuần hoàn.

\textbf{Bước 5 — Kiểm tra góc và đường ghép.} Biên Dirichlet không thuần nhất cần $h(0)=f(0)$, $h'(0)=g(0)$, $h''(0)=a^2f''(0)$ để có nghiệm $C^2$ cho phương trình thuần nhất. So sánh hai vết trên $x=at$, gồm $u,u_x,u_t$ rồi đạo hàm cấp hai. Nếu không tương thích, nêu rõ mức trơn còn lại thay vì khẳng định nghiệm cổ điển trên toàn miền.
\subsection*{Lời giải áp dụng cho bài}
\textbf{Bước 1. Dùng đối xứng của dữ liệu.} Với biên cố định hai đầu và $f(L-x)=f(x)$, $g(L-x)=g(x)$, hệ số sine $A_n=(2/L)\int_0^Lf(x)\sin(n\pi x/L)dx$. Đổi $y=L-x$ cho $A_n=(-1)^{n+1}A_n$; do đó $A_n=0$ khi $n$ chẵn. Cùng lập luận cho $B_n$.

\textbf{Bước 2. Dịch thời gian từng mode còn lại.} Với $n$ lẻ, $\omega_n=n\pi c/L$ và $\omega_nL/c=n\pi$, nên cả cos và sin đổi dấu khi dịch $L/c$. Vì vậy $u(x,t+L/c)=-u(x,t)$ và chu kỳ $2L/c$.

\textbf{Bước 3. Chứng minh bằng phản xạ nếu dữ liệu chưa trơn.} Mở rộng lẻ chu kỳ $2L$ của một hàm đối xứng giữa đoạn thỏa $f_e(s+L)=-f_e(s)$; tương tự $g_e$. Với $H'=g_e$, $H(s+L)=C-H(s)$. Thế vào d'Alembert, hằng số $C$ triệt tiêu, thu cùng phản chu kỳ. Lập luận này áp dụng trực tiếp cho dữ liệu chia đoạn, không cần đạo hàm chuỗi Fourier.

\clearpage\noindent\texttt{\detokenize{BT_3_2_3_LT_03}}\par
\section*{Lời giải bổ sung — Bài tập 3 · Mục 3.2.3.LT}

\subsection*{Kiến thức và trình tự giải}
\textbf{Bước 1 — Đọc đúng điều kiện biên.} Biên $u(0,t)=0$ dùng mở rộng lẻ; biên $u_x(0,t)=0$ dùng mở rộng chẵn. Với dữ liệu trên $x\ge0$, đặt $\widetilde f(-s)=-f(s)$ hoặc $f(s)$ tương ứng, làm tương tự cho $g$. Mở rộng lẻ là $C^2$ nếu $f(0)=f''(0)=0$, và mở rộng lẻ của $g$ là $C^1$ nếu $g(0)=0$. Mở rộng chẵn đòi $f'(0)=g'(0)=0$ cho mức trơn ấy.

\textbf{Bước 2 — Lập nguyên hàm sau khi phản xạ.} Đặt $H(s)=\int_0^s\widetilde g(r)dr$. Nếu $\widetilde g$ lẻ thì $H$ chẵn; nếu $\widetilde g$ chẵn thì $H$ lẻ. Chính nguyên hàm, chứ không phải vận tốc, xuất hiện trong


\[
u=\frac{\widetilde f(x-at)+\widetilde f(x+at)}2+\frac{H(x+at)-H(x-at)}{2a}.
\]


\textbf{Bước 3 — Tách miền trước và sau khi chạm biên.} Với $u(0,t)=h(t)$, ở $x\ge at$ dùng d'Alembert với dữ liệu gốc. Ở $0<x<at$:


\[
u=h(t-x/a)+\frac{f(x+at)-f(at-x)}2+\frac1{2a}\int_{at-x}^{at+x}g(s)ds.
\]


Với $u_x(0,t)=h(t)$, dùng nghiệm của các mở rộng chẵn rồi trừ $a\int_0^{t-x/a}h(s)ds$ trên miền $x<at$. Tính riêng đạo hàm theo $x$ của hạng này để kiểm tra biên.

\textbf{Bước 4 — Trên đoạn hữu hạn, phản xạ ở cả hai đầu.} Dirichlet ở hai đầu cho mở rộng lẻ tại $0,L$, chu kỳ không gian $2L$. Neumann ở hai đầu cho mở rộng chẵn tại $0,L$, cũng chu kỳ $2L$. Một đầu lẻ, một đầu chẵn cho $\widetilde f(s+2L)=-\widetilde f(s)$ và chu kỳ $4L$, tương tự $g$. Vì thế chu kỳ thời gian thường là $2L/a$ hoặc $4L/a$. Với Neumann hai đầu còn phải xét trung bình $g$: trung bình khác không sinh hạng tuyến tính theo $t$, nên nghiệm không tuần hoàn.

\textbf{Bước 5 — Kiểm tra góc và đường ghép.} Biên Dirichlet không thuần nhất cần $h(0)=f(0)$, $h'(0)=g(0)$, $h''(0)=a^2f''(0)$ để có nghiệm $C^2$ cho phương trình thuần nhất. So sánh hai vết trên $x=at$, gồm $u,u_x,u_t$ rồi đạo hàm cấp hai. Nếu không tương thích, nêu rõ mức trơn còn lại thay vì khẳng định nghiệm cổ điển trên toàn miền.
\subsection*{Lời giải áp dụng cho bài}
\textbf{Bước 1. Viết hàm đường thẳng dưới dạng affine.} $f(x)=m_1x+b_1$, $g(x)=m_2x+b_2$ trên đoạn xét. Với hai hệ số $a,b$,
\[af+bg=(am_1+bm_2)x+ab_1+bb_2.\]
Đây vẫn là hàm affine, có thể là hằng nếu hệ số góc bằng $0$.

\textbf{Bước 2. Co giãn và tịnh tiến đối số.} $f(ax)=m_1ax+b_1$, $f(x-h)=m_1x+b_1-m_1h$. Phải giới hạn $x$ để $ax$ hoặc $x-h$ nằm trong đoạn định nghĩa ban đầu; nếu $a<0$, chiều của đoạn đảo; nếu $a=0$, cần $0$ thuộc miền của $f$ và kết quả là hằng.

\textbf{Bước 3. Hiệu của hai dịch chuyển bậc hai.} Với $f(s)=As^2+Bs+C$, khai triển hai bình phương và triệt tiêu $Aa^2x^2,C$:
\[f(ax+h_1)-f(ax+h_2)=2Aa(h_1-h_2)x+A(h_1^2-h_2^2)+B(h_1-h_2).\]
Kết quả affine trên giao hai miền hợp lệ. Vì vậy tích phân một vận tốc affine cho nguyên hàm bậc hai, nhưng hiệu hai cận dịch cùng tốc độ có thể trở lại affine; đây là cơ sở của các đoạn thẳng trong đồ thị sóng ở mục 3.2.
\end{document}